% Options for packages loaded elsewhere
\PassOptionsToPackage{unicode}{hyperref}
\PassOptionsToPackage{hyphens}{url}
%
\documentclass[
  lang=en,
  cyan,
  device=normal]{elegantbook}
\usepackage{amsmath,amssymb}
\usepackage{iftex}
\ifPDFTeX
  \usepackage[T1]{fontenc}
  \usepackage[utf8]{inputenc}
  \usepackage{textcomp} % provide euro and other symbols
\else % if luatex or xetex
  \usepackage{unicode-math} % this also loads fontspec
  \defaultfontfeatures{Scale=MatchLowercase}
  \defaultfontfeatures[\rmfamily]{Ligatures=TeX,Scale=1}
\fi
\usepackage{lmodern}
\ifPDFTeX\else
  % xetex/luatex font selection
\fi
% Use upquote if available, for straight quotes in verbatim environments
\IfFileExists{upquote.sty}{\usepackage{upquote}}{}
\IfFileExists{microtype.sty}{% use microtype if available
  \usepackage[]{microtype}
  \UseMicrotypeSet[protrusion]{basicmath} % disable protrusion for tt fonts
}{}
\makeatletter
\@ifundefined{KOMAClassName}{% if non-KOMA class
  \IfFileExists{parskip.sty}{%
    \usepackage{parskip}
  }{% else
    \setlength{\parindent}{0pt}
    \setlength{\parskip}{6pt plus 2pt minus 1pt}}
}{% if KOMA class
  \KOMAoptions{parskip=half}}
\makeatother
\usepackage{xcolor}
\usepackage{color}
\usepackage{fancyvrb}
\newcommand{\VerbBar}{|}
\newcommand{\VERB}{\Verb[commandchars=\\\{\}]}
\DefineVerbatimEnvironment{Highlighting}{Verbatim}{commandchars=\\\{\}}
% Add ',fontsize=\small' for more characters per line
\usepackage{framed}
\definecolor{shadecolor}{RGB}{255,255,255}
\newenvironment{Shaded}{\begin{snugshade}}{\end{snugshade}}
\newcommand{\AlertTok}[1]{\textcolor[rgb]{0.75,0.01,0.01}{\textbf{\colorbox[rgb]{0.97,0.90,0.90}{#1}}}}
\newcommand{\AnnotationTok}[1]{\textcolor[rgb]{0.79,0.38,0.79}{#1}}
\newcommand{\AttributeTok}[1]{\textcolor[rgb]{0.00,0.34,0.68}{#1}}
\newcommand{\BaseNTok}[1]{\textcolor[rgb]{0.69,0.50,0.00}{#1}}
\newcommand{\BuiltInTok}[1]{\textcolor[rgb]{0.39,0.29,0.61}{\textbf{#1}}}
\newcommand{\CharTok}[1]{\textcolor[rgb]{0.57,0.30,0.62}{#1}}
\newcommand{\CommentTok}[1]{\textcolor[rgb]{0.54,0.53,0.53}{#1}}
\newcommand{\CommentVarTok}[1]{\textcolor[rgb]{0.00,0.58,1.00}{#1}}
\newcommand{\ConstantTok}[1]{\textcolor[rgb]{0.67,0.33,0.00}{#1}}
\newcommand{\ControlFlowTok}[1]{\textcolor[rgb]{0.12,0.11,0.11}{\textbf{#1}}}
\newcommand{\DataTypeTok}[1]{\textcolor[rgb]{0.00,0.34,0.68}{#1}}
\newcommand{\DecValTok}[1]{\textcolor[rgb]{0.69,0.50,0.00}{#1}}
\newcommand{\DocumentationTok}[1]{\textcolor[rgb]{0.38,0.47,0.50}{#1}}
\newcommand{\ErrorTok}[1]{\textcolor[rgb]{0.75,0.01,0.01}{\underline{#1}}}
\newcommand{\ExtensionTok}[1]{\textcolor[rgb]{0.00,0.58,1.00}{\textbf{#1}}}
\newcommand{\FloatTok}[1]{\textcolor[rgb]{0.69,0.50,0.00}{#1}}
\newcommand{\FunctionTok}[1]{\textcolor[rgb]{0.39,0.29,0.61}{#1}}
\newcommand{\ImportTok}[1]{\textcolor[rgb]{1.00,0.33,0.00}{#1}}
\newcommand{\InformationTok}[1]{\textcolor[rgb]{0.69,0.50,0.00}{#1}}
\newcommand{\KeywordTok}[1]{\textcolor[rgb]{0.12,0.11,0.11}{\textbf{#1}}}
\newcommand{\NormalTok}[1]{\textcolor[rgb]{0.12,0.11,0.11}{#1}}
\newcommand{\OperatorTok}[1]{\textcolor[rgb]{0.12,0.11,0.11}{#1}}
\newcommand{\OtherTok}[1]{\textcolor[rgb]{0.00,0.43,0.16}{#1}}
\newcommand{\PreprocessorTok}[1]{\textcolor[rgb]{0.00,0.43,0.16}{#1}}
\newcommand{\RegionMarkerTok}[1]{\textcolor[rgb]{0.00,0.34,0.68}{\colorbox[rgb]{0.88,0.91,0.97}{#1}}}
\newcommand{\SpecialCharTok}[1]{\textcolor[rgb]{0.24,0.68,0.91}{#1}}
\newcommand{\SpecialStringTok}[1]{\textcolor[rgb]{1.00,0.33,0.00}{#1}}
\newcommand{\StringTok}[1]{\textcolor[rgb]{0.75,0.01,0.01}{#1}}
\newcommand{\VariableTok}[1]{\textcolor[rgb]{0.00,0.34,0.68}{#1}}
\newcommand{\VerbatimStringTok}[1]{\textcolor[rgb]{0.75,0.01,0.01}{#1}}
\newcommand{\WarningTok}[1]{\textcolor[rgb]{0.75,0.01,0.01}{#1}}
\usepackage{longtable,booktabs,array}
\usepackage{calc} % for calculating minipage widths
% Correct order of tables after \paragraph or \subparagraph
\usepackage{etoolbox}
\makeatletter
\patchcmd\longtable{\par}{\if@noskipsec\mbox{}\fi\par}{}{}
\makeatother
% Allow footnotes in longtable head/foot
\IfFileExists{footnotehyper.sty}{\usepackage{footnotehyper}}{\usepackage{footnote}}
\makesavenoteenv{longtable}
\usepackage{graphicx}
\makeatletter
\def\maxwidth{\ifdim\Gin@nat@width>\linewidth\linewidth\else\Gin@nat@width\fi}
\def\maxheight{\ifdim\Gin@nat@height>\textheight\textheight\else\Gin@nat@height\fi}
\makeatother
% Scale images if necessary, so that they will not overflow the page
% margins by default, and it is still possible to overwrite the defaults
% using explicit options in \includegraphics[width, height, ...]{}
\setkeys{Gin}{width=\maxwidth,height=\maxheight,keepaspectratio}
% Set default figure placement to htbp
\makeatletter
\def\fps@figure{htbp}
\makeatother
\setlength{\emergencystretch}{3em} % prevent overfull lines
\providecommand{\tightlist}{%
  \setlength{\itemsep}{0pt}\setlength{\parskip}{0pt}}
\setcounter{secnumdepth}{5}
\ifLuaTeX
\usepackage[bidi=basic]{babel}
\else
\usepackage[bidi=default]{babel}
\fi
\babelprovide[main,import]{russian}
% get rid of language-specific shorthands (see #6817):
\let\LanguageShortHands\languageshorthands
\def\languageshorthands#1{}
% ===== Настройки ElegantBook для русского community-издания =====
% Accent color comes from ElegantBook's cyan theme: \structurecolor = RGB(31,186,190)

% ── Шрифты с полной поддержкой кириллицы ─────────────────────────
\usepackage{unicode-math}
\setmainfont{PT Serif}
\usepackage{etoolbox}
\linespread{1}
\apptocmd{\normalsize}{\fontsize{11}{15}\selectfont}{}{}
\AtBeginDocument{\normalsize}
% Allow looser spacing only when normal paragraph breaking would overflow.
\setlength{\emergencystretch}{8em}
\setsansfont{DejaVu Sans}[Scale=0.9]
\setmathfont{texgyrepagella-math}[Extension = .otf]

% ── Unicode glyph fallback (the text/mono fonts lack these) ──────
\usepackage{newunicodechar}
\newfontfamily\fallbackfont{DejaVu Sans}
\newunicodechar{→}{{\fallbackfont →}}
\newunicodechar{↔}{{\fallbackfont ↔}}
\newunicodechar{Σ}{{\fallbackfont Σ}}
\newunicodechar{★}{{\fallbackfont ★}}
\newunicodechar{✓}{{\fallbackfont ✓}}
\newunicodechar{✗}{{\fallbackfont ✗}}
\newunicodechar{₀}{{\fallbackfont ₀}}
\newunicodechar{₁}{{\fallbackfont ₁}}
\newunicodechar{₂}{{\fallbackfont ₂}}
\newunicodechar{₃}{{\fallbackfont ₃}}
\newunicodechar{₄}{{\fallbackfont ₄}}
\newunicodechar{₅}{{\fallbackfont ₅}}
\newunicodechar{₆}{{\fallbackfont ₆}}
\newunicodechar{₇}{{\fallbackfont ₇}}
\newunicodechar{₈}{{\fallbackfont ₈}}
\newunicodechar{₉}{{\fallbackfont ₉}}
\newunicodechar{ⓐ}{{\fallbackfont ⓐ}}
\newunicodechar{①}{{\fallbackfont ①}}
\newunicodechar{②}{{\fallbackfont ②}}
\newunicodechar{③}{{\fallbackfont ③}}
\newunicodechar{④}{{\fallbackfont ④}}
\newunicodechar{⑤}{{\fallbackfont ⑤}}
\newunicodechar{ⓑ}{{\fallbackfont ⓑ}}
\newunicodechar{ⓒ}{{\fallbackfont ⓒ}}
\newunicodechar{ⓓ}{{\fallbackfont ⓓ}}
\newunicodechar{ⓔ}{{\fallbackfont ⓔ}}
% Math / symbols absent from the Latin text font
\newunicodechar{≈}{{\fallbackfont ≈}}
\newunicodechar{♠}{{\fallbackfont ♠}}
\newunicodechar{♣}{{\fallbackfont ♣}}
\newunicodechar{♥}{{\fallbackfont ♥}}
\newunicodechar{♦}{{\fallbackfont ♦}}
% Greek letters used inline in body text (Latin Modern lacks them here)
\newunicodechar{α}{{\fallbackfont α}}
\newunicodechar{β}{{\fallbackfont β}}
\newunicodechar{γ}{{\fallbackfont γ}}
\newunicodechar{δ}{{\fallbackfont δ}}
\newunicodechar{ε}{{\fallbackfont ε}}
\newunicodechar{η}{{\fallbackfont η}}
\newunicodechar{θ}{{\fallbackfont θ}}
\newunicodechar{λ}{{\fallbackfont λ}}
\newunicodechar{μ}{{\fallbackfont μ}}
\newunicodechar{π}{{\fallbackfont π}}
\newunicodechar{ρ}{{\fallbackfont ρ}}
\newunicodechar{σ}{{\fallbackfont σ}}
\newunicodechar{τ}{{\fallbackfont τ}}
\newunicodechar{φ}{{\fallbackfont φ}}
\newunicodechar{ω}{{\fallbackfont ω}}

% ── Code/monospace font (covers box-drawing glyphs used in code) ──
\setmonofont[Scale=0.86]{DejaVu Sans Mono}

% ── Tables: retain readable columns after Russian text expansion ──
\usepackage{ragged2e}
\AtBeginEnvironment{longtable}{\let\raggedright\RaggedRight\sloppy}

% ── tcolorbox (ElegantBook already loads it; just add libraries) ──
\tcbuselibrary{skins,breakable}

% ── TikZ (for the full-bleed image cover in cover.tex) ───────────
\usepackage{tikz}
\usetikzlibrary{fadings}
% Vertical fade: invisible at the top, solid at the bottom — lets the cover
% art melt smoothly into the navy title band.
\tikzfading[name=titlefade, top color=transparent!100, bottom color=transparent!0]

% Theme accent color: override ElegantBook's bright cyan with a deep navy
% (classic, authoritative technical-book tone). Drives headings, rules, links,
% code chrome and the cover. Tweak this one line to restyle the whole book.
\definecolor{structurecolor}{RGB}{30,58,107}

% Accent shade for code/figure chrome (follows the theme color above)
\colorlet{accent}{structurecolor}

% Internal cross-reference links (图N-M / 第N章) — clickable + subtle teal,
% independent of the document class's global linkcolor. Used by crossref.lua.
\newcommand{\crossreflink}[2]{\hyperref[#1]{\textcolor{structurecolor}{#2}}}
\colorlet{accentbg}{structurecolor!6}
\colorlet{accentrule}{structurecolor!35}

% ── Wrap long lines in code blocks ───────────────────────────────
% pandoc emits fancyvrb's Verbatim (via `Highlighting` for ```lang blocks
% and plain `verbatim` for un-tagged blocks); neither wraps by default, so
% long lines overflow the page. fvextra adds breaklines; re-define both
% environments to wrap, breaking mid-token when a chunk has no spaces (URLs,
% long identifiers, CJK comments). A teal ↳ marks each continued line.
\usepackage{fvextra}
\fvset{breaklines=true,breakanywhere=true,%
  breaksymbolleft={\footnotesize\textcolor{accentrule}{$\hookrightarrow$}\,}}
\DefineVerbatimEnvironment{Highlighting}{Verbatim}{commandchars=\\\{\},breaklines,breakanywhere}

% ── Code blocks: subtle O'Reilly-style framed box ────────────────
% One shared box style for ALL code blocks so highlighted (```lang → Shaded)
% and plain (``` → verbatim) blocks look identical. The plain verbatim nests
% a fvextra Verbatim inside the box; \VerbatimEnvironment makes it end at
% \end{verbatim} (the outer env name) rather than \end{Verbatim}.
\newtcolorbox{codebox}{%
  breakable, enhanced,
  colback=black!3, colframe=accentrule,
  boxrule=0.5pt, leftrule=2.5pt,
  left=0.7em, right=0.6em, top=0.4em, bottom=0.4em,
  arc=0.6mm,
  before skip=0.7em, after skip=0.7em
}
\renewenvironment{Shaded}{\begin{codebox}}{\end{codebox}}
\renewenvironment{verbatim}%
  {\VerbatimEnvironment\begin{codebox}\begin{Verbatim}}%
  {\end{Verbatim}\end{codebox}}

% ── Experiment & question callouts (used by experiment_box.lua) ───
\newtcolorbox{experimentbox}{
  breakable, enhanced,
  colback=accentbg, colframe=accent,
  boxrule=0.6pt, left=0.8em, right=0.8em, top=0.6em, bottom=0.6em,
  arc=1mm, before skip=0.9em, after skip=0.9em
}

\newtcolorbox{questionbox}{
  breakable, enhanced,
  colback=accentbg, colframe=accent,
  boxrule=0.8pt, left=0.9em, right=0.9em, top=0.9em, bottom=0.7em,
  arc=1.2mm, before skip=0.9em, after skip=0.9em,
  title={\sffamily\bfseries Вопросы для размышления},
  fonttitle=\normalsize\color{white},
  coltitle=white,
  attach boxed title to top left={xshift=1em, yshift=-\tcboxedtitleheight/2},
  boxed title style={colback=accent, colframe=accent, arc=0.8mm, outer arc=0.8mm}
}

% ── Figures: force [H] so they can sit inside callout/quote boxes ─
\usepackage{float}
\let\origfigure\figure
\let\origendfigure\endfigure
\renewenvironment{figure}[1][htbp]{\origfigure[H]}{\origendfigure}
\setkeys{Gin}{width=\linewidth,height=0.82\textheight,keepaspectratio}

% ── Figure captions: manual numbering in text, so suppress label ──
\usepackage{caption}
% Sans, bold, centered. Figure numbers live in the caption text itself
% (crossref.lua), so the automatic "Figure N:" label is suppressed.
\captionsetup{font={normalsize,sf,bf}, labelformat=empty, skip=4pt, justification=centering}

% ── Blockquotes: teal left bar (experiments render as quotes) ────
\renewenvironment{quote}{%
  \begin{tcolorbox}[
    blanker, breakable,
    borderline west={2.5pt}{0pt}{accent},
    colback=accent!4,
    left=1em, right=0.6em, top=0.5em, bottom=0.5em,
    before skip=0.6em, after skip=0.6em
  ]%
}{%
  \end{tcolorbox}%
}

% ── Inline code: keep it tight ───────────────────────────────────
\setlength{\parindent}{2em}
\ifLuaTeX
  \usepackage{selnolig}  % disable illegal ligatures
\fi
\IfFileExists{bookmark.sty}{\usepackage{bookmark}}{\usepackage{hyperref}}
\IfFileExists{xurl.sty}{\usepackage{xurl}}{} % add URL line breaks if available
\urlstyle{same}
\hypersetup{
  pdftitle={AI-агенты изнутри: принципы проектирования и инженерная
практика},
  pdflang={ru-RU},
  hidelinks,
  pdfcreator={LaTeX via pandoc}}

\author{Bojie Li}
\date{}

\newunicodechar{℃}{{\fallbackfont ℃}}\newunicodechar{κ}{{\fallbackfont κ}}\newunicodechar{⁵}{{\fallbackfont ⁵}}\IfFontExistsTF{Noto Sans CJK SC}{\newcommand{\cjkmain}{Noto Sans CJK SC}}{%
  \IfFontExistsTF{Microsoft YaHei}{\newcommand{\cjkmain}{Microsoft YaHei}}{%
    \IfFontExistsTF{SimSun}{\newcommand{\cjkmain}{SimSun}}{%
      \IfFontExistsTF{Songti SC}{\newcommand{\cjkmain}{Songti SC}}{}}}}
\IfFileExists{xeCJK.sty}{}{\let\cjkmain\undefined}
\ifdefined\cjkmain
  \usepackage{xeCJK}
  \setCJKmainfont{\cjkmain}
  \setCJKsansfont{\cjkmain}
  \setCJKmonofont{\cjkmain}
\fi

% ── Юникод-глифы, которых нет в основном шрифте ───────────────────

% Minimal TeX installations may have a CJK font but not xeCJK.
\ifdefined\cjkmain\else
  \IfFontExistsTF{Droid Sans Fallback}{%
    \newfontfamily\bookhanfont{Droid Sans Fallback}
    \newcommand{\bookhanglyph}[1]{\newunicodechar{#1}{{\bookhanfont #1}}}
    \forcsvlist{\bookhanglyph}{李,博,杰,陈,耳,东,嘉,奖,怡,然,佳,宜}
  }{}
\fi

\begin{document}
\frontmatter

% Cover page — a hand-drawn "agent motif": a central AI core with radiating
% arms ending in tool glyphs (the "one agent + many tools" idea). Pure vector
% TikZ, navy line art on white — no external image, fully reproducible. The
% build date is stamped near the bottom so successive versions are distinguishable.
\begin{titlepage}
\thispagestyle{empty}
% Thin navy top band + footer rule (series look; no tagline text).
\begin{tikzpicture}[remember picture, overlay]
  \fill[structurecolor] (current page.north west) rectangle ([yshift=-0.85cm]current page.north east);
  \fill[structurecolor] (current page.south west) rectangle ([yshift=0.5cm]current page.south east);
\end{tikzpicture}
\centering
\vspace*{2.5cm}

{\fontsize{31}{39}\selectfont\rmfamily\bfseries AI-агенты изнутри:\par}
\vspace{0.5cm}
{\Large\sffamily\color{structurecolor} принципы проектирования и инженерная практика\par}

\vspace{1.5cm}

% ── Agent motif: central core + radiating arms ending in tool glyphs ──
\begingroup
\definecolor{ink}{RGB}{30,58,107}
\begin{tikzpicture}[line join=round, line cap=round,
    arm/.style={ink, line width=1.1pt},
    ring/.style={ink, line width=1.1pt, fill=white},
    ic/.style={ink, line width=0.8pt}]
  \def\R{3.15}    % arm length (hub centre → tool node)
  \def\rh{0.98}   % hub radius

  % arms first, so tool nodes sit on top of their ends
  \foreach \a in {90,45,0,-45,-90,-135,180,135}{
    \draw[arm] (\a:\rh) to[bend left=8] (\a:\R);
  }

  % hub (the "agent core") + a 4-point AI spark inside
  \fill[ink!7] (0,0) circle (\rh);
  \draw[ink, line width=1.3pt] (0,0) circle (\rh);
  \fill[ink] (0,0.52) -- (0.13,0.13) -- (0.52,0) -- (0.13,-0.13) --
             (0,-0.52) -- (-0.13,-0.13) -- (-0.52,0) -- (-0.13,0.13) -- cycle;

  % ── tool nodes (r=0.52) with a simple glyph each ──
  % 90° — search (magnifier)
  \begin{scope}[shift={(90:\R)}]
    \draw[ring] (0,0) circle (0.52);
    \draw[ic] (-0.05,0.06) circle (0.15);
    \draw[ic] (0.06,-0.05) -- (0.20,-0.19);
  \end{scope}
  % 45° — code </>
  \begin{scope}[shift={(45:\R)}]
    \draw[ring] (0,0) circle (0.52);
    \draw[ic] (-0.06,0.17) -- (-0.21,0) -- (-0.06,-0.17);
    \draw[ic] (0.06,0.17) -- (0.21,0) -- (0.06,-0.17);
  \end{scope}
  % 0° — terminal
  \begin{scope}[shift={(0:\R)}]
    \draw[ring] (0,0) circle (0.52);
    \draw[ic] (-0.24,-0.17) rectangle (0.24,0.17);
    \draw[ic] (-0.15,0.07) -- (-0.07,0) -- (-0.15,-0.07);
    \draw[ic] (0.01,-0.08) -- (0.15,-0.08);
  \end{scope}
  % -45° — document
  \begin{scope}[shift={(-45:\R)}]
    \draw[ring] (0,0) circle (0.52);
    \draw[ic] (-0.15,-0.21) -- (-0.15,0.21) -- (0.07,0.21) -- (0.17,0.11) -- (0.17,-0.21) -- cycle;
    \draw[ic] (0.07,0.21) -- (0.07,0.11) -- (0.17,0.11);
    \draw[ic] (-0.08,0.05) -- (0.10,0.05);
    \draw[ic] (-0.08,-0.05) -- (0.10,-0.05);
  \end{scope}
  % -90° — gear
  \begin{scope}[shift={(-90:\R)}]
    \draw[ring] (0,0) circle (0.52);
    \draw[ic] (0,0) circle (0.13);
    \foreach \g in {0,45,90,135,180,225,270,315}{ \draw[ic] (\g:0.14) -- (\g:0.22); }
    \fill[ink] (0,0) circle (0.035);
  \end{scope}
  % -135° — globe / web
  \begin{scope}[shift={(-135:\R)}]
    \draw[ring] (0,0) circle (0.52);
    \draw[ic] (0,0) circle (0.20);
    \draw[ic] (0,0) ellipse (0.08 and 0.20);
    \draw[ic] (-0.20,0) -- (0.20,0);
  \end{scope}
  % 180° — database
  \begin{scope}[shift={(180:\R)}]
    \draw[ring] (0,0) circle (0.52);
    \draw[ic] (0,0.14) ellipse (0.18 and 0.07);
    \draw[ic] (-0.18,0.14) -- (-0.18,-0.14);
    \draw[ic] (0.18,0.14) -- (0.18,-0.14);
    \draw[ic] (-0.18,0.0) arc[start angle=180, end angle=360, x radius=0.18, y radius=0.07];
    \draw[ic] (-0.18,-0.14) arc[start angle=180, end angle=360, x radius=0.18, y radius=0.07];
  \end{scope}
  % 135° — chat bubble
  \begin{scope}[shift={(135:\R)}]
    \draw[ring] (0,0) circle (0.52);
    \draw[ic, rounded corners=2pt] (-0.20,-0.03) rectangle (0.20,0.21);
    \draw[ic] (-0.11,-0.03) -- (-0.05,-0.16) -- (0.02,-0.03);
    \fill[ink] (-0.10,0.09) circle (0.022);
    \fill[ink] (0,0.09) circle (0.022);
    \fill[ink] (0.10,0.09) circle (0.022);
  \end{scope}
\end{tikzpicture}
\endgroup

\vfill
{\LARGE\sffamily Bojie Li\par}
\vspace{0.45cm}
{\small\sffamily\color{structurecolor!80} Русский перевод: community edition\par}
\vspace{0.3cm}
{\small\sffamily\color{structurecolor!80} Редакция 3.1 · 40f4ce57 · \today\par}
\vspace{1.2cm}
\end{titlepage}

{
\setcounter{tocdepth}{2}
\tableofcontents
}
\mainmatter
\hypertarget{ux43fux440ux435ux434ux438ux441ux43bux43eux432ux438ux435}{%
\chapter*{Предисловие}\label{ux43fux440ux435ux434ux438ux441ux43bux43eux432ux438ux435}}
\addcontentsline{toc}{chapter}{Предисловие}

С августа 2025 года до начала 2026 года автор проводил у издательства
«Тьюринг» онлайн-видеокурс «Практикум по агентам», посвящённый созданию
AI-агентов с помощью объяснимых, проверяемых и переиспользуемых
инженерных методов. Курс начинался с запуска агента, после чего
разбиралось, почему он устроен именно так: какую проблему решает каждое
проектное решение, какой ценой и для каких задач оно подходит. Эта книга
подготовлена на основе материалов курса и сопутствующих практических
работ.

\begin{quote}
\textbf{Как создавалась эта книга: опыт сотрудничества человека и
машины}

При подготовке к занятиям я сначала надиктовывал план голосовому агенту
Pine, а тот искал материалы и составлял черновик. Такой формат
совместной работы с помощью диктовки называется \textbf{whisper coding}.
После занятий, если слушатели застревали на каком-то шаге или результат
эксперимента оказывался неожиданным, я возвращался к соответствующим
примерам и аргументам и перепроверял их.
\end{quote}

С 2025 года автор наблюдает важное изменение: фокус исследований и
прикладных разработок в области AI расширяется --- от повышения
возможностей базовых моделей до выполнения агентами реальных задач.
Благодаря этому обучение моделей и их практическое применение становятся
всё теснее связаны. Модели рассуждений, такие как DeepSeek-R1,
продемонстрировали, как обучение с подкреплением повышает способности в
математике, программировании и рассуждениях.\footnote{DeepSeek-AI,
  \emph{DeepSeek-R1: Incentivizing Reasoning Capability in LLMs via
  Reinforcement Learning}, 2025.} Обучение агентов с подкреплением
(Agentic Reinforcement Learning) включает в процесс обучения вызовы
инструментов и взаимодействие со средой, благодаря чему модель учится не
только программировать, но и работать с графическими интерфейсами.
Задачи второго типа часто называют Computer Use: модель определяет
состояние страницы по снимку экрана, а затем выполняет необходимые
действия с помощью щелчков мыши, ввода текста и прокрутки. Такие
продукты, как Manus, Claude Code и OpenClaw, позволяют пользователям
напрямую поручать задачи, а также показывают, как сочетание «генерации
кода и файловой системы» может служить основой для универсального
агента.

В сентябре 2026 года Meta представила Muse. В тот период продукты с
персональными агентами начали появляться один за другим как в Китае, так
и за его пределами. Со многими проблемами, рассматриваемыми в этой
книге, мы с командой столкнулись ещё раньше, работая над Pine. Я занимал
должность главного научного сотрудника в Pine AI и вместе с коллегами
разрабатывал персонального агента, который выполняет поручения
потребителей. Уже в начале января 2025 года Pine был запущен в США:
пользователи могли попросить его позвонить компании, договориться об
уменьшении счёта, запросить возврат средств или отменить подписку.

Если представитель компании по телефону соглашался списать часть
начислений, Pine мог сразу сообщить пользователю только одно: «Сотрудник
службы поддержки согласился». Чтобы сказать: «Деньги уже вернули», нужно
было дождаться, пока возврат появится в выписке. Для выполнения таких
поручений модель должна помнить, на что пользователь дал разрешение,
понимать, что следует проверить дальше, и различать этапы выполнения
задачи. Позднее в отрасли начали говорить о Skill, Harness и инженерии
Loop, однако многие из соответствующих подходов мы нащупали именно во
время таких телефонных разговоров.

\begin{itemize}
\tightlist
\item
  \textbf{Skill.} Мы загружали соответствующие промпты только тогда,
  когда они требовались для текущей задачи. Когда инструментов стало
  больше, мы также научили модель вызывать их через командную строку,
  чтобы не перечислять определения всех инструментов в каждом запросе.
\item
  \textbf{Harness.} Мы выводили среду выполнения, текущее время и
  рабочее состояние пользователя в строку состояния, чтобы модель могла
  учитывать их при принятии решений. Иногда модель выдумывала
  результаты, вызывала не тот инструмент или упускала инструкции, а
  опасные действия могли выйти за пределы полномочий, предоставленных
  пользователем. Поэтому мы ввели ограничения для процесса выполнения;
  позднее похожие решения появились и в Claude Code.
\item
  \textbf{Инженерия Loop.} Иногда модель объявляла задачу выполненной,
  сделав лишь половину работы. Мы поручали одному агенту выполнить
  задачу, а другому --- проверить счёт или заказ. Если проверяющий
  обнаруживал упущение, задача возвращалась первому агенту на доработку.
  Так устроено разделение ролей «исполнитель --- проверяющий»
  (proposer-reviewer).
\end{itemize}

Другие команды, разрабатывавшие модели и агентов и решавшие похожие
задачи, независимо пришли к близким методам. Это ещё сильнее убедило
автора в том, что многие трудности, с которыми столкнулась наша команда,
являются общими для агентных систем, переходящих к практическому
применению. Помимо «Практикума по агентам», в 2024--2026 годах автор
также регулярно преподавал практический курс по агентам в Университете
Китайской академии наук. Вопросы слушателей обоих курсов помогли понять,
какие аспекты накопленного опыта требуют более ясного объяснения.
Публикуя книгу в открытом доступе, автор также надеется, что больше
исследователей и инженеров смогут применять, проверять и
совершенствовать эти методы.

\textbf{Сначала практика, потом название}. Для команд разработчиков это
означает ещё и следующее: когда какой-либо термин становится
общепринятым в отрасли, ведущие команды зачастую уже давно исследуют
стоящее за ним направление. Чтобы накопить собственные инженерные
преимущества, необходимо заранее выявлять проблемы в реальном бизнесе и
пробовать их решать. По мнению автора, для этого нужны два ключевых
условия.

\textbf{Во-первых, нужно выбрать реальный бизнес, предъявляющий
достаточно высокие требования к возможностям агентов, и постоянно
получать обратную связь.} Одна задача Pine могла длиться от нескольких
часов до нескольких недель, требовать общения с несколькими сторонами,
продолжительных телефонных разговоров, заполнения сложных форм и
многоэтапной переписки по электронной почте. При этом необходимо было
обеспечивать точность чисел, управляемость действий и защиту интересов
пользователя. Такие сценарии постоянно выявляют разрыв между
возможностями модели и требованиями бизнеса, побуждая команду
совершенствовать Harness. Если основная ценность продукта заключается
лишь в компенсации недостатков конкретного поколения моделей, следующее
обновление модели может устранить эти недостатки, и команде будет трудно
накопить устойчивые преимущества.

\textbf{Во-вторых, необходимо создать системный механизм оценки
(Evaluation).} Если после изменения системы переговоры об уменьшении
счёта завершились успешно, причиной могло быть как то, что помощник
задал правильные вопросы, так и то, что именно в этот раз сотрудник
службы поддержки случайно обладал полномочиями списать начисления.
Многократное сравнение новой и старой версий на одном и том же наборе
примеров позволяет понять, устойчивы ли различия. В седьмой главе
подробно рассматривается, как задавать критерии успеха, создавать среду
оценки, анализировать неудачи и проверять изменения.

Работая над Pine, мы извлекли из задач по уменьшению счетов один метод:
проверяющий должен посмотреть, появилась ли соответствующая сумма в
выписке. Если применить тот же метод к обмену билета, нужно проверить,
выпущен ли новый билет. Конкретные подтверждения зависят от бизнеса, но
за ними стоит общий принцип: чтобы определить, выполнена ли задача,
нужно проверить, какой результат действие оставило во внешнем мире.
\textbf{Хорошие принципы сохраняют силу при смене поколений моделей}.

\hypertarget{ux43cux438ux441ux441ux438ux44f-ux44dux442ux43eux439-ux43aux43dux438ux433ux438}{%
\section*{Миссия этой
книги}\label{ux43cux438ux441ux441ux438ux44f-ux44dux442ux43eux439-ux43aux43dux438ux433ux438}}
\addcontentsline{toc}{section}{Миссия этой книги}

Представьте, что человек передаёт помощнику счёт за мобильную связь и
просит связаться с оператором, чтобы договориться о снижении суммы или
освобождении от оплаты. Прежде чем начать, помощник должен выяснить,
какое именно начисление оспаривается; во время разговора --- выслушать
объяснения и уточнить условия; по завершении --- сообщить пользователю,
что именно пообещала служба поддержки. На каждом шаге помощнику
необходимо отвечать на три вопроса: что ему известно в данный момент?
Какие действия он может предпринять? На каком основании он выбирает
следующий шаг? Отталкиваясь от этих трёх вопросов, книга знакомит со
структурой AI-агента, методами его построения и границами возможностей.
Основная формула: \textbf{агент = LLM + контекст + инструменты}. LLM
отвечает за понимание, планирование и принятие решений, контекст
содержит информацию, необходимую для текущего решения, а инструменты
связывают агента с внешней средой, позволяя получать наблюдения или
выполнять действия.

Если воспользоваться наглядной аналогией, LLM --- это «мозг», контекст
--- «поле зрения при принятии текущего решения», а инструменты ---
«органы чувств, руки и ноги». В поле зрения входят сведения о среде,
история диалога, память и знания, а также системные инструкции,
определения инструментов и состояние задачи. Поиск, чтение и создание
снимков экрана приносят информацию, тогда как нажатия, запись и вызовы
бизнес-интерфейсов изменяют внешнее состояние. На
\crossreflink{fig:0-1}{рисунке 0-1} показан процесс одного запроса по
счёту. При первом решении модель знает лишь, что пользователь
сомневается в правильности начисления, а при втором уже может указать, о
какой именно сумме следует спросить.

\begin{figure}
\hypertarget{fig:0-1}{%
\centering
\includegraphics{images/revised/fig0-1.pdf}
\caption{Рисунок 0-1. LLM, контекст и инструменты при обработке запроса
по счёту}\label{fig:0-1}
}
\end{figure}

Читатели, знакомые с обучением с подкреплением (Reinforcement Learning,
RL), могут также осмыслить это разделение через политику, наблюдения и
действия. LLM выбирает следующий шаг на основе контекста, поэтому её
функция принятия решений близка к политике (Policy); контекст содержит
конкретные наблюдения и историю, используемые при этом решении;
инструменты предоставляют интерфейсы наблюдений и действий. Пространство
наблюдений (Observation Space) описывает, какие наблюдения можно
получить, а пространство действий (Action Space) --- какие действия
разрешено выполнять; вместе они очерчивают границы взаимодействия агента
со средой.\footnote{Раздел «Spaces» документации Gymnasium.} В первой
главе эти понятия будут раскрыты на примерах вызова инструментов, а
также будет рассмотрено распределение обязанностей между Model и Harness
в производственной реализации.

\begin{quote}
\textbf{Дополнительное размышление: от жизни к спроектированным
сущностям}

Помимо инженерной практики агенты поднимают вопросы об эволюции
интеллектуальных систем. Лауреат премии Тьюринга и один из
основоположников обучения с подкреплением Richard Sutton разделяет
эволюцию Вселенной на четыре этапа: пыль, звёзды, жизнь и
«спроектированные сущности» (designed entities).\footnote{Интервью
  Ричарда Саттона и Дваркеша Пателя от 26 сентября 2025 года, раздел
  «Succession to AI».} Биологические организмы эволюционируют
посредством случайных мутаций и естественного отбора, тогда как
спроектированные сущности могут целенаправленно создаваться и
совершенствоваться уже существующими интеллектуальными системами.

Агент может генерировать код, создавать новых агентов, затем
анализировать их поведение во время выполнения, предлагать изменения и
проверять результат, постепенно реализуя самозагрузку (bootstrap) и
непрерывную эволюцию. Это похоже на то, как один программист создаёт
другого программиста, который затем продолжает участвовать в построении
новых систем. Возможность понимать и совершенствовать себя --- один из
вопросов, которые больше всего интересовали автора при написании глав о
генерации кода и непрерывной эволюции.
\end{quote}

\hypertarget{ux441ux442ux440ux443ux43aux442ux443ux440ux430-ux43aux43dux438ux433ux438}{%
\section*{Структура
книги}\label{ux441ux442ux440ux443ux43aux442ux443ux440ux430-ux43aux43dux438ux433ux438}}
\addcontentsline{toc}{section}{Структура книги}

На \crossreflink{fig:0-2}{рис. 0-2} показаны связи между главами.
Допустим, во время телефонного разговора пользователь внезапно меняет
требования, но ассистент продолжает действовать по старым. Тогда можно
начать с \crossreflink{chap:6}{главы 6} в левой части рисунка и
посмотреть, как ассистент обрабатывает реплики, прерывающие его работу.
А как после исправления убедиться, что в других разговорах он тоже будет
учитывать новые требования? Для этого применяются методы оценки из
\crossreflink{chap:7}{главы 7} в правой части рисунка. Три стрелки,
исходящие из \crossreflink{chap:7}{главы 7}, показывают, что после
получения результатов оценки можно выбрать разные способы улучшения в
зависимости от причин неудачи.

\begin{figure}
\hypertarget{fig:0-2}{%
\centering
\includegraphics{images/revised/fig0-2.pdf}
\caption{Рисунок 0-2. Структура книги: создание AI-агентов и
совершенствование их возможностей}\label{fig:0-2}
}
\end{figure}

\begin{itemize}
\tightlist
\item
  \textbf{\crossreflink{chap:1}{Глава 1} (Введение в AI-агенты)} на
  примерах продуктов и вызовов инструментов знакомит с основными
  компонентами, рабочим циклом, Harness и способами оркестрации,
  формируя понятийную основу для понимания всей книги.
\item
  \textbf{\crossreflink{chap:2}{Глава 2} (Контекстная инженерия)} ---
  ключевая глава книги, посвящённая организации информации, которая
  нужна модели в текущий момент. Рассматриваются проектирование
  промптов, загрузка по требованию и управление историей, а на примере
  KV-кэша объясняется, как эти методы влияют на качество ответов и
  стоимость работы.
\item
  \textbf{\crossreflink{chap:3}{Глава 3} (Пользовательская память и база
  знаний)} посвящена сохранению и использованию информации между
  сеансами. Рассматриваются пользовательская память, RAG,
  структурированные индексы и обновление знаний, помогающие читателю
  освоить методы поиска и сопровождения постоянных знаний.
\item
  \textbf{\crossreflink{chap:4}{Глава 4} (Инструменты)} рассматривает,
  как агент получает наблюдения, выполняет действия и запрашивает
  совместную работу. Классификация инструментов, проектирование
  интерфейсов, MCP и управление выполнением помогают понять, как
  подключать к системе внешние возможности.
\item
  \textbf{\crossreflink{chap:5}{Глава 5} (Coding Agent и агенты общего
  назначения)} предлагает одну из ключевых архитектурных идей для агента
  общего назначения: использовать Coding Agent как исполнительное ядро,
  а файловую систему --- как рабочее пространство. Сначала
  рассматривается процесс разработки, а затем на примерах мультимедиа и
  бизнес-систем показывается, как агент пишет программы для решения
  задач, создаёт инструменты и строит новых агентов.
\item
  \textbf{\crossreflink{chap:6}{Глава 6} (Взаимодействие: расширение
  пространств наблюдений и действий)} рассматривает взаимодействие
  агента в двух измерениях --- модальности и времени. Асинхронные
  события, голосовое взаимодействие, Computer Use и робототехника
  помогают понять непрерывное наблюдение, быстрый и медленный пути, а
  также управление выполнением.
\item
  \textbf{\crossreflink{chap:7}{Глава 7} (Оценка AI-агента)} отвечает на
  вопрос, как оценивать работу системы и эффект отдельного изменения.
  Сначала определяются критерии успеха, среда и набор данных, затем
  рассматриваются верификаторы, выставление оценок и статистический
  анализ, а в завершение причины неудач анализируются для поиска
  направлений улучшения.
\item
  \textbf{\crossreflink{chap:8}{Глава 8} (Постобучение моделей)}
  посвящена обновлению возможностей модели с помощью данных и среды,
  причём приоритет отдаётся качеству данных и среды обучения. На примере
  дообучения с учителем, обучения с подкреплением и проектирования
  вознаграждений объясняется, почему обучение по демонстрациям и
  исследование дают разные результаты обобщения и почему изменение
  одного лишь алгоритма часто приносит мало пользы.
\item
  \textbf{\crossreflink{chap:9}{Глава 9} (Непрерывная эволюция
  AI-агента)} рассматривает использование опыта эксплуатации для
  улучшения последующих версий агента. Сравниваются четыре типа
  обновлений --- знаний, инструкций, программ и параметров, --- после
  чего обсуждается управление долгосрочными изменениями посредством
  верификации, выпуска и восстановления.
\item
  \textbf{\crossreflink{chap:10}{Глава 10} (Совместная работа нескольких
  AI-агентов)} рассматривает, когда несколько участников создают
  дополнительную ценность и как распределять между ними задачи,
  организовывать коммуникацию и координацию. Сначала объясняется
  распределение контекста и полномочий управления, затем на примерах
  разбираются механизмы совместной работы и причины неудач, а в
  завершение обсуждаются возможные формы общества и экономики агентов.
\end{itemize}

\hypertarget{ux43aux430ux43a-ux447ux438ux442ux430ux442ux44c-ux44dux442ux443-ux43aux43dux438ux433ux443}{%
\section*{Как читать эту
книгу}\label{ux43aux430ux43a-ux447ux438ux442ux430ux442ux44c-ux44dux442ux443-ux43aux43dux438ux433ux443}}
\addcontentsline{toc}{section}{Как читать эту книгу}

\begin{itemize}
\tightlist
\item
  \textbf{Создание продукта на основе AI-агента с нуля}: рекомендуется
  прочитать всю книгу. Сначала с помощью глав с первой по шестую
  создайте помощника, способного выполнять задачи, а затем в седьмой
  главе проверьте, в каких задачах он чаще ошибается. Последние три
  главы посвящены тому, как на основе этих неудач улучшать модель,
  систему или распределение обязанностей между участниками.
\item
  \textbf{Сначала сформировать целостное представление}: последовательно
  прочитайте первую, вторую, четвёртую и седьмую главы. Проследите один
  вызов инструмента от начала до конца: сначала разберитесь в рабочем
  цикле агента, затем посмотрите, какую информацию получает модель и что
  в действительности делает инструмент, а в конце с помощью оценки
  определите, соответствует ли результат требованиям пользователя. При
  первом чтении раздела второй главы о KV-кэше можно сначала усвоить три
  вывода, приведённые в начале, а к низкоуровневым принципам вернуться
  после проведения экспериментов.
\item
  \textbf{Сосредоточиться на обучении моделей}: сначала с помощью первой
  и второй глав разберитесь в принципах работы агента, а затем читайте
  восьмую главу вместе с методами оценки из седьмой. Это позволит
  связать цели обучения с результатами на реальных задачах.
\item
  \textbf{Сопровождать уже запущенного агента}: выберите в седьмой главе
  набор реальных задач и проведите оценку. Например, если помощник
  сообщил пользователю устаревшую цену, сначала выясните, какую версию
  ценового предложения он получил. Если устарели знания или процесс,
  изучите эту неудачу с помощью девятой главы. Если же и ценовое
  предложение, и инструменты были предоставлены в полном объёме, но
  модель даже после изменения процесса продолжает раз за разом выбирать
  неверный вариант, прочитайте восьмую главу и определите, стоит ли
  готовить данные для обучения.
\end{itemize}

В каждой главе предусмотрены эксперименты и вопросы для размышления.
Эксперименты пронумерованы в формате «Эксперимент X-Y», где X --- номер
главы, а Y --- порядковый номер внутри неё. Сложность обозначается
звёздами: ★ --- начальный уровень, ★★ --- требуется определённый
инженерный опыт, ★★★ --- затрагиваются открытые вопросы или
проектирование сложных систем. Для большинства экспериментов
предоставлен исполняемый код, собранный в открытом репозитории:

\begin{quote}
\textbf{Репозиторий с сопутствующим кодом}:
\texttt{bojieli/ai-agent-book}
\end{quote}

Если вы впервые выполняете эксперименты, можно начать с «Эксперимента
1-1: Ключевая роль контекста». Если вы используете Git, выполните в
терминале:

\begin{Shaded}
\begin{Highlighting}[]
\FunctionTok{git}\NormalTok{ clone https://github.com/bojieli/ai{-}agent{-}book.git}
\BuiltInTok{cd}\NormalTok{ ai{-}agent{-}book}
\end{Highlighting}
\end{Shaded}

Тот же каталог можно получить, выбрав \textbf{Code → Download ZIP} на
странице репозитория и распаковав архив. Найдите номер 1-1 в указателе
экспериментов первой главы (\texttt{chapter1/README.md}) и откройте
\texttt{chapter1/context/README.md}. В дальнейшем проекты и инструкции
по запуску для любого «Эксперимента X-Y» можно найти в файле
\texttt{chapterX/README.md} соответствующей главы. В «руководствах по
воспроизведению» указаны внешние проекты и способы их получения, а
команды установки и описание каталогов собраны в README.

\hypertarget{ux43fux440ux435ux434ux432ux430ux440ux438ux442ux435ux43bux44cux43dux44bux435-ux437ux43dux430ux43dux438ux44f}{%
\section*{Предварительные
знания}\label{ux43fux440ux435ux434ux432ux430ux440ux438ux442ux435ux43bux44cux43dux44bux435-ux437ux43dux430ux43dux438ux44f}}
\addcontentsline{toc}{section}{Предварительные знания}

Книга предназначена для читателей, обладающих базовыми навыками
программирования и желающих понять, как устроены AI-агенты, и научиться
создавать их. Для практических экспериментов потребуются следующие
базовые знания.

\textbf{Обязательно: основа для чтения всей книги}

\begin{itemize}
\tightlist
\item
  \textbf{Программирование на Python}: необходимо знать синтаксис,
  распространённые структуры данных, инструменты управления пакетами,
  такие как pip, а также уметь читать и изменять код экспериментов.
  Большинство экспериментов в книге основано на Python.
\item
  \textbf{Опыт использования LLM}: необходимо иметь опыт работы с
  ChatGPT, Claude или аналогичными продуктами, уметь формулировать
  запросы, читать ответы и задавать уточняющие вопросы.
\item
  \textbf{Основы программной инженерии}: необходимо знать командную
  строку, Git, JSON и REST API, чтобы уметь запускать проекты, проверять
  изменения и читать данные API.
\end{itemize}

Для практических экспериментов автор настоятельно рекомендует выбрать
один из следующих инструментов: Claude Code, Codex, Cursor или TRAE.
Попросите его помочь с отладкой фрагмента кода и посмотрите, какие файлы
он прочитает и какие инструменты вызовет. Тогда при дальнейшем чтении о
цикле выполнения и управлении контекстом у вас уже будет собственный
опыт для сопоставления.

\textbf{Рекомендуется: знания, дополняемые по главам}

\begin{itemize}
\tightlist
\item
  \textbf{Основы машинного обучения (\crossreflink{chap:8}{глава 8})}:
  понимание обучения и инференса, функции потерь, градиентного спуска и
  переобучения поможет разобраться в процессе обновления модели.
\item
  \textbf{Основы математики (\crossreflink{chap:2}{главы 2}, 3, 7 и 8)}:
  знание векторов и матриц поможет понять эмбеддинги и механизм
  внимания, а основы теории вероятностей и статистики --- метрики оценки
  и ожидаемое вознаграждение. В книге основное внимание уделяется
  наглядным объяснениям и практическому применению.
\item
  \textbf{Основы веб-разработки (\crossreflink{chap:6}{глава 6})}:
  знание HTTP, WebSocket и разделения задач между фронтендом и бэкендом
  поможет понять архитектуру асинхронных событий и голосовую
  коммуникацию в реальном времени.
\item
  \textbf{Основы Transformer (\crossreflink{chap:2}{главы 2} и 8)}:
  понимание механизма внимания и структуры модели поможет разобраться в
  вычислениях с контекстом и постобучении. Читателям, желающим изучить
  эту тему системно, рекомендуется книга «Большие модели в иллюстрациях:
  принципы и практика генеративного AI» Джея Аламмара и Мартена
  Гротендорста, переведённая Ли Боцзе и выпущенная издательством
  «Тьюринг». В ней с помощью наглядных схем объясняются Transformer,
  предобучение и дообучение, что дополняет инженерный взгляд,
  представленный в этой книге.
\end{itemize}

За исключением постобучения в \crossreflink{chap:8}{главе 8}, в
остальных главах требуется лишь небольшой объём знаний математики и
машинного обучения, поэтому начать чтение можно, имея базовый опыт
программирования.

\hypertarget{ux431ux43bux430ux433ux43eux434ux430ux440ux43dux43eux441ux442ux438}{%
\section*{Благодарности}\label{ux431ux43bux430ux433ux43eux434ux430ux440ux43dux43eux441ux442ux438}}
\addcontentsline{toc}{section}{Благодарности}

Благодарю редакторов издательства Turing Мэн Гэ и Лю Мэйин за работу над
этой книгой и усилия по организации «Практического интенсива по
AI-агентам»; благодарю Лю Цзюньмина за проведение практического курса по
агентам в Университете Китайской академии наук. Благодарю слушателей и
студентов обоих курсов: ваши вопросы, отзывы об экспериментах и
рекомендации помогали автору постоянно совершенствовать объяснения и
глубже понимать эти концепции.

Благодарю всех коллег из Pine AI. Вызовы реального бизнеса и обсуждения
в команде стали источником многих методов и примеров из этой книги.
Благодаря вашему практическому опыту и идеям эти представления постоянно
совершенствовались.

Благодарю друзей из AI-индустрии. Во время наших многочисленных бесед вы
откровенно высказывали свои замечания, делились разным опытом и
суждениями, помогая автору глубже понять модели и агентов.

Но больше всего автор благодарен своей семье, особенно жене Мэн Цзяин.
Она неизменно поддерживала его работу над этой книгой и высказала
множество ценных замечаний к рукописи.

\hypertarget{ux432ux432ux435ux434ux435ux43dux438ux435-ux432-ai-ux430ux433ux435ux43dux442ux43eux432}{%
\chapter{\texorpdfstring{Введение в
AI-агентов\label{chap:1}}{Введение в AI-агентов}}\label{ux432ux432ux435ux434ux435ux43dux438ux435-ux432-ai-ux430ux433ux435ux43dux442ux43eux432}}

При исследовании темы с помощью Deep Research он многократно выполняет
поиск и читает материалы, а затем составляет отчёт. Manus может
управлять браузером для выполнения онлайн-задач, мобильный помощник
Doubao --- бронировать билеты и отправлять сообщения с телефона, а Pine
AI --- звонить оператору связи от имени пользователя и договариваться о
снижении счёта. Всё это --- примеры применения AI-агентов.

Когда вы поручаете Cursor исправить функцию, он находит её, редактирует
файл, а затем запускает тесты. Если тесты по-прежнему не проходят,
информация об ошибке становится основой для следующего изменения; только
после успешного прохождения тестов у него появляются доказательства,
позволяющие сообщить вам о результате. Последовательность трёх этапов
--- поиска, редактирования и тестирования --- можно заранее заложить в
программу, но то, какую именно строку изменить и к какому этапу
вернуться после неудачи, зависит от информации, полученной во время
выполнения. Такой процесс, при котором действия выбираются в
соответствии с целью, а затем план корректируется с учётом результатов,
происходит и при выполнении таких задач, как исследование темы и
переговоры по телефону.

\hypertarget{ux441ux43eux432ux440ux435ux43cux435ux43dux43dux44bux439-ai-ux430ux433ux435ux43dux442-llm-ux43aux43eux43dux442ux435ux43aux441ux442-ux438ux43dux441ux442ux440ux443ux43cux435ux43dux442ux44b}{%
\section{Современный AI-агент = LLM + контекст +
инструменты}\label{ux441ux43eux432ux440ux435ux43cux435ux43dux43dux44bux439-ai-ux430ux433ux435ux43dux442-llm-ux43aux43eux43dux442ux435ux43aux441ux442-ux438ux43dux441ux442ux440ux443ux43cux435ux43dux442ux44b}}

В этой книге три типа функциональных компонентов современного агента
представлены формулой: \textbf{агент = LLM (большая языковая модель,
Large Language Model) + контекст + инструменты}.

\begin{itemize}
\tightlist
\item
  \textbf{LLM --- мозг агента}, отвечающий за понимание задачи,
  планирование шагов и выбор действий. Предобучение даёт модели языковые
  способности и знания о мире, а постобучение дополнительно формирует
  такие аспекты поведения, как следование инструкциям, рассуждение и
  вызов инструментов. В \crossreflink{chap:8}{главе 8} будут рассмотрены
  методы постобучения, включая дообучение с учителем и обучение с
  подкреплением.
\item
  \textbf{Контекст --- текущее поле принятия решений агента}, содержащее
  информацию, необходимую для конкретного решения, включая инструкции,
  определения инструментов, наблюдения за окружением, память о
  пользователе, знания предметной области, историю сообщений,
  собственное состояние и ход выполнения задачи.
\item
  \textbf{Инструменты --- органы чувств, руки и ноги агента},
  позволяющие получать информацию или выполнять операции через
  интерфейсы. Система инструментов включает определения, протоколы
  вызова и адаптеры и поддерживает как операции наблюдения --- поиск,
  чтение и создание снимков экрана, --- так и запись, выполнение кода,
  делегирование задач подагентам и общение с пользователем.
\end{itemize}

Один и тот же инструмент может выполнять несколько функций. Например,
когда модель читает файл через командную строку, возвращённое содержимое
становится основой для следующего решения; запуск программы через
командную строку также может изменить состояние файлов или процессов.

Предположим, вы только что распаковали посылку и хотите спросить службу
поддержки: «Можно ли вернуть товар после вскрытия упаковки?» Система
может с помощью \textbf{генерации с дополнением извлечёнными данными}
(Retrieval-Augmented Generation, RAG) найти действующие правила возврата
магазина, а затем попросить модель ответить на основании их
положений.\footnote{Patrick Lewis и др. Retrieval-Augmented Generation
  for Knowledge-Intensive NLP Tasks. NeurIPS, 2020.} На примере одного
запроса хорошо видна роль каждого компонента. Исходный текст правил и
индекс хранятся во внешней базе знаний; инструмент извлечения находит
относящиеся к вопросу фрагменты, после чего программа, вызывающая
модель, помещает их вместе с указанием источника в текущий запрос.
Увидев пункт о вскрытой упаковке, модель может разъяснить
соответствующее положение; если правила также различаются по категориям
товаров, сначала потребуется уточнить, что именно вы купили.

В этом запросе о возврате запуск извлечения --- действие агента, а
полученные положения правил --- наблюдение, поступившее из окружения
(Environment). Выбор между ответом и уточняющим вопросом зависит от
содержания правил. Так образуется замкнутый контур, известный из
обучения с подкреплением и теории управления и показанный на
\crossreflink{fig:1-1}{рисунке 1-1}.

\begin{figure}
\hypertarget{fig:1-1}{%
\centering
\includegraphics{images/revised/fig1-1.pdf}
\caption{Рисунок 1-1. Замкнутый контур взаимодействия агента и
Environment, а также внутренняя структура Model--Harness
агента}\label{fig:1-1}
}
\end{figure}

Внешний уровень рисунка показывает границу между агентом и окружением.
Файлы, базы данных, веб-страницы, пользователи, другие агенты, а также
физический или моделируемый мир могут служить окружением, с которым он
взаимодействует. Внутренний уровень показывает разделение
ответственности в производственной реализации: \textbf{Model отвечает за
понимание, рассуждение и принятие решений; Harness находится внутри
границ агента, но за пределами Model, и отвечает за организацию
контекста, интерфейсы инструментов, цикл выполнения и состояние, а также
реализует ограничения, проверку и корректировку.} Harness может
создавать, изолировать или проксировать окружение; содержимое файлов,
состояние процессов и другие элементы окружения изменяются в результате
действий.\footnote{См. «The Anatomy of an Agent Harness» от LangChain и
  «Effective harnesses for long-running agents» от Anthropic. В
  некоторых реализациях песочница также считается частью Harness; на
  \crossreflink{fig:1-1}{рисунке 1-1} окружение показано отдельно, чтобы
  было проще проследить, что именно изменяют действия.}

Таким образом, в книге используются два взаимодополняющих способа
декомпозиции: \textbf{«LLM + контекст + инструменты» описывает
функциональные компоненты, а «Model + Harness» --- производственную
реализацию.} Harness организует контекст и передаёт его Model; Model
принимает решение о вызове инструмента, Harness проверяет и
диспетчеризует вызов, после чего инструмент или окружение выполняет его.
Создание простой системы можно начать с формирования контекста и
диспетчеризации инструментов, а затем постепенно дополнить
ограничениями, проверкой и механизмами корректировки, необходимыми для
производственной эксплуатации.

Если вы знакомы с обучением с подкреплением, следующую таблицу можно
читать через три вопроса: «кто решает, что делать дальше», «на основании
чего принимается решение» и «как оно воплощается в действие». Это
приблизительное функциональное сопоставление.

\begin{longtable}[]{@{}
  >{\raggedright\arraybackslash}p{(\columnwidth - 6\tabcolsep) * \real{0.2145}}
  >{\raggedright\arraybackslash}p{(\columnwidth - 6\tabcolsep) * \real{0.2039}}
  >{\raggedright\arraybackslash}p{(\columnwidth - 6\tabcolsep) * \real{0.2852}}
  >{\raggedright\arraybackslash}p{(\columnwidth - 6\tabcolsep) * \real{0.2964}}@{}}
\toprule\noalign{}
\begin{minipage}[b]{\linewidth}\raggedright
\hspace{0pt}Интуитивный образ
\end{minipage} & \begin{minipage}[b]{\linewidth}\raggedright
\hspace{0pt}Инженерный компонент
\end{minipage} & \begin{minipage}[b]{\linewidth}\raggedright
\hspace{0pt}Соответствующее понятие в обучении с подкреплением
\end{minipage} & \begin{minipage}[b]{\linewidth}\raggedright
\hspace{0pt}Значение
\end{minipage} \\
\midrule\noalign{}
\endhead
\bottomrule\noalign{}
\endlastfoot
\hspace{0pt}\textbf{Мозг} & \hspace{0pt}LLM &
\hspace{0pt}\textbf{Стратегия} (Policy) & \hspace{0pt}Определяет
следующее действие на основании текущей информации \\
\hspace{0pt}\textbf{Текущее поле принятия решений} &
\hspace{0pt}Контекст & \hspace{0pt}\textbf{Наблюдения и история} &
\hspace{0pt}Организует возвращённую окружением информацию и имеющуюся
историю в основу для принятия решения \\
\hspace{0pt}\textbf{Органы чувств, руки и ноги} &
\hspace{0pt}Инструменты и адаптеры & \hspace{0pt}\textbf{Интерфейсы
наблюдений и действий} & \hspace{0pt}Определяют способы чтения
информации из окружения и выполнения действий, а также форматы вызовов и
возвращаемых результатов \\
\end{longtable}

\hypertarget{ux43fux440ux43eux441ux442ux440ux430ux43dux441ux442ux432ux43e-ux43dux430ux431ux43bux44eux434ux435ux43dux438ux439-ux438-ux43fux440ux43eux441ux442ux440ux430ux43dux441ux442ux432ux43e-ux434ux435ux439ux441ux442ux432ux438ux439-ux438ux43dux442ux435ux440ux444ux435ux439ux441-ux43cux435ux436ux434ux443-ux43cux43eux434ux435ux43bux44cux44e-ux438-ux43cux438ux440ux43eux43c}{%
\subsection{Пространство наблюдений и пространство действий: интерфейс
между моделью и
миром}\label{ux43fux440ux43eux441ux442ux440ux430ux43dux441ux442ux432ux43e-ux43dux430ux431ux43bux44eux434ux435ux43dux438ux439-ux438-ux43fux440ux43eux441ux442ux440ux430ux43dux441ux442ux432ux43e-ux434ux435ux439ux441ux442ux432ux438ux439-ux438ux43dux442ux435ux440ux444ux435ux439ux441-ux43cux435ux436ux434ux443-ux43cux43eux434ux435ux43bux44cux44e-ux438-ux43cux438ux440ux43eux43c}}

\textbf{Пространство наблюдений описывает, какие наблюдения может
получать AI-агент, а пространство действий --- какие действия он может
выполнять.} В рамках одного вызова конкретная информация передаётся
через контекст, а доступные инструменты предоставляют интерфейс для
взаимодействия со средой. Наблюдения могут поступать не только от
инструментов, но и из пользовательского ввода, уведомлений о событиях и
других источников; диапазон действий совместно ограничивают параметры
вызовов инструментов, права доступа и среда выполнения.

Представим два разных сбоя. В первом случае модель способна объяснить
каждую статью расходов в счёте, но не может получить счёт за текущий
месяц. Даже после того как пользователь разрешит чтение, счёт ещё нужно
найти и передать модели --- только тогда она сможет определить, какие
расходы можно оспорить. Во втором случае пользователь просит ассистента
сменить тарифный план. Ассистент уже нашёл подходящий вариант, но
интерфейс учётной записи позволяет только просматривать данные. Даже
если ещё раз прочитать счёт, изменить тариф не получится: для этого
нужны соответствующие права на выполнение операций. Задача Harness ---
поместить прочитанные материалы в контекст, принять выбранное моделью
действие, проверить параметры и права, а затем направить его на
выполнение.

Поэтому \textbf{при неизменной базовой модели настройка пространства
наблюдений и пространства действий --- важный инженерный способ повысить
качество выполнения задач.} Если модель упускает важный факт,
соответствующие данные можно включить в контекст. Если модель часто
выбирает не тот инструмент из нескольких похожих, дублирующие интерфейсы
можно объединить либо загружать только инструменты, необходимые для
текущей задачи.

При просмотре видео возникает и обратная проблема: изображение доступно
постоянно, но изменения, требующие повторной оценки, происходят редко.
Предположим, ассистент ждёт загрузки файла. Если вызывать основную
модель при каждом движении индикатора прогресса, это приведёт к
множеству лишних вычислений. Если же проверять состояние лишь раз в
десять секунд, можно пропустить кратковременно появившееся сообщение об
ошибке. \textbf{По мнению автора, разрежение пространства наблюдений ---
одна из фундаментальных проблем проектирования агентов.} Система может
непрерывно наблюдать за средой, но передавать основной модели вместе с
необходимым контекстом только те изменения, которые способны повлиять на
следующее действие.

\textbf{Manus: сочетание исследования, программирования и управления
компьютером в рамках одной задачи.} Deep Research хорошо справляется с
поиском и чтением, агент для программирования может писать и выполнять
код, а агент для управления компьютером --- взаимодействовать с
интерфейсом. Manus объединяет браузер, файловую систему, выполнение кода
и командную строку в облачной песочнице, позволяя одному агенту сначала
искать материалы, затем обрабатывать файлы и запускать код, а в конце
выполнять действия через браузер. Браузер и чтение файлов расширяют
объём доступной информации, а выполнение кода и команд --- набор
действий, которые можно комбинировать. \textbf{По мнению автора, такое
сочетание стало важным прорывом для универсальных агентов: один и тот же
агент может пересекать прежде разделённые границы продуктов, объединяя
получение материалов, обработку данных и выполнение действий в единую
задачу.}

\textbf{OpenClaw: расширение интерфейса до цифровой жизни пользователя.}
OpenClaw получает задачи и возвращает результаты через такие каналы
обмена сообщениями, как WhatsApp, Telegram, Slack, Discord и iMessage, а
работающий на устройстве пользователя Gateway связывает каналы обмена
сообщениями, модель и инструменты. После настройки соответствующих
инструментов и получения разрешений агент может обрабатывать локальные
файлы и подключаться к таким облачным сервисам, как Google Drive и
Notion. Благодаря этому данные пользователя, распределённые между
учётными записями и устройствами, можно включить в процесс выполнения
одной задачи. Ранние версии Manus работали преимущественно в облачной
песочнице, поэтому пользователю приходилось загружать файлы. Появившиеся
позднее коннектор Google Drive и функция My Computer для настольных
компьютеров дополнительно расширили доступ к облачным файлам и локальным
ресурсам. Эта эволюция показывает, что возможности продукта часто растут
благодаря расширению интерфейсов наблюдений и действий\footnote{Структура
  продуктов описана в официальных материалах Manus о Sandbox, Google
  Drive Connector и My Computer, а также в официальном репозитории
  OpenClaw и документации по его инструментам. Конкретный набор
  доступных ресурсов определяется коннекторами, конфигурацией
  инструментов и разрешениями пользователя.}.

Задачи у разных продуктов различаются, но каждый из них можно
рассматривать с трёх сторон: основания для принятия решений, доступные
операции и стратегия выполнения.

\begin{longtable}[]{@{}
  >{\raggedright\arraybackslash}p{(\columnwidth - 6\tabcolsep) * \real{0.2414}}
  >{\raggedright\arraybackslash}p{(\columnwidth - 6\tabcolsep) * \real{0.2283}}
  >{\raggedright\arraybackslash}p{(\columnwidth - 6\tabcolsep) * \real{0.2080}}
  >{\raggedright\arraybackslash}p{(\columnwidth - 6\tabcolsep) * \real{0.3223}}@{}}
\toprule\noalign{}
\begin{minipage}[b]{\linewidth}\raggedright
\hspace{0pt}Продукт на основе агентов
\end{minipage} & \begin{minipage}[b]{\linewidth}\raggedright
\hspace{0pt}Основная информация в контексте
\end{minipage} & \begin{minipage}[b]{\linewidth}\raggedright
\hspace{0pt}Доступные инструменты и операции
\end{minipage} & \begin{minipage}[b]{\linewidth}\raggedright
\hspace{0pt}Стратегия выполнения
\end{minipage} \\
\midrule\noalign{}
\endhead
\bottomrule\noalign{}
\endlastfoot
\hspace{0pt}\textbf{Cursor и другие агенты для программирования} &
\hspace{0pt}Требования, фрагменты кода, списки каталогов, вывод
терминала & \hspace{0pt}Поиск по коду, чтение и запись файлов,
выполнение команд & \hspace{0pt}Понять требования→найти связанный
код→отредактировать код→проверить тестами→найти и исправить ошибки \\
\hspace{0pt}\textbf{Deep Research и другие поисковые агенты} &
\hspace{0pt}Результаты поиска, содержимое
веб-\allowbreak{}\hspace{0pt}страниц, аннотации и цитаты из научных
работ & \hspace{0pt}Поисковые запросы, чтение
веб-\allowbreak{}\hspace{0pt}страниц, создание отчётов &
\hspace{0pt}Корректировать направление поиска с учётом уже полученной
информации, постепенно объединяя её в полный отчёт \\
\hspace{0pt}\textbf{Browser Use и другие агенты для управления
компьютером} & \hspace{0pt}Снимки экрана, DOM или дерево специальных
возможностей, результаты операций & \hspace{0pt}Нажатия, ввод,
прокрутка, снимки экрана, выполнение кода & \hspace{0pt}Прочитать
интерфейс→определить целевой элемент→выполнить действие→проверить
результат \\
\hspace{0pt}\textbf{Doubao и другие мобильные
ассистенты-\allowbreak{}\hspace{0pt}агенты} & \hspace{0pt}Снимки экрана
телефона, состояние интерфейса приложения, обратная связь от системы &
\hspace{0pt}Нажатия, жесты прокрутки, ввод, открытие приложений &
\hspace{0pt}Понять задачу→найти нужное приложение→выполнить
действие→подтвердить завершение \\
\hspace{0pt}\textbf{Pine AI, Meta Muse и другие персональные агенты для
выполнения поручений} & \hspace{0pt}Полученные с разрешения записи
учётных записей, счета, сведения о поставщиках услуг &
\hspace{0pt}Телефонные звонки, электронные письма, заполнение форм,
уточнения у пользователя & \hspace{0pt}Собрать информацию→разработать
стратегию переговоров→связаться с поставщиком услуг→провести
переговоры→сообщить результат \\
\end{longtable}

\textbf{Как поручение превращается в исполнимую задачу?} Работая в Pine
над персональным ассистентом для выполнения поручений, автор заметил,
что пользователи часто формулируют только цель, например: «Забронируй
мне авиабилет для командировки, желательно не дороже 1200 юаней».
Сначала ассистенту нужно выяснить даты, пункт назначения и главные
требования к поездке, затем собрать недостающие сведения и согласовать с
пользователем способ выполнения поручения и условия, по которым допустим
компромисс. Только после этого можно обращаться к поставщику услуг.
Новая информация, полученная в ходе выполнения поручения, может вернуть
задачу на предыдущий этап.

Предположим, сотрудник службы поддержки нашёл билет по минимальной цене
1250 юаней. Модель будет принимать решение с учётом общих знаний,
особенностей этой поездки и запомненных предпочтений пользователя. Если
пользователь находится в командировке и обычно готов немного доплатить
за удобный маршрут, Pine может сразу забронировать билет. Если же речь
идёт о семейной поездке, а пользователь ранее явно подчёркивал важность
общих расходов, ассистент может сначала сообщить ему о превышении
бюджета и попросить сотрудника службы поддержки поискать другие рейсы.
Если впоследствии пользователь добавит: «В следующий раз, если билет
окажется дороже бюджета, сначала спроси меня», эта обратная связь даст
для следующего решения более конкретное основание, чем общий опыт.

Рассмотрим другой сценарий: сотрудник банковской службы поддержки
предлагает отправить документ по домашнему адресу, но ассистент помнит,
что пользователь находится в зарубежной поездке. На этот раз необходимо
изменить способ доставки: во время разговора можно сразу сообщить о
невозможности получить письмо и попытаться договориться о доставке по
электронной почте. Теперь рассмотрим уменьшение суммы по счёту. Если
сотрудник службы поддержки устно согласился уменьшить сумму, ассистент
может точно сообщить об обещании компании. Но чтобы подтвердить, что
сумму действительно уменьшили, понадобится новый счёт как очередное
наблюдение. Пространство действий ассистента может включать реальные
телефонные звонки, однако его выводы по-прежнему ограничены
пространством наблюдений.

С точки зрения поручений в этой книге рассматриваются три типа
приложений: выполнение бытовых поручений, совместная офисная работа и
научные исследования. Один и тот же ассистент может выполнять задачи
разных типов; различия заключаются в том, кто предоставляет ему
полномочия, что он может изменить и кто в итоге принимает результат.
Если в офисе коллега вернул результаты проверки, можно ли закрывать
рабочую заявку? Предположим, руководитель проекта упомянул ассистента в
комментарии к заявке Multica и попросил проверить расхождение в
бюджетной таблице. Этот комментарий запустит проверку, но заявка
останется закреплена за прежним ответственным, а её статус не изменится.
Ассистент вернёт две суммы и укажет их источники, тем самым ответив на
вопрос коллеги. Однако решение о том, можно ли отправлять весь бюджет,
должен принять ответственный с учётом остальных незавершённых
задач.\footnote{Multica: упоминание агента в комментарии к рабочей
  заявке. В десятой главе рассматривается, как в совместной работе
  различать временную помощь и передачу ответственности.}

К научной задаче можно применить тот же вопрос: на каком основании
считать, что достигнут прогресс? Когда AlphaEvolve от Google DeepMind
ищет алгоритмы умножения матриц, модель предлагает программы-кандидаты,
а оценщик запускает их и проверяет правильность вычисления произведения
и количество выполненных умножений. Затем поиск может продолжить
улучшать программы с хорошими результатами. Исследовательская группа
сообщила, что один из алгоритмов позволяет перемножить две комплексные
матрицы размером \(4\times 4\) с помощью 48 скалярных умножений. У
утверждения «потребовалось на несколько умножений меньше» есть чётко
определённый объект подсчёта, а алгоритм при этом должен вычислять
правильное произведение матриц.\footnote{Google DeepMind, описание
  исследования AlphaEvolve, май 2025 года. Под количеством умножений
  здесь понимается число скалярных умножений в алгоритме; фактическое
  время выполнения также зависит от реализации и оборудования.}

Пространства действий этих систем носят \textbf{открытый} характер.
Невозможно заранее полностью перечислить поисковые фразы, вводимый
текст, создаваемый код и комбинации операций, однако каждое действие
по-прежнему ограничено инструментами, параметрами, правами доступа и
средой выполнения. Именно поэтому исполнимый план должен допускать
изменение следующего шага с учётом новых наблюдений.

\hypertarget{ux438ux43dux441ux442ux440ux443ux43cux435ux43dux442ux44b-ux43eux440ux433ux430ux43dux44b-ux447ux443ux432ux441ux442ux432-ux440ux443ux43aux438-ux438-ux43dux43eux433ux438-ai-ux430ux433ux435ux43dux442ux430}{%
\subsection{Инструменты: органы чувств, руки и ноги
AI-агента}\label{ux438ux43dux441ux442ux440ux443ux43cux435ux43dux442ux44b-ux43eux440ux433ux430ux43dux44b-ux447ux443ux432ux441ux442ux432-ux440ux443ux43aux438-ux438-ux43dux43eux433ux438-ai-ux430ux433ux435ux43dux442ux430}}

Инструменты воплощают решения модели в чтении внешней среды или
воздействии на неё. Инструменты поиска и чтения возвращают наблюдения, а
инструменты записи файлов, выполнения команд и другие подобные средства
реализуют действия. По полученному результату модель узнаёт об эффекте
операции, а затем решает, продолжить работу, повторить попытку или
завершить задачу.

По основному назначению в этой книге инструменты делятся на пять
категорий. Один интерфейс может одновременно выполнять несколько
функций: например, голосовой звонок позволяет как получать информацию,
так и передавать запросы.

\textbf{Инструменты восприятия} отвечают за получение информации:
поисковые системы предоставляют материалы из интернета, инструменты
чтения файлов открывают локальные документы, а API и интерфейсы баз
данных обеспечивают доступ к внешним сервисам и корпоративным данным.

\textbf{Инструменты выполнения} отвечают за совершение операций:
запускают код, изменяют файлы, выполняют системные команды или вызывают
внешние API, воплощая решения в среде.

\textbf{Инструменты совместной работы} поддерживают разделение и
координацию работы: делегируют специализированные задачи дочерним
агентам, собирают результаты и координируют работу нескольких агентов.
Если для принятия решения необходимо участие человека, система также
может запросить подтверждение через интерфейс общения с пользователем.

\textbf{Инструменты запуска по событиям} подключают внешние события к
процессу выполнения. Новое письмо, наступление заданного времени или
обратный вызов Webhook могут запустить задачу. Адаптер событий
преобразует такие входные данные в наблюдения, которые способен
обработать агент, а механизм планирования запускает или возобновляет
выполнение. Правила запуска могут задаваться разработчиком либо
регистрироваться посредством разрешённых вызовов инструментов.

\textbf{Инструменты общения с пользователем} передают пользователю
информацию и собирают обратную связь, например отправляют текстовые
сообщения и электронные письма или инициируют голосовые звонки. Они
используются для отчётов о ходе работы и запросов подтверждения, а также
могут сами напоминать пользователю о необходимых действиях.

В четвёртой главе подробно рассматривается проектирование этих пяти
категорий инструментов. От определения интерфейса зависит, сможет ли
модель правильно выбрать инструмент и заполнить параметры; обработка
ошибок определяет, можно ли продолжить работу после сбоя; контроль прав
ограничивает область воздействия операций. MCP (Model Context Protocol,
протокол контекста модели) предоставляет единый способ обнаружения и
вызова инструментов, сокращая объём адаптации между различными
системами.

\textbf{Вызов инструмента} (Tool Calling) --- способность модели
формировать запрос на выполнение операции в согласованном формате.
Распространённый вызов функций (Function Calling) выражает запрос через
имя инструмента и структурированные параметры. В этой книге везде
используется термин «вызов инструмента».

Рассмотрим запрос погоды в Пекине. Модель может указать название города,
но фактическая температура появится в контексте только после того, как
Harness выполнит запрос. Ниже показаны четыре шага обмена сообщениями
при этом вызове:

\begin{Shaded}
\begin{Highlighting}[]
\NormalTok{Шаг 1: объявление инструмента          Шаг 2: модель решает выполнить вызов}
\NormalTok{tools: [\{                              assistant: \{}
\NormalTok{  name: "get\_weather",                   tool\_calls: [\{}
\NormalTok{  parameters: \{                            id: "call\_1",}
\NormalTok{    city: "string"                         function: "get\_weather",}
\NormalTok{  \}                                        arguments: \{city: "Пекин"\}}
\NormalTok{\}]                                       \}]}
\NormalTok{                                       \}}

\NormalTok{Шаг 3: результат добавляется в контекст Шаг 4: модель отвечает по результату}
\NormalTok{tool: \{                                assistant: \{}
\NormalTok{  tool\_call\_id: "call\_1",                content: "Сегодня в Пекине 28°C, ясно."}
\NormalTok{  content: \textquotesingle{}\{"temp":28,"sky":"ясно"\}\textquotesingle{}   \}}
\NormalTok{\}}
\end{Highlighting}
\end{Shaded}

Здесь \texttt{call\_1} связывает возвращённое инструментом значение с
запросом на вызов. Разработчик описывает способ вызова инструмента,
реализует функцию запроса погоды и передаёт результат запроса обратно
модели. Модель решает, следует ли выполнять вызов, какой инструмент
выбрать и какие параметры ему передать. Во второй главе структура
сообщений API рассматривается подробнее.

При проектировании инструмента сначала можно предоставить ровно те
возможности, которых достаточно для задачи. Для четырёх арифметических
действий подойдёт калькулятор с понятными параметрами. Если задача
расширяется до чтения таблиц, очистки пропущенных значений, вычисления
статистик и построения графиков, ограниченный интерпретатор Python
позволит удобнее комбинировать эти шаги. С расширением универсальных
возможностей необходимо усиливать изоляцию и управление ресурсами: код
должен выполняться в песочнице, сетевой доступ по умолчанию должен быть
отключён, доступ к файлам --- ограничен разрешённым рабочим каталогом, а
для времени выполнения, CPU, памяти и размера вывода следует установить
предельные значения.

Для записи отдельного процесса выполнения можно использовать единый
инструмент журналирования. Для задач, продолжающихся несколько часов или
даже дней, лучше использовать управляемый виртуальный рабочий каталог, в
котором сохраняются план, промежуточные результаты, журналы выполнения и
итоговые материалы. Тогда агент сможет продолжать работу между
отдельными запусками. Для каталога следует ограничить доступные для
чтения и записи пути, объём и типы файлов, а также проверять, не выходит
ли путь за разрешённые границы.

Платежи, удаление данных, отправка электронной почты, развёртывание в
промышленной среде и другие операции сопряжены с повышенным риском или
регулируются явными бизнес-правилами. Их целесообразно оформлять как
специализированные инструменты с понятными параметрами, ограниченными
правами и полной аудитируемостью, при необходимости добавляя
предварительный просмотр и подтверждение человеком.
\textbf{Универсальные базовые возможности поддерживают комбинирование и
исследование, а специализированные инструменты обеспечивают соблюдение
ограничений и бизнес-правил при выполнении высокорисковых операций.}

\hypertarget{llm-ux43cux43eux437ux433-ai-ux430ux433ux435ux43dux442ux430}{%
\subsection{LLM: мозг
AI-агента}\label{llm-ux43cux43eux437ux433-ai-ux430ux433ux435ux43dux442ux430}}

LLM --- ядро принятия решений агентом. Получив запрос, она сначала
определяет цель и ограничения из формулировки пользователя, а затем
разбивает сложную задачу на выполнимые шаги. В процессе выполнения ей
также необходимо решать, что делать дальше, следует ли вызвать
инструмент, какой инструмент выбрать и какими значениями заполнить
параметры. Эти возможности опираются на знания и языковые способности,
накопленные в ходе предварительного обучения, а постобучение формирует
на их основе подходящее для решения задач поведение при принятии
решений. Такие решения необходимы как рабочим процессам, так и
автономным агентам, а практические действия выполняются посредством
Harness, инструментов и среды.

Планирование и моделирование перед выполнением действий помогают агенту
сравнивать варианты, прогнозировать результаты и сокращать число
бесполезных попыток. Например, модель может проверять этапы вычислений с
помощью математических закономерностей, анализировать сбои на основе
причинно-следственных связей, а затем выстраивать порядок операций с
учётом зависимостей. Сначала модель сравнивает различные варианты в ходе
рассуждений, а выбрав действие, запрашивает его выполнение у
инструмента. Традиционные системы обучения с подкреплением также могут
организовывать исследование с помощью выученной стратегии, функции
ценности, модели среды или алгоритма планирования. Особенность агентов
на основе LLM заключается в том, что ещё до начала конкретной задачи они
благодаря масштабному предварительному обучению получают обширные знания
и шаблоны решения задач, а постобучение дополнительно совершенствует
способы применения этих возможностей.

\hypertarget{ux432ux441ux442ux440ux43eux435ux43dux43dux44bux435-ux432-ux43cux43eux434ux435ux43bux44c-ux432ux44bux437ux43eux432ux44b-ux438ux43dux441ux442ux440ux443ux43cux435ux43dux442ux43eux432}{%
\subsubsection{Встроенные в модель вызовы
инструментов}\label{ux432ux441ux442ux440ux43eux435ux43dux43dux44bux435-ux432-ux43cux43eux434ux435ux43bux44c-ux432ux44bux437ux43eux432ux44b-ux438ux43dux441ux442ux440ux443ux43cux435ux43dux442ux43eux432}}

Выражение «модель как агент» (Model as Agent) описывает отраслевую
тенденцию: всё больше решений, связанных с планированием и вызовами
инструментов, принимает сама модель. Благодаря постобучению, особенно
обучению с подкреплением, модель может научиться определять, когда
вызывать инструмент, какой именно инструмент выбрать, какие параметры
передать и как продолжить работу после получения результата. Разработчик
задаёт цель, интерфейсы и ограничения, а модель организует конкретные
шаги с учётом хода выполнения задачи.

Для самостоятельного принятия решений модели необходим Harness.
Изначально английское слово harness обозначало упряжь и другие элементы
конской сбруи, передающие силу лошади повозке и позволяющие управлять
направлением движения. В агенте Harness организует возможности модели в
непрерывно выполняемый процесс решения задачи и отвечает за управление
контекстом, интерфейсы инструментов, ограничения, проверку и
исправление. Во второй половине этой главы мы последовательно рассмотрим
все пять функций.

Чем больше операций модель может организовывать самостоятельно, тем шире
возможные последствия ошибок. Поэтому модель и Harness необходимо
оптимизировать совместно: модель повышает качество планирования и выбора
инструментов, а Harness управляет правами доступа, проверяет результаты
и обрабатывает сбои. Производители моделей могут одновременно
совершенствовать обе составляющие, и такая согласованная итеративная
разработка становится важным инженерным преимуществом.

\begin{quote}
\textbf{Эволюция возможностей модели и Harness}

В работе «Горький урок» (The Bitter Lesson) Рич Саттон рассматривает
долгосрочную тенденцию в исследованиях AI: методы поиска и обучения,
способные использовать всё больше вычислительных ресурсов, постепенно
превосходят подходы, в значительной мере опирающиеся на спроектированные
вручную предметно-ориентированные решения\footnote{Sutton, Rich. The
  Bitter Lesson, 2019.}. В русле этой тенденции решения о выборе
инструментов, порядке шагов и других аспектах, прежде возлагавшиеся на
внешние процессы, могут постепенно усваиваться моделью в ходе обучения.

Позиция этой книги такова: \textbf{мы согласны с направлением, но
придерживаемся прагматичного темпа.} Модель будет последовательно
принимать на себя часть функций Harness по принятию решений, однако
каждый цикл обучения и развёртывания требует времени, тогда как
бизнес-правила и предпочтения пользователей постоянно меняются. Кроме
того, все ограничения реальных бизнес-сценариев трудно за один раз
зафиксировать в параметрах модели. Поэтому те этапы, с которыми модель
пока не справляется достаточно надёжно, сначала дополняет Harness.
\textbf{Текущая граница возможностей модели определяет текущую ценность
Harness.} Это разделение обязанностей меняется вместе с возможностями
модели: осваивая очередной класс решений, модель позволяет Harness
сократить соответствующую оркестрацию и переключиться на новые
сложности. Обеспечение соблюдения прав доступа, сохранение состояния и
проверка результатов по-прежнему остаются задачами системы выполнения.
\end{quote}

\hypertarget{ux43cux435ux445ux430ux43dux438ux437ux43cux44b-ux43eux431ux443ux447ux435ux43dux438ux44f-ux430ux433ux435ux43dux442ux430-ux43eux442-ux430ux434ux430ux43fux442ux430ux446ux438ux438-ux43a-ux43aux43eux43dux442ux435ux43aux441ux442ux443-ux434ux43e-ux443ux441ux442ux43eux439ux447ux438ux432ux43eux433ux43e-ux43eux431ux43dux43eux432ux43bux435ux43dux438ux44f}{%
\subsubsection{Механизмы обучения агента: от адаптации к контексту до
устойчивого
обновления}\label{ux43cux435ux445ux430ux43dux438ux437ux43cux44b-ux43eux431ux443ux447ux435ux43dux438ux44f-ux430ux433ux435ux43dux442ux430-ux43eux442-ux430ux434ux430ux43fux442ux430ux446ux438ux438-ux43a-ux43aux43eux43dux442ux435ux43aux441ux442ux443-ux434ux43e-ux443ux441ux442ux43eux439ux447ux438ux432ux43eux433ux43e-ux43eux431ux43dux43eux432ux43bux435ux43dux438ux44f}}

После ошибки инструмента агент в этот раз научился искать файл по
другому пути. Но сохранится ли этот урок завтра, в новом сеансе? Если
проследить по \crossreflink{fig:1-2}{рисунку 1-2}, где хранится этот
опыт, можно различить три вида обновлений.

\begin{figure}
\hypertarget{fig:1-2}{%
\centering
\includegraphics{images/revised/fig1-2.pdf}
\caption{Рисунок 1-2. Три уровня обновления возможностей
агента}\label{fig:1-2}
}
\end{figure}

\textbf{Адаптация к контексту} происходит в процессе выполнения задачи.
Прочитав примеры, состояние задачи или результаты поиска, модель может
скорректировать последующее поведение. Например, если инструмент
сообщает, что файл не существует, модель может заново найти путь и
продолжить работу. Этот подход обеспечивает быструю реакцию при
небольших затратах, но его эффективность зависит от ёмкости контекстного
окна и способа организации информации. Чтобы использовать накопленный
опыт в последующих задачах, соответствующую информацию необходимо
сохранить и снова предоставить модели. Во второй главе будут рассмотрены
способы организации контекста.

\textbf{Внешние артефакты} (artifact) позволяют сохранять опыт между
задачами. Факты и опыт записываются в документы базы знаний, стратегии,
которые можно выразить на естественном языке, --- в промпты или Skill, а
детерминированные процессы и ограничения --- в программы или логику
Harness. Такие артефакты удобно аудировать и редактировать; во время
выполнения они участвуют в принятии решений и осуществлении действий
через контекст или интерфейсы инструментов. В
\crossreflink{chap:3}{главах 3}--5 рассматриваются соответствующие
основы знаний и программ, а в \crossreflink{chap:9}{главе 9} обсуждается
создание, проверка и публикация обновлений на основе оценённых
траекторий выполнения.

\textbf{Параметры модели} обновляются посредством обучения. Такие
возможности, как понимание медицинских изображений, воспроизведение
языкового стиля и неявные стратегии принятия решений, связаны с
множеством закономерностей, которые трудно по отдельности сформулировать
в виде правил. Модель может усвоить их из данных и обратной связи в ходе
постобучения. Обновление параметров требует затрат на обучение, оценку и
развёртывание, а его цель --- позволить модели применять усвоенные
возможности и в новых задачах. Соответствующие методы подробно
рассматриваются в \crossreflink{chap:8}{главе 8}.

Один и тот же сбой можно сначала использовать для исправления текущей
операции, а затем оформить как проектную документацию, доступную
следующему агенту. Если ошибки одного типа неоднократно возникают во
множестве проектов, соответствующие случаи также можно использовать для
обучения.

\hypertarget{ux43aux43eux43dux442ux435ux43aux441ux442-ux442ux435ux43aux443ux449ux435ux435-ux43fux43eux43bux435-ux43fux440ux438ux43dux44fux442ux438ux44f-ux440ux435ux448ux435ux43dux438ux439-ai-ux430ux433ux435ux43dux442ux430}{%
\subsection{Контекст: текущее поле принятия решений
AI-агента}\label{ux43aux43eux43dux442ux435ux43aux441ux442-ux442ux435ux43aux443ux449ux435ux435-ux43fux43eux43bux435-ux43fux440ux438ux43dux44fux442ux438ux44f-ux440ux435ux448ux435ux43dux438ux439-ai-ux430ux433ux435ux43dux442ux430}}

Принимая решение, человек должен одновременно учитывать описание задачи,
справочники, историю общения и свежие данные. Агенту тоже нужен такой
набор информации --- это и есть контекст. Объём информации, который
можно передать за один вызов LLM, ограничен контекстным окном, а Harness
отвечает за отбор и организацию его содержимого. С точки зрения API
базовая реализация в этой книге включает пять типов информации:

\begin{itemize}
\tightlist
\item
  \textbf{Системный промпт} (System Prompt): создаётся разработчиком или
  Harness и определяет роль агента, границы задачи и правила поведения.
  Промпт-инжиниринг формирует рабочий стиль агента посредством изменения
  этих инструкций. Системный промпт может оставаться относительно
  стабильным в течение сеанса, а конфигурация продукта может добавлять в
  него сохраняемую между сеансами память о пользователе, например
  предпочтения, историю поведения и фоновые сведения, либо внедрять
  текущее состояние окружения. Фактические права доступа обеспечиваются
  Harness и средствами контроля доступа в окружении. Память о
  пользователе будет рассмотрена в \crossreflink{chap:3}{главе 3}.
\item
  \textbf{Определения инструментов} (Tool Definitions): объявляют имена,
  функции и форматы параметров доступных инструментов, помогая модели
  выбирать действия и формировать запросы. В базовой реализации книги
  определения инструментов вместе с системным промптом располагаются в
  начале запроса и образуют относительно стабильный префикс. Если
  инструментов много, полную схему также можно загружать по требованию,
  сохраняя уже существующий префикс; этот подход будет подробно
  рассмотрен в \crossreflink{chap:2}{главах 2} и 4.
\item
  \textbf{Сообщения пользователя} (User Messages): содержат задачу
  пользователя и дополнительные входные данные. К ним также могут
  прилагаться внешние знания, полученные посредством RAG (генерации с
  дополнением извлечёнными данными, Retrieval-Augmented Generation),
  например информация, появившаяся после даты отсечения обучающих
  данных, или закрытые материалы предметной области. Извлечение и
  организация знаний будут рассмотрены в \crossreflink{chap:3}{главе 3}.
\item
  \textbf{Ответы модели} (Assistant Messages): включают предназначенное
  пользователю текстовое \texttt{content}, запросы на вызов инструментов
  \texttt{tool\_calls}, а также предоставляемую сервисом информацию о
  рассуждениях. В книге для такой информации используется сокращённое
  обозначение \texttt{reasoning}; разные сервисы могут возвращать
  содержимое рассуждений, их сводку или непрозрачное состояние для
  продолжения. При вызове инструмента ответ может содержать информацию о
  рассуждениях и запрос на вызов, а при завершении задачи --- информацию
  о рассуждениях и окончательный текст. Конкретные поля и их видимость
  определяются сервисом модели.
\item
  \textbf{Результаты выполнения инструментов} (Tool Results): содержат
  возвращённые после выполнения значения, ошибки или состояние. Они
  добавляются в историю сообщений и служат основанием для следующего
  решения. По ним модель определяет, успешно ли выполнена операция,
  требуется ли скорректировать действия, и сокращает число повторных
  попыток в рамках текущей задачи.
\end{itemize}

При каждом вызове модели Harness должен передавать ей относительно
стабильный префикс --- системный промпт и определения инструментов ---
вместе с постоянно растущей историей сообщений. Благодаря этому, даже
если инструмент конвертации валют вернул результат в предыдущем раунде,
модель по-прежнему может прочитать его при вычислении общей суммы в
текущем.

Роль разных типов информации в принятии решений можно проверить с
помощью \textbf{абляционного эксперимента} (Ablation Study):
зафиксировать задачу и условия выполнения, при каждом запуске удалять
один тип информации и сравнивать изменения в поведении. В эксперименте
1-1 сохраняются системный промпт и текущая пользовательская задача,
чтобы зафиксировать роль, ограничения и требования задачи, а абляции по
отдельности подвергаются определения инструментов, результаты
инструментов, предыдущая информация о рассуждениях и история сообщений.

\begin{quote}
\textbf{Эксперимент 1-1 ★★: ключевая роль контекста}

Предложим агенту суммировать квартальные доходы компании: 2,5 миллиона
долларов США в первом квартале, 2,1 миллиона евро во втором, 1,8
миллиона фунтов стерлингов в третьем и 380 миллионов иен в четвёртом. По
условиям задачи доходы необходимо привести к долларам с помощью
инструмента конвертации валют, а затем вычислить годовой и средний
квартальный доход с помощью вычислительного инструмента, округлив
результаты до двух знаков после запятой. Таким образом, модели
необходимо не только выбирать инструменты, но и помнить уже выполненные
шаги и возвращённые значения.

Эксперимент включает полную базовую конфигурацию и четыре группы с
абляциями; каждая группа выполняется не более пяти раундов. При удалении
предыдущей информации о рассуждениях модели по-прежнему разрешается
рассуждать заново в каждом раунде. При удалении определений инструментов
запрос больше не содержит доступных инструментов. При удалении
результатов инструментов записи о вызовах сохраняются, но возвращённое
содержимое скрывается. При удалении истории сообщений остаются только
системный промпт, определения инструментов и текущая пользовательская
задача. Сначала попробуйте предсказать: какой этап вычислений можно
выполнить, зная лишь, что «курс иены уже запрашивали», а какой --- зная,
сколько долларов получилось после конвертации 380 миллионов иен? На
\crossreflink{fig:1-3}{рисунке 1-3} показаны результаты пяти групп.

\begin{figure}
\hypertarget{fig:1-3}{%
\centering
\includegraphics{images/revised/fig1-3.pdf}
\caption{Рисунок 1-3. Полная базовая конфигурация и четыре абляции в
эксперименте 1-1}\label{fig:1-3}
}
\end{figure}

При запуске Kimi K3 25 августа 2026 года \textbf{полная базовая
конфигурация} выполнила задачу за три раунда и четыре вызова
инструментов, получив годовой доход 9 602 895,73 доллара и средний
квартальный доход 2 400 723,93 доллара. После \textbf{удаления
предыдущей информации о рассуждениях} правильный результат также был
получен за то же число раундов и вызовов: результаты конвертации
сохранились в записях инструментов, поэтому модель смогла по ним
восстановить информацию, необходимую для следующего вычисления.

После \textbf{удаления определений инструментов} модель не инициировала
ни одного вызова и сообщила об отсутствии инструмента конвертации валют.
После \textbf{удаления результатов инструментов} модель инициировала
девять вызовов, но так и не завершила задачу в пределах пяти раундов:
существующие записи о вызовах не сообщали ей, сколько долларов
получилось при конвертации. После \textbf{удаления истории сообщений}
модель инициировала в общей сложности пятнадцать вызовов, многократно
повторяя конвертацию валют, и также исчерпала бюджет из пяти раундов.
Ранее выполненная работа не попадала в контекст следующего раунда,
поэтому прогресс задачи не мог накапливаться.

Повторим ту же задачу с другой моделью. 1 сентября DeepSeek V4 Flash
после удаления предыдущей информации о рассуждениях также выполнила три
раунда и четыре вызова инструментов, но получила результат 9 600 000,00
доллара. Куда исчезли недостающие 2 895,73 доллара? Подсказку даёт
первый вызов конвертации валют: доход за второй квартал составлял 2,1
миллиона евро, но модель указала сумму \texttt{2.1}, намереваясь считать
её в миллионах. Инструмент получил 2,1 евро и после конвертации с
округлением до двух знаков вернул 2,28 доллара. Затем модель включила в
общую сумму 2,28 миллиона долларов, тем самым увеличив ошибку округления
при конвертации в миллион раз. Для двух остальных сумм использовалась та
же запись. Даже если в конце вывести итог с двумя знаками после запятой,
ранее отброшенные разряды уже не вернуть.\footnote{В каталоге
  \texttt{chapter1/context/validation/} репозитория сохранены запросы и
  ответы каждого раунда для обоих запусков. Запись августовского запуска
  находится в \texttt{real\_20260825T025409Z/evidence.json},
  сентябрьского --- в \texttt{real\_20260901T163146Z/evidence.json}.}

Это сравнение показывает роль разных типов информации: определения
инструментов предоставляют доступные для вызова возможности, результаты
инструментов --- факты из окружения, а история сообщений сохраняет
прогресс задачи. Ценность предыдущей информации о рассуждениях зависит
от того, можно ли восстановить основания текущего решения по другим
записям.
\end{quote}

\hypertarget{ux446ux438ux43aux43b-react}{%
\subsection{Цикл ReAct}\label{ux446ux438ux43aux43b-react}}

\textbf{ReAct} (Reasoning + Acting) --- типичная реализация такого
цикла, в которой рассуждение чередуется с действиями\footnote{Yao,
  Shunyu, et al.~ReAct: Synergizing Reasoning and Acting in Language
  Models, ICLR 2023.}. Сначала модель анализирует текущую информацию и
выбирает действие, а затем корректирует план с учётом
\textbf{наблюдения}, полученного от инструмента. В последующих примерах
мы будем кратко указывать основания для решений на каждой итерации в
поле \texttt{reasoning}, чтобы помочь читателю понять связь между
действиями.

Последовательность решений, действий и наблюдений, накопленная в ходе
выполнения задачи, называется \textbf{траекторией} (trajectory). Один из
способов записи траектории --- история сообщений: пользовательские
сообщения содержат задачу, ответы модели --- решения и запросы на
вызовы, а сообщения инструментов --- результаты выполнения. Полная
запись запуска может также включать время, состояние среды и результаты
проверки для последующей отладки и оценки.

В базовой реализации из этого раздела \textbf{контекст каждой итерации =
относительно стабильный префикс + накопленная история сообщений}.
Префикс содержит системный промпт и определения инструментов, а история
сообщений растёт по мере взаимодействия. На
\crossreflink{fig:1-4}{рисунке 1-4} показаны сообщения трёх итераций
эксперимента 1-1. Посмотрите на три запроса конвертации валют на первой
итерации: какой результат соответствует 2,1 миллиона евро за второй
квартал?

\begin{figure}
\hypertarget{fig:1-4}{%
\centering
\includegraphics{images/revised/fig1-4.pdf}
\caption{Рисунок 1-4. Траектория агента --- цикл ReAct в задаче сведения
данных в нескольких валютах}\label{fig:1-4}
}
\end{figure}

ID запроса для евро --- \texttt{convert\_currency\_0}. Harness проверяет
параметры и выполняет инструмент, а затем записывает тот же ID в поле
\texttt{tool\_call\_id} сообщения инструмента. Так можно найти результат
2 282 608,70 доллара и проверить, правильно ли он использован при
суммировании на второй итерации. Эту сумму можно пересчитать: в
фиксированной таблице курсов эксперимента 1 доллар равен 0,92 евро,
поэтому деление 2 100 000 на 0,92 с округлением до двух знаков после
запятой даёт именно это число. Обратный курс 1,087, возвращённый
инструментом, уже округлён: если напрямую умножить его на сумму в евро,
результат будет отличаться на 91,30 доллара\footnote{В каталоге
  \texttt{chapter1/context/validation/} репозитория сохранены запросы и
  ответы каждого раунда для обоих запусков. Запись августовского запуска
  находится в \texttt{real\_20260825T025409Z/evidence.json},
  сентябрьского --- в \texttt{real\_20260901T163146Z/evidence.json}.}.

Цикл выполнения можно представить следующим псевдокодом в стиле Python.
Model определяет следующий шаг, Harness формирует контекст, проверяет и
диспетчеризирует вызовы, а инструмент выполняется в среде и возвращает
наблюдение. Здесь условием завершения служит прекращение запросов
инструментов со стороны модели; механизмы проверки и завершения в
промышленных системах будут рассмотрены далее.

\begin{Shaded}
\begin{Highlighting}[]
\NormalTok{history }\OperatorTok{=}\NormalTok{ [user\_request]}

\NormalTok{repeat:}
\NormalTok{    context }\OperatorTok{=}\NormalTok{ stable\_prefix }\OperatorTok{+}\NormalTok{ history}
\NormalTok{    decision }\OperatorTok{=}\NormalTok{ Model(context)}
\NormalTok{    history.append(decision)}

    \ControlFlowTok{if}\NormalTok{ decision has no tool call:}
        \ControlFlowTok{return}\NormalTok{ decision.answer}

    \ControlFlowTok{for}\NormalTok{ call }\KeywordTok{in}\NormalTok{ decision.tool\_calls:  }\CommentTok{\# Взаимно независимые вызовы можно выполнять параллельно}
\NormalTok{        validated\_call }\OperatorTok{=}\NormalTok{ Harness.validate(call)}
\NormalTok{        observation }\OperatorTok{=}\NormalTok{ Harness.dispatch(validated\_call)}
\NormalTok{        history.append(observation)}
\end{Highlighting}
\end{Shaded}

Ниже приведена выдержка из сообщений в порядке запуска августовской
базовой версии. Исходные ID вызовов и имена параметров сохранены; для
результатов конвертации указаны только сумма и валюта, а вычислительный
код сокращён. Дойдя до второй итерации, можно найти источник каждого
слагаемого в выражении: квартальных доходов четыре --- почему же
инструмент вернул только три результата конвертации?

\begin{Shaded}
\begin{Highlighting}[]
\NormalTok{История сообщений = [}
\NormalTok{  \{role: "user",}
\NormalTok{   content: "Первый квартал — 2,5 миллиона долларов США, второй — 2,1 миллиона евро, третий — 1,8 миллиона фунтов стерлингов, четвёртый — 380 миллионов иен. С помощью инструментов конвертируй суммы в доллары США и рассчитай общий годовой и средний квартальный доход."\},}

\NormalTok{  \# Первая итерация: три запроса конвертации валют}
\NormalTok{  \{role: "assistant",}
\NormalTok{   reasoning: "Сначала конвертировать в USD доходы в EUR, GBP и JPY.",}
\NormalTok{   tool\_calls: [}
\NormalTok{     \{id: "convert\_currency\_0", name: "convert\_currency",}
\NormalTok{      args: \{amount: 2100000, from\_currency: "EUR", to\_currency: "USD"\}\},}
\NormalTok{     \{id: "convert\_currency\_1", name: "convert\_currency",}
\NormalTok{      args: \{amount: 1800000, from\_currency: "GBP", to\_currency: "USD"\}\},}
\NormalTok{     \{id: "convert\_currency\_2", name: "convert\_currency",}
\NormalTok{      args: \{amount: 380000000, from\_currency: "JPY", to\_currency: "USD"\}\}}
\NormalTok{   ]\},}
\NormalTok{  \{role: "tool", tool\_call\_id: "convert\_currency\_0", content: "2282608.70 USD"\},}
\NormalTok{  \{role: "tool", tool\_call\_id: "convert\_currency\_1", content: "2278481.01 USD"\},}
\NormalTok{  \{role: "tool", tool\_call\_id: "convert\_currency\_2", content: "2541806.02 USD"\},}

\NormalTok{  \# Вторая итерация: вычисления по результатам конвертации}
\NormalTok{  \{role: "assistant",}
\NormalTok{   reasoning: "Доход за первый квартал уже указан в долларах США; прибавить к нему остальные три результата конвертации, затем разделить сумму на 4.",}
\NormalTok{   tool\_calls: [}
\NormalTok{     \{id: "code\_interpreter\_3", name: "code\_interpreter",}
\NormalTok{      args: \{code: """}
\NormalTok{        total = 2500000 + 2282608.70 + 2278481.01 + 2541806.02}
\NormalTok{        average = total / 4}
\NormalTok{        print(round(total, 2), round(average, 2))}
\NormalTok{      """\}\}}
\NormalTok{   ]\},}
\NormalTok{  \{role: "tool", tool\_call\_id: "code\_interpreter\_3",}
\NormalTok{   content: "9602895.73 2400723.93"\},}

\NormalTok{  \# Третья итерация: выдача результата}
\NormalTok{  \{role: "assistant",}
\NormalTok{   reasoning: "Вычисления завершены; представить результаты с двумя знаками после запятой.",}
\NormalTok{   content: "Общий годовой доход составляет 9 602 895,73 доллара США, а средний квартальный доход — 2 400 723,93 доллара США."\}}
\NormalTok{]}
\end{Highlighting}
\end{Shaded}

2,5 миллиона долларов за первый квартал уже указаны в сообщении
пользователя, поэтому эту сумму можно сразу включить в сложение.
Остальные три суммы взяты из результатов инструментов. Сложив все четыре
значения и разделив результат на 4, можно проверить итоговый средний
квартальный доход. Всего здесь модель вызывается три раза, а инструменты
--- четыре: на первой итерации модель одним запросом инициирует три
независимые операции конвертации валют, на второй запрашивает выполнение
кода для сведения результатов, а на третьей напрямую отвечает
пользователю.

Запись процесса выполнения в виде структурированной траектории также
помогает разработчикам улучшать систему. По цепочке из пользовательского
запроса, решений модели, параметров вызовов и результатов инструментов
можно локализовать этап сбоя. Сведение нескольких траекторий позволяет
выявлять повторяющиеся операции, совершенствовать интерфейсы
инструментов и оптимизировать пути принятия решений. Оценённый и
отобранный опыт можно записывать в базу знаний или использовать как
данные для постобучения модели. В \crossreflink{chap:7}{главах 7}--9
будут последовательно рассмотрены оценка, постобучение и непрерывная
эволюция.

В задаче конвертации валют понадобилось лишь несколько вызовов
инструментов. Задачи поиска и анализа обычно требуют более длинных
циклов и чаще сталкиваются с нехваткой информации и сбоями инструментов.
Следующие два эксперимента покажут, как такие задачи выполняются шаг за
шагом и за что отвечают клиент и сервер.

\begin{quote}
\textbf{Эксперимент 1-2 ★: многоэтапный поиск и корректировка запросов в
Kimi K3}

Исследовательская задача обычно требует нескольких поисковых запросов:
сначала устанавливаются основные факты, затем разыскиваются недостающие
доказательства и, наконец, сводятся результаты. В этом эксперименте Kimi
K3 проверяет членство стран в АСЕАН и статус столицы Индонезии. После
первоначального поиска модель также должна продолжить поиск ещё не
подтверждённых фактов.

Kimi K3 использует архитектуру смеси экспертов (MoE, Mixture of
Experts), содержит около 2,8 триллиона параметров, изначально
поддерживает визуальные данные и контекст объёмом в миллион токенов. MoE
можно сравнить с командой специалистов: при обработке каждого токена
механизм маршрутизации выбирает для вычислений только часть экспертов,
ограничивая вычислительные затраты. Постобучение модели включает
обучение с подкреплением для универсальных задач, программирования и
задач агентов с упором на повышение способности к длительному
выполнению\footnote{Kimi Team. Kimi K3: Open Frontier Intelligence,
  2026. Архитектура модели, длина контекста и методы постобучения
  описаны в техническом отчёте; подключение сервиса --- в официальной
  документации по инструментам Kimi.}.

В этой траектории модель сначала ищет официальные материалы АСЕАН о
членстве, затем дополнительно проверяет список участников и даты их
вступления, после чего сужает запрос до нормативных актов, президентских
указов и материалов судебных органов, связанных с переносом столицы
Индонезии. Задача завершается структурированным ответом со ссылками на
источники.

Отслеживание этого процесса позволяет увидеть распределение обязанностей
между тремя сторонами. \textbf{Модель} определяет содержание запросов,
параметры вызовов и следующий поисковый шаг. \textbf{Клиентский Harness}
получает определения инструментов, вызывает модель, передаёт запросы
Formula и добавляет результаты в контекст. \textbf{Сервис инструментов
Formula} выполняет поиск в интернете и возвращает результаты. Таким
образом, результаты предыдущей итерации попадают в следующий запрос, и
модель может соответствующим образом скорректировать поисковый
запрос\footnote{Благодарим читателя asdlem, который в задаче № 30
  сопутствующего репозитория указал на необходимость раздельно описать
  принятие решений о вызове инструментов и механизм их выполнения.}.

Стратегию вызова инструментов можно внедрить в параметры модели
посредством обучения, тогда как реализация поиска и возврат его
результатов относятся к системе выполнения. Такое распределение
обязанностей применимо и к другим инструментам, например исполнению
кода: модель решает, что следует выполнить, сервис инструментов
предоставляет среду выполнения, а Harness организует вызовы и обратную
связь.
\end{quote}

\begin{quote}
\textbf{Эксперимент 1-3 ★: управляемый сервером замкнутый
исследовательский цикл Responses API}

После поиска часто требуются дополнительные вычисления. Например, чтобы
найти ближайшую пару среди группы столиц, сначала нужно получить их
координаты, а затем перебрать пары городов и вычислить расстояния. В
этом эксперименте используется \texttt{qwen3.7-plus} от Bailian, который
через Responses API вызывает два управляемых сервером инструмента:
\texttt{web\_search} получает данные и источники, а
\texttt{code\_interpreter} запускает Python в песочнице для вычислений и
построения графиков\footnote{Способ подключения управляемых инструментов
  описан в документации Bailian по интерпретатору кода. Модель, порядок
  вызовов и результаты этого эксперимента приведены по сопутствующей
  записи запуска.}.

В первом задании этого запуска вычисления выполняются для десяти столиц:
Брунея, Камбоджи, Индонезии, Лаоса, Малайзии, Мьянмы, Филиппин,
Сингапура, Таиланда и Вьетнама. Получив координаты с помощью поиска,
модель пишет код, перебирающий \(10\times9/2=45\) пар городов, и
рассчитывает расстояния по большой окружности с помощью формулы
гаверсинусов, используя сферическую модель Земли с радиусом 6 371
километр. При координатах, использованных в этом запуске, ближайшей
оказалась пара Куала-Лумпур --- Сингапур с расстоянием около 316
километров. Поиск предоставляет входные данные, код выполняет
воспроизводимые вычисления, а модель объединяет источники, допущения и
результаты в ответ.

Второе задание: «Найти динамику биткоина за последний месяц и выполнить
технический анализ». Системный промпт предписывает сначала запросить у
пользователя уточнения, если не определены источник данных или
индикаторы. Поэтому модель сначала спрашивает об источнике данных и
показателях анализа. После того как пользователь указывает дневные цены
закрытия BTC/USD с CoinGecko, а также MA7, MA20, RSI14, MACD(12,26,9),
доходность за период и максимальную просадку, модель приступает к поиску
и вычислениям.

В квитанции выполнения кода этого запуска видна деталь, которую легко
упустить. После нескольких тайм-аутов при прямом обращении к интерфейсу
рыночных данных модель в последнем фрагменте Python-кода генерирует 31
последовательную дату и напрямую записывает 31 цену в список. Код
рассчитывает по этим значениям индикаторы и сохраняет график в
песочнице. В отчёте модели говорится, что цены взяты из поисковых
материалов. Квитанция выполнения позволяет понять, как рассчитаны
индикаторы, однако источник списка цен остаётся неясным.

В обеих задачах поиск и исполнение кода управляются сервером. Модель
выбирает дальнейшие действия на основе возвращённых результатов, а
серверный Harness связывает последовательные вызовы. Клиент отправляет
задачу и получает результат; если возникает уточняющий вопрос, он
передаёт дополнительную информацию пользователя, продолжая работу в
существующем контексте. На \crossreflink{fig:1-5}{рисунке 1-5} эта
архитектура сопоставляется с клиентским циклом из эксперимента 1-2, а
различия в распределении обязанностей показаны на конкретных шагах
поиска и вычислений.

\begin{figure}
\hypertarget{fig:1-5}{%
\centering
\includegraphics{images/revised/fig1-5.pdf}
\caption{Рисунок 1-5. Управляемые моделью вызовы инструментов:
клиентский цикл и цикл под управлением сервера}\label{fig:1-5}
}
\end{figure}
\end{quote}

\begin{quote}
\textbf{Дополнение об интерфейсах: вызовы инструментов в свободном
формате}

Параметры инструмента также можно передавать в виде текста свободного
формата. Например, после объявления пользовательского инструмента в
OpenAI Responses API с помощью \texttt{type:\ "custom"} модель может
напрямую отправлять код Python или SQL-запрос, что сокращает объём
экранирования при размещении многострочного текста в строке
JSON\footnote{OpenAI, Function calling; свободный формат входных данных
  и возврат результатов вызова описаны в разделе Custom tools.}. Обычные
инструменты-функции, напротив, принимают структурированные параметры
согласно схеме JSON.

Формат параметров и место выполнения --- два независимых проектных
решения. Для пользовательского инструмента со свободным форматом сторона
приложения по-прежнему принимает входные данные, проверяет их, выполняет
вызов и возвращает результат, тогда как для управляемого интерпретатора
кода песочницу и цикл выполнения предоставляет сервер.
\end{quote}

\hypertarget{ux438ux43dux436ux435ux43dux435ux440ux438ux44f-harness-ux43aux43eux43dux43aux443ux440ux435ux43dux442ux43eux441ux43fux43eux441ux43eux431ux43dux43eux441ux442ux44c-ux437ux430-ux43fux440ux435ux434ux435ux43bux430ux43cux438-ux43cux43eux434ux435ux43bux438}{%
\section{Инженерия Harness: конкурентоспособность за пределами
модели}\label{ux438ux43dux436ux435ux43dux435ux440ux438ux44f-harness-ux43aux43eux43dux43aux443ux440ux435ux43dux442ux43eux441ux43fux43eux441ux43eux431ux43dux43eux441ux442ux44c-ux437ux430-ux43fux440ux435ux434ux435ux43bux430ux43cux438-ux43cux43eux434ux435ux43bux438}}

Два приложения службы поддержки используют одну и ту же модель и оба
отвечают пользователю: «Возврат уже выполнен». Первое сохраняет только
эту фразу, а второе позволяет также получить сумму возврата и статус
транзакции, возвращённые платёжным сервисом. Чтобы определить, какой
возврат действительно состоялся, можно пройти от ответа назад: что
выполнил платёжный сервис? Какую квитанцию увидела модель перед ответом?
Если вызов завершился по тайм-ауту, кто зафиксировал требующую проверки
транзакцию и как было определено следующее действие?

Генерация фразы моделью и изменение состояния платёжной системой --- две
разные операции, каждую из которых необходимо проверять отдельно.
Современные LLM используют коннекционистский подход на основе нейронных
сетей: на данных учатся определять условные вероятности токенов и шаг за
шагом генерируют ответ или запрос к инструменту.\footnote{Консультативная
  исследовательская группа по искусственному интеллекту академических
  отделений Китайской академии наук. «Искусственный интеллект на службе
  научных исследований: система дисциплины искусственного интеллекта».
  Science Press, 2026, ISBN 978-7-03-084786-7,
  \crossreflink{chap:1}{глава 1}, раздел 1.1.} Даже если при каждом
декодировании получается один и тот же ответ, из этого нельзя узнать,
были ли возвращены деньги. Приложение должно передать запрос на
выполнение операции, записать результат, возвращённый инструментом, а
затем передать этот результат модели, чтобы на его основе проверить
окончательный ответ. Организация этих шагов и есть задача инженерии
Harness.

Выше функциональные компоненты AI-агента объяснялись через LLM, контекст
и инструменты. При реализации в коде и сервисах Model принимает решения
на основе входных данных, Harness организует входные данные,
диспетчеризует вызовы инструментов и обрабатывает обратную связь, а
инструменты или среда выполняют реальные операции. В этой книге
разделение обязанностей выражается следующими формулами, а обязанности
Harness сводятся к пяти функциям:

\begin{quote}
\textbf{Agent = Model + Harness}

\textbf{Harness = управление контекстом + интерфейс инструментов +
ограничения + верификация + исправление}

\textbf{Agent \(\leftrightarrow\) Environment}
\end{quote}

Управление контекстом и интерфейс инструментов обеспечивают работу
агента, а ограничения, верификация и исправление повышают надёжность
выполнения. Например, для агента возврата Harness формирует контекст из
политики возврата, заказа и причины обращения, через интерфейс
инструментов запрашивает заказ и инициирует возврат; перед выполнением
проверяет права и правила для суммы, а после выполнения запрашивает
статус транзакции. Если запрос завершается по тайм-ауту, сначала
выясняется, произошёл ли возврат, а затем с учётом идемпотентности
выбирается повторная попытка, компенсация, откат или передача человеку.

Такое разделение обязанностей также помогает определить границы системы.
Определения инструментов, адаптеры вызовов, разрешения песочницы и
механизмы сброса относятся к Harness; изменяющиеся при выполнении
операций файлы и процессы внутри песочницы, внешние базы данных,
веб-страницы, пользователи и физический мир относятся к среде. Даже если
имитационная среда и агент развёрнуты в одном процессе, они по-прежнему
отдельно отвечают за эволюцию среды и принятие решений о действиях.

\begin{longtable}[]{@{}
  >{\raggedright\arraybackslash}p{(\columnwidth - 6\tabcolsep) * \real{0.2145}}
  >{\raggedright\arraybackslash}p{(\columnwidth - 6\tabcolsep) * \real{0.3473}}
  >{\raggedright\arraybackslash}p{(\columnwidth - 6\tabcolsep) * \real{0.2971}}
  >{\raggedright\arraybackslash}p{(\columnwidth - 6\tabcolsep) * \real{0.1411}}@{}}
\toprule\noalign{}
\begin{minipage}[b]{\linewidth}\raggedright
\hspace{0pt}Функция
\end{minipage} & \begin{minipage}[b]{\linewidth}\raggedright
\hspace{0pt}Обязанности и ключевые аспекты проектирования
\end{minipage} & \begin{minipage}[b]{\linewidth}\raggedright
\hspace{0pt}Практические примеры
\end{minipage} & \begin{minipage}[b]{\linewidth}\raggedright
\hspace{0pt}Подробнее
\end{minipage} \\
\midrule\noalign{}
\endhead
\bottomrule\noalign{}
\endlastfoot
\hspace{0pt}\textbf{Context (управление контекстом)} & \hspace{0pt}Выбор
и организация информации, необходимой для текущего решения; сохранение
прогресса задачи и ключевых оснований & \hspace{0pt}Системные промпты,
база знаний, строка состояния агента, запросы Sidecar по отдельному
каналу & \hspace{0pt}\crossreflink{chap:2}{Главы 2} и 3 \\
\hspace{0pt}\textbf{Tools (интерфейс инструментов)} &
\hspace{0pt}Предоставление интерфейсов наблюдения и действий; понятные
названия, примеры параметров, ясные возвращаемые значения и ошибки &
\hspace{0pt}Инструменты MCP, интерпретатор кода, инструменты поиска &
\hspace{0pt}\crossreflink{chap:4}{Глава 4} \\
\hspace{0pt}\textbf{Constrain (ограничения)} &
\hspace{0pt}Предоставление возможностей по принципу минимальных
привилегий и проверка операции, параметров и границ авторизации перед
выполнением & \hspace{0pt}Область доступа к файлам, верхний предел суммы
возврата, подтверждение вызовов инструментов &
\hspace{0pt}\crossreflink{chap:4}{Глава 4} \\
\hspace{0pt}\textbf{Verify (верификация)} & \hspace{0pt}Проверка
операций с помощью проверяемых правил, результатов инструментов и
состояния среды; одновременная проверка источника и действительности
данных & \hspace{0pt}Линтер, проверка типов, тесты, проверка статуса
транзакции & \hspace{0pt}\crossreflink{chap:5}{Главы 5} и 6 \\
\hspace{0pt}\textbf{Correct (исправление)} & \hspace{0pt}Выбор повторной
попытки, продолжения, компенсации, отката или передачи человеку в
зависимости от типа ошибки; ограничение числа попыток восстановления &
\hspace{0pt}Незаметная повторная попытка, продолжение генерации,
передача человеку после срабатывания автоматического выключателя &
\hspace{0pt}\crossreflink{chap:2}{Главы 2} и 5 \\
\end{longtable}

Ограничения можно начинать с отключения чувствительных возможностей по
умолчанию: доступ к чувствительным ресурсам или возможности с внешними
побочными эффектами открываются только после явного разрешения ---
подобно тому, как мобильное приложение запрашивает доступ к камере или
геопозиции. Верификация должна связывать утверждения модели с
проверяемыми доказательствами. Например, успешный возврат должен
подтверждаться состоянием в транзакционной системе; если инструмент
возвращает JSON, необходимо также проверять типы полей, значения и
источник. Как вывод модели, так и содержимое внешних инструментов могут
подвергаться воздействию ошибочной информации или инъекции в промпт,
поэтому доверенная логика верификации должна работать в независимой
границе выполнения.

Механизм исправления должен учитывать и восстановление задачи, и
удобство пользователя. При кратковременном сбое можно автоматически
повторить попытку в пределах бюджета, не отвлекая пользователя слишком
часто; если операция имеет внешние побочные эффекты, сначала необходимо
проверить её состояние. При устойчивом сбое следует прекратить
повторение операции и сообщить о текущем прогрессе и действиях,
требующих вмешательства человека. Если добавить эти пять функций в
рассмотренный выше цикл, получится следующий псевдокод:

\begin{Shaded}
\begin{Highlighting}[]
\NormalTok{history }\OperatorTok{=}\NormalTok{ [user\_request, Environment.observe()]}

\ControlFlowTok{while}\NormalTok{ Harness.within\_budget():}
\NormalTok{    decision }\OperatorTok{=}\NormalTok{ Model(Harness.build\_context(history))}
\NormalTok{    history.append(decision)}

    \ControlFlowTok{if}\NormalTok{ decision has no tool call:}
\NormalTok{        verdict }\OperatorTok{=}\NormalTok{ Harness.verify\_completion(decision, Environment)}
        \ControlFlowTok{if}\NormalTok{ verdict.passed:}
            \ControlFlowTok{return}\NormalTok{ decision.answer}
\NormalTok{        history.append(verdict.feedback)}
        \ControlFlowTok{continue}

    \ControlFlowTok{for}\NormalTok{ call }\KeywordTok{in}\NormalTok{ decision.tool\_calls:}
\NormalTok{        permission }\OperatorTok{=}\NormalTok{ Harness.constrain(call)}
        \ControlFlowTok{if}\NormalTok{ permission.denied:}
\NormalTok{            history.append(permission.feedback)}
            \ControlFlowTok{continue}

\NormalTok{        result }\OperatorTok{=}\NormalTok{ Harness.dispatch(call)}
\NormalTok{        verdict }\OperatorTok{=}\NormalTok{ Harness.verify(result, Environment)}
        \ControlFlowTok{if} \KeywordTok{not}\NormalTok{ verdict.passed:}
\NormalTok{            result }\OperatorTok{=}\NormalTok{ Harness.correct(call, result, verdict)}
\NormalTok{        history.append(result)}

\ControlFlowTok{return}\NormalTok{ Harness.report\_incomplete(history)}
\end{Highlighting}
\end{Shaded}

В этом цикле ограничения применяются перед выполнением, верификация ---
после возврата результата инструментом и при завершении задачи, а
результат исправления возвращается в контекст. Во второй главе подробно
рассматривается цикл сообщений API, а в четвёртой и пятой --- выполнение
инструментов и автоматическая верификация.

\textbf{При переходе от демонстрации к промышленной эксплуатации
инженерный акцент смещается с подключения контекста и инструментов на
ограничения, верификацию и исправление.} Возможности системы можно
продемонстрировать на нескольких успешно выполненных задачах; продукт же
должен стабильно выдавать результат при изменении входных данных, сбоях
сервисов и длительной работе. Поэтому права, состояние, верификация и
восстановление становятся основными направлениями инженерных усилий.
Конкретный пример даёт Claude Code: помимо чтения и записи файлов,
выполнения команд и поиска, он предоставляет правила разрешений, сжатие
контекста и контрольные точки файлов. Правила разрешений определяют
необходимость подтверждения с учётом типа инструмента, области операции
и выбранного режима; контрольные точки позволяют восстановить
зафиксированные ими изменения файлов\footnote{Официальная документация
  Claude Code: настройка разрешений, контрольные точки, интерактивные
  команды и сжатие контекста. Область восстановления файловой
  контрольной точки определяется данными, зафиксированными инструментом;
  состояние во внешних сервисах требует соответствующего механизма
  восстановления.}. В автоматическом режиме Claude Code также передаёт
вызовы инструментов на проверку независимому классификатору. Именно
такие задачи подходят специализированным моделям принятия решений вроде
Jev: основная модель должна продумать способ выполнения задачи, а
вспомогательной достаточно ответить на один чётко ограниченный вопрос,
например: «Выходит ли эта операция за рамки того, что только что поручил
пользователь?» В четвёртой главе разделение обязанностей будет показано
на примере подтверждения вызовов инструментов.

На основе этих продуктовых механизмов можно выделить несколько
направлений работы, которые необходимо учитывать при создании Harness:

\begin{itemize}
\tightlist
\item
  \textbf{Управление состоянием процесса}: фиксировать текущий этап
  выполнения, завершённые операции и результаты, которые ещё требуют
  подтверждения.
\item
  \textbf{Сжатие контекста}: по мере роста объёма информации сохранять
  ключевые факты, ограничения и прогресс, освобождая место для
  последующих решений.
\item
  \textbf{Классификация разрешений}: различать чтение, изменение,
  внешние вызовы и другие операции, применяя правила в соответствии с
  риском и границами авторизации.
\item
  \textbf{Автоматический выключатель} (Circuit Breaker): прекращать
  повторные попытки после достижения порога последовательных сбоев,
  чтобы одна и та же неисправность не продолжала расходовать ресурсы и
  увеличивать ущерб; этот механизм подобен срабатыванию защитного
  автомата при неисправности электрической цепи.
\item
  \textbf{Восстановление после ошибок}: перехватывать исключения и с
  учётом контрольных точек и состояния среды выбирать откат к
  стабильному состоянию, повторную попытку или передачу управления
  человеку.
\end{itemize}

\hypertarget{ux43eux442-ux43fux440ux43eux43cux43fux442-ux438ux43dux436ux438ux43dux438ux440ux438ux43dux433ux430-ux43a-ux433ux440ux430ux444ux43eux432ux43eux43cux443-ux438ux43dux436ux438ux43dux438ux440ux438ux43dux433ux443-ux440ux430ux441ux448ux438ux440ux435ux43dux438ux435-ux441ux444ux435ux440ux44b-ux438ux43dux436ux435ux43dux435ux440ux43dux43eux433ux43e-ux432ux43dux438ux43cux430ux43dux438ux44f}{%
\subsection{От промпт-инжиниринга к графовому инжинирингу: расширение
сферы инженерного
внимания}\label{ux43eux442-ux43fux440ux43eux43cux43fux442-ux438ux43dux436ux438ux43dux438ux440ux438ux43dux433ux430-ux43a-ux433ux440ux430ux444ux43eux432ux43eux43cux443-ux438ux43dux436ux438ux43dux438ux440ux438ux43dux433ux443-ux440ux430ux441ux448ux438ux440ux435ux43dux438ux435-ux441ux444ux435ux440ux44b-ux438ux43dux436ux435ux43dux435ux440ux43dux43eux433ux43e-ux432ux43dux438ux43cux430ux43dux438ux44f}}

Прежде чем менять промпт для неудачно выполненной задачи, стоит
определить, на каком этапе произошёл сбой. Если модель неверно поняла
одно из ограничений, сначала следует проверить инструкции; если весь
счёт так и не попал в модель --- организацию информации; если после
прерывания задачи потерялись уже выполненные шаги --- механизм
выполнения. Когда задача растягивается на несколько дней и требует
совместной работы нескольких участников, на результат также влияют
непрерывность работы и структура выполнения. Ниже представлены пять
инженерных подходов, каждый из которых даёт свою отправную точку для
диагностики.

\textbf{Промпт-инжиниринг} (Prompt Engineering) сосредоточен на
формулировке инструкций. От того, как описаны цель задачи, ограничения и
примеры, напрямую зависит, как модель поймёт запрос и сформирует
результат.

\textbf{Контекстный инжиниринг} (Context Engineering) идёт дальше и
организует всю доступную модели информацию, включая системные
инструкции, определения инструментов, историю диалога и внешние знания.
Промпт --- лишь одна из её частей: на принятие решений также влияют
поиск, сжатие и организация состояния.

\textbf{Harness-инжиниринг} сосредоточен на выполнении модели и её
взаимодействии с окружением. Он объединяет формирование контекста,
диспетчеризацию инструментов, разрешения, проверку и восстановление
после ошибок в исполняемую систему, позволяющую модели последовательно
выполнять задачу.

\textbf{Инжиниринг циклов} (Loop Engineering) сосредоточен на
непрерывной автономной работе: кто после завершения одного раунда
определяет следующую задачу, когда перезапускать выполнение, когда
проверять результат и по каким условиям определять, что работа
действительно завершена. В десятой главе этот подход будет рассмотрен на
примере систем совместной работы нескольких AI-агентов.

\textbf{Графовый инжиниринг} (Graph Engineering) сосредоточен на
объединении нескольких исполнительных единиц. Циклы агентов,
детерминированные программы и ручное подтверждение могут выступать
узлами; рёбра задают маршрутизацию и зависимости; структурированное
состояние передаётся между узлами, а на критических границах сохраняются
контрольные точки. В обсуждениях, появившихся в 2026 году, это название
использовалось как обобщение подобных задач оркестрации\footnote{Josh C.
  Simmons, We Are Entering the Graph Engineering Phase, 4 июля 2026
  года. В статье рассматриваются узлы, типизированные соединения и
  состояние с возможностью сохранения контрольных точек. Связанные
  практические подходы также представлены в официальных материалах
  LangGraph, рабочих процессов Microsoft Agent Framework и Google ADK.}.

Harness обеспечивает управление контекстом и цикл выполнения отдельной
задачи, а внешний механизм диспетчеризации может многократно запускать
задачу. Граф выполнения может включать циклы агентов, а цикл, в свою
очередь, может вызывать графовый рабочий процесс для выполнения
подзадач.

\textbf{Когда все рассматриваемые модели уже способны справиться с
задачей, различия между продуктами зачастую определяются инженерными
практиками за пределами модели.} Улучшение результатов LangChain в
Terminal Bench 2.0 служит конкретным примером: при неизменном
использовании GPT-5.2-Codex за счёт доработки Harness результат вырос с
52,8\% до 66,5\%, а позиция в рейтинге на тот момент поднялась с места
за пределами первой тридцатки в первую пятёрку. Доработки включали
самопроверку перед завершением, обнаружение повторяющихся операций,
дополнение сведений об окружении и распределение бюджета рассуждений.
Этот прирост был получен благодаря анализу сценариев сбоев в траекториях
выполнения и последующим многократным изменениям и оценке
Harness\footnote{Vivek Trivedy, Improving Deep Agents with harness
  engineering, 17 февраля 2026 года. Результаты и позиции в рейтинге
  соответствуют модели, бенчмарку и конфигурации тестирования, указанным
  в статье.}.

\hypertarget{ux43eux441ux43dux43eux432ux43dux44bux435-ux43fux440ux438ux43dux446ux438ux43fux44b-ux43fux43eux441ux442ux440ux43eux435ux43dux438ux44f-ux44dux444ux444ux435ux43aux442ux438ux432ux43dux43eux433ux43e-ai-ux430ux433ux435ux43dux442ux430}{%
\subsection{Основные принципы построения эффективного
AI-агента}\label{ux43eux441ux43dux43eux432ux43dux44bux435-ux43fux440ux438ux43dux446ux438ux43fux44b-ux43fux43eux441ux442ux440ux43eux435ux43dux438ux44f-ux44dux444ux444ux435ux43aux442ux438ux432ux43dux43eux433ux43e-ai-ux430ux433ux435ux43dux442ux430}}

Прежде чем приступать к проектированию, выясните, какой именно этап
пользователь хочет передать агенту. Автор провёл эксперимент с
покупками: помощник подобрал десять товаров, но сначала покупательнице
не понравился ни один. Она хотела, чтобы товар был красивым, недорогим и
желательно без отрицательных отзывов. Только самостоятельно просмотрев
предложения на Taobao и сравнив варианты, она постепенно определилась, в
чём готова уступить, и в итоге выбрала один из тех самых десяти товаров.
Если смотреть лишь на совершённую покупку, может показаться, что первой
рекомендации было достаточно. Но, впервые увидев список, она ещё не
сформировала предпочтения, достаточные для принятия решения. Помощник
мог бы проверить для неё характеристики, сравнить цены и перечислить
показательные различия, чтобы она могла самостоятельно просматривать
предложения и выбирать.

Для оспаривания счёта ответ будет другим. Пользователь хочет разобраться
с одним списанием, но мог отказаться от этого из-за необходимости час
ждать в очереди. Помощник ждёт и проверяет вместо него, экономя именно
то время, которое пользователь хотел сберечь. Есть и задачи, которые
человек изначально не способен выполнить самостоятельно: например, не
умеющий программировать пользователь с помощью помощника создаёт нужную
ему небольшую утилиту. Определяя границы делегирования, можно спросить
пользователя: если бы кто-то сделал работу за вас, от какого этапа вы
больше всего хотели бы избавиться? А в каком этапе вы хотели бы
участвовать лично, чтобы получить от него пользу?

Anthropic свела опыт построения эффективных агентов к трём принципам:
сохраняйте простоту, обеспечивайте прозрачность и тщательно проектируйте
интерфейсы между агентом и инструментами\footnote{Anthropic. «Building
  effective agents», декабрь 2024 года.}.

\textbf{Сохраняйте простоту.} Начинайте с самого простого решения,
способного выполнить задачу, и добавляйте новые механизмы лишь после
того, как оценка выявит устойчивую проблему. Прямые вызовы API зачастую
позволяют яснее видеть входные данные, результаты и поток управления.
При использовании фреймворка также необходимо понимать, как он
организует контекст и вызывает инструменты.

\textbf{Обеспечивайте прозрачность.} Чтобы выяснить, почему возврат
средств завершился неудачей, разработчик должен иметь возможность
проследить по журналу выполнения ответ платёжного интерфейса.
Пользователю же важно знать, возвращены ли деньги и кто отвечает за
следующий шаг. План задачи, журналы инструментов и ключевые решения
можно записывать, а пользователю сообщать о ходе работы и вопросах,
требующих подтверждения.

\textbf{Качественно проектируйте интерфейсы инструментов.} ACI
(Agent-Computer Interface, интерфейс «агент --- компьютер») предполагает
проектирование инструментов с учётом того, как ими пользуется агент:
названия должны быть интуитивно понятными, параметры --- сопровождаться
конкретными примерами, возвращаемые значения --- облегчать последующие
решения, а места, где легко допустить ошибку, --- защищаться исполнимыми
ограничениями. Чем яснее интерфейс, тем проще модели преобразовать
намерение в правильное действие.

Здесь можно воспользоваться производственным методом \textbf{защиты от
ошибок} (Poka-yoke). Скошенный угол SIM-карты помогает пользователю
определить правильное направление установки, а блокировка дверцы
микроволновой печи останавливает нагрев при её открытии. В
проектировании инструментов этому соответствуют перечисления,
ограничивающие тип операции, проверка формата параметров и управление
разрешениями, ограничивающее область действия. Модель понимает интерфейс
через определение инструмента и возвращаемый результат, поэтому
неоднозначное описание будет приводить к неправильному использованию на
протяжении нескольких циклов вызовов. Ясный интерфейс и исполнимые
ограничения позволяют предотвращать ошибки в самом их источнике.

Далее проектные компромиссы рассматриваются в трёх аспектах: выбор
модели, паттерны оркестрации, защитные барьеры и безопасность. Выбор
модели определяет доступные возможности принятия решений, оркестрация
организует путь выполнения, а защитные барьеры воплощают ограничения,
проверку и исправление в конкретных операциях.

\hypertarget{ux43aux430ux43a-ux432ux44bux431ux440ux430ux442ux44c-ux43cux43eux434ux435ux43bux44c}{%
\subsection{Как выбрать
модель}\label{ux43aux430ux43a-ux432ux44bux431ux440ux430ux442ux44c-ux43cux43eux434ux435ux43bux44c}}

\textbf{Правильный выбор модели зачастую эффективнее многократной
доработки промпта.} Модель определяет, насколько сложные задачи способен
понять AI-агент, может ли он правильно выбрать инструмент и как
корректирует свои действия после ошибки. Если информации о задаче
достаточно и инструменты доступны, но ошибки планирования или
использования инструментов продолжают повторяться, следует сначала
сравнить возможности моделей и только затем решать, сколько усилий
вкладывать в промпты и оркестрацию.

\textbf{Закрытые модели подходят для быстрой интеграции через сервисные
интерфейсы.} В этой книге в качестве примеров рассматриваются серии GPT
от OpenAI и Claude от Anthropic. При оценке необходимо учитывать не
только возможности в решении задач, но и предоставляемые API способы
вызова инструментов, размер контекстного окна, ограничения сервиса и
стоимость.

\textbf{Модели с открытыми весами предоставляют больше возможностей для
развёртывания и настройки.} Такие модели, как DeepSeek, Kimi и GLM, в
зависимости от условий лицензии конкретной версии можно развёртывать
самостоятельно или использовать для последующего обучения. Развёртывание
в собственной инфраструктуре помогает контролировать движение данных, а
дообучение позволяет адаптировать модель к конкретной задаче.
\textbf{Для бизнес-сценариев, в которых возможностей модели уже
достаточно и объём вызовов стабилен, а также для сценариев, требующих
строгого контроля движения данных, автор в первую очередь рассматривает
решения с открытыми весами.}

На момент написания книги, по эмпирической оценке автора, модели с
открытыми весами сократили отставание до менее чем полугода.

\textbf{Одновременно оценивайте возможности модели, интерфейсы продукта
и правила использования.} Когда одна и та же модель предоставляется
через чат-продукт, Coding Agent или API, доступные инструменты,
разрешения и ограничения среды выполнения могут различаться. Для задач,
связанных с кибербезопасностью, дистилляцией и извлечением моделей,
конфиденциальными данными или высокорисковыми операциями, необходимо
также отдельно проверять правила использования сервиса.

\textbf{Для многоэтапных агентных задач автор предпочитает начинать с
reasoning-модели.} Планирование, выбор инструментов, проверка
результатов и восстановление после неудач требуют последовательных
решений, поэтому выделение большего объёма вычислений на рассуждения на
этих этапах обычно помогает повысить качество выполнения задачи. Для
одношаговой классификации, простого преобразования формата или нажатия в
заранее известной фиксированной точке в Computer Use может быть
достаточно и более лёгкой модели.

\textbf{Оценивайте скорость в масштабе полной задачи.} Агент часто
выполняет несколько последовательных раундов вызовов, поэтому время
генерации модели в каждом раунде накапливается в сквозной задержке. Если
задача требует двадцати раундов рассуждений и каждый из них занимает на
две секунды больше, общее время увеличивается на сорок секунд. В него
входят ожидание первого токена, генерация модели и выполнение
инструментов.

\textbf{Проверяйте мультимодальные возможности с учётом формата входных
данных.} Если требуется понимать изображения, аудио или видео,
соответствующие возможности следует оценивать на реальных материалах
задачи: например, на мелком тексте в снимках интерфейса, перебиваниях в
речи или непрерывных действиях в видео.

\hypertarget{ux440ux435ux436ux438ux43cux44b-ux43eux440ux43aux435ux441ux442ux440ux430ux446ux438ux438-ux440ux430ux431ux43eux447ux438ux435-ux43fux440ux43eux446ux435ux441ux441ux44b-ux438-ux430ux432ux442ux43eux43dux43eux43cux43dux43eux441ux442ux44c}{%
\subsection{Режимы оркестрации: рабочие процессы и
автономность}\label{ux440ux435ux436ux438ux43cux44b-ux43eux440ux43aux435ux441ux442ux440ux430ux446ux438ux438-ux440ux430ux431ux43eux447ux438ux435-ux43fux440ux43eux446ux435ux441ux441ux44b-ux438-ux430ux432ux442ux43eux43dux43eux43cux43dux43eux441ux442ux44c}}

Кто решает, что делать дальше после получения запроса пользователя?
Однократный вызов модели подходит для задач, где уже достаточно исходных
данных и чётко определён результат. Если постоянно повторяются одни и те
же шаги, программа может организовать их в рабочий процесс. Если же путь
нужно менять с учётом обратной связи от среды, выбор можно поручить
модели. Такое разделение определяет, как контекст передаётся между
вызовами и когда выполняются инструменты.

Рассмотрим задачу «извлечь личные воспоминания из миллиона сообщений».
Один из вариантов архитектуры предполагает сначала разделить диалоги, а
затем поручить нескольким агентам разные задачи: извлекать факты,
сверять их с исходными репликами, распознавать упомянутых в разговорах
людей и систематизировать воспоминания. После этого ещё один агент
проверяет и объединяет результаты. Система также должна хранить в графе
связи между фактами, отмечать уже обработанные сообщения и сохранять
каждую версию результата как доступный только для чтения снимок. Каждый
из этих механизмов решает практическую проблему, однако их одновременное
внедрение с самого начала повышает нагрузку, связанную с передачей
информации, синхронизацией состояния и локализацией ошибок. Если при
каждом новом сценарии продолжать добавлять ветвления, семантические
решения, которые модель могла бы принимать с учётом контекста,
постепенно превратятся в трудно сопровождаемые фиксированные правила.

Лучше начать с достаточно способного агента, предоставив ему ясную цель,
необходимый контекст и комбинируемые инструменты, чтобы он мог проверять
сведения, сопоставлять утверждения и разрешать противоречия. Программа
при этом обеспечивает правила, которые должны соблюдаться всегда:
полномочия, защиту исходных материалов, атомарную публикацию и другие
гарантии. Если оценка покажет, что система постоянно ошибается в задачах
одного типа, можно добавить специализированный валидатор, независимого
агента или детерминированный процесс. \textbf{Сложность структуры должна
определяться бизнес-ограничениями и реальными неудачами.} Manus
описывает этот компромисс фразой «Less structure, more intelligence».

\hypertarget{ux440ux435ux436ux438ux43c-ux440ux430ux431ux43eux447ux435ux433ux43e-ux43fux440ux43eux446ux435ux441ux441ux430-ux43fux440ux435ux434ux43eux43fux440ux435ux434ux435ux43bux451ux43dux43dux430ux44f-ux43eux440ux43aux435ux441ux442ux440ux430ux446ux438ux44f}{%
\subsubsection{Режим рабочего процесса: предопределённая
оркестрация}\label{ux440ux435ux436ux438ux43c-ux440ux430ux431ux43eux447ux435ux433ux43e-ux43fux440ux43eux446ux435ux441ux441ux430-ux43fux440ux435ux434ux43eux43fux440ux435ux434ux435ux43bux451ux43dux43dux430ux44f-ux43eux440ux43aux435ux441ux442ux440ux430ux446ux438ux44f}}

\textbf{Рабочий процесс} (Workflow) организует модели и инструменты
посредством заранее определённых путей в коде. Разработчик задаёт узлы,
связи и условия ветвления, а внутри узлов может использоваться LLM для
понимания или генерации содержимого. При разных входных данных система
переходит в соответствующую ветвь согласно программно заданным условиям.

Например, упрощённый бизнес-процесс бронирования авиабилетов можно
спроектировать в виде четырёх узлов:

\begin{enumerate}
\def\labelenumi{\arabic{enumi}.}
\tightlist
\item
  \textbf{Проверить личность пользователя}: подтвердить учётную запись и
  полномочия на выполнение операции через службу идентификации.
\item
  \textbf{Найти и подтвердить рейс}: найти доступные рейсы и попросить
  пользователя подтвердить маршрут, стоимость, а также условия обмена и
  возврата.
\item
  \textbf{Выполнить оплату}: после получения разрешения пользователя
  оплатить заказ через платёжный интерфейс.
\item
  \textbf{Подтвердить бронирование}: оформить билет или подтвердить
  бронирование согласно процессу поставщика и отправить пользователю
  результат.
\end{enumerate}

В узле поиска модель может преобразовать требования на естественном
языке в дату, пункт назначения и условия фильтрации, а программа ---
проверить статус идентификации, ценовое предложение и наличие
разрешения, необходимые для перехода к оплате. Реальный интерфейс
бронирования может требовать сначала зарезервировать место, а затем
выполнить оплату или оформить билет. Рабочий процесс должен учитывать
порядок, установленный протоколом поставщика, и предусматривать
компенсирующие пути на случай сбоев.

Преимущество рабочего процесса --- \textbf{явное управление ходом
выполнения}: ключевые шаги, ветви и контрольные точки можно проверять, а
такие правила, как «платить можно только после подтверждения суммы», ---
реализовывать в коде. Кроме того, для каждого узла удобно ограничивать
полномочия и область доступа к данным.

Чтобы фиксированный процесс правильно обрабатывал исключения,
разработчик должен заранее спроектировать соответствующие ветви. Если
пользователь изменит маршрут на этапе оплаты или после отмены рейса
потребуется заново сравнить варианты, система сможет действовать только
по предусмотренной ветви либо передать управление человеку. Чем больше в
задаче исключений, тем дороже сопровождение таких ветвей.

Другой наглядный пример рабочего процесса --- \textbf{генерация
изображений по тексту}. Пользователь может сказать лишь: «Нарисуй, как
будут работать программисты после создания AGI». Рабочий процесс может
состоять из двух узлов:

\begin{enumerate}
\def\labelenumi{\arabic{enumi}.}
\tightlist
\item
  \textbf{Переработка промпта}: LLM разворачивает общее требование,
  описывая главный объект, сцену, композицию, освещение и стиль.
  Например, представляет сцену как «программист отдыхает на берегу моря
  и управляет AI-сотрудниками через нейроинтерфейс», а затем оформляет
  её в виде промпта, подходящего целевой модели.
\item
  \textbf{Генерация изображения}: подготовленный промпт передаётся
  модели генерации изображений, которая возвращает картинку.
\end{enumerate}

В некоторых рабочих процессах Stable Diffusion разработчики используют
разделённые запятыми английские теги, обозначения качества и негативные
промпты. Узел переработки может приводить входные данные к единому
формату и добавлять подробности композиции. Stable Diffusion также
принимает описания на естественном языке\footnote{Stability AI,
  справочник по Stable Diffusion API с описанием текстовых промптов,
  негативных промптов, стилей и других входных данных; официальные
  примеры применения демонстрируют использование описаний на
  естественном языке.}. Здесь узел LLM преобразует требования
пользователя во входные данные, подходящие модели генерации изображений,
и добавляет творческие детали, поэтому его можно рассматривать как
\textbf{адаптационный слой} Harness.

Если сервис генерации изображений способен непосредственно понимать
требования на естественном языке, \textbf{автор предпочитает передавать
исходное описание напрямую, не добавляя вызов для переработки
исключительно ради адаптации формата входных данных.} Каждая переработка
создаёт ещё одну возможность упустить или изменить ограничение. Узел
переработки стоит задействовать для явно определённых творческих задач:
исследования композиции при общих требованиях или обеспечения единого
стиля серии изображений.

\begin{quote}
\textbf{Эксперимент 1-4 ★: сравнение рабочего процесса генерации
изображений по тексту с прямой генерацией}

Передайте одно и то же требование по двум маршрутам. В \textbf{маршруте
рабочего процесса} Kimi K3 сначала перерабатывает промпт, после чего
изображение генерирует Tongyi Wanxiang. В \textbf{маршруте прямой
генерации} исходное требование передаётся сервису генерации изображений.
В данном эксперименте для прямой генерации использовались две модели:
Gemini для изображений и GPT-Image 2.

Сравнение охватывает требования двух типов. К конкретным относятся
программист, работающий сверхурочно в выходные, комнатные растения на
подоконнике в утреннем свете и рекламный плакат наушников с
шумоподавлением с текстом «В одиночестве глубокой ночью тоже спокойно».
К общим --- рабочая обстановка программиста после создания AGI и утро в
городе будущего.

В случае плаката наушников узел переработки добавил матово-чёрные
наушники, тёмно-синий фон, лунный контровой свет и свободное
пространство. Одновременно он поместил \texttt{text} в негативный промпт
и предложил добавить китайский рекламный текст позднее. Таким образом,
рабочий процесс разделил первоначальную задачу создания готового плаката
на генерацию изображения и последующую вёрстку.
\end{quote}

По мере развития моделей часть адаптационной работы можно передать самой
модели или интерфейсу сервиса. Дообучение на инструкциях и модели
рассуждения снизили потребность в некоторых few-shot-примерах и промптах
с требованием «рассуждать пошагово». Структурированный вывод и нативные
вызовы инструментов сократили объём работы по исправлению формата и
отказоустойчивому разбору JSON. Улучшение понимания языка моделями
генерации изображений упростило переработку промптов. \textbf{После
обновления модели следует целенаправленно удалять адаптационный код,
утративший своё назначение.} Ранее необходимые шаги преобразования и
исправления формата могут превратиться в лишние вызовы, дополнительную
задержку и новые точки отказа.

\hypertarget{ux430ux432ux442ux43eux43dux43eux43cux43dux44bux439-ai-ux430ux433ux435ux43dux442-ux434ux438ux43dux430ux43cux438ux447ux435ux441ux43aux43eux435-ux43fux440ux438ux43dux44fux442ux438ux435-ux440ux435ux448ux435ux43dux438ux439}{%
\subsubsection{Автономный AI-агент: динамическое принятие
решений}\label{ux430ux432ux442ux43eux43dux43eux43cux43dux44bux439-ai-ux430ux433ux435ux43dux442-ux434ux438ux43dux430ux43cux438ux447ux435ux441ux43aux43eux435-ux43fux440ux438ux43dux44fux442ux438ux435-ux440ux435ux448ux435ux43dux438ux439}}

\textbf{Автономный AI-агент} (Autonomous Agent) выбирает следующий шаг с
учётом хода задачи и обратной связи от среды, поэтому подходит для
проблем, где невозможно заранее перечислить все пути выполнения. Если
пользователь говорит: «Забронируй мне билет в Шанхай на следующую
среду», агент может сначала найти рейсы, при появлении требования войти
в систему --- проверить личность, а затем продолжить поиск. Обнаружив,
что самый дешёвый вариант предусматривает пересадку, он спрашивает
пользователя, приемлемо ли это, и после ответа «только прямой рейс»
корректирует условия. Конкретные шаги он выбирает в пределах
предоставленных полномочий, а критические операции, такие как оплата,
по-прежнему ограничиваются бизнес-правилами.

Если запрос завершается неудачей, модель может изменить параметры и
повторить попытку, использовать другой способ запроса либо попросить
пользователя предоставить дополнительную информацию. Harness
регистрирует каждую попытку и ограничивает время или ресурсы, доступные
задаче.

Оба цикла на \crossreflink{fig:1-6}{рисунке 1-6} проходят проверку
бюджета. Слева инструмент вернул результат, справа модель собирается
завершить работу, но при приёмке обнаруживается, что чего-то не хватает.
Когда модель вызывает инструмент окончательного вывода, например
\texttt{final\_answer}, либо возвращает ответ без вызова инструмента,
Harness проверяет наличие всех требуемых результатов. Результат
передаётся пользователю только после успешной проверки. Если
обязательного артефакта не хватает, Harness сообщает модели, что именно
отсутствует, и предлагает продолжить выполнение. При невосстановимой
ошибке, превышении допустимого числа ошибок или максимального количества
раундов система сохраняет достигнутый прогресс, сообщает о незавершённой
части либо передаёт задачу человеку.

\begin{figure}
\hypertarget{fig:1-6}{%
\centering
\includegraphics{images/revised/fig1-6.pdf}
\caption{Рисунок 1-6. Цикл выполнения автономного
AI-агента}\label{fig:1-6}
}
\end{figure}

Автономный цикл подходит для сценариев, где количество шагов зависит от
задачи: например, когда Coding Agent исправляет реальную проблему в
коде, Computer Use Agent управляет интерфейсом либо исследовательская
задача требует многократного поиска и анализа. SWE-bench оценивает
способность исправлять код на задачах из реальных программных
репозиториев. В таких задачах следующий шаг часто определяется только
что полученным выводом тестов, состоянием интерфейса или результатами
поиска.

\hypertarget{ux432ux44bux431ux43eux440-ux438-ux441ux43eux447ux435ux442ux430ux43dux438ux435-ux434ux432ux443ux445-ux440ux435ux436ux438ux43cux43eux432}{%
\subsubsection{Выбор и сочетание двух
режимов}\label{ux432ux44bux431ux43eux440-ux438-ux441ux43eux447ux435ux442ux430ux43dux438ux435-ux434ux432ux443ux445-ux440ux435ux436ux438ux43cux43eux432}}

Рабочий процесс можно сочетать с автономным агентом: программа задаёт
критические этапы, условия проверки и полномочия, а внутри узлов,
требующих гибких решений, работает автономный агент. По сплошной линии в
верхней части \crossreflink{fig:1-7}{рисунка 1-7} видно, что сообщение
Telegram сначала проходит определение типа, голосовое сообщение
предварительно преобразуется в текст, а затем ответ отправляется через
узел Telegram. Пунктирная линия под узлом AI Agent показывает, что к
нему подключены такие инструменты, как Gmail и Google Calendar, а
конкретный инструмент модель выбирает в зависимости от содержимого
сообщения\footnote{n8n, документация узла AI Agent.}.

\begin{figure}
\hypertarget{fig:1-7}{%
\centering
\includegraphics{images/revised/n8n-workflow-grayscale.pdf}
\caption{Рисунок 1-7. Интерфейс редактора рабочих процессов
n8n}\label{fig:1-7}
}
\end{figure}

Другой подход --- \textbf{сначала поручить автономному агенту создать
рабочий процесс, а затем выполнить проверенный рабочий процесс}. Агент
проектирует топологию задачи и генерирует код оркестрации. Исполняющая
система проверяет код, полномочия и входные данные, после чего выполняет
сгенерированный путь. Для повторяющихся данных или шагов это сокращает
затраты на вызов модели при каждой диспетчеризации, сохраняя гибкость
динамического проектирования процесса в начале задачи. Эта форма будет
подробнее рассмотрена в \crossreflink{chap:10}{главе 10}.

\begin{quote}
\textbf{Справочник инструментов: фреймворки и платформы агентов}

Фреймворки различаются степенью инкапсуляции цикла выполнения,
управления состоянием и подключения инструментов.\footnote{Возможности и
  интерфейсы фреймворков описаны в официальных материалах: Claude Agent
  SDK, LangGraph, узлы Agent в Dify, CrewAI, OpenClaw, DeepSeek Harness
  и Pi. Узел Agent в n8n описан в приведённом выше источнике.}
\end{quote}

\begin{longtable}[]{@{}
  >{\raggedright\arraybackslash}p{(\columnwidth - 8\tabcolsep) * \real{0.1864}}
  >{\raggedright\arraybackslash}p{(\columnwidth - 8\tabcolsep) * \real{0.2243}}
  >{\raggedright\arraybackslash}p{(\columnwidth - 8\tabcolsep) * \real{0.1690}}
  >{\raggedright\arraybackslash}p{(\columnwidth - 8\tabcolsep) * \real{0.1900}}
  >{\raggedright\arraybackslash}p{(\columnwidth - 8\tabcolsep) * \real{0.2303}}@{}}
\toprule\noalign{}
\begin{minipage}[b]{\linewidth}\raggedright
\hspace{0pt}Фреймворк/\allowbreak{}платформа
\end{minipage} & \begin{minipage}[b]{\linewidth}\raggedright
\hspace{0pt}Основное назначение
\end{minipage} & \begin{minipage}[b]{\linewidth}\raggedright
\hspace{0pt}Способ оркестрации
\end{minipage} & \begin{minipage}[b]{\linewidth}\raggedright
\hspace{0pt}Способ разработки
\end{minipage} & \begin{minipage}[b]{\linewidth}\raggedright
\hspace{0pt}Типичное применение
\end{minipage} \\
\midrule\noalign{}
\endhead
\bottomrule\noalign{}
\endlastfoot
\hspace{0pt}\textbf{Codex Harness} & \hspace{0pt}Встраиваемая среда
выполнения агента & \hspace{0pt}Автономный цикл & \hspace{0pt}Интеграция
через код и интерфейсы & \hspace{0pt}Coding Agent, встраиваемый в
продукт агент \\
\hspace{0pt}\textbf{Claude Agent SDK} & \hspace{0pt}Переиспользование
инструментов и цикла выполнения Claude Code & \hspace{0pt}Автономный
цикл & \hspace{0pt}Python, TypeScript & \hspace{0pt}Задачи
программирования, сложные задачи с инструментами \\
\hspace{0pt}\textbf{LangChain /\allowbreak{} LangGraph} &
\hspace{0pt}Интеграция моделей и инструментов, оркестрация графов
состояния & \hspace{0pt}Рабочие процессы и автономные циклы &
\hspace{0pt}Код & \hspace{0pt}Многошаговые задачи, постоянное состояние
и восстановление \\
\hspace{0pt}\textbf{n8n} & \hspace{0pt}Платформа визуальной
автоматизации & \hspace{0pt}Рабочие процессы и узлы Agent &
\hspace{0pt}Визуальная разработка, узлы с кодом &
\hspace{0pt}Автоматизация бизнеса, подключение сервисов \\
\hspace{0pt}\textbf{Dify} & \hspace{0pt}Платформа разработки приложений
на основе LLM & \hspace{0pt}Рабочие процессы, диалоги и узлы Agent &
\hspace{0pt}Визуальная разработка и API & \hspace{0pt}RAG, базы знаний и
бизнес-\allowbreak{}\hspace{0pt}приложения \\
\hspace{0pt}\textbf{CrewAI} & \hspace{0pt}Команды из нескольких агентов
и оркестрация процессов & \hspace{0pt}Crews и Flows & \hspace{0pt}Код &
\hspace{0pt}Разделение ролей, совместная работа и
бизнес-\allowbreak{}\hspace{0pt}процессы \\
\hspace{0pt}\textbf{OpenClaw} & \hspace{0pt}Самостоятельно размещаемый
персональный помощник & \hspace{0pt}Автономный цикл и событийное
управление & \hspace{0pt}Конфигурация и расширения & \hspace{0pt}Обмен
сообщениями на разных платформах, личные задачи и интеграция
инструментов \\
\hspace{0pt}\textbf{DeepSeek Harness} & \hspace{0pt}Расширяемая
плагинами среда выполнения агента & \hspace{0pt}Комбинируемые циклы и
плагины & \hspace{0pt}Код и конфигурация & \hspace{0pt}Настройка
механизма выполнения, разработка и исследование агентов \\
\hspace{0pt}\textbf{Pi} & \hspace{0pt}Минималистичный набор инструментов
для агентов и Coding Agent & \hspace{0pt}Автономный цикл и расширения &
\hspace{0pt}Код & \hspace{0pt}Пользовательские инструменты, циклы
выполнения и помощники программиста \\
\end{longtable}

\begin{quote}
На примере Codex: разработчик может ставить задачи непосредственно в
настольном приложении, командной строке или редакторе. Для встраивания
таких возможностей в собственный продукт необходимо решить, как
запускать среду выполнения, принимать события и обрабатывать
согласования. Существует три распространённых способа интеграции:
\texttt{codex\ exec} подходит для неинтерактивных задач в скриптах и CI;
Codex SDK позволяет запускать и возобновлять задачи из кода приложения,
а также обрабатывать их события; app-server поддерживает сеансы, потоки
событий и интерактивные согласования посредством двунаправленной связи
по JSON-RPC, поэтому подходит для создания собственных
клиентов\footnote{Официальная документация OpenAI: Codex SDK, Codex App
  Server.}. Claude Agent SDK, в свою очередь, предоставляет приложениям
доступ к инструментам, циклу выполнения и управлению контекстом Claude
Code; через SDK разработчик может настраивать полномочия, сеансы и хуки.
\end{quote}

Определив, как должна продвигаться задача, необходимо также проверять,
можно ли выполнить каждый шаг и соответствует ли результат требованиям.
Рассматриваемые далее защитные механизмы обеспечивают соблюдение
ограничений, проверку и исправление именно в этих точках.

\hypertarget{ux437ux430ux449ux438ux442ux43dux44bux435-ux431ux430ux440ux44cux435ux440ux44b-ux438-ux431ux435ux437ux43eux43fux430ux441ux43dux43eux441ux442ux44c}{%
\subsection{Защитные барьеры и
безопасность}\label{ux437ux430ux449ux438ux442ux43dux44bux435-ux431ux430ux440ux44cux435ux440ux44b-ux438-ux431ux435ux437ux43eux43fux430ux441ux43dux43eux441ux442ux44c}}

\textbf{Защитные барьеры} (Guardrails) воплощают ограничения, проверку и
исправление на границах системы, помогая контролировать раскрытие
конфиденциальной информации, несанкционированные операции и вывод, не
соответствующий требованиям продукта. При проектировании сначала следует
выявить риски задачи: какие данные будет читать модель, какие
инструменты сможет вызывать и какие действия повлекут внешние
последствия, --- а затем установить в этих точках соответствующие
проверки и средства контроля.

При оценке защитных барьеров необходимо учитывать как пропуск опасных
запросов, так и \textbf{ложные отказы}. В первом случае проверяется, не
разрешены ли опасные запросы, а во втором --- не блокируются ли по
ошибке явно разрешённые задачи, например авторизованное тестирование
безопасности или исследование дистилляции моделей, соответствующее
правилам сервиса.

\hypertarget{ux442ux438ux43fux44b-ux437ux430ux449ux438ux442ux43dux44bux445-ux431ux430ux440ux44cux435ux440ux43eux432}{%
\subsubsection{Типы защитных
барьеров}\label{ux442ux438ux43fux44b-ux437ux430ux449ux438ux442ux43dux44bux445-ux431ux430ux440ux44cux435ux440ux43eux432}}

По месту действия защиты в этой книге выделяются \textbf{контекстный
уровень, исполнительный уровень и уровень данных}. Инженерная ценность
такого разделения состоит в постепенном снижении зависимости
безопасности от решений модели: поступающая в контекст информация
проверяется по источнику и содержимому, реальные действия дополнительно
ограничиваются исполнительными разрешениями, а итоговые операции чтения
и записи проходят через правила уровня данных. Совместная работа трёх
уровней позволяет обнаружить на последующих границах проблему,
пропущенную одной из проверок.

\textbf{Контекстный уровень} определяет, что читает модель и как
различаются источники информации. К распространённым механизмам
относятся:

\begin{itemize}
\tightlist
\item
  \textbf{Классификация релевантности}: выявление запросов,
  отклоняющихся от задач продукта. Например, если помощнику по
  программированию задают вопрос «Какова высота Эмпайр-стейт-билдинг?»,
  он отвечает или предлагает пользователю вернуться к задаче в
  зависимости от области применения продукта.
\item
  \textbf{Классификация безопасности}: выявление джейлбрейк-запросов,
  побуждающих модель нарушить правила безопасности, а также инъекций в
  промпт, которые пытаются превратить входные данные в управляющие
  инструкции. Инъекция в промпт может содержаться как во вводе
  пользователя, так и косвенно попадать в контекст через веб-страницы,
  документы или результаты инструментов.
\item
  \textbf{Модерация контента}: проверка содержимого на чувствительные
  темы, включая насилие и дискриминацию, с учётом назначения и правил, а
  также различий между обсуждением, анализом и недопустимым применением.
\item
  \textbf{Проверка правил и маркировка источников}: ограничение длины
  ввода, проверка формата полей, маркировка недоверенных источников и
  разделение инструкций и данных в контексте. Для устранения риска
  SQL-инъекций при обращении к базе данных следует также применять
  механизмы исполнительного уровня, например параметризованные запросы.
\end{itemize}

\textbf{Критически важные для безопасности проверки должны выполняться
за пределами контекста модели.} Если недоверенное содержимое уже
повлияло на решения модели, самопроверка той же моделью в том же
контексте с большой вероятностью продолжит опираться на
скомпрометированные предпосылки. Маркировка источников и классификаторы
помогают уменьшить этот риск, а независимые проверки разрешений и
выполнения ограничивают его реальные последствия. Организация контекста
и защита от инъекций в промпт будут подробно рассмотрены во второй
главе.

\textbf{Constitutional Classifiers} от Anthropic демонстрируют другой
способ организации классификаторов: синтетические данные генерируются по
правилам на естественном языке, после чего на них обучаются механизмы
проверки входа и вывода. В последующей системе вопрос и ответ
оцениваются совместно, благодаря чему классификатор может понимать смысл
с учётом диалога. Например, нейтральное само по себе выражение «пищевая
приправа» в контексте может фактически обозначать химический реагент.
Чтобы снизить вычислительные затраты, система сначала выполняет отсев с
помощью лёгкого зонда, считывающего внутренние активации, а затем
передаёт подозрительные диалоги на повторную проверку набору более
мощных классификаторов. Первый уровень отвечает за отбор, а
окончательное решение принимается на этапе повторной проверки, что
помогает сбалансировать стоимость и долю ложных отказов\footnote{Anthropic,
  «Next-generation Constitutional Classifiers», 9 января 2026 года;
  Cunningham et al., «Constitutional Classifiers++».}.

\textbf{Исполнительный уровень} определяет, могут ли предложенные
моделью операции вступить в силу. Вызовы инструментов можно разделить на
операции низкого, среднего и высокого риска с учётом обратимости, объёма
разрешений, финансовых последствий и чувствительности данных, а затем
динамически проверять их фактические параметры. Чтение общедоступных
материалов и изменение данных учётной записи требуют разных мер
контроля; для операций высокого риска можно предусмотреть дополнительную
независимую проверку или подтверждение человеком.

Исполнительные проверки должны обеспечиваться механизмами с независимыми
границами разрешений: например, доверенным процессом проверки, учётными
данными с минимальными привилегиями, песочницей или ручным
согласованием. Даже если на модель повлияло недоверенное содержимое,
фактический вызов должен пройти эти проверки. Ответ пользователю также
является действием: фильтр PII --- персональных идентификационных данных
--- может выявлять такие конфиденциальные поля, как номер удостоверения
личности или телефона, а проверка вывода --- контролировать соответствие
содержимого назначению продукта и требованиям бренда. Риски инструментов
и исполнительные разрешения будут рассмотрены в четвёртой главе.

\textbf{Уровень данных} непосредственно контролирует доступ к ресурсам и
их изменение. Политики построчной безопасности базы данных определяют,
какие записи доступны текущему субъекту; ограничения и валидаторы
проверяют допустимость записи; контролируемые представления и хранимые
процедуры ограничивают разрешённые операции; доверенная среда выполнения
привязывает идентичность и контекст доступа. Если фактические операции
чтения и записи проходят через эти механизмы, а разрешения учётных
данных среды выполнения ограничены, уровень данных сможет отклонить
несанкционированный доступ, даже если сгенерированный на верхнем уровне
код пропустил проверку. В пятой главе на примере динамически
создаваемого программного обеспечения будет показано, как перенести
границу доверия на этот уровень.

\hypertarget{ux443ux447ux430ux441ux442ux438ux435-ux447ux435ux43bux43eux432ux435ux43aux430}{%
\subsubsection{Участие
человека}\label{ux443ux447ux430ux441ux442ux438ux435-ux447ux435ux43bux43eux432ux435ux43aux430}}

\textbf{Человек в контуре} (human-in-the-loop, HITL) означает, что на
определённых этапах система просит человека принять решение или передаёт
ему управление. Например, помощник может самостоятельно сверить счёт,
отправить пользователю два тарифных плана, между которыми необходимо
выбрать, и разницу в цене, а после получения ответа оформить выбранный
вариант. Каждый дополнительный запрос требует внимания пользователя;
необходимость такой остановки зависит от последствий неверного выбора и
от того, предоставил ли пользователь достаточно однозначное разрешение.
На ранних этапах развёртывания обратная связь от людей также помогает
команде выявлять режимы отказа, дополнять набор пограничных случаев и
совершенствовать оценку. Система клиентского обслуживания может передать
обращение оператору, а Coding Agent --- вернуть разработчику нерешённую
проблему и текущие изменения.

Вмешательство человека особенно уместно в двух случаях. Первый ---
\textbf{превышение порога неудач}, например достижение предельного числа
повторных попыток или операций. Второй --- \textbf{операции высокого
риска}, например возврат крупной суммы, платёж или изменение данных,
которое трудно отменить. Условия срабатывания следует закрепить в
правилах выполнения, чтобы система могла своевременно остановиться,
сохранить состояние и передать управление.

При передаче управления следует сообщить исходную цель, выполненные
шаги, текущий результат, причину неудачи и вопросы, требующие решения.
Если пользователь временно недоступен, задачу можно сохранить и
перевести в режим ожидания либо продолжить самостоятельную работу
низкого риска по заранее установленным правилам; операции, требующие
разрешения, должны оставаться в состоянии ожидания. Благодаря этому
человек сможет сразу определить следующий шаг, а выполнение задачи можно
будет возобновить из сохранённого состояния.

\hypertarget{ux441ux43eux43eux442ux432ux435ux442ux441ux442ux432ux438ux435-ux43fux44fux442ux438-ux444ux443ux43dux43aux446ux438ux439-harness-ux441ux442ux440ux443ux43aux442ux443ux440ux435-ux43aux43dux438ux433ux438}{%
\subsection{Соответствие пяти функций Harness структуре
книги}\label{ux441ux43eux43eux442ux432ux435ux442ux441ux442ux432ux438ux435-ux43fux44fux442ux438-ux444ux443ux43dux43aux446ux438ux439-harness-ux441ux442ux440ux443ux43aux442ux443ux440ux435-ux43aux43dux438ux433ux438}}

\textbf{«LLM + контекст + инструменты» описывает функциональные
компоненты, а «Model + Harness» --- производственную реализацию.} Первая
формула показывает, откуда берутся возможности, вторая --- как работает
система.

В таблице ниже перечислены функции Harness, рассматриваемые в
соответствующих главах, и связанные с ними риски.

\begin{longtable}[]{@{}
  >{\raggedright\arraybackslash}p{(\columnwidth - 6\tabcolsep) * \real{0.2059}}
  >{\raggedright\arraybackslash}p{(\columnwidth - 6\tabcolsep) * \real{0.1400}}
  >{\raggedright\arraybackslash}p{(\columnwidth - 6\tabcolsep) * \real{0.3238}}
  >{\raggedright\arraybackslash}p{(\columnwidth - 6\tabcolsep) * \real{0.3303}}@{}}
\toprule\noalign{}
\begin{minipage}[b]{\linewidth}\raggedright
\hspace{0pt}Функция Harness
\end{minipage} & \begin{minipage}[b]{\linewidth}\raggedright
\hspace{0pt}Основные главы
\end{minipage} & \begin{minipage}[b]{\linewidth}\raggedright
\hspace{0pt}Техническое содержание
\end{minipage} & \begin{minipage}[b]{\linewidth}\raggedright
\hspace{0pt}Аспекты безопасности и надёжности
\end{minipage} \\
\midrule\noalign{}
\endhead
\bottomrule\noalign{}
\endlastfoot
\hspace{0pt}Управление контекстом &
\hspace{0pt}\crossreflink{chap:2}{Глава 2} & \hspace{0pt}Проектирование
промптов, строка состояния, сжатие контекста, Agent Skills &
\hspace{0pt}Инъекция в промпт, загрязнение контекста, потеря
прогресса \\
\hspace{0pt}Управление контекстом между сессиями &
\hspace{0pt}\crossreflink{chap:3}{Глава 3} & \hspace{0pt}Память
пользователя, RAG, структурированные индексы, агентный RAG &
\hspace{0pt}Раскрытие конфиденциальной информации, источники и
приватность \\
\hspace{0pt}Интерфейсы и ограничения инструментов &
\hspace{0pt}\crossreflink{chap:4}{Глава 4} & \hspace{0pt}Классификация
инструментов, управление разрешениями, MCP, активное обнаружение
инструментов & \hspace{0pt}Ошибочные действия, несанкционированный
доступ, необратимые операции \\
\hspace{0pt}Верификация & \hspace{0pt}\crossreflink{chap:5}{Главы 5} и 7
& \hspace{0pt}Разработка через тестирование, правила в виде кода, оценка
результатов работы & \hspace{0pt}Границы идентичности и ответственности,
ошибочная оценка результатов \\
\hspace{0pt}Коррекция & \hspace{0pt}\crossreflink{chap:2}{Главы 2} и 5 &
\hspace{0pt}Восстановление контекста, безопасные повторные попытки,
проверка состояния, откат и передача управления человеку &
\hspace{0pt}Повторное выполнение, сбой восстановления, распространение
ошибок \\
\end{longtable}

\crossreflink{chap:6}{Глава 6} дополнительно расширяет модальности и
временные режимы взаимодействия, включая изображения, речь, работу с
интерфейсом и асинхронные события. Для реализации этих возможностей
необходимы соответствующие контекст и интерфейсы инструментов. В
\crossreflink{chap:7}{главе 7} оценивается работа агента в целом, в
\crossreflink{chap:8}{главах 8} и 9 рассматривается улучшение его
возможностей посредством обновления соответственно параметров и внешних
результатов работы, а \crossreflink{chap:10}{глава 10} посвящена
распределению задач, коммуникации и координации нескольких агентов.

Безопасность проходит сквозной темой через все эти главы. Управление на
уровне контекста и выполнения в основном организует Harness, а контроль
доступа на уровне данных действует со стороны окружения. В каждой главе
на примере конкретных технологий объясняется, где следует размещать эти
средства управления и как организовать их совместную работу.

Пример Anthropic с длительно работающим AI-агентом демонстрирует
взаимодействие пяти функций: \textbf{агент инициализации} настраивает
окружение и разбивает работу на список задач, после чего \textbf{агент
выполнения} последовательно решает их, а в конце каждой сессии оставляет
запись о прогрессе и результаты, с которыми можно продолжить работу.
Новая сессия считывает эти записи и продолжает работу, а тесты и список
задач помогают системе определить, какие функции ещё не реализованы.
Такая организация смягчает проблемы передачи работы между контекстами и
преждевременного объявления о завершении\footnote{Anthropic, Effective
  harnesses for long-running agents, 26 ноября 2025 года.}.
\crossreflink{chap:2}{Глава 2} начинается с организации контекста, а в
\crossreflink{chap:5}{главе 5} подробнее рассматривается полный Harness
агента для программирования.

\hypertarget{ux441ux43aux432ux43eux437ux43dux44bux435-ux43fux430ux442ux442ux435ux440ux43dux44b-ux43fux440ux43eux435ux43aux442ux438ux440ux43eux432ux430ux43dux438ux44f}{%
\section{Сквозные паттерны
проектирования}\label{ux441ux43aux432ux43eux437ux43dux44bux435-ux43fux430ux442ux442ux435ux440ux43dux44b-ux43fux440ux43eux435ux43aux442ux438ux440ux43eux432ux430ux43dux438ux44f}}

Представим, что AI-агент должен обновить политику службы поддержки.
Проверяющий должен иметь возможность найти исходный текст изменённых
положений; при работе модели ей достаточно сначала увидеть каталог
политик; после обновления также должна сохраняться возможность найти
старую версию и убедиться, что вопросы, на которые система прежде
отвечала правильно, по-прежнему обрабатываются верно. Уже в этой задаче
проявляются несколько приведённых ниже паттернов проектирования.

\textbf{Предлагающий --- проверяющий (Proposer--Reviewer)} поручает
создание результата и его оценку двум независимым ролям. В примере с
обновлением политики службы поддержки один агент готовит проект
изменений, а другой сверяет их с исходными положениями. Способ проверки
зависит от представленного результата: при проверке слайдов можно
изучить итог рендеринга, при проверке кода --- вывод тестов, а при
проверке вызовов инструментов --- их параметры. Проверяющий должен
понимать контекст задачи и выносить решение в независимом контексте,
чтобы выводы создающей стороны меньше влияли на проверку. При
самопроверке модель склонна сохранять прежние предположения, тогда как
независимая проверка повышает вероятность обнаружения слепых зон;
границы полномочий и детерминированные проверки при этом продолжают
ограничивать выполнение. В \crossreflink{chap:3}{главе 3} этот паттерн
используется для обновления знаний, в \crossreflink{chap:4}{главе 4} ---
для проверки вызовов инструментов; там же представлен его вариант
Sidecar, работающий только для чтения. Этот паттерн также применяется в
примерах с PPT, видео и журналами в \crossreflink{chap:5}{главе 5}, при
оценке UI в \crossreflink{chap:7}{главе 7}, при пересмотре правил в
\crossreflink{chap:9}{главе 9} и при равноправном взаимодействии в
\crossreflink{chap:10}{главе 10}.

\textbf{Постепенное раскрытие (Progressive Disclosure)} сначала
предоставляет доступный для поиска каталог или сводку, а затем по мере
необходимости загружает подробности. Это позволяет контролировать
затраты на контекст и помогает модели сначала сузить область выбора. В
Agent Skills из \crossreflink{chap:2}{главы 2} метаданные постоянно
находятся в контексте, а основное содержимое загружается по требованию;
аналогичный подход используется при многоуровневом поиске в
\crossreflink{chap:3}{главе 3}, активном обнаружении инструментов и
усечении при постраничной выдаче в \crossreflink{chap:4}{главе 4}, а
также при обнаружении агентов в \crossreflink{chap:10}{главе 10}.

\textbf{Только добавление (Append-only)} выражает изменения посредством
добавления новых записей, сохраняя историю для кэширования, повторного
воспроизведения и аудита. В \crossreflink{chap:2}{главе 2} стабильный
префикс позволяет повторно использовать KV-кэш; чем ближе изменение к
началу, тем больше последующих кэшированных данных оно обычно
затрагивает. Событийная память из \crossreflink{chap:3}{главы 3}
фиксирует изменения состояния посредством новых событий, а добавление
схемы инструмента в последующий контекст по мере необходимости в
\crossreflink{chap:4}{главе 4} также использует идею сохранения
существующего префикса. В \crossreflink{chap:6}{главе 6} события из
разных источников добавляются в единый журнал приёма данных для
последующей диспетчеризации, повторного воспроизведения и расследования
проблем. Когда требуется сжатие или обновление состояния, можно создать
новую версию, сохранив записи, необходимые для отслеживания истории.

\textbf{Граничный набор + набор сохранения (Boundary Set + Retention
Set)} проверяет изменение с двух сторон: один набор содержит проблемы,
которые требуется исправить этим изменением, а другой --- задачи, с
которыми предыдущая версия уже справлялась правильно. Например, после
исправления определённого типа ошибок в параметрах инструментов
необходимо как повторно проверить примеры с этими ошибками, так и
убедиться, что ранее выполнявшиеся задачи по-прежнему работают. Если
проверять только целевые примеры, переобучение на них легко принять за
улучшение; если проверять лишь прежние возможности, можно одобрить
изменение, не дающее никакого эффекта. Подобное сопоставление
используется при регрессионной оценке в \crossreflink{chap:7}{главе 7},
разделении обучения и оценки в \crossreflink{chap:8}{главе 8} и проверке
обновлений в \crossreflink{chap:9}{главе 9}, позволяя одновременно
контролировать улучшения и ухудшения.

\textbf{Минимальный diff + возможность отката} требует, чтобы область
каждого изменения была чётко определена, само изменение имело указанный
источник, а систему можно было вернуть к предыдущей версии. Небольшие
последовательные изменения позволяют связать эффект с конкретной правкой
и упрощают её отмену в случае неудачи. Этому принципу следуют обновление
знаний в \crossreflink{chap:3}{главе 3}, исправления кода в
\crossreflink{chap:5}{главе 5}, а также обновления промптов и программ в
\crossreflink{chap:9}{главе 9}. Для контекста, внешних артефактов и
параметров модели предусмотрены свои способы восстановления: сохранение
снимка входных данных, версии артефакта или контрольной точки модели
позволяет восстановить соответствующее состояние; уже возникшие внешние
побочные эффекты дополнительно требуют механизма компенсации.

\hypertarget{ux438ux442ux43eux433ux438-ux433ux43bux430ux432ux44b}{%
\section{Итоги
главы}\label{ux438ux442ux43eux433ux438-ux433ux43bux430ux432ux44b}}

\textbf{Возможности AI-агента возникают благодаря совместной работе LLM,
контекста и инструментов.} Вернёмся к примеру с авиабилетом: цена в 1250
юаней --- это наблюдение, полученное от оператора, а цель командировки и
предыдущие отзывы пользователя образуют текущий контекст. На их основе
модель решает, принять предложение или продолжить уточнения. Ассистент
продолжает торговаться по телефону; чтобы сообщить пользователю «билет
забронирован», ему также необходимо получить подтверждение успешного
бронирования.

\textbf{Расширение контекста и инструментов --- один из основных
инженерных способов повышения возможностей агента.} При фиксированной
модели добавление ключевых материалов и предоставление необходимых
операций часто позволяет сделать выполнимыми задачи, которые прежде
выполнить было невозможно. Примеры Manus и OpenClaw демонстрируют такое
расширение; инструменты, параметры, права доступа и среда выполнения
вместе задают границы допустимых действий.

\textbf{Организация контекста влияет на каждое решение.} Базовая
реализация передаёт модели относительно стабильный префикс и
накапливаемую историю сообщений. В экспериментах с абляцией определения
инструментов сообщают модели, какие возможности ей доступны, результаты
работы инструментов предоставляют факты, необходимые для вычислений, а
история сообщений сохраняет уже выполненные шаги. Если имеющихся записей
достаточно для восстановления оснований принятого решения, модель может
продолжить выполнение задачи и после удаления истории рассуждений. ReAct
связывает рассуждения, действия и новые наблюдения, позволяя шаг за
шагом продвигаться к выполнению задачи.

\textbf{Harness --- ключ к накоплению инженерного конкурентного
преимущества продукта на основе агента.} Когда все рассматриваемые
модели удовлетворяют базовым требованиям к возможностям, именно
управление контекстом, интерфейсы инструментов, ограничения, проверка и
коррекция определяют, сможет ли продукт стабильно и надёжно обеспечивать
результат.

\textbf{Оркестрацию следует начинать с простого решения и усложнять по
мере необходимости.} В первую очередь стоит попробовать однократный
вызов; если требуется несколько фиксированных шагов, следует
использовать рабочий процесс, а автономного агента вводить только тогда,
когда путь необходимо корректировать на основе обратной связи. Эти три
подхода также можно комбинировать.

Что делать, если ассистент начинает обрабатывать заказ, на который у
него изначально не было прав? Сначала нужно провести проверку по трём
направлениям: источники данных, права инструментов и изменения
состояния. Из каких записей модель узнала о заказе? Почему уровень
выполнения разрешил эту операцию? Остановил ли её контроль доступа в
среде? Столкнувшись с высокорисковой операцией или действием за
пределами предоставленных полномочий, ассистент может остановиться и
передать задачу человеку. Предположим, проверка показала, что одно из
правил авторизации сформулировано слишком широко. Тогда предлагающий
может изменить только это правило. Проверяющий формирует граничный набор
из заказов, обрабатываемых с превышением полномочий, и набор сохранения
из обычных заказов, а затем проверяет, блокируются ли после изменения
недопустимые операции и продолжают ли выполняться допустимые. Правила
можно загружать по мере необходимости, а изменения следует добавлять в
журнал истории, чтобы при необходимости можно было восстановить прежнюю
версию. Эти приёмы будут неоднократно использоваться в последующих
главах.

Вернёмся к разговору о бронировании билета. Когда оператор называет цену
в 1250 юаней, ассистент должен одновременно видеть предложение, бюджет
пользователя и ранее согласованные условия. Если разговор длится уже
полчаса, какие сообщения следует сохранить для следующего шага, а какую
информацию можно сохранить отдельно и извлечь при необходимости? Во
второй главе ответ на этот вопрос начинается с рассмотрения фактических
входных данных одного вызова модели. Связь между агентом в обучении с
подкреплением и рассматриваемым здесь агентом на основе LLM будет
подробнее раскрыта в восьмой главе.

\begin{questionbox}

\hypertarget{ux432ux43eux43fux440ux43eux441ux44b-ux434ux43bux44f-ux440ux430ux437ux43cux44bux448ux43bux435ux43dux438ux44f}{%
\section*{Вопросы для
размышления}\label{ux432ux43eux43fux440ux43eux441ux44b-ux434ux43bux44f-ux440ux430ux437ux43cux44bux448ux43bux435ux43dux438ux44f}}
\addcontentsline{toc}{section}{Вопросы для размышления}

\begin{enumerate}
\def\labelenumi{\arabic{enumi}.}
\tightlist
\item
  ★★ Если приоритетно можно улучшить только что-то одно --- модель,
  контекст или инструменты, --- как определить, куда выгоднее всего
  вложить усилия? Какие результаты оценки заставят вас изменить выбор?
\item
  ★★★ Предположим, что в каждом раунде добавляется примерно одинаковый
  объём сообщений и модели каждый раз передаётся полная предшествующая
  история. Почему совокупный объём входных данных растёт примерно
  пропорционально квадрату числа раундов? Какие составляющие затрат
  изменяют кэширование, сжатие и обрезка истории?
\item
  ★★ Модель принимает всё больше решений о вызове инструментов, а
  Harness должен поддерживать всё более длительные и сложные задачи. Как
  будут меняться обязанности модели и Harness? Какие механизмы можно
  упростить, а какие необходимо усилить?
\item
  ★★ В эксперименте отсутствие результатов инструментов приводило к
  повторным вызовам. Какие ещё причины могут вызывать подобные циклы в
  промышленной системе? Как их обнаруживать, восстанавливать работу и
  задавать условия завершения?
\item
  ★ Выберите AI-продукт, которым пользуетесь ежедневно, и отдельно
  перечислите основные сведения в его контексте, доступные операции и
  стратегию выполнения. На каком этапе он наиболее подвержен сбоям? Как
  бы вы его улучшили?
\item
  ★★ Разработайте процесс для службы поддержки по бронированию
  авиабилетов: какие этапы лучше реализовать в виде рабочего процесса, а
  какие поручить автономному агенту? Как эти две части должны
  взаимодействовать, если пользователь внезапно изменит маршрут?
\item
  ★★★ Один и тот же инструмент \texttt{delete\_file} несёт разный риск
  при удалении временного и критически важного системного файла. Как бы
  вы выполняли динамическую оценку риска с учётом параметров, области
  затрагиваемых ресурсов и возможности восстановления?
\item
  ★★ Открытое пространство действий допускает сочетания содержимого и
  операций, которые трудно исчерпывающе перечислить заранее. В каких
  задачах уместнее ограничить действия небольшим набором
  предопределённых вариантов? Какие преимущества и издержки создаёт
  такое ограничение?
\item
  ★★ Когда агенту требуется подтверждение человека, пользователь может
  быть не в сети, отвечать медленно или давать неоднозначные указания.
  Какое состояние должна сохранять система, какую работу можно
  продолжить, а какие операции обязательно должны ожидать?
\item
  ★★★ Приведите пример инженерного метода, который будет упрощаться по
  мере совершенствования моделей. Объясните, какую проблему он решал
  изначально и как лежащие в его основе принципы проектирования систем
  продолжат приносить пользу.
\end{enumerate}

\textbf{Практикум к \crossreflink{chap:1}{главе 1}.}
\href{https://github.com/bojieli/ai-agent-book/tree/40f4ce5778e6cc93e3606c3079bf8cf13a0c4d04/chapter1}{Каталог
экспериментов} ·
\href{https://github.com/bojieli/ai-agent-book/blob/40f4ce5778e6cc93e3606c3079bf8cf13a0c4d04/chapter1/README.ru.md}{Путеводитель
на русском}.

\end{questionbox}

\hypertarget{ux438ux43dux436ux435ux43dux435ux440ux438ux44f-ux43aux43eux43dux442ux435ux43aux441ux442ux430}{%
\chapter{\texorpdfstring{Инженерия
контекста\label{chap:2}}{Инженерия контекста}}\label{ux438ux43dux436ux435ux43dux435ux440ux438ux44f-ux43aux43eux43dux442ux435ux43aux441ux442ux430}}

В первой главе контекст сравнивался с «текущим полем принятия решений»
AI-агента: всякий раз, когда модель принимает решение, она опирается на
информацию, полученную в данный момент. \textbf{Инженерия контекста}
(Context Engineering) отвечает за отбор, организацию и обновление этой
информации, включая системные инструкции, определения инструментов,
пользовательские запросы, историю сообщений и состояние задачи. Перед
каждым вызовом модели Harness должен подготовить сведения о задаче и
последних изменениях в ходе её выполнения в качестве входных данных. На
\crossreflink{fig:2-1}{рисунке 2-1} на конкретном содержимом показана
структура контекстного окна.

\begin{figure}
\hypertarget{fig:2-1}{%
\centering
\includegraphics{images/revised/fig2-1.pdf}
\caption{Рисунок 2-1. Общая структура контекстного окна}\label{fig:2-1}
}
\end{figure}

\hypertarget{ux43aux43eux43dux442ux435ux43aux441ux442-ux43fux440ux438ux43cux435ux43dux435ux43dux438ux435-ux432ux43eux437ux43cux43eux436ux43dux43eux441ux442ux435ux439-ux43cux43eux434ux435ux43bux438-ux43a-ux43aux43eux43dux43aux440ux435ux442ux43dux43eux439-ux437ux430ux434ux430ux447ux435}{%
\section{Контекст: применение возможностей модели к конкретной
задаче}\label{ux43aux43eux43dux442ux435ux43aux441ux442-ux43fux440ux438ux43cux435ux43dux435ux43dux438ux435-ux432ux43eux437ux43cux43eux436ux43dux43eux441ux442ux435ux439-ux43cux43eux434ux435ux43bux438-ux43a-ux43aux43eux43dux43aux440ux435ux442ux43dux43eux439-ux437ux430ux434ux430ux447ux435}}

Представьте, что к команде только что присоединился опытный инженер и
получил первое задание: «Изменить конфигурацию в соответствии с нашими
правилами». Он умеет программировать, но ещё не знает, какой файл служит
источником конфигурации, какие поля можно изменять и как проверять
результат. Даже получив весь репозиторий, он может изменить не тот файл.
AI-агент, приступающий к работе с незнакомым проектом, испытывает такой
же недостаток исходной информации.

Например, для выполнения просьбы «Помоги исправить этот баг» агенту для
работы с кодом обычно требуются три вида исходной информации:

\begin{itemize}
\tightlist
\item
  \textbf{Актуальный контекст кода}: структура каталогов, назначение
  модулей, основные структуры данных и правила оформления кода.
  Например, если изменить автоматически сгенерированный конфигурационный
  файл, изменение сработает локально, но будет перезаписано при
  следующей генерации. Только зная, какой модуль создаёт этот файл,
  можно найти правильное место для изменения.
\item
  \textbf{Правила процесса}: стратегия ветвления Git, требования к
  коммитам, процесс ревью кода и требования CI/CD. На их основе агент
  планирует изменение, тестирование и коммит, а исполняющая система
  выполняет эти операции в пределах предоставленных прав.
\item
  \textbf{Информация об окружении}: конфигурация среды разработки,
  способ подключения к тестовой базе данных, порядок развёртывания
  тестового окружения и правила управления ключами API. Эти сведения
  помогают выбрать правильное окружение и понять различия между
  локальной и тестовой средами.
\end{itemize}

Код проекта и тестовая база данных находятся за пределами контекста
модели. В примере с автоматически сгенерированной конфигурацией агент
лишь после чтения генерирующего скрипта с помощью инструмента поиска
узнаёт, с чего следует начинать изменение. \textbf{Качество контекста
определяет, сможет ли модель применить имеющиеся у неё способности к
текущей задаче.} \textbf{В задачах, требующих обширной исходной
информации, достаточно способная модель с полным и точным контекстом
может показать лучший результат, чем более мощная модель без необходимых
материалов.}

\textbf{Создание AI-нативной команды начинается с документирования
ключевых знаний.} Если причины выбора определённой базы данных помнят
только старые сотрудники, а изменения бизнес-правил обсуждались лишь в
личной переписке, новые коллеги и агенты столкнутся с одной и той же
проблемой: из кода видно, что сделала команда, но невозможно понять,
почему она это сделала.

\textbf{Команды, хорошо приспособленные к удалённой работе, зачастую
удобны и для агентов.} Разработчики ядра Linux обсуждают патчи и
обмениваются результатами ревью в списках рассылки, а документация
процесса разработки объясняет, как отправлять, проверять и объединять
патчи\footnote{Официальная документация ядра Linux, How the development
  process works.}. Обнаружив неожиданное изменение, новый участник может
найти соответствующий патч и проследить ход состоявшегося тогда
обсуждения; эта запись пригодится и агенту при проверке исходной
информации.

Прослушав интервью Вэн Цзяи на WhynotTV, автор сформулировал одну из его
мыслей так: \textbf{и люди, и модели зависят от контекста.} Человек
может задать дополнительный вопрос коллеге, участвовавшему в принятии
решения. То, какую исходную информацию сможет прочитать модель, зависит
от оставленных командой записей и от того, какие материалы Harness
включит в текущий запрос.\footnote{WhynotTV, «Вэн Цзяи: OpenAI, GPT,
  обучение с подкреплением, Infra, постобучение, природный дар, tuixue,
  открытый исходный код, CMU, Университет Цинхуа», Podcast \#4.}

В статье о ReAct с помощью формального выражения описывается, как
история взаимодействия участвует в принятии решений\footnote{Yao,
  Shunyu, et al., ReAct: Synergizing Reasoning and Acting in Language
  Models, ICLR 2023.}. В момент времени \(t\) агент получает из
окружения наблюдение \(o_t\in\mathcal{O}\) и в соответствии с политикой
\(\pi(a_t\mid c_t)\) выбирает действие \(a_t\in\mathcal{A}\), где
контекст имеет вид:

\[
c_t=(o_1,a_1,\ldots,o_{t-1},a_{t-1},o_t)
\]

Это выражение показывает, что \textbf{следующее действие агента зависит
от всей истории взаимодействия к текущему моменту.} Например, на выбор
следующего шага влияют сведения о том, какой файл уже проверял агент и
какую ошибку вернул инструмент. Для агента на основе LLM
пользовательский ввод и результаты инструментов могут служить
наблюдениями, а вызовы инструментов и исходящие ответы фиксируют
действия; системные инструкции и определения инструментов дополнительно
описывают правила задачи и доступные интерфейсы.

В этой главе сначала используется способ отправки запросов без
сохранения состояния: в каждом раунде Harness формирует сообщения и
определения инструментов, после чего вызывает модель. Полная история
сообщений позволяет сохранять ход выполнения задачи; по мере роста
истории с помощью поиска, записей состояния и сжатия можно удерживать в
контекстном окне информацию, необходимую для последующих решений.
Компоновка контекста, панель состояния и методы сжатия посвящены одной и
той же проблеме: как поместить в ограниченный контекст сведения, которые
модели нужны на текущем шаге в составе \(c_t\)?

\hypertarget{ux43aux430ux43a-ai-ux430ux433ux435ux43dux442-ux432ux44bux437ux44bux432ux430ux435ux442-ux431ux43eux43bux44cux448ux443ux44e-ux44fux437ux44bux43aux43eux432ux443ux44e-ux43cux43eux434ux435ux43bux44c-ux441ux442ux440ux443ux43aux442ux443ux440ux430-ux43aux43eux43dux442ux435ux43aux441ux442ux430-api}{%
\section{Как AI-агент вызывает большую языковую модель: структура
контекста
API}\label{ux43aux430ux43a-ai-ux430ux433ux435ux43dux442-ux432ux44bux437ux44bux432ux430ux435ux442-ux431ux43eux43bux44cux448ux443ux44e-ux44fux437ux44bux43aux43eux432ux443ux44e-ux43cux43eux434ux435ux43bux44c-ux441ux442ux440ux443ux43aux442ux443ux440ux430-ux43aux43eux43dux442ux435ux43aux441ux442ux430-api}}

Сначала на примере простого приветствия рассмотрим формат сообщений, а
затем разберём, как вопрос «Который сейчас час в Ванкувере и какая там
погода?» приводит к двум вызовам инструментов. В этом разделе
используются интерфейс, совместимый с Chat Completions, и локально
запущенная модель Qwen3-0.6B, а в завершение с помощью Python запросы
всех ходов связываются в единую последовательность\footnote{OpenAI,
  справочник по Chat Completions API с описанием ролей сообщений,
  определений инструментов и полей результатов вызова; Qwen, Function
  Calling с описанием вызовов инструментов и шаблонов в совместимых
  сервисах.}.

\hypertarget{ux447ux435ux442ux44bux440ux435-ux440ux43eux43bux438-ux441ux43eux43eux431ux449ux435ux43dux438ux439-ux432-ux43fux440ux438ux43cux435ux440ux435}{%
\subsection{Четыре роли сообщений в
примере}\label{ux447ux435ux442ux44bux440ux435-ux440ux43eux43bux438-ux441ux43eux43eux431ux449ux435ux43dux438ux439-ux432-ux43fux440ux438ux43cux435ux440ux435}}

\textbf{Список сообщений} (\texttt{messages}) в запросе хранит сведения
о взаимодействии в хронологическом порядке. Каждое сообщение содержит
\textbf{роль} (\texttt{role}), обозначающую его источник и назначение. В
этом разделе используются четыре роли:

\begin{itemize}
\tightlist
\item
  \textbf{system}: системный промпт. В этом разделе он содержит заданные
  разработчиком сведения о роли, правила поведения и ограничения и
  размещается в начале списка сообщений. В интерфейсах некоторых моделей
  инструкции разработчика передаются через роль \texttt{developer}.
\item
  \textbf{user}: сообщение пользователя, содержащее задачу, уточняющие
  вопросы и дополнительные требования.
\item
  \textbf{assistant}: ответ модели, включая текст и запросы на вызов
  инструментов. Если добавить эти сообщения во входные данные следующего
  раунда, модель сможет прочитать свои предыдущие ответы и запросы на
  выполнение операций.
\item
  \textbf{tool}: результат выполнения инструмента. Получив возвращённое
  значение, Harness записывает его как сообщение инструмента и связывает
  с конкретным вызовом посредством \texttt{tool\_call\_id}.
\end{itemize}

Определения инструментов передаются в поле верхнего уровня
\texttt{tools} запроса и содержат названия, назначение и формат
параметров доступных инструментов.

Модель формирует запрос на вызов на основе определений инструментов.
Возможность выполнить эту операцию определяется Harness и средой
выполнения.

\hypertarget{ux43eux434ux43dux43eux445ux43eux434ux43eux432ux44bux439-ux434ux438ux430ux43bux43eux433-ux43fux440ux43eux441ux442ux435ux439ux448ux438ux439-ux432ux44bux437ux43eux432-api}{%
\subsection{Одноходовый диалог: простейший вызов
API}\label{ux43eux434ux43dux43eux445ux43eux434ux43eux432ux44bux439-ux434ux438ux430ux43bux43eux433-ux43fux440ux43eux441ux442ux435ux439ux448ux438ux439-ux432ux44bux437ux43eux432-api}}

Сначала рассмотрим один обмен вопросом и ответом. Пользователь
спрашивает: «Hello, who are you?», Harness отправляет модели системную
инструкцию вместе с вопросом пользователя, а затем считывает ответ. На
\crossreflink{fig:2-2}{рисунке 2-2} показано соответствие между запросом
и ответом.

\begin{figure}
\hypertarget{fig:2-2}{%
\centering
\includegraphics{images/revised/fig2-2.pdf}
\caption{Рисунок 2-2. Структура запроса и ответа при одноходовом вызове
API}\label{fig:2-2}
}
\end{figure}

Ниже запрос и пример ответа представлены в формате интерфейса локальной
модели. Комментарии в коде указывают источник каждой части:

\begin{Shaded}
\begin{Highlighting}[]
\CommentTok{// Harness формирует запрос}
\NormalTok{\{}
  \StringTok{"model"}\OperatorTok{:} \StringTok{"Qwen3{-}0.6B"}\OperatorTok{,}
  \StringTok{"messages"}\OperatorTok{:}\NormalTok{ [}
\NormalTok{    \{}
      \StringTok{"role"}\OperatorTok{:} \StringTok{"system"}\OperatorTok{,}                           \CommentTok{// Инструкция от разработчика}
      \StringTok{"content"}\OperatorTok{:} \StringTok{"You are a helpful coding assistant. Follow user instructions."}
\NormalTok{    \}}\OperatorTok{,}
\NormalTok{    \{}
      \StringTok{"role"}\OperatorTok{:} \StringTok{"user"}\OperatorTok{,}                              \CommentTok{// Ввод пользователя}
      \StringTok{"content"}\OperatorTok{:} \StringTok{"Hello, who are you?"}
\NormalTok{    \}}
\NormalTok{  ]}
\NormalTok{\}}
\end{Highlighting}
\end{Shaded}

\begin{Shaded}
\begin{Highlighting}[]
\CommentTok{// API возвращает ответ}
\NormalTok{\{}
  \StringTok{"choices"}\OperatorTok{:}\NormalTok{ [\{}
    \StringTok{"message"}\OperatorTok{:}\NormalTok{ \{}
      \StringTok{"role"}\OperatorTok{:} \StringTok{"assistant"}\OperatorTok{,}                         \CommentTok{// Сгенерировано моделью}
      \StringTok{"content"}\OperatorTok{:} \StringTok{"Hi! I\textquotesingle{}m a coding assistant. I can help you write code, debug issues, and explain technical concepts. How can I help?"}
\NormalTok{    \}}
\NormalTok{  \}]}
\NormalTok{\}}
\end{Highlighting}
\end{Shaded}

Запрос содержит одно сообщение \texttt{system} и одно сообщение
\texttt{user}, а \texttt{choices{[}0{]}.message} в ответе хранит
сообщение \texttt{assistant}, сгенерированное моделью. В используемом
здесь интерфейсе без сохранения состояния \textbf{Harness отвечает за
передачу необходимого контекста при каждом вызове}. Чтобы следующий
запрос продолжил этот диалог, в него потребуется включить
соответствующую историю.

\hypertarget{ux43cux43dux43eux433ux43eux445ux43eux434ux43eux432ux43eux435-ux432ux437ux430ux438ux43cux43eux434ux435ux439ux441ux442ux432ux438ux435-ux441-ux432ux44bux437ux43eux432ux430ux43cux438-ux438ux43dux441ux442ux440ux443ux43cux435ux43dux442ux43eux432-ux43eux441ux43dux43eux432ux43dux43eux439-ux446ux438ux43aux43b-ux430ux433ux435ux43dux442ux430}{%
\subsection{Многоходовое взаимодействие с вызовами инструментов:
основной цикл
агента}\label{ux43cux43dux43eux433ux43eux445ux43eux434ux43eux432ux43eux435-ux432ux437ux430ux438ux43cux43eux434ux435ux439ux441ux442ux432ux438ux435-ux441-ux432ux44bux437ux43eux432ux430ux43cux438-ux438ux43dux441ux442ux440ux443ux43cux435ux43dux442ux43eux432-ux43eux441ux43dux43eux432ux43dux43eux439-ux446ux438ux43aux43b-ux430ux433ux435ux43dux442ux430}}

Рассмотрим ещё один вопрос, требующий инструментов: ``What's the current
time and weather in Vancouver?'' Чтобы ответить на него, сначала нужно
узнать местные время и погоду. На \crossreflink{fig:2-3}{рисунке 2-3}
первый вызов модели со стороны Harness возвращает два запроса на вызов
инструментов. Harness обращается к сервисам времени и погоды, добавляет
оба результата в историю сообщений и только после второго вызова модели
получает ответ для пользователя.

\begin{figure}
\hypertarget{fig:2-3}{%
\centering
\includegraphics{images/revised/fig2-3.pdf}
\caption{Рисунок 2-3. Полная последовательность взаимодействия при двух
вызовах API модели}\label{fig:2-3}
}
\end{figure}

Какой из этих двух инструментов следует выполнить первым? Инструменту
времени требуется часовой пояс, а инструменту погоды --- город и единица
измерения температуры. Пользователь указал только Ванкувер; на основании
этого модель в примере подставляет \texttt{America/Vancouver} и
\texttt{Vancouver}, а в качестве единицы измерения выбирает градусы
Цельсия. Для запроса погоды не нужно ждать результата инструмента
времени, поэтому оба запроса можно отправить одновременно.

\textbf{Первый вызов API --- фреймворк агента отправляет исходный
запрос:}

\begin{Shaded}
\begin{Highlighting}[]
\CommentTok{// Первый вызов модели: Harness формирует запрос}
\NormalTok{\{}
  \StringTok{"model"}\OperatorTok{:} \StringTok{"Qwen3{-}0.6B"}\OperatorTok{,}
  \StringTok{"messages"}\OperatorTok{:}\NormalTok{ [}
\NormalTok{    \{}
      \StringTok{"role"}\OperatorTok{:} \StringTok{"system"}\OperatorTok{,}                           \CommentTok{// Инструкции, предоставленные разработчиком}
      \StringTok{"content"}\OperatorTok{:} \StringTok{"You are a helpful assistant. Use the provided tools to get real{-}time information when needed."}
\NormalTok{    \}}\OperatorTok{,}
\NormalTok{    \{}
      \StringTok{"role"}\OperatorTok{:} \StringTok{"user"}\OperatorTok{,}                              \CommentTok{// Ввод пользователя}
      \StringTok{"content"}\OperatorTok{:} \StringTok{"What\textquotesingle{}s the current time and weather in Vancouver?"}
\NormalTok{    \}}
\NormalTok{  ]}\OperatorTok{,}
  \StringTok{"tools"}\OperatorTok{:}\NormalTok{ [                                       }\CommentTok{// Инструменты, определённые разработчиком}
\NormalTok{    \{}
      \StringTok{"type"}\OperatorTok{:} \StringTok{"function"}\OperatorTok{,}
      \StringTok{"function"}\OperatorTok{:}\NormalTok{ \{}
        \StringTok{"name"}\OperatorTok{:} \StringTok{"get\_current\_time"}\OperatorTok{,}
        \StringTok{"description"}\OperatorTok{:} \StringTok{"Get the current date and time in a specific timezone"}\OperatorTok{,}
        \StringTok{"parameters"}\OperatorTok{:}\NormalTok{ \{}
          \StringTok{"type"}\OperatorTok{:} \StringTok{"object"}\OperatorTok{,}
          \StringTok{"properties"}\OperatorTok{:}\NormalTok{ \{}
            \StringTok{"timezone"}\OperatorTok{:}\NormalTok{ \{ }\StringTok{"type"}\OperatorTok{:} \StringTok{"string"}\OperatorTok{,} \StringTok{"description"}\OperatorTok{:} \StringTok{"Timezone name, e.g. America/Vancouver"}\NormalTok{ \}}
\NormalTok{          \}}\OperatorTok{,}
          \StringTok{"required"}\OperatorTok{:}\NormalTok{ [}\StringTok{"timezone"}\NormalTok{]}
\NormalTok{        \}}
\NormalTok{      \}}
\NormalTok{    \}}\OperatorTok{,}
\NormalTok{    \{}
      \StringTok{"type"}\OperatorTok{:} \StringTok{"function"}\OperatorTok{,}
      \StringTok{"function"}\OperatorTok{:}\NormalTok{ \{}
        \StringTok{"name"}\OperatorTok{:} \StringTok{"get\_weather"}\OperatorTok{,}
        \StringTok{"description"}\OperatorTok{:} \StringTok{"Get the current weather for a specific city"}\OperatorTok{,}
        \StringTok{"parameters"}\OperatorTok{:}\NormalTok{ \{}
          \StringTok{"type"}\OperatorTok{:} \StringTok{"object"}\OperatorTok{,}
          \StringTok{"properties"}\OperatorTok{:}\NormalTok{ \{}
            \StringTok{"city"}\OperatorTok{:}\NormalTok{ \{ }\StringTok{"type"}\OperatorTok{:} \StringTok{"string"}\OperatorTok{,} \StringTok{"description"}\OperatorTok{:} \StringTok{"City name"}\NormalTok{ \}}\OperatorTok{,}
            \StringTok{"unit"}\OperatorTok{:}\NormalTok{ \{ }\StringTok{"type"}\OperatorTok{:} \StringTok{"string"}\OperatorTok{,} \StringTok{"enum"}\OperatorTok{:}\NormalTok{ [}\StringTok{"celsius"}\OperatorTok{,} \StringTok{"fahrenheit"}\NormalTok{] \}}
\NormalTok{          \}}\OperatorTok{,}
          \StringTok{"required"}\OperatorTok{:}\NormalTok{ [}\StringTok{"city"}\OperatorTok{,} \StringTok{"unit"}\NormalTok{]}
\NormalTok{        \}}
\NormalTok{      \}}
\NormalTok{    \}}
\NormalTok{  ]}
\NormalTok{\}}
\end{Highlighting}
\end{Shaded}

Этот список \texttt{tools} заранее определяется разработчиком и содержит
имена, описания и схемы параметров. В данном примере в запросах каждого
раунда передаются одни и те же два инструмента. Прочитав список и вопрос
пользователя, модель выбирает инструменты для вызова и заполняет их
параметры.

\textbf{Модель возвращает два запроса на вызов инструментов:}

\begin{Shaded}
\begin{Highlighting}[]
\CommentTok{// Модель возвращает запросы на вызов инструментов}
\NormalTok{\{}
  \StringTok{"choices"}\OperatorTok{:}\NormalTok{ [\{}
    \StringTok{"message"}\OperatorTok{:}\NormalTok{ \{}
      \StringTok{"role"}\OperatorTok{:} \StringTok{"assistant"}\OperatorTok{,}                         \CommentTok{// Сгенерировано моделью}
      \StringTok{"content"}\OperatorTok{:} \KeywordTok{null}\OperatorTok{,}                             \CommentTok{// В этом раунде текстовый ответ отсутствует}
      \StringTok{"tool\_calls"}\OperatorTok{:}\NormalTok{ [                              }\CommentTok{// Запрашиваются вызовы двух инструментов}
\NormalTok{        \{}
          \StringTok{"id"}\OperatorTok{:} \StringTok{"call\_abc123"}\OperatorTok{,}
          \StringTok{"type"}\OperatorTok{:} \StringTok{"function"}\OperatorTok{,}
          \StringTok{"function"}\OperatorTok{:}\NormalTok{ \{}
            \StringTok{"name"}\OperatorTok{:} \StringTok{"get\_current\_time"}\OperatorTok{,}
            \StringTok{"arguments"}\OperatorTok{:} \StringTok{"\{}\SpecialCharTok{\textbackslash{}"}\StringTok{timezone}\SpecialCharTok{\textbackslash{}"}\StringTok{: }\SpecialCharTok{\textbackslash{}"}\StringTok{America/Vancouver}\SpecialCharTok{\textbackslash{}"}\StringTok{\}"}
\NormalTok{          \}}
\NormalTok{        \}}\OperatorTok{,}
\NormalTok{        \{}
          \StringTok{"id"}\OperatorTok{:} \StringTok{"call\_def456"}\OperatorTok{,}
          \StringTok{"type"}\OperatorTok{:} \StringTok{"function"}\OperatorTok{,}
          \StringTok{"function"}\OperatorTok{:}\NormalTok{ \{}
            \StringTok{"name"}\OperatorTok{:} \StringTok{"get\_weather"}\OperatorTok{,}
            \StringTok{"arguments"}\OperatorTok{:} \StringTok{"\{}\SpecialCharTok{\textbackslash{}"}\StringTok{city}\SpecialCharTok{\textbackslash{}"}\StringTok{: }\SpecialCharTok{\textbackslash{}"}\StringTok{Vancouver}\SpecialCharTok{\textbackslash{}"}\StringTok{, }\SpecialCharTok{\textbackslash{}"}\StringTok{unit}\SpecialCharTok{\textbackslash{}"}\StringTok{: }\SpecialCharTok{\textbackslash{}"}\StringTok{celsius}\SpecialCharTok{\textbackslash{}"}\StringTok{\}"}
\NormalTok{          \}}
\NormalTok{        \}}
\NormalTok{      ]}
\NormalTok{    \}}
\NormalTok{  \}]}
\NormalTok{\}}
\end{Highlighting}
\end{Shaded}

Поле \texttt{tool\_calls} в ответе содержит два запроса, каждый со своим
\texttt{id}, именем инструмента и параметрами. \texttt{arguments} ---
строка с данными в формате JSON; исполнитель должен сначала разобрать
её, а затем проверить поля и их значения. \textbf{Модель определяет, что
вызвать и какие параметры передать; Harness отвечает за проверку и
диспетчеризацию; реализация инструмента обращается к сервису времени или
погоды.}

\textbf{Фреймворк агента выполняет инструменты, а затем делает второй
вызов API:}

После проверки запросов Harness запускает оба инструмента, записывает
возвращённые значения в сообщения \texttt{tool}, а затем передаёт модели
имеющуюся историю вместе с новыми результатами:

\begin{Shaded}
\begin{Highlighting}[]
\CommentTok{// Второй вызов модели: добавление результатов инструментов}
\NormalTok{\{}
  \StringTok{"model"}\OperatorTok{:} \StringTok{"Qwen3{-}0.6B"}\OperatorTok{,}
  \StringTok{"messages"}\OperatorTok{:}\NormalTok{ [}
\NormalTok{    \{}
      \StringTok{"role"}\OperatorTok{:} \StringTok{"system"}\OperatorTok{,}                           \CommentTok{// Сообщение сохранено из первого запроса}
      \StringTok{"content"}\OperatorTok{:} \StringTok{"You are a helpful assistant. Use the provided tools to get real{-}time information when needed."}
\NormalTok{    \}}\OperatorTok{,}
\NormalTok{    \{}
      \StringTok{"role"}\OperatorTok{:} \StringTok{"user"}\OperatorTok{,}                              \CommentTok{// Сообщение сохранено из первого запроса}
      \StringTok{"content"}\OperatorTok{:} \StringTok{"What\textquotesingle{}s the current time and weather in Vancouver?"}
\NormalTok{    \}}\OperatorTok{,}
\NormalTok{    \{}
      \StringTok{"role"}\OperatorTok{:} \StringTok{"assistant"}\OperatorTok{,}                         \CommentTok{// Запросы модели из предыдущего раунда сохранены}
      \StringTok{"content"}\OperatorTok{:} \KeywordTok{null}\OperatorTok{,}
      \StringTok{"tool\_calls"}\OperatorTok{:}\NormalTok{ [}
\NormalTok{        \{ }\StringTok{"id"}\OperatorTok{:} \StringTok{"call\_abc123"}\OperatorTok{,} \StringTok{"type"}\OperatorTok{:} \StringTok{"function"}\OperatorTok{,} \StringTok{"function"}\OperatorTok{:}\NormalTok{ \{ }\StringTok{"name"}\OperatorTok{:} \StringTok{"get\_current\_time"}\OperatorTok{,} \StringTok{"arguments"}\OperatorTok{:} \StringTok{"\{}\SpecialCharTok{\textbackslash{}"}\StringTok{timezone}\SpecialCharTok{\textbackslash{}"}\StringTok{: }\SpecialCharTok{\textbackslash{}"}\StringTok{America/Vancouver}\SpecialCharTok{\textbackslash{}"}\StringTok{\}"}\NormalTok{ \} \}}\OperatorTok{,}
\NormalTok{        \{ }\StringTok{"id"}\OperatorTok{:} \StringTok{"call\_def456"}\OperatorTok{,} \StringTok{"type"}\OperatorTok{:} \StringTok{"function"}\OperatorTok{,} \StringTok{"function"}\OperatorTok{:}\NormalTok{ \{ }\StringTok{"name"}\OperatorTok{:} \StringTok{"get\_weather"}\OperatorTok{,} \StringTok{"arguments"}\OperatorTok{:} \StringTok{"\{}\SpecialCharTok{\textbackslash{}"}\StringTok{city}\SpecialCharTok{\textbackslash{}"}\StringTok{: }\SpecialCharTok{\textbackslash{}"}\StringTok{Vancouver}\SpecialCharTok{\textbackslash{}"}\StringTok{, }\SpecialCharTok{\textbackslash{}"}\StringTok{unit}\SpecialCharTok{\textbackslash{}"}\StringTok{: }\SpecialCharTok{\textbackslash{}"}\StringTok{celsius}\SpecialCharTok{\textbackslash{}"}\StringTok{\}"}\NormalTok{ \} \}}
\NormalTok{      ]}
\NormalTok{    \}}\OperatorTok{,}
\NormalTok{    \{}
      \StringTok{"role"}\OperatorTok{:} \StringTok{"tool"}\OperatorTok{,}                              \CommentTok{// Harness записывает результат выполнения инструмента}
      \StringTok{"tool\_call\_id"}\OperatorTok{:} \StringTok{"call\_abc123"}\OperatorTok{,}
      \StringTok{"content"}\OperatorTok{:} \StringTok{"\{}\SpecialCharTok{\textbackslash{}"}\StringTok{timezone}\SpecialCharTok{\textbackslash{}"}\StringTok{: }\SpecialCharTok{\textbackslash{}"}\StringTok{America/Vancouver}\SpecialCharTok{\textbackslash{}"}\StringTok{, }\SpecialCharTok{\textbackslash{}"}\StringTok{datetime}\SpecialCharTok{\textbackslash{}"}\StringTok{: }\SpecialCharTok{\textbackslash{}"}\StringTok{2025{-}09{-}13T05:18:47}\SpecialCharTok{\textbackslash{}"}\StringTok{, }\SpecialCharTok{\textbackslash{}"}\StringTok{day\_of\_week}\SpecialCharTok{\textbackslash{}"}\StringTok{: }\SpecialCharTok{\textbackslash{}"}\StringTok{Saturday}\SpecialCharTok{\textbackslash{}"}\StringTok{\}"}
\NormalTok{    \}}\OperatorTok{,}
\NormalTok{    \{}
      \StringTok{"role"}\OperatorTok{:} \StringTok{"tool"}\OperatorTok{,}                              \CommentTok{// Harness записывает результат выполнения инструмента}
      \StringTok{"tool\_call\_id"}\OperatorTok{:} \StringTok{"call\_def456"}\OperatorTok{,}
      \StringTok{"content"}\OperatorTok{:} \StringTok{"\{}\SpecialCharTok{\textbackslash{}"}\StringTok{city}\SpecialCharTok{\textbackslash{}"}\StringTok{: }\SpecialCharTok{\textbackslash{}"}\StringTok{Vancouver}\SpecialCharTok{\textbackslash{}"}\StringTok{, }\SpecialCharTok{\textbackslash{}"}\StringTok{temperature}\SpecialCharTok{\textbackslash{}"}\StringTok{: 13.2, }\SpecialCharTok{\textbackslash{}"}\StringTok{unit}\SpecialCharTok{\textbackslash{}"}\StringTok{: }\SpecialCharTok{\textbackslash{}"}\StringTok{celsius}\SpecialCharTok{\textbackslash{}"}\StringTok{, }\SpecialCharTok{\textbackslash{}"}\StringTok{conditions}\SpecialCharTok{\textbackslash{}"}\StringTok{: }\SpecialCharTok{\textbackslash{}"}\StringTok{clear}\SpecialCharTok{\textbackslash{}"}\StringTok{, }\SpecialCharTok{\textbackslash{}"}\StringTok{humidity}\SpecialCharTok{\textbackslash{}"}\StringTok{: 93\}"}
\NormalTok{    \}}
\NormalTok{  ]}\OperatorTok{,}
  \StringTok{"tools"}\OperatorTok{:}\NormalTok{ [ }\OperatorTok{...}\NormalTok{ ]                                 }\CommentTok{// Используются приведённые выше определения инструментов; здесь они опущены}
\NormalTok{\}}
\end{Highlighting}
\end{Shaded}

Проследим во втором запросе идентификатор \texttt{call\_def456}: в
предыдущем сообщении \texttt{assistant} он обозначает запрос погоды в
Ванкувере, а последующее сообщение \texttt{tool} с тем же
идентификатором возвращает температуру 13,2 ℃ и ясную погоду. Даже если
погода вернётся раньше времени, Harness сможет сопоставить каждый
результат с соответствующим запросом. В этом раунде модель одновременно
получает исходный вопрос пользователя, два сформированных ею запроса и
данные, возвращённые инструментами. Теперь она может ответить
пользователю.

\textbf{На основе результатов инструментов модель формирует
окончательный ответ:}

\begin{Shaded}
\begin{Highlighting}[]
\CommentTok{// Модель формирует окончательный ответ}
\NormalTok{\{}
  \StringTok{"choices"}\OperatorTok{:}\NormalTok{ [\{}
    \StringTok{"message"}\OperatorTok{:}\NormalTok{ \{}
      \StringTok{"role"}\OperatorTok{:} \StringTok{"assistant"}\OperatorTok{,}                         \CommentTok{// Сгенерировано моделью}
      \StringTok{"content"}\OperatorTok{:} \StringTok{"It\textquotesingle{}s currently 5:18 AM on Saturday, September 13, 2025 in Vancouver.}\SpecialCharTok{\textbackslash{}n\textbackslash{}n}\StringTok{Weather: 13.2°C with clear skies and 93\% humidity. It\textquotesingle{}s quite cool this morning {-} you might want to grab a jacket."}
\NormalTok{    \}}
\NormalTok{  \}]}
\NormalTok{\}}
\end{Highlighting}
\end{Shaded}

Второй ответ содержит окончательный текст, а \texttt{tool\_calls} пуст,
поэтому цикл в примере завершается. Процесс «вызвать модель → выполнить
инструменты → вернуть результаты → снова вызвать модель» представляет
собой конкретную реализацию цикла выполнения из первой главы на уровне
API.

Если пользователь затем спросит: «А что в Токио?», Harness добавит это
сообщение в конец истории и снова вызовет модель. Опираясь на контекст
задачи «узнать текущие время и погоду», модель сможет сформировать новые
запросы для Токио, а после получения результатов инструментов ---
подготовить новый ответ.

\textbf{Как вернуть результаты инструментов при использовании Responses
API?} В приведённом выше локальном эксперименте используется интерфейс,
совместимый с Chat Completions. Для новых проектов на сервисах OpenAI
официально рекомендуется Responses API. Вернёмся к примеру с Ванкувером:
приложение помещает вопрос пользователя в поле \texttt{input}, а в
\texttt{tools} объявляет две функции --- времени и погоды. Возвращаемое
поле \texttt{output} представляет собой набор типизированных элементов.
Если среди них есть два элемента \texttt{function\_call}, приложение
считывает из каждого имя функции, параметры и \texttt{call\_id}, после
чего отдельно выполняет каждый инструмент.\footnote{OpenAI, документация
  по переходу на Responses API, вызову функций и параметрам создания
  ответа.}

Предположим, что идентификатор вызова функции погоды ---
\texttt{call\_weather}. Тогда в \texttt{input} следующего запроса можно
поместить такой результат; результат функции времени также необходимо
вернуть с её собственным идентификатором.

\begin{Shaded}
\begin{Highlighting}[]
\FunctionTok{\{}\DataTypeTok{"type"}\FunctionTok{:} \StringTok{"function\_call\_output"}\FunctionTok{,} \DataTypeTok{"call\_id"}\FunctionTok{:} \StringTok{"call\_weather"}\FunctionTok{,} \DataTypeTok{"output"}\FunctionTok{:} \StringTok{"\{}\CharTok{\textbackslash{}"}\StringTok{temperature}\CharTok{\textbackslash{}"}\StringTok{: 13.2, }\CharTok{\textbackslash{}"}\StringTok{unit}\CharTok{\textbackslash{}"}\StringTok{: }\CharTok{\textbackslash{}"}\StringTok{celsius}\CharTok{\textbackslash{}"}\StringTok{\}"}\FunctionTok{\}}
\end{Highlighting}
\end{Shaded}

\texttt{call\_id} связывает результат с конкретным вызовом функции.
\texttt{previous\_response\_id} указывает на предыдущий ответ модели,
позволяя сервису продолжить контекст того раунда. Эти два идентификатора
относятся к разным объектам. При таком способе продолжения в текущем
раунде всё равно необходимо передать инструменты, которые должны
оставаться доступными, и поле верхнего уровня \texttt{instructions}:
инструкции верхнего уровня из предыдущего запроса не наследуются
автоматически через идентификатор ответа. Модель может снова запросить
инструменты, поэтому приложение должно прочитать следующий
\texttt{output} и только затем решить, завершать ли цикл. В эксперименте
1-3 из первой главы использовался управляемый сервисом поиск; две
пользовательские функции в этом примере фактически выполняет приложение.

\hypertarget{ux440ux435ux430ux43bux438ux437ux430ux446ux438ux44f-ux43eux441ux43dux43eux432ux43dux43eux433ux43e-ux446ux438ux43aux43bux430-ux430ux433ux435ux43dux442ux430-ux432-ux43aux43eux434ux435}{%
\subsection{Реализация основного цикла агента в
коде}\label{ux440ux435ux430ux43bux438ux437ux430ux446ux438ux44f-ux43eux441ux43dux43eux432ux43dux43eux433ux43e-ux446ux438ux43aux43bux430-ux430ux433ux435ux43dux442ux430-ux432-ux43aux43eux434ux435}}

Два вызова \texttt{messages.append} в цикле сохраняют соответственно
ответ модели и результат выполнения инструмента. И время, и температура
относятся к Ванкуверу; если поменять местами два ID вызовов, в следующем
раунде модель прочитает, что «запрос времени» вернул температуру 13,2
градуса Цельсия.

В примере для подключения к совместимому сервису используется OpenAI
Python SDK. По умолчанию указан локальный адрес
\texttt{http://localhost:8000/v1}, подходящий для уже запущенного
локального сервиса без аутентификации. Эти значения по умолчанию
соответствуют бэкенду vLLM из эксперимента 2-1. При использовании
бэкенда Ollama из того же эксперимента задайте адрес
\texttt{http://localhost:11434/v1} и измените имя модели на
\texttt{qwen3:0.6b}. Для подключения к другим сервисам при необходимости
задайте переменные среды \texttt{LLM\_BASE\_URL} и
\texttt{LLM\_API\_KEY}, а также приведите указанный ниже параметр
\texttt{model} в соответствие с фактическим именем развёрнутой модели.
Инструменты времени и погоды пока используют фиксированные данные из
предыдущего примера, чтобы сообщения можно было сопоставить; при
подключении реальных данных замените реализации этих двух инструментов.

\begin{Shaded}
\begin{Highlighting}[]
\ImportTok{import}\NormalTok{ json}
\ImportTok{import}\NormalTok{ os}

\ImportTok{from}\NormalTok{ openai }\ImportTok{import}\NormalTok{ OpenAI}

\NormalTok{client }\OperatorTok{=}\NormalTok{ OpenAI(}
\NormalTok{    base\_url}\OperatorTok{=}\NormalTok{os.getenv(}\StringTok{"LLM\_BASE\_URL"}\NormalTok{, }\StringTok{"http://localhost:8000/v1"}\NormalTok{),}
\NormalTok{    api\_key}\OperatorTok{=}\NormalTok{os.getenv(}\StringTok{"LLM\_API\_KEY"}\NormalTok{, }\StringTok{"EMPTY"}\NormalTok{),}
\NormalTok{)}

\CommentTok{\# Определения инструментов: имена, назначение и структура параметров}
\NormalTok{tools }\OperatorTok{=}\NormalTok{ [}
\NormalTok{    \{}
        \StringTok{"type"}\NormalTok{: }\StringTok{"function"}\NormalTok{,}
        \StringTok{"function"}\NormalTok{: \{}
            \StringTok{"name"}\NormalTok{: }\StringTok{"get\_current\_time"}\NormalTok{,}
            \StringTok{"description"}\NormalTok{: }\StringTok{"Get the current date and time in a specific timezone"}\NormalTok{,}
            \StringTok{"parameters"}\NormalTok{: \{}
                \StringTok{"type"}\NormalTok{: }\StringTok{"object"}\NormalTok{,}
                \StringTok{"properties"}\NormalTok{: \{}
                    \StringTok{"timezone"}\NormalTok{: \{}\StringTok{"type"}\NormalTok{: }\StringTok{"string"}\NormalTok{, }\StringTok{"description"}\NormalTok{: }\StringTok{"Timezone name, e.g. America/Vancouver"}\NormalTok{\}}
\NormalTok{                \},}
                \StringTok{"required"}\NormalTok{: [}\StringTok{"timezone"}\NormalTok{],}
\NormalTok{            \},}
\NormalTok{        \},}
\NormalTok{    \},}
\NormalTok{    \{}
        \StringTok{"type"}\NormalTok{: }\StringTok{"function"}\NormalTok{,}
        \StringTok{"function"}\NormalTok{: \{}
            \StringTok{"name"}\NormalTok{: }\StringTok{"get\_weather"}\NormalTok{,}
            \StringTok{"description"}\NormalTok{: }\StringTok{"Get the current weather for a specific city"}\NormalTok{,}
            \StringTok{"parameters"}\NormalTok{: \{}
                \StringTok{"type"}\NormalTok{: }\StringTok{"object"}\NormalTok{,}
                \StringTok{"properties"}\NormalTok{: \{}
                    \StringTok{"city"}\NormalTok{: \{}\StringTok{"type"}\NormalTok{: }\StringTok{"string"}\NormalTok{, }\StringTok{"description"}\NormalTok{: }\StringTok{"City name"}\NormalTok{\},}
                    \StringTok{"unit"}\NormalTok{: \{}\StringTok{"type"}\NormalTok{: }\StringTok{"string"}\NormalTok{, }\StringTok{"enum"}\NormalTok{: [}\StringTok{"celsius"}\NormalTok{, }\StringTok{"fahrenheit"}\NormalTok{]\},}
\NormalTok{                \},}
                \StringTok{"required"}\NormalTok{: [}\StringTok{"city"}\NormalTok{, }\StringTok{"unit"}\NormalTok{],}
\NormalTok{            \},}
\NormalTok{        \},}
\NormalTok{    \},}
\NormalTok{]}

\CommentTok{\# В примере инструменты используют фиксированные данные, а параметры соответствуют приведённому выше запросу для Ванкувера.}
\KeywordTok{def}\NormalTok{ execute\_tool(name, arguments):}
\NormalTok{    args }\OperatorTok{=}\NormalTok{ json.loads(arguments)}
    \ControlFlowTok{if}\NormalTok{ name }\OperatorTok{==} \StringTok{"get\_current\_time"} \KeywordTok{and}\NormalTok{ args }\OperatorTok{==}\NormalTok{ \{}\StringTok{"timezone"}\NormalTok{: }\StringTok{"America/Vancouver"}\NormalTok{\}:}
\NormalTok{        result }\OperatorTok{=}\NormalTok{ \{}
            \StringTok{"timezone"}\NormalTok{: }\StringTok{"America/Vancouver"}\NormalTok{,}
            \StringTok{"datetime"}\NormalTok{: }\StringTok{"2025{-}09{-}13T05:18:47"}\NormalTok{,}
            \StringTok{"day\_of\_week"}\NormalTok{: }\StringTok{"Saturday"}\NormalTok{,}
\NormalTok{        \}}
    \ControlFlowTok{elif}\NormalTok{ name }\OperatorTok{==} \StringTok{"get\_weather"} \KeywordTok{and}\NormalTok{ args }\OperatorTok{==}\NormalTok{ \{}\StringTok{"city"}\NormalTok{: }\StringTok{"Vancouver"}\NormalTok{, }\StringTok{"unit"}\NormalTok{: }\StringTok{"celsius"}\NormalTok{\}:}
\NormalTok{        result }\OperatorTok{=}\NormalTok{ \{}
            \StringTok{"city"}\NormalTok{: }\StringTok{"Vancouver"}\NormalTok{, }\StringTok{"temperature"}\NormalTok{: }\FloatTok{13.2}\NormalTok{,}
            \StringTok{"unit"}\NormalTok{: }\StringTok{"celsius"}\NormalTok{, }\StringTok{"conditions"}\NormalTok{: }\StringTok{"clear"}\NormalTok{, }\StringTok{"humidity"}\NormalTok{: }\DecValTok{93}\NormalTok{,}
\NormalTok{        \}}
    \ControlFlowTok{else}\NormalTok{:}
\NormalTok{        result }\OperatorTok{=}\NormalTok{ \{}\StringTok{"error"}\NormalTok{: }\StringTok{"В примере доступны только время и погода в градусах Цельсия для Ванкувера"}\NormalTok{\}}
    \ControlFlowTok{return}\NormalTok{ json.dumps(result, ensure\_ascii}\OperatorTok{=}\VariableTok{False}\NormalTok{)}

\CommentTok{\# Исходные сообщения: инструкция и задача пользователя}
\NormalTok{messages }\OperatorTok{=}\NormalTok{ [}
\NormalTok{    \{}\StringTok{"role"}\NormalTok{: }\StringTok{"system"}\NormalTok{, }\StringTok{"content"}\NormalTok{: }\StringTok{"You are a helpful assistant. Use tools to get real{-}time information when needed."}\NormalTok{\},}
\NormalTok{    \{}\StringTok{"role"}\NormalTok{: }\StringTok{"user"}\NormalTok{, }\StringTok{"content"}\NormalTok{: }\StringTok{"What\textquotesingle{}s the current time and weather in Vancouver?"}\NormalTok{\},}
\NormalTok{]}

\CommentTok{\# Выполняем цикл: не более восьми раундов вызова модели}
\NormalTok{rounds }\OperatorTok{=} \DecValTok{0}
\ControlFlowTok{while}\NormalTok{ rounds }\OperatorTok{\textless{}} \DecValTok{8}\NormalTok{:}
\NormalTok{    rounds }\OperatorTok{+=} \DecValTok{1}
\NormalTok{    response }\OperatorTok{=}\NormalTok{ client.chat.completions.create(}
\NormalTok{        model}\OperatorTok{=}\StringTok{"Qwen3{-}0.6B"}\NormalTok{, messages}\OperatorTok{=}\NormalTok{messages, tools}\OperatorTok{=}\NormalTok{tools}
\NormalTok{    )}
\NormalTok{    assistant\_message }\OperatorTok{=}\NormalTok{ response.choices[}\DecValTok{0}\NormalTok{].message}

    \CommentTok{\# Сохраняем ответ модели, включая запросы на вызов инструментов в текущем раунде}
\NormalTok{    messages.append(assistant\_message.model\_dump(exclude\_none}\OperatorTok{=}\VariableTok{True}\NormalTok{))}

    \CommentTok{\# В этом примере цикл завершается, когда модель возвращает ответ без вызовов инструментов}
    \ControlFlowTok{if} \KeywordTok{not}\NormalTok{ assistant\_message.tool\_calls:}
        \BuiltInTok{print}\NormalTok{(assistant\_message.content)}
        \ControlFlowTok{break}

    \CommentTok{\# Поочерёдно выполняем инструменты и добавляем результаты; независимые вызовы можно выполнять параллельно}
    \ControlFlowTok{for}\NormalTok{ tool\_call }\KeywordTok{in}\NormalTok{ assistant\_message.tool\_calls:}
        \ControlFlowTok{try}\NormalTok{:}
\NormalTok{            result }\OperatorTok{=}\NormalTok{ execute\_tool(tool\_call.function.name, tool\_call.function.arguments)}
        \ControlFlowTok{except}\NormalTok{ json.JSONDecodeError:}
\NormalTok{            result }\OperatorTok{=}\NormalTok{ json.dumps(\{}\StringTok{"error"}\NormalTok{: }\StringTok{"Параметры инструмента должны содержать корректный JSON"}\NormalTok{\})}
\NormalTok{        messages.append(\{}
            \StringTok{"role"}\NormalTok{: }\StringTok{"tool"}\NormalTok{,}
            \StringTok{"tool\_call\_id"}\NormalTok{: tool\_call.}\BuiltInTok{id}\NormalTok{,}
            \StringTok{"content"}\NormalTok{: result,}
\NormalTok{        \})}
    \CommentTok{\# В следующем раунде используется список сообщений с добавленными результатами}
\ControlFlowTok{else}\NormalTok{:}
    \BuiltInTok{print}\NormalTok{(}\StringTok{"Достигнуто максимальное количество раундов, но задача ещё не завершена."}\NormalTok{)}
\end{Highlighting}
\end{Shaded}

Предположим, что в первом раунде модель одновременно запросила время и
погоду. Этот цикл \texttt{for} сначала получит время и добавит результат
в \texttt{messages}, а затем получит погоду и добавит второй результат.
Только после выполнения обеих операций программа вернётся к началу
\texttt{while} и отправит второй запрос модели. Если развернуть
содержимое \texttt{messages} на этих шагах, можно увидеть, какие данные
модель получала при каждом вызове.

\textbf{Исходное состояние (до 1-го вызова):}

\begin{Shaded}
\begin{Highlighting}[]
\NormalTok{messages = [}
\NormalTok{  \{ role: "system",  content: "You are a helpful assistant..." \},     \# Написано разработчиком}
\NormalTok{  \{ role: "user",    content: "What\textquotesingle{}s the current time and weather in Vancouver?" \},  \# Ввод пользователя}
\NormalTok{]}
\end{Highlighting}
\end{Shaded}

\textbf{После завершения 1-го вызова модели и выполнения инструментов:}

\begin{Shaded}
\begin{Highlighting}[]
\NormalTok{messages = [}
\NormalTok{  \{ role: "system",    content: "..." \},}
\NormalTok{  \{ role: "user",      content: "What\textquotesingle{}s the current time..." \},}
\NormalTok{  \{ role: "assistant", tool\_calls: [get\_current\_time, get\_weather] \},  \# Запросы на вызов, сгенерированные моделью}
\NormalTok{  \{ role: "tool",      tool\_call\_id: "call\_abc123", content: "\{time...\}" \},  \# Записано Harness после выполнения инструмента}
\NormalTok{  \{ role: "tool",      tool\_call\_id: "call\_def456", content: "\{weather...\}" \},  \# Записано Harness после выполнения инструмента}
\NormalTok{]}
\end{Highlighting}
\end{Shaded}

\textbf{После 2-го вызова (модель возвращает окончательный ответ, цикл
завершается):}

\begin{Shaded}
\begin{Highlighting}[]
\NormalTok{messages = [}
\NormalTok{  \{ role: "system",    content: "..." \},}
\NormalTok{  \{ role: "user",      content: "What\textquotesingle{}s the current time..." \},}
\NormalTok{  \{ role: "assistant", tool\_calls: [get\_current\_time, get\_weather] \},}
\NormalTok{  \{ role: "tool",      tool\_call\_id: "call\_abc123", content: "\{time...\}" \},}
\NormalTok{  \{ role: "tool",      tool\_call\_id: "call\_def456", content: "\{weather...\}" \},}
\NormalTok{  \{ role: "assistant", content: "It\textquotesingle{}s currently Saturday, Sep 13, 2025 in Vancouver..." \},  \# Окончательный ответ}
\NormalTok{]}
\end{Highlighting}
\end{Shaded}

\hypertarget{ux441ux43eux441ux442ux430ux432-ux43aux43eux43dux442ux435ux43aux441ux442ux430-ux441-ux442ux43eux447ux43aux438-ux437ux440ux435ux43dux438ux44f-api}{%
\subsection{Состав контекста с точки зрения
API}\label{ux441ux43eux441ux442ux430ux432-ux43aux43eux43dux442ux435ux43aux441ux442ux430-ux441-ux442ux43eux447ux43aux438-ux437ux440ux435ux43dux438ux44f-api}}

На \crossreflink{fig:2-4}{рисунке 2-4} запрос разделён на две части:
префикс и историю сообщений. Сразу после отправки приведённого выше
запроса о Ванкувере история содержит только вопрос пользователя и
предшествующий диалог. После завершения обоих запросов в ней появляются
предложенные моделью вызовы инструментов и два результата их выполнения.

\begin{figure}
\hypertarget{fig:2-4}{%
\centering
\includegraphics{images/revised/fig2-4.pdf}
\caption{Рисунок 2-4. Состав контекста при каждом обращении AI-агента к
модели}\label{fig:2-4}
}
\end{figure}

Таким образом, общий начальный фрагмент двух последовательных запросов
может распространяться и на историю сообщений. Объём старых вычислений,
который удастся повторно использовать при следующем вызове, зависит от
длины этого префикса.

В приведённом ниже Harness редко изменяющиеся системные инструкции и
определения инструментов размещаются в начале, а в каждом цикле текущее
состояние лишь добавляется в конец истории сообщений. Если история
становится слишком длинной, старые сообщения сначала сжимаются, а затем
собирается запрос.

\begin{Shaded}
\begin{Highlighting}[]
\NormalTok{stable\_prefix }\OperatorTok{=}\NormalTok{ system\_message}
\NormalTok{stable\_tools }\OperatorTok{=}\NormalTok{ core\_tool\_schemas}
\NormalTok{history }\OperatorTok{=}\NormalTok{ load\_message\_history(session)}
\NormalTok{status\_message }\OperatorTok{=}\NormalTok{ make\_status\_message(derive\_current\_state(history))}

\ControlFlowTok{if}\NormalTok{ estimated\_tokens(stable\_prefix, stable\_tools, history, status\_message) }\OperatorTok{\textgreater{}}\NormalTok{ budget:}
\NormalTok{    history }\OperatorTok{=}\NormalTok{ compress\_old\_evidence(}
\NormalTok{        history,}
\NormalTok{        preserve }\OperatorTok{=}\NormalTok{ [decisions, constraints, failures, citations]}
\NormalTok{    )}

\NormalTok{request.messages }\OperatorTok{=}\NormalTok{ [stable\_prefix] }\OperatorTok{+}\NormalTok{ history }\OperatorTok{+}\NormalTok{ [status\_message]}
\NormalTok{request.tools }\OperatorTok{=}\NormalTok{ stable\_tools}
\NormalTok{response }\OperatorTok{=}\NormalTok{ call\_model(request)}
\end{Highlighting}
\end{Shaded}

При сжатии сохраняется необходимая для последующей работы информация,
включая решения, ограничения, сбои и ссылки на источники, а также
поддерживается связь между запросами на вызов инструментов и
результатами их выполнения. Сжатие может запускаться не только при
превышении бюджета длины, но и при наличии повторяющейся или устаревшей
информации. Конкретные стратегии будут рассмотрены во второй половине
этой главы.

\begin{quote}
\textbf{Эксперимент 2-1 ★: локальное развёртывание LLM-сервиса и вызов
инструментов}

В этом эксперименте Qwen3-0.6B запускается локально. На примере задачи
«узнать время и погоду в Ванкувере» рассматривается прохождение
сообщений через модель, парсер и исполнитель инструментов. На
\crossreflink{fig:2-5}{рисунке 2-5} показано разделение обязанностей
между локальным сервисом и клиентом: сервер загружает модель,
обрабатывает входные данные и генерирует ответ, а клиент ведёт историю
сообщений, диспетчеризует инструменты и возвращает результаты.

\begin{figure}
\hypertarget{fig:2-5}{%
\centering
\includegraphics{images/revised/fig2-5.pdf}
\caption{Рисунок 2-5. Архитектура вызова инструментов локальной
LLM}\label{fig:2-5}
}
\end{figure}

\textbf{Сначала проследим за выполнением задачи целиком.} В одном цикле
вывода модель запрашивает время и температуру, указывая соответственно
параметры \texttt{America/Vancouver}, \texttt{Vancouver} и
\texttt{celsius}. Клиент параллельно запускает два инструмента. Первый
результат содержит местное время 11:11:28, второй --- температуру 17,9
°C. Получив оба результата, модель сообщает время с точностью до минуты
и температуру 17,9 °C, после чего завершает текущий цикл. Этот запуск
показывает, что модель с 600 миллионами параметров уже способна выбрать
для такой задачи два правильных инструмента, заполнить их параметры и
ответить пользователю на основе возвращённых значений.

\textbf{Теперь рассмотрим соответствие между исходным выводом и полями
API.} В режиме рассуждений Qwen3 может генерировать содержимое
рассуждений с тегами \texttt{\textless{}think\textgreater{}}, а для
запросов к инструментам используются соответствующие маркеры вызова.
После разбора на сервере клиент может получить отдельное поле
\texttt{thinking} или поле вызова инструмента. Порядок и способ
сохранения отдельных частей вывода совместно определяются шаблоном
модели и протоколом сервиса.

\textbf{Затем рассмотрим потоковую обработку.} По мере поступления
событий интерфейс может обновлять состояния «Выполняется рассуждение»,
«Вызывается инструмент» и «Генерируется ответ». Параметры инструмента
обычно поступают несколькими фрагментами, которые необходимо накопить до
получения полного запроса. Среда выполнения может приступить к
диспетчеризации после того, как протокол подтвердит завершение вызова, а
параметры, разрешения и зависимости пройдут проверку. Независимые
инструменты можно выполнять параллельно, а их работу --- объединять с
последующей обработкой вывода в конвейер.

\textbf{Наконец, сравним скорость и кэширование.} В этом запуске на Mac
с процессором Apple средняя скорость декодирования при выполнении задачи
с инструментами составила около 107 токенов в секунду. Токен --- базовая
единица текста, обрабатываемая моделью; конкретное разбиение
определяется токенизатором, поэтому одному слову или иероглифу может
соответствовать один или несколько токенов. Скорость декодирования
характеризует скорость вывода после начала генерации, но длительность
ожидания пользователя также зависит от задержки до первого токена, длины
рассуждений и времени работы инструментов.

При сравнении кэширования использовался один и тот же длинный префикс: в
одних запросах он сохранялся, а в других изменялось его начало. В
четырёх из пяти пар запрос с сохранённым префиксом оказался быстрее.
Средняя задержка до первого токена в двух группах составила
приблизительно 1,30 и 1,51 секунды соответственно. В следующем разделе
объясняется, чем эти запросы различаются с точки зрения вычислений.

Этот эксперимент также укрепил уверенность автора в перспективах
агентов, работающих на устройстве: небольшие модели уже способны
выполнять подобные задачи с вызовом инструментов непосредственно на
устройстве пользователя. \textbf{Эпоха агентов, работающих на
устройстве, ближе, чем многие предполагают.}
\end{quote}

\hypertarget{ux43fux440ux43eux435ux43aux442ux438ux440ux43eux432ux430ux43dux438ux435-ux43aux43eux43dux442ux435ux43aux441ux442ux430-ux441-ux443ux447ux451ux442ux43eux43c-kv-ux43aux44dux448ux430}{%
\section{Проектирование контекста с учётом
KV-кэша}\label{ux43fux440ux43eux435ux43aux442ux438ux440ux43eux432ux430ux43dux438ux435-ux43aux43eux43dux442ux435ux43aux441ux442ux430-ux441-ux443ux447ux451ux442ux43eux43c-kv-ux43aux44dux448ux430}}

Соседние запросы AI-агента обычно содержат много общей информации:
системные инструкции, определения инструментов и историю предшествующих
взаимодействий. Если движок инференса может повторно использовать
промежуточные результаты вычислений для этих данных, больше
вычислительных ресурсов останется на обработку новой информации.
\textbf{KV-кэш} сохраняет уже вычисленные на каждом слое внимания ключи
(Key) и значения (Value), чтобы использовать их для последующих токенов;
\textbf{кэш промптов (Prompt Cache)} дополнительно позволяет повторно
использовать между запросами кэш подходящих префиксов\footnote{Hugging
  Face, How caching works --- описание механизма послойного сохранения и
  повторного использования K и V.}.

Рассмотрим агента службы поддержки. Чтобы модель знала текущее время,
разработчик добавил в начало системного промпта строку
\texttt{Current\ time:\ \{\{now\}\}}, которая обновляется при каждом
запросе. Последующие бизнес-правила и определения инструментов остаются
неизменными, однако фактические входные данные начинают различаться уже
с временной метки. В сервисах, повторно использующих KV для одинаковых
префиксов, это приводит к повторному вычислению временной метки и всей
следующей за ней части, увеличивая время и стоимость обработки входных
данных.

Если добавлять время отдельным сообщением в конец либо при необходимости
вызывать инструмент времени, существующий префикс можно сохранить. Обе
компоновки предоставляют примерно одинаковую информацию, но требуют
повторных вычислений разного объёма.

Сначала запомним три практических правила, а далее последовательно
разберём лежащие в их основе принципы и эксперименты:

\begin{enumerate}
\def\labelenumi{\arabic{enumi}.}
\tightlist
\item
  \textbf{Сохраняйте повторно используемый префикс стабильным.}
  Используйте детерминированные содержимое и порядок системных
  инструкций и определений часто применяемых инструментов.
\item
  \textbf{Добавляйте изменяющиеся во время выполнения данные в конец.}
  Время, баланс, число вызовов и ход выполнения задачи можно передавать
  в новых сообщениях или результатах инструментов. Для последнего
  состояния следует указывать момент вступления в силу, чтобы модель
  могла определить, какая запись применима к текущему решению.
\item
  \textbf{Организуйте сообщения в соответствии с протоколом модели.}
  Используйте правильные роли, идентификаторы вызовов и формат
  результатов инструментов, а входные данные модели формируйте с помощью
  соответствующего шаблона чата. При сопоставлении кэша учитываются
  обработанный префикс и связанные с ним параметры конфигурации, а
  протокол сообщений также определяет, как модель распознаёт инструкции,
  пользовательские запросы и возвращаемые инструментами значения.
\end{enumerate}

В используемых в этом разделе моделях с причинным вниманием добавление
последующих токенов не изменяет уже вычисленные состояния предыдущих
позиций. При изменении существующих входных данных требуется повторно
вычислить состояния начиная с первой изменившейся позиции; префикс до
неё по-прежнему можно использовать повторно. Фактический диапазон
попадания в кэш на стороне сервера также определяется размером блоков
кэша, временем их хранения и конфигурацией среды выполнения.

\begin{quote}
\textbf{Эксперимент 2-2 ★: визуализация механизма внимания}

Сначала разберём одно вычисление внимания на примере фразы «Какова
погода в Пекине?». Для наглядности на
\crossreflink{fig:2-6}{рисунке 2-6} исходная китайская фраза разделена
на четыре смысловые единицы: «Пекин», показатель связи, «погода» и
«какова»; в реальной модели входные данные разбиваются токенизатором.
При обработке позиции «какова» слой внимания должен сформировать текущее
представление из уже имеющейся информации. Query, Key и Value отвечают
соответственно за запрос, сопоставление и агрегирование информации.

\begin{longtable}[]{@{}
  >{\raggedright\arraybackslash}p{(\columnwidth - 4\tabcolsep) * \real{0.2277}}
  >{\raggedright\arraybackslash}p{(\columnwidth - 4\tabcolsep) * \real{0.3593}}
  >{\raggedright\arraybackslash}p{(\columnwidth - 4\tabcolsep) * \real{0.4130}}@{}}
\toprule\noalign{}
\begin{minipage}[b]{\linewidth}\raggedright
\hspace{0pt}Вектор
\end{minipage} & \begin{minipage}[b]{\linewidth}\raggedright
\hspace{0pt}Назначение
\end{minipage} & \begin{minipage}[b]{\linewidth}\raggedright
\hspace{0pt}Интуитивное понимание в этом примере
\end{minipage} \\
\midrule\noalign{}
\endhead
\bottomrule\noalign{}
\endlastfoot
\hspace{0pt}\textbf{Query (запрос)} & \hspace{0pt}Вектор текущей
позиции, используемый для сопоставления с другими позициями &
\hspace{0pt}Позиции «какова» нужно найти релевантную информацию в уже
имеющемся содержимом \\
\hspace{0pt}\textbf{Key (ключ)} & \hspace{0pt}Вектор каждой позиции,
участвующий в сопоставлении & \hspace{0pt}Позиции «Пекин», «погода» и
другие предоставляют представления для сопоставления \\
\hspace{0pt}\textbf{Value (значение)} & \hspace{0pt}Информационный
вектор, агрегируемый с учётом весов & \hspace{0pt}Информация о месте,
погоде и других аспектах включается в текущее представление в
соответствующих пропорциях \\
\end{longtable}

Key можно представить как метку для поиска, а Value --- как извлекаемое
содержимое. В модели эти векторы являются числовыми величинами,
вычисленными из скрытых представлений, а способ их сопоставления
формируется в процессе обучения.

Одно вычисление состоит из трёх шагов. Сначала для текущей позиции
создаётся Query. Затем вычисляются его скалярные произведения с Key всех
позиций, после чего применяются масштабирование, причинная маска и
softmax, формирующий веса внимания с суммой 1. Наконец, вычисляется
взвешенная сумма Value с этими весами. Скалярное произведение --- это
сумма попарных произведений двух наборов чисел, а softmax преобразует
оценки видимых позиций в нормализованные веса. Векторы Value всех
позиций участвуют в агрегировании с соответствующими весами и вместе
формируют выход данной головки внимания.

Прочитаем по горизонтали последнюю строку
\crossreflink{fig:2-6}{рисунка 2-6}: позиция «какова» присваивает
позиции «погода» вес 0,55. Чтобы получить выход данной головки внимания
в позиции «какова», это число нужно умножить на Value позиции «погода»,
а затем прибавить вклады остальных трёх позиций. Почему в строке
«погода» справа стоит дополнительный крестик? Потому что при вычислении
этой позиции последовательность ещё не дошла до «какова».

\begin{figure}
\hypertarget{fig:2-6}{%
\centering
\includegraphics{images/revised/fig2-6.pdf}
\caption{Рисунок 2-6. Интуитивное представление механизма
внимания}\label{fig:2-6}
}
\end{figure}

Расположив веса всех позиций в виде матрицы, мы получим \textbf{тепловую
карту внимания}. Каждая строка соответствует одному Query, каждый
столбец --- одному Key, а интенсивность цвета обозначает величину веса.
Причинная маска позволяет текущей позиции читать саму себя и
предшествующие позиции. Заблокированные позиции обозначены на рисунке
символом ×; после softmax их веса равны нулю. Такая структура
применяется как при обработке входных данных, так и при последующей
генерации.

Отсюда также видно назначение кэша K и V: по мере продолжения генерации
новая позиция по-прежнему обращается к Key и Value уже существующих
позиций. Если сохранить эти векторы, новая позиция сможет
непосредственно использовать их для сопоставления и агрегирования, без
повторного вычисления исторических позиций.

Разобравшись с этими иллюстративными весами, исследуем распределение
внимания, фактически вычисляемое моделью Qwen3-0.6B. В эксперименте
отдельно записываются короткая фраза и сгенерированная
последовательность, содержащая рассуждение и окончательный ответ, после
чего сравниваются матрицы внимания на разных слоях. На
\crossreflink{fig:2-7}{рисунке 2-7} показаны слои 0, 13 и 27: в левом
столбце находится короткая фраза из девяти токенов, а в правом ---
последовательность вычисления 17 × 23, где отмечены диапазоны входных
данных, рассуждения и ответа. Для градаций серого используется
логарифмическая шкала, позволяющая одновременно различать сильные и
слабые веса внимания.

\begin{figure}
\hypertarget{fig:2-7}{%
\centering
\includegraphics{images/revised/fig2-7.png}
\caption{Рисунок 2-7. Визуализация тепловых карт
внимания}\label{fig:2-7}
}
\end{figure}

При чтении можно анализировать четыре аспекта:

\begin{enumerate}
\def\labelenumi{\arabic{enumi}.}
\tightlist
\item
  \textbf{Сток внимания (Attention Sink).} Некоторые слои присваивают
  большой вес началу последовательности, даже если первые токены несут
  мало смысловой информации. Рассмотрим короткую фразу в левом столбце
  \crossreflink{fig:2-7}{рисунка 2-7}. На каждом изображении сначала
  усреднены все головки внимания одного слоя. Если затем усреднить
  девять чисел в первом столбце матрицы, можно определить, какую долю
  внимания девять позиций в среднем присваивают первому токену. Для
  слоёв 0, 13 и 27 получаются соответственно 14,7\%, 73,9\% и 83,2\%. В
  более длинной сгенерированной последовательности в правом столбце этот
  показатель для слоя 13 снижается до 28,0\%. Softmax требует, чтобы
  сумма весов всех позиций равнялась 1, поэтому модель может обучиться
  использовать относительно стабильную принимающую позицию. Явление и
  его роль в потоковом инференсе подробно исследованы в
  StreamingLLM\footnote{Xiao et al., Efficient Streaming Language Models
    with Attention Sinks, ICLR 2024.}.
\item
  \textbf{Область рассуждения.} Отметьте диапазон генерации,
  соответствующий \texttt{\textless{}think\textgreater{}}, и проследите
  по его строкам, каким словам промпта, описаниям инструментов и
  предшествующим рассуждениям модель уделяет внимание. Треугольная форма
  возникает из-за причинной маски, а различия интенсивности внутри
  треугольника отражают веса отдельных позиций.
\item
  \textbf{Область ответа.} Затем отметьте диапазон генерации
  окончательного ответа и сравните распределение его весов между
  исходными входными данными и предшествующим рассуждением. Это помогает
  определить, какая уже имеющаяся информация вошла в текущее
  агрегирование внимания.
\item
  \textbf{Различия между позициями.} В работе Lost in the Middle
  наблюдалось, что в задачах ответов на вопросы с длинным контекстом и
  поиска по ключам и значениям информацию из средней части контекста
  сложнее эффективно использовать\footnote{Liu et al., Lost in the
    Middle: How Language Models Use Long Contexts, TACL 2024.}. Поэтому
  на практике постоянные правила следует размещать ближе к началу, а
  текущую цель и ход выполнения --- в конце.
\end{enumerate}
\end{quote}

\hypertarget{ux43eux442-ux441ux43eux43eux431ux449ux435ux43dux438ux439-api-ux43a-ux442ux43eux43aux435ux43dux430ux43c-ux43cux43eux434ux435ux43bux438-chat-template}{%
\subsection{От сообщений API к токенам модели: Chat
Template}\label{ux43eux442-ux441ux43eux43eux431ux449ux435ux43dux438ux439-api-ux43a-ux442ux43eux43aux435ux43dux430ux43c-ux43cux43eux434ux435ux43bux438-chat-template}}

В приведённом выше примере API использовался набор сообщений с ролями, а
входные данные на карте внимания представляли собой последовательность
токенов. Преобразование между ними выполняет \textbf{Chat Template
(шаблон чата)}: он добавляет к сообщениям маркеры ролей и границ, в
соответствии с соглашением организует описания инструментов и
возвращаемые значения, а затем передаёт результат токенизатору для
кодирования. Сравните два столбца на
\crossreflink{fig:2-8}{рисунке 2-8}: слева переданы только системная
инструкция и вопрос пользователя, но почему справа появился ещё и
\texttt{assistant}? Это начальный маркер, добавленный шаблоном для
следующего ответа. Эти маркеры ещё понадобятся нам при обсуждении
вызовов инструментов, сохранения сведений о рассуждениях и попаданий в
кэш.

\begin{figure}
\hypertarget{fig:2-8}{%
\centering
\includegraphics{images/revised/fig2-8.pdf}
\caption{Рисунок 2-8. Структура токенов Chat Template}\label{fig:2-8}
}
\end{figure}

У каждого семейства моделей --- Qwen, Llama, Gemma и других --- свой
шаблон чата. Поэтому один и тот же набор сообщений API при передаче
разным моделям может преобразовываться в разные маркеры ролей и границ.
В размещаемых у провайдера API это преобразование обычно выполняется на
стороне сервера, а при самостоятельном развёртывании для сервиса
инференса необходимо настроить шаблон, соответствующий модели.

Предположим, пользователь отправил фразу «Ванкувер 17,9 °C», и
инструмент погоды вернул точно такой же текст. Если по отдельности
передать эти два сообщения шаблону Qwen3-0.6B, получатся два входных
фрагмента, показанных на \crossreflink{fig:2-9}{рисунке 2-9}. Оба
начинаются с
\texttt{\textless{}\textbar{}im\_start\textbar{}\textgreater{}user}.
Какой маркер указывает, что текст получен от инструмента?

\begin{figure}
\hypertarget{fig:2-9}{%
\centering
\includegraphics{images/revised/fig2-9.pdf}
\caption{Рисунок 2-9. Пользовательское сообщение и результат инструмента
в шаблоне Qwen3}\label{fig:2-9}
}
\end{figure}

В правом столбце добавлена пара тегов
\texttt{\textless{}tool\_response\textgreater{}}. Кроме того, шаблон
использует роли сообщений, чтобы найти последний вопрос пользователя: в
раундах до этого вопроса сохраняется только содержимое ответов, а после
него взаимодействие с инструментами может также включать сведения о
рассуждениях модели\footnote{В конфигурационном файле токенизатора
  Qwen3-0.6B \texttt{tokenizer\_config.json} хранится шаблон чата,
  содержащий логику обработки пользовательских запросов, содержимого
  рассуждений и ответов инструментов.}.

Например, после проверки файла модель должна продолжить анализ. Если
вернуть результат с ролью \texttt{tool}, шаблон распознает его как
обратную связь от инструмента в рамках текущей задачи. Если же записать
тот же результат как обычное сообщение \texttt{user}, шаблон может
распознать его как новый запрос и изменить область сохранения предыдущих
рассуждений. Это похоже на необходимость сохранить черновик вычислений,
чтобы продолжить решать задачу: Harness должен сохранять требуемые
сервисом поля рассуждений и продолжать вызовы с правильной структурой
сообщений.

\textbf{По мере перехода от однораундового инференса к долгосрочной
работе AI-агента история рассуждений становится состоянием задачи,
которое необходимо поддерживать.} Ранний шаблон чата DeepSeek-R1 удалял
из исторических ответов часть с рассуждениями и сохранял только
окончательный ответ\footnote{Официальный шаблон чата DeepSeek-R1
  сохраняет в обычных исторических ответах содержимое после
  \texttt{\textless{}/think\textgreater{}}; правила возврата рассуждений
  при последующих вызовах инструментов см. в документации DeepSeek
  Thinking Mode.}. Более поздний протокол вызова инструментов DeepSeek,
напротив, требует возвращать сведения о рассуждениях. После завершения
отдельного вопроса и ответа окончательный ответ уже может содержать
информацию, необходимую для следующего раунда. Но если задача с
инструментами выполнена лишь наполовину, сведения о том, «зачем
проверяется этот файл» и «какие гипотезы уже исключены», могут
присутствовать только в промежуточных рассуждениях. Их сохранение
помогает в следующем раунде продолжить уже начатый ход мысли.

\textbf{Способ возврата сведений о рассуждениях определяется протоколом
модели.} При интеграции необходимо одновременно обрабатывать
\texttt{content}, \texttt{tool\_calls} и предоставляемые сервисом поля
рассуждений. Следующие варианты отражают разные правила продолжения:

\begin{itemize}
\tightlist
\item
  \textbf{Режим рассуждений DeepSeek} возвращает рассуждения через
  \texttt{reasoning\_content}. Согласно протоколу вызова инструментов,
  если запрос содержит \texttt{tools}, последующие запросы должны
  полностью возвращать исторические рассуждения, включая ответы модели
  без вызова инструментов; отсутствие поля приводит к ошибке 400. В
  обычных запросах чата без \texttt{tools} сервис игнорирует
  возвращённые исторические рассуждения\footnote{DeepSeek, Thinking
    Mode.}.
\item
  \textbf{Рассуждения и взаимодействие с инструментами в Kimi} также
  требуют управлять рассуждениями, запросами к инструментам и их
  результатами внутри рабочего цикла. Например, в примере чередующихся
  рассуждений Kimi K2.5 модель может продолжать рассуждать между
  несколькими вызовами инструментов, а клиент сохраняет и возвращает
  сообщения модели в соответствии с протоколом\footnote{Moonshot AI,
    официальный репозиторий Kimi K2.5, раздел Interleaved Thinking and
    Multi-Step Tool Call.}.
\item
  \textbf{Сохранение рассуждений в GLM} можно включить с помощью
  \texttt{thinking.\allowbreak{}clear\_\allowbreak{}thinking:\ false}, после чего необходимо
  полностью возвращать \texttt{reasoning\_content} в исходном порядке.
  Настройки по умолчанию различаются между Coding Plan и стандартным
  API, поэтому клиент должен явно выбрать требуемый режим\footnote{Z.ai,
    Thinking Mode --- описание чередующихся рассуждений, сохранения
    рассуждений и настроек конечных точек по умолчанию.}.
\item
  \textbf{Thinking block в Claude} необходимо сохранять вместе с
  подписью. Модели с поддержкой preserved thinking также проверяют
  источник и префикс блока; конкретные правила приведены далее, в
  разделе «Кэш как архитектурное ограничение».
\end{itemize}

Эти правила влияют как на последовательные вызовы инструментов, так и на
смену модели во время выполнения. При переключении слой адаптации должен
преобразовать запросы к инструментам, их результаты и состояние задачи в
форму, принимаемую целевым интерфейсом, а затем обработать сведения о
рассуждениях согласно правилам целевой модели. Такое продолжение
траектории подробно рассматривается в эксперименте 5-1
\crossreflink{chap:5}{главы 5}.

\textbf{Шаблон также определяет входные данные, сопоставляемые кэшем.}
Одно и то же содержимое сообщений после обработки разными шаблонами, при
разном порядке инструментов или разных способах сериализации может
образовать разные префиксы токенов. Для повторного использования
существующего KV-кэша необходимо сохранять согласованность
соответствующей конфигурации модели и префикса.

\hypertarget{ux43fux440ux438ux43dux446ux438ux43f-ux440ux430ux431ux43eux442ux44b-ux438-ux43eux433ux440ux430ux43dux438ux447ux435ux43dux438ux44f-kv-ux43aux44dux448ux430}{%
\subsection{Принцип работы и ограничения
KV-кэша}\label{ux43fux440ux438ux43dux446ux438ux43f-ux440ux430ux431ux43eux442ux44b-ux438-ux43eux433ux440ux430ux43dux438ux447ux435ux43dux438ux44f-kv-ux43aux44dux448ux430}}

При обработке запроса модель сначала выполняет \textbf{prefill
(предварительное заполнение)}, вычисляя представления уже имеющихся
токенов, а затем переходит к \textbf{decode (декодированию)} и
последовательно генерирует новые токены. Во время декодирования KV-кэш
хранит K и V для уже обработанных позиций, а межзапросное кэширование
префиксов позволяет дополнительно использовать одинаковый префикс,
сохранившийся после предыдущего запроса.

Например, к шестому раунду AI-агент уже накопил 2 000 входных токенов.
Если сервис использует KV-кэш только в пределах текущего запроса, ему
необходимо сначала обработать эти 2 000 токенов и лишь затем продолжить
генерацию с помощью кэша. Если же в межзапросном кэше уже имеется часть
того же префикса, prefill можно продолжить с первой ещё не
закэшированной позиции. Такое повторное использование между запросами
показано на \crossreflink{fig:2-10}{рисунке 2-10}.

\begin{figure}
\hypertarget{fig:2-10}{%
\centering
\includegraphics{images/revised/fig2-10.pdf}
\caption{Рисунок 2-10. Prompt Cache: повторное использование KV-кэша
префикса между запросами}\label{fig:2-10}
}
\end{figure}

\textbf{Проследим за ростом кэша на примере четырёх токенов.}
Предположим, на вход подаётся \texttt{{[}A,\ B,\ C,\ D{]}}. Обрабатывая
позицию D, модель сопоставляет Query токена D с Key токенов A, B, C и D,
а затем вычисляет взвешенную агрегацию их Value. После вычислений во
всех слоях итоговое представление в позиции D используется для
предсказания следующего токена E. К этому моменту K и V токенов от A до
D уже сохранены для каждого слоя.

После семплирования E на следующем шаге он передаётся в модель, которая
вычисляет его Query, Key и Value для каждого слоя. Query токена E
обращается к ранее сохранённому кэшу, а также к K и V текущей позиции,
после чего модель предсказывает шестой токен. Одновременно K и V токена
E добавляются в кэш. Таким образом, на каждом шаге генерации можно
использовать ранее вычисленное состояние и рассчитывать состояние только
для нового входного токена.

Если полностью отказаться от KV-кэша, то для предсказания E придётся
заново обработать токены от A до D, а для предсказания шестого токена
--- токены от A до E. Даже если учитывать только число позиций, для
которых повторно генерируются K и V, совокупный объём работы будет расти
как \(1+2+\cdots+N\). Кэш позволяет повторно использовать послойные
вычисления для исторических позиций, однако новая позиция по-прежнему
должна считывать все доступные ей K и V. При стандартном декодировании с
полным вниманием объём такого чтения и сопоставления линейно растёт с
длиной контекста, поэтому длинный контекст всё равно требует больше
видеопамяти и пропускной способности памяти.

\textbf{Почему изменение входных данных влияет на последующий кэш?} Слои
Transformer последовательно обрабатывают скрытые представления, причём у
каждого слоя есть собственные K и V. После вычисления внимания
представление определённой позиции уже содержит информацию из
предшествующих позиций и в таком виде передаётся следующему слою.
Поэтому при изменении \(k\)-го токена влияние распространяется от этой
позиции вперёд и на более высокие слои.

При неизменных модели и конфигурации вычислений позиции до \(k\) зависят
только от более ранних входных данных и могут использоваться повторно.
Начиная с \(k\), затронутое состояние необходимо создать заново. На
практике сервисы обычно сопоставляют префиксы блоками кэша, поэтому
граница повторного вычисления зависит также от размера блока. Чем ближе
обновление к началу, тем больше входных данных обычно приходится
вычислять заново. Размещение стабильных инструкций в начале и добавление
нового состояния в конец позволяет использовать эту зависимость.

\begin{quote}
\textbf{Эксперимент 2-3 ★★: сравнение способов управления контекстом}

В этом эксперименте агент выполняет одну и ту же задачу с помощью
инструментов обнаружения, чтения и поиска файлов. Сравниваются базовый
вариант со стабильным добавлением данных и пять вариантов изменения
контекста. В базовом варианте порядок инструкций и инструментов остаётся
неизменным, а сообщения модели и результаты инструментов последовательно
добавляются после каждого раунда. В остальных вариантах изменяются
инструкции, состояние, порядок инструментов или представление истории,
чтобы проследить за изменениями статистики кэша и процесса выполнения
задачи.

\textbf{Динамический системный промпт} обновляет текущее время в
системных инструкциях каждого раунда. Часть до изменившегося фрагмента
всё ещё может оставаться общим префиксом, однако входные данные начиная
с информации о времени необходимо сопоставить или вычислить заново.
Время можно добавлять как новое сообщение либо получать по необходимости
с помощью инструмента времени, сохраняя возможность повторного
использования постоянных бизнес-правил.

\textbf{Динамическая пользовательская конфигурация} обновляет в одном из
начальных сообщений оставшееся число вызовов или баланс аккаунта. Само
состояние необходимо для выполнения задачи, поэтому главное проектное
решение заключается в выборе места его обновления. Добавление
актуального состояния в конец с указанием времени вступления в силу
позволяет сохранить и достигнутый прогресс, и префикс. В рассматриваемой
далее панели состояния агента будет использоваться именно такая
организация.

\textbf{Перестановка определений инструментов} меняет порядок списка
инструментов в каждом раунде. Хотя имена и параметры остаются прежними,
порядок токенов после сериализации уже отличается. В отчёте о выполнении
Kimi K2.6 значение \texttt{Hit\%} как для базового варианта, так и для
варианта с перестановкой составляет 100\%. Действительно ли оба варианта
повторно использовали одинаковый объём входных данных? За три раунда
базовый вариант получил в общей сложности 7 567 входных токенов, из
которых 768 были найдены в кэше. Вариант с перестановкой также выполнил
три раунда, но из 7 568 входных токенов в кэше нашлись лишь 256. Если
разделить число закэшированных токенов на число входных, получатся
значения 10,1\% и 3,4\% в столбце \texttt{Cache\%}. Другой столбец,
\texttt{Hit\%}, подсчитывает запросы, для которых зарегистрировано
попадание: в базовом варианте отмечены два попадания, а в варианте с
перестановкой --- только одно; при этом в обоих вариантах
зарегистрировано ноль промахов. По одному этому отчёту невозможно
определить наличие попаданий в остальных запросах. Фиксированный порядок
инструментов позволяет избежать необоснованных изменений префикса.

\textbf{Скользящее окно} сохраняет только последний фрагмент истории.
Например, если хранить десять последних раундов взаимодействия, то в
пятнадцатом раунде могут оказаться недоступны важные данные из файла,
полученные во втором раунде. Обрезка не только сокращает общий префикс,
но и меняет доступные модели факты и сведения о ходе выполнения задачи.
В сопутствующей реализации отправной точкой служат шесть последних
сообщений истории, к которым спереди добавляются запросы на вызов
инструментов, соответствующие сохранённым результатам. Вариант со
скользящим окном повторно использовал 1 510 из 2 224 входных токенов,
достигнув максимальной среди шести вариантов доли кэша --- 67,9\%.
Однако журнал инструментов этого варианта в эксперименте 2-3 показывает,
что модель трижды искала файлы, затем прочитала два файла, но так и не
предоставила запрошенное пользователем резюме из трёх предложений.

\textbf{Текстовое форматирование} преобразует историю с ролями в текст
вида \texttt{USER:\ ...}, \texttt{ASSISTANT:\ ...} и т. п. Это изменяет
границы сообщений и структуру инструментов в шаблоне, поэтому модели
приходится восстанавливать исходные связи из обычного текста. В данном
эксперименте текстовый вариант всё же сформировал итоговый ответ и
получил попадания в кэш. Предоставляемые интерфейсом роли сообщений, ID
вызовов и поля результатов инструментов помогают стабильно передавать
эти связи.

Даже если программа заново формирует список сообщений, его префикс
по-прежнему может использоваться повторно при условии, что отправляемый
префикс и конфигурация кэша остаются неизменными.

Чтобы определить, насколько сократилось время ожидания пользователя,
сначала необходимо выяснить, в какой момент останавливается таймер. В
первом раунде варианта со скользящим окном значение \texttt{ttft}
составило около 2,23 секунды. Хотя программа называет это поле TTFT ---
задержкой до первого токена, --- таймер останавливается лишь после того,
как нестриминговый запрос возвращает ответ целиком.
\end{quote}

\hypertarget{kv-ux43aux44dux448-ux438-prompt-cache-ux434ux432ux430-ux443ux440ux43eux432ux43dux44f-ux43aux44dux448ux438ux440ux43eux432ux430ux43dux438ux44f}{%
\subsection{KV-кэш и Prompt Cache: два уровня
кэширования}\label{kv-ux43aux44dux448-ux438-prompt-cache-ux434ux432ux430-ux443ux440ux43eux432ux43dux44f-ux43aux44dux448ux438ux440ux43eux432ux430ux43dux438ux44f}}

Эти два уровня кэширования можно различить по двум моментам времени.
Когда модель генерирует второй токен ответа, она может прочитать
состояния ключей и значений уже существующих позиций этого запроса ---
здесь используется \textbf{KV-кэш}. Позднее пользователь задаёт
уточняющий вопрос, и Harness отправляет новый запрос; чтобы сократить
вычисления для этой части входных данных, сервис сначала должен найти
ранее сохранённый идентичный префикс --- это относится к \textbf{Prompt
Cache}. Первый объясняет, как продолжается генерация в рамках одного
запроса, а второй дополнительно требует механизма сохранения и
сопоставления между запросами.

В сервисах API на кэширование также влияют минимальная длина префикса,
размер блока, срок действия, конфигурация модели и область изоляции
арендаторов. Повторно используемый префикс необходимо искать в той же
области кэширования. Например, Claude API изолирует кэш по организациям
и рабочим пространствам; если одинаковый текст встречается в разных
областях изоляции, для каждой из них создаётся отдельный кэш\footnote{Anthropic,
  «Prompt caching», описание сопоставления кэша, изоляции, срока
  действия и тарификации.}.

Стоимость записи и чтения следует учитывать отдельно. Для некоторых
моделей Claude цена чтения из кэша составляет одну десятую от обычной
цены обработки входных данных, а за первоначальную запись взимается
дополнительная плата. Если один и тот же префикс многократно
используется в течение срока действия кэша, запись оплачивается только
один раз, а все последующие запросы с попаданием в кэш тарифицируются по
цене чтения.

\hypertarget{ux43aux44dux448-ux43aux430ux43a-ux430ux440ux445ux438ux442ux435ux43aux442ux443ux440ux43dux43eux435-ux43eux433ux440ux430ux43dux438ux447ux435ux43dux438ux435}{%
\subsection{Кэш как архитектурное
ограничение}\label{ux43aux44dux448-ux43aux430ux43a-ux430ux440ux445ux438ux442ux435ux43aux442ux443ux440ux43dux43eux435-ux43eux433ux440ux430ux43dux438ux447ux435ux43dux438ux435}}

\textbf{В промышленной системе кэш уже стал архитектурным ограничением.}
Порядок фрагментов промпта, начальная точка дочернего агента и способ
восстановления сеанса влияют на объём вычислений, который можно
использовать повторно. \textbf{Границы кэша следует определять с самого
начала проектирования.} Если сначала установить, какое содержимое может
быть общим и какое состояние будет изменяться, а затем организовать
промпт и структуру сеанса, можно сократить повторные вычисления во время
работы. Если перестраивать их лишь после того, как задержка и стоимость
долгосрочных задач стали проблемой, изменения зачастую затронут
управление состоянием, дочерних агентов и восстановление сеансов.

\textbf{Располагайте префикс с учётом стабильности содержимого.} Общие
для нескольких задач инструкции и определения инструментов можно
поместить в начало, а затем добавить профиль пользователя, состояние
сеанса и новые наблюдения, сгруппировав их по области применимости.
Например, если операционная система выбирается из macOS и Linux, режим
работы --- из обычного режима и режима отладки, а язык --- из китайского
и английского, три независимых условия создадут не более
\(2\times2\times2=8\) комбинаций префикса. Размещая чаще изменяющиеся
поля после общей части, можно увеличить вероятность повторного
использования общего префикса. При совместном использовании общей части
также необходимо соблюдать установленные сервисом границы организации,
рабочего пространства и конфиденциальности.

\textbf{Выбирайте для дочернего агента начальную точку, допускающую
повторное использование.} Если дочерний агент наследует контекст задачи
родительского агента, он может повторно использовать совпадающий у обоих
префикс, а затем добавить собственную подзадачу. Область повторного
использования совместно определяется промптом, порядком инструментов,
моделью и соответствующей конфигурацией. Для дочернего агента, которому
требуется отдельная роль или изолированный контекст, входные данные
формируются согласно его задаче, после чего следует проверить, какие
общие части допустимо использовать совместно.

\textbf{Сохраняйте фактически отправленное содержимое.} После сжатия
большого вывода инструмента в сводку или предварительное представление
следует сохранять версию, фактически переданную модели. Её
воспроизведение при восстановлении сеанса позволяет воссоздать
согласованную последовательность сообщений. Если при каждом
восстановлении генерировать сводку заново, могут появиться новые
формулировки и изменения префикса. При обновлении сводки можно
опубликовать новую версию и явно указать, с какого места далее следует
заново формировать кэш.

\textbf{Сохранение рассуждений также требует соблюдения ограничений
протокола.} Конкретным примером служит механизм preserved thinking в
Claude. Согласно документации, он применяется для защиты от дистилляции:
сервис проверяет подпись thinking block, модель, которая его
сгенерировала, и соответствующий префикс. Для моделей с привязкой к
префиксу, таких как Claude Fable 5.1 и Claude Opus 5.5, область проверки
образуют предшествующие блоку верхнеуровневые \texttt{system},
\texttt{tools} и сообщения. Изменение их содержимого делает
недействительными этот и последующие thinking block\footnote{Anthropic,
  Preserved thinking: модели, привязка к префиксу, обработка
  недействительных блоков и правила обрезки блоков рассуждений;
  интерфейсы для обновлений во время сеанса описаны в Mid-conversation
  system messages and tool changes.}.

Например, перезапись системного времени в каждом раунде, повторный
рендеринг сведений о среде в первом сообщении и обрезка старого
результата инструмента на месте изменяют префикс, к которому привязаны
последующие блоки рассуждений. После срабатывания проверки API может
вернуть ошибку 400 либо, в соответствии с настройкой
\texttt{prefix\_mismatch\_behavior}, отбросить недействительные блоки и
продолжить обработку запроса. Во втором случае модель по-прежнему может
использовать сохранённые сведения о задаче и результаты инструментов,
однако недействительные рассуждения больше не входят в текущие входные
данные. Для аккаунтов, созданных начиная с 31 августа 2026 года,
проверка префикса включена по умолчанию; способ её включения для более
ранних аккаунтов приведён в описании протокола.

При интеграции такого протокола верхнеуровневые инструкции и набор
инструментов сеанса можно сохранять стабильными, добавляя изменения
состояния в конец. Если инструкции или инструменты необходимо обновить
во время сеанса, используйте предоставляемые сервисом интерфейсы для
внутрисеансовых системных сообщений, блоков изменения инструментов и
подобных операций. При обрезке thinking block также необходимо сохранять
непрерывный интервал в исходном порядке: блоки можно удалить с начала, с
конца либо удалить полностью, но разрыв в середине повлияет на
последующие блоки. Блоки, удалённые намеренно, также должны оставаться
удалёнными.

При смене модели во время выполнения сервис также проверяет, может ли
целевая модель читать блоки рассуждений исходной модели. Клиенту следует
сохранять исходную историю, позволяя сервису обрабатывать её по правилам
совместимости. Если клиент самостоятельно удалит эти записи,
восстановить их при последующем возврате к исходной модели уже не
удастся.

\hypertarget{ux440ux430ux441ux448ux438ux440ux435ux43dux438ux435-ux438ux441ux441ux43bux435ux434ux43eux432ux430ux43dux438ux44f-ux440ux435ux434ux430ux43aux442ux438ux440ux443ux435ux43cux44bux439-ux438-ux43aux43eux43cux43fux43eux43dux443ux435ux43cux44bux439-kv-ux43aux44dux448}{%
\subsection{Расширение исследования: редактируемый и компонуемый
KV-кэш}\label{ux440ux430ux441ux448ux438ux440ux435ux43dux438ux435-ux438ux441ux441ux43bux435ux434ux43eux432ux430ux43dux438ux44f-ux440ux435ux434ux430ux43aux442ux438ux440ux443ux435ux43cux44bux439-ux438-ux43aux43eux43cux43fux43eux43dux443ux435ux43cux44bux439-kv-ux43aux44dux448}}

\begin{quote}
Во всех описанных выше решениях вычисления повторно используются за счёт
совпадения префиксов. Автор предлагает исследовать следующий вопрос:
можно ли напрямую редактировать и компоновать уже сохранённую в KV-кэше
информацию, чтобы кэш адаптировался к изменениям контекста? В работе
«Models Take Notes at Prefill» автор сравнивает такую информацию с
«заметками» модели\footnote{Li, Bojie. Models Take Notes at Prefill: KV
  Cache Can Be Editable and Composable. 2026. В экспериментах с
  редактированием, композицией и онлайн-кэшированием использовались
  разные задачи и метрики; подробные условия приведены в работе.}: после
того как на этапе prefill выполняются многослойные вычисления внимания,
влияние некоторого поля уже оказывается встроено в представления
последующих позиций. Например, поле «Город пользователя: Пекин» может
влиять на последующие решения о местоположении. Если заменить только K и
V самого поля, последующие состояния всё равно могут сохранять выводы,
сделанные на основе прежнего города.

С помощью вмешательств в работе разграничивается влияние KV самого поля
и последующих состояний. В исследованных задачах принятия решений
непосредственный вклад KV самого поля составлял менее 1\%, а основная
часть влияния сохранялась в последующих представлениях. Эти оставленные
далее по цепочке «заметки» позволяют исследовать два способа
использования кэша.

\textbf{Редактирование (Editing)} обновляет уже сделанные выводы в
последующих вычислениях посредством явного исправления информации. В
работе также показано, что в исследованных задачах для моделей размера
8B при наличии явной цепочки рассуждений редактирования одного поля
достаточно для восстановления целевого решения, а вычислительные затраты
составляют около 1\% от полного перерасчёта. Без такого процесса
рассуждения прежние последующие состояния могут продолжать влиять на
решение, из-за чего изменение поля не даёт нужного эффекта.

\textbf{Композиция (Composition)} подключает заранее вычисленный кэш
навыка к новому контексту. Для моделей с вращательным позиционным
кодированием (RoPE) перед объединением необходимо скорректировать
представления позиций в кэше. Это можно представить как перенумерацию
уже готовой страницы заметок перед её добавлением к другому материалу. В
работе сравнивались двенадцать моделей: косинусное сходство выходных
логитов после объединения и полного перерасчёта составляло 0,90--0,999.
Перепозиционирование и само объединение кэша можно выполнять за линейное
по длине время, сокращая объём повторных вычислений внимания для всего
контекста.

Добавление исправляющей информации также можно сочетать с существующим
кэшированием префиксов. В онлайн-бенчмарке vLLM из этой работы такая
комбинация сохранила около 98,5\% попаданий в кэш префиксов, а задержка
до первого токена на уровне p90 составила примерно от 1/398 до 1/53
относительно базового варианта.

\textbf{Автор считает это направление исследований перспективным,
поскольку оно может позволить сохранять больше уже выполненных
вычислений при обновлении контекста.} AI-агент часто обновляет поля
памяти, добавляет состояние или заменяет определения инструментов. Если
соответствующие состояния удастся исправлять и компоновать подобно
заметкам, затраты на поддержку длинного контекста можно будет
значительно снизить.
\end{quote}

\hypertarget{ux43eux442-ux43cux435ux445ux430ux43dux438ux437ux43cux430-ux43aux44dux448ux438ux440ux43eux432ux430ux43dux438ux44f-ux43a-ux43fux440ux43eux435ux43aux442ux438ux440ux43eux432ux430ux43dux438ux44e-ux441ux43eux434ux435ux440ux436ux438ux43cux43eux433ux43e-ux43aux43eux43dux442ux435ux43aux441ux442ux430}{%
\subsection{От механизма кэширования к проектированию содержимого
контекста}\label{ux43eux442-ux43cux435ux445ux430ux43dux438ux437ux43cux430-ux43aux44dux448ux438ux440ux43eux432ux430ux43dux438ux44f-ux43a-ux43fux440ux43eux435ux43aux442ux438ux440ux43eux432ux430ux43dux438ux44e-ux441ux43eux434ux435ux440ux436ux438ux43cux43eux433ux43e-ux43aux43eux43dux442ux435ux43aux441ux442ux430}}

Разобравшись, как кэш позволяет повторно использовать результаты
вычислений, нужно решить, что включать во входные данные на каждом
раунде и как обновлять это содержимое. Дальнейший материал состоит из
трёх частей:

\begin{itemize}
\tightlist
\item
  \textbf{Промпт-инжиниринг и Agent Skills}: сначала проектируются
  системные инструкции и определения инструментов, описывающие цели,
  правила и способы работы; затем предметные знания загружаются по мере
  необходимости с помощью постепенного раскрытия. Защита от инъекций в
  промпт сосредоточена на источниках внешнего содержимого и границах
  инструкций при его добавлении в контекст.
\item
  \textbf{Строка состояния агента}: предоставляет модели такие сведения,
  как ход выполнения задачи, сводка наблюдений за окружением и
  количество вызовов, сокращая затраты на их повторное выведение из
  длинной истории. Она подобна индикаторам времени, заряда и состояния
  сети на телефоне, благодаря которым текущие условия работы остаются
  видимыми.
\item
  \textbf{Сжатие контекста}: по мере роста истории сохраняет информацию,
  необходимую для последующих решений, определяет, когда выполнять
  сжатие, как организовать сводку, а также как обрабатывать повторные
  вычисления кэша и продолжение задачи.
\end{itemize}

\hypertarget{ux43fux440ux43eux43cux43fux442-ux438ux43dux436ux438ux43dux438ux440ux438ux43dux433-ux43eux43fux442ux438ux43cux438ux437ux430ux446ux438ux44f-ux441ux438ux441ux442ux435ux43cux43dux43eux433ux43e-ux43fux440ux43eux43cux43fux442ux430}{%
\section{Промпт-инжиниринг: оптимизация системного
промпта}\label{ux43fux440ux43eux43cux43fux442-ux438ux43dux436ux438ux43dux438ux440ux438ux43dux433-ux43eux43fux442ux438ux43cux438ux437ux430ux446ux438ux44f-ux441ux438ux441ux442ux435ux43cux43dux43eux433ux43e-ux43fux440ux43eux43cux43fux442ux430}}

Промпт-инжиниринг с помощью инструкций, процессов и примеров направляет
модель на выполнение конкретных задач. В этом разделе основное внимание
уделяется \textbf{системному промпту (System Prompt)}: для AI-агента это
своего рода руководство сотрудника, описывающее его роль, цели, правила
поведения и рабочие процессы. В примерах API из этой главы такое
содержимое передаётся в сообщении \texttt{system}.

\hypertarget{ux442ux43eux43d-ux438-ux441ux442ux438ux43bux44c-ux43aux430ux43a-ux43eux431ux440ux430ux449ux430ux442ux44cux441ux44f-ux43a-ux43fux43eux43bux44cux437ux43eux432ux430ux442ux435ux43bux44e}{%
\subsection{Тон и стиль: как обращаться к
пользователю}\label{ux442ux43eux43d-ux438-ux441ux442ux438ux43bux44c-ux43aux430ux43a-ux43eux431ux440ux430ux449ux430ux442ux44cux441ux44f-ux43a-ux43fux43eux43bux44cux437ux43eux432ux430ux442ux435ux43bux44e}}

Требования к тону должны воплощаться в конкретных формулировках.
Например, для кратких ответов можно потребовать: «Вы ДОЛЖНЫ отвечать
кратко, менее чем в 4 строках». Если выполнение задачи заблокировано,
можно указать: «Ограничьте ответ 1--2 предложениями», --- чтобы сообщить
текущее состояние и следующий шаг, сразу показав пользователю, что
требуется сделать.

Выделение прописными буквами, например «НИКОГДА не делайте X», можно
использовать для подчёркивания ключевых требований. Но ещё важнее
пояснить, в каких ситуациях действует требование и какие действия
следует предпринять.

\hypertarget{ux441ux442ux440ux443ux43aux442ux443ux440ux438ux440ux43eux432ux430ux43dux43dux44bux435-ux43fux440ux43eux43cux43fux442ux44b-ux438ux435ux440ux430ux440ux445ux438ux44f-ux438-ux433ux440ux430ux43dux438ux446ux44b-ux441ux43eux434ux435ux440ux436ux438ux43cux43eux433ux43e}{%
\subsection{Структурированные промпты: иерархия и границы
содержимого}\label{ux441ux442ux440ux443ux43aux442ux443ux440ux438ux440ux43eux432ux430ux43dux43dux44bux435-ux43fux440ux43eux43cux43fux442ux44b-ux438ux435ux440ux430ux440ux445ux438ux44f-ux438-ux433ux440ux430ux43dux438ux446ux44b-ux441ux43eux434ux435ux440ux436ux438ux43cux43eux433ux43e}}

Заголовки, списки и теги помогают организовать длинные промпты.
Заголовки Markdown подходят для обозначения разделов и шагов, а
XML-подобные теги --- для разделения блоков разного назначения.
Например, информацию о рабочем каталоге можно заключить в
\texttt{\textless{}working\_directory\textgreater{}}. Название тега
должно прямо указывать назначение содержимого, чтобы модели было легче
распознавать связи между блоками.

В одном промпте можно сочетать оба подхода:

\begin{Shaded}
\begin{Highlighting}[]
\NormalTok{\# Правила использования инструментов}

\NormalTok{\#\# Операции с файлами}
\NormalTok{\textless{}file\_operation\textgreater{}}
\NormalTok{{-} Перед чтением файла проверить, существует ли путь}
\NormalTok{{-} Перед изменением существующего файла сохранить резервную копию}
\NormalTok{\textless{}/file\_operation\textgreater{}}

\NormalTok{\#\# Сетевые запросы}
\NormalTok{\textless{}network\_request\textgreater{}}
\NormalTok{{-} Установить тайм{-}аут 30 секунд}
\NormalTok{{-} При кратковременных сбоях повторить попытку не более 3 раз}
\NormalTok{{-} Для запросов с внешними побочными эффектами перед повторной попыткой сначала проверить состояние выполнения}
\NormalTok{\textless{}/network\_request\textgreater{}}
\end{Highlighting}
\end{Shaded}

Markdown создаёт удобочитаемую иерархию, а
\texttt{\textless{}file\_operation\textgreater{}} и
\texttt{\textless{}network\_request\textgreater{}} разграничивают
области применения правил. Например, правило «повторить попытку не более
3 раз» относится только к сетевым запросам, тогда как при изменении
файлов следует соблюдать правила резервного копирования из предыдущего
блока.

\hypertarget{ux43eux440ux433ux430ux43dux438ux437ux430ux446ux438ux44f-ux43fux440ux430ux432ux438ux43b-ux432-ux432ux438ux434ux435-ux43fux440ux43eux446ux435ux441ux441ux430}{%
\subsection{Организация правил в виде
процесса}\label{ux43eux440ux433ux430ux43dux438ux437ux430ux446ux438ux44f-ux43fux440ux430ux432ux438ux43b-ux432-ux432ux438ux434ux435-ux43fux440ux43eux446ux435ux441ux441ux430}}

\textbf{Если задача состоит из чётко определённых этапов, промпт лучше
организовать в виде процесса.} Например, чтобы оформить возврат, сначала
нужно найти заказ и лишь затем на основании правил возврата определить,
можно ли вернуть деньги. Если заказ не найден, следует запросить у
пользователя номер заказа и другие сведения; если заказ соответствует
условиям --- инициировать возврат. Если те же правила разбросаны по
списку из сотни пунктов, на каждом шаге приходится заново искать
применимые положения. Когда входные условия, последовательность действий
и выходы при ошибках описаны на соответствующих этапах, модель может
определить, где находится процесс и что следует делать дальше.

Ниже на примере обработки файлов проверка, классификация,
предварительная обработка, выполнение и проверка результата объединены в
стандартную операционную процедуру (SOP):

\begin{Shaded}
\begin{Highlighting}[]
\NormalTok{Процесс обработки файлов}

\NormalTok{Шаг 1: Проверка}
\NormalTok{  Убедиться, что файл существует и доступен}
\NormalTok{  Файл не существует → записать ошибку в журнал и остановить обработку}
\NormalTok{  ↓}
\NormalTok{Шаг 2: Классификация}
\NormalTok{  Определить тип по расширению и содержимому файла}
\NormalTok{  ↓}
\NormalTok{Шаг 3: Предварительная обработка}
\NormalTok{  Файл конфигурации → сначала создать резервную копию}
\NormalTok{  Большой файл (\textgreater{}1 MB) → использовать потоковую обработку}
\NormalTok{  ↓}
\NormalTok{Шаг 4: Выполнение}
\NormalTok{  Выполнить соответствующую обработку с учётом типа файла}
\NormalTok{  ↓}
\NormalTok{Шаг 5: Проверка результата}
\NormalTok{  Проверить целостность обработанного файла}
\end{Highlighting}
\end{Shaded}

На каждом шаге задаётся локальная цель, а относящиеся к ней правила
размещаются непосредственно под ним. Если файл не существует, процесс
завершается на этапе проверки; если требуется изменить файл
конфигурации, резервная копия сохраняется на этапе предварительной
обработки. Сведения о ходе выполнения модель получает из истории
сообщений или записи состояния.

\hypertarget{ux434ux435ux442ux430ux43bux438ux437ux430ux446ux438ux44f-ux431ux438ux437ux43dux435ux441-ux43fux440ux430ux432ux438ux43b-ux43eux434ux43dux43eux437ux43dux430ux447ux43dux43eux435-ux43eux43fux438ux441ux430ux43dux438ux435-ux43fux440ux43eux434ux443ux43aux442ux43eux432ux44bux445-ux434ux43eux433ux43eux432ux43eux440ux451ux43dux43dux43eux441ux442ux435ux439}{%
\subsection{Детализация бизнес-правил: однозначное описание продуктовых
договорённостей}\label{ux434ux435ux442ux430ux43bux438ux437ux430ux446ux438ux44f-ux431ux438ux437ux43dux435ux441-ux43fux440ux430ux432ux438ux43b-ux43eux434ux43dux43eux437ux43dux430ux447ux43dux43eux435-ux43eux43fux438ux441ux430ux43dux438ux435-ux43fux440ux43eux434ux443ux43aux442ux43eux432ux44bux445-ux434ux43eux433ux43eux432ux43eux440ux451ux43dux43dux43eux441ux442ux435ux439}}

К какой категории отнести запрос на основании одних и тех же сведений о
заказе, какую услугу предоставить и какую плату взимать --- прежде всего
бизнес-правила, которые должны однозначно определить продуктовая и
операционная команды. Затем обязанности разделяются на уровне инженерной
реализации: промпт сообщает модели, какие поля нужно получить и к какому
сервису обратиться, а инструменты и бизнес-системы выполняют
детерминированные расчёты и применяют ограничения.

Рассмотрим агента, который звонит от имени пользователя, чтобы снизить
стоимость подписки или запросить возврат. Команда хочет снизить барьер
для попытки с помощью правила «если результата нет, деньги
возвращаются», одновременно установив чёткий порядок оплаты особых
задач. Можно предусмотреть три модели тарификации:

\begin{itemize}
\tightlist
\item
  \textbf{Комиссия с сэкономленной суммы}: если в результате переговоров
  существующий счёт снижается, взимается определённая доля фактически
  сэкономленной суммы, например 20\%.
\item
  \textbf{Фиксированная плата за услугу}: за сервисные задачи, такие как
  бронирование ресторана, взимается плата с учётом типа или сложности
  задачи.
\item
  \textbf{Предоплата за сложную задачу}: для сложных задач, отвечающих
  определённым условиям, заранее сообщается размер предоплаты и условие
  о её невозврате, а задача принимается после подтверждения
  пользователя.
\end{itemize}

Формулировка «выберите подходящий тип тарификации с учётом задачи»
по-прежнему оставляет неоднозначность классификации. Запрос «Помоги мне
вернуть одежду, купленную в прошлом месяце» связан с возвратом уже
уплаченных средств, а «Помоги мне отменить подписку Netflix» --- с
прекращением будущих платежей. В широком смысле оба результата можно
считать экономией, однако команда может ограничить комиссионную модель
переговорами о снижении существующего счёта, а для возвратов и отмены
услуг установить единую фиксированную плату.

Соответствующую границу категорий можно прямо сформулировать в промпте:
«НИКОГДА не используйте percentage\_based\_one\_time для возвратов и
отмены услуг. Вместо этого используйте fixed\_fee». Иными словами, для
возвратов и отмены услуг применяется \texttt{fixed\_fee}, а комиссионная
модель --- только для снижения существующего счёта путём переговоров.

Команде также необходимо оценивать вероятность успешного выполнения
задачи. На способ принятия задачи влияют условия, на которых компания
удовлетворяет обращения, соответствие пользователя этим условиям, состав
начислений в счёте и результаты выполнения аналогичных поручений в
прошлом. Если правила устанавливают пороги «выше 60\% --- режим с
возвратом средств, ниже 30\% --- отказ от задачи», необходимо также
определить порядок обработки промежуточного диапазона и самих граничных
значений.

При расчёте сумм следует явно задать единицы, базу расчёта и порядок
округления. Например, стоимость телефонного разговора накапливается из
расчёта 0,05 доллара за минуту, а итоговая сумма округляется до целого
числа долларов; экономия рассчитывается относительно существующего
счёта. Если текущий счёт составляет 150 долларов, а компания может в
следующем году повысить его до 180 долларов, предотвращённое сохранением
текущей цены потенциальное повышение следует учитывать по заранее
согласованным правилам, а не автоматически включать в сумму экономии от
снижения существующего счёта.

Автор рекомендует, чтобы \textbf{проектированием бизнес-правил руководил
ответственный за продукт}, поддерживая определения на основе данных
промышленной эксплуатации, отзывов пользователей и операционного опыта.
Затем инженеры реализуют эти правила в промптах, параметрах
инструментов, логике тарификации и правилах проверки.

\hypertarget{few-shot-ux43fux440ux438ux43cux435ux440ux44b-ux434ux435ux43cux43eux43dux441ux442ux440ux430ux446ux438ux44f-ux43eux436ux438ux434ux430ux435ux43cux43eux433ux43e-ux43fux43eux432ux435ux434ux435ux43dux438ux44f}{%
\subsection{Few-shot-примеры: демонстрация ожидаемого
поведения}\label{few-shot-ux43fux440ux438ux43cux435ux440ux44b-ux434ux435ux43cux43eux43dux441ux442ux440ux430ux446ux438ux44f-ux43eux436ux438ux434ux430ux435ux43cux43eux433ux43e-ux43fux43eux432ux435ux434ux435ux43dux438ux44f}}

Если требования к результату трудно выразить несколькими правилами,
\textbf{лучше привести два-три качественных примера, чем продолжать
нагромождать абстрактные объяснения.} Такой подход называется
\textbf{малочисленными примерами (few-shot examples)}. Тексты
определённого стиля, структурированные отчёты и тон ответов службы
поддержки удобно показывать в виде пар из входных данных и ожидаемого
результата. Модель распознаёт шаблон по примерам из текущего контекста и
применяет его к новым запросам.

Примеры можно поместить в системный промпт как общее содержимое для
множества запросов. Их также можно оформить как демонстрационные
сообщения \texttt{user}/\texttt{assistant}, загружаемые в начале
диалога. \textbf{При многократном выполнении однотипных задач автор
предпочитает использовать отобранный фиксированный набор примеров.}
Фиксированные примеры также сохраняют стабильность префикса. Если задачи
существенно различаются и фиксированный набор не обеспечивает
достаточного покрытия, можно рассмотреть поиск примеров с учётом
конкретной задачи. При динамическом выборе общие инструкции можно
оставить в начале, а выбранные примеры добавлять ближе к концу.

Для пояснения легко смешиваемых границ в описанных выше правилах
тарификации достаточно двух примеров. Если пользователь просит
договориться о снижении ежемесячной платы с 80 до 60 долларов, можно
взимать комиссию с сэкономленной суммы. Если же пользователь просит
отменить услугу стоимостью 80 долларов в месяц, взимается фиксированная
плата за отмену.

\hypertarget{ux43fux440ux43eux435ux43aux442ux438ux440ux43eux432ux430ux43dux438ux435-ux43eux43fux440ux435ux434ux435ux43bux435ux43dux438ux439-ux438ux43dux441ux442ux440ux443ux43cux435ux43dux442ux43eux432}{%
\subsection{Проектирование определений
инструментов}\label{ux43fux440ux43eux435ux43aux442ux438ux440ux43eux432ux430ux43dux438ux435-ux43eux43fux440ux435ux434ux435ux43bux435ux43dux438ux439-ux438ux43dux441ux442ux440ux443ux43cux435ux43dux442ux43eux432}}

Определение инструмента --- это руководство по его использованию для
агента. По названию и описанию модель должна определить, когда вызывать
инструмент, а затем заполнить запрос в соответствии со структурой
параметров. Понятное определение должно сообщать модели, что может
сделать инструмент, какие входные данные ему нужны, что он вернёт и
какие условия должны быть выполнены перед вызовом.

Эти сведения можно формулировать очень конкретно. Например, для
инструмента поиска можно указать: «При поиске содержимого файлов
используйте специализированный инструмент поиска»; для инструмента
времени --- привести \texttt{timezone:\ "America/New\_York"} как пример
фактической записи параметра; для инструмента чтения --- подсказать:
«Независимые друг от друга операции чтения можно вызывать пакетно»; для
инструмента редактирования --- потребовать: «Сначала прочитайте целевой
файл, затем отредактируйте его с учётом текущего содержимого». Эти
четыре типа указаний соответственно описывают границы применения, формат
параметров, эффективность выполнения и предварительные зависимости.
Проектирование интерфейсов инструментов будет рассмотрено на
дополнительных примерах в четвёртой главе.

Если инструментов немного, вместе с запросом можно предоставлять полные
определения \texttt{tools}, сохраняя их названия, содержимое и порядок в
последующих раундах. При увеличении числа инструментов полные
определения занимают значительную часть контекста и усложняют для модели
различение похожих инструментов. В таком случае можно использовать
\textbf{загрузку по запросу}: сначала сообщить модели о доступных
возможностях, а при необходимости найти и загрузить полную schema
соответствующих инструментов. Рассматриваемые далее Skills организованы
по тому же принципу: сначала предоставляется каталог, а затем
загружаются подробности.

OpenAI Responses API помечает инструменты с отложенной загрузкой с
помощью \texttt{defer\_loading:\ true}. В управляемом клиентом процессе
поиска инструментов модель выдаёт \texttt{tool\_search\_call}, а клиент
возвращает выбранные определения инструментов через
\texttt{tool\_search\_output}; часто используемые инструменты можно
предоставлять напрямую и применять вместе с загружаемыми по
запросу\footnote{OpenAI, Tool search, документация Responses API.}.
Claude Code в конфигурациях модели и сервиса с поддержкой поиска
инструментов по умолчанию загружает инструменты MCP по запросу: в начале
сеанса предоставляются названия инструментов и описания серверов, а
полные определения загружаются после поиска. Также доступен пороговый
режим, включающий отложенную загрузку, когда определения инструментов
достигают заданной доли контекста\footnote{Anthropic, Scale with MCP
  tool search, документация Claude Code.}. Шаблон обнаружения
инструментов Apps/Connectors в Codex использует поиск BM25 по метаданным
инструментов и предоставляет найденные инструменты следующему вызову
модели\footnote{OpenAI Codex, шаблон обнаружения инструментов
  Apps/Connectors.}.

Загрузку по запросу необходимо согласовать с кэшированием префикса. В
интерфейсах, позволяющих добавлять определения инструментов в ходе
сеанса, существующее содержимое остаётся на месте, а новые определения
вместе с результатами обнаружения добавляются в последующий контекст.
При причинном внимании K и V каждого токена зависят только от него
самого и предшествующего ему содержимого, поэтому добавление данных в
конец сохраняет условия для повторного использования предыдущего
префикса. В последующих раундах сообщения продолжают добавляться после
этой истории, а уже загруженные определения остаются на исходных
позициях. Некоторые интерфейсы добавляют новые инструменты в
расположенный в начале запроса набор \texttt{tools}. В таком случае
исторические сообщения, находящиеся после новых определений, приходится
вычислять заново, даже если они не изменились.

\begin{quote}
\textbf{Эксперимент 2-4 ★★: абляционный эксперимент в
промпт-инжиниринге}

Какую роль в промпте играют тон, организация информации и описания
инструментов? В этом эксперименте они сопоставляются на задачах
обслуживания клиентов авиакомпании из τ-bench. τ-bench создаёт
интерактивные сценарии с помощью симуляции пользователя,
бизнес-инструментов и правил, а затем проверяет целевое состояние после
завершения задачи; авиация и розничная торговля входят в число его
предметных областей\footnote{Yao et al., τ-bench: A Benchmark for
  Tool-Agent-User Interaction in Real-World Domains, 2024.}.

В этом запуске Kimi K3 использовалась как исполняющая модель и модель
симуляции пользователя. Один и тот же набор из десяти задач обслуживания
клиентов авиакомпании выполнялся в шести конфигурациях. Базовая
конфигурация сохраняла структурированные правила, полные описания
инструментов и профессиональный нейтральный тон. В остальных
конфигурациях по отдельности изменялись тон, организация правил и
наличие описаний инструментов, а ещё в одной группе эти изменения
сочетались.

В \textbf{сравнении тона} правила задачи сохранялись, но менялся способ
выражения. Стиль Трампа использовал преувеличенные и уверенные
формулировки, например: «Я забронирую тебе лучший рейс в истории».
Непринуждённый стиль использовал свободный тон и эмодзи.

В \textbf{сравнении организации информации} содержание правил
сохранялось, но заголовки и порядок перемешивались.

В \textbf{сравнении описаний инструментов} сохранялись названия,
сигнатуры функций и структура параметров, но удалялся описательный
текст. Название и schema инструмента также передают информацию, поэтому
результат после удаления описания зависит от того, достаточно ли
оставшихся сведений для текущей задачи.

Из десяти задач базовая конфигурация выполнила семь, конфигурации с
тоном Трампа и непринуждённым тоном --- соответственно шесть и девять, с
перемешанной организацией --- восемь, без описаний --- девять, а с
сочетанием изменений --- восемь.
\end{quote}

\hypertarget{ux438ux43dux44aux435ux43aux446ux438ux44f-ux432-ux43fux440ux43eux43cux43fux442-ux43eux441ux43dux43eux432ux43dux430ux44f-ux443ux433ux440ux43eux437ux430-ux431ux435ux437ux43eux43fux430ux441ux43dux43eux441ux442ux438-ux43aux43eux43dux442ux435ux43aux441ux442ux430}{%
\subsection{Инъекция в промпт: основная угроза безопасности
контекста}\label{ux438ux43dux44aux435ux43aux446ux438ux44f-ux432-ux43fux440ux43eux43cux43fux442-ux43eux441ux43dux43eux432ux43dux430ux44f-ux443ux433ux440ux43eux437ux430-ux431ux435ux437ux43eux43fux430ux441ux43dux43eux441ux442ux438-ux43aux43eux43dux442ux435ux43aux441ux442ux430}}

Для выполнения задач агенту необходимо читать веб-страницы, электронные
письма, документы и результаты инструментов. В этом содержимом могут
находиться инструкции, пытающиеся изменить цель задачи. \textbf{Инъекция
в промпт} (Prompt Injection) использует тот факт, что инструкции и
обрабатываемый материал вместе попадают в контекст модели, и побуждает
её принять содержимое из источника с низким уровнем доверия за основание
для действий. Например, пользователь просит кратко изложить
веб-страницу, а на ней агенту предлагается выполнить не связанную с
изложением операцию. Система должна сохранить цель пользовательской
задачи, а требования с веб-страницы рассматривать как анализируемое
содержимое.

Вызовы инструментов переносят такое отклонение во внешнюю среду.
Ошибочное краткое изложение может ввести пользователя в заблуждение, а
несанкционированное изменение файлов, отправка писем или раскрытие
данных непосредственно изменяют состояние системы. При анализе риска
можно последовательно проверить цепочку «точка поступления содержимого →
организация контекста → решение о вызове → фактическое выполнение»:
через какой инструмент поступил внешний контент, сохранилась ли
информация об источнике, какую операцию запросила модель и разрешает ли
её уровень выполнения\footnote{Anthropic, How we contain Claude across
  products.}.

При организации контекста прежде всего необходимо сохранять границы
источников и полномочий:

\begin{itemize}
\tightlist
\item
  \textbf{Маркировка источника}. При добавлении внешних материалов в
  контекст сохраняйте сведения о том, откуда получены веб-страница,
  письмо или документ, и заключайте содержимое в чётко обозначенные
  границы. Например,
  \texttt{\textless{}external\_content\ source="webpage"\textgreater{}...\textless{}/external\_content\textgreater{}}
  указывает, что внутри находятся материалы веб-страницы. Маркеры и их
  атрибуты создаются Harness, а внешний текст помещается в
  контролируемую область данных.
\item
  \textbf{Структурированные роли}. Системные правила, пользовательские
  запросы и результаты инструментов передаются отдельно в соответствии с
  соглашениями интерфейса. Результат инструмента связывается с ID
  вызова, что позволяет проследить его происхождение. При суммаризации
  или передаче другому агенту эту связь с источником также необходимо
  сохранять.
\item
  \textbf{Проверка входных данных}. Следует обнаруживать инструктивное
  содержимое, отклоняющееся от задачи, и удалять посторонние данные,
  попавшие в ходе разбора. Результаты проверки могут использоваться
  моделью и уровнем выполнения, однако решения об авторизации
  по-прежнему принимаются на основании запроса пользователя и системных
  полномочий.
\end{itemize}

Источниками необходимо управлять и для Skills, и для состояния среды
выполнения. При установке стороннего Skill следует проверить его
инструкции, документы, на которые он ссылается, и входящие в комплект
скрипты, чтобы убедиться, что требуемые им операции соответствуют
назначению. Ход выполнения задачи, права и счётчики в строке состояния
должны поддерживаться средой выполнения; факты, извлечённые с
веб-страницы, должны сохранять метку внешнего источника. Тогда после
суммаризации, сохранения или повторной загрузки информации система
по-прежнему сможет распознать исходные границы доверия.

Уровень выполнения обеспечивает соблюдение этих границ: область доступа
инструментов ограничивается правами, операции с файлами и процессами ---
песочницей, а высокорисковые действия подтверждаются согласно правилам.
Даже если модель предложит вызов с превышением полномочий, Harness и
среда должны заблокировать его до выполнения. Эти механизмы будут
подробно рассмотрены в четвёртой и пятой главах, а управление
источниками извлечённого содержимого и долговременной памяти --- в
третьей.

\begin{quote}
\textbf{Эксперимент 2-5 ★★: эксперимент по атаке и защите от инъекций в
промпт}

В этом эксперименте проверяется, как система сохраняет цель задачи и
полномочия на выполнение операций после попадания внешнего содержимого в
контекст. Эксперимент проводится в изолированном рабочем пространстве с
синтетическими данными, а чтение веб-страниц, запись файлов и отправка
электронной почты фиксируются в доступных для проверки журналах.

Три категории сценариев охватывают прямые запросы, содержимое
веб-страниц и межсессионную память. Сценарий прямого запроса проверяет,
соблюдает ли модель правила вывода синтетической защищённой информации.
Сценарий с веб-страницей проверяет, приводит ли задача суммаризации к
несанкционированной записи. Сценарий с памятью проверяет, влияет ли
внешний материал на операции в новом сеансе после его преобразования в
запись памяти.

В четырёх конфигурациях защита добавляется последовательно. Базовая
конфигурация содержит основные правила защиты информации и подтверждения
операций. Во второй добавляются указания о границах источников и
инструкций. В третьей дополнительно маркируются источники внешнего
содержимого. В четвёртой добавляется проверка полномочий в среде
выполнения для записи файлов и отправки электронной почты.

В этом запуске Kimi K3 были охвачены три категории сценариев и четыре
конфигурации; каждая комбинация повторялась пять раз --- всего
шестьдесят запусков. Ни в одном из них заданная цель атаки достигнута не
была.
\end{quote}

\hypertarget{ux434ux438ux43dux430ux43cux438ux447ux435ux441ux43aux438ux435-ux43fux440ux43eux43cux43fux442ux44b-ux438-agent-skills}{%
\section{Динамические промпты и Agent
Skills}\label{ux434ux438ux43dux430ux43cux438ux447ux435ux441ux43aux438ux435-ux43fux440ux43eux43cux43fux442ux44b-ux438-agent-skills}}

Если AI-агент поддерживает одновременно возврат средств, изменение кода
и упорядочивание документов, а к каждому запросу прикладываются три
полных руководства, то даже при оформлении одного возврата контекстное
окно будет занято нерелевантными стандартами кода и правилами
форматирования. Какие инструкции следует передавать каждый раз, а какие
можно загружать лишь после того, как модель определит задачу? Этот
вопрос приводит нас к динамическим промптам и Skills.

При оформлении возврата модель может сначала найти в каталоге
руководство по возвратам, прочитать его и только затем решить, нужно ли
обращаться к упомянутой в нём специальной таблице тарифов. На
\crossreflink{fig:2-11}{рисунке 2-11} тот же подход применяется к
созданию презентации: модель выбирает PPTX в каталоге и, прочитав
процесс создания новой презентации, обращается к руководству по
преобразованию HTML. Если требуется отредактировать существующую
презентацию, она может пойти другим путём и обратиться к справочнику по
XML.

\begin{figure}
\hypertarget{fig:2-11}{%
\centering
\includegraphics{images/revised/fig2-11.pdf}
\caption{Рисунок 2-11. Механизм постепенного раскрытия
Skills}\label{fig:2-11}
}
\end{figure}

\hypertarget{skills-ux43aux43eux43cux43fux43eux43dux443ux435ux43cux44bux435-ux435ux434ux438ux43dux438ux446ux44b-ux43fux440ux435ux434ux43cux435ux442ux43dux44bux445-ux43aux43eux43cux43fux435ux442ux435ux43dux446ux438ux439}{%
\subsection{Skills: компонуемые единицы предметных
компетенций}\label{skills-ux43aux43eux43cux43fux43eux43dux443ux435ux43cux44bux435-ux435ux434ux438ux43dux438ux446ux44b-ux43fux440ux435ux434ux43cux435ux442ux43dux44bux445-ux43aux43eux43cux43fux435ux442ux435ux43dux446ux438ux439}}

\textbf{Agent Skills} объединяют методы выполнения определённого класса
задач, справочные материалы, скрипты и шаблоны в многократно
используемые пакеты. Новому сотруднику необязательно в первый рабочий
день заучивать все руководства компании; модель также может подробно
прочитать нужный раздел только при выполнении соответствующей задачи.
Такой способ загрузки называется \textbf{постепенным раскрытием
(Progressive Disclosure)}\footnote{Anthropic, Equipping Agents for the
  Real World with Agent Skills, 2025; Agent Skills, How to add skills
  support to your agent.}.

Вернёмся к десяти товарам из \crossreflink{chap:1}{главы 1} и рассмотрим
на их примере, какую информацию стоит оформлять как Skill. Предположим,
пользователь хочет сравнить итоговую стоимость с доставкой. Для этого
достаточно добавить в текущий контекст сведения о товарах и адресе
доставки. Если команда регулярно сравнивает цены у разных продавцов,
можно создать Skill для сравнения цен, описав в нём, как сопоставлять
характеристики, различать цену за единицу и комплект и отмечать
неизвестную стоимость доставки. Когда в следующий раз пользователь
захочет лишь посмотреть доступные варианты, каталог останется видимым, а
полный процесс проверки цен можно будет загрузить только при
необходимости сравнить цены.

\textbf{Первый уровень --- метаданные.} Каждый Skill описывается в файле
\texttt{SKILL.md}, в начале которого находится YAML frontmatter,
отделённый символами \texttt{-\/-\/-} и содержащий \texttt{name} и
\texttt{description}. Сначала среда выполнения формирует из названий и
описаний каталог, чтобы модель могла определить, какие Skills требуются
для текущей задачи. Порядком и представлением каталога управляет
Harness, а полный процесс загружается после выбора.

\textbf{Описание Skill прежде всего следует формулировать как условие
маршрутизации.} Модель должна решить, использовать ли Skill, ещё до
загрузки его основного текста, поэтому в \texttt{description} сначала
следует указать, когда он применим и когда его нужно пропустить, а затем
кратко описать его возможности. Например, ``help with backend'' имеет
слишком широкую область действия и может соответствовать почти любой
задаче, связанной с бэкендом; формулировка «написание интеграционных
тестов для существующего HTTP-интерфейса с проверкой аутентификации,
валидации параметров и ответов с ошибками» задаёт чёткие границы. Для
задач, которые легко перепутать, можно также добавить условия
исключения, например: «не использовать, если обсуждается только проект
интерфейса».

\textbf{Второй уровень --- основной процесс.} После выбора Skill среда
выполнения передаёт полные инструкции в контекст. Если пользователь явно
вводит команду вроде \texttt{/pptx}, клиент может напрямую найти и
загрузить Skill. Если выбор делает сама модель, она получает содержимое
через инструмент Skill или инструмент чтения файлов. Первый вариант
избавляет модель от необходимости инициировать вызов активации, а во
втором выбор и загрузка остаются в траектории выполнения.

Например, основные инструкции PPTX Skill, использованного в
экспериментах этой книги, описывают несколько распространённых путей:
извлечение текста презентации с помощью markitdown от Microsoft,
распаковку PPTX для просмотра лежащего в его основе XML и выбор процесса
создания или редактирования презентации. Прочитав инструкции этого
уровня, модель может в зависимости от текущей задачи решить, к какому
материалу обратиться дальше и какой тип скриптов вызвать\footnote{Anthropic,
  PPTX Skill. В экспериментах этой книги зафиксирована версия,
  содержащая процесс преобразования HTML; конкретная версия и способ
  запуска приведены в README сопутствующих экспериментов.}.

\textbf{Третий уровень --- подробные материалы.} Сложные операции
распределены по специализированным документам, а основной файл
предоставляет точки входа. В использованной для этой книги версии
\texttt{html2pptx.md} объясняет, как создавать презентации с помощью
HTML, а другие справочные документы дополняют сведения о форматировании
и редактировании. Агент продолжает читать нужные материалы в
соответствии с задачей и выполняет её с помощью входящих в пакет
скриптов, ресурсов и шаблонов.

\hypertarget{ux43aux430ux43a-ux43dux430ux43fux438ux441ux430ux442ux44c-ux43fux440ux438ux433ux43eux434ux43dux44bux439-ux43a-ux438ux441ux43fux43eux43bux44cux437ux43eux432ux430ux43dux438ux44e-skill}{%
\subsection{Как написать пригодный к использованию
Skill}\label{ux43aux430ux43a-ux43dux430ux43fux438ux441ux430ux442ux44c-ux43fux440ux438ux433ux43eux434ux43dux44bux439-ux43a-ux438ux441ux43fux43eux43bux44cux437ux43eux432ux430ux43dux438ux44e-skill}}

Пригодный к использованию Skill должен отвечать на четыре вопроса: для
каких задач он применим, в какой последовательности их выполнять, в
каких случаях необходимо подтверждение и как определить, что работа
завершена. Автору нужно преобразовать свой опыт в исполняемые шаги и
снабдить ключевые решения примерами.

При создании Skill для письма можно опереться на подход Баоюя к
улучшению текстов AI на основе личных особенностей изложения\footnote{Баоюй,
  «Хватит устранять привкус AI с помощью промптов --- сам подход
  неверен», 2026-02-14.}, а затем разделить содержимое на четыре части:

\begin{itemize}
\tightlist
\item
  \textbf{Роль и читатели}: указать целевую аудиторию, типы задач и
  требования к готовому тексту.
\item
  \textbf{Основные принципы}: сначала сформулировать от трёх до пяти
  часто применяемых правил принятия решений, сопроводив каждое
  конкретными примерами до и после редактирования.
\item
  \textbf{Распространённые ошибки и границы}: зафиксировать часто
  возникающие проблемы изложения, действия с превышением полномочий и
  исключения.
\item
  \textbf{Справочные материалы}: включить глоссарий, шаблоны, образцы и
  специализированные инструкции для использования во время работы.
\end{itemize}

Правила можно формулировать по схеме «область действия, действие,
исключение, способ проверки». Например, правило сокращения длинных
предложений должно указывать, к каким абзацам оно применяется, как
разделять предложения и как после разделения проверить, что условия и
выводы остались полными.

Для начала выберите от трёх до пяти собственных оригинальных статей,
которыми вы довольны, и попросите агента обобщить лексику, структуру
предложений и абзацев, а также тон, сформировав первоначальный вариант
примерно из двадцати строк. Затем примените его к новой теме, после чего
автор должен отредактировать текст предложение за предложением. Различия
между версиями до и после редактирования дадут более конкретную обратную
связь: какие слова избыточны, какие предложения нужно разделить, где не
хватает фактов, а где следует вернуть технические подробности.
Преобразуйте повторяющиеся исправления в правила, сохранив
соответствующие примеры и области применения.

Skill также может включать исполняемые инструменты и шаблоны. Например,
PPTX Skill объединяет инструкции по работе, скрипты разбора и генерации
и ресурсы макетов, позволяя модели вызывать соответствующие ресурсы
согласно инструкциям. Каждый Skill можно независимо тестировать,
обновлять и версионировать, а участники команды могут постоянно
дополнять общий процесс своим опытом. \textbf{Долгосрочная ценность
Skills заключается в том, что они позволяют распределённо накапливать,
независимо публиковать и повторно использовать предметный опыт.}
Расширяя возможности агента, команды могут поддерживать Skills в
собственных областях компетенции, постепенно формируя экосистему
пакетов, подобную pip и npm. Эти Skills могут использовать общий
механизм обнаружения, загрузки и выполнения, предоставляемый основной
системой.

\textbf{При проектировании способа взаимодействия следует отдавать
предпочтение режимам, на которых модель уже обучалась и к которым была
адаптирована.} Распределение ролей, формат результатов инструментов и
способ загрузки Skill влияют на то, как модель продолжает выполнение
задачи. Стандартные процессы поставщика подходят в качестве отправной
точки и снижают необходимость заново выводить смысл незнакомого формата.

\hypertarget{ux43cux435ux441ux442ux43e-skills-ux432-ux43aux43eux43dux442ux435ux43aux441ux442ux435}{%
\subsection{Место Skills в
контексте}\label{ux43cux435ux441ux442ux43e-skills-ux432-ux43aux43eux43dux442ux435ux43aux441ux442ux435}}

Затраты контекста на Skills состоят из двух частей: каталога и основного
текста. Каталог позволяет модели обнаруживать доступные возможности, а
основной текст после выбора предоставляет процесс работы. Спецификация
требует загрузки по мере необходимости, тогда как конкретные роли
сообщений, способ обёртки и обновления определяются реализацией Harness.

Claude Code добавляет инструкции Skill в сеанс как сообщение
\texttt{user} в точке вызова\footnote{Claude Code, How Claude Code uses
  prompt caching, Invoking skills and commands; Agent Skills, How to add
  skills support to your agent, User-explicit activation.}. На
\crossreflink{fig:2-12}{рисунке 2-12} показан один из вариантов
расположения сообщений при самостоятельной активации моделью: сначала
отправляется вызов \texttt{Skill(skill:\ "pptx")}, затем поступает ответ
инструмента, после чего считывается добавленный основной текст Skill.
При использовании специализированного инструмента активации ответ
инструмента может подтверждать вызов, а основной текст отдельно
внедряется Harness. При использовании обычного инструмента чтения файлов
основной текст может непосредственно входить в результат инструмента.
Если пользователь явно вводит \texttt{/pptx}, клиент может сразу
развернуть инструкции, избавив модель от необходимости инициировать
вызов активации.

\begin{figure}
\hypertarget{fig:2-12}{%
\centering
\includegraphics[width=\textwidth,height=0.55\textheight]{images/revised/fig2-12.pdf}
\caption{Рисунок 2-12. Полная структура траектории агента при включённых
Skills}\label{fig:2-12}
}
\end{figure}

В реализации Skills для Codex каталог возможностей и выбранные
инструкции также организованы отдельно. В общедоступном исходном коде
часть каталога входит в контекст \texttt{developer}, а другие части
предоставляются вместе с текущим ходом и средой выполнения; явно
выбранный Skill добавляется через фрагмент пользовательского контекста.
Источник каталога и доступная среда выполнения влияют на конкретное
расположение элементов\footnote{OpenAI, Build skills; реализация Skills
  в общедоступном репозитории OpenAI Codex.}.

На \crossreflink{fig:2-13}{рисунке 2-13} показан рост кэша при загрузке
по мере необходимости, начиная с запроса, для которого ещё не создан
повторно используемый кэш. Первое включение каталога в запрос требует
обработки, а первая загрузка основного текста также увеличивает объём
вычислений. Если последующие запросы сохраняют тот же префикс, эти
данные можно будет считывать из кэша. Добавление нового содержимого в
точке вызова позволяет сохранить предыдущие сообщения неизменными. Уже
загруженные инструкции остаются в истории, а при необходимости
обновляются в соответствии со стратегией управления контекстом среды
выполнения.

\begin{figure}
\hypertarget{fig:2-13}{%
\centering
\includegraphics{images/revised/fig2-13.pdf}
\caption{Рисунок 2-13. Эволюция KV-кэша по мере роста траектории
агента}\label{fig:2-13}
}
\end{figure}

\hypertarget{ux441ux432ux44fux437ux44c-skills-ux441-ux438ux43dux441ux442ux440ux443ux43cux435ux43dux442ux430ux43cux438}{%
\subsection{Связь Skills с
инструментами}\label{ux441ux432ux44fux437ux44c-skills-ux441-ux438ux43dux441ux442ux440ux443ux43cux435ux43dux442ux430ux43cux438}}

Skill предоставляет знания о задаче и процесс работы, а инструменты
отвечают за чтение, выполнение и получение результатов. Семь базовых
инструментов, представленных в \crossreflink{chap:5}{главе 5}, можно
сочетать с различными Skills: один и тот же инструмент чтения файлов
способен прочитать как руководство по созданию презентаций, так и
стандарты кода, а универсальный исполнитель --- запустить
соответствующие скрипты. Благодаря этому профессиональные возможности
расширяются посредством процессов и ресурсов, тогда как базовый набор
инструментов остаётся относительно стабильным.

Специализированные инструменты также можно использовать совместно со
Skill. Рассмотренный выше поиск инструментов предоставляет их
определения по мере необходимости, а Skill объясняет, когда их вызывать,
как сочетать и как проверять результаты. Оба механизма должны
ограничивать объём постоянно присутствующей информации и отводить
подходящее место новому загружаемому содержимому. В
\crossreflink{chap:4}{главе 4} рассматривается выбор между универсальным
исполнителем и специализированными инструментами, а в
\crossreflink{chap:9}{главе 9} --- вопрос о том, в какой форме лучше
сохранять опыт: как знания, инструкции, программу или параметры модели.

\begin{quote}
\textbf{Эксперимент 2-6 ★★: создание презентации из научной статьи с
помощью Agent Skills}

В этом эксперименте на материале статьи «Attention Is All You Need»
демонстрируется, как агент загружает по мере необходимости
специализированный процесс создания презентации. Задача состоит в том,
чтобы создать из PDF-файла статьи презентацию на десять--пятнадцать
слайдов, охватывающую постановку проблемы, метод, ключевые результаты и
выводы, и использовать для пояснений иллюстрации и таблицы из статьи.

В этом запуске использовались Kimi Code и официальный PPTX Skill от
Anthropic. Сначала агент нашёл PPTX Skill в каталоге и прочитал
инструкции по работе, затем обратился к руководству по преобразованию
HTML, вызвал входящие в комплект инструменты генерации и в завершение
проверил макет с помощью миниатюр и постраничного рендеринга. Все
материалы, прочитанные на каждом шаге, служили следующей операции, а
полный процесс создания раскрывался постепенно по мере выполнения
задачи.

Готовая презентация состояла из тринадцати слайдов и включала схему
архитектуры, схему механизма внимания, таблицу результатов экспериментов
и визуализацию внимания из статьи.

Проверка готового результата также входит в процесс Skill. Во время
этого запуска были обнаружены отсутствующие нижние индексы, некорректные
маркеры списков, неверные пути к иллюстрациям и проблемы с точностью
числовых подписей. Агент последовательно исправил их по результатам
рендеринга и проверки структуры. Итоговый файл прошёл проверку структуры
и был повторно открыт и отрендерен с помощью LibreOffice.
\end{quote}

\begin{quote}
\textbf{Эксперимент 2-7 ★★: создание Skill для письма «без привкуса AI»
на основе личных образцов}

Это упражнение преобразует личный опыт письма в многократно используемые
инструкции. Подготовьте от трёх до пяти оригинальных статей, попросите
агента обобщить особенности изложения и создать первоначальный вариант
Skill с условиями активации, тремя--пятью принципами и соответствующими
примерами.

Выберите новую тему для пробного текста, после чего автор должен
отредактировать его предложение за предложением. Сопоставьте две версии
и найдите повторяющиеся проблемы: например, расплывчатые формулировки,
чрезмерно длинные предложения, пропущенные причинно-следственные связи
или удалённые конкретные примеры, необходимые для понимания.
Повторяющиеся способы исправления можно добавить в Skill.
\end{quote}

\hypertarget{ux43fux430ux43dux435ux43bux44c-ux441ux43eux441ux442ux43eux44fux43dux438ux44f-ai-ux430ux433ux435ux43dux442ux430-ux443ux43bux443ux447ux448ux435ux43dux438ux435-ux443ux43fux440ux430ux432ux43bux435ux43dux438ux44f-ux442ux440ux430ux435ux43aux442ux43eux440ux438ux435ux439-ux430ux433ux435ux43dux442ux430-ux441-ux43fux43eux43cux43eux449ux44cux44e-ux43cux435ux442ux430ux434ux430ux43dux43dux44bux445}{%
\section{Панель состояния AI-агента: улучшение управления траекторией
агента с помощью
метаданных}\label{ux43fux430ux43dux435ux43bux44c-ux441ux43eux441ux442ux43eux44fux43dux438ux44f-ai-ux430ux433ux435ux43dux442ux430-ux443ux43bux443ux447ux448ux435ux43dux438ux435-ux443ux43fux440ux430ux432ux43bux435ux43dux438ux44f-ux442ux440ux430ux435ux43aux442ux43eux440ux438ux435ux439-ux430ux433ux435ux43dux442ux430-ux441-ux43fux43eux43cux43eux449ux44cux44e-ux43cux435ux442ux430ux434ux430ux43dux43dux44bux445}}

Пользователь торопит: «Можно позвонить ещё раз?» Слева на
\crossreflink{fig:2-14}{рисунке 2-14} зафиксированы три звонка. При
первом ответа ждали 45 минут, но так и не дождались; лишь во время
третьего удалось договориться о цене 59 долларов в месяц. По правилам
этого примера каждый звонок расходует одну попытку. Хотя соединиться
удалось только дважды, лимит из трёх попыток уже исчерпан. Harness может
собрать такие счётчики, ход выполнения и вопросы, ожидающие
подтверждения, в обновляемую вместе с сообщениями сводку ---
\textbf{панель состояния агента (Agent Status Bar)}.

\begin{figure}
\hypertarget{fig:2-14}{%
\centering
\includegraphics{images/revised/fig2-14.pdf}
\caption{Рисунок 2-14. Архитектура панели состояния
агента}\label{fig:2-14}
}
\end{figure}

Смартфон показывает время, заряд аккумулятора и уровень сигнала в
верхней части экрана, позволяя человеку одним взглядом оценить текущее
состояние. Сводка справа на рисунке выполняет похожую функцию: увидев
«лимит достигнут» и «ожидается подтверждение пользователя», помощник
может сообщить пользователю согласованную цену и попросить решить,
принимать ли предложение. Harness может прикреплять сводку к результату
инструмента или добавлять её отдельным сообщением о состоянии.

\hypertarget{ux43eux442-ux438ux441ux445ux43eux434ux43dux44bux445-ux437ux430ux43fux438ux441ux435ux439-ux43a-ux44fux432ux43dux43eux43cux443-ux441ux43eux441ux442ux43eux44fux43dux438ux44e}{%
\subsection{От исходных записей к явному
состоянию}\label{ux43eux442-ux438ux441ux445ux43eux434ux43dux44bux445-ux437ux430ux43fux438ux441ux435ux439-ux43a-ux44fux432ux43dux43eux43cux443-ux441ux43eux441ux442ux43eux44fux43dux438ux44e}}

Рассмотрим задачу телефонного обслуживания: система разрешает звонить
каждой компании не более трёх раз в рамках одной задачи. Между тремя
разговорами могут выполняться поиск, чтение материалов и обработка
дополнительных вопросов пользователя. Чтобы решить, можно ли звонить
снова, модели необходимо найти в истории все соответствующие записи,
разделить их по компаниям и подсчитать общее количество.

Среда выполнения может обновлять счётчик после каждого вызова и сразу
возвращать: «Это третий звонок в Xfinity; лимит для этой задачи
достигнут». При следующем решении модель сможет использовать этот
результат. Программа ведёт счёт вызовов, а модель с учётом хода
выполнения решает, что делать дальше: ждать, искать другие каналы или
объяснить пользователю текущую ситуацию. Исполнительный слой
одновременно проверяет последующие запросы на звонок по этому счётчику.

Инженерная ценность такого подхода заключается в
\textbf{заблаговременной подготовке состояния, которое многократно
потребуется для последующих решений}. Полный журнал звонков по-прежнему
хранит содержание разговоров и историю обработки, а панель состояния
предоставляет часто используемые при принятии решений статистику и
прогресс. \textbf{Многократно используемую статистику следует в первую
очередь поддерживать инкрементально с помощью кода.} Найти одну запись в
контексте и вычислить итог по всем записям --- разные задачи. Сохранение
исходной траектории само по себе не создаёт точную и всегда доступную
статистическую сводку. Если каждый раз заставлять модель заново
просматривать историю и подсчитывать записи, она будет расходовать
токены инференса, а вероятность пропусков и двойного учёта возрастёт.

Автор сравнивает контекст с поисковой системой, которой необходим
дополнительный «слой извлечения сути»: исходные записи удобны для
ретроспективного анализа, а извлечённое состояние --- для
непосредственного использования. Счётчики, лимиты и проверенные
состояния завершения задачи целесообразно поддерживать кодом.
Информацию, требующую семантической оценки, можно сначала извлекать по
отдельным записям, а затем агрегировать и проверять. Так при каждом
решении модель одновременно видит текущее состояние и необходимые
доказательства.

Длинная траектория также затрудняет поиск ключевой информации. Если
сосредоточить последние цели, ограничения и прогресс ближе к концу,
модели будет проще найти сведения, необходимые для текущего решения.

\begin{quote}
\textbf{Эксперимент 2-8 ★★: наблюдение за эффектом панели состояния
агента посредством визуализации внимания}

В эксперименте используется Qwen3-0.6B и моделируется траектория
возврата средств через службу поддержки: агент уже трижды звонил в
Xfinity, между звонками выполнял поиск в интернете, после чего
пользователь спросил: «Можно ещё раз позвонить и поторопить их?»
Системные правила разрешают звонить каждой компании не более трёх раз в
рамках одной задачи.

Обе группы получают одинаковую историю и дополнительный вопрос
пользователя. Группа A использует исходную траекторию без изменений;
группе B в конец добавляется следующее состояние:

\begin{Shaded}
\begin{Highlighting}[]
\NormalTok{\textless{}}\KeywordTok{agent\_status}\NormalTok{\textgreater{}}
\NormalTok{Current State:}
\NormalTok{{-} Tool call summary: \textquotesingle{}phone\_call\textquotesingle{} has been invoked 3 times (Xfinity: 3 times)}
\NormalTok{{-} Constraint check: Maximum calls to Xfinity reached (3/3)}
\NormalTok{\textless{}/}\KeywordTok{agent\_status}\NormalTok{\textgreater{}}
\end{Highlighting}
\end{Shaded}

Для каждой группы выполняется по три запуска. В группе A модель дважды
явно отказалась звонить снова, а один ответ был отнесён к другим
результатам. В группе B модель явно отказалась звонить во всех трёх
запусках. Ни в одной группе не было ответов, классифицированных как
нарушающее правила продолжение звонков. Чтобы воспроизвести это
сравнение, перейдите в \texttt{chapter2/attention\_visualization} и
выполните
\texttt{python\ run\_status\_bar\_experiment.py\ -\/-output\ runs/status-bar-reproduction}.
Входные данные и результаты каждого запуска сохранены в файле
\texttt{runs/exp2-7-qwen3-0.6b-20260730-v2/comparison.json} каталога
эксперимента.

Диаграмма внимания показывает, как веса внимания области генерации
распределяются по истории звонков, результатам поиска, дополнительному
вопросу пользователя и панели состояния.
\end{quote}

Исследование ContextDistill-Bench, проведённое автором совместно с Noah
Shi, дополнительно рассматривает преимущества заранее извлечённого
состояния\footnote{Li, Bojie and Noah Shi. \emph{Distill, Don't
  Retrieve: Inference-Time Context Distillation for LLM Agent
  Reasoning.} 2026.}. В охваченных этим бенчмарком запросах состояния
небольшие модели также достигают высокой точности благодаря явному
состоянию, а объём рассуждений передовых моделей, задержка и стоимость
запросов могут сократиться примерно на порядок. Например, стоимость
тысячи запросов к Claude 4 Sonnet снизилась приблизительно с 35,46 до
3,51 доллара. \textbf{Такой подход обновляет состояние в момент
возникновения событий и устраняет необходимость при каждом запросе
заново просматривать историю.} После фиксации полей сводки запросу
достаточно читать только эти поля, даже если история продолжает расти.

\hypertarget{ux441ux43eux441ux442ux430ux432-ux43fux430ux43dux435ux43bux438-ux441ux43eux441ux442ux43eux44fux43dux438ux44f-ux430ux433ux435ux43dux442ux430}{%
\subsection{Состав панели состояния
агента}\label{ux441ux43eux441ux442ux430ux432-ux43fux430ux43dux435ux43bux438-ux441ux43eux441ux442ux43eux44fux43dux438ux44f-ux430ux433ux435ux43dux442ux430}}

Панель состояния обычно содержит три типа информации, отвечающие на
вопросы «что нужно сделать», «что только что произошло» и «каково
текущее состояние».

\textbf{Планирование задачи} фиксирует цель, этапы и состояние их
выполнения. Сложную задачу можно разбить на список TODO, сохраняя для
каждого пункта состояние и необходимые зависимости. Работая над текущим
этапом, модель по-прежнему видит исходные требования, ключевые
ограничения и дальнейшие действия; после завершения этапа список
обновляется.

\textbf{Вспомогательная информация о событиях (Side-channel
Information)} дополняет контекст, в котором произошло событие: например,
временем, географическим положением и интервалом с момента последнего
ответа. Эти метаданные сохраняются вместе с событием и помогают модели
понимать последовательность и обстоятельства. Одна и та же фраза «ответа
пока нет», сказанная через две минуты или через два дня после подачи
заявки, может требовать разных действий.

\textbf{Сводка наблюдений за текущим состоянием среды} предоставляет
текущее время, рабочий каталог, количество вызовов инструментов и
предупреждения об аномалиях. Она выполняет функцию, похожую на
отображение текущего каталога или состояния выполнения в CLI и GUI: даёт
явные основания для следующего действия. Например, предупреждение о
повторных вызовах может побудить модель проверить причину неудачи, а
сведения об оставшемся бюджете помогают решить, продолжать ли поиск.

Метаданные события после добавления обычно остаются неизменными, тогда
как ход задачи и состояние среды продолжают меняться. Эти виды
информации следует записывать в историю в соответствии с их способами
обновления, которые далее рассматриваются на примере структуры
сообщений.

\hypertarget{ux43aux43eux43dux43aux440ux435ux442ux43dux43eux435-ux43fux43eux43bux43eux436ux435ux43dux438ux435-ux43fux430ux43dux435ux43bux438-ux441ux43eux441ux442ux43eux44fux43dux438ux44f-ux430ux433ux435ux43dux442ux430-ux432-ux43aux43eux43dux442ux435ux43aux441ux442ux435}{%
\subsection{Конкретное положение панели состояния агента в
контексте}\label{ux43aux43eux43dux43aux440ux435ux442ux43dux43eux435-ux43fux43eux43bux43eux436ux435ux43dux438ux435-ux43fux430ux43dux435ux43bux438-ux441ux43eux441ux442ux43eux44fux43dux438ux44f-ux430ux433ux435ux43dux442ux430-ux432-ux43aux43eux43dux442ux435ux43aux441ux442ux435}}

В примере этого раздела Harness добавляет в конец сообщений сообщение
\texttt{user}, передающее сводку состояния. Здесь \texttt{user} --- роль
сообщения в выбранном API; содержимое состояния генерируется средой
выполнения, а его назначение обозначается тегами. При использовании
других интерфейсов можно применять поддерживаемые ими блоки контекста
или передавать состояние вместе с результатом соответствующего
инструмента. Если используется описанный выше шаблон, который обрезает
исторические рассуждения по реальным пользовательским запросам,
состояние среды выполнения можно передавать вместе с результатами
инструментов, чтобы шаблон продолжал сохранять сведения о рассуждениях
для текущей задачи.

Проследите по \crossreflink{fig:2-15}{рисунку 2-15} от исходного
поручения пользователя до раунда, в котором записано последнее
состояние. Сохранение исходного запроса позволяет после завершения
звонков одновременно проверить две вещи: чего хотел пользователь и до
какого этапа дошло выполнение.

\begin{figure}
\hypertarget{fig:2-15}{%
\centering
\includegraphics{images/revised/fig2-15.pdf}
\caption{Рисунок 2-15. Место вставки панели состояния агента в списке
сообщений API}\label{fig:2-15}
}
\end{figure}

Последовательность с рисунка можно представить следующим списком
сообщений. Для удобства сопоставления повторяющиеся раунды и часть полей
вызовов опущены.

\begin{Shaded}
\begin{Highlighting}[]
\NormalTok{messages: [}
\NormalTok{  \{ role: "system", content: "You are a customer service assistant..." \}}
\NormalTok{  \{ role: "user", content: "Help me cancel my Xfinity plan" \}}
\NormalTok{  \{ role: "assistant", tool\_calls: [\{ id: "call\_1", ... \}] \}}
\NormalTok{  \{ role: "tool", tool\_call\_id: "call\_1", content: "Call log..." \}}
\NormalTok{  \{ role: "assistant", tool\_calls: [\{ id: "call\_2", ... \}] \}}
\NormalTok{  \{ role: "tool", tool\_call\_id: "call\_2", content: "Call log..." \}}
\NormalTok{  ...остальные раунды...}
\NormalTok{  \{ role: "user", content: "Can you call them again to follow up?" \}}
\NormalTok{  \{ role: "user", content: "\textless{}agent\_status\textgreater{}}
\NormalTok{      Current State:}
\NormalTok{      {-} phone\_call invoked 3 times (Xfinity: 3/3 max)}
\NormalTok{      {-} Current time: 2025{-}09{-}14 10:30:45}
\NormalTok{      {-} TODO: [1] Cancel plan (in\_progress)}
\NormalTok{    \textless{}/agent\_status\textgreater{}" \}}
\NormalTok{]}
\end{Highlighting}
\end{Shaded}

Последнее сообщение объединяет количество вызовов, лимит, время и ход
выполнения задачи. Модель может напрямую использовать эти поля, а при
необходимости обратиться к конкретным записям по предшествующим
идентификаторам вызовов.

После добавления этого сообщения предшествующая последовательность
токенов остаётся неизменной и по-прежнему удовлетворяет базовым условиям
повторного использования префикса. Способ обработки старого состояния в
следующем раунде приходится выбирать с учётом компромисса между расходом
контекста и стоимостью кэширования.

\hypertarget{ux434ux432ux430-ux441ux43fux43eux441ux43eux431ux430-ux43eux431ux43dux43eux432ux43bux435ux43dux438ux44f-ux441ux43eux441ux442ux43eux44fux43dux438ux44f-ux438-ux438ux445-ux432ux43bux438ux44fux43dux438ux435-ux43dux430-ux43aux44dux448}{%
\subsection{Два способа обновления состояния и их влияние на
кэш}\label{ux434ux432ux430-ux441ux43fux43eux441ux43eux431ux430-ux43eux431ux43dux43eux432ux43bux435ux43dux438ux44f-ux441ux43eux441ux442ux43eux44fux43dux438ux44f-ux438-ux438ux445-ux432ux43bux438ux44fux43dux438ux435-ux43dux430-ux43aux44dux448}}

В следующем раунде может завершиться пункт TODO или увеличиться счётчик
вызовов инструмента. Harness должен обновить состояние; обычно применяют
один из двух способов.

\textbf{Замена в каждом раунде}: при формировании нового запроса
предыдущее состояние удаляется, а в конец добавляется актуальное. В
контексте остаётся только одна версия текущего состояния, удобная для
использования моделью. Например, предыдущий запрос состоял из «описания
задачи, состояния A», а новый --- из «описания задачи, результата
инструмента, состояния B». Описание задачи по-прежнему может быть взято
из кэша; результат инструмента отправляется впервые, поэтому его в любом
случае необходимо вычислить. Если после состояния A в предыдущем запросе
находились другие сообщения, удаление A лишит их прежнего совпадения с
префиксом. В текущем запросе модель видит только B. Если потребуется
просмотреть A, соответствующий снимок можно получить из журнала среды
выполнения.

\textbf{Постоянное добавление}: существующие записи состояния
сохраняются, а каждая новая версия добавляется в конец. Исторический
префикс остаётся стабильным, а изменения состояния удобно отслеживать.
Старые состояния постепенно занимают контекст, поэтому необходимо явно
указывать время или версию, чтобы модель определяла текущую ситуацию по
последней записи. Способ, которым Claude Code добавляет в сеанс
уведомления об изменении файлов, также отражает этот подход.

Если состояние невелико, а после него уже закэширован большой объём
результатов инструментов, сохранение старого состояния может оказаться
дешевле повторного вычисления длинного суффикса. Если же состояние
велико и часто обновляется, замена позволяет ограничить накопление
исторических состояний.

Граничное значение можно оценить с помощью приближённой модели. Пусть
каждое состояние содержит \(S\) токенов, средний размер закэшированного
суффикса после старого состояния, который при каждой замене необходимо
вычислить заново, равен \(R\) токенам --- без учёта содержимого, впервые
отправленного в текущем раунде, --- всего состояние внедряется \(N\)
раз, а стоимость кэшированного ввода составляет \(\alpha\) от стоимости
обычного ввода. Если не учитывать общие для двух подходов затраты, то
\(C_{\text{замена}} \approx (N-1)(1-\alpha)R\), а
\(C_{\text{добавление}} \approx \alpha S N(N-1)/2\). При \(N>1\)
оценочная стоимость постоянного добавления ниже, если
\(\alpha SN/2 < (1-\alpha)R\); в противном случае дешевле замена в
каждом раунде. Здесь стоимость обычного ввода принята за единицу: при
замене дополнительно оплачивается разница между чтением суффикса из кэша
и его повторной обработкой, а при добавлении --- повторное чтение
исторических состояний.

Например, состояние объёмом 100 токенов внедряется десять раз, а чтение
из кэша стоит одну десятую от обычного ввода. При добавлении последующие
запросы дополнительно прочитают соответственно 1, 2, \ldots, 9 старых
состояний, что эквивалентно стоимости 450 токенов обычного ввода. Если
замена состояния приводит к повторной обработке в среднем 1000
закэшированных токенов после него, девять замен обойдутся в эквивалент
8100 токенов обычного ввода, поэтому добавление будет дешевле. Если же
старое состояние и так находилось в самом конце кэша и после него не
было закэшированного содержимого, то \(R=0\), и непосредственная замена,
напротив, сэкономит эквивалент стоимости этих 450 токенов.

\begin{quote}
\textbf{Эксперимент 2-9 ★★: несколько полезных методов построения панели
состояния агента}

В эксперименте реализованы пять методов построения панели состояния,
которые можно включать отдельно или сочетать. Каждый метод предоставляет
определённый вид информации, помогающий модели выбрать следующее
действие.

\textbf{Отслеживание временно́й метки} записывает время события в
сообщениях пользователя и ответах инструментов, например
\texttt{{[}2025-09-14\ 10:30:45{]}}. Время добавляется в историю вместе
с событием и сохраняет исходное значение в последующих раундах, тогда
как текущее время обновляется с каждым новым сообщением о состоянии.
Эксперимент также поддерживает моделирование времени, позволяющее
проверить, различает ли модель «вчерашний файл» и «сегодняшнее
изменение» и правильно ли понимает длительность ожидания.

\textbf{Счётчик вызовов инструментов} хранит в словаре общее количество
вызовов каждого инструмента и добавляет в ответ отметку
\texttt{Tool\ call\ \#3\ for\ \textquotesingle{}read\_file\textquotesingle{}}.
С учётом содержимого ошибки модель может решить, продолжать ли попытки.
Например, после первой неудачи чтения она проверяет путь, а затем
выводит содержимое каталога; после нескольких неудач ищет альтернативный
файл или сообщает пользователю об отсутствии нужного файла. Если число
повторных попыток необходимо ограничить, среда выполнения также должна
принудительно применять лимит по счётчику.

\textbf{Управление списком TODO} предоставляет два инструмента:
\texttt{rewrite\_todo\_list} и \texttt{update\_todo\_status}. Каждый
пункт содержит уникальный идентификатор, описание, временную метку и
одно из четырёх состояний: \texttt{pending}, \texttt{in\_progress},
\texttt{completed} или \texttt{cancelled}. Список централизованно хранит
сведения о том, «что сделано и что ещё осталось», что напоминает подход
Manus, сохраняющий видимость цели посредством повторного представления
плана задачи.

\textbf{Подробные сведения об ошибках} включают тип и описание ошибки,
JSON с параметрами вызова, стек вызовов и рекомендации по исправлению.
Например, при \texttt{FileNotFoundError} модели можно предложить
проверить имя файла, текущий каталог и абсолютный путь. Параметры
помогают выявить ошибочные входные данные, стек вызовов --- найти место
сбоя, а рекомендации --- определить направление следующей проверки.
Перед передачей этих данных модели необходимо удалить учётные данные и
другие конфиденциальные поля, а также ограничить длину чрезмерно
подробных стеков.

\textbf{Восприятие состояния системы} предоставляет текущее время,
рабочий каталог, операционную систему, Shell и версию Python. Командный
инструмент в эксперименте распознаёт отдельные операции смены каталога,
обновляет рабочий каталог, поддерживаемый средой выполнения, и
использует его в последующих вызовах. Изменение каталога подпроцессом в
составной команде действует в пределах соответствующего процесса. Поля
среды также помогают модели выбирать доступные инструменты, например
использовать \texttt{apt} или \texttt{brew} с учётом фактической системы
и установленного программного обеспечения.

В этом сравнении на Kimi K3 для каждого типа задач использовалось по
пять примеров. Конфигурации с временными метками и с сочетанием методов
прошли все пять примеров, тогда как соответствующие контрольные группы
--- по три. При включении и отключении каждого из двух методов ---
счётчика вызовов инструментов и восприятия состояния системы --- были
пройдены все пять примеров. Конфигурация TODO не прошла ни одного
примера, а контрольная группа прошла один; конфигурация с подробными
ошибками прошла четыре примера, а контрольная группа --- пять.
\end{quote}

\textbf{Панель состояния следует по возможности поддерживать с помощью
кода.} Модель напрямую использует приведённые в ней статистику и
ограничения, поэтому ошибочное утверждение «выполнено три звонка» может
изменить последующие действия. Счётчики, время и переходы состояний
следует поручить программе. Если требуется семантическое понимание,
модель должна сначала извлечь сведения из каждой записи отдельно, после
чего код агрегирует и проверяет результаты, не заставляя модель
единовременно подсчитывать всю длинную историю. Между исходными
событиями, извлечёнными результатами и итоговым состоянием следует
сохранять соответствие, чтобы обнаруживать двойной учёт, пропуски и
смешение источников.

Необходимо заранее выбрать, какие сведения агрегирует панель состояния.
Например, утверждения «выполнено три звонка» достаточно, чтобы
определить, можно ли звонить снова, но оно не отвечает на вопрос «что
сотрудник службы поддержки пообещал во время второго разговора». Если
последующим задачам нужны только статистические значения, историю можно
заменить сводкой состояния, а исходные записи --- архивировать. Если
требуется проверить содержание разговора, исходный текст можно повторно
прочитать по индексу. Состав полей сводки определяет, на какие вопросы
можно ответить напрямую и когда необходимо обращаться к исходным
записям.

\hypertarget{ux441ux442ux440ux430ux442ux435ux433ux438ux438-ux441ux436ux430ux442ux438ux44f-ux43aux43eux43dux442ux435ux43aux441ux442ux430}{%
\section{Стратегии сжатия
контекста}\label{ux441ux442ux440ux430ux442ux435ux433ux438ux438-ux441ux436ux430ux442ux438ux44f-ux43aux43eux43dux442ux435ux43aux441ux442ux430}}

Предположим, что в ходе исследовательской задачи были прочитаны десять
длинных веб-страниц, а в окне осталось место лишь для двух материалов.
Сначала нужно спросить модель, что предстоит сделать на следующем этапе:
если требуется сравнить два предложения, следует сохранить цены, условия
их действия и источники; если проверить цитату --- место, по которому
можно получить исходный текст. Содержимое уже просмотренных, но
нерелевантных веб-страниц можно архивировать, освободив место для нового
ввода и вывода. При этом нужно рассчитать не только объём освобождённого
пространства, но и цену восстановления: сколько будет стоить повторно
найти удалённые материалы, если они понадобятся позже?

\hypertarget{ux442ux440ux438-ux446ux435ux43bux438-ux441ux436ux430ux442ux438ux44f}{%
\subsection{Три цели
сжатия}\label{ux442ux440ux438-ux446ux435ux43bux438-ux441ux436ux430ux442ux438ux44f}}

\textbf{Контроль длины и стоимости.} Например, длинные веб-страницы и
большие файлы могут занять значительную часть окна размером 128K
токенов. Harness должен резервировать бюджет для последующего ввода,
вывода модели и взаимодействия с инструментами.

\textbf{Подготовка знаний, необходимых для последующих решений.} В
результатах более десятка поисковых раундов одно событие может
упоминаться многократно, а некоторые утверждения могут противоречить
друг другу. Если привести их к виду «A и B подтверждены, для C всё ещё
не хватает доказательств», модель сможет продолжить работу над
нерешёнными вопросами. Резюме должно сохранять связи между фактами,
источники и открытые вопросы, чтобы последующие решения опирались на
ясные основания.

\textbf{Поддержка непрерывной работы над длительными задачами.}
Некоторые модели преждевременно завершают работу, когда считают, что
окно скоро будет исчерпано. Это явление называется \textbf{контекстной
тревожностью (Context Anxiety)}. При тестировании Sonnet 4.5 на
длительных задачах Anthropic обнаружила, что такое поведение может
сохраняться даже после сжатия, поэтому для продолжения работы после
сброса контекста применяются структурированные записи для передачи
состояния\footnote{Prithvi Rajasekaran, Harness design for long-running
  application development, Anthropic Engineering, 2026.}.

\hypertarget{ux43fux440ux435ux43eux431ux440ux430ux437ux43eux432ux430ux43dux438ux435-ux440ux435ux437ux443ux43bux44cux442ux430ux442ux43eux432-ux43fux43eux432ux442ux43eux440ux43dux44bux445-ux432ux44bux447ux438ux441ux43bux435ux43dux438ux439-ux432-ux43fux440ux438ux433ux43eux434ux43dux44bux435-ux434ux43bux44f-ux438ux441ux43fux43eux43bux44cux437ux43eux432ux430ux43dux438ux44f-ux437ux43dux430ux43dux438ux44f}{%
\subsection{Преобразование результатов повторных вычислений в пригодные
для использования
знания}\label{ux43fux440ux435ux43eux431ux440ux430ux437ux43eux432ux430ux43dux438ux435-ux440ux435ux437ux443ux43bux44cux442ux430ux442ux43eux432-ux43fux43eux432ux442ux43eux440ux43dux44bux445-ux432ux44bux447ux438ux441ux43bux435ux43dux438ux439-ux432-ux43fux440ux438ux433ux43eux434ux43dux44bux435-ux434ux43bux44f-ux438ux441ux43fux43eux43bux44cux437ux43eux432ux430ux43dux438ux44f-ux437ux43dux430ux43dux438ux44f}}

Панель состояния заранее поддерживает счётчики и прогресс, а резюме
дополнительно упорядочивает факты и связи из длинных текстов. В обоих
случаях разрозненная информация заблаговременно приводится к виду,
пригодному для непосредственного использования при последующих решениях.

Рассмотрим пример проверки зоомагазина:

\begin{quote}
Клетка 1: чёрная кошка. Клетка 2: белая кошка. Клетка 3: чёрная кошка.
Клетка 4: чёрная кошка. Клетка 5: белая кошка. \ldots{} (всего 100
клеток: 90 чёрных кошек и 10 белых)
\end{quote}

Чтобы ответить на вопрос «Какая кошка находится в клетке 37?», нужно
найти соответствующую запись; чтобы ответить «Сколько чёрных и сколько
белых кошек?», необходимо обобщить все записи. Если общие количества
впоследствии используются многократно, можно заранее вычислить и
сохранить: «90 чёрных кошек, 10 белых», --- чтобы каждый раз не начинать
подсчёт заново. Программа хорошо подходит для такого детерминированного
агрегирования, а модель может продолжить отвечать на прикладные вопросы
на основе полученной статистики.

Содержимое резюме следует выбирать в соответствии с задачей. Общих
количеств кошек двух типов достаточно, чтобы отвечать на вопросы о
численности; но для ответа о клетке 37 потребуются записи по каждой
клетке или соответствующий индекс.

Даже если в окне ещё есть место, чрезмерно длинный или неоднородный
контекст может мешать использованию информации: модель упускает ключевые
факты, повторно обрабатывает уже решённые вопросы или опирается на
устаревшее состояние при принятии текущих решений. Такая деградация
результатов по мере роста контекста обычно называется
\textbf{деградацией контекста (Context Rot)}. Она проявляется в
ухудшении выполнения задачи; если же запрос превышает предельную
ёмкость, происходит переполнение окна.

\textbf{Даже если в окне ещё есть место, многократно используемые знания
следует упорядочивать заранее.} Если ключевые выводы уже можно обобщить,
но модель по-прежнему вынуждена снова и снова искать их в длинной
истории, это продолжает расходовать вычислительные ресурсы и повышает
вероятность пропусков. В поисковой задаче можно сохранить подтверждённые
факты, источники и вопросы для дальнейшего поиска; в задаче
программирования --- архитектурные решения, область изменений,
результаты тестов и причины сбоев.

Примеры в контексте также влияют на то, как модель решает новые задачи.
В исследовании Dherin и соавторов определённые вычисления прямого
прохода представлены как зависящие от контекста низкоранговые обновления
весов, возникающие из сочетания слоёв внимания и MLP в Transformer. Это
предлагает одно из теоретических объяснений контекстного
обучения\footnote{Benoit Dherin et al., Learning without training: The
  implicit dynamics of in-context learning, 2025.}. Во время
фактического инференса веса модели остаются неизменными, а примеры
влияют на вывод через текущий ввод. Чтобы повторно использовать такой
опыт в последующих сеансах, его можно оформить как память, примеры или
Skill и загружать в зависимости от задачи; способы накопления опыта
подробнее рассматриваются в \crossreflink{chap:9}{главе 9}.

\hypertarget{ux43aux43eux43cux43fux440ux43eux43cux438ux441ux441ux44b-ux43cux435ux436ux434ux443-ux441ux436ux430ux442ux438ux435ux43c-ux438-kv-ux43aux44dux448ux435ux43c}{%
\subsection{Компромиссы между сжатием и
KV-кэшем}\label{ux43aux43eux43cux43fux440ux43eux43cux438ux441ux441ux44b-ux43cux435ux436ux434ux443-ux441ux436ux430ux442ux438ux435ux43c-ux438-kv-ux43aux44dux448ux435ux43c}}

Сжатие обычно выполняется между двумя вызовами модели. Harness с учётом
оставшегося бюджета и прогресса задачи подготавливает сообщения,
необходимые для следующего запроса, а затем отправляет их модели. Замена
истории изменяет префикс, поэтому необходимо учитывать одновременно
пользу сжатия и стоимость перестройки кэша.

Во-первых, системные инструкции и определения инструментов, которые
продолжают действовать, следует сохранять неизменными. Они находятся в
начале, поэтому при сжатии последующей истории эту часть префикса можно
использовать повторно. Во-вторых, нужно выбрать конкретные объекты
сжатия: длинные результаты инструментов можно заменить резюме, а
завершённое многораундовое взаимодействие --- объединить в запись об
этапе. Кэш до точки замены по-прежнему можно использовать, а содержимое
в точке замены и после неё необходимо обработать с учётом нового
префикса.

\textbf{Историю следует по возможности сжимать пакетно, не переписывая
один и тот же участок префикса в каждом раунде.} После сжатия ввод
каждого раунда сокращается, однако вызов сжатия и повторное вычисление
суффикса также расходуют ресурсы. Момент активации можно выбирать с
учётом заполненности окна, ожидаемого количества последующих раундов и
этапа задачи: пакетное упорядочивание после появления группы длинных
результатов инструментов и сохранение резюме для передачи состояния
после завершения этапа сокращают затраты на частое изменение истории.
При сохранении нескольких последних раундов следует убедиться, что
вызовы инструментов и их результаты образуют пары, а состояние
незавершённых операций записано в резюме.

При сжатии необходимо также проверить протокол обратной передачи
информации о рассуждениях. В некоторых моделях исторический thinking
привязан к префиксу, на основе которого был создан. Если заменить
историю перед такими блоками, сохранённый thinking может перестать
соответствовать новому префиксу. При таком протоколе сеанс можно
преобразовать в резюме, удалить затронутый прежний thinking и позволить
модели продолжить рассуждение на основе резюме; также можно использовать
поддерживаемые сервером механизмы сжатия или редактирования контекста.

Сначала рассмотрим индивидуальное резюмирование на
\crossreflink{fig:2-16}{рисунке 2-16}: после обработки осталось лишь
10,9\% исходного количества символов, однако совокупный расход токенов
основного AI-агента оказался самым высоким среди шести групп. Чтобы
объяснить, как эти два показателя могут одновременно быть верны, нужно
проследить, какие материалы были сжаты и в скольких последующих запросах
эти материалы снова отправлялись модели.

\begin{figure}
\hypertarget{fig:2-16}{%
\centering
\includegraphics{images/revised/fig2-16.pdf}
\caption{Рисунок 2-16. Сравнение стратегий сжатия
контекста}\label{fig:2-16}
}
\end{figure}

\begin{quote}
\textbf{Эксперимент 2-10 ★★★: сравнение стратегий сжатия контекста}

В этом эксперименте Kimi K3 используется для исследования
профессионального статуса соучредителей OpenAI. Для выполнения задачи
необходимо поэтапно составить список людей, найти сведения об их карьере
и изменениях должностей, а затем сопоставить соответствующие факты с
источниками. Чтобы наблюдать сжатие и обработку переполнения, бюджет
контекста на уровне приложения ограничен 128 000 токенов. В качестве
базового варианта используется режим без сжатия, который сравнивается
ещё с пятью подходами; каждый вариант выполняется не более 15
раундов\footnote{Файл результатов эксперимента 2-10
  \texttt{kimi\_k3\_real\_20260718.json} содержит статистические поля,
  итоговый ответ и временную метку для каждой стратегии, что позволяет
  напрямую повторно рассчитать коэффициент по символам. Запуск 18 июля
  2026 года использовал прежнее условие защиты; в изменении кода от 22
  июля --- коммите \texttt{aea4d72} --- совокупный ввод был заменён
  вводом последнего запроса.}.

\textbf{Стратегия 1: без сжатия.} Результаты поиска и чтения страниц
полностью включаются в историю. Семь вызовов инструментов дали 367 246
исходных символов. На пятом раунде программа активировала защиту и
остановилась, а совокупный расход основного агента составил 166 043
токена. В прежней версии защита активировалась по сумме входных токенов
всех запросов, поэтому число 165 227 в сообщении об ошибке обозначает
совокупный ввод. Например, если один и тот же фрагмент истории
отправлялся в трёх раундах, он учитывался трижды.

\textbf{Стратегия 2: индивидуальное резюмирование.} Для каждого
результата поиска отдельно создаётся резюме из двух-трёх абзацев,
сохраняющее основное содержимое соответствующей страницы. Такой подход
легко распараллелить, и он позволяет сохранить границы страниц; однако
если несколько страниц описывают одно событие, повторяющаяся информация
всё равно попадает в историю. В этой конфигурации итоговый ответ был
получен за двенадцать раундов, основной агент совокупно использовал 276
608 токенов, а коэффициент сжатия по символам составил 10,9\%.

\textbf{Стратегия 3: объединённое резюмирование.} Группа результатов
поиска объединяется, после чего создаётся общее резюме, позволяющее
упорядочить повторяющиеся факты и связи между разными источниками. В
этой конфигурации итоговый ответ был получен за десять раундов, основной
агент совокупно использовал 93 449 токенов, а коэффициент сжатия по
символам составил 4,3\%.

\textbf{Стратегия 4: контекстно-зависимое резюмирование.} При сжатии
одновременно передаются текущий запрос и уже известная информация:
например, в промпт добавляются
\texttt{Given\ the\ search\ query:\ \{query\}} и
\texttt{Current\ context:\ \{context\}}. На этой основе модель сохраняет
факты, относящиеся к текущему вопросу, и в первую очередь восполняет
пробелы существующего резюме. В этой конфигурации итоговый ответ был
получен за семь раундов, основной агент совокупно использовал 40 157
токенов, а коэффициент сжатия по символам составил 3,0\%. В одном случае
сжатия около 150 тысяч символов было преобразовано примерно в две тысячи
с сохранением имён людей и изменений их должностей.

\textbf{Стратегия 5: контекстно-зависимое резюмирование с цитированием.}
В относящемся к задаче резюме сохраняются ссылки на источники
соответствующих фактов. В этой конфигурации итоговый ответ был получен
за десять раундов, а основной агент совокупно использовал 222 992
токена. Согласно сводке записей инструментов, обработанные материалы
насчитывали 52 891 символ, а исходный текст --- 1 295 044 символа; их
отношение составляет около 4,1\%.

\textbf{Стратегия 6: адаптивное оконное сжатие.} На первых этапах
исходные материалы сохраняются, а затем пакетно упорядочиваются после
достижения вводом порогового значения. В эксперименте использовались три
механизма: сжатие активировалось при превышении 80\% бюджета; обобщались
ещё не сжатые результаты инструментов; обработанные результаты
помечались маркером \texttt{{[}COMPRESSED{]}}, чтобы исключить повторное
резюмирование. При бюджете эксперимента в 128 000 токенов порог
активации составлял 102 400 токенов. В этой конфигурации итоговый ответ
был получен за семь раундов, а основной агент совокупно использовал 174
601 токен. В исторической версии порог активировался по совокупному
вводу; позднее программу изменили так, чтобы она проверяла размер ввода
последнего запроса.

На \crossreflink{fig:2-17}{рисунке 2-17} сопоставлены ввод и вывод шести
стратегий.

\begin{figure}
\hypertarget{fig:2-17}{%
\centering
\includegraphics{images/revised/fig2-17.pdf}
\caption{Рисунок 2-17. Процессы обработки для шести стратегий
сжатия}\label{fig:2-17}
}
\end{figure}

При сравнении чисел на \crossreflink{fig:2-16}{рисунке 2-16} необходимо
различать три показателя: ввод отдельного запроса определяет достижение
бюджета окна; совокупное количество токенов отражает расход
многораундовых вызовов; коэффициент сжатия по символам описывает
отношение длины учтённых материалов после обработки к их исходной длине.
Здесь совокупный расход включает ввод и вывод всех запросов основного
агента, но не отдельные запросы модели, создававшие резюме.

Чтобы воспроизвести эксперимент, настройте модель и реальную поисковую
службу согласно сопроводительным инструкциям, а затем выполните в
каталоге эксперимента
\texttt{python\ experiment.py\ -m\ kimi-k3\ -n\ 15\ -o\ results/reproduction.json}.
Прочитайте \texttt{metrics} каждой стратегии: повторно рассчитайте
коэффициент по символам как
\texttt{total\_compressed\_size\ /\ total\_original\_size}, а затем
сравните \texttt{total\_tokens} и \texttt{iterations}. В базовом
варианте без сжатия к текстам результатов поиска также добавляется
форматирование, например заголовки и URL, поэтому количество символов
после обработки может немного превышать исходное. Получив итоговый
ответ, программа эксперимента считает соответствующую группу
завершённой. Попробуйте найти одного и того же человека в двух готовых
ответах: группа с цитированием указала местом работы Pamela Vagata
компанию Pebblebed, а в сводной таблице группы с оконным сжатием указана
Stripe. В первом ответе упомянут сайт организации, но открываемая ссылка
на веб-страницу не приведена. Обе группы предоставили сводные таблицы,
однако сведения о человеке в них противоречат друг другу.
\end{quote}

\hypertarget{ux43cux43dux43eux433ux43eux443ux440ux43eux432ux43dux435ux432ux44bux439-ux43cux435ux445ux430ux43dux438ux437ux43c-ux441ux436ux430ux442ux438ux44f-ux43fux440ux43eux43cux44bux448ux43bux435ux43dux43dux43eux433ux43e-ux443ux440ux43eux432ux43dux44f}{%
\subsection{Многоуровневый механизм сжатия промышленного
уровня}\label{ux43cux43dux43eux433ux43eux443ux440ux43eux432ux43dux435ux432ux44bux439-ux43cux435ux445ux430ux43dux438ux437ux43c-ux441ux436ux430ux442ux438ux44f-ux43fux440ux43eux43cux44bux448ux43bux435ux43dux43dux43eux433ux43e-ux443ux440ux43eux432ux43dux44f}}

Полные результаты инструментов иногда требуется временно сохранять,
завершённые шаги можно архивировать, а текущие ограничения должны
оставаться видимыми. Приведённые ниже пять типов механизмов
предназначены для разных частот использования и сроков хранения.

\begin{enumerate}
\def\labelenumi{\arabic{enumi}.}
\tightlist
\item
  \textbf{Управление бюджетом результатов инструментов}. Объёмный вывод
  сохраняется во внешнем хранилище, а в контексте остаются краткий
  предварительный обзор и точка входа для чтения. После включения обзора
  в историю его содержимое сохраняется неизменным; если нужны
  дополнительные подробности, соответствующий фрагмент читается
  отдельно, а результат добавляется в историю.
\item
  \textbf{Очистка нерелевантного содержимого}. Навигационные панели
  веб-страниц, повторяющиеся нижние колонтитулы и подобные элементы
  можно отфильтровать до попадания результата в контекст. Если требуется
  удалить уже включённое в историю содержимое, сначала следует
  проверить, не зависит ли от него дальнейшая задача или цитирование, а
  затем рассчитать стоимость изменения истории для кэша.
\item
  \textbf{Редактирование контекста на уровне API}. Службы,
  поддерживающие такую возможность, могут очищать старые результаты
  инструментов по заданным правилам. Клиент должен настроить условия
  активации и обработать возвращаемую службой информацию о
  редактировании. Область очистки и поведение кэша определяются
  контрактом интерфейса; этот механизм удобно подключать вместе с
  мониторингом бюджета ввода.
\item
  \textbf{Архивное резюмирование}. Цели, действия, результаты и
  незавершённые вопросы записываются по раундам или этапам задачи с
  сохранением порядка выполнения. Это можно рассматривать как журнал
  изменений задачи: впоследствии можно выяснить, на каком шаге было
  принято определённое решение и на чём оно основывалось.
\item
  \textbf{Сжатие всего сеанса}. Если локальной очистки по-прежнему
  недостаточно для соблюдения бюджета, текущая цель, ограничения,
  подтверждённые факты, ключевые решения и дальнейшие действия
  оформляются как резюме для передачи состояния. Сначала можно
  упорядочить существующие заметки сеанса, а затем обобщить всю историю.
  Если сжатие неоднократно завершается неудачей, механизм ограничения
  повторных попыток или автоматический выключатель прекращает повторные
  вызовы, сохраняет восстанавливаемое состояние, после чего можно
  уменьшить ввод или перейти к ручной обработке.
\end{enumerate}

Эти механизмы можно сочетать по уровням: контроль на входе уменьшает
разовые объёмные результаты, локальное упорядочивание сдерживает
повседневный рост, а резюме этапов и сжатие всего сеанса сохраняют
информацию для передачи длительных задач.

При чтении dsh-superagent автор обратил внимание на одну деталь: для
событий и результатов инструментов сначала можно передавать резюме и ID
объекта, а при необходимости получать подробности по этому ID.
Предположим, в фоне завершился набор тестов, а переднему плану сообщили
лишь: «Два теста завершились неудачей». Передний план может сообщить о
прогрессе, но не сможет ответить на последующий вопрос пользователя:
«Это те же два теста, что и вчера?» Если к резюме приложен идентификатор
записи текущего запуска тестов, передний план сможет получить имена
неудачных тестов и сопоставить их со вчерашней записью.

\hypertarget{ux43fux440ux438ux43dux446ux438ux43fux44b-ux43fux440ux43eux435ux43aux442ux438ux440ux43eux432ux430ux43dux438ux44f-ux441ux442ux440ux430ux442ux435ux433ux438ux438-ux441ux436ux430ux442ux438ux44f}{%
\subsection{Принципы проектирования стратегии
сжатия}\label{ux43fux440ux438ux43dux446ux438ux43fux44b-ux43fux440ux43eux435ux43aux442ux438ux440ux43eux432ux430ux43dux438ux44f-ux441ux442ux440ux430ux442ux435ux433ux438ux438-ux441ux436ux430ux442ux438ux44f}}

При выборе информации, сохраняемой после сжатия, можно проверить
следующие четыре аспекта. Процесс накопления опыта между задачами
рассматривается в \crossreflink{chap:9}{главе 9}.

\begin{itemize}
\tightlist
\item
  \textbf{Распределяйте объём по ценности для задачи}. Списки людей,
  связи между идентичностями и текущие вопросы для поиска обычно следует
  сохранять непосредственно; для ключевых доказательств, подтверждающих
  вывод, нужно оставлять источник и необходимые подробности;
  навигационные панели, рекламу и повторяющееся содержимое можно
  удалять. Если факт оспаривается, различия между свидетельствами также
  представляют важную информацию.
\item
  \textbf{Сохраняйте семантическую целостность}. Например, в утверждении
  «Sutskever покинул OpenAI в мае 2024 года» человек, время, действие и
  организация вместе образуют единый факт. При сжатии необходимо
  сохранить все эти элементы, не оставляя труднообъяснимый фрагмент
  вроде «Sutskever ушёл».
\item
  \textbf{Организуйте содержимое вокруг последующих вопросов}. Для
  задачи «найти список соучредителей» требуется полный охват людей, а
  для задачи «изучить биографию человека» --- опыт, хронология и
  источники. Один и тот же материал можно резюмировать по-разному в
  зависимости от задачи: при поиске важен охват, при анализе --- связи и
  основания, а при создании содержимого также необходимо сохранять
  полезные подробности и примеры.
\item
  \textbf{Проверяйте, поддерживает ли резюме выполнение задачи}.
  Семантическое резюме должно учитывать задачу, устранять повторы и
  сохранять условия и исключения. \textbf{Для семантического сжатия
  сложных материалов предпочтительно использовать модель, способность к
  пониманию которой близка к основной модели.} Выполняющая сжатие модель
  должна определить, какие условия, исключения и доказательства повлияют
  на последующую работу; если они потеряны, даже мощной основной модели
  будет трудно их восстановить.
\end{itemize}

Ранние архитектурные решения, причины ограничений и неудачные пути легко
теряются в длинной истории. \textbf{При выполнении длительных задач
необходимо регулярно обновлять документ о прогрессе}, сохраняя сведения
о том, «что решено, почему принято такое решение, что проверено и что
делать дальше». Промпт или Skill может задавать моменты обновления, а
Harness --- загружать соответствующие записи при передаче задачи. Такой
документ также помогает исполнителям, проверяющим и последующим агентам
понять уже выполненную работу.

\hypertarget{ux438ux437ux43eux43bux44fux446ux438ux44f-ux43aux43eux43dux442ux435ux43aux441ux442ux430-ux434ux43eux447ux435ux440ux43dux435ux433ux43e-ux430ux433ux435ux43dux442ux430}{%
\subsection{Изоляция контекста дочернего
агента}\label{ux438ux437ux43eux43bux44fux446ux438ux44f-ux43aux43eux43dux442ux435ux43aux441ux442ux430-ux434ux43eux447ux435ux440ux43dux435ux433ux43e-ux430ux433ux435ux43dux442ux430}}

\textbf{Для задач с чёткими границами результата и большим объёмом
промежуточных материалов автор предпочитает изоляцию посредством
дочернего агента.} При широком поиске читается множество материалов,
полезных лишь на этапе исследования; их предварительная изоляция
уменьшает продолжительное влияние этих материалов на основную задачу.
Главный агент передаёт задачу дочернему агенту с независимым контекстом;
дочерний агент изучает промежуточные результаты, а затем возвращает
выводы, местоположение доказательств и нерешённые вопросы. Благодаря
этому в контексте главного агента остаётся только информация,
необходимая для последующих решений.

Например, чтобы «найти функцию, обрабатывающую обратные вызовы платёжной
системы», главному агенту при самостоятельном поиске может потребоваться
прочитать более десятка файлов. Если поручить задачу поисковому
дочернему агенту, можно потребовать, чтобы он вернул имя функции,
модуль, точки вызова и основания для вывода. Его ответ может выглядеть
так: «За обработку обратных вызовов отвечает \texttt{handle\_callback} в
платёжном модуле; также существуют ещё две точки вызова», --- с
приложенными ссылками, позволяющими найти соответствующий код. На основе
этих результатов главный агент решает, нужно ли продолжать чтение
относящегося к задаче кода.

Описание задачи должно быть достаточно полным: какова цель, какова
область поиска, какие ограничения необходимо соблюдать, в каком формате
вернуть результат и как поступать при неоднозначности. Если запрос на
делегирование и возвращённый результат добавляются в историю, прежний
префикс главного агента по-прежнему можно использовать повторно.

В \crossreflink{chap:4}{главе 4} рассматривается проектирование
интерфейса дочернего агента как инструмента, а в
\crossreflink{chap:10}{главе 10} --- распределение контекста в
многоагентных системах.

\hypertarget{ux438ux442ux43eux433ux438-ux433ux43bux430ux432ux44b-1}{%
\section{Итоги
главы}\label{ux438ux442ux43eux433ux438-ux433ux43bux430ux432ux44b-1}}

Как помощнику службы поддержки действовать после третьего звонка? В
запросе следующего раунда нужно передать результат инструмента для этого
звонка, сопоставив его с операцией набора номера по ID вызова. В строке
состояния следует обновить число использованных попыток до трёх и
перечислить условия, которые ещё не удалось согласовать. Системные
правила и определения инструментов остаются в стабильном префиксе, а
соответствующий Skill загружается только тогда, когда нужно выполнить
конкретную операцию. Так происхождение каждого поля остаётся понятным,
как и то, какие данные изменятся при следующем вызове.

Если разговор затягивается и результатов инструментов становится всё
больше, можно ограничить их длину, свести уже завершённые части в
краткое резюме с указанием источников или передать независимо
выполняемую работу субагенту, получив от него только выводы и
подтверждающие материалы.

После завершения задачи счётчик трёх выполненных звонков можно
архивировать вместе с ней. Но если пользователь специально попросил: «В
дальнейшем связывайтесь со мной во второй половине дня по будням», а эта
фраза осталась только в резюме разговора, помощник, который примет
следующее поручение, может её вообще не увидеть. Как сохранять
информацию, полезную в долгосрочной перспективе, и как извлекать её в
подходящих задачах? В третьей главе рассматриваются сохраняющаяся между
задачами память пользователя и общая база знаний.

\begin{questionbox}

\hypertarget{ux432ux43eux43fux440ux43eux441ux44b-ux434ux43bux44f-ux440ux430ux437ux43cux44bux448ux43bux435ux43dux438ux44f-1}{%
\section*{Вопросы для
размышления}\label{ux432ux43eux43fux440ux43eux441ux44b-ux434ux43bux44f-ux440ux430ux437ux43cux44bux448ux43bux435ux43dux438ux44f-1}}
\addcontentsline{toc}{section}{Вопросы для размышления}

\begin{enumerate}
\def\labelenumi{\arabic{enumi}.}
\tightlist
\item
  ★★★ Изучите журнал работы со скользящим окном из эксперимента 2-3: что
  делал каждый из пяти вызовов инструментов? Какие дополнительные данные
  нужно записать и какое контрольное сравнение провести, чтобы
  определить, вызваны ли повторные вызовы обрезкой истории? Разработайте
  схему обрезки, сохраняющую ключевое состояние, и объясните её влияние
  на префиксное кэширование. Журнал находится в
  \texttt{chapter2/kv-cache/result\_\allowbreak{}sliding\_\allowbreak{}window\_\allowbreak{}20260718\_\allowbreak{}kimi\_\allowbreak{}k2\_\allowbreak{}6.\allowbreak{}json}.
\item
  ★★ Chat Template модели Qwen3 обрабатывает предыдущие рассуждения с
  учётом позиции последнего настоящего пользовательского запроса. Как бы
  вы управляли постепенно растущим объёмом рассуждений, если один цикл
  ReAct продолжается более ста итераций? Сравните обрезку предыдущих
  рассуждений, их повторную передачу и протокол привязки к префиксу, а
  также объясните, как Harness должен к ним адаптироваться.
\item
  ★★ Предположим, нужно сжать примерно 150 тысяч символов поисковых
  материалов до двух тысяч символов. Какие факты, связи и источники вы
  сохраните? Как проверить, достаточно ли краткого содержания для
  последующих вопросов, и как восстановить опущенные подробности?
\item
  ★★ Счётчик инструментов не учёл один вызов, поэтому строка состояния
  показывает неверное количество оставшихся вызовов. Разработайте
  механизмы регистрации событий, проверки сводных данных и контроля на
  уровне выполнения. Объясните, как с их помощью обнаружить и устранить
  такое несоответствие.
\item
  ★★ В эксперименте 2-4 сравнивались тон, организация информации и
  описания инструментов. Как бы вы управляли версиями системного
  промпта, который длительное время поддерживают несколько человек,
  проверяли конфликты между правилами и оценивали каждое изменение с
  помощью регрессионных задач?
\item
  ★★★ В примере с зоомагазином требуются как поиск отдельных записей,
  так и агрегированная статистика. Разработайте структуру исходных
  записей и статистических сводок, а также маршрутизацию запросов, чтобы
  система могла эффективно отвечать на вопросы обоих типов. Объясните,
  когда необходимо обновлять сводки.
\item
  ★★★ Skill уже находится в каталоге, но модель не обнаружила, что он
  применим к текущей задаче. Как бы вы искали причину проблемы с точки
  зрения описания, организации каталога, явной активации и примеров для
  оценки?
\item
  ★★ Skill успешно загружен, но модель пропустила предусмотренный в нём
  этап проверки. Разработайте контрольный эксперимент, позволяющий
  различить проблемы с активацией, пониманием, использованием
  инструментов и определением завершённости, а также сравните результаты
  разных моделей.
\item
  ★★★ Система содержит множество инструментов, а их набор часто
  меняется. Как бы вы расположили фиксированные определения, каталог
  инструментов, результаты загрузки по требованию и состояние во время
  выполнения? Учтите точность обнаружения, бюджет контекста и долю
  попаданий в кэш.
\end{enumerate}

\textbf{Практикум к \crossreflink{chap:2}{главе 2}.}
\href{https://github.com/bojieli/ai-agent-book/tree/40f4ce5778e6cc93e3606c3079bf8cf13a0c4d04/chapter2}{Каталог
экспериментов} ·
\href{https://github.com/bojieli/ai-agent-book/blob/40f4ce5778e6cc93e3606c3079bf8cf13a0c4d04/chapter2/README.ru.md}{Путеводитель
на русском}.

\end{questionbox}

\hypertarget{ux43fux430ux43cux44fux442ux44c-ux43fux43eux43bux44cux437ux43eux432ux430ux442ux435ux43bux44f-ux438-ux431ux430ux437ux430-ux437ux43dux430ux43dux438ux439}{%
\chapter{\texorpdfstring{Память пользователя и база
знаний\label{chap:3}}{Память пользователя и база знаний}}\label{ux43fux430ux43cux44fux442ux44c-ux43fux43eux43bux44cux437ux43eux432ux430ux442ux435ux43bux44f-ux438-ux431ux430ux437ux430-ux437ux43dux430ux43dux438ux439}}

Если на прошлой неделе пользователь сообщил помощнику, что предпочитает
места у окна, а сегодня сказал лишь: «Забронируй ещё один билет до
Токио», --- откуда взять сведения о его прежнем предпочтении для
выполнения нового запроса? А если вчера авиакомпания изменила правила
провоза багажа, где найти новые условия? В предыдущей главе
рассматривалась организация контекста одной задачи. В этой главе мы
разберём, откуда извлекать сказанное пользователем, где искать новые
правила авиакомпании и когда следует обновлять каждый из этих двух видов
данных.

\textbf{Память пользователя} организует сведения вокруг конкретного
пользователя: его предпочтения, привычки, потребности и опыт. Благодаря
ей AI-агент в последующих сеансах продолжает обслуживать пользователя с
учётом его особенностей. \textbf{База знаний}, напротив, доступна
нескольким пользователям в пределах определённых прав доступа: например,
она может содержать отраслевые нормативные требования, рабочие процессы
компании и техническую документацию. Память помогает агенту понимать,
кого он обслуживает, а база знаний предоставляет предметные знания,
необходимые для выполнения задачи.

Обе системы должны решать, как сохранять информацию и находить её при
необходимости. Обе сталкиваются с противоречиями между фактами,
устареванием содержимого и пропусками при поиске. Такие технологии, как
векторный поиск, структурированные индексы и сжатие знаний, применимы к
системам обоих типов. На \crossreflink{fig:3-1}{рисунке 3-1} показана
структура главы: сначала рассматриваются представление памяти
пользователя и управление ею, затем --- генерация с дополнением данными
из поиска (RAG), а в завершение методы поиска применяются к памяти между
сеансами и более сложной организации знаний.

\begin{figure}
\hypertarget{fig:3-1}{%
\centering
\includegraphics{images/revised/fig3-1.pdf}
\caption{Рисунок 3-1. Структура знаний этой главы}\label{fig:3-1}
}
\end{figure}

\hypertarget{ux441ux438ux441ux442ux435ux43cux430-ux43fux430ux43cux44fux442ux438-ux43fux43eux43bux44cux437ux43eux432ux430ux442ux435ux43bux44f}{%
\section{Система памяти
пользователя}\label{ux441ux438ux441ux442ux435ux43cux430-ux43fux430ux43cux44fux442ux438-ux43fux43eux43bux44cux437ux43eux432ux430ux442ux435ux43bux44f}}

Разберём фразу: «На следующей неделе я буду за границей, поэтому
банковские документы пришлите по электронной почте». Пользователь
сообщил о поездке с определённым сроком и указал способ получения
документов на время этой поездки. Если банк предложит отправить
документы обычной почтой, ассистент сможет на основании этой информации
запросить электронную доставку. Когда поездка закончится, при следующем
вопросе об отправке письма потребуется заново проверить текущие условия.
Если сохранить лишь запись «пользователь не любит бумажные письма», в
будущем эта информация найдётся в памяти, но будет применена за
пределами нужного контекста.

Рассмотрим также авиабилеты. Если пользователь скажет: «В прошлый раз
билет за 1250 оказался дороже, чем я ожидал», система может сохранить
как минимум предложенную тогда цену, имевшийся бюджет, обстоятельства
поездки и эту обратную связь. В следующий раз в похожей ситуации модель
будет осторожнее с предложениями, превышающими бюджет. Если пользователь
отдельно уточнит: «В дальнейшем при семейных поездках всегда сначала
спрашивайте меня», это задаст явное условие для всех будущих семейных
поездок. Таким образом, можно различить три вида информации: что именно
сказал пользователь, при каких условиях это произошло и какие выводы
модель сделала из этого случая.

В конце сеанса можно поручить LLM упорядочить эти записи и сохранить
факты, которые могут пригодиться в дальнейшем, в постоянном хранилище.
При начале новой задачи в контекст следует добавлять только относящиеся
к ней и по-прежнему применимые записи.

Например, при бронировании авиабилета пользователь предоставил следующую
информацию:

\begin{Shaded}
\begin{Highlighting}[]
\NormalTok{Пользователь: Помогите мне забронировать билет до Токио на следующую пятницу.}
\NormalTok{              Я предпочитаю места у окна и придерживаюсь вегетарианской диеты,}
\NormalTok{              поэтому мне понадобится специальное питание.}
\NormalTok{Агент: Я поищу рейсы до Токио на следующую пятницу...}
\NormalTok{       [вызывает инструмент flight\_search, получает 3 варианта]}
\NormalTok{Агент: Вот доступные варианты. С учётом ваших предпочтений я отобрал рейсы,}
\NormalTok{       где доступны места у окна. Забронировать прямой рейс ANA?}
\NormalTok{Пользователь: Да, и укажите мой номер United MileagePlus: 12345678.}
\end{Highlighting}
\end{Shaded}

После завершения сеанса модуль извлечения записей для памяти может
сформировать четыре записи:

\begin{Shaded}
\begin{Highlighting}[]
\NormalTok{Извлечённые записи для памяти:}
\NormalTok{{-} Пользователь предпочитает места у окна (предпочтение)}
\NormalTok{{-} Пользователь придерживается вегетарианской диеты и нуждается в специальном питании на рейсах (ограничение питания)}
\NormalTok{{-} Номер пользователя в программе United MileagePlus: 12345678 (программа лояльности)}
\NormalTok{{-} Пользователь планирует поездку в Токио (недавняя активность)}
\end{Highlighting}
\end{Shaded}

Этот пример включает три операции. \textbf{Избирательное извлечение}
сохраняет предпочтение места, требования к питанию и номер участника
программы лояльности, отбрасывая временные подробности вроде «поиск
вернул три варианта». \textbf{Абстрагирование} преобразует явно
выраженное предпочтение места у окна в правило, пригодное для будущих
бронирований. \textbf{Структурирование} присваивает фактам категории,
упрощая их поиск и обновление. Для поездки в Токио следует также
сохранить дату отправления и текущее состояние; выражение «следующая
пятница» необходимо преобразовать с учётом даты сеанса, а после
завершения или отмены поездки --- обновить состояние. Так долгосрочные
предпочтения и недавнюю активность можно поддерживать с учётом их разных
сроков действия.

В описании персонального ассистента Today приводился следующий пример.
Пользователь сообщил, что тренируется каждый вторник и четверг, готовясь
к полумарафону. Шесть недель спустя подключённые с разрешения
пользователя данные о сне показали, что минувшей ночью он спал очень
мало. Ассистент сопоставил эти сведения и предложил перенести
запланированную на этот день высокоинтенсивную интервальную тренировку.
Как долго следует хранить такую информацию? План тренировок актуален на
всём протяжении подготовки к забегу, а запись о сне за минувшую ночь
должна содержать дату. Даже если пользователь согласится перенести эту
тренировку, в следующий четверг она по-прежнему должна пройти по
исходному расписанию, если пользователь не попросит изменить его и в
дальнейшем. Today также разделяет предложение и изменение календаря на
два этапа: для второго требуется подтверждение пользователя.\footnote{Команда
  Today в статье «Что такое Today» на примере подготовки к полумарафону
  объясняет, как память сочетается с актуальной информацией за текущий
  день. В статье «Знакомство с Today» основатель объясняет механизм
  подтверждения перед изменением календаря.}

\hypertarget{ux43eux446ux435ux43dux43aux430-ux43fux430ux43cux44fux442ux438-ux442ux440ux451ux445ux443ux440ux43eux432ux43dux435ux432ux430ux44f-ux441ux438ux441ux442ux435ux43cux430}{%
\subsection{Оценка памяти: трёхуровневая
система}\label{ux43eux446ux435ux43dux43aux430-ux43fux430ux43cux44fux442ux438-ux442ux440ux451ux445ux443ux440ux43eux432ux43dux435ux432ux430ux44f-ux441ux438ux441ux442ux435ux43cux430}}

Открытый бенчмарк \textbf{LoCoMo} (Long-term Conversational Memory,
долговременная память диалогов) содержит длинные диалоги, в среднем
состоящие примерно из 300 реплик и охватывающие до 35 сеансов. Память
оценивается с помощью вопросов и ответов, составления сводок событий и
генерации мультимодальных диалогов. Вопросы делятся на пять категорий.
Ответ на одношаговый вопрос можно найти в одном сеансе, а для
многошагового требуется объединить сведения из нескольких сеансов.
Вопросы на временно́е рассуждение дополнительно требуют определить
последовательность событий. Вопросы на знания из открытого домена
предполагают объединение высказываний пользователя с внешними знаниями.
Вернёмся к примеру с бронированием билетов: пользователь сообщил лишь,
что «на следующей неделе едет в Токио». Чтобы ответить на вопрос «какая
валюта используется в месте назначения», необходимо и помнить о Токио, и
знать, что там используют иены. Если же спрашивается номер рейса,
который пользователь никогда не называл, помощник должен указать на
недостаток информации. Способность распознавать такие вопросы, на
которые невозможно дать обоснованный ответ, проверяется состязательными
заданиями.\footnote{Определения пяти категорий вопросов и ответов
  приведены в разделе 4.1 статьи Maharana и соавторов о LoCoMo
  (arXiv:2402.17753). Данные и примеры диалогов доступны на странице
  проекта авторов.}

Опираясь на подобные бенчмарки и практический опыт создания продуктов с
памятью, автор выделяет восемь способностей, которым необходимо уделять
внимание:

\begin{itemize}
\tightlist
\item
  \textbf{Сохранение персональных данных}: точное сохранение и
  извлечение долговременной информации, например сведений о личности
  пользователя.
\item
  \textbf{Отслеживание предпочтений}: фиксация предпочтений пользователя
  и их обновление при появлении новых явно выраженных пожеланий.
\item
  \textbf{Переключение контекста}: продолжение работы с релевантной
  исходной информацией после смены темы.
\item
  \textbf{Обновление памяти}: обработка дополнений и исправлений
  пользователя к ранее сохранённым фактам.
\item
  \textbf{Непрерывность между сеансами}: использование ранее сохранённой
  информации в новом сеансе.
\item
  \textbf{Сложные рассуждения}: связывание нескольких фрагментов памяти.
  Например, если пользователь ранее упоминал об аллергии на арахис, при
  обсуждении тайской кухни следует учесть это и напомнить о
  необходимости проверить состав блюд.
\item
  \textbf{Восприятие времени}: сохранение дат, интерпретация
  относительных указаний времени и вычисление интервалов между
  событиями.
\item
  \textbf{Разрешение противоречий}: обнаружение несогласованной
  информации и принятие решения с учётом источника, времени и объекта, к
  которому она относится.
\end{itemize}

Для практических задач AI-агента в этой книге дополнительно применяется
трёхуровневая система оценки. В экспериментах 3-1 и 3-2 с её помощью
сначала сравниваются представления памяти, а в экспериментах 3-9 и 3-11
по той же системе оцениваются методы поиска.

\textbf{Первый уровень: базовое воспроизведение.} Точное сохранение и
извлечение явно сформулированной информации, непосредственно
предоставленной пользователем. Например, если пользователь сказал: «Мой
номер участника программы --- 12345», при последующем бронировании
билета система должна вернуть номер полностью. На этом уровне главным
образом проверяется, не была ли информация потеряна, перепутана или
искажена.

\textbf{Второй уровень: поиск по нескольким сеансам.} Необходимо найти
всю релевантную информацию в относящихся к разным периодам и объектам
сеансах, а затем определить, кого касается текущий запрос и какие записи
нужны. Если у пользователя две машины, но он просит лишь «записать мою
машину на техническое обслуживание», система должна сначала извлечь
информацию об обеих машинах и уточнить, о какой из них идёт речь. При
запросе статуса кредита необходимо отличить действующий договор от
предложения, полученного во время прошлой консультации. При отмене
«поездки в Лос-Анджелес» требуется связать билеты, гостиницу и другие
бронирования с одной и той же поездкой.

\textbf{Третий уровень: проактивное обслуживание.} Исходя из текущей
задачи, необходимо выявить связанные с ней обстоятельства, которые
пользователь не перечислил явно. Например, при бронировании
международного рейса следует найти сохранённые несколько месяцев назад
паспортные данные и проверить, достаточен ли срок действия паспорта для
поездки. Если сломался телефон, нужно собрать сведения о гарантии на
устройство, дополнительной гарантии по кредитной карте и страховке
оператора, а затем перечислить доступные варианты действий. При
подготовке документов для налоговой декларации --- найти среди записей
за весь год сведения об операциях с акциями, доходах от внештатной
работы и документы по налогу на недвижимость.

\begin{quote}
\textbf{Эксперимент 3-1 ★: оценка системы памяти по трёхуровневой
системе}

Набор для оценки содержит по 20 примеров на каждом уровне --- всего 60.
На первом уровне преимущественно используются одиночные сеансы. Второй и
третий уровни включают несколько разнесённых во времени сеансов,
относящихся к разным объектам, и проверяют связывание сущностей,
обновление информации и проактивные напоминания.

Каждый пример выполняется в два этапа. Сначала диалоги подаются
последовательно: модель формирует память по первому фрагменту, а при
каждой последующей итерации получает только существующую память и новый
фрагмент диалога, после чего создаёт обновлённую память. Затем исходный
текст предыдущих диалогов становится недоступен модели, поэтому для
ответа на обсуждавшиеся в них вопросы она может использовать только
сведения, записанные в память. После обработки всех диалогов модели
задаётся новый вопрос, на который агент должен ответить с помощью
памяти.

Для оценки применяется LLM-as-a-judge: независимая модель-судья
сопоставляет ответ с эталоном и критериями оценивания, проверяя полноту
охвата фактов и наличие ошибок. В этом запуске Doubao Seed 1.6
использовалась для формирования памяти и ответов, а модель Moonshot ---
для выставления оценок. Модель-судья видит исходный диалог, тогда как
отвечающей модели доступна только память, сохранившаяся в конце.

Как читать ведомость результатов? Например, пользователь хочет лишь
найти номер учётной записи для услуги широкополосного доступа. В группе
с расширенными заметками помощник назвал правильный номер. Модель-судья
выставила по 4 балла за точность, полноту и рассуждение, но только 1
балл за проактивность --- и всё же задание было пройдено. Согласно
заданному в программе условию, для прохождения необходимо получить не
менее 3 из 4 баллов по каждому из первых трёх критериев, а модель-судья
не должна обнаружить существенных фактов, не имеющих оснований или
противоречащих исходным высказываниям. Проактивность учитывается в
совокупной оценке, поэтому одно лишь количество пройденных заданий не
показывает, насколько проактивно помощник оказывал помощь.
\end{quote}

\hypertarget{ux438ux435ux440ux430ux440ux445ux438ux44f-ux43fux430ux43cux44fux442ux438}{%
\subsection{Иерархия
памяти}\label{ux438ux435ux440ux430ux440ux445ux438ux44f-ux43fux430ux43cux44fux442ux438}}

Чтобы ответить на вопрос «Какой рейс был забронирован в прошлый раз?»,
нужно найти отдельное событие; чтобы ответить на вопрос «Какое место
выбирать в первую очередь при бронировании билета?», нужно извлечь
предпочтение, сохранённое по результатам нескольких взаимодействий. Оба
вопроса относятся к одному человеку, но требуют разной гранулярности
записей и разных способов поиска.

\textbf{Траектория (trajectory)} в хронологическом порядке фиксирует
события одного запуска AI-агента, включая сообщения пользователя, ответы
модели, вызовы инструментов и результаты их выполнения. При
использовании журнала событий, доступного только для добавления
(append-only), новые события дописываются в конец, а исходные записи
сохраняются для последующего отслеживания. Harness может выбирать из
траектории историю, извлекать состояние или формировать обобщение,
создавая контекст, который фактически отправляется модели при каждом
ходе.

\textbf{Долговременная пользовательская память} сохраняет полезную между
сеансами информацию, такую как предпочтения, отношения между людьми,
обобщения предыдущих действий и подтверждённые факты, и закрепляет её
принадлежность с помощью ID пользователя и других идентификаторов. Новые
формулировки могут дополнять, исправлять или заменять старые записи,
поэтому долговременную память необходимо объединять и обновлять. Агент
может явно читать и обслуживать её посредством вызовов инструментов,
либо Harness может запускать эти операции при загрузке или завершении
сеанса, а также во время фонового упорядочивания.

\textbf{Бизнес-состояние} описывает логический этап задачи, например
«требуется уточнение», «в обработке», «ожидание оплаты» или «выполнено».
Программа может проверять эти состояния, чтобы определить, какого
события ожидать дальше и какую операцию выполнить. Управление состоянием
в событийно-ориентированной архитектуре будет подробнее рассмотрено в
шестой главе.

Эти три типа записей могут работать совместно: траектория служит
основанием для выполнения, долговременная память предоставляет контекст
между сеансами, а бизнес-состояние показывает, до какого этапа дошло
выполнение текущей задачи. Далее основное внимание будет уделено
представлению, управлению и поиску в долговременной памяти.

\hypertarget{ux447ux435ux442ux44bux440ux435-ux444ux43eux440ux43cux430ux442ux430-ux445ux440ux430ux43dux435ux43dux438ux44f-ux43fux43eux43bux44cux437ux43eux432ux430ux442ux435ux43bux44cux441ux43aux43eux439-ux43fux430ux43cux44fux442ux438}{%
\subsection{Четыре формата хранения пользовательской
памяти}\label{ux447ux435ux442ux44bux440ux435-ux444ux43eux440ux43cux430ux442ux430-ux445ux440ux430ux43dux435ux43dux438ux44f-ux43fux43eux43bux44cux437ux43eux432ux430ux442ux435ux43bux44cux441ux43aux43eux439-ux43fux430ux43cux44fux442ux438}}

Если нужно изменить только адрес электронной почты, лучше иметь
возможность сразу найти соответствующую запись. Если же требуется
понять, почему адрес когда-то изменили, необходимо прочитать
предысторию. В левой верхней части \crossreflink{fig:3-2}{рисунка 3-2}
электронная почта, компания и должность хранятся в отдельных записях, а
в правой верхней части должность, команда и проект записаны вместе.
Теперь рассмотрим две нижние карточки: в каком поле следует искать
информацию, чтобы в будущем обновить должность или выяснить, кто именно
посещал кардиолога?

\begin{figure}
\hypertarget{fig:3-2}{%
\centering
\includegraphics{images/revised/fig3-2.pdf}
\caption{Рисунок 3-2. Сравнение четырёх стратегий памяти}\label{fig:3-2}
}
\end{figure}

\textbf{Simple Notes} сохраняют каждое воспоминание как краткий факт,
например: «Электронная почта пользователя:
\href{mailto:john@example.com}{\nolinkurl{john@example.com}}». Отдельные
записи легко читать и записывать. Если для поиска по стабильному ключу
используется хеш-таблица, средняя временная сложность поиска и
обновления обычно составляет O(1). Для семантического поиска по большому
количеству заметок дополнительно требуется поисковый индекс. Однако
чрезмерное дробление приводит к утрате связей. Если запись «Работает
старшим инженером в TechCorp и отвечает за разработку рекомендательной
системы» разделить на три не связанные между собой записи --- о
компании, должности и проекте, --- впоследствии придётся заново
определять, относятся ли они к одному месту работы.

\textbf{Enhanced Notes} сохраняют факты вместе с их контекстом в виде
полных абзацев. Например: «Пользователь работает старшим
инженером-программистом в TechCorp, уже три года занимается машинным
обучением и сейчас руководит проектом рекомендательной системы с
командой из пяти человек». Связи между компанией, должностью, опытом и
командой сохраняются в одном абзаце. Недостаток этого подхода состоит в
том, что отдельные факты могут повторяться в нескольких абзацах. После
смены должности необходимо найти и обновить все связанные абзацы, чтобы
новые и устаревшие записи не существовали одновременно.

\textbf{JSON Cards} организуют информацию по схеме «категория ---
подкатегория --- пары ключ-значение», например
\texttt{personal.contact.email} и \texttt{work.position.title}. При
изменении поля должности можно сохранить без изменений другие поля,
такие как \texttt{work.company.name}, поэтому этот формат удобен для
локальных обновлений. При проектировании категорий необходимо также
учитывать информацию, относящуюся сразу к нескольким из них. Например,
запись «По выходным разрабатывает личный проект на Python» одновременно
связана со временем, техническими предпочтениями и деятельностью. Можно
сохранить одну запись о деятельности и добавить к ней несколько тегов
либо связать соответствующие категории ссылками, сохранив эти отношения.

\textbf{Advanced JSON Cards} дополняют факты контекстом источника
\texttt{backstory}, субъектом \texttt{person}, отношением к пользователю
\texttt{relationship} и временной меткой. Эти поля помогают системе
определить, к кому относится информация, когда она была получена и в
каких ситуациях применима. Например, «доктор Чжан» может быть
стоматологом пользователя или кардиологом его отца. Если вместе с именем
сохранить пациента, родственную связь и медицинскую специальность,
вероятность путаницы снизится.

Когда пользователь говорит: «Помоги запланировать ежегодный медицинский
осмотр для членов моей семьи», система может найти родственников по
полям \texttt{person} и \texttt{relationship}, а затем прочитать
связанные сведения из \texttt{backstory}. Чем больше полей, тем больше
связей необходимо обрабатывать при извлечении и обновлении. Автор
рекомендует сочетать разные форматы хранения с учётом важности
информации и потребности в связях: для ключевых людей, договоров и
долгосрочных предпочтений использовать Advanced JSON Cards, сохраняя
источник и отношения, а для многочисленных простых фактов с небольшим
количеством связей --- Simple Notes, снижая затраты на извлечение и
сопровождение. Записи двух типов связываются идентификаторами или
ссылками и при необходимости извлекаются вместе.

\begin{quote}
\textbf{Эксперимент 3-2 ★★: сравнение четырёх форматов памяти}

Сопутствующая реализация предоставляет единый интерфейс создания,
обновления и извлечения воспоминаний для всех четырёх форматов. Итоговые
ответы сравниваются на том же наборе из 60 примеров, что и в
эксперименте 3-1.

В этом запуске Simple Notes, Enhanced Notes, JSON Cards и Advanced JSON
Cards прошли соответственно 48, 52, 51 и 49 примеров. На втором уровне
Enhanced Notes прошли 16 примеров, а Advanced JSON Cards --- 12; на
третьем уровне оба формата прошли по 19 примеров. С учётом сопровождения
памяти, формирования ответов и оценивания Enhanced Notes потребовали
около 990 тысяч токенов, а Advanced JSON Cards --- около 1,55 миллиона
токенов; среднее время вызова модели у последнего формата также
оказалось больше.

Рассмотрим в качестве примера задачу второго уровня о подписках на
стриминговые сервисы. Пользователь хочет сократить расходы и спрашивает,
на какие сервисы он подписан и сколько стоит каждый из них. Группа
Enhanced Notes перечислила ежемесячную стоимость Netflix, Disney+ и Max:
22,99, 19,99 и 19,99 доллара соответственно. Группа Advanced JSON Cards
сообщила только стоимость Max, дополнительно перечислив два более
дешёвых тарифа этого сервиса. Какими бы конкретными ни были
рекомендации, они не восполнят сведения о двух других счетах.

После какого обновления возникла проблема? Если просмотреть
воспоминания, последовательно сохранённые группой Advanced JSON Cards,
то после первого диалога среди них ещё присутствует Netflix, после
второго остаются только записи, связанные с Disney+, а после третьего
--- только с Max. К моменту вопроса сведения о первых двух подписках уже
исчезли из памяти.
\end{quote}

\hypertarget{ux440ux430ux441ux448ux438ux440ux435ux43dux43dux430ux44f-ux444ux43eux440ux43cux430-ux43fux440ux435ux434ux441ux442ux430ux432ux43bux435ux43dux438ux44f-ux437ux43dux430ux43dux438ux439-ux438ux441ux43fux43eux43bux43dux44fux435ux43cux44bux439-ux43aux43eux434}{%
\subsection{Расширенная форма представления знаний: исполняемый
код}\label{ux440ux430ux441ux448ux438ux440ux435ux43dux43dux430ux44f-ux444ux43eux440ux43cux430-ux43fux440ux435ux434ux441ux442ux430ux432ux43bux435ux43dux438ux44f-ux437ux43dux430ux43dux438ux439-ux438ux441ux43fux43eux43bux43dux44fux435ux43cux44bux439-ux43aux43eux434}}

После сохранения фактов необходимо также вычислять сводную статистику,
проверять конфликты и определять следующий шаг по бизнес-правилам. В
проекте \textbf{User as Code} (UaC) автор попытался организовать
информацию о пользователе в виде типизированных объектов Python, а
операции над ней оформить как функции\footnote{Li, Bojie, User as Code:
  Executable Memory for Personalized Agents, 2026.}. Модель извлекает
факты из диалога, а программа выполняет чётко определённые вычисления.

Для сравнения дат, суммирования денежных значений и проверки явных
бизнес-ограничений автор предпочитает оформлять вычисления как функции,
чтобы один и тот же набор фактов давал воспроизводимые и тестируемые
результаты. Благодаря этому модель может сосредоточиться на понимании
смысла, не тратя силы на повторные вычисления в уме и сверку правил в
длинном диалоге; представление и использование памяти также можно
проверять в рамках одной кодовой базы.

UaC использует двухэтапный процесс: сначала извлечённые факты
добавляются в журнал, а затем периодически создаётся контрольная точка,
восстанавливающая из журнала типизированное состояние. Журнал сохраняет
источники фактов и историю их изменений, а объект состояния
предоставляет актуальные поля для запросов.

Например, состояние с данными паспорта и поездок можно определить так:

\begin{Shaded}
\begin{Highlighting}[]
\ImportTok{from}\NormalTok{ dataclasses }\ImportTok{import}\NormalTok{ dataclass}
\ImportTok{from}\NormalTok{ datetime }\ImportTok{import}\NormalTok{ date}

\AttributeTok{@dataclass}
\KeywordTok{class}\NormalTok{ PassportInfo:}
\NormalTok{    number: }\BuiltInTok{str}
\NormalTok{    country: }\BuiltInTok{str}
\NormalTok{    expiry\_date: date}

\AttributeTok{@dataclass}
\KeywordTok{class}\NormalTok{ Trip:}
\NormalTok{    destination: }\BuiltInTok{str}
\NormalTok{    departure\_date: date}
\NormalTok{    is\_international: }\BuiltInTok{bool}

\AttributeTok{@dataclass}
\KeywordTok{class}\NormalTok{ UserState:}
\NormalTok{    passport: PassportInfo}
\NormalTok{    trips: }\BuiltInTok{list}\NormalTok{[Trip]}

\NormalTok{state }\OperatorTok{=}\NormalTok{ UserState(}
\NormalTok{    passport}\OperatorTok{=}\NormalTok{PassportInfo(}\StringTok{"AB1234567"}\NormalTok{, }\StringTok{"US"}\NormalTok{, date(}\DecValTok{2025}\NormalTok{, }\DecValTok{2}\NormalTok{, }\DecValTok{18}\NormalTok{)),}
\NormalTok{    trips}\OperatorTok{=}\NormalTok{[}
\NormalTok{        Trip(}\StringTok{"Tokyo"}\NormalTok{, date(}\DecValTok{2025}\NormalTok{, }\DecValTok{1}\NormalTok{, }\DecValTok{15}\NormalTok{), }\VariableTok{True}\NormalTok{),}
\NormalTok{        Trip(}\StringTok{"Paris"}\NormalTok{, date(}\DecValTok{2025}\NormalTok{, }\DecValTok{6}\NormalTok{, }\DecValTok{10}\NormalTok{), }\VariableTok{True}\NormalTok{),}
\NormalTok{    ],}
\NormalTok{)}
\end{Highlighting}
\end{Shaded}

\textbf{Агрегированная статистика} вычисляется прямым перебором записей.
Например, количество зарегистрированных международных поездок в 2025
году:

\begin{Shaded}
\begin{Highlighting}[]
\NormalTok{count }\OperatorTok{=} \BuiltInTok{sum}\NormalTok{(}
    \DecValTok{1} \ControlFlowTok{for}\NormalTok{ trip }\KeywordTok{in}\NormalTok{ state.trips}
    \ControlFlowTok{if}\NormalTok{ trip.is\_international }\KeywordTok{and}\NormalTok{ trip.departure\_date.year }\OperatorTok{==} \DecValTok{2025}
\NormalTok{)}
\CommentTok{\# count == 2}
\end{Highlighting}
\end{Shaded}

\textbf{Проверка конфликтов} сопоставляет два набора полей по заранее
заданным правилам. Ниже на примере совпадения классов лекарственных
препаратов показано, как поместить записи, требующие дополнительной
проверки, в очередь обработки:

\begin{Shaded}
\begin{Highlighting}[]
\KeywordTok{def}\NormalTok{ find\_allergy\_matches(profile):}
\NormalTok{    matches }\OperatorTok{=}\NormalTok{ []}
    \ControlFlowTok{for}\NormalTok{ medication }\KeywordTok{in}\NormalTok{ profile.current\_medications:}
        \ControlFlowTok{for}\NormalTok{ allergy }\KeywordTok{in}\NormalTok{ profile.allergies:}
            \ControlFlowTok{if}\NormalTok{ medication.drug\_class }\OperatorTok{==}\NormalTok{ allergy.drug\_class:}
\NormalTok{                matches.append((medication, allergy))}
    \ControlFlowTok{return}\NormalTok{ matches}
\end{Highlighting}
\end{Shaded}

Каждое совпадение сохраняет соответствующие записи о принимаемом
препарате и аллергии. Последующий процесс может сверить их, а затем по
правилам предметной области определить, следует ли выдать
предупреждение, передать случай на проверку или запросить дополнительные
сведения.

\textbf{Проверка ограничений} может автоматически запускаться при
обновлении состояния. Например, если по правилам поездки на дату
отправления паспорт должен оставаться действительным ещё не менее 180
дней, количество дней можно передать как параметр:

\begin{Shaded}
\begin{Highlighting}[]
\KeywordTok{def}\NormalTok{ passport\_alerts(state, required\_days):}
\NormalTok{    alerts }\OperatorTok{=}\NormalTok{ []}
    \ControlFlowTok{for}\NormalTok{ trip }\KeywordTok{in}\NormalTok{ state.trips:}
        \ControlFlowTok{if}\NormalTok{ trip.is\_international:}
\NormalTok{            days }\OperatorTok{=}\NormalTok{ (state.passport.expiry\_date }\OperatorTok{{-}}\NormalTok{ trip.departure\_date).days}
            \ControlFlowTok{if}\NormalTok{ days }\OperatorTok{\textless{}}\NormalTok{ required\_days:}
\NormalTok{                alerts.append((trip.destination, days))}
    \ControlFlowTok{return}\NormalTok{ alerts}

\NormalTok{alerts }\OperatorTok{=}\NormalTok{ passport\_alerts(state, required\_days}\OperatorTok{=}\DecValTok{180}\NormalTok{)}
\end{Highlighting}
\end{Shaded}

При расчёте с указанными в примере 180 днями на дату поездки в Токио до
окончания срока действия паспорта остаётся лишь 34 дня, а к дате
отправления в Париж паспорт уже просрочен. После замены паспорта
программа может заново проверить обе поездки с новой датой окончания его
действия.

\hypertarget{ux43aux43eux433ux43dux438ux442ux438ux432ux43dux43e-ux43dux430ux443ux447ux43dux44bux435-ux43eux441ux43dux43eux432ux44b-ux43fux43eux43bux44cux437ux43eux432ux430ux442ux435ux43bux44cux441ux43aux43eux439-ux43fux430ux43cux44fux442ux438}{%
\subsection{Когнитивно-научные основы пользовательской
памяти}\label{ux43aux43eux433ux43dux438ux442ux438ux432ux43dux43e-ux43dux430ux443ux447ux43dux44bux435-ux43eux441ux43dux43eux432ux44b-ux43fux43eux43bux44cux437ux43eux432ux430ux442ux435ux43bux44cux441ux43aux43eux439-ux43fux430ux43cux44fux442ux438}}

Когнитивная наука различает типы памяти по её содержимому и способу
использования. Эти понятия также помогают организовать долговременную
информацию AI-агента:

\begin{itemize}
\tightlist
\item
  \textbf{Эпизодическая память (episodic memory)} сохраняет конкретные
  события, их время, участников и ход. Пример для человека: «В прошлую
  среду ужинал с коллегами в итальянском ресторане». У агента ей может
  соответствовать запись «Пользователь планирует в следующую пятницу
  отправиться в Токио рейсом ANA» с указанием того, удалось ли
  впоследствии забронировать билет и менялось ли затем бронирование.
\item
  \textbf{Семантическая память (semantic memory)} сохраняет факты и
  понятия. Пример для человека: «Столица Италии --- Рим». У агента ей
  могут соответствовать записи «Пользователь --- вегетарианец» и
  «Пользователь предпочитает места у окна». Эти сведения могут быть как
  явно сообщены пользователем, так и обобщены на основе нескольких
  событий; сохранение источника помогает различать эти случаи.
\item
  \textbf{Процедурная память (procedural memory)} связана с навыками и
  порядком действий. Пример для человека: умение ездить на велосипеде. В
  инженерии агентов её аналогом могут служить многократно используемые
  процессы, например «найти прямые рейсы → подтвердить предпочтение по
  месту → использовать номер участника программы для часто летающих
  пассажиров → заказать питание».
\end{itemize}

Вернёмся к истории покупок из первой главы. То, что пользователь
«просмотрела десять товаров и купила один из них», позволяет сохранить
событие отдельной покупки, но почти ничего не говорит о том, что ей
понравится в будущем. Возможно, ей понравился этот фасон, а возможно,
товар просто можно было получить в тот же день. Если затем она скажет:
«При покупке подобных товаров сначала всегда проверяй, доступна ли
доставка в тот же день», появится ещё одно долгосрочное предпочтение с
определённой областью применимости. При следующей рекомендации
аналогичных товаров можно в первую очередь проверить срок доставки и
объяснить ей, из какой обратной связи взялся этот критерий отбора.

В таблице 3-1 эта классификация сопоставлена с рассмотренными ранее
уровнями памяти и форматами хранения, чтобы показать, на какой вопрос
отвечает каждый из трёх аспектов.

Таблица 3-1. Три аспекта проектирования памяти

\begin{longtable}[]{@{}
  >{\raggedright\arraybackslash}p{(\columnwidth - 4\tabcolsep) * \real{0.3006}}
  >{\raggedright\arraybackslash}p{(\columnwidth - 4\tabcolsep) * \real{0.3041}}
  >{\raggedright\arraybackslash}p{(\columnwidth - 4\tabcolsep) * \real{0.3954}}@{}}
\toprule\noalign{}
\begin{minipage}[b]{\linewidth}\raggedright
\hspace{0pt}Аспект проектирования
\end{minipage} & \begin{minipage}[b]{\linewidth}\raggedright
\hspace{0pt}Вопрос
\end{minipage} & \begin{minipage}[b]{\linewidth}\raggedright
\hspace{0pt}Конкретные категории
\end{minipage} \\
\midrule\noalign{}
\endhead
\bottomrule\noalign{}
\endlastfoot
\hspace{0pt}Уровень записи & \hspace{0pt}Какую функцию выполняет
информация? & \hspace{0pt}Траектория, долговременная пользовательская
память, бизнес-\allowbreak{}\hspace{0pt}состояние \\
\hspace{0pt}Формат хранения & \hspace{0pt}В какой структуре
представлять? & \hspace{0pt}Simple Notes, Enhanced Notes, JSON Cards,
Advanced JSON Cards \\
\hspace{0pt}Содержимое памяти & \hspace{0pt}Какой тип информации
сохранять? & \hspace{0pt}Эпизодическая память, семантическая память,
процедурная память \\
\end{longtable}

Например, «Пользователь предпочитает места у окна» относится к
семантической памяти и может храниться в долговременной памяти в формате
Simple Notes. Процесс бронирования билетов относится к процедурной
памяти и может храниться в карточке с указанием источника и условий
применимости.

\hypertarget{ux43fux440ux438ux43cux435ux440ux44b-ux444ux440ux435ux439ux43cux432ux43eux440ux43aux43eux432-ux43fux430ux43cux44fux442ux438}{%
\subsection{Примеры фреймворков
памяти}\label{ux43fux440ux438ux43cux435ux440ux44b-ux444ux440ux435ux439ux43cux432ux43eux440ux43aux43eux432-ux43fux430ux43cux44fux442ux438}}

Пользователь говорит: «Недавно я пересел к проходу», хотя в старой
записи указано: «Предпочитает места у окна». Система может обновить
предпочтение по месту при записи или сохранить оба высказывания и при
следующем бронировании определить, какое из них применимо. Сначала
рассмотрим, на каком этапе Mem0 обрабатывает такие конфликты, а затем
--- как Memobase раздельно записывает профиль пользователя и конкретные
события.

\textbf{Mem0: изменение этапа обработки конфликтов.} В статье 2025 года
и процессе v2 новые факты сравнивались со старыми при записи, тогда как
в последующей версии v3 применяется добавление фактов и гибридный поиск.
На \crossreflink{fig:3-3}{рисунке 3-3} сопоставлены эти два подхода.

\begin{figure}
\hypertarget{fig:3-3}{%
\centering
\includegraphics{images/revised/fig3-3.pdf}
\caption{Рисунок 3-3. Архитектура управления памятью
Mem0}\label{fig:3-3}
}
\end{figure}

В раннем процессе LLM сначала извлекает из диалога потенциальные факты,
затем выполняется поиск похожих существующих воспоминаний и принимается
решение \textbf{ADD}, \textbf{UPDATE}, \textbf{DELETE} или
\textbf{NOOP}. Если пользователь сначала говорит: «Я живу в Пекине», а
позднее: «Я переехал в Шанхай», система может обновить текущее место
жительства на Шанхай. Благодаря этому при запросе проще получить уже
упорядоченный факт, однако при обслуживании памяти необходимо правильно
различать исправления, изменения во времени и разные объекты, а также
сохранять основания обновлений, чтобы при обнаружении ошибки можно было
вернуться к ним для проверки. В статье также предложен вариант графовой
памяти \textbf{Mem0-g}, представляющий связанные знания с помощью
сущностей и отношений\footnote{Chhikara и др. \emph{Mem0: Building
  Production-Ready AI Agents with Scalable Long-Term Memory}. 2025.}.

В v3 автоматическое извлечение выполняется одним вызовом LLM, а записи
только добавляются посредством \textbf{ADD}. Сведения о Пекине и Шанхае
можно сохранить вместе с их временным контекстом, а текущее место
жительства определить при запросе с учётом контекста. При поиске
одновременно используются семантическое сходство, ключевые слова BM25 и
сопоставление сущностей; действия, завершение которых уже подтверждено,
также могут извлекаться как факты. Согласно руководству по миграции,
показатель LoCoMo вырос с 71.4 до 91.6, а LongMemEval --- с 67.8 до
93.4. В этой версии с открытым исходным кодом связи между сущностями
используются для внутреннего поиска, а внешний интерфейс графового
хранилища и возвращаемое значение \texttt{relations} удалены; графовая
память предоставляется управляемой платформой\footnote{Mem0. \emph{Open
  Source: Migrating to the New Memory Algorithm}.}.

Эта эволюция показывает, что система памяти может перенести разрешение
конфликтов с этапа записи на этап чтения. Если при записи ошибочно
объединить опыт двух людей или принять дополнительную информацию за
исправление, последующие ответы могут постоянно опираться на неверное
текущее состояние. Хранение только с добавлением позволяет учитывать при
принятии решения как новые, так и старые свидетельства, а также
сокращает число вызовов модели для последовательного принятия решений
при записи. Для часто изменяющейся пользовательской информации,
конфликтность которой трудно определить сразу, автор предпочитает
сначала сохранять факты вместе с их временем и источником, а затем
формировать актуальное состояние при ответе на конкретный вопрос.

\textbf{Memobase: профиль пользователя и память событий.} Memobase
организует долговременную информацию в виде двух частей. \textbf{Профиль
пользователя (profile)} сохраняет относительно устойчивые свойства по
темам и подтемам, например имя в \texttt{basic\_info}, интересы в
\texttt{interest} и должность в \texttt{work}. Разработчик может
настроить извлекаемые слоты, управляя охватом и гранулярностью профиля.
\textbf{Память событий (event memory)} записывает события из жизни
пользователя и время, когда они произошли, позволяя системе отвечать на
вопросы вроде «Когда мы в прошлый раз обсуждали бюджет?».

Memobase сначала помещает новый диалог в буфер. Когда в буфере
накапливается заданный объём данных или сеанс остаётся неактивным в
течение установленного времени, запускается пакетная обработка; её также
можно явно запустить при завершении сеанса. Благодаря этому несколько
сообщений обрабатываются в рамках одного процесса упорядочивания, а при
запросе читаются уже сформированные профиль и события. Асинхронное
обновление сокращает время ожидания извлечения памяти в ходе текущего
взаимодействия; если требуется немедленно использовать новейшие факты,
можно дождаться завершения текущей обработки\footnote{Memobase,
  официальное описание проекта, включая поля профиля, память событий,
  буфер и интерфейс синхронного обновления.}.

Если помощнику необходимо помнить как обстоятельства предыдущего
бронирования, так и предпочтения пользователя по месту и порядок
возврата билета, события, факты и процессы можно хранить раздельно, как
показано на \crossreflink{fig:3-4}{рисунке 3-4}. При получении каждой
задачи нужные сведения загружаются в рабочую память.

\begin{figure}
\hypertarget{fig:3-4}{%
\centering
\includegraphics{images/revised/fig3-4.pdf}
\caption{Рисунок 3-4. Эталонная архитектура совместной работы нескольких
типов памяти}\label{fig:3-4}
}
\end{figure}

Эпизодическая память сохраняет последовательности событий и такие
метаданные, как временные метки, эмоциональные маркеры и идентификаторы
задач, позволяя комбинировать поиск по времени и теме. Например, для
запроса «Когда мы в прошлый раз обсуждали бюджет?» одновременно
требуются тематическое сопоставление и временная сортировка.
Семантическая память сохраняет извлечённые из событий факты и
предпочтения, а процедурная --- многократно используемые процессы.

\textbf{Рабочая память (working memory)} организует текущую цель,
состояние задачи и извлечённую релевантную информацию. Harness выбирает
содержимое из траектории и долговременного хранилища и загружает его в
текущий контекст. Новые факты, появившиеся при выполнении задачи, после
упорядочивания записываются в долговременную память. Траектория
фиксирует, как происходили события, а рабочая память сосредоточена на
информации, необходимой для принятия следующего решения; они связаны
посредством идентификаторов событий и указаний источников.

При чтении пользовательской памяти с помощью RAG сначала по текущему
запросу находят набор потенциально релевантных воспоминаний, а затем
решают, какие из них следует поместить в рабочую память. Предположим,
сегодня помощник должен выбрать пользователю место и при поиске
одновременно находит две фразы: «В командировках я предпочитаю места у
окна» и «В прошлый раз я повредил колено, на обратном пути выбери мне
место у прохода». Обе тесно связаны с выбором места, но во второй
требование относится только к тому конкретному обратному рейсу. Даже
если при повторном ранжировании по релевантности поставить её на первое
место, это не отвечает на вопрос, применима ли она сегодня.

В такой ситуации автор считает более перспективным отбирать
потенциальные воспоминания с помощью модели принятия решений наподобие
Jev. Ей следует передать сведения о текущей поездке, актуальные
требования пользователя, а также исходную формулировку, дату и условия
применимости каждого воспоминания, чтобы по каждому из них вынести одно
из решений: «применимо в этом случае», «только для справки» или
«недостаточно оснований». Если сегодня пользователь тоже отправляется в
командировку, первое воспоминание можно использовать как основание для
выбора места. Если же пользователь только что сказал, что колено ещё не
зажило, стоит вернуться ко второму воспоминанию и задать уточняющий
вопрос.

\hypertarget{ux43cux435ux445ux430ux43dux438ux437ux43cux44b-ux441ux436ux430ux442ux438ux44f-ux438-ux443ux43fux43eux440ux44fux434ux43eux447ux438ux432ux430ux43dux438ux44f-ux43fux430ux43cux44fux442ux438}{%
\subsection{Механизмы сжатия и упорядочивания
памяти}\label{ux43cux435ux445ux430ux43dux438ux437ux43cux44b-ux441ux436ux430ux442ux438ux44f-ux438-ux443ux43fux43eux440ux44fux434ux43eux447ux438ux432ux430ux43dux438ux44f-ux43fux430ux43cux44fux442ux438}}

Пользователь, часто бронирующий авиабилеты, может десять раз упомянуть
«место у окна». Возможно, девять из этих упоминаний относятся к одному и
тому же бронированию. Если считать количество упоминаний
непосредственной мерой силы предпочтения, значимость одного события
окажется завышенной. При упорядочивании следует сначала определить
источник и событие, затем объединить повторяющиеся записи, проверить
наличие более поздних исправлений и лишь после этого решать, какие
закономерности можно использовать в будущих задачах. Встретив фразу «В
прошлый раз то место у окна было хорошим», можно поручить модели оценки
вроде Jev прочитать её вместе с соответствующим бронированием и
определить только одно: относится ли она к категории «воспоминание о
предыдущем перелёте», «требование к текущему бронированию» или
«выражение долгосрочного предпочтения». Если не удаётся сопоставить
фразу с бронированием или правдоподобны несколько толкований, следует
сохранить исходную формулировку и уточнить её при следующем
бронировании. Такое локальное решение позволяет не учитывать один и тот
же опыт многократно как новые свидетельства.

\begin{enumerate}
\def\labelenumi{\arabic{enumi}.}
\tightlist
\item
  \textbf{Отбор записей для упорядочивания.} Сигналами для ранжирования
  могут служить частота обращений, время создания записи, эмоциональный
  маркер и уникальность информации. Например, запись, созданную три дня
  назад, к которой обращались пять раз, с сильным эмоциональным маркером
  и без дубликатов, целесообразнее сохранить отдельно. Запись, созданную
  90 дней назад, к которой обращались лишь один раз и которая сильно
  дублирует три другие записи, можно в первую очередь объединить с
  соответствующими записями. Система может сначала по оценке отобрать
  записи, требующие упорядочивания, а затем на основе бизнес-правил
  решить, следует ли их сжать, объединить или удалить. Долгосрочные
  ограничения и действующие договоры необходимо сохранять до окончания
  срока их действия, даже если к ним обращаются редко.
\item
  \textbf{Кластеризация и формирование обобщений.} Похожие воспоминания
  объединяются в группы, а их общее содержание сохраняется в виде
  обобщения. Например, несколько разговоров о погоде можно упорядочить
  до «Пользователь часто спрашивает о погоде и особенно интересуется
  дождём». Конкретные даты и исходные запросы сохраняются в архиве, а
  обобщение связывается с исходными записями, чтобы к ним можно было
  вернуться позднее.
\item
  \textbf{Извлечение закономерностей и процессов.} Из нескольких записей
  о покупках можно вывести: «Предпочитает товары с хорошим соотношением
  цены и качества, придаёт большое значение отзывам пользователей», ---
  либо выделить последовательность выбора товаров.
\end{enumerate}

\hypertarget{ux437ux430ux449ux438ux442ux430-ux43aux43eux43dux444ux438ux434ux435ux43dux446ux438ux430ux43bux44cux43dux43eux441ux442ux438-ux43eux431ux435ux437ux43bux438ux447ux438ux432ux430ux43dux438ux435-ux436ux443ux440ux43dux430ux43bux43eux432}{%
\subsection{Защита конфиденциальности: обезличивание
журналов}\label{ux437ux430ux449ux438ux442ux430-ux43aux43eux43dux444ux438ux434ux435ux43dux446ux438ux430ux43bux44cux43dux43eux441ux442ux438-ux43eux431ux435ux437ux43bux438ux447ux438ux432ux430ux43dux438ux435-ux436ux443ux440ux43dux430ux43bux43eux432}}

Пользовательская память часто содержит идентификационные данные,
контактную информацию и сведения об учётных записях. Система должна с
учётом задачи определять, какие поля передавать на вход модели, а с
учётом эксплуатационных потребностей --- какие данные записывать в
журналы. Обезличивание можно выполнять перед записью в журнал, заменяя
конфиденциальные фрагменты метками типов или устойчивыми заполнителями и
сохраняя структуру, необходимую для диагностики проблем.

\begin{quote}
\textbf{Эксперимент 3-3 ★★: обезличивание журналов с помощью локальной
модели}

В этом эксперименте через Ollama вызывается локальная модель Qwen3-0.6B,
которая распознаёт конфиденциальные фрагменты в журналах непосредственно
на компьютере пользователя. Малые модели подходят для развёртывания на
потребительских устройствах; также можно переключаться на варианты 1.7B,
4B и другие, сравнивая качество распознавания и время обработки. Если
исходные журналы должны оставаться локальными, распознавание и замену
также следует выполнять на локальном компьютере, передавая последующим
сервисам уже обезличенный результат.

Тестовые данные включают структурированные номера, например номера
удостоверений личности и банковских карт, полуструктурированную
информацию, например адреса, а также конфиденциальные сведения,
раскрытые на естественном языке, например пароли. Согласно JSON Schema
модель возвращает тип информации, положение текста и уровень
уверенности, после чего программа заменяет соответствующие фрагменты
исходного текста. В контрольной группе используются регулярные
выражения, а в гибридной группе сначала выполняется поиск по регулярным
выражениям, после чего модель обрабатывает оставшийся текст.

Сначала рассмотрим две записи журнала. Первая содержит номер
удостоверения личности и номер мобильного телефона, но регулярное
выражение скрывает только номер телефона, поэтому полнота составляет
\(1/2\). Во второй записи регулярные выражения находят номер социального
страхования США и номер банковской карты, но также скрывают похожий по
формату номер заявки в службу поддержки. Из трёх отмеченных фрагментов
обезличить требовалось только два, поэтому точность составляет \(2/3\).
Если проверять лишь то, остались ли в журнале конфиденциальные значения,
вторая запись пройдёт проверку. Потеря важной зацепки обнаружится лишь
тогда, когда служба поддержки попытается найти по этой записи
соответствующую заявку.

Если вычислить полноту отдельно для каждого из 12 примеров, а затем
усреднить результаты, регулярные выражения, модель и гибридный подход
показывают соответственно около 50\%, 69\% и 78\%. После замены во всех
примерах остаётся соответственно 9, 3 и 2 конфиденциальных значения. При
выборе подхода необходимо учитывать и время обработки: на этой машине
регулярные выражения обрабатывают одну запись в среднем примерно за 0,1
миллисекунды, а модель и гибридный подход --- почти за 3 секунды. При
большом объёме журналов поля с устойчивым форматом можно сначала
обрабатывать регулярными выражениями, поручая модели только
дополнительную проверку фрагментов, для оценки которых требуется
учитывать контекст.
\end{quote}

По мере роста числа записей следующей задачей становится поиск среди
тысяч и десятков тысяч элементов информации, необходимых для текущей
задачи. Рассматриваемая далее технология RAG применима как к общим базам
знаний, так и к поиску по пользовательской памяти.

\hypertarget{ux43eux441ux43dux43eux432ux44b-rag-ux43fux43eux441ux442ux440ux43eux435ux43dux438ux435-ux43aux43eux43dux432ux435ux439ux435ux440ux430-ux43fux43eux43bux443ux447ux435ux43dux438ux44f-ux437ux43dux430ux43dux438ux439-ai-ux430ux433ux435ux43dux442ux43eux43c}{%
\section{Основы RAG: построение конвейера получения знаний
AI-агентом}\label{ux43eux441ux43dux43eux432ux44b-rag-ux43fux43eux441ux442ux440ux43eux435ux43dux438ux435-ux43aux43eux43dux432ux435ux439ux435ux440ux430-ux43fux43eux43bux443ux447ux435ux43dux438ux44f-ux437ux43dux430ux43dux438ux439-ai-ux430ux433ux435ux43dux442ux43eux43c}}

Магазин изменил правила возврата, и пользователь спрашивает: «Как мне
сейчас подать заявку на возврат товара, купленного в прошлом месяце?»
Чтобы ответить правильно, сначала нужно найти условия, применимые к
этому заказу, а затем ясно изложить требования и порядок действий.
\textbf{Генерация с дополнением из поиска} (Retrieval-Augmented
Generation, RAG) связывает эти два этапа, позволяя модели использовать
внутренние материалы магазина, а также правила, опубликованные уже после
завершения её обучения.

Начнём с ситуации, когда существует только одна версия правил. Ниже
процедура возврата в том же магазине запрашивается через упрощённый
интерфейс. Посмотрим, как результаты поиска попадают в ответ модели:

\begin{Shaded}
\begin{Highlighting}[]
\NormalTok{query }\OperatorTok{=} \StringTok{"процедура возврата"}
\NormalTok{results }\OperatorTok{=}\NormalTok{ retriever.search(query, top\_k}\OperatorTok{=}\DecValTok{2}\NormalTok{)}
\CommentTok{\# results = [}
\CommentTok{\#   "Правила возврата: полный возврат средств можно запросить в течение 7 дней после получения заказа; необходимо указать номер заказа. Средства поступят в течение 3–5 рабочих дней.",}
\CommentTok{\#   "Порядок возврата: откройте «Мои заказы», выберите заказ и нажмите «Запросить возврат»."}
\CommentTok{\# ]}
\NormalTok{answer }\OperatorTok{=}\NormalTok{ llm.generate(system}\OperatorTok{=}\StringTok{"Ты — помощник службы поддержки."}\NormalTok{, context}\OperatorTok{=}\NormalTok{results, question}\OperatorTok{=}\NormalTok{query)}
\CommentTok{\# Ответ содержит срок подачи заявки, необходимые сведения, порядок действий и срок поступления средств.}
\end{Highlighting}
\end{Shaded}

Сопоставьте \crossreflink{fig:3-5}{рисунок 3-5} и два результата поиска
в коде: порядок подачи заявки взят из «Порядка возврата», а поступления
средств нужно ждать 3--5 рабочих дней --- это следует из «Правил
возврата». Если в результатах поиска останется только порядок действий,
будут ли у модели основания сообщить пользователю, когда поступят
деньги? От ответа на этот вопрос зависит, как следует разбивать
руководство на фрагменты.

\begin{figure}
\hypertarget{fig:3-5}{%
\centering
\includegraphics{images/revised/fig3-5.pdf}
\caption{Рисунок 3-5. Процесс выполнения запроса RAG: поиск, дополнение
и генерация}\label{fig:3-5}
}
\end{figure}

\hypertarget{ux440ux430ux437ux431ux438ux435ux43dux438ux435-ux434ux43eux43aux443ux43cux435ux43dux442ux43eux432-ux43dux430-ux444ux440ux430ux433ux43cux435ux43dux442ux44b-chunking}{%
\subsection{Разбиение документов на фрагменты
(Chunking)}\label{ux440ux430ux437ux431ux438ux435ux43dux438ux435-ux434ux43eux43aux443ux43cux435ux43dux442ux43eux432-ux43dux430-ux444ux440ux430ux433ux43cux435ux43dux442ux44b-chunking}}

Предположим, в руководстве по послепродажному обслуживанию сначала
указано: «Товар можно вернуть в течение семи дней», а в следующем абзаце
перечислены исключения. Если извлечь только первое предложение, можно
пообещать возврат средств и за товар, на который это правило не
распространяется. Если же использовать всё руководство как единую
поисковую единицу, вместе с нужной информацией в неё попадут инструкции
по установке, настройке и устранению неполадок, из-за чего вектору будет
сложно выделить тему возврата средств. \textbf{Разбиение на фрагменты
(chunking)} позволяет выбрать промежуточный размер: текста во фрагменте
должно быть достаточно для объяснения всех относящихся к нему условий,
но при этом конкретный вопрос должен легко его находить. Длина фрагмента
также должна укладываться в бюджет входных данных модели эмбеддингов и
оставлять в контексте место для последующей генерации.

Существуют три распространённых метода.

При \textbf{разбиении на фрагменты фиксированного размера} текст делится
по количеству токенов: например, по 512 токенов на фрагмент с
перекрытием соседних фрагментов на 50--100 токенов. Благодаря перекрытию
содержимое возле границы может целиком попасть в другой фрагмент. Этот
метод упрощает контроль длины фрагментов, но граница может пройти внутри
абзаца, блока кода или таблицы.

При \textbf{рекурсивном разбиении или разбиении с учётом структуры}
текст прежде всего делится по естественным границам, таким как разделы и
абзацы, а слишком длинные части затем разбиваются на более мелкие
единицы, например предложения. Заголовки Markdown и элементы HTML могут
служить структурными ориентирами. Для блоков кода, таблиц и текста,
относящегося к заголовкам, также можно применять специальные правила
обработки. Для технических документов с чёткой структурой автор
рекомендует начинать именно с этого метода: сначала использовать
существующее деление на разделы и абзацы, а затем дополнительно
разбивать слишком длинные части. Обычно так удаётся лучше сохранить
целостный смысл, чем при простом отсечении текста.

При \textbf{семантическом разбиении} вычисляется сходство эмбеддингов
соседних предложений или групп предложений, а текст разделяется там, где
тема заметно меняется. Этот метод требует дополнительных вычислений
эмбеддингов, но может помочь сохранить материал по одной теме в одном
фрагменте.

Предположим, модуль поиска вернул только одно предложение: «Доход
компании вырос на 3\%». Какой компании? За какой квартал? Число найдено
верно, но основание для ответа на вопрос осталось в предыдущем абзаце.
Если же поместить во фрагмент весь финансовый отчёт, название компании и
квартал сохранятся, но при запросе о доходе вместе с ними попадут абзацы
о найме сотрудников и закупке оборудования. Всё это содержимое также
придётся представить одним вектором, поэтому тема дохода может стать в
нём менее заметной. На практике при выборе размера можно начать с
256--1024 токенов на фрагмент и перекрытия соседних фрагментов на
10--20\%.

Вместе с фрагментом можно сохранять метаданные, например заголовок,
источник и время. Рассматриваемый далее контекстно-зависимый поиск
дополнительно снабжает каждый фрагмент сведениями о его месте в исходном
тексте, помогая при независимом поиске разрешать ссылки и определять
область действия.

\hypertarget{ux43fux43bux43eux442ux43dux44bux435-ux44dux43cux431ux435ux434ux434ux438ux43dux433ux438-ux43fux440ux435ux434ux441ux442ux430ux432ux43bux435ux43dux438ux435-ux442ux435ux43aux441ux442ux430-ux432ux435ux43aux442ux43eux440ux430ux43cux438}{%
\subsection{Плотные эмбеддинги: представление текста
векторами}\label{ux43fux43bux43eux442ux43dux44bux435-ux44dux43cux431ux435ux434ux434ux438ux43dux433ux438-ux43fux440ux435ux434ux441ux442ux430ux432ux43bux435ux43dux438ux435-ux442ux435ux43aux441ux442ux430-ux432ux435ux43aux442ux43eux440ux430ux43cux438}}

\textbf{Эмбеддинг (embedding)} отображает слово, предложение или
фрагмент текста в числовой вектор, например
\texttt{{[}0.2,\ -0.5,\ 0.8,\ ...{]}}. Поисковая модель обучается так,
чтобы релевантные запросы и документы имели высокое сходство в векторном
пространстве. «Плотность» означает, что большинство измерений вектора
имеют ненулевые значения и совместно представляют информацию текста; в
рассматриваемом далее разреженном представлении ненулевыми остаются лишь
немногие измерения.

Векторное пространство можно представить как многомерную карту: каждому
фрагменту текста соответствует определённая точка, а поиск заключается в
нахождении содержимого, близкого к запросу. Одна из распространённых мер
--- \textbf{косинусное сходство}:

\[
\operatorname{cos}(A,B)=\frac{A\cdot B}{\|A\|\,\|B\|}
\]

Скалярное произведение \(A\cdot B\) вычисляется умножением
соответствующих измерений с последующим суммированием, а норма \(\|A\|\)
--- как квадратный корень из суммы квадратов всех измерений. После
деления на произведение норм результат отражает только соотношение
направлений. Для ненулевых векторов косинусное сходство лежит в
диапазоне \([-1,1]\): чем оно ближе к 1, тем сильнее совпадают
направления. В разных поисковых моделях также применяются скалярное
произведение или евклидово расстояние. Если оба вектора нормализованы до
единичной длины, их скалярное произведение равно косинусному сходству.

\begin{quote}
\textbf{Расчёт вручную: косинусное сходство трёх предложений}

Представим фразы «Как ухаживать за кошкой» вектором \(A=(0.9,0.5,0.1)\),
«Руководство по уходу за кошками» --- вектором \(B=(0.8,0.6,0.1)\), а
«Стратегия инвестирования в акции» --- вектором \(C=(0.1,0.1,0.9)\).

\(A\cdot B=0.9\times0.8+0.5\times0.6+0.1\times0.1=1.03\). Нормы этих
векторов равны соответственно \(\sqrt{1.07}\) и \(\sqrt{1.01}\), поэтому
сходство составляет приблизительно \textbf{0,991}.

\(A\cdot C=0.9\times0.1+0.5\times0.1+0.1\times0.9=0.23\), норма \(C\)
равна \(\sqrt{0.83}\), а сходство составляет приблизительно
\textbf{0,244}. В этом наборе векторов запрос об уходе за кошкой ближе к
руководству по уходу.
\end{quote}

На практике модели обычно используют представления размерностью от
нескольких сотен до более чем тысячи.

\hypertarget{ux43eux442-word2vec-ux43a-ux43aux43eux43dux442ux435ux43aux441ux442ux43dux43e-ux437ux430ux432ux438ux441ux438ux43cux44bux43c-ux43fux440ux435ux434ux441ux442ux430ux432ux43bux435ux43dux438ux44fux43c}{%
\subsubsection{От Word2Vec к контекстно-зависимым
представлениям}\label{ux43eux442-word2vec-ux43a-ux43aux43eux43dux442ux435ux43aux441ux442ux43dux43e-ux437ux430ux432ux438ux441ux438ux43cux44bux43c-ux43fux440ux435ux434ux441ux442ux430ux432ux43bux435ux43dux438ux44fux43c}}

\crossreflink{fig:3-6}{Рисунок 3-6} начинает обзор с фиксированных
векторов слов. \textbf{Word2Vec} обучает для каждого слова фиксированный
вектор на основе совместной встречаемости слов. Некоторые отношения
проявляются в векторных операциях, например в классическом выражении
«король − мужчина + женщина ≈ королева». Ограничение фиксированных
представлений также очевидно: для \texttt{bank} в выражениях
\texttt{river\ bank} («берег реки») и \texttt{investment\ bank}
(«инвестиционный банк») используется один и тот же вектор слова.

\textbf{GloVe} подсчитывает совместную встречаемость слов во всём
корпусе, а затем обучает векторы слов с помощью взвешенной целевой
функции наименьших квадратов. Word2Vec ориентирован на задачу
предсказания в локальном окне, тогда как GloVe напрямую включает в
обучение глобальную статистику совместной встречаемости; оба метода
обучают фиксированные представления слов\footnote{Pennington и др.,
  GloVe: Global Vectors for Word Representation, 2014.}.

\begin{figure}
\hypertarget{fig:3-6}{%
\centering
\includegraphics{images/revised/fig3-6.pdf}
\caption{Рисунок 3-6. Эволюция технологий плотных
эмбеддингов}\label{fig:3-6}
}
\end{figure}

В контекстных моделях наподобие \textbf{BERT} представление одного и
того же токена меняется в зависимости от предложения. Благодаря
механизму самовнимания модель может по контекстам «выпустила новый
продукт» и «купил два килограмма» вычислить разные представления слова
«яблоко» в этих двух предложениях. Модели эмбеддингов для поиска также
должны агрегировать представления токенов в вектор запроса или документа
и обучаться так, чтобы эти векторы можно было сравнивать.

\textbf{Sentence-BERT} сочетает BERT с пулингом и сиамской либо
триплетной сетевой архитектурой, обучая векторы предложений, пригодные
для непосредственного сравнения. Документы можно заранее закодировать и
сохранить в индексе. При поступлении запроса достаточно закодировать
только его, а затем вычислить сходство векторов. Предварительное
кодирование устраняет затраты на повторное вычисление представлений всех
документов при каждом запросе\footnote{Reimers и Gurevych,
  Sentence-BERT: Sentence Embeddings using Siamese BERT-Networks, 2019.}.

\textbf{BGE-M3} одновременно предоставляет плотные, разреженные и
многовекторные представления для поиска, а также поддерживает несколько
языков и длинные входные данные; в этом разделе его плотное
представление используется для объяснения одновекторного
поиска\footnote{В статье Chen et al.~о M3-Embedding (arXiv:2402.03216)
  описано совместное обучение трёх видов поисковых представлений. См.
  также описание модели BGE-M3 от BAAI, включающее способы
  представления, размерность, длину входа и методы использования.}.
Допустимая для модели длина входа влияет на стратегию разбиения на
фрагменты: типичный предел исходной BERT составляет 512 токенов, а
BGE-M3 поддерживает до 8192 токенов.

\begin{quote}
\textbf{Эксперимент 3-4 ★★: сравнение двух индексов приближённого поиска
ближайших соседей}

После построения векторов документов необходимо быстро находить векторы,
близкие к запросу. Попарное вычисление позволяет найти точных ближайших
соседей, но с ростом объёма данных индексы \textbf{приближённого поиска
ближайших соседей} (Approximate Nearest Neighbor, ANN) обменивают
контролируемое снижение полноты на скорость выполнения запросов. В этом
эксперименте ANNOY и HNSW сравниваются на одном наборе реальных
эмбеддингов и одинаковых запросах, а результаты точного поиска
используются как эталон. Например, если точный поиск для некоторого
запроса определил десять ближайших соседей, а среди десяти результатов
приближённого индекса оказалось восемь из них, то
\href{mailto:Recall@10}{\nolinkurl{Recall@10}} для этого запроса равен
0,8. Эта метрика показывает, какую долю точных ближайших соседей
сохранил индекс.

ANNOY сужает область кандидатов с помощью множества деревьев случайных
проекций. HNSW использует многоуровневый граф ближайших соседей: сначала
быстро приближается к целевой области на разреженных верхних уровнях, а
затем переходит к поиску на более плотных нижних уровнях. На
\crossreflink{fig:3-7}{рисунке 3-7} показана многоуровневая структура
HNSW.

\begin{figure}
\hypertarget{fig:3-7}{%
\centering
\includegraphics{images/revised/fig3-7.pdf}
\caption{Рисунок 3-7. Структура индекса HNSW}\label{fig:3-7}
}
\end{figure}

Таблица 3-2. Компромиссы при реализации ANNOY и HNSW

\begin{longtable}[]{@{}
  >{\raggedright\arraybackslash}p{(\columnwidth - 4\tabcolsep) * \real{0.2894}}
  >{\raggedright\arraybackslash}p{(\columnwidth - 4\tabcolsep) * \real{0.3596}}
  >{\raggedright\arraybackslash}p{(\columnwidth - 4\tabcolsep) * \real{0.3510}}@{}}
\toprule\noalign{}
\begin{minipage}[b]{\linewidth}\raggedright
\hspace{0pt}Аспект сравнения
\end{minipage} & \begin{minipage}[b]{\linewidth}\raggedright
\hspace{0pt}ANNOY
\end{minipage} & \begin{minipage}[b]{\linewidth}\raggedright
\hspace{0pt}HNSW (на примере hnswlib)
\end{minipage} \\
\midrule\noalign{}
\endhead
\bottomrule\noalign{}
\endlastfoot
\hspace{0pt}Структура индекса & \hspace{0pt}Множество деревьев случайных
проекций & \hspace{0pt}Многоуровневый граф ближайших соседей \\
\hspace{0pt}Затраты на построение & \hspace{0pt}Зависят от числа
деревьев, числа векторов и размерности & \hspace{0pt}Зависят от числа
связей, диапазона поиска при построении и объёма данных \\
\hspace{0pt}Память и файлы & \hspace{0pt}Статический индекс поддерживает
отображение в память и совместное использование несколькими процессами &
\hspace{0pt}Хранит векторы и связи графа; число связей влияет на
затраты \\
\hspace{0pt}Инкрементальное обновление & \hspace{0pt}Добавление
элементов после построения индекса требует его перестроения &
\hspace{0pt}Поддерживает добавление, обновление и пометку элементов как
удалённых; требуется управление вместимостью \\
\hspace{0pt}Полнота и скорость & \hspace{0pt}Настраиваются числом
деревьев и бюджетом поиска при запросе & \hspace{0pt}Настраиваются
параметрами построения и диапазоном поиска при запросе \\
\hspace{0pt}Сценарий использования & \hspace{0pt}Статический индекс,
удобный для распространения и совместного использования &
\hspace{0pt}Динамический индекс с постоянной записью и выполнением
запросов \\
\end{longtable}

В этом запуске индекс сначала строился по 240 фрагментам текста. На
одних и тех же 20 запросах оба индекса нашли все 10 ближайших соседей,
полученных точным поиском. При первоначальном построении базы ANNOY
создавал индекс в контейнере Linux приблизительно за 13 миллисекунд, а
HNSW --- в локальной среде приблизительно за 78 миллисекунд; размеры
сохранённых индексов составили около 1,40 и 1,02 МБ соответственно.
Затем были добавлены оставшиеся 60 фрагментов: индекс ANNOY перестроили,
а в HNSW новые векторы вставили напрямую. После этого база содержала 300
фрагментов, поэтому исходный набор ближайших соседей также мог
измениться. После повторного точного поиска по новой базе средняя
полнота обоих индексов по-прежнему составляла 100\%.
\end{quote}

\textbf{Индекс следует выбирать с учётом способа обновления базы
знаний.} Если корпус периодически публикуется пакетами, а несколько
процессов совместно используют один индекс только для чтения, автор
прежде всего рассматривает ANNOY. Если новые документы поступают
постоянно и требуется их добавление в реальном времени, предпочтение
отдаётся HNSW. Модель эмбеддингов определяет семантическое
представление, а индекс --- возможность использовать эти представления с
приемлемыми задержкой и стоимостью обновления. Вместе они определяют
качество поискового сервиса.

Изменяя число деревьев и бюджет поиска ANNOY либо число связей и
диапазон поиска HNSW, можно наблюдать изменения полноты и
задержки\footnote{В документации проекта Spotify Annoy описаны
  построение индекса только для чтения, его загрузка посредством
  отображения в память и настройка числа проверяемых при запросе узлов с
  помощью \texttt{search\_k}. В разделе документации hnswlib о
  Python-интерфейсе приведены метод \texttt{add\_items} для добавления
  векторов и метод \texttt{set\_ef} для настройки диапазона поиска при
  запросе.}. Полнота индекса показывает, были ли найдены целевые
ближайшие соседи. Способны ли эти соседи дать ответ на вопрос
пользователя, зависит уже от представления эмбеддингов и содержания
материалов.

\hypertarget{ux440ux430ux437ux440ux435ux436ux435ux43dux43dux43eux435-ux43fux440ux435ux434ux441ux442ux430ux432ux43bux435ux43dux438ux435-ux43eux442-ux43aux43bux44eux447ux435ux432ux44bux445-ux441ux43bux43eux432-ux43a-ux43eux446ux435ux43dux43aux435-ux440ux435ux43bux435ux432ux430ux43dux442ux43dux43eux441ux442ux438}{%
\subsection{Разреженное представление: от ключевых слов к оценке
релевантности}\label{ux440ux430ux437ux440ux435ux436ux435ux43dux43dux43eux435-ux43fux440ux435ux434ux441ux442ux430ux432ux43bux435ux43dux438ux435-ux43eux442-ux43aux43bux44eux447ux435ux432ux44bux445-ux441ux43bux43eux432-ux43a-ux43eux446ux435ux43dux43aux435-ux440ux435ux43bux435ux432ux430ux43dux442ux43dux43eux441ux442ux438}}

\textbf{Разреженное представление} использует векторы высокой
размерности, в которых лишь небольшое число измерений имеет ненулевые
значения. В традиционном поиске по ключевым словам каждому терму
соответствует отдельное измерение, а вес присваивается словам,
встречающимся в документе. В основе этого подхода лежит \textbf{модель
мешка слов} (Bag of Words, BoW): она сохраняет слова и количество их
вхождений, игнорируя порядок. Поэтому фразы «кошка гонится за собакой» и
«собака гонится за кошкой» имеют одинаковое представление в виде мешка
слов.

Такое представление хорошо подходит для сопоставления названий, номеров
и технических терминов. Начнём с TF-IDF и рассмотрим, как частота терма
и частота документов преобразуются в оценку для ранжирования.

\hypertarget{ux43eux442-tf-idf-ux43a-bm25}{%
\subsubsection{От TF-IDF к BM25}\label{ux43eux442-tf-idf-ux43a-bm25}}

\textbf{TF-IDF} (Term Frequency--Inverse Document Frequency, частота
терма --- обратная частота документов) объединяет два сигнала: сколько
раз слово встречается в текущем документе и насколько часто оно
встречается во всём корпусе. Предположим, что из 100 статей 60 содержат
слово «модель» и только 3 --- слово «дистилляция». Тогда в запросе
«дистилляция модели» слово «дистилляция» обладает большей различительной
способностью.

Один из базовых вариантов формулы выглядит так:

\[\operatorname{TF\!\text{-}\!IDF}(t,d)=\operatorname{TF}(t,d)\times\operatorname{IDF}(t),\qquad \operatorname{IDF}(t)=\ln\frac{N}{\operatorname{DF}(t)}\]

\(\operatorname{TF}(t,d)\) обозначает число вхождений терма \(t\) в
документ \(d\), \(\operatorname{DF}(t)\) --- число документов,
содержащих этот терм, а \(N\) --- общее число документов. Согласно этой
формуле, вклад терма, встретившегося десять раз, вдвое выше вклада при
пяти вхождениях; длинные документы также могут получать более высокие
оценки лишь потому, что содержат больше слов. Существуют различные
варианты TF-IDF с нормализацией, а рассматриваемый далее BM25 учитывает
оба этих фактора с помощью насыщения частоты терма и нормализации по
длине.

\textbf{BM25} суммирует вклады всех термов запроса:

\[\operatorname{Score}(Q,D)=\sum_i\operatorname{IDF}_{\mathrm{BM25}}(q_i)\cdot\frac{\operatorname{TF}(q_i,D)(k_1+1)}{\operatorname{TF}(q_i,D)+k_1\left(1-b+b\frac{|D|}{\operatorname{avgdl}}\right)}\]

Здесь \(q_i\) --- терм запроса, \(|D|\) --- длина документа, а
\(\operatorname{avgdl}\) --- средняя длина документов. Параметр \(k_1\)
управляет степенью насыщения частоты терма, а \(b\) --- силой
нормализации по длине; при значении 0 эта коррекция отключается. На
\crossreflink{fig:3-8}{рисунке 3-8} эти компоненты вычисляются по
отдельности. Сначала зафиксируем длину документа и увеличим число
вхождений терма запроса с одного до трёх, чтобы увидеть, насколько
возрастает оценка при каждом дополнительном вхождении. Затем зафиксируем
частоту терма и лишь удвоим длину документа. В нижней строке
сравнивается распространённость двух термов в коллекции документов.

\begin{figure}
\hypertarget{fig:3-8}{%
\centering
\includegraphics{images/revised/fig3-8.pdf}
\caption{Рисунок 3-8. Механизм оценки BM25}\label{fig:3-8}
}
\end{figure}

В первой строке рисунка второе вхождение увеличивает множитель частоты
терма примерно на 0,43, а третье --- лишь примерно на 0,24. Это и есть
насыщение частоты терма. Во второй строке один и тот же терм в обоих
случаях встречается лишь один раз, но в более длинном документе он менее
заметен, поэтому коррекция по длине снижает его множитель частоты. Чтобы
получить вклад терма в итоговую оценку, этот множитель необходимо также
умножить на IDF.

Остаётся определить ещё один компонент формулы: почему при одинаковых
трёх вхождениях редкий терм обычно полезнее распространённого? В этом
разделе для взвешивания термов используется следующая обратная частота
документов:

\[\operatorname{IDF}_{\mathrm{BM25}}(t)=\ln\frac{N-\operatorname{DF}(t)+0.5}{\operatorname{DF}(t)+0.5}\]

Числитель соответствует числу документов без данного терма, а
знаменатель --- числу документов с ним. Добавление 0,5 к обоим значениям
обеспечивает сглаживание, благодаря которому вычисление остаётся
возможным, когда частота документов равна 0 или \(N\). Редкие термы
получают больший вес; если терм содержится более чем в половине
документов, эта формула даёт отрицательное значение. В сопутствующей
реализации нижняя граница фактически используемого при оценке веса
установлена равной \(10^{-6}\), а исходное значение сохраняется для
проверки.

\begin{quote}
\textbf{Эксперимент 3-5 ★★: реализация поисковой системы BM25 с нуля}

В этом эксперименте последовательно реализуются токенизация, обработка
стоп-слов, инвертированный индекс и вычисление оценок. Обработка
стоп-слов позволяет отфильтровать часто встречающиеся служебные слова,
например «и» и «в». \textbf{Инвертированный индекс} хранит сведения о
том, в каких документах встречается каждый терм, и соответствующую
частоту терма. Он похож на предметный указатель в конце книги: по
запросу «TCP» можно найти страницы 45, 112 и 203, а затем прочитать
содержимое в этих местах.

Вручную рассчитаем один запрос «дистилляция модели». Предположим, корпус
содержит 10 документов, средняя длина равна 250 словам, \(k_1=1.5\), а
\(b=0.75\). Слово «модель» встречается в 3 документах, и его IDF
составляет примерно 0,762; слово «дистилляция» встречается в 2
документах, и его IDF составляет примерно 1,224. В таблице для каждого
терма перечислены совпавшие документы, частота терма и длина документа.
Подставив эти значения в формулу, можно вычислить вклад каждого терма в
оценку документа.

\begin{longtable}[]{@{}
  >{\raggedright\arraybackslash}p{(\columnwidth - 8\tabcolsep) * \real{0.2684}}
  >{\raggedright\arraybackslash}p{(\columnwidth - 8\tabcolsep) * \real{0.2017}}
  >{\raggedleft\arraybackslash}p{(\columnwidth - 8\tabcolsep) * \real{0.1686}}
  >{\raggedleft\arraybackslash}p{(\columnwidth - 8\tabcolsep) * \real{0.2142}}
  >{\raggedleft\arraybackslash}p{(\columnwidth - 8\tabcolsep) * \real{0.1471}}@{}}
\toprule\noalign{}
\begin{minipage}[b]{\linewidth}\raggedright
\hspace{0pt}Терм запроса
\end{minipage} & \begin{minipage}[b]{\linewidth}\raggedright
\hspace{0pt}Документ
\end{minipage} & \begin{minipage}[b]{\linewidth}\raggedleft
\hspace{0pt}Частота терма TF
\end{minipage} & \begin{minipage}[b]{\linewidth}\raggedleft
\hspace{0pt}Длина документа
\end{minipage} & \begin{minipage}[b]{\linewidth}\raggedleft
\hspace{0pt}Вклад BM25
\end{minipage} \\
\midrule\noalign{}
\endhead
\bottomrule\noalign{}
\endlastfoot
\hspace{0pt}модель & \hspace{0pt}doc\_1 & \hspace{0pt}5 &
\hspace{0pt}200 & \hspace{0pt}1.52 \\
\hspace{0pt}модель & \hspace{0pt}doc\_3 & \hspace{0pt}2 &
\hspace{0pt}500 & \hspace{0pt}0.82 \\
\hspace{0pt}модель & \hspace{0pt}doc\_7 & \hspace{0pt}8 &
\hspace{0pt}150 & \hspace{0pt}1.68 \\
\hspace{0pt}дистилляция & \hspace{0pt}doc\_1 & \hspace{0pt}3 &
\hspace{0pt}200 & \hspace{0pt}2.15 \\
\hspace{0pt}дистилляция & \hspace{0pt}doc\_5 & \hspace{0pt}1 &
\hspace{0pt}250 & \hspace{0pt}1.22 \\
\end{longtable}

Сумма вкладов двух термов для doc\_1 составляет примерно 3,67, поэтому
документы ранжируются следующим образом: \textbf{doc\_1 (3,67)
\textgreater{} doc\_7 (1,68) \textgreater{} doc\_5 (1,22) \textgreater{}
doc\_3 (0,82)}. В doc\_1 частота терма «дистилляция» ниже частоты терма
«модель», но из-за большей редкости его итоговый вклад выше; совпадение
сразу по двум термам дополнительно повышает общую оценку.

Теперь выполним поиск по слову \texttt{cat} в англоязычном корпусе
сопутствующего эксперимента. В двух коротких текстах, отмеченных как
релевантные, используются соответственно слова \texttt{kitten} и
\texttt{feline}, однако терм \texttt{cat} отсутствует в инвертированном
индексе, поэтому список кандидатов пуст. Как бы сильно мы ни увеличивали
описанный выше вес частоты терма, оценивать всё равно будет нечего. Если
повторить поиск со словами из этих двух текстов, можно увидеть, что
каждый из них попадает в список кандидатов.
\end{quote}

\hypertarget{ux433ux438ux431ux440ux438ux434ux43dux44bux439-ux43fux43eux438ux441ux43a-ux43eux431ux44aux435ux434ux438ux43dux435ux43dux438ux435-ux440ux435ux437ux443ux43bux44cux442ux430ux442ux43eux432-ux434ux432ux443ux445-ux43cux435ux442ux43eux434ux43eux432-ux441-ux43fux43eux441ux43bux435ux434ux443ux44eux449ux438ux43c-ux43fux435ux440ux435ux440ux430ux43dux436ux438ux440ux43eux432ux430ux43dux438ux435ux43c}{%
\subsection{Гибридный поиск: объединение результатов двух методов с
последующим
переранжированием}\label{ux433ux438ux431ux440ux438ux434ux43dux44bux439-ux43fux43eux438ux441ux43a-ux43eux431ux44aux435ux434ux438ux43dux435ux43dux438ux435-ux440ux435ux437ux443ux43bux44cux442ux430ux442ux43eux432-ux434ux432ux443ux445-ux43cux435ux442ux43eux434ux43eux432-ux441-ux43fux43eux441ux43bux435ux434ux443ux44eux449ux438ux43c-ux43fux435ux440ux435ux440ux430ux43dux436ux438ux440ux43eux432ux430ux43dux438ux435ux43c}}

Для некоторых вопросов трудно заранее перечислить все синонимы. В одной
из задач эксперимента 3-6 спрашивается, «как автоматически освобождать
неиспользуемую память в куче», однако в целевом фрагменте сказано:
«Автоматическое управление памятью избавляет разработчиков от
необходимости вручную освобождать объекты». В английской формулировке
запроса и фрагменте нет общих слов, поэтому поиск по ключевым словам не
нашёл ни одного результата, тогда как векторный поиск поставил целевой
фрагмент на первое место. Но для запроса по модели XR-7003 поиск по
ключевым словам сразу нашёл соответствующую модель, а векторный поиск
также вывел несколько похожих моделей.

\textbf{Гибридный поиск} сначала отдельно ищет релевантные фрагменты по
ключевым словам и векторам, а затем определяет, какие из них стоит
передать модели. Автор рекомендует начинать с этого подхода, если база
знаний одновременно содержит точные идентификаторы и описания на
естественном языке. Проследим по \crossreflink{fig:3-9}{рисунку 3-9} за
XR-7003: оба метода ставят его на первое место, поэтому при объединении
он получает вклад от обоих рейтингов. Переранжирование заменяет
результат на втором месте, но пользователь по-прежнему сначала видит
правильную модель.

\begin{figure}
\hypertarget{fig:3-9}{%
\centering
\includegraphics{images/revised/fig3-9.pdf}
\caption{Рисунок 3-9. Конвейер гибридного поиска и
переранжирования}\label{fig:3-9}
}
\end{figure}

Первый этап --- \textbf{параллельный поиск}: запрос отправляется в
плотный и разреженный поисковые движки, каждый из которых возвращает
набор потенциально релевантных документов. Оба метода могут использовать
один и тот же исходный текст либо предварительно обрабатывать его в
соответствии с требованиями своих моделей.

Второй этап --- \textbf{объединение результатов}: один и тот же документ
из результатов обоих методов объединяется, после чего формируется общий
рейтинг. Числовые шкалы косинусного сходства и BM25 различаются, поэтому
при непосредственном сложении результатов может начать доминировать
метод с более высокими числовыми значениями. Один из вариантов ---
сначала откалибровать или нормализовать оценки, а затем объединить их с
весами. Другой вариант --- \textbf{объединение по обратному рангу}
(Reciprocal Rank Fusion, RRF), при котором учитывается место документа в
каждом рейтинге:

\[\operatorname{RRF}(d)=\sum_{j:\,d\in R_j}\frac{1}{k+\operatorname{rank}_j(d)}\]

\(R_j\) --- результаты \(j\)-го метода, места нумеруются начиная с 1.
Если документ отсутствует в результатах одного из методов, вклад этого
метода равен нулю. Константу сглаживания \(k\) часто принимают равной
60: чем она больше, тем меньше разница в оценках между первыми местами.
RRF удобно объединяет разные способы оценки, но сохраняет только места в
рейтингах; различия в величине исходных оценок приходится учитывать
другими методами.

Рассмотрим пример с двумя методами поиска. Пользователь спрашивает:
«Когда поступят деньги после возврата?» Документ А занимает первое место
в поиске по ключевым словам, но отсутствует среди первых результатов
векторного поиска. Документ Б занимает второе место в обоих рейтингах.
При \(k=60\) документ А получает \(1/61\), то есть примерно 0,0164, а
документ Б --- \(1/62+1/62\), то есть примерно 0,0323. После объединения
документ Б оказывается выше документа А. Этот результат говорит лишь о
том, что документ Б найден обоими методами. Чтобы определить,
действительно ли он отвечает на вопрос о сроке поступления денег, нужно
проверить, о чём в нём говорится: «как подать заявку на возврат» или
«когда поступят деньги после возврата».

На третьем этапе можно добавить \textbf{нейросетевое переранжирование}.
Система выбирает документы, занимающие верхние места после объединения,
например первые 50, и модель переранжирования оценивает релевантность
каждого документа запросу. Переранжирование способно изменить порядок
этих 50 документов. Если действительно полезный документ вообще не попал
в выборку, необходимо вернуться к предыдущим этапам и улучшить поиск.

Для этого часто применяется \textbf{кросс-энкодер (cross-encoder)}:
запрос и документ-кандидат вместе передаются модели, две части
взаимодействуют внутри неё, а на выходе формируется оценка
релевантности. На этапе поиска \textbf{би-энкодер (bi-encoder)} кодирует
запрос и документы по отдельности: векторы документов можно вычислить
заранее, а при поступлении запроса выполнять только векторный поиск.
Первый подходит для точного ранжирования небольшого количества
документов, второй --- для первичного отбора из большого количества
документов. Это можно сравнить с отбором резюме и собеседованием:
сначала круг кандидатов сужают по уже известным характеристикам, а затем
подробно сравнивают их с учётом конкретных требований. Одна из моделей,
пригодных для этого этапа, --- BGE Reranker v2 M3\footnote{BAAI,
  описание модели BGE Reranker v2 M3, включая ввод пар «запрос ---
  документ» и оценку релевантности.}.

Если релевантные материалы уже попадают в пул кандидатов, но часто
располагаются ниже нерелевантных, автор рекомендует прежде всего оценить
переранжирование с помощью кросс-энкодера. Он совместно обрабатывает
вопрос и документ, поэтому может непосредственно сопоставлять условия,
ограничения и подробности документа с требованиями вопроса. В таблице
3-3 перечислены три распространённые метрики.

Таблица 3-3. Три метрики качества поиска

\begin{longtable}[]{@{}
  >{\raggedright\arraybackslash}p{(\columnwidth - 2\tabcolsep) * \real{0.4165}}
  >{\raggedright\arraybackslash}p{(\columnwidth - 2\tabcolsep) * \real{0.5835}}@{}}
\toprule\noalign{}
\begin{minipage}[b]{\linewidth}\raggedright
\hspace{0pt}Метрика
\end{minipage} & \begin{minipage}[b]{\linewidth}\raggedright
\hspace{0pt}Расчёт и значение
\end{minipage} \\
\midrule\noalign{}
\endhead
\bottomrule\noalign{}
\endlastfoot
\hspace{0pt}\href{mailto:recall@k}{\nolinkurl{recall@k}} (полнота) &
\hspace{0pt}Количество релевантных документов среди первых k
результатов, делённое на общее количество релевантных документов для
данного запроса; показывает, какая доля релевантного содержимого
найдена. Для тестового набора обычно дополнительно вычисляется среднее
значение по всем запросам \\
\hspace{0pt}MRR (средний обратный ранг) & \hspace{0pt}Для каждого
запроса берётся величина, обратная месту первого релевантного документа,
после чего вычисляется среднее значение: первое место даёт 1, десятое
--- 0,1, отсутствие результата --- 0 \\
\hspace{0pt}\href{mailto:nDCG@k}{\nolinkurl{nDCG@k}} (нормализованная
дисконтированная кумулятивная полезность) & \hspace{0pt}Результатам
присваиваются оценки в соответствии со степенью релевантности, позиции
ниже в рейтинге дисконтируются, а полученное значение делится на оценку
идеального ранжирования; показывает общее качество ранжирования первых k
результатов \\
\end{longtable}

Также можно учитывать \textbf{долю попаданий}: для какой доли запросов
среди первых k результатов найден хотя бы один релевантный документ?
Например, для ответа на один вопрос нужны три документа, но среди первых
двадцати результатов найдены только два из них. Для этого запроса уже
есть попадание, поэтому он не учитывается как случай полного отсутствия
результатов, однако его \href{mailto:recall@20}{\nolinkurl{recall@20}}
составляет лишь \(2/3\): ещё \(1/3\) материалов не найдена. Если для
каждого вопроса существует только один целевой документ, вопрос с
найденным документом получает 1 балл, а с пропущенным --- 0. При
усреднении по одному и тому же набору результатов доля попаданий будет
равна средней полноте.

\begin{quote}
\textbf{Эксперимент 3-6 ★★: сравнение разреженного и плотного поиска,
объединения и переранжирования}

В этом эксперименте на одном и том же наборе из 12 размеченных запросов
сравниваются пять конфигураций: разреженный поиск, плотный поиск,
объединение посредством RRF, взвешенное объединение и переранжирование
после объединения. Запросы охватывают семантическое сходство, точные
названия, разные языки и технические коды. В двух задачах на английском
и китайском языках спрашивается, как растения преобразуют солнечный свет
в пищу. Для плотного поиска используется модель Qwen3-Embedding-0.6B, а
для переранжирования --- кросс-энкодер MS MARCO MiniLM. Каждый метод
поиска возвращает не более 10 документов, а модель переранжирования
обрабатывает не более 10 документов, занимающих верхние места в рейтинге
RRF. Полнота и nDCG вычисляются только по первым трём документам
итогового рейтинга, тогда как для MRR первый релевантный документ ищется
во всём возвращённом списке.

Сначала откройте ранжирование по отдельным задачам в журнале выполнения
эксперимента 3-6 и сравните результаты поиска для двух вопросов о
фотосинтезе. Для вопроса на английском языке целевой документ
по-прежнему занимает первое место. Для вопроса на китайском языке
сначала возвращаются документы о языке программирования и вулканах, а о
фотосинтезе рассказывает лишь третий документ. Для каждого вопроса
релевантным размечен только один документ, и в обоих случаях он найден
среди первых трёх результатов, поэтому
\href{mailto:recall@3}{\nolinkurl{recall@3}} равен 1. Однако для вопроса
на китайском языке первый релевантный документ занимает третье место, а
его обратный ранг составляет лишь \(1/3\). При принятом в этом
эксперименте способе подсчёта полезность третьего места равна
\(1/\log_2(3+1)=0.5\), а первого --- 1, поэтому
\href{mailto:nDCG@3}{\nolinkurl{nDCG@3}} для вопроса на китайском языке
составляет 0,5.

Теперь рассмотрим все 12 задач. После переранжирования, помимо вопроса о
фотосинтезе на китайском языке, целевой документ ещё в одной задаче ---
об освобождении памяти в куче --- опустился с первого места на второе. В
остальных десяти задачах целевые документы сохранили первое место,
поэтому MRR после переранжирования равен
\((10+1/2+1/3)/12\approx0.903\). Для сравнения, MRR плотного поиска и
RRF равен 1. Их \href{mailto:nDCG@3}{\nolinkurl{nDCG@3}} также равен 1,
а после переранжирования снижается примерно до 0,928. Если смотреть
только на среднюю полноту, эти изменения мест останутся незаметными:
\href{mailto:recall@3}{\nolinkurl{recall@3}} разреженного поиска равен
0,833, а во всех остальных четырёх конфигурациях --- 1.

Предположим, теперь для этого набора задач нужно выбрать одну
конфигурацию. Что дала дополнительная задержка, затраченная на
переранжирование? На этом CPU RRF от начала поиска до выдачи рейтинга
занимает в среднем около 194 миллисекунд на вопрос, а с подключённым
переранжированием --- около 237 миллисекунд. Дополнительные примерно 43
миллисекунды не позволили найти новые целевые документы, но отодвинули
правильные ответы в двух задачах.
\end{quote}

\hypertarget{ux43eux440ux433ux430ux43dux438ux437ux430ux446ux438ux44f-ux438-ux43fux43eux438ux441ux43a-ux437ux43dux430ux43dux438ux439}{%
\section{Организация и поиск
знаний}\label{ux43eux440ux433ux430ux43dux438ux437ux430ux446ux438ux44f-ux438-ux43fux43eux438ux441ux43a-ux437ux43dux430ux43dux438ux439}}

Способ организации знаний определяет, что содержит каждая запись, а
механизм поиска --- что именно будет извлечено из этих записей. Абзацы в
технических руководствах, юридических документах и научных работах
связаны иерархией разделов, зависимостями между понятиями и
перекрёстными ссылками между документами. Сохранение этих структур
позволяет при обработке запроса одновременно обращаться к конкретным
доказательствам и обобщениям более высокого уровня.

В этом разделе последовательно рассматриваются несколько способов
организации: RAPTOR и GraphRAG создают иерархические индексы или индексы
связей; OpenViking организует контекстные ресурсы с помощью каталогов;
инкрементальные обновления и периодическая систематизация поддерживают
актуальность знаний; агентный RAG позволяет AI-агенту продолжать поиск
на основе уже полученных результатов; контекстно-зависимый поиск
дополняет фрагменты фоновой информацией; наконец, из наборов данных
извлекаются знания, пригодные для повторного использования. Эти подходы
воздействуют соответственно на процессы представления, индексирования,
обновления и обработки запросов, а также могут совместно применяться для
организации пользовательской памяти.

Сначала рассмотрим два примера, чтобы понять разрыв между информацией,
необходимой для ответа на запрос, и исходными записями.

\textbf{Пример 1: подсчёт чёрных и белых кошек.} В предыдущей главе
программа пересчитала кошек во всех 100 клетках и получила 90 чёрных и
10 белых кошек. На этот раз механизму поиска разрешено вернуть только
первые 20 записей. Предположим, что 18 из них содержат достоверные
сведения о цвете: в 15 записях указаны чёрные кошки, а в 3 --- белые.
Полученные 83,3\% (15/18) --- это доля чёрных кошек среди этих
достоверных записей, тогда как доля чёрных кошек во всём магазине
составляет 90\% (90/100). Механизм поиска отбирает записи по
релевантности, а не формирует репрезентативную выборку по всему
магазину.

Если долю по всему магазину спрашивают часто, можно заранее подсчитать
кошек во всех клетках, сохранить количество чёрных и белых кошек, а
также указать охват подсчёта и время обновления. Если кошка в клетке
сменилась, соответствующую подробную запись обновляют и пересчитывают
итоговые значения. Если кто-то спрашивает, какая кошка находится в
клетке № 37, по-прежнему нужно получить именно эту подробную запись.

\textbf{Пример 2: может ли медсестра получить скидку, предоставляемую
врачу?} Рассмотрим учебную ситуацию, построенную вокруг консультаций по
скидкам Xfinity: в обращениях ветеран John успешно прошёл проверку, врач
Sarah получила скидку, а учителю Mike сообщили, что он не соответствует
требованиям. Каждое обращение фиксирует один результат обработки. Если
медсестра спросит, может ли она получить скидку, поиск только по похожим
обращениям столкнётся с тремя проблемами.

\begin{itemize}
\tightlist
\item
  \textbf{Близкие профессии не обязательно дают право на одну и ту же
  льготу.} Понятия «медсестра» и «врач» семантически близки, поэтому
  запись Sarah может оказаться среди первых результатов; однако
  применимость скидки также зависит от того, кому условия акции
  разрешают подавать заявку и какие подтверждающие документы необходимо
  предъявить.
\item
  \textbf{Отдельные обращения не содержат полного перечня условий
  применимости.} Успешная проверка конкретного человека непосредственно
  свидетельствует лишь о результате для этого человека при
  существовавших на тот момент условиях; увеличение числа извлекаемых
  записей не гарантирует, что найдутся ранее не встречавшиеся категории
  или исключения.
\item
  \textbf{Область применимости материалов неясна.} Обращения могут
  относиться к разным акциям или периодам. Сначала необходимо
  определить, какие записи применимы и каких нормативных оснований не
  хватает.
\end{itemize}

При индексировании можно связать обращения, относящиеся к одной акции,
затем изучить соответствующие официальные условия и оформить сведения о
целевых категориях, исключениях, сроке действия и источнике положений в
виде карточки правил. Например, если условия в этой учебной ситуации
явно распространяются на действующих военнослужащих и ветеранов, а также
на лицензированных медицинских работников, но исключают учителей, каждое
из этих условий следует сохранить отдельно; принадлежность медсестёр к
категории медицинских работников также должна подтверждаться условиями.
Обращения продолжают храниться, чтобы объяснять фактический процесс
оформления и выявлять расхождения в применении правил.

Эти два примера отражают разные информационные потребности: для подсчёта
кошек требуется статистика по всем клеткам магазина, а для определения
права на скидку --- полные условия и исключения. Для таких запросов
автор предлагает перенести подготовку статистики и правил на этап
создания базы знаний. Простое увеличение числа извлекаемых фрагментов
всё равно может не дать модели полного охвата статистики или положений,
определяющих, будет ли одобрена конкретная заявка. Только если заранее
создавать статистические сводки, карточки правил и связи между
источниками с учётом задачи, механизм поиска сможет за один раз получить
достаточно информации для принятия решения и при необходимости
обратиться к исходным материалам.

\hypertarget{ux441ux442ux440ux443ux43aux442ux443ux440ux438ux440ux43eux432ux430ux43dux43dux44bux439-ux438ux43dux434ux435ux43aux441-ux43fux43eux441ux442ux440ux43eux435ux43dux438ux435-ux438ux435ux440ux430ux440ux445ux438ux438-ux438-ux441ux432ux44fux437ux435ux439}{%
\subsection{Структурированный индекс: построение иерархии и
связей}\label{ux441ux442ux440ux443ux43aux442ux443ux440ux438ux440ux43eux432ux430ux43dux43dux44bux439-ux438ux43dux434ux435ux43aux441-ux43fux43eux441ux442ux440ux43eux435ux43dux438ux435-ux438ux435ux440ux430ux440ux445ux438ux438-ux438-ux441ux432ux44fux437ux435ux439}}

При чтении длинного отчёта ответ на вопрос «Какую графу заполнять в
заявке на ремонт?» зачастую можно найти в одном фрагменте исходного
текста, тогда как вопрос «Почему несколько типов неисправностей
возникают снова и снова?» требует анализа нескольких разделов. Чтобы
поддерживать оба типа запросов, при создании базы можно генерировать
сводки, строить иерархию либо извлекать сущности и связи.
Рассматриваемый ниже RAPTOR формирует иерархию с помощью рекурсивного
обобщения, а GraphRAG организует материал посредством связей между
сущностями и сводок сообществ.

\textbf{RAPTOR} (Recursive Abstractive Processing for Tree-Organized
Retrieval) начинает с текстовых блоков: каждый блок преобразуется в
вектор, затем блоки объединяются в группы по степени сходства, после
чего для каждой группы создаётся сводка. Если считать сводки новой
партией текстовых блоков и повторять эти шаги, получится дерево,
показанное на \crossreflink{fig:3-10}{рисунке 3-10}. Если двигаться по
стрелкам вверх от текстовых блоков 3, 4 и 5, видно, что вместе они
образуют сводку B; ещё выше сводки трёх групп A, B и C объединяются в
общую сводку.\footnote{Sarthi et al.~RAPTOR: Recursive Abstractive
  Processing for Tree-Organized Retrieval, 2024.}

\begin{figure}
\hypertarget{fig:3-10}{%
\centering
\includegraphics{images/revised/fig3-10.pdf}
\caption{Рисунок 3-10. Иерархический древовидный индекс
RAPTOR}\label{fig:3-10}
}
\end{figure}

Например, утверждения «SSE2 поддерживает целочисленные SIMD-операции в
128-битных регистрах XMM» и «В SSE4.2 добавлены инструкции сравнения
строк» можно включить в сводку «Развитие наборов инструкций x86 SIMD».
Вопросы о деталях могут сопоставляться с блоками нижнего уровня, а
обзорные вопросы --- со сводками верхних уровней. При поиске можно как
последовательно углубляться по иерархии, так и искать среди узлов всех
уровней одновременно, выбирая наиболее релевантную группу.

\textbf{GraphRAG} организует знания с помощью сущностей, связей и сводок
сообществ. Одна связь записывается в виде триплета «субъект ---
отношение --- объект», например (Китай, столица, Пекин) или (Чжан Сань,
работает в, Tencent). Если две записи упоминают одну сущность, их можно
связать в графе. Оба врача на \crossreflink{fig:3-11}{рисунке 3-11}
носят фамилию Чжан. По какой цепочке связей должна пройти система, если
пользователь спрашивает: «Каков адрес больницы, в которой работает мой
стоматолог?»

\begin{figure}
\hypertarget{fig:3-11}{%
\centering
\includegraphics{images/revised/fig3-11.pdf}
\caption{Рисунок 3-11. Граф знаний GraphRAG «сущность ---
отношение»}\label{fig:3-11}
}
\end{figure}

\begin{enumerate}
\def\labelenumi{\arabic{enumi}.}
\tightlist
\item
  \textbf{Различайте одноимённые сущности при построении графа.} Один
  врач Чжан лечит зубы пользователя в стоматологической клинике
  «Жэньай», а другой лечит заболевание сердца у отца пользователя в
  больнице «Хуашань». При извлечении сущностей и объединении записей
  система может на основании этих признаков назначить двум врачам разные
  идентификаторы. Аналогичную функцию выполняют рассмотренные ранее поля
  \texttt{person} и \texttt{relationship}. Определение того, означает ли
  \texttt{bank} в предложении берег или банк, называется разрешением
  лексической неоднозначности; определение того, о каком именно враче
  Чжане говорится в записи, --- разрешением неоднозначности сущностей.
\item
  \textbf{При выполнении запроса ищите по цепочке связей.} Если начать с
  пользователя и последовательно пройти по цепочке «пользователь → врач
  → больница → адрес», можно объединить сведения из нескольких записей в
  единый ответ.
\end{enumerate}

В процессе индексирования Microsoft GraphRAG сначала извлекает сущности
и связи из текстовых единиц, а затем с помощью обнаружения сообществ
выявляет тесно связанные группы сущностей и создаёт иерархию сообществ
со сводками. При запросе о конкретной сущности локальный поиск
объединяет связанные с ней сведения и исходный текст; при запросе о
тематике всего корпуса глобальный поиск организует ответ с помощью
сводок сообществ.\footnote{Microsoft GraphRAG, обзор индексирования и
  выполнения запросов, локальный и глобальный поиск.}

Представление связей также должно учитывать условия и время. Например,
высказывание «Если на следующей неделе дождь не прекратится, я отменю
поездку на море и вместо этого пойду в музей» содержит условие, период
действия и альтернативное действие. Если извлечь только (пользователь,
планирует, поездка на море) и (пользователь, имеет альтернативный план,
поход в музей), условие переключения будет потеряно. Можно сохранить
исходное предложение целиком, а ограничение «если на следующей неделе
будет дождь» представить с помощью узла условия, узла события или
структурированных полей, связав отношение с исходным текстом.

Для пользовательской памяти автор рекомендует \textbf{многоуровневое
сочетание естественного языка и структурированного индекса}: сохранять
полное высказывание на естественном языке, записывать объект, время и
источник в метаданных, а для содержимого, которое требуется часто
связывать, создавать графовый индекс.

\begin{quote}
\textbf{Эксперимент 3-7 ★★★: сравнение иерархических сводок и индекса
связей}

В этом эксперименте используются 14 страниц руководства по архитектуре
Intel, посвящённых SSE, AVX и управлению состоянием, и создаются две
учебные реализации. В ветви иерархических сводок сначала генерируется
сводка для каждой страницы, затем сводки объединяются в три группы,
после чего создаются групповые сводки и корневая сводка. В ветви индекса
связей извлекаются сущности и отношения, после чего формируются сводки
сообществ. Обе группы используют одинаковые материалы и восемь вопросов
для сравнения поиска деталей понятий и рассуждений о связях.

При запросе «Объясните среду программирования SSE» иерархические сводки
позволяют одновременно найти обзор SSE и сведения о регистрах. В этой
реализации поиск выполняется сразу среди узлов всех уровней, после чего
выбираются пять наиболее релевантных. Ветвь связей расширяет поиск от
релевантных сущностей к их соседям и организует контекст с помощью
сводок сообществ. Цепочка связей «SSE → регистры XMM → операции с
плавающей точкой → ADDPS» позволяет соединить понятие с конкретной
инструкцией.

Более сложный вопрос: «Если процессор заявляет о поддержке AVX, означает
ли это, что программа уже может использовать AVX?» Для ответа необходимо
связать возможности процессора с управлением состоянием операционной
системы: проверить флаги AVX и OSXSAVE в CPUID, затем прочитать XCR0 и
убедиться, что состояния XMM и YMM включены. Для этого вопроса
подсказками о конечных точках графового поиска служат CPUID и YMM.
Программа находит соединяющий их кратчайший путь через XSAVE,
CR4.OSXSAVE и XCR0, а затем передаёт этот путь модели для подготовки
ответа.

Среднее покрытие страниц-источников по восьми вопросам в обеих группах
составляет около 95,8\%.
\end{quote}

Автор рекомендует сначала качественно реализовать разбиение на блоки и
гибридный поиск, а к созданию структурированного индекса переходить при
регулярной потребности в обобщении материалов из нескольких документов и
поиске по связям. Чтобы найти конкретный фрагмент политики возврата
средств, обычно достаточно разбиения на блоки и гибридного поиска; для
сравнения SSE и AVX или перехода от общей архитектуры к конкретной
инструкции в большей степени требуются межфрагментные сводки и индекс
связей.

\hypertarget{ux444ux430ux439ux43bux43eux432ux430ux44f-ux43fux430ux440ux430ux434ux438ux433ux43cux430-ux43eux440ux433ux430ux43dux438ux437ux430ux446ux438ux44f-ux437ux43dux430ux43dux438ux439-ux441-ux43fux43eux43cux43eux449ux44cux44e-ux441ux442ux440ux443ux43aux442ux443ux440ux44b-ux43aux430ux442ux430ux43bux43eux433ux43eux432}{%
\subsection{Файловая парадигма: организация знаний с помощью структуры
каталогов}\label{ux444ux430ux439ux43bux43eux432ux430ux44f-ux43fux430ux440ux430ux434ux438ux433ux43cux430-ux43eux440ux433ux430ux43dux438ux437ux430ux446ux438ux44f-ux437ux43dux430ux43dux438ux439-ux441-ux43fux43eux43cux43eux449ux44cux44e-ux441ux442ux440ux443ux43aux442ux443ux440ux44b-ux43aux430ux442ux430ux43bux43eux433ux43eux432}}

\textbf{OpenViking} организует память, ресурсы и навыки посредством
виртуальной файловой системы. Каждый элемент доступен по URI
\texttt{viking://}: AI-агент может просматривать каталоги, читать
обзоры, а затем обращаться к необходимому содержимому\footnote{OpenViking,
  Viking URI и архитектура системы.}. По назначению информации каталоги
можно разделить на три группы: внешние ресурсы, данные пользователя, а
также используемые агентом навыки и опыт. Идентификаторы пользователя и
агента дополнительно ограничивают их пространства имён.

\begin{longtable}[]{@{}
  >{\raggedright\arraybackslash}p{(\columnwidth - 2\tabcolsep) * \real{0.3940}}
  >{\raggedright\arraybackslash}p{(\columnwidth - 2\tabcolsep) * \real{0.6060}}@{}}
\toprule\noalign{}
\begin{minipage}[b]{\linewidth}\raggedright
\hspace{0pt}Содержимое
\end{minipage} & \begin{minipage}[b]{\linewidth}\raggedright
\hspace{0pt}Типичное назначение
\end{minipage} \\
\midrule\noalign{}
\endhead
\bottomrule\noalign{}
\endlastfoot
\hspace{0pt}Resources & \hspace{0pt}Внешние материалы: техническая
документация, репозитории кода, веб-\allowbreak{}\hspace{0pt}страницы и
т. д. \\
\hspace{0pt}User memories & \hspace{0pt}Предпочтения, привычки и опыт
конкретного пользователя \\
\hspace{0pt}Skills и Agent memories & \hspace{0pt}Повторно используемые
процессы, записи о задачах и опыт, требующий проверки \\
\end{longtable}

URI обеспечивают единый способ адресации, а физическим хранилищем и
индексом управляет сервис. Интерфейс каталогов позволяет агенту
последовательно просматривать материалы по темам и ограничивать
семантический поиск соответствующим поддеревом. Доступ в стиле файловой
системы, хранилище содержимого и векторный индекс могут совместно
образовывать единый сервис.

OpenViking разделяет содержимое на три уровня: \textbf{L0, L1 и L2}. L0
--- краткая сводка для быстрой оценки релевантности каталога. L1 ---
обзор каталога, описывающий основное содержимое и точки входа для
навигации. L2 --- основной текст, требующий подробного изучения. Сводка
и обзор уровня каталога хранятся соответственно в \texttt{.abstract.md}
и \texttt{.overview.md}. Бюджеты на генерацию текста для них составляют
соответственно 256 и 4000 символов\footnote{OpenViking, Context Layers
  (L0/L1/L2). L0 и L1 --- представления, генерируемые на уровне
  каталога; L2 может содержать исходный файл или основной текст после
  разбора.}.

При выполнении запроса можно сначала прочитать сводку и определить
нужную область, затем изучить обзор и выбрать файл, а в конце загрузить
исходный текст. Это перекликается с прогрессивным раскрытием Skills из
второй главы: облегчённые каталоги помогают находить материалы, а полное
содержимое загружается по мере необходимости для конкретной задачи. Если
требуются точные числовые значения, условия или исходные формулировки,
необходимо дополнительно получить соответствующий основной текст.

Для знаний, которые должны совместно поддерживать человек и агент, автор
предпочитает использовать \textbf{обычный текст Markdown и управление
версиями}. Человек может непосредственно читать, редактировать и
исправлять знания, а агент --- работать через тот же файловый интерфейс,
предлагать обновления в рабочей ветви и публиковать их после проверки.
Система управления версиями сохраняет различия каждого изменения и
позволяет откатывать ошибочные обновления. В конце сеанса
пользовательские предпочтения добавляются в память соответствующего
пользователя, а история операций --- в записи о задачах. После оценки
результатов, обобщения между траекториями и последующей проверки записи
о задачах могут стать опытом, рассматриваемым в девятой главе.

При файловой организации знаний одновременно необходимо
\textbf{создавать ссылки и индексы для связанных элементов}. Набор
файлов без взаимных ссылок позволяет заново обнаруживать связи только
посредством полнотекстового сканирования или семантического поиска. Чем
больше файлов, тем сложнее агенту продолжать исследование, отталкиваясь
от уже найденного содержимого. Ссылки обеспечивают явные пути поиска:
например, политика возврата средств может ссылаться на положения об
исключениях, сводка договора --- на запись о его подписании, а элемент о
человеке --- на связанные события. Начальные страницы описывают основные
темы, индексные страницы обеспечивают навигацию, а перекрёстные ссылки
между элементами поддерживают горизонтальный поиск. Такой подход
напоминает организацию энциклопедических статей и позволяет сравнительно
простым способом сохранять часть отношений.

При создании базы знаний полезно указать, какие записи можно применять
вместе, а какие инструкции исключают друг друга. Отвечая на просьбу
отменить заказ, сотрудник поддержки должен соблюдать и правила общения,
и правила конфиденциальности: оба документа уместны в контексте. А если
у склада две инструкции по отмене отправки, в индексе нужны условия
выбора. Пока посылка на складе, следует отменить распоряжение на
отгрузку; если посылка уже у перевозчика --- запросить её перехват.
Получив заявку «Отмените заказ» без других подробностей, агент должен
сначала проверить статус посылки и только затем загрузить подходящую
инструкцию.

Поддержку ссылок следует встроить в процесс обновления: перед
добавлением нового элемента сначала искать существующие записи и
определять, нужно ли объединить данные или создать новый элемент; после
записи добавлять связанные ссылки и обновлять индекс соответствующего
каталога. Ссылка может привести читателя к нужному положению, но
применимо ли оно к конкретному заказу, зависит также от даты его
вступления в силу.

\hypertarget{ux43aux430ux43a-ux434ux43eux43bux436ux43dux44b-ux43eux431ux43dux43eux432ux43bux44fux442ux44cux441ux44f-ux437ux43dux430ux43dux438ux44f}{%
\subsection{Как должны обновляться
знания}\label{ux43aux430ux43a-ux434ux43eux43bux436ux43dux44b-ux43eux431ux43dux43eux432ux43bux44fux442ux44cux441ux44f-ux437ux43dux430ux43dux438ux44f}}

Сегодня политика изменилась, а завтра пользователь может спросить: «Что
теперь делать?» --- или: «Почему в прошлом месяце мой заказ обработали
именно так?» Если просто заменить старые положения новыми, во втором
случае исчезнут основания, действовавшие на тот момент.
\textbf{Инкрементальные обновления} позволяют своевременно записывать
новые свидетельства, а \textbf{периодическая реорганизация} --- выявлять
в более широком масштабе дубликаты, конфликты и структурные изменения. В
обоих процессах необходимо сохранять версии и ссылки на свидетельства, а
также выполнять проверку перед публикацией, чтобы выбирать период,
применимый к конкретному вопросу.

\hypertarget{ux438ux43dux43aux440ux435ux43cux435ux43dux442ux430ux43bux44cux43dux44bux435-ux43eux431ux43dux43eux432ux43bux435ux43dux438ux44f-ux43fux43eux43bux44cux437ux43eux432ux430ux442ux435ux43bux44cux441ux43aux43eux439-ux43fux430ux43cux44fux442ux438-ux438-ux431ux430ux437ux44b-ux437ux43dux430ux43dux438ux439}{%
\subsubsection{Инкрементальные обновления пользовательской памяти и базы
знаний}\label{ux438ux43dux43aux440ux435ux43cux435ux43dux442ux430ux43bux44cux43dux44bux435-ux43eux431ux43dux43eux432ux43bux435ux43dux438ux44f-ux43fux43eux43bux44cux437ux43eux432ux430ux442ux435ux43bux44cux441ux43aux43eux439-ux43fux430ux43cux44fux442ux438-ux438-ux431ux430ux437ux44b-ux437ux43dux430ux43dux438ux439}}

Автор рекомендует \textbf{рассматривать базу знаний как репозиторий
кода, а каждое изменение знаний --- как Pull Request (PR)}. Каждое
обновление должно иметь чётко определённую область изменений и
сопровождаться подтверждающими свидетельствами, доступными проверяющему.
В этот процесс можно включить документы Markdown, файлы правил и
исполняемое состояние User as Code. История версий помогает
сопровождающим просматривать изменения, выяснять их причины и
восстанавливать данные после ошибочной публикации.

Можно использовать представленную в первой главе схему
\textbf{предлагающий --- проверяющий}, организовав обновление следующим
образом:

\begin{enumerate}
\def\labelenumi{\arabic{enumi}.}
\tightlist
\item
  \textbf{Предлагающий формирует изменение.} Proposer Agent читает новые
  свидетельства, выполняет поиск по существующим знаниям, находит
  элементы, которые нужно добавить, исправить или признать устаревшими,
  и предлагает в рабочей ветви изменение с чётко очерченной областью.
  Одновременно он поддерживает связанные ссылки, индекс каталога,
  временные сведения и ссылки на свидетельства.
\item
  \textbf{Проверяющий выполняет независимую проверку.} Reviewer Agent
  сопоставляет содержимое до и после изменения с исходными
  свидетельствами и проверяет, имеют ли новые утверждения источник,
  полностью ли указаны условия применимости, согласуются ли формулировки
  в разных файлах и не осталось ли полезной информации в удалённых
  фрагментах. Обнаружив проблему, он указывает конкретный элемент и
  основание, а также объясняет, что следует дополнить, восстановить или
  исправить.
\item
  \textbf{Стороны итеративно дорабатывают изменение и принимают
  решение.} Предлагающий вносит правки с учётом замечаний, после чего
  проверяющий выполняет повторную проверку. Одобренное изменение
  передаётся в процесс объединения. Если по достижении предельного числа
  итераций или бюджета затрат разногласия сохраняются, изменение
  остаётся в рабочей ветви и передаётся человеку на проверку.
\item
  \textbf{Изменение проверяется и публикуется.} Автоматические проверки
  охватывают формат, ссылки, метаданные и метки разрешений; исполняемые
  знания также должны пройти проверку типов и соответствующие тесты.
  После объединения из зафиксированной версии знаний перестраиваются
  затронутые фрагменты, сводки и векторные индексы, а затем публикуется
  версия сервиса.
\end{enumerate}

Этот процесс охватывает три уровня данных. \textbf{Уровень исходных
свидетельств} хранит диалоги, траектории выполнения, бизнес-документы и
результаты инструментов. \textbf{Уровень знаний} хранит прошедший
обработку и допускающий дальнейшее редактирование текст или код.
\textbf{Сервисный уровень} хранит индексы, сформированные из конкретной
версии знаний. Журнал изменений связывает идентификаторы свидетельств,
версию знаний, замечания проверки и окончательное решение, позволяя
сопровождающим выяснить, почему определённое знание было записано и
когда оно вступило в силу. Исходные свидетельства архивируются согласно
политике хранения, а индексы можно повторно сформировать из версии
знаний. Рассмотренная ранее память с добавлением только новых данных
подходит для хранения новых и старых фактов вместе с их источниками;
описанный здесь процесс проверки преобразует эти свидетельства в
пригодные для использования версии знаний. Если при запросе
обнаруживается ещё не разрешённый конфликт, можно вернуться к
свидетельствам и вынести решение с учётом условий задачи.

В такой архитектуре \textbf{и предлагающий, и проверяющий должны быть
агентами, способными вызывать инструменты и продолжать поиск}.
Предлагающему может потребоваться изучить множество записей, чтобы найти
затронутые элементы. Проверяющий также может переходить по ссылкам,
продолжать поиск, сравнивать версии или запускать проверки. Инструменты
поиска по файлам и свидетельствам, просмотра различий и запуска тестов
позволяют выполнять эти дополнительные действия, поэтому такую работу
может взять на себя Coding Agent с соответствующими возможностями.

Область поиска обоих агентов должна охватывать знания и исходные
свидетельства, относящиеся к текущему изменению, с учётом ограничений
прав пользователя, арендатора и задачи. Траектория проверки, ссылки на
результаты инструментов и замечания архивируются вместе, чтобы
впоследствии можно было установить, что именно проверяющий фактически
сверил.

При выборе моделей автор рекомендует отдавать предпочтение сочетаниям
сопоставимой мощности из разных семейств, например Claude и GPT либо
DeepSeek и Kimi. Различия в обучающих данных и привычках рассуждения
могут приводить к различным суждениям. Проверяющий должен самостоятельно
искать свидетельства и сверять с ними утверждения и условия в изменении.
Разрешения должны обеспечиваться исполняющей системой: предлагающий
записывает данные в рабочую ветвь, проверяющий читает свидетельства и
отправляет заключение, а разрешения на обновление основного репозитория
и онлайн-индекса принадлежат сервисам объединения и публикации.

\hypertarget{ux43fux435ux440ux438ux43eux434ux438ux447ux435ux441ux43aux430ux44f-ux440ux435ux43eux440ux433ux430ux43dux438ux437ux430ux446ux438ux44f-ux43fux43eux43bux44cux437ux43eux432ux430ux442ux435ux43bux44cux441ux43aux43eux439-ux43fux430ux43cux44fux442ux438-ux438-ux431ux430ux437ux44b-ux437ux43dux430ux43dux438ux439}{%
\subsubsection{Периодическая реорганизация пользовательской памяти и
базы
знаний}\label{ux43fux435ux440ux438ux43eux434ux438ux447ux435ux441ux43aux430ux44f-ux440ux435ux43eux440ux433ux430ux43dux438ux437ux430ux446ux438ux44f-ux43fux43eux43bux44cux437ux43eux432ux430ux442ux435ux43bux44cux441ux43aux43eux439-ux43fux430ux43cux44fux442ux438-ux438-ux431ux430ux437ux44b-ux437ux43dux430ux43dux438ux439}}

После множества локальных обновлений один и тот же факт может оказаться
разбросан по нескольким файлам, а старые и новые формулировки ---
существовать одновременно. Периодическая реорганизация выявляет такие
накопившиеся проблемы, последовательно обрабатывая знания по каталогам,
времени или темам. Её рабочая схема совпадает с периодическим
пересмотром, рассмотренным в девятой главе: на переднем плане
накапливаются новые свидетельства, а в фоновом режиме они
централизованно систематизируются, образуя записи для последующих задач.

Реорганизация включает три задачи.

\begin{enumerate}
\def\labelenumi{\arabic{enumi}.}
\tightlist
\item
  \textbf{Удаление дубликатов, объединение и изменение структуры.}
  Выявляются дублирующиеся, устаревшие или чрезмерно раздробленные
  элементы, связанное содержимое объединяется, слишком большие файлы
  разделяются, а ссылки, начальные и индексные страницы обновляются. При
  удалении из сервисного представления старой формулировки фиксируются
  заменивший её элемент и исходное основание.
\item
  \textbf{Повторная проверка по исходным материалам.} Если при каждой
  реорганизации переписывать только сводку, созданную на предыдущей
  итерации, ранние пропуски и неверные интерпретации будут передаваться
  дальше. Поэтому необходимо сопоставлять знания с диалогами,
  траекториями выполнения, бизнес-документами и результатами
  инструментов, проверяя, не пропущены ли в сводке ключевые факты,
  отрицания, временные сведения или условия и не представлено ли
  предположение как определённый вывод. Крупную базу знаний можно
  обрабатывать частями по каталогам, времени или темам, фиксируя в
  списке покрытия уже проверенные, ожидающие проверки и спорные области,
  пока не будет проверена вся база.
\item
  \textbf{Разрешение конфликтов и уточнение областей применимости.}
  Сравниваются источники противоречащих друг другу записей, а также их
  время, объекты, регионы, задачи и предварительные условия. Если две
  записи применимы к разным сценариям, следует явно указать область
  каждой. Если новое свидетельство однозначно заменяет старое правило,
  необходимо сохранить связь между версиями. Если оснований
  недостаточно, фиксируются вопрос, ожидающий подтверждения, и
  недостающие свидетельства.
\end{enumerate}

Результат реорганизации также отправляется на проверку как изменение.
Крупную перестройку можно разделить на несколько PR по каталогам или
темам, использующих общий план реорганизации и список покрытия. После
одобрения связанных изменений перестраиваются производные индексы и
повторно выполняются типичные сценарии поиска и вопросов-ответов, чтобы
проверить, остались ли ранее доступные знания доступными для поиска.
Реорганизацию можно запускать еженедельно, ежемесячно или с другой
периодичностью, а также при превышении порогов числа новых элементов,
конфликтов или изменения качества поиска.

\textbf{Вывод из эксплуатации недействительного содержимого.} Фрагменты
и сводки должны сохранять версию, время вступления в силу и время утраты
силы. После замены старой политики текущие бизнес-запросы
отфильтровывают недействительные версии, а исторические запросы получают
правила, действовавшие в указанный период. Это позволяет одновременно
использовать актуальные знания и объяснять результаты обработки в
прошлом.

\textbf{Разрешения и изоляция арендаторов.} Фильтрация по правам доступа
должна выполняться на уровне поиска: поисковый сервис сначала определяет
набор документов, доступных вызывающей стороне, а затем ищет в нём
релевантные материалы. Как только конфиденциальные данные другого
пользователя попадут в контекст, модель может раскрыть их в ответе или
при вызове инструмента. Исходные тексты, сводки, рёбра отношений и
кэшированные результаты должны наследовать соответствующую область
доступа, чтобы производное содержимое не раскрывало сведения других
пользователей или арендаторов. Сервис может использовать отдельные
индексы либо общий индекс с принудительной фильтрацией; главное, чтобы
ограничения разрешений применялись к каждому запросу.

\hypertarget{ux430ux433ux435ux43dux442ux43dux44bux439-rag-ux443ux43fux440ux430ux432ux43bux435ux43dux438ux435-ux43fux440ux43eux446ux435ux441ux441ux43eux43c-ux43fux43eux438ux441ux43aux430-ux441-ux43fux43eux43cux43eux449ux44cux44e-ai-ux430ux433ux435ux43dux442ux430}{%
\subsection{Агентный RAG: управление процессом поиска с помощью
AI-агента}\label{ux430ux433ux435ux43dux442ux43dux44bux439-rag-ux443ux43fux440ux430ux432ux43bux435ux43dux438ux435-ux43fux440ux43eux446ux435ux441ux441ux43eux43c-ux43fux43eux438ux441ux43aux430-ux441-ux43fux43eux43cux43eux449ux44cux44e-ai-ux430ux433ux435ux43dux442ux430}}

Сначала следует определить, достаточно ли одного поиска. Например, ответ
на вопрос «Какой номер телефона у службы послепродажного обслуживания?»
зачастую можно найти в одном актуальном руководстве по обслуживанию,
поэтому автор рекомендует придерживаться простого сценария с однократным
поиском и последующей генерацией. Но если пользователь спрашивает:
«Можно ли обменять купленную мной модель по действовавшим в прошлом
месяце условиям гарантии?», первый поиск может найти только серию
продукта, чего недостаточно для определения действовавших тогда условий
обмена. Если предоставить поиск как инструмент, агент сможет
использовать сведения из предыдущего результата для составления
следующего запроса. Подход, при котором агент решает, нужен ли
дополнительный поиск и что именно искать дальше, называется
\textbf{агентным RAG (Agentic RAG)}. Сравните два пути на
\crossreflink{fig:3-12}{рисунке 3-12}: результаты первого поиска
сведений о продукте полностью совпадают. Какие дополнительные сведения
использует правый путь, чтобы найти условия, действовавшие в тот момент?

\begin{figure}
\hypertarget{fig:3-12}{%
\centering
\includegraphics{images/revised/fig3-12.pdf}
\caption{Рисунок 3-12. Сравнение агентного и неагентного
RAG}\label{fig:3-12}
}
\end{figure}

Рассмотрим ReAct. Сначала модель анализирует вопрос, выбирает поисковый
запрос и вызывает \texttt{knowledge\_base\_search}. Получив наблюдение
от инструмента, модель определяет, достаточно ли доказательств. Если не
хватает условий или связанных объектов, она переформулирует запрос и
продолжает поиск. Если материалов уже достаточно, модель формирует ответ
с учётом источников. Harness ведёт историю, отслеживает вызовы
инструментов и бюджет, а при достижении лимита завершает поиск или
переходит к уточнению.

Теперь рассмотрим \crossreflink{fig:3-13}{рисунок 3-13}. Даже если
заменить в бэкенде базы знаний векторный индекс на граф связей, модель
по-прежнему сможет выполнять поиск через тот же интерфейс инструмента.
Но если в найденном документе окажется какая-либо инструкция, может ли
модель её выполнить?

\begin{figure}
\hypertarget{fig:3-13}{%
\centering
\includegraphics{images/revised/fig3-13.pdf}
\caption{Рисунок 3-13. Архитектура системы агентного
RAG}\label{fig:3-13}
}
\end{figure}

\textbf{Граница доверия к найденному содержимому.} Внешние документы
могут содержать текст, пытающийся изменить поведение модели, либо в них
ещё до добавления в базу знаний могли быть внесены ложные факты. Первый
случай относится к косвенной инъекции в промпт, второй --- к отравлению
базы знаний. Для защиты необходимо одновременно контролировать
происхождение содержимого и права на выполнение действий: результаты
поиска должны включать источник, версию и место цитирования и
использоваться как доказательства для ответа, а такие действия, как
перевод средств, удаление данных или отправка внешних сообщений, должны
проверяться Harness и средой выполнения на соответствие разрешениям
пользователя. Проверка материалов перед добавлением в базу, отслеживание
их происхождения и проверка фактов дополнительно ограничивают
распространение ошибочных знаний. Здесь совместно применяются
рассмотренная в \crossreflink{chap:2}{главе 2} маркировка источников и
ограничения инструментов из \crossreflink{chap:4}{главы 4}.

\begin{quote}
\textbf{Эксперимент 3-8 ★★: сравнение однократного и многоэтапного
поиска}

В этом эксперименте на семи вопросах сравниваются два режима с
использованием одного и того же корпуса китайских юридических
материалов, одного поискового механизма BM25 и одной модели. В базовом
режиме поиск выполняется один раз по исходному вопросу, а в режиме
агента модель может корректировать запрос и выполнять до четырёх
поисков.

Для запросов вроде «Как регулируется необходимая оборона?» основные
положения можно найти напрямую. Более сложный пример: «Как назначается
наказание за причинение тяжкого вреда здоровью по неосторожности в
состоянии опьянения при наличии судимости за кражу?». В нём одновременно
затрагиваются текущее деяние, состояние опьянения и предшествующая
судимость, поэтому при поиске необходимо отдельно найти соответствующие
материалы, а затем проверить условия их применимости.

При фактическом запуске последовательно выполнялись поисковые запросы
«состояние опьянения причинение тяжкого вреда здоровью по неосторожности
судимость за кражу назначение наказания», «состояние опьянения
преступление назначение наказания», «влияние судимости за кражу на
назначение наказания уголовное право» и «влияние судимости на назначение
наказания уголовное право». Затем модель предложила запрос «признание
рецидива уголовное право», но к этому моменту поисковый бюджет был
исчерпан, и Harness вернул сообщение об исчерпании бюджета.

Для четырёх сложных вопросов, выделенных в эксперименте, охват целевых
правовых норм вырос примерно с 8 до 92\%, среднее количество поисков
увеличилось с одного до 3,5, а среднее время этапа формирования ответа
--- примерно с 8 до 66 секунд. Для трёх простых вопросов охват
доказательств составил 100\%, и обе группы выполняли только один поиск.
\end{quote}

Тот же поисковый цикл можно применять к истории пользователя. Далее
используется трёхуровневый набор для оценки из эксперимента 3-1, чтобы
проследить, как агент извлекает факты из записей разных сеансов,
различает объекты и обрабатывает обновления.

\begin{quote}
\textbf{Эксперимент 3-9 ★★: чтение пользовательской памяти с помощью
агентного поиска}

В этом эксперименте полная история диалогов используется как база
знаний: она разбивается на окна, после чего для неё создаётся индекс, а
агенту предоставляется инструмент \texttt{search\_user\_memory}. Размер
окна можно изменять. В данном запуске каждый фрагмент содержит восемь
раундов диалога, соседние фрагменты перекрываются на два раунда, а также
сохраняются идентификатор сеанса, время и бизнес-метаданные. Допускается
не более трёх поисков, каждый из которых возвращает четыре фрагмента.

Для базового припоминания можно искать по конкретным полям, например по
запросу «Какой номер у моего чекового счёта?». При работе с несколькими
сеансами также необходимо проверять полноту набора объектов: у
пользователя есть и Honda Accord, и Tesla Model 3, поэтому в ответе на
вопрос о запланированном обслуживании автомобилей система должна
различать обе машины и состояние записи каждой из них.

На этой основе можно построить следующий маршрут поиска: сначала найти
сведения об автомобилях и обслуживании; обнаружив упоминание второй
машины, отдельно запросить планы по Tesla; затем сообщить о существующей
записи для Honda и уточнить, требуется ли пользователю также записать
Tesla. Однако в фактическом запуске был выполнен только запрос
\texttt{scheduled\ car\ services}, поэтому ответ группы с обычным
поиском сосредоточился на Honda и упустил Tesla.

Ещё одна категория трудностей связана с последовательными изменениями. В
примере с переводом средств жена создаёт распоряжение, затем муж меняет
сумму и дату, после чего жена снова вносит корректировку. Получив
записи, необходимо восстановить последовательность изменений с учётом
времени, объекта и состояния подтверждения, чтобы определить
окончательно вступившее в силу распоряжение. В текущей схеме разбиения
на фрагменты уже сохраняется информация о времени, однако поиск всё
равно может пропустить критически важное изменение, поэтому фрагменты
можно дополнительно снабдить контекстом событий.

Из 60 примеров в 18 поиск выполнялся один раз, в 6 --- дважды, а в 36
был полностью использован лимит из трёх поисков. Группа с обычным
поиском по фрагментам прошла 47 примеров: 19 на базовое припоминание, 15
на поиск по нескольким сеансам и 13 на проактивное обслуживание. В
следующем разделе на тех же примерах сравниваются префикс контекста и
структурированная память, а также анализируется, какие сведения
добавляет каждый из этих подходов.
\end{quote}

\hypertarget{ux43aux43eux43dux442ux435ux43aux441ux442ux43dux43e-ux437ux430ux432ux438ux441ux438ux43cux44bux439-ux43fux43eux438ux441ux43a-ux434ux43eux431ux430ux432ux43bux435ux43dux438ux435-ux444ux43eux43dux43eux432ux43eux439-ux438ux43dux444ux43eux440ux43cux430ux446ux438ux438-ux43a-ux444ux440ux430ux433ux43cux435ux43dux442ux430ux43c}{%
\subsection{Контекстно-зависимый поиск: добавление фоновой информации к
фрагментам}\label{ux43aux43eux43dux442ux435ux43aux441ux442ux43dux43e-ux437ux430ux432ux438ux441ux438ux43cux44bux439-ux43fux43eux438ux441ux43a-ux434ux43eux431ux430ux432ux43bux435ux43dux438ux435-ux444ux43eux43dux43eux432ux43eux439-ux438ux43dux444ux43eux440ux43cux430ux446ux438ux438-ux43a-ux444ux440ux430ux433ux43cux435ux43dux442ux430ux43c}}

Вернёмся к финансовому отчёту. Чтобы найти результаты компании ACME за
второй квартал 2025 года, в запросе указывают название компании и год,
однако во фрагменте со словами «выручка компании выросла на 3\%» как раз
нет этих двух элементов. Помимо увеличения размера блоков, можно
сохранить этот короткий фрагмент, добавив перед ним недостающую фоновую
информацию. При \textbf{контекстно-зависимом поиске (Contextual
Retrieval)} во время построения индекса для фрагмента создаётся краткий
контекстный префикс, после чего префикс и исходный текст вместе
используются для поиска\footnote{Anthropic. \emph{Introducing Contextual
  Retrieval}. Сентябрь 2024 года.}. Этот процесс показан на
\crossreflink{fig:3-14}{рисунке 3-14}.

\begin{figure}
\hypertarget{fig:3-14}{%
\centering
\includegraphics{images/revised/fig3-14.pdf}
\caption{Рисунок 3-14. Контекстно-зависимый поиск}\label{fig:3-14}
}
\end{figure}

Например, в префиксе можно указать: «Этот фрагмент взят из раздела с
ключевыми показателями отчёта компании ACME за второй квартал 2025
года». При генерации модели одновременно передаются исходный текст и
целевой блок, чтобы она добавила подтверждённую текстом информацию о
референтах, времени и разделе. Благодаря этому BM25 получает доступные
для сопоставления термы, такие как название компании и квартал, а
плотные эмбеддинги могут представлять фрагмент с учётом этого контекста.

Этот подход можно сочетать с контекстно-зависимым сжатием из второй
главы: на этапе индексирования к блокам знаний добавляется контекст,
повышающий вероятность обнаружения релевантной информации; во время
выполнения уже извлечённая история упорядочивается с учётом текущей
задачи, чтобы контролировать длину входных данных. Первый подход
направлен на доступность фрагментов для поиска, второй --- на информацию
и бюджет, необходимые для принятия текущего решения.

\begin{quote}
\textbf{Эксперимент 3-10 ★★: сравнение исходных блоков с блоками,
дополненными контекстом}

В этом эксперименте для одного набора материалов создаются два индекса:
в одном используются исходные фрагменты, а в другом к ним добавляются
сгенерированные моделью префиксы. Запросы и разметка целевых фрагментов
остаются одинаковыми; отдельно сравниваются BM25, плотный и гибридный
поиск.

Если воспользоваться предыдущим примером с финансовым отчётом и запросом
«Недавний рост выручки компании ACME», можно проверить, содержит ли
найденный фрагмент одновременно целевую компанию, период и данные о
росте. Одно лишь словосочетание «рост выручки» может привести к
совпадению с другой компанией; после добавления префикса название
компании и время могут учитываться при ранжировании.

Для 15 запросов этого эксперимента после добавления префиксов
\href{mailto:recall@1}{\nolinkurl{recall@1}} у BM25 вырос с 0,60 до
0,80, а у плотного поиска --- приблизительно с 0,867 до 0,933;
\href{mailto:recall@3}{\nolinkurl{recall@3}} у них сохранился на уровне
0,867 и 1,0 соответственно. \href{mailto:Recall@5}{\nolinkurl{Recall@5}}
у гибридного поиска вырос с 14/15 до 15/15. Для какого запроса удалось
дополнительно найти целевой фрагмент? В запросе о полномочиях
Постоянного комитета Всекитайского собрания народных представителей
целевой фрагмент перечислял полномочия, но не содержал подлежащего.
После добавления префикса BM25 поднял его с 16-го на 12-е место, однако
он по-прежнему не вошёл в первую пятёрку; гибридный поиск поднял его с
6-го на 5-е место, благодаря чему фрагмент оказался среди первых пяти.
Если обратиться к исходному тексту, первый фрагмент содержит полномочия
с пункта 1 по пункт 4, а пятый продолжает перечисление с пункта 4 по
пункт 22. При оценке проверяется лишь попадание размеченного пятого
фрагмента в результаты, однако для полного перечисления полномочий
необходимо объединить оба фрагмента.

Генерация префиксов увеличивает число вызовов на этапе индексирования.
При создании контекста для нескольких блоков одного документа можно
повторно использовать кэш префикса документа. В отчёте Anthropic за 2024
год при размере блока 800 токенов, документа --- 8000 токенов и заданной
длине промпта расчётная стоимость однократной обработки миллиона токенов
документов составляла около 1,02 доллара.

Согласно этому отчёту, сочетание контекстных эмбеддингов с BM25 снизило
долю неудач поиска в top-20 с 5,7\% до 2,9\%, то есть приблизительно на
49\% относительно исходного значения; после добавления переранжирования
она снизилась до 1,9\%, то есть приблизительно на 67\%. В отчёте доля
неудач определяется как \(1-\mathrm{recall@20}\).
\end{quote}

Для относительно стабильных баз материалов, к которым требуется
обращаться многократно, автор рекомендует в первую очередь дополнять
фрагменты контекстом при построении индекса: однократная обработка
сможет обслуживать множество последующих запросов. Если при разбиении
теряются название компании, субъект события или время, простое
увеличение количества результатов всё равно может вернуть набор
фрагментов с неясными референтами. Если добавить эти ориентиры при
генерации префиксов, последующий поиск сможет использовать их для
нахождения материалов.

\begin{quote}
\textbf{Эксперимент 3-11 ★★★: расширение памяти с помощью контекстных
префиксов и структурированных карточек}

Реплику «Хорошо, тогда бронируй этот вариант» в диалоге необходимо
связать с ранее упомянутым товаром или маршрутом. Если в исходном тексте
обсуждается билет в одну сторону из Шанхая в Сиэтл стоимостью 500
долларов, префикс может добавить этот контекст, а также сохранить сеанс
и время подтверждения.

В этом эксперименте к обычному разбиению на блоки добавляются две
конфигурации: группа с контекстными префиксами и группа с двухуровневой
памятью, в которую также загружаются Advanced JSON Cards. Сначала модель
формирует последовательность поисковых запросов на основе обычного
индекса, после чего тот же набор запросов применяется к двум другим
индексам. Двухуровневая группа использует те же результаты расширенного
поиска и дополнительно получает структурированные карточки.

В примере с переводом средств префиксы соответственно указывают, что
Patricia инициировала перевод, James изменил договорённость, а затем
Patricia снова внесла изменения, помогая установить последовательность
событий. Согласно исходной записи, окончательное состояние --- 95 000
долларов и 1 октября; прежняя договорённость была заменена. Все три
группы в этом эксперименте правильно определили это состояние.

Двухуровневая структура также может поддерживать проактивное выявление
связей. В карточках сохраняются такие ключевые факты, как «поездка в
Токио» и «срок действия паспорта», а подробности диалога извлекаются по
мере необходимости. Например, если поездка запланирована на 15--22
января, а срок действия паспорта истекает 18 февраля, система может
сначала обнаружить близость дат, затем найти записи о продлении паспорта
и требованиях для поездки и наконец напомнить пользователю проверить
незавершённые приготовления. В этом примере с путешествием ответы всех
трёх групп прошли проверку; различия между карточками и поиском
проявились главным образом в рассмотренных ниже примерах с автомобилями
и несколькими сеансами.

Из 60 примеров обычный поиск, контекстные префиксы и двухуровневая
память прошли соответственно 47, 52 и 57. В группе с префиксами число
пройденных примеров проактивного обслуживания выросло с 13 до 18;
двухуровневая группа прошла все 20 примеров на уровне поиска по
нескольким сеансам и также 18 примеров на уровне проактивного
обслуживания. В примере с автомобилями карточка добавила сведения о
второй машине, благодаря чему ответ охватил ранее пропущенный объект.
Рассмотрим и одно ухудшение: когда пользователь спросил, какое покрытие
распространяется на ущерб недвижимости от воды, обычная группа упомянула
собственное и сдаваемое в аренду жильё, а группа с префиксами пропустила
сдаваемое жильё. Сопоставление по каждому заданию показывает, что группа
с префиксами дала на 8 правильных ответов больше обычной группы, но при
этом ошиблась в 3 заданиях, на которые та ответила верно.
\end{quote}

\textbf{Когда требуется проактивное обслуживание, автор рекомендует
двухуровневую архитектуру, в которой ключевые факты постоянно находятся
в контексте, а подробные записи извлекаются по мере необходимости.} Если
полагаться только на поиск, модели зачастую сначала нужно знать, какой
вопрос задать, и лишь затем она сможет найти соответствующие записи;
паспорт, автомобиль или родственная связь, ещё не попавшие в её поле
зрения, поэтому легко останутся незамеченными. Небольшой объём ключевых
фактов, заранее помещённый в контекст, помогает модели обнаруживать
связи, формировать запросы и затем извлекать исходные подробности. При
выборе постоянно доступной информации необходимо учитывать отношения
между людьми, текущие ограничения и незавершённые дела; в результатах
поиска следует сохранять источник, время и основание для изменений.

\hypertarget{ux438ux437ux432ux43bux435ux447ux435ux43dux438ux435-ux437ux43dux430ux43dux438ux439-ux438ux437-ux43dux430ux431ux43eux440ux430-ux434ux430ux43dux43dux44bux445-ux43eux442-ux43fux440ux438ux43cux435ux440ux43eux432-ux43a-ux441ux442ux440ux443ux43aux442ux443ux440ux438ux440ux43eux432ux430ux43dux43dux43eux43cux443-ux43fux440ux435ux434ux441ux442ux430ux432ux43bux435ux43dux438ux44e}{%
\subsection{Извлечение знаний из набора данных: от примеров к
структурированному
представлению}\label{ux438ux437ux432ux43bux435ux447ux435ux43dux438ux435-ux437ux43dux430ux43dux438ux439-ux438ux437-ux43dux430ux431ux43eux440ux430-ux434ux430ux43dux43dux44bux445-ux43eux442-ux43fux440ux438ux43cux435ux440ux43eux432-ux43a-ux441ux442ux440ux443ux43aux442ux443ux440ux438ux440ux43eux432ux430ux43dux43dux43eux43cux443-ux43fux440ux435ux434ux441ux442ux430ux432ux43bux435ux43dux438ux44e}}

Знания также можно получать путём систематизации большого количества
примеров. В судебных решениях, обращениях в службу поддержки и записях о
ремонте множество полезных сведений рассеяно по конкретным описаниям.
Когда задача требует выявить закономерности в большом массиве примеров,
сравнить результаты или определить, какие уточняющие вопросы следует
задать, автор рекомендует сначала извлечь унифицированные поля, а затем
провести статистический анализ и сформировать прототипы. Извлечение лишь
нескольких похожих примеров не позволяет увидеть распределение во всём
наборе данных. После унификации полей можно собрать статистику по всему
массиву записей и получить общий обзор для последующего поиска отдельных
случаев.

Этот процесс можно разделить на два этапа. \textbf{Извлечение знаний}
преобразует исходные описания в структурированные записи, где явно
определены смысл полей, их значения и источники. Затем \textbf{анализ
паттернов} с помощью статистики и кластеризации выявляет похожие случаи
и их общие признаки. На \crossreflink{fig:3-15}{рисунке 3-15} показан
процесс от извлечения данных из примеров и кодирования признаков до
поиска прототипов.

\begin{figure}
\hypertarget{fig:3-15}{%
\centering
\includegraphics{images/revised/fig3-15.pdf}
\caption{Рисунок 3-15. Конвейер извлечения структурированных
знаний}\label{fig:3-15}
}
\end{figure}

\begin{quote}
\textbf{Эксперимент 3-12 ★★★: извлечение факторов и прототипов из
судебных дел}

В этом эксперименте используются материалы уголовных дел из CAIL2018. Из
них выбраны примеры трёх категорий --- кража, умышленное причинение
вреда здоровью и мошенничество, --- на основе которых построен учебный
процесс «выявление факторов --- извлечение полей --- кластеризация ---
объяснение прототипов»\footnote{CAIL2018 --- официальный проект набора
  данных. Для сопутствующего эксперимента были отобраны примеры из
  общедоступных материалов, а моделирование и оценка проводились с
  разделением внутри выборки.}. Для моделирования взято по 120 записей
каждой категории, а ещё по 20 оставлено отдельно. Итого получено 360
примеров для моделирования и 60 отложенных примеров.

\textbf{Выявление и систематизация полей.} Сначала модель пакетами
читает примеры для моделирования и предлагает потенциально полезные
факторы, а затем объединяет их в модульную схему. В общей части хранятся
поля, которые необходимо учитывать для нескольких видов преступлений,
например явка с повинной и возмещение ущерба. В расширениях сохраняются
сведения, относящиеся к конкретному виду преступления, например сумма по
делу или характер причинённого вреда. В результате было сформировано 16
общих полей и три группы расширений, содержащие соответственно 12, 12 и
11 полей. Для каждого поля также необходимо определить допустимые
значения и состояние отсутствующих данных, чтобы обеспечить
единообразное извлечение.

\textbf{Кодирование полей в векторы.} Для категориальных полей можно
использовать унитарное кодирование. Например, категории «кража»,
«грабёж» и «мошенничество» кодируются соответственно как
\texttt{{[}1,0,0{]}}, \texttt{{[}0,1,0{]}} и \texttt{{[}0,0,1{]}},
причём каждая позиция обозначает только принадлежность к соответствующей
категории. Значения 1, 2 и 3 привнесли бы в некоторые алгоритмы порядок
и расстояние между категориями. Для таких полей, как «была ли явка с
повинной» и «был ли возмещён ущерб», значениями 1 и 0 можно обозначить
два подтверждённых состояния. Случаи, не упомянутые в исходном тексте,
следует обрабатывать отдельно, чтобы отсутствие данных автоматически не
считалось отрицательным ответом.

Для неотрицательных числовых значений с большим диапазоном, например
денежных сумм, в этой реализации сначала применяется \(\ln(1+x)\) для
сжатия масштаба, после чего выполняется стандартизация. Если поле
допускает одновременное наличие нескольких категорий --- например, одна
запись содержит несколько принятых мер, --- для каждой категории можно
выделить отдельный индикатор, сформировав многозначное бинарное
кодирование.

\textbf{Построение прототипов дел.} Закодированные признаки
стандартизируются, после чего дела кластеризуются отдельно по каждому
виду преступления. В этой реализации сравниваются варианты с числом
кластеров от двух до пяти, а группировка выбирается по коэффициенту
силуэта. В результате сформировано 11 прототипов. Для каждого прототипа
сохраняются количество примеров, основные признаки, а также медиана,
квартили и диапазон сроков лишения свободы в примерах. Например, на
основе причины конфликта, способа совершения преступления и последствий
причинённого вреда можно сравнить такие сочетания признаков, как
начавшаяся после ссоры драка без оружия и заранее спланированное
нападение с оружием. Конкретные группы совместно определяются
результатами извлечения и кластеризацией.

\textbf{Объяснение признаков и результатов сопоставления.} На основе
различий между стандартизированными центрами кластеров программа
вычисляет веса признаков при различении прототипов. При сопоставлении
новый случай отображается в то же пространство признаков, после чего
извлекаются статистическая сводка по ближайшей группе и основания для
совпадения, чтобы модель могла сформировать объяснение.

Из 60 отложенных примеров были выбраны 12, для которых сгенерировали
объяснения, а затем передали их независимой модели для оценки.
Рассмотрим один пример кражи. Соответствующий ему прототип содержал
только одну запись для моделирования со сроком лишения свободы четыре
месяца, поэтому минимальное, медианное и максимальное значения
составляли четыре месяца. Генератор достоверно сообщил это число и также
указал, что для сравнения доступна лишь одна запись. За верность чисел
данным прототипа оценщик поставил максимальные 4 балла, хотя фактическая
метка отложенного примера составляла восемь месяцев.

Если поиск по прототипам используется в диалоге, сначала можно уточнить,
какие отсутствующие факты способны повлиять на сопоставление, а затем
повторить поиск с учётом ответа.
\end{quote}

При переходе от примеров к прототипам база знаний получает
дополнительный слой доступного для поиска статистического представления.
Исходный текст служит доказательством, поля обеспечивают возможность
вычислений, а прототипы обобщают общие признаки похожих примеров. Связи
между всеми тремя уровнями необходимо сохранять, чтобы при обнаружении
ошибки извлечения или смещения выборки можно было вернуться к исходным
данным и внести исправления.

\hypertarget{ux43fux435ux440ux435ux434ux43dux438ux439-ux43aux440ux430ux439-ux438ux441ux441ux43bux435ux434ux43eux432ux430ux43dux438ux439-ux43cux443ux43bux44cux442ux438ux43cux43eux434ux430ux43bux44cux43dux430ux44f-ux438-ux43fux430ux440ux430ux43cux435ux442ux440ux438ux447ux435ux441ux43aux430ux44f-ux43fux430ux43cux44fux442ux44c}{%
\subsection{Передний край исследований: мультимодальная и
параметрическая
память}\label{ux43fux435ux440ux435ux434ux43dux438ux439-ux43aux440ux430ux439-ux438ux441ux441ux43bux435ux434ux43eux432ux430ux43dux438ux439-ux43cux443ux43bux44cux442ux438ux43cux43eux434ux430ux43bux44cux43dux430ux44f-ux438-ux43fux430ux440ux430ux43cux435ux442ux440ux438ux447ux435ux441ux43aux430ux44f-ux43fux430ux43cux44fux442ux44c}}

Представьте, что пользователь указывает на ранее загруженную фотографию
и говорит: «Нужен именно этот цвет». Если система сохранила лишь слова
«синяя куртка», в следующий раз она может не различить два оттенка
синего. При преобразовании в текст также теряются детали тембра голоса и
пространственных связей на диаграммах. Чтобы при будущих запросах эти
сведения можно было восстановить, можно сохранять исходные материалы,
представления, пригодные для использования моделью, либо исследовать
постоянную память внутри самой модели.

\textbf{Первый подход: сочетание исходных материалов с текстовым
индексом.} Например, в авторизованной пользователем задаче по
систематизации фотоальбома можно вырезать и сохранять отдельные области,
такие как лица людей, ссылаться на изображения в записях Markdown, а
также фиксировать время съёмки, обстановку и предоставленные
пользователем метки людей. При поиске сначала по описанию находят
несколько потенциально релевантных фотографий, а затем загружают
оригиналы или вырезанные области и сравнивают детали внешности.
Фрагменты аудио также можно связать со временем, источником и
расшифровкой, чтобы при необходимости воспроизводить исходную запись.
Текст служит точкой входа, а исходные материалы сохраняют детали,
которые трудно выразить словами.

\textbf{Второй подход: хранение эмбеддингов и их загрузка по мере
необходимости.} Визуальный или аудиоэнкодер может преобразовать
материалы в векторы и сохранить их в индексе для поиска по сходству.
Например, после кодирования отдельных областей изображения или
фрагментов аудио их можно связать с соответствующими записями памяти.
Если модель поддерживает обученный проекционный слой и интерфейс памяти,
в контексте также можно зарезервировать набор слотов, загрузить в них
эти непрерывные векторы и предоставить модели доступ к ним через
механизм внимания. Для этого необходимо согласовать представления
энкодера с внутренними представлениями языковой модели.

В предварительном исследовании автора и Noah Shi эта задача разделена на
два этапа. На групповой фотографии изображены два человека, и
пользователь спрашивает: «Тот, кто стоит у окна, --- это человек,
которого мы видели в прошлый раз?» Сначала по указанию «у окна» нужно
определить, кого имеет в виду пользователь, а затем сравнить изображение
этого человека с ранее сохранённым образцом. Если на первом этапе выбран
не тот человек, даже безошибочное распознавание на втором этапе даст
ответ на другой вопрос. В статье такие факты, как имена, которые можно
точно записать, также оставлены в текстовой памяти, чтобы для каждого из
двух типов информации использовался подходящий способ
доступа.\footnote{Боцзе Ли, Noah Shi, Parametric Multimodal User Memory;
  сопутствующий проект \texttt{19PINE-AI/multimodal-user-memory}
  содержит реализацию и эксперименты PerceptMem. Вопрос о групповой
  фотографии приведён в качестве учебного примера.}

\textbf{Третий подход: постоянная память внутри модели.} Обучение
отдельной LoRA для каждого пользователя --- один из способов записать
персонализированные факты в модель. Прямой вопрос «Где живёт
пользователь?» и планирование поездки с учётом его адреса предъявляют
разные требования к использованию памяти: в первом случае нужно извлечь
факт, а во втором --- ещё и задействовать его в последующих
рассуждениях. Если обучать модель только воспроизводить факты, она может
помнить ответ, но не использовать его самостоятельно в сложных задачах.
Поэтому обучающие примеры должны включать как прямые вопросы, так и
задачи, решение которых зависит от этих фактов.

В исследовании автора \textbf{User as Engram} текстовые факты
записываются в локальную таблицу памяти модели Engram, а общий адаптер
обучается использовать эти факты в рассуждениях\footnote{Li, Bojie, User
  as Engram: Internalizing Per-User Memory as Local Parametric Edits,
  2026.}. Поиск N-gram с хешированными ключами отвечает за обнаружение
записей таблицы, а механизм гейтинга передаёт извлечённое содержимое в
модель; для каждого пользователя сохраняются локальные записи, тогда как
компонент рассуждений может быть общим. В экспериментах с текстовыми
фактами в статье сравниваются прямое воспроизведение и косвенное
рассуждение, а также показано, как раздельно управлять содержимым памяти
каждого пользователя и общей способностью к рассуждению.

Для распространения этого подхода на мультимодальные данные также
необходимо разработать механизмы адресации, записи и чтения
представлений изображений и аудио.

\hypertarget{ux438ux442ux43eux433ux438-ux433ux43bux430ux432ux44b-2}{%
\section{Итоги
главы}\label{ux438ux442ux43eux433ux438-ux433ux43bux430ux432ux44b-2}}

Вернёмся к поручению забронировать билет из начала главы. «Я предпочитаю
место у окна» --- это предпочтение пользователя, а допустимое количество
багажа нужно уточнять по правилам авиакомпании, действующим на дату
поездки. Если сохранить оба сведения одним предложением, легко упустить
важнейшее различие: к кому относится первое и с какого момента действует
второе. При сохранении нужно фиксировать источник, время и условия
применимости --- только тогда после извлечения можно обоснованно решить,
подходит ли информация в данном случае. Отдельный факт удобно хранить в
виде краткой заметки; если требуется описать отношения между людьми или
сохранить контекст происхождения сведений, можно использовать карточку
или фрагмент исходного текста; даты и количества также можно хранить в
типизированных полях, чтобы программа могла выполнять вычисления.

Если помощник ошибся в норме провоза багажа, сначала следует проверить,
загружены ли действующие правила в базу знаний, затем выяснить, в какой
фрагмент они попали, нашёл ли их поисковый компонент и какое место они
заняли после переранжирования, а в завершение --- какие материалы
получила модель. Если исходного текста много, структурированный обзор
можно оставить в контексте, а подробности извлекать по мере
необходимости. Резюме глав помогает искать по темам, а индекс связей ---
находить взаимосвязанную информацию; если во фрагменте потерялась
отсылка, к нему можно добавить сведения об источнике и контексте; если
после одного поиска доказательств всё ещё недостаточно, агент может
изменить запрос и продолжить поиск.

После публикации авиакомпанией новых правил сначала можно выполнить
небольшое обновление, сохранив даты вступления в силу и источники как
новых, так и прежних правил; накопившиеся дубликаты, противоречия и
изменения структуры затем следует периодически упорядочивать. Инициатор
предлагает изменение, проверяющий сверяет доказательства, а после
одобрения перестраивается рабочий индекс. История версий позволяет
повторно воспроизвести вопросы, на которые раньше были даны неверные
ответы, а также установить, на каких материалах основывался прежний
ответ. В девятой главе будет отдельно рассмотрено, как накапливать опыт
на основе результатов выполнения.

Сведения о том, что можно зарегистрировать два места багажа, отвечают
лишь на вопрос: «Допускают ли это правила?» Чтобы добавить багаж в заказ
пользователя, помощнику также нужно знать, какой инструмент можно
вызвать, какие параметры ему требуются и какое подтверждение означает,
что операция выполнена. С этого начинается следующая глава.

\begin{questionbox}

\hypertarget{ux432ux43eux43fux440ux43eux441ux44b-ux434ux43bux44f-ux440ux430ux437ux43cux44bux448ux43bux435ux43dux438ux44f-2}{%
\section*{Вопросы для
размышления}\label{ux432ux43eux43fux440ux43eux441ux44b-ux434ux43bux44f-ux440ux430ux437ux43cux44bux448ux43bux435ux43dux438ux44f-2}}
\addcontentsline{toc}{section}{Вопросы для размышления}

\begin{enumerate}
\def\labelenumi{\arabic{enumi}.}
\tightlist
\item
  ★★ Пользователь говорит: «1 мая я переехал из A в B, с 10 по 20 мая
  живу в гостинице C, а в этот период счета присылайте по электронной
  почте». Как 15 мая помощник должен ответить на вопросы о домашнем
  адресе, временном месте проживания и способе доставки счетов?
  Разработайте правила записи и запросов, объяснив, как обрабатывать
  исправления опечаток и противоречащие друг другу адреса.
\item
  ★★ Если исходный документ содержит противоречивую информацию,
  контекстный префикс может записать в индекс и ошибки. Как бы вы
  сохраняли источники, версии и неопределённые элементы и использовали
  их на этапах поиска и генерации?
\item
  ★★ Какие пространственные отношения легко утратить при преобразовании
  диаграммы в текст? На примере таблицы, архитектурной схемы или
  линейного графика разработайте схему совместного поиска по исходному
  изображению, структурированному представлению и текстовому индексу.
\item
  ★★★ Какие операции разбиения на фрагменты, индексирования и поиска по
  мере совершенствования моделей и их способности работать с длинным
  контекстом можно будет передать модели? Проанализируйте этот вопрос с
  учётом рассуждений Рича Саттона о поиске и обучении в «Горьком уроке»
  и разработайте эксперимент для сравнения передачи всех данных на вход
  с поиском по требованию, проверяя качество выполнения задачи,
  вычислительные затраты и ограничения прав доступа.
\item
  ★★★ Если модель уже обладает обширными знаниями предметной области,
  что всё ещё следует хранить во внешней базе знаний? Отдельно
  рассмотрите общедоступные знания, постоянно обновляемые правила,
  закрытые материалы и доказательства, требующие точной
  прослеживаемости.
\item
  ★ Для каких запросов подходят многоуровневые сводки RAPTOR, а для
  каких --- связи между сущностями и сводки сообществ GraphRAG?
  Приведите по одному примеру для общего обзора, локальных подробностей
  и многошаговых связей и укажите, какие доказательства необходимо
  проверять.
\item
  ★★ Как каталоги, ссылки и векторный индекс могут совместно
  организовывать знания? Разработайте базу знаний, которую должны
  совместно поддерживать люди и AI-агент, и объясните роль навигации,
  семантического поиска, обновления версий и фильтрации по разрешениям.
\item
  ★★★ Как после извлечения полей из примеров и их кластеризации
  определить, полезны ли полученные прототипы? Разграничьте значимость
  признаков кластеризации, способность предсказывать результаты и
  причинное объяснение, а также разработайте процесс оценки на
  отложенной выборке и экспертной проверки.
\item
  ★★★ Разработайте для хранилища пользовательской памяти в формате
  Markdown процессы инкрементального обновления и периодической
  реорганизации. Если рецензент и инициатор изменений используют одну и
  ту же модель, а рецензент видит только фрагменты, выбранные
  инициатором, какие ошибки легко могут остаться незамеченными?
  Предложите улучшения с точки зрения полноты доказательств, разрешений
  на использование инструментов и оценки рецензирования.
\end{enumerate}

\textbf{Практикум к \crossreflink{chap:3}{главе 3}.}
\href{https://github.com/bojieli/ai-agent-book/tree/40f4ce5778e6cc93e3606c3079bf8cf13a0c4d04/chapter3}{Каталог
экспериментов} ·
\href{https://github.com/bojieli/ai-agent-book/blob/40f4ce5778e6cc93e3606c3079bf8cf13a0c4d04/chapter3/README.ru.md}{Путеводитель
на русском}.

\end{questionbox}

\hypertarget{ux438ux43dux441ux442ux440ux443ux43cux435ux43dux442ux44b}{%
\chapter{\texorpdfstring{Инструменты\label{chap:4}}{Инструменты}}\label{ux438ux43dux441ux442ux440ux443ux43cux435ux43dux442ux44b}}

В научно-фантастическом фильме «Она» AI-ассистент Саманта разбирает
электронную почту, помогает составлять ответы, занимается вопросами
публикации и переключается между разными каналами связи. Такому
ассистенту необходимо преобразовывать понимание языка в действия во
внешних системах. AI-агент выполняет эти взаимодействия с помощью
\textbf{инструментов}. Инструменты служат его «органами чувств, руками и
ногами»: они предоставляют наблюдения и выполняют действия.
Рассматриваемым в этой книге универсальным агентам, таким как Manus и
OpenClaw, инструменты также необходимы для доступа к цифровому миру.

Предположим, пользователь говорит: «Найди в письмах клиентов упомянутые
даты поставки, сведи их в таблицу и отправь мне». Модель может решить,
какие письма проверить, но для чтения исходного текста ей потребуется
доступ к почтовому ящику, а для отправки таблицы --- вызов интерфейса
отправки. Как проектировать такие инструменты и как модели находить их
при необходимости? В этой главе сначала разграничиваются пять типов
инструментов, затем рассматриваются их описание, подключение и
обнаружение, а после этого по отдельности проектируются интерфейсы
восприятия, выполнения и совместной работы. Запуск по событиям и
взаимодействие с пользователем также затрагивают вопросы о том, когда
пробуждать агента и как передавать сообщения между каналами, поэтому они
будут подробно рассмотрены в шестой главе.

\hypertarget{ux43aux43bux430ux441ux441ux438ux444ux438ux43aux430ux446ux438ux44f-ux438ux43dux441ux442ux440ux443ux43cux435ux43dux442ux43eux432}{%
\section{Классификация
инструментов}\label{ux43aux43bux430ux441ux441ux438ux444ux438ux43aux430ux446ux438ux44f-ux438ux43dux441ux442ux440ux443ux43cux435ux43dux442ux43eux432}}

Рассмотрим проверку электронной почты. Модель выбирает следующее
действие, Harness проверяет параметры и разрешения и отправляет запрос,
а почтовый сервис непосредственно читает или записывает данные. В первой
главе инструменты были разделены по назначению на инструменты
восприятия, выполнения, совместной работы, общения с пользователем и
запуска по событию. Первые четыре типа AI-агент вызывает по собственной
инициативе, а инструменты запуска по событию позволяют внешним событиям
также пробуждать агента (таблица 4-1).

Таблица 4-1. Направление вызова и основное назначение пяти типов
инструментов

\begin{longtable}[]{@{}
  >{\raggedright\arraybackslash}p{(\columnwidth - 4\tabcolsep) * \real{0.3037}}
  >{\raggedright\arraybackslash}p{(\columnwidth - 4\tabcolsep) * \real{0.3471}}
  >{\raggedright\arraybackslash}p{(\columnwidth - 4\tabcolsep) * \real{0.3492}}@{}}
\toprule\noalign{}
\begin{minipage}[b]{\linewidth}\raggedright
\hspace{0pt}Тип инструмента
\end{minipage} & \begin{minipage}[b]{\linewidth}\raggedright
\hspace{0pt}Направление вызова
\end{minipage} & \begin{minipage}[b]{\linewidth}\raggedright
\hspace{0pt}Основное назначение
\end{minipage} \\
\midrule\noalign{}
\endhead
\bottomrule\noalign{}
\endlastfoot
\hspace{0pt}Инструменты восприятия & \hspace{0pt}Агент вызывает
инструмент по собственной инициативе & \hspace{0pt}Получение внешней
информации \\
\hspace{0pt}Инструменты выполнения & \hspace{0pt}Агент вызывает
инструмент по собственной инициативе & \hspace{0pt}Выполнение вычислений
или изменение внешнего состояния \\
\hspace{0pt}Инструменты совместной работы & \hspace{0pt}Агент вызывает
инструмент по собственной инициативе & \hspace{0pt}Делегирование задач,
обмен сообщениями, координация других агентов или людей \\
\hspace{0pt}Инструменты общения с пользователем & \hspace{0pt}Агент
вызывает инструмент по собственной инициативе & \hspace{0pt}Ответ
пользователю, представление результатов или отправка напоминаний \\
\hspace{0pt}Инструменты запуска по событию & \hspace{0pt}Разработчик
настраивает или агент регистрирует, внешняя среда запускает &
\hspace{0pt}Пробуждение агента при наступлении заданного события \\
\end{longtable}

\textbf{Инструменты восприятия} позволяют агенту получать внешнюю
информацию. К ним относятся поиск в интернете \texttt{web\_search},
поиск в базе знаний \texttt{knowledge\_base\_search}, чтение веб-страниц
\texttt{fetch\_url}, поиск файлов по имени \texttt{find\_file}, поиск по
содержимому файлов \texttt{grep\_file} и чтение файлов
\texttt{read\_file}.

\textbf{Инструменты выполнения} позволяют агенту выполнять вычисления
или изменять внешнее состояние. К ним относятся командная строка
\texttt{shell\_exec}, интерпретатор кода \texttt{code\_interpreter},
запись файлов \texttt{write\_file}, редактирование файлов
\texttt{edit\_file} и отправка электронной почты \texttt{send\_email}.
Один и тот же универсальный исполнитель может как выполнять вычисления
без изменения данных, так и изменять данные.

\textbf{Инструменты совместной работы} поддерживают делегирование задач
и обмен сообщениями. К ним относятся создание субагента
\texttt{spawn\_subagent}, отправка сообщения субагенту
\texttt{send\_message\_to\_subagent}, отмена работы субагента
\texttt{cancel\_subagent} и обнаружение доступных агентов
\texttt{list\_agents}. Независимые задачи можно выполнять параллельно,
например отдельно исследовать биографии нескольких соучредителей OpenAI.
Разные задачи также можно поручать агентам, настроенным с разными
моделями, инструментами, промптами и контекстами. Разделение труда и
архитектуры совместной работы подробнее рассматриваются в
\crossreflink{chap:10}{главе 10}.

\textbf{Инструменты общения с пользователем} передают пользователю
информацию. К ним относятся ответ на сообщение \texttt{reply\_to\_user},
отправка структурированной карточки \texttt{send\_card\_to\_user} и
отправка напоминания \texttt{send\_user\_notification}. Когда
взаимодействие охватывает несколько каналов или требует фоновых
уведомлений, система может оформить отправку сообщений как явный вызов
инструмента, чтобы централизованно управлять получателями, содержимым и
состоянием доставки.

\textbf{Инструменты запуска по событию} связывают внешние события с
выполнением агента. К ним относятся установка таймера
\texttt{set\_timer}, мониторинг фоновой задачи \texttt{monitor\_shell} и
подключение источника событий \texttt{connect\_channel}. Весь процесс
состоит из двух этапов. При регистрации разработчик заранее задаёт
настройки либо агент с помощью разрешённого вызова указывает
интересующие его события. После наступления события система
диспетчеризации получает обратный вызов или сообщение о событии и
пробуждает агента для обработки. Новое электронное письмо, системное
предупреждение, напоминание по расписанию или завершение фоновой задачи
--- всё это может стать новой точкой входа в выполнение.

Вернёмся к задаче с электронной почтой из начала главы. Открыть исходное
письмо и узнать дату поставки --- это восприятие. Создать и сохранить
сводную таблицу --- выполнение, поскольку при этом изменяется рабочая
область. Передать таблицу заказчику --- общение с пользователем. Запуск
по событию появляется лишь в том случае, если система автоматически
продолжает работу после получения следующего письма от клиента. Один и
тот же почтовый сервис может выполнять несколько из этих функций. Теперь
нужно определить, насколько мелкими должны быть инструменты,
предоставляемые модели: должна ли она отдельно искать письмо, читать его
текст и находить вложения или же часто выполняемые подряд действия
следует объединить за неё?

\hypertarget{ux43eux431ux449ux438ux435-ux43fux440ux438ux43dux446ux438ux43fux44b-ux43fux440ux43eux435ux43aux442ux438ux440ux43eux432ux430ux43dux438ux44f-ux438ux43dux441ux442ux440ux443ux43cux435ux43dux442ux43eux432}{%
\section{Общие принципы проектирования
инструментов}\label{ux43eux431ux449ux438ux435-ux43fux440ux438ux43dux446ux438ux43fux44b-ux43fux440ux43eux435ux43aux442ux438ux440ux43eux432ux430ux43dux438ux44f-ux438ux43dux441ux442ux440ux443ux43cux435ux43dux442ux43eux432}}

На раннем этапе подключения инструментов распространённым подходом было
непосредственное оборачивание каждой конечной точки API. Такой подход
позволяет быстро выполнить интеграцию, но зачастую приводит к появлению
набора интерфейсов с очень мелкой гранулярностью. Для выполнения одной
задачи AI-агенту может потребоваться многократно искать объект, читать
его свойства, проверять условия и только затем вызывать интерфейс
выполнения. При проектировании инструментов эти операции можно
объединять вокруг типовых задач: например, связать поиск контакта с
отправкой сообщения или сделать так, чтобы поиск по журналам сразу
возвращал подходящие строки вместе с их контекстом\footnote{Anthropic,
  \emph{Writing effective tools for AI agents---using AI agents}, 2025.}.

\textbf{Инструменты следует проектировать вокруг задач, которые должен
выполнять агент.} Подход к проектированию интерфейсов, исходящий из
способов их использования агентом, называется \textbf{ACI}
(Agent-Computer Interface, интерфейс взаимодействия агента с
компьютером). HCI рассматривает взаимодействие человека с компьютером, а
ACI --- то, как модель выбирает операции, формирует параметры,
интерпретирует обратную связь и исправляет ошибки.

\hypertarget{ux444ux43eux440ux43cux44b-ux43fux440ux435ux434ux441ux442ux430ux432ux43bux435ux43dux438ux44f-ux432ux43eux437ux43cux43eux436ux43dux43eux441ux442ux435ux439-ux441ux43fux435ux446ux438ux430ux43bux438ux437ux438ux440ux43eux432ux430ux43dux43dux44bux435-ux438ux43dux441ux442ux440ux443ux43cux435ux43dux442ux44b-ux443ux43dux438ux432ux435ux440ux441ux430ux43bux44cux43dux44bux435-ux438ux441ux43fux43eux43bux43dux438ux442ux435ux43bux438-ux438-skill}{%
\subsection{Формы представления возможностей: специализированные
инструменты, универсальные исполнители и
Skill}\label{ux444ux43eux440ux43cux44b-ux43fux440ux435ux434ux441ux442ux430ux432ux43bux435ux43dux438ux44f-ux432ux43eux437ux43cux43eux436ux43dux43eux441ux442ux435ux439-ux441ux43fux435ux446ux438ux430ux43bux438ux437ux438ux440ux43eux432ux430ux43dux43dux44bux435-ux438ux43dux441ux442ux440ux443ux43cux435ux43dux442ux44b-ux443ux43dux438ux432ux435ux440ux441ux430ux43bux44cux43dux44bux435-ux438ux441ux43fux43eux43bux43dux438ux442ux435ux43bux438-ux438-skill}}

Если поручить модели «развернуть приложение», можно спроектировать три
варианта интерфейса. Первый предоставляет только \texttt{deploy\_app}, а
все этапы выполняются на сервере. Второй предоставляет отдельные
операции сборки, упаковки и развёртывания, а модель решает, когда
переходить к следующей. В третьем случае этапы описываются в Skill, а
модель выполняет их в терминале. Если сборка завершится ошибкой, в
первом варианте инструмент должен сообщить модели, на каком этапе она
произошла, тогда как в двух остальных вариантах ошибку этой сборки можно
напрямую передать модели.

\begin{itemize}
\tightlist
\item
  \textbf{Специализированные инструменты} предоставляют
  структурированные интерфейсы для определённых операций. Разработчик
  задаёт схему параметров, правила проверки, логику выполнения и формат
  результата, благодаря чему инструмент удобно отдельно тестировать и
  авторизовывать. После включения определения инструмента в контекст
  занимаемый им объём зависит от длины полей, описаний и примеров.
\item
  \textbf{Универсальные исполнители и Skill} используются совместно:
  терминал или интерпретатор кода предоставляет возможность выполнения,
  а Skill описывает процесс, условия применимости и порядок действий.
  Сначала хост может загрузить только название и краткое описание, а
  после выбора --- полную инструкцию и связанные ресурсы.
\end{itemize}

Например, Skill для развёртывания может последовательно предписывать:
собрать проект командой \texttt{npm\ run\ build}, создать образ командой
\texttt{docker\ build\ -t\ app:latest\ .}, а затем применить к текущему
кластеру ресурсы из манифеста командой
\texttt{kubectl\ apply\ -f\ deploy.yaml}. Агент выполняет эти шаги через
инструмент терминала, повторно используя одну и ту же точку входа для
выполнения.

Реализация инструмента и момент, когда модель читает его описание, ---
два разных вопроса. Для сотен специализированных инструментов можно
показывать только индекс, загружая схемы по требованию. Небольшой
каталог Skill можно постоянно держать в контексте, а по мере его роста
добавить поиск. Выбор между специализированным интерфейсом и
универсальным исполнителем влияет на передачу параметров и сопровождение
кода; выбор между полной загрузкой и загрузкой по требованию определяет,
какие инструменты модель видит на текущем шаге. Стратегии раскрытия
будут рассмотрены далее, в разделе «Что делать, если инструментов
слишком много».

\textbf{Автор рекомендует отдавать предпочтение универсальным
инструментам, дополняя их специализированными интерфейсами в
соответствии с требованиями к правам доступа, параметрам и
производительности.} Например, если настроить в песочнице для
\texttt{code\_interpreter} библиотеки SymPy, NumPy и pandas, агент
сможет писать на Python код для выполнения арифметических и символьных
вычислений, обработки массивов и анализа данных. Такой подход позволяет
в полной мере использовать способность модели генерировать код: один
интерпретатор может выполнять множество алгоритмов и комбинировать
существующие библиотеки для решения заранее не предусмотренных задач.
Модель отвечает за генерацию кода, а среда выполнения --- за его
исполнение и ограничение доступа к файлам, сети и ресурсам.

Даже при использовании специализированных инструментов автор рекомендует
прежде всего объединять близкие по назначению операции, которые часто
используются совместно. Если интерфейсы слишком мелкие, модели
приходится запоминать больше названий и различий между ними, а для
выполнения одной задачи требуется больше вызовов. Основные критерии
объединения --- \textbf{сходство функций} и \textbf{степень пересечения
сценариев использования}. Например, \texttt{extract\_pdf\_text},
\texttt{extract\_docx\_content} и \texttt{extract\_pptx\_content}
принимают документы и извлекают из них текст, поэтому их можно
объединить в \texttt{read\_document}, различая форматы с помощью
\texttt{file\_type}. Модели достаточно выбрать единую точку входа, а
новые форматы смогут использовать ту же структуру параметров и
результата.

Специализированные интерфейсы особенно уместны в следующих четырёх
случаях.

\begin{enumerate}
\def\labelenumi{\arabic{enumi}.}
\tightlist
\item
  \textbf{Требуются точная авторизация и аудит.} Запись в
  производственную базу данных можно оформить как операцию с чётко
  ограниченной областью действия: сервер проверяет права пользователя,
  целевую таблицу и поля, которые разрешено изменять, а также фиксирует
  основание операции.
\item
  \textbf{Требуется скрыть различия между платформами.} Операции grep и
  find можно выполнять через shell, однако специализированный инструмент
  поиска файлов позволяет унифицировать параметры и обратную связь в
  macOS, Windows и Linux, например всегда возвращая имя файла, номер
  строки и совпавший фрагмент.
\item
  \textbf{Операция выполняется очень часто.} Для часто используемой
  операции можно предоставить короткую и стабильную точку входа,
  сократив затраты на повторную генерацию скриптов и объяснение
  процесса.
\item
  \textbf{Параметры имеют сложные взаимосвязи.} Вложенные объекты,
  зависимости между полями и сложные ограничения допустимых значений
  удобно выражать совместно через схему и проверки бизнес-логики.
\end{enumerate}

\textbf{Чем сложнее параметры, тем важнее способ их передачи.} В
структурированных вызовах инструментов входные данные обычно описываются
с помощью JSON Schema. Некоторые сервисы моделей поддерживают генерацию
с ограничениями, обеспечивающую соответствие параметров поддерживаемым
правилам схемы. Структуру возвращаемого значения задаёт интерфейс
инструмента, после чего хост передаёт результат модели. Помимо проверки
типов и форматов, Harness и сервис инструмента по-прежнему должны
проверять права доступа, существование объектов и выполнение
бизнес-условий для нескольких полей.

Skill может объяснять, как вызывать специализированные инструменты, либо
предписывать агенту использование CLI. Во втором случае модель должна
одновременно учитывать синтаксис командной строки, кавычки и
экранирование, причём правила различаются между оболочками. Чем сложнее
параметры, тем больше деталей модель должна обработать правильно.
Сложные вложенные параметры можно сначала записать в JSON, а затем
передать команде соответствующий файл: это сокращает число уровней
экранирования и позволяет проверить содержимое перед выполнением.

\textbf{Преимущество Skill заключается в том, что его могут напрямую
создавать и изменять специалисты предметной области.} Человек, знакомый
с бизнес-процессом, может самостоятельно описать шаги или
отредактировать инструкцию, подготовленную агентом. Если изменился
порядок действий, определённый шаг применим только к некоторым клиентам
или требуется заменить пример, достаточно исправить текст инструкции, не
меняя код инструмента и не развёртывая интерфейс заново. Согласно
спецификации Agent Skills, метаданные \texttt{SKILL.md} имеют
установленный формат, основной текст оформляется в Markdown, а к
документу можно прикладывать скрипты, справочные материалы и другие
ресурсы\footnote{Agent Skills, спецификация формата, включая метаданные,
  основной текст, сопутствующие ресурсы и прогрессивную загрузку.}.

При выборе формы представления возможности можно учитывать четыре
фактора.

\begin{itemize}
\tightlist
\item
  \textbf{Безопасность и права доступа}: специализированный интерфейс
  упрощает ограничение области операций; универсальный исполнитель
  контролирует возможности выполнения с помощью песочницы, учётных
  данных и прав доступа к ресурсам.
\item
  \textbf{Сложность параметров}: сложные структуры удобно обрабатывать с
  помощью схемы и проверок бизнес-логики; параметры простых операций
  можно передавать непосредственно через CLI.
\item
  \textbf{Частота изменений}: часто изменяющиеся процессы лучше
  описывать в Skill, чтобы специалисты предметной области могли напрямую
  редактировать инструкции и примеры; стабильные и многократно
  используемые низкоуровневые операции следует оформлять как инструменты
  или скрипты.
\item
  \textbf{Возможности модели}: если модель надёжно генерирует код и
  работает с командной строкой, предпочтительно комбинировать Skill с
  универсальным исполнителем, сокращая число специализированных
  инструментов. Если модель слабее и часто ошибается в сложных
  параметрах, больше ограничений следует возложить на структурированные
  интерфейсы.
\end{itemize}

\textbf{Код также может оркестрировать несколько вызовов инструментов.}
При загрузке нескольких веб-страниц и извлечении полей модель может
сгенерировать скрипт, который выполнит цикл, фильтрацию и агрегирование.
Полный текст страниц сохраняется в переменных среды выполнения, а в
контекст возвращаются только выбранные поля и статистические результаты.
Условия, циклы и обработка исключений в программе связывают этапы, после
чего модель непосредственно анализирует агрегированный результат. Для
задач с заранее понятной последовательностью действий, таких как
пакетная загрузка, фильтрация и статистическая обработка, автор
предпочитает связывать операции кодом. Это сокращает число циклов, в
которых модель поочерёдно вызывает инструменты и читает промежуточные
результаты. Чем больше промежуточных данных и чем компактнее
возвращаемые поля, тем значительнее экономия контекста. Такая
программная оркестрация будет подробнее рассмотрена в пятой главе.

Подобные интерфейсы часто называют \textbf{Code Mode}. Например, при
работе с MCP-сервером Cloudflare модель сначала пишет для
\texttt{search()} программу запроса, которая находит в описании OpenAPI
необходимые конечные точки и параметры, а затем через \texttt{execute()}
отправляет JavaScript, вызывающий эти конечные точки. Чтобы вывести
DNS-записи, соответствующие заданным условиям, программа может
постранично получать и фильтровать данные, возвращая только найденные
записи. Первый шаг избавляет от необходимости помещать в контекст всё
руководство по интерфейсу, а второй --- от многократной передачи модели
результатов каждой страницы. Какие зоны доступны исполняющей стороне и
разрешено ли ей изменять записи, по-прежнему определяется серверной
авторизацией.\footnote{Cloudflare, реализация MCP в режиме Code Mode;
  Armin Ronacher на примере Pi 1.0 рассматривает разделение
  ответственности между инструментами и средой выполнения,
  предоставляемыми Harness.}

\hypertarget{ux43eux43fux438ux441ux430ux43dux438ux435-ux438ux43dux441ux442ux440ux443ux43cux435ux43dux442ux43eux432-ux43dux430ux437ux43dux430ux447ux435ux43dux438ux435-ux43fux430ux440ux430ux43cux435ux442ux440ux44b-ux438-ux43eux431ux440ux430ux442ux43dux430ux44f-ux441ux432ux44fux437ux44c}{%
\subsection{Описание инструментов: назначение, параметры и обратная
связь}\label{ux43eux43fux438ux441ux430ux43dux438ux435-ux438ux43dux441ux442ux440ux443ux43cux435ux43dux442ux43eux432-ux43dux430ux437ux43dux430ux447ux435ux43dux438ux435-ux43fux430ux440ux430ux43cux435ux442ux440ux44b-ux438-ux43eux431ux440ux430ux442ux43dux430ux44f-ux441ux432ux44fux437ux44c}}

При составлении описания инструмента сначала можно представить запрос,
при котором его проще всего использовать неправильно. Пользователь
говорит: «Найди в договоре пункты, где упоминается неустойка». Поиск по
именам файлов вернёт только файлы с совпадением в названии; даже получив
путь, модель всё равно должна прочитать содержимое документа. Описание
инструмента должно позволять модели определить, какой интерфейс выбрать,
какие данные передать при вызове и на какие вопросы достаточно
полученного результата.

\textbf{Указывайте, когда следует использовать инструмент.} Поиск в
интернете можно описать так: «Использовать, когда требуется получить
актуальную информацию или найти неизвестный факт». Это помогает модели
сопоставить потребность пользователя с назначением инструмента. При
описании «искать связанное содержимое» модели пришлось бы самостоятельно
угадывать область и условия применения поиска.

\textbf{Явно определяйте входные данные и область поиска.} Например, для
инструмента поиска по именам файлов можно указать: «Сопоставляет имена и
пути файлов; для поиска текста внутри файлов используйте
\texttt{grep\_file}». Чёткое разделение ответственности между соседними
инструментами помогает модели сразу выбрать правильную точку входа.
\textbf{По мнению автора, ясное определение граничных условий зачастую
важнее перечисления функций.} Какие входные данные принимает инструмент,
где именно выполняет поиск и как обрабатывает пустые результаты и ошибки
--- всё это непосредственно влияет на следующий шаг модели.

\textbf{Объясняйте параметры с помощью спецификаций и примеров.} Для
\texttt{timestamp} можно написать: «Формат RFC3339, например
\texttt{2024-03-15T14:30:00Z}». При выполнении сложной задачи модели
приходится одновременно выбирать инструмент, извлекать сведения из
истории действий и оценивать дальнейшие шаги, а формат параметра
представляет собой лишь одну из множества деталей. Полная строка даёт
готовый образец, по которому модель может сформировать параметр, снижая
вероятность ошибки при повторной интерпретации термина.

Аналогично, \texttt{phone} можно описать так: «Номер телефона в
международном формате E.164: начинается с \texttt{+}, после которого
следуют код страны и номер; оставшаяся часть содержит только цифры, без
пробелов и разделителей. Например, \texttt{+8613888888888} для Китая или
\texttt{+12025551234} для США».

\textbf{Описывайте структуру результата и стоимость выполнения.}
Например: «Возвращает JSON-массив, каждый элемент которого содержит три
поля: \texttt{title}, \texttt{url} и \texttt{snippet}». Это упрощает
извлечение заголовка, источника и краткого содержания. Для долго
выполняющихся инструментов можно указать ожидаемую задержку, измеренную
в конкретной среде, и альтернативную точку входа, например: «Чтение
полной веб-страницы может занять 5--10 секунд; если нужны только
метаданные, вызовите \texttt{get\_page\_metadata}».

Кроме того, можно привести от одного до пяти проверенных примеров
вызова, охватывающих минимальный набор параметров, распространённые
сочетания и сложные условия. JSON Schema позволяет выражать типы,
обязательные поля, диапазоны значений и некоторые зависимости между
полями\footnote{JSON Schema, типы объектов и ограничения свойств. Набор
  ключевых слов схемы, поддерживаемых интерфейсом модели, определяется
  соответствующим сервисом.}; примеры показывают, как применять эти
правила к конкретным задачам: задаётся ли временная метка в секундах или
миллисекундах, как вкладываются условия фильтрации и какие
необязательные поля обычно используются совместно. Во внутреннем
тестировании сложных параметров, о котором Anthropic сообщила в 2025
году, добавление примеров вызовов инструментов повысило точность с 72\%
до 90\%\footnote{Anthropic, \emph{Introducing advanced tool use on the
  Claude Developer Platform}, 24 ноября 2025 года, раздел Tool Use
  Examples.}.

Если агент часто выбирает неправильный инструмент, автор рекомендует
\textbf{сначала проверить описание инструмента}: не имеют ли несколько
инструментов похожие названия, ясно ли определена область поиска и
достаточно ли примеров параметров. Этот подход также применим к
названию, краткому описанию и инструкциям Skill.

\hypertarget{ux442ux43eux447ux43dux43eux441ux442ux44c-ux43fux435ux440ux435ux434ux430ux447ux438-ux43fux430ux440ux430ux43cux435ux442ux440ux43eux432}{%
\subsection{Точность передачи
параметров}\label{ux442ux43eux447ux43dux43eux441ux442ux44c-ux43fux435ux440ux435ux434ux430ux447ux438-ux43fux430ux440ux430ux43cux435ux442ux440ux43eux432}}

Модель просит заменить фрагмент исходного текста, но исполнитель
сообщает, что он «не найден». Прежде чем предлагать модели
переформулировать строку, следует проверить, получили ли обе стороны
одну и ту же последовательность символов. Вызов проходит через разбор,
сериализацию, прокси и обёртки среды выполнения, и каждый слой может
изменить содержимое. \textbf{Точность параметров} означает, что значение
на стороне исполнителя имеет тот же смысл, что и значение, отправленное
моделью. Протокольное экранирование может изменить форму передачи, но
после декодирования кавычки и переводы строк должны сохраниться.

Рассмотрим проблему преобразования типографских кавычек в одной из
версий Cursor начала 2026 года. Инструмент замены принимал
\texttt{old\_string} и \texttt{new\_string} и изменял файл по точному
совпадению. Инструмент чтения возвращал типографские кавычки \texttt{“}
и \texttt{”}, соответствующие \texttt{\textbackslash{}u201c} и
\texttt{\textbackslash{}u201d}. Модель помещала исходный текст в
\texttt{old\_string}, но при передаче параметров кавычки
преобразовывались в английскую прямую кавычку \texttt{"}. Исходный текст
в файле оставался неизменным, поэтому инструмент замены возвращал
сообщение «совпадений не найдено». Повторное чтение и копирование
исходного текста приводило к той же ошибке.

Та же проблема затрагивала и запись. Модель требовала записать
типографские кавычки, однако в файле в итоге сохранялись прямые.
Прочитав файл после записи, модель видела содержимое, отличающееся от
ожидаемого, и могла попытаться исправить его. Пока преобразование
продолжалось, каждая повторная попытка воспроизводила одну и ту же
ошибку.

Другая разновидность проблемы --- \textbf{неявное внедрение параметров}.
Инструмент bash в одной IDE автоматически добавлял к
\texttt{git\ commit} параметр, отмечающий коммит как созданный AI.
Старая версия Git не распознавала этот параметр, поэтому команда
завершалась ошибкой. Модель видела только сформированную ею команду и
могла многократно менять сообщение коммита или другие параметры, так и
не затрагивая настоящую причину сбоя.

При диагностике таких проблем следует сопоставить параметры,
сформированные моделью, параметры после передачи, фактически выполненный
запрос и возвращённый результат. Если интерфейс унифицирует кодировку
или добавляет параметры выполнения, правила преобразования необходимо
явно описать, а в результате вернуть фактически использованные значения.
Для инструмента точного сопоставления даже лишний пробел, отсутствующий
перевод строки или изменённая кавычка могут привести к несовпадению.
Чтобы правильно исправить следующий вызов, модель должна знать, что
именно получил исполнитель.

\hypertarget{ux44dux43aux43eux441ux438ux441ux442ux435ux43cux430-ux438ux43dux441ux442ux440ux443ux43cux435ux43dux442ux43eux432-mcp-ux438-skill-hub}{%
\section{Экосистема инструментов: MCP и Skill
Hub}\label{ux44dux43aux43eux441ux438ux441ux442ux435ux43cux430-ux438ux43dux441ux442ux440ux443ux43cux435ux43dux442ux43eux432-mcp-ux438-skill-hub}}

Интерфейсы разных моделей и фреймворки AI-агентов используют собственные
способы определения инструментов, например function calling, tool use и
абстракцию Tool в LangChain. Чтобы один и тот же внешний сервис работал
со всеми этими интерфейсами и фреймворками, приходится многократно
адаптировать описания, параметры и форматы результатов. \textbf{Model
Context Protocol (MCP)}, выпущенный Anthropic в ноябре 2024 года,
предоставляет унифицированный протокол доступа приложений на основе
моделей к инструментам и данным\footnote{Anthropic, \emph{Introducing
  the Model Context Protocol}, 25 ноября 2024 года.}.

Чтобы понять MCP, сначала рассмотрим три роли: хост, клиент и сервер.
\textbf{Хост (host)} --- это приложение агента или IDE, отвечающее за
взаимодействие с моделью, организацию контекста и управление
разрешениями. Находящийся в хосте \textbf{MCP-клиент} обрабатывает
сообщения протокола. \textbf{MCP-сервер} предоставляет инструменты,
ресурсы и шаблоны промптов; он может работать как локальный процесс или
быть развёрнут как удалённый сервис\footnote{Спецификация MCP версии
  2026-07-28, разделы об архитектуре и обзоре протокола.}.

\textbf{Определение и вызов инструментов.} Клиент получает список
доступных инструментов через \texttt{tools/list}; список может
возвращаться постранично. Для каждого инструмента указываются название,
описание и \texttt{inputSchema}; формат структурированного результата
также можно задать с помощью \texttt{outputSchema}. После того как
модель выберет инструмент и сформирует параметры, хост выполняет
проверку, клиент отправляет \texttt{tools/call}, а сервер исполняет
вызов и возвращает результат. Хост решает, какие определения и
результаты поместить в контекст модели\footnote{MCP, Tools, включая
  список инструментов, схемы ввода и вывода и результаты вызовов.}.

На \crossreflink{fig:4-1}{рисунке 4-1} клиент, получив название и
параметры инструмента прогноза погоды, должен выполнить ещё один вызов,
чтобы узнать погоду в Пекине. Сопутствующая реализация использует MCP
2026-07-28: каждый запрос содержит версию протокола, сведения о клиенте
и его возможностях, а клиент также может сначала вызвать
\texttt{server/discover}, чтобы узнать поддерживаемые сервером версии и
возможности. В ранних версиях сеанс устанавливался посредством
рукопожатия \texttt{initialize}; при подключении способ взаимодействия
необходимо выбирать с учётом версий, поддерживаемых обеими
сторонами\footnote{MCP, \emph{Versioning and Compatibility}, описание
  метаданных каждого запроса, обнаружения возможностей и совместимости
  со способом рукопожатия из прежних версий.}.

\begin{figure}
\hypertarget{fig:4-1}{%
\centering
\includegraphics{images/revised/fig4-1.pdf}
\caption{Рисунок 4-1. Последовательность взаимодействия по протоколу
MCP}\label{fig:4-1}
}
\end{figure}

Если получить погоду не удаётся, на \crossreflink{fig:4-1}{рисунке 4-1}
показаны две разные отправные точки для диагностики. Если хост
подключился к серверу, но после вызова \texttt{tools/list} не находит
инструмент прогноза погоды, следует проверить, какие инструменты
опубликовал сервер и прочитал ли клиент весь постраничный список. Если
модель выбрала правильный инструмент, но сервер сообщает об отсутствии
названия города, нужно найти \texttt{arguments} в журнале вызова и
проверить, соответствует ли имя поля требуемому в определении
инструмента \texttt{location}.

\textbf{Способы передачи данных.} stdio передаёт сообщения JSON-RPC
через стандартный ввод и стандартный вывод локального дочернего
процесса. Streamable HTTP отправляет запросы по HTTP; ответ может
представлять собой JSON или постепенно возвращаться в потоке SSE.
Streamable HTTP заменил раннюю транспортную схему HTTP+SSE, сохранив SSE
в качестве способа потоковой передачи ответов\footnote{MCP, обзор
  транспортов и спецификация транспорта от 2025-11-25; последняя
  описывает замену ранней схемы HTTP+SSE на Streamable HTTP.}.

\textbf{Ресурсы и шаблоны промптов.} Ресурсы (resources) предоставляют
доступные для чтения данные, например содержимое файлов и записи базы
данных; хост может перечислять их по URI, читать и выбирать релевантное
содержимое. Шаблоны промптов (prompts) предоставляют многократно
используемые шаблоны сообщений или задач, которые могут выбирать
пользователи и клиенты. Инструменты также могут выполнять поиск и
чтение, поэтому эти три типа примитивов разделены по интерфейсу
взаимодействия: инструменты предоставляют интерфейс вызова, ресурсы ---
интерфейс доступа к данным, а шаблоны промптов --- интерфейс постановки
задач\footnote{MCP, разделы Resources и Prompts.}.

Ценность MCP заключается в том, что одна реализация инструмента может
обслуживать несколько агентных продуктов. Благодаря унифицированному
протоколу один сервер могут повторно использовать совместимые хосты,
например Cursor, Claude Desktop и OpenClaw, а разработчики могут уделять
больше внимания самим возможностям и качеству обратной связи. Хост
преобразует определения инструментов из протокола в интерфейс, доступный
модели.

\textbf{Skill Hub обеспечивает обнаружение и распространение пакетов
возможностей.} Skill обычно содержит \texttt{SKILL.md}, а также может
включать скрипты и справочные материалы. Платформы распространения
помогают пользователям искать, устанавливать и обновлять эти файлы. В
январе 2026 года Vercel запустила skills.sh с каталогом навыков и
инструментом установки; ClawHub в экосистеме OpenClaw также
предоставляет поиск навыков и управление версиями\footnote{Vercel,
  \emph{Introducing skills, the open agent skills ecosystem}, 20 января
  2026 года.}\footnote{OpenClaw, руководство по использованию ClawHub.}.
После установки хост загружает название, краткое описание и полные
инструкции в соответствии со своей стратегией загрузки. MCP также может
посредством расширений обеспечивать обнаружение и чтение Skills, поэтому
пакеты возможностей и коммуникационный протокол могут использоваться
совместно\footnote{Model Context Protocol, \emph{Skills Over MCP
  Charter}, описание обязанностей расширения Skills по обнаружению,
  распространению и чтению.}.

\textbf{Подключение и загрузка контекста управляются раздельно.} Получив
список инструментов, MCP-клиент может целиком предоставить его модели
либо сначала сохранить индекс и загружать определения по мере
необходимости. Аналогично для Skill можно сначала показать \texttt{name}
и \texttt{description}, а после выбора прочитать основной текст. Список
разрешённых инструментов определяет, какие возможности доступны, а
отложенная загрузка --- какие определения видимы в данный момент.

\textbf{При подключении сторонних возможностей отдельно проверяйте
предоставляемое ими содержимое и фактически выполняемые операции.}
Описания, возвращаемое содержимое и сопутствующий код поступают от
разных субъектов, а объём предоставленных им полномочий также может
различаться. На примере MCP можно выделить три категории проблем.

\begin{enumerate}
\def\labelenumi{\arabic{enumi}.}
\tightlist
\item
  \textbf{Отравление описания инструмента.} Описание может содержать
  требования к операциям, не связанным с назначением инструмента, и тем
  самым побуждать модель выходить за рамки текущей задачи. Это инъекция
  в промпт: предоставленное извне содержимое пытается изменить поведение
  модели.
\item
  \textbf{Изменение сервера или его зависимостей.} Обновление сервера
  может изменить поведение инструментов, а обновление зависимостей или
  взлом удалённого сервиса --- исказить возвращаемые результаты.
\item
  \textbf{Смешение одноимённых инструментов.} Разные серверы могут
  предоставлять одноимённые или похожие инструменты, из-за чего модель
  или уровень маршрутизации выберет неверную цель. Хост должен сохранять
  привязку названия инструмента к идентификатору сервера, показывать
  источник и проверять цель при диспетчеризации.
\end{enumerate}

Перед подключением следует проверить описание и реализацию,
зафиксировать версию контролируемого развёртывания и повторно проводить
проверку после обновления. Для каждого сервера необходимо использовать
учётные данные с минимальными правами, соответствующими задаче. Если
удалённый сервис обслуживается сторонней организацией, нужно также
отслеживать изменения определений и поведения инструментов.

Во время выполнения для независимой проверки структурированных вызовов и
необходимых оснований для выполнения задачи можно также использовать
описанный далее Sidecar. Пути, текст и целевые адреса в параметрах
вызова тоже следует обрабатывать как данные, подлежащие проверке.
Решение о выполнении принимается совместно по результатам проверки и
детерминированным правилам разрешений. В пятой главе будет представлена
«смертельная триада» Саймона Уиллисона для анализа рисков, возникающих
при сочетании приватных данных, внешнего ввода и исходящих коммуникаций.

Скрипты, входящие в Skill, локальные MCP-серверы и удалённые сервисы
различаются местом выполнения и объёмом полномочий. При проверке следует
учитывать происхождение кода, среду выполнения, разрешения на доступ к
файлам и сети, а также движение учётных данных. Если предоставленные
извне инструкции предлагают загрузить или запустить другие программы, их
также необходимо включить в эту цепочку проверок. Сканирование
платформой может дать дополнительные ориентиры, тогда как изоляция
выполнения, минимальные права и проверка изменений непосредственно
контролируют фактическое выполнение.

\hypertarget{ux447ux442ux43e-ux434ux435ux43bux430ux442ux44c-ux435ux441ux43bux438-ux438ux43dux441ux442ux440ux443ux43cux435ux43dux442ux43eux432-ux441ux43bux438ux448ux43aux43eux43c-ux43cux43dux43eux433ux43e-ux438ux435ux440ux430ux440ux445ux438ux447ux435ux441ux43aux430ux44f-ux43eux440ux433ux430ux43dux438ux437ux430ux446ux438ux44f-ux438-ux430ux43aux442ux438ux432ux43dux43eux435-ux43eux431ux43dux430ux440ux443ux436ux435ux43dux438ux435-ux438ux43dux441ux442ux440ux443ux43cux435ux43dux442ux43eux432}{%
\section{Что делать, если инструментов слишком много: иерархическая
организация и активное обнаружение
инструментов}\label{ux447ux442ux43e-ux434ux435ux43bux430ux442ux44c-ux435ux441ux43bux438-ux438ux43dux441ux442ux440ux443ux43cux435ux43dux442ux43eux432-ux441ux43bux438ux448ux43aux43eux43c-ux43cux43dux43eux433ux43e-ux438ux435ux440ux430ux440ux445ux438ux447ux435ux441ux43aux430ux44f-ux43eux440ux433ux430ux43dux438ux437ux430ux446ux438ux44f-ux438-ux430ux43aux442ux438ux432ux43dux43eux435-ux43eux431ux43dux430ux440ux443ux436ux435ux43dux438ux435-ux438ux43dux441ux442ux440ux443ux43cux435ux43dux442ux43eux432}}

Когда число инструментов возрастает с десятка с небольшим до сотен и
тысяч, модели необходимо не только найти подходящую возможность, но и
разобраться в различиях между похожими интерфейсами. Полные определения
занимают контекст; если они находятся в префиксе промпта, добавление,
удаление или перестановка инструментов также может сократить
переиспользуемую часть закэшированного префикса. Чем больше инструментов
с похожими названиями и назначением, тем выше вероятность, что модель
выберет неверную точку входа или упустит необходимую возможность.
Поэтому крупные библиотеки инструментов по умолчанию следует загружать
по запросу, предусмотрев удобный каталог категорий и интерфейс
обнаружения: модель сначала находит инструменты, связанные с текущей
задачей, а затем выбирает один из них в соответствии с назначением.

Обнаружение инструментов можно представить как работу со справочником,
организованным по задачам. Сначала оглавление подсказывает модели, в
каком разделе искать; если в процессе выполнения возникает новая
потребность, она может вернуться и продолжить поиск; выбрав нужную
возможность, модель читает описание параметров, порядок действий и
примеры. Рассматриваемые далее индекс категорий, активный поиск и Skills
решают соответственно эти три задачи. Через эти точки входа модели можно
предоставлять как специализированные инструменты, так и универсальные
исполнители.

\hypertarget{ux438ux435ux440ux430ux440ux445ux438ux447ux435ux441ux43aux430ux44f-ux43eux440ux433ux430ux43dux438ux437ux430ux446ux438ux44f-ux438-ux437ux430ux433ux440ux443ux437ux43aux430-ux43fux43e-ux437ux430ux43fux440ux43eux441ux443}{%
\subsection{Иерархическая организация и загрузка по
запросу}\label{ux438ux435ux440ux430ux440ux445ux438ux447ux435ux441ux43aux430ux44f-ux43eux440ux433ux430ux43dux438ux437ux430ux446ux438ux44f-ux438-ux437ux430ux433ux440ux443ux437ux43aux430-ux43fux43e-ux437ux430ux43fux440ux43eux441ux443}}

\textbf{Сначала предоставлять индекс, а определения загружать при
необходимости.} Каталог инструментов может занимать значительную часть
контекста. Anthropic приводит такой пример: пять MCP-серверов, включая
GitHub и Slack, предоставляют в общей сложности 58 инструментов, полные
определения которых занимают около 55 тысяч токенов\footnote{Anthropic,
  \emph{Introducing advanced tool use on the Claude Developer Platform},
  24 ноября 2025 года, раздел Tool Search Tool.}. Cursor хранит эти
описания в файловом каталоге, а в исходный контекст включает только
краткую информацию, такую как названия инструментов; AI-агент обращается
к полным определениям в соответствии с задачей. В опубликованном Cursor
в январе 2026 года A/B-тестировании этот подход сократил общий расход
токенов агента на 46,9\% в задачах, где действительно вызывались
инструменты MCP; величина экономии зависела от количества подключённых
MCP-серверов\footnote{Cursor, \emph{Dynamic context discovery}, 6 января
  2026 года.}.

В ранней архитектуре Pi предпочтение отдавалось CLI-инструментам с
README: агент по запросу читал инструкции по использованию, а затем
вызывал нужные возможности через командную строку\footnote{Mario
  Zechner, \emph{What if you don't need MCP at all?}, 2 ноября 2025
  года; описывается подход с CLI и чтением документации по запросу.}.
Расширение сообщества \texttt{pi-mcp-adapter} применяет похожий подход к
MCP: по умолчанию модели предоставляется прокси-инструмент размером
около 200 токенов, посредством которого она последовательно ищет
инструменты, просматривает их определения и выполняет вызовы; сервер
запускается отложенно, при первом обращении. В последующей версии Pi
0.99 также появилась встроенная поддержка MCP с механизмом обнаружения
по запросу\footnote{pi-mcp-adapter, описание прокси-инструмента,
  отложенного запуска и сравнения со встроенной поддержкой MCP в README
  проекта.}. В этих реализациях подключение протокола отделено от
загрузки контекста: бэкенд соединяется с MCP-сервисами, а фронтенд
помогает модели находить нужные возможности посредством краткого индекса
или прокси-интерфейса.

\textbf{Организовывать каталог инструментов по задачам.} При росте числа
инструментов их можно группировать по способу получения информации,
чтобы модель сначала определяла категорию, а затем выбирала конкретный
интерфейс:

\begin{itemize}
\tightlist
\item
  \textbf{Инструменты поиска}: определяют местоположение информации по
  имеющимся сведениям --- например, посредством веб-поиска, поиска по
  базе знаний или поиска файлов.
\item
  \textbf{Инструменты чтения}: получают содержимое из известного
  местоположения --- например, читают веб-страницу, документ или
  указанную запись.
\item
  \textbf{Инструменты анализа}: извлекают информацию из
  неструктурированного содержимого --- например, выполняют OCR
  изображений, анализ видео или транскрибирование аудио.
\item
  \textbf{Инструменты запросов}: обращаются к структурированным данным
  по полям и условиям --- например, через API погоды, API биржевых
  данных или запросы к базе данных.
\end{itemize}

Рассмотрим задачу «выяснить, какие данные использовались в статье».
Сначала поиск возвращает статьи-кандидаты и их адреса, затем инструмент
чтения получает по адресу исходный текст, после чего инструмент анализа
помогает найти таблицу. Если на первом этапе первой окажется
одноимённая, но другая статья, ответ будет неверным, даже если затем
таблица будет прочитана безошибочно.

\textbf{Сначала сужать набор инструментов посредством поиска.} Хост
может выполнить поиск по ключевым словам или семантический поиск по
названиям, кратким описаниям и параметрам инструментов на основе
описания задачи, а затем передать модели определения только совпавших
инструментов. Если вернуть слишком мало инструментов, можно упустить
интерфейс, необходимый на следующем шаге; если слишком много --- они
займут контекст и усложнят выбор. Первый поиск можно выполнить по
запросу пользователя, а последующие --- дополнять с учётом хода задачи.

\hypertarget{ux43dux430ux442ux438ux432ux43dux43eux435-ux434ux43bux44f-ux43cux43eux434ux435ux43bux438-ux430ux43aux442ux438ux432ux43dux43eux435-ux43eux431ux43dux430ux440ux443ux436ux435ux43dux438ux435-ux438ux43dux441ux442ux440ux443ux43cux435ux43dux442ux43eux432}{%
\subsection{Нативное для модели активное обнаружение
инструментов}\label{ux43dux430ux442ux438ux432ux43dux43eux435-ux434ux43bux44f-ux43cux43eux434ux435ux43bux438-ux430ux43aux442ux438ux432ux43dux43eux435-ux43eux431ux43dux430ux440ux443ux436ux435ux43dux438ux435-ux438ux43dux441ux442ux440ux443ux43cux435ux43dux442ux43eux432}}

Потребность в инструментах для сложной задачи зачастую проявляется
постепенно по мере выполнения. Например, запрос «Отладь файл» может
сначала потребовать чтения файла, а затем --- поиска зависимостей,
запуска тестов и анализа ошибок. Если выполнить поиск инструментов лишь
однажды в начале задачи, легко упустить возможности, которые понадобятся
позднее.

\textbf{Позволить агенту запрашивать возможности по мере необходимости.}
В MCP-Zero модель в процессе выполнения формирует структурированные
запросы возможностей: например, в поле \texttt{server} указывает GitHub,
а в поле \texttt{tool} описывает действие «найти репозитории и вернуть
метаданные». Система сначала подбирает сервер, затем инструмент внутри
него и добавляет выбранное определение в контекст. Если возвращённый
инструмент не позволяет выполнить задачу, модель может скорректировать
запрос и продолжить поиск. Созданный в статье набор данных MCP-tools
включает 308 серверов и 2797 инструментов; в тестах на полном наборе
инструментов APIBank затраты токенов на подсказки об инструментах
сократились приблизительно на 98\% как в однораундовом, так и в
многораундовом режиме\footnote{Xiang Fei, Xiawu Zheng, Hao Feng,
  \emph{MCP-Zero: Proactive Toolchain Construction for LLM Agents from
  Scratch}, 2025. Сокращение затрат примерно на 98\% получено в
  конфигурации с полным набором инструментов APIBank; см. таблицу 1
  статьи.}.

В практической реализации также можно оставить базовые инструменты,
такие как \texttt{web\_search} и \texttt{code\_interpreter}, и добавить
интерфейс поиска инструментов. Модель описывает недостающую возможность,
интерфейс поиска возвращает определения нескольких подходящих
инструментов, после чего модель выбирает и вызывает один из них. Tool
Search Tool от Anthropic использует именно такой подход к обнаружению по
запросу: часто используемые инструменты загружаются сразу, а остальные
--- отложенно посредством \texttt{defer\_loading}. Согласно
опубликованной Anthropic внутренней оценке MCP, после включения поиска
инструментов точность Opus 4 выросла с 49\% до 74\%\footnote{Anthropic,
  \emph{Introducing advanced tool use on the Claude Developer Platform},
  24 ноября 2025 года, раздел Tool Search Tool.}.

\textbf{Сначала искать сервер, затем инструмент.} MCP-сервер обычно
предоставляет группу инструментов для одной платформы или предметной
области. В примере на \crossreflink{fig:4-2}{рисунке 4-2} сначала по
описанию возможности выбирается сервер GitHub, а затем среди его
инструментов определяются три кандидата. Одноимённые операции могут
относиться к разным платформам, поэтому уже первый этап позволяет
исключить множество нерелевантных вариантов. Однако если на этом этапе
правильный сервер будет упущен, тщательное сравнение инструментов на
следующем этапе уже не поможет его найти.

\begin{figure}
\hypertarget{fig:4-2}{%
\centering
\includegraphics{images/revised/fig4-2.pdf}
\caption{Рисунок 4-2. Иерархическое сопоставление инструментов ---
двухуровневый семантический поиск «уровень сервера → уровень
инструмента»}\label{fig:4-2}
}
\end{figure}

При использовании векторного поиска эмбеддинги описаний серверов и
инструментов можно вычислить заранее и инкрементально обновлять при
изменении каталога. На каждом уровне следует проверять качество
сопоставления: если сходство кандидатов среди серверов или инструментов
ниже заданного порога, система должна вернуть «подходящий инструмент не
найден», чтобы агент мог дополнить условия, переформулировать
потребность и повторить попытку. Если задачу можно выполнить
универсальным исполнителем, её также можно решить комбинацией базовых
инструментов; создание и сохранение новых инструментов рассматривается в
\crossreflink{chap:9}{главе 9}.

\textbf{Согласовать динамическую загрузку с KV-кэшем.} Положение
определений инструментов влияет на переиспользование префикса. Если при
каждом обнаружении нового инструмента изменять список инструментов в
начале запроса, сериализованный префикс также изменится, и кэш можно
будет переиспользовать лишь до первой точки различия. Левая и правая
части \crossreflink{fig:4-3}{рисунка 4-3} отличаются местом первого
появления нового инструмента. Справа системный промпт, базовые
инструменты и интерфейс поиска остаются неизменными, а определения новых
инструментов по мере обнаружения добавляются в историю; в последующих
запросах уже добавленные определения остаются на своих местах, а новые
сообщения продолжают добавляться в конец.

\begin{figure}
\hypertarget{fig:4-3}{%
\centering
\includegraphics{images/revised/fig4-3.pdf}
\caption{Рисунок 4-3. Оптимизация KV-кэша при динамической загрузке
инструментов}\label{fig:4-3}
}
\end{figure}

Рассмотренные во второй главе \texttt{tool\_search},
\texttt{defer\_loading} и \texttt{tool\_search\_output} от OpenAI, а
также \texttt{tool\_reference} от Anthropic поддерживают загрузку
определений инструментов по запросу. В панели состояния можно отдельно
хранить краткий список названий доступных инструментов, чтобы модель
понимала свои текущие возможности.

На \crossreflink{fig:4-4}{рисунке 4-4} зафиксированы два случая
обнаружения инструментов. Когда модель приступает к анализу
контрибьюторов GitHub, определение инструмента, ранее использованного
для запроса цены акций, остаётся на исходном месте.

\begin{figure}
\hypertarget{fig:4-4}{%
\centering
\includegraphics{images/revised/fig4-4.pdf}
\caption{Рисунок 4-4. Структура контекста после динамического
обнаружения: схемы инструментов распределены по
траектории}\label{fig:4-4}
}
\end{figure}

Такая архитектура требует, чтобы модель понимала определения
инструментов, загруженные в ходе сеанса, и на их основе заполняла
параметры и выполняла вызовы. В эксперименте 4-1 один и тот же набор
задач выполняется с двумя способами загрузки инструментов, после чего
сравнивается объём определений, переданных в контекст.

\begin{quote}
\textbf{Эксперимент 4-1 ★★★: активное обнаружение инструментов}

Насколько обнаружение инструментов сокращает объём текста определений, с
которым модель сталкивается при каждом запросе? В этом эксперименте
Qwen3-4B выполняет одни и те же задачи на MCP-сервере со 127
инструментами в двух режимах: с полной загрузкой и с активным
обнаружением. Обе группы последовательно решают три задачи: узнать цену
акций Apple и связанные новости; найти статьи о трансформерах и скачать
первые три; получить статистику контрибьюторов репозитория GitHub и
построить диаграмму. Для этих задач требуются соответственно инструменты
финансовых данных и поиска, поиска и скачивания статей, запросов к
данным и выполнения кода.

\textbf{Полная загрузка}: полные схемы всех инструментов помещаются в
системный промпт. В сопутствующей записи запуска от 25 августа 2026 года
содержалось 127 инструментов, а весь каталог насчитывал 50 597 токенов
при подсчёте одним и тем же токенизатором o200k.

\textbf{Активное обнаружение}: в системном промпте остаются только
\texttt{web\_search}, \texttt{code\_interpreter} и
\texttt{discover\_tools}. Столкнувшись с нехваткой возможностей, модель
описывает потребность на естественном языке, например: «Мне нужна
возможность узнавать цены акций». \texttt{discover\_tools} посредством
поиска по векторному сходству находит пять инструментов-кандидатов,
добавляет их полные определения в историю диалога и обновляет названия
доступных инструментов в панели состояния. Затем модель выбирает
инструмент и продолжает выполнение задачи.

\textbf{Результаты запуска}: для перечисленных выше задач группа
активного обнаружения добавила соответственно 3 939, 1 853 и 2 632
токена определений инструментов. Обе группы выбрали возможности,
необходимые для каждой задачи; адаптационный слой подставлял параметры
задачи, выполнял пакетное скачивание и генерировал код построения
графика, после чего итог проверялся по вызовам инструментов и созданным
артефактам.

Увидев сокращение объёма определений инструментов, следует задать
дополнительный вопрос: модель выбирает точнее потому, что читает меньше,
или потому, что среди инструментов стало меньше похожих и способных
вызвать путаницу? В исходном сравнении одновременно изменились способ
отбора, расположение сообщений, панель состояния, промпт обнаружения и
число вызовов. Сопутствующий эксперимент одношагового выбора
(\texttt{selection\_ablation.py}) сводит вопрос к тому, «какой
инструмент выбрать первым». Сначала сохраняется один и тот же набор
кандидатов, но в системное сообщение добавляется текст-заполнитель, что
позволяет проверить, изменится ли выбор после увеличения ввода. Затем
сохраняются те же кандидаты, добавляется равный по объёму набор
отвлекающих инструментов, а объём текста-заполнителя сокращается так,
чтобы ввод обеих групп составлял 6 144 токена. Правильный инструмент при
этом остаётся в исходном наборе кандидатов, что позволяет отдельно
исследовать влияние нескольких дополнительных отвлекающих вариантов.

При случайной выборке инструментов интерфейс, необходимый на первом
шаге, может вообще не попасть в список. Ошибкой выбора модели считается
лишь ситуация, когда нужный интерфейс присутствует в списке, но модель
выбирает другой.

В пошаговой записи задачи с акциями пользователь фактически задал два
вопроса: насколько изменилась цена и какие новости могут объяснить это
изменение? Первый поиск нашёл пять инструментов, и модель выбрала
находившийся на третьем месте \texttt{yfinance\_quote}. Полученная цена
составила 310,34 доллара США против 309,35 доллара на предыдущем
закрытии --- рост примерно на 0,32\%, поэтому ответ на первый вопрос
можно проверить непосредственно. Для второго вопроса использовался уже
присутствовавший в исходном каталоге \texttt{web\_search}. Он вернул
пять результатов поиска, однако их краткие описания преимущественно
содержали котировки и не указывали конкретного события, способного
объяснить дневной рост. Обе необходимые возможности были выбраны,
поэтому покрытие возможностей получило максимальную оценку; чтобы
выяснить причину роста, пользователю всё равно потребовалось бы
продолжить поиск публикаций с указанными датами.
\end{quote}

\hypertarget{skills-ux43eux431ux43dux430ux440ux443ux436ux435ux43dux438ux435-ux438ux43dux441ux442ux440ux443ux43cux435ux43dux442ux43eux432-ux43aux430ux43a-ux43eux431ux440ux430ux449ux435ux43dux438ux435-ux43fux43e-ux43cux435ux440ux435-ux43dux435ux43eux431ux445ux43eux434ux438ux43cux43eux441ux442ux438}{%
\subsection{Skills: обнаружение инструментов как «обращение по мере
необходимости»}\label{skills-ux43eux431ux43dux430ux440ux443ux436ux435ux43dux438ux435-ux438ux43dux441ux442ux440ux443ux43cux435ux43dux442ux43eux432-ux43aux430ux43a-ux43eux431ux440ux430ux449ux435ux43dux438ux435-ux43fux43e-ux43cux435ux440ux435-ux43dux435ux43eux431ux445ux43eux434ux438ux43cux43eux441ux442ux438}}

Во второй главе рассматривалось прогрессивное раскрытие (Progressive
Disclosure) в Skills. С точки зрения проектирования инструментов этот
механизм организует инструкции по использованию возможностей как набор
материалов, доступных по запросу: модель сначала выбирает Skill по
краткому описанию, затем читает порядок действий и, наконец, вызывает
инструменты или выполняет скрипты в соответствии с задачей.

\textbf{Прогрессивное раскрытие состоит из трёх этапов.} При запуске
загружаются \texttt{name} и \texttt{description} каждого Skill, чтобы
модель знала, какие возможности доступны и для каких задач они
предназначены; после выбора читается полный файл \texttt{SKILL.md}; в
процессе выполнения дополнительно открываются указанные в нём документы,
шаблоны или скрипты\footnote{Agent Skills, спецификация формата, включая
  метаданные, основной текст, сопутствующие ресурсы и прогрессивную
  загрузку.}. Например, при работе с PDF сначала определяется, требуется
ли извлечь текст или заполнить форму, а затем читается соответствующая
инструкция. Так необходимые для конкретных операций подробности хранятся
в пакете возможностей и постепенно входят в контекст по мере выполнения
задачи. Общие затраты каталога зависят от количества Skills и длины их
описаний.

Этот процесс похож на работу со справочником: сначала читатель
просматривает оглавление, затем открывает нужную статью, а при
возникновении конкретных вопросов обращается к примерам. Для небольших
библиотек возможностей, состоящих главным образом из рабочих процедур и
справочных материалов, автор предпочитает каталог Skills в сочетании с
чтением файлов. Такой подход повторно использует способность модели
читать документацию и вызывать универсальные инструменты: на начальном
этапе достаточно поддерживать описания, ссылки и инструкции, не создавая
отдельный сервис поиска инструментов. Когда библиотека разрастается и
каталог становится трудно просматривать, можно добавить поиск по
ключевым словам или семантический поиск.

Специализированные инструменты также можно загружать по запросу. Хост
может на основе пространства имён сервера, списка разрешений и поиска
определять, какие определения доступны в текущей задаче, а затем
добавлять их в контекст через интерфейс поиска инструментов. Skills, в
свою очередь, объединяют шаги, примеры и справочные материалы в одной
файловой структуре, удобной для написания и сопровождения авторами.

Если одни и те же инструкции по заполнению PDF многократно используются
в разных сеансах, можно ли переиспользовать и вычисления модели,
выполненные при их чтении? Сначала следует различить две ситуации. В
пределах одного сеанса новые инструкции можно добавить в конец, и
совпадающий предшествующий префикс сможет попасть в кэш. В другом сеансе
даже при полном совпадении текста инструкций пользовательский запрос и
материалы перед инструкциями, а также положение самих инструкций обычно
отличаются, поэтому обычный префиксный кэш не позволяет непосредственно
переиспользовать вычисления для данного фрагмента. Рассмотренное во
второй главе исследование компонуемого KV-кэша предлагает отдельно
предварительно вычислить этот фрагмент, а затем скорректировать
вращательное позиционное кодирование (RoPE) и встроить его в новый
контекст. Вычислительные затраты на изменение позиции и объединение
линейно зависят от длины фрагмента.

MCP и Skill Hub также могут совместно обеспечивать распространение.
Skill Hub отвечает за поиск, установку и обновление пакетов
возможностей, а расширение Skills для MCP позволяет клиенту обнаруживать
и читать Skills, предоставляемые сервером\footnote{Model Context
  Protocol, \emph{Skills Over MCP Charter}, описание обязанностей
  расширения Skills по обнаружению, распространению и чтению.}. Один и
тот же набор практических знаний можно загрузить как локальный пакет
возможностей либо получить через MCP-сервис, поддерживающий это
расширение.

Инструмент уже найден --- но что именно следует вернуть модели, чтобы
результат оказался полезным? Вернёмся к примеру с письмом клиента:
одного имени файла достаточно, чтобы найти исходное письмо, однако для
проверки даты поставки необходимо прочитать исходное предложение. С
этого различия начинается проектирование инструментов восприятия.

\hypertarget{ux438ux43dux441ux442ux440ux443ux43cux435ux43dux442ux44b-ux432ux43eux441ux43fux440ux438ux44fux442ux438ux44f}{%
\section{Инструменты
восприятия}\label{ux438ux43dux441ux442ux440ux443ux43cux435ux43dux442ux44b-ux432ux43eux441ux43fux440ux438ux44fux442ux438ux44f}}

Предположим, пользователю нужно проверить итоговую сумму в отчёте.
Поисковый инструмент возвращает краткое описание найденного результата:
его достаточно, чтобы найти отчёт, но по нему не всегда можно определить
единицы измерения сумм. Только если модель перейдёт от описания к
исходной таблице и увидит названия столбцов, единицы измерения и строку
с итогом, она сможет выполнить проверку. При проектировании инструментов
восприятия здесь возникают два взаимосвязанных вопроса: какие материалы
вернуть в этот раз и где модель сможет продолжить чтение, если их
окажется недостаточно?

Один поиск может выдать десятки тысяч символов, а PDF-документ ---
насчитывать более ста страниц. Можно использовать описанное во второй
главе \textbf{сжатие с учётом контекста} и при превышении заданного
порога извлекать содержимое, относящееся к текущему вопросу. Например,
можно установить порог в 10000 символов, возвращать сводку и сохранять
возможность прочитать исходный результат, чтобы модель могла продолжить
проверку цитат, таблиц или деталей.

\textbf{Сначала перечислите результаты поиска, затем прочитайте полный
текст выбранной моделью записи.} Возвращаемый список должен содержать
заголовок, расположение и фрагмент описания, помогая AI-агенту
определить, какие результаты стоит прочитать. Если результатов много,
используйте пагинацию или курсор (\texttt{cursor}): по умолчанию
возвращайте несколько первых записей и указывайте способ получения
следующей страницы; если можно получить общее количество результатов,
также возвращайте его. На основе уже имеющихся доказательств модель
сможет решить, продолжать ли просмотр страниц или открыть один из
результатов и прочитать исходный текст.

\textbf{Читайте большие файлы фрагментами и при усечении указывайте
диапазон.} Интерфейс чтения должен поддерживать такие параметры, как
\texttt{offset} и \texttt{limit}, и пояснять, что используется в
качестве единицы измерения: строки, байты или записи. Например, вместе с
сообщением «Показаны строки 1--200 из 5000» следует вернуть параметры
для продолжения чтения. Информация о диапазоне помогает агенту понять,
является ли текущее содержимое полным, а также повторно обращаться к
нему по расположению. Неявное усечение опасно: модель может решить, что
уже прочитала весь материал, и принять вывод «в этом фрагменте ничего
нет» за вывод «во всём материале ничего нет». Если содержимое в середине
пропущено, следует указать место пропуска и способ чтения оставшейся
части.

\textbf{Используйте режим только для чтения для организации кэширования
и параллельного выполнения.} Результаты доступа к бизнес-данным только
для чтения обычно можно использовать повторно, а сами операции удобно
выполнять параллельно, например одновременно читать пять файлов или
запускать три независимых поисковых запроса. Ключ кэша должен включать
параметры запроса, идентификационные данные пользователя и область его
полномочий; для погоды, цен акций и результатов поиска также необходимо
задавать подходящий срок действия (TTL). При параллельном чтении нужно
ограничивать частоту запросов и учитывать влияние обновлений внешних
данных: для задач, которым требуется согласованное представление,
несколько операций чтения можно согласовывать с помощью снимка, номера
версии или времени получения данных.

\textbf{Выбирайте мультимодальный вывод в соответствии с задачей.}
Снимки экрана, диаграммы и сканированные документы можно напрямую
передавать модели с поддержкой визуального ввода либо сначала извлекать
из них текст с помощью OCR, анализа документов и других инструментов.
Извлечённый текст удобен для поиска, цитирования и сжатия; сохранение
изображения помогает понимать компоновку страницы, соответствия строк и
столбцов таблицы и расположение элементов интерфейса. При обработке
преимущественно текстового содержимого можно сначала возвращать
результат извлечения; если задача связана с UI, сложными таблицами или
дизайн-макетами, модель должна иметь возможность просмотреть
соответствующее изображение. Обе формы также можно использовать
совместно.

Инструменты восприятия также могут преобразовывать то, что человек
различает невооружённым глазом, в наблюдения, которые можно проверить.
Исследователи Coscientist поместили красный, жёлтый и синий образцы в
три известные лунки планшета и поручили помощнику определить, в какой
лунке находится каждый цвет. Планшетный считыватель вернул данные об
оптическом поглощении на разных длинах волн, а помощник с помощью Python
нашёл пики поглощения. На этом этапе было установлено соответствие между
длинами волн и лунками, но для присвоения цвета каждой лунке ещё
требовалось интерпретировать эти пики. В статье отмечена одна деталь:
помощник правильно распознал все три образца лишь после того, как
исследователь предложил ему подумать, как разные цвета поглощают
свет.\footnote{Boiko и др., \emph{Autonomous chemical research with
  large language models}, Nature, 2023. Схему эксперимента см. на рис. 4
  и в опубликованной авторами упрощённой реализации.}

\begin{quote}
\textbf{Эксперимент 4-2 ★★: MCP-сервер инструментов восприятия}

В этом эксперименте распространённые средства получения информации
подключаются к одному MCP-серверу, а интерфейсы организуются по пяти
типам сценариев:

\begin{itemize}
\tightlist
\item
  \textbf{Поиск и получение}: поиск в интернете, поиск в локальной базе
  знаний и скачивание файлов.
\item
  \textbf{Мультимодальное понимание}: чтение веб-страниц, извлечение
  данных из документов PDF, Word и PPT, OCR и визуальный анализ
  изображений, а также транскрибирование и анализ аудио и видео.
\item
  \textbf{Файловая система}: чтение файлов, поиск содержимого и просмотр
  каталогов; сервер также предоставляет операции копирования,
  перемещения и удаления, разрешения для которых управляются как для
  инструментов выполнения.
\item
  \textbf{Открытые источники данных}: погода, цены акций, валютные
  курсы, Wikipedia и статьи arXiv.
\item
  \textbf{Частные источники данных}: чтение календаря и содержимого
  Notion с разрешения пользователя.
\end{itemize}

В этом запуске были охвачены первые четыре типа сценариев. Сначала
рассмотрим чтение документов: инструмент открыл одностраничный PDF и
вернул количество страниц и извлечённый текст, который можно сверить с
исходной страницей. Затем рассмотрим запрос к календарю: он также вернул
ответ MCP, во внешней оболочке которого признак ошибки имел значение
false, но внутри были указаны \texttt{success:\ false} и
\texttt{missing\_credentials}. На основании этого помощник может
сообщить пользователю: «Разрешение на чтение календаря ещё не получено».
Ответ «На этот день ничего не запланировано» ошибочно представил бы
невыполненный запрос как результат запроса. Запрос к Notion также
остановился на этапе настройки учётных данных.
\end{quote}

\hypertarget{ux43cux443ux43bux44cux442ux438ux43cux43eux434ux430ux43bux44cux43dux43eux435-ux432ux43eux441ux43fux440ux438ux44fux442ux438ux435}{%
\subsection{Мультимодальное
восприятие}\label{ux43cux443ux43bux44cux442ux438ux43cux43eux434ux430ux43bux44cux43dux43eux435-ux432ux43eux441ux43fux440ux438ux44fux442ux438ux435}}

Зададим два вопроса по одному PDF с диаграммой: «Какое решение авторы
рекомендуют в заключении?» и «Какая кривая первой пересекает порог?»
Ответ на первый вопрос часто можно найти в извлечённом тексте, а для
второго также необходимы легенда, координаты и пространственные
отношения между кривыми. Изображения, аудио, видео и PDF можно напрямую
передавать основной модели, поддерживающей соответствующую модальность,
сначала преобразовывать в текст либо анализировать мультимодальным
инструментом, который модель вызывает по мере необходимости.

\hypertarget{ux432ux441ux442ux440ux43eux435ux43dux43dux430ux44f-ux43cux443ux43bux44cux442ux438ux43cux43eux434ux430ux43bux44cux43dux430ux44f-ux43eux431ux440ux430ux431ux43eux442ux43aux430}{%
\subsubsection{Встроенная мультимодальная
обработка}\label{ux432ux441ux442ux440ux43eux435ux43dux43dux430ux44f-ux43cux443ux43bux44cux442ux438ux43cux43eux434ux430ux43bux44cux43dux430ux44f-ux43eux431ux440ux430ux431ux43eux442ux43aux430}}

\textbf{Встроенная мультимодальная обработка} позволяет основной модели
непосредственно использовать визуальный, аудио- и другой ввод при ответе
на вопрос. Типичная архитектура визуально-языковых моделей, таких как
Qwen-VL и LLaVA, включает визуальный кодировщик, соединительный модуль и
языковую модель: визуальный кодировщик извлекает признаки из
изображения, соединительный модуль преобразует их в представление,
которое может принять языковая модель, после чего они вместе с текстовой
инструкцией участвуют в рассуждении.\footnote{Bai и др., \emph{Qwen-VL:
  A Versatile Vision-Language Model for Understanding, Localization,
  Text Reading, and Beyond}, 2023; Liu и др., \emph{Improved Baselines
  with Visual Instruction Tuning}, 2023. В работах соответственно
  описаны визуальный кодировщик и структура его соединения с языковой
  моделью.}

При использовании \textbf{Vision Transformer (ViT)} изображение
разделяется на блоки (patch), которые кодируются в визуальные признаки.
Затем модель получает эти признаки через проекцию, повторную выборку или
другой соединительный механизм и с помощью механизма внимания
обрабатывает связи между содержимым изображения и текстовым вопросом.
Например, чтобы ответить на вопрос «Какое число находится в правом
нижнем углу таблицы?», модель должна одновременно учитывать текст,
расположение и структуру таблицы.

Если задача требует неоднократно проверять связи между компоновкой,
диаграммами и текстом, автор предпочитает предоставлять исходный
материал непосредственно основной модели с соответствующими
возможностями. Визуальная информация продолжает участвовать в
последующем рассуждении, что снижает вероятность потери деталей при
многократном пересказе.

PDF обычно сначала разбирается слоем обработки ввода либо его страницы
визуализируются как изображения и передаются модели. Сохранение
изображений страниц помогает анализировать компоновку, диаграммы и
соответствия между изображениями и текстом; для длинных документов
необходимо контролировать количество страниц, разрешение и способ
разбиения, чтобы нужные для задачи детали попадали в доступный модели
диапазон.

\hypertarget{ux438ux437ux432ux43bux435ux447ux435ux43dux438ux435-ux432-ux442ux435ux43aux441ux442}{%
\subsubsection{Извлечение в
текст}\label{ux438ux437ux432ux43bux435ux447ux435ux43dux438ux435-ux432-ux442ux435ux43aux441ux442}}

\textbf{Извлечение в текст (Extract to Text)} разделяет обработку на два
этапа: сначала текст получают с помощью анализа документов, OCR или
инструмента транскрибирования аудио, а затем языковая модель выполняет
поиск, обобщение и отвечает на вопросы. Текстовые модели, такие как GLM
5.2 и DeepSeek V4 Flash, могут таким способом обрабатывать текстовую
информацию из документов, аудио и видео.\footnote{Z.AI, документация
  модели GLM-5.2, входная модальность --- текст; DeepSeek, карточка
  модели DeepSeek-V4-Flash, описывающая модель генерации текста и
  способы её использования.}

Для PDF, состоящих преимущественно из текста, непосредственное
извлечение текста обычно компактнее постраничной передачи снимков
страниц, а также удобнее для поиска и цитирования. Количество токенов
для изображения зависит от разрешения, разбиения и способа кодирования
модели; объём текстового ввода --- от извлечённого содержимого и способа
токенизации.

Результат извлечения должен сохранять связи, от которых зависит задача.
Обычный OCR может распознать числа на диаграмме, но потерять их
соответствие подписям осей; чтение только текстового слоя PDF также
может пропустить диаграммы во встроенных изображениях. Анализ компоновки
и разбор таблиц позволяют дополнительно сохранить порядок абзацев,
структуру строк и столбцов или маркеры расположения.

\hypertarget{ux43cux443ux43bux44cux442ux438ux43cux43eux434ux430ux43bux44cux43dux44bux439-ux430ux43dux430ux43bux438ux437-ux43aux430ux43a-ux438ux43dux441ux442ux440ux443ux43cux435ux43dux442}{%
\subsubsection{Мультимодальный анализ как
инструмент}\label{ux43cux443ux43bux44cux442ux438ux43cux43eux434ux430ux43bux44cux43dux44bux439-ux430ux43dux430ux43bux438ux437-ux43aux430ux43a-ux438ux43dux441ux442ux440ux443ux43cux435ux43dux442}}

\textbf{Мультимодальный анализ как инструмент} позволяет основной модели
вызывать средства анализа изображений, аудио или документов в
зависимости от текущего вопроса. Например, \texttt{analyze\_image},
\texttt{analyze\_pdf} и \texttt{analyze\_audio} принимают ссылку на файл
и вопрос на естественном языке, а возвращают результат анализа. Основная
модель решает, когда вызвать инструмент, о чём спросить и как
использовать результат; модель внутри инструмента сосредоточена на
обработке соответствующей модальности.

Если основная модель поддерживает только текст, а задача зависит от
визуальных деталей, автор предпочитает дополнять её такими инструментами
анализа.

Например, агент сначала читает основной текст отчёта и обнаруживает, что
квартальная выручка, о которой спрашивает пользователь, приведена только
на диаграмме. Тогда он вызывает инструмент анализа изображений с
вопросом: «Какова выручка в каждом квартале?» Получив значения и
соответствующие им кварталы, основная модель продолжает вычислять рост
или разницу. В контексте основной модели достаточно сохранить вопрос и
результат анализа, а исходное изображение обрабатывается инструментом.

\begin{quote}
\textbf{Эксперимент 4-3 ★★: извлечение мультимодальной информации ---
сравнительный анализ трёх технических парадигм}

В этом эксперименте используются столбчатая диаграмма квартальной
выручки и содержащий её PDF-отчёт, чтобы сравнить три способа передачи
информации. Выручка за четыре квартала на диаграмме последовательно
составляет 120, 150, 95 и 180 миллионов долларов; в основном тексте
отчёта описана лишь динамика показателя. Для изображения и PDF задаются
два вопроса: «В каком квартале выручка была максимальной и сколько
именно она составила?» и «В каком квартале выручка была минимальной,
сколько она составила и насколько выручка четвёртого квартала была
выше?»

В режиме \textbf{непосредственного просмотра изображения} исходное
изображение или изображение страницы PDF передаётся модели вместе с
вопросом. В режиме \textbf{извлечения в текст} ответы формируются
соответственно по результатам OCR и извлечения текста из PDF. В режиме
\textbf{вызова инструмента по мере необходимости} модель сначала читает
извлечённый текст, обнаруживает, что чисел или соответствий
недостаточно, вызывает \texttt{inspect\_visual} для просмотра исходного
изображения, а затем формирует ответ. Во всех трёх путях используется
одна и та же модель Doubao, что позволяет сравнить влияние способа
передачи информации на ответы.

Во всех четырёх комбинациях «материал --- вопрос» этого запуска
непосредственный просмотр изображения и вызов инструмента по мере
необходимости дали правильные ответы, а группа извлечения текста ни разу
не ответила полностью правильно. Один из ответов легко ошибочно признать
правильным: модель, получившая только результат OCR изображения,
ответила: «Минимальная выручка была в первом квартале и составила 95
миллионов долларов; в четвёртом квартале она была на 85 миллионов
долларов выше». Сумма и разница верны, но квартал указан неправильно. В
переданном модели тексте суммы над столбцами смешались с делениями
вертикальной оси, а подписи кварталов находились отдельно в последней
строке. Из такого текста можно извлечь два числа, но невозможно
определить, какой квартал соответствует какому столбцу. Только просмотр
исходного изображения визуальным инструментом позволил сопоставить 95
миллионов долларов с третьим кварталом.

При переходе к PDF с той же диаграммой причина ошибки изменилась. В
основном тексте PDF упоминались спад в середине года и восстановление в
последнем квартале, но все точные суммы находились во встроенном
изображении. На вопрос о максимальной выручке модель, читавшая только
основной текст, указала четвёртый квартал, но также сообщила, что не
может назвать точное значение. В такой ситуации наиболее полезный
следующий шаг --- передать визуальному инструменту страницу с
диаграммой.

При вызове визуального инструмента по мере необходимости добавился ещё
один раунд запроса и возврата результата, поэтому задержка в этом
запуске также увеличилась.
\end{quote}

\hypertarget{ux438ux43dux441ux442ux440ux443ux43cux435ux43dux442ux44b-ux432ux44bux43fux43eux43bux43dux435ux43dux438ux44f}{%
\section{Инструменты
выполнения}\label{ux438ux43dux441ux442ux440ux443ux43cux435ux43dux442ux44b-ux432ux44bux43fux43eux43bux43dux435ux43dux438ux44f}}

Инструменты выполнения позволяют AI-агенту записывать файлы, запускать
программы и изменять внешние системы. Инструменты восприятия
ориентированы на получение полезной информации, тогда как инструменты
выполнения должны ещё и управлять изменениями, вызванными операциями.
При проектировании необходимо одновременно ответить на три вопроса:
соответствует ли операция полномочиям, предоставленным для задачи, как
ограничить область её выполнения и как подтвердить результат после
завершения. Ошибочное удаление файлов, остановка сервисов, повторная
отправка запросов и другие подобные ошибки связаны именно с этими
границами.

Рассмотрим изменение адреса доставки в существующем заказе. Получение
нового адреса от пользователя, отправка изменения и повторное чтение
нового адреса из заказа --- три наблюдаемых момента. Если после отправки
соединение прервалось, клиент знает лишь, что не получил подтверждение,
но не может на этом основании заключить, что сервер ничего не записал.
После восстановления соединения можно сначала запросить текущее
состояние заказа и лишь затем решать, следует ли повторять операцию.
Если при повторном чтении по-прежнему отображается старый адрес,
необходимо выяснить, не применяет ли сервис изменения с задержкой. При
наличии идентификатора запроса следует прежде всего проверить по нему
состояние именно этого изменения. Повторная отправка только на основании
одного чтения всё ещё может привести к двукратному выполнению одной
операции.

Иногда для обнаружения полезной задачи и её фактического выполнения
требуются разные полномочия. При выпуске \textbf{dots} OpenAI рассказала
об одном из первых пользователей: его dot заметил, что пользователь
забыл выставить счёт одному изданию, подготовил счёт и отправил его
только после согласия пользователя. Согласно описанию продукта, когда
пользователь не поручал новую работу, dot может самостоятельно искать
полезные сведения в подключённых приложениях. На этом этапе инструменты
работают только на чтение: они не могут отправлять сообщения, изменять
содержимое приложений или управлять браузером.\footnote{OpenAI,
  \emph{Introducing dots}, опубликовано 29 сентября 2026 года; пример со
  счётом и ограничения фонового самостоятельного поиска режимом только
  для чтения приведены в описании выпуска.}

\textbf{Многоуровневая проверка параметров и прав доступа.}

Сначала проверяются входные данные. Файловый инструмент должен разобрать
целевой путь, убедиться, что фактическое место обращения находится в
разрешённой области, и учесть символические ссылки и другие механизмы,
способные изменить фактическое назначение пути. Инструмент команд должен
различать программу, аргументы и скрипт, требующий интерпретации.
Внешний API должен проверять типы, форматы и диапазоны значений полей.
Если параметры недопустимы, следует вернуть конкретное поле и причину
ошибки, чтобы модель могла исправить запрос и отправить его заново.

Затем проверяются права доступа. Доступ к файлам следует ограничивать
разрешёнными областями чтения и записи, процессы запускать от имени
учётной записи с ограниченными правами, а для внешних API использовать
полномочия и учётные данные, соответствующие задаче. Квоты и ограничения
частоты запросов сдерживают потребление ресурсов. Чёрный список команд
позволяет блокировать известные опасные шаблоны, а разбор команд и
проверка моделью помогают определить смысл операции. Базовые права
файловой системы, сети и процессов ограничивают фактическую область
воздействия. В пятой главе подробнее рассматривается совместное
применение этих механизмов.

\textbf{Предлагающий --- проверяющий: независимая проверка плана
действий и результата выполнения.}

Описанную в первой главе \textbf{парадигму «предлагающий ---
проверяющий» (Proposer--Reviewer)} можно применять и к инструментам
выполнения: предлагающий формирует план действий, а проверяющий
оценивает его относительно цели задачи, предоставленных полномочий и
критериев приёмки. В зависимости от момента проверки различают
предварительное одобрение и последующую проверку.

\textbf{Предварительное одобрение} выполняется до отправки операции.
Например, перед отправкой электронного письма проверяющий сверяет
получателей, текст, вложения и область рассылки с требованиями
пользователя; перед изменением критически важной конфигурации ---
целевую среду, содержимое изменения и план восстановления. На основании
результата проверки Harness разрешает выполнение, требует исправлений
или передаёт решение человеку.

Для сложных проверок автор предпочитает выбирать модели с близкими
возможностями, но обученные в разных условиях: например, перекрёстно
проверять семейства GPT и Claude либо Kimi и DeepSeek. Сопоставимые
возможности позволяют проверяющему понять сложный план, а различия в
обучении дают шанс сократить число общих слепых зон. Модели одного
семейства могут исходить из сходных предположений, поэтому обе проверки
пропустят одну и ту же проблему. Если же возможности моделей различаются
слишком сильно, проверяющий может не заметить ключевую ошибку в плане.
Требуемый уровень возможностей определяется задачей проверки: для
понимания сложного плана необходимы развитые способности к рассуждению,
а проверку простых полей по явным правилам можно поручить модели
меньшего размера.

Обе стороны должны исходить из одной и той же цели задачи, одинаковых
границ полномочий и критериев приёмки, а также получать доказательства,
необходимые для принятия решения. Промпты могут распределять
обязанности: предлагающий организует действия, а проверяющий уделяет
особое внимание рискам, упущениям и соблюдению правил.

При отказе конкретную причину следует добавить в траекторию как
результат инструмента, например: «Получатель не входит в подтверждённую
область рассылки; проверьте список адресатов». На основании этого модель
может изменить план. После последовательных отказов или при неясных
полномочиях решение передаётся пользователю. Для списания средств,
отправки электронной почты, изменения критически важных конфигураций,
создания внешних ресурсов и других операций можно установить уровни
одобрения в зависимости от области воздействия. Обычные операции с
чёткими границами и возможностью восстановления можно обрабатывать по
заранее настроенным правилам.

\textbf{Последующая проверка} оценивает созданные результаты и состояние
системы после выполнения. Автор рекомендует прежде всего менять способ
проверки, позволяя проверяющему непосредственно исследовать результат
работы или его итоговое представление. При повторном чтении одного и
того же кода легко сохранить предположения, использованные при его
создании. Переход к наблюдаемому, исполняемому результату помогает
выявить ошибки, незаметные в самом коде. Например, после программного
создания документа следует визуализировать страницы и проверить вёрстку;
после изменения конфигурации --- запустить программу в изолированной
среде и убедиться, что изменение вступило в силу; после записи данных
--- повторно прочитать целевую запись и сверить содержимое. В пятой
главе этот цикл будет показан на примере презентации: предлагающий
генерирует код, проверяющий изучает снимки визуализированного результата
и сообщает, что необходимо исправить.

\textbf{Sidecar: модуль проверки, сопровождающий основной цикл.}

Название Sidecar заимствует метафору «коляски» из микросервисной
архитектуры и обозначает вспомогательный модуль, который сопровождает
работу основного агента и выполняет специализированные задачи. При
проверке безопасности он оценивает конкретное действие до вызова
инструмента; при обработке контекста может параллельно отбирать сведения
из памяти или сжимать результаты инструментов.

Auto Mode в Claude Code автоматически определяет, можно ли разрешить
конкретный вызов инструмента. В реализации, представленной Anthropic в
марте 2026 года, часть вызовов сначала обрабатывается обычными
разрешающими правилами, а остальные передаются модельному
классификатору. Классификатор сначала выполняет быструю оценку, а затем
подробнее анализирует заблокированные вызовы. Он читает сообщения
пользователя и вызовы инструментов, исключая пояснительный текст
основной модели и результаты инструментов, чтобы сопоставлять намерение
пользователя с фактическим действием.\footnote{Anthropic, \emph{How we
  built Claude Code auto mode: a safer way to skip permissions}, 25
  марта 2026 года; описаны многоуровневая система прав, классификатор
  вызовов и способ изоляции входных данных.}

При проектировании собственного модуля проверки ему можно передавать
название инструмента, целевой ресурс, параметры, тип операции и
относящиеся к ней полномочия. Например, перед удалением файла следует
проверить, находится ли цель в разрешённой области, относится ли она к
результатам, подлежащим очистке, и требуется ли дополнительное
подтверждение.

Проверка может выполняться параллельно с генерацией основной моделью
последующего содержимого или обновлением интерфейса, сокращая
дополнительное ожидание, однако проверяемый вызов обязан дождаться
разрешения. Например, пользователь просит «посмотреть конфигурацию», а
основная модель готовится записать файл обратно. Путь может быть
совершенно допустимым, но смысл операции изменился с просмотра на
редактирование. При проектировании такой проверки можно поручить
модели-классификатору вроде Jev сопоставить исходные слова пользователя
с параметрами инструмента и различить варианты «соответствует запросу»,
«выходит за рамки запроса» и «требует дополнительного подтверждения».
Результат классификации используется процессом одобрения, а права
доступа к файлам по-прежнему проверяются уровнем выполнения при каждой
операции.

Парадигма «предлагающий --- проверяющий» описывает распределение
обязанностей, а Sidecar --- связь вспомогательного модуля с основным
циклом во время выполнения. Функции проверяющего может выполнять
Sidecar, однако Sidecar способен решать и другие задачи. В таблице 4-2
эти подходы сопоставлены на примере типичных применений, рассмотренных в
этом разделе.

Таблица 4-2. Сравнение механизма «предлагающий --- проверяющий» и
механизма Sidecar

\begin{longtable}[]{@{}
  >{\raggedright\arraybackslash}p{(\columnwidth - 4\tabcolsep) * \real{0.2374}}
  >{\raggedright\arraybackslash}p{(\columnwidth - 4\tabcolsep) * \real{0.3592}}
  >{\raggedright\arraybackslash}p{(\columnwidth - 4\tabcolsep) * \real{0.4034}}@{}}
\toprule\noalign{}
\begin{minipage}[b]{\linewidth}\raggedright
\hspace{0pt}Аспект
\end{minipage} & \begin{minipage}[b]{\linewidth}\raggedright
\hspace{0pt}Предлагающий --- проверяющий
\end{minipage} & \begin{minipage}[b]{\linewidth}\raggedright
\hspace{0pt}Sidecar
\end{minipage} \\
\midrule\noalign{}
\endhead
\bottomrule\noalign{}
\endlastfoot
\hspace{0pt}Основная задача & \hspace{0pt}Независимая проверка
предложения или результата & \hspace{0pt}Выполнение специализированных
задач совместно с основным циклом \\
\hspace{0pt}Время выполнения & \hspace{0pt}Одобрение до действия или
проверка после действия & \hspace{0pt}Запуск по мере необходимости, в
том числе параллельно основному циклу \\
\hspace{0pt}Объект проверки & \hspace{0pt}План, операция и её результат
& \hspace{0pt}Вызовы инструментов, память, вывод инструментов и
другое \\
\hspace{0pt}Организация входных данных & \hspace{0pt}Общие критерии
задачи и доказательства, необходимые для независимой оценки &
\hspace{0pt}Выбор необходимых полей и контекста в соответствии с
обязанностями \\
\hspace{0pt}Типичные применения & \hspace{0pt}Одобрение критически
важных операций, проверка качества документов, проверка конфигураций &
\hspace{0pt}Шлюзовая проверка вызовов, отбор сведений из памяти,
суммаризация вывода \\
\end{longtable}

При дополнении контекста Sidecar может одновременно с обработкой
текущего шага основной моделью отбирать релевантные сведения из памяти
пользователя, обобщать длинный вывод инструментов или запрашивать
данные, необходимые на последующих шагах. Основной цикл получает
результат в подходящий момент, сокращая последовательное ожидание.

\textbf{Автоматическая проверка и замкнутый контур обратной связи.}

\textbf{Автор рекомендует включать в процесс инструмента проверку всех
результатов операций, которые можно проверить автоматически.} Например,
после записи кода с помощью \texttt{write\_file} можно вызвать
соответствующий языку синтаксический анализатор или linter,
преобразовать диагностические сообщения в сведения о файле, номере
строки, типе и описании ошибки и вернуть их вместе с состоянием записи.
Увидев сообщение «строка 10: переменная \texttt{result} не определена»,
агент сможет найти и исправить проблему.

Инструмент должен раздельно сообщать состояние операции и состояние
проверки: файл уже может быть записан, хотя проверка завершилась с
ошибкой. Синтаксическая проверка обнаруживает код, который невозможно
разобрать; статическая проверка дополнительно выявляет некоторые ошибки
имён, типов и правил. Для проверки функциональной корректности требуются
соответствующие тесты или результаты выполнения.

\textbf{Усечение и сохранение длинного вывода.}

Журналы терминала и вывод программ могут быть очень длинными. Можно
установить порог, например 200 строк или 10000 символов, после
превышения которого полный результат сохраняется как доступный для
последующего чтения артефакт, а в контекст возвращаются лишь фрагменты
из начала и конца:

\begin{itemize}
\tightlist
\item
  Сохранить первые 50 строк с информацией о запуске и исходным
  контекстом.
\item
  Сохранить последние 50 строк с итоговым состоянием или последней
  ошибкой.
\item
  Указать объём пропущенного содержимого, например: «В середине
  пропущено 8523 строки; полный вывод сохранён».
\item
  Вернуть ссылку на артефакт и способ чтения, чтобы агент мог продолжить
  просмотр или поиск с помощью \texttt{read\_file}.
\end{itemize}

Каждому вызову следует присваивать отдельный идентификатор артефакта,
чтобы последующий вывод не перезаписывал предыдущую запись. Возвращаемое
инструментом значение должно одновременно содержать статус выхода,
признак усечения и местонахождение артефакта, позволяя сопоставить
текущую сводку с полной записью.

\textbf{Изоляция и песочница среды выполнения.}

Универсальные инструменты, такие как интерпретатор Python и Shell,
позволяют выполнять программы, созданные моделью. Область их доступа
совместно определяется учётной записью процесса, файловой системой и
сетевыми политиками. Виртуальное окружение Python (venv) изолирует
зависимости пакетов, но код по-прежнему использует права процесса на
доступ к файлам, сети и системе. Для ограничения этих возможностей
необходимы механизмы операционной системы или виртуализации.

Распространённые способы развёртывания:

\begin{itemize}
\tightlist
\item
  \textbf{Ограниченный локальный процесс}: выполняется на локальной
  машине, а область доступа контролируется ограниченной учётной записью,
  правилами чтения и записи рабочего пространства и сетевыми политиками.
  Локальное выполнение Coding Agent обычно относится к этому типу
  развёртывания; конкретные права определяются конфигурацией хоста.
\item
  \textbf{Контейнер}: пространства имён изолируют процессы, файловую
  систему и представление сети, а средства управления ресурсами
  ограничивают их потребление. Контейнер использует общее с хостом ядро,
  поэтому фактические границы также зависят от подключённых каталогов,
  прав процессов и конфигурации контейнера.\footnote{Docker,
    \emph{Docker Engine security}; описано влияние пространств имён,
    управления ресурсами, ядра и прав процессов на изоляцию контейнеров.}
\item
  \textbf{microVM или виртуальная машина}: виртуализация предоставляет
  задаче отдельное ядро. Такие microVM, как Firecracker, подходят для
  сценариев, требующих более строгой изоляции при одновременном учёте
  времени запуска и затрат ресурсов.\footnote{Firecracker, Design;
    описана архитектура microVM на основе KVM.}
\end{itemize}

При локальной разработке, доверенных входных данных и чётких границах
задачи можно начать с ограниченного процесса и правил рабочего
пространства. Если требуется выполнять код неизвестного происхождения
или работать с чувствительными учётными данными, уровень изоляции
следует повысить, отдавая предпочтение контейнеру или microVM с
отдельным ядром и ограничивая раскрытие учётных данных. Если задача
требует только чтения материалов, соответствующие каталоги можно
подключить в режиме только для чтения. Для записи результатов следует
предоставить отдельное пространство вывода. Среда выполнения также
должна ограничивать CPU, память, диск, число процессов и доступ к сети,
а также устанавливать максимальное время работы, чтобы неконтролируемую
задачу можно было остановить и освободить её ресурсы.

После выбора способа изоляции возникает вопрос стоимости: для кого
предназначена эта среда и как долго её сохранять? Когда автор создавал
своего первого агента, один сервис на Go обслуживал параллельные запросы
нескольких пользователей. Для доступа к веб-страницам он вызывал
безголовый браузер, а загруженные файлы можно было хранить в GCS или
Google Drive и извлекать только при необходимости. По тогдашней оценке
автора, постоянно работающая облачная виртуальная машина для каждого
пользователя могла обходиться в десятки долларов за песочницу
ежемесячно. Даже если пользователь обращался к помощнику лишь один раз
за месяц, среда всё это время ожидала его.

Эти затраты можно рассчитать, разделив поручение на этапы. Предположим,
сегодня пользователь загрузил счёт, а попросил помощника проверить его
только через неделю. В течение недели необходимо сохранять лишь счёт и
запись о задаче. Если для проверки потребуется выполнить новый
сгенерированный скрипт, изолированную среду можно выделить только тогда,
предоставить ей файлы, разрешённые для чтения в рамках этой операции,
сохранить результат и освободить среду. Предположим, тарификация ведётся
по времени работы экземпляра, характеристики машины одинаковы, а сама
проверка занимает 10 минут. Если та же машина семь дней простаивает, а
затем выполняет проверку, придётся оплатить 168 часов и 10 минут
вычислений. Если же запуск по требованию занимает 2 минуты на подготовку
среды, суммарное оплачиваемое время составит 12 минут.

Если через неделю пользователь хочет продолжить работу с уже открытым
рабочим столом, одного восстановления файлов может оказаться
недостаточно. Сколько времени потребуется, чтобы вернуть приложение на
прежнюю страницу? Можно ли восстановить несохранённую работу? Только
ответив на эти два вопроса, можно определить, стоит ли сохранять среду.
Это различие снова понадобится в пятой главе при обсуждении сеансов
командной строки.

\textbf{Наблюдаемость выполнения инструментов.}

Записи выполнения должны отвечать на вопрос: «Кто инициировал какую
операцию, на какие ресурсы она повлияла и каков результат?» Для каждого
вызова следует записывать ID вызова, время, название инструмента,
необходимые параметры, целевые ресурсы, состояние выполнения и
длительность, а также связывать его с задачей и записями об одобрении.
Учётные данные и чувствительное содержимое необходимо маскировать либо
ограничивать доступ к ним. Агрегированные показатели частоты вызовов,
доли успешных вызовов, тайм-аутов и длительности позволяют выявлять
медленные инструменты, повторяющиеся сбои и узкие места ресурсов. При
достижении аномалиями заданного порога можно отправлять оповещение или
приостанавливать соответствующее выполнение.

\textbf{Идемпотентность и семантика отмены.}

Предположим, помощник отправил заказ, продавец уже создал его, но
подтверждение потерялось при передаче. Помощник видит только тайм-аут
запроса. Если он отправит тот же заказ ещё раз, как продавцу понять, что
пользователь хотел купить только один экземпляр?

Один из подходов --- сохранить перед первой отправкой \textbf{ключ
идемпотентности}. При повторной попытке оформить тот же заказ для того
же пользователя помощник использует прежний ключ и прежние параметры.
Сервер связывает ключ с типом операции, параметрами и состоянием
обработки, атомарно регистрирует запрос и устраняет дубликаты. Поэтому
даже при получении двух запросов создаётся только один заказ. Это и есть
\textbf{идемпотентность}. При этом каждый сетевой запрос может иметь
собственный ID вызова, позволяющий выяснить, какой запрос завершился по
тайм-ауту, а какой получил подтверждение.

При интеграции внешнего инструмента необходимо внимательно изучить ещё
два соглашения: что произойдёт, если один ключ используется с разными
параметрами, и как долго сервер хранит ключ? Например, Stripe API v1
отклоняет повторный запрос, если его параметры изменились. После
хранения в течение не менее 24 часов ключ может быть удалён; если старый
ключ уже удалён, новый запрос с ним будет считаться новым. Помощник,
возобновляющий задачу через несколько дней, должен сначала проверить
предыдущую операцию по бизнес-идентификатору.\footnote{Stripe,
  \emph{Idempotent requests}; описаны ключи идемпотентности, повторное
  использование результата запроса, согласованность параметров и срок
  хранения.}

Перед повторной попыткой можно также запросить состояние цели: например,
найти заказ по бизнес-идентификатору или прочитать файл, чтобы
убедиться, что ожидаемое содержимое уже записано. Параллельные запросы
могут одновременно инициировать изменение после проверки, поэтому для
предотвращения повторных операций дополнительно требуются ограничения
уникальности, условные обновления или серверная поддержка
идемпотентности.

Отправка электронной почты, телефонные звонки, денежные переводы и
другие операции воздействуют на внешний мир, поэтому для них можно
применять процесс «предварительная проверка --- подтверждение ---
отправка»: сначала проверить цель, содержимое, сумму и другие
бизнес-условия, затем получить подтверждение в соответствии с
требованиями к полномочиям и только после этого отправить операцию.
Одобрение проверяет соответствие действия требованиям, а идемпотентность
и запрос состояния обеспечивают обработку повторных попыток после
отправки. Если сервер временно не может подтвердить результат, задачу
следует перевести в состояние «результат ожидает подтверждения»,
сохранить идентификатор запроса и продолжить проверку либо передать её
человеку.

Для отмены также необходимы явные соглашения. Если задача ещё не
началась, её можно удалить из очереди. Выполняемый локальный процесс
можно завершить и освободить его ресурсы. Если внешний сервис уже
получил запрос, необходимо выяснить, поддерживает ли он отмену, откат
или компенсирующую операцию. Инструмент должен возвращать текущее
состояние, чтобы агент мог различать «ожидание прекращено», «выполнение
уже прекращено» и «внешняя операция выполнена».

\begin{quote}
\textbf{Эксперимент 4-4 ★★: MCP-сервер инструментов выполнения}

В этом эксперименте операции с файлами, программами, внешними системами
и графическим интерфейсом подключаются к MCP, а для вызовов каждой
категории организуются проверка и обратная связь:

\begin{itemize}
\tightlist
\item
  \textbf{Запись и редактирование файлов}: проверка синтаксиса Python и
  JavaScript, возврат структурированных диагностических сообщений,
  проверка результата редактирования и границ рабочего пространства.
\item
  \textbf{Выполнение в терминале}: запись команды и результата,
  установка тайм-аута и вызов независимой модели для операций, требующих
  проверки.
\item
  \textbf{Интерпретатор кода}: выполнение Python в контейнере с
  ограничениями сети, файловой системы и потребления ресурсов.
\item
  \textbf{Обработка длинного вывода}: возврат усечённых фрагментов,
  долговременное сохранение полного вывода и предоставление способа
  последующего чтения.
\item
  \textbf{Операции с данными}: чтение и запись Excel, применение формул
  и создание снимков визуализированного результата для проверки его
  отображения.
\item
  \textbf{Внешние системы}: вызов Webhook, создание GitHub PR, а также
  интерфейсы для событий календаря и отправки электронной почты.
\item
  \textbf{Графический интерфейс}: навигация и снимки экрана в браузере,
  управление клавиатурой на виртуальном рабочем столе и управление
  интерфейсом эмулятора Android.
\end{itemize}

В примере для рабочего стола веб-страница открывается с помощью
клавиатуры, после чего создаётся снимок экрана. В примере для телефона
открываются настройки Wi-Fi, проверяется текущий экран, а затем
выполняется возврат на главный экран.

Рассмотрим операцию с таблицей. Инструмент умножает количество на цену
за единицу: две суммы составляют 25 и 21, а итог --- 46. В
экспортированном PDF на первой странице перечислены позиции, количество
и цена за единицу, но сумма видна только на второй странице. Инструмент
сделал снимок лишь первой страницы, поэтому рядом с надписью «Grand
total» на изображении нет числа. Если задача требует предоставить счёт,
на котором итоговую сумму можно проверить с первого взгляда, после
правильного расчёта формул необходимо ещё настроить макет печати и
повторно просмотреть предварительный результат.
\end{quote}

\hypertarget{ux438ux43dux441ux442ux440ux443ux43cux435ux43dux442ux44b-ux441ux43eux432ux43cux435ux441ux442ux43dux43eux439-ux440ux430ux431ux43eux442ux44b}{%
\section{Инструменты совместной
работы}\label{ux438ux43dux441ux442ux440ux443ux43cux435ux43dux442ux44b-ux441ux43eux432ux43cux435ux441ux442ux43dux43eux439-ux440ux430ux431ux43eux442ux44b}}

Как текущему AI-агенту передать работу, когда требуется поручить другому
агенту проверку материалов или возникает бизнес-вопрос, решение по
которому должен принять человек? Инструменты совместной работы должны не
только объяснять исполнителю, что нужно сделать, но и определять, какой
результат он должен вернуть. Например, при делегировании поиска следует
указать, какие материалы нужно найти; при делегировании проверки ---
приложить проверяемое решение; при обращении к человеку --- перечислить
вопросы, требующие решения.

Когда несколько агентов приступают к работе одновременно, распределение
задач должно предусматривать и распределение вычислительных ресурсов.
Автор наблюдал на своих машинах два вида сбоев: параллельные задачи по
отдельности создавали полные копии кода и за один день заполняли диск
объёмом 3 ТБ; другая группа задач занимала все ресурсы сервера с сотней
ядер, из-за чего даже человеку было трудно войти в систему. Для каждой
задачи можно создать отдельное рабочее пространство, хранящее только
различия, после завершения удалить временные файлы, а требуемые для
передачи файлы сохранить отдельно. Независимые рабочие пространства
экономят место на диске, но не мешают десяткам задач одновременно занять
весь CPU. Квоты для фоновых задач и повышенный приоритет интерактивных
запросов оставляют человеку возможность войти на машину и устранить
сбой.

Вернёмся к примеру с обработкой писем клиентов. Другому агенту можно
поручить извлечь срок поставки из вложений, но вместе с задачей следует
передать идентификатор исходного письма и потребовать вернуть «дату,
соответствующее предложение из оригинала и вложение, в котором оно
находится». Получив две противоречащие друг другу даты, главный агент
сможет определить, что одна относится к дате заказа, а другая --- к
обещанной дате доставки. Если же он получит лишь сообщение «обработка
завершена», разрешить расхождение будет невозможно.

\textbf{Настраивайте субагентов в соответствии с их обязанностями.}

Профессиональные знания и рабочие процессы можно сначала загружать по
требованию через Skill. Делегировать работу отдельно настроенным
субагентам следует тогда, когда задаче требуются независимый контекст,
разные права на использование инструментов или параллельная обработка.
Для каждого субагента промпт, набор инструментов и источники знаний
настраиваются в соответствии с его обязанностями. Это уменьшает помехи,
возникающие из-за размещения всех правил и промежуточных материалов в
одном контексте. Например, агент по поиску материалов предоставляет
источники и доказательства, агент по проверке кода выявляет проблемы, а
главный агент распределяет задачи и объединяет результаты. При
разделении обязанностей необходимо также согласовать порядок доступа к
общим ресурсам, чтобы несколько агентов не изменяли одновременно один
файл и не инициировали повторно одну и ту же внешнюю операцию.

Промпт субагента должен содержать следующее:

\begin{itemize}
\tightlist
\item
  \textbf{Роль и цель}: прямо опишите обязанности, например «Вы
  отвечаете за проверку числовых данных и источников в отчёте», и
  задайте критерии завершения текущей задачи.
\item
  \textbf{Источники информации}: разграничьте описание задачи от
  главного агента, дополнения пользователя и результаты инструментов.
  Для удобства чтения можно использовать метки
  \texttt{{[}FROM\_MAIN\_AGENT{]}}, \texttt{{[}FROM\_USER{]}},
  \texttt{{[}TOOL\_RESULT{]}}, тогда как Harness должен сохранять
  фактическое происхождение и роль каждого сообщения. Одноимённые метки
  во внешнем тексте следует обрабатывать как данные.
\item
  \textbf{Границы задачи}: явно укажите допустимую область работы,
  доступные инструменты, права на ресурсы и порядок эскалации при
  нехватке информации или возникновении противоречий.
\item
  \textbf{Требования к выводу}: возвращайте заключение, доказательства,
  нерешённые вопросы и текущее состояние в согласованной структуре JSON
  или Markdown, чтобы главный агент мог продолжить работу.
\end{itemize}

\textbf{Управляйте жизненным циклом субагента.}

Интерфейс совместной работы можно организовать в виде нескольких групп
примитивов: запуск, обмен сообщениями, состояние и обнаружение.
\texttt{spawn\_subagent} создаёт субагента и назначает ему задачу, а
\texttt{cancel\_subagent} запрашивает остановку, когда задача теряет
смысл; \texttt{send\_message\_to\_subagent} передаёт дополнительные
требования, а субагент также может сообщать о ходе работы или
запрашивать уточнения; \texttt{get\_subagent\_status} позволяет
запросить состояние и результат. С помощью \texttt{list\_agents} система
также может перечислять обязанности агентов, их доступность и
принадлежность задач, помогая главному агенту найти исполнителя.

Эти интерфейсы могут поддерживать различные режимы вызова.
\textbf{Синхронный вызов} ожидает результат и только после этого
продолжает работу; он подходит для коротких задач или случаев, когда
следующий шаг зависит от результата. \textbf{Асинхронный вызов} сначала
возвращает ID задачи, после чего главный агент может заняться другой
работой, а затем получить результат посредством запроса состояния или
события завершения. \textbf{Потоковое взаимодействие} непрерывно
возвращает инкрементальные сообщения, позволяя заранее использовать
промежуточные результаты. \textbf{Многораундовое взаимодействие} даёт
обеим сторонам возможность задавать уточняющие вопросы и дополнять
условия.

ID задачи связывает создание, сообщения, результат и отмену. Harness
должен управлять такими состояниями, как ожидание в очереди, выполнение,
завершение, сбой и отмена, явно определять, в каких состояниях ещё можно
принимать сообщения и как после отмены обрабатывать уже выполненные
внешние операции. Передача контекста, топология взаимодействия и
распределение обязанностей между несколькими агентами рассматриваются в
\crossreflink{chap:10}{главе 10}.

\textbf{Запрашивайте вмешательство, когда требуется решение человека.}

HITL (Human-in-the-Loop, «человек в контуре») подходит для этапов,
требующих выбора пользователя, авторизации или дополнительного
профессионального суждения. В запросе следует вместе представить вопрос,
требующий решения, возможные действия и их последствия, чтобы человек
мог принять обоснованное решение. Например, когда решение уже
подготовлено, но для его исполнения требуется явная авторизация, следует
отправить запрос на одобрение; если для задачи не хватает обязательных
условий --- запросить дополнительную информацию.

Ответ человека может поступить не сразу, поэтому для запроса необходимо
задать срок действия и поведение при тайм-ауте. Например, после пяти
минут ожидания можно сохранить задачу в состоянии ожидания,
приостановить зависящие от ответа операции или применить заранее
разрешённый консервативный вариант. Если срок запроса истёк, он отменён
или уже обработан, повторный ответ должен возвращать текущее состояние;
если требуется новое решение, следует создать новый запрос. Уведомления
можно отправлять в соответствии с приоритетом: срочные запросы --- через
настроенные каналы мгновенной связи, обычные --- в очередь электронной
почты или задач.

Одобрения, отказы и причины внесённых человеком изменений также могут
стать данными обратной связи. Сначала следует сохранить задачу,
доступные на тот момент доказательства и принятое решение, а затем
оценить, какой опыт можно использовать повторно: явные бизнес-правила
можно включить в базу знаний или Skill, а суждения о предпочтениях ---
преобразовать в примеры для оценки или постобучения. В
\crossreflink{chap:9}{главе 9} рассматривается, как отбирать такие
траектории и выбирать способ обновления.

\begin{quote}
\textbf{Эксперимент 4-5 ★★: MCP-сервер инструментов совместной работы}

В этом эксперименте реализуются интерфейсы управления субагентами,
привлечения человека и отправки уведомлений, а также отрабатываются два
способа передачи контекста и синхронный и асинхронный вызовы.

Сначала рассмотрим запрос на возврат средств. Товар приобретён три дня
назад за 80 сингапурских долларов. Правила магазина разрешают возврат в
течение семи дней после покупки, если сумма меньше 100 сингапурских
долларов. Главный агент передаёт субагенту задачу «определить по
правилам магазина, соответствует ли этот возврат условиям». Какие факты
следует отправить вместе с задачей?

При синхронном делегировании через \texttt{mcp\_spawn\_subagent} вариант
с передачей только необходимого контекста упустил сведения о покупке и
отправил лишь правила магазина. Внешнее состояние квитанции ---
\texttt{completed}, поскольку субагент уже ответил; однако внутри ответа
указано состояние \texttt{need\_info}: не хватает суммы и времени
покупки. Главный агент с помощью
\texttt{mcp\_send\_message\_to\_subagent} добавил сведения о том, что
товар не использовался, но это по-прежнему не позволяло ответить на два
указанных вопроса. Только вариант с созданным моделью резюме передал
одновременно сведения о трёх днях, 80 сингапурских долларах и правилах
магазина, после чего было получено заключение: «Возврат соответствует
условиям».

При переходе к асинхронному делегированию квитанция о запуске
одновременно содержит \texttt{task\_id} и \texttt{subagent\_id}. В
интерфейсе этого эксперимента при запросе состояния, добавлении
сообщения или запросе отмены используется параметр
\texttt{subagent\_id}. По этому идентификатору можно найти историю
одного и того же субагента от \texttt{running} до \texttt{completed}.
Для обращения к человеку предусмотрен отдельный \texttt{request\_id}:
\texttt{mcp\_request\_admin\_approval} запрашивает одобрение, а
\texttt{mcp\_request\_admin\_input} --- дополнительную информацию. При
фактическом запуске администратор одобрил публикацию отчёта об этом
эксперименте. Для другого запроса тайм-аут наступил через одну секунду,
и в квитанции одновременно были указаны \texttt{approved:\ false} и
\texttt{timeout:\ true}; на этом основании главному агенту можно лишь
сообщить, что одобрение ещё не получено.

Сопутствующий сервер также предоставляет интерфейсы для отправки
сообщений по электронной почте, через Telegram и Slack.
\end{quote}

\hypertarget{ux438ux442ux43eux433ux438-ux433ux43bux430ux432ux44b-3}{%
\section{Итоги
главы}\label{ux438ux442ux43eux433ux438-ux433ux43bux430ux432ux44b-3}}

При подключении к модели новой возможности необходимо отдельно принять
два решения. Первое: как дать модели использовать эту возможность?
\textbf{Автор обычно сначала выбирает универсальную среду выполнения},
позволяя модели писать код, который комбинирует операции.
Специализированный интерфейс стоит создавать для операций, требующих
отдельной авторизации, сложных параметров, в которых легко ошибиться,
часто повторяющихся вызовов или различий между платформами. Второе:
когда модель должна увидеть эту возможность? Даже если инструмент уже
подключён через MCP, модель может сначала просмотреть каталог и
прочитать полное определение лишь тогда, когда оно понадобится.
Инструкции по выполнению операций, скрипты и справочные материалы,
распространяемые через Skill Hub, также можно загружать по мере
необходимости.

Проверьте на примере таблицы сроков поставки из начала главы, достаточно
ли возможностей интерфейса. Когда помощник \textbf{собирает
свидетельства} из писем и вложений, длинные результаты можно читать
постранично или по диапазонам, но места усечения должны быть отмечены, а
при сжатии таблиц или изображений необходимо сохранять соответствие
между датами и заказами. Перед \textbf{выполнением} отправки помощник
должен проверить получателей и вложения. Если во время отправки
соединение прервалось, сначала следует по идентификатору сообщения
проверить, получил ли его сервер, и лишь затем решать, нужна ли
повторная попытка. Скрипты можно выполнять в песочнице, а внешние
операции записи требуют проверки идемпотентности. Для задачи, отмена
которой уже запрошена, необходимо также получить и прочитать
подтверждение, чтобы удостовериться, какие операции действительно
остановлены. Если \textbf{делегировать} проверку вложений другому
AI-агенту, должна быть возможность по идентификатору задачи проверять
ход выполнения, отправлять дополнительные сообщения и отменять задачу, а
при обнаружении двух противоречащих друг другу записей --- запрашивать
решение человека.

Теперь инструменты позволяют помощнику узнать срок поставки
определённого заказа. Но если ни почтовая система, ни система
электронных таблиц не предоставляют готового интерфейса для сводки,
сможет ли помощник самостоятельно написать такой интерфейс, чтобы
воспользоваться им в следующий раз? С этого начинается пятая глава. О
том, как будить помощника при поступлении нового письма и как доставлять
результат, когда пользователь не в сети, пойдёт речь в шестой главе.

\begin{questionbox}

\hypertarget{ux432ux43eux43fux440ux43eux441ux44b-ux434ux43bux44f-ux440ux430ux437ux43cux44bux448ux43bux435ux43dux438ux44f-3}{%
\section*{Вопросы для
размышления}\label{ux432ux43eux43fux440ux43eux441ux44b-ux434ux43bux44f-ux440ux430ux437ux43cux44bux448ux43bux435ux43dux438ux44f-3}}
\addcontentsline{toc}{section}{Вопросы для размышления}

\begin{enumerate}
\def\labelenumi{\arabic{enumi}.}
\tightlist
\item
  ★★ MCP отделяет подключение инструментов от фреймворка AI-агента.
  Выберите одно из требований: потоковая передача результатов,
  двунаправленное взаимодействие или управление состоянием длительных
  задач. Укажите, какие возможности протокола и хоста необходимы для его
  реализации, и обсудите, чем потребуется дополнить существующий MCP.
\item
  ★★ Как следует выбирать между несколькими серверами MCP,
  предоставляющими сходные инструменты? На примере, когда один
  возвращает краткое содержание, а другой --- полный текст, разработайте
  описания инструментов, обозначения источников и критерии выбора,
  позволяющие агенту использовать эти различия.
\item
  ★★ В каких ещё сценариях, помимо запуска проверок после написания
  кода, можно применять цикл «выполнение --- проверка --- обратная
  связь»? Выберите два примера, опишите методы проверки, обратную связь
  при неудаче и затраты, а также обсудите, когда уместны немедленная,
  выборочная и асинхронная проверка.
\item
  ★★ Как ещё, помимо активного обнаружения, можно организовать
  возможности при наличии тысяч инструментов? Опишите условия
  применимости и компромиссы одного из подходов, использующего каталог с
  категориями, шаблоны рабочих процессов или универсальный исполнитель.
\end{enumerate}

\textbf{Практикум к \crossreflink{chap:4}{главе 4}.}
\href{https://github.com/bojieli/ai-agent-book/tree/40f4ce5778e6cc93e3606c3079bf8cf13a0c4d04/chapter4}{Каталог
экспериментов} ·
\href{https://github.com/bojieli/ai-agent-book/blob/40f4ce5778e6cc93e3606c3079bf8cf13a0c4d04/chapter4/README.ru.md}{Путеводитель
на русском}.

\end{questionbox}

\hypertarget{coding-agent-ux438-ux443ux43dux438ux432ux435ux440ux441ux430ux43bux44cux43dux44bux439-ux430ux433ux435ux43dux442}{%
\chapter{\texorpdfstring{Coding Agent и универсальный
агент\label{chap:5}}{Coding Agent и универсальный агент}}\label{coding-agent-ux438-ux443ux43dux438ux432ux435ux440ux441ux430ux43bux44cux43dux44bux439-ux430ux433ux435ux43dux442}}

Предположим, пользователь передал вам три таблицы продаж с записями о
заказах и возвратах и попросил определить товар с самой высокой долей
возвратов, а затем оформить обоснование в виде отчёта. Чтобы AI-агент
взялся за такую задачу открытого типа, ему сначала нужно разобраться в
заголовках столбцов и единицах измерения, написать программу для
объединения данных, а затем сверить аномальные значения с исходными
записями. Если при сверке обнаружится, что названия товаров в двух
таблицах различаются, ему придётся вернуться и изменить программу.
Контекст, память, база знаний и инструменты, рассмотренные по
отдельности в предыдущих главах, пригодятся на разных этапах такого
итеративного процесса. В этой главе мы начнём с задач программирования и
постепенно распространим тот же подход на исследования, создание
контента и анализ данных.

Автор видит архитектуру универсального агента для задач открытого типа
следующим образом: \textbf{исполнительным ядром служит Coding Agent, а
рабочим пространством --- файловая система.} Coding Agent способен
писать, изменять и выполнять код, а файловая система хранит код, данные,
память и промежуточные результаты. Вместе они позволяют агенту по мере
необходимости писать программы и сохранять результаты в файлах, чтобы
использовать их на последующих этапах.

Генерация кода --- это \textbf{метавозможность}: агент может с помощью
кода организовывать работу существующих инструментов, а также создавать
новые инструменты под новые требования. Кроме того, код позволяет явно
выражать правила и этапы рассуждения. Например, если записать условие
«возраст старше 18 лет и личность подтверждена по настоящему имени» как
\texttt{age\ \textgreater{}\ 18\ and\ verified}, это правило можно
проверить на нескольких наборах входных данных. Из двух людей с
подтверждённой личностью --- 18-летнего и 19-летнего --- условию должен
соответствовать только второй. Если личность этого 19-летнего заявителя
ещё не подтверждена по настоящему имени, его заявку также следует
отклонить.

Сначала в этой главе будут рассмотрены базовые инструменты Coding Agent,
а на примере OpenClaw будет показано, как код и файловая система
обеспечивают выполнение универсальных задач. Затем мы разберём рабочий
процесс, Harness и конкретную реализацию. После этого применение
генерации кода будет рассмотрено в шести направлениях: рассуждение,
бизнес-правила, мультимедиа, адаптация систем, генеративный UI и
самозагрузка агента.

\hypertarget{coding-ux430ux433ux435ux43dux442}{%
\section{Coding-агент}\label{coding-ux430ux433ux435ux43dux442}}

\hypertarget{ux43dux430ux43fux438ux441ux430ux43dux438ux435-ux43aux43eux434ux430-ux431ux430ux437ux43eux432ux430ux44f-ux441ux43fux43eux441ux43eux431ux43dux43eux441ux442ux44c-ai-ux430ux433ux435ux43dux442ux430}{%
\subsection{Написание кода --- базовая способность
AI-агента}\label{ux43dux430ux43fux438ux441ux430ux43dux438ux435-ux43aux43eux434ux430-ux431ux430ux437ux43eux432ux430ux44f-ux441ux43fux43eux441ux43eux431ux43dux43eux441ux442ux44c-ai-ux430ux433ux435ux43dux442ux430}}

\textbf{Автор считает генерацию кода базовой способностью универсального
агента.} Математические вычисления можно поручить решателю, для
обработки данных можно вызвать программную библиотеку, а повторяющиеся
операции --- оформить в виде скрипта. Столкнувшись с новым форматом или
интерфейсом, агент также может сгенерировать логику разбора и адаптации.
Базовые инструменты предоставляют точки выполнения, а код объединяет их
возможности в последовательность шагов, необходимую для решения задачи.

Рассмотрим задачу: «Собрать все оставшиеся в репозитории комментарии
TODO, классифицировать их по приоритету и создать issue». Сначала нужно
найти TODO, прочитать код до и после каждого комментария и определить,
какие из них служат лишь личными напоминаниями, а какие требуют
отдельной задачи. Подготовив имена файлов, номера строк и обоснования
решений, можно после получения разрешения создать issue. Записывать файл
потребуется лишь в том случае, если результаты классификации также нужно
сохранить в репозитории. В подобных задачах могут понадобиться и такие
операции, как выполнение команд и редактирование кода. В учебных целях
мы используем следующие семь базовых инструментов:

\begin{enumerate}
\def\labelenumi{\arabic{enumi}.}
\tightlist
\item
  \textbf{Code Interpreter (интерпретатор кода)}: выполняет код Python и
  других языков для вычислений, обработки данных и построения графиков;
  песочница ограничивает доступ к файлам, сети и ресурсам.
\item
  \textbf{Bash Shell (терминал командной строки)}: выполняет команды
  сборки, тестирования, обработки файлов и другие команды, возвращая
  вывод и статус завершения.
\item
  \textbf{Инструмент чтения файлов}: читает код, конфигурации,
  документацию и журналы, поддерживая получение заданного диапазона
  содержимого.
\item
  \textbf{Инструмент записи файлов}: создаёт файлы или полностью
  заменяет их содержимое.
\item
  \textbf{Инструмент редактирования файлов}: вносит локальные изменения
  в существующие файлы и подходит для сопровождения и итеративной
  разработки.
\item
  \textbf{Инструмент поиска по именам файлов (Glob)}: находит файлы по
  шаблону, например \texttt{**/*.py} позволяет найти файлы Python.
\item
  \textbf{Инструмент поиска по содержимому файлов (Grep)}: ищет текст
  или шаблон и возвращает имя файла, номер строки и совпавшее
  содержимое, например позволяет найти места вызова определённой
  функции.
\end{enumerate}

Эти семь инструментов главным образом обеспечивают возможности
восприятия и выполнения из четвёртой главы. Средства совместной работы,
событийного запуска и взаимодействия с пользователем можно подключать в
соответствии с потребностями приложения. Количество инструментов и
гранулярность их интерфейсов определяются реализацией: Shell позволяет
просматривать каталоги, выполнять поиск и читать файлы, тогда как
отдельные файловые инструменты предоставляют единообразные параметры и
структурированную обратную связь. Например, Codex может читать и искать
с помощью инструмента выполнения команд, а изменения файлов выражать
через \texttt{apply\_patch}. Патч объединяет целевой файл, место
изменения, добавляемое и удаляемое содержимое, что упрощает проверку и
воспроизведение изменений.

Сузим задачу до «составить список TODO в файлах Python» и посмотрим, как
взаимодействуют инструменты:

\begin{longtable}[]{@{}
  >{\raggedright\arraybackslash}p{(\columnwidth - 4\tabcolsep) * \real{0.2678}}
  >{\raggedright\arraybackslash}p{(\columnwidth - 4\tabcolsep) * \real{0.3142}}
  >{\raggedright\arraybackslash}p{(\columnwidth - 4\tabcolsep) * \real{0.4181}}@{}}
\toprule\noalign{}
\begin{minipage}[b]{\linewidth}\raggedright
\hspace{0pt}Шаг
\end{minipage} & \begin{minipage}[b]{\linewidth}\raggedright
\hspace{0pt}Операция
\end{minipage} & \begin{minipage}[b]{\linewidth}\raggedright
\hspace{0pt}Ответ или результат
\end{minipage} \\
\midrule\noalign{}
\endhead
\bottomrule\noalign{}
\endlastfoot
\hspace{0pt}Поиск комментариев & \hspace{0pt}Grep ищет \texttt{TODO},
область поиска ограничена шаблоном \texttt{**/*.py} &
\hspace{0pt}\texttt{api.py}, строка 42: добавить ограничение частоты
запросов; \texttt{db.py}, строка 15: перейти на PostgreSQL;
\texttt{test\_}\allowbreak{}\texttt{api.py}, строка 8: добавить тесты
граничных случаев \\
\hspace{0pt}Составление списка & \hspace{0pt}Write записывает
\texttt{TODO\_}\allowbreak{}\texttt{LIST.md} & \hspace{0pt}Три задачи с
соответствующими файлами, номерами строк и пояснениями \\
\hspace{0pt}Сообщение результата & \hspace{0pt}Сообщить пользователю
расположение списка и область охвата & \hspace{0pt}Всего найдено три
TODO в файлах Python \\
\end{longtable}

Если пользователь также попросит подсчитать количество TODO по модулям и
построить столбчатую диаграмму, агент может сгенерировать программу
Python и передать её Code Interpreter для вычислений и построения
графика. Если же необходимо определить приоритеты, агент продолжит
читать связанную реализацию, тесты или существующие issue, чтобы найти
основания для оценки.

\hypertarget{ux43fux440ux438ux43cux435ux440-ux43eux442-manus-ux434ux43e-openclaw-ux43fux440ux43eux433ux440ux430ux43cux43cux43dux43eux435-ux44fux434ux440ux43e-ux443ux43dux438ux432ux435ux440ux441ux430ux43bux44cux43dux43eux433ux43e-ai-ux430ux433ux435ux43dux442ux430}{%
\subsection{Пример: от Manus до OpenClaw --- программное ядро
универсального
AI-агента}\label{ux43fux440ux438ux43cux435ux440-ux43eux442-manus-ux434ux43e-openclaw-ux43fux440ux43eux433ux440ux430ux43cux43cux43dux43eux435-ux44fux434ux440ux43e-ux443ux43dux438ux432ux435ux440ux441ux430ux43bux44cux43dux43eux433ux43e-ai-ux430ux433ux435ux43dux442ux430}}

Универсальные агенты, такие как Manus и OpenClaw, могут в рамках одной
задачи искать информацию, управлять компьютером и выполнять код.
Несколько видов работы часто связываются одной программой: она получает
данные через интерфейс, упорядочивает результаты поиска, создаёт
диаграммы и документы. При повторном выполнении аналогичной задачи эту
программу можно использовать снова.

Например, программные библиотеки позволяют напрямую создавать и изменять
документы Word и презентации PowerPoint в формате OOXML (Office Open
XML). Отчёты PDF можно формировать из Markdown, HTML или LaTeX, а анализ
и визуализацию данных --- выполнять с помощью Python. Стабильные
операции с веб-страницами также можно оформлять как повторно
используемые скрипты. В девятой главе подробнее рассматривается
извлечение повторно используемых шагов из единожды успешно выполненной
операции.

\textbf{Если для операции уже существует подходящий API или интерфейс
командной строки, автор предпочитает вызывать её программно.}
Структурированные параметры избавляют от затрат на повторные снимки
экрана, визуальный поиск элементов и щелчки, а скрипт можно сохранить,
протестировать и использовать повторно. Computer Use расширяет круг
программ, которыми можно управлять; закрепление в коде уже отлаженных и
регулярно повторяемых шагов дополнительно снижает стоимость и задержку.
Файловая система на \crossreflink{fig:5-1}{рисунке 5-1} связывает
отдельные вызовы. Файл с данными, сохранённый одним вызовом, может стать
входом для следующего скрипта анализа.

\begin{figure}
\hypertarget{fig:5-1}{%
\centering
\includegraphics{images/revised/fig5-1.pdf}
\caption{Рисунок 5-1. Coding-агент как ядро архитектуры
OpenClaw}\label{fig:5-1}
}
\end{figure}

Например, процесс выполнения задачи «проанализировать данные о продажах
за прошлый квартал и подготовить краткий отчёт» можно разделить на пять
шагов:

\begin{enumerate}
\def\labelenumi{\arabic{enumi}.}
\tightlist
\item
  \textbf{Чтение памяти}: получить сведения о предпочитаемом
  пользователем формате отчётов PDF и о том, что данные о продажах
  хранятся в Google Sheets.
\item
  \textbf{Получение данных}: изучить использование Google Sheets API и
  загрузить необходимую таблицу через уже авторизованный интерфейс.
\item
  \textbf{Написание кода анализа}: агрегировать данные о продажах с
  помощью pandas, построить диаграммы с помощью matplotlib и проверить
  методику расчётов и результаты.
\item
  \textbf{Создание результатов}: объединить анализ и диаграммы в PDF,
  сохранить исходные данные, промежуточные результаты и скрипт генерации
  для последующей проверки или обновления.
\item
  \textbf{Обновление памяти}: записать подтверждённый источник данных и
  идентификатор таблицы, чтобы найти их при выполнении следующей задачи.
\end{enumerate}

Если поручить ту же задачу офисному сервису Qwen, пользователь сможет
напрямую загрузить детализацию продаж, после чего сервис проанализирует
данные и сформирует визуальный отчёт.\footnote{Справочный центр Qwen
  Office, обзор Qwen Office с описанием загрузки файлов данных, анализа
  данных, создания визуальных отчётов и других возможностей.}
Предположим, по заказу получена оплата в размере 100 юаней, а затем
оформлен возврат 20 юаней. В чистые продажи за квартал должны войти
только 80 юаней. Если на диаграмме в отчёте указано 100 юаней, по
цепочке от диаграммы к сводной таблице и исходной детализации можно
определить, какой возврат не был учтён.

Файловая система связывает эти шаги: результат, сохранённый на
предыдущем шаге, может стать входом для следующего, а программа может
многократно читать одни и те же данные. Если задача была прервана, агент
может определить достигнутый прогресс по сохранённым результатам и
продолжить обработку незавершённой части.

Промежуточные результаты полезны и в другом случае: если пользователь
задаст связанный вопрос, работу можно продолжить с подходящего шага.
Исследовательское подразделение Microsoft продемонстрировало это в Data
Formulator на данных о производстве электроэнергии в разных странах.
Исходная таблица содержала объёмы выработки для различных видов энергии.
Сначала аналитик рассчитал долю возобновляемой энергетики по странам и
годам, а затем построил на основе этих данных рейтинг. Если требуется
узнать, насколько выросла доля в определённой стране, одних мест на
диаграмме недостаточно: они не показывают разницу между двумя годами в
процентных пунктах. Сохранившийся в интерфейсе ход вычислений позволяет
вернуться к таблице долей, взять значения за два года и вычесть одно из
другого.\footnote{Исследовательское подразделение Microsoft, пример
  анализа энергетических данных в Data Formulator и реализованная в
  версии 0.7.0 связь между диаграммами и таблицами данных.}

\textbf{Сохраняйте память, знания и возможности в файлах.} Рабочее
пространство OpenClaw хранит в файлах Markdown несколько видов
информации: \texttt{MEMORY.md} содержит долговременные факты и решения,
организованные по датам журналы --- сведения о ходе работы,
\texttt{USER.md} может использоваться для предпочтений пользователя, а
файлы Skill описывают способы применения возможностей\footnote{OpenClaw,
  \emph{Agent workspace} и \emph{Memory overview} с описанием рабочего
  пространства, файловой памяти и механизмов поиска.}. \textbf{Для
памяти, которую должны совместно поддерживать человек и агент, автор
предпочитает непосредственно читаемые файлы.} Пользователь может открыть
записи и исправить ошибочный факт, а Git позволяет проследить источник,
сравнить версии и выполнить откат. Организация журналов по датам
сохраняет явную хронологию событий; когда файлов становится много,
поисковый индекс ускоряет их нахождение.

Способность агента записывать файлы также позволяет накапливать опыт
выполнения задач. Например, если при обращении в банк выяснилось, что
для подтверждения личности требуется адрес отделения, в котором был
открыт счёт, вместе с этим фактом можно записать применимые операции,
источник информации и время события. После проверки такие записи можно
преобразовать в знания, инструкции или программы для повторного
использования в последующих задачах. Этот процесс непрерывного развития
рассматривается в девятой главе.

\textbf{Для каких задач агенту подходит самостоятельное написание
программ?} В открытых исследованиях, генерации контента и обработке
данных часто встречаются новые форматы входных данных, интерфейсы и
требования к результатам, поэтому генерация кода хорошо подходит на роль
исполнительного ядра. Системы с более стабильными процессами, например
специализированная служба поддержки, могут организовать работу вокруг
бизнес-состояния, предметных инструментов и диалоговой стратегии,
используя код для вычислений, преобразований и проверки правил. Оба типа
систем могут применять возможности Coding, но их роль меняется в
зависимости от задачи.

\hypertarget{ux43eux431ux449ux438ux439-ux43fux440ux43eux446ux435ux441ux441-ux440ux430ux431ux43eux442ux44b-coding-ux430ux433ux435ux43dux442ux430}{%
\subsection{Общий процесс работы
Coding-агента}\label{ux43eux431ux449ux438ux439-ux43fux440ux43eux446ux435ux441ux441-ux440ux430ux431ux43eux442ux44b-coding-ux430ux433ux435ux43dux442ux430}}

Рассмотрим \crossreflink{fig:5-2}{рисунок 5-2} на примере небольшого
дефекта. Пользователь сообщает: «При сортировке по общей стоимости 100
юаней оказываются перед 20 юанями». Сначала нужно подобрать входные
данные, воспроизводящие проблему, а затем проследить движение данных и
определить, где сумма превратилась в строку. После исправления логики
сравнения следует повторно выполнить проверку с исходными входными
данными и убедиться, что значения с дробной частью и одинаковые суммы
по-прежнему обрабатываются правильно.

\begin{figure}
\hypertarget{fig:5-2}{%
\centering
\includegraphics{images/revised/fig5-2.pdf}
\caption{Рисунок 5-2. Процесс работы Coding-агента}\label{fig:5-2}
}
\end{figure}

\textbf{Понимание проекта и дополнение документации.}

Получив задачу, агент сначала читает описание проекта, а затем
сопоставляет его с реализацией, местами вызова и тестами, чтобы
определить область влияния изменений. Если в документации недостаточно
ясно описаны назначение соответствующих модулей, используемые ключевые
абстракции, взаимодействие модулей или порядок тестирования, агент
дополняет её сведениями, подтверждёнными по коду, чтобы облегчить работу
следующим разработчикам.

\textbf{Файлы проектных инструкций} размещают часто используемые
соглашения там, где агенту легко их получить. В \texttt{CLAUDE.md},
\texttt{AGENTS.md} и правилах проекта Cursor можно зафиксировать команды
сборки, стиль кода и границы изменений, например: «используй
\texttt{pnpm\ test} для запуска тестов», «объявляй конкретные типы
переменных и избегай \texttt{any}», «изменения миграций базы данных
требуют отдельной проверки». Разные среды загружают эти инструкции в
зависимости от каталога, конфигурации или условий применимости; ранний
файл Cursor \texttt{.cursorrules} также воплощал этот подход\footnote{Anthropic,
  \emph{How Claude remembers your project}; Cursor, Rules --- описание
  проектных инструкций и правил, загружаемых в соответствии с областью
  действия.}.

У файлов инструкций и файлов OpenClaw \texttt{SOUL.md} и
\texttt{MEMORY.md} разные назначения: первые описывают правила работы
внутри проекта, \texttt{SOUL.md} определяет личность и поведение агента,
а файлы памяти сохраняют сведения, которые можно использовать в
последующих задачах. Стабильные и краткие инструкции помогают сохранять
неизменным префикс контекста, а загрузка правил каталогов по мере
необходимости позволяет контролировать объём входных данных текущей
задачи. Разрешения на выполнение обеспечиваются Harness и средой
выполнения.

Как отмечалось во второй главе, команды, хорошо приспособленные к
удалённой совместной работе, обычно также упрощают участие агента. На
уровне репозитория кода это означает, что решения фиксируются в
проектной документации, контекст задачи --- в issue и PR, а внутренний
опыт --- в руководстве для разработчиков. Полезный проверочный вопрос:
\textbf{сможет ли новый удалённый сотрудник, располагая только
репозиторием и документацией, найти точку входа в задачу, запустить
тесты и понять требования к изменениям?}

\textbf{Определение цели и критериев приёмки.}

Если bug уже локализован или требуется изменить параметр, правку можно
выполнить непосредственно на основе существующего поведения и тестов.
Более широкую цель сначала нужно конкретизировать. Например,
«оптимизировать производительность системы» может означать снижение
времени отклика, уменьшение потребления памяти или повышение пропускной
способности. Агент должен определить текущее узкое место, измерить
базовые показатели и выяснить допустимую область изменений, подтвердить
такие ограничения, как совместимость интерфейсов и удобство
использования, и лишь затем выбрать способ оптимизации.

\textbf{Подготовка проектного решения.}

Описание решения должно отвечать на следующие вопросы: какие модули
будут изменены, какой подход будет применён, почему он подходит для
текущей задачи, появятся ли новые зависимости, а также как проверить
результат и выполнить восстановление. Для небольшого изменения
достаточно краткого плана; сложное изменение следует оформить как
проектный документ, доступный для проверки. Если проект или пользователь
требуют одобрения решения, сначала необходимо представить конкретный
план и только после его утверждения приступать к реализации.

\textbf{Реализация, тестирование и проверка.}

При реализации агент соблюдает соглашения проекта, повторно использует
существующие абстракции и инструменты и ограничивает изменения областью,
относящейся к задаче. Тесты проектируются вокруг ожидаемого поведения и
охватывают нормальные пути, граничные условия и исключительные ситуации.
При исправлении bug регрессионный тест должен воспроизводить исходную
проблему и подтверждать результат исправления.

После сбоя теста агент изучает диагностические сведения, локализует
причину, изменяет код и повторно запускает соответствующие проверки.
\textbf{Критерием завершения задачи программирования должно быть
прохождение приёмки.} Создание патча --- лишь один из этапов: необходимо
также сверить связанные тесты, требования пользователя и результаты
работы, подлежащие передаче. Цикл «тестирование --- исправление ---
проверка» позволяет агенту продолжать устранение проблем и предоставляет
Loop-инженерии чёткое условие остановки.

Проверка кода дополнительно охватывает аспекты, которые могут не
проверяться тестами, например читаемость, совместимость интерфейсов,
производительность, разрешения и использование зависимостей. Для этого
можно сочетать просмотр диффа, статические проверки и привлечение
субагента проверки кода. Обнаруженные проблемы следует исправить, а
затронутые проверки --- запустить повторно.

\textbf{Синхронизация документации и передача результата.}

Добавление модулей, изменение зависимостей и ключевого поведения должны
отражаться в соответствующей документации. В описании результата следует
кратко изложить решённую проблему, изменения поведения, результаты
проверки и порядок использования, а при необходимости --- добавить
инструкции по восстановлению. Код, тесты и документация вместе сохраняют
результат выполненной задачи.

Разные модели также по-разному организуют исследование и внесение
изменений. Одни сначала подробно изучают реализацию, места вызова и
тесты; другие начинают с наиболее релевантных файлов, пробуют небольшой
патч, а затем продолжают локализацию по результатам компиляции или
тестирования. \textbf{Автор прежде всего рассматривает такой ритм
исследования и действий как усвоенную моделью стратегию поведения.} Если
после смены Harness сохраняется схожий ритм, а при смене модели внутри
одного Harness он заметно меняется, различие указывает на саму модель.
Промпты, инструменты и бюджет Harness могут усиливать или подавлять эту
склонность. В седьмой главе рассматривается оценка этих различий, а в
восьмой --- влияние постобучения на стратегию поведения.

\hypertarget{ux43fux440ux430ux43aux442ux438ux43aux430-harness-ux438ux43dux436ux435ux43dux435ux440ux438ux438-ux432-coding-ux430ux433ux435ux43dux442ux430ux445}{%
\subsection{Практика Harness-инженерии в
Coding-агентах}\label{ux43fux440ux430ux43aux442ux438ux43aux430-harness-ux438ux43dux436ux435ux43dux435ux440ux438ux438-ux432-coding-ux430ux433ux435ux43dux442ux430ux445}}

В первой главе производственная реализация описывалась формулой
\textbf{Agent = Model + Harness}, а также были представлены пять функций
Harness: управление контекстом, интерфейсы инструментов, ограничения,
проверка и исправление. Coding-агент может задействовать тестирование,
статический анализ, непрерывную интеграцию и контроль версий из
программной инженерии, чтобы воплотить эти функции в исполняемом
процессе.

Вернёмся к рассмотренному выше дефекту сортировки. Передача модели
ошибочного ввода, соответствующей функции и вывода с ошибкой относится к
управлению контекстом. Возможность читать и изменять код, а также
запускать тесты обеспечивается интерфейсами инструментов и
ограничениями. Проверка начинается лишь тогда, когда изменение
тестируется на том же вводе. Если числа по-прежнему сортируются как
символы, пригодной для исправления обратной связью станет совместная
передача модели фактического и ожидаемого порядков.

Модель отвечает за понимание проблемы и создание решения и кода. Harness
предоставляет необходимые для работы сведения и интерфейсы, организует
проверки, а по их результатам продолжает или восстанавливает выполнение.
На уровне инженерной системы следует подготовить четыре компонента:

\begin{itemize}
\tightlist
\item
  \textbf{Базовые критерии приёмки}: комплект тестов, CI-конвейер и
  стандарты проверки кода определяют, что считается завершённой работой.
  CI, то есть конвейер непрерывной интеграции, может автоматически
  запускать согласованные проверки после коммита.
\item
  \textbf{Границы выполнения}: явно определяют обязанности модулей,
  правила зависимостей, доступ к данным и разрешения на операции.
\item
  \textbf{Сигналы обратной связи}: результаты linter, проверки типов и
  тестов преобразуются в диагностические сообщения, по которым модель
  может локализовать и понять проблему, а затем принять меры.
\item
  \textbf{Средства восстановления}: сохраняются версии Git, снимки
  рабочего пространства и состояние задачи, а для внешних изменений
  подготавливаются соответствующие процессы восстановления.
\end{itemize}

\textbf{Проектируйте цель вместе со способом проверки.}

Способ работы агента зависит как от ясности задачи, так и от возможности
автоматически проверить результат. В таблице 5-1 представлены четыре
сочетания этих факторов и дополнительные инженерные меры, необходимые в
каждом случае. Уточнение требований, добавление примеров и проверок
позволяют постепенно преобразовать изначально расплывчатую задачу в
исполняемую цель.

Таблица 5-1. Четыре квадранта ясности задачи и автоматизации проверки

\begin{longtable}[]{@{}
  >{\raggedright\arraybackslash}p{(\columnwidth - 4\tabcolsep) * \real{0.2609}}
  >{\raggedright\arraybackslash}p{(\columnwidth - 4\tabcolsep) * \real{0.3747}}
  >{\raggedright\arraybackslash}p{(\columnwidth - 4\tabcolsep) * \real{0.3644}}@{}}
\toprule\noalign{}
\begin{minipage}[b]{\linewidth}\raggedright
\end{minipage} & \begin{minipage}[b]{\linewidth}\raggedright
\hspace{0pt}Результат можно проверить автоматически
\end{minipage} & \begin{minipage}[b]{\linewidth}\raggedright
\hspace{0pt}Результат требует ручной проверки
\end{minipage} \\
\midrule\noalign{}
\endhead
\bottomrule\noalign{}
\endlastfoot
\hspace{0pt}\textbf{Цель ясна} & \hspace{0pt}Исправление ошибки с
регрессионным тестом: итерации по обратной связи &
\hspace{0pt}Рефакторинг кода: сочетание тестов и архитектурной
проверки \\
\hspace{0pt}\textbf{Цель неоднозначна} & \hspace{0pt}«Повысить качество
кода»: сначала определить цели качества помимо требований linter &
\hspace{0pt}«Сделать UI красивее»: уточнить предпочтения с помощью
примеров и сравнений \\
\end{longtable}

\textbf{Накопленная программной инженерией инфраструктура проверки ---
одна из важных причин, по которым Coding-агенты сравнительно рано
достигли зрелости.} Компиляторы, системы типов, комплекты тестов и
контроль версий позволяют локализовать ошибки, проверять изменения и
выполнять откат после неудачи, естественным образом образуя мощный
Harness. Этот вывод также определяет направление вложений для других
агентных сценариев: прежде всего следует перемещать задачи в область
«цель ясна, результат можно проверить автоматически». Если приёмка
требует каждый раз ручной проверки, скорость человека становится верхней
границей пропускной способности. Если же цель неоднозначна, даже самая
быстрая автоматизация может непрерывно оптимизировать неверный
показатель.

\textbf{Извлекайте методы из инженерных примеров.}

Крупномасштабная миграция кода позволяет понять, как взаимодействуют эти
компоненты. Правила миграции и связи между модулями записываются в
документации репозитория, ограничения зависимостей обеспечиваются
статическими проверками, каждая группа изменений проходит тесты и CI, а
при сбое возвращается диагностическое сообщение, позволяющее
локализовать проблему. Благодаря этому агент может изучить требования к
миграции, получать обратную связь после каждой группы изменений и на её
основе продолжать работу.

Практика LangChain показывает, как улучшать Harness по траекториям
неудач: зафиксировать модель, проанализировать, на каком шаге задача
остановилась или завершилась ошибкой, а затем изменить системный промпт,
middleware инструментов, цикл самопроверки и бюджет рассуждений. Пример
с Terminal Bench 2.0 уже рассматривался в первой главе; здесь важен
итеративный метод «зафиксировать неудачу --- предложить изменение ---
провести повторную оценку»\footnote{LangChain, \emph{Improving Deep
  Agents with harness engineering}, 17 февраля 2026 года.}.

Подход Anthropic к длительным задачам разделяет инициализацию среды и
непрерывное выполнение. Агент инициализации создаёт список задач и
рабочую среду, а последующие Coding-агенты последовательно выполняют
пункты, сохраняя код, результаты тестов и прогресс. Новая сессия читает
эти результаты и продолжает работу, что сокращает потери информации при
передаче между контекстами\footnote{Anthropic, \emph{Effective harnesses
  for long-running agents}, 26 ноября 2025 года.}.

\textbf{Преобразуйте накопленный опыт в принципы проектирования
системы.}

\begin{enumerate}
\def\labelenumi{\arabic{enumi}.}
\tightlist
\item
  \textbf{Исполняемые ограничения важнее указаний в промпте.}
  Зависимости модулей, типы полей и стандарты кода можно проверять
  посредством статического анализа и CI, а права доступа должны
  принудительно обеспечиваться средой выполнения. Промпт и документация
  описывают правила и их причины, а механизмы проверки обнаруживают и
  блокируют соответствующие нарушения.
\item
  \textbf{В первую очередь вкладывайтесь в автоматическую проверку.}
  Количество повторяющихся проверок растёт вместе с объёмом задач,
  поэтому постоянное увеличение ручной проверки становится препятствием
  для масштабирования. Эту работу берут на себя тесты, проверки качества
  и мониторинг поведения, а ручная проверка сосредотачивается на поиске
  компромиссов между целями и сложных решениях.
\item
  \textbf{Чем своевременнее и структурированнее обратная связь, тем
  легче выполнить исправление.} Сообщение об ошибке должно содержать
  соответствующее место, условие сбоя и фактический результат, помогая
  модели определить следующий шаг. Строка состояния из второй главы
  может сводить последние диагностические сообщения, число вызовов и
  прогресс задачи.
\item
  \textbf{Проектируйте пути восстановления для изменений.} Ветки Git
  сохраняют изменения кода, песочницы ограничивают последствия проб и
  ошибок, а снимки восстанавливают управляемое состояние. Операции с
  внешними API обрабатываются с помощью механизмов идемпотентности,
  отмены и компенсации из четвёртой главы.
\end{enumerate}

\textbf{Проверяйте и результат, и процесс.} При устранении сбоя базы
данных очистка данных с последующим пересозданием может восстановить
сервис, но нарушит требование о сохранности данных. При исправлении
ошибки компиляции удаление соответствующей реализации может обеспечить
успешную компиляцию, но приведёт к потере функциональности. Эти примеры
показывают, что приёмка должна охватывать бизнес-цели, а процесс
выполнения --- подчиняться ограничениям. Harness может проверять
разрешение и область воздействия таких операций, как удаление
производственных данных и перезапись файлов, а также вызывать
представленный в четвёртой главе механизм независимой проверки.

При обучении уже выявленные нарушения также можно наказывать.
Рассматриваемый в восьмой главе RLVP --- штраф за путь проверки ---
включает ограничения процесса в обучающую обратную связь и направляет
модель к соблюдению требований выполнения при достижении результата. Во
время работы разрешения и проверки продолжают ограничивать доступные
действия, а обучающая обратная связь влияет на то, какие действия
предлагает модель.

\textbf{Управляйте распространением сбоев с учётом зависимостей.} Если
параллельно читаются три независимых файла и один из них отсутствует,
следует вернуть ошибку для этого файла, сохранив результаты чтения
остальных. Если какой-либо вызов зависит от вывода неудачной задачи,
соответствующую цепочку зависимостей следует остановить или
перепланировать. Harness должен определять область отмены по связям
между задачами, сообщать родительской операции состояние каждого
элемента и не отбрасывать вместе с ошибкой результаты независимой
работы. Реализация параллельных вызовов, потокового разбора и каскадной
отмены рассматривается далее.

\hypertarget{ux432ux43eux441ux441ux442ux430ux43dux43eux432ux43bux435ux43dux438ux435-ux43fux43eux441ux43bux435-ux441ux431ux43eux435ux432-ux438-ux43eux448ux438ux431ux43eux43a}{%
\subsection{Восстановление после сбоев и
ошибок}\label{ux432ux43eux441ux441ux442ux430ux43dux43eux432ux43bux435ux43dux438ux435-ux43fux43eux441ux43bux435-ux441ux431ux43eux435ux432-ux438-ux43eux448ux438ux431ux43eux43a}}

При выполнении длительной задачи соединение может оборваться между двумя
запросами к модели или после исполнения инструмента, но до сохранения
результата. Если Harness знает лишь о том, что «задача завершилась с
ошибкой», повторный запуск с начала может привести к дублированию уже
выполненных операций. Чтобы восстановить работу, сначала нужно
определить место прерывания, выяснить, какие результаты уже сохранены, а
затем решить, следует ли повторить попытку, исправить вызов, продолжить
с контрольной точки или завершить выполнение.

Первый вопрос при восстановлении: «Какие свидетельства остались после
предыдущего шага?» Если соединение оборвалось при чтении файла, обычно
его можно прочитать повторно. Если это произошло после отправки
удалённой задачи, она уже могла быть принята удалённой системой. При
наличии идентификатора задачи следует сначала запросить её состояние и
лишь затем решать, отправлять ли её повторно. Раздельная фиксация
состояний «запрос отправлен», «удалённая система подтвердила приём» и
«результат сохранён» создаёт отправную точку для восстановления.

\textbf{Выбирайте способ обработки в зависимости от места сбоя.}

\begin{longtable}[]{@{}
  >{\raggedright\arraybackslash}p{(\columnwidth - 4\tabcolsep) * \real{0.2430}}
  >{\raggedright\arraybackslash}p{(\columnwidth - 4\tabcolsep) * \real{0.3690}}
  >{\raggedright\arraybackslash}p{(\columnwidth - 4\tabcolsep) * \real{0.3880}}@{}}
\toprule\noalign{}
\begin{minipage}[b]{\linewidth}\raggedright
\hspace{0pt}Уровень
\end{minipage} & \begin{minipage}[b]{\linewidth}\raggedright
\hspace{0pt}Типичные сбои
\end{minipage} & \begin{minipage}[b]{\linewidth}\raggedright
\hspace{0pt}Основное направление действий
\end{minipage} \\
\midrule\noalign{}
\endhead
\bottomrule\noalign{}
\endlastfoot
\hspace{0pt}Уровень API & \hspace{0pt}Ограничение частоты HTTP 429,
перегрузка сервиса, тайм-\allowbreak{}\hspace{0pt}аут, разрыв
соединения, достижение предельной длины вывода & \hspace{0pt}Проверить
причину ошибки и состояние запроса, повторить попытку с задержкой или
скорректировать запрос \\
\hspace{0pt}Уровень инструментов & \hspace{0pt}Инструмент с таким именем
отсутствует, недопустимые параметры, ошибка исполнения, многократное
получение одной и той же ошибки & \hspace{0pt}Вернуть диагностические
данные, позволяющие локализовать проблему, исправить входные данные или
условия исполнения \\
\hspace{0pt}Уровень контекста & \hspace{0pt}Переполнение контекстного
окна, сбой сжатия, отсутствие пары между вызовом и результатом &
\hspace{0pt}Скорректировать объём входных данных, восстановить
фактическую запись, проверить структуру сообщений \\
\hspace{0pt}Уровень потока управления & \hspace{0pt}Повторение операций
без прогресса, рекурсивный сбой логики восстановления &
\hspace{0pt}Вести счётчики, ограничить глубину восстановления и
инициировать завершение или эскалацию \\
\end{longtable}

При ограничении частоты и временной перегрузке можно подождать и
повторить попытку. Ошибка параметров требует их изменения, недостаток
прав --- получения авторизации, а отсутствие инструмента --- повторного
выбора доступной возможности. Сопоставление типов ошибок с конкретными
действиями восстановления позволяет избежать многократного сбоя одного и
того же запроса при неизменных условиях.

\textbf{Одновременно отслеживайте сбои и прогресс.}

Harness может вычислять отпечаток вызова по сочетанию «имя инструмента +
нормализованные параметры» и выявлять повторяющиеся действия с учётом
результата и состояния задачи. Одинаковые вызовы могут быть как
обоснованным опросом состояния, так и циклом без прогресса. При оценке
необходимо также учитывать условие ожидания, внешнее состояние и
продолжительность происходящего. Для каждого пути восстановления следует
вести отдельный счётчик сбоев, ограничивая всю задачу общим бюджетом.

Для потоковых запросов необходимо отдельно задавать тайм-аут соединения,
тайм-аут простоя при чтении и общий предел времени. Тайм-аут соединения
ограничивает ожидание его установки, тайм-аут простоя при чтении ---
интервал между двумя порциями полученных данных, а общий предел ---
продолжительность всего запроса. Если соединение сохраняется, но данные
долго не поступают, сторожевой таймер простоя может закрыть запрос и
запустить процесс восстановления. \textbf{Каждое длительное соединение
должно контролироваться на активность; одного тайм-аута установки
соединения недостаточно.} Пороговые значения тайм-аутов должны учитывать
способ вывода модели: видимый текст, инкрементальные данные рассуждения
и сигналы активности могут показывать, продолжается ли запрос, а
продолжительное внутреннее рассуждение увеличивает время ожидания.

Проверка целостности контекста строится вокруг сопоставления ID вызовов
и результатов. Если результат отсутствует, сначала следует запросить
журнал исполнения и восстановить уже сохранённый фактический результат.
Если состояние исполнения неизвестно, это нужно явно зафиксировать и
действовать в соответствии с описанными в \crossreflink{chap:4}{главе 4}
механизмами идемпотентности и запроса состояния.

\textbf{Восстанавливайтесь поэтапно и сообщайте о текущем состоянии.}

\begin{enumerate}
\def\labelenumi{\arabic{enumi}.}
\tightlist
\item
  \textbf{Повторная попытка с задержкой.} Для допускающих повтор ошибок
  применяйте экспоненциальную задержку со случайным джиттером и
  соблюдайте указания сервера о времени ожидания. Выделяйте отдельные
  бюджеты повторных попыток для основной и вспомогательных задач.
  Фоновые операции, такие как генерация заголовков и подсказок ввода,
  после сбоя можно пропустить, чтобы они не отнимали ресурсы у основной
  задачи.
\item
  \textbf{Корректировка или переключение.} При достижении предельной
  длины вывода сначала проверьте разрешённый интерфейсом предел и бюджет
  задачи, а затем решите, следует ли увеличить лимит, разделить задачу
  или продолжить вывод. Если сервис остаётся недоступным, можно
  переключиться на резервную модель. Перед переключением преобразуйте
  сообщения, определения инструментов и специфичные для модели поля.
\item
  \textbf{Информирование о ходе восстановления.} Во время
  автоматического восстановления можно показывать сообщения «Повторная
  попытка» или «Переключение сервиса», сохраняя в журнале каждую ошибку.
  После успешного восстановления продолжите задачу. Если пользователю
  необходимо предоставить авторизацию или дополнительные данные либо
  выбрать следующий шаг, сообщите о текущем прогрессе, причине сбоя и
  уже предпринятых действиях.
\end{enumerate}

\textbf{Стратегия восстановления должна обеспечивать дальнейшее
продвижение всей задачи.} Восстановимые промежуточные ошибки следует
сохранять в журнале выполнения, а в интерфейсе --- показывать ход
восстановления. Сообщать пользователю и нижестоящим системам об
окончательном сбое нужно лишь после исчерпания способов восстановления.
Это позволяет нижестоящим системам различать состояния «выполняется
повторная попытка» и «задача завершилась с ошибкой», не запуская
компенсацию или повторную отправку преждевременно.

Ошибки инструментов могут передаваться в контекст как структурированные
результаты. Например, можно вернуть имя отсутствующего инструмента,
ошибочное поле и допустимые значения, чтобы модель могла исправить
вызов. Преобразования формата, разрешённые контрактом интерфейса и
имеющие однозначную семантику, можно выполнить программно и
зафиксировать в журнале. Если смысл параметров неоднозначен, следует
поручить модели сгенерировать их заново, чтобы исправление не изменило
цель операции.

\textbf{Подхватывайте незавершённую траекторию.}

При переключении модели необходимо передать принимающей стороне уже
подтверждённые факты задачи, результаты инструментов и незавершённые
действия. Роли сообщений, формат вызовов инструментов и поля рассуждения
различаются между интерфейсами, поэтому содержимое и состояние протокола
следует преобразовывать раздельно. Некоторые сервисы возвращают
доступное для чтения содержимое рассуждения или его сводку, а также
непрозрачные поля, предназначенные для последующих запросов к тому же
сервису. Такие поля могут быть прикреплены к блокам рассуждения или
вызовам инструментов. Требования к их дословной передаче в последующих
запросах определяются соответствующим интерфейсом.

\textbf{Системы, которым требуется продолжение работы между разными
поставщиками, должны сохранять рабочую траекторию в нейтральном
формате.} Способность модели понять текст другого поставщика не
означает, что целевой API примет непрозрачные протокольные поля
исходного поставщика. Отделение фактов задачи от протокола поставщика
позволяет использовать одну запись для подхвата задачи, воспроизведения
при оценке и обучающих данных. Нейтральная запись включает задачу и
сообщения, имена и параметры инструментов, стабильные идентификаторы
вызовов, состояния исполнения, результаты и их источники. Исходные
ответы поставщика сохраняются отдельно для продолжения работы через тот
же интерфейс и расследования проблем. При преобразовании в запрос к
целевой модели следует создать отображение ID вызовов и сохранить пары
вызовов с результатами. Доступные для чтения сводки передаются как
обычный контекст с указанием источника, а непрозрачные поля исходного
сервиса остаются в исходной записи.

Если целевой интерфейс не позволяет выразить прежний вызов, выполненные
шаги можно оформить как явную запись о состоянии задачи, например:
«Авиабилеты и проживание уже найдены; далее необходимо получить
стоимость питания и обменный курс», приложив фактические результаты.
После этого целевая модель должна создать новый допустимый вызов. Когда
принимающая сторона начнёт работу, возвращаемые ею протокольные поля
следует сохранять в соответствии с требованиями её собственного
интерфейса. Такая нейтральная запись также может использоваться для
воспроизведения при оценке в \crossreflink{chap:7}{главе 7},
формирования обучающих примеров в \crossreflink{chap:8}{главе 8} и
извлечения опыта в \crossreflink{chap:9}{главе 9}.

\begin{quote}
\textbf{Эксперимент 5-1 ★★★: подхват траектории между поставщиками}

В этом эксперименте AI-агент запрашивает стоимость авиабилетов,
проживания и питания, а также обменный курс, после чего рассчитывает
общую сумму в юанях. После первых двух вызовов инструментов скрипт в
заданной контрольной точке переключает сервис, имитируя передачу
управления после недоступности исходного сервиса. Принимающая сторона
продолжает выполнение оставшихся шагов. В эксперименте используются три
модели --- Moonshot, Anthropic и Google --- и сравниваются шесть
направленных комбинаций переключения.

Сравниваются три способа обработки: передача исходных сообщений с
отображением полей; преобразование сообщений после удаления содержимого
рассуждения и непрозрачных полей; сохранение состояния задачи в
нейтральной траектории с последующим восстановлением контекста согласно
целевому интерфейсу. Место переключения задаётся скриптом, а ошибки
формата принимающей стороны получены из фактических ответов сервисов.

В этой серии испытаний нейтральная траектория обеспечила успешный
подхват и правильный расчёт общей суммы во всех шести комбинациях. Два
других способа успешно завершили соответственно три и четыре комбинации.
Прогресс задачи был сохранён в результатах инструментов, поэтому ни в
одной из трёх групп повторных вызовов инструментов не возникло.
\end{quote}

\begin{quote}
\textbf{Эксперимент 5-2 ★★: продолжение после обрыва вывода}

В этом эксперименте предпринимается попытка прервать поток в трёх
местах: при выводе рассуждения, основного текста и параметров
инструмента. Сравниваются три способа: повторная отправка всего раунда,
продолжение с уже полученным основным текстом в качестве префикса и
добавление инструкции о продолжении. Для исполнения вызова инструмента
необходимы полные параметры. Усечённый фрагмент сначала сохраняется в
журнале выполнения, а после восстановления повторно разбирается и
проверяется.

Конкретная ошибка возникла в параметре города: изначально требовалось
сгенерировать «Токио», но в контрольной точке остался лишь фрагмент
\texttt{\{"city":"То}, а продолжение закрыло его как
\texttt{\{"city":"То"\}}. JSON стал допустимым, однако город остался
указан неверно.

В этом испытании с длинным основным текстом продолжение по префиксу
сократило объём последующего вывода в токенах. При прерывании параметров
результаты отличались: во всех трёх испытаниях Moonshot продолжение по
префиксу не смогло восстановить правильный город.
\end{quote}

\textbf{Задайте условия завершения цикла восстановления.}

Для каждого пути восстановления необходим верхний предел. Например,
после нескольких последовательных сбоев сжатия следует перейти к другой
стратегии работы с контекстом; если проверку прав невозможно завершить
--- передать задачу человеку; по достижении предельного числа раундов
или бюджета стоимости продолжения вывода --- остановиться. \textbf{Порог
повторных попыток следует определять по распределению восстановлений в
производственной среде.} Если подавляющее большинство временных сбоев
устраняется за первые несколько попыток, а вероятность успеха
последующих однотипных попыток приближается к нулю, следует остановиться
в этой точке перегиба, чтобы бесполезные вызовы не продолжали
расходовать квоту.

Сама логика восстановления также может завершиться с ошибкой. Например,
переполнение контекстного окна запускает хук остановки, который
обращается к модели для генерации описания коммита, снова переполняет
окно и повторно запускает тот же хук. Эту цепочку можно разорвать с
помощью признака повторного входа и ограничения глубины рекурсии,
разделив очистку после ошибки и вспомогательные задачи, требующие вызова
модели, а также выделив последним отдельный бюджет. Для задачи также
следует задать общие пределы числа итераций, времени и стоимости, чтобы
каждый путь восстановления оставался в рамках общего бюджета задачи.

\hypertarget{ux43fux440ux438ux451ux43cux44b-ux440ux435ux430ux43bux438ux437ux430ux446ux438ux438-coding-ux430ux433ux435ux43dux442ux430}{%
\subsection{Приёмы реализации
Coding-агента}\label{ux43fux440ux438ux451ux43cux44b-ux440ux435ux430ux43bux438ux437ux430ux446ux438ux438-coding-ux430ux433ux435ux43dux442ux430}}

Сначала по записи выполнения задачи по исправлению падающего теста
оценим, какой участок стоит оптимизировать. В докладе автора «The Road
to JARVIS», представленном в Порту в ноябре 2026 года, одна и та же
сессия воспроизводилась с разной скоростью инференса. При расчёте по
времени работы доступной тогда онлайн-модели общее время составило 24,2
секунды. При неизменном времени работы инструментов и расчётной скорости
инференса ещё не созданного специализированного чипа результат составил
около 8,1 секунды. Из них 7,5 секунды занял запуск тестов. Даже если
устранить оставшиеся примерно 0,6 секунды работы модели, всё равно
придётся ждать завершения тестов.

Какие операции из этих 7,5 секунды можно выполнять одновременно? Это
зависит от связей между вызовами инструментов.

\textbf{Параллельные вызовы инструментов, потоковое выполнение и
каскадное прерывание.}

Поиск в коде, чтение конфигурации и просмотр журналов зачастую не
зависят друг от друга. Когда модель генерирует такие вызовы в одном
ответе, Harness может совместить генерацию вызовов с выполнением
инструментов: как только интерфейс подтвердит, что первый вызов выведен
полностью, а его параметры и разрешения прошли проверку, Harness
немедленно ставит его на выполнение и продолжает принимать последующие
вызовы. Полноту следует определять по событию завершения вызова,
предусмотренному интерфейсом, а не по тому, что временно полученный JSON
уже удаётся разобрать: иначе вызов может быть выполнен преждевременно.

Независимые вызовы можно выполнять параллельно. Например, одновременно
искать имя функции и читать конфигурационный файл, а решение об
изменениях принимать после получения обоих результатов. Если выбор
читаемого файла зависит от результатов поиска, его чтение переносится на
следующий шаг. Планировщик также должен учитывать общее состояние:
одновременное изменение одного файла двумя инструментами приведёт к
конфликту записи, даже если их входные данные не зависят друг от друга.

В определении инструмента можно объявить возможности конкурентного
выполнения, используемые ресурсы и способ отмены. В учебной реализации
инструменты без таких объявлений можно выполнять последовательно. Если
один вызов завершается с ошибкой, следует остановить последующие
операции, зависящие от его результата, сохранить результаты независимых
вызовов и вернуть родительской операции состояние каждого из них. При
отмене уже запущенной задачи также необходимо зафиксировать, сработала
ли отмена и успели ли произойти изменения файлов или внешние операции.

\textbf{Точное управление контекстом.}

Большую кодовую базу можно добавлять в контекст постепенно посредством
многоуровневого поиска: сначала найти файлы, затем прочитать нужные
функции, а при необходимости перейти к местам вызова и тестам. Каждое
чтение дополняет сведения, необходимые для текущего решения, чтобы
модель одновременно видела код, ограничения и обратную связь от
выполнения.

Инструмент чтения должен поддерживать возврат фрагмента по диапазону
строк, например со 100-й по 150-ю, и указывать для каждой строки её
фактический номер в файле. Благодаря этому модель может точно сослаться
на «строку 142», а инструмент последующего редактирования --- найти
нужное место. В возвращаемом содержимом также следует указывать исходный
файл и прочитанный диапазон. Если файл изменился, позицию редактирования
необходимо подтвердить заново.

Компиляция или тестирование могут вывести тысячи строк. Здесь можно
применить описанный в четвёртой главе механизм обрезки и сохранения:
вернуть начало и конец вывода с указанием количества пропущенных строк,
а полную версию сохранить для чтения по мере необходимости. Ошибка может
находиться в середине, поэтому агент сможет дополнительно найти её по
ключевым словам или прочитать соответствующий диапазон.

\textbf{Динамическое добавление информации о среде.}

Строка состояния агента из второй главы применительно к задачам
программирования может объединять четыре типа сведений:

\begin{itemize}
\tightlist
\item
  \textbf{Текущий рабочий каталог}: задаёт базу для относительных путей.
\item
  \textbf{Ветка Git}: показывает, в какой ветке разработки находятся
  текущие изменения.
\item
  \textbf{Последние коммиты}: помогают понять недавние изменения,
  связанные с задачей.
\item
  \textbf{Сводка проиндексированных и непроиндексированных изменений}:
  позволяет различать подготовленное к коммиту содержимое и изменения
  рабочего дерева.
\end{itemize}

Harness может обновлять сводку после операций, изменяющих это состояние,
и добавлять её в конец контекста при следующем вызове модели. Реже
меняющиеся сведения, например последние коммиты, можно обновлять по
необходимости, сокращая число повторных запросов и объём входных данных.
Стабильность начальной части контекста способствует повторному
использованию кэша префикса.

\textbf{Сохранение состояния среды выполнения команд.}

Переход в другой каталог, активация виртуального окружения, установка
переменных среды и запуск фоновых сервисов влияют на последующие
команды. Если агент выполняет \texttt{cd} в shell, изменение каталога
естественным образом сохраняется только тогда, когда следующие команды
продолжают использовать ту же сессию. При выполнении в отдельном
процессе рабочий каталог определяется конфигурацией нового процесса.
Переменные среды, создаваемые при активации виртуального окружения,
подчиняются тем же границам сессии.

\textbf{Для интерактивной отладки и задач, требующих длительно
работающих сервисов, автор рекомендует по умолчанию использовать
постоянный терминал.} Он сохраняет рабочий каталог, переменные среды и
состояние сессии, соответствуя привычке разработчика последовательно
работать в одном терминале. Harness назначает сессии идентификатор и
отслеживает рабочий каталог, работающие процессы и статусы их
завершения, чтобы модель знала, в какой среде будет выполнена следующая
команда. Другой вариант реализации --- явно передавать каждой команде
рабочий каталог и конфигурацию среды, поддерживая согласованность
посредством воспроизводимой конфигурации.

Для параллельных задач следует использовать отдельные сессии или явные
снимки среды, чтобы смена каталога в одной задаче не повлияла на другую.
После завершения задачи фоновые процессы необходимо централизованно
остановить, а при возобновлении --- повторно проверить, жива ли сессия.
Например, если на следующий день нужно продолжить запуск тестов, код и
список зависимостей можно восстановить из рабочего дерева, но сначала
следует проверить, работает ли ещё сервисный процесс, запущенный вчера в
терминале. Именно эти два вида состояния необходимо различать при оценке
стоимости песочницы в четвёртой главе.

\textbf{Механизм немедленной синтаксической обратной связи.}

После записи файла слой инструментов может автоматически запускать
соответствующий языку синтаксический анализатор или linter и передавать
диагностику модели как часть результата инструмента. Например, при
отсутствии скобки возвращаются файл, номера строки и столбца и тип
ошибки, чтобы модель могла исправить её уже в следующем цикле, получая
обратную связь, подобную мгновенной подсветке в IDE.

Область проверки можно расширять вместе с областью изменений: сначала
проверить изменённые файлы, затем запустить тесты затронутых модулей и,
наконец, выполнить регрессионные проверки, предусмотренные требованиями
к передаче результата. Легковесные проверки помогают раньше обнаруживать
локальные проблемы, а интеграционные тесты проверяют взаимодействие
между модулями.

\hypertarget{ux438ux43dux441ux442ux440ux443ux43cux435ux43dux442ux44b-ux43fux43eux438ux441ux43aux430-ux432-coding-ux430ux433ux435ux43dux442ux435}{%
\subsection{Инструменты поиска в
Coding-агенте}\label{ux438ux43dux441ux442ux440ux443ux43cux435ux43dux442ux44b-ux43fux43eux438ux441ux43aux430-ux432-coding-ux430ux433ux435ux43dux442ux435}}

Если сообщение об ошибке содержит текст «отсутствует поле суммы», можно
сначала найти исходную формулировку, а затем от места возникновения
ошибки проследить вызывающий код. Если же пользователь лишь говорит, что
«сортировка работает неправильно», а в коде используются английские
имена, буквальный поиск этой фразы почти бесполезен. Сначала можно найти
интерфейс, отображающий результат, а затем по полям, связям вызовов и
тестовым данным добраться до логики сортировки. Четыре типа инструментов
поиска на \crossreflink{fig:5-3}{рисунке 5-3} используют разные
подсказки; результаты первого поиска зачастую определяют, какой
инструмент следует выбрать для следующего.

\begin{figure}
\hypertarget{fig:5-3}{%
\centering
\includegraphics{images/revised/fig5-3.pdf}
\caption{Рисунок 5-3. Сравнение инструментов поиска
Coding-агента}\label{fig:5-3}
}
\end{figure}

\textbf{Сопоставление содержимого по регулярным выражениям.} grep и
ripgrep ищут содержимое по текстовым шаблонам и подходят для задач, в
которых известны имя функции, имя переменной или сообщение об ошибке.
Например, \texttt{def\ handle.*} может найти строки текста с
соответствующей формой определения функции. Фильтрация по типам файлов и
путям позволяет также искать только в файлах Python, исключать каталоги
тестов и сосредоточивать результаты на релевантной реализации.

Текстовое сопоставление непосредственно соответствует содержимому
файлов, поэтому результаты легко проверить. Поисковый запрос необходимо
связать с формулировками в коде: при поиске «аутентификации
пользователя» можно дополнительно попробовать такие имена, как
\texttt{authenticate}, \texttt{login} или \texttt{check\_credentials}.
Когда модель расширяет запрос с учётом предыдущих результатов, поиск и
понимание кода образуют единый итеративный процесс.

\textbf{Сопоставление шаблонов имён файлов.} glob ищет файлы по шаблону
пути: например, \texttt{**/*.test.ts} находит тестовые файлы TypeScript
в каталогах любого уровня, а \texttt{src/components/**/Button.tsx} ---
файл Button.tsx в каталогах компонентов. Сопоставление путей не требует
чтения содержимого и подходит для первоначального изучения структуры
проекта с последующим выбором файлов, которые нужно открыть.

\textbf{Семантический поиск по коду.} Если задача описана только через
функциональность, семантически связанный код можно находить с помощью
векторного поиска. Например, «аутентификация пользователя» может
соответствовать функции \texttt{check\_credentials}. Системам такого
типа обычно необходимо решать две задачи:

\begin{itemize}
\tightlist
\item
  \textbf{Разбиение по структуре кода.} Основными единицами служат
  функции, классы и методы; вместе с ними сохраняются связанные
  сигнатуры, комментарии и сведения о модуле. При дальнейшем разделении
  длинных единиц необходимо добавлять контекст, достаточный для
  понимания фрагмента.
\item
  \textbf{Сочетание нескольких поисковых сигналов.} Векторный поиск
  находит семантически связанный код, поиск по ключевым словам --- код с
  точным совпадением имён, после чего две группы результатов
  объединяются. Описанный в третьей главе перекрёстный кодировщик
  позволяет дополнительно выполнить ранжирование: он одновременно
  получает запрос и фрагмент кода, а затем оценивает релевантность этого
  фрагмента.
\end{itemize}

При исследовательских запросах вроде «взаимодействует с базой данных»
или «обрабатывает проверку пользовательского ввода» можно сначала найти
модуль семантическим поиском, а затем прочитать реализацию.

\textbf{При создании поисковых возможностей Coding-агента автор
рекомендует сначала качественно реализовать Grep, Glob и чтение
заданного диапазона.} Модель может на основании результатов поиска
переформулировать запрос, открыть вызывающий код и тесты и постепенно
сузить область. Эти инструменты непосредственно читают текущее рабочее
пространство, поэтому не требуют поддерживать индекс представлений и
синхронизировать его с кодом. Семантический индекс для улучшения полноты
поиска стоит добавлять при многократном исследовании крупных кодовых баз
или когда задача описана только через функциональность и точку входа
найти трудно. Cursor ранее публично рассказывал об индексировании
кодовой базы на основе представлений, а в последующей документации о
поиске описал локальный индекс Instant Grep и исследовательских
субагентов\footnote{Cursor, \emph{Securely indexing large codebases} и
  Search --- материалы соответственно об индексировании на основе
  представлений и последующих механизмах локального поиска.}. Эта
эволюция также показывает, что «использование интерфейса, подобного
grep» и «создание индекса» --- два независимых архитектурных измерения:
один и тот же поисковый интерфейс может использовать сканирование или
индекс.

\textbf{Поиск определений и ссылок на уровне символов.} Языковой сервер
может использовать сведения о синтаксисе, областях видимости и типах для
поддержки функций «перейти к определению» и «найти ссылки». Например,
если \texttt{authenticate} определена в строке 42 и вызывается в строке
189, инструмент может вернуть связь между этими местами и отличить их от
одноимённых символов в других областях видимости. Плагины
интеллектуального анализа кода Claude Code предоставляют один из
способов такой интеграции: через протокол языковых серверов (LSP) они
обеспечивают навигацию по символам и диагностику после
редактирования\footnote{Anthropic, \emph{Code intelligence plugins} ---
  материал о предоставлении через LSP поиска определений и ссылок, а
  также диагностики в реальном времени.}.

Эффективность инструментов работы с символами зависит от поддержки
языка, конфигурации проекта и состояния анализа. Агент может дополнять
их текстовым поиском строк, комментариев или файлов, которые языковой
сервер ещё не разобрал.

\hypertarget{ux438ux43dux441ux442ux440ux443ux43cux435ux43dux442ux44b-ux440ux435ux434ux430ux43aux442ux438ux440ux43eux432ux430ux43dux438ux44f-ux444ux430ux439ux43bux43eux432-ux432-coding-ux430ux433ux435ux43dux442ux435}{%
\subsection{Инструменты редактирования файлов в
Coding-агенте}\label{ux438ux43dux441ux442ux440ux443ux43cux435ux43dux442ux44b-ux440ux435ux434ux430ux43aux442ux438ux440ux43eux432ux430ux43dux438ux44f-ux444ux430ux439ux43bux43eux432-ux432-coding-ux430ux433ux435ux43dux442ux435}}

Представьте: после того как модель прочитала файл, кто-то добавил
комментарий в самое его начало. Будут ли теперь указания «изменить
строку 28» и «заменить этот уникальный фрагмент исходного текста»
по-прежнему относиться к одному и тому же месту? А если исходный текст
встречается по одному разу в двух разных функциях? Пять подходов к
редактированию на \crossreflink{fig:5-4}{рисунке 5-4} по-разному
определяют место изменения: одни опираются на номера строк, другие ---
на исходный текст, а третьи поручают объединение другой модели.

\begin{figure}
\hypertarget{fig:5-4}{%
\centering
\includegraphics{images/revised/fig5-4.pdf}
\caption{Рисунок 5-4. Сравнение пяти подходов к редактированию
файлов}\label{fig:5-4}
}
\end{figure}

\textbf{Описание различий + Apply Model.} Основная модель передаёт
намерение изменения и фрагменты кода: например, описание удалений и
добавлений, похожее на Git diff, либо каркас кода, где неизменённые
части пропущены с помощью маркеров многоточия. Затем специализированная
модель применения объединяет это с исходным файлом и генерирует полное
изменённое содержимое. Ранний Fast Apply в Cursor --- характерный пример
такого подхода\footnote{Cursor, \emph{Editing Files at 1000 Tokens per
  Second}, описание специализированной модели, применяющей описание
  изменений к исходному файлу.}. Он распределяет «формулирование
изменения» и «применение изменения» между двумя моделями. По мере того
как базовые модели стали лучше генерировать точные локальные правки, у
этого двухмодельного процесса начали появляться более простые
альтернативы.

Такое разделение позволяет основной модели сосредоточиться на выражении
логики кода, а модели применения --- на объединении локальных изменений
с файлом. Соответственно, основное внимание при проверке уделяется
результату объединения: несколько похожих фрагментов могут создать
неоднозначность позиционирования, маркеры пропусков могут быть
истолкованы неверно, а неизменённые области --- случайно переписаны.

\textbf{Старая строка → новая строка (Old String → New String).}
Инструмент Edit в Claude Code использует точную замену строк; в версии
0.50 за май 2025 года Cursor добавил Search \& Replace, который сначала
использовался для редактирования длинных файлов моделями
Anthropic\footnote{Anthropic, \emph{Tools reference}, инструмент Edit;
  журнал обновления Cursor 0.50 от 15 мая 2025 года, где представлено
  добавление Search \& Replace; в журнале обновления 1.3 также
  зафиксированы улучшения Search \& Replace и Apply.}. При таком подходе
основание для позиционирования и содержимое замены передаются
непосредственно программе, что избавляет от дополнительного вызова
модели на этапе объединения. Модель предоставляет \texttt{old\_string} и
\texttt{new\_string}, а инструмент находит исходный текст в файле и
заменяет его. Интерфейс может требовать, чтобы при единичной замене
исходный текст был уникален. Если совпадений нет или их несколько, он
возвращает позиции и диагностическую информацию, чтобы модель добавила
контекст. Для пакетной замены область действия задаётся явным
параметром.

Например, в конфигурационном файле для двух сервисов указано
\texttt{timeout\ =\ 30}. Требуется изменить на 60 секунд только тайм-аут
сервиса загрузки. Если инструмент требует уникального совпадения,
передача одного лишь \texttt{timeout\ =\ 30} найдёт два совпадения, и
файл должен остаться без изменений. В следующем запросе модель включает
в старую строку и заголовок раздела сервиса загрузки, и строку
настройки, а в новой строке сохраняет заголовок, изменяя только
значение.

Сопоставление по содержимому наглядно и удобно для проверки. При вызове
инструмента модель должна вывести старую строку целиком: при удалении
большого блока кода параметр становится длинным, а небольшие различия в
пробелах, переносах строк или кавычках также влияют на совпадение.
Пример с фигурными кавычками Unicode из четвёртой главы показывает, как
подобные сбои возникают при передаче параметров.

\textbf{Позиционирование по номерам строк (Old Line Numbers → New
String).} Модель указывает удалить строки с X по Y, а затем вставить
новый текст. Если инструмент чтения уже предоставляет фактические номера
строк, для удаления большого блока достаточно двух позиций --- начальной
и конечной. Значение номеров строк зависит от версии файла на момент
чтения: предыдущее изменение может сместить последующие позиции.

Поэтому при пакетном редактировании все правки могут ссылаться на один
исходный снимок, а инструмент --- проверять пересечение диапазонов и
применять изменения целиком. Другой вариант --- вносить изменения
последовательно, каждый раз заново читать файл и только затем определять
позиции. Инструмент также может сверять версию файла или ожидаемое
содержимое и после конкурентного изменения возвращать конфликт, чтобы
старые номера строк не применялись к новому содержимому файла.

\textbf{Команды редактирования в стиле Vim.} Команды редактора позволяют
перемещать, копировать, вырезать и вставлять содержимое, поэтому хорошо
подходят для реорганизации кода, например переноса функции из одного
места в другое. Они зависят от текущего положения курсора, буфера и
режима команд, причём каждый шаг изменяет точку отсчёта для последующих
операций.

\textbf{Интерактивные привычки человека при работе с редактором не
обязаны становиться архитектурой инструментов модели по умолчанию.}
Человек обычно смотрит на текущее состояние, а затем выполняет одну
небольшую операцию. Модель же может после длительного рассуждения за
один раз сгенерировать изменения на десятки или даже сотни строк.
Разбиение такой пакетной правки на множество перемещений курсора и
переключений режимов повышает стоимость вызовов и требует от модели
прогнозировать, как предыдущие команды повлияют на состояние. Поэтому
для задач генерации кода автор отдаёт предпочтение интерфейсам, напрямую
выражающим цель и содержимое изменения. При использовании команд
редактора необходимо явно задавать состояние, служащее точкой отсчёта, и
своевременно возвращать обратную связь.

\textbf{Сопоставление начала и конца старой строки (Old String Start +
End → New String).} Модель передаёт только несколько строк в начале и
конце заменяемой области, а инструмент находит обе границы и содержимое
между ними. При удалении сотен строк кода такой способ позволяет
сократить объём повторно выводимого исходного текста, сохранив основание
для позиционирования по содержимому.

Инструмент должен проверить порядок появления двух фрагментов исходного
текста и в соответствии с контрактом интерфейса определить, включаются
ли начальный и конечный фрагменты в замену. Изменение выполняется только
тогда, когда условиям соответствует ровно один диапазон. Если совпадений
несколько, инструмент сообщает модели их позиции и просит добавить
контекст. Для вложенных конструкций или повторяющегося кода определить
диапазон помогают дополнительные соседние якоря, например имена и
сигнатуры функций.

\textbf{При создании нового Coding-агента автор рекомендует начать с
простого и детерминированного интерфейса редактирования наподобие замены
старой строки новой.} Позиционирование по исходному тексту однозначно,
причины сбоя легко объяснить, а различия удобно проверять. Когда модель
способна надёжно генерировать локальные изменения, Harness может
передать больше работы детерминированной программе и отказаться от
дополнительной модели применения. Если при удалении крупных блоков
потребуется сократить объём параметров, можно добавить начальные и
конечные якоря либо редактирование по номерам строк на основе снимка.
Инструмент \texttt{apply\_patch} в Codex демонстрирует ещё один способ
прямого редактирования: несколько изменений выражаются с помощью патча с
ограниченным синтаксисом.

\hypertarget{ux431ux435ux437ux43eux43fux430ux441ux43dux43eux441ux442ux44c-coding-ux430ux433ux435ux43dux442ux430}{%
\subsection{Безопасность
Coding-агента}\label{ux431ux435ux437ux43eux43fux430ux441ux43dux43eux441ux442ux44c-coding-ux430ux433ux435ux43dux442ux430}}

Coding-агент может читать проект, изменять файлы, запускать программы, а
иногда и обращаться к внешним сервисам по сети. При проектировании
разрешений необходимо проследить поток данных и проверить три аспекта:
какое содержимое попадает в контекст, какие операции может запрашивать
модель и какие механизмы в итоге одобряют и выполняют эти операции.

Simon Willison сформулировал риск одного из видов инъекции в промпт в
виде «смертельной триады»: доступ к приватным данным, контакт с
недоверенным содержимым и возможность внешнего
взаимодействия\footnote{Simon Willison, \emph{The lethal trifecta for AI
  agents: private data, untrusted content, and external communication},
  16 июня 2025 года.}. Например, если помощник может читать приватные
файлы, обрабатывать внешние письма и отправлять сообщения, между
чтением, интерпретацией и отправкой необходимо установить явные границы
авторизации. Если инструкции из внешнего содержимого будут приняты за
требования пользователя, это может привести к несанкционированной
передаче данных.

Помимо этих трёх аспектов автор рекомендует проверять и четвёртый:
\textbf{постоянная память усиливает последствия инъекции в промпт}. Если
внешнее содержимое попадёт в долговременную память, его влияние может
сохраниться между сессиями и вновь проявиться в последующих задачах.
Поэтому при записи в память следует фиксировать источник, область
применимости и статус подтверждения, а при чтении сохранять эти
атрибуты. Достоверность текста, сохранённого в \texttt{MEMORY.md},
по-прежнему зависит от процесса его появления. После обнаружения
ошибочной записи также должна существовать возможность найти, исправить
и отозвать связанное с ней содержимое.

Три уровня защитных барьеров из первой главы выполняют здесь разные
задачи: уровень контекста обозначает источники информации и требования к
поведению, уровень выполнения проверяет операции и авторизацию, а
уровень данных ограничивает чтение и запись ресурсов в соответствии с
доверенной идентичностью. Песочница Claude Code ограничивает выполнение
команд границами файловой системы и сети, а Cowork также выполняет
локальные задачи в изолированной среде\footnote{Anthropic, \emph{Making
  Claude Code more secure and autonomous with sandboxing} и \emph{How we
  contain Claude across products} --- материалы о границах файловой
  системы, сети и изоляции продуктов.}. Ниже эти механизмы
рассматриваются подробнее.

\textbf{Ограничение ресурсов и взаимодействий в песочнице.}

\begin{itemize}
\tightlist
\item
  \textbf{Сетевой выход.} Это одна из границ песочницы для кода, которой
  легко пренебречь, но которую следует реализовать в первую очередь.
  Автор рекомендует начинать с отключённого по умолчанию внешнего
  доступа и по мере необходимости разрешать источники пакетов, сайты
  документации и бизнес-API. Список разрешений должен соответствовать
  конкретным адресатам и области допустимых операций; для общих сервисов
  также необходимо учитывать принадлежность учётной записи, арендатора и
  ресурсов. Сетевая политика должна распространяться на запущенные
  инструментами программы и их дочерние процессы, а передаваемые наружу
  данные следует отдельно проверять на соответствие области авторизации.
\item
  \textbf{Файловая система.} Исходный код можно подключить только для
  чтения, а его копию поместить в доступное для записи рабочее
  пространство. Другой вариант --- создавать патчи через контролируемый
  интерфейс редактирования и применять их после проверки. Результаты и
  промежуточные файлы размещаются в отдельном каталоге, а файлы учётных
  данных остаются за пределами песочницы. Если требуется доступ к
  сервису аутентификации, контролируемый интерфейс должен предоставлять
  только необходимые для задачи возможности, чтобы не передавать
  сгенерированной программе все учётные данные хоста.
\item
  \textbf{Ресурсы и время.} Следует ограничить CPU, память, число
  процессов, дисковое пространство и время выполнения, чтобы уменьшить
  последствия бесконечных циклов, неконтролируемого роста числа
  процессов и непрерывной записи на диск. При превышении лимита
  возвращается структурированный статус, например «выполнение превысило
  120 секунд и было остановлено», вместе с сохранённым выводом и
  результатами очистки. По ним модель сможет решить, следует ли сузить
  задачу или изменить реализацию.
\end{itemize}

Эти границы должны охватывать всю цепочку выполнения. Даже после
одобрения команды фактически запущенная программа должна работать в
пределах тех же разрешений; соответствующий контроль также должен
распространяться на фоновые процессы, вспомогательные инструменты и
внешние интерфейсы.

\textbf{Разбор структуры команд и назначения аргументов.}

Shell-команда состоит из имени программы, аргументов, перенаправлений,
конвейеров и подвыражений. При проверке необходимо разобрать эти
структуры, а затем определить связанные с ними операции чтения, записи,
выполнения и взаимодействия. \textbf{Чёрного списка команд недостаточно
для окончательного решения о безопасности.} Одну и ту же операцию можно
выразить разными сочетаниями аргументов и косвенными вызовами. Проверка
должна учитывать структуру команды и назначение аргументов, а разрешения
среды выполнения --- ограничивать фактическую область выполнения.

Например, инструмент поиска файлов может иметь аргумент, запускающий
последующую операцию для каждого совпадения, а параметр вывода
инструмента загрузки определяет место записи. Валидатор должен
распознавать, «какой параметр использует следующий аргумент как своё
значение», отслеживать назначение аргументов, а затем проверять целевой
путь, разрешения доступа и допустимые операции. Тогда результат проверки
можно конкретизировать: «какие файлы будут прочитаны», «куда будет
выполнена запись» и «какая программа будет запущена».

Из-за подстановки переменных, динамических скриптов и путей, возникающих
только во время работы, часть фактического поведения можно определить
лишь при выполнении. Для неопределённых частей Harness может потребовать
более явную форму вызова, передать их на проверку или выполнить в
ограниченной среде. Проверка с помощью модели обеспечивает контекстную
оценку, а детерминированное управление доступом ограничивает фактические
последствия. Описанный в четвёртой главе Sidecar может работать
параллельно с подготовительными операциями, а действия с внешними
последствиями отправляются на выполнение только после успешной проверки.

\textbf{Явное определение обслуживаемой стороны при многостороннем
делегировании.}

Лояльность принципалу (principal loyalty) отвечает на вопрос: как агент,
действующий от имени одной стороны и одновременно общающийся с другими
участниками, может последовательно соблюдать цели и полномочия
доверителя. Например, пользователь поручает агенту вести переговоры о
цене с поставщиком и конфиденциально сообщает предельный бюджет в 12 000
юаней. Агент должен понимать предложения поставщика и объяснять позицию
своей стороны, одновременно защищая это приватное ограничение.

При оценке многосторонних взаимодействий автор и его соавторы наблюдали
два типа ошибок. В первом случае агент раскрывал приватные сведения
доверителя или принимал неправомерные требования другой стороны. Во
втором он проявлял чрезмерную осторожность и не мог выполнить даже
обычные запросы доверителя. Оценённые стратегии промптов и методы
обучения продемонстрировали компромисс между защитой информации и
нормальным взаимодействием\footnote{Li, Bojie and Noah Shi, \emph{Whose
  Side Is Your Agent On? Multi-Party Principal Loyalty in LLM Agents},
  2026.}. \textbf{Лояльность необходимо оценивать одновременно по защите
информации и выполнению задачи.} Если стремиться лишь к увеличению числа
отказов, будет заблокировано и нормальное взаимодействие; если
стремиться только к беспрепятственному продолжению диалога, другая
сторона сможет повлиять на интересы доверителя.

Для Coding-агента файлы репозитория, вывод инструментов и ответы
сторонних MCP-сервисов могут содержать как сведения, относящиеся к
задаче, так и текст с требованиями изменить поведение. Harness должен
различать инструкции доверителя и внешние материалы по их
действительному источнику, снабжать сообщения сведениями об источнике и
роли и преобразовывать область делегирования в разрешения на выполнение.
Если внешний материал утверждает, что содержит «требование
пользователя», это необходимо подтвердить через доверенный канал.

В промпте можно дополнительно сформулировать правила взаимодействия:
защищать приватные сведения доверителя, причём иногда конфиденциальным
должен оставаться даже сам факт существования определённой записи;
раскрывать внешним сторонам только необходимые сведения; различать
внутренние пределы и публичную позицию; продвигать задачу в пределах
предоставленных полномочий; сохранять единые стандарты авторизации при
повторяющихся запросах. Расширение полномочий или смена обслуживаемой
стороны должны подтверждаться через доверенный канал и сопровождаться
соответствующим обновлением разрешений среды выполнения. Так модель
получает явные основания для поведения, а инструменты и среда
обеспечивают фактические границы доступа.

\hypertarget{ux43aux43eux434-ux43cux435ux442ux430ux441ux43fux43eux441ux43eux431ux43dux43eux441ux442ux44c-ux443ux43dux438ux432ux435ux440ux441ux430ux43bux44cux43dux43eux433ux43e-ai-ux430ux433ux435ux43dux442ux430}{%
\section{Код: метаспособность универсального
AI-агента}\label{ux43aux43eux434-ux43cux435ux442ux430ux441ux43fux43eux441ux43eux431ux43dux43eux441ux442ux44c-ux443ux43dux438ux432ux435ux440ux441ux430ux43bux44cux43dux43eux433ux43e-ai-ux430ux433ux435ux43dux442ux430}}

В предыдущих разделах требования пользователя преобразовывались в
изменения кода. Теперь вернёмся к таблице продаж из начала главы: даже
если пользователю в итоге нужен только отчёт, агент всё равно может
написать программу, которая унифицирует идентификаторы товаров,
группирует данные и рассчитывает статистику, а затем повторно вычисляет
долю возвратов для одного товара по исходной таблице. Сначала нужно
уточнить, следует ли вести расчёт по количеству товаров или по
количеству заказов. Предположим, определённый товар продан в количестве
десяти штук в двух заказах, причём из одного заказа вернули одну штуку.
При расчёте по количеству товаров результат составит 1/10, а при расчёте
по заказам, в которых был возврат, --- 1/2. Оба алгоритма успешно
выполняются, но только предварительное уточнение методики позволяет
понять, какое из значений отвечает на вопрос пользователя. Такую
программу можно будет снова запустить, когда поступят данные за
следующий период. Этот цикл «генерация, выполнение, проверка» применим
также к созданию контента, интеграции бизнес-систем и разработке
агентов.

\begin{quote}
\textbf{Метаспособность.} В этой книге метаспособностью
(meta-capability) называется способность создавать или комбинировать
другие способности. Генерация кода позволяет создавать новые
инструменты, например скрипты обработки данных и процессы вызова API;
выражать новые ограничения, например утверждения и правила проверки; а
также формировать новые способы представления, например HTML-формы,
презентации и видео.
\end{quote}

Начнём с простейшего случая. Даже если конечным результатом должен быть
только ответ на одну задачу, модель может написать небольшую программу
для проверки собственных вычислений.

\hypertarget{ux43aux43eux434-ux43aux430ux43a-ux438ux43dux441ux442ux440ux443ux43cux435ux43dux442-ux43cux44bux448ux43bux435ux43dux438ux44f}{%
\subsection{Код как инструмент
мышления}\label{ux43aux43eux434-ux43aux430ux43a-ux438ux43dux441ux442ux440ux443ux43cux435ux43dux442-ux43cux44bux448ux43bux435ux43dux438ux44f}}

\textbf{Для точных вычислений и формальных ограничений автор
предпочитает сначала поручить модели составить исполняемое выражение, а
затем передать его интерпретатору или решателю.} При пословной генерации
ответа легко допустить ошибку при переписывании или вычислении, а также
упустить условие; программа позволяет зафиксировать эти шаги и выполнять
их многократно. Задачи по математике и логике обычно состоят из двух
этапов: сначала условие преобразуется в отношения и ограничения, затем
вычисляется результат. LLM может помочь понять задачу, предложить
решение и написать программу, а интерпретатор или решатель --- выполнить
заданные операции.

Например, в классе 40 учеников: 60 \% выбрали математику, 45 \% ---
физику, а 25 \% --- оба предмета. Чтобы найти число учеников, выбравших
только физику, нужно из числа выбравших физику вычесть число выбравших
оба предмета: 18 − 10 = 8. Если по ошибке использовать число выбравших
математику, получится 24 − 10 = 14. Ошибка возникает на этапе выбора
уменьшаемого.

Если выразить количественные отношения с помощью осмысленных переменных,
каждое из них можно проверить отдельно:

\begin{Shaded}
\begin{Highlighting}[]
\NormalTok{math }\OperatorTok{=} \DecValTok{40} \OperatorTok{*} \DecValTok{60} \OperatorTok{//} \DecValTok{100}       \CommentTok{\# 24 ученика выбрали математику}
\NormalTok{phys }\OperatorTok{=} \DecValTok{40} \OperatorTok{*} \DecValTok{45} \OperatorTok{//} \DecValTok{100}       \CommentTok{\# 18 учеников выбрали физику}
\NormalTok{both }\OperatorTok{=} \DecValTok{40} \OperatorTok{*} \DecValTok{25} \OperatorTok{//} \DecValTok{100}       \CommentTok{\# 10 учеников выбрали оба предмета}
\NormalTok{only\_phys }\OperatorTok{=}\NormalTok{ phys }\OperatorTok{{-}}\NormalTok{ both     }\CommentTok{\# 8 учеников выбрали только физику}
\end{Highlighting}
\end{Shaded}

Символьные вычисления дополнительно расширяют такое разделение
обязанностей. Они позволяют сохранять точные формы вроде \(\sqrt{2}\) и
преобразовывать выражения по правилам алгебры. Если сразу возвести это
значение в квадрат, результатом будет целое число 2. Если же сначала
округлить его до 1,414, а затем возвести в квадрат, получится 1,999396.
Когда задача требует точного упрощения, следует сохранять символьную
форму; если нужно сообщить результат измерения, численное значение можно
вычислить в конце с заданной точностью. Рассматривая интеграцию LLM с
Wolfram\textbar Alpha, Wolfram подчёркивал взаимодополняемость понимания
языка и вычислительных возможностей\footnote{Stephen Wolfram,
  \emph{Wolfram\textbar Alpha as the Way to Bring Computational
  Knowledge Superpowers to ChatGPT}, январь 2023 года.}. AI-агент может
сначала преобразовать задачу на естественном языке в выражение Wolfram
Language или Python, затем вызвать вычислительную систему и использовать
полученный результат в дальнейших рассуждениях.

Для разных задач нужны разные инструменты. SymPy подходит для символьных
выражений, уравнений и математического анализа; NumPy предоставляет
операции с массивами и матрицами; SciPy --- алгоритмы численной
оптимизации и другие методы.

\begin{quote}
\textbf{Эксперимент 5-3 ★★: решение математических задач с помощью кода}

В этом эксперименте одна и та же модель Doubao решает 30 задач AIME
2024, чтобы сравнить два подхода: непосредственное рассуждение и
использование инструментов Python. В среде инструментов доступны SymPy,
NumPy и SciPy; модель может сгенерировать программу, прочитать результат
её выполнения и затем дать окончательный ответ.

В этом запуске при непосредственном рассуждении были правильно решены 11
задач, а с помощью кода --- 16.
\end{quote}

Логические задачи можно решать аналогичным образом. Например, в задачах
о рыцарях и лжецах каждому жителю сопоставляется булева переменная:
значение «истина» означает рыцаря, а «ложь» --- лжеца. Затем каждое
высказывание преобразуется в булево выражение. Для говорящего X задаётся
ограничение: «X --- рыцарь» эквивалентно утверждению «высказывание X
истинно». Например, если A говорит: «B --- рыцарь», получается
\texttt{A\ ==\ B}; если A говорит: «B --- лжец», получается
\texttt{A\ ==\ (not\ B)}.

Решатель ищет значения переменных, удовлетворяющие всем ограничениям,
после чего модель интерпретирует их как роли персонажей. Если
высказывания противоречат друг другу, удовлетворяющего условиям набора
значений не существует; возможна и ситуация, когда заданные высказывания
допускают сразу несколько различных комбинаций ролей.

\begin{quote}
\textbf{Эксперимент 5-4 ★★: логическое рассуждение с помощью решения
ограничений}

В этом эксперименте из набора данных K\&K с разбиением по числу
персонажей выбраны 84 задачи, чтобы сравнить непосредственное
рассуждение одной и той же модели Doubao с вызовом
\texttt{python-constraint}. При использовании кода необходимо определить
переменные персонажей, преобразовать высказывания и вызвать решатель для
поиска комбинации ролей, удовлетворяющей всем условиям.

В этом запуске при непосредственном рассуждении были правильно решены 63
задачи, а с помощью кода --- 33. Во всех 84 траекториях с кодом
встречалась инструкция импорта библиотеки для решения ограничений. В
одной из задач модель даже исправила программу, но так и не предоставила
ответ.

Рассмотрим задачу с двумя персонажами, Джеймсом и Софией. Джеймс
говорит, что София --- лжец; София утверждает: «Джеймс --- рыцарь тогда
и только тогда, когда Джеймс --- лжец». Последнее высказывание ложно
независимо от роли Джеймса, следовательно, София --- лжец, а
высказывание Джеймса истинно. В первой программе, отправленной моделью,
бикондиционал был записан в виде \texttt{james\ ==\ not\ james}, который
Python не может разобрать. К восьмому вызову инструмента в программе уже
использовалась запись \texttt{james\ ==\ (not\ james)}, однако итоговый
ответ всё равно остался пустым, поэтому задача вновь не была засчитана.
\end{quote}

\textbf{Необходимая сложность Harness зависит от границ возможностей
модели в конкретной задаче.} Если модель уже стабильно справляется с
рассуждением, дополнительный этап решения может лишь увеличить накладные
расходы; те же части, в которых модель часто ошибается при вычислениях,
стоит передать детерминированным средствам. Однако использование
инструментов также требует, чтобы модель правильно выполнила
формализацию: ошибка при преобразовании ограничений превратит исходную
задачу в другую. Даже после исправления ограничений помощник может
исчерпать число шагов и не успеть предоставить ответ.

\hypertarget{ux43aux43eux434-ux43aux430ux43a-ux441ux440ux435ux434ux441ux442ux432ux43e-ux43eux431ux435ux441ux43fux435ux447ux435ux43dux438ux44f-ux431ux438ux437ux43dux435ux441-ux43fux440ux430ux432ux438ux43b}{%
\subsection{Код как средство обеспечения
бизнес-правил}\label{ux43aux43eux434-ux43aux430ux43a-ux441ux440ux435ux434ux441ux442ux432ux43e-ux43eux431ux435ux441ux43fux435ux447ux435ux43dux438ux44f-ux431ux438ux437ux43dux435ux441-ux43fux440ux430ux432ux438ux43b}}

Бизнес-правила можно закрепить в процессе выполнения с помощью кода.
Например, для правила «возврат возможен в течение 7 дней после покупки»
необходимо дополнительно определить, идёт ли речь о календарных или
рабочих днях, с какого момента начинается отсчёт --- оформления заказа,
оплаты или отправки --- и включается ли граничный момент. Только после
того, как бизнес-подразделение уточнит смысл этих условий, их можно
выразить в виде проверок и тестовых сценариев, обеспечив правилам
стабильную основу для выполнения.

Например, ответственный за правило подтвердил формулировку: «с момента
оплаты прошло не более 168 часов». Тогда можно составить три соседних
сценария: 167 часов 59 минут после оплаты, ровно 168 часов и 168 часов 1
минута. Из серверного времени подачи заявки вычитается время оплаты,
причём обе временные метки предварительно приводятся к единой временной
базе. Согласно этому правилу первые два сценария проходят проверку, а
третий отклоняется. Если ответственный изменит правило и потребует
считать календарные дни по месту оплаты, необходимо заново проверить как
существующие сценарии, так и код.

Описание на естественном языке помогает модели понять политику,
объяснить пользователю причину и в пределах полномочий найти другие
варианты, например проверить возможность изменения бронирования.
Программная проверка перед критически важными операциями контролирует
конкретные условия и хорошо подходит для сочетаний множества условий,
расчётов времени и сверки данных из разных источников. Модель может
помочь сгенерировать код проверки, а ответственный за правило
подтверждает его смысл посредством ревью и граничных сценариев. При
изменении правила следует одновременно обновить описание, реализацию и
тесты.

\textbf{Главная ценность программных правил --- предотвращение ошибочных
операций, способных причинить реальный ущерб.} Отмена заказа, перевод
средств и удаление данных изменяют состояние бизнеса, а исправление
последствий может стоить дорого. Автор рекомендует помещать проверку
фактических данных непосредственно в исполняющий интерфейс, чтобы каждая
отправка проходила через один и тот же набор правил и модель не могла
пропустить вызов отдельного инструмента проверки. Отдельный интерфейс
запросов по-прежнему можно использовать для предварительного объяснения
политики, но при выполнении следует заново считать необходимые факты,
проверить, не изменилось ли состояние, и управлять последующими
операциями с помощью описанных в \crossreflink{chap:4}{главе 4}
транзакций, идемпотентности и механизмов компенсации.

\textbf{Параметры направляют проверку, а решение принимается по
серверным фактам.}

τ-bench исследует вызов инструментов и соблюдение политик на примере
авиации, электронной коммерции и других сценариев\footnote{Sierra
  Research, τ-bench: A Benchmark for Tool-Agent-User Interaction in
  Real-World Domains --- политики, инструменты и интерактивные задачи
  для сценариев авиации и электронной коммерции.}. Ниже в качестве
учебного сценария используется бронирование авиабилетов и задаются пять
правил отмены: сначала исключаются заказы с уже использованными
сегментами маршрута, затем проверяются время бронирования, отмена рейса
авиакомпанией, бизнес-класс и наличие страховки.
\texttt{cancellation\_reason} фиксирует указанную пользователем причину
заявки, а \texttt{expected\_*} --- факты, которые модель ожидает
увидеть; фактическое решение принимается на основе данных из бизнес-базы
и серверных часов.

Этот код демонстрирует ветви политики. Входной слой отвечает за
подтверждение личности инициатора запроса, принадлежности заказа и
полномочий на операцию. Исполняющий слой в управляемой транзакции или
эквивалентном механизме согласованности считывает состояние, проверяет
условия, выполняет отмену и сохраняет запись аудита.

\begin{Shaded}
\begin{Highlighting}[]
\KeywordTok{def}\NormalTok{ cancel\_reservation(}
\NormalTok{    reservation\_id: }\BuiltInTok{str}\NormalTok{,}
\NormalTok{    cancellation\_reason: }\BuiltInTok{str}\NormalTok{,        }\CommentTok{\# "change\_of\_plan", "airline\_cancelled", "other"}
\NormalTok{    expected\_cabin\_class: }\BuiltInTok{str} \OperatorTok{|} \VariableTok{None} \OperatorTok{=} \VariableTok{None}\NormalTok{,  }\CommentTok{\# Необязательно: класс обслуживания, ожидаемый моделью}
\NormalTok{    expected\_has\_insurance: }\BuiltInTok{bool} \OperatorTok{|} \VariableTok{None} \OperatorTok{=} \VariableTok{None}  \CommentTok{\# Необязательно: статус страховки, ожидаемый моделью}
\NormalTok{) }\OperatorTok{{-}\textgreater{}} \BuiltInTok{dict}\NormalTok{:}
    \CommentTok{"""}
\CommentTok{    Отменяет бронирование авиабилета.}

\CommentTok{    Политика отмены (принудительно применяется сервером по фактам из базы данных):}
\CommentTok{    {-} Правило 1: заказ нельзя отменить, если использован хотя бы один сегмент маршрута}
\CommentTok{    {-} Правило 2: после прохождения правила 1 заказ можно отменить в течение 24 часов после бронирования}
\CommentTok{    {-} Правило 3: после прохождения правила 1 можно отменить билет на рейс, отменённый авиакомпанией}
\CommentTok{    {-} Правило 4: после прохождения правила 1 можно отменить билет бизнес{-}класса}
\CommentTok{    {-} Правило 5: билеты базового экономкласса и экономкласса можно отменить только при наличии туристической страховки}

\CommentTok{    Перед вызовом сначала запросите сведения о заказе и последовательно сверьте их}
\CommentTok{    с приведённой политикой. Параметры expected\_* описывают основания вашего решения}
\CommentTok{    и используются сервером только для сверки и аудита, не влияя на решение по политике.}
\CommentTok{    """}
    \CommentTok{\# Считать из бизнес{-}базы факты, необходимые для проверки политики}
\NormalTok{    r }\OperatorTok{=}\NormalTok{ db.get\_reservation(reservation\_id)}
\NormalTok{    now }\OperatorTok{=}\NormalTok{ server\_clock.now()  }\CommentTok{\# Использовать серверные часы}

    \CommentTok{\# Зафиксировать расхождения ожидаемых значений с текущими фактами для диагностики и аудита}
    \ControlFlowTok{if}\NormalTok{ expected\_cabin\_class }\KeywordTok{is} \KeywordTok{not} \VariableTok{None} \KeywordTok{and}\NormalTok{ expected\_cabin\_class }\OperatorTok{!=}\NormalTok{ r.cabin\_class:}
\NormalTok{        log\_mismatch(reservation\_id, }\StringTok{"cabin\_class"}\NormalTok{, expected\_cabin\_class, r.cabin\_class)}
    \ControlFlowTok{if}\NormalTok{ expected\_has\_insurance }\KeywordTok{is} \KeywordTok{not} \VariableTok{None} \KeywordTok{and}\NormalTok{ expected\_has\_insurance }\OperatorTok{!=}\NormalTok{ r.has\_insurance:}
\NormalTok{        log\_mismatch(reservation\_id, }\StringTok{"has\_insurance"}\NormalTok{, expected\_has\_insurance, r.has\_insurance)}

    \ControlFlowTok{if}\NormalTok{ r.any\_segment\_used:}
        \ControlFlowTok{return}\NormalTok{ \{}\StringTok{"success"}\NormalTok{: }\VariableTok{False}\NormalTok{, }\StringTok{"reason"}\NormalTok{: }\StringTok{"Cannot cancel with used segments"}\NormalTok{\}}

\NormalTok{    hours\_since\_booking }\OperatorTok{=}\NormalTok{ (now }\OperatorTok{{-}}\NormalTok{ r.booking\_time).total\_seconds() }\OperatorTok{/} \DecValTok{3600}
    \ControlFlowTok{if}\NormalTok{ hours\_since\_booking }\OperatorTok{\textless{}} \DecValTok{0}\NormalTok{:}
        \ControlFlowTok{return}\NormalTok{ \{}\StringTok{"success"}\NormalTok{: }\VariableTok{False}\NormalTok{, }\StringTok{"reason"}\NormalTok{: }\StringTok{"Booking time is in the future"}\NormalTok{\}}
    \ControlFlowTok{if}\NormalTok{ hours\_since\_booking }\OperatorTok{\textless{}=} \DecValTok{24}\NormalTok{:}
\NormalTok{        execute\_cancellation(reservation\_id)}
        \ControlFlowTok{return}\NormalTok{ \{}\StringTok{"success"}\NormalTok{: }\VariableTok{True}\NormalTok{, }\StringTok{"reason"}\NormalTok{: }\StringTok{"Cancelled within 24{-}hour window"}\NormalTok{\}}

    \ControlFlowTok{if}\NormalTok{ r.flight\_status }\OperatorTok{==} \StringTok{"cancelled\_by\_airline"}\NormalTok{:}
\NormalTok{        execute\_cancellation(reservation\_id)}
        \ControlFlowTok{return}\NormalTok{ \{}\StringTok{"success"}\NormalTok{: }\VariableTok{True}\NormalTok{, }\StringTok{"reason"}\NormalTok{: }\StringTok{"Airline cancelled flight"}\NormalTok{\}}

    \ControlFlowTok{if}\NormalTok{ r.cabin\_class }\OperatorTok{==} \StringTok{"business"}\NormalTok{:}
\NormalTok{        execute\_cancellation(reservation\_id)}
        \ControlFlowTok{return}\NormalTok{ \{}\StringTok{"success"}\NormalTok{: }\VariableTok{True}\NormalTok{, }\StringTok{"reason"}\NormalTok{: }\StringTok{"Business class cancellation"}\NormalTok{\}}

    \ControlFlowTok{if}\NormalTok{ r.cabin\_class }\KeywordTok{in}\NormalTok{ [}\StringTok{"basic\_economy"}\NormalTok{, }\StringTok{"economy"}\NormalTok{]:}
        \ControlFlowTok{if}\NormalTok{ r.has\_insurance:}
\NormalTok{            execute\_cancellation(reservation\_id)}
            \ControlFlowTok{return}\NormalTok{ \{}\StringTok{"success"}\NormalTok{: }\VariableTok{True}\NormalTok{, }\StringTok{"reason"}\NormalTok{: }\SpecialStringTok{f"}\SpecialCharTok{\{}\NormalTok{r}\SpecialCharTok{.}\NormalTok{cabin\_class}\SpecialCharTok{\}}\SpecialStringTok{ with insurance"}\NormalTok{\}}
        \ControlFlowTok{return}\NormalTok{ \{}\StringTok{"success"}\NormalTok{: }\VariableTok{False}\NormalTok{, }\StringTok{"reason"}\NormalTok{: }\SpecialStringTok{f"}\SpecialCharTok{\{}\NormalTok{r}\SpecialCharTok{.}\NormalTok{cabin\_class}\SpecialCharTok{\}}\SpecialStringTok{ requires insurance"}\NormalTok{\}}

    \ControlFlowTok{return}\NormalTok{ \{}\StringTok{"success"}\NormalTok{: }\VariableTok{False}\NormalTok{, }\StringTok{"reason"}\NormalTok{: }\StringTok{"Does not meet cancellation policy"}\NormalTok{\}}
\end{Highlighting}
\end{Shaded}

\textbf{Параметры помогают модели последовательно сверить каждый пункт.}
В описании инструмента перечисляется политика отмены и содержится
требование сначала запросить заказ. Модель может передать класс
обслуживания и статус страховки через \texttt{expected\_cabin\_class} и
\texttt{expected\_has\_insurance}, чтобы основания решения попали в
журнал вызовов. Необязательные параметры предоставляют точку для сверки,
однако факт выполнения запроса и самой сверки всё равно необходимо
проверять по траектории. Если процесс требует обязательного
предварительного запроса, Harness может проверить наличие записи о
запросе и срок её действия.

Например, если с момента бронирования прошло более 24 часов, рейс
выполняется по расписанию, билет относится к экономклассу и страховка не
приобретена, модель может на основании этих данных объяснить, почему
билет нельзя отменить по политике из примера, и найти разрешённые
варианты, такие как изменение бронирования.

\textbf{Сервер применяет правила по текущим фактам.} Класс обслуживания,
наличие страховки, время бронирования, использование сегментов маршрута
и статус рейса считываются из бизнес-источников данных, а временное окно
рассчитывается по серверным часам. Расхождения \texttt{expected\_*}
используются для диагностики, а \texttt{cancellation\_reason} --- для
фиксации намерения заявителя. Такие факты, как отмена рейса
авиакомпанией, проверяются сервером: если пользователь укажет в заявке
«отмена авиакомпанией», это не сделает изначально невозвратный билет
возвратным.

\emph{Уточнение русского издания: В демонстрационном коде проверяется
наличие страховки, но не причина отмены. В правилах τ-bench страховка
даёт право на отмену, только если покрывает её причину.
\href{https://github.com/sierra-research/tau-bench/blob/main/tau_bench/envs/airline/wiki.md}{Источник}.}

Независимость источников данных необходимо обеспечить и на уровне
контроля доступа: достоверность проверки зависит от того, кто вправе
изменять факты о заказе, когда они обновляются и как сохраняется
история. Параллельные операции могут изменить состояние после запроса,
поэтому при отправке возврата сервер заново считывает состояние заказа,
подтверждает, что билет всё ещё можно вернуть, и обеспечивает одинаковый
результат обработки повторных запросов.

\begin{quote}
\textbf{Эксперимент 5-5 ★★: малая модель и программные бизнес-правила}

В этом эксперименте используется Qwen3-4B, а две реализации сравниваются
на 60 фиксированных сценариях возврата авиабилетов. В обеих группах
модель получает политику на естественном языке. В контрольной группе
инструмент отмены после вызова сразу возвращает деньги, а в группе с
программными правилами перед выполнением считываются имитационная база
данных и серверные часы, после чего проверяется право на возврат.

В эксперименте используется упрощённая политика: билеты гибкого
экономкласса и бизнес-класса подлежат возврату; билеты базового
экономкласса можно вернуть в течение 24 часов после бронирования, а
также при отмене рейса авиакомпанией или значительной задержке.
Параметры сверки \texttt{expected\_refundable} и
\texttt{expected\_reason} фиксируют соответственно оценку моделью права
на возврат и его основания.

При оценке сначала проверяется, были ли возвращены деньги, а затем ---
соответствовал ли возврат политике для данного билета. По этому критерию
контрольная группа правильно обработала 57 сценариев, а группа с
программными правилами --- 55. Ни одна из групп не вернула деньги
пассажирам, не имевшим на это права. Все потерянные баллы связаны с
заказами, по которым возврат полагался, но не был выполнен: 3 сценария в
контрольной группе и 5 в группе с программными правилами. Пассажирам
также важно, объяснила ли служба поддержки причину и предложила ли
дальнейшие варианты после отказа. В эксперименте отдельно проверялись
объяснение политики и альтернативные варианты в итоговом ответе, однако
эти проверки не учитывались в приведённой выше оценке возвратов.
\end{quote}

\hypertarget{ux43fux440ux43eux433ux440ux430ux43cux43cux43dux430ux44f-ux433ux435ux43dux435ux440ux430ux446ux438ux44f-ux43cux443ux43bux44cux442ux438ux43cux435ux434ux438ux430}{%
\subsection{Программная генерация
мультимедиа}\label{ux43fux440ux43eux433ux440ux430ux43cux43cux43dux430ux44f-ux433ux435ux43dux435ux440ux430ux446ux438ux44f-ux43cux443ux43bux44cux442ux438ux43cux435ux434ux438ux430}}

Подготовка презентаций, технических отчётов и интерактивных страниц
включает две задачи: организацию содержимого и визуальное представление.
AI-агент может сначала описать структуру в виде кода или языка разметки,
а затем вызвать инструмент рендеринга для создания доступного для
просмотра результата. Например, HTML задаёт структуру страницы, CSS
управляет компоновкой и стилями, а JavaScript реализует интерактивность.
Положение и стили элементов становятся изменяемыми параметрами, что
упрощает массовую корректировку, повторную генерацию и сравнение версий.

Программная генерация подходит для структурированного содержимого,
которое требуется неоднократно изменять. GUI предоставляет возможность
непосредственно просматривать результат и выполнять локальные операции.
Оба подхода можно использовать совместно.

\textbf{Агент для создания PPT.}

Для научного доклада необходимо разложить проблему, метод, эксперименты
и выводы статьи на последовательное изложение по слайдам. Агент может
сначала определить вопрос, на который должен отвечать каждый слайд,
затем выбрать основные тезисы и иллюстрации и, наконец, сгенерировать
исходный код презентации. Slidev использует Markdown как основной формат
содержимого и поддерживает HTML и компоненты Vue, поэтому хорошо
подходит для организации этого процесса в исполняемый конвейер
генерации\footnote{Slidev, \emph{Getting Started}: создание, показ и
  экспорт презентаций на основе Markdown и веб-технологий.}.

Если организаторы предоставили доклад прошлых лет в качестве образца,
можно напрямую переиспользовать его структуру слайдов. PPTAgent сначала
анализирует назначение каждого слайда в референсной презентации, затем
подбирает подходящие слайды для нового плана и изменяет расположенные на
них объекты с помощью кода. Например, для представления результатов
нового эксперимента можно выбрать слайд, изначально предназначенный для
параллельного показа двух групп результатов, и заменить на нём
иллюстрации и пояснения. Такое переиспользование требует дополнительной
проверки: если обе иллюстрации уже заменены результатами текущего
эксперимента, не остались ли в соседних абзацах пояснения к
прошлогоднему? В версии из статьи текстовые поля изменяются по номерам
элементов и абзацев. Если заменить только иллюстрации, прежние пояснения
останутся на слайде.\footnote{Zheng и др., статья о PPTAgent, EMNLP
  2025. В сопутствующем коде статьи интерфейс для редактирования
  отдельных абзацев реализован в \texttt{pptagent/apis.py}.}

Даже если текст и изображения заменены правильно, вёрстка всё равно
может нарушиться. Например, после добавления к иллюстрации двух строк
пояснений компиляция может пройти успешно, но текст перекроет оси
координат. Как показано на \crossreflink{fig:5-5}{рисунке 5-5}, только
увидев снимок отрендеренного слайда, проверяющий сможет передать
подобное замечание автору, отвечающему за вёрстку.

\begin{figure}
\hypertarget{fig:5-5}{%
\centering
\includegraphics{images/revised/fig5-5.pdf}
\caption{Рисунок 5-5. Механизм «автор --- проверяющий» при создании
PPT}\label{fig:5-5}
}
\end{figure}

\begin{itemize}
\tightlist
\item
  \textbf{Автор (Proposer)} читает статью, планирует структуру
  изложения, пишет исходный код Slidev и исправляет слайды с учётом
  обратной связи.
\item
  \textbf{Проверяющий (Reviewer)} рендерит слайды с помощью
  инструментов, использует визуальную модель для проверки плотности
  содержимого, читаемости, компоновки и соответствия текста
  иллюстрациям, а затем возвращает структурированную обратную связь с
  номером слайда, типом проблемы, серьёзностью и рекомендацией по
  исправлению.
\end{itemize}

Обратная связь должна позволять сразу приступить к исправлению.
Например: «На слайде 3 слишком много тезисов: разделите его на обзор
метода и этапы вычислений»; «На слайде 7 слишком мелкий шрифт кода:
сначала увеличьте его до 14 pt, затем снова проверьте читаемость при
проецировании». В первую очередь исправляются фактические ошибки и
неверная интерпретация графиков, а затем --- компоновка и визуальные
детали.

При локальном изменении необходимо также указать конкретный объект.
Функция Doubao Work в Feishu Slides позволяет пользователю выбрать
элемент слайда и сформулировать требование к редактированию через
плавающую панель инструментов.\footnote{Справочный центр Feishu,
  использование Doubao Work в презентациях.} Предположим, на одном
слайде расположены два текстовых поля: в одном находится исходная цитата
из интервью, а в другом --- резюме докладчика. Пользователь выбирает
резюме и просит: «Сократи до двух строк». Выбранное текстовое поле
ограничивает область изменения. Резюме можно переформулировать, но
соседняя цитата должна сохранить исходные слова респондента.

После исправлений автор повторно рендерит слайды, а проверяющий изучает
новую версию и подтверждает, устранены ли проблемы предыдущего цикла.
Для цикла задаются условия выхода двух типов: достижение явных
требований к содержимому и оформлению либо исчерпание бюджета по числу
итераций, времени или стоимости --- например, не более пяти циклов
исправлений. При передаче результата сохраняются текущая версия и список
нерешённых проблем, чтобы облегчить последующую проверку.

Здесь используется тот же шаблон «автор --- проверяющий», что и при
утверждении операций в \crossreflink{chap:4}{главе 4}. При утверждении
операций в основном определяется, допустимо ли выполнить конкретное
действие; при проверке содержимого результат улучшается в ходе
нескольких циклов обратной связи. В обоих случаях необходимо совместно
использовать требования задачи, приводить основания, позволяющие
локализовать проблему, и поручать Harness определение дальнейшего
процесса. Использование моделей из разных семейств позволяет получить
ещё одну точку зрения при проверке.

\textbf{Автор считает важным преимуществом разделения обязанностей между
двумя агентами более эффективное управление контекстом.} Каждый из двух
участников может хранить только необходимую ему информацию: проверяющий
в каждом цикле получает актуальные снимки слайдов, критерии проверки и
необходимый список проблем; автор сохраняет исходный код, цели
исправлений и структурированную текстовую обратную связь. Изображения
слайдов остаются в контексте проверяющего, а основная траектория
продолжает работу с краткой текстовой обратной связью. Каждый цикл
проверки сосредоточен на последней версии слайдов, что одновременно
уменьшает помехи от старых снимков и сдерживает рост контекста.

\begin{quote}
\textbf{Эксперимент 5-6 ★★: автоматическое создание PPT по научной
статье}

В этом эксперименте по PDF-файлу статьи «Attention Is All You Need»
создаётся презентация Slidev из 20 слайдов. Агент извлекает структуру и
основные тезисы статьи, использует три исходные иллюстрации из неё с
указанием источников и последовательно организует изложение по слайдам.
После рендеринга проверяются выход текста за границы, читаемость
графиков и соответствие текста иллюстрациям, а затем вносятся
исправления с учётом обратной связи.

Сравниваются самостоятельная проверка одним агентом и разделение
обязанностей между автором и проверяющим. В этом запуске обе итоговые
презентации получили по 95 баллов в рамках одной независимой процедуры
оценки. Пиковый объём входных данных при разделении обязанностей
составил около 24 тысяч токенов, а в варианте с одним агентом и
сохранением прежних снимков --- около 93 тысяч токенов.
\end{quote}

\begin{quote}
\textbf{Эксперимент 5-7 ★★: автоматическое создание видеолекции по
научной статье}

На основе презентации агент может написать для каждого слайда текст
устного объяснения, описывающий представленные на иллюстрации объекты,
связи и выводы, а затем синтезировать речь с помощью TTS. Текст
выступления передаётся визуальному проверяющему вместе со слайдами,
чтобы убедиться, что упомянутые элементы присутствуют в кадре, а
описание соответствует графикам.

При создании видео в эксперименте использовался процесс генерации из
предыдущего раздела и более ранняя версия презентации из 22 слайдов
(v2). После последовательного отбора 12 слайдов последний кадр видео
соответствовал слайду 21 исходной презентации. Для сравнения оценок и
объёма ресурсов в предыдущем разделе использовалась более поздняя версия
из 20 слайдов (v9). После выбора слайдов Kimi сгенерировала текст
выступления, визуальная модель Qwen проверила соответствие текста
изображениям, а Fish Audio синтезировала речь для каждого слайда. FFmpeg
установил длительность показа слайдов по фактической продолжительности
каждого аудиофрагмента и собрал видеолекцию длительностью около 8,6
минуты.

Как показано на \crossreflink{fig:5-6}{рисунке 5-6}, слайды, текст
выступления и речь объединяются в видео только на последнем этапе.
Некоторые проблемы можно услышать лишь тогда: не обрывается ли фраза при
смене слайда, не усиливается ли к концу рассинхронизация звука и
изображения.

\begin{figure}
\hypertarget{fig:5-6}{%
\centering
\includegraphics{images/revised/fig5-6.pdf}
\caption{Рисунок 5-6. Сквозной конвейер от научной статьи до
видеолекции}\label{fig:5-6}
}
\end{figure}
\end{quote}

\textbf{Агент для монтажа видео.}

\textbf{Операции видеомонтажа, которые можно выразить скриптом, следует
преимущественно выполнять через программный интерфейс.} Временная шкала,
слои, субтитры и звуковые дорожки распределены по сложным панелям GUI, а
повторный поиск и перетаскивание элементов увеличивают задержку и
вероятность ошибок. Программный интерфейс позволяет напрямую описывать
операции временными интервалами, дорожками и параметрами эффектов.
Python API Blender и FFmpeg позволяют оформить часть задач монтажа в
виде повторяемых скриптов.

Агент может сначала интерпретировать требование вроде «вырежи фрагмент с
сёрфингом», затем найти целевой интервал в исходном материале и
сгенерировать код для обрезки, компоновки и обработки звуковых дорожек.
Автор отвечает за план монтажа и скрипты, а проверяющий изучает
результат рендеринга. Если обнаружены пропущенные или посторонние
фрагменты либо перекрытие субтитрами, он указывает точную временную
отметку и проблему, после чего автор вносит исправления. Ключевые кадры
подходят для проверки сцен и компоновки, но непрерывность движения и
синхронизацию звука с изображением необходимо проверять также по
соседним кадрам или непрерывным фрагментам.

\begin{quote}
\textbf{Эксперимент 5-8 ★★: интеллектуальный видеомонтаж на основе API}

В этом эксперименте поиск целевого фрагмента, генерация скрипта и
проверка эффектов объединены в единый процесс. Сначала ссылка на видео,
целевой вопрос и интервал извлечения кадров передаются дочернему агенту
для анализа видео. Он обрабатывает большое количество кадров и
возвращает основному агенту временной интервал вместе с обоснованием.

Поиск выполняется в два этапа. На первом этапе извлекается один кадр
каждые 10 секунд, чтобы определить приблизительные границы целевого
фрагмента, например: «сёрфинг находится между 40-й и 110-й секундами».
На втором этапе вокруг найденного интервала и соседних границ
извлекается по одному кадру в секунду, чтобы точнее определить начало и
конец. Сцены короче интервала грубой выборки могут оказаться между двумя
кадрами. В этом случае можно повысить частоту выборки с учётом
продолжительности искомого фрагмента или просмотреть изображения между
двумя извлечёнными кадрами.

После определения интервала автор генерирует скрипт Blender Python,
который выполняет обрезку и необходимые операции с субтитрами, скоростью
воспроизведения и другими параметрами. Проверяющий просматривает
предварительный результат и указывает границы и эффекты, требующие
корректировки. Полная версия видео экспортируется только после успешного
прохождения проверки.

В действительности вырезался фрагмент, в котором большой белый кролик в
одиночестве идёт по траве. При просмотре одного кадра каждые 10 секунд
модель определила интервал от 10-й до 20-й секунды. Если начать обрезку
непосредственно с 10-й секунды, начало сцены будет потеряно. Поэтому
программа расширила диапазон на 5 секунд в обе стороны, просмотрела по
одному кадру в секунду между 5-й и 25-й секундами и сузила интервал до
9--12 секунд. Затем проверяющий обнаружил, что во фрагмент попали
соседние кадры, и исправленный интервал составил с 9,25-й по 11,00-ю
секунду. По сравнению с ручной разметкой интервала с 9,08-й по 11,20-ю
секунду начало запаздывает на 0,17 секунды, а конец наступает на 0,20
секунды раньше. Обе границы соответствуют требованию о погрешности не
более 3 секунд.

Временные коды позволяют проверить лишь место обрезки. Но сможет ли
проверяющий обнаружить ошибку, если наложить те же субтитры и замедление
на намеренно неверно выбранный фрагмент? В эксперименте ему сначала
предложили проверить контрольный фрагмент с 0-й по 3-ю секунду, а затем
--- исправленный результат рендеринга. Проверялись большой белый кролик
в кадре, субтитры в нижней части и длительность после изменения
скорости.
\end{quote}

\textbf{3D и промышленные детали: выбор подходящего способа генерации.}

\textbf{Для промышленных деталей, у которых размеры и сборочные
отношения являются жёсткими ограничениями, предпочтительна генерация
параметрическим кодом; для объектов, где важнее естественность формы и
внешний вид, следует в первую очередь рассматривать генеративные
модели.} Выбор зависит от того, можно ли точно описать результат
небольшим числом параметров и какие отношения должны сохраняться после
изменений. CadQuery, OpenSCAD и скрипты Blender предоставляют
редактируемый процесс построения, а такие модели, как Hunyuan 3D,
используют усвоенные априорные представления о формах для создания
сложного внешнего вида.

\textbf{Пригодность описания для параметризации.} Основную геометрию
упрощённого фланца можно выразить через внешний диаметр, толщину,
диаметр окружности расположения отверстий, диаметр и число отверстий.
Код связывает параметры с этапами построения и поэтому подходит для
многократной корректировки. Растение в горшке, камень Тайху или
человеческое лицо содержат множество нерегулярных форм. Генеративная
модель может помочь создать их внешний вид, а программные методы ---
сформировать детали с помощью правил роста, случайных процессов и
комбинации ассетов.

\textbf{Проверка точности.} Для требования вроде «диаметр отверстия 5
mm, допуск ±0,05 mm» необходимо сначала явно определить единицы
измерения, систему координат и метод измерения, а затем сопоставить
размеры модели со спецификацией. В параметрической модели можно
непосредственно проверять геометрические объекты. Сетку можно измерять
по сечениям, отверстиям, габаритным размерам и другим характеристикам с
учётом погрешности дискретизации.

\textbf{Сохранение и изменение.} B-rep --- граничное представление ---
описывает границу твёрдого тела гранями, рёбрами и вершинами, а
треугольная сетка аппроксимирует поверхность гранями. CadQuery может
экспортировать STEP для обмена данными CAD и сетки STL, а редактируемые
параметры и логика построения сохраняются в исходном коде
CadQuery\footnote{CadQuery, \emph{Importing and Exporting Files}:
  описание форматов STEP, STL и других, а также указание на то, что
  полная параметрическая информация сохраняется в исходном коде CadQuery
  на Python.}. Для перехода к обработке на станках с ЧПУ дополнительно
требуются планирование технологического процесса, построение траекторий
инструмента и адаптация к станку.

Например, чтобы вместо M5 использовать крепёжные отверстия под M6, можно
изменить параметр диаметра отверстий, повторно сгенерировать модель и
проверить, что внешний диаметр, толщина и расположение отверстий
остались неизменными. Способ изменения результата генеративной модели
зависит от возможностей интерфейса: модель можно сгенерировать заново
либо применить локальное редактирование и последующую геометрическую
обработку.

При сочетании двух подходов также следует явно разделять обязанности:
размеры и сборочные отношения поручить параметрическому коду, а текстуры
и свободную форму --- генеративной модели. \textbf{Если ключевые размеры
требуется неоднократно и точно изменять, редактируемый процесс
построения ценнее однократно сгенерированного внешнего вида.}

\begin{quote}
\textbf{Эксперимент 5-9 ★★: два способа создания одной детали --- код и
генеративная модель}

В этом эксперименте необходимо создать фланец внешним диаметром 80 mm и
толщиной 10 mm с четырьмя равномерно расположенными крепёжными
отверстиями на окружности диаметром 60 mm. Чтобы однозначно определить
требования к сборке, первоначальная спецификация задаёт сквозные
отверстия диаметром 5,5 mm под M5, а запрос на изменение требует
заменить их сквозными отверстиями диаметром 6,5 mm под M6.

В программном варианте модель генерирует программу CadQuery, после
выполнения которой экспортируются STEP и STL. В варианте с генеративной
моделью сначала создаётся референсное изображение изделия, которое затем
передаётся Hunyuan 3D для генерации сетки. Проверяются внешний диаметр,
толщина, число и диаметр отверстий, их расположение и монтажная
поверхность; также фиксируются единицы измерения и направление оси.

Как проверить, что модель изменила только указанный размер? При
сопоставлении экспортированных STL до и после изменения средний диаметр
четырёх крепёжных отверстий увеличился с 5,499 mm до 6,499 mm. Внешний
диаметр обеих деталей составляет 80 mm, толщина --- 10 mm, а диаметр
окружности расположения отверстий при округлении до трёх знаков равен
60,000 mm. Исходная сетка, созданная генеративной моделью, не имеет явно
заданного физического масштаба. Программа измерения построила сечение
посередине по оси Z, но не обнаружила внутренних контуров отверстий.
Чтобы выяснить, проходят ли четыре отверстия через всё тело диска,
необходимо определить монтажную поверхность и построить сечения на
разной глубине по её нормали. После повторной генерации изменились как
общая форма, так и ориентация.

Наконец, рассмотрим другую задачу: как принимать изображение растения в
горшке, которое должно выглядеть естественно? В эксперименте одно
изображение растения было нарисовано с помощью matplotlib, а другое
сгенерировано Gemini. Затем одна и та же визуальная модель оценила их
естественность по десятибалльной шкале. Изображения получили
соответственно 3 и 8 баллов. Эта оценка показывает, какое из них
модель-проверяющий считает более похожим на настоящее растение. Если
вернуться к фланцу, даже когда на изображении изделия видны четыре
отверстия, для проверки сборки всё равно необходимо измерить их
положение и диаметр на трёхмерной детали.
\end{quote}

\hypertarget{ux43aux43eux434-ux43aux430ux43a-ux441ux438ux441ux442ux435ux43cux43dux44bux439-ux430ux434ux430ux43fux442ux435ux440}{%
\subsection{Код как системный
адаптер}\label{ux43aux43eux434-ux43aux430ux43a-ux441ux438ux441ux442ux435ux43cux43dux44bux439-ux430ux434ux430ux43fux442ux435ux440}}

Интерфейсы, способы аутентификации и форматы данных внешних сервисов
различаются. AI-агент может изучить документацию и примеры ответов,
сгенерировать HTTP-клиент, настроить заголовки аутентификации, разобрать
возвращаемые значения и преобразовать поля вышестоящей системы в
структуру, необходимую нижестоящей. Код здесь выполняет роль системного
адаптера, объединяя существующие сервисы в новый процесс выполнения
задач.

Например, платёжный сервис может возвращать сумму в фэнях, а отчёт ---
отображать её в юанях. Успешный разбор JSON говорит лишь о совместимости
форматов. Транзакция на 25 юаней представлена в вышестоящей системе
целым числом 2500. Если напрямую поместить его в отчёт, 25 юаней
превратятся в 2500 юаней.

Для адаптера необходимо явно определить входные и выходные данные,
обработку ошибок и диапазон поддерживаемых версий. Примеры помогают
модели понять поля, а спецификация интерфейса и тесты проверяют
отсутствующие значения, пагинацию, ошибочные ответы и изменения версий.
Учётными данными управляет среда выполнения, а сгенерированный код
использует авторизованные возможности через контролируемые интерфейсы.
После успешной проверки адаптер можно сохранить и использовать повторно,
а при последующих изменениях формата --- снова запустить процесс
модификации.

Аналогичный подход применим и к системам, доступным только через
графический интерфейс. Сначала агент выполняет задачу с помощью Computer
Use, представленного в \crossreflink{chap:6}{главе 6}, а затем
закрепляет стабильную последовательность действий в виде инструмента
RPA. При изменении элементов интерфейса, состояния входа или
последовательности страниц инструмент должен обнаруживать отклонения и
заново находить нужные элементы. В \crossreflink{chap:9}{главе 9}
подробнее рассматривается извлечение, проверка и повторное использование
рабочих процессов из успешных траекторий.

\textbf{Разбор и визуализация журналов агента.}

Сложная задача может включать несколько запросов к модели, выполнений
инструментов и взаимодействий с дочерними агентами. Обзор журнала
показывает взаимосвязи шагов, их состояния и длительность, а в
развёрнутом представлении отображаются конкретные входные и выходные
данные, позволяя проследить, что произошло на определённом шаге. На
разбор и отображение влияют новые поля, вложенные структуры и версии
формата.

Адаптивный разбор обрабатывает такие изменения посредством цикла
«обнаружить изменение → сгенерировать парсер → проверить результат». При
сбое разбора система сохраняет исходный пример, позицию ошибки и
ожидаемые поля, а затем сопоставляет их со спецификацией интерфейса и
версией производителя данных, чтобы определить, было ли изменение
ожидаемым. Агенту поручают написать новый парсер только после
подтверждения обновления формата. Если ошибка возникла в программе,
записывающей журнал, исходный журнал сохраняется, а о сбое сообщается,
чтобы новый парсер не принял ошибочные данные за нормальный формат.
Новый код сначала запускается в ограниченной среде: проверяются поля и
их смысл, затем выполняется проверка на обычных примерах, примерах с
отсутствующими полями и регрессионных сценариях старого формата. После
успешной проверки новый парсер регистрируется с указанием версии и
применимого формата. Неудачные примеры продолжают храниться для
последующего отслеживания и повторных попыток.

Для разбора и отображения используются разные методы проверки. Программа
проверяет временные метки, состояния, числовые типы и связи между
вызовами. Рендеринг в браузере проверяет работоспособность страницы, а
визуальная модель помогает оценить понятность таблиц, временных шкал и
меток. При горячем обновлении сохраняются предыдущая версия и
возможность отката, чтобы поддержку нового формата можно было вводить в
работающую систему постепенно.

\begin{quote}
\textbf{Эксперимент 5-10 ★★★: адаптивная система разбора журналов}

В этом эксперименте сначала запускаются два процесса формирования
журналов, выводящие соответственно формат с разделителями-вертикальными
чертами и формат с вложенными скобками. Исходный парсер не способен
обработать эти примеры, после чего система запрашивает у модели две
функции разбора, а после успешного прохождения тестов загружает их без
перезапуска и повторно обрабатывает ранее не разобранные журналы.

Из первого формата извлекаются \texttt{timestamp}, \texttt{level},
\texttt{module}, \texttt{step} и \texttt{message}; второй также содержит
имя инструмента, длительность и состояние. В журнале запуска успешно
разобраны по три записи каждого типа, а вновь созданный движок способен
загрузить сохранённые парсеры и продолжить обработку смешанного потока
журналов.

Chromium отображает результаты разбора в виде таблицы, а визуальная
проверка оценивает названия столбцов, содержимое строк и
удобочитаемость.
\end{quote}

\textbf{Автоматический анализ журналов выполнения агента и диагностика
проблем.}

После извлечения структурированных траекторий агент может сопоставить
ожидаемый процесс с фактическими действиями на основе архитектурной
документации и PRD --- документа с требованиями к продукту. Например,
проверялась ли возможность возврата средств по заказу до самого
возврата? Переключилась ли программа на чтение из кэша после истечения
времени ожидания сервиса складского учёта?

Отчёт о проблеме может содержать приоритет, модуль, проявление,
доказательства и рекомендации по улучшению. На его основе агент также
может сгенерировать регрессионные тесты, зафиксировав входные данные и
условия приёмки, соответствующие проблеме. В контролируемой среде
тестовый фреймворк повторно воспроизводит запросы или считывает
сохранённые ответы инструментов, сопоставляя поведение дефектной и
исправленной версий. Для операций с внешней записью используется
изолированная среда или замещающие интерфейсы, чтобы последствия
воспроизведения оставались контролируемыми.

Утверждения должны соответствовать проверяемым свойствам. Например,
наличие в журнале проверки возможности возврата средств говорит лишь о
том, что программа её выполнила. Чтобы определить правильность возврата,
необходимо также проверить, что проверка предшествовала возврату, её
результат разрешал возврат и каков итоговый статус заказа.

После проверки диагностики агент может, получив соответствующее
разрешение, создать GitHub Issue через MCP и предложить ответственного с
учётом модуля. Если для того же дефекта уже существует рабочая задача,
новые доказательства следует добавить в неё.

\begin{quote}
\textbf{Эксперимент 5-11 ★★★: интеллектуальная система диагностики
производственных журналов}

В этом эксперименте локальный HTTP-сервис формирует две
бизнес-траектории. Первая показывает, что система вернула деньги, не
проверив возможность возврата по заказу. Вторая --- что после истечения
времени ожидания запроса к складскому учёту система не переключилась на
получение данных из кэша. Модель сопоставляет траектории с архитектурной
документацией и описанием требований, формирует отчёт о проблемах со
ссылками на траектории и ходы, а затем создаёт три исполняемых
регрессионных теста.

На \crossreflink{fig:5-7}{рисунке 5-7} диагностический отчёт расположен
между траекторией и регрессионными тестами. В отчёте указано: «Система
пропустила проверку и сразу вернула деньги». Вернитесь к траектории и
найдите доказательство: какого шага не хватает в журнале вызовов? Затем
посмотрите на тесты: обнаружат ли они повторный пропуск этого шага в
будущем?

\begin{figure}
\hypertarget{fig:5-7}{%
\centering
\includegraphics{images/revised/fig5-7.pdf}
\caption{Рисунок 5-7. Конвейер интеллектуальной диагностики
производственных журналов}\label{fig:5-7}
}
\end{figure}

Все три теста завершаются неудачей на дефектной версии и успешно
проходят после её замены на исправленную. Один из тестов проверяет
только факт вызова программой инструмента проверки возможности возврата.
Этот тест пройдёт, даже если программа сначала вернёт деньги и лишь
затем проверит возможность возврата по заказу. В журнале запуска также
сохранено подтверждение создания рабочей задачи через GitHub MCP, по
которому можно найти задачу, соответствующую переданному разработчику
диагностическому отчёту.
\end{quote}

\hypertarget{ux43aux43eux434-ux43aux430ux43a-ux433ux435ux43dux435ux440ux430ux442ux438ux432ux43dux44bux439-ui}{%
\subsection{Код как генеративный
UI}\label{ux43aux43eux434-ux43aux430ux43a-ux433ux435ux43dux435ux440ux430ux442ux438ux432ux43dux44bux439-ui}}

Если пользователь также должен самостоятельно изменять материалы, над
которыми работает AI-агент, интерфейс можно превратить в общее рабочее
пространство для человека и агента. Например, в маршруте путешествия
пользователь может видеть, какой пункт сейчас изменяет помощник,
перехватывать редактирование, подтверждать или отменять изменения. Если
пользователь только что перенёс дату отправления на пятницу, помощник
должен продолжить работу уже с обновлённым маршрутом. Одновременно
просматривать и изменять данные можно лишь тогда, когда обе стороны
совместно используют одно видимое состояние и заранее определено, кто в
текущий момент управляет редактируемым полем.

Представьте, что пользователь хочет сравнить три варианта поездки. Пока
место назначения не определено, удобнее уточнять его в диалоге. После
выбора направления пользователю проще непосредственно указать даты и
число путешественников, чем многократно сверять несколько текстовых
фрагментов. Наконец, изменение общей стоимости в зависимости от дат
нагляднее всего представить на графике. Агент генерирует подходящие для
задачи формы, диаграммы или интерфейсы приложений, в которых
пользователь может непосредственно вводить данные, выбирать варианты и
выполнять действия. Это и называется генеративным UI (Generative UI).

\textbf{Описывать интерфейс и взаимодействие посредством протокола.}

Существует два способа генерировать интерфейсы: выводить код HTML,
JavaScript и других языков для выполнения в контролируемой среде либо
выводить декларативную структуру, которую визуализируют компоненты
клиента. Первый способ упрощает создание нестандартных макетов и
взаимодействий, второй ограничивает доступные возможности набором
компонентов и действий, предоставленных клиентом.

A2UI (Agent-to-User Interface) использует декларативное описание. Агент
создаёт структурированные сообщения, описывающие компоненты, привязку
данных и обновления, а клиент отображает интерфейс с помощью собственных
компонентов\footnote{Официальная документация A2UI и \emph{AG-UI
  Overview} описывают соответственно декларативное представление
  интерфейса и обмен событиями между агентом и фронтендом.}. Например,
приложение может предоставлять компоненты карточки, кнопки, текстового
поля и таблицы, а агент --- описать с их помощью таблицу с заголовком
«Данные о продажах», тремя строками и двумя столбцами. Каталог
компонентов также можно расширить диаграммами, картами и другими
возможностями.

AG-UI (Agent--User Interaction) обеспечивает обмен событиями между
агентом и фронтендом, передавая сообщения, вызовы инструментов и
обновления состояния, а также может переносить сообщения интерфейса
A2UI. При совместном использовании A2UI описывает отображаемое
содержимое, а AG-UI связывает выполнение на бэкенде с взаимодействием на
фронтенде.

Клиент может получать обновления интерфейса поэтапно и отображать их в
веб-, мобильной или настольной среде с помощью адаптированного средства
визуализации. Клиент проверяет свойства компонентов, привязку данных и
запросы действий, а сервер продолжает проверять бизнес-права. Для
ссылок, форматированного текста и доступа к данным внутри компонентов
также необходимы отдельные правила обработки.

При непосредственной генерации кода необходимо проверять зависимости,
доступ к ресурсам и разрешения на выполнение, а также изолировать
сгенерированный интерфейс от конфиденциального состояния
хост-приложения. Инъекция в промпт из внешних материалов может повлиять
на процесс генерации, а риск выполнения скриптов в браузере зависит ещё
и от способа загрузки и изоляции страницы.

\textbf{Предоставлять интерактивные результаты в HTML.}

\textbf{Для результатов исследований, включающих данные, демонстрации и
непрерывные обновления, автор отдаёт предпочтение HTML-сайтам.} Длинный
линейный отчёт в Markdown заставляет читателя постоянно искать контекст.
Интерактивная веб-страница позволяет собрать пояснения, данные и
демонстрации в одной точке входа: читатель просматривает диаграммы,
фильтрует интересующие его примеры, раскрывает подробности отдельных
результатов, а затем взаимодействует с демонстрацией, чтобы понять
поведение системы. По мере выполнения задачи агент также может обновлять
страницу, непрерывно фиксируя на ней актуальные результаты и
доказательства.

Например, в исследовательской практике автора каждому проекту
соответствует интерактивный сайт, который одновременно служит живым
документом в ходе исследования и итоговым результатом. По мере
проведения экспериментов агент постоянно дополняет его содержимое. Такие
сайты решают три основные задачи.

Первая --- \textbf{прослеживаемость экспериментальных данных}. Страница
связывает примеры, промпты, ответы модели и оценки, позволяя проверять
их по отдельности. Например, если ответы одного типа постоянно получают
низкую оценку, по истории оценивания можно найти исходные ответы и
промпты.

Вторая --- \textbf{мониторинг показателей обучения}. Сигналы, отражающие
внутреннее состояние обучения, автор называет «клиническими
показателями»: потери на обучающей и валидационной выборках, нормы
градиентов, скорость обучения, а при обучении с подкреплением ---
вознаграждение, дивергенция KL, энтропия стратегии и другие значения.
Для языковой модели также можно использовать перплексию, чтобы измерять
качество предсказания заданной последовательности текста. Подобно
физиологическим показателям при медицинском обследовании, эти кривые
иногда обнаруживают проблемы раньше, чем итоговые результаты задачи. Их
совместный просмотр с точностью выполнения задач помогает определить, на
каком этапе возникли отсутствие сходимости, аномалии градиентов или
деградация обучения.

Третья --- \textbf{демонстрация принципа работы}. При объяснении
взаимодействия нескольких модулей страница может позволить читателю
изменить входные данные, просмотреть промежуточное состояние, а затем
сопоставить изменение результата с записями эксперимента.

Если работа продолжается несколько часов, как пользователю время от
времени подключаться к ней? При использовании Coding Agent автор хочет,
чтобы короткая устная команда точно передавалась агенту, ведущему
задачу, чтобы результаты можно было открыть в любой момент, а
уведомления приходили только при появлении новой версии или
необходимости принять решение. Например, одновременно создаются два
отчёта. Пользователь открывает отчёт о продажах, указывает на вторую
диаграмму и говорит: «Перестрой её по чистой выручке». Интерфейс должен
передать вместе с командой идентификаторы текущего отчёта и диаграммы.
Когда диаграмма будет готова, уведомление должно позволить сразу открыть
новую версию. Если для «чистой выручки» доступны два способа расчёта,
помощник должен предложить пользователю выбрать один из них.

\textbf{Уточнять намерение пользователя.}

Когда пользователь говорит: «Я хочу купить билет на самолёт до Пекина»,
место назначения уже известно, но ещё требуется уточнить город
отправления, дату и тип поездки --- в одну сторону или туда и обратно.
Агент может собрать все вопросы в одном раунде диалога либо
сгенерировать форму, чтобы связанные условия были представлены на одной
странице.

Элементы управления формы соответствуют разным типам информации:
текстовые поля собирают открытые ответы, раскрывающиеся меню и
переключатели перечисляют доступные варианты, флажки поддерживают
множественный выбор, а выбор даты помогает указать время. JavaScript
позволяет описывать зависимости между полями: например, после выбора
варианта «туда и обратно» показывать дату возвращения и проверять, что
она не предшествует дате отправления. При изменении выбора также
необходимо очищать или игнорировать поля, утратившие актуальность.

\crossreflink{fig:5-8}{Рисунок 5-8} взят из одного запуска следующего
эксперимента. Проследите эту отправку данных: среди четырёх полей,
возвращённых браузером, нет места назначения. Откуда модель знает, что
пассажир направляется в Пекин?

\begin{figure}
\hypertarget{fig:5-8}{%
\centering
\includegraphics{images/revised/fig5-8.pdf}
\caption{Рисунок 5-8. От формы покупки авиабилета к полному
маршруту}\label{fig:5-8}
}
\end{figure}

\begin{quote}
\textbf{Эксперимент 5-12 ★★: система уточнения намерения посредством
динамически генерируемой формы}

Эксперимент начинается с запроса «Я хочу купить билет на самолёт до
Пекина». Модель должна сгенерировать HTML-форму: текстовое поле для
города отправления, выбор даты отправления, варианты поездки в одну
сторону и туда и обратно, а также дату возвращения, которая отображается
при выборе поездки туда и обратно.

Программа передаёт модели исходное поручение вместе с данными формы.
Если сохранить только отправленные поля, место назначения всё равно
будет отсутствовать, даже когда каждый элемент управления заполнен
правильно.

Выполните ещё одну проверку с той же формой: заполните дату возвращения,
а затем выберите поездку в одну сторону. Страница скроет поле даты
возвращения, но отправленный JSON сохранит и вариант поездки в одну
сторону, и дату возвращения 18 августа. При переключении на поездку в
одну сторону клиент также должен очистить ставшую недействительной дату
возвращения.
\end{quote}

\textbf{Генерировать SQL-запросы.}

Пользователь может сформулировать запрос к базе данных на естественном
языке, а агент --- сгенерировать SQL на основе структуры таблиц и
бизнес-правил расчёта. \textbf{Для отображения полной таблицы
результатов автор отдаёт предпочтение режиму Artifact, в котором данные
поступают из базы непосредственно в интерфейс.} Если заставлять модель
читать и построчно пересказывать большую таблицу, возрастут расход
токенов и время ожидания, а также появится риск ошибок при переносе
данных. SQL сохраняется как доступный для проверки артефакт запроса.
Если требуется проанализировать причины или обобщить тенденции, модели
передаются лишь необходимые агрегированные результаты.

На \crossreflink{fig:5-9}{рисунке 5-9} пользователь хочет узнать число
работающих сотрудников в каждом отделе, и страница может непосредственно
отобразить статистическую таблицу, возвращённую базой данных. Если затем
пользователь спросит: «На сколько человек в отделе маркетинга больше,
чем в отделе разработки?», модели достаточно взять из таблицы числа 9 и
5 и вычислить разницу в 4 человека. Если же последует вопрос: «В прошлом
месяце разница была такой же?», текущей таблицы уже недостаточно ---
потребуется повторно выполнить запрос со статистической датой прошлого
месяца.

\begin{figure}
\hypertarget{fig:5-9}{%
\centering
\includegraphics{images/revised/fig5-9.pdf}
\caption{Рисунок 5-9. Процесс выполнения SQL-запроса
агентом}\label{fig:5-9}
}
\end{figure}

Сервис выполнения определяет доступные таблицы, столбцы и строки на
основе доверенной идентификационной информации, использует ограниченную
учётную запись базы данных и проверяет структуру SQL и тип операции. В
сценариях выполнения запросов можно разрешить только одобренные запросы
на чтение, ограничить вызовы функций, бюджет времени, число возвращаемых
строк и диапазон данных, а введённые пользователем значения передавать
посредством параметризованной привязки.

Агент также может создавать визуальные артефакты, например представлять
сводные данные по отделам в виде столбчатой диаграммы. Сервис выполнения
передаёт результаты запроса в контролируемую среду визуализации, а
интерфейс отображает только разрешённые данные и взаимодействия.

\begin{quote}
\textbf{Эксперимент 5-13 ★★: ERP-агент с взаимодействием на естественном
языке}

На примере запросов к данным сотрудников и заработной платы в
ERP-системе --- системе планирования ресурсов предприятия --- этот
эксперимент показывает, как генерировать SQL из бизнес-вопросов и
передавать результаты базы данных непосредственно на страницу. В
PostgreSQL создаются две таблицы. Таблица сотрудников содержит
идентификатор сотрудника, имя, отдел, уровень, дату приёма и дату
увольнения; пустая дата увольнения означает, что сотрудник продолжает
работать. Таблица заработной платы содержит идентификатор сотрудника,
дату выплаты и сумму зарплаты; для каждого сотрудника предусмотрена одна
запись в месяц.

На основе этих двух таблиц агент обрабатывает следующие десять вопросов:

\begin{enumerate}
\def\labelenumi{\arabic{enumi}.}
\tightlist
\item
  Какова средняя продолжительность работы одного сотрудника?
\item
  Сколько работающих сотрудников в каждом отделе?
\item
  В каком отделе самый высокий средний уровень сотрудников?
\item
  Сколько новых сотрудников было принято в каждый отдел в этом и прошлом
  году?
\item
  Какова средняя зарплата в отделе A с марта позапрошлого года по май
  прошлого года?
\item
  В каком отделе средняя зарплата за прошлый год была выше: A или B?
\item
  Какова средняя зарплата сотрудников каждого уровня в этом году?
\item
  Какова средняя зарплата за последний месяц у сотрудников со стажем до
  года, от одного до двух лет и от двух до трёх лет?
\item
  Какие десять сотрудников получили наибольшее повышение зарплаты с
  прошлого по текущий год?
\item
  Есть ли задолженность по зарплате --- месяцы, в которые сотрудник
  работал, но не получил выплату?
\end{enumerate}

Перед выполнением запросов необходимо определить правила расчёта.
Например, стаж работающих сотрудников рассчитывается до базовой даты
запроса, а уволенных --- до даты увольнения; для «средней зарплаты»
нужно решить, усреднять ли записи о выплатах или средние значения по
сотрудникам; при ранжировании повышения зарплаты необходимо учесть
сотрудников без записей за прошлый год.

В этом запуске были сгенерированы и выполнены десять SQL-запросов,
результаты которых совпали с независимым эталонным расчётом, после чего
браузер непосредственно отобразил таблицы. Модель получает структуру
таблиц, задачу и необходимую обратную связь о выполнении, а строки базы
данных передаются приложением в интерфейс.
\end{quote}

\textbf{Динамически генерировать программное обеспечение.}

Агент может на основе требований сгенерировать интерфейс и логику
взаимодействия, а затем продолжить их изменять в ответ на действия
пользователя. Представленная Anthropic в 2025 году временная
исследовательская предварительная версия Imagine with Claude
демонстрировала такое взаимодействие: пользователь исследует интерфейс,
а модель непрерывно генерирует соответствующий программный
продукт\footnote{Anthropic, \emph{Introducing Claude Sonnet 4.5} ---
  описание выпущенной вместе с моделью временной исследовательской
  предварительной версии Imagine with Claude.}.

\textbf{Для реальных продуктов автор считает более перспективной
настройку поверх существующего каркаса.} При непрерывной генерации всего
приложения с нуля приходится снова и снова оплачивать генерацию,
визуализацию и проверку. Стабильный каркас с возможностью изменять
важные для пользователя элементы интерфейса и процессы позволяет быстрее
предоставлять надёжно персонализированный продукт. Пользователь просит:
«Сделай кнопку синей», «Добавь в боковую панель меню быстрых действий»
или «Выбери более читаемый шрифт». Агент находит соответствующие
компоненты, изменяет код и проверяет результат в режиме предварительного
просмотра. Базовый каркас предоставляет стабильную структуру страниц и
бизнес-интерфейсы, а область настройки можно постепенно расширять от
цветов и шрифтов до макета и компонентов.

HMR (Hot Module Replacement, горячая замена модулей) в среде разработки
позволяет обновлять затронутые фронтенд-модули после изменения кода.
Сохранение состояния зависит от интеграции с фреймворком, границ
обновления и типа изменений. Если необходима полная перезагрузка,
приложение должно предварительно сохранить требуемое состояние. В свою
очередь, reload в режиме разработки FastAPI перезагружает серверную
программу: эти механизмы воздействуют соответственно на фронтенд-модули
и бэкенд-процесс\footnote{\emph{Vite HMR API} и \emph{FastAPI CLI}
  описывают соответственно горячее обновление фронтенда и автоматическую
  перезагрузку сервера разработки.}.

\begin{quote}
\textbf{Эксперимент 5-14 ★★: система диалоговой настройки интерфейса}

В этом эксперименте создаётся чат-приложение с фронтендом на React и
бэкендом на FastAPI. В среде разработки включаются HMR для фронтенда и
reload для бэкенда. Модель изменяет цвет, шрифт и заголовок по
требованиям трёх последовательных раундов, а браузер после каждого
изменения проверяет фактический интерфейс.

В этом запуске были зарегистрированы события HMR от Vite. Ни в одном из
трёх раундов не происходило переходов с полной загрузкой страницы,
существующее состояние чата сохранилось, а итоговая сборка прошла
успешно.
\end{quote}

\textbf{Отделять границы прав доступа от динамического бизнес-кода.}

\textbf{Чем проще агенту динамически переписывать бизнес-код, тем ниже,
ближе к стабильному уровню доступа к данным, должны располагаться
окончательные границы прав.} В традиционном приложении проверки прав
записываются в проверенном бизнес-коде. Если генерируемый код постоянно
изменяется, новый путь может пропустить такую же проверку. Сохранить
ограничения прав можно лишь в том случае, если каждая операция чтения и
записи проходит через независимый сервис доступа к данным. Например,
приложение для подбора персонала может динамически генерировать списки
кандидатов и формы операций, однако принадлежность кандидата организации
и поля, которые разрешено изменять рекрутеру, должен постоянно проверять
доверенный сервис доступа к данным.

Цель такого проектирования: \textbf{даже если в сгенерированном коде
возникнет ошибка, доверенный уровень доступа к данным продолжит
принудительно применять ограничения прав.} Это развивает принцип
защитных механизмов уровня данных из первой главы: правила, обязательные
для исполнения, следует размещать там, где проходят все соответствующие
операции. Политики безопасности на уровне строк базы данных могут
ограничивать доступ к записям, ограничения и валидаторы --- проверять
допустимость состояния, а контролируемые представления, хранимые
процедуры или сервисы данных --- предоставлять разрешённые операции.
Прикладной уровень может заранее выполнить проверку и объяснить её
пользователю, но окончательное чтение и запись всё равно контролируются
доверенным сервисом данных и правами базы данных.

Для каждой операции также необходим доверенный \textbf{контекст доступа}
(access context), содержащий сведения о пользователе, арендаторе, роли
или идентичности агента. Сервис аутентификации и доверенная среда
выполнения устанавливают и привязывают идентичность, а сервис данных
выполняет авторизацию на её основе. Динамический код получает
ограниченные интерфейсы, необходимые для выполнения задачи. Учётные
данные и полномочия на управление политиками находятся под контролем
независимых компонентов. Все подключения к базе данных, вспомогательные
запросы и фоновые процессы должны соблюдать одну область доступа.

Например, безопасность на уровне строк PostgreSQL должна использоваться
совместно с настройкой ролей базы данных. Приложение подключается под
ролью, на которую распространяются политики, а его учётная запись не
имеет права изменять политики доступа\footnote{PostgreSQL, \emph{Row
  Security Policies} --- описание настройки политик безопасности на
  уровне строк в зависимости от ролей и операций, а также области их
  применения.}. Механизмы привязки правил к идентичностям проходят
независимую проверку, тогда как бизнес-интерфейс может непрерывно
меняться в заданных границах.

Права доступа и целостность данных выполняют разные задачи. После
подтверждения права рекрутера изменять запись кандидата конечный автомат
дополнительно проверяет, разрешён ли переход из состояния «прошёл отбор»
в состояние «прошёл собеседование». Валидатор заработной платы читает
связанную вакансию и проверяет, находится ли ожидаемая зарплата в
установленном для неё диапазоне. Проверка ссылок подтверждает
действительность связи между вакансией и кандидатом.

\begin{quote}
\textbf{Эксперимент 5-15 ★★★: объект данных со встроенными правами для
динамически генерируемого программного обеспечения}

В этом эксперименте создаётся прототип объектного хранилища на Python
поверх PostgreSQL, в котором правила доступа, валидаторы, отношения
между объектами и последующие реакции объявляются внутри типов данных.
При чтении проверяются права доступа. При создании и обновлении сначала
проверяются права, бизнес-правила и ссылочные отношения, затем данные
сохраняются, после чего объявленные последующие реакции добавляются в
очередь с ограничением глубины цепочки.

Пример системы подбора персонала содержит объекты вакансий, кандидатов,
собеседований и оценок. Кандидат последовательно проходит этапы подачи
заявки, отбора, собеседования, направления предложения о работе и найма;
на соответствующих этапах также возможен переход в состояние отказа.
Диапазон заработной платы берётся из связанной вакансии. В примере
проверяются обычное обновление состояния, недопустимый переход
состояния, выход за диапазон заработной платы вакансии и доступ из
другого арендатора. Это помогает понять, какие задачи решают
соответственно авторизация и проверка бизнес-правил.

Прототип представляет пользователя, роль и атрибуты организации
посредством \texttt{AccessContext}; учебный скрипт напрямую создаёт
такие контексты для имитации разных идентичностей. При развёртывании
контекст создаётся в точке аутентификации, а объектное хранилище и
сгенерированный код размещаются по разные стороны явно определённой
границы сервиса или прав доступа.
\end{quote}

\hypertarget{ux43aux43eux434-ux441ux43eux437ux434ux430ux451ux442-ux43aux43eux434-ux441ux430ux43cux43eux437ux430ux433ux440ux443ux437ux43aux430-ai-ux430ux433ux435ux43dux442ux430}{%
\subsection{Код создаёт код: самозагрузка
AI-агента}\label{ux43aux43eux434-ux441ux43eux437ux434ux430ux451ux442-ux43aux43eux434-ux441ux430ux43cux43eux437ux430ux433ux440ux443ux437ux43aux430-ai-ux430ux433ux435ux43dux442ux430}}

Возможности генерации кода можно использовать и для разработки самих
агентов. Создатель получает требования к задаче, генерирует или изменяет
промпт, инструменты и цикл выполнения нового агента, а затем проверяет
результат с помощью тестов и реальных взаимодействий. На
\crossreflink{fig:5-10}{рисунке 5-10} показан этот цикл самозагрузки:
существующая реализация служит отправной точкой, обратная связь по
задаче направляет изменения, а результаты проверки определяют, какие из
них можно сохранить.

\begin{figure}
\hypertarget{fig:5-10}{%
\centering
\includegraphics{images/revised/fig5-10.pdf}
\caption{Рисунок 5-10. Цикл самозагрузки агента}\label{fig:5-10}
}
\end{figure}

Как отличить изменение, улучшившее общие способности, от изменения,
которое случайно подстроилось под только что увиденную задачу? Оставьте
одну задачу, которую создателю не показывали. Например, разработайте
агента мониторинга цен на данных одного магазина, а затем передайте ему
пример из другого магазина, одновременно повторно запустив исходную
задачу.

\textbf{Обслуживайте среду выполнения агента.}

Работа нового агента зависит от конфигурации, учётных данных, состояния
сеанса, сервисов и песочницы. Команда \texttt{doctor} в OpenClaw служит
примером инструмента обслуживания: она проверяет миграцию старых
конфигураций, состояние учётных данных OAuth, блокировки сеансов,
установку плагинов, работоспособность шлюза, образы песочницы и другие
возможные проблемы, после чего предлагает рекомендации по исправлению
или сама выполняет соответствующие исправления\footnote{OpenClaw, Doctor
  --- описание функций проверки конфигурации и состояния, миграции и
  исправления.}.

Эти детерминированные проверки могут стать основой процесса
самообслуживания. Правила обнаружения охватывают известные сбои, а
действия по исправлению имеют чётко определённые входные данные и
ожидаемые результаты. При дальнейшем проектировании процесса диагностики
агента можно передавать модели нерешённые проблемы вместе с журналами,
конфигурацией и результатами проверок, чтобы она объяснила, какую
конфигурацию и каким образом намерена изменить. Затем исполнительный
слой проверяет права доступа и область изменений, после чего запускает
проверку работоспособности и регрессионные тесты. Снимки конфигурации,
журнал изменений и механизм отката помогают восстановить работу системы
после неудачного исправления.

\textbf{Закрепляйте инженерный опыт в эталонной реализации.}

Создание специализированного агента одновременно затрагивает логику
предметной области и общую инфраструктуру. При проверке сгенерированного
результата можно начать с четырёх категорий вопросов:

\begin{enumerate}
\def\labelenumi{\arabic{enumi}.}
\tightlist
\item
  \textbf{Контекст и цикл выполнения.} Сохраняйте связи между ролями,
  идентификаторами вызовов инструментов и результатами инструментов;
  проверяйте состояние в многоходовых взаимодействиях, условия
  завершения, бюджет и восстановление после исключений.
  Структурированные сообщения упрощают поддержку связей протокола, а
  стабильный префикс создаёт условия для кэширования промпта.
\item
  \textbf{Интерфейсы инструментов.} Описание инструмента должно
  объяснять его назначение, условия применимости и границы
  использования, а параметры --- содержать типы, смысл и конкретные
  примеры. Исполнительный слой по-прежнему должен проверять параметры,
  идентичность и состояние бизнес-процесса.
\item
  \textbf{Выбор модели и API.} Модель склонна использовать старые модели
  и интерфейсы, часто встречавшиеся в обучающих данных, поэтому даже
  хорошо написанный код может опираться на устаревший выбор технологий.
  Предоставьте создателю актуализируемые материалы по выбору технологий
  или разрешите ему обращаться к официальной документации, чтобы сверять
  возможности моделей, версии интерфейсов и способы вызова. После
  генерации выполните реальные вызовы, чтобы проверить работоспособность
  конфигурации.
\item
  \textbf{Совместимость зависимостей и экосистемы.} Проверяйте состояние
  сопровождения библиотек, их версии и изменения интерфейсов, а также
  подтверждайте работоспособность собранной программы установкой
  зависимостей, проверкой компиляции и тестами.
\end{enumerate}

\textbf{При создании специализированного агента автор в первую очередь
предлагает создателю скопировать проверенную эталонную реализацию, а
затем адаптировать её к предметной области.} В промпте трудно
исчерпывающе описать все детали протокола сообщений, управления
состоянием и механизмов восстановления, тогда как эталонный код напрямую
переносит этот инженерный опыт в новую систему. Даже модель, хорошо
умеющая программировать, может упустить неявные условия выполнения
задачи; повторное использование фреймворка сокращает число таких
повторяющихся ошибок. Создатель копирует подходящий базовый фреймворк,
корректирует системный промпт под новую роль, добавляет или удаляет
инструменты предметной области, изменяет бизнес-логику и дополняет
соответствующие тесты. Фреймворк обеспечивает обработку сообщений,
связывание вызовов, управление состоянием и пути завершения, а изменения
предметной области проверяются новыми тестовыми сценариями.

Если шаблонный подход оказался намного быстрее, как исключить объяснение
«в нём просто реализовали меньше функций»? Поручим двум группам
создателей разработать один и тот же тип агента и проверим их по одному
набору критериев.

\begin{quote}
\textbf{Эксперимент 5-16 ★★★: разработка агента, способного создавать
агентов}

Попросите одну и ту же модель-создателя разработать агента проверки
релизов: он должен прочитать набор записей о проверках развёртывания и,
если хотя бы одна обязательная проверка не пройдена, отклонить релиз с
указанием основания. При создании с нуля модель должна реализовать цикл
выполнения, протокол инструментов, инструменты предметной области,
интерфейс командной строки и тесты. При шаблонном подходе сначала
копируется агент, умеющий проверять структурированные записи, а модели
остаётся лишь заполнить спецификацию предметной области. Например: в
каком поле хранится состояние и какие значения означают успешное
прохождение? Скрипт создания генерирует по спецификации системный
промпт, определения инструментов и документацию. Существующий инструмент
на Python читает спецификацию и выполняет проверку. Оба подхода
демонстрируют метапрограммирование: одна программа генерирует или
изменяет другие программы.

На \crossreflink{fig:5-11}{рисунке 5-11} оба подхода сходятся на одном
наборе критериев приёмки.

\begin{figure}
\hypertarget{fig:5-11}{%
\centering
\includegraphics{images/revised/fig5-11.pdf}
\caption{Рисунок 5-11. Конвейер агента, способного создавать
агентов}\label{fig:5-11}
}
\end{figure}

Для обоих подходов используются одинаковые требования к задаче и
критерии приёмки. Сначала проверяются структура результата, компиляция и
тесты, затем --- связь вызовов инструментов с их результатами в
стандартных сообщениях. Наконец, запускаются реальные задачи модели и
многоходовые взаимодействия, чтобы проверить обрезку контекста,
управление состоянием и ограничение количества вызовов.

В сохранённом контрольном эксперименте с Kimi K3 обе группы прошли 39 из
39 детерминированных проверок, а также успешно выполнили реальную задачу
модели и проверку состояния в многоходовом взаимодействии. При создании
шаблонным способом было использовано 1 181 токен и около 25 секунд, а
при создании с нуля --- 381 814 токенов и около 90 минут.
\end{quote}

\hypertarget{ux438ux442ux43eux433ux438-ux433ux43bux430ux432ux44b-4}{%
\section{Итоги
главы}\label{ux438ux442ux43eux433ux438-ux433ux43bux430ux432ux44b-4}}

\textbf{Агенты для программирования достигли зрелости сравнительно рано,
и у этого есть одна причина, которую легко упустить из виду: в
программной инженерии уже существуют системы типов, наборы тестов,
системы контроля версий и инструменты отладки.} Если код исправлен
неверно, тесты зачастую показывают, на каком наборе входных данных был
получен неожиданный результат; после другого исправления код можно снова
проверить на тех же данных. Чтобы перенести этот опыт на другие задачи,
сначала нужно задать три конкретных вопроса: как обнаруживать ошибки,
как возвращаться к известному состоянию и когда следует остановиться и
передать задачу человеку? Например, если соединение оборвалось во время
вызова возврата средств, Harness должен сначала проверить проводки и
лишь затем решать, следует ли повторять операцию; иначе восстановление
после одного сбоя может привести ко второму возврату.

Код также позволяет сохранить однажды принятое решение в форме,
пригодной для многократного использования. Перебор в математической
задаче и ограничения в логической задаче можно передать программе для
поэлементной проверки; правила возврата средств можно реализовать в виде
проверки на входе в операцию, чтобы сервер считывал фактическое
состояние заказа и только после этого решал, разрешать ли возврат. Если
изменился формат журнала, можно исправить парсер, а вызвавший ошибку
образец сохранить для регрессионного теста. Во всех трёх случаях можно
выполнять код, но критерии различаются: удовлетворяет ли решение
условиям задачи, соответствует ли возврат средств политике и способен ли
новый парсер прочитать ранее пропущенное поле.

Если результатом работы служит документ, видео или геометрическая
модель, необходимо также проверить качество рендеринга, синхронизацию
звука с изображением или точность размеров и попросить проверяющего
просмотреть фактический результат. Если результатом служит форма, запрос
или визуализация, помимо проверки того, может ли пользователь понять
результат и работать с ним, нужно проверить, получает ли исполняющий
сервис данные в соответствии с правами доступа. Если код создаёт другого
AI-агента, автор предпочитает сначала повторно использовать проверенный
фреймворк, сгенерировать инструменты предметной области и бизнес-логику,
а затем проверить протокол вызовов, состояние между ходами
взаимодействия и фактический результат задачи. Накопленные таким образом
инструменты, валидаторы и тесты можно повторно использовать в следующей
задаче.

Код позволяет выполнять задачу снова и снова, однако среда не
обязательно дождётся завершения выполнения. Пока запрос ещё выполняется,
пользователь уже может передумать. В шестой главе рассматривается, как
такое изменение включается в уже идущее взаимодействие.

\begin{questionbox}

\hypertarget{ux432ux43eux43fux440ux43eux441ux44b-ux434ux43bux44f-ux440ux430ux437ux43cux44bux448ux43bux435ux43dux438ux44f-4}{%
\section*{Вопросы для
размышления}\label{ux432ux43eux43fux440ux43eux441ux44b-ux434ux43bux44f-ux440ux430ux437ux43cux44bux448ux43bux435ux43dux438ux44f-4}}
\addcontentsline{toc}{section}{Вопросы для размышления}

\begin{enumerate}
\def\labelenumi{\arabic{enumi}.}
\tightlist
\item
  ★★ Спроектируйте песочницу для Coding Agent, которому необходимо
  читать файлы проекта, устанавливать зависимости и запускать тесты. Как
  бы вы настроили права доступа к файлам, сети, процессам и ресурсам?
  Если для задачи требуется расширить полномочия, как сформулировать
  запрос, проверить его область и сохранить запись?
\item
  ★★★ Когда создатель последовательно генерирует и изменяет несколько
  поколений агентов, дефекты фреймворка и ошибки предметной области
  могут постепенно накапливаться. Как отслеживать эти изменения с
  помощью эталонных версий, регрессионных тестов, оценки поведения и
  механизмов отката?
\item
  ★★ Изменение формата журналов может быть следствием как штатного
  обновления, так и дефекта системы, которая их формирует. Какие
  свидетельства должна сохранять система адаптивного синтаксического
  анализа и на основании каких условий она должна добавлять новый
  парсер, сообщать об аномалии или приостанавливать обработку?
\item
  ★★ В цикле «автор предложения --- проверяющий» проверяющий считает
  плотность информации на странице приемлемой, а пользователь ---
  чрезмерной. Как преобразовать предпочтения пользователя в примеры,
  критерии оценки и интерактивную обратную связь, а затем проверить,
  соответствуют ли последующие изменения требованиям?
\item
  ★★ Coding Agent может обновлять базу знаний, архитектурную
  документацию и инструкции проекта, а также закреплять
  последовательности операций в коде. Как выявлять избыточные или
  устаревшие правила и задавать для них область применимости, методы
  проверки и условия прекращения действия? Какие механизмы накопления и
  оценки необходимы, чтобы перейти от единичного успешного изменения к
  непрерывной эволюции, рассматриваемой в
  \crossreflink{chap:9}{главе 9}?
\item
  ★ В командах, удобно организованных для удалённой совместной работы,
  агентам обычно также проще получать рабочую информацию. Изучите
  требования, архитектуру, журнал решений и критерии приёмки знакомого
  проекта, перечислите документацию, которую важнее всего дополнить, и
  объясните, какие конкретные задачи она поможет выполнить агенту.
\item
  ★★★ Как организовать авторизацию на основе идентичности, маркировку
  источников, проверку потоков данных и контроль записи в память для
  агента, который одновременно обращается к закрытым данным,
  обрабатывает недоверенное содержимое, взаимодействует с внешними
  системами и использует долговременную память? Покажите расположение
  этих проверок на примере одного штатного бизнес-процесса.
\item
  ★★ В режимах Artifact для SQL и HTML модель генерирует артефакт,
  сервис выполнения обрабатывает запрос или выполняет рендеринг, а
  интерфейс непосредственно отображает результат. Какую информацию
  целесообразно снова передавать модели для объяснения? Как проверять
  права на выполнение запросов, бюджет ресурсов, изоляцию страницы и
  согласованность отображаемого результата?
\item
  ★★ Инструмент проверяет правила по бизнес-состоянию в базе данных, а
  описание параметров помогает модели подготовить вызов. На примере
  отмены бронирования объясните, за что отвечают правила на естественном
  языке, примеры параметров, проверка на уровне выполнения и тесты
  граничных случаев. Что необходимо синхронно обновить после изменения
  правил?
\end{enumerate}

\textbf{Практикум к \crossreflink{chap:5}{главе 5}.}
\href{https://github.com/bojieli/ai-agent-book/tree/40f4ce5778e6cc93e3606c3079bf8cf13a0c4d04/chapter5}{Каталог
экспериментов} ·
\href{https://github.com/bojieli/ai-agent-book/blob/40f4ce5778e6cc93e3606c3079bf8cf13a0c4d04/chapter5/README.ru.md}{Путеводитель
на русском}.

\end{questionbox}

\hypertarget{ux432ux437ux430ux438ux43cux43eux434ux435ux439ux441ux442ux432ux438ux435-ux440ux430ux441ux448ux438ux440ux435ux43dux438ux435-ux43fux440ux43eux441ux442ux440ux430ux43dux441ux442ux432ux430-ux43dux430ux431ux43bux44eux434ux435ux43dux438ux439-ux438-ux434ux435ux439ux441ux442ux432ux438ux439}{%
\chapter{\texorpdfstring{Взаимодействие: расширение пространства
наблюдений и
действий\label{chap:6}}{Взаимодействие: расширение пространства наблюдений и действий}}\label{ux432ux437ux430ux438ux43cux43eux434ux435ux439ux441ux442ux432ux438ux435-ux440ux430ux441ux448ux438ux440ux435ux43dux438ux435-ux43fux440ux43eux441ux442ux440ux430ux43dux441ux442ux432ux430-ux43dux430ux431ux43bux44eux434ux435ux43dux438ux439-ux438-ux434ux435ux439ux441ux442ux432ux438ux439}}

\textbf{Когда базовая модель остаётся неизменной, реорганизация или
расширение интерфейсов наблюдений и действий зачастую оказывается одним
из наиболее эффективных инженерных способов повысить качество выполнения
задач AI-агентом.} В предыдущих главах агент обычно ждал, пока
пользователь закончит говорить, затем искал информацию, вызывал
инструменты и давал ответ. Однако пока он ожидает результаты работы
инструментов, пользователь и внешняя среда продолжают меняться.

Пользователь поручает помощнику переоформить авиабилет. Пока помощник
уточняет у службы поддержки доступные рейсы, пользователь присылает ещё
одно сообщение: «Мне также нужно сдать одно место багажа»; представитель
службы поддержки всё ещё говорит, а отображаемая на экране стоимость
билета тем временем может измениться. Помощнику нужно своевременно
уточнить, включён ли зарегистрированный багаж в стоимость, а также
проверить, остаётся ли в силе только что названная цена. Здесь возникают
два вопроса проектирования: по каким каналам помощник может узнавать об
изменениях и в какие моменты он может на них реагировать? В этой главе
мы сначала рассмотрим асинхронно поступающие электронные письма и
результаты работы инструментов, а затем --- непрерывную речь, изменения
на экране и обратную связь от действий робота.

\hypertarget{ux440ux430ux441ux448ux438ux440ux435ux43dux438ux435-ux43cux43eux434ux430ux43bux44cux43dux43eux441ux442ux435ux439-ux438-ux43cux43eux43cux435ux43dux442ux43eux432-ux441ux440ux430ux431ux430ux442ux44bux432ux430ux43dux438ux44f}{%
\section{Расширение модальностей и моментов
срабатывания}\label{ux440ux430ux441ux448ux438ux440ux435ux43dux438ux435-ux43cux43eux434ux430ux43bux44cux43dux43eux441ux442ux435ux439-ux438-ux43cux43eux43cux435ux43dux442ux43eux432-ux441ux440ux430ux431ux430ux442ux44bux432ux430ux43dux438ux44f}}

\textbf{Модальность} описывает форму информации и действий. AI-агент
может читать текст, слушать звуки, видеть экран и получать сигналы
крутящего момента, а также выводить текст, синтезировать речь, нажимать
элементы интерфейса или приводить в движение суставы.

\textbf{Момент срабатывания} касается того, когда поступает информация,
когда начинается действие и когда система отвечает. Агент может
самостоятельно запрашивать информацию или получать её от внешних
сервисов; инструмент может немедленно вернуть результат либо сначала
запустить задачу и позднее сообщить о её завершении. При поступлении
нового события система также должна решить, следует ли прервать текущую
работу и как затем возобновить прерванную задачу.

\begin{longtable}[]{@{}
  >{\raggedright\arraybackslash}p{(\columnwidth - 4\tabcolsep) * \real{0.3056}}
  >{\raggedright\arraybackslash}p{(\columnwidth - 4\tabcolsep) * \real{0.3339}}
  >{\raggedright\arraybackslash}p{(\columnwidth - 4\tabcolsep) * \real{0.3605}}@{}}
\toprule\noalign{}
\begin{minipage}[b]{\linewidth}\raggedright
\hspace{0pt}Измерение
\end{minipage} & \begin{minipage}[b]{\linewidth}\raggedright
\hspace{0pt}Интерфейсы получения информации и наблюдения
\end{minipage} & \begin{minipage}[b]{\linewidth}\raggedright
\hspace{0pt}Интерфейсы действий и исполнения
\end{minipage} \\
\midrule\noalign{}
\endhead
\bottomrule\noalign{}
\endlastfoot
\hspace{0pt}Организация содержимого (\crossreflink{chap:2}{главы 2}--5)
& \hspace{0pt}Контекст, память и база знаний организуют информацию для
принятия решений & \hspace{0pt}Инструменты и код выражают исполняемые
операции \\
\hspace{0pt}Модальность (эта глава) & \hspace{0pt}Речь, экран,
физические датчики & \hspace{0pt}Речь, нажатия, движения суставов \\
\hspace{0pt}Момент срабатывания (эта глава) & \hspace{0pt}Активные
запросы, доставка событий, непрерывный ввод & \hspace{0pt}Исполнение в
течение нескольких ходов, отмена, вытеснение и возобновление \\
\end{longtable}

Контекст содержит информацию, необходимую для принятия конкретного
решения, а пространство наблюдений описывает возможный диапазон
информации, которую способна получить система. Инструменты предоставляют
конкретные интерфейсы, а пространство действий описывает диапазон
действий, разрешённых системе. При расширении интерфейсов эти диапазоны
по-прежнему ограничиваются возможностями устройств, параметрами,
разрешениями и средой выполнения.

Многие интерфейсы вызова инструментов организуют сообщения по ходам:
модель предлагает вызов, система добавляет результат инструмента, после
чего запускает следующий раунд инференса. Событийно-ориентированная
система должна управлять внешними изменениями поверх этого протокола.
Если интерфейс нативно поддерживает асинхронность, во время ожидания
инструмента можно принимать новые требования, а после получения
результата --- продолжать обработку. При использовании синхронного
интерфейса взаимодействие организует Harness с помощью очередей,
дескрипторов задач и механизмов планирования. Далее эти два подхода
рассматриваются отдельно.

Темп обратной связи сильно различается в разных сценариях. В следующей
таблице приведены распространённые временные масштабы.

\begin{longtable}[]{@{}
  >{\raggedright\arraybackslash}p{(\columnwidth - 6\tabcolsep) * \real{0.2226}}
  >{\raggedright\arraybackslash}p{(\columnwidth - 6\tabcolsep) * \real{0.3188}}
  >{\raggedright\arraybackslash}p{(\columnwidth - 6\tabcolsep) * \real{0.1906}}
  >{\raggedright\arraybackslash}p{(\columnwidth - 6\tabcolsep) * \real{0.2680}}@{}}
\toprule\noalign{}
\begin{minipage}[b]{\linewidth}\raggedright
\hspace{0pt}Типичный масштаб
\end{minipage} & \begin{minipage}[b]{\linewidth}\raggedright
\hspace{0pt}Сценарий
\end{minipage} & \begin{minipage}[b]{\linewidth}\raggedright
\hspace{0pt}Изменения в наблюдениях
\end{minipage} & \begin{minipage}[b]{\linewidth}\raggedright
\hspace{0pt}Обработка действий
\end{minipage} \\
\midrule\noalign{}
\endhead
\bottomrule\noalign{}
\endlastfoot
\hspace{0pt}Секунды --- дни & \hspace{0pt}Асинхронная и
событийно-\allowbreak{}\hspace{0pt}ориентированная работа &
\hspace{0pt}Внешнее событие активирует агента & \hspace{0pt}Сначала
запускается задача, затем последующие события сообщают о ходе выполнения
и завершении \\
\hspace{0pt}Десятки миллисекунд --- секунды & \hspace{0pt}Речь &
\hspace{0pt}Аудио поступает непрерывно & \hspace{0pt}Потоковый вывод с
обработкой перебиваний и остановки воспроизведения \\
\hspace{0pt}Доли секунды --- секунды & \hspace{0pt}Computer Use &
\hspace{0pt}Экран непрерывно меняется & \hspace{0pt}После действия
состояние интерфейса запрашивается заново \\
\hspace{0pt}Цикл управления миллисекундного масштаба &
\hspace{0pt}Низкоуровневое управление роботом & \hspace{0pt}Данные от
датчиков поступают непрерывно & \hspace{0pt}Выполняются короткие
фрагменты действий с реакцией на сигналы управления и остановки \\
\end{longtable}

\hypertarget{ux430ux441ux438ux43dux445ux440ux43eux43dux43dux43eux441ux442ux44c-ux438-ux441ux43eux431ux44bux442ux438ux439ux43dux43e-ux43eux440ux438ux435ux43dux442ux438ux440ux43eux432ux430ux43dux43dux430ux44f-ux430ux440ux445ux438ux442ux435ux43aux442ux443ux440ux430-ux43aux43eux433ux434ux430-ux43cux438ux440-ux43eux431ux440ux430ux449ux430ux435ux442ux441ux44f-ux43a-ux430ux433ux435ux43dux442ux443-ux441ux430ux43c}{%
\section{Асинхронность и событийно-ориентированная архитектура: когда
мир обращается к агенту
сам}\label{ux430ux441ux438ux43dux445ux440ux43eux43dux43dux43eux441ux442ux44c-ux438-ux441ux43eux431ux44bux442ux438ux439ux43dux43e-ux43eux440ux438ux435ux43dux442ux438ux440ux43eux432ux430ux43dux43dux430ux44f-ux430ux440ux445ux438ux442ux435ux43aux442ux443ux440ux430-ux43aux43eux433ux434ux430-ux43cux438ux440-ux43eux431ux440ux430ux449ux430ux435ux442ux441ux44f-ux43a-ux430ux433ux435ux43dux442ux443-ux441ux430ux43c}}

В \crossreflink{chap:4}{главе 4} были представлены инструменты
восприятия, выполнения и совместной работы, а в этом разделе
рассматриваются инструменты запуска по событию и взаимодействия с
пользователем. Инструменты запуска по событию передают изменения среды
AI-агенту, а инструменты взаимодействия позволяют агенту связываться с
пользователем в ходе выполнения задачи. Harness соотносит эти входные и
выходные данные с выполняемой в данный момент задачей, формируя
непрерывно работающий событийный цикл.

\hypertarget{ux437ux430ux447ux435ux43c-ux43dux443ux436ux43dux430-ux430ux441ux438ux43dux445ux440ux43eux43dux43dux43eux441ux442ux44c}{%
\subsection{Зачем нужна
асинхронность}\label{ux437ux430ux447ux435ux43c-ux43dux443ux436ux43dux430-ux430ux441ux438ux43dux445ux440ux43eux43dux43dux43eux441ux442ux44c}}

При синхронном вызове вызывающая сторона может продолжить текущий
процесс только после получения результата. Асинхронный вызов сначала
возвращает дескриптор задачи --- идентификатор для запроса её состояния
и управления ею --- либо предусматривает последующее уведомление о
завершении, позволяя вызывающей стороне во время ожидания заниматься
другой работой. Благодаря этому система может организовать несколько
выполняющихся задач, а затем определять, какие операции могут
выполняться параллельно, исходя из ресурсов и зависимостей.

AI-агента можно представить как секретаря, который работает с
электронной почтой, телефонными звонками и посетителями: в ожидании
ответа он продолжает готовить материалы, при поступлении срочного звонка
сохраняет текущий прогресс, а после разговора возвращается к прежней
задаче. В инженерной реализации необходимо решить три задачи:

\begin{itemize}
\tightlist
\item
  \textbf{Отделение длительных задач от взаимодействия.} Анализ файлов,
  телефонные звонки и сборка выполняются в фоне, а пользователь
  по-прежнему может запрашивать прогресс и дополнять требования.
\item
  \textbf{Приоритет событий.} Срочная остановка своевременно поступает в
  контур управления, обычные обновления добавляются в очередь, а
  независимые запросы можно передать параллельной сессии.
\item
  \textbf{Прерывание и возобновление.} Необходимо сохранять состояние
  задачи, фиксировать уже выполненные операции, а при возобновлении
  повторно проверять, какие условия по-прежнему действительны.
\end{itemize}

Исследовательская предварительная версия Claude Duet предлагает такой
пример: в два часа дня у пользователя назначена встреча, и он хочет
заранее получить справку о компании другой стороны. Он может оставить
голосовое сообщение, а затем закрыть приложение. Подготовка продолжится
в облаке, и пользователь сможет в любой момент открыть страницу задачи и
посмотреть прогресс; если потребуются дополнительные сведения или
справка будет готова, помощник вернётся в основной диалог и уведомит
его. Продолжение задачи после ухода пользователя --- это асинхронное
выполнение. Для полнодуплексного телефонного разговора, который будет
рассмотрен далее, требуется ещё одна возможность: когда одна сторона
говорит, она всё равно может услышать, как другая её
перебивает.\footnote{Anthropic, \emph{Get started with Claude Duet}:
  вводные материалы перед встречей, страницы задач и уведомления в
  основном диалоге.}

API моделей также задают порядок сообщений. Одни интерфейсы принимают
запрос следующего цикла только после полного сопоставления вызовов
инструментов с их результатами, другие позволяют оставить вызов
инструмента незавершённым и получать последующие обновления. Асинхронные
функции приложения управляют ожиданием и выполнением, а протокол модели
определяет, когда можно отправлять сообщения.

На \crossreflink{fig:6-1}{рисунке 6-1} новые сообщения и результаты
инструментов могут поступать в очередь, однако способ их доставки
зависит от источника. Если сервис поддерживает push-уведомления, он
может сообщить об изменении состояния сразу; если сервис предоставляет
только интерфейс запросов, адаптер должен периодически опрашивать его и
при обнаружении изменений создавать событие. Событие также может
возникнуть по истечении таймера. Затем цикл событий передаёт эти входные
данные агенту.

\begin{figure}
\hypertarget{fig:6-1}{%
\centering
\includegraphics{images/revised/fig6-1.pdf}
\caption{Рисунок 6-1. Событийно-ориентированная асинхронная архитектура
агента}\label{fig:6-1}
}
\end{figure}

\hypertarget{ux440ux435ux430ux43bux438ux437ux430ux446ux438ux44f-ux441ux43eux431ux44bux442ux438ux439ux43dux43e-ux43eux440ux438ux435ux43dux442ux438ux440ux43eux432ux430ux43dux43dux43eux433ux43e-ux43cux435ux445ux430ux43dux438ux437ux43cux430-ux432-openclaw}{%
\subsection{Реализация событийно-ориентированного механизма в
OpenClaw}\label{ux440ux435ux430ux43bux438ux437ux430ux446ux438ux44f-ux441ux43eux431ux44bux442ux438ux439ux43dux43e-ux43eux440ux438ux435ux43dux442ux438ux440ux43eux432ux430ux43dux43dux43eux433ux43e-ux43cux435ux445ux430ux43dux438ux437ux43cux430-ux432-openclaw}}

OpenClaw принимает сообщения из разных каналов через Gateway и
маршрутизирует события в среду выполнения агента. Механизмы
автоматизации включают внутренние Hooks, планирование по времени,
Heartbeat, внешние Webhook и другие точки входа\footnote{OpenClaw,
  Automation: такие точки входа для автоматизации, как внутренние Hooks,
  запуск по расписанию, Heartbeat и внешние Webhook.}. Эти механизмы
можно классифицировать по условиям срабатывания:

\begin{itemize}
\tightlist
\item
  \textbf{Hooks (событийные хуки)} реагируют на внутренние события
  фреймворка, например создание или сброс сессии либо изменение
  жизненного цикла выполнения.
\item
  \textbf{Cron (планирование по времени)} выполняет работу в заданное
  время или с заданной периодичностью. Например, \texttt{0\ 9\ *\ *\ 5}
  означает срабатывание каждую пятницу в 9:00 в настроенном часовом
  поясе.
\item
  \textbf{Heartbeat (сердцебиение)} периодически пробуждает агента,
  чтобы проверить, нет ли требующих внимания дел.
\item
  \textbf{Webhook (внешний callback)} принимает события от внешних
  сервисов и запускает их обработку после аутентификации и
  маршрутизации.
\end{itemize}

Еженедельные отчёты удобно создавать по расписанию, требующие
постоянного внимания задачи можно проверять периодически, а
push-доставка новых писем и обратных вызовов сервисов позволяет
сократить задержку обнаружения.

\textbf{События звонков PineClaw.} Pine AI совершает телефонные звонки
вместо пользователя, решая такие задачи, как переговоры по счетам,
отмена подписок и урегулирование страховых случаев. PineClaw подключает
подобные телефонные задачи к OpenClaw. Звонок обычно длится некоторое
время, и в процессе может потребоваться участие пользователя:

\begin{itemize}
\tightlist
\item
  \textbf{Проверка личности.} Сотрудник службы поддержки просит
  владельца учётной записи пройти проверку, например ввести секретный
  код или одноразовый пароль.
\item
  \textbf{Трёхсторонний звонок.} Сотрудник службы поддержки требует
  напрямую поговорить с владельцем учётной записи, поэтому пользователя
  необходимо своевременно уведомить, чтобы он ответил на звонок.
\item
  \textbf{Подтверждение решения.} Когда сотрудник службы поддержки
  предлагает альтернативный вариант, клиентская часть с учётом текущего
  поручения и прежних предпочтений пользователя определяет, можно ли его
  принять. Если решение должен принять пользователь, ему отправляются
  конкретные варианты, а параллельно продолжается уточнение остальных
  условий у сотрудника поддержки.
\end{itemize}

Например, сотруднику службы поддержки требуется подключить владельца
учётной записи к трёхстороннему звонку. Чем раньше будет доставлено
уведомление, тем меньше времени обеим сторонам придётся ждать на линии.
В PineClaw используется механизм Channel, передающий события звонка
между Gateway и Pine API: установление соединения, необходимость
пользовательского ввода и завершение разговора поступают в поток как
отдельные события, на основании которых агент обновляет задачу и
уведомляет пользователя.

\textbf{По мнению автора, для своевременного проактивного обслуживания
лучше всего, чтобы источник события уведомлял агента при изменении
состояния.} После доставки уведомления через канал Harness находит
соответствующую задачу и проверяет, не является ли событие повторным или
просроченным, а затем модель решает, что делать дальше. Если
пользователю необходимо подтвердить личность, интерфейс должен открыть
доверенную точку входа для аутентификации, после чего её результат
возвращается в исходную задачу как новое событие.

\hypertarget{ux438ux43dux441ux442ux440ux443ux43cux435ux43dux442ux44b-ux437ux430ux43fux443ux441ux43aux430-ux43fux43e-ux441ux43eux431ux44bux442ux438ux44e}{%
\subsection{Инструменты запуска по
событию}\label{ux438ux43dux441ux442ux440ux443ux43cux435ux43dux442ux44b-ux437ux430ux43fux443ux441ux43aux430-ux43fux43e-ux441ux43eux431ux44bux442ux438ux44e}}

Если поручить помощнику «напомнить мне после ответа клиента», как он
проснётся завтра точно в нужный момент? Сначала следует проверить, может
ли почтовая система отправить событие при получении письма; если система
допускает только активные запросы, необходимо запомнить время следующей
проверки. Три приведённых ниже условных интерфейса передают эти условия
ожидания на хранение среде выполнения, позволяя модели завершить текущую
итерацию работы.

\textbf{Таймер (\texttt{set\_timer})} используется для пробуждения
задачи в назначенное время. Одноразовый таймер подходит для задачи «в
следующий понедельник в 10:00 позвонить в банк и узнать статус
оформления ипотеки», а периодический --- для ежечасной проверки
состояния сервиса. С помощью таймера также можно запланировать ожидание
ответа после отправки письма или повторный звонок после наступления
рабочего времени. Интерфейс должен сохранять часовой пояс, условие
срабатывания и идентификатор задачи, а также предоставлять возможность
отмены или изменения. Если внешний сервис поддерживает только запросы,
задача по расписанию может также выполнять опрос статуса.

\textbf{Мониторинг фоновых задач (\texttt{monitor\_shell})} используется
для преобразования длительно выполняющейся команды в события о ходе
работы. При запуске возвращается идентификатор задачи, во время
выполнения среда уведомляет о новом выводе, появлении ключевых слов или
изменении состояния, а после завершения сообщает статус выхода и
результат.

Если заставлять модель постоянно запрашивать прогресс, каждая такая
проверка будет требовать вызова модели, даже когда состояние не
изменилось. Если же ждать только уведомления о завершении, можно
пропустить сбой, возникший в процессе. Поэтому автор рекомендует
поручить мониторинг среде выполнения и пробуждать модель лишь при
полезном прогрессе, ошибке или застое. Это позволяет продолжать общение
во время сборки и своевременно вмешиваться при возникновении проблем. В
правилах мониторинга можно настроить фильтрацию, объединение событий и
ограничение частоты, чтобы каждая строка журнала не приводила к
отдельному вызову модели.

Например, подписавшись на обновления состояния ожидающей обработки
заявки, модель может завершить текущую итерацию инференса и дождаться,
пока уведомление с номером заявки не пробудит её. Если другая сторона
предоставляет только интерфейс запросов, модель может самостоятельно
запланировать следующую проверку. Если в назначенный момент выяснится,
что заявка «всё ещё обрабатывается», она назначит ещё одну проверку;
если заявка уже закрыта --- отменит последующие. В обоих случаях за
ожидание отвечает среда выполнения, а модель принимает следующее решение
только после получения уведомления или срабатывания таймера.

\textbf{Каналы внешних событий (\texttt{connect\_channel})} используются
для подписки на электронную почту, мгновенные сообщения или обратные
вызовы API. Адаптер канала проверяет источник, сопоставляет собственные
идентификаторы сервиса с ветвью диалога и связанным объектом, а затем
отправляет событие в едином формате; подробнее это рассматривается далее
в разделе «Уровень приёма». К этой категории относятся уведомления о
состоянии телефонного звонка в PineClaw.

Для всех трёх видов инструментов необходимо явно определять условия
срабатывания, правила фильтрации и полезную нагрузку события. Полезная
нагрузка должна содержать идентификатор события, время его
возникновения, источник, связанную задачу и необходимые данные, чтобы
AI-агент мог определить, остаётся ли событие актуальным и какая
информация потребуется на следующем шаге.

\hypertarget{ux438ux43dux441ux442ux440ux443ux43cux435ux43dux442ux44b-ux432ux437ux430ux438ux43cux43eux434ux435ux439ux441ux442ux432ux438ux44f-ux441-ux43fux43eux43bux44cux437ux43eux432ux430ux442ux435ux43bux435ux43c}{%
\subsection{Инструменты взаимодействия с
пользователем}\label{ux438ux43dux441ux442ux440ux443ux43cux435ux43dux442ux44b-ux432ux437ux430ux438ux43cux43eux434ux435ux439ux441ux442ux432ux438ux44f-ux441-ux43fux43eux43bux44cux437ux43eux432ux430ux442ux435ux43bux435ux43c}}

Вернёмся к задаче изменения бронирования. Задать вопрос службе
поддержки, отправить пользователю сообщение с запросом подтверждения, на
которое можно ответить напрямую, и записать в журнале задачи «ожидается
ответ» --- это три операции, данные о которых сохраняются в разных
местах. Пользователь может временно отвлечься и вернуться, когда
потребуется уточнить условия, подтвердить решение или получить
результат.

В интерфейсе сессии можно показывать окончательный ответ, ход работы
инструментов и результаты выполнения; инструмент взаимодействия, в свою
очередь, явно задаёт получателя, канал, содержимое сообщения и вложения.
Модель решает, что необходимо сообщить, Harness проверяет адресата и
разрешения, а служба сообщений выполняет доставку. Среда выполнения
также может автоматически направить окончательный ответ в исходную
сессию. Эти два способа можно использовать совместно.

Многоканальное взаимодействие OpenClaw позволяет пользователю отправлять
сообщения, ждать новостей о ходе работы и получать напоминания так же,
как при общении с секретарём. Такое «ощущение живого человека» возникает
благодаря непрерывному ведению состояния задачи и своевременной
коммуникации. Пользователь может продолжать задачу в привычном канале, а
агент должен сохранять согласованность идентичности, сессии и привязки
сообщений.

Форма коммуникации также может меняться вместе с задачей: краткие
сведения о ходе работы передаются текстом, к сложным результатам
прикладываются изображения или файлы, запрашиваемая информация
собирается с помощью структурированных карточек, а интерактивные
результаты представляются через генеративный UI, описанный в
\crossreflink{chap:5}{главе 5}. Если канал поддерживает уведомления о
доставке или прочтении, система сохраняет соответствующий статус;
повторные попытки после сбоев, дублирующиеся сообщения и межканальная
синхронизация также заносятся в журнал доставки.

\textbf{Многоканальные уведомления и повторное вовлечение пользователя.}

Мгновенные сообщения, SMS, электронная почта, телефонные звонки и
push-уведомления приложений подходят для разных степеней срочности и
форм содержимого. Агент может выбирать канал с учётом срочности,
пользовательских предпочтений, текущего состояния и длины сообщения.
Событие телефонного разговора, требующее немедленного подтверждения,
сопровождается своевременным напоминанием, а полный отчёт можно
доставить по электронной почте или через ссылку в приложении.

Однажды автор поручил помощнику следить за стоимостью гостиничных
номеров на время конференции ATC. Каждый день в одно и то же время
помощник присылал сообщение о том, что цена ещё не опустилась до
целевого уровня. Однажды цена ненадолго снизилась, но это окно обнаружил
написанный самим автором скрипт. Кратковременное снижение могло
произойти между двумя проверками. Даже если помощник своевременно
обнаружил бы его, из-за поздней доставки сообщения пользователь всё
равно мог упустить возможность забронировать номер. В дни без снижения
цены плановые проверки могут выполняться как обычно, но пользователю
незачем ежедневно читать одно и то же сообщение. В
\crossreflink{chap:7}{главе 7} работа таких уведомлений будет проверена
на примере пятиминутного окна снижения цены.

После завершения длительной задачи уведомление возвращает пользователя
на страницу с результатом; ежедневные сводки и еженедельные отчёты
отправляются с согласованной периодичностью. Многоканальная стратегия
также должна предусматривать дедупликацию, периоды тишины и условия
эскалации --- например, следует ли перейти на другой канал, если на
первое напоминание не последовало ответа. Журнал уведомлений связывается
с состоянием задачи, чтобы пользователь мог продолжить работу
непосредственно из полученного сообщения.

\hypertarget{ux432ux438ux440ux442ux443ux430ux43bux44cux43dux430ux44f-ux438ux434ux435ux43dux442ux438ux447ux43dux43eux441ux442ux44c-ux438-ux438ux437ux43eux43bux438ux440ux43eux432ux430ux43dux43dux430ux44f-ux441ux440ux435ux434ux430-ux432ux44bux43fux43eux43bux43dux435ux43dux438ux44f}{%
\subsection{Виртуальная идентичность и изолированная среда
выполнения}\label{ux432ux438ux440ux442ux443ux430ux43bux44cux43dux430ux44f-ux438ux434ux435ux43dux442ux438ux447ux43dux43eux441ux442ux44c-ux438-ux438ux437ux43eux43bux438ux440ux43eux432ux430ux43dux43dux430ux44f-ux441ux440ux435ux434ux430-ux432ux44bux43fux43eux43bux43dux435ux43dux438ux44f}}

В \crossreflink{chap:4}{главе 4} на примере Саманты из фильма «Она» было
показано, как универсальный помощник с помощью инструментов входит в
цифровой мир. Для реализации такого помощника необходимо чётко
определить, под какой идентичностью он работает и к каким ресурсам имеет
доступ. \textbf{Автор склоняется к тому, чтобы предоставить
универсальному помощнику отдельную виртуальную идентичность.} Выделенные
помощнику служебный почтовый ящик, номер телефона, хранилище и
вычислительная среда позволяют AI-агенту, подобно секретарю с
собственной рабочей учётной записью, выполнять задачи от имени
пользователя под явно обозначенной идентичностью представителя.
Независимая идентичность отделяет повседневную работу помощника от
личных учётных записей пользователя, упрощая отдельную выдачу, аудит и
отзыв разрешений; при ошибочных действиях также легче ограничить масштаб
последствий. Для задач, требующих доступа к учётной записи пользователя,
необходимые разрешения предоставляются с ограниченной областью действия.

От того, кому принадлежит номер телефона, зависит, кто сможет продолжить
работу при обратном звонке. Cue от Manus предоставляет помощнику
отдельный номер, позволяет принимать входящие звонки и оставлять
пользователю краткое содержание разговора. Если организация согласится
использовать оставленный помощником номер для обратной связи, при звонке
сотрудник службы поддержки сможет напрямую связаться с помощником,
который ранее занимался этим вопросом.\footnote{Manus, \emph{Introducing
  Manus 2.0}: отдельный номер телефона Cue, приём входящих звонков и
  краткое содержание разговоров.}

Например, для получения скидки или отмены комиссии помощник может
связаться с организацией как уполномоченный представитель. Если
сотрудник службы поддержки потребует устного подтверждения личности,
пользователя можно подключить к трёхстороннему звонку, а после проверки
помощник продолжит разговор. Если организация собирается перезвонить
пользователю, ему заранее передаются номер обращения, уже согласованные
условия и вопросы, которые нужно уточнить. Отчитываясь перед
пользователем, помощник также должен ясно указывать имеющиеся
доказательства: если сотрудник службы поддержки устно согласился
отменить комиссию, следует зафиксировать именно это обещание; сообщать о
фактической отмене комиссии можно лишь после последующей проверки
обновлённого счёта.

Когда автор, находясь в Пекине, перевыпускал кредитную карту
американского банка, Pine связался с банком, организовал срочный
перевыпуск карты и добился отмены комиссии за срочность в размере 15
долларов. По окончании разговора помощник уже мог сообщить, какое
обращение принял банк и какую именно комиссию согласился отменить. Если
через несколько дней пользователь спросит: «Где моя новая карта?»,
потребуется дополнительно проверить доставку или попросить получателя
подтвердить её получение.

Виртуальный компьютер и виртуальный телефон предоставляют постоянно
работающую операционную среду. На виртуальной машине можно запускать
настольную операционную систему, контейнеры подходят для изоляции
процессов и зависимостей, а эмулятор Android служит средой для мобильных
приложений. Они поддерживают фоновую работу и уменьшают влияние
автоматических операций на текущую работу пользователя. Границы изоляции
совместно определяются настройками сети, подключённых томов, учётных
данных и разрешений среды выполнения; эти настройки также влияют на то,
к каким данным и внешним системам может обращаться агент.

Предположим, вы хотите, чтобы корпоративная версия Tencent WorkBuddy
создала еженедельный отчёт после того, как вы закончите работу и
выключите компьютер, но таблица продаж за текущую неделю хранится только
на рабочем компьютере. Если выбрать локальную задачу, помощник сможет
прочитать таблицу, однако компьютер должен оставаться включённым, а
клиент --- подключённым к сети. Если выбрать облачную задачу, помощник
сможет продолжить работу после выключения компьютера при условии, что
разрешённые для использования в отчёте данные заранее загружены в задачу
или помещены в доступное ей облачное хранилище материалов. Если таблица
продаж должна оставаться на локальном компьютере, её обработку следует
запланировать на время, когда он подключён к сети.

Веб-сайты и мобильные приложения также могут требовать аутентификацию,
проверочный код или ручное подтверждение. При подключении следует в
первую очередь использовать предоставляемые сервисом API и способы
авторизации. Если требуется аутентификация через графический интерфейс,
пользователь может лично войти в учётную запись через удалённый рабочий
стол VNC/RDP, а автоматические операции на время аутентификации
приостанавливаются. После завершения входа система проверяет состояние
сессии и возобновляет задачу. Токены и сессии сохраняются в пределах
срока действия, а после отзыва разрешения или завершения задачи
удаляются согласно установленной политике.

Агент и виртуальная среда могут обмениваться данными через \textbf{общую
файловую систему}. Например, пользователь загружает CSV-файл,
виртуальный компьютер читает его, выполняет анализ и записывает графики
в общее хранилище; среда выполнения передаёт агенту расположение файлов
с графиками, после чего тот предоставляет пользователю доступную ссылку
для скачивания. Содержимое больших файлов остаётся в хранилище, а модель
при необходимости читает краткие сводки или отдельные фрагменты.

\hypertarget{ux443ux440ux43eux432ux435ux43dux44c-ux43fux440ux438ux451ux43cux430-ux440ux430ux437ux43dux44bux435-ux438ux441ux442ux43eux447ux43dux438ux43aux438-ux435ux434ux438ux43dux43eux435-ux441ux43eux431ux44bux442ux438ux435}{%
\subsection{Уровень приёма: разные источники, единое
событие}\label{ux443ux440ux43eux432ux435ux43dux44c-ux43fux440ux438ux451ux43cux430-ux440ux430ux437ux43dux44bux435-ux438ux441ux442ux43eux447ux43dux438ux43aux438-ux435ux434ux438ux43dux43eux435-ux441ux43eux431ux44bux442ux438ux435}}

AI-агент службы поддержки интернет-магазина обрабатывает возврат
средств. Покупатель подаёт заявку в чате поддержки в App. Сумма
превышает полномочия агента, поэтому он просит руководителя подтвердить
операцию в группе послепродажного обслуживания корпоративного WeChat.
После успешного возврата платёжный сервис отправляет Webhook. Затем
система заказов публикует изменение состояния в Kafka. Все четыре
сообщения относятся к одному заказу, но поступают разными способами. Чем
больше источников, тем сильнее среда выполнения распадается на отдельные
тракты, каждый из которых по-своему обрабатывает повторные попытки,
дедупликацию и порядок. \textbf{Уровень приёма унифицирует форму событий
после их поступления в систему, сохраняя различия возможностей
источников до момента поступления.} Такая многоуровневая архитектура
показана на \crossreflink{fig:6-2}{рисунке 6-2}.

\begin{figure}
\hypertarget{fig:6-2}{%
\centering
\includegraphics{images/revised/fig6-2.pdf}
\caption{Рисунок 6-2. Уровень приёма: адаптеры, унифицированный конверт
и журнал приёма}\label{fig:6-2}
}
\end{figure}

Распространённые источники делятся на четыре категории.
\textbf{Диалоговые каналы} включают собственный диалоговый интерфейс
продукта, системы мгновенного обмена сообщениями --- корпоративный
WeChat, Feishu, Slack и другие, --- а также электронную почту; агент
может ответить по тому же каналу. \textbf{Обратные вызовы} --- это
HTTP-интерфейсы, вызываемые внешними сервисами при изменении состояния.
Обычно они подписаны и доставляют только уведомление. После ошибки
доставки некоторые сервисы автоматически повторяют попытку, а в GitHub
доставку приходится повторять вручную или через API. \textbf{Шины
сообщений}, такие как SQS и Kafka, переносят бизнес-события внутренних
сервисов и обладают явно определённой семантикой подтверждения и
хранения. \textbf{Таймеры, опросы и внутренние события} создаются самой
системой, и их достаточно записывать в тот же журнал. В таблице ниже
сравниваются первые три категории.\footnote{Amazon Web Services,
  ``Amazon SQS standard queues''; ``Amazon SQS visibility timeout'';
  ``Using the message deduplication ID in Amazon SQS''; ``Using the
  message group ID with Amazon SQS FIFO Queues''. Документация Apache
  Kafka 4.4, разделы ``Consumer Position'' и ``Message Delivery
  Semantics'' главы ``Design''. Дата проверки: 2026-11-24.}\footnote{Открытая
  платформа Feishu, ``Обзор событий''; ``Получение истории сообщений
  диалога''. Центр разработчиков корпоративного WeChat, ``Обзор'' в
  разделе ``Получение сообщений и событий''. GitHub Docs, ``Best
  practices for using webhooks''; ``Handling failed webhook
  deliveries''. Дата проверки: 2026-11-24.}

\begin{longtable}[]{@{}
  >{\raggedright\arraybackslash}p{(\columnwidth - 6\tabcolsep) * \real{0.2471}}
  >{\raggedright\arraybackslash}p{(\columnwidth - 6\tabcolsep) * \real{0.2331}}
  >{\raggedright\arraybackslash}p{(\columnwidth - 6\tabcolsep) * \real{0.2705}}
  >{\raggedright\arraybackslash}p{(\columnwidth - 6\tabcolsep) * \real{0.2493}}@{}}
\toprule\noalign{}
\begin{minipage}[b]{\linewidth}\raggedright
\end{minipage} & \begin{minipage}[b]{\linewidth}\raggedright
\hspace{0pt}Диалог в App
\end{minipage} & \begin{minipage}[b]{\linewidth}\raggedright
\hspace{0pt}Корпоративный WeChat /\allowbreak{} Feishu
\end{minipage} & \begin{minipage}[b]{\linewidth}\raggedright
\hspace{0pt}Электронная почта
\end{minipage} \\
\midrule\noalign{}
\endhead
\bottomrule\noalign{}
\endlastfoot
\hspace{0pt}Способ поступления & \hspace{0pt}Долговременное соединение &
\hspace{0pt}Обратный вызов или долговременное соединение &
\hspace{0pt}Push-\allowbreak{}\hspace{0pt}доставка или опрос \\
\hspace{0pt}Область упорядочивания & \hspace{0pt}В пределах диалога &
\hspace{0pt}Для большинства событий не гарантируется & \hspace{0pt}Не
гарантируется \\
\hspace{0pt}Встроенное основание дедупликации & \hspace{0pt}ID сообщения
клиента & \hspace{0pt}ID сообщения, ID события &
\hspace{0pt}Message-\allowbreak{}\hspace{0pt}ID \\
\hspace{0pt}Возможность повторного воспроизведения & \hspace{0pt}Зависит
от собственных записей & \hspace{0pt}Сам обратный вызов воспроизвести
нельзя; в Feishu сообщения можно повторно получить через интерфейс
истории & \hspace{0pt}Почтовый ящик можно прочитать повторно \\
\hspace{0pt}Способ подтверждения & \hspace{0pt}Квитанция клиента &
\hspace{0pt}Успешный ответ в установленный срок &
\hspace{0pt}Самостоятельный учёт позиции чтения \\
\hspace{0pt}Канал ответа & \hspace{0pt}Исходный диалог &
\hspace{0pt}Исходный диалог & \hspace{0pt}Ответное письмо \\
\end{longtable}

\begin{longtable}[]{@{}
  >{\raggedright\arraybackslash}p{(\columnwidth - 6\tabcolsep) * \real{0.2471}}
  >{\raggedright\arraybackslash}p{(\columnwidth - 6\tabcolsep) * \real{0.2098}}
  >{\raggedright\arraybackslash}p{(\columnwidth - 6\tabcolsep) * \real{0.3364}}
  >{\raggedright\arraybackslash}p{(\columnwidth - 6\tabcolsep) * \real{0.2066}}@{}}
\toprule\noalign{}
\begin{minipage}[b]{\linewidth}\raggedright
\end{minipage} & \begin{minipage}[b]{\linewidth}\raggedright
\hspace{0pt}Webhook
\end{minipage} & \begin{minipage}[b]{\linewidth}\raggedright
\hspace{0pt}SQS
\end{minipage} & \begin{minipage}[b]{\linewidth}\raggedright
\hspace{0pt}Kafka
\end{minipage} \\
\midrule\noalign{}
\endhead
\bottomrule\noalign{}
\endlastfoot
\hspace{0pt}Способ поступления & \hspace{0pt}Обратный вызов &
\hspace{0pt}Получение потребителем & \hspace{0pt}Получение
потребителем \\
\hspace{0pt}Область упорядочивания & \hspace{0pt}Не гарантируется &
\hspace{0pt}Для стандартной очереди не гарантируется, для FIFO --- в
пределах группы сообщений & \hspace{0pt}В пределах раздела \\
\hspace{0pt}Встроенное основание дедупликации & \hspace{0pt}ID доставки
& \hspace{0pt}ID сообщения; у FIFO также есть
\texttt{Message\allowbreak{}Deduplication\allowbreak{}Id} & \hspace{0pt}Тема, раздел и смещение \\
\hspace{0pt}Возможность повторного воспроизведения & \hspace{0pt}Зависит
от сервиса & \hspace{0pt}После удаления невозможно &
\hspace{0pt}Возможно в течение срока хранения \\
\hspace{0pt}Способ подтверждения & \hspace{0pt}Ответ 2xx в установленный
срок & \hspace{0pt}Удаление сообщения & \hspace{0pt}Фиксация смещения \\
\hspace{0pt}Канал ответа & \hspace{0pt}Нет & \hspace{0pt}Нет &
\hspace{0pt}Нет \\
\end{longtable}

\textbf{Унифицированный конверт.}

Адаптер преобразует исходное сообщение в \textbf{конверт} (envelope)
унифицированного формата. В конверте записываются только факты, которые
можно определить в момент поступления. Приоритет и срок обработки
добавляются правилами маршрутизации, уровень клиента и состояние заказа
--- бизнес-системой, а состояние обработки хранится отдельно. Заявку
покупателя на возврат средств можно представить так:

\begin{Shaded}
\begin{Highlighting}[]
\FunctionTok{\{}
  \DataTypeTok{"id"}\FunctionTok{:} \StringTok{"evt\_7f3a"}\FunctionTok{,}
  \DataTypeTok{"source"}\FunctionTok{:} \FunctionTok{\{}\DataTypeTok{"type"}\FunctionTok{:} \StringTok{"app\_chat"}\FunctionTok{,} \DataTypeTok{"sender"}\FunctionTok{:} \StringTok{"user:u\_5521"}\FunctionTok{,} \DataTypeTok{"verified"}\FunctionTok{:} \KeywordTok{true}\FunctionTok{\},}
  \DataTypeTok{"dedupe\_key"}\FunctionTok{:} \StringTok{"app:msg:msg\_1029"}\FunctionTok{,}
  \DataTypeTok{"thread"}\FunctionTok{:} \StringTok{"app:conv:conv\_88"}\FunctionTok{,}
  \DataTypeTok{"kind"}\FunctionTok{:} \StringTok{"message"}\FunctionTok{,}
  \DataTypeTok{"correlation"}\FunctionTok{:} \OtherTok{[}\StringTok{"order:12345"}\OtherTok{]}\FunctionTok{,}
  \DataTypeTok{"content"}\FunctionTok{:} \FunctionTok{\{}\DataTypeTok{"type"}\FunctionTok{:} \StringTok{"text"}\FunctionTok{,} \DataTypeTok{"text"}\FunctionTok{:} \StringTok{"Я хочу вернуть деньги за этот заказ"}\FunctionTok{\},}
  \DataTypeTok{"raw\_ref"}\FunctionTok{:} \StringTok{"raw/2026{-}11{-}23/app/msg\_1029.json"}\FunctionTok{,}
  \DataTypeTok{"reply"}\FunctionTok{:} \FunctionTok{\{}\DataTypeTok{"via"}\FunctionTok{:} \StringTok{"app\_chat"}\FunctionTok{,} \DataTypeTok{"target"}\FunctionTok{:} \StringTok{"app:conv:conv\_88"}\FunctionTok{\},}
  \DataTypeTok{"occurred\_at"}\FunctionTok{:} \StringTok{"2026{-}11{-}23T14:02:09+08:00"}\FunctionTok{,}
  \DataTypeTok{"expires\_at"}\FunctionTok{:} \KeywordTok{null}
\FunctionTok{\}}
\end{Highlighting}
\end{Shaded}

\begin{itemize}
\tightlist
\item
  \texttt{source}: тип источника и проверенный отправитель. Признак
  проверки устанавливается адаптером на основании подписи, токена или
  учётной записи канала; самопредставление в тексте сообщения проверкой
  не считается. Среда выполнения определяет полномочия по этому полю.
\item
  \texttt{dedupe\_key}: ключ дедупликации; критерий повтора определяется
  семантикой источника. ID доставки, ID сообщения либо сочетание раздела
  и смещения обозначают одну и ту же доставку и блокируют повторную
  отправку на транспортном уровне. Если бизнес-сообщение содержит ключ
  идемпотентности (idempotency key), например номер заявки на возврат
  вместе с состоянием, адаптер должен извлечь его, чтобы блокировать
  повторы на бизнес-уровне. Для пользовательского диалога действует
  обратное правило: если покупатель дважды говорит «верните деньги», это
  два разных сообщения, поэтому дедупликация выполняется только по ID
  сообщения клиента.
\item
  \texttt{raw\_ref}: исходная полезная нагрузка, сохранённая без
  изменений. Только по ней можно расследовать ошибки адаптера.
\item
  \texttt{reply}: исходящий адаптер и целевой адрес ответа. У обратного
  вызова платёжного сервиса нет канала ответа, поэтому это поле пусто.
\item
  \texttt{occurred\_at} --- указанное источником время возникновения
  события; \texttt{expires\_at} --- срок действия, заполняемый только
  для таких событий, как ценовые предложения и коды подтверждения.
\end{itemize}

\textbf{Потоки и связи.}

Покупатель сначала говорит: «Верните деньги», а сразу затем добавляет:
«Нет, лучше обменяйте товар». Порядок этих сообщений нельзя менять.
Сообщения другого покупателя с ними не связаны и могут обрабатываться
параллельно. \textbf{Уровень приёма гарантирует порядок только в
пределах потока.} Адаптер сопоставляет собственную единицу
упорядочивания источника с полем \texttt{thread}: диалог в App, группа
сообщений SQS FIFO (\texttt{MessageGroupId}), ключ раздела Kafka, номер
Issue или PR в обратном вызове GitHub. Если источник не гарантирует
порядок, среда выполнения выявляет опоздавшее старое состояние по
\texttt{occurred\_at} или предоставленному источником номеру версии ---
например, когда событие «возврат обрабатывается» приходит после события
«возврат выполнен».

Четыре сообщения об этом возврате относятся к четырём разным потокам.
Объединять их в один поток затратно: не связанное с возвратом обсуждение
в группе послепродажного обслуживания тоже пришлось бы поставить после
сообщения покупателя. Кроме того, идентификатор потока часто
одновременно служит ключом раздела, поэтому его последующее изменение
потребует миграции данных. \textbf{Связи выражаются отдельным полем без
изменения потоков.} В \texttt{correlation} перечисляются бизнес-объекты,
которых касается событие, а задача подписывается на связанные с этими
объектами потоки.

\textbf{Разделение подтверждения и прогресса.}

\textbf{Адаптер должен подтверждать источнику получение после записи
события в журнал приёма, а не после завершения обработки агентом.}
Обратный вызов событий Feishu требует ответа в течение 3 секунд,
корпоративный WeChat --- в течение 5 секунд; при превышении срока
событие отправляется повторно. Если сообщение SQS не удалено до
истечения тайм-аута видимости, оно снова становится доступным для
получения. Ожидание завершения обработки перед подтверждением приведёт к
повторной доставке при любом достаточно долгом вызове модели. Если
сначала записывать событие в журнал и только затем подтверждать его,
гарантии доставки разных источников сводятся к одному правилу:
\textbf{каждое событие попадает в журнал приёма не менее одного раза
(at-least-once)}, а повторная доставка блокируется ключом дедупликации.

Журнал приёма играет роль папки входящих сообщений (inbox), следует
изложенному в первой главе принципу записи только добавлением и
разделяется по потокам. Если журнал приёма реализован на Kafka, для
дедупликации необходимо отдельное хранилище с ограничением уникальности
либо фильтрация на стороне потребителя по ключу дедупликации. Kafka не
умеет отклонять запись по заданному приложением ключу: идемпотентный
производитель защищает только от повторных попыток самого производителя.

Система должна раздельно хранить три вида прогресса:

\begin{longtable}[]{@{}
  >{\raggedright\arraybackslash}p{(\columnwidth - 4\tabcolsep) * \real{0.3502}}
  >{\raggedright\arraybackslash}p{(\columnwidth - 4\tabcolsep) * \real{0.3527}}
  >{\raggedright\arraybackslash}p{(\columnwidth - 4\tabcolsep) * \real{0.2971}}@{}}
\toprule\noalign{}
\begin{minipage}[b]{\linewidth}\raggedright
\hspace{0pt}Вид прогресса
\end{minipage} & \begin{minipage}[b]{\linewidth}\raggedright
\hspace{0pt}Что записывается
\end{minipage} & \begin{minipage}[b]{\linewidth}\raggedright
\hspace{0pt}Кто ведёт
\end{minipage} \\
\midrule\noalign{}
\endhead
\bottomrule\noalign{}
\endlastfoot
\hspace{0pt}Прогресс источника & \hspace{0pt}Смещение Kafka, дескриптор
получения SQS, позиция чтения почтового ящика & \hspace{0pt}Адаптер \\
\hspace{0pt}Прогресс диспетчеризации & \hspace{0pt}Каким задачам уже
переданы события из журнала & \hspace{0pt}Маршрутизатор \\
\hspace{0pt}Состояние задачи & \hspace{0pt}До какого шага дошла
обработка, отправлен ли ответ & \hspace{0pt}Среда выполнения \\
\end{longtable}

Если считать смещение прогрессом задачи, при повторном воспроизведении
журнала уже отправленные ответы будут отправлены снова. После разделения
повторное воспроизведение затрагивает только прогресс диспетчеризации.
Исходящему направлению соответствует папка исходящих сообщений (outbox):
перед отправкой записываются намерение отправки и ключ идемпотентности,
а после получения квитанции канала операция помечается завершённой. Как
выполнять сверку после перезапуска, рассматривается далее в разделе
«Построение базового цикла на основе почтовых событий».

\textbf{Срок действия, объединение и обратное давление.}

Если в один поток поступают сотни событий в секунду, не следует
передавать модели каждое по отдельности. Перед диспетчеризацией
маршрутизатор обрабатывает накопившиеся события:

\begin{itemize}
\tightlist
\item
  \textbf{Срок действия.} Ценовые предложения, коды подтверждения и уже
  снятые предупреждения после определённого момента теряют смысл.
  События с истёкшим \texttt{expires\_at} только записываются в
  траекторию и больше не активируют модель.
\item
  \textbf{Объединение.} \texttt{kind} различает два типа событий.
  Событие состояния описывает текущее значение некоторого объекта,
  например состояние возврата в обратном вызове платёжного сервиса или
  состояние заказа в Kafka. Если в одном потоке накопилось несколько
  таких событий, модели передаётся только последнее, а промежуточные
  значения сохраняются в журнале. Событие-сообщение содержит
  произнесённую кем-то фразу или отданную команду, например сообщение
  покупателя или нажатие руководителя. Такие события можно передавать
  модели пакетами, но нельзя отбрасывать.
\item
  \textbf{Обратное давление.} Источники, из которых данные извлекаются,
  можно опрашивать медленнее, оставляя накопившиеся события в источнике.
  Источники с push-доставкой нельзя попросить приостановиться, поэтому
  адаптер продолжает записывать события в журнал и подтверждать их, а
  очередь переносится в журнал приёма и обрабатывается по двум
  предыдущим правилам.
\end{itemize}

\textbf{Один возврат проходит через четыре источника.}

\begin{enumerate}
\def\labelenumi{\arabic{enumi}.}
\tightlist
\item
  Покупатель запрашивает возврат средств в диалоге App. Адаптер
  выполняет дедупликацию по ID сообщения клиента, записывает событие в
  журнал и отправляет клиенту квитанцию. Маршрутизатор создаёт по
  \texttt{order:12345} задачу возврата, которая подписывается на поток
  этого диалога.
\item
  Сумма превышает полномочия, поэтому агент отправляет в группу
  послепродажного обслуживания корпоративного WeChat карточку
  подтверждения, а идентификатор карточки записывается в задаче.
  Руководитель нажимает «Одобрить». Адаптер записывает событие в журнал,
  возвращает успешный ответ и добавляет связанный объект по
  идентификатору карточки.
\item
  После вызова агентом интерфейса возврата платёжный сервис дважды
  отправляет обратный вызов «возврат выполнен» с одинаковым ID доставки.
  Вторая доставка отбрасывается при дедупликации.
\item
  Система заказов публикует в Kafka событие «средства за заказ
  возвращены», описывающее то же самое, что и обратный вызов. Если оба
  адаптера извлекают из полезной нагрузки номер заявки на возврат вместе
  с состоянием и используют их как ключ дедупликации, событие
  блокируется уже на этапе приёма. Если один адаптер выполняет
  дедупликацию по ID доставки, а другой --- по разделу и смещению, оба
  события попадут в журнал. Среда выполнения обнаружит, что возврат уже
  подтверждён в состоянии задачи, и только запишет второе событие в
  траекторию.
\item
  Ни у обратного вызова, ни у события Kafka нет канала ответа. Среда
  выполнения использует \texttt{reply} потока диалога из шага 1 и
  уведомляет покупателя через исходящий адаптер App.
\end{enumerate}

Шаг 4 показывает границу дедупликации на уровне приёма: какие повторы
удастся заблокировать, зависит от ключа дедупликации, объявленного
адаптером. Остальные повторы определяются по состоянию задачи. После
диспетчеризации события задаче выбор между продолжением существующей
сессии и созданием новой относится к управлению контекстом. Информация
между сессиями передаётся через память, описанную в третьей главе, а
возможность совместного использования контекста рассматривается в
десятой главе.

\hypertarget{ux43cux435ux445ux430ux43dux438ux437ux43c-ux43eux431ux440ux430ux431ux43eux442ux43aux438-ux441ux43eux431ux44bux442ux438ux439}{%
\subsection{Механизм обработки
событий}\label{ux43cux435ux445ux430ux43dux438ux437ux43c-ux43eux431ux440ux430ux431ux43eux442ux43aux438-ux441ux43eux431ux44bux442ux438ux439}}

Пользовательские сообщения, результаты инструментов, таймеры и запросы
на совместную работу могут одновременно поступать одному агенту. Среда
выполнения должна сохранять эти события, определять порядок их обработки
и затем передавать модели соответствующую информацию. \textbf{Цикл
событий} (event loop) непрерывно продвигает этот процесс: извлекает
события из очереди, обновляет состояние задачи, формирует запросы к
модели, планирует выполнение инструментов, а затем ожидает новый ввод
или уведомление о завершении. Аналогичный подход применяется в
программах на Go, где несколько потоков сообщений организуются с помощью
goroutine, channel и \texttt{select}.

Рассмотрим прототип dsh-superagent. В одном проекте GitHub обсуждение в
Issue и обсуждение в запросе на включение изменений (PR) относятся к
двум разным ветвям диалога (Thread). Когда системе нужно решить, следует
ли принять новую работу, она считывает состояние ветви и передаёт его
модели для оценивания (Eval). В одной ветви одновременно может
выполняться не более одного оценивания, однако на неё могут быть
подписаны несколько заявок (Ticket). Предположим, что заявка на
исправление дефекта и заявка на дополнение документации подписаны на
один и тот же PR. Если ревьюер оставит комментарий в PR, обе заявки
получат его в качестве контекста. Если затем кто-то дополнит Issue
условиями воспроизведения, сообщение будет доставлено заявкам,
подписанным на этот Issue. Получатели сообщения определяются ветвью и
связями подписки. Доступные после получения сообщения действия
по-прежнему ограничиваются исходными полномочиями каждой
заявки.\footnote{DeepSeek, доменная модель прототипа dsh-superagent, а
  также правила авторизации и подписки.}

В dsh-superagent взаимодействие с людьми и непосредственное выполнение
работы также возложены на разные сессии. Основная сессия (Main Session)
только принимает заявки, отвечает и фиксирует полномочия: она не читает
код и не останавливается в ожидании результата. У каждой заявки есть
собственная исполнительная сессия (Executor Session), которая изменяет
код и запускает тесты в отдельном рабочем каталоге. Только когда
исполнительная сессия завершает исправление, оказывается заблокированной
или должна задать человеку вопрос, она отправляет основной сессии
сообщение со сводкой и адресом для чтения --- то есть использует
описанный в \crossreflink{chap:2}{главе 2} принцип «сводка плюс ID
объекта». Благодаря этому основная сессия всегда может принять следующий
запрос, а вопросы человеку, уведомления и остановка также проходят через
неё, поэтому обязательства выдаются только через одну точку. Быстрая и
медленная модели, рассматриваемые далее в разделе о голосовом
взаимодействии, представляют собой тот же принцип разделения
обязанностей применительно к разговору в реальном времени.\footnote{DeepSeek,
  прототип dsh-superagent, версия \texttt{8c32a9d}. Обязанности двух
  типов сессий описаны в \texttt{docs/004-prompts.md}, сессия выполнения
  и рабочий каталог --- в \texttt{docs/003-execution.md}, а формат
  сообщений, передаваемых в основную сессию, --- в
  \texttt{docs/007-architecture.md}.}

Синхронный интерфейс модели обычно принимает следующую группу сообщений
после завершения одной генерации или получения всех результатов
инструментов. Здесь можно установить \textbf{безопасную точку}: прежде
чем перейти к новому шагу обработки, среда выполнения записывает
завершённые сообщения, всё ещё выполняющиеся вызовы и текущий прогресс
задачи. Нативный асинхронный интерфейс, напротив, предоставляет протокол
для получения обновлений во время выполнения, а конкретный момент
продолжения совместно определяется интерфейсом и средой выполнения.

Для отмены также необходимо содействие исполнителя. После остановки
генерации модели фоновые инструменты могут продолжать работать;
исполнитель должен получить сигнал отмены и сообщить фактическое
состояние, подобно тому как задача в Go проверяет \texttt{ctx.Done()} и
затем выполняет очистку. Уже совершённая запись, отправка или оплата
должна отменяться либо компенсироваться соответствующим
бизнес-процессом.

\textbf{Три стратегии обработки.}

На \crossreflink{fig:6-3}{рисунке 6-3} сравниваются три варианта
действий при поступлении нового события. Если поставить его в очередь,
время ответа будет зависеть от того, когда встретится следующая
безопасная точка. При отмене текущей задачи необходимо выяснить, какие
операции уже завершены и откуда можно возобновить работу. При запуске
отдельного цикла следует проверить, не будут ли две задачи конкурировать
за один ресурс.

\begin{figure}
\hypertarget{fig:6-3}{%
\centering
\includegraphics{images/revised/fig6-3.pdf}
\caption{Рисунок 6-3. Три стратегии обработки асинхронных
событий}\label{fig:6-3}
}
\end{figure}

\textbf{Обработка с отменой (Cancellation-Based)} подходит для
пользовательских команд остановки и управляющих событий, требующих
немедленной обработки. Сначала среда выполнения прекращает дальнейшую
диспетчеризацию связанных операций, отправляет сигнал остановки
выполняющему генерацию запросу и допускающим отмену инструментам, а
затем сохраняет завершённые сообщения, состояния инструментов и
подтверждения отмены. После этого новое событие включается в контекст
принятия решения. Например, если пользователь говорит: «Стой! Я ошибся»,
система сначала ограничивает дальнейшее выполнение, а затем повторно
уточняет задачу с учётом фактического текущего состояния. Необработанные
элементы очереди сохраняют свои идентификаторы и состояния и
переупорядочиваются с учётом области их связи.

\textbf{Обработка через очередь (Queued)} подходит для обычных
дополнений и обновлений прогресса. Событие поступает в очередь и в
следующей безопасной точке передаётся модели вместе с результатами
инструментов. Например, если во время поиска пользователь уточняет:
«Учитывай результаты только за последний месяц», то после завершения
поиска модель одновременно читает результаты и новое условие, а затем
решает, отфильтровать их или выполнить поиск заново.

\textbf{Параллельная обработка (Parallel)} подходит для независимых от
основной задачи запросов, требующих быстрого ответа. Например, пока
выполняется анализ большого объёма данных, пользователь может спросить о
погоде, и система создаст для этого отдельную сессию.

\textbf{Срочность и маршрутизация.}

Тип события даёт первоначальную подсказку: \texttt{user.interrupt},
доверенные инструкции надзора и системные тревоги могут направляться в
высокоприоритетный канал; для \texttt{user.input}, \texttt{agent.input},
\texttt{tool.result} и \texttt{timer.trigger} решение принимается также
с учётом содержимого и состояния задачи. Даже если отправитель называет
сообщение «срочным» в его тексте, система определяет полномочия по
проверенной идентичности и установленным правилам.

Для маршрутизации, требующей понимания смысла сообщения, автор
рекомендует использовать лёгкую классификационную модель: «немедленно
остановись» соответствует отмене, «какая сегодня погода?» может
подходить для параллельной обработки, а «отчёт должен быть на китайском
языке» обычно попадает в очередь дополнений к текущей задаче. Явный
механизм остановки обрабатывается непосредственно средой выполнения, а
семантическая классификация применяется к остальным сообщениям, которые
необходимо интерпретировать.

Универсальная модель тоже может присваивать сообщениям метки, однако
даже для вывода одной лишь метки «требует обработки» ей необходимо
прочитать описание задачи и содержимое сообщения. Если выполнять такую
оценку для каждого нового сообщения, стоимость вызовов будет расти
вместе с их количеством, а действительно срочному сообщению придётся
ждать, пока модель выдаст метку. Jev от TypeSafe специально предназначен
для подобных задач принятия решения с заранее заданными вариантами.
Получив сообщение и текущее состояние задачи, он может с меньшими
затратами и задержкой параллельно вернуть вероятности для каждого
варианта.\footnote{В примечаниях к выпуску Jev компания TypeSafe описала
  заранее определённые варианты вывода, параллельный возврат
  вероятностей и метод обучения RLCD, нацеленный на калибровку решений.}

Рассмотренное в \crossreflink{chap:1}{главе 1} разрежение пространства
наблюдений здесь определяет, какие сообщения запускают новый инференс.
Например, для сообщений Feishu и комментариев к GitHub Issue программа
может сначала отфильтровать повторные уведомления по номеру события, а
затем поручить Jev определить, дополняет ли новое сообщение существующую
задачу или содержит новый запрос, требующий основной модели. Прочитав
фразу «дополнение к описанным вчера шагам воспроизведения», модель
должна лишь определить, что это дополнительное пояснение. К какому
именно Issue оно относится, программа может узнать непосредственно из
идентификатора, которым снабжён комментарий.

Даже высокая вероятность, выданная моделью, может отражать завышенную
уверенность. Например, все 100 сообщений могут быть отнесены к категории
«дополнение существующей задачи» с вероятностью около 90\%, хотя в
действительности к ней относятся лишь 60 сообщений. Если автоматически
добавить всю эту группу к существующим задачам, исходное назначение
остальных 40 сообщений будет упущено.

\hypertarget{ux43fux43eux441ux442ux440ux43eux435ux43dux438ux435-ux431ux430ux437ux43eux432ux43eux433ux43e-ux446ux438ux43aux43bux430-ux43dux430-ux43eux441ux43dux43eux432ux435-ux43fux43eux447ux442ux43eux432ux44bux445-ux441ux43eux431ux44bux442ux438ux439}{%
\subsection{Построение базового цикла на основе почтовых
событий}\label{ux43fux43eux441ux442ux440ux43eux435ux43dux438ux435-ux431ux430ux437ux43eux432ux43eux433ux43e-ux446ux438ux43aux43bux430-ux43dux430-ux43eux441ux43dux43eux432ux435-ux43fux43eux447ux442ux43eux432ux44bux445-ux441ux43eux431ux44bux442ux438ux439}}

Предположим, что процесс обработки новых писем внезапно перезапустился.
Идентификатор письма поможет найти исходный текст, но не покажет, был ли
отправлен ответ на предыдущее письмо. Лучше регистрировать каждую
операцию записи и полученную квитанцию под идентификатором
соответствующего письма --- от постановки в очередь, чтения письма и
вызова инструментов до обновления почтового ящика или уведомления
пользователя. Если после перезапуска обнаружится запрос без квитанции,
сначала следует проверить фактическое состояние почтового ящика и лишь
затем решать, нужно ли отправлять запрос повторно.

\begin{quote}
\textbf{Эксперимент 6-1 ★★★: событийная обработка почты с помощью
AI-агента}

В этом эксперименте тестовый почтовый ящик служит точкой входа для
отработки регистрации событий, работы с очередью и обработки задач.
Сначала новые письма преобразуются в структурированные события и
обрабатываются в порядке поступления, затем подключаются инструменты для
работы с историей переписки, чтения вложений, поиска знаний и отправки
уведомлений. Проследите по \crossreflink{fig:6-4}{рисунку 6-4} обработку
одного приглашения на встречу: найдите в цикле событий этапы проверки
конфликтов в календаре и составления черновика, а затем найдите в
постоянном хранилище запись об обработке этого письма.

\begin{figure}
\hypertarget{fig:6-4}{%
\centering
\includegraphics{images/revised/fig6-4.pdf}
\caption{Рисунок 6-4. Архитектура событийного агента в эксперименте
6-1}\label{fig:6-4}
}
\end{figure}

Одну очередь можно расширить для поддержки разных источников:
поступления электронной почты, сообщений в IM или SMS, обновлений GitHub
PR и Issue, таймеров, внешних Webhook, а также внутренних событий, таких
как переход пользователя в автономный режим, истечение времени ожидания
процесса и предупреждения о ресурсах. Каждый адаптер проверяет источник
и дополняет событие связанной информацией, а результат обработки
записывается по идентификатору события.

Для приёмки используются три тестовых письма: приглашение на встречу
должно запускать проверку конфликтов в календаре и создавать черновик
ответа с принятием или отклонением; жалоба клиента --- извлечение
ключевой информации, установку высокого приоритета и уведомление
пользователя; рекламное письмо --- архивирование по заданным правилам.

На текущем этапе в сопутствующем журнале зафиксированы проверка
документации интерфейсов и исследование аутентификации; для подключения
тестового почтового ящика по-прежнему требуются действительные учётные
данные.

\textbf{Расширение (необязательно): подключение нескольких источников и
внедрение сбоев.} Помимо тестового почтового ящика подключите локальную
конечную точку Webhook и очередь SQS, которую можно эмулировать
локально. Каждый адаптер должен проверить источник, сформировать
конверт, записать его в журнал приёма и лишь затем подтвердить получение
источнику. Среда выполнения должна читать события из журнала и
обрабатывать их по порядку внутри каждой цепочки. Затем внедрите три
типа сбоев:

\begin{enumerate}
\def\labelenumi{\arabic{enumi}.}
\tightlist
\item
  \textbf{Повторная доставка.} Конечная точка Webhook дважды получает
  запрос с одним идентификатором доставки, но в журнале появляется
  только одно событие.
\item
  \textbf{Сбой до подтверждения.} Адаптер принудительно завершается
  после записи в журнал, но до удаления сообщения SQS. После того как
  сообщение снова становится видимым и считывается повторно, в журнале
  по-прежнему остаётся только одно событие.
\item
  \textbf{Сбой до записи квитанции.} Среда выполнения принудительно
  завершается после отправки ответа, но до записи квитанции, а после
  перезапуска воспроизводит события из журнала. Если канал поддерживает
  ключ идемпотентности либо позволяет найти уже отправленное сообщение,
  например в папке отправленных писем, один и тот же ответ не
  отправляется повторно. Если канал не предоставляет ни одной из этих
  возможностей, ответ можно лишь пометить как ожидающий ручного
  подтверждения.
\end{enumerate}

Все три типа сбоев должны оставлять записи в журнале: какая доставка
была дедуплицирована, какое сообщение было считано повторно и какой
ответ после проверки не стали отправлять ещё раз.

Сопутствующий код этого расширения можно запускать автономно, без API
Key. Тестовый почтовый ящик и очередь SQS эмулируются локальными файлами
и отделены от реального почтового ящика, для которого выше ожидаются
учётные данные. Эмулятор очереди сохраняет семантику тайм-аута
видимости, дескрипторов получения и порядка сообщений внутри группы. Для
имитации сбоя родительский процесс напрямую завершает дочерний процесс в
заданной точке. Для каждого канала ответа демонстрируется отдельный
подход: для электронной почты проверяется папка отправленных сообщений,
для диалога в приложении используется ключ идемпотентности, а SMS не
поддерживает ни одну из этих возможностей. Для воспроизведения перейдите
в \texttt{chapter6/agent-with-event-trigger} и выполните
\texttt{python\ multisource\_demo.py}: терминал последовательно выведет
оставшиеся в журнале записи для всех трёх типов сбоев. В сопутствующем
журнале запуска успешно пройдены проверки всех трёх типов сбоев.
\end{quote}

Почтовая очередь создаёт основу для последовательной обработки. Далее мы
добавим длительно работающие инструменты и рассмотрим, как во время
ожидания принимать новые требования, отменять задачи и возобновлять
работу.

\hypertarget{ux441ux43eux432ux43cux435ux441ux442ux438ux43cux43eux441ux442ux44c-ux43fux440ux438-ux43eux442ux441ux443ux442ux441ux442ux432ux438ux438-ux432ux441ux442ux440ux43eux435ux43dux43dux43eux439-ux43fux43eux434ux434ux435ux440ux436ux43aux438-ux430ux441ux438ux43dux445ux440ux43eux43dux43dux43eux441ux442ux438}{%
\subsection{Совместимость при отсутствии встроенной поддержки
асинхронности}\label{ux441ux43eux432ux43cux435ux441ux442ux438ux43cux43eux441ux442ux44c-ux43fux440ux438-ux43eux442ux441ux443ux442ux441ux442ux432ux438ux438-ux432ux441ux442ux440ux43eux435ux43dux43dux43eux439-ux43fux43eux434ux434ux435ux440ux436ux43aux438-ux430ux441ux438ux43dux445ux440ux43eux43dux43dux43eux441ux442ux438}}

Предположим, AI-агент составляет письмо и уже вызвал
\texttt{search\_contacts} для поиска контакта. Пока результат ещё не
получен, пользователь просит: «Подождите, сначала узнайте погоду на
завтра». Если текущий интерфейс модели требует сначала предоставить
результат для каждого вызова инструмента, Harness должен преобразовать
фоновое состояние в последовательность сообщений, допустимую для этого
интерфейса.

\textbf{Представление фонового состояния в синхронном формате
сообщений.}

\textbf{При работе с синхронным интерфейсом следует по возможности
сохранять соответствие между обычными вызовами и реальными результатами,
добавляя ответ о состоянии только при переключении задачи в процессе
выполнения.} Если слишком рано отвечать на каждый вызов «всё ещё
выполняется», траектория переполнится незавершёнными состояниями, и
модели станет сложнее отслеживать реальные результаты. На практике можно
придерживаться пяти правил:

\begin{enumerate}
\def\labelenumi{\arabic{enumi}.}
\tightlist
\item
  \textbf{Записывать полный вызов.} Записывать в траекторию завершённое
  сообщение assistant и элементы вызовов инструментов, возвращённые API,
  сохраняя идентификаторы вызовов. Состояние серверного инференса
  сохраняется предусмотренным поставщиком способом продолжения.
\item
  \textbf{Записывать фактическое завершение.} После завершения
  инструмента записывать реальный результат, идентификатор задачи и
  состояние выполнения; пока задача выполняется, диспетчер задач
  постоянно сохраняет её незавершённое состояние.
\item
  \textbf{При переключении сообщать текущее состояние.} Если интерфейс
  допускает подобную полезную нагрузку результата, для незавершённого
  вызова можно вернуть явную информацию о состоянии, например: «Задача
  всё ещё выполняется в фоновом режиме, результат пока не получен», а
  затем добавить событие пользователя. Идентификатор задачи в ответе о
  состоянии связывает его с последующим уведомлением о завершении.
\item
  \textbf{Упорядочивать сообщения после прерывания.} Если продолжение в
  процессе выполнения не поддерживается, следует отменить незавершённую
  генерацию, сохранить подтверждённые завершённые сообщения и состояния
  инструментов, а затем отправить новый запрос. Незавершённый вывод
  обрабатывается согласно требованиям интерфейса, чтобы фрагмент
  параметров не был ошибочно принят за исполняемый вызов.
\item
  \textbf{Принимать обычные события пакетами.} Дополнительная информация
  поступает в очередь и передаётся единым пакетом по завершении текущего
  цикла или в заданной безопасной точке; при этом регистрируется, какие
  события уже обработаны.
\end{enumerate}

В примере с поиском контакта полный вызов сначала записывается в
траекторию. Когда пользователь переключает задачу, Harness добавляет
ответ о том, что поиск всё ещё выполняется, а затем передаёт запрос
погоды. Позже результат поиска контакта поступает в контекст как новое
событие с указанием источника и идентификатора задачи, после чего агент
продолжает составлять письмо. Если исходный вызов уже получил ответ о
состоянии, последующий результат передаётся в предусмотренной адаптером
форме события, сохраняя соответствие между вызовом и результатом.

Ответ о состоянии сообщает только о текущем ходе выполнения. Шаги,
зависящие от данных контакта, должны дождаться реального результата.
Диспетчер задач может хранить состояния \texttt{running},
\texttt{completed}, \texttt{cancelled}, \texttt{failed} и перед
отправкой ответа или вызовом следующего инструмента проверять, выполнены
ли зависимости.

\textbf{Проектирование асинхронных инструментов с помощью дескрипторов
задач.}

Для новых инструментов, выполняющих длительные задачи, автор
предпочитает определять «запуск» как законченную операцию, которая
немедленно возвращает реальный дескриптор задачи. Благодаря этому модель
различает состояния «уже запущено» и «уже завершено», а последующие
запросы, отмена и события с результатами получают единое основание для
связывания. Например, \texttt{initiate\_phone\_call} запускает задачу
телефонного звонка и возвращает идентификатор задачи вместе с состоянием
«идёт набор номера», а \texttt{phone\_call\_connected} и
\texttt{phone\_call\_ended} сообщают соответственно об установлении и
завершении соединения.

В описании инструмента следует явно указать: «Этот инструмент запускает
задачу телефонного звонка и после успешного запуска возвращает ID
задачи. Вы можете продолжить работу над другими вопросами; сведения о
ходе звонка и его результате возвращаются отдельными событиями». Имя
инструмента показывает, что это операция запуска, а структура ответа и
последующие события объясняют, как задача выполняется и завершается.
Система также предоставляет интерфейсы для запроса состояния, отмены и
возобновления по ID задачи, чтобы модель понимала, какие действия ещё
доступны.

\textbf{Поэлементная обработка пакета событий.}

Когда в контекст одновременно поступает несколько событий, модель может
ответить только на последнее и пропустить ранее сформулированные
требования. Среда выполнения должна хранить состояние обработки каждого
события, чтобы условия, на которые ещё не дан ответ, оставались
видимыми. Промпт может требовать от модели проверять события по одному,
а строка состояния --- перечислять ещё не обработанную информацию:

\begin{Shaded}
\begin{Highlighting}[]
\NormalTok{[Необработанное событие 1/4] Результат database\_query: …}
\NormalTok{[Необработанное событие 2/4] Дополнение пользователя: учитывать только данные по Пекину}
\NormalTok{[Необработанное событие 3/4] Системное напоминание: до срока сдачи отчёта осталось 30 минут}
\NormalTok{[Необработанное событие 4/4] Вопрос пользователя: как продвигается работа?}
\end{Highlighting}
\end{Shaded}

В сводке пакета можно указать: «1 результат инструмента, 2 сообщения
пользователя, 1 системное напоминание». Среда выполнения дополнительно
связывает события с ограничениями задачи: условие по региону изменяет
область запроса, срок сдачи влияет на планирование, а вопрос о ходе
работы требует ответа. После обработки регистрируется состояние
«завершено», «ожидает» или «требует уточнения», а на следующей итерации
нерешённые пункты продолжают сохраняться.

Один и тот же номер задачи связывает запуск, отчёты о ходе выполнения и
окончательный результат. Если в процессе поступит требование остановить
задачу, необходимо дополнительно проверить, получено ли для этого номера
подтверждение отмены.

\begin{quote}
\textbf{Эксперимент 6-2 ★★★: асинхронный агент с поддержкой
параллельного выполнения и прерывания}

В этом эксперименте сначала проверяется асинхронная среда выполнения, а
затем подключается модель для проверки решений. В сопутствующем журнале
выполнения четыре сценария реализованы с помощью настоящих подпроцессов
операционной системы, а диспетчеризацией событий и формированием отчётов
управляет детерминированный скрипт.

\begin{enumerate}
\def\labelenumi{\arabic{enumi}.}
\tightlist
\item
  \textbf{Фоновое выполнение и немедленный ответ.} Запустите процесс
  анализа длительностью несколько секунд и немедленно верните
  идентификатор задачи; во время ожидания ответьте на вопрос о текущем
  времени, а после завершения примите результат анализа.
\item
  \textbf{Очередь и пакетная обработка.} Во время выполнения длительной
  задачи последовательно добавьте требования «ответить по-японски» и
  «оформить в виде веб-страницы», а после завершения задачи обработайте
  результат и оба требования одним пакетом. Скрипт создаёт HTML на
  японском языке.
\item
  \textbf{Отмена и возобновление.} Продемонстрируйте, что при получении
  команды «Остановить анализ» выполняющийся ход модели прерывается, а
  все активные фоновые задачи получают требование остановиться. В этом
  сценарии прерывание инициирует скрипт; на тот момент выполнялся только
  один процесс анализа --- T1. На 0,803 секунды событие прерывания
  зафиксировало прогресс T1 на уровне 18\%; лишь на 0,806 секунды
  система зарегистрировала его код завершения −15. Сразу после этого
  была запущена новая задача T2. На 4,436 секунды был получен результат
  T2; T1 не вызвала обратный вызов завершения.
\item
  \textbf{Параллельные задачи и решения на основе прогресса.}
  Одновременно запустите три процесса анализа, продвигающихся на каждом
  логическом шаге времени соответственно на 3\%, 2\% и 1\%. В журнале
  параллельного выполнения этого эксперимента к моменту завершения самой
  быстрой задачи T1 ответ на запрос состояния показывает, что T2
  достигла 68\%, а T3 --- 34\%. Согласно
  \crossreflink{fig:6-5}{рисунку 6-5}, задачи с прогрессом не более 50\%
  отменяются. Скрипт отменил только T3, а в итоге получил результаты T1
  и T2. Это показывает, что одна среда выполнения способна как
  отреагировать на пакетную остановку, так и сохранить полезную
  параллельную работу.
\end{enumerate}

\begin{figure}
\hypertarget{fig:6-5}{%
\centering
\includegraphics{images/revised/fig6-5.pdf}
\caption{Рисунок 6-5. Прерывание и возобновление асинхронного агента в
эксперименте 6-2}\label{fig:6-5}
}
\end{figure}
\end{quote}

\hypertarget{ux43dux430ux442ux438ux432ux43dux430ux44f-ux430ux441ux438ux43dux445ux440ux43eux43dux43dux43eux441ux442ux44c-ux43cux43eux434ux435ux43bux438-gpt-6-astra}{%
\subsection{Нативная асинхронность модели: GPT-6
Astra}\label{ux43dux430ux442ux438ux432ux43dux430ux44f-ux430ux441ux438ux43dux445ux440ux43eux43dux43dux43eux441ux442ux44c-ux43cux43eux434ux435ux43bux438-gpt-6-astra}}

В описанной выше схеме Harness адаптирует синхронный формат сообщений.
Нативный асинхронный интерфейс напрямую позволяет модели продолжать
независимую работу, пока инструмент ещё не вернул результат, и получить
этот результат позднее. По состоянию на сентябрь 2026 года GPT-6 Astra
поддерживает асинхронные вызовы инструментов (Async tool calling) и
корректировку в середине хода (Mid-turn steering), которые позволяют
соответственно обрабатывать ожидание инструментов и обновления от
пользователя. На \crossreflink{fig:6-6}{рисунке 6-6} сопоставлены потоки
сообщений при адаптации синхронного интерфейса и нативной
асинхронности.\footnote{OpenAI, «Async tool calling»; «Using GPT-6
  Astra», дата проверки: 2026-09-05.}\footnote{OpenAI, «Mid-turn
  steering», дата проверки: 2026-09-05.}

\begin{figure}
\hypertarget{fig:6-6}{%
\centering
\includegraphics{images/revised/fig6-6.pdf}
\caption{Рисунок 6-6. Совместимость с синхронным интерфейсом и нативная
асинхронность модели}\label{fig:6-6}
}
\end{figure}

\textbf{Продолжение независимой работы во время ожидания инструмента.}
Например, после запуска поиска помещения для встречи AI-агент может
сначала подготовить повестку и контрольный список, а после поступления
результатов сравнить цены и вместимость. Поле \texttt{async:\ true} в
определении инструмента объявляет поддержку асинхронного выполнения.
Получив полный элемент вызова, приложение запускает выполнение и
сохраняет идентификатор вызова \texttt{call\_id}. Когда результат
поступает, приложение возвращает его с исходным идентификатором и
продолжает работу с последнего состояния ответа. Модель планирует
последующую работу, а Harness управляет сопоставлением вызовов и
жизненным циклом задач.\footnote{OpenAI, «Async tool calling»; «Using
  GPT-6 Astra», дата проверки: 2026-09-05.}

\textbf{Получение новых требований во время выполнения задачи.} Пока
модель рассуждает или формирует ответ, пользователь может добавить
требования: «сократить бюджет до 1000» или «изменить число участников на
10». Приложение отправляет событие корректировки через то же соединение,
а после подтверждения его получения сервером продолжает обработку с
текущего состояния ответа. После разрыва и восстановления соединения
приложение может возобновить задачу на основе сохранённых событий и
состояния ответа.\footnote{OpenAI, «Mid-turn steering», дата проверки:
  2026-09-05.}

Эти две возможности можно использовать совместно: поиск помещения
продолжает выполняться, новые требования пользователя учитываются в
последующих решениях, а после поступления результатов поиска завершается
выбор. Исполнитель при этом поддерживает фактическое состояние фоновых
задач. Если поиск нужно остановить или бронирование отменить, Harness
вызывает соответствующий интерфейс и проверяет подтверждение; уже
завершённые бизнес-операции передаются в процесс отмены или компенсации.

Если модель и интерфейс уже поддерживают нативную асинхронность, автор
предпочитает использовать нативный протокол, чтобы приостановка
инструментов, обновления от пользователя и продолжение ответа
обрабатывались в соответствии с контрактом интерфейса. Слой
совместимости с синхронным форматом предназначен прежде всего для
существующих моделей и старых интерфейсов. После обновления модели можно
постепенно удалить логику совместимости, служащую только для
восстановления пар сообщений, и сосредоточить инженерные усилия на
состоянии задач, отмене и восстановлении.

\hypertarget{ux43dux430ux434ux451ux436ux43dux430ux44f-ux43eux431ux440ux430ux431ux43eux442ux43aux430-ux430ux441ux438ux43dux445ux440ux43eux43dux43dux44bux445-ux437ux430ux434ux430ux447}{%
\subsection{Надёжная обработка асинхронных
задач}\label{ux43dux430ux434ux451ux436ux43dux430ux44f-ux43eux431ux440ux430ux431ux43eux442ux43aux430-ux430ux441ux438ux43dux445ux440ux43eux43dux43dux44bux445-ux437ux430ux434ux430ux447}}

Для надёжного асинхронного взаимодействия необходимо одновременно
проверять сообщения, задачи и состояние выполнения.

\begin{enumerate}
\def\labelenumi{\arabic{enumi}.}
\tightlist
\item
  \textbf{Принадлежность результата.} Задержанный результат содержит
  идентификатор исходного вызова, по которому приложение сопоставляет
  его с соответствующей задачей; решения, зависящие от этого результата,
  принимаются после поступления данных.
\item
  \textbf{Возобновление и отмена.} После изменения требований
  пользователем система сохраняет работу, которая остаётся актуальной,
  обновляет последующие шаги и явно указывает, какие фоновые задачи
  продолжают выполняться, а какие уже остановлены.
\item
  \textbf{Несколько требований.} Бюджет, число участников, регион и срок
  выполнения по отдельности сохраняются в состоянии задачи, а итоговый
  результат проверяется на соответствие каждому из этих условий.
\end{enumerate}

\begin{quote}
\textbf{Эксперимент 6-3 ★★★: нативная асинхронность модели и
корректировка в середине хода}

В этом эксперименте агент выбирает для условной встречи самое дешёвое
помещение, соответствующее требованиям по бюджету и числу участников.
Исходный бюджет составляет 2000, число участников --- 20; инструмент
поиска возвращает цены и вместимость трёх помещений только через 8
секунд.

\begin{longtable}[]{@{}
  >{\raggedright\arraybackslash}p{(\columnwidth - 4\tabcolsep) * \real{0.3588}}
  >{\raggedleft\arraybackslash}p{(\columnwidth - 4\tabcolsep) * \real{0.2079}}
  >{\raggedleft\arraybackslash}p{(\columnwidth - 4\tabcolsep) * \real{0.4334}}@{}}
\toprule\noalign{}
\begin{minipage}[b]{\linewidth}\raggedright
\hspace{0pt}Помещение
\end{minipage} & \begin{minipage}[b]{\linewidth}\raggedleft
\hspace{0pt}Цена
\end{minipage} & \begin{minipage}[b]{\linewidth}\raggedleft
\hspace{0pt}Вместимость
\end{minipage} \\
\midrule\noalign{}
\endhead
\bottomrule\noalign{}
\endlastfoot
\hspace{0pt}A & \hspace{0pt}1800 & \hspace{0pt}30 \\
\hspace{0pt}B & \hspace{0pt}900 & \hspace{0pt}12 \\
\hspace{0pt}C & \hspace{0pt}600 & \hspace{0pt}8 \\
\end{longtable}

При исходных условиях следует выбрать A. Группы с корректировкой во
время выполнения задачи получают два обновления: бюджет изменён на 1000,
а число участников --- на 10; теперь следует выбрать B. Вместе с
результатом инструмент возвращает новую сгенерированную квитанцию,
которую необходимо процитировать в итоговом ответе, чтобы проверить,
использовала ли модель фактически поступивший результат.

Эксперимент включает пять групп по три запуска: синхронный инструмент,
нативный асинхронный инструмент, корректировка во время инференса,
корректировка во время ожидания асинхронного инструмента и контрольная
модель без поддержки корректировки. В первых двух группах проверяется,
можно ли во время ожидания составить независимый контрольный список
подготовки; в двух группах с корректировкой --- одновременно ли
применяются требования по бюджету и числу участников; в контрольной
группе --- сообщает ли интерфейс явно об отсутствии поддержки.

В сохранённых записях запусков группа с асинхронным инструментом во
время ожидания составила независимый контрольный список, а обе группы с
корректировкой выбрали B согласно новым условиям. Во всех запусках с
инструментом запрос выполнялся только один раз, а фактический результат
возвращался исходному вызову; интерфейс контрольной модели явно ответил,
что она не поддерживает получение новых требований в процессе генерации
одного хода. Общее время выполнения в синхронной и асинхронной группах
было близким; основное различие заключалось в возможности продвигать
другую работу во время ожидания.

Что выберет модель, если запомнит только изменение числа участников на
10, но продолжит использовать прежний бюджет 2000? И A, и B вмещают
достаточно людей, однако самым дешёвым всё равно остаётся B. Поэтому
правильный выбор помещения ещё не доказывает, что обновление бюджета
было учтено.

Если число участников изменилось, нужно ли повторять уже отправленный
запрос? В этом эксперименте инструмент не принимает параметры фильтрации
и возвращает все три помещения. Агент может повторно использовать уже
выполняющийся запрос, дождаться результата и затем отфильтровать
варианты по условиям: 10 участников и бюджет 1000. Если бы инструмент
возвращал только помещения вместимостью не менее 20 человек, в старом
результате оказалось бы лишь A, а B вместимостью 12 человек вообще
отсутствовало бы. В таком случае простое изменение числа участников на
10 при работе со старым списком всё равно привело бы к пропуску B. Если
исходный запрос уже отфильтровал B, запрос придётся выполнить повторно.
\end{quote}

Асинхронные механизмы позволяют внешним изменениям поступать в систему
непосредственно в ходе выполнения задачи. Голосовое взаимодействие,
работа с экраном и взаимодействие с роботами дополнительно сокращают
цикл обратной связи: система должна непрерывно принимать входные данные
и в подходящий момент корректировать вывод и действия.

\hypertarget{ux433ux43eux43bux43eux441-ux441ux430ux43cux44bux439-ux435ux441ux442ux435ux441ux442ux432ux435ux43dux43dux44bux439-ux447ux435ux43bux43eux432ux435ux43aux43e-ux43cux430ux448ux438ux43dux43dux44bux439-ux438ux43dux442ux435ux440ux444ux435ux439ux441}{%
\section{Голос: самый естественный человеко-машинный
интерфейс}\label{ux433ux43eux43bux43eux441-ux441ux430ux43cux44bux439-ux435ux441ux442ux435ux441ux442ux432ux435ux43dux43dux44bux439-ux447ux435ux43bux43eux432ux435ux43aux43e-ux43cux430ux448ux438ux43dux43dux44bux439-ux438ux43dux442ux435ux440ux444ux435ux439ux441}}

Голос позволяет пользователю сообщать AI-агенту о своих потребностях,
дополнять условия и подтверждать результаты, даже когда его руки или
зрительное внимание заняты другой работой. Голосовой ввод преобразует
диктовку в текст, а голосовой агент дополнительно управляет диалогом,
вызовами инструментов и воспроизведением ответов; оба подхода
поддерживают описанный в предисловии whisper coding.

В этом разделе рассматриваются два варианта использования:
непосредственный диалог пользователя с агентом и общение агента с
внешними людьми от имени пользователя. В первом случае основное внимание
уделяется распознаванию, ответам и прерываниям, а во втором во время
разговора также необходимо подтверждать личность, выполнять
бизнес-запросы и принимать решения.

\hypertarget{ux43fux43eux441ux43bux435ux434ux43eux432ux430ux442ux435ux43bux44cux43dux43eux441ux442ux44c-ux432ux437ux430ux438ux43cux43eux434ux435ux439ux441ux442ux432ux438ux44f-ux43eux442-ux43aux430ux441ux43aux430ux434ux430-ux43a-ux43fux43eux43bux43dux43eux43cux443-ux434ux443ux43fux43bux435ux43aux441ux443}{%
\subsection{Последовательность взаимодействия: от каскада к полному
дуплексу}\label{ux43fux43eux441ux43bux435ux434ux43eux432ux430ux442ux435ux43bux44cux43dux43eux441ux442ux44c-ux432ux437ux430ux438ux43cux43eux434ux435ux439ux441ux442ux432ux438ux44f-ux43eux442-ux43aux430ux441ux43aux430ux434ux430-ux43a-ux43fux43eux43bux43dux43eux43cux443-ux434ux443ux43fux43bux435ux43aux441ux443}}

В описании GPT-Live OpenAI рассматривает три типа голосового
взаимодействия: каскадные системы, аудиомодели с разделением на реплики
и полнодуплексные системы\footnote{OpenAI, \emph{Introducing GPT-Live},
  2026-07-08, описание каскадных систем, аудиомоделей с разделением на
  реплики и полнодуплексного взаимодействия.}. По мере развития этого
направления решение о том, «когда начать говорить, а когда продолжить
слушать», всё больше передаётся модели, которая принимает его
непрерывно. В каскадных системах передачей права говорить в основном
управляет логика определения конечной точки, тогда как полнодуплексная
модель может обрабатывать наложение речи, паузы и прерывания с учётом
непрерывного звука и состояния диалога.

\begin{longtable}[]{@{}
  >{\raggedright\arraybackslash}p{(\columnwidth - 6\tabcolsep) * \real{0.2372}}
  >{\raggedright\arraybackslash}p{(\columnwidth - 6\tabcolsep) * \real{0.2159}}
  >{\raggedright\arraybackslash}p{(\columnwidth - 6\tabcolsep) * \real{0.2748}}
  >{\raggedright\arraybackslash}p{(\columnwidth - 6\tabcolsep) * \real{0.2721}}@{}}
\toprule\noalign{}
\begin{minipage}[b]{\linewidth}\raggedright
\hspace{0pt}Парадигма
\end{minipage} & \begin{minipage}[b]{\linewidth}\raggedright
\hspace{0pt}Основная структура
\end{minipage} & \begin{minipage}[b]{\linewidth}\raggedright
\hspace{0pt}Возможности, которые проще реализовать
\end{minipage} & \begin{minipage}[b]{\linewidth}\raggedright
\hspace{0pt}Основные вопросы проектирования
\end{minipage} \\
\midrule\noalign{}
\endhead
\bottomrule\noalign{}
\endlastfoot
\hspace{0pt}Каскадная & \hspace{0pt}Разделение функций между VAD, ASR,
LLM и TTS & \hspace{0pt}Независимая замена модулей, поэтапная проверка
ввода и вывода & \hspace{0pt}Задержки между модулями, потеря
акустической информации при расшифровке \\
\hspace{0pt}Сквозная аудиомодель с разделением на реплики &
\hspace{0pt}Модель напрямую принимает и генерирует аудио &
\hspace{0pt}Формирование ответа с учётом интонации, эмоций и другой
информации & \hspace{0pt}Определение окончания речи, смена реплик и
выбор момента ответа \\
\hspace{0pt}Полнодуплексная & \hspace{0pt}Непрерывное получение аудио и
генерация вывода & \hspace{0pt}Одновременная речь, реплики согласия,
прерывания и непрерывное взаимодействие & \hspace{0pt}Непрерывное
состояние, управление выводом и оценка взаимодействия \\
\end{longtable}

Прямая обработка аудио моделью, потоковый режим ввода и вывода и
способность одновременно слушать и говорить --- отдельные возможности,
которые следует проверять независимо. В каскадную систему также можно
добавить потоковые модули и управление прерываниями; темп взаимодействия
сквозной модели зависит от конкретного устройства модели.

\hypertarget{ux43fux430ux440ux430ux434ux438ux433ux43cux430-1-ux43aux430ux441ux43aux430ux434ux43dux44bux439-ux43aux43eux43dux432ux435ux439ux435ux440-cascading}{%
\subsection{Парадигма 1 · Каскадный конвейер
(Cascading)}\label{ux43fux430ux440ux430ux434ux438ux433ux43cux430-1-ux43aux430ux441ux43aux430ux434ux43dux44bux439-ux43aux43eux43dux432ux435ux439ux435ux440-cascading}}

Каскадная система разделяет голосовой диалог на несколько модулей.
Детектор голосовой активности (VAD) распознаёт участки речи в аудио, а
логика определения конечной точки решает, когда передать такой фрагмент
дальше; автоматическое распознавание речи (ASR) преобразует аудио в
текст; LLM понимает запрос и генерирует ответ; затем синтез речи (TTS)
преобразует ответ в аудио, готовое к воспроизведению. На
\crossreflink{fig:6-7}{рисунке 6-7} показана базовая структура с
последовательным выполнением этапов.

\begin{figure}
\hypertarget{fig:6-7}{%
\centering
\includegraphics{images/revised/fig6-7.pdf}
\caption{Рисунок 6-7. Последовательный конвейер голосового
AI-агента}\label{fig:6-7}
}
\end{figure}

Акустический VAD находит участки, где говорит человек, поэтому
последующее распознавание может пропускать длительные периоды тишины.
Однако одна и та же пауза имеет разное значение после фраз «Я хочу
забронировать\ldots» и «Есть билеты на утро четверга?». Семантический
VAD --- семантическое определение конечной точки --- в первом случае
может продолжить запись и лишь во втором передать реплику диалоговой
модели. Так голосовые наблюдения становятся разреженными во времени:
непрерывный звук запускает инференс только в отдельные моменты. Решение
о том, нужно ли отвечать, принимает модель, услышавшая реплику; она
также может продолжить слушать или сначала проверить рейсы.

\begin{longtable}[]{@{}
  >{\raggedright\arraybackslash}p{(\columnwidth - 4\tabcolsep) * \real{0.2960}}
  >{\raggedright\arraybackslash}p{(\columnwidth - 4\tabcolsep) * \real{0.3195}}
  >{\raggedright\arraybackslash}p{(\columnwidth - 4\tabcolsep) * \real{0.3846}}@{}}
\toprule\noalign{}
\begin{minipage}[b]{\linewidth}\raggedright
\hspace{0pt}Модуль
\end{minipage} & \begin{minipage}[b]{\linewidth}\raggedright
\hspace{0pt}Назначение
\end{minipage} & \begin{minipage}[b]{\linewidth}\raggedright
\hspace{0pt}Распространённые источники ожидания или ошибок
\end{minipage} \\
\midrule\noalign{}
\endhead
\bottomrule\noalign{}
\endlastfoot
\hspace{0pt}VAD и логика определения конечной точки &
\hspace{0pt}Обнаружение речи и выбор момента передачи &
\hspace{0pt}Порог тишины, шум и паузы в середине реплики \\
\hspace{0pt}ASR & \hspace{0pt}Преобразование аудио в текст &
\hspace{0pt}Время декодирования, распознавание имён и имён
собственных \\
\hspace{0pt}LLM & \hspace{0pt}Понимание, рассуждение и генерация ответа
& \hspace{0pt}Очередь запросов, генерация первого фрагмента ответа \\
\hspace{0pt}TTS и воспроизведение & \hspace{0pt}Синтез аудио и его
передача на динамик & \hspace{0pt}Синтез первого пакета, передача по
сети и буферизация воспроизведения \\
\end{longtable}

В полностью последовательной реализации после окончания реплики
пользователь должен поочерёдно дождаться подтверждения конечной точки,
транскрибирования, генерации ответа и первого фрагмента аудио. В примере
на \crossreflink{fig:6-8}{рисунке 6-8} первый звук ответа слышен
пользователю лишь через 1,4 секунды. Самая продолжительная часть этого
времени --- ожидание подтверждения того, что пользователь закончил
говорить: 500 миллисекунд. Даже если сократить транскрибирование со 100
до 50 миллисекунд, пользователь будет ждать лишь на 50 миллисекунд
меньше.

\begin{figure}
\hypertarget{fig:6-8}{%
\centering
\includegraphics{images/revised/fig6-8.pdf}
\caption{Рисунок 6-8. Каскад задержек: последовательное накопление
общего времени ответа}\label{fig:6-8}
}
\end{figure}

Нагрузка на сервис дополнительно создаёт ожидание в очереди. На
\crossreflink{fig:6-9}{рисунке 6-9} эта зависимость показана с помощью
одноканальной модели M/M/1: поступление запросов описывается
пуассоновским процессом, время обслуживания --- экспоненциальным
распределением, а при среднем времени обслуживания \(S\) и коэффициенте
использования \(\rho<1\) среднее общее время обслуживания и ожидания в
очереди составляет \(W=S/(1-\rho)\). Если \(S=1\) секунда, то при росте
коэффициента использования с 0,5 до 0,8 среднее общее время
увеличивается с 2 до 5 секунд. Само обслуживание по-прежнему занимает в
среднем 1 секунду, а дополнительные 3 секунды целиком уходят на ожидание
в очереди. Когда поток запросов приближается к пропускной способности
системы, кратковременные колебания нагрузки приводят к накоплению
очереди.

\begin{figure}
\hypertarget{fig:6-9}{%
\centering
\includegraphics{images/revised/fig6-9.pdf}
\caption{Рисунок 6-9. Коэффициент использования одноканальной системы и
среднее общее время, включая ожидание в очереди и
обслуживание}\label{fig:6-9}
}
\end{figure}

\begin{quote}
\textbf{Эксперимент 6-4 ★: создание традиционного голосового агента}

Предположим, помощник перечисляет рейсы на четверг, а вы перебиваете его
поправкой: «Подождите, я хочу забронировать на пятницу». Браузер
передаёт звук с микрофона серверу через WebSocket. Как только Silero VAD
обнаруживает начало вашей речи, сервер отправляет
\texttt{speech\_start}, после чего браузер сразу останавливает
воспроизведение и очищает очередь старого аудио. В этот момент помощник
ещё не знает, на какой день вы изменили бронирование. Лишь когда вы
закончите говорить и тишина продлится дольше стандартных 500
миллисекунд, сервер отправит \texttt{speech\_end} и передаст фрагмент
речи локальной модели Whisper tiny для транскрибирования.

Doubao начинает возвращать текст. Сколько ещё пройдёт до того, как
пользователь услышит ответ? В записи одного прогона с воспроизведением
транскрибирование заняло 2,710 секунды. Doubao вернула первый токен
через 0,749 секунды после отправки запроса, но полный ответ поступил
лишь через 0,792 секунды. Программа воспроизведения дождалась всего
текста и только затем передала его Fish S1 для синтеза; ещё через 2,819
секунды сервер получил первую порцию аудио. Если отсчитывать от
определения конечной точки, сумма длительностей этих трёх
последовательных этапов составляет около \textbf{6,321 секунды}. В
архивном отчёте указано 6,279 секунды, поскольку вместо времени
генерации полного ответа там использовано время получения первого
токена.

Эти 6,321 секунды охватывают только обработку на сервере. Сначала
пользователь должен дождаться, пока детектор тишины определит, что
реплика завершена, а в конце --- пока аудио поступит в браузер и начнёт
воспроизводиться.
\end{quote}

\hypertarget{ux43eux442-ux43fux43eux441ux43bux435ux434ux43eux432ux430ux442ux435ux43bux44cux43dux43eux439-ux43eux431ux440ux430ux431ux43eux442ux43aux438-ux43a-ux43fux43eux442ux43eux43aux43eux432ux43eux43cux443-ux432ux43eux441ux43fux440ux438ux44fux442ux438ux44e}{%
\subsubsection{От последовательной обработки к потоковому
восприятию}\label{ux43eux442-ux43fux43eux441ux43bux435ux434ux43eux432ux430ux442ux435ux43bux44cux43dux43eux439-ux43eux431ux440ux430ux431ux43eux442ux43aux438-ux43a-ux43fux43eux442ux43eux43aux43eux432ux43eux43cux443-ux432ux43eux441ux43fux440ux438ux44fux442ux438ux44e}}

Последовательная схема удобна для поэтапной проверки, но выявляет три
проблемы. Ожидание тишины увеличивает задержку определения конечной
точки; в транскрипции могут теряться интонация, эмоции и окружающие
звуки; при разбиении аудио по паузам адреса электронной почты, имена
людей и имена собственные могут лишиться целостного контекста. Для
систем, уже использующих каскадную архитектуру, автор рекомендует
сначала перейти к потоковой обработке: каждый этап должен как можно
раньше предоставлять инкрементальные результаты, чтобы распознавание,
рассуждение и воспроизведение перекрывались по времени, а подтверждённые
и ещё изменяемые части учитывались отдельно. Это позволяет повторно
использовать существующие модули и постепенно определять источники
ожидания.

\begin{itemize}
\tightlist
\item
  \textbf{Потоковое ASR.} Непрерывно принимать аудио и выдавать
  предварительную транскрипцию; по мере поступления последующего аудио
  пересматривать нестабильную часть, а после подтверждения конечной
  точки выдавать окончательный текст.
\item
  \textbf{Спекулятивное выполнение LLM.} Заранее готовить ответ на
  основе предварительной транскрипции. После подтверждения
  окончательного текста проверять, совпадают ли входные данные и
  остаются ли применимыми уже выполненные действия, а затем решать,
  сохранить, скорректировать или заново сгенерировать результат.
  Отправку данных наружу, запись бизнес-данных и другие операции
  выполнять лишь после подтверждения необходимых условий.
\item
  \textbf{Сегментированный вывод LLM.} Передавать фрагмент текста в TTS
  сразу после того, как он становится пригоден для озвучивания; при
  разбиении сохранять смысл предложений, пунктуацию и интонационные
  подсказки.
\item
  \textbf{Инкрементальный синтез TTS.} Постепенно возвращать блоки
  аудио, чтобы генерация последующего текста, синтез и воспроизведение
  перекрывались по времени. При вмешательстве пользователя одновременно
  обрабатывать задачу генерации и ещё не воспроизведённое содержимое
  буфера.
\end{itemize}

Стандартная реализация транскрибирования Whisper кодирует и декодирует
аудио по окнам\footnote{OpenAI, Whisper, описание транскрибирования по
  аудиоокнам и реализации модели.}. Инкрементальный процесс
транскрибирования можно построить посредством внешнего разбиения на
блоки или многократной передачи накопленного префикса. Конкретная
реализация должна управлять повторным использованием состояния,
стоимостью повторного кодирования и пересмотром предшествующего текста.
Модели со встроенной поддержкой потокового распознавания обрабатывают
непрерывный ввод непосредственно в своей архитектуре и процессе
обучения.

Системы восприятия на основе аудиоязыковых моделей также могут совместно
выводить текст и акустические события. История аудио, транскрипция и
лексические подсказки помогают распознавать имена людей, бренды и имена
собственные, а объём сохраняемой истории зависит от устройства контекста
и кэша. Определяя, закончил ли пользователь говорить, модель может
учитывать как смысловую завершённость, так и тишину. Во время паузы
после фразы «Я хочу забронировать\ldots» система реального времени
услышала только эти слова и ещё не может знать, что следом будет названа
дата.

Помимо транскрипции, интерфейс может предусматривать следующие события:
\texttt{speak\_start} и \texttt{speak\_end} обозначают начало и
окончание речи, \texttt{interrupt} --- прерывание, \texttt{emotion}
описывает эмоциональное состояние, а \texttt{laugh}, \texttt{sigh} и
\texttt{noise} фиксируют смех, вздохи и окружающий шум. Эти события
можно передавать агенту вместе с текстом и временными метками.

\begin{quote}
\textbf{Эксперимент 6-5 ★: имитация потокового восприятия речи с помощью
Qwen2-Audio}

В этом эксперименте префикс аудио постепенно увеличивается: при каждом
запуске в Qwen2-Audio повторно передаётся аудио от начала записи до
текущего момента, чтобы проследить, как с изменением входных данных
меняются транскрипция и акустические события. В контрольном варианте
используются конечная точка после 600 миллисекунд тишины и Whisper.
Тесты включают обычную речь, длинное предложение с паузой в 900
миллисекунд и речь с примесью розового шума.

В этом запуске локальной квантованной модели сохранены результаты
инференса для 13 префиксов; каждый запуск занимал около 8--11 секунд.
Пауза в середине реплики привела к разбиению в контрольном варианте;
аудиомодель обнаружила шум, но также ошибочно распознала кашель и смех.
\end{quote}

\hypertarget{ux43fux430ux440ux430ux434ux438ux433ux43cux430-2-ux441ux43aux432ux43eux437ux43dux430ux44f-ux432ux441ux435ux43cux43eux434ux430ux43bux44cux43dux430ux44f-ux43cux43eux434ux435ux43bux44c-omni}{%
\subsection{Парадигма 2 · Сквозная всемодальная модель
(Omni)}\label{ux43fux430ux440ux430ux434ux438ux433ux43cux430-2-ux441ux43aux432ux43eux437ux43dux430ux44f-ux432ux441ux435ux43cux43eux434ux430ux43bux44cux43dux430ux44f-ux43cux43eux434ux435ux43bux44c-omni}}

Сквозная аудиомодель непосредственно принимает звук, генерирует ответ на
основе языковой и акустической информации, а затем синтезирует речь
через ветвь аудиовывода. Непосредственная обработка аудио позволяет
учитывать темп речи, высоту голоса, паузы и окружающие звуки; подход с
предварительной расшифровкой и последующим ответом упрощает проверку
промежуточного текста и независимую замену модулей.

Эти акустические признаки влияют на понимание. Например, одна и та же
фраза «Пожалуйста, пришлите отчёт до обеда», произнесённая быстро или
медленно, воспринимается слушателем по-разному. При выводе модель также
может организовывать смысл с помощью пауз, ударений и интонации.

Система, непосредственно обрабатывающая аудио, также может дождаться,
пока пользователь закончит говорить, и только затем ответить. Например,
если пользователь сделал паузу, диктуя номер телефона, система всё равно
должна определить, закончил ли он перечислять цифры. Определение
конечной точки задаёт момент ответа; если собеседник вмешивается во
время воспроизведения, системе также приходится согласовывать генерацию
модели с выводом через динамик.\footnote{OpenAI, \emph{Voice activity
  detection}: представлены два способа определения границ хода --- по
  голосовой активности и по семантике.}

\begin{quote}
\textbf{Эксперимент 6-6 ★★: локальный запуск MiniCPM-o 4.5 и сравнение
сквозного режима с самокаскадом}

В этом эксперименте используется одна и та же модель MiniCPM-o 4.5 с
отключённым режимом мышления и сравниваются два пути: прямой ответ на
вопрос по аудио и предварительная расшифровка моделью с последующим
ответом только на основе этой расшифровки. Входные данные включают две
устные арифметические задачи и две аудиозаписи с одинаковым текстом, но
разным темпом речи.

На каждом из двух путей модель правильно ответила на три задания, однако
ошибки возникли в разных местах. При прямой обработке аудио модель
расслышала в одной из задач английскую фразу ``twelve boxes''
(«двенадцать коробок») как число 8, что привело к ошибке в вычислениях;
при расшифровке число 12 было сохранено правильно. В задании на темп
речи обе аудиозаписи содержали фразу ``Please send the report before
lunch.'' («Пожалуйста, пришлите отчёт до обеда»). После расшифровки
текст оказался полностью одинаковым, поэтому информация о темпе была
потеряна; при прямой обработке аудио модель правильно различила два
темпа речи.

В ходе эксперимента также была сгенерирована аудиозапись с частотой
дискретизации 24 kHz, объединившая понимание входных данных и
аудиовывод.
\end{quote}

\hypertarget{ux43fux430ux440ux430ux434ux438ux433ux43cux430-3-ux43fux43eux43bux43dux43eux434ux443ux43fux43bux435ux43aux441ux43dux430ux44f-ux43cux43eux434ux435ux43bux44c-ux432ux437ux430ux438ux43cux43eux434ux435ux439ux441ux442ux432ux438ux44f}{%
\subsection{Парадигма 3 · Полнодуплексная модель
взаимодействия}\label{ux43fux430ux440ux430ux434ux438ux433ux43cux430-3-ux43fux43eux43bux43dux43eux434ux443ux43fux43bux435ux43aux441ux43dux430ux44f-ux43cux43eux434ux435ux43bux44c-ux432ux437ux430ux438ux43cux43eux434ux435ux439ux441ux442ux432ux438ux44f}}

Полнодуплексная система одновременно говорит и слушает, а самое заметное
улучшение --- ощущение разговора с живым собеседником. Попробуйте
произнести: «Хочу забронировать билет на среду\ldots{} нет, на четверг».
Пауза посередине позволяет исправить себя; если ассистент уже начал
перечислять рейсы на среду, то, услышав «нет, на четверг», он должен
остановиться. Возможен и другой момент: пока собеседник ещё объясняет
маршрут, короткая реплика согласия даст ему понять, что ассистент
слушает. Только непрерывный приём аудио позволяет выбирать между
говорением, слушанием, паузой и вызовом инструмента. Синхронный перевод
также требует наложения ввода и вывода.

Выпущенная Kyutai в 2024 году \textbf{Moshi} раздельно моделирует
аудиопотоки пользователя и модели, представляя одновременную речь и
прерывания с помощью параллельных потоков\footnote{Kyutai, \emph{Moshi:
  a speech-text foundation model for real-time dialogue}, 2024.
  Описывает параллельные аудиопотоки пользователя и модели.}. Если
сохранить эти временны́е связи во входных данных, модель сможет научиться
определять, когда отвечать, а когда продолжать слушать.

Модель взаимодействия Thinking Machines Lab разделяет непрерывные ввод и
вывод на короткие блоки одинаковой длины, каждый из которых охватывает
200 миллисекунд и называется микроходом (micro-turn)\footnote{Thinking
  Machines Lab, \emph{Interaction Models: A Scalable Approach to
  Human-AI Collaboration}. Описывает микроходы, многопоточное
  взаимодействие и делегирование фоновых задач.}. Посмотрите на второй
блок временно́й шкалы на \crossreflink{fig:6-10}{рисунке 6-10}: ассистент
ещё произносит «четыре», но уже слышит, как пользователь исправляет дату
словом «нет». В третьем блоке ассистент замолкает и продолжает слушать,
пока пользователь не договорит «пятница».

\begin{figure}
\hypertarget{fig:6-10}{%
\centering
\includegraphics{images/revised/fig6-10.pdf}
\caption{Рисунок 6-10. Ввод, вывод и 200-миллисекундные микроходы в
полнодуплексной голосовой системе}\label{fig:6-10}
}
\end{figure}

Теперь рассмотрим нижнюю половину рисунка. Участники разговора слышат
два способных накладываться друг на друга голосовых потока, однако
модель может в одной последовательности попеременно обрабатывать блоки
ввода и генерировать блоки вывода. Тишина также занимает отдельный блок,
поэтому ассистент может сначала некоторое время послушать и лишь затем
заговорить. Указанные на рисунке 200 миллисекунд --- длина блока в этой
реализации; полное время от момента, когда пользователь перебивает
ассистента, до фактического прекращения звука из динамика включает также
затраты на обработку и буферизацию воспроизведения.

\hypertarget{ux43aux43eux433ux43dux438ux442ux438ux432ux43dux430ux44f-ux43fux43eux441ux43bux435ux434ux43eux432ux430ux442ux435ux43bux44cux43dux43eux441ux442ux44c-ux432ux437ux430ux438ux43cux43eux434ux435ux439ux441ux442ux432ux438ux435-ux432-ux440ux435ux430ux43bux44cux43dux43eux43c-ux432ux440ux435ux43cux435ux43dux438-ux438-ux433ux43bux443ux431ux43eux43aux43eux435-ux43cux44bux448ux43bux435ux43dux438ux435}{%
\subsection{Когнитивная последовательность: взаимодействие в реальном
времени и глубокое
мышление}\label{ux43aux43eux433ux43dux438ux442ux438ux432ux43dux430ux44f-ux43fux43eux441ux43bux435ux434ux43eux432ux430ux442ux435ux43bux44cux43dux43eux441ux442ux44c-ux432ux437ux430ux438ux43cux43eux434ux435ux439ux441ux442ux432ux438ux435-ux432-ux440ux435ux430ux43bux44cux43dux43eux43c-ux432ux440ux435ux43cux435ux43dux438-ux438-ux433ux43bux443ux431ux43eux43aux43eux435-ux43cux44bux448ux43bux435ux43dux438ux435}}

Когда во время звонка Pine в службу поддержки сотрудник ещё не закончил
фразу, ассистенту иногда следует подождать, а иногда --- негромко дать
понять, что он слушает. Даже когда ассистент говорит сам, он должен
услышать исправление, с которым собеседник его перебивает. Именно
поэтому автор считает перспективным направление Moshi: лишь правильно
выбирая момент, можно создать у собеседника ощущение, что на другом
конце провода его действительно слушают. Если для ответа требуется найти
информацию, возникает ещё один вопрос времени: сколько передач между
компонентами произойдёт от понимания вопроса до сообщения найденного
ответа?

Если передняя голосовая модель может сама вызвать инструмент, пусть
вызывает. Для поиска рейса достаточно двух обращений к модели: сначала
она запрашивает рейсы по условиям пользователя, затем сообщает
результат. При делегировании обращений четыре: передняя модель ставит
задачу, фоновая запрашивает данные, затем обрабатывает результат, а
передняя пересказывает его. Есть и другая точка зрения: модель реального
времени может ошибиться в параметрах, поэтому инструменты лучше поручить
фоновой модели, а передней оставить выбор момента, когда слушать и когда
говорить.

У передней модели GPT-Live-1 нет интерфейса для прямого вызова
прикладных функций: она может только делегировать задачу фоновому
агенту. GPT-Realtime в OpenAI Realtime поддерживает прямые вызовы
функций. Голосовая модель Doubao 3.0 (Seeduplex) тоже поддерживает их в
полностью сквозном режиме. При гибридной организации работы функции
вызывает подключённая языковая модель.\footnote{В описаниях моделей
  OpenAI GPT-Live-1 и GPT-Realtime рассматриваются соответственно
  фоновое делегирование и вызовы функций. Способ подключения модели для
  голосового взаимодействия в реальном времени Doubao 3.0 описан в
  документации Volcengine «Function Calling».}

Пусть одно обращение к передней модели занимает \(L_f\), к фоновой ---
\(L_b\), а работа инструмента --- \(T\). Тогда от конца реплики
пользователя до озвучивания результата при прямом вызове пройдёт
примерно \(2L_f+T\), а при делегировании --- \(2L_f+2L_b+T\).
Дополнительные \(2L_b\) не зависят от времени работы инструмента \(T\).
Сведения, которые могут понадобиться в любой реплике, например данные из
памяти пользователя, среда выполнения может сама найти и поместить в
контекст. Тогда достаточно одного обращения к передней модели, а
ожидание составит примерно \(L_f+T\). Если сведения найдены до того, как
пользователь заговорил, из времени ожидания исчезает и \(T\).

Задачи с несколькими шагами рассуждения, например сравнение вариантов
обмена билета, удобно поручить фоновой модели. Бронирование связано со
списанием денег и оформлением билета: кто бы ни вызывал инструмент,
делать заказ можно только после подтверждения пользователя. Результаты
поиска можно кратко озвучить, а полную таблицу показать на экране. Это
не зависит от того, кто вызвал инструмент.

\begin{quote}
\textbf{Эксперимент 6-7 ★★: кто вызывает инструменты голосового агента
--- передняя или фоновая модель?}

Для диалога используйте голосовую модель реального времени Doubao 3.0,
для фоновой работы --- DeepSeek V4 Pro на платформе Volcengine Ark.
Подготовьте пять имитаций инструментов: погода, чтение памяти
пользователя, веб-поиск, поиск рейсов и бронирование. Задайте задержки
ответа; для поиска рейсов и бронирования --- от 2 до 5 секунд. Заранее
синтезируйте реплики пользователя и при каждом повторении сценария
воспроизводите одну и ту же запись.

Сравните четыре схемы. В первой все инструменты вызывает передняя
модель, во второй --- фоновая, а передняя может только делегировать. В
двух остальных погоду запрашивает передняя модель, среда выполнения в
начале сеанса загружает память пользователя в контекст, а веб-поиск и
бронирование поручены фоновой. Различие только в том, кто ищет рейсы.
Передняя модель слышит голос пользователя, а фоновая видит только текст.
После каждой реплики пользователя интерфейс Doubao возвращает также
результат распознавания речи --- расшифровку. При делегировании
составьте из неё и произнесённых реплик помощника запись разговора в
хронологическом порядке. Передайте её фоновой модели вместе со списком
уже выполненных за этот разговор вызовов инструментов и заданием от
передней модели. Для каждой из четырёх схем повторите каждый сценарий 10
раз.

\begin{enumerate}
\def\labelenumi{\arabic{enumi}.}
\tightlist
\item
  По отдельности запросите погоду и рейсы. Засекайте время от конца
  реплики пользователя до момента, когда он услышит результат; «Сейчас
  проверю» ответом не считается. Сравните медианное время ожидания при
  прямом вызове и при делегировании. У автора для погоды получилось 2,81
  и 6,96 секунды соответственно.
\item
  Найдите рейсы, выберите билет и подтвердите заказ. Сверьте каждое из
  13 полей бронирования с результатами предыдущего поиска, последней
  репликой пользователя и его профилем. Проверьте, что заказ оформлен
  после подтверждения. У автора во всех четырёх схемах все поля были
  заполнены верно и ни один заказ не был оформлен раньше времени.
\item
  Теперь забронируйте билет для коллеги. Включите запись, в которой
  пользователь уточняет написание имени: «Чэнь Цзяи: Чэнь --- это 陈, из
  частей 耳 и 东; Цзя --- 嘉, как в 嘉奖 (``награждать''); И --- 怡, как
  в 怡然 (``безмятежно'')». Сопоставьте запись, расшифровку и имя в
  заказе при прямом вызове и при делегировании. У автора передняя модель
  верно записала 陈嘉怡 во всех десяти попытках. В расшифровке же стояло
  созвучное 陈佳宜, и при делегировании часть заказов оформили на
  неверное имя.
\item
  Пока не вернулись результаты поиска рейсов, включите запись с вопросом
  о погоде. Ответит ли передняя модель на новый вопрос? Сообщит ли затем
  о рейсах? У автора вопрос о погоде не остался без ответа ни разу, а
  вот о найденных рейсах модель почти всегда забывала
  сообщить.\footnote{\textbf{Уточнение русского издания.} В тексте автор
    пишет, что помощник отвечал на каждый попутный вопрос о погоде. В
    \href{https://github.com/bojieli/ai-agent-book/blob/40f4ce5778e6cc93e3606c3079bf8cf13a0c4d04/chapter6/voice-tool-routing/validation/runs/doubao-seeduplex-20261125-v1/tables.md}{таблице
    4 с результатами эксперимента} указано 39 ответов из 40: в схеме,
    где рейсы ищет передняя модель, --- 9 из 10. Результат поиска рейсов
    после ответа о погоде был озвучен всего один раз.}
\end{enumerate}
\end{quote}

Если делегирование необходимо, передней модели стоит прямо в задании
указать имена и числа, которые она уверенно расслышала, а также
последние исправления пользователя. Тогда ошибка распознавания не
попадёт в параметры инструмента. Среде выполнения следует запоминать ещё
не озвученные результаты, чтобы, ответив на попутный вопрос, передняя
модель вернулась к ним.

Фраза быстрой модели «Сейчас проверю» сообщает сотруднику поддержки, что
ассистент правильно понял запрос. Когда можно будет назвать цену билета,
зависит от того, как поиск и рассуждение согласованы с разговором. На
\crossreflink{fig:6-11}{рисунке 6-11} показаны три варианта: сначала
отреагировать, а затем дать развёрнутый ответ; поддерживать
взаимодействие передней моделью, получающей рекомендации от фоновой;
параллельно продвигать формирование мысли и её выражение по фрагментам.

\begin{figure}
\hypertarget{fig:6-11}{%
\centering
\includegraphics{images/revised/fig6-11.pdf}
\caption{Рисунок 6-11. Архитектура быстрого и медленного мышления и
сравнение подходов}\label{fig:6-11}
}
\end{figure}

\hypertarget{ux43fux43eux434ux445ux43eux434-1-ux431ux44bux441ux442ux440ux43eux435-ux43cux44bux448ux43bux435ux43dux438ux435-ux440ux435ux430ux433ux438ux440ux443ux435ux442-ux43cux435ux434ux43bux435ux43dux43dux43eux435-ux43eux442ux432ux435ux447ux430ux435ux442}{%
\subsubsection{Подход 1: быстрое мышление реагирует, медленное
отвечает}\label{ux43fux43eux434ux445ux43eux434-1-ux431ux44bux441ux442ux440ux43eux435-ux43cux44bux448ux43bux435ux43dux438ux435-ux440ux435ux430ux433ux438ux440ux443ux435ux442-ux43cux435ux434ux43bux435ux43dux43dux43eux435-ux43eux442ux432ux435ux447ux430ux435ux442}}

Быстрая модель сначала кратко реагирует, медленная выполняет анализ в
фоновом режиме, после чего система даёт подробный ответ. Такая
архитектура позволяет как можно раньше подтвердить получение запроса, но
требует согласовывать два ответа. Например, быстрая модель сначала
рекомендует приобрести определённый тариф, а медленная затем
обнаруживает, что в нём отсутствует критически важная функция:
рекомендации противоречат друг другу. Причина конфликта в том, что обе
стороны отвечают независимо, а основания для решений и стадии
рассуждения могут различаться. Высказанная быстрой моделью рекомендация
уже влияет на решение пользователя, который даже может успеть ей
последовать. Последующее исправление от медленной модели потребует
объяснить изменение и разобраться с уже возникшими последствиями.

Поэтому передняя модель должна различать подтверждение задачи и
бизнес-вывод. Фразу «Я проверю, входит ли эта функция в тариф» можно
произнести сразу, а с рекомендацией о покупке следует дождаться
необходимых запросов и анализа. Когда фоновый результат готов, система
сопоставляет его с уже сказанным и формирует продолжение ответа. Простые
вопросы быстрая модель может обрабатывать самостоятельно, сокращая
дублирование работы.

\hypertarget{ux43fux43eux434ux445ux43eux434-2-ux431ux44bux441ux442ux440ux43eux435-ux43cux44bux448ux43bux435ux43dux438ux435-ux432ux437ux430ux438ux43cux43eux434ux435ux439ux441ux442ux432ux443ux435ux442-ux43cux435ux434ux43bux435ux43dux43dux43eux435-ux43fux43eux434ux441ux43aux430ux437ux44bux432ux430ux435ux442}{%
\subsubsection{Подход 2: быстрое мышление взаимодействует, медленное
подсказывает}\label{ux43fux43eux434ux445ux43eux434-2-ux431ux44bux441ux442ux440ux43eux435-ux43cux44bux448ux43bux435ux43dux438ux435-ux432ux437ux430ux438ux43cux43eux434ux435ux439ux441ux442ux432ux443ux435ux442-ux43cux435ux434ux43bux435ux43dux43dux43eux435-ux43fux43eux434ux441ux43aux430ux437ux44bux432ux430ux435ux442}}

Голос и экран также могут разделять функции. На слух пользователь сразу
поймёт главное из фразы: «В этом месяце в счёте появились дополнительные
15 евро за роуминг». Исходную сумму, новую стоимость и подробности
спорных начислений можно оставить на экране, чтобы пользователь сверил
их в удобном темпе. Конкретную операцию можно подтвердить нажатием, а
конфиденциальные данные пользователь самостоятельно вводит в одобренном
продавцом интерфейсе. Передней модели достаточно получить состояние
наподобие «проверка завершена», чтобы естественно продолжить разговор.

Предположим, путешественник просит Pine забронировать билет в Шанхай на
утро четверга с бюджетом 1200 юаней, чтобы успеть на место встречи до
полудня. Разговаривая со службой поддержки авиакомпании, Pine
одновременно передаёт фоновой модели «записку»: «Если прилететь в Хунцяо
в десять утра, можно ли успеть из аэропорта на место встречи?» Проверив
наличие билетов, сотрудник поддержки сообщает Pine: «Билеты на рейс с
прибытием в Хунцяо в десять часов распроданы. Есть другой рейс, который
прибывает в Пудун в половине одиннадцатого; билет стоит 1250 юаней».
Если фоновая модель видит только переданную передней моделью записку,
она продолжает рассчитывать маршрут из Хунцяо. Передней модели придётся
дать новое указание: «Пересчитай для Пудуна: посадка в половине
одиннадцатого, место и требуемое время прибытия не изменились».

Другой вариант --- последовательно синхронизировать с фоновой моделью и
расшифровку разговора. Фоновая модель уже прочитала в ней названные
сотрудником поддержки аэропорт и время посадки, поэтому для
делегирования той же задачи передней модели достаточно сказать:
«Рассчитай, можно ли успеть этим рейсом». В синхронизируемой записи
необходимо отмечать, кому принадлежат реплики. Недосказанную из-за
прерывания фразу передней модели также следует отличать от уже
произнесённого. В расшифровке важные сведения могут быть пропущены или
записаны неверно. Имена и числа, которые передняя модель уверенно
расслышала, а также исправления пользователя следует явно указать в
задании. При расхождении с расшифровкой фоновая модель руководствуется
заданием. Завершив поиск, фоновая модель должна лишь сообщить передней,
во сколько путешественник предположительно доберётся до места встречи и
насколько жёстким будет график. Проверенные по пути маршруты и поля,
возвращённые картографическим интерфейсом, можно сохранить в фоновой
системе для справки. Передняя модель использует результат, чтобы
продолжить разговор со службой поддержки.

Деловой путешественник из первой главы хотел уложиться в 1200 юаней. В
предыдущих поездках он был готов немного доплатить за подходящий рейс, а
на этот раз не установил иной предельной суммы. Услышав, что минимальная
цена составляет 1250 юаней, ассистент может забронировать билет от его
имени. Но если речь идёт о бережливом пользователе, отправляющемся в
семейную поездку, о дополнительных 50 юанях следует сначала спросить.
Поэтому передней модели приходится одновременно вести два диалога:

\begin{itemize}
\tightlist
\item
  Сообщить пользователю: «Минимальная цена этого рейса --- 1250 юаней,
  на 50 юаней выше бюджета. Вы согласны его забронировать?»
\item
  Затем спросить сотрудника поддержки: «Можно ненадолго зафиксировать
  эту цену? Есть ли в тот же день более поздний рейс стоимостью до 1200
  юаней?» Пока сотрудник проверяет билеты, передняя модель может также
  уточнить норму провоза багажа и условия возврата и обмена.
\end{itemize}

Когда пользователь вернётся, он увидит несколько вариантов, а сотрудник
поддержки будет знать, что искать дальше. Если пользователь впоследствии
пояснит: «В прошлый раз я спешил на встречу; в дальнейшем перед
превышением бюджета в командировке тоже сначала спрашивай меня», при
следующем бронировании нужно будет следовать этому новому требованию.

Во время ожидания подтверждения личности ритм разговора меняется. Если
сотрудник поддержки сообщает, что код подтверждения отправлен, передняя
модель может ответить: «Хорошо, пожалуйста, подождите: я попрошу
путешественника посмотреть телефон». Если путешественник напишет, что
код ещё не пришёл, можно уточнить у сотрудника поддержки последние цифры
номера, на который отправлено сообщение, и через какое время его
разрешается отправить повторно. После прохождения проверки в указанном
продавцом канале передняя модель сможет продолжить оформление. В таком
разговоре лучше немного подождать и оформить билет позднее, чем торопить
пользователя с немедленным решением.

В этот момент фоновая модель присылает сообщение: «Исходное ценовое
предложение обновилось». Передняя модель сначала уточняет новую цену у
сотрудника поддержки, а затем сообщает её пользователю. Если
пользователь именно в это время отвечает «Можно бронировать», его ответ
относится к ранее отправленному предложению на 1250 юаней. Передняя
модель должна определить по новой цене и имеющемуся разрешению,
требуется ли повторное подтверждение. Фоновая модель лишь предоставляет
информацию; разговор со службой поддержки и принятие обязательств всегда
остаются за передней моделью. Чтобы распознать устаревшее предложение,
передней модели необходимы указанные в сообщении рейс, условия
применимости цены и время.

Даже для одного и того же фонового сообщения «Цена изменилась,
проверьте» способ продолжить разговор зависит от момента его
поступления. Если сотрудник поддержки только что сказал: «Минимальная
цена составляет\ldots», передняя модель может сначала дослушать сумму.
Если цена уже названа и сотрудник ждёт ответа, уместно уточнить новую
стоимость. Если путешественник только что вмешался в телефонный разговор
со словами «Подождите, я хочу изменить дату», сначала нужно расслышать
новую дату, а затем решить, требуется ли ещё проверять исходный рейс.
Одно и то же «хорошо» во время объяснения условий может лишь показывать,
что собеседник слушает, но после вопроса «Оформить вам билет прямо
сейчас?» его могут воспринять как согласие на покупку. По одной лишь
расшифровке текста такие различия можно не заметить.

\hypertarget{ux43fux43eux434ux445ux43eux434-3-ux441ux43eux433ux43bux430ux441ux43eux432ux430ux43dux43dux43eux435-ux444ux43eux440ux43cux438ux440ux43eux432ux430ux43dux438ux435-ux438-ux432ux44bux440ux430ux436ux435ux43dux438ux435-ux43cux44bux441ux43bux438}{%
\subsubsection{Подход 3: согласованное формирование и выражение
мысли}\label{ux43fux43eux434ux445ux43eux434-3-ux441ux43eux433ux43bux430ux441ux43eux432ux430ux43dux43dux43eux435-ux444ux43eux440ux43cux438ux440ux43eux432ux430ux43dux438ux435-ux438-ux432ux44bux440ux430ux436ux435ux43dux438ux435-ux43cux44bux441ux43bux438}}

Серия Step-Audio предлагает два связанных примера. \textbf{Дистилляция
рассуждения с привязкой к модальности (MGRD)} сосредоточена на
информации, лежащей в основе рассуждения: данные для рассуждения
создаются и отбираются с учётом таких акустических признаков, как высота
голоса, ритм и интонация, а понимание аудио совершенствуется посредством
итеративной дистилляции. На этапе обучения с подкреплением Step-Audio R1
оценивает задания с аудио по двум составляющим: 0,8 балла за правильный
ответ и ещё 0,2 балла за приведённый процесс рассуждения. Вторая
составляющая побуждает модель указывать основания ответа; использует ли
рассуждение акустические признаки записи, проверяется на предшествующем
этапе отбора примеров.\footnote{StepFun, \emph{Step-Audio-R1 Technical
  Report} --- описание дистилляции рассуждения с привязкой к
  модальности.}

\textbf{Mind-Paced Speaking (MPS)} рассматривает временное согласование
рассуждения и выражения. В этом методе используются модуль формирования
мысли (Formulation Brain) и модуль выражения (Articulation Brain),
причём процесс формирования мысли направляет генерацию речи. Этот метод
применяется в решении реального времени Step-Audio R1.1.\footnote{StepFun,
  \emph{Mind-Paced Speaking: A Dual-Brain Approach to Real-Time
  Reasoning in Spoken Language Models} и карточка модели Step-Audio-R1.1
  --- описание модулей формирования и выражения мысли и их применения в
  системах реального времени.}

Два модуля обмениваются информацией по мере работы: формирование мысли
непрерывно продвигается, а модуль выражения организует речь на основе
уже доступного содержания и предшествующего вывода. Это позволяет
частично совмещать выражение с последующим рассуждением, одновременно
отслеживая уже сказанное.

\hypertarget{ux440ux430ux437ux434ux435ux43bux435ux43dux438ux435-ux444ux443ux43dux43aux446ux438ux439-ux43cux435ux436ux434ux443-ux43cux43eux434ux443ux43bux44fux43cux438-ux438-ux441ux43fux43eux441ux43eux431-ux43eux431ux43dux43eux432ux43bux435ux43dux438ux44f}{%
\subsubsection{Разделение функций между модулями и способ
обновления}\label{ux440ux430ux437ux434ux435ux43bux435ux43dux438ux435-ux444ux443ux43dux43aux446ux438ux439-ux43cux435ux436ux434ux443-ux43cux43eux434ux443ux43bux44fux43cux438-ux438-ux441ux43fux43eux441ux43eux431-ux43eux431ux43dux43eux432ux43bux435ux43dux438ux44f}}

\textbf{Автор считает важным инженерным преимуществом разделения
быстрого и медленного мышления возможность независимо развивать
интерактивность и способность к рассуждению.} После выпуска более
сильной модели рассуждения система может заменить фоновую модель.
Благодаря этому фоновая часть следует за развитием моделей рассуждения,
а передняя продолжает совершенствовать слушание, реагирование и
обработку прерываний. Передняя модель должна понимать задачу и
дополнения пользователя, тогда как фоновая сосредоточена на более
длительном рассуждении, планировании и выполнении.

\textbf{Когда модели рассуждения быстро развиваются, а команде
необходимо непрерывно совершенствовать бизнес-возможности, автор отдаёт
предпочтение заменяемой архитектуре с разделением быстрого и медленного
мышления.} Такое разделение устраняет необходимость переобучать всю
голосовую систему при каждом обновлении модели рассуждения. Более тесную
координацию следует применять тогда, когда промежуточное рассуждение
должно непосредственнее управлять выражением и для этого имеются
необходимые ресурсы обучения.

Голосовой агент Pine AI использует архитектуру, сочетающую быстрое
мышление передней модели с медленным мышлением фоновой. В опубликованной
компанией статье сообщается о первом месте в голосовом треке τ²-bench и
описывается сочетание непрерывного разговора, бизнес-рассуждения и
работы с инструментами\footnote{Pine AI, \emph{The Most Natural
  Human-Computer Interface Is Your Voice} --- описание архитектуры
  быстрого и медленного мышления и представленных компанией результатов
  оценки голосовых задач.}.

Для понимания таких сочетаний можно выделить три измерения архитектуры:
\textbf{аудиотракт} определяет, обрабатывает ли модель звук
непосредственно или передаёт данные через расшифровку; \textbf{временная
организация слушания и речи} определяет, выполняются ли ввод и вывод по
очереди или параллельно; \textbf{когнитивное разделение функций}
определяет, как взаимодействие в реальном времени и глубокое рассуждение
совместно используют состояние и обмениваются результатами. Система
может непосредственно обрабатывать аудио и сохранять полный дуплекс,
одновременно передавая сложные задачи отдельной фоновой модели. Подобные
сочетания демонстрируют модель взаимодействия Thinking Machines Lab и
GPT-Live.

\hypertarget{ux431ux43eux43bux435ux435-ux447ux435ux43bux43eux432ux435ux447ux43dux44bux439-ux441ux438ux43dux442ux435ux437-ux440ux435ux447ux438}{%
\subsection{Более человечный синтез
речи}\label{ux431ux43eux43bux435ux435-ux447ux435ux43bux43eux432ux435ux447ux43dux44bux439-ux441ux438ux43dux442ux435ux437-ux440ux435ux447ux438}}

Паузы, темп и логическое ударение помогают слушателю понять структуру
предложения и намерение говорящего. Например, при сообщении номера
телефона цифры следует правильно группировать, при объяснении сложных
условий можно замедлить речь, а во время ожидания результатов запроса
--- сначала сообщить, что именно сейчас обрабатывается. По мнению
автора, уместные паузы и изменения ритма делают речь естественнее и
позволяют пользователю понимать, как продвигается выполнение задачи.
Управление речью можно проектировать с учётом этих конкретных задач.

Основная LLM может добавлять в текст управляющие метки, например
\texttt{{[}THINKING{]}}, \texttt{{[}EMO:happy{]}} и
\texttt{{[}SPEED:0.8x{]}}. Исполнительный слой разбирает метки,
разделяет текст на фрагменты речи и паузы, а затем сопоставляет их с
поддерживаемыми интерфейсом синтеза эмоциями, стилями или эталонными
аудиозаписями. Значение меток определяется приложением: \texttt{0.8x}
может соответствовать медленному темпу; паузы может вставлять
проигрыватель, а смех и вздохи можно выражать с помощью поддерживаемых
звуковых меток или средств генерации.

При использовании эталонных аудиозаписей различные варианты речи одного
и того же говорящего объединяются в библиотеку. Для каждого фрагмента
текста выбирается соответствующее сочетание эмоции, темпа и стиля, после
чего синтезированные фрагменты последовательно объединяются. Для
управляющих состояний необходимо определить область действия: например,
эмоция сохраняется до следующего изменения, а пауза действует только в
текущей позиции.

\begin{quote}
\textbf{Эксперимент 6-8 ★★: TTS на основе Fish Audio, управляемый
метками}

В этом эксперименте с помощью Fish Audio S1 создаётся голосовая
библиотека из 24 эталонных аудиозаписей, сочетающих четыре эмоции, три
варианта темпа и два стиля. Парсер преобразует ответ с метками во
фрагменты речи и тишины, а исполнительный слой выбирает для каждого
фрагмента эталонную аудиозапись и выполняет синтез.

Сравните три конфигурации: без управляющих меток, с одной эталонной
аудиозаписью и с несколькими эталонными аудиозаписями. Сохраняя
содержимое ответа неизменным, проверьте, обеспечивают ли метки
маршрутизацию в ожидаемую конфигурацию, а затем оцените на слух
изменения пауз, эмоций и тембра.

В сохранённом тесте на слух используется модель понимания аудио. Порядок
представления трёх конфигураций чередуется, а оценка учитывает
естественность, соответствие эмоции, паузы для размышления,
согласованность тембра и впечатление от общения со службой поддержки.
Группа с несколькими эталонными аудиозаписями получила наивысшую
совокупную оценку, а группа с одной эталонной аудиозаписью уступила
группе без меток.
\end{quote}

\hypertarget{computer-use-ai-ux430ux433ux435ux43dux442-ux430ux432ux442ux43eux43cux430ux442ux438ux437ux430ux446ux438ux438-gui}{%
\section{Computer Use: AI-агент автоматизации
GUI}\label{computer-use-ai-ux430ux433ux435ux43dux442-ux430ux432ux442ux43eux43cux430ux442ux438ux437ux430ux446ux438ux438-gui}}

Предположим, вы просите агента изменить адрес доставки на веб-странице.
Сначала ему нужно найти поле адреса, ввести новый адрес, а затем по
отклику интерфейса убедиться, что результат сохранён. Если перед щелчком
появится окно входа, то даже при неизменных координатах кнопки на
предыдущем снимке экрана щелчок уже придётся по другому объекту.
Computer Use управляет браузером, электронной таблицей или системными
настройками с помощью экрана, мыши и клавиатуры, а его основная задача
--- согласовывать каждое действие с текущим состоянием интерфейса.

Проследим по циклу \textbf{восприятие --- размышление --- действие} на
\crossreflink{fig:6-12}{рисунке 6-12}, что модель должна проверить перед
вводом нового адреса. На основе снимка экрана и хода выполнения задачи
она выбирает место щелчка, Harness проверяет параметры и разрешения,
после чего передаёт действие исполнительному слою. Успешный щелчок
означает лишь то, что соответствующее действие было отправлено.
Безопасность последующего ввода адреса определяется тем, где окажется
фокус ввода на следующем снимке экрана.

\begin{figure}
\hypertarget{fig:6-12}{%
\centering
\includegraphics{images/revised/fig6-12.pdf}
\caption{Рисунок 6-12. Цикл «восприятие --- размышление --- действие»
агента Computer Use}\label{fig:6-12}
}
\end{figure}

Как после заполнения адреса ассистент должен сообщить исполнителю, по
какой именно кнопке «Сохранить» нужно щёлкнуть? Цель можно указать как
экранными координатами, так и номером элемента, однако эти способы
предъявляют разные требования к модели и инструменту.

\hypertarget{ux43fux440ux43eux435ux43aux442ux438ux440ux43eux432ux430ux43dux438ux435-ux43fux440ux43eux441ux442ux440ux430ux43dux441ux442ux432ux430-ux434ux435ux439ux441ux442ux432ux438ux439}{%
\subsection{Проектирование пространства
действий}\label{ux43fux440ux43eux435ux43aux442ux438ux440ux43eux432ux430ux43dux438ux435-ux43fux440ux43eux441ux442ux440ux430ux43dux441ux442ux432ux430-ux434ux435ux439ux441ux442ux432ux438ux439}}

В эталонной реализации Computer Use от Anthropic на
\crossreflink{fig:6-13}{рисунке 6-13} предусмотрены три категории
инструментов: графический интерфейс, выполнение команд и редактирование
файлов. Даже у одной лишь категории графического интерфейса данные
направляются по двум путям: действие щелчка передаётся исполнителю, а
полученный после него снимок экрана возвращается модели как наблюдение
для принятия решения на следующем шаге. Формат инструментов задаётся
нативным API; при использовании других моделей адаптационный слой должен
преобразовывать снимки экрана, запросы действий и результаты выполнения
в сообщения, принимаемые соответствующим интерфейсом.\footnote{Anthropic,
  документация и эталонная реализация инструмента Computer Use с
  описанием версионируемого интерфейса действий, исполнительного слоя и
  обработки скриншотов.}

\begin{figure}
\hypertarget{fig:6-13}{%
\centering
\includegraphics{images/revised/fig6-13.pdf}
\caption{Рисунок 6-13. Пространство действий Computer
Use}\label{fig:6-13}
}
\end{figure}

\textbf{Инструменты управления GUI.} Действия мыши включают перемещение,
щелчок левой, правой или средней кнопкой, двойной и тройной щелчок, а
также перетаскивание. Более детализированные интерфейсы позволяют
отдельно задавать нажатие и отпускание кнопки. Для прокрутки необходимо
указать направление и величину; некоторые реализации также поддерживают
клавиши-модификаторы. Клавиатурный интерфейс позволяет вводить текст,
использовать сочетания клавиш и удерживать клавиши. Снимки экрана и
положение указателя служат для получения состояния, а интерфейс ожидания
даёт время на загрузку страницы или выполнение анимации. Конкретный
набор действий определяется версией инструмента, а Harness проверяет
параметры в соответствии с используемым интерфейсом.

\textbf{Инструменты выполнения команд.} Постоянный терминал позволяет
сохранять рабочий каталог и состояние окружения между вызовами.
Например, если один вызов изменил каталог, следующий продолжит работу в
нём. Эталонная реализация может распознавать завершение команды по
строке-маркеру и задавать тайм-аут; после его истечения исполнитель всё
равно должен проверить состояние процесса.

\textbf{Инструменты редактирования файлов.} Интерфейс замены строк
определяет изменяемый диапазон по исходному тексту; также могут быть
доступны просмотр, создание, вставка и отмена изменений. Перед
выполнением проверяются позиции и количество совпадений, а после него
--- различия. Если файл изменился между двумя вызовами, сначала следует
заново прочитать соответствующий фрагмент. Здесь можно повторно
использовать стратегии редактирования, представленные в пятой главе.

\begin{quote}
\textbf{Эксперимент 6-9 ★: запуск эталонной реализации Anthropic
Computer Use}

В этом эксперименте используются рабочий стол Ubuntu, браузер и слой
выполнения инструментов из официального примера с контейнером. После
получения задачи от пользователя модель анализирует снимок экрана и
возвращает действие, слой выполнения управляет рабочим столом, а затем
передаёт модели новый снимок.

В сопутствующем задании необходимо узнать погоду в Сан-Франциско.
Сохранённая траектория начинается с обращения к Google, но после
появления CAPTCHA в соответствии с инструкцией по восстановлению во
время выполнения переключается на общедоступные данные Open-Meteo и в
итоге сообщает температуру и погодные условия. Запись содержит 15
действий на рабочем столе и соответствующие снимки экрана, позволяя шаг
за шагом проследить, как щелчки, нажатия клавиш, ввод текста и ожидание
изменяют страницу.

Способ управления браузером с помощью открытой модели представлен в
эксперименте 6-10.
\end{quote}

\hypertarget{ux432ux438ux437ux443ux430ux43bux44cux43dux430ux44f-ux43bux43eux43aux430ux43bux438ux437ux430ux446ux438ux44f-grounding}{%
\subsection{Визуальная локализация
(Grounding)}\label{ux432ux438ux437ux443ux430ux43bux44cux43dux430ux44f-ux43bux43eux43aux430ux43bux438ux437ux430ux446ux438ux44f-grounding}}

Задачу щёлкнуть по полю поиска можно поставить перед моделью в двух
формах. В первом случае ей показывают скриншот и просят указать
координаты клика. Во втором --- нумеруют возможные цели и предлагают
выбрать поле поиска; номера могут быть нанесены на области-кандидаты
непосредственно на скриншоте либо взяты из предоставленного браузером
списка элементов. Ниже по отдельности рассмотрены три соответствующих
подхода. При выборе номера исполнительный слой преобразует его в
фактическую цель. При непосредственном указании координат модель должна
дополнительно оценить положение цели, что предъявляет более высокие
требования к обучению локализации и адаптации к разрешению экрана.

\textbf{Set-of-Mark: визуальная разметка.}

Set-of-Mark (SoM) накладывает на изображение доступные для указания
номера, связывая визуальные области с выводом модели. В исходной работе
области-кандидаты создавались методами сегментации: например, области
выделялись с помощью SAM или SEEM, после чего на них размещались метки.
Когда модель выбирает номер, исполнительный слой определяет место
действия по соответствующей области. Этот метод применим к интерфейсам
настольных приложений и игр, где отсутствует DOM. Необходимый элемент
управления должен попасть в пронумерованную область, а модель ---
правильно распознать соответствующий ей номер.\footnote{Microsoft,
  \emph{Set-of-Mark Prompting}, 2023; реализация DOM-сервиса browser-use
  служит примером обработки структурированных элементов.}

\textbf{Индекс структурированных элементов.}

Браузер может предоставить и другой тип ориентиров. Объектная модель
документа (DOM) описывает элементы страницы и их иерархию, данные о
компоновке --- отображаемое положение элементов, а дерево доступности
(accessibility tree) сохраняет роли и доступные имена, необходимые
вспомогательным технологиям. Например, если кнопке, содержащей только
значок лупы, на странице присвоено доступное имя «Поиск», AI-агент
получает возможность найти её по роли и имени. В решениях наподобие
browser-use система может получать сведения об элементах через
интерфейсы отладки браузера, такие как Chrome DevTools Protocol (CDP),
распознавать поля ввода, кнопки и ссылки, а затем нумеровать
интерактивные элементы текущей страницы. На
\crossreflink{fig:6-14}{рисунке 6-14} модель выбрала элемент {[}2{]} в
правом верхнем углу страницы. Как этот номер преобразуется в клик по
кнопке Submit?

\begin{figure}
\hypertarget{fig:6-14}{%
\centering
\includegraphics{images/revised/fig6-14.pdf}
\caption{Рисунок 6-14. Set-of-Mark и индекс структурированных элементов
(реализация browser-use)}\label{fig:6-14}
}
\end{figure}

Процесс можно разделить на четыре этапа: прочитать DOM, данные о
компоновке и доступные сведения из дерева доступности; отфильтровать
интерактивные элементы; создать для текущей страницы отображение
номеров; передать модели описания элементов и при необходимости
размеченный скриншот. Например, модель может увидеть следующие четыре
пронумерованных элемента:

\begin{Shaded}
\begin{Highlighting}[]
\NormalTok{На скриншоте отмечены [1], [2], [3], [4]}

\NormalTok{[1] \textless{}input type="text" placeholder="Search" aria{-}label="Search" /\textgreater{}}
\NormalTok{[2] \textless{}button id="submit{-}btn" aria{-}label="Submit form" /\textgreater{}}
\NormalTok{[3] \textless{}input type="text" placeholder="Enter your name" value="" /\textgreater{}}
\NormalTok{[4] \textless{}a href="/docs" aria{-}label="Documentation" /\textgreater{}}
\end{Highlighting}
\end{Shaded}

После выбора номера моделью исполнительный слой преобразует его в ссылку
на элемент или координаты, проверяет, что цель по-прежнему видима и
интерактивна, а затем выполняет действие. После обновления или прокрутки
страницы либо динамической вставки элементов отображение необходимо
обновить, чтобы не использовать устаревшие номера. Структурированный
индекс позволяет задействовать метки, роли и текст, тогда как Canvas,
перекрытия и неполная семантическая разметка требуют дополнительного
анализа скриншота. Чем больше перечислено элементов и чем крупнее
скриншот, тем больше информации передаётся в контекст.

\textbf{Непосредственное предсказание координат.}

При координатном подходе модель получает скриншот и описание цели, а в
ответ непосредственно указывает позицию: например, сопоставляет команду
«нажать кнопку отправки» с определённой точкой. Координатный вывод
применяется в исследованиях локализации в GUI, таких как SeeClick, а
также в Claude Computer Use. Исполнительный слой преобразует координаты
изображения в фактические координаты рабочего стола.

После масштабирования скриншота обе системы координат должны
использовать одно и то же преобразование. Например, если рабочий стол
имеет разрешение 2560×1440, а модели передано изображение размером
1366×768, указанные моделью координаты (683, 384) соответствуют:

\[
x = 683 \times \frac{2560}{1366} = 1280,\qquad
y = 384 \times \frac{1440}{768} = 720
\]

На \crossreflink{fig:6-15}{рисунке 6-15} соотношения сторон для двух
размеров близки. Если для сохранения пропорций к скриншоту добавлены
чёрные поля, указанные моделью координаты будут учитывать эти пустые
области. Перед обратным преобразованием в координаты исходного экрана
необходимо сначала вычесть соответствующее смещение полей, а затем
применить коэффициент масштабирования. Обрезка, масштабирование дисплея
и начала координат нескольких мониторов также изменяют отображение,
поэтому инструментальный слой должен сохранять фактические параметры
преобразования.

\begin{figure}
\hypertarget{fig:6-15}{%
\centering
\includegraphics{images/revised/fig6-15.pdf}
\caption{Рисунок 6-15. Согласование разрешений и двунаправленное
масштабирование координат}\label{fig:6-15}
}
\end{figure}

В эталонной реализации использовались, среди прочего, целевые разрешения
XGA (1024×768), WXGA (1280×800) и FWXGA (1366×768). Принудительное
изменение соотношения сторон искажает форму элементов управления, а
неверное обратное преобразование переносит правильно предсказанную точку
в ошибочное место.\footnote{Anthropic, документация и эталонная
  реализация инструмента Computer Use с описанием версионируемого
  интерфейса действий, исполнительного слоя и обработки скриншотов.}

\textbf{При наличии полной структурированной информации автор отдаёт
предпочтение индексированию по DOM или дереву доступности.} Имена, роли
и ссылки на элементы можно непосредственно использовать для локализации,
уменьшая объём работы модели по ситуативной оценке положения в пикселях.
При отсутствии структуры универсальную визуальную модель можно сначала
использовать в сочетании с SoM, чтобы она выбирала одну из размеченных
областей. Модель, уже обученная локализации в GUI, может предсказывать
координаты напрямую, без отдельной разметки.

\begin{quote}
\textbf{Эксперимент 6-10 ★: автоматизация работы в браузере с помощью
browser-use}

В этом эксперименте Qwen3-VL 32B Instruct управляет browser-use и
генерирует действия на основе состояния браузера, сведений об элементах
и скриншотов. Задача состоит в поиске погоды в Сан-Франциско: найти поле
поиска, ввести «San Francisco weather today», отправить запрос, а затем
извлечь температуру и погодные условия.

Код и инструкции по запуску эксперимента 6-10 предусматривают два
варианта. При наличии ключа OpenRouter можно использовать размещённую у
провайдера визуальную модель. Если уже доступен совместимый сервис,
например vLLM или SGLang, переменные \texttt{OPEN\_MODEL\_BASE\_URL} и
\texttt{OPEN\_MODEL\_MODEL} можно направить на собственную конечную
точку и модель, а ключ задать в соответствии с требованиями сервиса. В
обоих вариантах сначала выполните в каталоге эксперимента команду
\texttt{python\ main.py\ -\/-dry-run}, чтобы просмотреть конфигурацию, а
затем запустите фактический запрос согласно инструкции.

Сохранённый запуск выполнил 16 шагов взаимодействия и сохранил 15
скриншотов. После появления CAPTCHA в Google для получения результата
задача перешла на weather.com.
\end{quote}

\hypertarget{ux430ux433ux435ux43dux442-computer-use-ux441ux43fux43eux441ux43eux431ux43dux44bux439-ux432ux438ux434ux435ux442ux44c-ux430ux43dux438ux43cux430ux446ux438ux44e-ux438-ux441ux43bux44bux448ux430ux442ux44c-ux437ux432ux443ux43a}{%
\subsection{Агент Computer Use, способный видеть анимацию и слышать
звук}\label{ux430ux433ux435ux43dux442-computer-use-ux441ux43fux43eux441ux43eux431ux43dux44bux439-ux432ux438ux434ux435ux442ux44c-ux430ux43dux438ux43cux430ux446ux438ux44e-ux438-ux441ux43bux44bux448ux430ux442ux44c-ux437ux432ux443ux43a}}

Лишь на короткое время после сохранения страница показывает сообщение
«Сбой синхронизации», однако на двух снимках экрана, сделанных с
интервалом в несколько секунд, отображается одна и та же форма.
Просматривая только эти два изображения, модель, скорее всего, не сможет
восстановить, что происходило между ними. Та же проблема возникает с
видео, анимацией, кратковременными уведомлениями и речью участников
встречи. \textbf{Для непрерывно изменяющихся интерфейсов автор
предлагает сосредоточить улучшения на интерфейсе наблюдения.} Чтобы
выбирать следующий шаг с учётом изменений, необходимо сохранять также
события и изображения, появлявшиеся между действиями.

Предложенный автором и его соавтором \textbf{Agent--Computer Observation
Interface (AOI, интерфейс наблюдения «агент--компьютер»)} преобразует
непрерывный поток изображения с экрана и звука в изображения и текст,
которые может принимать модель\footnote{Bojie Li, Noah Shi,
  \emph{Agent-Computer Observation Interfaces Enable Dynamic Computer
  Use}, 2026. В работе представлены AOI и результаты абляционного
  исследования его компонентов.}. Интерфейс включает три компонента:

\begin{itemize}
\tightlist
\item
  \textbf{Захват ключевых кадров между действиями.} Во время работы
  основного цикла отслеживаются изменения изображения; различия между
  пикселями или лёгкая визуальная модель используются для отбора
  кандидатных кадров, после чего полезные изображения передаются для
  принятия следующего решения.
\item
  \textbf{Транскрибирование с управлением по уровню громкости.} При
  обнаружении звука вызывается распознавание речи, а транскрипция вместе
  с временной информацией добавляется в журнал наблюдений, помогая
  агенту использовать содержание встречи или устные инструкции.
\item
  \textbf{Сохраняемое описание изображения.} Ключевое изображение кратко
  описывается текстом; после удаления исходного изображения из контекста
  текст продолжает хранить сведения о произошедших изменениях, чтобы к
  ним можно было обращаться на последующих шагах.
\end{itemize}

Например, в видео показана операция с меню, и модель, получив
промежуточные кадры и соответствующие описания, сможет обращаться к
этому шагу и после завершения воспроизведения. Результаты абляционного
исследования в статье показывают, что особенно важны сохраняемые
текстовые описания, тогда как тщательный отбор ключевых кадров даёт
меньший вклад; некоторые модели при добавлении слишком большого числа
изображений, напротив, работают хуже.

Вернёмся к примеру с сохранением формы. Даже если на ключевом кадре
запечатлено сообщение «Сбой синхронизации», после удаления исходного
изображения из контекста модель сможет опираться только на оставшееся
описание. Фраза «На странице появилось уведомление» может означать как
успешное сохранение, так и сбой; лишь описание «Сбой синхронизации,
адрес не сохранён» сохраняет информацию, необходимую для выбора
следующего шага.

\hypertarget{ux43cux43eux434ux435ux43bux44c-ux43cux438ux440ux430-ux434ux43bux44f-computer-use}{%
\subsection{Модель мира для Computer
Use}\label{ux43cux43eux434ux435ux43bux44c-ux43cux438ux440ux430-ux434ux43bux44f-computer-use}}

Своевременно получив наблюдения, AI-агент должен также выбрать действие
и дождаться его выполнения. В исследовании OSWorld-Human эффективность
работы агента сопоставляется с траекториями действий человека и
анализируются такие источники временных затрат, как избыточные шаги и
вызовы модели\footnote{Reyna Abhyankar и др., \emph{OSWorld-Human:
  Benchmarking the Efficiency of Computer-Use Agents} --- анализ
  траекторий выполнения задач и временных затрат.}.

Модель мира позволяет при планировании прогнозировать последствия
действий. Состояние рабочего стола может включать окна, фокус, положение
прокрутки, введённые данные и ход загрузки; сведения о разрешениях,
результатах сетевых запросов и другие данные необходимо получать через
соответствующие интерфейсы. На основании доступного текущего состояния и
возможных действий модель прогнозирует следующее состояние:

\begin{Shaded}
\begin{Highlighting}[]
\NormalTok{Текущее состояние + click / type / scroll / wait → прогнозируемое следующее состояние}
\end{Highlighting}
\end{Shaded}

Работая за компьютером, человек зачастую ещё до щелчка представляет себе
следующий шаг: если интерфейс соответствует ожиданиям, он продолжает
работу, а при неожиданном результате останавливается и проверяет
состояние. Модель мира может предоставить агенту аналогичную основу для
спекулятивного выполнения: планировщик заранее сравнивает возможные
действия, во время загрузки страницы подготавливает последующие шаги и
продолжает выполнение после подтверждения фактического состояния,
сокращая задержку из-за повторного планирования с нуля после каждого
щелчка. Например, при создании в VS Code нового файла Python и записи в
него hello world ожидаемые изменения включают появление нового файла в
дереве файлов, наличие кода в области редактирования и обновление
состояния сохранения. После выполнения проверяются эти изменения: если
они соответствуют ожиданиям, работа продолжается, а при отклонении
заново определяются положение элементов и план действий.

Для таких операций, как удаление файла, также требуются отдельные
условия выполнения: правильно ли выбран объект, разрешено ли удаление,
можно ли восстановить файл и требуется ли подтверждение пользователя.
Прогноз помогает понять возможные результаты, но фактические разрешения
и бизнес-условия проверяются на уровне выполнения.

\textbf{Прогнозирование латентного состояния в Photon-1.} В июле 2026
года Induction Labs представила Photon-1: видеокадры кодируются в
дискретные латентные токены, а изменения состояния изучаются посредством
прогнозирования последующих токенов. По данным компании, для
предварительного обучения было использовано около 30 тысяч часов работы
GPU H200. Состояние прогнозируется в пространстве латентных токенов, а
вспомогательный генератор изображений преобразует их обратно в доступные
для просмотра кадры, чтобы человек мог проверить прогноз. При
планировании и генерации действий можно напрямую использовать внутреннее
представление, избегая затрат на создание полного снимка экрана на
каждом шаге.\footnote{David Li, Jonathan Li, \emph{Scaling Video
  Pretraining with Imagination Models}, Induction Labs, 2026-07-23 ---
  описание предварительного обучения Photon-1, обучения действиям и
  симуляции состояний.}

Далее в этом подходе посредством дообучения осваиваются форматы
инструкций и действий: сначала прогнозируется целевое состояние, затем
генерируются ведущие к нему операции, а качество выполнения задач
повышается с помощью онлайн-обучения в виртуальной машине. Дообучение с
обусловливанием на действия, в свою очередь, используется для
прогнозирования изменений рабочего стола после заданной операции. Эти
два варианта применения связывают представление состояния соответственно
с планированием и симуляцией, показывая, как оно участвует в Computer
Use.

\hypertarget{ux43cux43eux431ux438ux43bux44cux43dux44bux435-ux443ux441ux442ux440ux43eux439ux441ux442ux432ux430-ux442ux435ux445ux43dux438ux447ux435ux441ux43aux438ux435-ux438ux43dux442ux435ux440ux444ux435ux439ux441ux44b-ux438-ux432ux437ux430ux438ux43cux43eux434ux435ux439ux441ux442ux432ux438ux435-ux432-ux44dux43aux43eux441ux438ux441ux442ux435ux43cux435}{%
\subsection{Мобильные устройства: технические интерфейсы и
взаимодействие в
экосистеме}\label{ux43cux43eux431ux438ux43bux44cux43dux44bux435-ux443ux441ux442ux440ux43eux439ux441ux442ux432ux430-ux442ux435ux445ux43dux438ux447ux435ux441ux43aux438ux435-ux438ux43dux442ux435ux440ux444ux435ux439ux441ux44b-ux438-ux432ux437ux430ux438ux43cux43eux434ux435ux439ux441ux442ux432ux438ux435-ux432-ux44dux43aux43eux441ux438ux441ux442ux435ux43cux435}}

На мобильных устройствах по-прежнему применяется цикл «наблюдение ---
принятие решения --- действие», но действия чаще выражаются в виде
жестов касания. Одна и та же координата может обозначать место
одиночного нажатия, место долгого нажатия или начальную точку свайпа,
поэтому интерфейсу также требуются тип жеста, его продолжительность и
траектория. Поворот экрана, появление экранной клавиатуры и системных
диалоговых окон тоже изменяют доступную для взаимодействия область.

Android \texttt{AccessibilityService} позволяет в пределах
предоставленных разрешений считывать сведения об интерфейсе и выполнять
соответствующие действия; для интерфейса жестов также необходимо заявить
соответствующую возможность\footnote{Android Developers,
  AccessibilityService: обзор возможностей доступа к интерфейсу,
  выполнения действий и жестов.}. Снимки экрана, дерево специальных
возможностей и интерфейсы системной автоматизации можно использовать
совместно, а конкретные доступные возможности определяются операционной
системой, реализацией приложения и предоставленными разрешениями. Такие
представленные в \crossreflink{chap:7}{главе 7} бенчмарки, как
AndroidWorld, проверяют на задачах в реальных приложениях, приводят ли
эти действия к ожидаемому результату.

\textbf{Автор считает, что при внедрении на мобильных устройствах
экосистемные барьеры зачастую преодолеть сложнее, чем технические
проблемы.} Платформе необходимо определить, какие операции можно
поручить AI-агенту, как предоставлять доступ к учётным записям, платежам
и данным, а также как отслеживать ответственность после возникновения
ошибки. Эти договорённости напрямую влияют на круг задач, которые может
выполнять агент.

Один из структурных конфликтов интересов связан с трафиком и вниманием
пользователей. Многие приложения получают доход за счёт показа рекламы,
рекомендаций контента и точек входа в транзакции. Когда пользователь
поручает агенту напрямую выполнить поиск или покупку, путь просмотра
страниц может сократиться, а объём внимания, получаемого рекомендациями
и рекламой, --- измениться. Для платформ, монетизирующих эти этапы
взаимодействия, действия агента затрагивают существующие источники
дохода. Поэтому платформам необходимо заново осмыслить открытие
интерфейсов, распределение коммерческих доходов и управление
идентичностью агентов.

\hypertarget{ux443ux43fux440ux430ux432ux43bux435ux43dux438ux435-ux440ux43eux431ux43eux442ux430ux43cux438-ux443ux431ux43eux440ux43aux430-ux441ux442ux43eux43bux430-ux441-ux43fux43eux43cux43eux449ux44cux44e-xlerobot}{%
\section{Управление роботами: уборка стола с помощью
XLeRobot}\label{ux443ux43fux440ux430ux432ux43bux435ux43dux438ux435-ux440ux43eux431ux43eux442ux430ux43cux438-ux443ux431ux43eux440ux43aux430-ux441ux442ux43eux43bux430-ux441-ux43fux43eux43cux43eux449ux44cux44e-xlerobot}}

Поместить красную чашку в поднос, выбросить жёлтую макулатуру в мусорный
контейнер и в конце подтвердить, что стол убран, --- эта задача
объединяет восприятие, планирование и управление. Когда захват
приближается к чашке, её положение на изображении с камеры меняется;
после захвата чашки манипулятор должен обойти другие предметы на столе;
если бумагу не удалось захватить с первой попытки, система должна
обнаружить сбой и повторить попытку. Каждое действие изменяет условия
следующего решения.

В этом разделе задача рассматривается в пяти экспериментах. Эксперименты
6-11 и 6-13 посвящены практическому проектированию телеуправления и
автономного управления настоящим роботом; эксперименты 6-12, 6-14 и 6-15
выполняются в локальной симуляции и исследуют соответственно базовый
уровень управления, восстановление после сбоев и изменение визуальных
условий.

Одним и тем же роботом может управлять человек, непрерывно наблюдающий
за изображением с камеры, либо модель, которая считывает изображение,
выбирает навык и передаёт его системе управления для выполнения. Если
сохранить неизменными задачу и оборудование, различия между двумя
подходами можно анализировать с помощью четырёх вопросов: что хочет
выполнить человек, какую подзадачу следует выполнить следующей, какие
действия выводит текущий навык и соответствует ли среда плану после
выполнения действия. Долгосрочное планирование определяет порядок
обработки чашки и макулатуры, VLA-стратегия или заранее заданные
примитивы действий отвечают за захват и размещение, мировая модель
помогает оценивать последствия действий, а перенос из симуляции в
реальность учитывает изменения изображения, контакта и характеристик
исполнительных механизмов.

\hypertarget{ux440ux430ux437ux434ux435ux43bux435ux43dux438ux435-ux444ux443ux43dux43aux446ux438ux439-ux43eux431ux43eux440ux443ux434ux43eux432ux430ux43dux438ux44f-ux438-ux430ux43bux433ux43eux440ux438ux442ux43cux43eux432}{%
\subsection{Разделение функций оборудования и
алгоритмов}\label{ux440ux430ux437ux434ux435ux43bux435ux43dux438ux435-ux444ux443ux43dux43aux446ux438ux439-ux43eux431ux43eux440ux443ux434ux43eux432ux430ux43dux438ux44f-ux438-ux430ux43bux433ux43eux440ux438ux442ux43cux43eux432}}

Если автономная уборка стола завершается неудачей, сначала с помощью
телеуправления можно проверить, способно ли оборудование выполнить
задачу. Сохраните неизменными камеру, манипулятор, захват, расположение
предметов на столе и критерии приёмки, передав оператору замкнутый
контур восприятия и управления. Если человек стабильно выполняет задачу
в нескольких попытках, а автономная система раз за разом терпит неудачу,
значит, имеющееся оборудование уже способно справиться с текущей
задачей, а узкое место преимущественно находится в алгоритмах и обратной
связи автономного управления. В этом случае в первую очередь следует
проверить локализацию объектов, выбор навыков, момент выполнения
действий и восстановление после сбоев, чтобы автономная система освоила
эффективные корректировки, применяемые при телеуправлении.

XLeRobot поддерживает телеуправление с помощью клавиатуры, контроллера
Xbox и Switch Joy-Con, а управление через VR сначала можно опробовать в
симуляционной среде.\footnote{XLeRobot, документация по Teleop.} При
приближении захвата к чашке оператор замедляет движение; если чашка
скользит, приходится корректировать точку захвата. Если захват не
удержал бумагу, необходимо снова осмотреть сцену; после размещения
объекта также следует проверить результат. Все эти корректировки на
основе наблюдений стоит записывать в демонстрационные данные. При
последующем автономном управлении робот также должен уметь обнаруживать
скольжение и неудачные захваты.

\begin{quote}
\textbf{Эксперимент 6-11 ★: уборка стола посредством телеуправления
настоящим XLeRobot}

Разместите красную чашку, поднос, жёлтую макулатуру и мусорный контейнер
в рабочей области XLeRobot, где выполнены калибровка, проверка
ограничителей и аварийной остановки. С помощью одного из откалиброванных
способов телеуправления выполните фиксированную задачу: «поместить
красную чашку в поднос, выбросить жёлтую макулатуру в мусорный
контейнер, затем подтвердить состояние стола». Повторите эксперимент
несколько раз и запишите изображение с камеры, управляющий ввод
оператора, состояние манипулятора, время выполнения действий, неудачные
захваты и количество повторных попыток.

В конце каждой попытки убедитесь, что чашка находится внутри подноса,
макулатура --- внутри мусорного контейнера, а манипулятор вернулся в
безопасную позу, и проверьте записи о столкновениях, выходе за границы и
вмешательстве человека на протяжении всего эксперимента. Совместный
анализ этапа возникновения сбоя, способа восстановления оператором и
итогового результата позволяет сформировать опорный уровень для
автономного управления.
\end{quote}

Симуляционный эксперимент позволяет многократно изменять количество и
исходное положение объектов и отдельно исследовать процесс управления
при известном состоянии. В 2D-модели стола приближение к объекту, его
перемещение к цели и переход к следующему объекту можно представить как
три этапа, а затраты на управление описать количеством шагов до
завершения и длиной пути.

\begin{quote}
\textbf{Эксперимент 6-12 ★: создание идеального базового уровня
управления для той же задачи в симуляторе}

Случайным образом генерируйте в 2D-пространстве стола положения объектов
и целей, позволяя контроллеру напрямую считывать координаты, поочерёдно
приближаться к объектам, перемещать их к соответствующим целям и
переходить к следующему объекту. Изменяйте исходное положение и
количество объектов, сравнивая долю успешных попыток, количество шагов
до завершения и длину пути.

Сначала проследите полную траекторию одного объекта: после достижения
объекта контроллер переходит к этапу перемещения, а число завершённых
объектов увеличивается только после достижения цели. Затем увеличьте
количество объектов и убедитесь, что для каждого из них выполняются оба
условия достижения. В сопутствующей 2D-симуляции для пяти случайных
начальных значений размещалось по 1, 2, 3 и 4 объекта, каждая комбинация
выполнялась в 512 эпизодах, и все 10 240 эпизодов завершились успешно.
Контроллер может напрямую считывать координаты, а перемещение объектов
не зависит от настоящего захвата.
\end{quote}

\hypertarget{ux431ux430ux437ux43eux432ux430ux44f-ux441ux442ux440ux443ux43aux442ux443ux440ux430-ux443ux43fux440ux430ux432ux43bux435ux43dux438ux44f-ux440ux43eux431ux43eux442ux43eux43c}{%
\subsection{Базовая структура управления
роботом}\label{ux431ux430ux437ux43eux432ux430ux44f-ux441ux442ux440ux443ux43aux442ux443ux440ux430-ux443ux43fux440ux430ux432ux43bux435ux43dux438ux44f-ux440ux43eux431ux43eux442ux43eux43c}}

Одна уборка стола может занимать несколько минут, однако управление
суставами должно обновляться непрерывно и своевременно. Система может
разделить работу по временным масштабам: высокоуровневый планировщик
задаёт порядок выполнения задачи, стратегия навыков генерирует движения,
а низкоуровневый контроллер отслеживает их выполнение:

\begin{longtable}[]{@{}
  >{\raggedright\arraybackslash}p{(\columnwidth - 6\tabcolsep) * \real{0.2314}}
  >{\raggedright\arraybackslash}p{(\columnwidth - 6\tabcolsep) * \real{0.2403}}
  >{\raggedright\arraybackslash}p{(\columnwidth - 6\tabcolsep) * \real{0.2771}}
  >{\raggedright\arraybackslash}p{(\columnwidth - 6\tabcolsep) * \real{0.2512}}@{}}
\toprule\noalign{}
\begin{minipage}[b]{\linewidth}\raggedright
\hspace{0pt}Уровень
\end{minipage} & \begin{minipage}[b]{\linewidth}\raggedright
\hspace{0pt}Основной вопрос
\end{minipage} & \begin{minipage}[b]{\linewidth}\raggedright
\hspace{0pt}Выход
\end{minipage} & \begin{minipage}[b]{\linewidth}\raggedright
\hspace{0pt}Временной масштаб и пример конфигурации
\end{minipage} \\
\midrule\noalign{}
\endhead
\bottomrule\noalign{}
\endlastfoot
\hspace{0pt}Цель задачи & \hspace{0pt}Что хочет выполнить человек &
\hspace{0pt}«Вернуть чашку и макулатуру на свои места» &
\hspace{0pt}Задача длительностью в минуты \\
\hspace{0pt}Долгосрочное планирование & \hspace{0pt}Что делать сначала,
а что потом & \hspace{0pt}Сначала обработать чашку, затем макулатуру и в
конце выполнить проверку & \hspace{0pt}От секунд до минут \\
\hspace{0pt}Базовый навык & \hspace{0pt}Какое изменение состояния
требуется сейчас &
\hspace{0pt}\texttt{pick(red\_}\allowbreak{}\texttt{cup)},
\texttt{place(red\_}\allowbreak{}\texttt{cup,\ tray)} & \hspace{0pt}Один
навык может длиться несколько секунд, например 1--3 секунды \\
\hspace{0pt}VLA /\allowbreak{} стратегия навыка & \hspace{0pt}Как именно
выполнить текущий навык & \hspace{0pt}Действия захвата, поза рабочего
органа или фрагмент траектории суставов & \hspace{0pt}Инференс можно
настроить на частоту 1--10 Гц \\
\hspace{0pt}Низкоуровневое управление и уровень безопасности &
\hspace{0pt}Как обеспечить стабильное и своевременное выполнение &
\hspace{0pt}Управляющие воздействия на суставы или рабочий орган,
ограничение скорости и аварийная остановка & \hspace{0pt}Контроллер
можно настроить на частоту от десятков до тысяч Гц \\
\end{longtable}

Планировщиком может служить программа на основе правил, VLM или
оптимизатор, а VLA также может брать на себя часть выбора навыков.
\textbf{Автор рекомендует разделять порядок выполнения задачи и текущие
действия.} Высокоуровневому мышлению требуется время на составление
плана, а низкоуровневое управление должно своевременно реагировать на
изменения положения и усилия; только работа каждого уровня в подходящем
ему темпе позволяет совместить планирование и управление.

В этом разделе высокоуровневой модели предоставляются пять типов
инструментов: \texttt{observe\_scene}, \texttt{pick}, \texttt{place},
\texttt{verify\_state} и \texttt{stop}, предназначенные соответственно
для наблюдения, захвата, размещения, проверки и остановки. Harness
проверяет параметры, рабочую область и текущее состояние, а исполнитель
преобразует прошедшие проверку навыки в движения манипулятора и
обрабатывает ограничения скорости и тайм-ауты. Углы суставов и
непрерывные траектории генерируются стратегией навыков и контроллером,
благодаря чему между выбором навыка моделью и фактическим движением
формируется чёткий интерфейс.

\hypertarget{ux434ux43eux43bux433ux43eux441ux440ux43eux447ux43dux43eux435-ux43fux43bux430ux43dux438ux440ux43eux432ux430ux43dux438ux435-ux438-ux434ux435ux43aux43eux43cux43fux43eux437ux438ux446ux438ux44f-ux437ux430ux434ux430ux447}{%
\subsection{Долгосрочное планирование и декомпозиция
задач}\label{ux434ux43eux43bux433ux43eux441ux440ux43eux447ux43dux43eux435-ux43fux43bux430ux43dux438ux440ux43eux432ux430ux43dux438ux435-ux438-ux434ux435ux43aux43eux43cux43fux43eux437ux438ux446ux438ux44f-ux437ux430ux434ux430ux447}}

Когда пользователь говорит «убери на столе», планировщик сначала
распознаёт объекты и соответствующие им целевые области, определяет
порядок обработки и явно задаёт для каждого шага начальные условия,
условия завершения и ограничения действий. Например, порядок уборки
может выглядеть так:

\begin{Shaded}
\begin{Highlighting}[]
\NormalTok{Обработать красную чашку → Убрать жёлтую бумагу → Проверить стол}
\end{Highlighting}
\end{Shaded}

«Обработать красную чашку» необходимо дополнительно разделить на два
действия и одну проверку:

\begin{Shaded}
\begin{Highlighting}[]
\NormalTok{pick(red\_cup) → place(red\_cup, tray) → verify\_state()}
\end{Highlighting}
\end{Shaded}

Рассмотрим задачу «взять чашку и поместить её в поднос». Завершение
программы захвата означает лишь окончание выполнения этой программы.
Прежде чем поднимать манипулятор, необходимо по обратной связи от
захвата или новому изображению определить, действительно ли чашка
удерживается. Если захват оказался пустым, следует повторно определить
положение чашки; если кто-то убрал поднос, успешный захват остаётся
действительным, а перепланировать нужно лишь действие размещения. Если
каждый навык соответствует проверяемому изменению состояния, имеет
ограничения по действиям и длительности, а затем возвращает состояние
выполнения и новое наблюдение, планировщик может восстановиться именно с
того шага, на котором произошла ошибка.

\begin{quote}
\textbf{Эксперимент 6-13 ★★: автономная уборка стола XLeRobot с помощью
Gemini Robotics-ER 1.5}

Используйте оборудование, расположение предметов на столе, задачу и
критерии приёмки из эксперимента 6-11, передав наблюдение и планирование
модели воплощённого мышления. Ориентируясь на способ оркестрации
RoboCrew, спроектируйте для этой задачи пять инструментов:
\texttt{observe\_scene}, \texttt{pick}, \texttt{place},
\texttt{verify\_state} и \texttt{stop}.\footnote{Google DeepMind,
  \emph{Gemini Robotics-ER 1.5}. XLeRobot, управление с помощью
  LLM-агента. Пример из исходного проекта XLeRobot демонстрирует способ
  оркестрации модели и вызовов инструментов; в этом разделе сохраняется
  тот же принцип оркестрации, но инструменты действий ограничены
  откалиброванными примитивами для работы на столе: захвата, размещения,
  проверки и остановки.}

Сначала модель считывает наблюдение стола и определяет порядок обработки
объектов, а затем вызывает откалиброванные навыки захвата и размещения.
После завершения каждого навыка система получает новое наблюдение и
проверяет постусловия; после сбоя она выбирает повторную попытку или
остановку в соответствии с текущим состоянием. Система выполнения
немедленно прекращает работу при получении сигнала аварийной остановки,
обнаружении выхода за границы или невозможности подтвердить критически
важное состояние; модель также может самостоятельно завершить задачу с
помощью \texttt{stop}.

В конце по новому наблюдению среды подтвердите, что чашка находится
внутри подноса, макулатура --- внутри мусорного контейнера, а
манипулятор вернулся в безопасную позу, и проверьте записи о
столкновениях и выходе за границы. Оператор отвечает за запуск,
аварийную остановку и наблюдение на месте, а любое вмешательство в ходе
выполнения записывается отдельно.
\end{quote}

Эксперимент 6-13 показывает, как подключить действия настоящего робота и
проверить результат уборки, однако в сопутствующих записях пока нет
испытаний этой задачи на реальном оборудовании. Поэтому сначала сравним
в симуляционной среде три подхода: разомкнутое выполнение по плану,
проверку результата после каждого шага и предварительное прогнозирование
последствий действия. Цель задачи и пять типов инструментов остаются
неизменными, а сбои можно многократно воспроизводить, чтобы наблюдать,
как три подхода справляются с одной и той же ошибкой.

\hypertarget{vla-ux443ux43fux440ux430ux432ux43bux435ux43dux438ux435}{%
\subsection{VLA-управление}\label{vla-ux443ux43fux440ux430ux432ux43bux435ux43dux438ux435}}

VLA --- сокращение от Vision-Language-Action, то есть модель «зрение ---
язык --- действие». Она генерирует следующее действие на основе текущего
изображения, инструкции задачи и доступного состояния робота:

\begin{Shaded}
\begin{Highlighting}[]
\NormalTok{Текущее наблюдение + инструкция навыка → действие или последовательность действий}
\end{Highlighting}
\end{Shaded}

В примере с уборкой стола высокоуровневый планировщик передаёт
\texttt{pick(red\_cup)}, а стратегия навыка по изображению определяет, с
какой стороны приблизиться к чашке, когда закрыть захват и по какой
траектории поднять манипулятор. После выполнения этого фрагмента
движения система получает новое наблюдение стола, подтверждает, что
чашка захвачена, и только затем переходит к
\texttt{place(red\_cup,\ tray)}. Интерфейс инструмента задаёт желаемое
изменение состояния, а VLA учится реализовывать его посредством
непрерывных действий.

Действия можно представлять разными способами. RT-2 и исходная OpenVLA
квантуют числовые значения каждого измерения действий в дискретные
токены и выводят их с помощью механизма генерации языковой модели; π₀
использует сопоставление потоков для генерации непрерывных блоков
действий. Первый подход позволяет повторно использовать
авторегрессионную архитектуру генерации, а второй непосредственно
моделирует распределение фрагмента непрерывных действий.\footnote{Moo
  Jin Kim et al., \emph{OpenVLA: An Open-Source Vision-Language-Action
  Model}, arXiv:2406.09246, 2024. См. также описание проекта RT-2 и
  публикацию Physical Intelligence \emph{Our First Generalist Policy}.}

Инференс модели и управление исполнительным механизмом обычно
выполняются с разной частотой обновления. Разбиение действий на блоки
(action chunking) позволяет модели за один раз генерировать короткий
фрагмент будущих действий, который управляющий поток выполняет с более
высокой частотой. Если дополнительно наложить инференс и выполнение друг
на друга, система сможет готовить следующий блок во время выполнения
текущего, сокращая ожидание. При стыковке блоков действий также
необходимо учитывать непрерывность траектории и актуальность наблюдений:
если чашка сдвинется, пока манипулятор тянется к ней, оставшиеся
действия могут по-прежнему быть направлены к старому положению.
Сокращение фактически выполняемого за раз фрагмента, скользящее
обновление по новым наблюдениям или прерывание действия при обнаружении
аномалии повышают скорость реакции, но также увеличивают затраты на
инференс и диспетчеризацию.

\hypertarget{ux434ux430ux43dux43dux44bux435-ux438-ux443ux441ux43bux43eux432ux438ux44f-ux440ux430ux437ux432ux451ux440ux442ux44bux432ux430ux43dux438ux44f-vla}{%
\subsection{Данные и условия развёртывания
VLA}\label{ux434ux430ux43dux43dux44bux435-ux438-ux443ux441ux43bux43eux432ux438ux44f-ux440ux430ux437ux432ux451ux440ux442ux44bux432ux430ux43dux438ux44f-vla}}

Чтобы стабильно преобразовывать команду «захватить чашку» в действия,
необходимо одновременно учитывать несколько типов различий:

\begin{itemize}
\tightlist
\item
  \textbf{Охват обучающих данных}: доступных демонстраций действий
  роботов намного меньше, чем текстов и изображений из интернета. Знания
  о том, что представляет собой «чашка», недостаточно: для определения
  усилия необходим практический опыт захвата. Демонстрации должны
  охватывать разные материалы, формы, положения и условия трения чашек.
  Одинаковое направление захвата может привести к разным результатам для
  керамической чашки и мягкого бумажного стаканчика.
\item
  \textbf{Цель обучения действий и их последствий}: клонирование
  поведения в основном учит тому, какое действие демонстратор предпримет
  в текущих условиях. Для сравнения последствий нескольких действий
  можно дополнительно использовать прогнозирование состояния или оценку
  результатов действий.
\item
  \textbf{Конструкция робота и интерфейс}: степени свободы, система
  координат, устройство захвата и задержки исполнительных механизмов
  определяют, как действие реализуется на оборудовании. При переносе
  стратегии необходимо согласовать представление действий, параметры
  калибровки и цикл управления.
\item
  \textbf{Актуальность наблюдений}: во время выполнения блока действий
  объект может переместиться, оказаться перекрыт или опрокинуться.
  Система должна определить, когда получать новое наблюдение и когда
  обновлять по нему оставшиеся действия.
\end{itemize}

Эти условия связывают визуальное распознавание, управление навыками и
изменения среды. Жидкость в чашке влияет на скорость перемещения, кабель
питания у края стола изменяет допустимый путь, а контакт бумаги со
столом влияет на способ захвата. Если планировщик способен заранее
сравнить эти последствия, он получает больше оснований для выбора
навыка.

\hypertarget{ux43cux438ux440ux43eux432ux430ux44f-ux43cux43eux434ux435ux43bux44c}{%
\subsection{Мировая
модель}\label{ux43cux438ux440ux43eux432ux430ux44f-ux43cux43eux434ux435ux43bux44c}}

В планировании робота мировую модель с условием на действие можно
понимать как «предиктор результата действия»: получив текущее состояние
и действие-кандидат, она оценивает, что произойдёт дальше.

\begin{Shaded}
\begin{Highlighting}[]
\NormalTok{Текущее состояние + действие{-}кандидат}
\NormalTok{    → прогноз следующего состояния или будущего фрагмента}
\NormalTok{    → сравнение результатов кандидатов}
\NormalTok{    → выбор действия, перепланирование или остановка}
\end{Highlighting}
\end{Shaded}

Этот процесс включает три задачи: представить текущее состояние,
спрогнозировать вызванное действием изменение состояния и предоставить
планировщику или контроллеру возможность выбрать действие на основе
прогноза. Состояние может быть представлено структурированными
переменными, такими как положение объекта и признак его удержания
захватом, либо внутренним представлением, извлечённым из изображения.
V-JEPA 2 сначала обучается по видео представлению и прогнозированию
состояния, а дополнительно обученная V-JEPA 2-AC вводит условие на
действие для планирования робота. Связывание действий с будущими
наблюдениями также относится к подходам мировых моделей действий
(World-Action Model).\footnote{Meta AI, \emph{Представляем мировую
  модель V-JEPA 2 и новые бенчмарки для рассуждений о физическом мире},
  11 июня 2025 года. Номер технического отчёта: arXiv:2506.09985.}

Мировая модель может помогать роботу на трёх этапах: до начала действий
отдельно прогнозировать результаты захвата, толкания или ожидания; во
время выполнения сопоставлять прогноз с фактическим наблюдением и при
расхождении корректировать действие или перепланировать; во время
обучения изучать изменения состояния по видео, симуляционным данным и
траекториям неудач, сокращая количество проб и ошибок на реальном
оборудовании.

Предположим, жёлтая макулатура частично закрыта красной чашкой.
Планировщик может сравнить варианты «сначала захватить бумагу», «сначала
убрать чашку» и «изменить направление захвата». Автор рекомендует прежде
всего использовать мировую модель для прогнозирования непосредственно
связанных с задачей и проверяемых изменений состояния. Если она способна
определить, станет ли бумагу проще захватить и легко ли опрокинуть
чашку, этого достаточно для упорядочивания действий-кандидатов.
Представление состояния и цель обучения модели следует проектировать
вокруг этих суждений. Если модель генерирует будущие изображения, автор
рекомендует сначала проверить, точно ли они отражают последствия
действий, и лишь затем добавлять визуальные подробности; ценность
реалистичности изображения в конечном счёте должна выражаться в
улучшении планирования и управления. После выполнения система проверяет
фактический результат по камере и состоянию робота, а затем решает,
продолжить работу, повторить попытку или сменить подход.

Чем длиннее горизонт прогнозирования, тем легче накапливаются ошибки. На
практике можно сначала прогнозировать краткосрочные последствия,
выполнять небольшой фрагмент, затем перепланировать по новому наблюдению
и учитывать неопределённость прогноза при выборе. Генеративные мировые
модели также могут предоставлять интерактивную симуляцию и визуализацию
сцен; при их использовании для управления необходимо дополнительно
проверить обусловленность действиями, закономерности контакта и ошибки
прогнозирования. Аппаратные ограничители, аварийная остановка и
управление в реальном времени по-прежнему остаются задачами системы
выполнения.\footnote{Google DeepMind, \emph{Genie 3: A new frontier for
  world models}, 2025-08-05; Zachary Lin и др., \emph{Cosmos World
  Foundation Model Platform for Physical AI}, 2025.}

\begin{quote}
\textbf{Эксперимент 6-14 ★★: сравнение трёх замкнутых контуров
автономной уборки стола в симуляторе}

Представьте красную чашку и жёлтую макулатуру с помощью
структурированного состояния: каждый объект находится на одном из трёх
этапов --- «на столе», «взят» или «размещён», --- а для целевой области
записывается её доступность. Симулятор предоставляет пять типов
инструментов: наблюдение, захват, размещение, проверку и остановку, а
при захвате макулатуры внедряет не более одного временного сбоя.

Сравните три способа управления. \textbf{Разомкнутое выполнение}
выполняет четыре действия в фиксированном порядке. \textbf{Пошаговая
проверка} считывает состояние после каждого действия, после сбоя снова
получает наблюдение и выбирает доступный для выполнения шаг.
\textbf{Выполнение с прогнозированием} использует небольшую нейронную
сеть для изучения отображения «состояние и действие --- следующее
состояние» и упорядочивает доступные действия по прогнозируемому
прогрессу размещения объектов.

Сравните завершение задач, количество вызовов инструментов и процесс
восстановления после сбоев. Например, после неудачного захвата
макулатура остаётся на столе, поэтому непосредственная попытка
размещения завершится ошибкой из-за отсутствия захваченного объекта;
продолжить размещение можно только после считывания состояния и
повторного захвата. Обе стратегии с обратной связью в конце вызывают
\texttt{verify\_state}, чтобы проверить размещение объектов.

Сначала рассмотрим эпизоды без внедрённых сбоев: все три стратегии
успешно завершают задачу. Затем установим вероятность внедрения сбоя в
50\% и запустим каждую стратегию всего в 384 эпизодах на трёх начальных
значениях --- по 128 эпизодов на каждое начальное значение. Разомкнутая
стратегия по-прежнему вызывает инструменты 4 раза за эпизод, но успешно
завершает лишь 205 эпизодов. Пошаговая проверка и выполнение с
прогнозированием завершают все 384 эпизода, в среднем выполняя примерно
по 6,93 вызова. Дополнительные вызовы позволяют контроллеру обнаружить
пустой захват и восстановиться. В этой задаче, где внедряется только
один восстанавливаемый сбой, прогнозирование следующего состояния не
позволило спасти ни одного дополнительного эпизода по сравнению с
непосредственной проверкой.
\end{quote}

\hypertarget{ux43eux442-ux441ux438ux43cux443ux43bux44fux446ux438ux43eux43dux43dux43eux439-ux441ux440ux435ux434ux44b-ux43a-ux43dux430ux441ux442ux43eux44fux449ux435ux43cux443-ux440ux43eux431ux43eux442ux443}{%
\subsection{От симуляционной среды к настоящему
роботу}\label{ux43eux442-ux441ux438ux43cux443ux43bux44fux446ux438ux43eux43dux43dux43eux439-ux441ux440ux435ux434ux44b-ux43a-ux43dux430ux441ux442ux43eux44fux449ux435ux43cux443-ux440ux43eux431ux43eux442ux443}}

При переходе от симулированного стола к настоящему роботу необходимо
последовательно учитывать различия между средами. С точки зрения
изображения одна и та же чашка и поднос могут оказаться на другом фоне,
при ином освещении, положении камеры и взаимном перекрытии; с точки
зрения движения на траекторию и захват влияют трение, шум датчиков и
задержки исполнительных механизмов.

Доменная рандомизация изменяет во время обучения фон, внешний вид
объектов, освещение или шум, чтобы модель выполняла одинаковые решения в
разных условиях. Эксперимент 6-15 исследует только один элемент этого
процесса: в какую сторону следует двигаться после обнаружения объекта и
целевой области?

\begin{quote}
\textbf{Эксперимент 6-15 ★★★: межсредовое RGB-тестирование одной и той
же задачи на столе}

Представьте объект кругом, а целевую область --- прямоугольником и
сгенерируйте RGB-изображения стола. По взаимному положению объекта и
цели модель выводит одно из четырёх направлений движения: влево, вправо,
вверх или вниз; когда две позиции оказываются достаточно близко, она
выводит определённую в эксперименте метку \texttt{grasp}. Эта метка
означает, что 2D-позиции удовлетворяют условию близости.

Обучите четыре визуальные стратегии с одинаковой архитектурой: с
фиксированными визуальными условиями, преимущественным изменением фона,
преимущественным изменением внешнего вида объекта и сочетанием
нескольких типов изменений. Примените стратегии каждой группы к исходным
условиям и тестовым изображениям с новыми цветами, текстурами, яркостью
и шумом, сравните точность определения направления и изучите
изображения, на которых допущены ошибки.

Каждый тестовый домен для каждого начального значения содержит 1024
синтетических изображения. Сначала сравним среднюю точность по трём
начальным значениям: стратегия, обученная в фиксированных визуальных
условиях, достигает 93,03\% в исходном домене, но при переходе в целевой
домен с новыми текстурами и шумом (\texttt{target\_realistic}) её
точность падает до 54,82\%; стратегия, обученная с сочетанием нескольких
типов изменений, достигает 87,50\% в исходном домене и 76,11\% в том же
целевом домене. Для сравнения, группа с рандомизацией только фона
достигает в этом целевом домене лишь 36,88\%.\footnote{LeRobot,
  руководство по Sim2Real.}
\end{quote}

\hypertarget{ux438ux442ux43eux433ux438-ux433ux43bux430ux432ux44b-5}{%
\section{Итоги
главы}\label{ux438ux442ux43eux433ux438-ux433ux43bux430ux432ux44b-5}}

Вернёмся к звонку об изменении бронирования из начала главы. Пока в
фоновом режиме выполняется поиск рейсов, пользователь уточняет
требования к багажу. Даже если результаты поиска поступят позже
сообщения пользователя, они могут отвечать на уже неактуальный вопрос.
Очередь событий должна сохранять источник сообщения, идентификатор
задачи и порядок событий, чтобы передний план мог приоритетно принять
новые требования и определить, можно ли продолжить исходный поиск,
следует ли его отменить или нужно повторить поиск с другими условиями.
Это проблема \textbf{момента выполнения}, и простой переход от
текстового запроса к голосовому её не решает.

Голосовое взаимодействие добавляет ещё один уровень требований. Система
должна слышать, когда пользователь перебивает её во время озвучивания
ценового предложения; если дата изменилась, неактуальное озвучивание
следует прекратить. Каскадная и сквозная архитектуры описывают, как звук
проходит через распознавание, инференс и генерацию, а полнодуплексная
--- могут ли пользователь и ассистент одновременно слушать и говорить.

При переходе к экрану или роботизированной руке проблема остаётся той же
--- меняется лишь \textbf{форма наблюдений и действий}. Координаты
кнопки берутся с предыдущего скриншота, поэтому после обновления
страницы её нужно найти заново. Роботизированная рука может за один раз
предсказывать небольшой непрерывный фрагмент действий, чтобы движения
суставов были более плавными, но чем длиннее такой фрагмент, тем сложнее
реагировать на изменения во время его выполнения. Высокоуровневое
планирование, навыки и низкоуровневое управление работают в разных
временных масштабах, а чтобы убедиться, что чашка действительно
поставлена на поднос, необходимо ещё раз осмотреть стол. Если объединить
эти взаимодействия, становится видна общая для них последовательность:

\begin{Shaded}
\begin{Highlighting}[]
\NormalTok{Получить наблюдение}
\NormalTok{  → Определить текущее состояние и момент выполнения}
\NormalTok{  → Выбрать ответ или действие}
\NormalTok{  → Выполнить его и получить обратную связь от среды}
\NormalTok{  → Продолжить, скорректировать, повторить, остановиться или перепланировать}
\end{Highlighting}
\end{Shaded}

Этот процесс помогает локализовать сбои. Если пользователь перебивает
систему, а прежнее ценовое предложение продолжает озвучиваться, следует
проверить, кто отвечает за пробуждение переднего плана, можно ли
вытеснить воспроизведение и до какой безопасной точки дошёл сигнал
отмены. Если озвучивание остановилось вовремя, но система всё равно
называет прежнюю дату, нужно проследить передачу результатов поиска
между фоновым и передним планами. Как после исправления доказать, что
система будет правильно работать и в большем числе звонков? С этого
момента вся книга переходит к улучшению возможностей. В
\crossreflink{chap:7}{главе 7} сначала рассматривается, как
проектировать задачи, регистрировать сбои и сравнивать изменения; эта
система оценки затем используется для постобучения, непрерывной эволюции
и совместной работы нескольких агентов.

\begin{questionbox}

\hypertarget{ux432ux43eux43fux440ux43eux441ux44b-ux434ux43bux44f-ux440ux430ux437ux43cux44bux448ux43bux435ux43dux438ux44f-5}{%
\section*{Вопросы для
размышления}\label{ux432ux43eux43fux440ux43eux441ux44b-ux434ux43bux44f-ux440ux430ux437ux43cux44bux448ux43bux435ux43dux438ux44f-5}}
\addcontentsline{toc}{section}{Вопросы для размышления}

\begin{enumerate}
\def\labelenumi{\arabic{enumi}.}
\tightlist
\item
  ★★ Когда асинхронный AI-агент получает новое событие, ему необходимо
  определить его приоритет относительно текущей задачи. Если такое
  решение требует понимания семантики --- например, распознать,
  дополняет ли пользователь условия или требует немедленно остановиться,
  --- как следует распределить обязанности между движком правил и LLM?
  Сравните время отклика, стоимость и последствия ошибочной
  классификации.
\item
  ★★ В очереди накопилось 20 событий: 10 результатов работы
  инструментов, 5 сообщений пользователей и 5 системных напоминаний. Как
  бы вы организовали их порядок, формат и состояния обработки, чтобы
  модель могла отслеживать зависимости между задачами и реагировать на
  важные изменения?
\item
  ★★★ Агент может использовать выделенный адрес электронной почты и
  выделенный номер телефона либо после авторизации работать с личными
  учётными записями пользователя. Обсудите, для каких сценариев подходит
  каждый из этих двух способов идентификации, учитывая фоновую работу,
  требования другой стороны к проверке личности, объём полномочий и
  аудит.
\item
  ★★ Сквозная речевая модель обрабатывает аудио напрямую, поэтому ошибки
  на промежуточных этапах трудно локализовать извне. Как организовать
  наблюдение и контрольные эксперименты, чтобы различать ошибки
  восприятия, рассуждения и выражения?
\item
  ★ Механизм MPS в Step-Audio R1.1 позволяет одновременно развивать
  замысел и его выражение. Как системе следует использовать паузы,
  сообщения о ходе работы и поэтапное изложение? Что делать, если
  последующие рассуждения изменили уже озвученный вывод?
\item
  ★★ И разметка визуальных областей в SoM, и индексация элементов DOM
  позволяют модели одним номером указать место выполнения операции. Как
  при нестандартных элементах управления, перекрытии или динамических
  изменениях сочетать повторную разметку, прогнозирование координат и
  проверку после операции?
\item
  ★★ Качество данных телеуправления зависит от навыков оператора. Как
  дрожание, движение обходным путём, паузы и неудачные попытки захвата
  влияют на обучение VLA? Как отбирать данные, сохраняя при этом ценные
  процессы восстановления после ошибок?
\item
  ★★★ Для речи, Computer Use и робототехники можно использовать единую
  модель, модульный каскад либо разделение на взаимодействие на переднем
  плане и фоновое рассуждение. Какому сочетанию этих подходов вы бы
  отдали предпочтение на ближайшие пять лет? Обоснуйте решение с учётом
  задержки, наблюдаемости, стоимости обновлений и рисков задач.
\item
  ★★ Стандартные веб-страницы могут предоставлять DOM или информацию о
  доступности, тогда как Canvas, WebGL и самостоятельно отрисованные
  элементы управления часто требуют визуальной локализации. Какие
  дополнительные затраты возникают при одновременной поддержке
  структурированного и визуального путей и какие преимущества это даёт?
\item
  ★★ Разбиение действий на блоки позволяет VLA за один раз сформировать
  последовательность будущих действий. Если во время её выполнения
  объект переместили, как согласовать длину блока, повторное наблюдение,
  фоновое рассуждение и прерывание действий, чтобы обеспечить как
  эффективность выполнения, так и скорость реакции?
\item
  ★★★ Для быстрого взаимодействия необходимо своевременно действовать.
  На примерах голосового ответа, отправки веб-формы и захвата объекта
  роботом разработайте градацию действий, подтверждение состояния и
  механизм остановки так, чтобы скорость реакции соответствовала
  последствиям задачи.
\item
  ★★★ Выберите один из примитивов: пробуждение, безопасная точка,
  отмена, вытеснение или разделение на быстрый и медленный контуры.
  Сравните его реализацию в событийно-ориентированной задаче и при
  разбиении действий робота на блоки. Как скорость изменений среды,
  обратимость действий и стоимость наблюдения влияют на проектирование?
\item
  ★★ Следует ли объединять чат с клиентом, обсуждение согласования и
  callback об оплате по одному заказу в один поток? Сравните объединение
  в один поток и подписку по связанным объектам с точки зрения порядка
  обработки, степени параллелизма, миграции ключа секционирования и
  изоляции прав доступа.
\item
  ★★ Голосовой помощник для путешествий умеет узнавать погоду, читать
  данные из памяти пользователя, искать в интернете, находить рейсы и
  бронировать билеты. Как распределить работу между передней моделью,
  фоновой моделью и средой выполнения? Пользователь диктует имя
  попутчика, а пока идёт поиск рейса, спрашивает о погоде. Как не
  ошибиться в имени при бронировании и не забыть сообщить о рейсе?
\end{enumerate}

\textbf{Практикум к \crossreflink{chap:6}{главе 6}.}
\href{https://github.com/bojieli/ai-agent-book/tree/40f4ce5778e6cc93e3606c3079bf8cf13a0c4d04/chapter6}{Каталог
экспериментов} ·
\href{https://github.com/bojieli/ai-agent-book/blob/main/chapter6/README.ru.md}{Путеводитель
на русском} ·
\href{https://github.com/bojieli/ai-agent-book/tree/40f4ce5778e6cc93e3606c3079bf8cf13a0c4d04/chapter6/voice-tool-routing}{Исходники
эксперимента 6-7: вызов инструментов голосового агента} ·
\href{https://github.com/bojieli/ai-agent-book/tree/40f4ce5778e6cc93e3606c3079bf8cf13a0c4d04/chapter6/agent-with-event-trigger}{Исходники
эксперимента 6-1: приём событий и проверка отказоустойчивости}.

\end{questionbox}

\hypertarget{ux43eux446ux435ux43dux43aux430-ai-ux430ux433ux435ux43dux442ux430}{%
\chapter{\texorpdfstring{Оценка
AI-агента\label{chap:7}}{Оценка AI-агента}}\label{ux43eux446ux435ux43dux43aux430-ai-ux430ux433ux435ux43dux442ux430}}

В первых шести главах было рассмотрено создание отдельного агента:
контекст, память, база знаний, инструменты, код, а также интерфейсы
наблюдений и действий. После создания системы необходимо с помощью
набора воспроизводимых задач определить, что она умеет, при каких
условиях чаще всего терпит неудачу и что следует изменить на следующем
этапе. Результаты оценки также обеспечивают обратную связь для
постобучения и непрерывного развития.

В процессе разработки многие решения необходимо сравнивать на реальных
задачах: какую модель использовать, какие инструменты сделать
доступными, какое содержимое и способ индексации выбрать для базы
знаний, как организовать память пользователя, промпты и Skills, а также
какие ограничения добавить в Harness. Эти решения могут влиять друг на
друга. Например, более подробное описание инструмента способно сократить
число ошибок в параметрах, но увеличить расход контекста; добавление
этапа проверки может повысить долю успешно выполненных задач,
одновременно увеличив время выполнения.

\textbf{Сравнительные эксперименты} сопоставляют разные варианты на
одинаковых задачах и в одинаковых условиях, фиксируя изменения доли
успешных выполнений, затраченного времени и стоимости. В ходе
\textbf{абляционных экспериментов} удаляются или отключаются такие
компоненты, как переранжирование результатов поиска, сводка состояния
или проверка результата, после чего оценивается изменение работы
системы. Если в каждом цикле явно фиксировать внесённое изменение,
сохранять остальные условия выполнения неизменными и повторять
измерения, можно постепенно установить связь между инженерными решениями
и результатами выполнения задач. Автор рекомендует включить
сравнительные и абляционные эксперименты в повседневную разработку. По
нескольким успешно завершённым траекториям легко принять случайный успех
за улучшение или заметить выигрыш только в одном типе задач, упустив
ухудшение в других.

\textbf{Объектом оценки должен быть полноценный агент, состоящий из
модели и Harness.} Модель определяет, как понимать задачу и выбирать
действия, а Harness организует контекст, управляет вызовами
инструментов, поддерживает состояние и выполняет проверки. Одна и та же
модель при подключении разных инструментов или процессов обратной связи
может давать разные результаты выполнения задач. Coding Harness,
представленный в пятой главе, как раз предоставляет условия для подобных
сравнений.

Взаимодействие модели и Harness можно исследовать с помощью
\textbf{эксперимента с заменой модели (model swap)}: сначала фиксируются
среда выполнения и задачи, затем заменяется модель и сравниваются
итоговые результаты и распределение ошибок. Если более сильная модель
по-прежнему терпит неудачу на тех же этапах, можно проверить, не
отсутствуют ли необходимые инструменты, своевременно ли поступает
обратная связь и охватывает ли оценивание поставленные цели. Если
результаты выполнения задач существенно меняются вместе с моделью,
следует проанализировать, какие способности к рассуждению, планированию
или следованию инструкциям повлияли на результат.

Замену модели можно сочетать с абляцией компонентов. Например, две
модели можно запустить в конфигурациях «со сводкой состояния» и «без
сводки состояния», чтобы проверить, помогает ли сводка обеим моделям
возобновлять длительные задачи и зависит ли получаемый выигрыш от
модели. Это позволяет не только анализировать различия между моделями,
но и определять, при каких условиях эффективен конкретный компонент
Harness.

Оценка также служит регрессионной проверкой при обновлении модели. Новая
модель может улучшить результаты на публичных бенчмарках, но ухудшить их
на отдельных бизнес-задачах; собственный набор задач позволяет выявлять
такие изменения. Автор рекомендует «разрабатывать продукт для будущих
моделей»: заранее подготовить инструменты, бизнес-процессы и набор для
оценки, а затем непрерывно отслеживать результаты новых моделей на одних
и тех же задачах. Когда доля успешных выполнений, время отклика и
стоимость будут соответствовать требованиям продукта, можно переходить к
запуску.

Как преобразовать требования к запуску в систему оценки? При изучении
\crossreflink{fig:7-1}{рисунка 7-1} можно проследить весь путь одного
теста: что требуется измерить, откуда взять задачи, по каким
свидетельствам определять успех и, наконец, как записывать и сравнивать
результаты?

\begin{figure}
\hypertarget{fig:7-1}{%
\centering
\includegraphics{images/revised/fig7-1.pdf}
\caption{Рисунок 7-1. Четыре этапа системы оценки агента}\label{fig:7-1}
}
\end{figure}

\hypertarget{ux43eux446ux435ux43dux43aux430-ux43dux430-ux43fux440ux438ux43cux435ux440ux435-ux43eux434ux43dux43eux439-ux437ux430ux434ux430ux447ux438-ux43eux431ux43bux430ux441ux442ux44c-telecom-ux432-ux3c4uxb2-bench}{%
\section{Оценка на примере одной задачи: область telecom в
τ²-bench}\label{ux43eux446ux435ux43dux43aux430-ux43dux430-ux43fux440ux438ux43cux435ux440ux435-ux43eux434ux43dux43eux439-ux437ux430ux434ux430ux447ux438-ux43eux431ux43bux430ux441ux442ux44c-telecom-ux432-ux3c4uxb2-bench}}

Сначала рассмотрим сбой мобильного интернета. John Smith пользуется
номером 555-123-2002 и находится во Франции, а в строке состояния его
телефона отображается «No Service». Он хочет восстановить подключение к
сети и считает проблему решённой, только если результат проверки
скорости достигает уровня excellent. При необходимости он готов докупить
2.0 GB трафика, но хочет сохранить текущий тариф.

Этот пример взят из области telecom бенчмарка τ²-bench от Sierra. Среда
одновременно моделирует системы оператора и пользовательское устройство,
а AI-агент службы поддержки и симулятор пользователя с помощью
инструментов читают и изменяют разные части состояния. При оценке
необходимо проследить взаимодействие обеих сторон: как обнаруживается
проблема, как выполняются действия и восстанавливается ли в итоге
подключение к сети.\footnote{В τ²-bench от Sierra Research
  пользовательские инструменты для области telecom определены в
  \texttt{user\_tools.py}, а точка входа для оценки находится в
  \texttt{evaluator.py}. Содержание примера и результаты подготовлены на
  основе сохранённых сопутствующих задач и траекторий запусков для пяти
  примеров.}

\hypertarget{ux447ux435ux442ux44bux440ux435-ux441ux43eux441ux442ux430ux432ux43bux44fux44eux449ux438ux435-ux43eux43fux440ux435ux434ux435ux43bux435ux43dux438ux44f-ux437ux430ux434ux430ux447ux438}{%
\subsection{Четыре составляющие определения
задачи}\label{ux447ux435ux442ux44bux440ux435-ux441ux43eux441ux442ux430ux432ux43bux44fux44eux449ux438ux435-ux43eux43fux440ux435ux434ux435ux43bux435ux43dux438ux44f-ux437ux430ux434ux430ux447ux438}}

Чтобы другой агент службы поддержки мог решить ту же задачу, необходимо
сначала воспроизвести состояние телефона и учётной записи, заставить
симулятор пользователя выполнять диагностику в соответствии с теми же
требованиями, а затем применить те же критерии оценки. В таблице ниже
перечислены четыре составляющие определения задачи. При чтении обратите
внимание на различие: автор задачи знает исходные настройки телефона, но
сколько информации в начале диалога доступно агенту и пользователю?

\begin{longtable}[]{@{}
  >{\raggedright\arraybackslash}p{(\columnwidth - 4\tabcolsep) * \real{0.2370}}
  >{\raggedright\arraybackslash}p{(\columnwidth - 4\tabcolsep) * \real{0.4577}}
  >{\raggedright\arraybackslash}p{(\columnwidth - 4\tabcolsep) * \real{0.3053}}@{}}
\toprule\noalign{}
\begin{minipage}[b]{\linewidth}\raggedright
\hspace{0pt}Составляющая
\end{minipage} & \begin{minipage}[b]{\linewidth}\raggedright
\hspace{0pt}Поле
\end{minipage} & \begin{minipage}[b]{\linewidth}\raggedright
\hspace{0pt}Содержимое в этом примере
\end{minipage} \\
\midrule\noalign{}
\endhead
\bottomrule\noalign{}
\endlastfoot
\hspace{0pt}Обращение & \hspace{0pt}\texttt{ticket} &
\hspace{0pt}Пользователь, номер, страна пребывания, невозможность
подключиться к интернету, требования к проверке скорости и ограничения
тарифа \\
\hspace{0pt}Пользовательский сценарий &
\hspace{0pt}\texttt{user\_}\allowbreak{}\texttt{scenario.instructions} &
\hspace{0pt}Что известно пользователю, как получать сведения в ответ на
вопросы и когда подтверждать решение проблемы \\
\hspace{0pt}Исходное состояние &
\hspace{0pt}\texttt{initial\_}\allowbreak{}\texttt{state.initialization\_}\allowbreak{}\texttt{actions}
& \hspace{0pt}Задать личность и местоположение пользователя, включить
авиарежим, отключить роуминг на устройстве и подключить роуминг на
стороне оператора \\
\hspace{0pt}Критерии оценки &
\hspace{0pt}\texttt{evaluation\_}\allowbreak{}\texttt{criteria} &
\hspace{0pt}Проверить действия и состояние среды, а в
\texttt{reward\_}\allowbreak{}\texttt{basis} указать оцениваемые
элементы \\
\end{longtable}

\textbf{Что известно пользователю.} Поле \texttt{known\_info} содержит
имя, номер и страну пребывания. Включённый авиарежим и отключённый
роуминг на устройстве хранятся в состоянии среды, поэтому симулятор
пользователя должен запросить их с помощью инструментов и лишь затем
сообщить агенту. Отдельное поле с известными пользователю сведениями
позволяет симулятору различать информацию, которую можно сообщить сразу,
и информацию, которую необходимо фактически проверить. Одного требования
в промпте «не рассказывать всё сразу» недостаточно, чтобы стабильно
контролировать, когда именно раскрывается каждый факт. Такая организация
реализует постепенное раскрытие информации (Progressive Information
Disclosure): диалог начинается со слов «У меня не работает интернет», а
агент с помощью вопросов и действий постепенно получает сведения,
необходимые для диагностики.

\textbf{Как отвечает пользователь.} Поле \texttt{task\_instructions}
предписывает проявить лёгкое недовольство после первой неудачной попытки
исправления, подтверждать решение проблемы только при результате
проверки скорости excellent и основывать ответы о состоянии устройства
на результатах инструментов. Уровни poor, fair и good не соответствуют
цели этой задачи. Привязка к фактам (Grounding) заставляет пользователя
отвечать в соответствии с фактическим состоянием моделируемой среды. Без
этого ограничения симулятор пользователя мог бы согласиться с
утверждением агента «Теперь всё должно работать»: диалог выглядел бы
успешным, хотя неисправность устройства осталась бы нерешённой.

\textbf{С какого состояния начинают обе стороны.} На стороне
пользователя \texttt{set\_user\_info} задаёт личность,
\texttt{set\_user\_location(abroad=true)} указывает, что он находится за
границей, после чего выполняются \texttt{turn\_airplane\_mode\_on} и
\texttt{turn\_roaming\_off}. На стороне оператора
\texttt{enable\_roaming} подключает роуминг для линии \texttt{L1002}
клиента \texttt{C1001}. Поэтому запрос к системам оператора показывает,
что роуминг подключён, но телефон всё равно не может выйти в сеть из-за
авиарежима и настройки роуминга на устройстве. Чтобы получить сведения,
которых нет в системах оператора, агент должен попросить пользователя
проверить устройство.

\textbf{Какие проверки определяют оценку.} В \texttt{actions}
перечислены выполняемые на стороне пользователя действия
\texttt{toggle\_airplane\_mode} и \texttt{toggle\_roaming}. В
\texttt{env\_assertions} входят проверка состояния мобильных данных и
проверка скорости с параметрами \texttt{expected\_speed=200} и
\texttt{expected\_desc="excellent"}. Формат задачи также предусматривает
\texttt{communicate\_info} и \texttt{nl\_assertions}, которые
соответственно описывают информацию, обязательную для сообщения
пользователю, и утверждения на естественном языке; в этой задаче оба
поля пусты.

В \texttt{reward\_basis} для этой задачи выбрано значение
\texttt{ENV\_ASSERTION}, поэтому итоговое вознаграждение зависит от
утверждений о состоянии среды. Проверки действий могут присутствовать в
результатах, но их участие в итоговой оценке определяется конфигурацией
подсчёта.

\hypertarget{ux43fux440ux43eux432ux435ux440ux43aux430-ux440ux435ux437ux443ux43bux44cux442ux430ux442ux43eux432-ux43fux43e-ux442ux440ux430ux435ux43aux442ux43eux440ux438ux438-ux432ux44bux43fux43eux43bux43dux435ux43dux438ux44f}{%
\subsection{Проверка результатов по траектории
выполнения}\label{ux43fux440ux43eux432ux435ux440ux43aux430-ux440ux435ux437ux443ux43bux44cux442ux430ux442ux43eux432-ux43fux43e-ux442ux440ux430ux435ux43aux442ux43eux440ux438ux438-ux432ux44bux43fux43eux43bux43dux435ux43dux438ux44f}}

В сопутствующем запуске были выбраны пять задач из области telecom,
каждая из которых выполнялась один раз; агент службы поддержки и
симулятор пользователя использовали GPT-4.1 mini. В итоге были пройдены
четыре задачи. Сначала разберём описанную выше задачу с роумингом во
Франции, а затем сопоставим её с неудачной.

\begin{quote}
\textbf{Эксперимент 7-1 ★: запуск τ²-bench и сравнение с развитием
τ-bench}

Сначала изучите пользовательский сценарий, действия инициализации и
условия оценки одной задачи, а затем выполните одно полное
взаимодействие. Отметьте на траектории запросы службы поддержки,
действия пользователя, изменения состояния и результат приёмки,
проверяя, на какой информации основывалось каждое решение.

При сравнении τ-bench и τ²-bench прежде всего выясните, может ли
пользователь изменять среду с помощью инструментов и какие требования к
координации со службой поддержки создают эти изменения.

При изучении \crossreflink{fig:7-2}{рисунка 7-2} сначала определите, кто
может отключить авиарежим на телефоне, а затем --- кто может изменить
конфигурацию роуминга у оператора. Возможности сторон различаются, что и
объясняет, почему после собственных действий агенту всё равно приходится
просить пользователя проверить скорость на телефоне.

\begin{figure}
\hypertarget{fig:7-2}{%
\centering
\includegraphics{images/revised/fig7-2.pdf}
\caption{Рисунок 7-2. Среда τ²-bench с двойным управлением и
многоуровневой проверкой}\label{fig:7-2}
}
\end{figure}
\end{quote}

При обработке задачи с роумингом во Франции агент сначала
идентифицировал учётную запись. По номеру он нашёл клиента
\texttt{C1001} и запросил сведения о трафике линий \texttt{L1001},
\texttt{L1002} и \texttt{L1003}, но на основании подробной информации о
\texttt{L1001} сделал ошибочный вывод: предоставленный пользователем
номер 555-123-2002 отсутствует среди активных линий, а наиболее близкий
номер --- 555-123-2001. После того как пользователь настоял, что номер
указан верно, агент продолжил проверку \texttt{L1002} и лишь тогда
подтвердил правильную линию.

Затем симулятор пользователя вызвал \texttt{check\_network\_status} и
\texttt{check\_status\_bar}. Результат содержал следующие состояния:
авиарежим включён, сотовая сеть недоступна, мобильные данные включены,
роуминг на устройстве отключён. На основании этого пользователь сообщил:
«На телефоне включён авиарежим, а роуминг данных отключён. Сначала нужно
отключить авиарежим?» Здесь запрос выполняется инструментом на стороне
пользователя, а его результат передаётся агенту через диалог.

Так работает механизм \textbf{двойного управления (Dual-Control)}: агент
запрашивает и изменяет данные в системах оператора, а симулятор
пользователя располагает отдельными инструментами устройства ---
например, может просматривать строку состояния, переключать авиарежим,
извлекать и повторно устанавливать SIM-карту и запускать проверку
скорости. Обе стороны координируют действия посредством общения, а среда
регистрирует вызванные каждой из них изменения состояния.

Затем агент попросил пользователя отключить авиарежим и включить роуминг
на устройстве, после чего в строке состояния появился сигнал 5G.
Пользователь запустил \texttt{run\_speed\_test}, получил результат 275
Mbps, Excellent и подтвердил решение проблемы. Оба утверждения о
состоянии среды были выполнены, поэтому вознаграждение за задачу
составило 1.0.

Эта успешная траектория также выявила проблему процесса. Политика
требует, чтобы служба поддержки вызывала только один инструмент за раз,
однако агент в одном сообщении инициировал вызовы
\texttt{get\_customer\_by\_phone} и \texttt{get\_customer\_by\_name}. В
итоговую оценку текущей задачи входят только утверждения о состоянии
среды, поэтому это нарушение не повлияло на вознаграждение.

В другой, неудачной траектории ошибочно выбранная линия так и не была
исправлена, из-за чего последующая диагностика также пошла по неверному
пути. Пользователь по-прежнему указывал номер 555-123-2002, но агент
выбрал \texttt{L1001} и продолжил диагностику на основании потребления
3.2/5 GB для этой линии. Инструмент уже сообщил, что \texttt{L1001}
соответствует номеру 555-123-2001, однако агент не исправил выбор и
после нескольких раундов нерелевантной диагностики перевёл обращение на
специалиста. На стороне пользователя удалось отключить режим экономии
трафика, но для правильной линии \texttt{L1002} не были приобретены
необходимые 2 GB трафика, поэтому ни одно из трёх утверждений о
состоянии среды не было выполнено.

Сравнение двух траекторий позволяет представить итоговую оценку в виде
конкретных направлений улучшения: проверять соответствие номера и линии,
корректировать выводы с учётом новых данных, различать состояние
устройства и систем оператора и подтверждать приобретение необходимого
пакета трафика. В последующем разделе об атрибуции сбоев мы продолжим
разбирать, как находить такие ключевые поворотные точки в длинных
траекториях.

\hypertarget{ux43cux435ux442ux440ux438ux43aux438-ux43eux446ux435ux43dux43aux438-ux43eux43fux440ux435ux434ux435ux43bux435ux43dux438ux435-ux443ux441ux43fux435ux445ux430}{%
\section{Метрики оценки: определение
успеха}\label{ux43cux435ux442ux440ux438ux43aux438-ux43eux446ux435ux43dux43aux438-ux43eux43fux440ux435ux434ux435ux43bux435ux43dux438ux435-ux443ux441ux43fux435ux445ux430}}

Четыре пройденных задания из пяти дают общий результат этого запуска.
Чтобы принимать на его основе продуктовые решения, необходимо также
учитывать тип задач и цену неудачи. Служба поддержки по возвратам должна
стабильно оформлять возвраты, соответствующие условиям; при написании
программы можно сгенерировать несколько реализаций и с помощью тестов
выбрать версию, которая работает и выдаёт правильный результат; при
проведении исследования также можно выдвинуть несколько гипотез, а затем
поручить исследователю проверить их. Метрики оценки должны
соответствовать такому способу использования.

\hypertarget{ux442ux435ux445ux43dux438ux447ux435ux441ux43aux43eux435-ux447ux443ux434ux43e-ux43eux446ux435ux43dux43aux430-ux432ux43eux437ux43cux43eux436ux43dux43eux441ux442ux435ux439-ux43fux440ux438-ux43dux435ux441ux43aux43eux43bux44cux43aux438ux445-ux43fux43eux43fux44bux442ux43aux430ux445-ux441-ux43fux43eux43cux43eux449ux44cux44e-passk}{%
\subsection{\texorpdfstring{Техническое чудо: оценка возможностей при
нескольких попытках с помощью
\href{mailto:Pass@k}{\nolinkurl{Pass@k}}}{Техническое чудо: оценка возможностей при нескольких попытках с помощью Pass@k}}\label{ux442ux435ux445ux43dux438ux447ux435ux441ux43aux43eux435-ux447ux443ux434ux43e-ux43eux446ux435ux43dux43aux430-ux432ux43eux437ux43cux43eux436ux43dux43eux441ux442ux435ux439-ux43fux440ux438-ux43dux435ux441ux43aux43eux43bux44cux43aux438ux445-ux43fux43eux43fux44bux442ux43aux430ux445-ux441-ux43fux43eux43cux43eux449ux44cux44e-passk}}

Некоторые впечатляющие демонстрации AI-агентов автор называет
«техническими чудесами»: если предоставить системе достаточно времени и
попыток, она выполняет задачу, которую прежде было трудно
автоматизировать. В исследовательской работе ценность может принести
одна проверенная успешная траектория. Метрика
\textbf{\href{mailto:Pass@k}{\nolinkurl{Pass@k}}} соответствует такому
способу использования: одно и то же задание выполняется \(k\) раз и
считается пройденным, если успешной была хотя бы одна попытка. Для
заданий с непрерывной оценкой также можно сообщать наивысший результат
среди \(k\) попыток, обозначаемый как
\textbf{\href{mailto:Best@k}{\nolinkurl{Best@k}}}.

Например, Anthropic организовала совместную работу нескольких
экземпляров Claude над созданием компилятора C, продемонстрировав, что
длительное выполнение, обратная связь от тестов и распределение задач
позволяют решать сложные инженерные задачи. При поиске математического
контрпримера исследователь также может разрешить модели многократно
повторять попытки, а затем независимо проверить справедливость одного из
найденных контрпримеров. При проверке безопасности ПО можно
перепроверить подозрительные проблемы, о которых сообщила модель, а в
творческих задачах с открытым результатом --- выбрать лучшую из
нескольких работ.\footnote{Anthropic, Building a C compiler with a team
  of parallel Claudes, 2026-02-05; материал посвящён совместной работе
  нескольких экземпляров, длительному выполнению и обратной связи от
  тестов.}

Прикладные продукты также демонстрируют свои возможности с помощью
сложных задач. Manus предоставляет агенту виртуальный компьютер,
позволяя пользователю наблюдать, как тот управляет интерфейсом, ищет
материалы и последовательно выполняет работу. OpenClaw принимает
поручения через приложения для обмена мгновенными сообщениями,
обращается к настроенным и разрешённым файлам и сервисам, сообщает о
ходе работы и запрашивает дополнительные сведения, а механизм расписания
может активировать его для проверки списка дел. Благодаря таким
конкретным действиям пользователю проще понять назначение агента.

По мнению автора, демонстрация способности выполнять задачи новых типов
--- один из важных факторов привлекательности этих продуктов для
пользователей.

\hypertarget{ux43dux430ux434ux451ux436ux43dux43eux441ux442ux44c-ux432-ux431ux438ux437ux43dux435ux441ux435-ux432ux43dux438ux43cux430ux43dux438ux435-ux43a-passk}{%
\subsection{Надёжность в бизнесе: внимание к
Pass\^{}k}\label{ux43dux430ux434ux451ux436ux43dux43eux441ux442ux44c-ux432-ux431ux438ux437ux43dux435ux441ux435-ux432ux43dux438ux43cux430ux43dux438ux435-ux43a-passk}}

Бизнес-процессы обычно требуют правильного результата при каждом из
нескольких запусков. \textbf{Pass\^{}k} измеряет случай, когда одно и то
же задание выполняется \(k\) раз и каждый запуск оказывается успешным;
эту метрику можно читать как «pass power k». Условия прохождения задания
могут одновременно учитывать итоговый результат, соблюдение политик и
ограничения прав доступа; каждый запуск проверяется по одним и тем же
условиям.

Если запуски независимы друг от друга, а вероятность успеха каждого из
них равна \(p\), то:

\[
\mathrm{Pass@k}=1-(1-p)^k,\qquad
\mathrm{Pass}^{k}=p^k.
\]

Например, при \(p=0.6\) и \(k=5\)
\href{mailto:Pass@5}{\nolinkurl{Pass@5}} составляет около 99,0\%, а
Pass\^{}5 --- около 7,8\%. Первая метрика показывает, что среди пяти
попыток обычно удаётся найти одну успешную, а вторая --- что вероятность
успеха всех пяти попыток по-прежнему очень мала.

Набор фактических результатов сначала можно рассмотреть без формул.
Предположим, что пять попыток одного задания дали результаты «успех,
неудача, успех, успех, неудача». Для этой группы
\href{mailto:Pass@5}{\nolinkurl{Pass@5}} равен 1, Pass\^{}5 --- 0, а
оценка доли успеха одной попытки составляет 3/5. Приведённые выше 99,0\%
и 7,8\% отвечают на другой вопрос: если известно, что вероятность успеха
каждой попытки равна 0,6, какая доля групп при многократном выполнении
серий по пять попыток будет соответствовать каждому из критериев?

Те же записи позволяют ответить на другой, более простой для проверки
вопрос: какая доля пар окажется успешной, если выбирать по две попытки
из этих пяти? Представим каждый запуск отдельной карточкой. При выборе
двух карточек без возвращения существует десять возможных пар. Обе
попытки будут неудачными только в одной паре, состоящей из двух карточек
с неудачами, поэтому хотя бы один успех будет в девяти парах, а
выборочная оценка \href{mailto:Pass@2}{\nolinkurl{Pass@2}} составит
9/10. Из трёх карточек с успехами можно составить три пары, поэтому
выборочная оценка Pass\^{}2 равна 3/10. Именно так τ²-bench рассчитывает
для каждого задания долю комбинаций, в которых успешны все попытки: если
задание выполнялось \(n\) раз и завершилось успехом \(s\) раз, то при
выборе \(k\) попыток доля комбинаций, где успешны все попытки, равна
\(\binom{s}{k}/\binom{n}{k}\), где \(1\le k\le n\).\footnote{В
  экспериментальной версии τ²-bench (ревизия 8d005b0) Pass\^{}k
  вычисляется отдельно для каждого задания в \texttt{agent\_metrics.py}.
  Если число успешных попыток меньше \(k\), количество комбинаций, в
  которых успешны все попытки, равно нулю.}

Теперь вернёмся к результату 4/5 из эксперимента 7-1. Пять траекторий
относятся к пяти разным заданиям, и у каждого задания есть лишь одна
карточка, поэтому составить даже одну пару невозможно. Чтобы оценить
Pass\^{}2 этих заданий, необходимо добавить независимые запуски для
каждого из них, а затем отдельно вычислить и агрегировать результаты.
Если просто возвести 0,8 в квадрат, это будет означать дополнительное
предположение, что вероятность успеха одинакова для всех этих заданий, а
отдельные попытки независимы друг от друга.

При вычислениях необходимо сохранять единообразие статистической
единицы. Пять выборок одного задания измеряют воспроизводимость
результата для этого задания; на пять последовательных разных задач в
промышленной среде также влияют сложность задач и распределение трафика.
Если вероятности успеха различаются между заданиями, можно сначала
вычислить для каждого из них метрику повторного успеха, а затем
агрегировать результаты с учётом весов заданий. Общие состояния, кэш или
одинаковые сбои также создают зависимость между запусками.

\textbf{Как выбирать знаменатель, если разрешён ответ «информации
недостаточно»?} Представим, что помощник по работе с литературой получил
100 вопросов и ответил на 80 из них, причём 68 ответов оказались
правильными. Для читателя, который полагается на предоставленные ответы,
важна доля 68/80, то есть 85\%; тому, кому необходимо выполнить весь
пакет заданий, важна также доля 68/100, то есть 68\%. Ещё для 20
вопросов требуется найти дополнительные доказательства или передать их
человеку.

В PaperQA2 при ответах на 248 вопросов о научной литературе из LitQA2
также можно было выбрать вариант «информации недостаточно». В статье
приведены средние результаты трёх запусков: когда система давала
содержательный ответ, доля правильных ответов составляла 85,2\%; при
расчёте по всем вопросам она составляла 66,0\%, а ещё в 21,9\% случаев
система выбирала ответ «информации недостаточно». Для девяти
участников-людей соответствующие показатели правильных ответов составили
73,8\% и 67,7\%. Согласно статье, различие по первому показателю было
статистически значимым, а по второму статистической значимости не
достигло.\footnote{Skarlinski и др., \emph{Language Agents Achieve
  Superhuman Synthesis of Scientific Knowledge}, 2024, см. раздел 2 и
  рисунок 2.}

Две траектории с одинаковым конечным состоянием могут иметь разную
стоимость. Рассмотрим проверку счёта: в первом случае правильная сумма
была получена одним запросом. Во втором тот же ответ был дан только
после двух запросов. Оба запроса содержали один и тот же идентификатор
запроса, а первый завершился по тайм-ауту, поэтому повторная попытка
помогла восстановить выполнение задания. Однако если результат первого
запроса уже был успешно доставлен, второй запрос лишь напрасно увеличил
время ожидания и расходы на инструменты.

Если выполнение связано с записью или другими внешними эффектами, перед
каждым испытанием следует восстанавливать заданное начальное состояние,
выполнять операцию в тестовой среде и регистрировать все попытки. Если
продукт допускает повторную попытку после неудачи, при оценке процесса
необходимо учитывать число повторов, действия по восстановлению и
стоимость выполнения; исходные записи о неудачах следует сохранять,
чтобы различать результат одной попытки и результат после
восстановления.

\hypertarget{ux441ux440ux435ux434ux430-ux43eux446ux435ux43dux43aux438}{%
\section{Среда
оценки}\label{ux441ux440ux435ux434ux430-ux43eux446ux435ux43dux43aux438}}

Предположим, нужно повторно проверить тот запуск, в котором ранее была
выбрана неверная линия связи. Проверяющий должен видеть того же
пользователя, те же две линии и тот же исходный тариф. Если при первом
запуске одна из линий уже была изменена, то при простом повторном
запуске невозможно определить, чем вызвано изменение результата: моделью
или базой данных. Среда оценки отвечает за создание исходной точки
задачи, предоставление интерфейсов взаимодействия, запись процесса и
проверку конечного состояния.

\hypertarget{ux43fux44fux442ux44c-ux43aux43eux43cux43fux43eux43dux435ux43dux442ux43eux432}{%
\subsection{Пять
компонентов}\label{ux43fux44fux442ux44c-ux43aux43eux43cux43fux43eux43dux435ux43dux442ux43eux432}}

На примере задачи telecom среду оценки можно разделить на пять частей.

\textbf{Датасет (Dataset)} хранит определение задачи: сведения о
пользователе, требования к поведению, начальное состояние и критерии
приёмки. Одна запись соответствует одному тестовому сценарию, а во время
выполнения информация распределяется между участниками согласно их
ролям.

\textbf{Состояние среды (Environment State)} хранит изменяемую
информацию в процессе выполнения задачи, например сведения о клиентах,
линиях связи, тарифах и счетах, а также об авиарежиме, роуминге и режиме
экономии трафика на устройстве. Перед началом нового запуска среда
сначала восстанавливает базовые данные, а затем выполняет
\texttt{initialization\_actions}, создавая неисправность,
предусмотренную текущей задачей. Очистка и инициализация должны быть
согласованы, чтобы пополнение пакета трафика или изменения настроек из
предыдущего запуска не переходили в следующий.

\textbf{Интерфейсы инструментов (Tools)} определяют, что участники могут
читать и изменять. Сотрудник поддержки может искать клиента, проверять
использование трафика, пополнять пакет трафика или переводить обращение
на оператора; пользователь может просматривать состояние устройства,
переключать настройки и измерять скорость. Гранулярность инструментов
должна соответствовать проверяемой способности: если требуется оценить
диагностику и планирование, AI-агент должен комбинировать запросы и
операции; если же напрямую предоставить инструмент «исправить проблему с
доступом в интернет», проверка будет сосредоточена на том, когда
вызывать эту способность.

\textbf{Критерии оценки (Rubric)} определяют, как оценивать результат. В
задаче telecom можно проверять действия, состояние среды, обязательную
коммуникацию и требования к естественному языку, а затем с помощью
\texttt{reward\_basis} указывать, какие проверки учитываются в итоговой
оценке.

\textbf{Протокол выполнения (Interaction Protocol)} определяет порядок
реплик, вызовов инструментов и завершения. Симулятор пользователя может
вывести \texttt{\#\#\#STOP\#\#\#}, чтобы завершить взаимодействие, а
раннер дополнительно задаёт бюджеты, например ограничение на число
раундов. Причина завершения и оценка задачи записываются отдельно: после
прекращения диалога пользователем проверки среды по-прежнему определяют,
достигнута ли цель; при исчерпании бюджета также необходимо сохранить
конечное состояние и незавершённые шаги.

\hypertarget{ux441ux440ux435ux434ux44b-ux43eux446ux435ux43dux43aux438-ux441-ux432ux437ux430ux438ux43cux43eux434ux435ux439ux441ux442ux432ux438ux435ux43c-ux447ux435ux43bux43eux432ux435ux43aux430-ux438-ux430ux433ux435ux43dux442ux430-ux438-ux441-ux432ux44bux437ux43eux432ux430ux43cux438-ux438ux43dux441ux442ux440ux443ux43cux435ux43dux442ux43eux432}{%
\subsection{Среды оценки с взаимодействием человека и агента и с
вызовами
инструментов}\label{ux441ux440ux435ux434ux44b-ux43eux446ux435ux43dux43aux438-ux441-ux432ux437ux430ux438ux43cux43eux434ux435ux439ux441ux442ux432ux438ux435ux43c-ux447ux435ux43bux43eux432ux435ux43aux430-ux438-ux430ux433ux435ux43dux442ux430-ux438-ux441-ux432ux44bux437ux43eux432ux430ux43cux438-ux438ux43dux441ux442ux440ux443ux43cux435ux43dux442ux43eux432}}

Задача telecom включает симулятор пользователя, который постоянно
отвечает и управляет устройством. Задачи исправления кода, анализа
данных и решения математических задач, напротив, часто запускаются
инструкцией, после чего агент взаимодействует преимущественно с
инструментами. Для сред обоих типов необходимы задача, состояние,
интерфейсы, критерии оценки и протокол выполнения. Различие состоит в
том, требуется ли дополнительно симулировать собеседника и каким образом
обе стороны согласуют действия.

В задачах программирования результат можно проверять тестами, в задачах
с базами данных --- запросами к изменённым записям, а открытый анализ
может потребовать оценки человеком или моделью. В задачах со
взаимодействием человека и агента нередко приходится объединять
доказательства двух типов. Например, восстановилась ли мобильная сеть,
определяется по состоянию устройства; уточнил ли сотрудник поддержки
номер и помог ли пользователю выполнить необходимые действия --- по
диалогу. То же относится к возврату средств на
\crossreflink{fig:7-3}{рисунке 7-3}: отмену заказа можно проверить в
базе данных, но сообщение сотрудника поддержки о сроке поступления
средств приходится искать в диалоге.

\begin{figure}
\hypertarget{fig:7-3}{%
\centering
\includegraphics{images/revised/fig7-3.pdf}
\caption{Рисунок 7-3. Среды оценки с вызовами инструментов и
взаимодействием с человеком}\label{fig:7-3}
}
\end{figure}

В Verifiers работу с задачами этих двух типов можно начинать через
разные интерфейсы среды. В таблице 7-1 перечислены интерфейсы версии
v0.1.9 для обработки однораундовых ответов, многораундовых вызовов
инструментов, внедрения состояния выполнения и управления ресурсами
песочницы, что упрощает их сопоставление с кодом.\footnote{Prime
  Intellect, \texttt{ToolEnv} и \texttt{StatefulToolEnv} из Verifiers
  v0.1.9, демонстрирующие многораундовые траектории, диспетчеризацию
  инструментов и внедрение параметров состояния.}

Таблица 7-1. Сравнение типов сред Verifiers

\begin{longtable}[]{@{}
  >{\raggedright\arraybackslash}p{(\columnwidth - 6\tabcolsep) * \real{0.2115}}
  >{\raggedright\arraybackslash}p{(\columnwidth - 6\tabcolsep) * \real{0.2586}}
  >{\raggedright\arraybackslash}p{(\columnwidth - 6\tabcolsep) * \real{0.2805}}
  >{\raggedright\arraybackslash}p{(\columnwidth - 6\tabcolsep) * \real{0.2494}}@{}}
\toprule\noalign{}
\begin{minipage}[b]{\linewidth}\raggedright
\hspace{0pt}Тип среды
\end{minipage} & \begin{minipage}[b]{\linewidth}\raggedright
\hspace{0pt}Основная задача
\end{minipage} & \begin{minipage}[b]{\linewidth}\raggedright
\hspace{0pt}Обработка инструментов и состояния
\end{minipage} & \begin{minipage}[b]{\linewidth}\raggedright
\hspace{0pt}Типичный сценарий
\end{minipage} \\
\midrule\noalign{}
\endhead
\bottomrule\noalign{}
\endlastfoot
\hspace{0pt}\texttt{SingleTurnEnv} & \hspace{0pt}Организация
однораундового ответа & \hspace{0pt}Оценка после одного ответа &
\hspace{0pt}Математические задачи, однораундовые вопросы и ответы \\
\hspace{0pt}\texttt{ToolEnv} & \hspace{0pt}Организация многораундового
взаимодействия с инструментами & \hspace{0pt}Сохранение траектории
выполнения, диспетчеризация инструментов и получение результатов &
\hspace{0pt}Поиск и обобщение информации \\
\hspace{0pt}\texttt{StatefulToolEnv} & \hspace{0pt}Подключение состояния
выполнения к работе инструментов & \hspace{0pt}Дополнение параметров
инструмента или обновление состояния с учётом текущего состояния &
\hspace{0pt}Привязка базы данных или идентификатора ресурса к текущей
задаче \\
\hspace{0pt}\texttt{SandboxEnv} & \hspace{0pt}Управление изолированной
средой выполнения & \hspace{0pt}Выделение песочницы для запуска и
использование доступных в ней инструментов & \hspace{0pt}Выполнение
кода, проверка файлов и результатов тестов \\
\end{longtable}

Например, если две задачи одновременно изменяют базу данных, раннер
может внедрять идентификатор соответствующей базы данных в вызовы
инструментов, чтобы каждая траектория работала со своими данными. При
ошибке выполнения инструмент должен возвращать информацию, позволяющую
принять решение о следующем шаге, например сообщать об отсутствии
объекта, недопустимых параметрах или тайм-ауте сервиса. Система оценки
сохраняет наблюдения, вызовы, результаты выполнения и вознаграждения,
благодаря чему при анализе ошибку можно локализовать до конкретного
шага.

\hypertarget{ux43fux440ux43eux435ux43aux442ux438ux440ux43eux432ux430ux43dux438ux435-ux43eux446ux435ux43dux43eux447ux43dux43eux433ux43e-ux43dux430ux431ux43eux440ux430-ux434ux430ux43dux43dux44bux445}{%
\section{Проектирование оценочного набора
данных}\label{ux43fux440ux43eux435ux43aux442ux438ux440ux43eux432ux430ux43dux438ux435-ux43eux446ux435ux43dux43eux447ux43dux43eux433ux43e-ux43dux430ux431ux43eux440ux430-ux434ux430ux43dux43dux44bux445}}

Теперь превратим пять задач из области telecom в оценочный набор,
способный служить ориентиром при выпуске. Если все новые задачи сводятся
к «подключить роуминг для указанного номера», система может получить
высокий балл, даже если так и не научилась различать две линии одного
аккаунта. Если по результатам оценки предстоит решать, можно ли
выпускать службу поддержки, задачи должны опираться на реальные
обращения: включать как стандартные операции, так и сложные случаи, в
которых легко перепутать линии. При оценке также необходимо проверять,
какую именно линию в итоге изменила служба поддержки. После выпуска эти
задачи нужно поддерживать, чтобы они продолжали отражать работу, которую
пользователи действительно поручают системе.

\hypertarget{ux441ux43eux43fux43eux441ux442ux430ux432ux43bux435ux43dux438ux435-ux43fux43eux434ux445ux43eux434ux43eux432-ux43a-ux43fux440ux43eux435ux43aux442ux438ux440ux43eux432ux430ux43dux438ux44e-ux431ux435ux43dux447ux43cux430ux440ux43aux43eux432}{%
\subsection{Сопоставление подходов к проектированию
бенчмарков}\label{ux441ux43eux43fux43eux441ux442ux430ux432ux43bux435ux43dux438ux435-ux43fux43eux434ux445ux43eux434ux43eux432-ux43a-ux43fux440ux43eux435ux43aux442ux438ux440ux43eux432ux430ux43dux438ux44e-ux431ux435ux43dux447ux43cux430ux440ux43aux43eux432}}

В таблице 7-2 шесть бенчмарков сопоставлены по оцениваемым способностям,
источникам задач, средам выполнения и способам проверки.

Таблица 7-2. Ключевые архитектурные решения нескольких бенчмарков
AI-агентов

\begin{longtable}[]{@{}
  >{\raggedright\arraybackslash}p{(\columnwidth - 8\tabcolsep) * \real{0.2200}}
  >{\raggedright\arraybackslash}p{(\columnwidth - 8\tabcolsep) * \real{0.1740}}
  >{\raggedright\arraybackslash}p{(\columnwidth - 8\tabcolsep) * \real{0.2026}}
  >{\raggedright\arraybackslash}p{(\columnwidth - 8\tabcolsep) * \real{0.1834}}
  >{\raggedright\arraybackslash}p{(\columnwidth - 8\tabcolsep) * \real{0.2200}}@{}}
\toprule\noalign{}
\begin{minipage}[b]{\linewidth}\raggedright
\hspace{0pt}Бенчмарк
\end{minipage} & \begin{minipage}[b]{\linewidth}\raggedright
\hspace{0pt}Оцениваемая способность
\end{minipage} & \begin{minipage}[b]{\linewidth}\raggedright
\hspace{0pt}Источник и организация задач
\end{minipage} & \begin{minipage}[b]{\linewidth}\raggedright
\hspace{0pt}Среда выполнения
\end{minipage} & \begin{minipage}[b]{\linewidth}\raggedright
\hspace{0pt}Способ проверки
\end{minipage} \\
\midrule\noalign{}
\endhead
\bottomrule\noalign{}
\endlastfoot
\hspace{0pt}τ²-\allowbreak{}\hspace{0pt}bench &
\hspace{0pt}Взаимодействие со службой поддержки и работа с инструментами
& \hspace{0pt}Сочетания предметных задач и неисправностей &
\hspace{0pt}Симулятор пользователя, устройства и
бизнес-\allowbreak{}\hspace{0pt}база данных & \hspace{0pt}Поэлементные
проверки с оценкой по \texttt{reward\_}\allowbreak{}\texttt{basis} \\
\hspace{0pt}SWE-\allowbreak{}\hspace{0pt}bench Verified &
\hspace{0pt}Исправление проблем в коде & \hspace{0pt}Реальные GitHub
issue, отобранные вручную & \hspace{0pt}Репозиторий кода и тестовая
среда & \hspace{0pt}Тесты
\texttt{FAIL\_}\allowbreak{}\texttt{TO\_}\allowbreak{}\texttt{PASS} и
\texttt{PASS\_}\allowbreak{}\texttt{TO\_}\allowbreak{}\texttt{PASS} \\
\hspace{0pt}AndroidWorld & \hspace{0pt}Работа с приложениями Android &
\hspace{0pt}Параметризованные шаблоны задач & \hspace{0pt}Эмулятор
Android с реальными приложениями & \hspace{0pt}Чтение состояния
устройства, файлов или внутренних данных приложения \\
\hspace{0pt}OSWorld & \hspace{0pt}Работа с рабочим столом и несколькими
приложениями & \hspace{0pt}Реальные сценарии использования с настроенной
начальной точкой задачи & \hspace{0pt}Виртуальная машина и настольные
приложения & \hspace{0pt}Проверка состояния выполнения задачи и
полученных артефактов \\
\hspace{0pt}Terminal-\allowbreak{}\hspace{0pt}Bench &
\hspace{0pt}Инженерные задачи в терминале & \hspace{0pt}Задачи и тесты,
разработанные вручную & \hspace{0pt}Контейнеры и необходимые для задач
зависимости & \hspace{0pt}Файлы, выполнение программ и результаты
тестов \\
\hspace{0pt}GAIA & \hspace{0pt}Обработка информации универсальным
помощником & \hspace{0pt}Составленные вручную вопросы, часть из которых
снабжена материалами & \hspace{0pt}Интернет, файлы и вспомогательные
инструменты & \hspace{0pt}Сопоставление ответов после нормализации
формата \\
\end{longtable}

\hypertarget{ux432ux430ux43bux438ux434ux430ux442ux43eux440ux44b}{%
\subsection{Валидаторы}\label{ux432ux430ux43bux438ux434ux430ux442ux43eux440ux44b}}

Валидатор связывает цель задачи с проверяемыми доказательствами. В
задачах с кодом необходимо запускать тесты, в задачах для рабочего стола
--- считывать состояние приложений или файлов, а в задачах поиска
информации --- сверять ответ и источники. \textbf{Приёмка должна
основываться на независимо полученных доказательствах.} Агент может
подготовить полный отчёт и заявить, что все шаги выполнены, однако
валидатор всё равно должен фактически запустить полученный артефакт,
считать состояние или проверить источники. Отчёт о выполнении помогает
определить объекты проверки, но оценка выставляется по результатам самих
проверок.

Представим, что пользователь открывает платёжное приложение и просит
помощника «одобрить все запросы на получение денег, отправленные соседом
по квартире в этом месяце». Допустим, ожидают обработки три запроса: от
соседа за текущий месяц, от того же соседа за прошлый месяц и от коллеги
за текущий месяц. Первый нужно одобрить, а два других должны остаться в
ожидании. Если проверять только первый запрос, полный балл получит даже
помощник, одобривший все три; если проверять только два последних,
задачу пройдёт помощник, который вообще ничего не сделал. Такая задача
представлена в AppWorld в симуляции платёжного приложения. При
проектировании валидатора можно сначала по данным до операции выделить
группу запросов, которую пользователь разрешил обработать, затем
проверить, что все они приняты, а в завершение выполнить обратную
проверку по записям после операции: какие ещё запросы затронул
помощник?\footnote{Trivedi и др., \emph{AppWorld}, ACL 2024; задача и
  условия приёмки приведены на рисунке 2.}

Научный поиск тоже может выдавать легко проверяемый артефакт. В
FunSearch требуется выбрать как можно больше точек в восьмимерном
пространстве, где каждая координата может принимать только значения 0, 1
или 2, а вычисления выполняются по модулю 3; среди выбранных точек
никакие три не должны лежать на одной прямой. Исследователи подготовили
каркас программы, пошагово выбирающей точки, поручили модели предлагать
функцию ранжирования точек-кандидатов, а результат оценивали по размеру
фактически построенного допустимого множества. На момент публикации
статьи известная конструкция содержала 496 точек, а новая программа
выдала список из 512 точек.\footnote{Romera-Paredes и др.,
  \emph{Mathematical discoveries from program search with large language
  models}, Nature, онлайн-публикация 2023 года. Авторы также отдельно
  опубликовали восьмимерное множество точек и вычислительный блокнот для
  проверки.}

Можно скачать только этот список и независимо проверить результат.
Сначала следует убедиться, что он содержит 512 различных восьмимерных
векторов и каждая координата находится в заданном диапазоне. Затем
возьмём любые две разные точки \(a,b\). На проходящей через них прямой
лежат всего три точки, причём третья равна \(c=-a-b \pmod{3}\).
Например, в двумерном случае для \((0,0)\) и \((1,2)\) третьей точкой
будет \((2,1)\). Если покоординатно применить ту же операцию в восьми
измерениях и перебрать 130 816 пар различных точек из списка, ни одна из
вычисленных третьих точек не окажется в списке. Этого достаточно, чтобы
доказать возможность выбрать не менее 512 точек. Чтобы доказать
невозможность выбора 513 точек, дополнительно требуется доказательство
верхней границы.

\textbf{Два вида тестов SWE-bench отдельно проверяют исправление и
регрессии.} \texttt{FAIL\_TO\_PASS} --- тесты, которые падали до
исправления и должны пройти после него; \texttt{PASS\_TO\_PASS} ---
тесты, которые проходили ранее и должны продолжать проходить после
исправления. Первые проверяют заявленную проблему, вторые --- сохранение
поведения, охваченного существующими тестами.\footnote{OpenAI,
  Introducing SWE-bench Verified, 2024; описание классификации тестов,
  стандартов разметки и ручного отбора.}

\textbf{OSWorld проверяет артефакты приложений посредством исполняемых
валидаторов.} В первоначальной версии предусмотрено 134 исполняемые
функции оценки, которые в сочетании с конфигурацией задачи считывают
файлы, данные приложений или состояние системы.\footnote{В описании
  проекта OSWorld перечислены инициализация задач и исполняемая оценка;
  134 функции оценки относятся к статистике первоначальной версии.}
Например, в задаче на изменение таблицы можно проверить заданные ячейки
и формулы, а в задаче на обновление настроек браузера --- считать
соответствующую конфигурацию.

\textbf{Terminal-Bench принимает инженерные задачи по результатам
выполнения.} Сценарий сборки ядра Linux объединяет изменение исходного
кода, компиляцию, создание initramfs и запуск QEMU. Конкретными
контрольными точками служат Linux 6.9, пользовательский \texttt{printk}
в \texttt{start\_kernel} и сообщение в журнале загрузки.

В пятой главе рассматривалась сверка выручки по заказам. Но с чего
начать проверку, если помощник сдаёт квартальный отчёт по двадцати
каналам? В Tencent WorkBuddy Bench есть публичная задача на сверку. Обе
рабочие таблицы содержат данные по каналам: расчётную сумму нужно
сравнивать с предыдущим кварталом, то есть по отношению к
предшествующему периоду, а сумму поступлений --- с тем же кварталом
предыдущего года, то есть год к году. Даже если в обоих случаях указать
в качестве базового периода предыдущий квартал, заголовки и число
каналов могут быть совершенно правильными. Например, для канала «Магазин
красивой жизни» в таблице расчётов после удаления повторяющихся строк и
записей без квартала и последующей агрегации по каналу и кварталу можно
вычислить, что общая сумма сделок выросла с 1 569 231,00 до 1 787
311,03, то есть на 218 080,03. Оценочная программа проверяет охват
каналов в обеих таблицах, а затем выборочно сверяет значения для двух
заданных каналов. Проверка охвата распространяется на все двадцать
каналов, а числовая проверка --- только на два выбранных.\footnote{Офисная
  задача Tencent WorkBuddy Bench \texttt{channel-period-compare-L4-017}.
  Пакет данных версии 3e8ad83 содержит исходную рабочую книгу, эталонный
  ответ и оценочную программу.}

Валидатор можно проверить с помощью задачи, чувствительной ко времени.
Предположим, пользователь просит уведомить его, когда цена номера в
определённом отеле опустится ниже 950 юаней. В 10:00 цена составляет
1000 юаней, в 10:10 снижается до 900, а в 10:15 снова возвращается к
1000. Если в 10:12 отправить уведомление о снижении цены с указанием
типа номера и ссылкой для бронирования, пользователь ещё успеет оформить
заказ. Если же каждый день сообщать «цена пока не снизилась», но
пропустить именно эти пять минут, поручение останется невыполненным даже
после недели непрерывной работы. При оценке следует отдельно учитывать
число окон, в которых условие выполнялось, количество своевременных
уведомлений внутри таких окон и число случаев, когда пользователя
беспокоили при отсутствии изменений. Если сам оценщик сохранял цену лишь
раз в день, число пропущенных окон определить вообще невозможно. Чтобы
измерять пропуски, сначала нужно подготовить полное воспроизведение цен
с временными метками, а затем позволить проверяемой системе опрашивать
их с собственной частотой.

\hypertarget{ux440ux430ux437ux434ux435ux43bux435ux43dux438ux435-ux437ux430ux434ux430ux447-ux43fux43e-ux441ux43bux43eux436ux43dux43eux441ux442ux438}{%
\subsection{Разделение задач по
сложности}\label{ux440ux430ux437ux434ux435ux43bux435ux43dux438ux435-ux437ux430ux434ux430ux447-ux43fux43e-ux441ux43bux43eux436ux43dux43eux441ux442ux438}}

Набор задач должен охватывать разные уровни сложности и объяснять, чем
обусловлена сложность. Базовые задачи проверяют, не деградировали ли
способности, а трудные оставляют пространство для измерения дальнейшего
прогресса моделей. Если все задачи быстро начинают проходить, оценочный
набор перестаёт помогать при выборе. Задачи можно разделять по видам
инструментов, длине последовательности действий, степени
распределённости информации, неопределённости среды или необходимым
профессиональным знаниям, а затем анализировать, при каких условиях
система заметно деградирует.

466 вопросов GAIA разделены на три уровня. Более простые задачи обычно
требуют небольшого числа инструментов и шагов, а более сложные ---
многоэтапной интеграции информации и сложного сочетания действий. В
исходной статье точность аннотаторов-людей на задачах трёх уровней
составляла соответственно 93,9\%, 91,8\% и 87,3\%, а результаты GPT-4 с
плагинами --- 30,3\%, 9,7\% и 0\%.\footnote{Grégoire Mialon и др., GAIA:
  A Benchmark for General AI Assistants; в исходной статье описаны
  уровни задач, правила оформления ответов и конфигурации для людей и
  моделей.}

Разброс задач Terminal-Bench столь же велик: регистрация модели требует
работы с MLflow, многокомпонентное развёртывание --- координации
Git-сервиса и веб-сервиса, а восстановление паролей и дифференциальный
анализ FEAL добавляют специализированные алгоритмические знания.

Бизнес-задачи также должны охватывать границы политик. Например, если
пользователь утверждает, что служба поддержки уже одобрила отмену
заказа, система должна проверить разрешение и состояние заказа, а затем
принять решение об отмене в соответствии с политикой. Такие сценарии
проверяют, как агент использует доказательства при обработке
противоречивой информации, и могут включаться в набор вместе со
стандартными задачами обслуживания.

\hypertarget{ux43fux440ux435ux434ux43eux442ux432ux440ux430ux449ux435ux43dux438ux435-ux443ux442ux435ux447ux43aux438-ux434ux430ux43dux43dux44bux445}{%
\subsection{Предотвращение утечки
данных}\label{ux43fux440ux435ux434ux43eux442ux432ux440ux430ux449ux435ux43dux438ux435-ux443ux442ux435ux447ux43aux438-ux434ux430ux43dux43dux44bux445}}

\textbf{Сочетание нескольких источников и специальные вложения.} Если
модель может найти в результатах поиска условие задачи и готовый ответ,
насколько единственный правильный ответ доказывает её способность к
поиску и рассуждению? GAIA снижает такую вероятность с помощью
многоэтапного поиска и специальных вложений. Например, сначала нужно
опознать астронавта на астрономическом изображении за определённую дату,
затем выяснить его группу астронавтов, сравнить продолжительность
пребывания её участников в космосе и наконец вывести фамилии и число
минут в требуемом формате. Точки с запятой, округление чисел и
разделители тысяч входят в соглашение о формате ответа. Специально
подготовленные PDF-файлы, аудиозаписи и изображения увеличивают
разнообразие сочетаний входных данных. Публичные задачи могут попасть в
обучающие материалы. Поэтому необходимо также сохранять отдельный набор
задач исключительно для оценки и фиксировать, какие задачи уже
использовались при обучении.\footnote{Grégoire Mialon и др., GAIA: A
  Benchmark for General AI Assistants; в исходной статье описаны уровни
  задач, правила оформления ответов и конфигурации для людей и моделей.}

\textbf{Параметризованные экземпляры.} AndroidWorld предоставляет для
задач логику инициализации, проверки выполнения и очистки, а различные
экземпляры создаёт посредством параметров.\footnote{Christopher Rawles и
  др., AndroidWorld: A Dynamic Benchmarking Environment for Autonomous
  Agents; описание параметризованных задач и приёмки по состоянию
  устройства.} Например, в задаче «Изменить телефон контакта
\texttt{{[}CONTACT\_NAME{]}} на \texttt{{[}NEW\_PHONE{]}}» можно менять
контакт и номер, сохраняя неизменной цель операции. Изменение параметров
позволяет проверить, действует ли модель в соответствии с текущей
задачей; фиксация части параметров при изменении остальных также
помогает исследовать влияние отдельных факторов.

\textbf{Идентификаторы-канарейки.} В файлы бенчмарка можно встраивать
уникальные canary GUID, помогающие обнаруживать признаки копирования
данных или их попадания в другие корпуса. Такие метки используются в
файлах задач Terminal-Bench. При анализе утечки необходимо также
проверить, встречала ли модель идентификатор в текущем контексте или
результатах поиска, и сопоставить это с происхождением данных.
Обнаружение метки может послужить поводом для расследования, но
принадлежность обучающих данных требует дополнительных доказательств.

\hypertarget{ux43aux43eux43dux442ux440ux43eux43bux44c-ux43aux430ux447ux435ux441ux442ux432ux430-ux438-ux434ux43eux43bux433ux43eux441ux440ux43eux447ux43dux43eux435-ux441ux43eux43fux440ux43eux432ux43eux436ux434ux435ux43dux438ux435}{%
\subsection{Контроль качества и долгосрочное
сопровождение}\label{ux43aux43eux43dux442ux440ux43eux43bux44c-ux43aux430ux447ux435ux441ux442ux432ux430-ux438-ux434ux43eux43bux433ux43eux441ux440ux43eux447ux43dux43eux435-ux441ux43eux43fux440ux43eux432ux43eux436ux434ux435ux43dux438ux435}}

Если одна и та же задача из области telecom вчера прошла, а сегодня нет,
помимо изменений агента нужно уметь ответить на следующие вопросы:
изменилось ли начальное состояние линии? Продолжает ли симулятор
пользователя действовать по тем же правилам? Проверяют ли условия оценки
прежнюю цель? Если проследить начальную точку, взаимодействие и конечное
состояние одного запуска, можно выделить пять направлений сопровождения.

\begin{enumerate}
\def\labelenumi{\arabic{enumi}.}
\tightlist
\item
  \textbf{Явно определить объём информации пользователя.} Отделить
  \texttt{known\_info} от требований к поведению, чтобы у запросов и
  уточнений была однозначная исходная информационная точка.
\item
  \textbf{Чётко сформулировать условия приёмки.} Выразить «сеть
  восстановлена» через состояние мобильной передачи данных, результат
  измерения скорости и оценку, проверяя эффект, которого действительно
  требует задача.
\item
  \textbf{Задать основания для ответов симулируемого пользователя.}
  Состояние устройства должно поступать из результатов инструментов, а
  правила эмоций и терпения --- определять организацию взаимодействия.
\item
  \textbf{Представить действия обеих сторон.} Пользователь и служба
  поддержки независимо изменяют среду, а для последующих решений
  получают новое состояние посредством доступных инструментов или
  диалога.
\item
  \textbf{Комбинировать задачи и неисправности.} Изменять сочетания
  неисправностей и параметры задач по явно заданным правилам генерации,
  а затем проверять разрешимость экземпляра и соответствие исходного
  состояния определению задачи.
\end{enumerate}

\textbf{SWE-bench Verified демонстрирует ручной отбор перед
публикацией.} Проект привлёк 93 разработчиков с опытом Python для
разметки 1699 случайно выбранных примеров. Основное внимание уделялось
тому, достаточно ли описания проблемы, не исключают ли тесты разумные
решения и насколько сложно исправление; каждый пример независимо
проверяли три аннотатора. Опубликованный набор Verified содержит 500
примеров.\footnote{OpenAI, Introducing SWE-bench Verified, 2024;
  описание классификации тестов, стандартов разметки и ручного отбора.}
Правила отбора и распределение сложности необходимо публиковать вместе,
чтобы пользователи могли понять, какие требования к способностям
сохранились после сокращения набора задач. Задачи агентов часто требуют
нескольких вызовов модели и операций со средой, поэтому неоднозначные
или неразрешимые задачи будут вновь и вновь расходовать бюджет при
каждом регрессионном прогоне. Автор рекомендует сначала устранять такой
шум и только затем увеличивать число задач. Небольшой качественный
оценочный набор лучше подходит для частых запусков и упрощает
интерпретацию неудач.

\textbf{OSWorld-Verified демонстрирует сопровождение после публикации.}
Примерно через 15 месяцев после выхода первоначальной версии проект
систематизировал и исправил более 300 замечаний; команда примерно из
десяти человек сосредоточенно работала над этим два месяца.\footnote{XLANG
  Lab, Introducing OSWorld-Verified, 2025; описание отзывов сообщества и
  изменений инициализации и валидаторов.}

Например, после изменения структуры веб-страницы необходимо обновить
извлечение состояния; перед началом задачи с документом --- убедиться,
что приложение загрузилось, файл открыт, а курсор находится в заданной
позиции; после изменения условий проверки --- выполнить регрессионную
проверку на правильных и неправильных артефактах.

\begin{quote}
\textbf{Эксперимент 7-2 ★: самостоятельное выполнение задач бенчмарков}

Выберите задачи из GAIA, AndroidWorld, SWE-bench Verified,
Terminal-Bench, OSWorld-Verified и τ²-bench: для каждого бенчмарка
возьмите по одной простой, средней и сложной задаче согласно принятому в
нём определению сложности. Сначала самостоятельно выполните задачи,
затем проверьте официальные оценки; зафиксируйте путь выполнения,
обнаруженные неоднозначности и итоговые артефакты.

В сопроводительных материалах сохранены 18 примеров работы Codex с этими
шестью видами бенчмарков, по которым можно сверить собственный подход. В
задаче GAIA о расстоянии между атомами формат округления разошёлся с
требованиями оценщика; в задаче AndroidWorld на умножение после
получения результата 840 страница всё ещё не перешла в требуемое
состояние успеха. В τ²-bench скрытый пользователь симулировался
локальной моделью. Из 18 примеров проверку прошли 13.

После выполнения проанализируйте: допускает ли задача несколько разумных
толкований, охватывает ли валидатор эти толкования и какие действия
после получения правильного промежуточного результата ещё необходимо
выполнить для прохождения приёмки.
\end{quote}

\hypertarget{ux442ux440ux438-ux438ux441ux442ux43eux447ux43dux438ux43aux430-ux43eux446ux435ux43dux43eux447ux43dux43eux433ux43e-ux43dux430ux431ux43eux440ux430}{%
\subsection{Три источника оценочного
набора}\label{ux442ux440ux438-ux438ux441ux442ux43eux447ux43dux438ux43aux430-ux43eux446ux435ux43dux43eux447ux43dux43eux433ux43e-ux43dux430ux431ux43eux440ux430}}

При производственной оценке можно сочетать публичные бенчмарки,
собственные бизнес-задачи и производственные траектории. Они выполняют
разные функции и вместе охватывают область применения системы.

\textbf{Публичные бенчмарки} подходят для предварительного сравнения
моделей, а также служат источником примеров проектирования задач,
параметризованной генерации и валидаторов. Автор использует их для
первичного отбора моделей-кандидатов, а окончательный выбор модели и
решение о выпуске принимает по результатам собственного бизнес-набора.

\textbf{Собственные бизнес-наборы} организуют задачи вокруг реальных
пользовательских целей, инструментов и политик и используются при выборе
модели и совершенствовании Harness. За основу можно взять структуру
τ²-bench: настроить предметные данные, пользовательские сценарии,
инструменты и проверки состояния, а затем откалибровать сложность задач
и поведение симулятора пользователя.

\textbf{Возврат производственных траекторий} помогает находить пробелы
по реальным запускам: исправления пользователей, отрицательные отзывы,
аномалии бизнес-состояния или ошибки, обнаруженные при проверке
траекторий. После обработки с надлежащим разрешением и обезличивания
человек проверяет причину неудачи, восстанавливает необходимые условия
задачи и добавляет случай в регрессионный набор. Пользовательская
обратная связь служит ориентиром для локализации проблемы, но оценочные
метки необходимо подтверждать по фактическому состоянию.

На начальном этапе можно использовать публичные бенчмарки и небольшое
число бизнес-задач, составленных вручную. После выпуска основным
источником новых регрессионных сценариев должны стать реальные ошибки и
новые бизнес-ситуации. Такие задачи непосредственно отражают проблемы
пользователей и позволяют постоянно восполнять условия, упущенные при
проектировании.

\hypertarget{ux43cux435ux442ux43eux434ux44b-ux430ux432ux442ux43eux43cux430ux442ux438ux437ux438ux440ux43eux432ux430ux43dux43dux43eux439-ux43eux446ux435ux43dux43aux438}{%
\section{Методы автоматизированной
оценки}\label{ux43cux435ux442ux43eux434ux44b-ux430ux432ux442ux43eux43cux430ux442ux438ux437ux438ux440ux43eux432ux430ux43dux43dux43eux439-ux43eux446ux435ux43dux43aux438}}

Автоматизированная оценка начинается с чёткого определения объекта
проверки. SWE-bench проверяет поведение кода с помощью набора тестов,
AndroidWorld запрашивает состояние приложения, а GAIA нормализует ответы
с учётом их типа и затем сравнивает их. Такие проверки можно выполнять
программно, поэтому они подходят для многократного запуска. В τ²-bench
сочетается несколько критериев: вызовы инструментов и состояние среды
можно проверить напрямую, а требования к естественно-языковому ответу
--- передать оценивающей модели. Одна и та же задача часто требует обоих
подходов.

Программные проверки также могут предоставлять диагностическую
информацию. Например, в задаче оператора связи из начала главы отдельная
проверка настроек роуминга, остатка трафика и операции пополнения
позволяет определить, какие условия ещё не выполнены. Чтобы установить
точный момент выбора неверной линии, необходимо связать результаты
проверок с траекторией вызовов.

Для открытых результатов требования к качеству необходимо формулировать
конкретнее. В ответе на жалобу следует изложить решение и учесть эмоции
и проблему клиента; исследовательский отчёт должен охватывать ключевые
факты, а его выводы --- подтверждаться источниками; ответ на основе
памяти должен правильно связывать людей, время и события. Для таких
требований можно с помощью \textbf{Rubric (критериев оценки)}
перечислить отдельные критерии, а затем поручить \textbf{LLM-as-a-Judge
(модели-оценщику)} выставить оценки на основе предоставленных
материалов.

\textbf{Условия, которые можно проверить программно, автор предпочитает
оформлять как программные проверки.} Формат, числовые значения, права
доступа и конечное состояние обычно имеют чёткие критерии; при
фиксированных входных данных и среде их можно проверять многократно, а
затраты легко оценить. В одном и том же ответе могут потребоваться оба
способа оценивания. Например, на \crossreflink{fig:7-4}{рисунке 7-4} для
проверки суммы возврата достаточно сравнить числа, но чтобы определить,
понятно ли сотрудник поддержки объяснил условия возврата, необходимо
прочитать весь текст. Этот аспект можно поручить оценивающей модели, а
затем откалибровать её оценки на выборке, проверенной людьми.

\begin{figure}
\hypertarget{fig:7-4}{%
\centering
\includegraphics{images/revised/fig7-4.pdf}
\caption{Рисунок 7-4. Спектр способов проверки: от детерминированной
проверки до оценки моделью}\label{fig:7-4}
}
\end{figure}

\hypertarget{llm-as-a-judge-ux43fux440ux435ux43eux431ux440ux430ux437ux43eux432ux430ux43dux438ux435-ux442ux440ux435ux431ux43eux432ux430ux43dux438ux439-ux43a-ux43aux430ux447ux435ux441ux442ux432ux443-ux432-ux43fux43eux44dux43bux435ux43cux435ux43dux442ux43dux443ux44e-ux43eux446ux435ux43dux43aux443}{%
\subsection{LLM-as-a-Judge: преобразование требований к качеству в
поэлементную
оценку}\label{llm-as-a-judge-ux43fux440ux435ux43eux431ux440ux430ux437ux43eux432ux430ux43dux438ux435-ux442ux440ux435ux431ux43eux432ux430ux43dux438ux439-ux43a-ux43aux430ux447ux435ux441ux442ux432ux443-ux432-ux43fux43eux44dux43bux435ux43cux435ux43dux442ux43dux443ux44e-ux43eux446ux435ux43dux43aux443}}

На \crossreflink{fig:7-5}{рисунке 7-5} служба поддержки правильно
указала сумму возврата и ожидаемый срок зачисления. Почему же за полноту
ответ получил лишь 2 балла? Критерии оценки также требовали сообщить,
куда будут возвращены средства, но в ответе было пропущено «тем же
способом, которым была произведена оплата». Чтобы определить, какой
именно пункт упущен, оценивающая модель должна одновременно видеть
задачу, критерии оценки и справочные материалы.

\begin{figure}
\hypertarget{fig:7-5}{%
\centering
\includegraphics{images/revised/fig7-5.pdf}
\caption{Рисунок 7-5. Конвейер LLM-as-a-Judge}\label{fig:7-5}
}
\end{figure}

Оценивающую модель необходимо предварительно откалибровать. Одна из
распространённых проблем --- \textbf{смещение в пользу длины}: при
сопоставимом содержании более длинный ответ может получить более высокую
оценку. Для проверки можно подготовить примеры с одинаковыми фактами и
выводами, но разной длиной изложения, и проверить, зависит ли оценка от
объёма. Если два ответа оцениваются вместе, можно менять их местами и
смотреть, не выберет ли оценивающая модель другой вариант. Если сама
задача требует краткости, это следует записать в Rubric и привести
пример ответа подходящей длины. Обнаруженная при массовом анализе
корреляция между высокой оценкой и длинным ответом может послужить
первым сигналом, после чего необходимо сопоставить объём информации,
правильность и сложность задач, чтобы выявить конкретную причину.

\textbf{Четыре принципа проектирования Rubric.} В работе Scale AI
«Rubrics as Rewards» предложены четыре принципа: опора на экспертные
знания, полное покрытие, взвешивание по важности и
самодостаточность.\footnote{Scale AI, «Использование Rubric для создания
  более качественных моделей», сентябрь 2025 года. Четыре критерия
  используются для организации стандартов оценки; условия безусловного
  отклонения и примеры, связанные с памятью пользователя, в этом разделе
  разработаны для конкретных задач.} При оценке AI-агентов их можно
реализовать следующим образом:

\begin{enumerate}
\def\labelenumi{(\arabic{enumi})}
\item
  \textbf{Опора на экспертные знания.} Сначала попросите специалиста,
  знакомого с предметной областью, составить эталонный ответ,
  перечислить ключевые факты и отношения, в которых легко допустить
  ошибку. Например, в задаче на память необходимо явно указать, кому
  принадлежит какой страховой полис и когда старый адрес был заменён
  новым.
\item
  \textbf{Полное покрытие.} Одновременно перечислите обязательную
  информацию, требования к определению отношений и форме ответа, а также
  ошибки высокого риска --- \textbf{ловушки (Pitfall)}. В ответе службы
  поддержки можно проверять, отреагировала ли она на запрос, указала ли
  следующий шаг и не выдумала ли обещание возврата средств.
\item
  \textbf{Взвешивание по важности.} Разделите критерии на обязательные
  (Essential), важные (Important) и необязательные (Optional), а для
  ловушек задайте правила вычитания баллов. Для ошибок, непосредственно
  приводящих к провалу задачи, можно дополнительно предусмотреть
  \textbf{вето (Veto)}. Например, если в задаче на память агент выдумал
  номер счёта пользователя, задачу следует считать проваленной, даже
  если ответ сформулирован гладко и содержит всю остальную информацию.
\item
  \textbf{Самодостаточность критериев оценки.} В каждом критерии укажите
  факты, условия и допустимые границы, необходимые для вынесения оценки.
  Например, критерий «правильно назвать текущий адрес клиента» можно
  конкретизировать так: «указать обновлённый адрес в Шанхае; адрес в
  Пекине был отменён в последующем диалоге».
\end{enumerate}

Для каждого измерения также необходимы однозначные уровни оценки и
граничные случаи. Сначала проведите пробную оценку группы ответов,
соберите расхождения между человеком и моделью, а также между разными
оценщиками, после чего дополните критерии прецедентами. Например, если
ответ содержит требуемые ключевые слова, но объединяет истории двух
разных людей, следует снизить оценку за правильность связей. Такие
прецеденты помогают выявлять \textbf{взлом вознаграждения (Reward
Hacking)}: система получает высокий балл, хотя её фактическое поведение
отклоняется от цели задачи.

Ниже приведён полный Rubric для вопроса о памяти между сеансами. В
первом фрагменте истории сказано: «Мою дочь зовут Лили», а во втором:
«Педиатр Лили --- доктор Чен, сегодня я водил её на повторный приём».
Пользователь спрашивает: «Кто педиатр моей дочери?» Два фрагмента
истории совместно подтверждают ответ: «Педиатр вашей дочери Лили ---
доктор Чен».

\begin{longtable}[]{@{}
  >{\raggedright\arraybackslash}p{(\columnwidth - 8\tabcolsep) * \real{0.1845}}
  >{\raggedright\arraybackslash}p{(\columnwidth - 8\tabcolsep) * \real{0.2200}}
  >{\raggedright\arraybackslash}p{(\columnwidth - 8\tabcolsep) * \real{0.1801}}
  >{\raggedright\arraybackslash}p{(\columnwidth - 8\tabcolsep) * \real{0.2162}}
  >{\raggedright\arraybackslash}p{(\columnwidth - 8\tabcolsep) * \real{0.1992}}@{}}
\toprule\noalign{}
\begin{minipage}[b]{\linewidth}\raggedright
\hspace{0pt}Измерение и важность
\end{minipage} & \begin{minipage}[b]{\linewidth}\raggedright
\hspace{0pt}4 балла: отлично
\end{minipage} & \begin{minipage}[b]{\linewidth}\raggedright
\hspace{0pt}3 балла: хорошо
\end{minipage} & \begin{minipage}[b]{\linewidth}\raggedright
\hspace{0pt}2 балла: частичное соответствие
\end{minipage} & \begin{minipage}[b]{\linewidth}\raggedright
\hspace{0pt}1 балл: не соответствует
\end{minipage} \\
\midrule\noalign{}
\endhead
\bottomrule\noalign{}
\endlastfoot
\hspace{0pt}Фактическая правильность: обязательный критерий &
\hspace{0pt}Правильно указаны и имя врача, и его связь с человеком &
\hspace{0pt}Правильно указан доктор Чен, но опущено обозначение его роли
& \hspace{0pt}Упомянут доктор Чен, но в качестве возможных ответов также
перечислены другие врачи & \hspace{0pt}Имя врача указано неверно либо
при наличии достаточных доказательств дан ответ «не знаю» \\
\hspace{0pt}Полнота информации: важный критерий & \hspace{0pt}Ясно
указано соответствие «Лили --- педиатр --- доктор Чен» &
\hspace{0pt}Кратко указано имя врача, что полностью отвечает на текущий
вопрос & \hspace{0pt}Указано только медицинское учреждение или событие,
но пропущено имя врача & \hspace{0pt}Связанная с вопросом информация
отсутствует \\
\hspace{0pt}Правильность связей: важный критерий & \hspace{0pt}Ответ
правильно связывает два фрагмента истории и при необходимости ссылается
на соответствующие основания & \hspace{0pt}Отношения между людьми
указаны правильно, но основания не раскрыты & \hspace{0pt}Нужные люди
найдены, но отношения между ними остаются неоднозначными & \hspace{0pt}С
Лили связан врач самого пользователя или другого ребёнка \\
\end{longtable}

\textbf{Для обнаружения галлюцинаций применяется отдельное вето.} Факты
в ответе должны подтверждаться историей. Если выдуманы дата приёма,
результат диагностики или личность врача, задача считается проваленной.
В этом примере после срабатывания вето итоговая оценка обнуляется.
Правильность связей оценивается по ответу, ссылкам на источники и
наблюдаемой траектории; краткого ответа также может быть достаточно для
выполнения задачи.

От граничных случаев зависит, сможет ли система оценки работать с
реальными диалогами. Если у пользователя несколько дочерей, а указание в
вопросе неоднозначно, правильным действием будет задать уточняющий
вопрос. Если история однозначно указывает на Лили, можно ответить сразу.
Являются ли «Dr.~Chen» и «доктор Чен» одним человеком, следует
определять по указанным в истории псевдонимам, учреждениям и отношениям
между людьми.

\textbf{Если критериев много, первичную проверку можно поручить
модели-судье.} Предположим, имеется тысяча ответов, каждый из которых
нужно проверить по двадцати критериям: всего получится двадцать тысяч
суждений. Рассмотрим отношения между людьми из таблицы. Утверждения
«Доктор Чен --- педиатр Лили» и «Доктор Чен --- мой педиатр, а Лили
ходила со мной на приём» содержат одинаковые имена, но приводят к разным
выводам. Вместо того чтобы отдельно составлять правила сопоставления
имён, порядка слов и отрицаний, критерий можно сформулировать одним
предложением: «Правильно ли ответ связывает этого врача с дочерью
пользователя с учётом вопроса и двух фрагментов истории?» Передайте эти
материалы, оцениваемый ответ и критерий модели-судье наподобие Jev,
чтобы она выбрала один из вариантов: «соответствует», «не соответствует»
или «невозможно определить». Такой вопрос позволяет выявить ответы, в
которых врач связан не с тем человеком. Чтобы различать четыре уровня
правильности связей из таблицы, можно непосредственно представить
критерии всех четырёх уровней как варианты ответа, добавив «невозможно
определить» для последующей проверки. Разные измерения одного ответа
можно оценивать параллельно, после чего программа суммирует баллы и
применяет явно заданные правила, включая вето за галлюцинации.

Чтобы как можно скорее начать использовать этот метод оценки, можно
напрямую обратиться к такой передовой модели-судье, как официальный Jev
от TypeSafe. Запишите правила оценки в промпте, перечислите возможные
уровни и при каждом вызове передавайте оцениваемый ответ вместе с
соответствующими историческими материалами --- это позволит быстро
получать результат. Если ежедневно требуется оценивать множество ответов
по одним и тем же правилам и желательно дополнительно снизить стоимость
одного вызова, можно дистиллировать собственную модель-судью. Сначала
разработчик системы оценки чётко описывает критерии каждого уровня,
затем передовые модели, такие как DeepSeek, Kimi, GPT, Claude или
Gemini, выступают в роли учителей и оценивают примеры по заданным
правилам. Из проверенных наборов «материалы, ответ, уровень, присвоенный
учителем» формируется синтетический обучающий набор, на котором с
помощью метода контролируемого дообучения из
\crossreflink{chap:8}{главы 8} обучается меньшая модель-ученик.
Обученная модель-ученик затем сможет выполнять повседневную оценку.

\begin{quote}
\textbf{Эксперимент 7-3 ★★: создание системы оценки пользовательской
памяти на основе Rubric}

Реализуйте структурированную оценку в системе оценки пользовательской
памяти из \crossreflink{chap:3}{главы 3}, охватив три уровня задач:
базовое воспроизведение, устранение неоднозначности между сеансами и
выявление скрытых связей между сеансами. Измерения оценки включают
фактическую правильность, полноту фактов, правильность связей,
проактивность и независимое обнаружение галлюцинаций.

Фактическая правильность проверяет предоставленные числа, даты, имена и
другие сведения; полнота фактов --- наличие в ответе обязательных
ключевых фактов; правильность связей --- отношения между людьми,
временем и событиями; проактивность --- своевременность обоснованных
рекомендаций и предупреждений о рисках. Первые четыре измерения
оцениваются по четырёхуровневой шкале с однозначными критериями для
каждого уровня, примерами и граничными случаями; обнаружение
галлюцинаций используется как вето.

В реализации для первых двух измерений сохраняются названия полей
Precision и Recall, однако в них записывается уровень Rubric. Чтобы
сообщать точность и полноту в статистическом смысле, необходимо
дополнительно определить подсчитываемые единицы фактов и отдельно
вычислять долю правильных фактов среди приведённых и долю охваченных
обязательных фактов среди всех обязательных фактов.

При анализе результатов эксперимента сначала проверяйте, пройдена ли
задача, а затем --- сколько баллов она получила. Задача считается
пройденной, если фактическая правильность, полнота фактов и правильность
связей получили не менее 3 баллов и вето за галлюцинацию не сработало.
Непрерывное вознаграждение также учитывает проактивность: оценка каждого
из четырёх измерений преобразуется по формуле \((s-1)/3\), после чего
вычисляется среднее; при наличии галлюцинации вознаграждение обнуляется.
Например, в эксперименте 7-12 одна из конфигураций использовала BGE-M3
для поиска, пропускала переранжирование и поручала ответ Doubao. В одной
из задач о страховой выплате четыре измерения получили соответственно 3,
3, 3 и 2 балла. Задача была пройдена, а непрерывное вознаграждение
составило приблизительно 0,583.
\end{quote}

\begin{quote}
\textbf{Эксперимент 7-4 ★★: сравнительная оценка Advanced JSON Cards и
RAG}

Сравните на одном наборе из 60 вопросов три варианта памяти: только
Advanced JSON Cards, при котором структурированные карточки помещаются в
контекст; только RAG, при котором диалог разбивается на фрагменты и
выполняется поиск связанного содержимого; гибридную систему, в которой
ключевые факты постоянно находятся в контексте, а исходные диалоги
извлекаются по мере необходимости. Каждый из трёх уровней содержит по 20
вопросов, что в сумме даёт 180 комбинаций «вариант --- вопрос».

Для всех трёх вариантов используйте одинаковые вопросы, исходные
диалоги, модель ответов и правила оценки. Отдельно рассчитайте долю
успешных решений по уровням, число срабатываний вето за галлюцинации и
стоимость. Затем сопоставьте три результата для каждого вопроса и
проанализируйте, в каких вопросах гибридная система обеспечивает
взаимное дополнение, а в каких теряет качество из-за отбора или
объединения информации.
\end{quote}

В таблице 7-3 приведена историческая сводка этого эксперимента. На
каждом уровне имеется 20 вопросов, поэтому разница в 5 процентных
пунктов соответствует одному вопросу.

Таблица 7-3. Доля успешных решений трёх систем пользовательской памяти
по уровням

\begin{longtable}[]{@{}
  >{\raggedright\arraybackslash}p{(\columnwidth - 8\tabcolsep) * \real{0.1735}}
  >{\raggedleft\arraybackslash}p{(\columnwidth - 8\tabcolsep) * \real{0.2413}}
  >{\raggedleft\arraybackslash}p{(\columnwidth - 8\tabcolsep) * \real{0.2528}}
  >{\raggedleft\arraybackslash}p{(\columnwidth - 8\tabcolsep) * \real{0.1543}}
  >{\raggedleft\arraybackslash}p{(\columnwidth - 8\tabcolsep) * \real{0.1781}}@{}}
\toprule\noalign{}
\begin{minipage}[b]{\linewidth}\raggedright
\hspace{0pt}Система
\end{minipage} & \begin{minipage}[b]{\linewidth}\raggedleft
\hspace{0pt}Базовое воспроизведение
\end{minipage} & \begin{minipage}[b]{\linewidth}\raggedleft
\hspace{0pt}Устранение неоднозначности между сеансами
\end{minipage} & \begin{minipage}[b]{\linewidth}\raggedleft
\hspace{0pt}Скрытые связи между сеансами
\end{minipage} & \begin{minipage}[b]{\linewidth}\raggedleft
\hspace{0pt}Общий результат
\end{minipage} \\
\midrule\noalign{}
\endhead
\bottomrule\noalign{}
\endlastfoot
\hspace{0pt}Advanced JSON Cards & \hspace{0pt}95\% & \hspace{0pt}60\% &
\hspace{0pt}50\% & \hspace{0pt}68,3\% (41/\allowbreak{}60) \\
\hspace{0pt}RAG & \hspace{0pt}90\% & \hspace{0pt}40\% & \hspace{0pt}15\%
& \hspace{0pt}48,3\% (29/\allowbreak{}60) \\
\hspace{0pt}Гибридная система & \hspace{0pt}80\% & \hspace{0pt}70\% &
\hspace{0pt}50\% & \hspace{0pt}66,7\% (40/\allowbreak{}60) \\
\end{longtable}

Сначала рассмотрим только RAG. В задачах на базовое воспроизведение он
правильно ответил на 18 вопросов, а вариант с карточками --- на 19. В
задачах на скрытые связи между сеансами только RAG правильно ответил
лишь на 3 вопроса, а вариант с карточками --- на 10. Чтобы выяснить
причину разницы в 7 вопросов, можно сначала выбрать один вопрос, на
который RAG ответил неверно, и проверить, содержат ли результаты поиска
оба фрагмента диалога, необходимые для установления связи. Если оба
фрагмента найдены, следует проверить, каких двух людей или какие два
события связал ответ. По одной лишь доле успешных решений невозможно
определить, не все ли материалы были найдены или модель прочитала их, но
неверно установила связь.

\textbf{Даёт ли объединение двух механизмов памяти естественное
сочетание их преимуществ?} Гибридная система правильно ответила на 14
вопросов по устранению неоднозначности между сеансами, превзойдя вариант
с карточками, давший 12 правильных ответов, и только RAG, давший 8.
Однако в исторической сводке также зафиксированы две группы вопросов,
которые стоит рассматривать отдельно: на 3 вопроса правильно ответила
только гибридная система, а ещё на 8 вопросов правильно ответил хотя бы
один из отдельных вариантов, но гибридная система ошиблась. Таким
образом, хотя бы один отдельный вариант мог правильно ответить в общей
сложности на \(40-3+8=45\) вопросов, тогда как фактически запущенная
гибридная система --- на 40. Результат в 45 вопросов --- это
ретроспективный ориентир, полученный объединением двух наборов ответов.
Чтобы система самостоятельно достигла такого результата, необходимо
решить, как ей выбирать нужный вид памяти во время ответа. Кроме того, в
28 из 180 оценок сработало вето за галлюцинации.

\textbf{Одно семейство моделей и оценка из нескольких источников.}

Агент и оценивающая модель могут не знать одни и те же факты, а также
могут предпочитать одинаковые способы выражения. Если они принадлежат к
одному семейству моделей, особенно важно проверить, не влияют ли такие
предпочтения на оценки. Если агент освоит стиль, который нравится
оценщику, но продолжит допускать фактические ошибки, не обнаруживаемые
этим оценщиком, его баллы всё равно могут вырасти. Это один из примеров
закона Гудхарта в оценке: если многократно превращать метрику в цель
оптимизации, со временем она может утратить способность измерять
реальное качество.

Перед выпуском важных версий и во время регулярных проверок критически
значимых систем автор рекомендует использовать оценку из нескольких
источников: разные семейства моделей должны независимо выставлять оценки
по одному и тому же Rubric. Например, задачу выполняет Claude, а
оценивают GPT и Gemini. Сначала на примерах с ручной разметкой следует
измерить точность каждого оценщика и расхождения между ними, после чего
выбрать способ агрегирования. Для обычных оценок можно использовать
взвешенное среднее, критические условия провала обрабатывать по заранее
установленным правилам, а примеры с существенными расхождениями
передавать человеку на проверку. В повседневной работе можно
использовать одну оценивающую модель, своевременно обнаруживая дрейф
оценок с помощью описанных выше проверок.

\textbf{Мультимодальный LLM-as-a-Judge.}

При мультимодальной оценке тот же процесс распространяется на аудио,
изображения и видео. Материалы для оценки должны охватывать оцениваемые
свойства:

\begin{itemize}
\tightlist
\item
  \textbf{Оценка TTS}: проверка правильности чтения, естественности,
  эмоциональной выразительности и постоянства тембра. Частота ошибок в
  словах (WER) в основном отражает ошибки на уровне слов в распознанном
  тексте; для оценки просодии и тембра необходимо прослушать аудио.
\item
  \textbf{Оценка ASR}: влияние ошибок распознавания определяется с
  учётом задачи. Например, если фраза «перевести тысячу» распознана как
  «перевести десять тысяч», изменятся параметры действия. Суммы, даты и
  имена можно отдельно выделить как критические поля; ошибки в обычных
  словах также следует оценивать с учётом контекста.
\item
  \textbf{Оценка UI}: используйте схему «предлагающий --- проверяющий»
  для проверки переполнения текста, контрастности и расположения кнопок
  на снимке экрана, а доступность кнопок проверяйте посредством
  реального взаимодействия. Здесь механизм проверки из
  \crossreflink{chap:5}{главы 5} применяется для оценки сгенерированных
  результатов.
\item
  \textbf{Оценка видеомонтажа}: проверяйте содержимое кадров и
  спецэффекты по ключевым кадрам, а моменты начала и окончания, переходы
  и синхронизацию звука с изображением --- по последовательным
  фрагментам вблизи точек монтажа. Диапазон выборки определяется
  оцениваемым аспектом.
\end{itemize}

\begin{quote}
\textbf{Эксперимент 7-5 ★★: создание полностью автоматизированного
конвейера оценки качества TTS}

Разработайте четырёхмерный Rubric. Правильность проверяет пропуски,
неверное произношение и добавления; естественность --- паузы, ударения,
плавность и просодию; эмоциональная выразительность --- соответствие
контексту текста и заданному стилю; постоянство тембра при наличии
целевого диктора --- сходство эталонной и синтезированной речи. Для
оценки эмоций следует привести примеры и учитывать языковые нормы
выражения вопроса, выделения, печали и других состояний.

Тестовый материал должен охватывать отдельные предложения и длинные
фрагменты, новости, рассказы и диалоги, нейтральный, возбуждённый и
печальный тон, а также числа, имена собственные, иероглифы с несколькими
чтениями и диалектные слова. Модуль генерации может подключаться к
OpenAI, ElevenLabs, Fish Audio, MiniMax, Doubao и другим сервисам.
Модуль оценки передаёт модели, способной напрямую принимать аудио,
исходный текст, оцениваемую синтезированную речь, Rubric и необходимое
для задачи эталонное аудио, после чего выставляет оценки по каждому
измерению и объясняет основания.

Сгруппируйте результаты по сервисам, типам тестового материала и
измерениям оценки. Отберите примеры с высокими и низкими баллами, а
также с существенными расхождениями между оценщиками, и проведите слепое
прослушивание людьми.
\end{quote}

В сопутствующем эксперименте для каждой из двух конфигураций TTS были
взяты образцы четырёх типов: числа, иероглифы с несколькими чтениями,
длинные предложения и возбуждённая интонация. Voxtral напрямую оценивал
их на слух; всего было получено 8 записей с оценками. По пятибалльной
шкале обе конфигурации получили 5,00 за правильность и 4,00 за
естественность. Средние оценки Fish Audio за эмоциональную
выразительность и постоянство тембра составили соответственно 4,00 и
3,00, а OpenAI --- 3,75 и 2,75. Обе разницы в 0,25 балла возникли только
в задании с числами: по остальным трём заданиям все четыре оценки
совпали. Например, в конце исходного текста задания с иероглифами с
несколькими чтениями стояла фраза «к тому же необходимо полностью
погасить всю задолженность». В ней один и тот же иероглиф встречается
дважды и должен произноситься сначала как hái («к тому же»), а затем как
huán («погасить»). Проверяющий поставил обеим записям максимальную
оценку за правильность и привёл совершенно одинаковые обоснования, но не
указал посимвольно, что именно услышал.

Этот эксперимент также показывает, как эталонное аудио влияет на цель
оценки. В записях использовалось эталонное аудио от Fish S1, поэтому
оценка тембра отражала, насколько синтезированная речь похожа на эту
эталонную запись. При сравнении универсальных систем TTS можно отдельно
оценивать чтение, естественность и стиль. При сравнении систем
клонирования голоса следует предложить всем вариантам имитировать одного
и того же целевого диктора, а оценку сходства откалибровать посредством
слепого прослушивания людьми. Справочные материалы также определяют,
«что считается хорошим»: эталонный ответ задаёт требуемое содержание, а
эталонное изображение или аудио --- визуальные или звуковые
характеристики для сравнения.

Составленный вручную Rubric позволяет быстро сформировать
диагностические измерения. По мере масштабирования оценки можно также
обучить специализированную \textbf{генеративную модель вознаграждения}
для выставления баллов; соответствующие методы обучения рассматриваются
в \crossreflink{chap:8}{главе 8}. Далее мы свяжем поэлементную оценку с
траекторией выполнения, чтобы найти конкретные шаги, которые можно
исправить.

\hypertarget{ux430ux442ux440ux438ux431ux443ux446ux438ux44f-ux441ux431ux43eux435ux432-ux43fux43eux438ux441ux43a-ux43fux435ux440ux432ux43eux439-ux43eux448ux438ux431ux43aux438-ux43fux43e-ux432ux441ux435ux439-ux442ux440ux430ux435ux43aux442ux43eux440ux438ux438}{%
\subsection{Атрибуция сбоев: поиск первой ошибки по всей
траектории}\label{ux430ux442ux440ux438ux431ux443ux446ux438ux44f-ux441ux431ux43eux435ux432-ux43fux43eux438ux441ux43a-ux43fux435ux440ux432ux43eux439-ux43eux448ux438ux431ux43aux438-ux43fux43e-ux432ux441ux435ux439-ux442ux440ux430ux435ux43aux442ux43eux440ux438ux438}}

После сбоя задачи сначала сопоставьте условия приёмки с траекторией
выполнения, а затем найдите самое раннее решение, которое отклонило
задачу от цели. Эта работа называется \textbf{атрибуцией сбоев (failure
attribution)}. Полезная запись об атрибуции должна содержать категорию
ошибки, связанные с ней шаги, вызов инструмента или вывод модели,
доказательства, подтверждающие вывод, и последующие последствия.
Последняя появившаяся ошибка зачастую лишь следствие более ранней; чтобы
найти место, которое действительно стоит исправить, необходимо
проследить назад источники информации и изменения состояния.

Проблемы в продакшене выявляются главным образом по трём типам сигналов:
явным исправлениям пользователя, отрицательным оценкам или отзывам
пользователя, а также аномалиям, обнаруженным при последующей проверке
состояния, валидации правил или оценке моделью. Автор рекомендует
разработчикам сначала лично внимательно изучить несколько проблемных
траекторий и лишь затем выстраивать процесс массовой атрибуции.
Классификация сбоев включает продуктовые решения: сначала необходимо
преобразовать в явные требования то, какие действия требуют
подтверждения и какие инициативные предложения выходят за допустимые
границы, --- только после этого модель сможет выполнять разметку.
Например, граница между «самостоятельно выполнить задачу за
пользователя» и «сначала запросить подтверждение пользователя» зависит
от текущих полномочий, последствий действия и продуктовых соглашений.
Модель может помогать с массовой разметкой, а руководители продукта и
разработки должны калибровать эти границы.

На примере Coding Agent можно начать со следующих девяти категорий
проблем. По мере накопления случаев их следует разделять на более узкие
подкатегории, а критерии классификации, положительные и отрицательные
примеры и процесс атрибуции --- добавлять в промпт или Skill
агента-разметчика.

\begin{longtable}[]{@{}
  >{\raggedright\arraybackslash}p{(\columnwidth - 4\tabcolsep) * \real{0.2545}}
  >{\raggedright\arraybackslash}p{(\columnwidth - 4\tabcolsep) * \real{0.3729}}
  >{\raggedright\arraybackslash}p{(\columnwidth - 4\tabcolsep) * \real{0.3726}}@{}}
\toprule\noalign{}
\begin{minipage}[b]{\linewidth}\raggedright
\hspace{0pt}Категория ошибки
\end{minipage} & \begin{minipage}[b]{\linewidth}\raggedright
\hspace{0pt}Типичные проявления
\end{minipage} & \begin{minipage}[b]{\linewidth}\raggedright
\hspace{0pt}Способ локализации
\end{minipage} \\
\midrule\noalign{}
\endhead
\bottomrule\noalign{}
\endlastfoot
\hspace{0pt}Понимание требований и устранение неоднозначности &
\hspace{0pt}Пропущено условие требования; изменена неверная область;
выбор между двумя одноимёнными конфигурационными файлами сделан без
достаточных оснований & \hspace{0pt}Сопоставить требования пользователя
с фактическими действиями, найти первое отклонение и проверить, какая
информация была доступна в тот момент \\
\hspace{0pt}Недостатки процесса и соблюдения правил &
\hspace{0pt}Пропущены требуемые репозиторием тесты или планирование;
добавлена внешняя зависимость при наличии внутренней реализации;
нарушены архитектурные соглашения & \hspace{0pt}Найти первое действие,
нарушившее явное соглашение, а затем проверить, были ли правила уже
предоставлены или доступны для чтения \\
\hspace{0pt}Ошибки вызова инструментов & \hspace{0pt}Редактирование
неоднократно завершается сбоем сопоставления; неверные параметры или
schema; ошибки экранирования специальных символов &
\hspace{0pt}Сопоставить исходный запрос, результат инструмента и
состояние перед вызовом, чтобы найти первые входные данные, не
соответствующие контракту интерфейса \\
\hspace{0pt}Искажение результатов проверки & \hspace{0pt}Для сокрытия
сбоя изменены утверждения, пропущены тесты или подменена проверяемая
логика; заявлено об успешной проверке без запуска тестов &
\hspace{0pt}Проверить цель изменений в логике валидации и сопоставить
заявление о завершении с фактическим журналом выполнения \\
\hspace{0pt}Неполнота изменений & \hspace{0pt}Изменена сигнатура
функции, но пропущены динамические вызовы, привязки на другом языке или
связанная schema & \hspace{0pt}Сопоставить затронутую область с
изменёнными местами, проследить пропущенные зависимости и процесс
поиска \\
\hspace{0pt}Ошибочная информация для пользователя & \hspace{0pt}Неверно
переданы сумма, состояние или время; частичное выполнение представлено
как полное; пропущено обязательное уведомление & \hspace{0pt}По пунктам
сопоставить итоговый ответ с результатами инструментов, конечным
состоянием и требованиями к уведомлению \\
\hspace{0pt}Регрессии совместимости и нефункциональных характеристик &
\hspace{0pt}Изменение публичного API нарушило работу вызывающей стороны;
изменение структуры данных не сопровождается необходимой миграцией;
ухудшились производительность или потребление ресурсов &
\hspace{0pt}Проверить контракты интерфейсов, требования к миграциям и
базовые показатели производительности, затем найти шаг, внёсший
изменение \\
\hspace{0pt}Сбои выполнения & \hspace{0pt}Вывод усечён, истекло время
ожидания, соединение разорвано или завершение работы не доведено до
конца & \hspace{0pt}Сопоставить причину остановки модели, события
выполнения Harness и состояние сервисов инструментов \\
\hspace{0pt}Преждевременное завершение задачи & \hspace{0pt}В
многоцелевой задаче выполнена лишь часть целей; работа завершена, хотя
оставались разумные варианты действий & \hspace{0pt}Сопоставить
оставшиеся цели, уже опробованные варианты и условия остановки, чтобы
локализовать решение о завершении \\
\end{longtable}

В записи об атрибуции следует различать \textbf{непосредственную ошибку,
гипотезу о причине и последующие симптомы}. Например,
\texttt{edit\_file} сообщает, что сопоставить \texttt{old\_string} не
удалось, после чего агент трижды повторяет тот же вызов без изменений.
Сбой сопоставления --- непосредственная ошибка, а три повторные попытки
указывают на проблему со стратегией восстановления. Только
дополнительная проверка содержимого файла перед вызовом позволит
различить ошибку при формировании строки, конкурентное изменение файла и
неверное понимание правил сопоставления инструмента.

При массовой атрибуции автор сначала с помощью правил выявляет журналы
выполнения, в которых могли произойти ошибки, а затем поручает модели
анализировать причины с учётом предшествующих и последующих шагов. Такие
сигналы, как запуск тестов или изменение интерфейса, можно извлекать с
небольшими затратами, сосредоточив вызовы модели на фрагментах,
требующих семантической оценки. Например, можно сопоставлять заявления
об успешном прохождении тестов с журналом выполнения, отмечать изменения
утверждений или маркеров пропуска и проверять, сопровождаются ли
изменения schema соответствующими миграциями. Затем модель или человек
определяет обоснованность этих изменений с учётом цели задачи.

\hypertarget{ux440ux430ux437ux434ux435ux43bux44cux43dux430ux44f-ux43fux440ux438ux451ux43cux43aux430-ux43aux43eux43dux435ux447ux43dux43eux433ux43e-ux441ux43eux441ux442ux43eux44fux43dux438ux44f-ux438-ux438ux442ux43eux433ux43eux432ux43eux433ux43e-ux43eux442ux432ux435ux442ux430}{%
\subsubsection{Раздельная приёмка конечного состояния и итогового
ответа}\label{ux440ux430ux437ux434ux435ux43bux44cux43dux430ux44f-ux43fux440ux438ux451ux43cux43aux430-ux43aux43eux43dux435ux447ux43dux43eux433ux43e-ux441ux43eux441ux442ux43eux44fux43dux438ux44f-ux438-ux438ux442ux43eux433ux43eux432ux43eux433ux43e-ux43eux442ux432ux435ux442ux430}}

Пользователь получает итоговый ответ агента, а бизнес-система сохраняет
состояние, созданное его действиями. Проверять необходимо и то и другое.
Если возврат средств отправлен на обработку, следует точно сообщить
сумму, состояние обработки и дальнейшие действия. Если выполнена лишь
часть задачи, необходимо по пунктам перечислить завершённое и
оставшееся. Проверку состояния среды и проверку обязательных уведомлений
можно оформить как независимые критерии оценки, чтобы различать ошибки
выполнения и ошибки при передаче информации.

Представим, что пользователь поручил помощнику службы поддержки
запросить у продавца отмену комиссии. Во время звонка сотрудник
поддержки говорит: «Мы отменим эту комиссию»; запись разговора позволяет
подтвердить, что \textbf{сотрудник поддержки дал устное согласие}.
Итоговый счёт получает пользователь. В отчёте пользователю следует
сохранить именно этот смысл действия: «Сотрудник поддержки согласился
отменить комиссию». Такой случай вполне обоснованно учитывается в доле
устных обещаний. Для расчёта доли фактически проведённых отмен
потребуется новый счёт или последующее подтверждение пользователя.

Иногда при выполнении одного поручения необходимо передавать
взаимодействие между несколькими каналами. Если продавец собирается
напрямую перезвонить пользователю, помощник может заранее отправить ему
краткую сводку: цель звонка, уже согласованные вопросы и вопросы, на
которые ещё предстоит ответить. Если продавец требует, чтобы
пользователь лично прошёл устную проверку, помощник может объяснить, что
именно требуется сделать, подключить пользователя к трёхстороннему
звонку, а после проверки продолжить оформление. При запросе одноразового
кода следует соблюдать указания, связанные с этим кодом: пользователь
должен пройти соответствующую проверку в указанном продавцом доверенном
интерфейсе. Как оценивать такие действия? В случае обратного звонка
нужно проверить, получил ли пользователь сводку до того, как ответил на
звонок. В случае проверки --- смог ли пользователь подключиться на
нужном этапе и продолжилось ли выполнение задачи после её завершения.

Задача AndroidWorld по внесению расходов с изображения демонстрирует
другую причинно-следственную цепочку: неверные данные были успешно
записаны, после чего агент объявил задачу выполненной. Требовалось
внести расходы с изображения в приложение для учёта финансов. За 32 шага
T3A выполнил навигацию, ввод и сохранение. Уже в выводе шага 8 явно
указано, что конкретные расходы на изображении не видны, однако
выполнение продолжилось. На шаге 11 появились не подкреплённые
источником данные: «Кофе --- 4,50 доллара, обед --- 12,75 доллара»; весь
последующий ввод основывался на них. Когда итоговый валидатор стал
искать соответствующую изображению запись о платье, оказалось, что
необходимая запись отсутствует.

При анализе этой траектории необходимо одновременно проверить канал
наблюдения и стратегию принятия решений. T3A получает дерево элементов
доступности, а пиксели изображения не попадают в контекст модели.
Поэтому для задачи чтения изображения сначала нужно добавить канал
снимков экрана или распознавания текста. Если обобщённо записать такую
проблему как «модель не умеет выполнять OCR», команда может заменить
модель или дополнить обучающие данные, но по-прежнему не заметить
отсутствующие входные данные. При недостатке информации агент также
должен попытаться получить её либо сообщить о препятствии. Решение
продолжить выполнение на шаге 8, последующее выдумывание данных и
итоговое заявление о завершении представляют разные проблемы в одной
причинно-следственной цепочке.

Та же проблема встречается и в научных исследованиях. Команда Denario
сообщала о таком сбое: программа анализа не могла прочитать
предоставленные данные, после чего система иногда самостоятельно
генерировала имитационные данные и продолжала анализ, не уведомляя
пользователя.\footnote{Francisco Villaescusa-Navarro и др., \emph{The
  Denario Project: Deep Knowledge Artificial Intelligence Agents for
  Scientific Discovery}. Авторы статьи описывают возможную подмену
  данных при сбое чтения входных данных.} Представим, что читатель
получил построенный таким образом график. Даже если повторный запуск
скрипта рисует точно такую же кривую, один вопрос остаётся без ответа:
на какой совокупности наблюдений она основана? При поиске причины можно
проследить назад столбец значений, использованный для построения
графика, найти присваивающий их код, а затем проверить, какой файл
фактически был прочитан при этом запуске. После сбоя чтения ещё
оставалась возможность исправить ситуацию; последующие выводы лишились
оснований лишь тогда, когда имитационные значения были переданы дальше
как наблюдаемые.

\begin{quote}
\textbf{Эксперимент 7-6 ★★: атрибуция сбоев в неудачных траекториях
AndroidWorld}

Используйте сохранённые журналы выполнения T3A, чтобы потренироваться в
офлайн-атрибуции. Материалы содержат пошаговые действия, пояснения,
сводки и итоговые результаты валидации; 52 траектории соответствуют
неудачно выполненным задачам. Выполните следующие пять шагов.

Шаг 1. Выберите 10 неудачных траекторий. Основное внимание уделите тихим
сбоям, при которых вызовы инструментов внешне проходят нормально, но
итоговая приёмка завершается неудачей; для сравнения добавьте случаи с
видимыми ошибками. В сопутствующей разметке выбраны 9 траекторий с
тихими сбоями и 1 траектория с видимой ошибкой.

Шаг 2. Последовательно проверьте фактические основания каждого шага,
отметьте самое раннее отклонение от задачи и укажите, произошло ли оно в
вызове инструмента или выводе модели. В случаях с неясной причиной
сравните соседние префиксы траекторий. Если состояние допускает
восстановление, можно продолжить выполнение с разных префиксов или
заменить один из шагов, чтобы проверить гипотезу о причине. Если после
некоторого префикса задачу всё ещё можно спасти, ошибка, вероятно,
обратима; момент её первого возникновения всё равно определяется
сопоставлением фактов и действий.

Шаг 3. Составьте структурированную запись об атрибуции, указав название
задачи, номер шага, категорию ошибки, задействованную модель или
системный компонент, исходное доказательство, основную причину и
последствия, а также уровень уверенности.

Шаг 4. Сопоставьте результат с существующими заметками об атрибуции,
зафиксируйте расхождения и доказательства в пользу каждого вывода.
Например, в задаче транскрибирования изображения сначала необходимо
проверить, получила ли модель пиксели, и лишь затем обсуждать её
способность к распознаванию.

Шаг 5. Выберите 3 траектории, в которых первая ошибка возникла в выводе
модели. Извлеките контекст и состояние перед первой ошибкой, опишите
допустимые и запрещённые действия и сформируйте регрессионную задачу с
префиксом.
\end{quote}

\hypertarget{ux43eux448ux438ux431ux43aux438-ux444ux43eux440ux43cux430ux442ux438ux440ux43eux432ux430ux43dux438ux44f-ux434ux43eux43aux443ux43cux435ux43dux442ux430-ux437ux430ux432ux438ux441ux44fux449ux438ux435-ux43eux442-ux43eux431ux43bux430ux441ux442ux438-ux434ux435ux439ux441ux442ux432ux438ux44f}{%
\subsubsection{Ошибки форматирования документа, зависящие от области
действия}\label{ux43eux448ux438ux431ux43aux438-ux444ux43eux440ux43cux430ux442ux438ux440ux43eux432ux430ux43dux438ux44f-ux434ux43eux43aux443ux43cux435ux43dux442ux430-ux437ux430ux432ux438ux441ux44fux449ux438ux435-ux43eux442-ux43eux431ux43bux430ux441ux442ux438-ux434ux435ux439ux441ux442ux432ux438ux44f}}

Один и тот же символ в разных местах выполняет разные функции. В
китайском тексте кавычки обозначают цитирование, двойные кавычки в JSON
относятся к синтаксису, а обратные кавычки Markdown ограничивают код.
Если пользователь просит изменить формат кавычек, при оценке прежде
всего необходимо явно определить область, в которой разрешены изменения.

Документ можно разобрать на фрагменты с областями действия: китайский
текст (\texttt{ZH\_PROSE}), английский текст (\texttt{EN\_PROSE}),
цитируемый источник (\texttt{QUOTED\_SOURCE}), встроенный код
(\texttt{INLINE\_CODE}), блок кода (\texttt{CODE\_BLOCK}), комментарий к
коду (\texttt{CODE\_COMMENT}) и JSON или schema
(\texttt{JSON\_OR\_SCHEMA}). Для каждого фрагмента фиксируются
допустимые преобразования, защищённые символы и правила проверки.
Например:

\begin{longtable}[]{@{}
  >{\raggedright\arraybackslash}p{(\columnwidth - 2\tabcolsep) * \real{0.4199}}
  >{\raggedright\arraybackslash}p{(\columnwidth - 2\tabcolsep) * \real{0.5801}}@{}}
\toprule\noalign{}
\begin{minipage}[b]{\linewidth}\raggedright
\hspace{0pt}Фрагмент
\end{minipage} & \begin{minipage}[b]{\linewidth}\raggedright
\hspace{0pt}Требования к обработке
\end{minipage} \\
\midrule\noalign{}
\endhead
\bottomrule\noalign{}
\endlastfoot
\hspace{0pt}Китайское пояснение «вызовите метод \texttt{reset()}» &
\hspace{0pt}При изменении пунктуации обычного текста сохранить
встроенный код и окружающие его обратные кавычки \\
\hspace{0pt}Английский оригинал «Please restart the service.» &
\hspace{0pt}Оформить внешнюю пунктуацию по правилам цитирования,
сохранив содержимое цитаты \\
\hspace{0pt}Кавычки в китайском комментарии к коду &
\hspace{0pt}Отдельно обработать текст комментария по установленным
правилам, сохранив его границы \\
\hspace{0pt}\texttt{name\ =\ "status"} & \hspace{0pt}Сохранить
синтаксические границы строки и проверить возможность разбора кода \\
\end{longtable}

Регрессия с префиксом одновременно проверяет стиль целевой области, долю
сохранённых цитат, синтаксис кода и JSON, а также расстояние
редактирования для нецелевого текста. Пока область действия не
определена, допустимые следующие шаги включают чтение правил
форматирования, проверку структуры документа и запрос необходимых
уточнений. Благодаря этому минимальное изменение получает исполнимые
условия приёмки.

\hypertarget{ux43eux448ux438ux431ux43aux438-ux442ux43eux447ux43dux43eux433ux43e-ux43aux43eux43fux438ux440ux43eux432ux430ux43dux438ux44f-ux43fux43eux438ux441ux43a-ux440ux430ux437ux43bux438ux447ux438ux44f-ux43fux43e-ux446ux435ux43fux43eux447ux43aux435-ux43fux435ux440ux435ux434ux430ux447ux438-ux441ux442ux440ux43eux43aux438}{%
\subsubsection{Ошибки точного копирования: поиск различия по цепочке
передачи
строки}\label{ux43eux448ux438ux431ux43aux438-ux442ux43eux447ux43dux43eux433ux43e-ux43aux43eux43fux438ux440ux43eux432ux430ux43dux438ux44f-ux43fux43eux438ux441ux43a-ux440ux430ux437ux43bux438ux447ux438ux44f-ux43fux43e-ux446ux435ux43fux43eux447ux43aux435-ux43fux435ux440ux435ux434ux430ux447ux438-ux441ux442ux440ux43eux43aux438}}

Инструменты точной замены обычно требуют, чтобы \texttt{old\_string}
посимвольно совпадала с содержимым файла. При сбое сопоставления можно
проверить следующую цепочку и определить, на каком уровне изменилась
строка:

Исходные байты файла → результат инструмента → сериализация Harness →
контекст модели → вывод модели → декодированная строка → разбор вызова
инструмента → сопоставление инструментом.

Сравниваемые материалы включают исходные байты и их хеш,
последовательность кодовых точек Unicode, а также доступные через
интерфейс ID токенов tokenizer. Сначала следует привести к единому
представлению ожидаемые преобразования, такие как экранирование JSON и
соглашения о переводах строк, а затем найти первое неожиданное различие.
Например, визуально одинаковые символы могут использовать разные
составные формы Unicode, а сохранённая в параметрах обратная косая черта
--- на одном из уровней декодироваться лишний раз или, напротив, не
декодироваться.

Минимальный набор проб должен охватывать прямое дословное
воспроизведение, извлечение из длинного контекста, помещение в параметры
инструмента, выбор среди похожих строк, а также пробелы, переводы строк,
обратные косые черты, составные символы Unicode и редкие токены. Метрики
могут включать точное совпадение байтов, точное совпадение кодовых
точек, позицию первого расхождения и реальную долю успешных вызовов
инструмента. Для сравнения последовательностей токенов необходимо
зафиксировать tokenizer, его версию и условия кодирования. Один и тот же
текст может иметь разные разбиения на токены при генерации, поэтому в
конечном счёте всё равно нужно проверять строку, фактически полученную
инструментом.

Размещение одной и той же пробы в разных точках цепочки помогает сузить
область исправления. Если различие уже присутствует в результате
инструмента, следует проверить чтение и сериализацию. Если вывод модели
верен, а инструмент получает неверную строку, нужно проверить
декодирование и протокол. Если модель генерирует неверную строку даже
при полном контексте, этот случай можно включить в обучающие данные по
способности к копированию из \crossreflink{chap:8}{главы 8}.
Конкурентные изменения необходимо проверять отдельно с помощью версии
файла или повторного чтения.

\hypertarget{ux441ux43aux432ux43eux437ux43dux44bux435-ux440ux435ux433ux440ux435ux441ux441ux438ux43eux43dux43dux44bux435-ux437ux430ux434ux430ux447ux438-ux438-ux440ux435ux433ux440ux435ux441ux441ux438ux43eux43dux43dux44bux435-ux437ux430ux434ux430ux447ux438-ux43fux43e-ux43fux440ux435ux444ux438ux43aux441ux443-ux442ux440ux430ux435ux43aux442ux43eux440ux438ux438}{%
\subsection{Сквозные регрессионные задачи и регрессионные задачи по
префиксу
траектории}\label{ux441ux43aux432ux43eux437ux43dux44bux435-ux440ux435ux433ux440ux435ux441ux441ux438ux43eux43dux43dux44bux435-ux437ux430ux434ux430ux447ux438-ux438-ux440ux435ux433ux440ux435ux441ux441ux438ux43eux43dux43dux44bux435-ux437ux430ux434ux430ux447ux438-ux43fux43e-ux43fux440ux435ux444ux438ux43aux441ux443-ux442ux440ux430ux435ux43aux442ux43eux440ux438ux438}}

После определения причины сбоя ожидаемое поведение оформляется в виде
многократно выполняемой \textbf{регрессионной задачи (regression task)}.
Сквозная регрессия начинается с исходного состояния и проверяет весь
рабочий процесс; регрессия по префиксу траектории фиксирует контекст и
состояние перед определённым решением и сосредоточивается на проверке
следующего действия. Они отвечают соответственно на вопросы «выполнена
ли задача целиком» и «соответствует ли требованиям конкретное решение».

\textbf{Сквозная регрессионная задача} задаёт пользовательский запрос,
исходную среду и бюджет выполнения, позволяет AI-агенту завершить
работу, а затем проверяет конечное состояние, обязательные результаты и
ограничения. Для такой регрессии можно использовать задачи OSWorld,
AndroidWorld, τ²-bench и другие. Сохранение промежуточных результатов
проверок и полной траектории также помогает локализовать сбои.

\textbf{Регрессионная задача по префиксу траектории} предоставляет
существующий диалог, наблюдения инструментов, состояние задачи и
сведения о версиях, после чего агент должен выдать следующее действие
или небольшую последовательность дальнейших действий. Она подходит для
частой проверки границ стратегии, формирования параметров и
восстановления после ошибок, а стоимость её выполнения обычно ниже. Если
действие необходимо выполнить фактически, следует также восстановить
соответствующий снимок среды; если предоставлен только текстовый
префикс, объектом приёмки становится предложенное моделью действие.

\textbf{Для производственных агентов, требующих высокой надёжности,
автор считает, что накопление регрессионных задач по префиксу траектории
зачастую важнее увеличения числа сквозных регрессионных задач.} Задачи
по префиксу фиксируют контекст перед ошибкой, позволяя непосредственно
проверить конкретное решение и часто запускать такую проверку. Для
постоянного накопления подобных примеров разработчикам необходимо
терпеливо выстраивать систему классификации и определения причин сбоев.
Основа регрессии по префиксу --- \textbf{множество допустимых действий}:
в зависимости от текущего состояния разумными следующими шагами могут
быть чтение правил репозитория, получение недостающей информации или
уточняющий вопрос пользователю; одновременно перечисляются действия,
способные привести к проблемам. Категории действий должны однозначно
соответствовать бизнес-семантике, чтобы синонимичные метки не
принимались за различия в поведении.

Между классификацией сбоев и регрессионными задачами можно установить
следующее соответствие:

\begin{longtable}[]{@{}
  >{\raggedright\arraybackslash}p{(\columnwidth - 2\tabcolsep) * \real{0.4338}}
  >{\raggedright\arraybackslash}p{(\columnwidth - 2\tabcolsep) * \real{0.5662}}@{}}
\toprule\noalign{}
\begin{minipage}[b]{\linewidth}\raggedright
\hspace{0pt}Тип сбоя
\end{minipage} & \begin{minipage}[b]{\linewidth}\raggedright
\hspace{0pt}Что должна проверять регрессионная задача
\end{minipage} \\
\midrule\noalign{}
\endhead
\bottomrule\noalign{}
\endlastfoot
\hspace{0pt}Пропущенный этап процесса & \hspace{0pt}Явные требования к
плану, тестированию и коммиту, а также порядок выполнения в полном
процессе \\
\hspace{0pt}Ошибка вызова инструмента & \hspace{0pt}Формат параметров,
экранирование специальных символов, стратегия чтения данных и повторных
попыток после сбоя \\
\hspace{0pt}Ошибка выполнения & \hspace{0pt}Восстановление состояния и
продолжение задачи после усечения, тайм-\allowbreak{}\hspace{0pt}аута
или сбоя инструмента \\
\hspace{0pt}Ошибки полноты и логики & \hspace{0pt}Список нескольких
целей, оставшиеся задачи, разумные границы исследования и условия
остановки \\
\hspace{0pt}Неоднозначность требований & \hspace{0pt}Допустимые
уточняющие вопросы или проверки при недостатке информации \\
\hspace{0pt}Искажение результатов проверки & \hspace{0pt}Защищённые
условия приёмки и фактические результаты выполнения, на которых основано
заявление о завершении \\
\hspace{0pt}Ошибка передачи информации & \hspace{0pt}Соответствие фактам
сумм, состояний, времени и обязательных уведомлений в итоговом ответе \\
\end{longtable}

Эти данные оценки также задают конкретные цели для постобучения модели в
\crossreflink{chap:8}{главе 8} и непрерывной эволюции агентов в
\crossreflink{chap:9}{главе 9}.

\begin{quote}
\textbf{Эксперимент 7-7 ★★: оценка границ по префиксу траектории ---
различные представления одного контекста}

Объедините в примеры уже предоставленные агенту пользовательскую память,
текущую задачу, префикс траектории и состояние среды, потребовав от
модели вывести следующее наблюдаемое действие. Сформируйте примеры по
типам производственных проблем, выявленных через исправления
пользователя, отрицательные оценки и последующий аудит. Охватите
различия в области действия между стилем научной статьи и публикацией в
X, конфликты предпочтений по worktree и PR с правилами репозитория,
приоритет текущих инструкций над старыми предпочтениями, выводы с низкой
уверенностью, подтверждение перед удалением и предварительный просмотр
перед внешней публикацией.

Представьте один и тот же набор семантических полей в трёх форматах:
JSON Cards, Markdown и Python-like. С помощью детерминированных правил
проверьте категорию решения, следующее действие, необходимые
доказательства, запрещённые действия и способ использования памяти.

В сопутствующем прогоне GPT 5.6 Sol обработала комбинации из 11 примеров
и трёх представлений --- всего 33 вызова. По исходным правилам
оценивания все три представления прошли 6 из 11 примеров, однако
результаты по отдельным задачам различались. Анализ ответов показывает
влияние самих правил оценивания: перед удалением модель решила запросить
подтверждение у пользователя, но \texttt{ask\_user} отсутствовал среди
допустимых для этой задачи меток действий; при конфликте старого
предпочтения с текущим требованием модель выбрала текущее требование,
однако соответствующая метка решения также не входила в допустимое
множество.
\end{quote}

Вернёмся к задаче о роуминге из начала главы. Если в текущей траектории
условиям соответствуют обе линии, уточнить номер --- разумное действие.
Если автор задачи принимает только изменение одной конкретной линии, то,
вероятно, в эталонный ответ включена информация, ставшая известной лишь
задним числом. В сквозной задаче можно продолжить работу с учётом
последующего ответа, тогда как задачу по префиксу необходимо оценивать
исключительно на основании доступных в тот момент сведений.

\hypertarget{ux43fux430ux440ux43dux44bux435-ux441ux440ux430ux432ux43dux435ux43dux438ux44f-ux438-ux440ux435ux439ux442ux438ux43dux433-ux43cux43eux434ux435ux43bux435ux439}{%
\subsection{Парные сравнения и рейтинг
моделей}\label{ux43fux430ux440ux43dux44bux435-ux441ux440ux430ux432ux43dux435ux43dux438ux44f-ux438-ux440ux435ux439ux442ux438ux43dux433-ux43cux43eux434ux435ux43bux435ux439}}

Прочитав два ответа на один и тот же открытый вопрос, люди обычно могут
сказать, какой из них им нравится больше. Собрав множество таких
\textbf{парных сравнений}, можно оценить относительные рейтинги моделей.
Рассмотрим два ответа о возврате средств на
\crossreflink{fig:7-6}{рисунке 7-6}: в одном указаны сумма и срок, а в
другом лишь сказано, что возврат будет обработан. Голосование
показывает, какой ответ читатели предпочли в этом случае. Чтобы
преобразовать этот голос в баллы, необходимо также знать разницу между
исходными рейтингами двух соперников.

\begin{figure}
\hypertarget{fig:7-6}{%
\centering
\includegraphics{images/revised/fig7-6.pdf}
\caption{Рисунок 7-6. Рейтинг Elo и ранжирование по парным
сравнениям}\label{fig:7-6}
}
\end{figure}

\textbf{Рейтинг Elo} изначально применялся в соревнованиях по играм на
доске. При использовании распространённой шкалы в 400 пунктов ожидаемый
результат модели A в сравнении с моделью B равен:

\[
E_A=\frac{1}{1+10^{(R_B-R_A)/400}}
\]

Например, если рейтинг A равен 1200, а B --- 1000, ожидаемый результат A
составляет примерно 0,76. В бинарной модели, где возможны только победа
и поражение, это также ожидаемая вероятность победы. Фактические
результаты победы, ничьей и поражения обозначаются соответственно как 1,
0,5 и 0, после чего рейтинг обновляется по следующему правилу:

\[
R_A'=R_A+K(S_A-E_A)
\]

Здесь \(S_A\) --- фактический результат, а \(K\) определяет влияние
отдельного сравнения. Если B неожиданно побеждает, фактический результат
A оказывается значительно ниже ожидаемого и рейтинг A заметно снижается.
Если A побеждает в соответствии с ожиданиями, изменение рейтинга
невелико. Рейтинг B обновляется симметрично на основе его фактического и
ожидаемого результатов. После последовательной обработки сравнений
рейтинги отражают отношения предпочтения в этом наборе данных.

\textbf{Модель Брэдли --- Терри} также объясняет вероятность победы
разностью скрытых уровней силы двух участников и позволяет подобрать
рейтинги сразу по всему набору сравнений методом максимального
правдоподобия. В онлайновом Elo рейтинг обновляется после каждого
сравнения, поэтому результат зависит от \(K\) и порядка поступления
данных. При пакетном подборе рейтинги всех участвующих моделей
оцениваются одновременно для фиксированных данных и заданной модели. В
2024 году Chatbot Arena использовала метод ранжирования Брэдли ---
Терри.\footnote{LMArena, «Набор данных рейтинга Arena», содержит историю
  использования оценивания Bradley--Terry в 2024 году. В эксперименте
  этого раздела фиксированный исторический снимок голосов используется
  для сравнения онлайн-рейтинга Elo с пакетной подгонкой модели.}

В Chatbot Arena пользователи анонимно сравнивают ответы двух моделей.
Чьи предпочтения отражает рейтинг и какие способности моделей он
оценивает, зависят от того, кто голосует, какие вопросы задают
пользователи и какие модели получают возможность встретиться в
сравнении. Если велика доля задач по программированию, итоговый рейтинг
будет в большей степени отражать качество программирования.

Для парных сравнений можно также применять оценивание моделями. В таком
случае A и B следует поменять местами и провести оценивание ещё раз,
чтобы проверить \textbf{позиционное смещение (Position Bias)}. Сначала
результаты двух оцениваний сопоставляются с соответствующими моделями и
фиксируется победитель либо ничья, после чего результаты объединяются по
заранее определённым правилам. При противоречии между результатами можно
отметить расхождение, провести дополнительное оценивание или передать
случай на проверку человеку.

\begin{quote}
\textbf{Эксперимент 7-8 ★★: построение рейтинга моделей по данным парных
сравнений}

Реализуйте онлайновый Elo на основе общедоступного снимка голосований
Chatbot Arena. В сопутствующем запуске из примерно 1,8 миллиона исходных
сравнений было отобрано 1 670 250. Начальный рейтинг каждой модели равен
1000. Сравнения считываются в хронологическом порядке, для каждого
вычисляется ожидаемый результат, после чего рейтинги обеих моделей
обновляются с фиксированным шагом \(K=4\). В 281 121 сравнении
пользователи выбрали ничью, а ещё в 295 254 --- вариант «оба ответа
плохие». В обоих случаях каждой модели начисляется по половине балла.
Если продуктовой команде важно знать, в скольких сравнениях пользователи
остались недовольны обоими ответами, исходную метку второго варианта
необходимо сохранить.

При тех же правилах фильтрации данных отдельно выполните пакетный подбор
рейтингов по модели Брэдли --- Терри. Сравните корреляцию двух
рейтингов, пересечение среди нескольких верхних позиций и модели с
наибольшими расхождениями. Изменяйте \(K\) и порядок сравнений, чтобы
проверить влияние онлайнового обновления на рейтинг. В этой серии
сравнений участвовало 129 моделей, а среди первых 20 позиций двух
рейтингов совпали 12.

Попробуйте интерпретировать пустую ячейку матрицы попарных результатов
(\texttt{win\_rate\_matrix.json}). В строке \texttt{chatgpt-4o-latest} и
столбце \texttt{bard-jan-24-gemini-pro} фактическая доля набранных
баллов отсутствует: в отфильтрованных данных эти модели не сравнивались
напрямую. Однако обе сравнивались с другими моделями, а их итоговые
онлайновые рейтинги Elo составили соответственно примерно 1126,69 и
1094,21. Подстановка этих значений в формулу даёт ожидаемый результат
первой модели примерно 0,547. Чтобы ответить на вопрос «какую из этих
двух моделей в действительности выбирали пользователи», всё равно
необходимо собрать голоса по их прямым сравнениям.

Сопутствующая реализация ежемесячно формирует накопительные снимки:
рейтинг за каждый месяц учитывает все сравнения, включённые к этому
моменту. Если вместо этого рассматривать только последние четыре недели,
рейтинги будут меняться ещё и из-за исключения старых сравнений из окна.
При повторении эксперимента сначала выберите один способ расчёта, а
затем объясняйте изменения позиций. Создайте с помощью D3.js анимацию
столбчатой диаграммы, где длина столбца обозначает рейтинг, а положение
по вертикали --- место в рейтинге. Объясняйте изменения рейтинга с
учётом времени добавления модели, количества сравнений, состава
соперников и распределения задач. Для позиций с заметными изменениями
вернитесь к конкретным данным и проверьте причины.
\end{quote}

\hypertarget{ux432ux44bux431ux43eux440-ux43cux43eux434ux435ux43bux438-ux43dux430-ux43eux441ux43dux43eux432ux435-ux43eux446ux435ux43dux43aux438}{%
\section{Выбор модели на основе
оценки}\label{ux432ux44bux431ux43eux440-ux43cux43eux434ux435ux43bux438-ux43dux430-ux43eux441ux43dux43eux432ux435-ux43eux446ux435ux43dux43aux438}}

При выборе модели необходимо ответить на конкретный вопрос: какая
конфигурация обеспечит лучший совокупный результат при текущем
распределении задач, эксплуатационном бюджете и требованиях к
надёжности? Для этого необходимо одновременно измерять качество
выполнения, скорость ответа, потребление ресурсов и стабильность работы.

\hypertarget{ux43aux43bux44eux447ux435ux432ux44bux435-ux43fux430ux440ux430ux43cux435ux442ux440ux44b-ux432ux44bux431ux43eux440ux430}{%
\subsection{Ключевые параметры
выбора}\label{ux43aux43bux44eux447ux435ux432ux44bux435-ux43fux430ux440ux430ux43cux435ux442ux440ux44b-ux432ux44bux431ux43eux440ux430}}

Чтобы разобраться в показателях скорости, следует начать с двух этапов
инференса. \textbf{Prefill (предварительное заполнение)} обрабатывает
входной контекст, а \textbf{Decode (декодирование)} последовательно
генерирует последующие токены. Более длинный ввод обычно требует больше
вычислений на этапе предварительного заполнения, однако на фактическое
время также влияют кэширование префикса, пакетная обработка и реализация
инференса. Скорость декодирования определяет время завершения вывода:
при стабильной скорости 50 токенов/с генерация 2000 токенов рассуждения
займёт около 40 секунд.

При измерениях необходимо различать отдельный запрос и сервис в целом, а
также явно определять начальную и конечную точки отсчёта времени.

\begin{longtable}[]{@{}
  >{\raggedright\arraybackslash}p{(\columnwidth - 4\tabcolsep) * \real{0.2551}}
  >{\raggedright\arraybackslash}p{(\columnwidth - 4\tabcolsep) * \real{0.3446}}
  >{\raggedright\arraybackslash}p{(\columnwidth - 4\tabcolsep) * \real{0.4003}}@{}}
\toprule\noalign{}
\begin{minipage}[b]{\linewidth}\raggedright
\hspace{0pt}Показатель
\end{minipage} & \begin{minipage}[b]{\linewidth}\raggedright
\hspace{0pt}Объект измерения
\end{minipage} & \begin{minipage}[b]{\linewidth}\raggedright
\hspace{0pt}Условия, учитываемые при интерпретации
\end{minipage} \\
\midrule\noalign{}
\endhead
\bottomrule\noalign{}
\endlastfoot
\hspace{0pt}Задержка до первого токена (TTFT) & \hspace{0pt}От отправки
запроса до получения первого содержательного выходного токена &
\hspace{0pt}Включает затраты на сеть, ожидание в очереди,
предварительное заполнение, генерацию первого токена и другие операции;
потоковый протокол также может буферизовать вывод \\
\hspace{0pt}Скорость вывода одного запроса & \hspace{0pt}Число токенов,
выводимых за единицу времени на этапе генерации & \hspace{0pt}Длина
ввода, фактическая длина вывода, параллелизм и планирование на стороне
сервера \\
\hspace{0pt}Совокупная пропускная способность & \hspace{0pt}Общее число
входных или выходных токенов, обработанных сервисом за определённое
время & \hspace{0pt}Уровень параллелизма, способ пакетной обработки,
доля успешных запросов; ввод и вывод учитываются раздельно \\
\hspace{0pt}Время, связанное с рассуждением & \hspace{0pt}Время,
затраченное рассуждающей моделью на наблюдаемых этапах &
\hspace{0pt}Бюджет рассуждения, доступность событий рассуждения в
сервисе и итоговая польза для задачи \\
\hspace{0pt}Сквозная задержка & \hspace{0pt}От начала запроса или задачи
до завершения & \hspace{0pt}Совместный вклад нескольких вызовов модели,
выполнения инструментов, повторных попыток и восстановления состояния \\
\hspace{0pt}Хвостовая задержка p95 &
\hspace{0pt}95-\allowbreak{}\hspace{0pt}й процентиль распределения
задержки при одинаковой нагрузке & \hspace{0pt}Около 95\% образцов не
превышают это значение; для анализа длинного хвоста его можно
рассматривать вместе со средним значением и p99 \\
\end{longtable}

У моделей со скрытым процессом рассуждения первый видимый пользователю
токен может поступить только после завершения внутреннего инференса. В
эксперименте метрики времени до «первого потокового содержимого» и до
«первого видимого ответа» следует называть раздельно. Одинаковая
скорость 50 токенов/с при разной длине рассуждения может означать
совершенно разное время ожидания для пользователя. В интерактивных
продуктах автор уделяет особое внимание p95, p99 и соответствующим им
траекториям: множество быстрых запросов способно снизить среднее
значение, но небольшое число надолго зависших задач всё равно серьёзно
ухудшит пользовательский опыт.

\textbf{Стоимость} следует рассчитывать на выполненную задачу.
Записывайте стоимость входных данных, кэшированного ввода, вывода и
инструментов, а затем рассчитывайте среднюю стоимость и её
распределение. Модель с дешёвым отдельным вызовом, но множеством
повторных попыток после неудач может обходиться дороже за одну
выполненную задачу.

\textbf{Качество выполнения задач} можно измерять с помощью
рассмотренных выше \href{mailto:Pass@1}{\nolinkurl{Pass@1}},
\href{mailto:Pass@k}{\nolinkurl{Pass@k}}, Pass\^{}k и
\href{mailto:Best@k}{\nolinkurl{Best@k}}. Для обычных одиночных запросов
важен \href{mailto:Pass@1}{\nolinkurl{Pass@1}}, для стабильности
повторных выполнений --- Pass\^{}k, а в исследовательских задачах с
несколькими допустимыми попытками можно сравнивать
\href{mailto:Pass@k}{\nolinkurl{Pass@k}}.
\href{mailto:Best@k}{\nolinkurl{Best@k}} фиксирует наивысшую оценку
среди этих \(k\) ответов; при фактическом развёртывании следует отдельно
оценивать результат, выбранный селектором. Определив метрики с учётом
бизнеса, можно выбирать число экспериментов и бюджет выполнения.

\textbf{Ограничения скорости и надёжность} влияют на пропускную
способность сервиса. RPM обозначает лимит числа запросов в минуту, TPM
--- лимит числа токенов в минуту. При выборе также необходимо учитывать
тайм-ауты, ограничения скорости, ошибки сервиса, сохранение цели в
длительных задачах и восстановление после ошибок.

\textbf{Кривая «бюджет --- возможности»} показывает, как изменяется
результат задачи при увеличении времени, числа токенов, вызовов
инструментов или вычислительных ресурсов. При сравнении людей и моделей
в RE-Bench с общим бюджетом в 2 часа на каждую среду результат лучшего
AI-агента примерно в четыре раза превосходил результат
человека-эксперта; при увеличении бюджета до 8 часов люди немного вышли
вперёд; за суммарные 32 часа нескольких попыток результат людей был
примерно вдвое выше результата лучшего агента.\footnote{Wijk, Hjalmar,
  et al.~\emph{RE-Bench: Evaluating Frontier AI R\&D Capabilities of
  Language Model Agents against Human Experts.} arXiv:2411.15114, 2025.}
Агент, лидировавший при двухчасовом бюджете, при более продолжительной
работе уступил людям.

Несколько моделей можно распределить по типам задач: лёгкая модель
обрабатывает подходящие ей запросы, более мощная --- сложное
оркестрирование, а специализированные модели --- понимание изображений
или генерацию кода. Сам маршрутизатор также необходимо оценивать.
Например, фраза «Что больше: 9.9 или 9.11?» коротка, но требует
понимания чисел; вопрос «До автомойки от дома 50 метров --- пойти пешком
или поехать на машине?» требует ещё и учитывать цель задачи. Отдельно
измеряйте правильную маршрутизацию, ошибочную маршрутизацию и
восстановление после передачи задачи, а затем сравнивайте качество,
задержку и стоимость полной задачи с вариантом на одной модели. Автор
использует более сложную комбинацию моделей, только если совокупная
польза покрывает ошибки маршрутизации и затраты на сопровождение.

\hypertarget{ux43fux43eux432ux435ux434ux435ux43dux447ux435ux441ux43aux430ux44f-ux441ux442ux440ux430ux442ux435ux433ux438ux44f-ux43cux43eux434ux435ux43bux438}{%
\subsection{Поведенческая стратегия
модели}\label{ux43fux43eux432ux435ux434ux435ux43dux447ux435ux441ux43aux430ux44f-ux441ux442ux440ux430ux442ux435ux433ux438ux44f-ux43cux43eux434ux435ux43bux438}}

Разные модели могут выбирать разные пути действий для одной и той же
задачи. Типичное различие между Coding Agent заключается в
\textbf{пороге перехода к действиям}: одни модели сначала широко
исследуют репозиторий и лишь затем начинают вносить изменения; другие по
локальной информации сначала создают патч, а потом дополняют своё
понимание с помощью тестов. Первые тратят больше времени на
исследование, вторые сильнее полагаются на быструю обратную связь для
исправления локальных суждений.

Эта склонность совместно формируется моделью и Harness. Системный промпт
может требовать сначала проверить правила, запустить тесты или
ограничить область исследования; постобучение также влияет на действия
модели по умолчанию. Траектории SFT демонстрируют, когда читать и когда
вносить изменения, вознаграждение за процесс оценивает конкретные пути
использования инструментов, а вознаграждение за результат закрепляет
сочетания действий в успешных траекториях. Чтобы наблюдать эти привычки,
необходимо сделать промежуточные действия объектом оценки.

\begin{quote}
\textbf{Эксперимент 7-9 ★★: измерение порога перехода модели к действиям
в фиксированном Coding Harness}

Предоставьте GPT-5.6-sol и Claude Sonnet 5 один и тот же системный
промпт, одну и ту же схему инструментов, один и тот же репозиторий
задачи, одну и ту же команду тестирования и одинаковое максимальное
число раундов и вызывайте их через одну и ту же конечную точку
совместимого сервиса. Нейтральный промпт позволяет модели самостоятельно
определять порядок исследования и внесения изменений. Три небольших
репозитория должны охватывать локальную ошибку, межмодульную
нормализацию идентичности и чувствительное к публичному контракту
исправление кэша. Каждая модель независимо выполняет каждую задачу
трижды --- всего 18 траекторий.

Запишите вызовы инструментов, число независимо прочитанных файлов и
затраченное время до первой попытки редактирования, а затем проверьте
первый протестированный патч, итоговые тесты, повторную работу и общую
стоимость. В сводных результатах до первой попытки редактирования
GPT-5.6-sol в среднем вызывала инструменты 6.89 раза и читала 4.67
файла, а Claude Sonnet 5 --- соответственно 4.56 раза и 3.56 файла. В
межмодульной задаче число прочитанных файлов оказалось близким: 7.00 и
6.67 соответственно. Первые протестированные патчи и итоговые тесты
обеих моделей прошли успешно. К моменту возврата из инструмента первого
редактирования прошло в среднем соответственно 15.01 и 14.48 секунды.
Среда запуска фиксировала время при возврате инструмента, поэтому эти
два значения также включают время обработки первого запроса на
редактирование инструментом.

Почему измеряется время до «попытки редактирования»? Рассмотрим
траекторию второй попытки Claude Sonnet 5 исправить задачу с кэшем:
модель один раз выполнила поиск, прочитала два файла, а затем попыталась
изменить существующий файл инструментом, способным только создавать
новые файлы. Инструмент отклонил запрос. Модель перешла на инструмент
замены текста, изменила код, после чего первый же тест прошёл.

Это сравнение демонстрирует различия в путях действий при фиксированных
условиях. Добавив условие промпта «сначала исследовать», можно провести
перекрёстный эксперимент «модель × стратегия промпта»: сравнить модели
при одном промпте и промпты для одной модели. Записывайте раздельно
число исследовательских действий, параллельные вызовы инструментов и
задержку модели, а затем оценивайте пользу дополнительного чтения или
раннего редактирования вместе с итоговым качеством.
\end{quote}

\hypertarget{ux430ux43dux430ux43bux438ux437-ux437ux430ux442ux440ux430ux442-ux430ux433ux435ux43dux442ux43dux43eux439-ux441ux438ux441ux442ux435ux43cux44b}{%
\subsection{Анализ затрат агентной
системы}\label{ux430ux43dux430ux43bux438ux437-ux437ux430ux442ux440ux430ux442-ux430ux433ux435ux43dux442ux43dux43eux439-ux441ux438ux441ux442ux435ux43cux44b}}

Затраты многораундового агента складываются из вызовов модели,
выполнения инструментов и инфраструктуры. Контекст непрерывно
накапливается, поэтому ранняя информация может многократно появляться в
последующих запросах и расходы необходимо декомпозировать по всей
траектории задачи.

\textbf{Стоимость инференса модели} включает тарифицируемые входные
данные, кэшированный ввод, вывод и другие составляющие. Предположим, в
первом, втором и третьем раундах вводится соответственно 1000, 2000 и
3000 токенов: суммарный объём ввода составит 6000 токенов. Если в каждом
раунде добавляется одинаковый объём и полностью сохраняется история,
суммарный ввод растёт приблизительно квадратично с числом раундов.
Кэширование префикса позволяет повторно использовать уже выполненные
вычисления, а скидки API рассчитываются по соответствующим правилам
тарификации.

\textbf{Стоимость вызовов инструментов} включает плату за внешние
сервисы, такие как поиск и доступ к базам данных, а также вычислительные
ресурсы песочницы для выполнения кода. Результаты инструментов также
становятся входными данными модели: один поиск может вернуть несколько
тысяч токенов, и их сохранение в нескольких последующих раундах
многократно увеличит стоимость ввода.

\textbf{Инфраструктурные затраты} включают векторные базы данных,
очереди сообщений, реляционные базы данных, а также хранилища журналов и
трассировок. После распределения этих расходов по задачам или
арендаторам можно дополнительно рассчитать совокупную стоимость
выполнения одной бизнес-операции.

Рассмотрим расчёт для выбора конфигурации. Предположим, две конфигурации
обрабатывают одну и ту же группу из 100 задач: общая стоимость варианта
А составляет 60 юаней при 60 успешных выполнениях, а варианта Б --- 72
юаня при 90 успешных выполнениях. Отдельная попытка в варианте Б дороже,
однако каждое успешное выполнение задачи обходится в 1 юань для варианта
А и в 0.8 юаня для варианта Б. Числитель здесь включает деньги,
потраченные на неудачные попытки. Если в двух группах используются
разные определения успеха --- например, в одной достаточно согласия
службы поддержки, а в другой требуется проверить счёт, --- напрямую
сравнивать эти значения стоимости нельзя.

В сопутствующем эксперименте используется фиксированный восьмираундовый
процесс возврата средств службой поддержки. На последовательных этапах
предоставляются заранее заданные результаты по заказу, доставке,
правилам, базе знаний, контролю рисков, возврату средств, уведомлению и
закрытию обращения, на основе которых модель генерирует последующие
ответы. В эксперименте независимо контролируются стабильность префикса и
сжатие истории, образуя четыре контрольные группы. В таблице 7-4
приведена стоимость вызовов модели, рассчитанная по зафиксированному
потреблению gpt-4o-mini и действовавшим на тот момент ценам.

Таблица 7-4. Стоимость вызовов модели в фиксированном восьмираундовом
процессе

\begin{longtable}[]{@{}
  >{\raggedright\arraybackslash}p{(\columnwidth - 8\tabcolsep) * \real{0.2559}}
  >{\raggedleft\arraybackslash}p{(\columnwidth - 8\tabcolsep) * \real{0.1439}}
  >{\raggedleft\arraybackslash}p{(\columnwidth - 8\tabcolsep) * \real{0.2033}}
  >{\raggedleft\arraybackslash}p{(\columnwidth - 8\tabcolsep) * \real{0.1790}}
  >{\raggedleft\arraybackslash}p{(\columnwidth - 8\tabcolsep) * \real{0.2180}}@{}}
\toprule\noalign{}
\begin{minipage}[b]{\linewidth}\raggedright
\hspace{0pt}Вариант
\end{minipage} & \begin{minipage}[b]{\linewidth}\raggedleft
\hspace{0pt}Входные токены
\end{minipage} & \begin{minipage}[b]{\linewidth}\raggedleft
\hspace{0pt}Кэшированные токены
\end{minipage} & \begin{minipage}[b]{\linewidth}\raggedleft
\hspace{0pt}Стоимость модели
\end{minipage} & \begin{minipage}[b]{\linewidth}\raggedleft
\hspace{0pt}Экономия относительно базового варианта
\end{minipage} \\
\midrule\noalign{}
\endhead
\bottomrule\noalign{}
\endlastfoot
\hspace{0pt}Нестабильный префикс, история сохраняется &
\hspace{0pt}20,700 & \hspace{0pt}0 & \hspace{0pt}\$0.003776 &
\hspace{0pt}--- \\
\hspace{0pt}Только стабильный префикс & \hspace{0pt}20,386 &
\hspace{0pt}13,568 & \hspace{0pt}\$0.002708 & \hspace{0pt}28.3\% \\
\hspace{0pt}Только сжатие истории & \hspace{0pt}16,177 & \hspace{0pt}0 &
\hspace{0pt}\$0.003115 & \hspace{0pt}17.5\% \\
\hspace{0pt}Стабильный префикс и сжатие включены одновременно &
\hspace{0pt}16,035 & \hspace{0pt}6,144 & \hspace{0pt}\$0.002643 &
\hspace{0pt}30.0\% \\
\end{longtable}

Столбец входных токенов включает кэшированную часть, которую при расчёте
стоимости следует отделять от некэшированного ввода. В базовой группе
объём ввода за раунд вырос с 1113 до 3668 токенов, а результаты
инструментов за восемь раундов внесли в совокупности 9544 входных
токена. При одновременном включении обеих оптимизаций вклад результатов
инструментов снизился до 5248 токенов.

При совместном включении двух оптимизаций измеренная экономия составила
30.0\%. Сжатие сокращает последующий ввод, но также может изменить длину
повторно используемого префикса; между запусками могут различаться и
длина вывода, и попадания в кэш. Поэтому доли экономии, измеренные по
отдельности для двух оптимизаций, нельзя просто складывать.

\textbf{Оптимизация затрат и управление выполнением.} Со стороны ввода в
первую очередь можно сравнить стабильный префикс и сжатие истории,
рассмотренные во второй главе; со стороны модели --- описанную выше в
этой главе многоуровневую маршрутизацию.

Для задач, не требующих выполнения в реальном времени, можно
использовать \textbf{асинхронную пакетную обработку}: объединять их для
планирования в пределах допустимого окна завершения и применять пакетные
тарифы сервиса. При локальном развёртывании ожидающие задачи можно
планировать на периоды меньшей загрузки ресурсов, повышая коэффициент
использования оборудования.

В промышленной системе также следует отслеживать затраты по типу задачи,
модели и пользователю, устанавливать бюджет задач и поэтапные
контрольные точки. При приближении к верхнему пределу бюджета Harness
может сохранить состояние, сузить область дальнейшего исследования или
передать решение об оставшейся части задачи пользователю.

\begin{quote}
\textbf{Эксперимент 7-10 ★: сквозной анализ стоимости задач агента}

Сначала освойте последовательный учёт на фиксированном восьмираундовом
процессе, а затем замените его собственной типичной рабочей нагрузкой. С
помощью LangSmith или собственной системы трассировки записывайте
входные, кэшированные входные, выходные токены и токены рассуждения,
число вызовов инструментов и размер их результатов, а также сквозную
задержку запросов и задач. Рассчитайте стоимость модели, инструментов и
инфраструктуры по правилам тарификации соответствующих сервисов.

Сначала попробуйте предположить: если кэширование префикса уже включено,
обязательно ли дополнительное сжатие истории снизит стоимость восьми
вызовов модели? В сопутствующей трассировке Kimi K2.5 при использовании
только стабильного префикса из 19,508 входных токенов 14,779
тарифицировались по цене кэшированных. После дополнительного сжатия
истории общий объём ввода снизился до 15,303 токенов, однако часть,
тарифицируемая по обычной цене, выросла с 4,729 до 7,180 токенов. С
учётом вывода стоимость восьми раундов, напротив, увеличилась примерно с
0.00825 до 0.00906 доллара.
\end{quote}

\hypertarget{ux43dux435ux43fux440ux435ux440ux44bux432ux43dux430ux44f-ux438ux442ux435ux440ux430ux446ux438ux44f-ux43dux430-ux43eux441ux43dux43eux432ux435-ux43eux446ux435ux43dux43aux438}{%
\subsection{Непрерывная итерация на основе
оценки}\label{ux43dux435ux43fux440ux435ux440ux44bux432ux43dux430ux44f-ux438ux442ux435ux440ux430ux446ux438ux44f-ux43dux430-ux43eux441ux43dux43eux432ux435-ux43eux446ux435ux43dux43aux438}}

Оценочный набор превращает смену модели в воспроизводимое инженерное
решение. Предположим, сейчас система использует Claude, а команда хочет
оценить недавно выпущенную Gemini. Сначала можно зафиксировать задачи,
инструменты, бюджет и правила приёмки, а затем сравнить долю успеха,
корректность вызовов, задержку и стоимость.

Результаты сравнения часто зависят от задачи. Например, новая модель
может быть дешевле и успешнее на простых задачах, но показывать снижение
доли успеха на 5 процентных пунктов в задачах со сложным многораундовым
оркестрированием. После подтверждения влияния с помощью повторных
запусков и интервала различий подходящие задачи можно перенести на новую
модель, сохранив прежнюю конфигурацию для остальных, а затем оценить
дополнительные затраты и ошибки маршрутизации. В следующем разделе
рассматривается, как определить, достаточно ли наблюдаемого различия для
принятия решения.

\begin{quote}
\textbf{Эксперимент 7-11 ★★: многомерное тестирование производительности
моделей}

Выберите актуальные для бизнеса конкретные версии моделей GPT, Claude,
Gemini, Doubao, Qwen, Kimi, DeepSeek и других и сформируйте матрицу
конфигураций; для одинаковых весов или одной версии модели можно также
сравнить разных поставщиков сервиса. Запишите доступные сведения о
конечной точке сервиса, квантовании, конфигурации инференса и других
параметрах и сравните результаты со сторонними данными мониторинга,
используя одинаковые определения измерений.

Зафиксируйте уровни ввода, например 8K, 32K и 128K токенов, и целевые
объёмы вывода, например 512 и 2048 токенов. Выберите сочетания в
пределах возможностей моделей и записывайте фактическое число токенов.
Выполните каждую конфигурацию не менее 100 раз и рассчитайте TTFT,
сквозную задержку, скорость вывода, стандартное отклонение и процентили.
Токены рассуждения и связанное с ними время записывайте отдельно в
пределах наблюдаемости интерфейса.

В течение недели проверяйте сервис каждый час, записывая успешные
выполнения, тайм-ауты, ограничения скорости и другие ошибки, а затем
обобщите интервалы непрерывных сбоев и время восстановления. Почасовые
проверки дают картину доступности с разрешением, соответствующим частоте
выборки; для точной оценки продолжительности сбоя необходимо повысить
частоту проверок или использовать записи о событиях сервиса. Отдельно
отмечайте не полностью наблюдавшиеся сбои, интервалы которых пересекают
начало или конец периода наблюдения.

Постепенно увеличивайте параллелизм в пределах согласованной для учётной
записи нагрузки, наблюдайте за долей успеха, совокупной пропускной
способностью и хвостовой задержкой и записывайте квоты RPM и TPM, а
также условия срабатывания ограничений скорости. Для расчёта стоимости
используйте датированные цены входных, кэшированных входных и выходных
токенов; при сравнении разных валют зафиксируйте обменный курс, а затем
рассчитайте стоимость типичной многораундовой задачи.
\end{quote}

\begin{quote}
\textbf{Эксперимент 7-12 ★★: сквозная оценка выбора системы
пользовательской памяти}

Повторно используйте поиск в контексте или агентный RAG из третьей главы
и на одном наборе из 60 вопросов о пользовательской памяти сравните три
точки выбора: модель эмбеддингов, reranker и основную модель агента.

В число моделей эмбеддингов можно включить соответствующие версии
BGE-M3, OpenAI, Qwen или Mistral; среди вариантов reranker должны быть
базовый вариант без переранжирования и конфигурация семантического
переранжирования; для основной модели при одинаковых условиях поиска
сравнивайте долю успеха и эффективность использования инструментов.
Отдельно записывайте долю попаданий top-5, полноту извлечения
необходимой информации, задержку поиска и стоимость, а затем --- оценку
итогового ответа по рубрике.

Сопутствующая матрица включает 4 варианта эмбеддингов, 3 конфигурации
переранжирования и 2 основные модели, образуя 24 конфигурации. При
анализе сначала фиксируйте две другие точки выбора и сравнивайте
изменение одного компонента, а затем проверяйте эффекты сочетаний: при
каких условиях поиска переранжирование приносит пользу, способна ли
основная модель использовать разрозненные свидетельства и как она
действует при отсутствии ключевого факта.

Рассмотрим одну ячейку сопутствующей матрицы: эмбеддинги BGE-M3,
семантическое переранжирование Doubao и ответы Doubao. В этой
конфигурации были даны ответы на 60 вопросов. Среди первых пяти
извлечённых материалов для каждого вопроса был хотя бы один релевантный,
однако из всех фрагментов, помеченных как релевантные, в первую пятёрку
в среднем попадало лишь около 47.8\% --- сначала показатель
рассчитывался для каждого вопроса, а затем усреднялся. Итоговые ответы
прошли проверку по правилам оценивания для 25 вопросов. Эти 25
результатов также можно разделить по уровням: на каждом из трёх уровней
было по 20 вопросов, и проверку прошли ответы соответственно на 18, 5 и
2 из них.

Например, в одном из вопросов первого уровня помощнику требовалось
перечислить все упомянутые пользователем хозяйственные расходы и их
суммы. Система поиска извлекла все фрагменты, помеченные как
релевантные, поэтому \href{mailto:recall@5}{\nolinkurl{recall@5}} для
этого вопроса равнялся 1; однако в записи оценки отмечено, что в ответе
отсутствовали расходы на домашний офис в размере 1,525 долларов. Здесь
необходимо различать два вида пропусков: первый показывает, сколько
фрагментов материалов было пропущено при поиске, второй --- какие
требуемые пункты ответа отсутствовали в итоговом ответе.
\end{quote}

\hypertarget{ux441ux442ux430ux442ux438ux441ux442ux438ux447ux435ux441ux43aux430ux44f-ux437ux43dux430ux447ux438ux43cux43eux441ux442ux44c-ux440ux435ux437ux443ux43bux44cux442ux430ux442ux43eux432-ux43eux446ux435ux43dux43aux438}{%
\section{Статистическая значимость результатов
оценки}\label{ux441ux442ux430ux442ux438ux441ux442ux438ux447ux435ux441ux43aux430ux44f-ux437ux43dux430ux447ux438ux43cux43eux441ux442ux44c-ux440ux435ux437ux443ux43bux44cux442ux430ux442ux43eux432-ux43eux446ux435ux43dux43aux438}}

В одном отчёте две конфигурации правильно решили соответственно 70 и 73
задачи. Сохранится ли преимущество в три задачи, если в следующий раз
взять другой набор? Оценка охватывает лишь ограниченное число задач, а
работа AI-агента подвержена случайности. Сначала оценим колебания
отдельной доли успехов, а затем сравним различия по результатам двух
конфигураций на одних и тех же задачах.

Если считать \(n\) задач независимой выборкой и для каждой фиксировать
успех или неудачу, то стандартную ошибку измеренной доли успехов
\(\hat p\) можно грубо оценить следующим образом:

\[
\mathrm{SE}(\hat p)\approx\sqrt{\frac{\hat p(1-\hat p)}{n}}
\]

Например, если успешно выполнены 70 из 100 примеров, то 95\%-й
доверительный интервал при нормальном приближении составляет примерно
\(70\%\pm9\) процентных пунктов. Этот интервал характеризует точность
оценки доли успехов одной группы; для сравнения новой и старой моделей
необходимо непосредственно оценить интервал разности между двумя
группами. При малом размере выборки или доле успехов, близкой к 0 или 1,
можно использовать интервал Уилсона либо подходящий метод повторной
выборки.

При сравнении двух конфигураций на одном наборе задач следует отдавать
предпочтение \textbf{парному анализу}. Для каждой задачи нужно
фиксировать один из четырёх исходов: «обе конфигурации успешны»,
«успешна только A», «успешна только B» или «обе неуспешны». Различие
проявляется главным образом в двух средних категориях. Если для каждой
задачи имеется один бинарный результат, несогласованные пары можно
сравнить с помощью критерия Мак-Немара; другой вариант --- выполнить
парный бутстрэп по задачам и оценить интервал средней разности. Вернёмся
к исходным 100 задачам. Предположим, B сохранила все 70 правильных
ответов A и дополнительно правильно решила ещё 3 задачи. Тогда следующим
шагом можно выяснить, что общего у этих трёх задач. Но возможен и другой
случай: B улучшила результат на 23 задачах, однако ошиблась в 20
задачах, которые раньше были решены правильно. В обоих отчётах будет
указано повышение с 70\% до 73\%, но во втором случае следует отдельно
исследовать каждую из 20 регрессий. Чтобы оценить надёжность этих трёх
процентных пунктов, также необходимы результаты по каждой задаче.

Повторные запуски позволяют наблюдать собственные колебания модели. Для
каждой конфигурации можно провести 3--5 независимых запусков, сохраняя
неизменными экземпляры задач, параметры среды, бюджет и правила приёмки.
Если сервис поддерживает начальные значения генератора случайных чисел,
их следует фиксировать; даже одно и то же значение для разных моделей
нужно считать отдельной настройкой выборки каждой модели. Повторные
запуски одной задачи имеют общий контекст, поэтому при статистическом
анализе результаты можно сначала агрегировать на уровне задач либо
применять повторную выборку с кластеризацией по задачам, а не считать
все запуски независимыми новыми задачами.

Конкретный процесс выглядит так: зафиксировать задачи и условия,
запустить на каждой задаче обе конфигурации, сохранить парные
результаты, агрегировать разности по задачам, а затем оценить общее
различие и его интервал. Если ожидаемый эффект составляет всего 2--3
процентных пункта, нескольких десятков задач обычно недостаточно для
точной оценки; при приближении независимой выборки стандартная ошибка
уменьшается с ростом размера выборки как \(1/\sqrt n\).

Ниже парный расчёт показан на четырёх задачах. Ключи словаря содержат
задачу и начальное значение генератора случайных чисел, а значения ---
результаты, полученные по одним и тем же правилам приёмки, где 1
означает успех. Данные специально составлены для демонстрации; код можно
непосредственно запустить в Python.

\begin{Shaded}
\begin{Highlighting}[]
\NormalTok{a }\OperatorTok{=}\NormalTok{ \{(}\StringTok{"проверить заказ"}\NormalTok{, }\DecValTok{7}\NormalTok{): }\DecValTok{1}\NormalTok{, (}\StringTok{"изменить адрес"}\NormalTok{, }\DecValTok{7}\NormalTok{): }\DecValTok{1}\NormalTok{,}
\NormalTok{     (}\StringTok{"проверить доставку"}\NormalTok{, }\DecValTok{7}\NormalTok{): }\DecValTok{0}\NormalTok{, (}\StringTok{"изменить время"}\NormalTok{, }\DecValTok{7}\NormalTok{): }\DecValTok{0}\NormalTok{\}}
\NormalTok{b }\OperatorTok{=}\NormalTok{ \{(}\StringTok{"проверить заказ"}\NormalTok{, }\DecValTok{7}\NormalTok{): }\DecValTok{1}\NormalTok{, (}\StringTok{"изменить адрес"}\NormalTok{, }\DecValTok{7}\NormalTok{): }\DecValTok{0}\NormalTok{,}
\NormalTok{     (}\StringTok{"проверить доставку"}\NormalTok{, }\DecValTok{7}\NormalTok{): }\DecValTok{1}\NormalTok{, (}\StringTok{"изменить время"}\NormalTok{, }\DecValTok{7}\NormalTok{): }\DecValTok{1}\NormalTok{\}}

\ControlFlowTok{if}\NormalTok{ a.keys() }\OperatorTok{!=}\NormalTok{ b.keys():}
    \ControlFlowTok{raise} \PreprocessorTok{ValueError}\NormalTok{(}\StringTok{"Обе конфигурации должны охватывать одинаковые задачи и начальные значения"}\NormalTok{)}

\NormalTok{changes }\OperatorTok{=}\NormalTok{ \{key: b[key] }\OperatorTok{{-}}\NormalTok{ a[key] }\ControlFlowTok{for}\NormalTok{ key }\KeywordTok{in}\NormalTok{ a\}}
\NormalTok{gains }\OperatorTok{=}\NormalTok{ [task }\ControlFlowTok{for}\NormalTok{ (task, seed), diff }\KeywordTok{in}\NormalTok{ changes.items() }\ControlFlowTok{if}\NormalTok{ diff }\OperatorTok{==} \DecValTok{1}\NormalTok{]}
\NormalTok{losses }\OperatorTok{=}\NormalTok{ [task }\ControlFlowTok{for}\NormalTok{ (task, seed), diff }\KeywordTok{in}\NormalTok{ changes.items() }\ControlFlowTok{if}\NormalTok{ diff }\OperatorTok{==} \OperatorTok{{-}}\DecValTok{1}\NormalTok{]}
\BuiltInTok{print}\NormalTok{(gains, losses)}
\BuiltInTok{print}\NormalTok{(}\SpecialStringTok{f"}\SpecialCharTok{\{}\BuiltInTok{sum}\NormalTok{(changes.values()) }\OperatorTok{/} \BuiltInTok{len}\NormalTok{(changes)}\SpecialCharTok{:+.0\%\}}\SpecialStringTok{"}\NormalTok{)}
\end{Highlighting}
\end{Shaded}

После запуска видно, что B впервые успешно выполнила задачи «проверить
доставку» и «изменить время», но показала регрессию в задаче «изменить
адрес». Чистый прирост составил одну задачу: общая доля успехов
увеличилась с 50\% до 75\%, то есть на 25 процентных пунктов. Если
смотреть только на итоговый результат, регрессия в задаче «изменить
адрес» останется незамеченной.

При одновременном отборе множества моделей, промптов и параметров
необходимо также учитывать \textbf{множественные сравнения}. Можно
заранее определить основной показатель и сравниваемые варианты,
скорректировать порог значимости либо повторно проверить победившую
конфигурацию на другом наборе задач, не участвовавшем в отборе. При
окончательном решении следует одновременно учитывать величину различия,
статистическую неопределённость и ценность для бизнеса: достаточно ли
снизились затраты, не стало ли больше критических сбоев и стабильно ли
воспроизводится результат на целевом распределении задач.

\hypertarget{ux43dux430ux431ux43bux44eux434ux430ux435ux43cux43eux441ux442ux44c-ai-ux430ux433ux435ux43dux442ux430}{%
\section{Наблюдаемость
AI-агента}\label{ux43dux430ux431ux43bux44eux434ux430ux435ux43cux43eux441ux442ux44c-ai-ux430ux433ux435ux43dux442ux430}}

Для оценивания необходимо знать, при каких условиях агент совершил то
или иное действие и к какому результату оно привело.
\textbf{Наблюдаемость (Observability)} фиксирует эти сигналы работы
системы с помощью логов, метрик и трассировки, помогая команде
диагностировать проблемы, анализировать затраты и непрерывно пополнять
набор примеров для оценивания. На \crossreflink{fig:7-7}{рисунке 7-7}
это показано на примере задачи, в которой нужно узнать погоду и
отправить сообщение. Пользователь ждал 3,2 секунды; чтобы определить,
какой этап следует оптимизировать, нужно проследить выполнение этого
запроса и найти шаг, занявший больше всего времени.

\begin{figure}
\hypertarget{fig:7-7}{%
\centering
\includegraphics{images/revised/fig7-7.pdf}
\caption{Рисунок 7-7. Стек технологий наблюдаемости}\label{fig:7-7}
}
\end{figure}

Путь выполнения агента меняется в зависимости от вывода модели,
результатов работы инструментов и состояния среды. Во время выполнения
следует сохранять текущий ввод, вызовы инструментов и их результаты,
переходы между состояниями и окончательный вывод. Предоставляемые
моделью пояснения помогают понять причины выбора действий; при
диагностике эти пояснения сопоставляют с фактическими действиями,
результатами инструментов и состоянием среды.

Одной задаче обычно соответствует одна \textbf{трасса (Trace)}, в
которой каждый вызов модели, вызов инструмента и операция извлечения
данных представлены отдельным \textbf{span}. В span записываются
название операции, время её начала и окончания, статус и ошибки; к нему
также можно прикрепить версию модели, расход токенов и обработанные
входные и выходные данные. «Основной цикл агента» может быть
родительским span, содержащим вызовы модели и выполнение инструментов;
связи с параллельными подзадачами и межсервисными обратными вызовами
также можно обозначать ссылками. Благодаря этому команда может как
просматривать сведения об отдельных шагах, так и восстанавливать
временную шкалу задачи.

\textbf{OpenTelemetry} предоставляет универсальные механизмы сбора и
передачи телеметрии, а спецификации вроде \textbf{OpenInference}
определяют для приложений на базе LLM поля вызовов модели, промптов,
извлечения данных и расхода ресурсов. Единообразные поля и
идентификаторы трассировки упрощают подключение средств сбора данных к
различным аналитическим бэкендам. Такие платформы, как LangSmith,
Langfuse и Arize Phoenix, могут обеспечивать просмотр трассировок,
управление наборами данных и процессы оценивания.

Сбор данных также требует инженерного проектирования. Асинхронная
буферизация и пакетная отправка позволяют снизить влияние на основной
процесс; кроме того, следует задать ёмкость очереди и правила
семплирования, а также вести статистику отброшенных данных, чтобы
контролировать накладные расходы на сбор. Для входных и выходных данных
выбирают, какие поля записывать, а чувствительное содержимое сохраняют
только после обезличивания или с применением контроля доступа. Состояние
среды, необходимое для воспроизведения, можно сохранять в виде версий,
ссылок на снимки или параметров, позволяющих восстановить это состояние.

Данные трассировки позволяют ответить на несколько конкретных вопросов:
какие задачи многократно выполняют извлечение, но так и не получают
необходимую информацию; какие инструменты часто завершаются с ошибкой;
какие вызовы занимают основную часть времени ожидания; какие задачи
расходуют значительно больше ресурсов, чем аналогичные. Если прикреплять
к трассировкам версии промптов, конфигурации функций и группы
экспериментов, можно также сравнивать поведение системы до и после
изменений.

Возврат данных в цикл образует процесс непрерывного совершенствования:
отбираются неудачные и подозрительные случаи, обрабатываются
персональные и чувствительные поля, специалисты вручную подтверждают
требования задачи и причину сбоя, после чего случаи преобразуются в
сквозные или префиксные регрессионные тесты. Сохранение типа задачи и
источника выборки позволяет одновременно поддерживать набор оценивания,
близкий к распределению в продакшене, и специализированные наборы для
границ с высоким уровнем риска.

\hypertarget{ux43eux442-ux43eux442ux447ux451ux442ux430-ux431ux435ux43dux447ux43cux430ux440ux43aux430-ux43a-ux443ux43bux443ux447ux448ux435ux43dux438ux44e-ux441ux438ux441ux442ux435ux43cux44b}{%
\section{От отчёта бенчмарка к улучшению
системы}\label{ux43eux442-ux43eux442ux447ux451ux442ux430-ux431ux435ux43dux447ux43cux430ux440ux43aux430-ux43a-ux443ux43bux443ux447ux448ux435ux43dux438ux44e-ux441ux438ux441ux442ux435ux43cux44b}}

AI-агент неоднократно застревает на странице настроек Wi-Fi в
AndroidWorld. Он не знает, куда нажать дальше, или вообще не видит
нужный элемент управления? На \crossreflink{fig:7-8}{рисунке 7-8}
показано, как команда по результатам одного эксперимента решает, что
проверять следующим.

\begin{figure}
\hypertarget{fig:7-8}{%
\centering
\includegraphics{images/revised/fig7-8.pdf}
\caption{Рисунок 7-8. Замкнутый цикл от бенчмарка к
улучшению}\label{fig:7-8}
}
\end{figure}

\textbf{Если результаты ухудшились, сначала проверьте систему оценки и
лишь затем изменяйте агента.} В частности, проверьте, успешно ли
инициализирована задача, достаточно ли ресурсов, правильно ли работает
валидатор и не изменились ли условия тестирования. Завершение процесса
из-за ограничения ресурсов и ошибочная оценка правильного ответа как
неверного могут одинаково проявляться в виде снижения результата. Только
отделив проблемы измерительного контура от проблем поведения модели,
можно определить, что именно следует исправлять дальше.

\hypertarget{ux43aux430ux43a-ux447ux438ux442ux430ux442ux44c-ux43eux442ux447ux451ux442-ux43fux43eux438ux441ux43a-ux441ux43aux43eux43dux446ux435ux43dux442ux440ux438ux440ux43eux432ux430ux43dux43dux44bux445-ux43fux430ux442ux442ux435ux440ux43dux43eux432-ux43dux435ux443ux434ux430ux447}{%
\subsection{Как читать отчёт: поиск сконцентрированных паттернов
неудач}\label{ux43aux430ux43a-ux447ux438ux442ux430ux442ux44c-ux43eux442ux447ux451ux442-ux43fux43eux438ux441ux43a-ux441ux43aux43eux43dux446ux435ux43dux442ux440ux438ux440ux43eux432ux430ux43dux43dux44bux445-ux43fux430ux442ux442ux435ux440ux43dux43eux432-ux43dux435ux443ux434ux430ux447}}

Исторический отчёт охватывает 116 задач, каждая из которых запускалась
один раз, а общая доля успешного выполнения составила около 88\%. При
этом в четырёх задачах настройки Wi-Fi сконцентрировались три неудачи, а
в траекториях наблюдались повторные переходы и неспособность подтвердить
состояние. На этом основании можно сформулировать два первоначальных
объяснения: модели недостаёт навигационной информации либо текущее
наблюдение недостаточно полно отображает интерфейс.

Команда сузила последующие эксперименты до четырёх задач Wi-Fi в
эмуляторе с API 35, чтобы с небольшими затратами проверить оба
объяснения. Общий показатель позволяет оценить результат в целом, а
результаты и траектории отдельных задач помогают найти кластеры неудач,
исправление которых стоит сделать приоритетным.

\hypertarget{ux43eux442-ux434ux430ux43dux43dux44bux445-ux43a-ux433ux438ux43fux43eux442ux435ux437ux430ux43c-ux43fux43eux441ux43bux435ux434ux43eux432ux430ux442ux435ux43bux44cux43dux43eux435-ux441ux443ux436ux435ux43dux438ux435-ux43eux431ux43bux430ux441ux442ux438-ux43fux440ux43eux431ux43bux435ux43cux44b}{%
\subsection{От данных к гипотезам: последовательное сужение области
проблемы}\label{ux43eux442-ux434ux430ux43dux43dux44bux445-ux43a-ux433ux438ux43fux43eux442ux435ux437ux430ux43c-ux43fux43eux441ux43bux435ux434ux43eux432ux430ux442ux435ux43bux44cux43dux43eux435-ux441ux443ux436ux435ux43dux438ux435-ux43eux431ux43bux430ux441ux442ux438-ux43fux440ux43eux431ux43bux435ux43cux44b}}

На первой итерации H1 добавляются подсказки по навигации в настройках
Wi-Fi и требование проверить итоговое состояние. Доля успеха осталась на
уровне 25\%, следовательно, для этих четырёх задач данный набор
подсказок не увеличил количество успешных выполнений.

На второй итерации H5 несовместимый с API 35 канал информации о
специальных возможностях заменяется на дерево элементов UIAutomator. Все
четыре задачи выполняются успешно, но при этом заметно увеличивается
количество токенов. На третьей итерации H5C из полного дерева элементов
удаляются избыточные контейнерные узлы, которые невидимы, не содержат
текста и не допускают взаимодействия, чтобы проверить, какую информацию
можно исключить.

Во всех сравниваемых итерациях неизменными оставались модель, параметры
задач, случайное начальное значение, ограничение количества шагов и
условия эмулятора, а порядок запусков чередовался.

\hypertarget{ux43eux442-ux440ux435ux437ux443ux43bux44cux442ux430ux442ux43eux432-ux43a-ux440ux435ux448ux435ux43dux438ux44e-ux441ux43eux432ux43cux435ux441ux442ux43dux43eux435-ux441ux440ux430ux432ux43dux435ux43dux438ux435-ux43aux430ux447ux435ux441ux442ux432ux430-ux438-ux441ux442ux43eux438ux43cux43eux441ux442ux438}{%
\subsection{От результатов к решению: совместное сравнение качества и
стоимости}\label{ux43eux442-ux440ux435ux437ux443ux43bux44cux442ux430ux442ux43eux432-ux43a-ux440ux435ux448ux435ux43dux438ux44e-ux441ux43eux432ux43cux435ux441ux442ux43dux43eux435-ux441ux440ux430ux432ux43dux435ux43dux438ux435-ux43aux430ux447ux435ux441ux442ux432ux430-ux438-ux441ux442ux43eux438ux43cux43eux441ux442ux438}}

В таблице 7-5 сведены результаты четырёх задач на каждой итерации. В H1
и H5 экспериментальный вариант сравнивается с исходной схемой
наблюдения, а в H5C контрольным вариантом служит полное дерево элементов
UIAutomator.

Таблица 7-5. Три эксперимента на подмножестве задач AndroidWorld с Wi-Fi

\begin{longtable}[]{@{}
  >{\raggedright\arraybackslash}p{(\columnwidth - 8\tabcolsep) * \real{0.1400}}
  >{\raggedright\arraybackslash}p{(\columnwidth - 8\tabcolsep) * \real{0.2090}}
  >{\raggedleft\arraybackslash}p{(\columnwidth - 8\tabcolsep) * \real{0.1866}}
  >{\raggedleft\arraybackslash}p{(\columnwidth - 8\tabcolsep) * \real{0.2778}}
  >{\raggedright\arraybackslash}p{(\columnwidth - 8\tabcolsep) * \real{0.1867}}@{}}
\toprule\noalign{}
\begin{minipage}[b]{\linewidth}\raggedright
\hspace{0pt}Эксперимент
\end{minipage} & \begin{minipage}[b]{\linewidth}\raggedright
\hspace{0pt}Изменение
\end{minipage} & \begin{minipage}[b]{\linewidth}\raggedleft
\hspace{0pt}Доля успеха контрольной группы → экспериментальной группы
\end{minipage} & \begin{minipage}[b]{\linewidth}\raggedleft
\hspace{0pt}Токены экспериментальной/\allowbreak{}контрольной группы
\end{minipage} & \begin{minipage}[b]{\linewidth}\raggedright
\hspace{0pt}Последующее решение
\end{minipage} \\
\midrule\noalign{}
\endhead
\bottomrule\noalign{}
\endlastfoot
\hspace{0pt}H1 & \hspace{0pt}Добавлены подсказки по навигации &
\hspace{0pt}25\%→25\% & \hspace{0pt}0.47× & \hspace{0pt}Перейти к
проверке канала наблюдения \\
\hspace{0pt}H5 & \hspace{0pt}Использовано дерево элементов UIAutomator &
\hspace{0pt}25\%→100\% & \hspace{0pt}2.498× & \hspace{0pt}Сохранить
эффективное наблюдение и продолжить снижать затраты на входные данные \\
\hspace{0pt}H5C & \hspace{0pt}Сокращено дерево элементов UIAutomator &
\hspace{0pt}100\%→100\% & \hspace{0pt}0.506× & \hspace{0pt}Расширить
набор задач для регрессионного тестирования \\
\end{longtable}

Почему сначала изменили наблюдение? При неудачах программа уже открывала
панель Internet, но исходный интерфейс часто возвращал только строку
состояния. Как бы внимательно модель ни читала подсказки по навигации,
она не могла найти элементы управления на странице настроек. После
перехода на полное дерево элементов все четыре задачи были выполнены
успешно, однако среднее количество токенов достигло 2,498-кратного
контрольного уровня, превысив установленный для этой итерации верхний
предел, равный 1,5-кратному контрольному уровню. Поэтому на следующей
итерации сохранили текст, состояния и доступные действия элементов
управления, попробовав удалить пустые контейнеры. В этом парном
сравнении все четыре задачи были пройдены как до, так и после
сокращения, а среднее количество токенов снизилось со 139 439,5 до 70
557,5. Значение 0.506 в таблице представляет собой отношение этих двух
чисел.

\hypertarget{ux43dux435ux43fux440ux435ux440ux44bux432ux43dux44bux435-ux438ux442ux435ux440ux430ux446ux438ux438-ux43fux435ux440ux435ux43dux43eux441-ux43bux43eux43aux430ux43bux44cux43dux43eux433ux43e-ux443ux43bux443ux447ux448ux435ux43dux438ux44f-ux432-ux43fux43eux43bux43dux43eux435-ux440ux435ux433ux440ux435ux441ux441ux438ux43eux43dux43dux43eux435-ux442ux435ux441ux442ux438ux440ux43eux432ux430ux43dux438ux435}{%
\subsection{Непрерывные итерации: перенос локального улучшения в полное
регрессионное
тестирование}\label{ux43dux435ux43fux440ux435ux440ux44bux432ux43dux44bux435-ux438ux442ux435ux440ux430ux446ux438ux438-ux43fux435ux440ux435ux43dux43eux441-ux43bux43eux43aux430ux43bux44cux43dux43eux433ux43e-ux443ux43bux443ux447ux448ux435ux43dux438ux44f-ux432-ux43fux43eux43bux43dux43eux435-ux440ux435ux433ux440ux435ux441ux441ux438ux43eux43dux43dux43eux435-ux442ux435ux441ux442ux438ux440ux43eux432ux430ux43dux438ux435}}

Последний шаг на \crossreflink{fig:7-8}{рисунке 7-8} --- вернуться к
стандартной среде Pixel 6 с API 33 и выполнить по пять запусков каждой
из 116 задач. Команда задала три цели для полного сравнения: сохранить
долю успеха, удержать количество токенов в пределах 75\% от контрольного
решения и время выполнения --- в пределах 1,5-кратного контрольного
значения. Чтобы проверить все три цели, сокращённое и полное деревья
необходимо запускать с одной моделью и на одном наборе задач.

Позднее в сопутствующем проекте завершили 580 тестов сокращённого дерева
с локальной моделью Qwen2.5-7B. В 77 запусках состояние приложения
получило максимальную оценку. Однако по более строгому критерию успеха
агент также должен был явно сообщить о завершении задачи, и оба условия
были выполнены лишь в 26 запусках --- приблизительно в 4,48\% случаев.
Иногда модель уже выполняла задачу, но не успевала своевременно
завершить работу.

\begin{quote}
\textbf{Эксперимент 7-13 ★★★: оценка и улучшение AndroidWorld}

Начните с исторического отчёта и результатов парных сравнений H1, H5 и
H5C, а затем выполните следующий процесс.

Шаг 1. Сопоставьте результаты отдельных задач с метками способностей,
чтобы найти кластер неудач. Затем изучите траектории и разделите
проблемы среды, наблюдения, принятия решений, выполнения и приёмки.

Шаг 2. Проследите путь от внешних симптомов к условиям процесса и
системным механизмам, сформулируйте гипотезу улучшения и укажите
ожидаемые изменения, измеряемые показатели и способности, которые могут
деградировать.

Шаг 3. На каждой итерации сохраняйте остальные условия неизменными,
сравнивайте одно конкретное изменение и одновременно регистрируйте долю
успеха, количество токенов, задержку и типы неудач.

Шаг 4. По результатам решите, следует ли сохранить, скорректировать или
отклонить изменение. Если качество улучшилось ценой дополнительных
затрат, определите, для каких задач его стоит включать.

Шаг 5. Перенесите локальное улучшение в полное регрессионное
тестирование в стандартной среде. Проверьте целевые показатели качества
и стоимости по схеме из 116 задач и пяти случайных начальных значений, а
затем определите область публикации и условия отката.
\end{quote}

\hypertarget{ux43eux442-ux432ux43dux435ux448ux43dux435ux439-ux43eux446ux435ux43dux43aux438-ux43a-ux432ux43dux443ux442ux440ux435ux43dux43dux435ux439-ux438ux43dux444ux440ux430ux441ux442ux440ux443ux43aux442ux443ux440ux430-ux43eux446ux435ux43dux43aux438-ai-ux430ux433ux435ux43dux442ux43eux432-ux43fux440ux43eux43cux44bux448ux43bux435ux43dux43dux43eux433ux43e-ux443ux440ux43eux432ux43dux44f}{%
\section{От внешней оценки к внутренней: инфраструктура оценки
AI-агентов промышленного
уровня}\label{ux43eux442-ux432ux43dux435ux448ux43dux435ux439-ux43eux446ux435ux43dux43aux438-ux43a-ux432ux43dux443ux442ux440ux435ux43dux43dux435ux439-ux438ux43dux444ux440ux430ux441ux442ux440ux443ux43aux442ux443ux440ux430-ux43eux446ux435ux43dux43aux438-ai-ux430ux433ux435ux43dux442ux43eux432-ux43fux440ux43eux43cux44bux448ux43bux435ux43dux43dux43eux433ux43e-ux443ux440ux43eux432ux43dux44f}}

Выше были представлены среда оценки, наборы данных и анализ отчётов. Для
включения этих методов в разработку продукта необходимо дополнительно
обеспечить управление конфигурацией, разделение на экспериментальные
группы, отслеживание версий и сбор данных. Тогда каждое изменение
модели, промпта или Harness можно будет связать с конкретными условиями
эксперимента и результатами выполнения.

\hypertarget{ux438ux43dux444ux440ux430ux441ux442ux440ux443ux43aux442ux443ux440ux430-ux430ux431ux43bux44fux446ux438ux438-ux441ux440ux430ux432ux43dux435ux43dux438ux435-ux432ux43aux43bux430ux434ux430-ux444ux443ux43dux43aux446ux438ux439}{%
\subsection{Инфраструктура абляции: сравнение вклада
функций}\label{ux438ux43dux444ux440ux430ux441ux442ux440ux443ux43aux442ux443ux440ux430-ux430ux431ux43bux44fux446ux438ux438-ux441ux440ux430ux432ux43dux435ux43dux438ux435-ux432ux43aux43bux430ux434ux430-ux444ux443ux43dux43aux446ux438ux439}}

\textbf{Абляционный эксперимент (Ablation Study)} позволяет наблюдать за
изменением поведения системы при отключении или замене отдельного
компонента. Для агента можно по отдельности управлять бюджетом
рассуждений, сжатием контекста, автоматической памятью и фоновыми
задачами либо отключить целую группу функций, создав более простой
инженерный базовый вариант.

Абляция одной функции помогает определить её предельный вклад, а
комбинированная абляция --- выявить зависимости. Например, сжатие
истории может снизить стоимость входных данных, но также повлиять на
долю попаданий в кэш; польза автоматической памяти, в свою очередь,
может меняться в зависимости от бюджета контекста.

Конфигурацию необходимо определить до того, как компонент её прочитает.
Если модуль при загрузке сохраняет конфигурацию в виде константы, её
нужно внедрить до загрузки; если компонент поддерживает конфигурацию
времени выполнения, её снимок можно зафиксировать при создании задачи
или сеанса.

Повторное проведение абляции перед обновлением модели и выпуском важных
версий позволяет выявить «долг функций»: изначально функция
компенсировала недостатки модели, но новая модель уже справляется без
неё, тогда как сама функция по-прежнему увеличивает количество вызовов,
задержку и объём работ по сопровождению.

\hypertarget{ab-ux442ux435ux441ux442ux438ux440ux43eux432ux430ux43dux438ux435-ux440ux430ux437ux434ux435ux43bux435ux43dux438ux435-ux43cux435ux445ux430ux43dux438ux437ux43cux430-ux438ux437ux43cux435ux43dux435ux43dux438ux44f-ux438-ux431ux438ux437ux43dux435ux441-ux446ux435ux43bux438}{%
\subsection{A/B-тестирование: разделение механизма изменения и
бизнес-цели}\label{ab-ux442ux435ux441ux442ux438ux440ux43eux432ux430ux43dux438ux435-ux440ux430ux437ux434ux435ux43bux435ux43dux438ux435-ux43cux435ux445ux430ux43dux438ux437ux43cux430-ux438ux437ux43cux435ux43dux435ux43dux438ux44f-ux438-ux431ux438ux437ux43dux435ux441-ux446ux435ux43bux438}}

При A/B-тестировании подходящие пользователи, сеансы или задачи
распределяются между разными версиями, после чего сравниваются
фактические результаты выполнения. Сначала нужно определить единицу
случайного распределения по группам, а затем сохранять конфигурацию
неизменной в пределах этой единицы, чтобы один и тот же пользователь или
последовательные задачи не переключались многократно между условиями
эксперимента.

\textbf{Создавайте несколько вариантов.} Помимо исходного и нового
вариантов можно предусмотреть изменения разной степени выраженности.
Например, для длины плана можно создать контрольную группу и три группы
с последовательно ужесточающимися ограничениями, чтобы выяснить, как
строгость ограничений влияет на поведение.

\textbf{Разделяйте метрики механизма и целевые метрики.} Длина файла
плана показывает, сработало ли «сокращение плана», а стоимость на уровне
сеанса измеряет конечную пользу. Более короткий план может уменьшить
объём вывода в текущем ходе, но также привести к дополнительным
исправлениям и проверкам. В эксперименте следует одновременно
фиксировать длину плана, последующие переделки, общее количество токенов
и долю успешно выполненных задач, оценивая ценность по результату задачи
в целом.

\textbf{Устанавливайте защитные метрики.} Даже при снижении стоимости
доля успешно выполненных задач, удовлетворённость пользователей и
частота критических ошибок должны оставаться в согласованных пределах.
Следует заранее зафиксировать допустимые изменения и условия остановки,
а при достижении границы прекратить расширение эксперимента или
выполнить откат.

\textbf{Что пользователь сделал бы сам?} Вспомним эксперимент с
покупками из первой главы. Из десяти рекомендованных товаров
покупательнице сначала не понравился ни один. Затем она самостоятельно
просмотрела товары на Taobao и в итоге купила именно один из
рекомендованных. Если смотреть только на заказ, рекомендация покажется
точной; если учитывать лишь первоначальную реакцию покупательницы,
покажется, что все десять товаров были выбраны неверно. Чтобы понять,
чем помогла рекомендация, можно попросить её вспомнить: «Что вы увидели
позже, что заставило вас передумать?»

Согласование счетов --- другой случай. Некоторые люди готовы
воспользоваться снижением платы, но не хотят ради этого час
разговаривать по телефону. Если формировать выборку только из тех, кто
уже звонил, в неё не попадут люди, готовые подать заявку лишь при
наличии услуги оформления от их имени. Пользователей можно распределить
по группам до того, как они решат, подавать ли заявку: одна группа
сможет воспользоваться такой услугой, а другая будет действовать прежним
способом. В обеих группах знаменателем служит общее количество
пользователей, соответствующих условиям; отдельно подсчитывается,
сколько из них подали заявку и у скольких в итоге подтвердилось снижение
платы. Разница в долях поданных заявок поможет оценить, скольким людям
услуга позволила сделать первый шаг; разница в долях фактического
снижения платы покажет, скольким людям этот шаг помог его получить. Что
касается создания слайдов, одни готовят презентацию, чтобы выступить
перед другими, а другие --- чтобы разобраться в вопросе, систематизируя
материал. Перед проектированием оценки нужно выяснить цель; человека из
второй группы можно попросить закрыть слайды и своими словами объяснить
то, чего он раньше не понимал.\footnote{Автор обсуждал различные мотивы
  покупок, согласования счетов и создания слайдов в докладе «The Road to
  JARVIS» в Порту. На основе этого в данном разделе разработаны примеры
  оценки с распределением по группам и последующими опросами.}

\hypertarget{ux434ux432ux443ux445ux443ux440ux43eux432ux43dux435ux432ux430ux44f-ux441ux438ux441ux442ux435ux43cux430-ux444ux43bux430ux433ux43eux432-ux444ux443ux43dux43aux446ux438ux439}{%
\subsection{Двухуровневая система флагов
функций}\label{ux434ux432ux443ux445ux443ux440ux43eux432ux43dux435ux432ux430ux44f-ux441ux438ux441ux442ux435ux43cux430-ux444ux43bux430ux433ux43eux432-ux444ux443ux43dux43aux446ux438ux439}}

\textbf{Автор рекомендует закладывать флаги функций в архитектуру с
самого начала разработки продукта.} Они одновременно служат проведению
экспериментов, поэтапному выпуску и устранению сбоев, позволяя команде
независимо сравнивать основные функции и своевременно восстанавливать
прежнюю конфигурацию при обнаружении регрессии. Флаги времени сборки и
времени выполнения выполняют разные задачи.

\textbf{Флаги времени сборки} определяют, какие компоненты войдут в
артефакт. Модули с функциями исключительно для внутреннего использования
можно исключить из внешней сборки, одновременно проверив, что связанные
зависимости и ресурсы также удалены.

\textbf{Флаги времени выполнения} определяют, какие функции будут
использоваться для текущего пользователя или задачи. Конфигурация может
поступать с сервера, кэшироваться на клиенте, а распределение по группам
--- выполняться экспериментальной платформой, например GrowthBook.
Представим, что флаг управляет отображением прогресса задачи, а в момент
открытия помощника пользователем сервис распределения по группам не
отвечает. Автор предпочёл бы сначала запустить помощника с последней
локально кэшированной конфигурацией, срок действия которой ещё не истёк,
а затем обновить её в фоновом режиме; при первом запуске или после
истечения срока действия кэша следует использовать заранее заданное
значение по умолчанию. Тогда пользователь сможет начать работу сразу
после открытия помощника.

\emph{Уточнение русского издания: Срок действия кэша и актуальность
запрета --- разные вещи. Если флаг управляет доступом, сохранённое
разрешение может уже быть отозвано. Это рекомендация русского издания, а
не описание реализации OpenClaw.
\href{https://github.com/bojieli/ai-agent-book/blob/c3352738f4b6fe42fe34e3cf6a79bcb424a133b8/book/chapter7.md\#L805}{Источник}.}

\hypertarget{ux43eux446ux435ux43dux43aux430-ux447ux443ux432ux441ux442ux432ux438ux442ux435ux43bux44cux43dux43eux441ux442ux438-ux43fux440ux43eux43cux43fux442ux43eux432}{%
\subsection{Оценка чувствительности
промптов}\label{ux43eux446ux435ux43dux43aux430-ux447ux443ux432ux441ux442ux432ux438ux442ux435ux43bux44cux43dux43eux441ux442ux438-ux43fux440ux43eux43cux43fux442ux43eux432}}

Системный промпт участвует в определении поведения агента, поэтому его
следует версионировать вместе с кодом. Помимо самого шаблона необходимо
сохранять итоговый текст после динамического рендеринга: список
инструментов, правила задачи, сведения об окружении и пользовательская
конфигурация могут изменять фактические входные данные.

Практичный подход --- предоставить отдельную точку входа для рендеринга:
она принимает версию кода, конфигурацию и снимок окружения, а возвращает
итоговый системный промпт. Если динамические поля, например время,
передавать как явные входные данные, один и тот же набор входных данных
будет давать одинаковый результат. Сохраняя версионируемые снимки, можно
установить, какое изменение и как именно изменило содержимое промпта.

\textbf{Автор рекомендует управлять изменениями промптов так же, как
изменениями кода: после каждого изменения запускать соответствующие
регрессионные проверки.} Для шаблонов с широкой областью влияния можно
дополнительно менять порядок абзацев или формулировки, измерять
чувствительность системы к форме выражения и выбирать устойчиво
работающий вариант.

\hypertarget{ux432ux43aux43bux44eux447ux435ux43dux438ux435-ux437ux430ux449ux438ux442ux44b-ux434ux430ux43dux43dux44bux445-ux432-ux430ux43dux430ux43bux438ux442ux438ux447ux435ux441ux43aux438ux435-ux438ux43dux442ux435ux440ux444ux435ux439ux441ux44b}{%
\subsection{Включение защиты данных в аналитические
интерфейсы}\label{ux432ux43aux43bux44eux447ux435ux43dux438ux435-ux437ux430ux449ux438ux442ux44b-ux434ux430ux43dux43dux44bux445-ux432-ux430ux43dux430ux43bux438ux442ux438ux447ux435ux441ux43aux438ux435-ux438ux43dux442ux435ux440ux444ux435ux439ux441ux44b}}

\textbf{Защиту данных следует проектировать одновременно со сбором
данных для оценки.} Сначала нужно определить, на какие вопросы предстоит
ответить, и лишь затем решать, какие поля собирать: это позволит
постоянно получать пригодные для диагностики материалы. Если заняться
чувствительными данными постфактум, уже собранные траектории зачастую
будет трудно безопасно использовать повторно, а всю цепочку сбора
придётся проектировать заново. Структурированных полей, таких как
категория задачи, продолжительность, объём использования и состояние,
часто уже достаточно для ответа на основные операционные вопросы;
подробное содержимое следует обезличивать и защищать контролем доступа в
соответствии с целями диагностики.

Система типов может помочь ограничить потоки данных. Например,
аналитический интерфейс может принимать только значения специального
типа-обёртки для «полей, разрешённых к сбору», а логику проверки или
обезличивания можно сосредоточить в точке их создания, запретив
произвольному тексту напрямую поступать в аналитическую систему.
Проверка типов ограничивает способы вызова, однако содержимое полей всё
равно необходимо проверять по соответствующим правилам обработки.

Все эти возможности вместе обеспечивают непрерывную оценку: абляция
проверяет вклад компонентов, A/B-тестирование сравнивает влияние на
продукт, флаги функций фиксируют условия эксперимента, снимки промптов
отслеживают входные данные, определяющие поведение, а сбор данных
предоставляет материал для выводов.

\hypertarget{ux441ux440ux435ux434ux44b-ux441ux438ux43cux443ux43bux44fux446ux438ux438-ux43cux43eux441ux442-ux43eux442-ux43eux446ux435ux43dux43aux438-ux43a-ux43fux43eux441ux442ux43eux431ux443ux447ux435ux43dux438ux44e}{%
\section{Среды симуляции: мост от оценки к
постобучению}\label{ux441ux440ux435ux434ux44b-ux441ux438ux43cux443ux43bux44fux446ux438ux438-ux43cux43eux441ux442-ux43eux442-ux43eux446ux435ux43dux43aux438-ux43a-ux43fux43eux441ux442ux43eux431ux443ux447ux435ux43dux438ux44e}}

Среда оценки предоставляет задачи, интерфейсы взаимодействия и способы
приёмки. При дальнейшем использовании для обучения тот же набор
механизмов должен также поддерживать многократное сэмплирование,
формирование обратной связи и накопление траекторий. \textbf{Среда
симуляции} позволяет AI-агенту взаимодействовать в условиях, допускающих
сброс, посредством управляемой модели среды или системы выполнения; её
можно использовать как для оценки, так и для обучения.

Валидаторы и критерии оценивания со стороны оценки могут участвовать в
формировании обучающего вознаграждения. Программные валидаторы проверяют
тесты, состояние и ограничения, поэтому подходят для предоставления
\textbf{проверяемого вознаграждения}; модельное оценивание на основе
критериев предоставляет соответствующий оценочный сигнал, но требует
калибровки смещений и защиты от манипулирования вознаграждением.
\textbf{Обучение с подкреплением с проверяемыми вознаграждениями (RLVR,
Reinforcement Learning with Verifiable Rewards)} означает применение
проверяемых вознаграждений в обучении с подкреплением; конкретные методы
обучения будут рассмотрены в \crossreflink{chap:8}{главе 8}.

Обучение дополнительно выдвигает на первый план три инженерных
требования. Первое --- надёжный \textbf{reset}: для каждого эпизода, то
есть полного цикла взаимодействия от начального состояния до завершения,
среду следует воссоздавать с заданными параметрами и устранять побочные
эффекты предыдущего запуска. Второе --- управляемая рандомизация:
обучение проводится на разных экземплярах задач, при этом для оценки
сохраняется отдельное распределение. Наконец, необходима достаточная
пропускная способность: по мере роста масштаба взаимодействий скорость
обучения совместно определяют параллелизм среды, планирование моделей и
накладные расходы отдельного экземпляра. Обратную связь можно
предоставлять на отдельных шагах либо вычислять целиком после завершения
эпизода --- это зависит от задачи и алгоритма обучения.

При выборе среды оценки прежде всего следует спросить, ошибки какого
типа требуется выявить. В задаче возврата средств слева на
\crossreflink{fig:7-9}{рисунке 7-9} для обнаружения повторного возврата
достаточно смоделированных заказов; если же ошибка возникает при нажатии
кнопки возврата, агенту потребуется работать с интерфейсом приложения.
Задача захвата справа зависит от деталей другого рода: как переместится
чашка после соприкосновения с захватом? В симуляции положение чашки
можно многократно изменять, а затем на реальном роботе проверить,
переносятся ли на него результаты таких упражнений.

\begin{figure}
\hypertarget{fig:7-9}{%
\centering
\includegraphics{images/revised/fig7-9.pdf}
\caption{Рисунок 7-9. Среды оценки для цифровых бизнес-процессов и
роботизированных манипуляций}\label{fig:7-9}
}
\end{figure}

\textbf{Цифровые среды} позволяют помещать сервисы инструментов,
файловую систему и состояние приложений в управляемую среду выполнения.
В примере инструментальной среды AWorld для GAIA предоставляются 26
MCP-серверов и 126 функций инструментов, а также используется
распределённое планирование. Время выполнения в этом примере сократилось
с 7695 секунд при последовательном выполнении до 525 секунд, что
соответствует ускорению примерно в 14,6 раза.

\textbf{Воплощённые среды} должны дополнительно моделировать физическое
состояние и процесс управления. RoboTwin2 предоставляет основанные на
физическом движке задачи по манипулированию двумя роботизированными
руками и позволяет рандомизировать положение, ориентацию и внешний вид
объектов, а также другие параметры среды. Модель формирует действия на
основе изображений с камер и состояния робота, а после их выполнения
среда возвращает новые наблюдения. Стратегия с \textbf{разбиением
действий на блоки (Action Chunking)} за один раз выдаёт несколько
последующих действий, которые затем выполняются согласно протоколу
управления с обновлением наблюдений; соответствующий механизм рассмотрен
в \crossreflink{chap:6}{главе 6}.

Снимки виртуальных машин в OSWorld и инициализация состояния приложений
в AndroidWorld также обеспечивают воспроизводимое взаимодействие. В
цифровых и воплощённых средах можно сочетать механизмы изолированного
выполнения из \crossreflink{chap:4}{главы 4} с независимой идентичностью
и передачей управления человеку из \crossreflink{chap:6}{главы 6}:
отдельно управлять вычислительными ресурсами, доступом к файлам,
внешними разрешениями и необходимыми этапами аутентификации, а также
выбирать для конкретной задачи общую файловую систему или независимые
экземпляры.

\begin{quote}
\textbf{Эксперимент 7-14 ★★: настройка среды воплощённого интеллекта
OpenVLA и RoboTwin2}

Изучите описание архитектуры и интерфейсов OpenVLA и SimpleVLA-RL, чтобы
понять, как соединяются визуальное кодирование, языковые условия и вывод
действий. Настройте RoboTwin2 и чётко определите используемую стратегию
и интерфейс робота: если наблюдения включают RGB-изображения с трёх
ракурсов и совокупное 14-мерное состояние суставов двух роботизированных
рук, а им соответствует 14-мерный вектор управления суставами,
поочерёдно проверьте порядок суставов, поля захватов и способ
нормализации.

Изучите рандомизацию объектов и пространственные ограничения задачи
\texttt{move\_can\_pot}, а затем настройте предварительно обученную или
адаптированную стратегию, соответствующую этим входным и выходным
данным. Зафиксируйте условия задачи и записывайте долю успешных
выполнений, время выполнения и типы сбоев. В сопутствующем эксперименте
для обоих вариантов используются одни и те же веса стратегии, а при
каждом предсказании формируются 25 действий. В одной группе после
выполнения только первого действия сразу получается новое наблюдение; в
другой перед получением нового наблюдения выполняется до 25 действий из
предсказанной последовательности. Таким образом, изменяется только
способ выполнения, а модель, задача и случайные начальные значения
остаются общими для обеих групп.

В каждой группе проведено по 128 тестов внутри распределения и 128
тестов вне распределения. Группа, выполнявшая только одно действие, не
достигла успеха ни разу; группа, выполнявшая длинные последовательности,
достигла успеха 26 раз --- по 13 раз в каждом из двух типов тестов. В
этой задаче длинные последовательности оказались эффективнее, однако
доля успешных выполнений составила лишь около 10\%. При просмотре
записей неудачных попыток сначала можно разделить их на два случая:
захват не был правильно наведён с самого начала или первые несколько
действий были правильными, но после смещения объекта система не смогла
своевременно скорректировать движение? В первом случае следует проверить
исходное наблюдение и предсказание действий; только второй случай
непосредственно связан с моментом получения нового наблюдения.
\end{quote}

На \crossreflink{fig:7-10}{рисунке 7-10} стратегия вычисляет действия по
изображениям с камер и состоянию робота, после чего адаптационный слой
передаёт их симулятору. Стратегия увидит последствия выполненного
действия лишь после получения следующего кадра. При замене стратегии
необходимо поочерёдно проверить количество камер, входные данные
состояния и выходную голову действий. Даже если программа выполняется
без ошибок, смысл действий уже может оказаться иным.

\begin{figure}
\hypertarget{fig:7-10}{%
\centering
\includegraphics{images/revised/fig7-10.pdf}
\caption{Рисунок 7-10. Среда воплощённого интеллекта OpenVLA и
RoboTwin2}\label{fig:7-10}
}
\end{figure}

\hypertarget{ux43aux43eux43cux43fux440ux43eux43cux438ux441ux441ux44b-ux442ux43eux447ux43dux43eux441ux442ux438-ux441ux438ux43cux443ux43bux44fux446ux438ux438-ux438-ux434ux43eux43cux435ux43dux43dux430ux44f-ux440ux430ux43dux434ux43eux43cux438ux437ux430ux446ux438ux44f}{%
\subsection{Компромиссы точности симуляции и доменная
рандомизация}\label{ux43aux43eux43cux43fux440ux43eux43cux438ux441ux441ux44b-ux442ux43eux447ux43dux43eux441ux442ux438-ux441ux438ux43cux443ux43bux44fux446ux438ux438-ux438-ux434ux43eux43cux435ux43dux43dux430ux44f-ux440ux430ux43dux434ux43eux43cux438ux437ux430ux446ux438ux44f}}

Точность симуляции следует выбирать с учётом целевой задачи. Для задачи
захвата необходимо правдоподобно моделировать столкновения, трение и
контакт с захватом, а для управления интерфейсом --- точно
воспроизводить состояния элементов управления и процесс отклика.
Повышение детализации обычно увеличивает и расход ресурсов, поэтому
сначала можно определить ключевые факторы, влияющие на стратегию, а
затем выбрать точность моделирования.

\textbf{Доменная рандомизация (Domain Randomization)} позволяет
стратегии столкнуться с более широким диапазоном условий за счёт
изменения физических параметров, визуального облика, положения объектов
и шума датчиков. Диапазон рандомизации должен охватывать ожидаемые
изменения при развёртывании, сохраняя при этом разрешимость задачи.
Например, изменение освещения и фона помогает проверить, не полагается
ли модель на случайные визуальные признаки, а изменение трения и массы
объектов --- адаптируются ли действия к изменениям динамики.

Разрыв между симуляцией и реальным выполнением существует и в цифровых
средах: например, могут различаться отрисовка интерфейса, задержка
сервисов и режимы отказа. В управляемую среду можно вносить задержки,
частичные сбои и изменения компоновки, а затем измерять качество
переноса посредством оценки в независимой целевой среде. Таким образом,
среда предоставляет воспроизводимые условия для оценки и разнообразные
взаимодействия для постобучения.

\hypertarget{ux438ux442ux43eux433ux438-ux433ux43bux430ux432ux44b-6}{%
\section{Итоги
главы}\label{ux438ux442ux43eux433ux438-ux433ux43bux430ux432ux44b-6}}

Если и служба поддержки, и пользователь говорят «исправлено», означает
ли это, что задача выполнена успешно? Вернёмся к проблеме с роумингом из
начала главы: при приёмке нужно проверить, достиг ли результат измерения
скорости на устройстве пользователя согласованного уровня excellent.
Если продукт предоставляет пользователю только одну попытку выполнить
операцию, лучший результат из пяти попыток не заменяет результат первой.
Если допускается \(k\) попыток, также следует различать три
формулировки: добиться успеха хотя бы один раз
(\href{mailto:Pass@k}{\nolinkurl{Pass@k}}), какую оценку получит лучшая
попытка (\href{mailto:Best@k}{\nolinkurl{Best@k}}) или добиться успеха в
каждой попытке (Pass\^{}k). Только после выбора критерия становится
понятно, какие задачи следует взять из открытых бенчмарков, собственных
наборов бизнес-задач и производственных траекторий, какие результаты
можно проверить программно, а какие требуют поэлементной оценки по
Rubric и калибровки с участием людей.

Некоторые ответы нужно прочитать даже после выставления оценки. Раньше
пользователь предпочитал краткие пояснения, но на этот раз явно
потребовал полностью описать этапы миграции. Если остановить траекторию
перед принятием решения ассистентом, можно отдельно спросить его, что
делать дальше. Он решил сразу написать полное описание, однако оценщик
засчитал неудачу, поскольку метка решения отсутствовала в допустимом
наборе. \textbf{Получить информацию и правильно применить её при
принятии текущего решения --- две разные способности.} По префиксу
траектории можно выяснить, не использовал ли агент память неправильно, а
также обнаружить, какое разумное поведение не учтено правилами оценки;
но качество итогового описания можно определить только после выполнения
задачи целиком.

При сравнении двух моделей или двух Harness нужно задать ещё один
вопрос: выполняли ли они один и тот же набор задач? Все четыре задачи
AndroidWorld, связанные с Wi-Fi, были успешно выполнены как до, так и
после упрощения дерева интерфейса, а средний расход токенов сократился
почти вдвое. Но если после расширения до полного набора задач
сохранились результаты только для варианта с упрощённым деревом, по ним
можно оценить его собственные качество, задержку и стоимость, однако
нельзя вычислить экономию относительно полного дерева.

Неудачи также указывают направление для следующего цикла обучения. Если
сотрудник службы поддержки всегда находит заказ, но раз за разом неверно
заполняет один и тот же параметр инструмента, можно сначала сохранить
вызов перед ошибкой как регрессионную задачу, а затем собрать правильные
демонстрации для похожих ситуаций. В восьмой главе рассматривается
следующий вопрос: как использовать эти материалы для улучшения модели и
на каких независимых задачах принимать новую версию? В девятой главе
обсуждается, как превратить такие изменения и проверки в повседневную
работу.

\begin{questionbox}

\hypertarget{ux432ux43eux43fux440ux43eux441ux44b-ux434ux43bux44f-ux440ux430ux437ux43cux44bux448ux43bux435ux43dux438ux44f-6}{%
\section*{Вопросы для
размышления}\label{ux432ux43eux43fux440ux43eux441ux44b-ux434ux43bux44f-ux440ux430ux437ux43cux44bux448ux43bux435ux43dux438ux44f-6}}
\addcontentsline{toc}{section}{Вопросы для размышления}

\begin{enumerate}
\def\labelenumi{\arabic{enumi}.}
\tightlist
\item
  ★★ Модель-судья может отдавать предпочтение более длинным ответам,
  определённому стилю или конкретной позиции в представлении. Как
  разработать контрольные образцы, тесты с перестановкой и наборы для
  ручной калибровки, чтобы отделить качество содержания от этих
  предпочтений?
\item
  ★★★ Открытые наборы для оценки могут попасть в обучающие данные. Как
  сочетать новые задачи, параметризованные среды, неопубликованные
  тестовые задания и производственные примеры, чтобы снизить влияние
  загрязнения данных на выводы? Какие признаки могут дать canary-метки?
\item
  ★★ Разработайте Rubric для оценки «полезности ответа» и «уместности
  тона» на основе четырёх критериев: экспертных рекомендаций, полноты
  охвата, взвешивания по важности и самодостаточности. Как с помощью
  пограничных случаев разрешать разногласия между оценщиками?
\item
  ★★ В τ²-bench модель имитирует пользователя. Как проверить, следует ли
  имитируемый пользователь условиям задачи и охватывает ли он реальные
  условия взаимодействия, включая эмоции, неоднозначность и нехватку
  информации?
\item
  ★★ Модель Bradley--Terry объясняет парные предпочтения одномерной
  латентной силой. При каких распределениях задач или группах
  пользователей могут возникать циклические предпочтения? Как
  использовать анализ по группам и ошибки подгонки модели для
  интерпретации рейтинга?
\item
  ★★ Вероятность успеха AI-агента за одну попытку равна 60\%. При
  условии независимости попыток сравните
  \href{mailto:Pass@k}{\nolinkurl{Pass@k}} и Pass\^{}k при \(k=5\).
  Затем с учётом стоимости повторных попыток, выбора результата и
  побочных эффектов действий объясните, как спроектировать оценку на
  практике.
\item
  ★★ Как при ограниченном бюджете распределить небольшие диагностические
  проверки, парные сравнения и повторные запуски, а затем провести
  последнюю проверку на задачах, не участвовавших в выборе решения,
  чтобы каждый раунд оценки приносил как можно больше пользы для
  принятия решений?
\item
  ★ В четырёх задачах AndroidWorld, связанных с Wi-Fi, полное дерево
  элементов повысило долю успешных выполнений с 25\% до 100\%, при этом
  расход токенов был в 2,498 раза больше, чем в исходном варианте. После
  сокращения дерева все четыре задачи по-прежнему выполнялись успешно, а
  расход токенов составил 0,506 от расхода для полного дерева. Как
  разработать правила сокращения и регрессионные задачи, чтобы сохранить
  состояния элементов управления, связи доступности и зависимости
  последующих операций?
\item
  ★★ Симулятор пользователя поэтапно раскрывает информацию в ответ на
  вопросы агента. Как сравнивать разные стратегии раскрытия информации и
  определять их влияние на число раундов взаимодействия, долю успешных
  выполнений и сложность задачи?
\item
  ★★ Пользователь передал исследовательскому ассистенту три результата
  наблюдений: 2, 5 и 8 --- и попросил построить график и сообщить
  среднее значение. Из журналов видно, что открыть файл не удалось,
  после чего программа построения графиков создала график по значениям
  4, 5 и 6, а в отчёте было указано среднее значение 5. Существующая
  проверка требует лишь, чтобы «график открывался, а среднее значение
  было правильным». Какой результат даст эта проверка? Как изменить
  критерии приёмки и регрессионные сценарии, чтобы обнаружить подмену
  данных?
\end{enumerate}

\textbf{Практикум к \crossreflink{chap:7}{главе 7}.}
\href{https://github.com/bojieli/ai-agent-book/tree/40f4ce5778e6cc93e3606c3079bf8cf13a0c4d04/chapter7}{Каталог
экспериментов} ·
\href{https://github.com/bojieli/ai-agent-book/blob/40f4ce5778e6cc93e3606c3079bf8cf13a0c4d04/chapter7/README.ru.md}{Путеводитель
на русском}.

\end{questionbox}

\hypertarget{ux43fux43eux441ux442ux43eux431ux443ux447ux435ux43dux438ux435-ux43cux43eux434ux435ux43bux438}{%
\chapter{\texorpdfstring{Постобучение
модели\label{chap:8}}{Постобучение модели}}\label{ux43fux43eux441ux442ux43eux431ux443ux447ux435ux43dux438ux435-ux43cux43eux434ux435ux43bux438}}

В первых семи главах мы в основном улучшали AI-агента с помощью
контекста, инструментов и механизмов выполнения. В этой главе мы
переходим к параметрам модели: как научить её стабильнее понимать
задачи, использовать инструменты, обрабатывать обратную связь и выбирать
подходящие действия в новых ситуациях? Эти способности можно постепенно
сформировать посредством обучения. Наборы задач, валидаторы и среды
симуляции из \crossreflink{chap:7}{главы 7} становятся источниками
необходимых для обучения данных и обратной связи.

Чтобы разработать схему обучения, сначала необходимо диагностировать
проблемы модели. Если модель не понимает материалы целевой предметной
области, нужно восполнить языковые способности и знания в этой области.
Если она понимает задачу, но постоянно генерирует неверные параметры
инструментов, следует с помощью демонстраций обучить её использованию
интерфейсов и порядку действий. Если модель иногда добивается успеха, но
часто выбирает неэффективные пути, можно с помощью вознаграждения
повысить вероятность эффективных стратегий. Разным проблемам
соответствуют разные данные, целевые функции и процессы обучения.

Развитие способностей современных моделей обычно включает четыре этапа:
\textbf{предобучение} закладывает универсальную языковую основу и
знания, \textbf{Mid-training} продолжает обучение на целевом
распределении, \textbf{дообучение с учителем (SFT)} обучает ответам и
действиям на демонстрациях, а \textbf{обучение с подкреплением (RL)}
оптимизирует стратегию с помощью обратной связи по результатам
выполнения задач. Сначала в этой главе мы разберём, как работают эти
четыре этапа, а затем на примерах из математики, навигации, вызова
инструментов и робототехники обсудим способы их сочетания.

Автор уделяет особое внимание двум вопросам. Первый --- как обучение на
демонстрациях и обучение с помощью вознаграждения влияют на обобщение. В
сравнительных экспериментах GeneralPoints и V-IRL ограниченный набор
демонстраций для SFT с большей вероятностью заставляет модель запоминать
паттерны обучающей выборки, тогда как RL обеспечивает лучшее обобщение
при изменяющихся условиях задачи. Далее эти различия будут объяснены с
точки зрения обучающих данных, целей оптимизации и тестового
распределения. Второй вопрос --- как результаты обучения зависят от
данных и среды. Инженерный вывод автора: \textbf{данные и среда важнее
алгоритма.} Разобравшись в основных механизмах таких алгоритмов, как PPO
и GRPO, следует сосредоточить основные усилия на охвате корпуса,
качестве демонстраций, обратной связи среды и проектировании
вознаграждения. Во многих задачах после формирования базовых
способностей и подготовки демонстрационных данных можно удовлетворить
требования практического применения с помощью одного лишь SFT.

\begin{quote}
\textbf{Рекомендации по чтению}

Если вы впервые знакомитесь с обучением моделей, читайте главу,
размышляя над одним вопросом: агент службы поддержки всегда может найти
заказ, но часто передаёт инструменту неверные параметры --- какие
обучающие материалы необходимо подготовить команде? Сначала прочитайте
разделы «Полная картина четырёх этапов» и «Когда выбирать Mid-training,
SFT и RL», запишите один демонстрационный пример, который вы бы собрали,
а затем проверьте свои выводы на трёх проблемных случаях в конце главы.
Читатели с опытом обучения моделей могут двигаться по другому маршруту:
начать с синтеза данных, вручную один раз рассчитать внутригрупповое
преимущество GRPO, а затем проверить, какие действия будет поощрять
вознаграждение в многоэтапной задаче. К разделам о классическом RL и
основах предобучения можно возвращаться при встрече с такими понятиями,
как действие, стратегия или потери.
\end{quote}

\hypertarget{ux43eux442-ux43fux440ux435ux434ux43eux431ux443ux447ux435ux43dux438ux44f-ux434ux43e-rl-ux43fux43eux43bux43dux430ux44f-ux43aux430ux440ux442ux438ux43dux430-ux447ux435ux442ux44bux440ux451ux445-ux44dux442ux430ux43fux43eux432}{%
\section{От предобучения до RL: полная картина четырёх
этапов}\label{ux43eux442-ux43fux440ux435ux434ux43eux431ux443ux447ux435ux43dux438ux44f-ux434ux43e-rl-ux43fux43eux43bux43dux430ux44f-ux43aux430ux440ux442ux438ux43dux430-ux447ux435ux442ux44bux440ux451ux445-ux44dux442ux430ux43fux43eux432}}

Один и тот же фрагмент диалога со службой поддержки можно использовать
как материал для предсказания моделью следующего токена, разметить в нём
достойные подражания вызовы инструментов или позволить модели
самостоятельно оформить заказ и выставить оценку по результату. Во всех
случаях корректируются параметры, но содержание обучения зависит от
того, кто формирует ответ, какие позиции учитываются при вычислении
потерь и как оценивается вывод. В таблице 8-1 сопоставлены материалы и
цели четырёх видов обучения.

Таблица 8-1. Четыре этапа развития возможностей модели

\begin{longtable}[]{@{}
  >{\raggedright\arraybackslash}p{(\columnwidth - 8\tabcolsep) * \real{0.1520}}
  >{\raggedright\arraybackslash}p{(\columnwidth - 8\tabcolsep) * \real{0.2194}}
  >{\raggedright\arraybackslash}p{(\columnwidth - 8\tabcolsep) * \real{0.2190}}
  >{\raggedright\arraybackslash}p{(\columnwidth - 8\tabcolsep) * \real{0.1856}}
  >{\raggedright\arraybackslash}p{(\columnwidth - 8\tabcolsep) * \real{0.2240}}@{}}
\toprule\noalign{}
\begin{minipage}[b]{\linewidth}\raggedright
\hspace{0pt}Этап
\end{minipage} & \begin{minipage}[b]{\linewidth}\raggedright
\hspace{0pt}Обучающие материалы
\end{minipage} & \begin{minipage}[b]{\linewidth}\raggedright
\hspace{0pt}Основная цель оптимизации
\end{minipage} & \begin{minipage}[b]{\linewidth}\raggedright
\hspace{0pt}Основное назначение
\end{minipage} & \begin{minipage}[b]{\linewidth}\raggedright
\hspace{0pt}Факторы, определяющие затраты ресурсов
\end{minipage} \\
\midrule\noalign{}
\endhead
\bottomrule\noalign{}
\endlastfoot
\hspace{0pt}Предобучение & \hspace{0pt}Крупномасштабные тексты, код и
данные других модальностей & \hspace{0pt}Предсказание следующего токена
авторегрессионной языковой моделью & \hspace{0pt}Формирование основ
языка, знаний и рассуждений & \hspace{0pt}Размер модели, число обучающих
токенов и эффективность системы \\
\hspace{0pt}Mid-\allowbreak{}\hspace{0pt}training & \hspace{0pt}Корпусы
на целевом языке, из предметной области, с кодом или длинным контекстом,
смешанные с общими корпусами для сохранения прежних возможностей &
\hspace{0pt}Продолжение языкового моделирования, обычно с обучающим
сигналом по большому числу текстовых токенов & \hspace{0pt}Дополнение
знаний и базовых возможностей в целевом распределении &
\hspace{0pt}Объём данных, длина последовательности и способ обновления
параметров \\
\hspace{0pt}SFT & \hspace{0pt}Ввод и ожидаемые ответы, вызовы
инструментов или полные траектории & \hspace{0pt}Повышение условной
вероятности демонстрационного вывода & \hspace{0pt}Обучение содержанию
задач, следованию инструкциям, формату и процессу действий &
\hspace{0pt}Число демонстраций, длина вывода и число эпох обучения \\
\hspace{0pt}RL & \hspace{0pt}Задачи, сгенерированные стратегией ответы
или траектории, вознаграждения & \hspace{0pt}Повышение ожидаемого
накопленного вознаграждения & \hspace{0pt}Оптимизация принятия решений,
исследования и использования обратной связи & \hspace{0pt}Число rollout,
длина генерации, стоимость среды и оценивания \\
\end{longtable}

\hypertarget{ux43fux440ux435ux434ux43eux431ux443ux447ux435ux43dux438ux435-ux438ux437ux443ux447ux435ux43dux438ux435-ux437ux430ux43aux43eux43dux43eux43cux435ux440ux43dux43eux441ux442ux435ux439-ux442ux435ux43aux441ux442ux430-ux43fux43e-ux441ux43bux435ux434ux443ux44eux449ux435ux43cux443-ux442ux43eux43aux435ux43dux443}{%
\subsection{Предобучение: изучение закономерностей текста по следующему
токену}\label{ux43fux440ux435ux434ux43eux431ux443ux447ux435ux43dux438ux435-ux438ux437ux443ux447ux435ux43dux438ux435-ux437ux430ux43aux43eux43dux43eux43cux435ux440ux43dux43eux441ux442ux435ux439-ux442ux435ux43aux441ux442ux430-ux43fux43e-ux441ux43bux435ux434ux443ux44eux449ux435ux43cux443-ux442ux43eux43aux435ux43dux443}}

Получив предшествующий текст, авторегрессионная языковая модель
распределяет вероятности между всеми токенами словаря. Если на вход
подано «Столица Китая ---», обучающей целью становится следующий символ
или фрагмент слова из обучающего текста. Разбиение на токены
определяется токенизатором: «Пекин» может соответствовать одному или
нескольким токенам.

В процессе обучения \textbf{функция потерь (Loss)} измеряет вероятность,
которую модель присвоила фактическому следующему токену: чем ниже
вероятность, тем выше соответствующие отрицательные логарифмические
потери. Затем вычисляется градиент потерь по параметрам, после чего
оптимизатор обновляет веса. Градиент показывает, какие изменения
параметров влияют на потери, а скорость обучения совместно с
оптимизатором определяет способ корректировки на данном шаге.
Посредством многократного предсказания в параметры модели переносятся
лексические и грамматические закономерности, факты и связи, необходимые
для рассуждений, из огромного объёма текстов.

Для понимания последующего обучения важно помнить, что вывод модели
представляет собой \textbf{условное распределение вероятностей}.
Обучение изменяет это распределение, повышая вероятность определённых
продолжений или действий в соответствующих ситуациях. Во время генерации
токены выбираются по правилам декодирования и шаг за шагом образуют
ответ.

Предобучающие тексты содержат вопросы и ответы, статьи, код и
разнообразные фрагменты, поэтому цель продолжения также меняется вместе
с контекстом. Получив вопрос, базовая модель может сгенерировать ответ,
а может продолжить текст ещё одним вопросом. Чтобы она стабильно
следовала инструкциям пользователя, ей необходимы дополнительные
соответствующие обучающие материалы.

Чтобы наглядно понять смысл «обучения предсказанию», достаточно
рассмотреть одну позицию. Предположим, обучающее предложение звучит так:
«У этого поезда конечная станция --- Шанхай», а токенизатор кодирует
слово «Шанхай» одним токеном. Если после фрагмента «конечная станция
---» вероятность токена «Шанхай» повышается с 0,25 до 0,5, отрицательное
логарифмическое правдоподобие в этой позиции снижается примерно с 1,39
до 0,69 при использовании натурального логарифма. Такое обновление
использует только слово, которое действительно следует далее в тексте.
Чтобы определить, оформлен ли билет пользователя, необходимо проверить
состояние заказа --- следующий токен обучающего текста такой обратной
связи не предоставляет.

\hypertarget{mid-training-ux43fux440ux43eux434ux43eux43bux436ux435ux43dux438ux435-ux43eux431ux443ux447ux435ux43dux438ux44f-ux43dux430-ux446ux435ux43bux435ux432ux43eux43c-ux440ux430ux441ux43fux440ux435ux434ux435ux43bux435ux43dux438ux438}{%
\subsection{Mid-training: продолжение обучения на целевом
распределении}\label{mid-training-ux43fux440ux43eux434ux43eux43bux436ux435ux43dux438ux435-ux43eux431ux443ux447ux435ux43dux438ux44f-ux43dux430-ux446ux435ux43bux435ux432ux43eux43c-ux440ux430ux441ux43fux440ux435ux434ux435ux43bux435ux43dux438ux438}}

Если общая модель почти не встречала профессиональных документов на
корейском языке или внутренних протоколов компании, можно продолжить
языковое моделирование на соответствующих корпусах, предоставив модели
больше возможностей изучить целевой язык, профессиональную терминологию,
шаблоны кода и связи в длинных документах. Добавление общего корпуса к
данным предметной области помогает модели сохранять качество выполнения
уже освоенных задач. Такое обучение часто называют \textbf{Mid-training}
или продолженным предобучением (Continued Pre-training, CPT); термины
«предобучение с адаптацией к предметной области» (DAPT) и «предобучение
с адаптацией к задаче» (TAPT) дополнительно уточняют область отбора
корпусов.

При Mid-training в вычислении потерь обычно участвует большое число
токенов из документов, кода и цепочек вывода. Одно понятие может
неоднократно встречаться в определениях, примерах, выводах и разных
документах, благодаря чему модель изучает множество связей и способов
его применения. При SFT то же содержимое обычно организуется как ввод и
ожидаемый вывод, а обучающий сигнал сосредоточен на требуемых задачей
ответах или действиях. Небольшой набор вопросов и ответов охватывает
лишь несколько способов формулировать вопросы и ответы. Даже если модель
запомнит содержащиеся в нём факты, при иной формулировке или
необходимости объединить несколько понятий она может оказаться
неспособна применить эти знания. Именно поэтому для крупномасштабных
знаний предметной области требуется более широкий охват корпусом.

При выборе подхода сначала можно оценить объём знаний и способ их
обновления: для крупного корпуса предметной области подходит продолжение
обучения; для материалов, которые требуется извлекать с учётом источника
и часто обновлять, подходит RAG; при наличии чётко определённой задачи и
качественных ответов посредством SFT можно научить модель использовать
эти знания. Подходы также можно сочетать, поручая им соответственно
развитие базовых возможностей, работу с внешними знаниями и выполнение
требований к поведению.

\hypertarget{sft-ux43eux431ux443ux447ux435ux43dux438ux435-ux43eux442ux432ux435ux442ux430ux43c-ux438-ux434ux435ux439ux441ux442ux432ux438ux44fux43c-ux43fux43e-ux434ux435ux43cux43eux43dux441ux442ux440ux430ux446ux438ux44fux43c}{%
\subsection{SFT: обучение ответам и действиям по
демонстрациям}\label{sft-ux43eux431ux443ux447ux435ux43dux438ux435-ux43eux442ux432ux435ux442ux430ux43c-ux438-ux434ux435ux439ux441ux442ux432ux438ux44fux43c-ux43fux43e-ux434ux435ux43cux43eux43dux441ux442ux440ux430ux446ux438ux44fux43c}}

\textbf{Дообучение с учителем (Supervised Fine-Tuning, SFT)} организует
пользовательский ввод и ожидаемый вывод в обучающие образцы, позволяя
модели повысить условную вероятность каждого токена демонстрационного
вывода. Для авторегрессионной языковой модели сохраняется цель в виде
отрицательного логарифмического правдоподобия, но обучающие материалы
гораздо точнее отражают требования задачи.

Например, демонстрация вызова инструмента содержит пользовательский
запрос, определение инструмента и ожидаемый вызов. Модель читает
предшествующий контекст и учится выбирать нужный инструмент, заполнять
его параметры и отвечать после получения результата. Многораундовые
демонстрации могут также включать уточнение, проверку, восстановление и
завершение задачи.

В демонстрации проверки заказа можно выделить два фрагмента, которые при
развёртывании должен генерировать ассистент: имя инструмента и параметры
для запуска проверки, а также ответ после получения сведений о заказе.
Если при вычислении потерь оцениваются только токены этих двух
фрагментов, это называется \textbf{loss masking (маскированием потерь)}.
Системная инструкция, вопрос пользователя и возвращённые инструментом
сведения о заказе по-прежнему входят в контекст и определяют, как должна
ответить модель. Если же разметить содержимое заказа как вывод
ассистента, модель будет учиться самой дописывать сведения, которые
должна предоставлять система заказов.

Эффективность SFT обеспечивается плотным обучающим сигналом по каждому
токену демонстрации. Для задач со стабильным форматом и чётко
ограниченной областью нескольких тысяч качественных образцов может быть
достаточно, чтобы модель освоила необходимое поведение; сложным задачам
требуется более широкий охват. Модель изучает не только знания из
демонстраций, но и формулировки, формат и принятый порядок действий. При
чрезмерно однородных данных вместе с ними могут усилиться и некоторые
поверхностные связи.

\begin{quote}
\textbf{Параметро-эффективное обучение: LoRA}

При полно-параметрическом дообучении необходимо хранить градиенты и
состояния оптимизатора для огромного числа весов. LoRA (Low-Rank
Adaptation, низкоранговая адаптация) замораживает исходную матрицу \(W\)
и представляет обновление как произведение двух меньших матриц:
\(W'=W+\gamma BA\). Здесь ранг (rank) определяет ёмкость обновления, а
\(\gamma\) --- масштабирование; число обучаемых параметров зависит от
rank, целевых слоёв и архитектуры модели. LoRA применима как к SFT, так
и к RL и позволяет загружать несколько адаптеров для одной базовой
модели. Автор обычно сначала создаёт базовый вариант постобучения с
LoRA, чтобы сравнить данные и схемы обучения с меньшими затратами
видеопамяти и хранилища; если ёмкости адаптера недостаточно, увеличивает
rank или сравнивает результат с полно-параметрическим обучением.

Эксперименты из работы «LoRA Without Regret» дают несколько полезных
ориентиров для конфигурации: включить в адаптацию основные веса, в том
числе MLP; постепенно повышать rank, проверяя достаточность ёмкости;
отдельно подбирать скорость обучения для LoRA. В этих тестах подходящая
для LoRA скорость обучения была примерно в десять раз выше, чем при
полно-параметрическом дообучении, а RL показало хорошие результаты и при
низком rank.\footnote{Schulman, John and Thinking Machines Lab, ``LoRA
  Without Regret'', 2025.} На практике в качестве кандидатных диапазонов
можно использовать rank 64--256 для SFT и 8--32 для RL, сравнивая их
также с меньшими значениями rank.
\end{quote}

\hypertarget{ux43fux43bux430ux43dux438ux440ux43eux432ux430ux43dux438ux435-sft-ux438-rl-ux441-ux443ux447ux451ux442ux43eux43c-ux438ux441ux445ux43eux434ux43dux43eux433ux43e-ux441ux43eux441ux442ux43eux44fux43dux438ux44f-ux43cux43eux434ux435ux43bux438}{%
\subsection{Планирование SFT и RL с учётом исходного состояния
модели}\label{ux43fux43bux430ux43dux438ux440ux43eux432ux430ux43dux438ux435-sft-ux438-rl-ux441-ux443ux447ux451ux442ux43eux43c-ux438ux441ux445ux43eux434ux43dux43eux433ux43e-ux441ux43eux441ux442ux43eux44fux43dux438ux44f-ux43cux43eux434ux435ux43bux438}}

RL оценивает посредством вознаграждения ответы или траектории, которые
модель генерирует самостоятельно. Процесс выполнения стратегии от
заданной исходной точки до завершения и сбора одной траектории
называется одним \textbf{rollout}. Возможность получить полезную
обратную связь при обучении зависит от протокола вывода и текущего
распределения возможностей.

Сначала следует проверить \textbf{протокол вывода}. Если модель часто
генерирует JSON, который невозможно разобрать, верификатор может
зарегистрировать ошибку формата, но такую обратную связь трудно
разделить по отдельным возможностям задачи. После стабилизации полей,
вызовов инструментов и основного процесса посредством демонстраций SFT
вознаграждение сможет в большей степени отражать качество содержимого
или действий. Это одно из распространённых применений схемы «сначала
SFT, затем RL».

Затем следует проверить \textbf{вероятность сэмплирования успешного
поведения}. На отложенных задачах в условиях, близких к обучению,
измеряются \href{mailto:Pass@1}{\nolinkurl{Pass@1}} и
\href{mailto:Pass@k}{\nolinkurl{Pass@k}}. Если вероятность успеха каждой
попытки одной задачи равна \(p\), а выборки приблизительно независимы,
то вероятность хотя бы одного успеха за \(k\) попыток составляет:

\[
\operatorname{Pass@}k=1-(1-p)^k
\]

Если \href{mailto:Pass@1}{\nolinkurl{Pass@1}} низок, но при увеличении
\(k\) вероятность успеха заметно возрастает, текущая стратегия уже
способна генерировать успешные пути. RL, сэмплирование с отбраковкой или
дистилляция могут использовать эти пути и повышать вероятность их
появления. Если же при разумном бюджете и охвате задач успешные попытки
по-прежнему крайне редки, сначала следует определить, на каких шагах
модели не хватает знаний, соблюдения формата или локальных возможностей.

Рассмотрим бинарное итоговое вознаграждение. Если все rollout одной
группы получают 0 баллов, относительные внутригрупповые преимущества в
GRPO также равны 0, поэтому по одним только вознаграждениям этой группы
невозможно определить, какая попытка была лучше. Оценка ценности,
исследование и другие ограничения в PPO могут и дальше приводить к
обновлениям, однако постоянное отсутствие успешных образцов всё равно
затрудняет обучение задаче. В приближении независимых попыток с
фиксированной вероятностью успеха до первого успеха в среднем требуется
\(1/p\) попыток; при низкой вероятности успеха стоимость сэмплирования
быстро растёт. Поэтому автор сначала добавляет демонстрации, базовые
знания или достижимую промежуточную обратную связь, повышая вероятность
получения полезного сигнала при сэмплировании, и лишь затем увеличивает
число rollout.

Обучение можно дополнять с учётом источника проблемы: при недостаточном
охвате языка, фактов и шаблонов кода следует добавить соответствующий
корпус; при нестабильности интерфейсов и основного процесса ---
демонстрации SFT; если модель способна выполнять отдельные шаги ---
спроектировать проверяемые частичные вознаграждения или обучение по
учебной программе. Модель с достаточно сильной исходной точкой можно
сразу обучать посредством RL. Такой путь продемонстрировала
DeepSeek-R1-Zero, тогда как DeepSeek-R1 дополнительно объединила данные
холодного старта с другими этапами обучения, улучшив представление
рассуждений, читаемость и языковую согласованность. В рассматриваемых
далее экспериментах GeneralPoints будет показан другой тип исходной
точки, для которой сначала необходимо стабилизировать формат.

\hypertarget{sft-ux438-rl-ux43aux430ux43a-ux434ux432ux430-ux43eux431ux443ux447ux430ux44eux449ux438ux445-ux441ux438ux433ux43dux430ux43bux430-ux432ux43bux438ux44fux44eux442-ux43dux430-ux43fux43eux432ux435ux434ux435ux43dux438ux435}{%
\subsection{SFT и RL: как два обучающих сигнала влияют на
поведение}\label{sft-ux438-rl-ux43aux430ux43a-ux434ux432ux430-ux43eux431ux443ux447ux430ux44eux449ux438ux445-ux441ux438ux433ux43dux430ux43bux430-ux432ux43bux438ux44fux44eux442-ux43dux430-ux43fux43eux432ux435ux434ux435ux43dux438ux435}}

SFT сопоставляет вывод с демонстрацией, а RL оценивает сгенерированное
поведение по вознаграждению. GeneralPoints предоставляет конкретный
пример для сравнения: если в ограниченных демонстрациях карты J, Q и K
всегда считаются равными 10, модель может выучить это фиксированное
соответствие; изменение значений карт при тестировании позволяет
проверить, действительно ли она выполняет вычисления по текущим
правилам.

Таблица 8-2. Сравнение механизмов обучения SFT и RL

\begin{longtable}[]{@{}
  >{\raggedright\arraybackslash}p{(\columnwidth - 4\tabcolsep) * \real{0.2833}}
  >{\raggedright\arraybackslash}p{(\columnwidth - 4\tabcolsep) * \real{0.3673}}
  >{\raggedright\arraybackslash}p{(\columnwidth - 4\tabcolsep) * \real{0.3493}}@{}}
\toprule\noalign{}
\begin{minipage}[b]{\linewidth}\raggedright
\hspace{0pt}Аспект
\end{minipage} & \begin{minipage}[b]{\linewidth}\raggedright
\hspace{0pt}SFT
\end{minipage} & \begin{minipage}[b]{\linewidth}\raggedright
\hspace{0pt}RL
\end{minipage} \\
\midrule\noalign{}
\endhead
\bottomrule\noalign{}
\endlastfoot
\hspace{0pt}Непосредственная цель оптимизации & \hspace{0pt}Повышение
условной вероятности демонстрационного вывода & \hspace{0pt}Повышение
ожидаемого вознаграждения стратегии \\
\hspace{0pt}Источник обучающего сигнала & \hspace{0pt}Токены заданного
ответа или действия & \hspace{0pt}Траектории текущей или ранее
существовавшей стратегии и обратная связь по ним \\
\hspace{0pt}Форма данных & \hspace{0pt}Образцы «ввод --- вывод»,
многораундовые демонстрации & \hspace{0pt}Задачи, состояния, действия,
вознаграждения и условия завершения \\
\hspace{0pt}Основная инженерная работа & \hspace{0pt}Создание и отбор
демонстраций, настройка позиций для вычисления потерь &
\hspace{0pt}Генерация rollout, выполнение среды, оценивание и обновление
стратегии \\
\hspace{0pt}Условия, от которых зависит обобщение & \hspace{0pt}Охват
демонстраций, базовые возможности и интенсивность обучения &
\hspace{0pt}Распределение задач, схема вознаграждения, исследование и
стабильность оптимизации \\
\hspace{0pt}Ресурсные особенности &
\hspace{0pt}Офлайн-\allowbreak{}\hspace{0pt}демонстрации можно
непосредственно использовать многократно & \hspace{0pt}Основные затраты
могут приходиться на онлайн-\allowbreak{}\hspace{0pt}генерацию и
выполнение среды \\
\hspace{0pt}Подходящая исходная точка & \hspace{0pt}Есть качественные
демонстрации, требуется сформировать поведение для задачи &
\hspace{0pt}Модель уже способна выдавать полезные ответы или действия, а
по результату можно вычислить вознаграждение \\
\end{longtable}

Эти две цели также можно понять через распределение вероятностей. Если у
демонстрации есть несколько правильных вариантов, обучение методом
максимального правдоподобия подталкивает модель охватить их, что часто
называют тенденцией к \textbf{mass-covering (охвату массы)}. RL
перераспределяет вероятности в соответствии с вознаграждением; при
наличии различий в вознаграждениях и соответствующей регуляризации
вероятность может сосредоточиться на нескольких вариантах с высоким
вознаграждением, проявляя тенденцию к \textbf{mode-seeking (поиску
мод)}. Итоговая степень разнообразия также зависит от данных,
ограничения KL, энтропийной регуляризации и настроек сэмплирования.

\textbf{Постобучение также влияет на момент начала действий.} Если
демонстрации многократно показывают «сначала широко исследовать
репозиторий, затем редактировать», модель осваивает такую
последовательность исследования. Если обратная связь постоянно поощряет
сочетание быстрой локализации, внесения изменений и тестирования,
соответствующее поведение усиливается. Момент перехода модели от чтения
к редактированию отражает выбор между двумя выгодами: насколько чтение
ещё одного файла повысит уверенность и сколько полезной обратной связи
можно получить, сначала изменив код, а затем запустив тесты. В
\crossreflink{chap:7}{главе 7} эта разница наблюдалась в одном и том же
Coding Harness: одни модели сначала читали множество файлов и лишь потом
приступали к работе, тогда как другие начинали вносить изменения гораздо
раньше.

Онлайн-взаимодействие также предоставляет три возможности для обучения.

\begin{itemize}
\tightlist
\item
  \textbf{Пробовать подходы, отсутствующие в демонстрациях.} Модель
  может комбинировать имеющиеся знания, предлагая новые пути, после чего
  среда проверяет результат. В робототехническом примере SimpleVLA-RL
  присутствуют стратегии действий, отсутствующие в демонстрациях; далее
  на основе описания задачи будет рассмотрено исследование способов
  формирования эффективного поведения.
\item
  \textbf{Использовать особенности задач, в которых проверка проще
  генерации.} Математические ответы, тесты кода и формальные
  доказательства можно проверять соответствующими верификаторами. Модель
  может самостоятельно решать задачи, писать код или приводить
  доказательства, а затем получать вознаграждение по результатам
  проверки. Условия, охватываемые верификатором, также определяют, в
  сторону какого поведения будет оптимизироваться модель.
\item
  \textbf{Обучаться на состояниях, реально посещаемых стратегией.}
  Отклонившись от демонстрации, модель попадает в новые состояния,
  вследствие чего меняются и последующие действия. Это называется
  \textbf{ковариатным сдвигом (covariate shift)}. При соответствующих
  допущениях классического последовательного имитационного обучения
  верхняя граница ошибки может расти с длиной траектории \(T\) как
  \(T^2\); онлайн-агрегация данных, охватывающая посещаемые стратегией
  состояния, позволяет улучшить эту границу до линейного роста с \(T\).
  Рассматриваемая далее On-Policy Distillation будет использовать
  контекст, сгенерированный текущей стратегией, а затем подключать
  плотный обучающий сигнал учителя.
\end{itemize}

Таким образом, процесс обучения можно разворачивать поэтапно: сначала
сформировать пригодные возможности и протокол поведения, а затем
повышать качество посредством цикла генерации, проверки и обновления.
Следующие разделы о классическом RL и основах предобучения дополняют
необходимый контекст; их также можно пропустить и сразу перейти к
разделу о Mid-training, где подробнее рассматривается выбор конкретных
данных.

\hypertarget{ux43eux442-ux43aux43bux430ux441ux441ux438ux447ux435ux441ux43aux43eux433ux43e-rl-ux430ux433ux435ux43dux442ux430-ux43a-ux441ux43eux432ux440ux435ux43cux435ux43dux43dux43eux43cux443-ai-ux430ux433ux435ux43dux442ux443-ux43dux435ux43eux431ux44fux437ux430ux442ux435ux43bux44cux43dux43eux435-ux447ux442ux435ux43dux438ux435}{%
\section{\texorpdfstring{От классического RL-агента к современному
AI-агенту
\texttt{{[}Необязательное\ чтение{]}}}{От классического RL-агента к современному AI-агенту {[}Необязательное чтение{]}}}\label{ux43eux442-ux43aux43bux430ux441ux441ux438ux447ux435ux441ux43aux43eux433ux43e-rl-ux430ux433ux435ux43dux442ux430-ux43a-ux441ux43eux432ux440ux435ux43cux435ux43dux43dux43eux43cux443-ai-ux430ux433ux435ux43dux442ux443-ux43dux435ux43eux431ux44fux437ux430ux442ux435ux43bux44cux43dux43eux435-ux447ux442ux435ux43dux438ux435}}

\hypertarget{ux432ux437ux430ux438ux43cux43eux434ux435ux439ux441ux442ux432ux438ux435-ux430ux433ux435ux43dux442ux430-ux441ux43e-ux441ux440ux435ux434ux43eux439}{%
\subsection{Взаимодействие агента со
средой}\label{ux432ux437ux430ux438ux43cux43eux434ux435ux439ux441ux442ux432ux438ux435-ux430ux433ux435ux43dux442ux430-ux441ux43e-ux441ux440ux435ux434ux43eux439}}

Обучение с подкреплением изучает, как выбирать действия на основе
доступной в текущий момент информации, чтобы повысить долгосрочное
накопленное вознаграждение. Например, в настольной игре стратегия
выбирает ход, после чего среда обновляет доску; по завершении партии
выдаётся вознаграждение в зависимости от победы или поражения. Первые
действия могут не приносить немедленных очков, но влиять на конечный
результат.

На \crossreflink{fig:8-1}{рисунке 8-1} связь между действиями,
наблюдениями и вознаграждениями показана на примере упрощённой игры с
поиском сокровищ. Игрок идёт из вестибюля на восток, и игра сообщает,
что он оказался в кладовой; вознаграждение за этот шаг равно 0. Затем
игрок подбирает красный ключ и получает 5 очков вознаграждения. Игровая
программа знает полное состояние мира, а объём доступной игроку
информации зависит от правил игры. Если наблюдение содержит полное
состояние, среда называется полностью наблюдаемой. Если видна лишь часть
состояния, игроку часто приходится помнить, что он видел и делал раньше,
чтобы решить, куда двигаться дальше.

\begin{figure}
\hypertarget{fig:8-1}{%
\centering
\includegraphics{images/revised/fig8-1.pdf}
\caption{Рисунок 8-1. Цикл взаимодействия агента со средой в обучении с
подкреплением}\label{fig:8-1}
}
\end{figure}

Непрерывное взаимодействие образует \textbf{траекторию}, включающую
наблюдения, действия, вознаграждения и информацию о завершении на каждом
шаге. Если обучающей системе доступно состояние среды, его также можно
записывать. Стратегия определяет способ действий, а \textbf{функция
ценности (Value Function)} оценивает, какую будущую отдачу можно
ожидать, если из текущего состояния продолжить выполнение определённой
стратегии.

Граница между агентом и средой определяется областью системы, выбранной
при моделировании. В рассматриваемых в этой книге системах модель и
Harness находятся внутри агента, а интерфейсы инструментов связывают его
с внешними приложениями, устройствами и сервисами. Например, агент
формирует запрос на изменение файла, файловая система выполняет запись,
а инструмент возвращает агенту результат выполнения.

В процессе обучения необходимо распределять возможности между
\textbf{исследованием и использованием (Exploration--Exploitation)}.
Исследование знакомит стратегию с ещё недостаточно изученными
действиями, а использование выбирает действия, которые в текущий момент
оцениваются как более выгодные. В задачах с отложенным вознаграждением
также требуется определить, какие ранние шаги способствовали конечному
результату. Это и есть проблема присвоения заслуг, рассматриваемая
далее.

Чтобы понять метод RL, сначала можно проверить пять аспектов:

\begin{itemize}
\tightlist
\item
  \textbf{Пространство действий}: что разрешено выдавать на выходе. В
  настольных играх используются дискретные ходы, а в управлении роботами
  --- непрерывные целевые значения суставов. Действия также могут
  ограничиваться состоянием, физическими пределами и разрешениями.
\item
  \textbf{Стратегия}: правило выбора действия на основе состояния,
  наблюдения или истории. Это может быть как таблица, так и нейронная
  сеть.
\item
  \textbf{Сигнал вознаграждения}: обратная связь, посредством которой
  выражается цель обучения. Вознаграждение каждого шага участвует в
  формировании долгосрочной отдачи, поэтому стратегия может принять
  текущие издержки ради будущей выгоды.
\item
  \textbf{Функция ценности}: оценивает последующую отдачу и может
  использоваться для сравнения действий или формирования сигнала
  обновления стратегии. Конкретный алгоритм определяет, обучается ли она
  явно.
\item
  \textbf{Модель среды}: предсказывает состояние, наблюдение или
  вознаграждение после действия. Методы на основе модели планируют с
  использованием известной или изученной модели среды; безмодельные
  методы обучают стратегию или ценность непосредственно на данных
  взаимодействия.
\end{itemize}

В таблице 8-3 для сопоставления нескольких знакомых систем из этих
понятий выбраны среда, действия и цель. При чтении первых двух строк
«сохранение равновесия» и «очищенную площадь» можно рассматривать как
возможные цели, выбранные при моделировании задачи.

Таблица 8-3. Сопоставление ключевых элементов разных систем агентов

\begin{longtable}[]{@{}
  >{\raggedright\arraybackslash}p{(\columnwidth - 6\tabcolsep) * \real{0.2140}}
  >{\raggedright\arraybackslash}p{(\columnwidth - 6\tabcolsep) * \real{0.2564}}
  >{\raggedright\arraybackslash}p{(\columnwidth - 6\tabcolsep) * \real{0.2644}}
  >{\raggedright\arraybackslash}p{(\columnwidth - 6\tabcolsep) * \real{0.2652}}@{}}
\toprule\noalign{}
\begin{minipage}[b]{\linewidth}\raggedright
\hspace{0pt}Тип агента
\end{minipage} & \begin{minipage}[b]{\linewidth}\raggedright
\hspace{0pt}Внешняя среда
\end{minipage} & \begin{minipage}[b]{\linewidth}\raggedright
\hspace{0pt}Представление действий
\end{minipage} & \begin{minipage}[b]{\linewidth}\raggedright
\hspace{0pt}Пример вознаграждения или цели
\end{minipage} \\
\midrule\noalign{}
\endhead
\bottomrule\noalign{}
\endlastfoot
\hspace{0pt}Новорождённая газель & \hspace{0pt}Рельеф, гравитация и
динамика тела & \hspace{0pt}Непрерывное управление мышцами &
\hspace{0pt}Сохранение равновесия, предотвращение падения \\
\hspace{0pt}Робот-\allowbreak{}\hspace{0pt}пылесос &
\hspace{0pt}Помещение, пол и зарядная инфраструктура &
\hspace{0pt}Перемещение, уборка, возврат на зарядку и управляющие
параметры & \hspace{0pt}Очищенная площадь, заряд аккумулятора и цена
столкновений \\
\hspace{0pt}Шахматный агент & \hspace{0pt}Доска, противник и правила
контроля времени & \hspace{0pt}Легальные ходы в текущем состоянии &
\hspace{0pt}Победа +1, поражение −1, ничья 0 \\
\hspace{0pt}Агент клиентской поддержки & \hspace{0pt}Пользователь,
бизнес-\allowbreak{}\hspace{0pt}база данных и интерфейсы сервисов &
\hspace{0pt}Текст ответа, запросы и
бизнес-\allowbreak{}\hspace{0pt}вызовы & \hspace{0pt}Решение проблемы,
соблюдение процедур, время обработки \\
\hspace{0pt}Агент --- помощник программиста & \hspace{0pt}Источник
требований, кодовая база и среда выполнения & \hspace{0pt}Поиск,
редактирование, запуск и ответ & \hspace{0pt}Результаты тестов,
выполнение задачи, регрессии и затраты \\
\end{longtable}

В этих системах используются разные представления действий. В настольных
играх обычно на каждом шаге выбирается одно дискретное действие, робот
выдаёт непрерывные управляющие величины заданной размерности, а LLM
сначала генерирует последовательность токенов, которую интерфейс затем
разбирает как ответ или вызов инструмента. Внутренние вычисления также
можно включить в модель стратегии и записывать отдельно от выполнения
инструментов, фактически изменяющего внешнее состояние.

\hypertarget{ux434ux432ux430-ux43fux440ux435ux434ux441ux442ux430ux432ux43bux435ux43dux438ux44f-ux434ux435ux439ux441ux442ux432ux438ux439-ux43aux43bux430ux441ux441ux438ux447ux435ux441ux43aux430ux44f-ux43fux43eux441ux442ux430ux43dux43eux432ux43aux430-rl-ux438-ux441ux442ux440ux430ux442ux435ux433ux438ux44f-llm-ux441-ux43fux435ux440ux435ux43cux435ux43dux43dux43eux439-ux434ux43bux438ux43dux43eux439}{%
\subsection{Два представления действий: классическая постановка RL и
стратегия LLM с переменной
длиной}\label{ux434ux432ux430-ux43fux440ux435ux434ux441ux442ux430ux432ux43bux435ux43dux438ux44f-ux434ux435ux439ux441ux442ux432ux438ux439-ux43aux43bux430ux441ux441ux438ux447ux435ux441ux43aux430ux44f-ux43fux43eux441ux442ux430ux43dux43eux432ux43aux430-rl-ux438-ux441ux442ux440ux430ux442ux435ux433ux438ux44f-llm-ux441-ux43fux435ux440ux435ux43cux435ux43dux43dux43eux439-ux434ux43bux438ux43dux43eux439}}

\textbf{Марковский процесс принятия решений (Markov Decision Process,
MDP)} описывает задачу принятия решений посредством состояний, действий,
переходов и вознаграждений. Марковское свойство требует, чтобы для
предсказания следующего состояния и вознаграждения при заданных текущем
состоянии и действии не требовалась дополнительная история: вся
относящаяся к прогнозу историческая информация уже должна содержаться в
текущем состоянии. На \crossreflink{fig:8-2}{рисунке 8-2} видно, почему
недостаточно записывать лишь комнату, в которой находится игрок: даже
если он дважды подходит к одной и той же двери, возможность открыть её
зависит от того, успел ли он подобрать ключ.

\begin{figure}
\hypertarget{fig:8-2}{%
\centering
\includegraphics{images/revised/fig8-2.pdf}
\caption{Рисунок 8-2. Схема марковского процесса принятия решений
(MDP)}\label{fig:8-2}
}
\end{figure}

Шахматы дают наглядный пример проектирования состояния. Помимо
расположения фигур необходимо записывать очерёдность хода, права на
рокировку и взятие на проходе, а также информацию, нужную для применения
правила пятидесяти ходов и определения повторения позиции. Только
включив эти данные в состояние, можно применять правила на основе
текущего состояния. Если агент видит лишь часть этой информации, ему
требуется восполнять её с помощью истории либо использовать частично
наблюдаемую модель.

MDP может включать дискретные, непрерывные и составные действия. В го
есть 361 пересечение доски, а также действие паса и ограничения
легальности; в играх Atari обычно доступно небольшое число дискретных
операций; управляющие величины робота определяются суставами, захватами
и физическими ограничениями. Если число состояний и действий невелико,
ценности можно хранить непосредственно в таблице. При больших масштабах
применяются соответственно аппроксимация функций, оптимизация стратегии
и поиск.

\textbf{Табличный Q-learning: обновление ценности действия по результату
одного шага.} \(Q(s,a)\) оценивает отдачу после выполнения действия
\(a\) в состоянии \(s\). Оптимальная ценность действия удовлетворяет
уравнению оптимальности Беллмана:

\[
Q^*(s,a)=\mathbb E\left[r+\gamma\max_{a'}Q^*(s',a')\mid s,a\right]
\]

Здесь \(r\) --- вознаграждение за текущий шаг, \(s'\) --- следующее
состояние, а \(\gamma\) --- коэффициент дисконтирования. Для
продолжающихся задач обычно выбирают \(0\leq\gamma<1\); в задачах с
явным завершением и конечной отдачей можно использовать и значение 1.
Чем больше коэффициент дисконтирования, тем больший вес сохраняют
отдалённые вознаграждения. Отдача, рассчитанная от текущего момента,
равна \(G_t=\sum_{k\geq0}\gamma^k r_{t+k+1}\).

Получив в результате одного взаимодействия \(r\) и \(s'\), Q-learning
выполняет следующее обновление:

\[
Q(s,a)\leftarrow Q(s,a)+\alpha\left[r+\gamma\max_{a'}Q(s',a')-Q(s,a)\right]
\]

Здесь \(\alpha\) --- скорость обучения, а выражение в квадратных скобках
представляет собой разность между «вознаграждением за один шаг плюс
дисконтированной последующей ценностью» и текущей оценкой. При
достижении конечного состояния последующая ценность принимается равной
0. Попробуйте выполнить это действие в прилагаемой игре: отключите
случайные возмущения, от входа идите на восток, а затем подберите
красный ключ в кладовой. За ключ выдаётся 5 очков, а каждый шаг стоит
0,5 очка, поэтому для этого действия \(r\) равняется 4,5. Предположим,
что прежняя табличная оценка действия «взять ключ в кладовой» равна
0,25, а максимальное Q-значение доступных после получения ключа действий
составляет 0,63. При \(\gamma=0.99\) и \(\alpha=0.2\) одношаговая цель
равна \(4.5+0.99\times0.63=5.1237\). Прежняя оценка приближается к этой
цели на одну пятую и обновляется до 1,22474. Такой метод коррекции
оценки по фактической обратной связи за один шаг называется
\textbf{обучением на основе временных различий (Temporal-Difference
Learning, TD)}.

Сопоставив Q-значения всех действий в одном состоянии, можно выбрать
действие с наилучшей текущей оценкой. На
\crossreflink{fig:8-3}{рисунке 8-3} такая стратегия представлена в виде
упрощённой сетки. Число в каждой ячейке обозначает нормализованную
ценность состояния. Если закрыть числа рукой и смотреть только на
стрелки, сможете ли вы добраться из левого верхнего угла до сокровища?

\begin{figure}
\hypertarget{fig:8-3}{%
\centering
\includegraphics{images/revised/fig8-3.pdf}
\caption{Рисунок 8-3. Q-learning в сеточном мире}\label{fig:8-3}
}
\end{figure}

Следующее исследование может изменить направления стрелок. На
\crossreflink{fig:8-4}{рисунке 8-4} разобран приведённый выше расчёт для
подбора красного ключа: полученное на этом шаге вознаграждение изменяет
оценку действия «взять ключ в кладовой», тогда как Q-значения других
состояний таблицы будут обновлены позднее, каждое в свой момент.

\begin{figure}
\hypertarget{fig:8-4}{%
\centering
\includegraphics{images/revised/fig8-4.pdf}
\caption{Рисунок 8-4. Визуализация обновления Q-значений}\label{fig:8-4}
}
\end{figure}

Поведенческая стратегия Q-learning может включать случайное
исследование, тогда как в цели обновления используется максимальная
ценность действия в следующем состоянии. Поэтому Q-learning относится к
обучению \textbf{вне стратегии (Off-Policy)}.

\begin{quote}
\textbf{Эксперимент 8-1 ★: поведение Q-learning в игре с поиском
сокровищ}

Игра с поиском сокровищ содержит скрытые правила, связанные с ключами и
дверями, эффектами оружия и объединением предметов. Для прохождения
необходимо выполнить несколько шагов в правильной последовательности;
оптимальный путь занимает 11 шагов. За ключевые действия и победу
выдаётся вознаграждение, а на остальных шагах обратная связь невелика.

В этом обучении скорость обучения равна 0,2, коэффициент дисконтирования
--- 0,99, а доля исследования ε постепенно снижается с 1,0 до 0,1.
Табличный Q-learning применяет ε-жадное исследование: на каждом шаге с
вероятностью ε случайно выбирается одно из доступных действий, а в
остальных случаях --- одно из действий с максимальным текущим
Q-значением. На каждом этапе записываются доля побед, число посещённых
состояний и число шагов до прохождения.

При чтении журналов обучения и оценки эксперимента с поиском сокровищ
сначала разделите два вопроса. Когда во время обучения стратегия со
случайным исследованием впервые находит сокровище? Способна ли она после
завершения обучения, при фиксированной Q-таблице и жадном выборе
действий, стабильно проходить игру? Первый вопрос показывает, как модель
обнаруживает путь, а второй проверяет, способна ли она выбирать действия
в соответствии с изученными ценностями.

В одном из предыдущих запусков за первые 1000 игровых сессий агент ни
разу не прошёл игру, а Q-таблица содержала 124 состояния. К 5000 сессиям
стабильных побед по-прежнему не было, а число состояний достигло 133.
Между 7000-й и 8000-й сессиями доля побед выросла с 34 до 96\%. К
10000-й сессии она достигла 100\%, Q-таблица содержала 145 состояний, а
агент нашёл путь из 11 шагов.

Почему вначале прогресса долго не было, а затем доля побед быстро
выросла? Сначала исследование обнаружило ключевые действия и путь, а
последующие посещения постепенно распространили отдачу конечной точки на
предшествующие состояния. Вся небольшая симуляция занимает около 10
секунд, поэтому допускает большое число повторных взаимодействий.

При анализе записывайте как время вычислений, так и число
взаимодействий. В игре одно действие стоит дёшево; при работе с
браузером, телефоном или физическим устройством каждое взаимодействие
также требует ожидания, расходов и изменения состояния. Кроме того,
таблица привязана к конкретному кодированию состояний: после изменения
правил или карты необходимо заново проверить применимость уже изученных
ценностей. В следующем эксперименте сравнивается влияние предварительно
обученных знаний на исходную точку исследования.
\end{quote}

\textbf{Предварительно обученная LLM: представление вычислений и
действий посредством языка.}

Перед вызовом инструмента LLM может сгенерировать фрагмент промежуточных
вычислений либо сразу выдать вызов. Оба варианта представлены
последовательностями токенов, но выполняют разные функции: промежуточные
вычисления изменяют контекст, доступный модели на следующих шагах, а
выполненный вызов инструмента изменяет или считывает внешнюю среду. Во
время обучения эти выходные данные можно записывать в соответствии с
протоколом и оценивать с учётом результата задачи и вычислительных
затрат. Таким образом, «как долго думать, о чём думать и когда вызывать
инструмент» становится набором вариантов, доступных для обучения одной
стратегии. Рассматриваемый далее AdaptThink посвящён именно бюджету
рассуждений.

Последовательности переменной длины образуют более крупное комбинаторное
пространство. Модель может писать запросы, программы, параметры JSON и
многошаговые рассуждения, тогда как длина последовательности
ограничивается бюджетом и контекстом. Предварительное обучение даёт
поиску исходную точку: математические материалы часто содержат анализ
условий и вычисления, а код и обсуждения --- поиск проблемы, внесение
изменений и проверку. Модель может использовать эти структуры для
генерации возможных путей, а затем корректировать их на основе обратной
связи от среды. Благодаря имеющимся знаниям полезные комбинации действий
с самого начала выбираются с большей вероятностью, что снижает
потребность в слепых пробах среди множества нерелевантных путей.

Например, задача «пересчитать сумму в иностранной валюте в доллары»
требует получить обменный курс и выполнить вычисление. Благодаря
предварительному обучению модель сравнительно легко генерирует такие
взаимосвязанные шаги. Постобучение может дополнительно скорректировать
вероятности путей с учётом источника курса, результата вычислений и
ограничений задачи. Языковые знания здесь задают исходную стратегию, а
вознаграждение предоставляет обратную связь по задаче.

На \crossreflink{fig:8-5}{рисунке 8-5} сопоставлены две постановки из
этого раздела: табличный Q-learning начинает с заданного кодирования
состояний и правил исследования, а LLM-агент --- с предварительно
обученной языковой стратегии и организует в контексте наблюдения от
инструментов и накопленный опыт.

\begin{figure}
\hypertarget{fig:8-5}{%
\centering
\includegraphics{images/revised/fig8-5.pdf}
\caption{Рисунок 8-5. Сравнение классического RL и современного
LLM-агента}\label{fig:8-5}
}
\end{figure}

Предварительное обучение также позволяет модели интерпретировать новые
инструкции и корректировать поведение в задаче на основе небольшого
числа демонстраций. Инструменты связывают языковой вывод с поиском,
выполнением кода и бизнес-операциями. При проектировании таких агентов
необходимо не только определить интерфейсы инструментов и разрешения на
выполнение, но и организовать обратную связь по результатам задачи,
затратам на выполнение и ключевым этапам процесса.

Как применить «предсказание следующего токена» к действиям, если модель
должна управлять роботизированной рукой? OpenVLA обучалась примерно на
970 тысячах реальных демонстраций роботов. Для одного шага демонстрации
на вход подаются текущий кадр с камеры и инструкция задачи, а целью
обучения служит действие, которое демонстратор выполняет следующим.
Например, горизонтальное перемещение роботизированной руки изначально
представляет собой непрерывное числовое значение. OpenVLA делит каждую
размерность действия на 256 диапазонов и одним токеном обозначает
диапазон, в который попало значение. Благодаря этому модель может
последовательно предсказывать токены действий и вычислять функцию потерь
только в этих позициях.\footnote{Kim, Moo Jin et al., ``OpenVLA: An
  Open-Source Vision-Language-Action Model'', 2024. arXiv:2406.09246.}

Этот пример также показывает, как различать клонирование поведения и
обучение с подкреплением: нужно определить, какой сигнал вызывает
обновление модели. OpenVLA обучается на записях демонстраций: чем выше
присвоенная моделью вероятность токенов демонстрируемого действия, тем
ниже функция потерь. Это относится к дообучению с учителем. В
SimpleVLA-RL из эксперимента 8-13 стратегия будет фактически пробовать
действия в среде, а затем корректировать будущий выбор на основе
вознаграждения за задачу.

\begin{quote}
\textbf{От алгоритма и среды к априорным знаниям}

В статье «The Second Half» Шуньюй Яо рассматривает развитие RL с точки
зрения алгоритмов, сред и априорных знаний.\footnote{Shunyu Yao, ``The
  Second Half'', 10 апреля 2025 года.} На этой основе автор обобщает
изменение исследовательского фокуса в виде трёх этапов: в 2015--2016
годах основное внимание уделялось совершенствованию алгоритмов, в
2016--2018 годах возросло значение построения сред, а начиная с 2018
года постепенно пришло осознание ценности априорных знаний, полученных
посредством предварительного обучения. Gym предоставил единый интерфейс
сред, Universe и World of Bits исследовали компьютерные и веб-задачи, а
такие проекты, как Dota 2, продемонстрировали возможность обучения в
сложных средах. Затем GPT-2 и GPT-3 показали накопленные посредством
предварительного языкового обучения знания и выразительные возможности,
а WebGPT, ChatGPT и другие системы стали в ещё большей степени
использовать эту исходную точку. Необходимые для задачи априорные знания
можно сначала получить посредством языкового моделирования, а затем
передать RL для оптимизации конкретного поведения.

Автор разделяет следующий порядок инженерных приоритетов:
\textbf{априорные знания, среда, алгоритм}. Сначала следует проверить,
какими знаниями и возможностями обладает модель, затем построить
подходящую среду и обратную связь и лишь после этого выбрать
соответствующий метод оптимизации. На
\crossreflink{fig:8-6}{рисунке 8-6} элементы расположены в порядке
появления исследовательского интереса к ним, однако при работе над
проектом можно двигаться в обратном порядке, начиная с уже имеющихся
возможностей.

\begin{figure}
\hypertarget{fig:8-6}{%
\centering
\includegraphics{images/revised/fig8-6.pdf}
\caption{Рисунок 8-6. Изменение направлений исследований обучения с
подкреплением}\label{fig:8-6}
}
\end{figure}
\end{quote}

\begin{quote}
\textbf{Эксперимент 8-2 ★★: сравнительное исследование традиционного RL
и LLM-агента}

В одной и той же игре с поиском сокровищ сравниваются табличный
Q-learning и LLM-агент. В конфигурации LLM используется Kimi K3, а в
контексте хранится не более 50 записей опыта. На каждом шаге оба
варианта выбирают одно из действий, предложенных средой. Два пути на
\crossreflink{fig:8-7}{рисунке 8-7} показывают, как этот выбор выполняют
соответственно Q-таблица и языковая модель.

\begin{figure}
\hypertarget{fig:8-7}{%
\centering
\includegraphics{images/revised/fig8-7.pdf}
\caption{Рисунок 8-7. Сравнение архитектур Q-learning и LLM-агента в
игре с поиском сокровищ}\label{fig:8-7}
}
\end{figure}

В прилагаемом журнале запуска Kimi K3 агент прошёл первую сессию за 17
шагов. Сначала он подобрал ржавый меч, обнаружил, что северная дверь
заперта, отправился в кладовую и получил красный ключ и магический
кристалл. На девятом шаге он объединил меч с кристаллом и создал
серебряный меч. Когда на двенадцатом шаге агент попытался использовать
ключ, среда сообщила: «При необходимости он будет использован
автоматически». Лишь на тринадцатом шаге, когда агент пошёл на север,
дверь действительно открылась. Победив стражника серебряным мечом на
четырнадцатом шаге, он ещё раз попытался использовать меч, но среда
сообщила, что сейчас это невозможно. На семнадцатом шаге агент получил
сокровище. Читая этот журнал, можно закрывать следующую строку и сначала
пытаться угадать по текущей обратной связи, какой шаг ещё необходим.

Сопоставив каждое действие с наблюдением в соответствующий момент, можно
проанализировать, как модель использует имеющиеся знания о том, что
«ключ открывает дверь», а «оружие нужно против стражника», и как
обратная связь от инструментов изменяет последующий план. Табличный
подход накапливает эти связи посредством посещения пар «состояние ---
действие».

В этом сравнении обучение Q-learning в течение 10000 игровых сессий
заняло 2,12 секунды. После этого отдельно провели оценку жадной
стратегии на 100 сессиях: все они завершились прохождением, в среднем за
12 шагов. Kimi за одну сессию 17 раз вызвала модель, затратив 6 минут 56
секунд. На основании этого автор использует следующий принцип выбора:
когда симуляция дешева и допускает большое число повторений, можно в
полной мере применять исследование традиционного RL; когда каждое
взаимодействие требует ожидания, расходов или реальных действий,
предпочтительнее стратегия, способная использовать уже имеющиеся знания
и тем самым повышать ценность каждого взаимодействия.
\end{quote}

Эти два эксперимента связывают настоящую главу с рассмотренными ранее
инженерными методами: контекст позволяет организовывать опыт в рамках
одного запуска, внешняя память --- сохранять его между задачами, а
постобучение --- изменять стратегию посредством обновления параметров.
Далее сначала будут рассмотрены основы предварительного обучения на
языковых и визуальных данных, а затем --- конкретные обучающие данные и
методы оптимизации.

\hypertarget{ux43eux441ux43dux43eux432ux44b-ux43fux440ux435ux434ux432ux430ux440ux438ux442ux435ux43bux44cux43dux43eux433ux43e-ux43eux431ux443ux447ux435ux43dux438ux44f-ux43cux43eux434ux435ux43bux435ux439-ux43dux435ux43eux431ux44fux437ux430ux442ux435ux43bux44cux43dux43eux435-ux447ux442ux435ux43dux438ux435}{%
\section{\texorpdfstring{Основы предварительного обучения моделей
\texttt{{[}необязательное\ чтение{]}}}{Основы предварительного обучения моделей {[}необязательное чтение{]}}}\label{ux43eux441ux43dux43eux432ux44b-ux43fux440ux435ux434ux432ux430ux440ux438ux442ux435ux43bux44cux43dux43eux433ux43e-ux43eux431ux443ux447ux435ux43dux438ux44f-ux43cux43eux434ux435ux43bux435ux439-ux43dux435ux43eux431ux44fux437ux430ux442ux435ux43bux44cux43dux43eux435-ux447ux442ux435ux43dux438ux435}}

Предварительное обучение формирует исходные представления и
распределение генерации модели, а на последующих этапах обучения её
параметры продолжают корректироваться относительно этой отправной точки.
Чтобы понять, как она возникает, сначала рассмотрим, как небольшая
языковая модель обучается на тексте, а затем --- как к ней подключается
визуальный ввод.

Сначала текст проходит \textbf{токенизацию (Tokenization)} и
преобразуется в дискретные токены. Например, фраза «Я люблю
программировать» условно может быть разделена на три единицы: «Я»,
«люблю» и «программировать»; фактическое разбиение определяется
токенизатором. Во время обучения предшествующие токены используются как
контекст для предсказания последующих, после чего вычисляются потери и
обновляются параметры. На \crossreflink{fig:8-8}{рисунке 8-8} обучающее
предложение сначала скрывается после слова «среда», предлагая угадать
следующее слово, а затем открывается исходный текст. При сдвиге этой
позиции вперёд получается следующая задача на предсказание.

\begin{figure}
\hypertarget{fig:8-8}{%
\centering
\includegraphics{images/revised/fig8-8.pdf}
\caption{Рисунок 8-8. Предсказание следующего токена при предварительном
обучении}\label{fig:8-8}
}
\end{figure}

После базового языкового моделирования можно добавить демонстрации
инструкций, чтобы обучить модель формату ответов и выполнению задач.
Можно также использовать данные о предпочтениях и с помощью DPO и других
методов повышать вероятность более одобряемых ответов.

\begin{quote}
\textbf{Эксперимент 8-3 ★★: обучение LLM с нуля --- сила алгоритмических
улучшений}

На примере MiniMind 2 примерно со 100 миллионами параметров на
потребительском GPU выполняются предварительное обучение, SFT и
оптимизация предпочтений. Сравните базовую конфигурацию с конфигурацией,
в которую добавлены QK Norm и оптимизатор Muon, и проследите за
снижением потерь, временем достижения целевого качества и результатами
генерации. Согласно сопутствующей сводке, скорость сходимости
увеличивается примерно втрое, а полное обучение занимает около 14 часов
и стоит примерно 34 доллара.

\emph{Уточнение русского издания: В отчёте указаны восемь RTX 4090:
около 6 часов на предобучение и ещё 8 часов на SFT. При QK Norm и Muon
порог функции потерь 3 достигнут за 12 шагов вместо 36. Это втрое меньше
шагов до выбранного порога, а не измеренное трёхкратное сокращение
времени всего обучения.
\href{https://github.com/bojieli/ai-agent-book/tree/c3352738f4b6fe42fe34e3cf6a79bcb424a133b8/chapter8/MiniMind-pretrain/validation/runs/exp8-3-training-report-20260731-v1}{Источник}.}

QK Norm корректирует числовой масштаб запросов и ключей в механизме
внимания, а Muon изменяет способ оптимизационного обновления
соответствующих параметров.

На разных этапах наблюдайте за ответами модели на один и тот же набор
вопросов: уже после предварительного обучения она может отвечать на
такие вопросы, как «Какая гора самая высокая в мире?»; после SFT
улучшаются следование инструкциям и организация вывода; оптимизация
предпочтений продолжает влиять на фактологические ответы и формулировки.
\end{quote}

\begin{quote}
\textbf{Эксперимент 8-4 ★★: самостоятельное обучение VLM}

На \crossreflink{fig:8-9}{рисунке 8-9} показано соединение «визуальный
энкодер --- проекционный слой --- языковая модель». Визуальный энкодер,
например CLIP, извлекает признаки изображения, проекционный слой
преобразует их к размерности и представлению, необходимым на входе
языковой модели, а языковая модель генерирует описание или ответ с
учётом этих входных данных.

\begin{figure}
\hypertarget{fig:8-9}{%
\centering
\includegraphics{images/revised/fig8-9.pdf}
\caption{Рисунок 8-9. Архитектура визуально-языковой модели
(VLM)}\label{fig:8-9}
}
\end{figure}

На первом этапе визуальный энкодер и LLM замораживаются, а новый
проекционный слой обучается таким образом, чтобы признаки изображения
могли участвовать в языковом предсказании. На втором этапе LLM
размораживается в соответствии с конфигурацией и проходит SFT на
высококачественных данных «изображение --- описание» или «вопрос ---
ответ». Отдельно проследите за изменениями в распознавании объектов,
описании деталей, определении отношений и языковых способностях.

Эта архитектура повторно использует существующие визуальную и языковую
модели и сначала сосредоточивает обучение на модуле межмодального
соединения, поэтому при ограниченном бюджете её стоит попробовать в
первую очередь. Замораживание весов сокращает объём обучаемого
состояния, а впоследствии соответствующие модули можно разморозить в
зависимости от требований задачи. Если описание содержит ошибки, сначала
проверьте, попала ли информация об объекте в визуальное представление.
Если информация уже присутствует, но модель всё равно не может её
использовать, основное внимание следует уделить выразительной
способности проекционного слоя, данным для выравнивания и надзору за
языковой генерацией. Если далее преобразовать вывод в представление
действий и согласовать его с интерфейсом управления, получится
архитектура VLA из \crossreflink{chap:6}{главы 6}.
\end{quote}

Эти два эксперимента показывают, как конфигурация оптимизации и
межмодальное соединение влияют на обучение. При ограниченном бюджете
автор рекомендует сначала сравнить преимущества алгоритмических и
архитектурных улучшений и лишь затем решать, следует ли увеличивать
масштаб модели. Когда существующую модель адаптируют к новому языку или
предметной области, обучение также можно продолжить с весов
предварительно обученной модели. В следующем разделе рассматривается
выбор данных для такого обучения.

\emph{Уточнение русского издания: Рекомендация автора не означает, что в
этом опыте измерялось преимущество перед увеличением масштаба.
Дополнительная слепая оценка относится к восьми сравнениям: семь
предпочтений в пользу QK Norm с Muon и одна ничья.
\href{https://github.com/bojieli/ai-agent-book/tree/c3352738f4b6fe42fe34e3cf6a79bcb424a133b8/chapter8/MiniMind-pretrain/validation/runs/exp8-3-training-report-20260731-v1}{Источник}.}

\hypertarget{mid-training-ux432ux43eux441ux43fux43eux43bux43dux435ux43dux438ux435-ux437ux43dux430ux43dux438ux439-ux438-ux431ux430ux437ux43eux432ux44bux445-ux441ux43fux43eux441ux43eux431ux43dux43eux441ux442ux435ux439}{%
\section{Mid-training: восполнение знаний и базовых
способностей}\label{mid-training-ux432ux43eux441ux43fux43eux43bux43dux435ux43dux438ux435-ux437ux43dux430ux43dux438ux439-ux438-ux431ux430ux437ux43eux432ux44bux445-ux441ux43fux43eux441ux43eux431ux43dux43eux441ux442ux435ux439}}

Под \textbf{Mid-training} в этой главе понимается дополнительный этап
обучения языковой модели на целевом распределении данных поверх уже
существующей модели. Большое количество токенов из документов, кода или
выкладок используется в качестве обучающих сигналов, причём корпус
целевой предметной области можно смешивать с общим корпусом.
Исследования DAPT/TAPT показали, что дополнительное предобучение на
корпусе предметной области или данных, связанных с задачей, улучшает
результаты в последующих задачах.\footnote{Gururangan, Suchin et al.,
  ``Don't Stop Pretraining: Adapt Language Models to Domains and
  Tasks'', ACL, 2020.}

При выборе данных следует исходить из пробелов двух типов. Первый ---
недостаточный охват знаний и языка: например, модель мало знакома с
внутренними протоколами компании, терминологией целевого языка или
определённой кодовой базой. Второй --- пробелы в базовых способностях,
необходимых для задачи: например, отслеживать состояние в длинном
контексте, объединять связи между документами или выполнять длинные
выкладки в коде. Сначала следует разделить проблемы понимания, знаний,
формата и действий по типам ошибок, а затем определить, какие материалы
необходимо добавить.

После снижения функции потерь на текстах предметной области следует
также проверить, способна ли модель применять усвоенные знания в
реальных задачах. На использование знаний влияют порядок дополнительного
предобучения и обучения на инструкциях, организация вопросов и ответов,
а также форма материалов.\footnote{Jiang, Zhengbao et al.,
  ``Instruction-tuned Language Models are Better Knowledge Learners'',
  ACL, 2024.} Если для задач предметной области модели явно недостаёт
базовых способностей, автор рекомендует сначала расширить охват с
помощью корпуса предметной области, а затем посредством примеров SFT
научить модель применять эту информацию в задачах. Если модель иногда
решает задачу правильно, но доля успехов остаётся низкой, можно
переходить к RL.

\hypertarget{ux43aux430ux43a-ux444ux43eux440ux43cux438ux440ux43eux432ux430ux442ux44c-ux434ux430ux43dux43dux44bux435-ux434ux43bux44f-mid-training}{%
\subsection{Как формировать данные для
Mid-training}\label{ux43aux430ux43a-ux444ux43eux440ux43cux438ux440ux43eux432ux430ux442ux44c-ux434ux430ux43dux43dux44bux435-ux434ux43bux44f-mid-training}}

\textbf{Шаг 1. Определить необходимые материалы по распределению
ошибок.} Отдельно оцените результаты по темам, языкам, типам документов,
шаблонам кода и длине контекста. Сопоставляя
\href{mailto:Pass@1}{\nolinkurl{Pass@1}},
\href{mailto:Pass@k}{\nolinkurl{Pass@k}} и примеры ошибок, определите,
при каких условиях модели не хватает применимого решения. Проблемы с
форматом параметров можно устранить с помощью примеров работы с
интерфейсом, тогда как пробелы в понятиях и связях предметной области
требуют соответствующего корпуса.

\textbf{Шаг 2. Организовать корпус высокой плотности.} Исходные
документы дают термины и связи между фактами, репозитории кода ---
структуру модулей и зависимости, а учебные выкладки и синтетические
пояснения раскрывают неявные шаги. Примеры, охватывающие несколько
документов, связывают разрозненную информацию. После фильтрации по
качеству, дедупликации и изоляции оценочного набора из этих материалов
формируются обучающие последовательности.

\textbf{Шаг 3. Смешивать данные в соответствии с целевыми
способностями.} Естественные длинные тексты можно брать из книг,
документов и репозиториев кода. Для развития поиска, многошагового
рассуждения, следования инструкциям, агрегирования и статистической
обработки можно добавить данные задач с выкладками. Для AI-агентов можно
добавить траектории планирования, выбора инструментов, составления
параметров, отслеживания состояния и восстановления после ошибок. Эти
материалы могут поступать из существующих наборов данных или
генерироваться моделью-учителем, но перед включением в обучение их
необходимо отфильтровать по требованиям к содержанию и поведению.

\textbf{Шаг 4. Позволить модели продолжать практиковаться в уже
освоенных задачах.} Одна группа примеров сохраняет короткие тексты,
общий корпус и исходные задачи. В другой группе те же задачи помещаются
в более длинный контекст, ключевая информация переносится в другие
позиции и добавляются отвлекающие элементы. При составлении смеси из
этих двух групп можно ориентироваться на распределение исходных
обучающих данных или испытать несколько соотношений на открытых
корпусах, таких как FineWeb-2, а затем скорректировать их по результатам
оценки.

\textbf{Шаг 5. Принимать решение о продолжении или остановке по
нескольким метрикам.} Одновременно отслеживайте обучающую функцию
потерь, результаты на отложенных задачах предметной области, общие
способности, сохранение исходного навыка следования инструкциям и долю
успехов в целевых задачах. Если результаты в предметной области
улучшаются, а общие способности снижаются, сравните разные пропорции
смеси, скорости обучения и объёмы обучения. Если функция потерь
снижается, а результаты в целевых задачах не меняются, проверьте,
способны ли обучающие материалы научить модель решать такие задачи и не
требуется ли добавить примеры применения знаний в реальных задачах.

Предположим, модель умеет находить заданный номер заказа в очень длинном
журнале заказов. Если вместо этого спросить: «Какой клиент неоднократно
оформлял заказы, а затем отменил их все?», ей придётся найти все заказы
одного клиента и по отдельности сопоставить каждый из них с записью об
отмене. В обеих задачах можно использовать один и тот же журнал, но
проверяемые способности будут различаться. RULER посредством
синтетических задач позволяет независимо управлять длиной входных данных
и сложностью задачи, отдельно проверяя поиск одной записи, многошаговое
отслеживание и агрегирование. Авторы оценили 17 моделей: многие из них
показали почти максимальный результат в простом поиске, однако по мере
увеличения длины входных данных совокупные результаты заметно
снижались.\footnote{Hsieh и др., RULER, 2024.}

Чтобы проверить, проявляется ли это различие на реальных материалах,
можно дополнительно использовать задачи LongBench v2 с несколькими
документами, длинными диалогами или репозиториями кода. Чтобы отдельно
проверить, способна ли модель запоминать требования к формату вывода и
другие инструкции, можно взять проверяемые инструкции из IFEval и
разместить их в разных позициях длинного материала. При сравнении двух
моделей позиции инструкций и отвлекающее содержимое должны быть
одинаковыми; кроме того, следует указать, что тест составлен
самостоятельно. Наконец, с помощью сквозных задач для агентов необходимо
проверить, могут ли эти способности работать совместно. При разработке
оценочного набора можно последовательно проверить следующие
аспекты.\footnote{Bai и др., LongBench v2, ACL 2025. В статье
  представлены шесть типов задач, включающих в общей сложности 503
  вопроса с вариантами ответа.}

\begin{itemize}
\tightlist
\item
  \textbf{Позиционирование и поиск}: способность находить одну или
  несколько ключевых частей информации, расположенных в начале, середине
  или конце.
\item
  \textbf{Связи и рассуждение}: отслеживание связей между абзацами и
  документами, разрешение противоречий и объединение доказательств.
\item
  \textbf{Агрегирование и статистика}: подсчёт, группировка, сортировка,
  сравнение и выявление тенденций в больших таблицах и длинных журналах.
\item
  \textbf{Следование инструкциям}: выполнение нескольких требований,
  разрешение конфликтов, соблюдение условий процесса и формата вывода.
\item
  \textbf{Длинные цепочки рассуждений}: длинные цепочки зависимостей в
  задачах по математике, логическому выводу и коду.
\item
  \textbf{Базовое поведение агента}: декомпозиция задач, выбор
  инструментов, формирование параметров, сохранение состояния и
  восстановление после неудач.
\end{itemize}

Длинные рассуждения и многораундовые вызовы инструментов постоянно
увеличивают контекст, однако ранние инструкции, факты и состояние задачи
должны по-прежнему учитываться в последующих решениях. Если модель не
может находить информацию в длинных текстах или соблюдать ограничения,
агент с длинным горизонтом не сможет работать надёжно. Перед увеличением
бюджета обучения следует сначала сравнить соотношения данных и объёмы
обучения на конфигурациях небольшого масштаба. Для часто обновляемых
фактов, требующих указания источников, автор отдаёт предпочтение RAG;
дополнительное обучение с записью в параметры следует рассматривать для
стабильных и масштабных знаний и базовых способностей предметной
области.

\begin{quote}
\textbf{Эксперимент 8-5 ★★: дополнительное предобучение для освоения
нового языка}

За основу берётся Mistral 7B v0.3, а способность работать с корейским
повышается с помощью материалов корейской Wikipedia и других источников.
Сначала обучение проводится на смеси примерно из 80 \% корейских и 20 \%
английских данных, а затем выполняется SFT на корейских данных с
инструкциями. Изменения в языковом моделировании и диалоговом поведении
наблюдаются отдельно.

Корейские данные расширяют охват целевого языка, а английские помогают
сохранить исходные способности. На каждом этапе модель проверяется на
одной и той же группе корейских и английских задач, не участвовавших в
обучении. В сопутствующих результатах слепая оценка на финальном этапе
для корейского языка улучшается, однако способности модели к английскому
снижаются. Это задаёт конкретное направление для следующей итерации:
увеличить долю английских обучающих материалов или снизить интенсивность
обновления.

Раздельный учёт целевого языка, исходного языка и качества фактов
позволяет оценить выгоды и издержки. На этапе SFT дополнительно
проверяется, способна ли модель формулировать ответы в соответствии с
инструкциями на корейском. Таким образом, влияние дополнительного
предобучения и обучения на инструкциях можно наблюдать раздельно.
\end{quote}

\hypertarget{sft-ux434ux43eux43eux431ux443ux447ux435ux43dux438ux435-ux441-ux443ux447ux438ux442ux435ux43bux435ux43c}{%
\section{SFT (дообучение с
учителем)}\label{sft-ux434ux43eux43eux431ux443ux447ux435ux43dux438ux435-ux441-ux443ux447ux438ux442ux435ux43bux435ux43c}}

Материалы предметной области позволяют модели несколько раз встретить
определённый интерфейс, однако, чтобы научить её исправлять параметры
после неудачного вызова, можно сразу показать соответствующий пример.
При SFT модели задают входные данные задачи и обучают её токен за
токеном генерировать ответ или действие из демонстрации. Рассматривая
\crossreflink{fig:8-10}{рисунок 8-10}, сначала определите, какому
поведению нужно научить модель, а затем решите, как собирать
демонстрации. Вместе с ними в обучение попадут факты, тон и формат
вывода.

\begin{figure}
\hypertarget{fig:8-10}{%
\centering
\includegraphics{images/revised/fig8-10.pdf}
\caption{Рисунок 8-10. Конвейер дообучения с учителем
(SFT)}\label{fig:8-10}
}
\end{figure}

Данные обычно поступают из трёх источников: эксперты определяют ключевые
варианты поведения и границы, модель-учитель массово создаёт
демонстрации, а валидатор отбирает корректные результаты из выборок
модели. Эти подходы можно сочетать. Например, сначала эксперты
составляют небольшое число демонстраций вызова инструментов, затем
учитель изменяет условия задачи и создаёт дополнительные примеры, после
чего проверяются параметры, результаты выполнения и бизнес-ограничения.

\textbf{Одно из применений SFT --- научить модель приёмам, которые
приходится многократно описывать в промптах и Skill.} Возьмём в качестве
примера оценивание памяти из \crossreflink{chap:7}{главы 7}. Допустим,
руководство по оцениванию вместе с прецедентами занимает десять тысяч
токенов. Опираясь на полное руководство, учитель оценивает набор ответов
и показывает, как обрабатывать совпадающие имена людей, перепутанные
отношения, недостаточность доказательств и другие ситуации. При обучении
ученика во входных данных оставляют только краткое название критерия,
материалы текущей истории и оцениваемый ответ, а целевым выводом служит
выбранная учителем категория. После развёртывания достаточно
предоставлять короткий промпт и текущие материалы, экономя при каждом
инференсе десять тысяч токенов неизменных инструкций. Тем же способом
для многократно вызываемого Skill можно показать, когда искать
информацию и как выбирать инструменты, а затем дистиллировать эти
устойчивые действия в модель.

При подготовке данных сначала необходимо определить распределение задач,
схему вывода и способ приёмки. После проверки демонстраций выполняются
удаление дубликатов, настройка пропорций разных категорий примеров и
разделение набора данных. Количество и охват примеров следует
проектировать совместно: если несколько тысяч раз повторить одну и ту же
формулировку вопроса, будет многократно закрепляться один паттерн;
помочь модели выбирать подходящее поведение для разных входных данных
могут лишь вариации ключевых условий. Автор сначала тщательно
дорабатывает набор демонстраций, охватывающих ключевые ситуации, и
только затем увеличивает его размер. При SFT целевой вывод изучается
токен за токеном, поэтому ошибочные факты, лишние шаги и неуместные
привычки также попадут в параметры.

\begin{quote}
\textbf{Эксперимент 8-6 ★★★: SFT речи --- от «копирования голоса» к
«моделированию паралингвистики» \texttt{{[}Расширенный\ эксперимент{]}}}

На примере Orpheus и Sesame CSM сравните, как в генерации речи участвуют
условия, задающие говорящего, контекст диалога и маркеры
выразительности. Orpheus представляет речь в виде токенов и может
использовать референсную пару «текст --- фрагмент речи» для генерации
соответствующего тембра, а также поддерживает маркеры выразительности
\texttt{\textless{}laugh\textgreater{}},
\texttt{\textless{}sigh\textgreater{}} и другие. Sesame CSM принимает
текст, идентификатор говорящего и имеющиеся фрагменты аудио, после чего
генерирует продолжение речи.\footnote{Документация Orpheus-TTS описывает
  маркеры экспрессии и контекстное обусловливание «текст --- речь»;
  интерфейс Sesame CSM генерирует продолжение речи на основе текста,
  данных о говорящем и аудиофрагментов.}

При создании обучающих примеров отдельно сохраняйте текст, говорящего,
аудио и условия выразительности. Для обучения согласованности
используйте материал одного говорящего из разных фраз, а для обучения
выразительности --- смех, вздохи и паузы в различных контекстах.
Сравните вывод без референса, с референсом и с заданными условиями
выразительности, а затем оцените точность, естественность,
эмоциональность и тембр.

Данные должны охватывать достаточное количество сочетаний говорящих,
типов предложений и эмоций. Если говорящих слишком мало, вывод может
тяготеть к одному тембру. Если определённый маркер всегда связан с одним
и тем же фрагментом звука, выражение может стать повторяющимся и
механическим.
\end{quote}

\begin{quote}
\textbf{Эксперимент 8-7 ★★★: многоязычное мышление --- генерация
процесса рассуждения на указанном языке
\texttt{{[}Расширенный\ эксперимент{]}}}

Проведите SFT модели gpt-oss-20b, задавая язык процесса рассуждения с
помощью инструкции. Например, \texttt{reasoning\ language:\ German}
требует генерировать текст рассуждения на немецком языке. Обучающие
примеры представлены на английском, испанском, французском и других
языках, а китайский оставлен для этапа тестирования.

В сопровождающем эксперименте после обучения значение \texttt{Chinese}
позволяет модели генерировать процесс рассуждения на китайском языке.
Можно дополнительно сравнить три результата: соблюдается ли требование к
языку, верен ли итоговый ответ, а также как изменяются длина генерации и
затраченное время. Такое межъязыковое поведение обеспечивается совместно
многоязычной основой, заложенной предварительным обучением, и освоенным
при SFT механизмом управления. Сравнение разных языков и базовых моделей
позволяет продолжить анализ их вклада.
\end{quote}

\begin{quote}
\textbf{Эксперимент 8-8 ★★: дистилляция промпта --- воспроизведение
полезных возможностей с меньшими затратами}

Устойчивая бизнес-задача может зависеть от системного промпта объёмом в
несколько тысяч токенов, а учитель при этом может генерировать длинные
рассуждения. При дистилляции промпта учитель сначала создаёт
высококачественный вывод в полных условиях, после чего пары
«пользовательский ввод --- итоговый ответ» образуют данные SFT для
ученика. Ученик использует более короткий промпт и сразу отвечает в
соответствии с целевой конфигурацией.

Уже из первого обучающего примера видно, как учитель Kimi K3 читает
арабский текст с вкраплениями латинских букв, объясняет основания своего
решения и в конце выдаёт код языка \texttt{ar}. Демонстрация для ученика
SmolLM2-135M-Instruct содержит только исходный текст и \texttt{ar}. В
одном и том же наборе из 80 тестовых записей на вход учителю в среднем
поступает около 386 токенов на запрос, а ученику --- около 97. Эти два
значения рассчитаны с помощью соответствующих токенизаторов моделей.

Сопроводительные записи сравнения при распознавании языка позволяют
разделить «правильность ответа» и «соответствие вывода формату». Учитель
создал демонстрации для 160 обучающих примеров и правильно ответил на
все 80 тестовых записей. Средству оценивания не удалось разобрать ни
один вывод ученика без дообучения, поэтому он получил нулевой результат.
После дообучения удалось разобрать все 80 ответов, причём 76 меток
оказались правильными. Сам по себе нулевой результат базовой модели не
позволяет определить, не распознала ли она язык или распознала, но
ответила в неверном формате. Этот тестовый набор содержит 75 разных
фрагментов текста; полностью совпадающих текстов в обучающем и тестовом
наборах нет.

При подготовке данных необходимо также сохранять условия, без которых
ученик не сможет выполнить задачу. Если учитель опирается на актуальные
результаты инструментов, эти результаты можно передать ученику во
входных данных либо сохранить интерфейс инструмента. Систематические
ошибки учителя следует отфильтровывать с помощью независимой проверки.

Когда правила продукта стабильны, распределение входных данных понятно,
а стоимость вызовов высока, автор рекомендует в первую очередь оценить
дистилляцию промпта, чтобы сжать многократно обрабатываемые инструкции и
рассуждения в параметры. Если продукт всё ещё находится на стадии
экспериментов, а требования задачи часто меняются, лучше сохранить
редактируемые промпты и явные рассуждения для быстрых экспериментов и
локализации ошибок.
\end{quote}

\begin{quote}
\textbf{Эксперимент 8-9 ★★★: дистилляция цепочки рассуждений (Chain of
Thought, CoT)}

При дистилляции CoT текст рассуждения учителя и итоговый ответ вместе
используются в качестве обучающей цели, чтобы ученик освоил процесс
разбиения задачи, вычислений, проверки и исправления. Характерный пример
этого подхода --- дистилляция DeepSeek-R1 в модели семейств Qwen и
Llama. Командам, которым нужна автономная и контролируемая модель при
ограниченных затратах на обучение, автор рекомендует в первую очередь
рассмотреть этот путь. Если учитель заметно превосходит ученика в
целевой задаче, он способен предоставить ценные демонстрации. Сначала
можно выбрать среди моделей с открытыми весами, позволяющих получить
полный вывод рассуждений, а затем сравнить возможности учителя и
ученика, а также качество данных. Если интерфейс предоставляет лишь
краткое изложение рассуждений, объём информации, доступной для
дистилляции процесса, также сокращается.

На первом этапе учитель предлагает несколько решений одной задачи. Для
математической задачи сначала проверяется ответ, а для задачи по
программированию запускаются тесты; сохраняются только решения,
прошедшие проверку. Такая фильтрация называется \textbf{выборкой с
отклонением (Rejection Sampling)}. Дообучение с учителем на отобранных
решениях называется \textbf{дообучением на основе выборки с отклонением
(Rejection Sampling Fine-Tuning, RFT)}. Если ученик должен освоить ещё и
процесс рассуждения, необходимо дополнительно проверить промежуточные
шаги и исключить демонстрации, в которых ошибочное рассуждение случайно
привело к правильному ответу.

Для задач с надёжным валидатором автор рекомендует сначала создать
базовый вариант RFT: генерации, фильтрации и SFT достаточно для
формирования процесса обучения, поэтому не требуется реализовывать
оптимизацию стратегии в режиме онлайн и оценку ценности. На основе
такого варианта можно решить, оправданы ли дополнительные затраты на RL.

На втором этапе обучается ученик. Пары вида «вопрос → текст рассуждения
+ итоговый ответ» оформляются в соответствии с диалоговым форматом
выбранной модели. При использовании границ
\texttt{\textless{}think\textgreater{}} и
\texttt{\textless{}/think\textgreater{}} необходимо явно определить,
какие токены участвуют в вычислении функции потерь. Можно начать с
ученика класса 7B и сравнить конфигурации с полными траекториями,
сокращёнными траекториями и обучением только на итоговых ответах.

На третьем этапе проверяется, чему научился ученик. В сопроводительном
эксперименте Kimi K3 решает 24 математические задачи AIME. Для
дообучения Qwen2.5-1.5B-Instruct сохраняются 23 траектории с правильными
ответами. При повторном решении того же набора число правильных ответов
ученика увеличивается с одного до двух, тогда как учитель правильно
решает 23 задачи. После обучения ученик иногда возвращается к проверке
вычислений, однако попарный тест по задачам не выявляет достоверного
повышения точности (\(p=1.0\)).
\end{quote}

\begin{quote}
\textbf{Практика подготовки книги: процесс разработки эксперимента по
дистилляции}

Для первоначальной реализации эксперимента по дистилляции автор
использовал Codex на базе GPT-5.6-Sol, однако после отказа перешёл на
Claude Code на базе Claude Opus 5, который также не смог продолжить
выполнение. Затем автор использовал Kimi K3 для завершения кода и
запуска эксперимента.
\end{quote}

Во всех четырёх примерах применяется обучение с учителем, но изучается
разное: тембр и условия выразительности речи, управление языком при
многоязычном рассуждении, прямые ответы при дистилляции промпта и
процесс решения задач при дистилляции CoT.

\hypertarget{ux441ux438ux43dux442ux435ux437-ux434ux430ux43dux43dux44bux445-ux434ux43bux44f-sft-ux43eux442-ux434ux435ux43cux43eux43dux441ux442ux440ux430ux446ux438ux439-ux43a-ux43eux431ux443ux447ux430ux44eux449ux438ux43c-ux442ux440ux430ux435ux43aux442ux43eux440ux438ux44fux43c}{%
\section{Синтез данных для SFT: от демонстраций к обучающим
траекториям}\label{ux441ux438ux43dux442ux435ux437-ux434ux430ux43dux43dux44bux445-ux434ux43bux44f-sft-ux43eux442-ux434ux435ux43cux43eux43dux441ux442ux440ux430ux446ux438ux439-ux43a-ux43eux431ux443ux447ux430ux44eux449ux438ux43c-ux442ux440ux430ux435ux43aux442ux43eux440ux438ux44fux43c}}

Какие сведения можно изменить, а какие взаимосвязи необходимо сохранить,
чтобы из одной реальной жалобы получить набор обучающих задач? Из
рабочих траекторий извлекаются намерение пользователя, начальное
состояние, доступные инструменты, бизнес-ограничения и типичные ошибки,
а затем из них формируется шаблон задачи. Идентификационные данные,
заказы, файлы и учётные данные следует заменять; если исходная проблема
связана с двумя похожими заказами, в сгенерированном примере также
необходимо сохранить эту требующую различения связь. Только запустив
пример в сбрасываемой среде, можно установить, действительно ли новая
задача сохранила сложность исходной.

Полный процесс выглядит так: \textbf{рабочий случай→шаблон
задачи→синтетический пример→несколько траекторий-кандидатов→проверка
задачи и траекторий→обучающие данные}. Два вида проверки отвечают за
разные аспекты: проверка задачи определяет, достаточно ли условий,
достижима ли цель и верен ли эталонный результат; проверка траектории
последовательно оценивает состояние на момент завершения задачи, порядок
вызовов инструментов и соответствие процесса выполнения
бизнес-ограничениям.

Для исполняемых условий используются тесты, утверждения базы данных или
проверка различий состояний, а открытые аспекты качества коммуникации
оцениваются по рубрике и моделью-оценщиком с калибровкой посредством
ручной выборочной проверки. Совместное использование графа навыков,
исполняемой среды и независимых валидаторов позволяет планомерно
расширять охват
задач.\footnote{Kulikov, Ilia, et al.~\emph{Autodata: An Agentic Data
  Scientist to Create High Quality Synthetic Data.} arXiv:2606.25996,
  2026.}\footnote{Tan, Zelin, et al.~``SKT: Skill-Use Training at Scale
  via Verified Synthetic Data Generation'', 2026. arXiv:2608.02287.}\footnote{Wei,
  Yifan, et al.~``Towards Compositional Generalization of LLMs via Skill
  Taxonomy Guided Data Synthesis'', 2026. arXiv:2601.03676.}\footnote{Zhu,
  Kaijie, et al.~``TermiGen: High-Fidelity Environment and Robust
  Trajectory Synthesis for Terminal Agents'', 2026. arXiv:2602.07274.}\footnote{Hua,
  Zhanbo, et al.~``CLI-Universe: Towards Verifiable Task Synthesis
  Engine for Terminal Agents'', 2026. arXiv:2606.22883.}

Подходящей демонстрацией может быть как успешно завершённая траектория,
так и полный процесс, включающий диагностику и исправление. Для
ошибочных действий следует явно указывать их роль в образце и выбирать
позиции вычисления функции потерь в соответствии с целью обучения.
Записи о неудачах также можно использовать для построения пар
предпочтений, выявления пробелов в покрытии и создания задач на
восстановление. При последующем переходе к RL текущая стратегия заново
сэмплирует траектории и обновляется с использованием обратной связи от
среды.

Способ разделения обучающего и оценочного наборов выбирается в
соответствии с целью: при оценке на том же распределении можно
изолировать конкретные примеры, а при оценке переноса дополнительно
изолируются шаблоны задач, группы клиентов или периоды времени. При
дедупликации проверяются приблизительные совпадения как входных данных,
так и способов решения; эталонные ответы, скрытые тесты и сведения,
предназначенные исключительно для оценки, в соответствии с протоколом
задачи остаются на стороне проверки.

Можно ли после завершения дообучения сверить обучающие материалы по
весам модели? Рассмотрим открытую модель рассуждений s1. Если
преобразовать открытые образцы в идентификаторы токенов и
дедуплицировать их с использованием применявшихся при обучении
токенизации и обработки текста, получится перечень идентификаторов.
Затем можно построчно сравнить матрицы входных эмбеддингов до и после
дообучения и зафиксировать, векторы каких токенов изменились. При
проверке автором все 16 573 идентификатора в двух перечнях попарно
совпали. Для такой проверки необходимо правильно определить базовую
модель и токенизатор, добавить системный промпт в соответствии с
конфигурацией обучения и обрезать длинные образцы; также нужно
убедиться, что при обучении действительно обновлялись входные эмбеддинги
и что они отделены от весов выходного слоя. Например, если LoRA
обучается только в слоях внимания, входные эмбеддинги остаются
неизменными, поэтому составить по ним перечень использованных токенов
невозможно.\footnote{В проекте training-forensics Ли Боцзе зафиксированы
  множества токенов и условия предварительной обработки для s1 и других
  моделей, а также рассмотрено влияние совместного использования весов
  эмбеддингов и сравнения LoRA на этот метод.}

Границы этих доказательств можно увидеть на примере трёх токенов.
Предположим, что после токенизации двух открытых образцов получены
последовательности \texttt{{[}10,\ 20{]}} и \texttt{{[}20,\ 30{]}}, а
после дедупликации --- множество \texttt{\{10,\ 20,\ 30\}}. Теперь
дополнительно добавим последовательность \texttt{{[}30,\ 10{]}} и
повторим вычисление: множество останется совершенно тем же. Даже если
при проверке совпадёт каждый идентификатор, по этому результату всё
равно нельзя определить, использовался ли третий образец, как нельзя и
исключить попадание какой-либо оценочной задачи в обучающий набор.

Рассмотрим описанный в седьмой главе случай «преждевременного
завершения»: в момент, когда агент готов объявить о выполнении задачи,
выделяется префикс траектории, после чего непроверенное заявление о
завершении помечается как \texttt{rejected}, а продолжение, в котором
сначала проверяются оставшиеся цели, запускаются соответствующие тесты и
только затем даётся ответ, --- как \texttt{chosen}. Пары предпочтений
можно использовать для DPO; продолжение, успешно прошедшее проверку,
также можно отдельно использовать как демонстрацию для SFT.

Приведённые далее эксперименты с китайскими фигурными кавычками и
специальными строками дополнительно показывают, как определять цели
обучения с учителем. Задача с китайскими фигурными кавычками исходит из
области действия документа: она требует изменять целевой основной текст,
защищая код, цитаты и JSON; задача со специальной строкой требует
сохранять \texttt{old\_string} в параметрах инструмента с точностью до
байта. Первая проверяет область редактирования и синтаксис документа,
вторая --- точность копирования и соответствие фактическому инструменту.
Общий процесс определения причин ошибок позволяет получить разные цели
обучения.

\hypertarget{ux43aux43eux433ux434ux430-ux432ux44bux431ux438ux440ux430ux442ux44c-mid-training-sft-ux438-rl}{%
\section{Когда выбирать Mid-training, SFT и
RL}\label{ux43aux43eux433ux434ux430-ux432ux44bux431ux438ux440ux430ux442ux44c-mid-training-sft-ux438-rl}}

Если модель не умеет правильно записывать даже формат входных данных
инструмента, то сколько вызовов действительно сработает, если уже сейчас
заставить её раз за разом действовать методом проб и ошибок? Сначала по
\crossreflink{fig:8-11}{рисунку 8-11} можно определить, на каком шаге
она застревает, а затем по таблице 8-4 выбрать обучающие материалы и
способ обратной связи.

\begin{figure}
\hypertarget{fig:8-11}{%
\centering
\includegraphics{images/revised/fig8-11.pdf}
\caption{Рисунок 8-11. Выбор и сочетание Mid-training, SFT и
RL}\label{fig:8-11}
}
\end{figure}

Таблица 8-4. Критерии выбора Mid-training, SFT и RL

\begin{longtable}[]{@{}
  >{\raggedright\arraybackslash}p{(\columnwidth - 6\tabcolsep) * \real{0.2577}}
  >{\raggedright\arraybackslash}p{(\columnwidth - 6\tabcolsep) * \real{0.2296}}
  >{\raggedright\arraybackslash}p{(\columnwidth - 6\tabcolsep) * \real{0.2515}}
  >{\raggedright\arraybackslash}p{(\columnwidth - 6\tabcolsep) * \real{0.2611}}@{}}
\toprule\noalign{}
\begin{minipage}[b]{\linewidth}\raggedright
\hspace{0pt}Наблюдаемое явление
\end{minipage} & \begin{minipage}[b]{\linewidth}\raggedright
\hspace{0pt}Что проверить в первую очередь
\end{minipage} & \begin{minipage}[b]{\linewidth}\raggedright
\hspace{0pt}Возможные методы
\end{minipage} & \begin{minipage}[b]{\linewidth}\raggedright
\hspace{0pt}Основные направления последующей оценки
\end{minipage} \\
\midrule\noalign{}
\endhead
\bottomrule\noalign{}
\endlastfoot
\hspace{0pt}Модель плохо владеет языком предметной области, её понятиями
или базовыми операциями, а многократное сэмплирование почти не даёт
прогресса & \hspace{0pt}Охват корпуса и базовые способности &
\hspace{0pt}Mid-\allowbreak{}\hspace{0pt}training; для доступных через
поиск фактов --- RAG & \hspace{0pt}Задачи предметной области, тестовая
выборка для проверки общих способностей и частично правильные
траектории \\
\hspace{0pt}Содержимое пригодно, но формат, параметры или фиксированный
процесс соблюдаются нестабильно & \hspace{0pt}Демонстрации, интерфейсы и
ограничения & \hspace{0pt}SFT, декодирование с ограничениями или
валидация инструментов & \hspace{0pt}Успешность разбора, результат
задачи и поведение при восстановлении \\
\hspace{0pt}Модель иногда выполняет задачу, но доля успеха мала или она
часто идёт обходным путём & \hspace{0pt}Качество вознаграждения,
исследование и распределение задач & \hspace{0pt}RL & \hspace{0pt}Польза
от выполнения задачи, различия в вознаграждении, независимое обобщение и
стоимость \\
\hspace{0pt}Есть демонстрации или данные о предпочтениях, но
интерактивная среда ещё не создана & \hspace{0pt}Охват
офлайн-\allowbreak{}\hspace{0pt}данных и способ валидации &
\hspace{0pt}SFT, RFT, офлайн-\allowbreak{}\hspace{0pt}оптимизация
предпочтений & \hspace{0pt}Качество базового уровня, смещения в данных и
польза от создания среды \\
\end{longtable}

Работу можно вести в четыре этапа. Сначала установите базовый уровень
для решения с помощью промпта, контекста, инструментов и ограничений.
Если эти методы уже удовлетворяют требованиям задачи, автор отдаёт
предпочтение внешним редактируемым решениям, чтобы снизить стоимость
сопровождения обучения, развёртывания и регрессионного тестирования.
Динамические факты предоставляются посредством поиска и управления
версиями; обучение параметров следует рассматривать для неявных
суждений, паттернов восприятия и стабильного поведения, которые
требуется усваивать многократно.

На втором этапе наблюдайте за успешностью, частичным прогрессом,
результатами разбора и типами ошибок на отложенных задачах. После
фиксации конфигурации сэмплирования
\href{mailto:Pass@1}{\nolinkurl{Pass@1}} и
\href{mailto:Pass@k}{\nolinkurl{Pass@k}} показывают, существуют ли в
текущей стратегии пути, которыми можно воспользоваться; неудачные
траектории помогают различить нехватку знаний, отсутствие обратной связи
и ошибки процесса.

На третьем этапе предоставьте демонстрации поведения, которое требуется
усвоить. JSON Schema, вызовы инструментов, употребление терминов,
уточнение и восстановление можно оформить в виде примеров SFT, а
необходимые знания --- предоставить через соответствующий корпус и
контекст задачи.

На четвёртом этапе оцените условия обратной связи для RL. Если можно
генерировать кандидатов, поддающихся оцениванию, вознаграждение обладает
достаточной различающей способностью, а среда допускает необходимое
количество взаимодействий, можно сравнить пользу от оптимизации
стратегии. Если вознаграждение слишком разрежено, сначала рассмотрите
учебную программу, частичные вознаграждения или более подходящую
начальную стратегию.

На каждом этапе дальнейшие вложения определяются результатами
оценивания: задачи с чётко заданным форматом можно развёртывать после
SFT, при достаточно сильной исходной модели можно сразу применять RL, а
после продолжения обучения на данных предметной области --- повторно
использовать уже имеющиеся способности к выполнению инструкций. Далее на
примере однораундовой задачи подробно рассмотрим различия в поведении,
возникающие при обучении на демонстрациях и на вознаграждении.

Вернёмся к задаче с заказом из рекомендаций по чтению. Передайте модели
тот же фрагмент реального описания инструмента и проверьте, сможет ли
она правильно указать параметры. Если поля в описании изменились,
сначала обновите описание, доступное во время выполнения. Если
назначение полей понятно, но модель раз за разом подставляет
пользовательский текст вместо номера заказа, можно предоставить
демонстрации SFT, включающие запрос данных, уточнение и правильное
заполнение параметров. Если модель уже иногда полностью выполняет
задачу, а разный порядок запросов приводит к разным показателям
успешности и стоимости, появляются условия для сравнения нескольких
траекторий стратегии.

\hypertarget{ux43eux434ux43dux43eux440ux430ux443ux43dux434ux43eux432ux43eux435-ux43eux431ux443ux447ux435ux43dux438ux435-ux441-ux43fux43eux434ux43aux440ux435ux43fux43bux435ux43dux438ux435ux43c-ux441ux43eux43fux43eux441ux442ux430ux432ux43bux435ux43dux438ux435-ux437ux430ux43fux43eux43cux438ux43dux430ux43dux438ux44f-ux438-ux43eux431ux43eux431ux449ux435ux43dux438ux44f}{%
\section{Однораундовое обучение с подкреплением: сопоставление
запоминания и
обобщения}\label{ux43eux434ux43dux43eux440ux430ux443ux43dux434ux43eux432ux43eux435-ux43eux431ux443ux447ux435ux43dux438ux435-ux441-ux43fux43eux434ux43aux440ux435ux43fux43bux435ux43dux438ux435ux43c-ux441ux43eux43fux43eux441ux442ux430ux432ux43bux435ux43dux438ux435-ux437ux430ux43fux43eux43cux438ux43dux430ux43dux438ux44f-ux438-ux43eux431ux43eux431ux449ux435ux43dux438ux44f}}

Однораундовая задача начинается с однократного ввода и заканчивается,
когда модель отправляет ответ и получает вознаграждение. При генерации
ответа модель по-прежнему может рассуждать в несколько шагов, а система
обучения --- позволять ей многократно пробовать решить одну и ту же
задачу. Ограничение взаимодействия одним ответом упрощает наблюдение за
тем, как метод обучения меняет поведение модели. Сначала на примере
AdaptThink рассмотрим, как модель учится распределять вычислительный
бюджет на рассуждения, а затем с помощью GeneralPoints сравним
результаты SFT и RL при изменении правил и визуального оформления.

Чтобы понять процесс обучения в последующих экспериментах, сначала
разберёмся, как алгоритмы RL используют вознаграждение для обновления
модели. Обучение RL в этой главе в основном опирается на
\textbf{градиент стратегии}. Возьмём несколько выборок ответов на один
вопрос: модель генерирует несколько ответов, после чего вознаграждение
каждого из них сравнивается с базовым уровнем. Ответы, превосходящие
базовый уровень, получают положительный сигнал, и обучение повышает
вероятность их генерации; ответы ниже базового уровня получают
отрицательный сигнал, и вероятность их генерации снижается.

Чтобы однократное обновление не увело модель в неверную сторону,
распространённый алгоритм \textbf{PPO} отсекает дополнительный выигрыш в
суррогатной целевой функции, когда отношение вероятностей выходит за
заданный диапазон; фактическая величина обновления также зависит от
скорости обучения и количества итераций обновления. Рассматриваемый
далее «PPO с сетью ценности» обучает оценку ценности, с помощью которой
формируется базовый уровень для сравнения и вычисляется преимущество
действий на каждом шаге. \textbf{GRPO} непосредственно сравнивает
несколько ответов на один вопрос и строит базовый уровень по
вознаграждениям внутри группы, благодаря чему отдельное обучение сети
ценности не требуется.

Следующий псевдокод в стиле Python связывает выборку, вознаграждение,
преимущество и обновление параметров, помогая понять, как информация
передаётся между этапами. В реальном обучении также требуется настроить
параллельную выборку, KL-регуляризацию и оптимизатор.

\begin{Shaded}
\begin{Highlighting}[]
\ControlFlowTok{for}\NormalTok{ prompt }\KeywordTok{in}\NormalTok{ batch:}
\NormalTok{    group }\OperatorTok{=}\NormalTok{ [rollout(policy, env.reset(prompt)) }\ControlFlowTok{for}\NormalTok{ \_ }\KeywordTok{in} \BuiltInTok{range}\NormalTok{(G)]}
\NormalTok{    rewards }\OperatorTok{=}\NormalTok{ [verify(trajectory) }\ControlFlowTok{for}\NormalTok{ trajectory }\KeywordTok{in}\NormalTok{ group]}
\NormalTok{    advantages }\OperatorTok{=}\NormalTok{ normalize\_within\_group(rewards)       }\CommentTok{\# базовая линия GRPO}
\NormalTok{    update(policy, group, advantages)}
\end{Highlighting}
\end{Shaded}

Как PPO обновляет стратегию по группе траекторий разной длины? Сначала
внутри каждой траектории в обратном порядке вычисляется дисконтированный
возврат для каждого шага. Накопление прекращается в конце траектории,
чтобы вознаграждение следующей задачи не попало в текущую. Затем
фактически выбранные моделью действия из всех траекторий объединяются и
для них вычисляются средние потери. Ниже показано одно обновление без
многократной оптимизации по мини-батчам.\footnote{Schulman et al.,
  \emph{Proximal Policy Optimization Algorithms}, 2017,
  arXiv:1707.06347.}

\begin{Shaded}
\begin{Highlighting}[]
\NormalTok{states, actions, returns }\OperatorTok{=}\NormalTok{ [], [], []}
\ControlFlowTok{for}\NormalTok{ trajectory }\KeywordTok{in}\NormalTok{ rollouts:}
\NormalTok{    states.extend(trajectory.states)       }\CommentTok{\# состояние перед каждым выбором действия}
\NormalTok{    actions.extend(trajectory.actions)     }\CommentTok{\# поэлементно соответствует состояниям}
\NormalTok{    returns.extend(discounted\_returns(trajectory.rewards))}

\NormalTok{returns }\OperatorTok{=}\NormalTok{ tensor(returns)}
\NormalTok{values }\OperatorTok{=}\NormalTok{ value\_model(states)}
\NormalTok{advantages }\OperatorTok{=}\NormalTok{ returns }\OperatorTok{{-}}\NormalTok{ stop\_gradient(values)}
\NormalTok{ratio }\OperatorTok{=}\NormalTok{ exp(policy.log\_prob(actions, given}\OperatorTok{=}\NormalTok{states)}
            \OperatorTok{{-}}\NormalTok{ old\_policy.log\_prob(actions, given}\OperatorTok{=}\NormalTok{states))}
\NormalTok{policy\_loss }\OperatorTok{=} \OperatorTok{{-}}\NormalTok{mean(}\BuiltInTok{min}\NormalTok{(}
\NormalTok{    ratio }\OperatorTok{*}\NormalTok{ advantages,}
\NormalTok{    clip(ratio, }\DecValTok{1} \OperatorTok{{-}}\NormalTok{ epsilon, }\DecValTok{1} \OperatorTok{+}\NormalTok{ epsilon) }\OperatorTok{*}\NormalTok{ advantages}
\NormalTok{))}
\NormalTok{value\_loss }\OperatorTok{=}\NormalTok{ mean((values }\OperatorTok{{-}}\NormalTok{ returns) }\OperatorTok{**} \DecValTok{2}\NormalTok{)}
\NormalTok{update(policy, value\_model, policy\_loss }\OperatorTok{+}\NormalTok{ value\_coef }\OperatorTok{*}\NormalTok{ value\_loss)}
\end{Highlighting}
\end{Shaded}

«Относительность» GRPO обеспечивается несколькими ответами на один
вопрос, тогда как здесь PPO использует сеть ценности для оценки
ожидаемого возврата из каждого состояния. \texttt{old\_policy} --- это
стратегия, зафиксированная на момент генерации данной группы траекторий.
Предположим, преимущество действия на некотором шаге равно 1, а текущая
стратегия назначает ему вероятность, в 1,4 раза превышающую вероятность
при выборке. При \texttt{epsilon\ =\ 0.2} вклад этого элемента в
суррогатную целевую функцию составит \texttt{min(1.4,\ 1.2)\ =\ 1.2}.
Дальнейшее повышение вероятности этого действия уже не принесёт
дополнительного выигрыша от данного элемента.

Если преимущество равно −1, направление меняется на противоположное.
Когда отношение вероятностей уже снизилось до 0,6, подстановка в ту же
строку кода даёт \texttt{min(-0.6,\ -0.8)\ =\ -0.8}, поэтому дальнейшее
снижение вероятности также не увеличивает выигрыш. И наоборот, если это
менее удачное действие стало встречаться чаще и отношение вероятностей
выросло до 1,4, выбор меньшего значения даёт
\texttt{min(-1.4,\ -1.2)\ =\ -1.4}: повышение его вероятности
по-прежнему влечёт штраф.

\begin{quote}
\textbf{Эксперимент 8-10 ★★: AdaptThink --- научиться понимать, «когда
не нужно рассуждать»}

Модели рассуждений часто решают задачи с помощью длинного процесса
рассуждения. Для простых вопросов прямой ответ позволяет сократить
затраты на генерацию. AdaptThink предоставляет модели два режима ---
Thinking и NoThinking: в моделях, поддерживающих этот протокол,
NoThinking завершает фрагмент рассуждения с помощью
\texttt{\textless{}think\textgreater{}\textless{}/think\textgreater{}},
после чего сразу генерируется ответ. Цель обучения --- научить модель
выбирать подходящий вычислительный бюджет на рассуждения в зависимости
от задачи.

Метод включает два ключевых компонента. \textbf{Целевая функция с
ограничениями} побуждает модель чаще использовать NoThinking, сохраняя
при этом общую результативность ответов; \textbf{выборка по значимости}
корректирует соотношение выборок двух режимов, а затем с помощью весов
компенсирует различие между распределением выборки и целевой функцией
оптимизации. Исходная модель обычно склоняется к Thinking, поэтому для
получения обучающего сигнала, позволяющего сравнивать два режима,
необходимо целенаправленно увеличить количество примеров NoThinking. Под
проблемой холодного старта здесь конкретно понимается нехватка примеров
NoThinking в начале обучения.

В эксперименте книги использовались восемь GPU H100 80GB. К шагу 300
точность MATH500 изменилась с 81,00\% до 81,80\%, а средняя длина ответа
сократилась примерно с 4 911 до 1 577 токенов; точность GSM8K выросла с
79,68\% до 81,88\%, а длина ответа сократилась примерно с 1 025 до 477
токенов; AIME \href{mailto:mean@16}{\nolinkurl{mean@16}} изменился с
31,46\% до 31,04\%, а длина ответа сократилась примерно с 12 120 до 6
402 токенов. Длина ответа в трёх наборах сократилась приблизительно на
68\%, 53\% и 47\% соответственно.

Доля NoThinking в трёх наборах данных составила соответственно 83,80\%,
84,15\% и 56,25\%. На AIME модель сохранила больший объём вычислений для
рассуждений, тогда как на MATH500 и GSM8K чаще отвечала напрямую. Для
каждой задачи AIME2024 выполнялось 16 выборок, а
\href{mailto:mean@16}{\nolinkurl{mean@16}} обозначает среднюю точность
этих 16 ответов. К шагу 300 средний ответ стал почти вдвое короче, а
точность снизилась примерно на 0,42 процентного пункта.

Шаг 300 был одной из промежуточных проверок во время обучения.
Изначально команда планировала провести 3 140 шагов, но позднее
выполнение аварийно завершилось, и последняя запись проверки осталась на
шаге 410. В этот момент средняя точность AIME2024 составляла 31,88\%, а
средняя длина ответа --- около 5 874 токенов.

AdaptThink можно сочетать с дистилляцией промптов: дистилляция
записывает в параметры многократно используемые инструкции и способы
решения задач, а AdaptThink затем учится определять, сколько вычислений
для рассуждения следует выделить каждой задаче. Вместе они влияют на
качество ответа, задержку и стоимость.
\end{quote}

\begin{quote}
\textbf{Эксперимент 8-11 ★★: GeneralPoints --- сопоставление
«запоминания и обобщения» в однораундовом RL}

Чу и соавторы использовали похожую на «24» карточную игру GeneralPoints
для проверки способности модели к обобщению\footnote{Chu, Tianzhe et
  al., ``SFT Memorizes, RL Generalizes: A Comparative Study of
  Foundation Model Post-training'', 2025. arXiv:2501.17161.}. Число на
каждой из четырёх карт необходимо использовать ровно один раз и с
помощью сложения, вычитания, умножения и деления получить 24. Если
модель научилась решать задачу, сможет ли она по-прежнему правильно
считать после изменения правил карт или мастей? На
\crossreflink{fig:8-12}{рисунке 8-12} показаны два варианта
представления задачи: в GP-L условие задаётся непосредственно текстом, а
в GP-VL модель должна распознать карты.

\emph{Уточнение русского издания: Одна карточная задача не означает одну
попытку: эксперимент допускает последовательные ответы после проверки.
Это важно при сравнении бюджетов вычислений.
\href{https://arxiv.org/html/2501.17161v1}{Источник}.}

\begin{figure}
\hypertarget{fig:8-12}{%
\centering
\includegraphics{images/revised/fig8-12.pdf}
\caption{Рисунок 8-12. Архитектура эксперимента GeneralPoints: схема
обучения и тестирования двух вариантов, GP-L и GP-VL}\label{fig:8-12}
}
\end{figure}

\textbf{Вариант правил}: при обучении J/Q/K имеют значение 10, а при
тестировании --- соответственно 11/12/13. Поэтому тестовый набор
содержит комбинации чисел, не встречавшиеся при обучении, что позволяет
проверить, способна ли модель выполнять вычисления с 11, 12 и 13 по
новым правилам. \textbf{Визуальный вариант}: при обучении используются
чёрные масти (♠♣), а при тестировании --- красные (♥♦), что позволяет
проверить способность модели адаптироваться к изменению внешнего вида
карт.

Эксперимент начинается с Llama-3.2-Vision-11B: сначала с помощью SFT
формируется базовая способность следовать инструкциям, после чего
обучение разделяется на две ветви --- дальнейшее SFT и RL, сравниваемые
при одинаковом вычислительном бюджете. В ветви RL применяется PPO с
сетью ценности. Обе ветви обучаются по правилу J/Q/K=10, после чего
оцениваются на тестовых наборах внутри распределения (ID) и вне
распределения (OOD).

Возьмём \texttt{{[}7,\ J,\ 3,\ 5{]}}: при J=10 можно использовать
\texttt{(J−7)×(5+3)=24}; после изменения значения J на 11 это выражение
даёт 32, поэтому его необходимо заменить на \texttt{(11−5)×(7−3)=24}.
Таким образом, одна и та же группа карт требует от модели прочитать
действующее правило и заново выстроить вычисление.

Два метода обучения показали заметно различающиеся результаты при
тестировании вне распределения. \textbf{OOD по правилам}: для GP-L доля
успеха RL выросла с 11,5\% до 15,0\%, то есть на 3,5 процентного пункта;
результат SFT снизился с 11,5\% до 3,4\%, то есть на 8,1 процентного
пункта. Для GP-VL результат RL вырос на 3,0 процентного пункта, а SFT
снизился на 5,6 процентного пункта. \textbf{Визуальный OOD}: для GP-VL
результат RL вырос с 23,6\% до 41,2\%, то есть на 17,6 процентного
пункта; результат SFT снизился с 23,6\% до 13,7\%, то есть на 9,9
процентного пункта.

Дополнительная проверка результатов распознавания карт позволяет связать
итоговую успешность решения задачи с визуальным восприятием: в ветви RL
точность распознавания росла вместе с долей успешно решённых задач,
тогда как в ветви SFT она снижалась. Исследователи связывают это
различие с целями обучения: подгонка по каждому токену может закреплять
текстовые паттерны демонстраций, тогда как вознаграждение за результат
постоянно требует от модели правильно распознавать карты и выполнять
вычисления. Точность восприятия служит промежуточным показателем для
анализа этого механизма.

Причину предварительного SFT также можно понять через формат вывода.
Использованная в эксперименте базовая модель Llama-3.2-Vision-11B не
могла стабильно генерировать требуемый структурированный вывод, поэтому
валидатор не мог вычислить по нему вознаграждение за задачу и попытка
непосредственного сквозного обучения RL провалилась. SFT сначала
формирует пригодный для синтаксического анализа и оценки вывод, и только
после этого RL может использовать обратную связь от задачи.
Рассмотренный ранее DeepSeek-R1-Zero, напротив, начинал с базовой
модели, уже способной генерировать ответы, пригодные для оценки, и сразу
переходил к RL.

Сколько модель может узнать из обратной связи валидатора? На рисунке 10
в статье отдельно сравнивается чисто текстовая карточная задача GP-L.
Исследователи задали 1, 3, 5 и 10 итераций проверки, отдельно обучив и
оценив соответствующие модели RL. При одинаковом вычислительном бюджете
обучения в группе с десятью проверками доля успеха на новых правилах
выросла относительно начального уровня обучения этой группы на 5,99
процентного пункта, тогда как в группе с одной проверкой --- на 0,48
процентного пункта.

Почему при смещении распределения в этом эксперименте производительность
SFT снижается, а RL показывает лучший результат? Одно из согласующихся с
наблюдениями объяснений состоит в том, что ограниченные данные SFT
закрепили фиксированный паттерн «считать J/Q/K равными 10». Когда при
тестировании J=11, модель по-прежнему считает его равным 10. Ветвь RL,
оптимизируемая по результату, с большей вероятностью закрепляет
стратегию «пересчитывать до получения правильного ответа», поэтому может
применять её и после изменения значения J на 11.

Изменения правил и мастей в GeneralPoints позволяют отдельно наблюдать,
научилась ли модель выполнять вычисления по новым условиям и правильно
ли она распознаёт новые визуальные входные данные. При использованных
здесь модели, обучающих данных и тестовом бюджете ветвь, продолжившая
SFT, ухудшила результаты при тестировании вне распределения, а ветвь с
последующим RL --- улучшила. Прежде чем переносить выводы на собственную
задачу, стоит спросить: освоил ли ученик формат вывода? Что именно
изменилось в новом тесте --- правила, внешний вид входных данных или и
то и другое? Только раздельные контролируемые эксперименты покажут,
следует ли добавить демонстрации, улучшить восприятие или расширить
исследование вариантов с обратной связью по результату.
\end{quote}

\hypertarget{ux430ux43bux433ux43eux440ux438ux442ux43c-rl-ux43eux442-16-rollout-ux434ux43e-ux43eux434ux43dux43eux433ux43e-ux43eux431ux43dux43eux432ux43bux435ux43dux438ux44f-ux43fux430ux440ux430ux43cux435ux442ux440ux43eux432}{%
\section{Алгоритм RL: от 16 rollout до одного обновления
параметров}\label{ux430ux43bux433ux43eux440ux438ux442ux43c-rl-ux43eux442-16-rollout-ux434ux43e-ux43eux434ux43dux43eux433ux43e-ux43eux431ux43dux43eux432ux43bux435ux43dux438ux44f-ux43fux430ux440ux430ux43cux435ux442ux440ux43eux432}}

\textbf{GRPO (Group Relative Policy Optimization)} предложен DeepSeek и
формирует обучающий сигнал на основе относительных результатов
нескольких попыток решения одной задачи. Рассмотрим задачу исправления
кода в стиле SWE-bench: при пустых входных данных \texttt{parser.py}
Python-проекта вызывает \texttt{IndexError}; AI-агент должен исправить
реализацию, не изменяя тестовые файлы. Одну итерацию обучения можно
разделить на четыре этапа.

\textbf{Первый этап: многократные попытки модели политики.} Модель
политики --- это языковая модель, которая обучается в данный момент.
Система помещает один и тот же исходный код и одно и то же описание
проблемы в 16 изолированных песочниц, позволяя модели независимо
выполнить задачу 16 раз. Каждая попытка начинается с чтения кода,
проходит через внесение изменений и тестирование и заканчивается
отправкой результата; вся эта траектория взаимодействия называется одним
\textbf{rollout}. Благодаря случайности семплирования из одной исходной
точки возникают разные пути: одни правильно добавляют проверку
граничного случая, другие лишь перехватывают исключение, третьи изменяют
не тот файл, а некоторые пытаются изменить тесты.

\textbf{Второй этап: вычисление вознаграждения.} После завершения
каждого rollout валидатор применяет патч в чистом окружении и запускает
тесты. Предположим, что из 16 попыток четыре прошли все тесты и не
изменили тестовые файлы, а остальные 12 завершились неудачно. Тогда
четыре успешные траектории получают вознаграждение 1, а остальные 12 ---
вознаграждение 0. На \crossreflink{fig:8-13}{рисунке 8-13} результаты
оценки всех 16 траекторий собраны в одной таблице. В подобных задачах
программирования вознаграждение можно вычислять непосредственно с
помощью тестов и правил. Для целей, которые сложнее выразить тестами,
например читаемости и качества объяснения, можно дополнительно
использовать человеческую оценку или модель вознаграждения.

\textbf{Третий этап: вычисление относительного преимущества.}
Вознаграждение показывает лишь успех или неудачу отдельной траектории, а
\textbf{относительное преимущество} --- насколько она хороша по
сравнению с другими попытками в той же группе. Средняя доля успешных
попыток в этой группе равна 4/16: четыре прошедшие тесты траектории
находятся выше среднего по группе и получают положительное преимущество;
12 неудачных --- ниже среднего и получают отрицательное преимущество.
Если выполнить нормализацию по внутригрупповому стандартному отклонению,
среднее вознаграждение составит 0,25, стандартное отклонение ---
приблизительно 0,433, преимущество успешной траектории --- около 1,73, а
неудачной --- около −0,58. GRPO формирует обучающий сигнал на основе
такого внутригруппового сравнения. Если вознаграждения всех 16
траекторий полностью совпадают, после вычитания среднего по группе
каждая из них получит ноль. Добавление в знаменатель небольшого
стабилизирующего слагаемого позволяет избежать деления на ноль, но после
нормализации результирующие вознаграждения в этой группе всё равно не
указывают, какой траектории следует отдавать предпочтение при обучении.
В таком случае можно проверить, позволяют ли ограничения пути,
процессные вознаграждения или частичный прогресс различить траектории
внутри группы. Например, из двух попыток, не исправивших код, только
одна прошла исходно падавший тест граничного случая. Если этот
промежуточный результат допускает независимую проверку, появляется ещё
один возможный сигнал для обучения.

\textbf{Четвёртый этап: обновление политики с помощью градиентного
спуска.} Программа обучения преобразует относительное преимущество в
функцию потерь, вычисляет градиенты, после чего оптимизатор, например
AdamW или Muon, выполняет градиентный спуск: повышает вероятность
решений, принятых моделью в траекториях с положительным преимуществом, и
снижает вероятность решений из траекторий с отрицательным преимуществом.
В процессе обучения вероятности этих решений многократно корректируются
на множестве задач и rollout. По мере действия оптимизации при
возникновении похожей ошибки последовательность «сначала воспроизвести
проблему, проверить граничные условия, изменить реализацию и запустить
тесты» будет появляться чаще, а действия «замаскировать исключение,
изменить тесты, отправить результат без проверки» --- реже.

\begin{figure}
\hypertarget{fig:8-13}{%
\centering
\includegraphics{images/revised/fig8-13.pdf}
\caption{Рисунок 8-13. Пример исправления кода: проверка и относительное
преимущество 16 rollout}\label{fig:8-13}
}
\end{figure}

Сколько \textbf{итераций обучения (step)} составляют 16 выборок одной
задачи на \crossreflink{fig:8-13}{рисунке 8-13}? Согласно принятому в
этом разделе способу подсчёта, все 16 rollout входят в один пакет. На
\(k\)-м раунде текущая политика выполняет семплирование, после чего
вычисляются вознаграждения, преимущества и градиенты и политика
обновляется один раз; повторное семплирование с новыми параметрами
начинается только на \(k+1\)-м раунде. Обучение в течение 100 steps
означает, что проводится около 100 таких итераций.

На этой основе можно оценить время обучения. Одна сложная задача может
включать десятки раундов вызовов инструментов. При синхронном обучении,
ожидающем завершения всего пакета, длительность этапа семплирования
определяется самой медленной из 16 параллельных траекторий. Предположим,
что она занимает около 2 000 секунд, а вычисление градиентов и
обновление оптимизатора --- ещё 600 секунд. Тогда один step потребует
примерно \(2{,}000+600=2{,}600\) секунд, то есть 43 минуты. Непрерывное
обучение в течение 100 steps с такими же затратами времени займёт около
72 часов.

PPO и GRPO следуют этому замкнутому циклу; основное различие заключается
в том, «с чем проводится сравнение». GRPO напрямую сравнивает несколько
rollout одной задачи и не требует отдельной модели ценности. PPO обучает
модель ценности, которая оценивает, «насколько хорошего результата
обычно можно достичь» на каждом шаге траектории, а затем определяет,
оказалось ли текущее действие лучше этого ожидания, предоставляя тем
самым базовую линию для пошагового распределения заслуг. Оба алгоритма
подавляют чрезмерно большие обновления с помощью таких механизмов, как
clipping суррогатной целевой функции. DPO устроен иначе: он напрямую
обучается на заранее собранных парах предпочтений «лучший ответ ---
худший ответ» и не заставляет текущую политику генерировать эту группу
rollout в режиме онлайн.

В примерах этой главы используются разные схемы обучения: AdaptThink
применяет собственную целевую функцию с ограничениями; GeneralPoints,
V-IRL и ReTool используют PPO, причём у первых двух модель ценности
участвует в оценке преимущества; SimpleVLA-RL и RLVP используют GRPO.

Предположим, что задачи А и Б выполняются по четыре раза: задача А
решена правильно во всех попытках, а задача Б --- неправильно. Все
четыре вознаграждения задачи А равны 1, а задачи Б --- 0. После
отдельного вычитания среднего по группе преимущества в обеих группах
превращаются в 0. Но если есть вычислительные ресурсы для дополнительных
попыток, на какую задачу перспективнее их потратить, чтобы получить
полезные обучающие примеры?

Never Give Up включает этот выбор в процесс семплирования. Если все
попытки задачи А успешны, группа пропускается; если все попытки задачи Б
неудачны, с некоторой вероятностью добавляется ещё одна группа.
Предположим, что во второй группе наконец появляется один правильный
ответ. Тогда при обучении объединяются все восемь ответов из двух групп:
один правильный и семь неправильных, благодаря чему для одной задачи
появляются разные уровни вознаграждения. Асинхронное семплирование
позволяет остальным задачам выполняться как обычно, не дожидаясь
результата задачи Б.\footnote{Michael Noukhovitch и др., \emph{Learning
  to Solve Hard Problems in RL for LLMs by Never Giving Up}, 2026.
  Описанное в статье дополнительное семплирование используется для
  обучения; уже решённые задачи по-прежнему могут снова выбираться в
  последующем обучении.}

\hypertarget{ux43fux43eux447ux435ux43cux443-ux432-rl-ux434ux43bux44f-llm-ux43eux431ux44bux447ux43dux43e-ux43fux440ux435ux434ux43fux43eux447ux438ux442ux430ux44eux442-on-policy}{%
\subsection{Почему в RL для LLM обычно предпочитают
On-Policy}\label{ux43fux43eux447ux435ux43cux443-ux432-rl-ux434ux43bux44f-llm-ux43eux431ux44bux447ux43dux43e-ux43fux440ux435ux434ux43fux43eux447ux438ux442ux430ux44eux442-on-policy}}

\textbf{Online (онлайн)} и \textbf{On-policy (на данных текущей
политики)} описывают два разных аспекта. При онлайн-обучении система во
время работы непрерывно взаимодействует со средой и создаёт данные.
Обучение на данных текущей политики относится к источнику данных и
требует, чтобы поведенческая политика \(\mu\), генерирующая траектории,
совпадала с оптимизируемой политикой \(\pi_\theta\) или была достаточно
близка к ней. Асинхронный кластер может непрерывно выполнять
онлайн-семплирование, но если worker отстаёт на несколько контрольных
точек, возвращаемые им данные характеризуются устареванием политики и
приобретают компонент off-policy (на данных другой политики).

Автор рекомендует сначала создать стабильную базовую конфигурацию с
обучением, близким к on-policy, и лишь затем вводить асинхронное
семплирование или повторное использование старых данных для повышения
пропускной способности. Рассматриваемые в этой главе схемы PPO/GRPO
обычно в начале каждого step выполняют семплирование с помощью последней
политики. Затем PPO выполняет несколько раундов обновлений по minibatch
на одном пакете данных, и текущая политика постепенно отклоняется от
\texttt{old\_policy}, с помощью которой эти данные были сгенерированы.
Отношение вероятностей измеряет это изменение, а целевая функция с
clipping подавляет чрезмерно большие обновления.

Цель градиента политики связана с ожидаемым вознаграждением для текущей
политики \(\pi_\theta\). При использовании данных, выбранных
поведенческой политикой \(\mu\), сначала можно вычислить отношение
вероятностей действия в одном и том же состоянии, а затем построить на
его основе коррекцию по важности:

\[
\rho_t=\frac{\pi_\theta(a_t\mid s_t)}{\mu(a_t\mid s_t)}
=\exp\left(\log\pi_\theta(a_t\mid s_t)-\log\mu(a_t\mid s_t)\right).
\]

Предположим, что после одного и того же префикса вероятность выбора
определённого токена политикой семплирования равна 0,2, а текущая
политика назначает ему вероятность 0,3. Тогда отношение вероятностей на
этом шаге равно 1,5. В старой выборке такой токен встречался реже,
поэтому при оценке по текущей политике вес подобных примеров
соответственно повышается. Если при семплировании и обучении
используется одна и та же политика, до обновления параметров отношение
вероятностей на каждом шаге равно 1.

Повторное использование старых траекторий позволяет выполнять
семплирование и обучение асинхронно, повышая пропускную способность.
Однако чем сильнее различаются две политики, тем вероятнее, что
небольшое число примеров получит очень большой вес. При перевзвешивании
всей траектории отношения вероятностей отдельных шагов также приходится
перемножать. PPO отсекает часть дополнительного выигрыша в суррогатной
целевой функции каждого действия, ослабляя влияние экстремальных
отношений вероятностей; однако состояния, которых старые примеры никогда
не посещали, можно увидеть только с помощью нового семплирования.
Поэтому автор отдаёт предпочтение свежим траекториям.\footnote{Zheng,
  Chujie et al., ``Stabilizing Reinforcement Learning with LLMs:
  Formulation and Practices'', 2025. arXiv:2512.01374.}

\hypertarget{ux447ux438ux441ux43bux435ux43dux43dux44bux435-ux440ux430ux437ux43bux438ux447ux438ux44f-ux43cux435ux436ux434ux443-ux43eux431ux443ux447ux435ux43dux438ux435ux43c-ux438-ux438ux43dux444ux435ux440ux435ux43dux441ux43eux43c}{%
\subsubsection{Численные различия между обучением и
инференсом}\label{ux447ux438ux441ux43bux435ux43dux43dux44bux435-ux440ux430ux437ux43bux438ux447ux438ux44f-ux43cux435ux436ux434ux443-ux43eux431ux443ux447ux435ux43dux438ux435ux43c-ux438-ux438ux43dux444ux435ux440ux435ux43dux441ux43eux43c}}

В крупномасштабном RL для LLM семплирование обычно выполняется движками
инференса вроде vLLM и SGLang, а логарифмические вероятности и градиенты
вычисляются обучающими системами вроде FSDP и Megatron. Даже при
загрузке одинаковых весов обе стороны могут вычислить немного разные
логарифмические вероятности одного и того же токена из-за различий в
точности вычислений с плавающей запятой, порядке редукции, тензорном
параллелизме, размере пакета, KV-кэше и объединённых вычислительных
ядрах. В результате параметры ещё не обновлены, а \(\rho_t\) уже
отклоняется от 1: между обучением и семплированием возникает численное
различие политик. Контролируемые эксперименты показывают, что даже
небольшие различия на уровне токенов могут вызвать коллапс
обучения.\footnote{Zhong, Tianle et al., ``Diagnosing Training Inference
  Mismatch in LLM Reinforcement Learning'', 2026. arXiv:2605.14220.}

Если распределение семплирования изменилось, наивный градиент политики
может отдать предпочтение одному из двух действий, даже когда их
вознаграждения полностью совпадают. Предположим, что при семплировании
действия А и Б выбираются с одинаковой вероятностью, но во время
обучения вероятность выбора А моделью уже равна 0,8. Для logit,
управляющего действием А, градиент логарифмической вероятности равен
\(1-0.8=0.2\), если выбрано А, и \(0-0.8=-0.8\), если выбрано Б. Если
оба действия получают 1 балл, средний градиент по распределению
семплирования равен \(-0.3\). Хотя обе альтернативы получают одинаковую
оценку, обновление снижает вероятность А.

Score centering вычитает из градиента логарифмической вероятности для
одного и того же префикса его среднее по распределению семплирования.
После вычитания \(-0.3\) в приведённом примере два значения становятся
равны соответственно \(0.5\) и \(-0.5\), а среднее обновление при
одинаковом вознаграждении оказывается точно нулевым. Выше GRPO вычитала
среднее из вознаграждений группы ответов; здесь центрируется
градиент.\footnote{Martin Marek, Max Ryabinin, \emph{Score Centering
  Stabilizes Off-policy Reinforcement Learning}, 2026. Если политика
  семплирования отстала слишком сильно, расстояние между обучаемой и
  поведенческой политиками также можно ограничить с помощью
  KL-регуляризации (arXiv:2611.21791).}

Численная ошибка между обучающей стороной и стороной семплирования также
усиливается вдоль длинной последовательности: \textbf{ошибка
логарифмической вероятности → отклонение отношения вероятностей →
накопление на длинном префиксе → изменение clipping и взвешивания
преимуществ → изменение градиента и эффективного размера выборки}.
Например, если логарифмические отношения вероятностей для 4 000 токенов
имеют отклонение \(10^{-3}\) одного знака, вычисленное отношение
вероятностей на уровне траектории будет в
\(e^{4000\times0.001}=e^4\approx54.6\) раза больше истинного значения.

Ранние изменения вероятностей также могут повлиять на результат
семплирования, направив последующие состояния по другому пути.
Наблюдаемые при обучении сигналы включают всплески отношений важности,
рост доли clipping, скачки градиента или длины ответа, а также снижение
вознаграждения и энтропии. При изменении размера пакета может меняться
процесс редукции и, следовательно, численный результат для одного и того
же входа, то есть нарушается инвариантность относительно размера пакета
(batch invariance). В соответствующих исследованиях наблюдался
возникающий из-за этого неявный эффект off-policy; согласование
вычислений при семплировании и обучении либо явная коррекция различий
между политиками повышали стабильность.\footnote{He, Horace and Thinking
  Machines Lab, ``Defeating Nondeterminism in LLM Inference'', 2025.}

Поэтому автор считает численную согласованность семплирования и обучения
ключевой проверкой перед запуском обучения и непрерывно отслеживает её
вместе с кривой вознаграждения. Обработку можно разделить на четыре
этапа: «измерить различия → согласовать конфигурацию → скорректировать
отклонение → ограничить устаревание».

\begin{itemize}
\tightlist
\item
  \textbf{До любого обновления параметров} сравнивайте логарифмические
  вероятности токенов, вычисленные стороной семплирования (sampler) и
  обучающей стороной (trainer) для одного пакета траекторий;
  отслеживайте среднее значение, квантили и максимум \(\rho_t\),
  приближённую KL-дивергенцию и долю значений, подвергнутых clipping.
  Это наиболее прямой модульный тест режима on-policy.
\item
  При синхронизации весов также согласуйте адаптеры LoRA, токенизатор,
  шаблон диалога, версию модели и конфигурацию позиционного кодирования.
  Rollout должен сохранять логарифмические вероятности, выданные
  политикой семплирования во время генерации, чтобы впоследствии на их
  основе вычислялось отношение вероятностей.
\item
  По возможности согласуйте точность, схему параллелизма и ключевые
  вычислительные ядра при семплировании и обучении. Если это невозможно,
  явно рассматривайте различие как off-policy, применяйте коррекцию по
  важности, отслеживайте эффективный размер выборки и оценивайте
  результат коррекции с учётом доли clipping.
\item
  Своевременно выполняйте повторное семплирование с обновлённой
  политикой, ограничивайте число раундов обновления на каждом пакете
  данных и число checkpoint, на которое асинхронное семплирование
  отстаёт от обучения.
\end{itemize}

\hypertarget{ux441ux440ux435ux434ux430-rl-ux43eux442-ux43eux446ux435ux43dux438ux432ux430ux43dux438ux44f-ux43a-ux441ux438ux43cux443ux43bux44fux446ux438ux438}{%
\section{Среда RL: от оценивания к
симуляции}\label{ux441ux440ux435ux434ux430-rl-ux43eux442-ux43eux446ux435ux43dux438ux432ux430ux43dux438ux44f-ux43a-ux441ux438ux43cux443ux43bux44fux446ux438ux438}}

Для обучения RL AI-агент должен многократно выполнять задачи.
Исправление кода можно повторить в изолированной копии, тогда как платёж
изменяет состояние реального счёта. Обучающая среда должна предоставлять
сбрасываемое начальное состояние, параллельно выполняемые экземпляры
задач и проверяемые результаты, чтобы модель могла исследовать разные
пути в контролируемых условиях. В седьмой главе уже были представлены
среды оценивания; при их использовании для обучения необходимо также
учитывать длительное выполнение, побочные эффекты действий и стоимость
большого числа взаимодействий. Автор рекомендует сначала выделить
специалистов и бюджет на создание среды и только затем планировать
крупномасштабное обучение. Реализация реалистичной обратной связи,
сброса состояния и параллельной изоляции зачастую требует больше
инженерных ресурсов, чем интеграция PPO или GRPO, а стоимость создания
среды может превышать стоимость самого обучения.

\hypertarget{ux441ux440ux435ux434ux430-ux442ux440ux435ux43dux438ux440ux43eux432ux43eux447ux43dux430ux44f-ux43fux43bux43eux449ux430ux434ux43aux430-ux43cux43eux434ux435ux43bux438}{%
\subsection{Среда: тренировочная площадка
модели}\label{ux441ux440ux435ux434ux430-ux442ux440ux435ux43dux438ux440ux43eux432ux43eux447ux43dux430ux44f-ux43fux43bux43eux449ux430ux434ux43aux430-ux43cux43eux434ux435ux43bux438}}

Среда обновляет состояние в соответствии с действиями и возвращает
модели новые наблюдения. При создании среды можно отталкиваться от трёх
вопросов.

\begin{itemize}
\tightlist
\item
  \textbf{Насколько взаимодействие близко к сценарию развёртывания?}
  Симулятор клиентской поддержки должен охватывать случаи, когда
  пользователь предоставляет дополнительные сведения, меняет требования
  или выражается неоднозначно; симулятор инструментов должен
  воспроизводить ошибки и задержки, распространённые в производственной
  среде. На основе этой обратной связи модель учится справляться с
  реальными проблемами. Слишком однообразная обратная связь приводит к
  тому, что стратегия начинает зависеть от фиксированных шаблонов
  симуляции.
\item
  \textbf{Можно ли стабильно повторять и параллельно выполнять задачи?}
  Для каждого rollout требуется независимое состояние заказа, файла или
  учётной записи с возможностью очистки и восстановления после
  завершения. Именно эти функции выполняют песочница интерпретатора кода
  ReTool и MCP-песочница AWorld. Для реальных API с ограничениями
  частоты запросов, платой за использование или побочными эффектами
  записи в зависимости от задачи можно применять учётные записи
  песочницы, воспроизведение данных или имитационные интерфейсы.
\item
  \textbf{Как определить результат?} Состояние среды служит важной
  основой для расчёта вознаграждения. Например, в задаче возврата
  средств можно проверить сумму, состояние заказа и условия авторизации,
  а в задаче по работе с кодом --- выполнить тесты. Проектирование
  вознаграждения будет отдельно рассмотрено далее.
\end{itemize}

\textbf{Прежде чем настраивать алгоритм, автор рекомендует проверить:
достаточно ли среда симуляции близка к реальной задаче? Это гораздо
важнее выбора между PPO и GRPO.} Реалистичность, возможность сброса и
проверяемость вместе определяют, чему модель сможет научиться из
взаимодействий. Как правило, сначала стоит отладить эти три составляющие
на небольшом числе задач и лишь затем увеличивать масштаб обучения: так
проще локализовать проблемы.

\hypertarget{ux441ux438ux43cux443ux43bux44fux446ux438ux44f-ux438ux43dux441ux442ux440ux443ux43cux435ux43dux442ux43eux432-ux438-ux441ux440ux435ux434ux44b-ux441-ux43fux43eux43cux43eux449ux44cux44e-ux43cux43eux434ux435ux43bux438}{%
\subsection{Симуляция инструментов и среды с помощью
модели}\label{ux441ux438ux43cux443ux43bux44fux446ux438ux44f-ux438ux43dux441ux442ux440ux443ux43cux435ux43dux442ux43eux432-ux438-ux441ux440ux435ux434ux44b-ux441-ux43fux43eux43cux43eux449ux44cux44e-ux43cux43eux434ux435ux43bux438}}

Результаты поиска, диалоги с пользователями и сложные состояния
бизнес-процессов можно синтезировать с помощью модели. Симулятор на
основе модели получает действия агента и генерирует следующую обратную
связь, снижая зависимость от внешних сервисов, доступных в реальном
времени. Этот подход можно расширить от отдельного инструмента до всей
среды выполнения задачи.

\textbf{Первый уровень: синтез возвращаемых значений инструментов.} В
ZeroSearch\footnote{Sun, Hao, et al.~``ZeroSearch: Incentivize the
  Search Capability of LLMs without Searching'', 2025. arXiv:2505.04588.}
LLM исполняет роль поискового движка: ученик отправляет запрос, а
симулятор генерирует результаты поиска. Это позволяет уменьшить
зависимость обучения от стоимости реальных поисковых API, ограничений
частоты запросов и колебаний результатов. Кроме того, применяется
учебный план: вначале возвращаются высококачественные и
высокорелевантные документы, а затем постепенно добавляется шум, чтобы
ученик учился извлекать полезную информацию из несовершенных материалов.
После обучения в симулируемой поисковой среде исследованная модель
сохранила высокое качество работы и при подключении к реальному
поисковому движку.

\textbf{Второй уровень: симуляция переходов состояния среды.}
DreamGym\footnote{``DreamGym: Scaling Agent Learning via Experience
  Synthesis'', 2025. arXiv:2511.03773.} дистиллирует динамику среды в
рассуждающую модель опыта. Модель считывает текущее состояние и действие
агента, выводит следующее состояние и обратную связь и массово
генерирует траектории взаимодействия для онлайн-RL, сокращая число
обращений к реальной среде. В задачах клиентской поддержки и продаж LLM
также может исполнять роль пользователя, имитируя уточняющие вопросы,
предоставление дополнительных сведений и изменение требований. В серии
τ-bench применяется симулятор пользователей; аналогичные интерактивные
системы можно использовать и для обучения, разделяя тренировочные задачи
и задачи оценивания.

Симуляторы на основе моделей также нуждаются в калибровке. Если
симулятор клиентской поддержки всегда правильно понимает неоднозначные
вопросы, у стратегии не будет возможности научиться задавать уточняющие
вопросы; если поисковый симулятор всегда возвращает правильные
документы, стратегия не получит опыта распознавания шума. Кроме того,
обучение может закрепить поведение, эксплуатирующее недостатки системы
оценивания симулятора. Автор рекомендует сочетать подходы: поручить
модели симуляцию основной массы взаимодействий, затем регулярно
проверять стратегию на задачах в реальной среде и использовать
обнаруженные расхождения для корректировки симулятора. Так образуется
итерационный процесс «синтетические взаимодействия → оценивание на
реальных задачах → корректировка симулятора».

Как понять, действительно ли симулятор способен научить пользоваться
инструментами? Можно провести небольшой тест: для одного и того же
исходного заказа сначала выполнить запрос, затем отправить изменение и
снова выполнить запрос, проверив, изменились ли возвращаемые поля заказа
ожидаемым образом. После этого следует сбросить среду и убедиться, что
снова возвращаются исходные значения. Среда, которая умеет лишь
сгенерировать правдоподобную фразу «изменение выполнено», но не может
отразить изменение в следующем запросе, будет вознаграждать стратегию за
убедительные слова, даже если заказ так и не был изменён.

\hypertarget{ux438ux437ux43eux43bux44fux446ux438ux44f-ux441ux440ux435ux434ux44b-ux440ux430ux441ux43fux440ux435ux434ux435ux43bux435ux43dux438ux44f-ux437ux430ux434ux430ux447-ux438-ux43eux446ux435ux43dux438ux432ux430ux43dux438ux44f}{%
\subsection{Изоляция среды, распределения задач и
оценивания}\label{ux438ux437ux43eux43bux44fux446ux438ux44f-ux441ux440ux435ux434ux44b-ux440ux430ux441ux43fux440ux435ux434ux435ux43bux435ux43dux438ux44f-ux437ux430ux434ux430ux447-ux438-ux43eux446ux435ux43dux438ux432ux430ux43dux438ux44f}}

Обучающая среда определяет, с чем может столкнуться модель. Сбрасываемое
состояние позволяет повторять попытки, изолированные экземпляры ---
выполнять задачи параллельно, а явные состояния и правила проверки ---
воспроизводить и оценивать результаты. Для создания задач можно
использовать описанный выше метод синтеза данных: извлечь из реальных
журналов бизнес-процессов шаблоны задач, удалить идентифицирующую
информацию, а затем сгенерировать вымышленных персонажей, заказы, файлы
и начальные состояния.

Обучение и оценивание могут использовать общий генератор задач и общий
код проверки, однако конкретные задачи следует разделить на независимые
наборы. Согласно подходам к проектированию сред SWE-Gym, τ²-bench и
AndroidWorld\footnote{Pan, Jiayi et al., ``Training Software Engineering
  Agents and Verifiers with SWE-Gym'', 2024. arXiv:2412.21139; Barres,
  Victor et al., ``\(\tau^2\)-Bench: Evaluating Conversational Agents in
  a Dual-Control Environment'', 2025. arXiv:2506.07982; Rawles,
  Christopher et al., ``AndroidWorld: A Dynamic Benchmarking Environment
  for Autonomous Agents'', 2024. arXiv:2405.14573.}, приёмочные тесты,
скрытые состояния, необходимые для оценивания, и эталонные решения
должны оставаться на стороне верификатора, а агент должен получать
наблюдения через предусмотренный задачей интерфейс. Прежде чем
увеличивать масштаб выборки, следует выполнить небольшое число rollout и
проверить, можно ли решить задачу и способен ли верификатор отличить
действительный результат от недействительного. Раннее обнаружение
смещений в оценивании уменьшает вероятность того, что стратегия научится
эксплуатировать уязвимости проверки.

Таким образом, среду следует разрабатывать в следующем порядке:
\textbf{шаблон задачи → сбрасываемая среда → проверка результата →
изоляция обучения и оценивания → калибровка по реальным
взаимодействиям}. Описанный выше синтез данных отбирает из этих сред
действительные траектории, формируя демонстрации для SFT. После перехода
к RL текущая стратегия продолжает учиться в среде методом проб и ошибок,
исследуя пути за пределами демонстраций.

Проверку также необходимо учитывать в бюджете пропускной способности.
Lean kernel, средство запуска тестов и выполнение в контейнере расходуют
ресурсы CPU и память. Если GPU генерирует результаты быстрее, чем
выполняется проверка, задачи начинают накапливаться в очереди
верификации. В таком случае следует увеличить число воркеров
верификатора, сократить задержку проверки либо скорректировать
соотношение ресурсов генерации и проверки\footnote{Li, Bojie, and Noah
  Shi, ``RLVP: Penalize the Path, Reward the Outcome'', 2026.
  arXiv:2607.07435.}.

\hypertarget{ux43eux442-ux43eux434ux43dux43eux43aux440ux430ux442ux43dux43eux433ux43e-ux43a-ux43cux43dux43eux433ux43eux44dux442ux430ux43fux43dux43eux43cux443-ux432ux437ux430ux438ux43cux43eux434ux435ux439ux441ux442ux432ux438ux44e-ux441ux446ux435ux43dux430ux440ux438ux438-ux437ux430ux434ux430ux447-ux438-ux440ux430ux441ux43fux440ux435ux434ux435ux43bux435ux43dux438ux435-ux437ux430ux441ux43bux443ux433}{%
\section{От однократного к многоэтапному взаимодействию: сценарии задач
и распределение
заслуг}\label{ux43eux442-ux43eux434ux43dux43eux43aux440ux430ux442ux43dux43eux433ux43e-ux43a-ux43cux43dux43eux433ux43eux44dux442ux430ux43fux43dux43eux43cux443-ux432ux437ux430ux438ux43cux43eux434ux435ux439ux441ux442ux432ux438ux44e-ux441ux446ux435ux43dux430ux440ux438ux438-ux437ux430ux434ux430ux447-ux438-ux440ux430ux441ux43fux440ux435ux434ux435ux43bux435ux43dux438ux435-ux437ux430ux441ux43bux443ux433}}

\hypertarget{ux43aux43bux44eux447ux435ux432ux44bux435-ux441ux43bux43eux436ux43dux43eux441ux442ux438-ux43cux43dux43eux433ux43eux44dux442ux430ux43fux43dux44bux445-ux437ux430ux434ux430ux447}{%
\subsection{Ключевые сложности многоэтапных
задач}\label{ux43aux43bux44eux447ux435ux432ux44bux435-ux441ux43bux43eux436ux43dux43eux441ux442ux438-ux43cux43dux43eux433ux43eux44dux442ux430ux43fux43dux44bux445-ux437ux430ux434ux430ux447}}

Слева на \crossreflink{fig:8-14}{рисунке 8-14} решается задача «24».
Модель может выполнить в ответе несколько шагов вычислений, но на всех
этих шагах использует четыре заданные карты и не получает по ходу новых
наблюдений от среды. Справа показана навигация по панорамам улиц. Только
повернув налево на перекрёстке, можно увидеть улицу за поворотом.
Дальнейший путь приходится заново выбирать по находящимся перед глазами
дорожным указателям. Организовать такое взаимодействие позволяет ReAct,
рассмотренный в первой и четвёртой главах: модель выбирает действие,
инструмент или среда выполняет его, а новое наблюдение поступает в
следующий раунд принятия решений.

\begin{figure}
\hypertarget{fig:8-14}{%
\centering
\includegraphics{images/revised/fig8-14.pdf}
\caption{Рисунок 8-14. Сравнение однократного и многоэтапного
RL}\label{fig:8-14}
}
\end{figure}

На \crossreflink{fig:8-15}{рисунке 8-15} AI-агент пропускает перекрёсток
на третьем шаге, на четвёртом замечает, что сбился с пути,
разворачивается и возвращается, а на пятом достигает цели.
Вознаграждение за достижение конечной точки подтверждает успешность всей
навигации, но не различает движение по неверному пути и своевременный
разворот. При следующей встрече с этим перекрёстком агент должен сразу
выбрать правильный путь, однако, если он снова ошибётся, следует
по-прежнему поощрять его за разворот. При обучении необходимо оценить,
как каждый предшествующий шаг повлиял на итоговую выгоду. Эта задача
называется \textbf{распределением заслуг (Credit Assignment)}. При
оценке ценности и преимущества можно учитывать последующие отдачи, а
если промежуточные действия допускают надёжную проверку --- добавлять
процессные вознаграждения. Разворот на четвёртом шаге получает ноль
баллов в рамках процессного вознаграждения, но всё же помогает достичь
конечной точки. Поэтому при оценке ценности этого шага необходимо
учитывать и последующую отдачу.

\begin{figure}
\hypertarget{fig:8-15}{%
\centering
\includegraphics{images/revised/fig8-15.pdf}
\caption{Рисунок 8-15. Распределение заслуг при многоэтапном
взаимодействии}\label{fig:8-15}
}
\end{figure}

Та же проблема возникает и при звонках в службу поддержки. Например,
пользователь оставил положительный отзыв на десятом раунде. Уточнение
спорного списания во втором раунде и объяснение комиссии в седьмом могли
оказаться полезными, но по одному лишь итоговому отзыву невозможно
определить вклад каждого из этих шагов. В следующих двух экспериментах
рассматриваются соответственно ситуации, когда обучающая система может
проверять навигационные действия шаг за шагом и когда об успехе робота
можно судить только по результату всей задачи.

\begin{quote}
\textbf{Эксперимент 8-12 ★★★: V-IRL-VL --- многоэтапная визуальная
навигация}

В V-IRL\footnote{Yang, Jihan et al., ``V-IRL: Grounding Virtual
  Intelligence in Real Life'', 2024. arXiv:2402.03310.} агент должен шаг
за шагом найти пункт назначения, ориентируясь по находящимся перед ним
панорамам улиц. Для обучения используются маршруты Нью-Йорка, а при
тестировании --- другие города, где меняются как формулировки указаний,
так и сами уличные виды. RL заметно превосходит SFT в обоих видах
тестирования вне распределения: при изменении правил и визуальных
данных. В эксперименте применяется PPO с сетью ценности. Обратная связь
после каждого шага помогает определить, какие решения на длинном
маршруте следует закрепить.

\emph{Уточнение русского издания: В исследовании перенос при смене
правил и перенос при изменении визуальной среды оцениваются в отдельных
условиях. Совместная смена города, изображения и правил сама по себе не
позволила бы разделить их эффекты.
\href{https://arxiv.org/html/2501.17161v1}{Источник}.}
\end{quote}

\begin{quote}
\textbf{Эксперимент 8-13 ★★★: SimpleVLA-RL --- открытое исследование с
вознаграждением за результат \texttt{{[}Расширенный\ эксперимент{]}}}

В робототехнических задачах LIBERO SimpleVLA-RL использует
вознаграждение за результат --- успех или неудачу. Для холодного запуска
посредством SFT на каждую задачу используется одна демонстрационная
траектория, после чего RL повышает долю успешных выполнений с 17,3\% до
91,7\% и вырабатывает отсутствующее в демонстрации действие «толкнуть и
разрезать». Сопоставим этот результат с навигацией по панорамам улиц:
если можно надёжно определить, на каком промежуточном шаге был выбран
неверный путь, обратную связь можно предоставить сразу и тем самым
ускорить обучение. Если же возможные способы выполнения задачи ещё
неизвестны, вознаграждение только за результат оставляет модели
пространство для испытания других путей. Если слишком рано превратить
последовательность действий из человеческой демонстрации в требования к
вознаграждению, новые стратегии наподобие «толкнуть и разрезать» могут
получить заниженную оценку.

\emph{Уточнение русского издания: Рост с 17,3 до 91,7 \% относится к
LIBERO-Long. Поведение pushcut наблюдалось в другом наборе задач ---
RoboTwin2.0. \href{https://arxiv.org/html/2509.09674v1}{Источник}.}
\end{quote}

\hypertarget{ux432ux44bux437ux43eux432ux44b-ux438ux43dux441ux442ux440ux443ux43cux435ux43dux442ux43eux432-ux43eux431ux443ux447ux435ux43dux438ux435-ux438ux441ux43fux43eux43bux44cux437ux43eux432ux430ux43dux438ux44e-ux432ux43dux435ux448ux43dux438ux445-ux432ux43eux437ux43cux43eux436ux43dux43eux441ux442ux435ux439-ux432ux43e-ux432ux440ux435ux43cux44f-ux432ux437ux430ux438ux43cux43eux434ux435ux439ux441ux442ux432ux438ux44f}{%
\subsection{Вызовы инструментов: обучение использованию внешних
возможностей во время
взаимодействия}\label{ux432ux44bux437ux43eux432ux44b-ux438ux43dux441ux442ux440ux443ux43cux435ux43dux442ux43eux432-ux43eux431ux443ux447ux435ux43dux438ux435-ux438ux441ux43fux43eux43bux44cux437ux43eux432ux430ux43dux438ux44e-ux432ux43dux435ux448ux43dux438ux445-ux432ux43eux437ux43cux43eux436ux43dux43eux441ux442ux435ux439-ux432ux43e-ux432ux440ux435ux43cux44f-ux432ux437ux430ux438ux43cux43eux434ux435ux439ux441ux442ux432ux438ux44f}}

После подключения внешних инструментов модель может выбирать поиск,
выполнение кода, изменение файлов, запросы к базам данных или
необходимую для задачи комбинацию нескольких API. Harness проверяет и
диспетчеризирует вызовы, а инструменты выполняются в среде и возвращают
наблюдения. На \crossreflink{fig:8-16}{рисунке 8-16} принятие решений в
рамках одной задачи отделено от обучения после её завершения. По
стрелкам справа результат инструмента возвращается в текущую задачу, и
модель использует его для выбора следующего шага. При движении вниз
видно, что только после завершения задачи обучающая программа вычисляет
вознаграждение по всей траектории и обновляет параметры, а новая
стратегия применяется уже при выборке следующего раунда. Если в рамках
одной задачи требуется последовательно вызвать несколько инструментов,
результат предыдущего вызова влияет на параметры следующего.

\begin{figure}
\hypertarget{fig:8-16}{%
\centering
\includegraphics{images/revised/fig8-16.pdf}
\caption{Рисунок 8-16. Цикл вознаграждения RL при вызовах
инструментов}\label{fig:8-16}
}
\end{figure}

Search-R1\footnote{Jin, Bowen et al., ``Search-R1: Training LLMs to
  Reason and Leverage Search Engines with Reinforcement Learning'',
  2025. arXiv:2503.09516.} обучает модель решать, когда выполнять поиск,
какой запрос отправить и как использовать результаты для продолжения
рассуждений. ReTool подключает к такому же циклу интерпретатор кода:
модель учится планировать выполнение кода, понимать возвращаемые
значения и исправлять решение с учётом ошибок. AWorld-train расширяет
этот подход до песочницы MCP с несколькими инструментами, где модель
должна выбирать инструменты и учитывать зависимости между вызовами, а
обучающая система --- уметь сбрасывать среду и повторно воспроизводить
траектории.

В траекториях использования инструментов необходимо различать токены,
сгенерированные моделью, и токены, возвращённые средой. При вычислении
потерь стратегии учитываются только такие действия модели, как
рассуждения, параметры вызовов и ответы, а позиции, соответствующие
обратной связи от среды, маскируются. При этом обратная связь остаётся
частью контекста для последующих вычислений и помогает модели определить
следующий шаг. Например, если интерпретатор возвращает сообщение об
ошибке, обучение должно закреплять то, как модель после этого исправляет
код. Если включить текст ошибки в потери стратегии, к обновлению
примешается цель «генерировать вывод среды», не относящаяся к решениям,
за которые в действительности отвечает модель.

\begin{quote}
\textbf{Эксперимент 8-14 ★★★: ReTool --- решение математических задач с
помощью интерпретатора кода}

На \crossreflink{fig:8-17}{рисунке 8-17} на примере квадратного
уравнения показано, как чередуются текстовые рассуждения, генерация кода
и обратная связь от интерпретатора.

\begin{figure}
\hypertarget{fig:8-17}{%
\centering
\includegraphics{images/revised/fig8-17.pdf}
\caption{Рисунок 8-17. Цикл ReTool, чередующий текстовые и программные
рассуждения с обратной связью от выполнения в песочнице}\label{fig:8-17}
}
\end{figure}

После предварительного прогрева посредством SFT ReTool обучается с
помощью PPO, чередуя текстовые рассуждения, выполнение кода и обратную
связь от интерпретатора. Модель постепенно учится определять, когда
стоит выполнить код, и как изменить его после ошибки. Здесь используются
обучающие данные DAPO-Math-17k, а стратегия по-прежнему обновляется с
помощью стандартного PPO\footnote{Feng, Jiazhan et al., ``ReTool:
  Reinforcement Learning for Strategic Tool Use in LLMs'', 2025.
  arXiv:2504.11536.}\footnote{Yu, Qiying et al., ``DAPO: An Open-Source
  LLM Reinforcement Learning System at Scale'', 2025. arXiv:2503.14476.}.

После обучения точность ответов модели на AIME 2024 выросла примерно с
25\% до 67,0\%. По сравнению с чисто текстовым RL обратная связь от кода
позволила модели быстрее научиться точным вычислениям и исправлению
ошибок.

\emph{Уточнение русского издания: Для Qwen2.5-32B-Instruct на AIME 2024
в статье указаны три отдельные точки: 26,7 \% для исходной модели, 40,9
\% после SFT и 67,0 \% после 400 шагов RL. Это не результат оценивания
на AIME 2025. \href{https://arxiv.org/html/2504.11536v1}{Источник}.}
\end{quote}

\begin{quote}
\textbf{Эксперимент 8-15 ★★★: AWorld-train --- обучение использованию
инструментов в песочнице}

При прохождении по верхней половине цикла на
\crossreflink{fig:8-18}{рисунке 8-18} модель получает результат вызова
инструмента и может продолжить попытки. Теперь пройдём по нижней
половине: эта партия задач уже завершена, обучающая программа получает
траектории с вознаграждениями и обновляет параметры, а новая стратегия
применяется только в следующей партии задач.

\begin{figure}
\hypertarget{fig:8-18}{%
\centering
\includegraphics{images/revised/fig8-18.pdf}
\caption{Рисунок 8-18. Циклы вызовов инструментов и обучения в
AWorld-train}\label{fig:8-18}
}
\end{figure}

Среда GAIA включает 26 MCP-серверов и 126 функций инструментов. Она
позволяет управлять веб-страницами посредством Playwright, извлекать
содержимое из PDF, документов Word и презентаций, транскрибировать
аудио, распознавать текст и выполнять код в песочнице E2B. Инструмент
Excel предоставляет 29 операций, включая работу с формулами, диаграммами
и сводными таблицами, а поиск знаний также обеспечивает доступ к
Wikipedia, arXiv и Wayback. После сброса песочницы попытку можно
повторить из того же состояния или воспроизвести предыдущую траекторию.

Показанные на рисунке измерения времени взяты из контрольного сравнения
при сборе данных, приведённого в статье. В сопроводительном эксперименте
также указана конфигурация обучения модели 4B: Qwen3-4B-Thinking-2507,
восемь A100, предел контекстного окна 131K, не более 32 раундов
взаимодействия, размер пакета 32, выборка 16 траекторий на задачу,
обновление посредством GRPO и скорость обучения 1e−6.
\end{quote}

Инструмент вернул результат --- но какой шаг должна вознаградить
обучающая программа? Успешный вызов означает лишь, что инструмент
выполнился без ошибки. Чтобы определить, приблизил ли этот шаг
выполнение задачи, необходимо с помощью валидатора проверить вызванные
им изменения и только затем решить, какое вознаграждение предоставить
модели.

\hypertarget{ux43fux440ux43eux435ux43aux442ux438ux440ux43eux432ux430ux43dux438ux435-ux432ux43eux437ux43dux430ux433ux440ux430ux436ux434ux435ux43dux438ux44f-ux43aux430ux43a-ux43fux440ux435ux432ux440ux430ux442ux438ux442ux44c-ux446ux435ux43bux44c-ux437ux430ux434ux430ux447ux438-ux432-ux43eux431ux443ux447ux430ux44eux449ux438ux439-ux441ux438ux433ux43dux430ux43b}{%
\section{Проектирование вознаграждения: как превратить цель задачи в
обучающий
сигнал}\label{ux43fux440ux43eux435ux43aux442ux438ux440ux43eux432ux430ux43dux438ux435-ux432ux43eux437ux43dux430ux433ux440ux430ux436ux434ux435ux43dux438ux44f-ux43aux430ux43a-ux43fux440ux435ux432ux440ux430ux442ux438ux442ux44c-ux446ux435ux43bux44c-ux437ux430ux434ux430ux447ux438-ux432-ux43eux431ux443ux447ux430ux44eux449ux438ux439-ux441ux438ux433ux43dux430ux43b}}

Представим две траектории возврата средств, в результате которых
возвращена одинаковая сумма, но в одной из них пропущена обязательная
проверка личности. Если учитывать только конечный баланс, обе получат
одинаковую оценку. Чтобы превратить полный набор требований пользователя
в обучающий сигнал, сначала определите, кто может проверить результат,
на каком шаге можно получить доказательства и какую информацию
необходимо передать модели, а затем проверьте, можно ли отдельно
верифицировать разрешения и процессные ограничения на всём пути
выполнения.

\hypertarget{ux43eux442ux43aux443ux434ux430-ux43fux43eux441ux442ux443ux43fux430ux435ux442-ux432ux43eux437ux43dux430ux433ux440ux430ux436ux434ux435ux43dux438ux435-ux43fux440ux430ux432ux438ux43bux430-ux447ux435ux43bux43eux432ux435ux447ux435ux441ux43aux438ux435-ux43fux440ux435ux434ux43fux43eux447ux442ux435ux43dux438ux44f-ux438-ux43eux446ux435ux43dux43aux430-ux43cux43eux434ux435ux43bux44cux44e}{%
\subsection{Откуда поступает вознаграждение: правила, человеческие
предпочтения и оценка
моделью}\label{ux43eux442ux43aux443ux434ux430-ux43fux43eux441ux442ux443ux43fux430ux435ux442-ux432ux43eux437ux43dux430ux433ux440ux430ux436ux434ux435ux43dux438ux435-ux43fux440ux430ux432ux438ux43bux430-ux447ux435ux43bux43eux432ux435ux447ux435ux441ux43aux438ux435-ux43fux440ux435ux434ux43fux43eux447ux442ux435ux43dux438ux44f-ux438-ux43eux446ux435ux43dux43aux430-ux43cux43eux434ux435ux43bux44cux44e}}

Для задач, допускающих автоматическую проверку, автор отдаёт
предпочтение \textbf{обучению с подкреплением на основе верифицируемого
вознаграждения (RLVR)}, при котором приёмочные требования
непосредственно выражаются программным кодом. В математической задаче
можно проверить ответ, в задаче по программированию --- запустить тесты,
а в задаче возврата средств --- проверить сумму и состояние заказа в
базе данных. За успешное прохождение можно начислить 1, за неудачу ---
0: так получается простейшее вознаграждение за результат. Чёткие правила
и воспроизводимое выполнение упрощают поиск причин расхождений в
оценках, а область покрытия верификатора определяет, какие именно
требования будут оптимизироваться при обучении.

\textbf{Обучение с подкреплением на основе обратной связи от людей
(RLHF)} преобразует человеческие предпочтения в обучающий сигнал. В
InstructGPT\footnote{Ouyang, Long et al., ``Training Language Models to
  Follow Instructions with Human Feedback'', OpenAI, 2022.} сначала
собираются результаты сравнения ответов людьми, затем обучается модель
вознаграждения, после чего стратегия оптимизируется с помощью PPO.
Модель вознаграждения усваивает предпочтения из ограниченной выборки,
поэтому постоянная погоня за более высокой оценкой модели может
усиливать её смещения\footnote{Gao, Leo, John Schulman, and Jacob
  Hilton, ``Scaling Laws for Reward Model Overoptimization'', OpenAI,
  2023.}. При обучении часто добавляют KL-регуляризацию относительно
эталонной SFT-модели, чтобы ограничить смещение стратегии.
DPO\footnote{Rafailov, Rafael et al., ``Direct Preference Optimization:
  Your Language Model is Secretly a Reward Model'', 2023.}, в свою
очередь, оптимизирует стратегию непосредственно по парам предпочтений
«лучший ответ --- худший ответ», устраняя отдельный этап обучения модели
вознаграждения.

В открытых задачах модель может выдавать обоснование оценки и балл.
Диагностику, создаваемую \textbf{генеративной моделью вознаграждения
(GRM)}, также можно использовать для отбора демонстраций, построения пар
предпочтений или подготовки последующих исправлений.
DeepSeek-GRM\footnote{Liu, Zijun et al., ``Inference-Time Scaling for
  Generalist Reward Modeling'', 2025. arXiv:2504.02495.} генерирует
принципы оценки с учётом задачи, а затем на их основе формирует
критический разбор и выставляет балл. Поведение оценщика улучшается
посредством онлайн-RL, а итоговое суждение агрегируется по нескольким
сэмплам и с помощью вспомогательной модели вознаграждения.

Разрыв между метрикой оценки и целью задачи также влияет на поведение
модели. \textbf{Reward hacking} означает получение стратегией высокой
оценки за счёт использования недостатков правил или их реализации. В
длительных задачах AI-агентов встречается и явление, называемое
\textbf{reward seeking}: модель действует в соответствии с собственным
предположением о способе оценки, например самостоятельно выбирает
проверку, которую легко пройти, и завершает задачу сразу после её
успешного прохождения\footnote{storm, «Самопроверка и ранняя остановка
  агентов на длинном горизонте: феномен Reward Seeking и меры по его
  смягчению», сообщество Qingke, 6 августа 2026 года.}. Этот пример
показывает, что критерии приёмки должны исходить из цели пользователя и
охватывать реальные результаты работы. Если проверять только показатели,
самостоятельно выбранные моделью, легко упустить ещё не выполненные
требования.

При проектировании вознаграждения для задачи возврата средств сначала
пропустите через верификатор две написанные человеком траектории. В
первой выясняются данные заказа, вызывается интерфейс возврата средств и
пользователю пересказывается полученное подтверждение; во второй лишь
говорится: «Возврат средств выполнен». Если обе траектории получают
одинаковую оценку, сначала выясните, что именно читает оценщик: итоговый
ответ, подтверждение интерфейса или состояние заказа. Этого небольшого
контрпримера достаточно, чтобы выявить проблему прокси-метрики и не
обнаруживать лишь после обучения на целом пакете данных, что модель
научилась преждевременно сообщать об успехе.

\hypertarget{ux43aux43eux433ux434ux430-ux432ux44bux434ux430ux451ux442ux441ux44f-ux432ux43eux437ux43dux430ux433ux440ux430ux436ux434ux435ux43dux438ux435-ux437ux430-ux440ux435ux437ux443ux43bux44cux442ux430ux442-ux438-ux437ux430-ux43fux440ux43eux446ux435ux441ux441}{%
\subsection{Когда выдаётся вознаграждение: за результат и за
процесс}\label{ux43aux43eux433ux434ux430-ux432ux44bux434ux430ux451ux442ux441ux44f-ux432ux43eux437ux43dux430ux433ux440ux430ux436ux434ux435ux43dux438ux435-ux437ux430-ux440ux435ux437ux443ux43bux44cux442ux430ux442-ux438-ux437ux430-ux43fux440ux43eux446ux435ux441ux441}}

\textbf{Вознаграждение за результат} оценивает выполнение задачи после
завершения эпизода; если для оценки результата применяется модель, её
называют \textbf{ORM (Outcome Reward Model)}. Вознаграждение
SimpleVLA-RL за успех или неудачу позволяет стратегии исследовать
различные последовательности действий, ориентируясь на конечную цель.
Такой сигнал легко определить, но после неудачи в длинной траектории
модели требуются дополнительные попытки, чтобы различить действия,
приведшие к сбою. Это одна из причин, по которым разреженная обратная
связь снижает эффективность по числу примеров\footnote{Silver, David and
  Richard S. Sutton, ``Welcome to the Era of Experience'', 2025.}.

Рассмотрим преждевременное завершение работы Coding Agent. После того
как модель заявляет о завершении, Harness запускает защищённые
приёмочные тесты в изолированном рабочем пространстве. Тесты читают
реальные файлы и состояние среды: за прохождение начисляется
положительное вознаграждение, за непрохождение --- отрицательное. Таким
образом, заявление о завершении инициирует внешнюю проверку, а результат
задачи определяется фактическими итогами работы. Для оценки отдельно
подготавливаются две группы задач: в одной работа ещё не завершена и
измеряется доля преждевременных завершений, а в другой всё уже выполнено
и проверяется, способна ли модель своевременно закончить работу.

\textbf{Вознаграждение за процесс} предоставляет обратную связь на
промежуточных шагах, например оценивает проверку личности, выбор
параметров, навигационные действия или шаги рассуждения. Модель,
оценивающая процесс, называется \textbf{PRM (Process Reward Model)}. В
работе OpenAI «Let's Verify Step by Step»\footnote{Lightman, Hunter et
  al., ``Let's Verify Step by Step'', OpenAI, 2023.} исследовалась
пошаговая проверка математических рассуждений. Процессный сигнал
помогает модели локализовать ошибки, но увеличивает стоимость разметки и
проверки. Если считать единственно правильным какой-либо конкретный
порядок решения задачи, это также снизит вознаграждение для других
эффективных стратегий. Поэтому автор добавляет процессные сигналы только
для действительно верифицируемых промежуточных событий, позволяя модели
продолжать исследование разных эффективных решений. V-IRL-VL
(эксперимент 8-12) использует пошаговую обратную связь при навигации, а
SimpleVLA-RL (эксперимент 8-13) --- вознаграждение в конечной точке,
демонстрируя разный выбор детализации обратной связи для двух задач.

На практике сначала можно построить базовую систему на вознаграждении за
результат, а затем добавить процессные сигналы для таких промежуточных
событий, как завершение проверки личности или прохождение тестов. В
многоходовой задаче ограниченной длины можно задать коэффициент
дисконтирования \(\gamma=1\), чтобы при распространении вознаграждения
из конечной точки назад его вес оставался неизменным. PPO связывает
обратную связь с действиями на каждом шаге посредством оценки ценности и
вычисления преимущества; GRPO с преимуществом на уровне траектории
обычно использует один и тот же вес преимущества при вычислении потерь
стратегии для каждого сгенерированного токена этой траектории.

\hypertarget{ux441ux43aux43eux43bux44cux43aux43e-ux438ux43dux444ux43eux440ux43cux430ux446ux438ux438-ux434ux43eux43bux436ux43dux43e-ux432ux44bux440ux430ux436ux430ux442ux44c-ux432ux43eux437ux43dux430ux433ux440ux430ux436ux434ux435ux43dux438ux435-ux441ux43aux430ux43bux44fux440-ux432ux435ux43aux442ux43eux440-ux438-ux433ux435ux43dux435ux440ux430ux442ux438ux432ux43dux430ux44f-ux434ux438ux430ux433ux43dux43eux441ux442ux438ux43aux430}{%
\subsection{Сколько информации должно выражать вознаграждение: скаляр,
вектор и генеративная
диагностика}\label{ux441ux43aux43eux43bux44cux43aux43e-ux438ux43dux444ux43eux440ux43cux430ux446ux438ux438-ux434ux43eux43bux436ux43dux43e-ux432ux44bux440ux430ux436ux430ux442ux44c-ux432ux43eux437ux43dux430ux433ux440ux430ux436ux434ux435ux43dux438ux435-ux441ux43aux430ux43bux44fux440-ux432ux435ux43aux442ux43eux440-ux438-ux433ux435ux43dux435ux440ux430ux442ux438ux432ux43dux430ux44f-ux434ux438ux430ux433ux43dux43eux441ux442ux438ux43aux430}}

Момент появления вознаграждения определяет плотность обратной связи, а
его содержимое --- форму представления. Скаляр обобщает общую
результативность; полускалярная форма «краткое обоснование + балл»
дополняет оценку её основанием; вектор раздельно выражает точность,
полноту, стоимость и другие измерения; генеративная диагностика
конкретно описывает проблему и направление исправления, а также
допускает агрегирование нескольких сэмплов.

\begin{itemize}
\tightlist
\item
  При наличии однозначного ответа или приёмочных тестов можно начать с
  бинарного скаляра.
\item
  При наличии нескольких целей качества их можно оценивать раздельно, а
  затем взвешивать или ограничивать в соответствии с требованиями
  бизнеса.
\item
  Если оценка требует контекстного объяснения, можно генерировать
  диагностику и калибровать её с помощью проверки фактов и выборочного
  контроля людьми.
\end{itemize}

Например, в задаче поддержки клиентов можно одновременно фиксировать,
решена ли проблема, число ходов диалога и факт завершения проверки
личности. Эти три измерения соответствуют результативности, стоимости и
процессным требованиям. Перед использованием в обучении следует
проверить стабильность оценок и их способность различать поведение
траекторий внутри одной группы. Автор начинает с небольшого числа
надёжных показателей и добавляет новые измерения лишь после
последовательной проверки каждого из них. Любая непроверенная оценочная
составляющая может заставить модель оптимизировать новую прокси-метрику,
не улучшая выполнение реальной задачи.

\hypertarget{ux43eux433ux440ux430ux43dux438ux447ux435ux43dux438ux44f-ux43fux443ux442ux438-ux438-rlvp}{%
\subsection{Ограничения пути и
RLVP}\label{ux43eux433ux440ux430ux43dux438ux447ux435ux43dux438ux44f-ux43fux443ux442ux438-ux438-rlvp}}

В процессе выполнения задачи могут действовать независимые требования к
поведению: связываться с пользователем только в рабочее время, проводить
проверку личности, не изменять защищённые файлы. RLVP (Reinforcement
Learning with Verified Penalty)\footnote{Li, Bojie, and Noah Shi,
  ``RLVP: Penalize the Path, Reward the Outcome'', 2026.
  arXiv:2607.07435.} объединяет вознаграждение за результат с
верифицируемым штрафом за путь по принципу \textbf{вознаграждать
результат, штрафовать путь}. Правила пути срабатывают на наблюдаемых
средой действиях, а полнота результата по-прежнему проверяется
приёмочными тестами задачи.

Оба исправления могут пройти приёмку, хотя в одном из них файл, который
сначала следовало прочитать, был сразу перезаписан. Это нарушение можно
точно определить по журналу инструментов; установить же вклад каждого
предшествующего шага в конечный успех значительно сложнее. Простейший
подход --- назначить всей траектории оценку \(R=O+\beta\Phi\), где \(O\)
--- оценка задачи, \(\Phi\) агрегирует штрафы за нарушения или
верифицируемый прогресс, а \(\beta\) определяет их относительный вес.
Это позволяет различать два исправления, но каждое действие внутри одной
траектории получает одинаковый сигнал.

Если известно, на каком ходу произошло нарушение, штраф можно отнести
только к этому ходу. Рассмотрим две тестовые траектории из сопутствующей
реализации: траектория A завершает задачу, но нарушает правила на втором
ходу; траектория B не завершает задачу, но соблюдает правила на обоих
ходах. Их оценки результата равны соответственно 1 и 0, а после
вычитания среднего по группе получаются +0,5 и −0,5. В каждой траектории
два хода, каждый из которых содержит по 10 токенов действий. При
конфигурации \texttt{c2} из реализации два сигнала учитываются
раздельно:

\begin{Shaded}
\begin{Highlighting}[]
\NormalTok{outcome\_weight }\OperatorTok{=}\NormalTok{ (}\DecValTok{1} \OperatorTok{{-}} \FloatTok{0.5}\NormalTok{) }\OperatorTok{/} \DecValTok{40}  \CommentTok{\# сигнал результата распределяется между всеми 40 токенами действий в пакете}
\NormalTok{penalty\_weight }\OperatorTok{=} \OperatorTok{{-}}\FloatTok{0.3} \OperatorTok{/} \DecValTok{10}       \CommentTok{\# штраф распределяется только между 10 токенами действий второго хода траектории A}
\NormalTok{first\_turn }\OperatorTok{=}\NormalTok{ outcome\_weight}
\NormalTok{second\_turn }\OperatorTok{=}\NormalTok{ outcome\_weight }\OperatorTok{+}\NormalTok{ penalty\_weight}
\end{Highlighting}
\end{Shaded}

На первом ходу траектории A каждый токен действия получает вес +0,0125,
а на втором --- −0,0175. Таким образом, два действия в рамках одного
успешного исправления могут получить противоположно направленные
обучающие сигналы. После включения таких весов в потери стратегии
оптимизатор соответствующим образом обновляет параметры. Тот же способ
распределения штрафа применим и к задаче поддержки клиентов, если журнал
инструментов позволяет определить ход, на котором была пропущена
проверка разрешения.

Роль сигнала пути можно понять через различия внутри группы. Если
вознаграждение за результат одинаково у всей группы траекторий GRPO,
рассчитанное по нему преимущество будет нулевым. Если же число нарушений
в траекториях различается, штраф за путь предоставляет новое основание
для сравнения. Вознаграждение за частичный прогресс подчиняется тому же
условию: чтобы оно создавало различия, модель должна сэмплировать
траектории с разной степенью прогресса.

При проектировании вознаграждения за путь можно следовать четырём
правилам: штрафовать за однозначно определяемые нарушения; сохранять
вознаграждение за результат задачи; предоставлять достижимый допустимый
путь для ограничиваемого поведения; использовать воспроизводимые правила
проверки. Например, отдельный штраф за каждое действие будет побуждать
модель оставаться на месте, поэтому необходимо сохранить вознаграждение
за успешное выполнение задачи, чтобы одновременно стимулировать её
завершение. Если базовой модели трудно генерировать допустимые пути,
сначала можно задать отправную точку небольшим числом демонстраций, а
затем в ходе обучения постепенно снижать зависимость от формирования
вознаграждения по пути.

\begin{quote}
\textbf{Эксперимент 8-16 ★★★: RLVP --- вознаграждать результат,
штрафовать путь}

В эксперименте к GRPO добавляются вознаграждение за результат \(O\) и
сигнал пути \(\Phi\), после чего система сравнивается с вариантом,
использующим только вознаграждение за результат. В агентных задачах
TerminalBench число нарушений снизилось с 3,71 до 0,66 при практически
неизменной доле успешно выполненных задач; в агентных задачах miniF2F
добавление достижимого вознаграждения за частичный прогресс сократило
число итераций, необходимых для достижения доли успеха 0,9, с 7,0 до
4,4.

Задача исправления программного обеспечения также демонстрирует роль
достижимости сигнала: если ни один роллаут не способен пройти хотя бы
один тест, вознаграждение за прогресс, вычисляемое по числу пройденных
тестов, всегда остаётся одинаковым и не увеличивает различия внутри
группы.
\end{quote}

Снижение среднего числа нарушений до 0,66 означает, что модель чаще
выбирает допустимый путь, но всё же иногда нарушает правила. Для
ограничений, обязательных к соблюдению, при вызове инструментов
необходимо дополнительно проверять права доступа.

\hypertarget{ux434ux438ux441ux442ux438ux43bux43bux44fux446ux438ux44f-ux43fux43eux432ux44bux448ux435ux43dux438ux435-ux44dux444ux444ux435ux43aux442ux438ux432ux43dux43eux441ux442ux438-ux432ux44bux431ux43eux440ux43aux438}{%
\section{Дистилляция: повышение эффективности
выборки}\label{ux434ux438ux441ux442ux438ux43bux43bux44fux446ux438ux44f-ux43fux43eux432ux44bux448ux435ux43dux438ux435-ux44dux444ux444ux435ux43aux442ux438ux432ux43dux43eux441ux442ux438-ux432ux44bux431ux43eux440ux43aux438}}

Одна траектория rollout AI-агента часто включает несколько раундов
генерации моделью и выполнения инструментов. Повысить
\textbf{эффективность выборки} --- значит достичь целевого качества с
меньшим числом взаимодействий со средой. Количество шагов обучения,
время работы GPU и число обращений к среде отражают разные виды затрат,
поэтому их необходимо рассматривать вместе с особенностями задачи.
Например, в рассмотренном в этой главе случае ReTool обучение SFT
занимает около 1 часа, а RL --- около 9 дней; значительная часть
ресурсов расходуется на многократное сэмплирование и выполнение
инструментов.

При вознаграждении за конечный результат целая траектория может получить
лишь одну оценку успеха или неудачи. Допустим, бизнес-инструмент уже
сообщил: «Нужны последние четыре цифры кредитной карты», однако агент
пропустил эту информацию и в итоге получил вознаграждение 0. В процессе
обучения ещё предстоит по множеству попыток определить, чем именно
вызвана неудача: отсутствующим полем, ошибочными параметрами инструмента
или ошибкой в последующей операции. Если учитель предоставляет пошаговый
надзор на этой траектории, одно и то же взаимодействие со средой создаёт
более детализированный обучающий сигнал.

\hypertarget{ux434ux438ux441ux442ux438ux43bux43bux44fux446ux438ux44f-ux43dux430-ux434ux430ux43dux43dux44bux445-ux442ux435ux43aux443ux449ux435ux439-ux43fux43eux43bux438ux442ux438ux43aux438-ux43fux43bux43eux442ux43dux44bux439-ux43dux430ux434ux437ux43eux440-ux438ux437-ux43eux434ux43dux43eux439-ux442ux440ux430ux435ux43aux442ux43eux440ux438ux438-rollout}{%
\subsection{Дистилляция на данных текущей политики: плотный надзор из
одной траектории
rollout}\label{ux434ux438ux441ux442ux438ux43bux43bux44fux446ux438ux44f-ux43dux430-ux434ux430ux43dux43dux44bux445-ux442ux435ux43aux443ux449ux435ux439-ux43fux43eux43bux438ux442ux438ux43aux438-ux43fux43bux43eux442ux43dux44bux439-ux43dux430ux434ux437ux43eux440-ux438ux437-ux43eux434ux43dux43eux439-ux442ux440ux430ux435ux43aux442ux43eux440ux438ux438-rollout}}

При \textbf{дистилляции на данных текущей политики (On-Policy
Distillation)} ученик генерирует траекторию, а учитель предоставляет
надзор на фактически созданных учеником префиксах. В исследовании
Thinking Machines Lab 2025 года этот способ обучения был систематически
продемонстрирован\footnote{Thinking Machines Lab, ``On-Policy
  Distillation'', 2025.}. Чтобы понять различия между тремя методами,
полезно отдельно рассматривать, «кто создаёт обучающие состояния» и «кто
предоставляет обучающий сигнал»:

\begin{longtable}[]{@{}
  >{\raggedright\arraybackslash}p{(\columnwidth - 4\tabcolsep) * \real{0.3178}}
  >{\raggedright\arraybackslash}p{(\columnwidth - 4\tabcolsep) * \real{0.2828}}
  >{\raggedright\arraybackslash}p{(\columnwidth - 4\tabcolsep) * \real{0.3994}}@{}}
\toprule\noalign{}
\begin{minipage}[b]{\linewidth}\raggedright
\hspace{0pt}Метод
\end{minipage} & \begin{minipage}[b]{\linewidth}\raggedright
\hspace{0pt}Кто формирует траекторию или состояние
\end{minipage} & \begin{minipage}[b]{\linewidth}\raggedright
\hspace{0pt}Основной обучающий сигнал
\end{minipage} \\
\midrule\noalign{}
\endhead
\bottomrule\noalign{}
\endlastfoot
\hspace{0pt}SFT /\allowbreak{} дистилляция на данных сторонней политики
& \hspace{0pt}Человек или учитель & \hspace{0pt}Надзор по целевому
ответу на уровне каждого токена \\
\hspace{0pt}On-\allowbreak{}\hspace{0pt}policy RL & \hspace{0pt}Текущий
ученик & \hspace{0pt}Вознаграждение за результат или процесс \\
\hspace{0pt}On-\allowbreak{}\hspace{0pt}policy Distillation &
\hspace{0pt}Текущий ученик & \hspace{0pt}Оценка учителем вероятностей
каждого токена на префиксах ученика \\
\end{longtable}

Например, в демонстрации учитель всегда сначала запрашивает заказ и
только затем оформляет возврат, тогда как ученик может сначала
попытаться оформить возврат и получить ошибку «Отсутствует номер
заказа». При дистилляции на данных текущей политики учитель оценивает
состояние, до которого уже дошёл ученик, и подсказывает, как запросить
недостающую информацию. Благодаря этому обучение охватывает собственные
ошибки ученика и пути восстановления после них.

Один из вариантов реализации минимизирует обратную KL-дивергенцию между
учеником и учителем на префиксах, сгенерированных учеником:

\[
D_{\mathrm{KL}}(\pi_\theta\Vert\pi_T)
=\mathbb{E}_{a\sim\pi_\theta(\cdot\mid h)}
\left[\log\pi_\theta(a\mid h)-\log\pi_T(a\mid h)\right],
\]

где \(h\) --- текущий префикс, \(a\) --- следующий токен, а
\(\pi_\theta\) и \(\pi_T\) обозначают распределения ученика и учителя
соответственно. Токены, которые учитель считает более подходящими,
получают более благоприятный обучающий сигнал. При реализации можно
сравнивать полные распределения либо запрашивать у учителя
логарифмические вероятности уже выбранных учеником токенов и на их
основе формировать сигнал оптимизации для каждого токена.

Рассмотрим пример «сначала запросить API, затем разобрать возвращённое
значение». Допустим, в определённой позиции вариантам «запросить» и
«вызвать» соответствует по одному токену, которым учитель назначает
вероятности 80\% и 15\%, а суммарная вероятность остальных вариантов
составляет 5\%. Это даёт гораздо более точное основание для выбора, чем
итоговый успех или неудача. Если траектория содержит \(T\) токенов,
сгенерированных моделью, надзор можно получить приблизительно в \(T\)
позициях. Результаты инструментов по-прежнему используются как контекст,
но маскируются при вычислении потерь стратегии.

Ниже этот процесс представлен в форме выравнивания полных распределений:

\begin{Shaded}
\begin{Highlighting}[]
\NormalTok{trajectory }\OperatorTok{=}\NormalTok{ rollout(student, task)}
\NormalTok{loss }\OperatorTok{=} \DecValTok{0}
\ControlFlowTok{for}\NormalTok{ prefix }\KeywordTok{in}\NormalTok{ model\_generated\_positions(trajectory):}
\NormalTok{    target }\OperatorTok{=}\NormalTok{ stop\_gradient(teacher.distribution(prefix))}
\NormalTok{    loss }\OperatorTok{+=}\NormalTok{ KL(student.distribution(prefix), target)}
\NormalTok{update\_student(loss)}
\end{Highlighting}
\end{Shaded}

Ученик отвечает за исследование, а учитель предоставляет надзор во
множестве позиций; основные дополнительные затраты связаны с инференсом
учителя. Когда взаимодействие со средой дорого, а учитель способен
надёжно направлять ученика, автор в первую очередь рассмотрел бы
дистилляцию на данных текущей политики, позволяющую извлекать больше
полезных обучающих сигналов из каждого обращения к среде. Если ученик
недостаточно владеет целевым языком, предметными знаниями или базовыми
действиями, сначала можно создать эффективную исходную точку с помощью
Mid-training и обучения на демонстрациях, а затем перейти к дистилляции
на данных текущей политики.

Не менее важно соответствие между сэмплированием и обучением: движок
сэмплирования формирует траектории по стратегии \(\mu\), а движок
обучения вычисляет потери по \(\pi_\theta\); различие между ними
изменяет распределение обучающих состояний.

Сравнение Thinking Machines Lab на математических задачах
продемонстрировало повышение эффективности обучения примерно на порядок.
В многораундовых агентах можно использовать тот же способ сравнения,
одновременно фиксируя долю успешных задач, количество обращений к среде
и вычислительные затраты учителя. Выгода совместно определяется тем, с
какими состояниями ученик сталкивается в среде и способен ли учитель
предоставить надёжную обратную связь.

Автор и его коллеги также сравнили DPO, четыре разновидности RL и
дистилляцию на данных текущей политики в задаче на «чувство
времени»\footnote{Li, Bojie, and Noah Shi, ``Agents That Sense Physical
  Time: Urgency, Persistence, and Vigilance as Missing Controls for LLM
  Agents'', 2026.}. Первые подходы сталкивались соответственно с
проблемами разреженного вознаграждения, рассогласования цели,
несоответствия формы rollout и коллапса политики. После перехода к
замороженному учителю Qwen3-32B и выравниванию на уровне отдельных
токенов на собственных многораундовых траекториях ученика обучение стало
плавно сходиться, а доля успешно пройденных сценариев в четырёх условиях
превысила базовый вариант SFT на той же исходной модели на 23--47
процентных пунктов. Этот опыт заставил автора придавать больше значения
плотности сигнала от каждого взаимодействия. Когда в обучении участвуют
все промежуточные решения одного взаимодействия, уже сформированные
траектории используются гораздо полнее.

Как сделать так, чтобы один ученик перенял сильные стороны и учителя
математики, хорошо умеющего рассуждать, и учителя программирования,
хорошо исправляющего программы? Дистилляция на данных текущей политики с
несколькими учителями (MOPD) использует описанный выше процесс обучения,
но заменяет фиксированного учителя выбором учителя в соответствии с
предметной областью задачи. Если при решении задачи по программированию
ученик на полпути свернул не туда, уже написанный им префикс передаётся
учителю программирования, чтобы получить распределение вероятностей для
следующего шага. При решении математической задачи ученика в ходе
решения направляет учитель математики. Учителей для разных областей
можно сначала обучить независимо, а на этапе дистилляции зафиксировать
их параметры. При развёртывании достаточно запускать только одного
ученика.\footnote{Wenhan Ma и др., \emph{MOPD: Multi-Teacher On-Policy
  Distillation for Capability Integration in LLM Post-Training}, 2026. В
  работе используются обученные по отдельным областям учителя, знания
  которых дистиллируются на собственных траекториях ученика.}

Здесь есть два практических условия. Во-первых, чтобы вычислять
KL-дивергенцию для каждого токена по приведённой выше формуле,
необходимо получать выравниваемые вероятности учителя на заданных
учеником префиксах. Несколько готовых полных ответов, написанных
учителем, можно использовать для SFT, но их недостаточно для вычисления
этих потерь. Во-вторых, при приёмочной оценке результаты необходимо
рассматривать отдельно для каждой области. Если качество ученика в
математике выросло, а в программировании снизилось, увеличение общего
среднего показателя легко скроет это ухудшение.

\hypertarget{ux441ux430ux43cux43eux434ux438ux441ux442ux438ux43bux43bux44fux446ux438ux44f-ux43dux430-ux434ux430ux43dux43dux44bux445-ux442ux435ux43aux443ux449ux435ux439-ux43fux43eux43bux438ux442ux438ux43aux438-ux441ux43eux437ux434ux430ux43dux438ux435-ux443ux447ux438ux442ux435ux43bux44f-ux441-ux43fux43eux43cux43eux449ux44cux44e-ux434ux43eux43fux43eux43bux43dux438ux442ux435ux43bux44cux43dux43eux433ux43e-ux43aux43eux43dux442ux435ux43aux441ux442ux430}{%
\subsection{Самодистилляция на данных текущей политики: создание учителя
с помощью дополнительного
контекста}\label{ux441ux430ux43cux43eux434ux438ux441ux442ux438ux43bux43bux44fux446ux438ux44f-ux43dux430-ux434ux430ux43dux43dux44bux445-ux442ux435ux43aux443ux449ux435ux439-ux43fux43eux43bux438ux442ux438ux43aux438-ux441ux43eux437ux434ux430ux43dux438ux435-ux443ux447ux438ux442ux435ux43bux44f-ux441-ux43fux43eux43cux43eux449ux44cux44e-ux434ux43eux43fux43eux43bux43dux438ux442ux435ux43bux44cux43dux43eux433ux43e-ux43aux43eux43dux442ux435ux43aux441ux442ux430}}

Если у команды есть только одна модель, хорошо владеющая данной
предметной областью, кто будет учителем? Можно дважды передать ей одну и
ту же задачу. В первый раз ей предоставляется вопрос пользователя в
условиях развёртывания. Во второй раз дополнительно передаются эталонный
ответ, демонстрация эксперта или предметная документация. Дополнительная
информация, доступная при втором ответе, позволяет направлять ученика,
работающего в условиях первого запуска.

\textbf{Самодистилляция на данных текущей политики (On-Policy
Self-Distillation, OPSD)}\footnote{Zhao, Siyan, et al.~``Self-Distilled
  Reasoner: On-Policy Self-Distillation for Large Language Models'',
  2026. arXiv:2601.18734.} возлагает на одну модель две роли: ученик
видит только входные данные задачи, а учитель дополнительно получает
эталонный ответ или уже проверенное решение. После того как ученик
сформирует траекторию, учитель с учётом этой привилегированной
информации оценивает каждый префикс и передаёт ученику сигнал надзора.

Например, ученик решает математическую задачу, а учителю уже известен
правильный ответ. Благодаря этому учитель может оценить текущее
направление рассуждений ученика, а ученик учится улучшать свою
генерацию, располагая только условием задачи. В следующем псевдокоде
показано, куда добавляется дополнительный контекст и где изолируются
градиенты:

\begin{Shaded}
\begin{Highlighting}[]
\NormalTok{trajectory }\OperatorTok{=}\NormalTok{ rollout(model, task\_without\_answer)}
\NormalTok{loss }\OperatorTok{=} \DecValTok{0}
\ControlFlowTok{for}\NormalTok{ prefix }\KeywordTok{in}\NormalTok{ model\_generated\_positions(trajectory):}
\NormalTok{    privileged\_prefix }\OperatorTok{=}\NormalTok{ add\_verified\_answer(prefix)}
\NormalTok{    target }\OperatorTok{=}\NormalTok{ stop\_gradient(model.distribution(privileged\_prefix))}
\NormalTok{    loss }\OperatorTok{+=}\NormalTok{ KL(model.distribution(prefix), target)}
\NormalTok{update(model, loss)}
\end{Highlighting}
\end{Shaded}

Привилегированная информация поступает только в ветвь учителя во время
обучения, тогда как входные данные ученика при обучении и оценке
остаются такими же, как при развёртывании. По псевдокоду можно проверить
два пути прохождения информации: эталонный ответ поступает только в
строку, создающую префикс учителя, и при развёртывании ученик его не
увидит; к выводу учителя применяется остановка градиента, поэтому на
этом шаге по различию двух распределений обновляются параметры только
ветви ученика. Когда в следующем раунде используется обновлённая модель,
суждения учителя также могут измениться. Таким образом, одна модель
может использовать более полную информацию, доступную на стороне
обучения, чтобы предоставлять надзор для собственной генерации в
реальных условиях.

Источником привилегированной информации также могут служить демонстрации
людей или прошедшая отбор предметная документация. При выборе следует
проверять, улучшает ли такая информация суждения учителя: чем лучше
учитель после её получения распознаёт ошибки ученика, тем полезнее
сигнал дистилляции. В процессе обучения также необходимо на заранее
отложенных тестовых образцах проверять исходные возможности и стиль
генерации. В исследованиях наподобие Purified OPSD дополнительно
применяются ограничения регуляризации, помогающие модели сохранить
первоначальный стиль рассуждений\footnote{Shen, Ziqi, et al.~``Purified
  OPSD: On-Policy Self-Distillation Without Losing How to Think'', 2026.
  arXiv:2607.02234.}.

\hypertarget{ux43eux442-ux43fux440ux43eux431ux43bux435ux43cux43dux44bux445-ux441ux43bux443ux447ux430ux435ux432-ux43a-ux43fux43eux441ux442ux43eux431ux443ux447ux435ux43dux438ux44e}{%
\section{От проблемных случаев к
постобучению}\label{ux43eux442-ux43fux440ux43eux431ux43bux435ux43cux43dux44bux445-ux441ux43bux443ux447ux430ux435ux432-ux43a-ux43fux43eux441ux442ux43eux431ux443ux447ux435ux43dux438ux44e}}

В \crossreflink{chap:7}{главе 7} проблемы из продакшена были
преобразованы в воспроизводимые задачи, результаты атрибуции сбоев и
критерии оценки. Эти материалы также можно использовать для
постобучения: полные задачи предоставляют среду взаимодействия, префиксы
траекторий до возникновения ошибки --- контекст принятия решений, а
валидаторы проверяют вновь сгенерированное поведение. В таблице 8-5
показано их соответствие методам обучения.

Таблица 8-5. Материалы для оценки из \crossreflink{chap:7}{главы 7} и их
применение при постобучении

\begin{longtable}[]{@{}
  >{\raggedright\arraybackslash}p{(\columnwidth - 2\tabcolsep) * \real{0.3699}}
  >{\raggedright\arraybackslash}p{(\columnwidth - 2\tabcolsep) * \real{0.6301}}@{}}
\toprule\noalign{}
\begin{minipage}[b]{\linewidth}\raggedright
\hspace{0pt}Материал из \crossreflink{chap:7}{главы 7}
\end{minipage} & \begin{minipage}[b]{\linewidth}\raggedright
\hspace{0pt}Способ применения при обучении в
\crossreflink{chap:8}{главе 8}
\end{minipage} \\
\midrule\noalign{}
\endhead
\bottomrule\noalign{}
\endlastfoot
\hspace{0pt}Сквозные задачи и валидаторы & \hspace{0pt}Задачи для
сэмплирования в RL, проверяемое вознаграждение, пул задач для выборки с
отклонением \\
\hspace{0pt}Префиксы траекторий и последующее поведение &
\hspace{0pt}Пары предпочтений для DPO, демонстрации решений для SFT,
входные данные учителя для
on-\allowbreak{}\hspace{0pt}policy-\allowbreak{}\hspace{0pt}дистилляции \\
\hspace{0pt}Место сбоя и категория ошибки & \hspace{0pt}Метки оценки
процесса, проверяемые правила штрафов за траекторию \\
\hspace{0pt}Оценки по отдельным критериям и эталонные ответы с ручной
разметкой & \hspace{0pt}Многомерное вознаграждение, обучение и
калибровка генеративной модели вознаграждения \\
\end{longtable}

Исторический префикс можно использовать как исходную точку задачи на
продолжение. При on-policy-дистилляции текущая модель-ученик сначала
продолжает генерацию с этой точки, после чего модель-учитель оценивает
фактически созданный учеником префикс; если требуется выполнение
инструментов, одновременно восстанавливается соответствующее состояние
среды. Так оценка отдельных решений из \crossreflink{chap:7}{главы 7}
связывается с описанным в этой главе процессом сэмплирования и обучения.

Обучающие задачи необходимо отделять от отложенных задач для оценки. Обе
стороны могут использовать общие интерфейсы и логику валидации, но
экземпляры задач должны быть независимыми; если требуется проверить
обобщение между задачами, разделение также можно выполнить по шаблонам
задач или типам бизнес-процессов. Ниже на трёх примерах показано, как
выполнить такое преобразование.

\hypertarget{ux441ux43bux443ux447ux430ux439-1-ux43fux440ux435ux436ux434ux435ux432ux440ux435ux43cux435ux43dux43dux43eux435-ux437ux430ux432ux435ux440ux448ux435ux43dux438ux435-ux440ux430ux431ux43eux442ux44b-coding-agent}{%
\subsection{Случай 1: преждевременное завершение работы Coding
Agent}\label{ux441ux43bux443ux447ux430ux439-1-ux43fux440ux435ux436ux434ux435ux432ux440ux435ux43cux435ux43dux43dux43eux435-ux437ux430ux432ux435ux440ux448ux435ux43dux438ux435-ux440ux430ux431ux43eux442ux44b-coding-agent}}

К преждевременному завершению относятся ситуации, когда агент объявляет
о завершении, ещё не запустив тесты, отправляет результат после
выполнения лишь двух из трёх требований или сдаётся после нескольких
неудач. В продакшене такие случаи можно собирать по исправлениям
пользователей, отрицательным оценкам и результатам аудита траекторий.
Например, если пользователь указывает: «Ты ещё не запускал тесты», по
траектории можно проверить состояние задачи на тот момент, имевшиеся
доказательства прохождения проверки и заявление модели о завершении.

При анализе причины сначала следует остановить траекторию в момент,
когда модель собирается сказать «готово», а затем проверить
предшествующие шаги: не пропустила ли она требования и не тот ли код она
изменила. В результате получаются два типа обучающих материалов: полная
задача используется для обучения модели доводить работу до конца, а
префикс перед решением --- для обучения выбору между продолжением работы
и завершением на основе доказательств.

\textbf{Полную задачу можно связать с вознаграждением за результат.}
Когда модель отправляет результат, запускаются защищённые приёмочные
тесты: за прохождение начисляется +1, за непрохождение --- −1. Тесты
охватывают фактические результаты работы; в процессе модель может
получать обратную связь, запуская обычные тесты для разработки, а
итоговую приёмку выполняет обучающая среда.

\textbf{Из префикса перед решением можно сформировать пару
предпочтений.} В одном и том же контексте отвергнутый пример
(\texttt{rejected}) представляет собой заявление о завершении без
достаточных доказательств, а выбранный (\texttt{chosen}) --- продолжение
работы с запуском необходимых тестов, проверкой пропущенных требований и
завершением оставшейся части. Можно поручить модели-учителю написать
несколько вариантов дальнейших действий, а затем с помощью валидатора
отобрать подходящие демонстрации. Если примеров недостаточно,
посредством замены типа задачи, пропущенного проверяемого пункта и
формулировки завершения можно сформировать несколько сотен пар
предпочтений. После этого выполняется DPO с LoRA, а в данные добавляется
достаточное количество общих задач для сохранения охвата поведения.

\textbf{При оценке одновременно проверяются незавершённые и завершённые
задачи.} В одной группе задач выполнена лишь часть работы --- она
позволяет проверить, продолжит ли модель выполнение; в другой группе всё
уже сделано --- она показывает, способна ли модель своевременно
завершить работу. Одновременно выборочно проверяются общие вызовы
инструментов и другие задачи, чтобы определить, повлияло ли направленное
обучение на существующие возможности. Если снижать только частоту
преждевременного завершения, модель может каждый раз выбирать
продолжение проверки и медлить с завершением даже уже выполненной
работы. В таком случае в группе завершённых задач появятся лишние вызовы
инструментов и ожидание.

\begin{quote}
\textbf{Эксперимент 8-17 ★★: исправление преждевременного завершения
посредством DPO на основе проблемных случаев}

Эксперимент объединяет в единый процесс атрибуцию сбоя, префикс перед
решением, пару предпочтений, обучение LoRA и оценку на отложенном
наборе; для обучения используется модель 7B.

Данные содержат 24 реалистичных проблемных случая четырёх типов:
объявление о завершении без запуска тестов, выполнение лишь части из
нескольких целей, невыполнение условий приёмки и отказ после ошибки. В
них также включён вариант манипулирования вознаграждением с удалением
неуспешного теста. Кроме того, отложены 12 задач, в которых ещё
требуется выполнить работу, и 8 уже завершённых задач; они используются
только для оценки.

Сначала рассмотрим одну отложенную задачу. Потребитель очереди сообщений
теряет сообщения после отключения и повторного подключения. Ассистент
изменил код повторного подключения, а инструмент сообщил лишь: «Файл
изменён». Следует ли теперь объявить проблему исправленной или
продолжить проверку? В эксперименте предлагаются два варианта ответа и
сравнивается, какому из них модель присваивает более высокую оценку. До
обучения модель выбрала правильный ответ в 3 из 12 незавершённых задач
--- 25,0\%; после обучения --- в 11 из 12, то есть в 91,7\%. Во всех 8
действительно завершённых задачах модель выбрала правильный ответ как
до, так и после обучения.

Если убрать варианты ответа и попросить модель самостоятельно написать
следующий шаг, результат меняется. Согласно правилам текстовой
классификации эксперимента, в незавершённых задачах число
преждевременных заявлений о завершении снизилось с 1/12 до 0/12; однако
в завершённых задачах число явных корректных завершений также снизилось
с 6/8 до 0/8. Если смотреть только на первую группу, легко решить, что
проблема устранена. Но вторая группа показывает: модель распознаёт
момент остановки среди двух готовых ответов, однако при самостоятельной
генерации всё ещё может неизменно требовать дополнительных проверок. В
примере с потерей сообщений ассистенту также необходимо действительно
выполнить тест с отключением и повторным подключением и проверить,
поступают ли сообщения после восстановления соединения. Сказать
«продолжу проверку» и получить результат проверки --- два разных шага.
\end{quote}

\hypertarget{ux441ux43bux443ux447ux430ux439-2-ux43aux430ux432ux44bux447ux43aux438-ux432-ux43aux438ux442ux430ux439ux441ux43aux43eux43c-ux442ux435ux43aux441ux442ux435}{%
\subsection{Случай 2: кавычки в китайском
тексте}\label{ux441ux43bux443ux447ux430ux439-2-ux43aux430ux432ux44bux447ux43aux438-ux432-ux43aux438ux442ux430ux439ux441ux43aux43eux43c-ux442ux435ux43aux441ux442ux435}}

Требование «заменить прямые кавычки в статье на китайском языке
типографскими» сначала требует определить область редактирования. В
цитатах на китайском естественном языке можно использовать \texttt{“”},
а вложенные цитаты следует оформлять по правилам китайской пунктуации;
кавычки же в английском оригинале, строковых литералах, JSON, путях и
идентификаторах выполняют соответствующие синтаксические функции. Цель
обучения --- научить модель распознавать эти области и выполнять
минимальное редактирование.

Сначала правила можно оформить как Skill, а затем на их основе
генерировать исходные и целевые тексты. Положительные примеры включают
абзацы на китайском языке, вложенные цитаты и комментарии на китайском
языке; защищаемые примеры --- цитаты на английском, строковые и
символьные литералы, JSON, встроенный код и блоки кода. Для комментариев
на китайском языке внутри кода необходимо одновременно распознавать
границы комментария и область цитирования; фрагменты с неопределимой
областью следует сохранять без изменений. Только сочетание этих случаев
в одном документе позволяет отрабатывать определение области при
реальном редактировании.

\begin{quote}
\textbf{Эксперимент 8-18 ★★: SFT для типографских кавычек в китайском
тексте с учётом области действия}

В эксперименте используется Qwen3-8B с LoRA SFT в формате bf16; обучение
длится 2 эпохи и включает 256 обновлений. Цель обучения ---
преобразовывать цитаты в естественном языке на китайском, сохраняя
синтаксические кавычки в коде и структурированных данных.

На основе 16 типов фрагментов, 10 жанров статей и 9 языков
программирования формируются 1 024 обучающих, 256 отложенных и 256
граничных примеров. В каждом примере сохраняются исходный и целевой
тексты, охватывающие цитаты, которые необходимо преобразовать, строки,
которые требуется защитить, а также случаи их вложения друг в друга.

После обучения доля точных совпадений на отложенном наборе составляет
96,9\%, а на граничном --- 97,7\%; доля сохранённых без изменений
защищённых областей равна 100\%.

Один и тот же набор правил областей действия используется для
формирования меток и проверки вывода. Модель учится определять область
редактирования, а парсер и синтаксические проверки отвечают за валидацию
структуры изменённого документа. Это позволяет раздельно проверять
качество редактирования текста и корректность форматирования.
\end{quote}

\hypertarget{ux441ux43bux443ux447ux430ux439-3-ux447ux430ux441ux442ux44bux435-ux441ux431ux43eux438-ux43fux440ux438-ux440ux435ux434ux430ux43aux442ux438ux440ux43eux432ux430ux43dux438ux438-ux444ux430ux439ux43bux43eux432}{%
\subsection{Случай 3: частые сбои при редактировании
файлов}\label{ux441ux43bux443ux447ux430ux439-3-ux447ux430ux441ux442ux44bux435-ux441ux431ux43eux438-ux43fux440ux438-ux440ux435ux434ux430ux43aux442ux438ux440ux43eux432ux430ux43dux438ux438-ux444ux430ux439ux43bux43eux432}}

В \crossreflink{chap:5}{главе 5} были представлены инструменты наподобие
\texttt{edit\_file(path,\ old\_string,\ new\_string)}. Модель должна
скопировать исходный текст в \texttt{old\_string}, после чего инструмент
выполняет точное сопоставление. Любое различие в пробелах, переводах
строк, обратных косых чертах или символах Unicode может привести к сбою
сопоставления.

Сначала следует последовательно сравнить каждый этап передачи
информации: \textbf{исходные байты файла → результат инструмента →
сериализация Harness → контекст модели → вывод токенов модели →
декодированная строка → разбор аргументов вызова → сопоставление
инструментом}. Если байты изменились на этапе чтения, следует исправить
инструмент; если содержимое изменилось при экранировании или
формировании контекста --- Harness; если изменение возникло при
кодировании с последующим декодированием --- проверить tokenizer. Только
если входные данные модели остались неизменными, а первое различие
появилось при генерации моделью, этот случай можно использовать для
постобучения способности к точному копированию.

Обучающие данные могут содержать три типа задач: дословно воспроизвести
указанное содержимое; выбрать полностью совпадающий вариант среди
нескольких похожих строк одинаковой длины; записать исходный текст в
поле \texttt{old\_string} вызова инструмента. В примеры включаются
пробелы, настоящие переводы строк, обратные косые черты и разные
представления Unicode, чтобы охватить места, наиболее подверженные
ошибкам при реальном редактировании.

\begin{quote}
\textbf{Эксперимент 8-19 ★★: SFT для точного копирования специальных
строк}

В качестве базовой модели в эксперименте используется Qwen3-8B, которая
обучается с bf16 LoRA в течение 2 эпох; промпт исключается из расчёта
функции потерь. В задачах на непосредственное воспроизведение
демонстрацией служит целевая строка, а в задачах с аргументами
инструмента --- полный JSON, содержащий только \texttt{old\_string};
имена полей, кавычки и скобки также включаются в надзор на уровне
каждого токена. Кроме того, отдельно проверяется согласованность
кодирования и декодирования tokenizer.

Доля побайтовых точных совпадений модели на отложенном наборе возросла с
37,5\% до 78,9\%, а на независимом граничном наборе составила 80,1\%;
средняя позиция первого расхождения байтов в этих двух наборах равна
соответственно 54,0 и 54,2. Дополнительно на 512 пробах из отложенного и
граничного наборов были сопоставлены три tokenizer; для Qwen3 и Qwen2.5
доля строк, прошедших round-trip без потерь, составила 80,1\%.

После дообучения модель правильно решила на 106 задач больше, но
остались ли правильными все задачи, которые она решала верно раньше? При
попарном сравнении двух ответов на каждую задачу выясняется, что в 107
случаях неправильный ответ стал правильным, а в одном случае правильный
стал неправильным. В первой задаче отложенного набора требовалось
вернуть JSON-объект. Базовая модель правильно скопировала
\texttt{old\_string}, но добавила лишние поля, включая путь к файлу;
после дообучения она вернула только запрошенное поле и получила балл. В
задаче с регрессией (\texttt{eval-095}) требовалось без изменений
воспроизвести две строки текста, однако ответ дообученной модели
объединил их, пропустив разделяющий перевод строки.
\end{quote}

Для каждого из этих трёх случаев следует составить диагностическую
карточку: сохранить входные данные, которые модель действительно видела
до сбоя, указать как минимум одно допустимое действие, а затем добавить
контрастный пример, который «выглядит похоже, но не должен выполняться
аналогично». Например, при конфликте редактирования следует повторно
прочитать файл, а если файл не изменился и правка уже проверена, работу
можно завершить. Обучение на первых примерах с последующей приёмкой на
не участвовавших в обучении однотипных и контрастных задачах позволяет
проверить, научилась ли модель принимать решения с учётом условий или
просто всегда выполняет ещё один шаг.

\hypertarget{ux43fux440ux430ux43aux442ux438ux447ux435ux441ux43aux438ux435-ux440ux435ux43aux43eux43cux435ux43dux434ux430ux446ux438ux438-ux43fux43e-ux43fux43eux441ux442ux43eux431ux443ux447ux435ux43dux438ux44e}{%
\section{Практические рекомендации по
постобучению}\label{ux43fux440ux430ux43aux442ux438ux447ux435ux441ux43aux438ux435-ux440ux435ux43aux43eux43cux435ux43dux434ux430ux446ux438ux438-ux43fux43e-ux43fux43eux441ux442ux43eux431ux443ux447ux435ux43dux438ux44e}}

Схему постобучения следует выбирать исходя из решаемой проблемы.
Недостаток предметных знаний, нестабильность протокола вывода и
неудачный выбор действий требуют разных данных и обратной связи. Автор
придаёт большее значение качеству данных и среды: они определяют, с
какими проблемами столкнётся модель и какие сигналы получит, а алгоритм
обучения уже преобразует эти сигналы в обновления параметров.

В проекте проверки можно проводить в следующем порядке.

\begin{enumerate}
\def\labelenumi{\arabic{enumi}.}
\tightlist
\item
  \textbf{Определите место хранения знаний и способностей.} Стабильные
  предметные корпуса можно использовать для Mid-training, а демонстрации
  способны научить модель отвечать и действовать; фактами, требующими
  своевременного обновления, цитирования, контроля доступа или удаления,
  следует управлять с помощью RAG.
\item
  \textbf{Проверьте, можно ли оценивать вывод.} На небольшом количестве
  задач измерьте долю успешного разбора JSON, полноту параметров и долю
  успешных вызовов инструментов. Если ошибки формата возникают часто,
  сначала улучшите вывод посредством SFT или декодирования с
  ограничениями и лишь затем увеличивайте объём выборки для RL.
\item
  \textbf{Проверьте эффективную длину контекста.} Даже если модель
  способна принимать 128K токенов, необходимо проверить, может ли она на
  такой длине выполнять поиск, рассуждать и планировать. Автор проводит
  поэтапную приёмку по матрице «способность × длина», затем расширяет
  окно, подмешивая короткие примеры и данные предыдущих этапов, чтобы
  сохранить способности на коротких контекстах.
\item
  \textbf{Проверьте вероятность выборки эффективных траекторий.}
  Совместно отслеживайте \texttt{pass@k}, частичный прогресс и различия
  вознаграждений внутри группы. Если успешных примеров крайне мало,
  лучшую отправную точку можно создать посредством Mid-training, SFT,
  дистилляции или учебных задач; частичное вознаграждение должно
  соответствовать прогрессу, достижимому для текущей стратегии.
\item
  \textbf{Обеспечьте охват целей задачи вознаграждением.} Если оценка
  зависит только от длины ответа, модель может начать добавлять
  многословный текст. Преобразуйте требования приёмки в тесты, проверки
  состояния и откалиброванные критерии оценки, после чего проверяйте
  фактическое качество траекторий с высокой оценкой.
\item
  \textbf{Откалибруйте обратную связь среды.} Сопоставьте формулировки
  пользователей, ошибки инструментов и изменения состояния в симуляции и
  при развёртывании, а обнаруженные различия добавьте в обучающие
  задачи. Стоимость создания стабильной среды также следует учитывать в
  бюджете обучения.
\item
  \textbf{Контролируйте обобщение и забывание на отложенных задачах.}
  Одновременно отслеживайте целевую способность и исходные общие
  способности. Если обучающие метрики продолжают улучшаться, а качество
  на отложенных задачах снижается, скорректируйте соотношение данных,
  количество циклов обучения или регуляризацию и выберите контрольную
  точку по результатам валидации.
\item
  \textbf{Наблюдайте за оценкой ценности и исследованием.} Ошибка оценки
  ценности в PPO влияет на вычисление преимущества; температура и
  случайность выборки влияют на разнообразие возможных путей. Совместное
  рассмотрение колебаний вознаграждения, ошибки ценности, энтропии
  вывода и различий между траекториями помогает понять, почему обучение
  застопорилось.
\item
  \textbf{Измеряйте численную согласованность выборки и обучения.} До
  обновления сравнивайте log probability на обеих сторонах, отслеживайте
  отношение вероятностей, приближённый KL, долю отсечения и устаревание
  стратегии, а также согласуйте веса, шаблоны и вычислительную
  конфигурацию.
\item
  \textbf{Оценивайте стоимость всей цепочки.} Стоимость RL включает
  генерацию моделью, выполнение инструментов, проверку, обновление
  параметров и ожидание. Сопоставление времени в ReTool показывает, что
  выборка может значительно увеличить продолжительность цикла обучения.
  Если SFT уже обеспечивает целевое качество, ресурсы можно в первую
  очередь направить на пополнение данных и оценку развёртывания.
\item
  \textbf{Проверяйте качество обучающих сигналов для каждого типа
  данных.} Ошибочные связи в корпусе Mid-training, шум в демонстрациях
  SFT и смещение верификатора RL попадут в обучающий сигнал. Сначала
  вручную проверьте небольшую партию данных и траекторий и лишь затем
  увеличивайте масштаб.
\end{enumerate}

Эти проверки можно объединить в недорогое пробное обучение: на небольшом
объёме предметного корпуса проверить прирост знаний и забывание, на
малом количестве демонстраций --- обучаемость поведения, а на небольшой
партии rollout --- различия вознаграждений и эксплуатационные расходы.
Результаты пробного обучения послужат основанием для дальнейших
вложений.

\textbf{Постобучение также можно сочетать с RAG и обучением в
контексте.} Предположим, вы создаёте помощника по денежным переводам.
Сначала можно показать модели несколько демонстрационных диалогов, чтобы
она попробовала понятно объяснить пользователю процедуру одобрения; чем
больше примеров, тем выше могут оказаться время и стоимость их обработки
при каждом запросе. Когда политика изменится, актуальные положения можно
получить из базы знаний и сослаться на источник в ответе. На этапе
отправки перевода сумма и полномочия должны проверяться программой, даже
если лимит одобрения не менялся много лет. При выборе метода следует
учитывать не только частоту изменения информации, но и возможность
выразить основания решения в виде правил. Естественный тон речи и
морфологические особенности на медицинских изображениях сложно
исчерпывающе описать правилами, поэтому обучение таким закономерностям
зачастую требует обучения параметров.

Система может сначала проверять стратегию с помощью контекста,
предоставлять доказательства посредством RAG, программно обеспечивать
жёсткие ограничения, а затем улучшать способности модели с помощью
Mid-training, SFT и RL. После достижения моделью целевого качества её
поведение также можно перенести в меньшую модель посредством
дистилляции, снизив стоимость развёртывания.

Как проверить, что «модель научилась принимать решения на основе
обратной связи»? В своём выступлении в Порту автор назвал открытой
проблемой непрерывное обучение модели принятию решений в её весах при
одновременном сохранении прежних способностей. Запоминание исправленного
пользователем остатка проверяет память на факты. Чтобы выяснить,
научилась ли модель принимать решения, следует задать вопрос в другом
сценарии.

Например, обучающие материалы описывают ситуацию, в которой для
пользователя в командировке допустим авиабилет, немного превышающий
бюджет. При тестировании замените его незнакомым отелем и новой суммой:
в одной задаче сохраните ограничение по времени, необходимое для
прибытия на конференцию, а в другой явно укажите, что речь идёт о
семейном отпуске при ограниченном бюджете. В обеих задачах предоставьте
доступные на тот момент варианты и заранее определите допустимые
действия и объяснения. Если в обеих задачах модель неизменно выбирает
более дорогой вариант или неизменно задаёт уточняющие вопросы, это может
указывать на то, что она запомнила лишь поверхностное действие. После
этого повторно выполните задачи, с которыми прежняя версия уже
справлялась, например задачу со строгим соблюдением явно заданного
пользователем предельного размера возмещения.

\hypertarget{ux438ux442ux43eux433ux438-ux433ux43bux430ux432ux44b-7}{%
\section{Итоги
главы}\label{ux438ux442ux43eux433ux438-ux433ux43bux430ux432ux44b-7}}

Столкнувшись с задачей, которую модель не может решить, сначала дайте ей
несколько попыток и посмотрите, где именно она терпит неудачу. Если
после разумного сэмплирования \texttt{pass@k} по-прежнему близок к нулю,
сначала проверьте, хватает ли модели предметных знаний и базовых
способностей, а затем решите, нужно ли дополнить корпус с помощью
Mid-training. Если модель не может прочитать длинный документ целиком,
поэтапно увеличивайте длину обучения и проверяйте сохранность прежних
способностей на коротких задачах. Если модель иногда отвечает правильно,
но часто записывает параметры инструментов в формате, не поддающемся
синтаксическому разбору, сначала стоит предоставить ей демонстрации с
помощью SFT. RL получает возможность повысить вероятность выбора хороших
стратегий лишь после того, как модель сможет создавать разнообразные
исполняемые траектории, а вознаграждение --- различать, какие из них
лучше.

\textbf{Автор в первую очередь вкладывался бы в данные и среду.} В
экспериментах GeneralPoints по переносу правил и визуальному переносу
обученная на ограниченном числе демонстраций модель SFT легко
переобучается, тогда как ветвь RL лучше справляется с задачами вне
обучающего распределения. При замене набора демонстраций или способа
верификации такое сравнение необходимо провести заново. В многошаговых
задачах наподобие ReTool, V-IRL и SimpleVLA-RL также приходится решать,
какие события при вызовах инструментов, навигации или управлении
действиями заслуживают обратной связи. Если реальная среда слишком
дорога, можно симулировать возвращаемые инструментами значения и
переходы между состояниями, однако затем необходимо проверить на
реальных задачах, не научилась ли модель эксплуатировать погрешности
симулятора. Для некоторых задач достаточно восполнить базовые знания и
предоставить демонстрации, чтобы уже выполнить требования к поставке.

\textbf{По мнению автора, главные ограничения современного RL для
AI-агентов --- эффективность использования выборок и согласованность
распределений.} Если одно взаимодействие с инструментом обходится
дорого, нельзя ли извлечь из него больше обучающего сигнала? On-Policy
Distillation позволяет учителю предоставлять поточечные указания на
уровне токенов в тех состояниях, до которых фактически дошёл ученик;
RLVP, в свою очередь, дополняет итоговое вознаграждение проверяемыми
промежуточными событиями. С другой стороны, если стратегия,
использованная при сэмплировании, уже отстаёт от обучаемой версии либо
сэмплер и обучающая сторона вычисляют заметно различающиеся вероятности
для одного и того же фрагмента вывода, сначала следует исследовать это
расхождение и лишь затем интерпретировать кривую вознаграждения.

Преждевременное завершение, редактирование неверного диапазона и
неспособность точно воспроизвести строку --- каждый такой сбой можно
проследить до конкретного обучающего примера. После устранения ошибки на
этой задаче необходимо провести регрессионную проверку на независимых
примерах и задачах, которые модель прежде решала правильно. Когда
появится следующая группа сбоев, какие из них стоят того, чтобы ещё раз
обучить модель, а для каких достаточно обновить знания, инструкции или
программу? Именно этот выбор рассматривается в следующей главе.

\begin{questionbox}

\hypertarget{ux432ux43eux43fux440ux43eux441ux44b-ux434ux43bux44f-ux440ux430ux437ux43cux44bux448ux43bux435ux43dux438ux44f-7}{%
\section*{Вопросы для
размышления}\label{ux432ux43eux43fux440ux43eux441ux44b-ux434ux43bux44f-ux440ux430ux437ux43cux44bux448ux43bux435ux43dux438ux44f-7}}
\addcontentsline{toc}{section}{Вопросы для размышления}

\begin{enumerate}
\def\labelenumi{\arabic{enumi}.}
\tightlist
\item
  ★★ Направленное дообучение улучшило результаты на целевой задаче, но
  снизило общую способность вызывать инструменты. Как бы вы разработали
  набор тестов для проверки исходных возможностей, скорректировали
  соотношение данных и объём обучения, а также сравнили результаты
  полного дообучения параметров и LoRA?
\item
  ★★ Некоторую способность в конкретной предметной области можно освоить
  как посредством обучения параметров, так и с помощью примеров в
  контексте и поиска знаний. На основании каких факторов вы бы выбрали
  подход?
\item
  ★★ Какие обучающие материалы и методы проверки необходимы для
  дистилляции Chat-модели, отвечающей напрямую, Reasoning-модели,
  генерирующей процесс рассуждения, и Agentic-модели, использующей
  инструменты в несколько раундов?
\item
  ★★★ AI-агент службы поддержки выполняет задачу за 10 раундов
  взаимодействия и в конце получает вознаграждение за успех. Как бы вы
  связали этот сигнал с промежуточными решениями? Для каких шагов
  уместно добавить вознаграждение за процесс?
\item
  ★★★ У вас есть бюджет в 10 000 долларов на повышение качества работы
  агента службы поддержки. Как бы вы распределили вложения между
  оцениванием, контекстом и знаниями, промптами и навыками, программными
  ограничениями и обучением параметров?
\item
  ★★★ У задачи нет явно заданного вознаграждения, а взаимодействие со
  средой обходится дорого. Из каких источников можно получить обучающие
  сигналы? Как бы вы проверили, что эти сигналы действительно улучшают
  выполнение задачи?
\item
  ★★ Какие персонализированные способности уместно реализовывать
  посредством отдельной LoRA для каждого пользователя или каждой
  компании? Какую информацию лучше хранить во внешней базе знаний?
\item
  ★★★ С учётом исследовательского подхода Weak-to-Strong Generalization
  может ли малая модель предоставлять полезный сигнал для обучения более
  сильного агента? Рассмотрите по отдельности локальную проверку, метки
  предпочтений и дополнительный контекст, а также разработайте метод
  выявления смещения в обучающем сигнале.
\item
  ★★ В одной траектории большинство промежуточных шагов правильны, но
  итоговый результат ошибочен. Другая траектория приводит к правильному
  результату, но нарушает ограничения процесса. Как оценить эти две
  траектории, совместив оценку процесса и результата?
\item
  ★★★ Среды и валидаторы SWE-bench, τ²-bench и AndroidWorld можно
  использовать для обучения. Как разделить экземпляры и шаблоны задач,
  чтобы задействовать обучающие материалы и при этом сохранить
  независимую оценку?
\item
  ★★★ У целевой задачи очень низкий показатель
  \href{mailto:pass@1}{\nolinkurl{pass@1}}. Как, объединив
  \href{mailto:pass@k}{\nolinkurl{pass@k}}, долю случаев успешного
  разбора формата, долю частичного прогресса и анализ причин неудач,
  решить, следует ли применять Mid-training, SFT или RL?
\item
  ★★★ В пакете роллаутов осталось лишь несколько незавершённых
  сверхдлинных ответов, тогда как остальные задачи уже выполнены. Как
  повысить утилизацию обучающего кластера, одновременно контролируя
  устаревание политики, вызванное асинхронным сэмплированием?
\item
  ★★★ Какие смещения симуляции могут вводить агента в заблуждение при
  использовании модели для имитации поисковой системы или пользователя?
  Как откалибровать обучение с помощью состояния задачи и реальных
  взаимодействий?
\end{enumerate}

\textbf{Практикум к \crossreflink{chap:8}{главе 8}.}
\href{https://github.com/bojieli/ai-agent-book/tree/40f4ce5778e6cc93e3606c3079bf8cf13a0c4d04/chapter8}{Каталог
экспериментов} ·
\href{https://github.com/bojieli/ai-agent-book/blob/40f4ce5778e6cc93e3606c3079bf8cf13a0c4d04/chapter8/README.ru.md}{Путеводитель
на русском}.

\end{questionbox}

\hypertarget{ux43dux435ux43fux440ux435ux440ux44bux432ux43dux430ux44f-ux44dux432ux43eux43bux44eux446ux438ux44f-ai-ux430ux433ux435ux43dux442ux430}{%
\chapter{\texorpdfstring{Непрерывная эволюция
AI-агента\label{chap:9}}{Непрерывная эволюция AI-агента}}\label{ux43dux435ux43fux440ux435ux440ux44bux432ux43dux430ux44f-ux44dux432ux43eux43bux44eux446ux438ux44f-ai-ux430ux433ux435ux43dux442ux430}}

Агент способен с первой попытки выполнить незнакомую сложную задачу, но
даже после обработки множества похожих задач может продолжать повторять
одну и ту же ошибку. Чтобы он, подобно новому сотруднику, постоянно
совершенствовался в ходе повседневной работы, система должна
преобразовывать опыт выполнения в возможности, пригодные для
последующего использования: фиксировать результаты, находить причины
проблем, формировать способы улучшения, а затем проверять их
эффективность на новых задачах. \textbf{По мнению автора, способность
агента учиться на опыте и непрерывно эволюционировать становится одной
из важнейших.} Эта способность требует инженерного развития прикладной
системы и заслуживает внимания как одно из приоритетных направлений
исследований моделей следующего поколения.

Как правило, после завершения одного инференса параметры развёрнутой
модели остаются неизменными. Рассмотренные во второй главе обучение в
контексте, поддержание состояния и сжатие помогают ей адаптироваться в
рамках текущей задачи, а описанные в третьей главе память и база знаний
позволяют сохранять информацию между задачами. Для дальнейшего обучения
на опыте необходимо также сопоставлять результаты множества выполнений:
какие шаги являются общими для успешных траекторий, чего не хватает
неудачным траекториям, для каких версий инструментов и условий задачи
подходит тот или иной способ действий.

Например, сохранив сто траекторий бронирования билетов, система сможет
при необходимости найти конкретный случай из прошлого. Но чтобы
сформировать правило «перед бронированием билета сначала проверять срок
заказа специального питания», ей потребуется сопоставить причины
успешных и неудачных бронирований, проверить условия авиакомпании и
ограничения по времени, а затем подтвердить этот способ действий на
новой задаче. Запись и поиск позволяют агенту восстановить прошлый опыт,
а сопоставление, обобщение и проверка помогают сформировать пригодный
для повторного использования способ действий. Лишь накапливая такой
опыт, агент может постепенно превращаться из просто умного в опытного.

\textbf{Непрерывное обучение (continual learning)} исследует, как
система может постоянно учиться по мере поступления новых задач и
данных, сохраняя при этом уже приобретённые возможности. Агент переносит
эту проблему в повседневное взаимодействие: обратная связь поступает от
пользователей, инструментов и среды, причём удовлетворённость
пользователя, выполнение задачи и соответствие процесса требованиям
могут приводить к разным выводам. По мнению автора, одна из главных
трудностей непрерывного обучения агента состоит в том, что в
повседневной работе почти никто не помечает пошагово, что было сделано
правильно, а что --- нет, поэтому система должна сама определять, какой
обратной связи можно доверять. Если непроверенная обратная связь будет
преобразована в правило или использована для обучения, ошибочная
атрибуция причин повлияет на последующие задачи. Такие ошибки могут
сохраняться и дальше: неверно записанное предпочтение повлияет на
следующее бронирование билета, а обучение модели на ошибочной обратной
связи может привести к неправильному выполнению даже тех задач, с
которыми она прежде справлялась.

В этой книге термином \textbf{непрерывная эволюция} обозначается весь
процесс обновления системы агента на основе опыта, охватывающий четыре
вида носителей: знания, инструкции, программы и параметры модели.
Периодическое обучение базовой модели обеспечивает общие возможности,
тогда как обновление агента на основе опыта учитывает конкретные
бизнес-правила, изменения инструментов и локальные стратегии. Такие
локальные изменения происходят постоянно, и ожидание следующего обучения
базовой модели зачастую не соответствует потребностям бизнеса.

На текущем этапе автор предлагает в первую очередь создавать вокруг
модели проверяемую обучающую систему: явно выражаемый опыт должен влиять
на последующее поведение через знания, инструкции и программы, а
материалы, подходящие для обучения параметров, следует направлять на
постобучение. Каждое обновление сначала проходит регрессионную оценку в
новой версии, и лишь после успешной проверки выпуска последующие задачи
начинают использовать эту версию.

В предыдущих главах уже были представлены основные компоненты этой
системы: вторая глава посвящена организации состояния задачи, третья ---
хранению и поиску знаний, пятая --- созданию инструментов с помощью кода
и изменению системы, седьмая --- оценке результатов, а восьмая ---
обновлению параметров модели. На \crossreflink{fig:9-1}{рис. 9-1} эти
компоненты объединены. Опыт, полученный при предыдущем бронировании,
когда не было заказано специальное питание, после оценки превращается в
правило, которое будет использовано при следующем бронировании билета.

\begin{figure}
\hypertarget{fig:9-1}{%
\centering
\includegraphics{images/revised/fig9-1.pdf}
\caption{Рисунок 9-1. Общий замкнутый цикл непрерывной эволюции
агента}\label{fig:9-1}
}
\end{figure}

\hypertarget{ux43fux43eux43bux443ux447ux435ux43dux438ux435-ux43eux431ux443ux447ux430ux44eux449ux438ux445-ux441ux438ux433ux43dux430ux43bux43eux432-ux438ux437-ux442ux440ux430ux435ux43aux442ux43eux440ux438ux439-ux432ux44bux43fux43eux43bux43dux435ux43dux438ux44f}{%
\section{Получение обучающих сигналов из траекторий
выполнения}\label{ux43fux43eux43bux443ux447ux435ux43dux438ux435-ux43eux431ux443ux447ux430ux44eux449ux438ux445-ux441ux438ux433ux43dux430ux43bux43eux432-ux438ux437-ux442ux440ux430ux435ux43aux442ux43eux440ux438ux439-ux432ux44bux43fux43eux43bux43dux435ux43dux438ux44f}}

Какие траектории пригодны для обучения? Если сотрудник поддержки
сообщает пользователю: «Возврат оформлен», сначала нужно проверить,
сколько денег возвращено и соблюдены ли правила при выполнении операции,
а затем --- сможет ли пользователь понять из ответа, когда поступят
деньги. На \crossreflink{fig:9-2}{рисунке 9-2} объединены все три
проверки. Для них используются соответственно данные заказа, журнал
действий и исходный ответ сотрудника поддержки.

\begin{figure}
\hypertarget{fig:9-2}{%
\centering
\includegraphics{images/revised/fig9-2.pdf}
\caption{Рисунок 9-2. Трёхуровневая проверка траектории: от результата в
среде до рубрики LLM}\label{fig:9-2}
}
\end{figure}

\textbf{Проверка результата определяет, что изменилось в среде.}
AI-агент для программирования может запускать тесты, проверку типов и
бенчмарки производительности; агент, оформляющий возврат средств, может
проверить состояние заказа и фактическую сумму возврата. Автор обычно
начинает с этого уровня проверки. Если модель заявляет о полном
возврате, а на счёт поступила лишь часть суммы, данные заказа позволяют
понять, какую именно операцию нужно проверить.

\textbf{Проверка процесса определяет, соответствуют ли действия правилам
задачи.} Даже после успешного прохождения тестов нужно убедиться, что
защищённые тестовые файлы не были изменены. Когда агент сообщает
пользователю: «Возврат поступит в течение семи дней», необходимо также
проверить, была ли подана заявка на возврат и обоснован ли указанный
срок. Таблицы прав доступа, структурированные политики и журналы
действий можно проверять программно; положения, требующие интерпретации
контекста, следует оценивать вместе с соответствующими доказательствами.

\textbf{Проверка качества определяет, соответствует ли результат
требованиям к использованию.} На итоговый эффект влияют ясность ответа
сотрудника поддержки, соответствие альтернативного решения цели
пользователя и наличие ключевых доказательств в исследовательском
отчёте. Здесь можно применять представленный в седьмой главе подход
LLM-as-a-Judge: заранее определить рубрику (Rubric) и потребовать от
оценивающей модели выставить по каждому пункту балл, привести
обоснование и указать доказательства.

В таблице 9-1 эти три уровня раскрыты на примере агента службы поддержки
через семь аспектов. Результат задачи относится к уровню результата;
соблюдение правил, конфиденциальность и обещания --- к уровню процесса;
качество изложения и допустимые альтернативы --- к уровню качества.
Достоверность фактов проверяется по-разному в зависимости от содержания.
Например, сумму возврата можно напрямую сопоставить с результатом
инструмента, а для интерпретации политики необходимо учитывать исходный
текст и условия её применимости.

Таблица 9-1. Аспекты оценки траектории агента службы поддержки

\begin{longtable}[]{@{}
  >{\raggedright\arraybackslash}p{(\columnwidth - 4\tabcolsep) * \real{0.3282}}
  >{\raggedright\arraybackslash}p{(\columnwidth - 4\tabcolsep) * \real{0.3552}}
  >{\raggedright\arraybackslash}p{(\columnwidth - 4\tabcolsep) * \real{0.3166}}@{}}
\toprule\noalign{}
\begin{minipage}[b]{\linewidth}\raggedright
\hspace{0pt}Аспект
\end{minipage} & \begin{minipage}[b]{\linewidth}\raggedright
\hspace{0pt}Проверочный вопрос
\end{minipage} & \begin{minipage}[b]{\linewidth}\raggedright
\hspace{0pt}Основные доказательства
\end{minipage} \\
\midrule\noalign{}
\endhead
\bottomrule\noalign{}
\endlastfoot
\hspace{0pt}Результат задачи & \hspace{0pt}Решена ли основная проблема
пользователя & \hspace{0pt}Итоговое состояние среды, результаты
инструментов \\
\hspace{0pt}Соблюдение правил & \hspace{0pt}Соблюдены ли политики, права
доступа и обязательные процедуры & \hspace{0pt}База политик, траектория
действий \\
\hspace{0pt}Границы конфиденциальности & \hspace{0pt}Находится ли
предоставленная информация в допустимых пределах & \hspace{0pt}Текст
ответа, журнал доступа \\
\hspace{0pt}Достоверность фактов & \hspace{0pt}Подтверждаются ли
утверждения знаниями или результатами инструментов & \hspace{0pt}Ссылки
на источники, результаты инструментов \\
\hspace{0pt}Согласованность обещаний и действий &
\hspace{0pt}Действительно ли были выполнены заявленные операции &
\hspace{0pt}Сопоставление ответа с журналом выполнения \\
\hspace{0pt}Качество изложения & \hspace{0pt}Является ли ответ
естественным, лаконичным и ясным & \hspace{0pt}Полный текст диалога,
языковая рубрика \\
\hspace{0pt}Допустимые альтернативы & \hspace{0pt}Найден ли разрешённый
альтернативный путь, если исходное решение оказалось недоступно &
\hspace{0pt}Цель пользователя, политики и последующие действия \\
\end{longtable}

Чтобы обеспечить последующее обновление, результаты оценки должны
содержать четыре элемента: общее состояние задачи, заключения по каждому
аспекту, местоположение доказательств для каждого заключения и тип
ошибки. Местоположение доказательств может указывать на конкретный ход
диалога или вызов инструмента. Аспекты, по которым доказательств
недостаточно, следует помечать как требующие подтверждения и передавать
на дополнительную проверку или ручной анализ.

Например, заключения «возврат не выполнен, ответ сформулирован ясно»
недостаточно, чтобы определить способ улучшения. Если дополнить его
сведениями «на четвёртом ходу интерфейс возврата был вызван без номера
заказа и отклонил операцию; на пятом ходу агент всё равно заявил, что
возврат выполнен», можно предложить отдельные изменения для проверки
параметров и валидации заявлений о завершении операции.

\begin{quote}
\textbf{Эксперимент 9-1 ★★: создание валидатора траекторий для агента
службы поддержки}

Сравните два вида вывода для одного набора траекторий службы поддержки:
только итоговую оценку и отдельные заключения, доказательства и уровень
уверенности по каждому аспекту. Проверьте, позволяет ли второй вариант
отдельно локализовать невыполнение задачи, нарушение правил, ложные
обещания и проблемы изложения, а затем предложить соответствующие
обновления.

Устаревшие знания требуют обновления материалов, при отсутствии
параметров можно добавить проверку перед вызовом, а повторяющиеся
формулировки --- исправить с помощью инструкций или обучения.
\end{quote}

В эксперименте 9-1 на основе вымышленных данных учётных записей были
созданы восемь траекторий: по две для нормального возврата, ложного
обещания, утечки конфиденциальных данных и чрезмерного отказа. В
\texttt{real-normal-refund-01} сотрудник поддержки сообщил, что возврат
оформлен, состояние заказа также изменилось на \texttt{refunded}, а
сумма возврата составила 480.

Каждая из восьми траекторий оценивалась по семи аспектам. Из 56
заключений 52 совпали с заданными для примеров экспертными метками, что
соответствует примерно 92,9 \%. Три из четырёх оставшихся расхождений
относились к одному и тому же случаю переоформления билета. Пользователь
попросил перенести невозвратный авиабилет на 3 октября. Сотрудник
поддержки проверил личность пользователя, трижды запросил данные заказа
и на десятом ходу сообщил, что всё выполнено. Однако состояние заказа
осталось \texttt{confirmed}, а в журнале вызовов не было операции
переоформления.

Кого следует перепроверить в такой ситуации? Модель, оценивавшая
«допустимые альтернативы», сослалась на ответ на десятом ходу и
поставила положительную оценку, не заметив, что допустимое
альтернативное решение можно было выполнить, но оно не было выполнено.
Программа, проверявшая «согласованность обещаний и действий»,
действительно сообщила об ошибке, но объяснила её отсутствием успешного
вызова для возврата средств. Пользователь просил переоформить билет,
поэтому проверять следовало запись о переоформлении и новую дату. При
этом заранее заданная метка для этого аспекта была положительной, что
также расходилось с фактически сгенерированным ответом.

В этом случае нужно исправить две вещи: сотрудник поддержки пропустил
одно действие, а программа оценки проверила не тот инструмент. При
бронировании билета внешне одинаковая ошибка также может возникать в
разных местах. Например, пользователь уже увеличил бюджет, но помощник
продолжает отбирать варианты по прежнему ограничению. Если проследить
траекторию назад, может оказаться, что новое предпочтение не было
сохранено, что в этот раз было извлечено старое предпочтение либо что
модель увидела новый бюджет, но не применила его при выборе билета.

\hypertarget{ux447ux435ux442ux44bux440ux435-ux441ux43fux43eux441ux43eux431ux430-ux43dux435ux43fux440ux435ux440ux44bux432ux43dux43eux439-ux44dux432ux43eux43bux44eux446ux438ux438-ai-ux430ux433ux435ux43dux442ux430}{%
\section{Четыре способа непрерывной эволюции
AI-агента}\label{ux447ux435ux442ux44bux440ux435-ux441ux43fux43eux441ux43eux431ux430-ux43dux435ux43fux440ux435ux440ux44bux432ux43dux43eux439-ux44dux432ux43eux43bux44eux446ux438ux438-ai-ux430ux433ux435ux43dux442ux430}}

После обнаружения проблемы необходимо решить, где именно вносить
изменения. Возьмём в качестве примера агента для анализа медицинских
изображений на \crossreflink{fig:9-3}{рисунке 9-3}. Если в заключении
приведена ссылка на устаревшую версию рекомендаций, сначала следует
проверить, какая версия хранится в базе знаний. Если рекомендации верны,
но в заключении не указан источник, нужно проверить инструкции по
составлению заключения. Даже если оба заключения были возвращены на
доработку, при их исправлении проверять нужно разные места.

\begin{figure}
\hypertarget{fig:9-3}{%
\centering
\includegraphics{images/revised/fig9-3.pdf}
\caption{Рисунок 9-3. Четыре области обновления агента для анализа
медицинских изображений}\label{fig:9-3}
}
\end{figure}

В системе поддержки клиентов также одновременно используются все четыре
способа. Обученная модель формулирует ответы, база знаний хранит
корпоративные политики, Skill объясняет, как обрабатывать исключения, а
серверный код проверяет права доступа и суммы. В таблице 9-2 перечислены
затраты и ограничения, которые необходимо учитывать при выборе способа
обновления.

Таблица 9-2. Границы применимости четырёх способов непрерывной эволюции

\begin{longtable}[]{@{}
  >{\raggedright\arraybackslash}p{(\columnwidth - 6\tabcolsep) * \real{0.1779}}
  >{\raggedright\arraybackslash}p{(\columnwidth - 6\tabcolsep) * \real{0.2432}}
  >{\raggedright\arraybackslash}p{(\columnwidth - 6\tabcolsep) * \real{0.2945}}
  >{\raggedright\arraybackslash}p{(\columnwidth - 6\tabcolsep) * \real{0.2845}}@{}}
\toprule\noalign{}
\begin{minipage}[b]{\linewidth}\raggedright
\hspace{0pt}Способ обновления
\end{minipage} & \begin{minipage}[b]{\linewidth}\raggedright
\hspace{0pt}Для чего подходит
\end{minipage} & \begin{minipage}[b]{\linewidth}\raggedright
\hspace{0pt}Инженерные преимущества
\end{minipage} & \begin{minipage}[b]{\linewidth}\raggedright
\hspace{0pt}Требующие решения проблемы
\end{minipage} \\
\midrule\noalign{}
\endhead
\bottomrule\noalign{}
\endlastfoot
\hspace{0pt}База знаний об опыте & \hspace{0pt}Факты, опыт, исключения и
источники & \hspace{0pt}Быстро обновляется, обеспечивает
прослеживаемость и поиск по требованию & \hspace{0pt}Качество поиска,
условия применимости и способность модели правильно использовать
найденное \\
\hspace{0pt}Промпт и Skill & \hspace{0pt}Принципы принятия решений,
требующие учёта контекста, исключений и приоритетов &
\hspace{0pt}Объяснимость, простота локальных изменений &
\hspace{0pt}Увеличение длины, конфликты инструкций и качество их
соблюдения \\
\hspace{0pt}Программа и Harness & \hspace{0pt}Поддающиеся разбору
процессы, исполняемые проверки, права доступа и строгие ограничения &
\hspace{0pt}Тестируемость, стабильность исполнения &
\hspace{0pt}Разработка и сопровождение, совместимость интерфейсов и
обработка исключений \\
\hspace{0pt}Параметры модели & \hspace{0pt}Восприятие, стиль генерации и
неявные паттерны принятия решений & \hspace{0pt}Повторное использование
усвоенных паттернов, сокращение повторяющегося контекста &
\hspace{0pt}Подготовка данных, стоимость обучения, обобщение и
забывание \\
\end{longtable}

Одно улучшение можно разделить на несколько обновлений. Автор в первую
очередь помещает факты о политиках в базу знаний, а точно исполняемые
ограничения, такие как проверка прав доступа, --- в программу; принципы,
требующие понимания контекста и исключений, выражает через Skill, а
многомерное восприятие и неявные стратегии улучшает посредством
обучения.

\hypertarget{ux43fux440ux435ux43eux431ux440ux430ux437ux43eux432ux430ux43dux438ux435-ux43eux43fux44bux442ux430-ux432-ux437ux43dux430ux43dux438ux44f}{%
\subsection{Преобразование опыта в
знания}\label{ux43fux440ux435ux43eux431ux440ux430ux437ux43eux432ux430ux43dux438ux435-ux43eux43fux44bux442ux430-ux432-ux437ux43dux430ux43dux438ux44f}}

Об одной и той же покупке билета можно сохранить информацию двух
уровней. Первая запись: «Эта авиакомпания требует заказывать специальное
питание за 24 часа». Вторая: «Перед бронированием проверить срок заказа
специального питания, чтобы не обнаружить после оплаты невозможность
выполнить требование». Первую следует сохранять вместе с названием
авиакомпании, применимыми маршрутами и источником политики. У другой
авиакомпании срок может отличаться, но правило «сначала проверить срок
для этого рейса» по-прежнему переносимо. Тогда при устаревании факта
потребуется заменить только конкретный срок, а опыт проверки останется
полезен при следующем бронировании. Оба вида информации можно сохранять
с помощью рассмотренных в третьей главе методов хранения, индексирования
и поиска.

Извлечение такого опыта из одного бронирования может стоить недорого, но
нельзя просто копировать всю запись о покупке билета. В ней содержится и
исходный текст политики, возвращённый инструментом, и обходные пути,
возникшие в ходе нескольких попыток AI-агента. Если один запрос
завершился по тайм-ауту, а результат удалось получить лишь после смены
инструмента, слепое воспроизведение этой последовательности может
превратить случайный сбой в основание для смены инструмента.

Материалы можно сохранять на трёх уровнях: \textbf{исходная траектория}
содержит полные доказательства выполнения, \textbf{анализ отдельного
запуска} фиксирует причины успеха или неудачи и черновик опыта, а
\textbf{итоговый документ об опыте} обобщает подходы, проверенные на
нескольких задачах. Исходные траектории сохраняются с указанием версий,
а анализ и знания пересматриваются по мере появления новых
доказательств. В итоговом документе указываются сценарии применения,
рекомендуемая стратегия, запрещённые действия, исключения, источники и
дата последней проверки, чтобы облегчить последующий поиск и контроль.

В подходе User as Code из третьей главы сначала пополняется журнал
фактов, а затем периодически перестраивается модель пользователя. Здесь
запись также отделяется от систематизации. Вернёмся к предыдущему
запросу. Допустим, другой запрос тоже завершился по тайм-ауту, но
повторная попытка с тем же инструментом оказалась успешной. Остаётся ли
обоснованной рекомендация «при тайм-ауте сменить инструмент»? На
\crossreflink{fig:9-4}{рисунке 9-4} показано, как на основе этих двух
записей сформулировать более точную рекомендацию, а затем проверить её
на других запросах.

\begin{figure}
\hypertarget{fig:9-4}{%
\centering
\includegraphics{images/revised/fig9-4.pdf}
\caption{Рисунок 9-4. От оценённых траекторий к документу знаний об
опыте}\label{fig:9-4}
}
\end{figure}

При использовании рассмотренных в третьей главе методов извлечения,
кластеризации и поиска можно сопоставлять положительные и отрицательные
примеры одного и того же опыта. В двух случаях бронирования
использовалась одна версия инструмента? Один запуск завершился успешно,
а другой --- неудачно потому, что был пропущен шаг, или потому, что
изначально различались условия бронирования? Преобразование исходных
записей в полезный опыт можно разделить на пять этапов:

\begin{enumerate}
\def\labelenumi{\arabic{enumi}.}
\tightlist
\item
  Сохранить исходную траекторию и результат среды, зафиксировав задачу и
  версии инструментов.
\item
  Сформировать анализ каждого запуска, записав тип задачи, необходимые
  способности, фактическую стратегию, ошибки и исключения.
\item
  Агрегировать запуски по семействам задач и для каждого элемента опыта
  перечислить подтверждающие и опровергающие его траектории.
\item
  Отобрать опыт с достаточной доказательной базой и оформить его как
  черновик записи знаний.
\item
  Протестировать его на новых задачах, не участвовавших в дистилляции, и
  проверить, улучшает ли рекомендуемая стратегия последующие результаты.
\end{enumerate}

GAIA\footnote{Mialon, G., et al.~\emph{GAIA: a benchmark for General AI
  Assistants.} arXiv:2311.12983, 2023.} предоставляет многошаговые
задачи, требующие поиска, чтения веб-страниц, обработки файлов и
вычислений, а AWorld\footnote{Yu, C., et al.~\emph{AWorld: Orchestrating
  the Training Recipe for Agentic AI.} arXiv:2508.20404, 2025.} ---
среду для выполнения действий агентом, вызова инструментов и записи
траекторий. Их сочетание позволяет сначала пометить успех, частичный
успех и неудачу по ответу или состоянию среды, а затем сравнить пути
внутри каждого семейства задач. В успешных записях можно найти подходы,
которые стоит повторить, а в неудачных --- условия, при которых этого
делать не следует. Если задача выполнена лишь частично, можно поэтапно
проверить, какие действия оказались эффективными и с какого места
начались проблемы.

Reflexion\footnote{Shinn, N., et al.~\emph{Reflexion: Language Agents
  with Verbal Reinforcement Learning.} arXiv:2303.11366, 2023.}
использует рефлексию на естественном языке для улучшения последующих
попыток. При включении этого подхода в дистилляцию знаний рефлексия
может применяться для выдвижения гипотез о причинах и создания
черновиков опыта, результаты среды --- для проверки, сопоставление
траекторий --- для определения условий применимости, а новые задачи ---
для проверки переноса. Решение о включении опыта в официальную базу
знаний принимается на основе всех четырёх видов материалов.

Вернёмся к сценарию покупки билета из шестой главы. Пользователь хотел
уложиться в 1200 юаней, чтобы в тот же день успеть на совещание, но в
итоге принял предложение за 1250 юаней. Ассистент может записать
исходный бюджет, цену покупки, цель поездки и дословные слова
пользователя. При следующих трёх бронированиях минимальная цена каждый
раз случайно составляла ровно 1250 юаней. Как следует использовать эту
запись в памяти?

\begin{itemize}
\tightlist
\item
  На следующей неделе пользователю снова нужно успеть на такое же
  деловое совещание. В нескольких предыдущих случаях он соглашался
  доплатить несколько десятков юаней, чтобы попасть на встречу, а теперь
  поручил ассистенту забронировать билет как обычно. Ассистент может
  купить этот билет и сообщить фактическую цену.
\item
  Пользователь едет в отпуск с ребёнком и отдельно просит: «В семейной
  поездке нужно немного сэкономить. Если дороже 1200 юаней, сначала
  спроси меня». Ассистент может отправить пользователю вариант за 1250
  юаней и одновременно проверить, есть ли в пределах бюджета рейсы на
  предыдущий или следующий день.
\item
  Позднее пользователь объясняет: «В прошлый раз я согласился только
  потому, что спешил, но вообще-то 1250 юаней для меня тоже дорого. В
  дальнейшем и в командировках сначала спрашивай меня, если цена выше
  бюджета». В следующий раз ассистент должен сначала спросить, даже если
  тот же человек снова летит на совещание в тот же город. Следует
  сохранить запись о той покупке, а это более новое и конкретное
  требование добавить в предпочтения для командировок: тогда обе записи
  в памяти останутся полезными.
\end{itemize}

Эти три заказа показывают уровни обновления. Одна пользовательская
поправка может немедленно изменить текущую задачу, а также служит
достаточным основанием для обновления личных предпочтений пользователя.
Предлагать изменение общего Skill бронирования, чтобы при сравнении цен
он извлекал соответствующие предпочтения, стоит лишь тогда, когда
множество независимых случаев показывает, что ассистент игнорирует цель
поездки. И только если явно заданные бюджет и предпочтения уже
загружены, но модель по-прежнему ошибается в оценке при иной
формулировке, стоит собирать такие примеры для обучения способности
модели принимать решения.

\hypertarget{ux43fux440ux435ux432ux440ux430ux449ux435ux43dux438ux435-ux43eux43fux44bux442ux430-ux432-ux438ux43dux441ux442ux440ux443ux43aux446ux438ux438}{%
\subsection{Превращение опыта в
инструкции}\label{ux43fux440ux435ux432ux440ux430ux449ux435ux43dux438ux435-ux43eux43fux44bux442ux430-ux432-ux438ux43dux441ux442ux440ux443ux43aux446ux438ux438}}

Если помощнику каждый раз передают правильные материалы, но он
по-прежнему ошибается на одном и том же этапе, а опытный человек может
одной фразой описать, как следует принимать решение, можно попробовать
сформулировать инструкцию. Документ с накопленным опытом служит
помощнику справочным материалом, системный промпт задаёт требования к
поведению в разных задачах, Skill загружает локальный процесс для
конкретной области или инструмента, а соблюдение жёстких ограничений,
например прав доступа, обеспечивает Harness. От правильного выбора места
для изменения зависит, на какие будущие задачи оно повлияет.

Andrej Karpathy использует термин \textbf{обучение системного промпта
(System Prompt Learning)} для описания способа обновления опыта:
столкнувшись с проблемой, модель записывает напоминание для
использования в будущих задачах.\footnote{Karpathy, A. ``We're missing
  (at least one) major paradigm for LLM learning \ldots{} system prompt
  learning?'' X, May 11, 2025.} DSPy\footnote{Khattab, O., et
  al.~\emph{DSPy: Compiling Declarative Language Model Calls into
  Self-Improving Pipelines.} arXiv:2310.03714, 2023.} оптимизирует
инструкции и примеры на наборе для разработки, OPRO\footnote{Yang, C.,
  et al.~\emph{Large Language Models as Optimizers.} arXiv:2309.03409,
  2023.} предлагает новые промпты на основе уже опробованных промптов и
их оценок, а GEPA\footnote{Agrawal, L., et al.~\emph{GEPA: Reflective
  Prompt Evolution Can Outperform Reinforcement Learning.}
  arXiv:2507.19457, 2025.} использует обратную связь по траекториям,
чтобы предлагать и отбирать изменения промпта. Все эти методы
переписывают инструкции на основе результатов выполнения, а затем
сравнивают новую и старую версии на задачах для разработки, поэтому они
подходят для пакетной офлайн-оптимизации.

При обновлениях в продакшене автор отдаёт предпочтение минимальному diff
с указанием источника, сохраняя вызвавший изменение пример, причину
изменения и старую версию. Локальный патч меньше затрагивает
существующие правила и позволяет при регрессии откатить отдельное
изменение.

Как убедиться, что новая инструкция действительно исправила предыдущую
ошибку? Предположим, пользователь указывает на ошибочно архивированное
письмо клиента и говорит: «Это клиент отвечает мне --- почему письмо
тоже посчитали рассылкой?» В открытом почтовом помощнике Inbox Zero
предусмотрен инструмент пробного запуска. После изменения правил
помощник может получить исходное письмо по его ID, повторно
классифицировать его согласно правилам, которые сейчас включены для этой
учётной записи, и показать, какие правила сработают, почему, какую метку
планируется назначить или в какую папку переместить письмо. На этом
этапе просматриваются только предполагаемые действия: почтовый ящик
остаётся без изменений, а черновик ответа не создаётся.\footnote{В
  проекте Inbox Zero инструмент для пробного запуска правил на указанном
  письме реализован в \texttt{test-rules-on-email-tool.ts}.} Затем
следует так же проверить настоящую рассылку: письмо клиента должно
остаться во входящих, а рассылка --- по-прежнему отправиться в архив.
Вместе эти два письма показывают, сохранило ли изменение прежнюю границу
классификации.

\hypertarget{ux43fux440ux438ux43cux435ux440-1.-ux447ux451ux442ux43aux438ux435-ux433ux440ux430ux43dux438ux446ux44b-ux43fux435ux440ux435ux434ux430ux447ux438-ux447ux435ux43bux43eux432ux435ux43aux443-ux438-ux438ux441ux43fux43eux43bux44cux437ux43eux432ux430ux43dux438ux44f-ux438ux43dux441ux442ux440ux443ux43cux435ux43dux442ux43eux432}{%
\subsubsection{Пример 1. Чёткие границы передачи человеку и
использования
инструментов}\label{ux43fux440ux438ux43cux435ux440-1.-ux447ux451ux442ux43aux438ux435-ux433ux440ux430ux43dux438ux446ux44b-ux43fux435ux440ux435ux434ux430ux447ux438-ux447ux435ux43bux43eux432ux435ux43aux443-ux438-ux438ux441ux43fux43eux43bux44cux437ux43eux432ux430ux43dux438ux44f-ux438ux43dux441ux442ux440ux443ux43cux435ux43dux442ux43eux432}}

Политика telecom в τ²-bench требует передавать запрос человеку, если он
выходит за пределы действий агента, а перед передачей приложить все
усилия для решения задачи. Из 20 задач, использованных для извлечения
правил, менее мощная модель прошла только одну. Как завершились неудачей
остальные 19? В 17 задачах вызывался инструмент передачи человеку, а ещё
две выполнялись до достижения предельного числа шагов.

Эти 19 неудачных траекторий передали модели, попросив её обобщить
правила и добавить их в конец исходной политики. При работе с теми же
114 новыми задачами исходная политика позволила пройти 14 задач
(12,3\%); с первой версией правил модель прошла 20 задач, а со второй
--- 22 (19,3\%). Сопоставление по каждой задаче показало, что вторая
версия сохранила успех во всех прежних 14 задачах и дополнительно решила
ещё восемь. Чем именно две версии правил помогли службе поддержки?

\textbf{Содержимое правил определяется входными материалами.} Когда
модели предоставили только тексты ошибок и сводки неудач, она вывела
правило: «Прекратить вызовы, если один и тот же инструмент неоднократно
возвращает ошибку». При втором извлечении модель также видела отдельные
списки инструментов службы поддержки и пользователя, а промпт предлагал
ей подумать, «кто должен выполнить этот шаг». На этот раз
сгенерированные правила явно разделили обязанности: операции на
устройстве, связанные с состоянием сети, SIM-картой и APN, должен
выполнять пользователь под руководством службы поддержки. Так опыт
превратился из общего напоминания о повторных попытках в конкретное
правило распределения операций.

\textbf{Обобщение также необходимо проверять по политике и результатам.}
В первой версии добавленных правил были два пункта: «Передать человеку
после трёх последовательных неудачных вызовов» и «Передать человеку,
если пользователь дважды не предоставил номер». При выполнении по этой
версии правил передача человеку произошла во всех 114 задачах набора
переноса, но пройти удалось лишь 20. Эти два правила усвоили и
преждевременную передачу человеку.

Одна траектория наглядно показывает разницу до и после исправления.
Пользователь сообщил, что на телефоне нет сигнала. В базовой версии
служба поддержки подставила в поле номера фразу «Пожалуйста,
предоставьте ваш правильный номер телефона», пять раз получила ошибку
при поиске, а затем передала запрос человеку. Со второй версией правил
первый запрос снова был отправлен без номера, но теперь служба поддержки
попросила пользователя сообщить номер в диалоге и после его получения
нашла учётную запись. Затем она попыталась проверить SIM-карту;
инструмент сообщил, что такая операция отсутствует, и служба поддержки
попросила пользователя вынуть и снова вставить SIM-карту. Пользователь
выполнил просьбу, а затем отключил на телефоне режим полёта --- только
после этого задача была пройдена. Рассмотрим и остальные семь задач,
которые вторая версия решила в отличие от базовой: лишь в одной из них
не было передачи человеку, а в остальных шести инструмент передачи
человеку вызывался. Следовательно, если считать критерием успеха
«отсутствие передачи человеку», эти шесть задач как раз будут пропущены.
Чтобы решить, нужна ли передача, необходимо учитывать, что ещё может
сделать служба поддержки и до какого этапа требуется довести задачу.

\begin{quote}
\textbf{Эксперимент 9-2 ★★. Извлечение правил передачи человеку и
использования инструментов из неудачных траекторий τ²-bench}

Используйте среду telecom из \crossreflink{chap:7}{главы 7}, разделив
набор для извлечения и набор переноса. Сначала запустите менее мощную
модель и соберите неудачные траектории, затем поручите модели
сгенерировать правила и добавьте их к исходной политике. На наборе
переноса сравните исходную и две эволюционировавшие версии; во всех трёх
группах заменяйте только файл политики, сохраняя один и тот же симулятор
пользователя.

В базовой версии передача человеку происходила в 104 задачах, а во
второй версии --- в 111; задача учитывается, если передача произошла в
ней хотя бы один раз. Ошибочные вызовы инструментов пользователя,
напротив, подсчитываются по числу вызовов: их количество снизилось с
1049 до 163. Ещё одна метрика подсчитывает вызовы, в параметрах которых
хотя бы одно значение было пустой строкой: их число снизилось с 904 до
174; вызовы, в которых параметры вообще не передавались, этой метрикой
не учитываются. Каждая стратегия запускалась на каждой задаче только
один раз.
\end{quote}

В службе поддержки авиакомпании также встречается преждевременная
передача человеку. Если пользователь возражает против сбора за возврат
билета, сбора за изменение бронирования или правил провоза багажа, агент
должен найти применимую политику, объяснить её и поискать допустимые
альтернативы. Согласно правилам следующего эксперимента, обычные споры о
политике агент продолжает обрабатывать самостоятельно, а при явном
требовании пользователя связать его с человеком или при возникновении
определённой ситуации, связанной с безопасностью, вызывает
\texttt{transfer\_to\_human}.

\begin{quote}
\textbf{Эксперимент 9-3 ★★. Оптимизация системного промпта службы
поддержки авиакомпании на основе неудачных траекторий}

Проанализируйте траектории преждевременной передачи человеку по трём
направлениям: соблюдение правил, решение задачи и поиск допустимых
обходных решений. На основе анализа сгенерируйте патч промпта с
указанием источников. Сравните автоматически извлечённую, исходную и
вручную оптимизированную версии в одинаковых условиях.

Ещё пять задач предназначены для проверки сохранения прежних
возможностей: в трёх служба поддержки должна действовать самостоятельно,
а в двух --- передать запрос человеку. Все три версии инструкций
получили 5/5. На пяти задачах с преждевременной передачей,
использованных для извлечения патча, одна и та же модель-рецензент
признала пройденными соответственно 0, 2 и 4 задачи для исходной,
автоматической и вручную оптимизированной версий. Автоматическая версия
исправила две задачи: возврат билета по специальному тарифу и
сверхнормативный багаж. На этом основании программа разрешила
ограниченное пробное применение.

Одна оценка вызывает сомнения. Пассажир попросил отменить сбор за
изменение бронирования в размере 200 юаней, но вручную оптимизированная
версия начала объяснять, можно ли вернуть билет по специальному тарифу.
Рецензент одобрил предложение «вернуть налоги» и засчитал задачу, хотя
служба поддержки так и не объяснила сбор за изменение бронирования, о
котором спрашивал пассажир.
\end{quote}

\hypertarget{ux43fux440ux438ux43cux435ux440-2.-ux443ux43bux443ux447ux448ux435ux43dux438ux435-skill-ux443ux442ux43eux447ux43dux435ux43dux438ux44f-ux442ux440ux435ux431ux43eux432ux430ux43dux438ux439-ux43dux430-ux43eux441ux43dux43eux432ux435-ux43fux43eux43bux44cux437ux43eux432ux430ux442ux435ux43bux44cux441ux43aux43eux439-ux43eux431ux440ux430ux442ux43dux43eux439-ux441ux432ux44fux437ux438}{%
\subsubsection{Пример 2. Улучшение Skill уточнения требований на основе
пользовательской обратной
связи}\label{ux43fux440ux438ux43cux435ux440-2.-ux443ux43bux443ux447ux448ux435ux43dux438ux435-skill-ux443ux442ux43eux447ux43dux435ux43dux438ux44f-ux442ux440ux435ux431ux43eux432ux430ux43dux438ux439-ux43dux430-ux43eux441ux43dux43eux432ux435-ux43fux43eux43bux44cux437ux43eux432ux430ux442ux435ux43bux44cux441ux43aux43eux439-ux43eux431ux440ux430ux442ux43dux43eux439-ux441ux432ux44fux437ux438}}

Существующий Skill уточнения требований может определять, когда задавать
вопросы, какие именно вопросы задавать и когда начинать исполнение. По
мере накопления пользовательской обратной связи система должна
корректировать эти условия, чтобы польза от уточнений соответствовала
цене отвлечения пользователя.

Например, если пользователь говорит: «Добавь на страницу входа поддержку
корпоративной аутентификации», перед реализацией может потребоваться
уточнить поставщика удостоверений, способ отката, совместимость с
существующими пользователями и охват развёртывания. Напротив,
исправление одной опечатки обычно можно выполнить сразу. Если вопросов
слишком мало, ключевые решения придётся пересматривать после сдачи
результата; если их слишком много, даже простое изменение превратится в
долгое интервью. Skill должен выбирать процесс с учётом неоднозначности,
риска и стоимости переделки.

Начальный процесс можно организовать следующим образом: небольшие
низкорисковые и легко обратимые изменения выполнять после указания
необходимых предположений; если изменение затрагивает архитектуру,
данные, права доступа, публичный интерфейс или имеет большой охват,
задать вместе несколько вопросов, ответы на которые способны изменить
решение. Затем составить краткую Spec или Plan с целями, нецелями,
компромиссами, предположениями и критериями приёмки. Получив ответы по
вопросам, требующим подтверждения, продолжить исполнение; если во время
исполнения изменятся исходные условия, определить затронутые решения и
запросить дополнительное подтверждение.

Для улучшения этого процесса необходимо фиксировать задачу, уточняющие
вопросы, версию Spec, пользовательские изменения, итоговый результат и
переделки. Если пользователь говорит: «Это не то, что я представлял»,
следует проверить, какие условия были упущены; если он говорит: «Ты
задаёшь слишком много вопросов», нужно выяснить, влияли ли ответы на эти
вопросы на выбранное решение.

Если в нескольких задачах аутентификации регулярно упускаются требования
совместимости со старым способом входа, в Skill можно добавить условие:
«Уточнить поставщика удостоверений, путь отката и границы
совместимости». Если множество исправлений опечаток задерживается из-за
вопросов, область срабатывания процесса уточнения, напротив, следует
сузить.

Здесь также проявляется разделение обязанностей между Skill и Harness.
Skill помогает модели понять ситуацию, задать вопросы и
систематизировать решение; Harness в соответствии с правилами проекта
блокирует операции, если для высокорисковой операции записи не получено
необходимое разрешение, предпринимается прямая попытка изменить
защищённую ветвь или пропускается процесс выпуска. Решение о содержании
требований по-прежнему принимает модель с учётом пользовательского
ввода. Проверенные траектории диалогов также могут стать обучающими
данными, описанными в \crossreflink{chap:8}{главе 8}.

\begin{quote}
\textbf{Эксперимент 9-4 ★★. Эволюция Skill уточнения требований и
подтверждения Spec на основе пользовательской обратной связи}

В этом упражнении необходимо самостоятельно разработать сравнение.
Сначала можно подготовить две задачи: исправить опечатку на кнопке входа
и добавить на ту же страницу корпоративную аутентификацию. Критерии
приёмки для первой задачи можно записать сразу; для второй, возможно,
сначала придётся уточнить поставщика удостоверений. Затем добавьте
задачи, затрагивающие архитектуру, права доступа, данные или публичные
интерфейсы, и испытайте три процесса: прямое исполнение, исполнение
после уточнения и исполнение после подтверждения Spec. Фиксируйте ответы
пользователя, изменения Spec, результаты работы и обратную связь о
переделках, а затем на их основе сгенерируйте локальное обновление
Skill.
\end{quote}

\hypertarget{ux43fux440ux435ux432ux440ux430ux449ux435ux43dux438ux435-ux43eux43fux44bux442ux430-ux432-ux43fux440ux43eux433ux440ux430ux43cux43cux443}{%
\subsection{Превращение опыта в
программу}\label{ux43fux440ux435ux432ux440ux430ux449ux435ux43dux438ux435-ux43eux43fux44bux442ux430-ux432-ux43fux440ux43eux433ux440ux430ux43cux43cux443}}

Стабильные, повторяемые и проверяемые операции можно на основе опыта
выполнения оформлять в виде рабочих процессов, инструментов или кода
Harness. Автор рекомендует в первую очередь превращать в программы
именно такие операции: после того как модель найдёт работоспособный
путь, последующие задачи смогут использовать его повторно, избавляя от
необходимости снова и снова читать инструкции, разбираться в страницах и
выбирать следующий шаг. В \crossreflink{chap:5}{главе 5} были
представлены возможности чтения и записи файлов, запуска тестов и
создания инструментов; в этом разделе они применяются в процессе
обновления «диагностика сбоя → изменение кода → регрессионная проверка».

Изменения могут вноситься на разных уровнях. На операционном уровне
создаются параметризованные рабочие процессы или адаптеры API; на уровне
управления настраиваются маршрутизация инструментов, повторные попытки,
автоматическое отключение при сбоях и сжатие контекста; на уровне
проверки добавляются проверки параметров, утверждения о состоянии и
регрессионные сценарии; на архитектурном уровне изменяются потоки
информации между планированием, выполнением и Reviewer.

Рассмотрим отправку электронной почты. При первом выполнении
мультимодальный AI-агент посредством наблюдений и действий находит
элементы управления «Написать», «Получатель», «Тема», «Текст» и
«Отправить». В следующей задаче обычно меняются лишь получатель и
содержимое, поэтому последовательность операций можно повторно
использовать подобно макросу электронной таблицы. Система превращает
исследовательскую траекторию в программу и сохраняет параметры, проверки
состояния и применимые версии.

Жизненный цикл браузерного рабочего процесса можно разделить на шесть
этапов:

\begin{enumerate}
\def\labelenumi{\arabic{enumi}.}
\tightlist
\item
  \textbf{Захват траектории.} Записываются переходы, щелчки, ввод и
  выбор значений из раскрывающихся списков; сохраняются параметры
  действий, URL страницы и сведения для поиска элементов, например XPath
  и CSS. Эти данные позволяют впоследствии снова находить элементы.
\item
  \textbf{Извлечение параметров.} Конкретные значения получателя, темы и
  текста заменяются переменными \texttt{\{recipient\}},
  \texttt{\{subject\}}, \texttt{\{content\}} и т. п., а стабильные
  действия сохраняются. Например, впервые введённый адрес
  \texttt{test@example.com} становится одним из значений параметра
  получателя.
\item
  \textbf{Добавление проверок состояния.} Перед выполнением проверяется
  видимость кнопки и заполненность обязательных полей, после выполнения
  --- изменение страницы; в конце задачи считывается фактическое
  состояние страницы или бэкенда и подтверждается, что письмо
  отправлено.
\item
  \textbf{Независимое воспроизведение.} Учётная запись в песочнице или
  тестовый сайт возвращается в новое начальное состояние, после чего
  программа выполняется полностью с проверкой условий каждого этапа и
  итогового результата.
\item
  \textbf{Сопоставление с новой задачей.} Рабочий процесс находится в
  библиотеке возможностей по намерению и ключевым словам, из текущей
  задачи извлекаются параметры, после чего выполнение передаётся
  Playwright. Программа отвечает за ожидание элементов и проверку
  состояния; при штатном воспроизведении можно обойтись без пошаговых
  вызовов LLM.
\item
  \textbf{Обработка утраты работоспособности.} Если перестал работать
  поиск элемента, не прошла проверка состояния, изменилась schema API
  или итоговый результат оказался аномальным, текущее воспроизведение
  останавливается. Сначала проверяется, какие операции уже выполнены,
  затем полный процесс агента заново исследует путь из текущего
  состояния и записывает новые успешно пройденные шаги в скрипт. Если
  письмо уже отправлено, но результат ещё не появился на странице,
  сначала следует запросить состояние отправки, чтобы повторное
  воспроизведение с начала не привело к повторной отправке.
\end{enumerate}

Экономия времени зависит от того, сколько времени первоначально уходило
на принятие решений моделью. В экспериментах PreAct с повторяющимися
задачами сообщалось о сквозном ускорении в 8,5--13 раз; в приведённом
ниже локальном примере с электронной почтой исключены все вызовы модели,
однако браузеру по-прежнему приходится ждать ответа страницы.\footnote{Li,
  Bojie. \emph{PreAct: Computer-Using Agents that Get Faster on Repeated
  Tasks.} arXiv:2606.17929, 2026.}

\begin{quote}
\textbf{Эксперимент 9-5 ★★★: создание проверяемого рабочего процесса из
браузерной траектории}

Создайте из успешной траектории рабочий процесс с параметрами и
проверками состояния, а затем сравните его долю завершённых задач, число
вызовов модели и затраченное время с базовым вариантом, который каждый
раз заново исследует путь. После этого измените структуру страницы или
создайте состояние, при котором щелчок выполнен успешно, но задача не
отправлена, и проверьте, способна ли программа обнаружить утрату
работоспособности и перейти к резервному сценарию.

В этом эксперименте с локальным почтовым сайтом первоначальное
исследование вызвало модель 4 раза, а параметризованное воспроизведение
--- 0 раз; выполнение заняло около 5,31 и 4,45 секунды соответственно.
При воспроизведении также возникли два типа сбоев: в одном случае сервер
отказался отправлять письмо с пустым текстом; в другом запрос был
принят, но сообщение не сохранилось. Версия, проверявшая лишь факт
выполнения щелчка, в обоих случаях ошибочно сообщила об успехе, тогда
как версия с проверкой состояния после отправки оба раза распознала
сбой.
\end{quote}

\textbf{Сначала изменяйте код в отдельной ветви.} Создайте из текущей
стабильной версии изолированную ветвь, поручите Coding Agent подготовить
локальный патч и последовательно выполните связанные с изменением
статические проверки, модульные тесты, проверки безопасности,
воспроизведение неудачной траектории и регрессионные проверки старых
задач, а затем переходите к поэтапному выпуску. Работающая система
продолжает использовать стабильную версию; переключение выполняется лишь
после прохождения патчем всех проверок.

Этот процесс можно организовать с помощью Git worktree и Pull Request.
Skill предписывает модели создать отдельную рабочую область, уточнить
требования, завершить реализацию и тестирование, а в PR описать
проблему, изменения, проверенные сценарии и оставшиеся риски, которые
должен проверить принимающий работу специалист. Harness, в свою очередь,
обеспечивает соблюдение проектных ограничений: например, защищает
\texttt{main}, проверяет обязательные статусы ревью и перед передачей
результата подтверждает полноту артефактов выпуска.

Чтобы определить, как конкретное изменение влияет на результат, к
каждому изменению следует прилагать проверяемое описание: доказательства
сбоя, предполагаемая первопричина, затронутые компоненты, предлагаемый
патч, исправляемая ошибка, изначально правильное поведение и сценарии,
отдельно проверяющие ошибку и правильное поведение. Agentic Harness
Engineering распределяет соответствующую информацию по трём уровням
наблюдаемости: компоненты, опыт и решения. Компоненты хранятся в виде
доступных для редактирования файлов, траектории преобразуются в
доказательства, которые можно просматривать уровень за уровнем, перед
редактированием объявляется ожидаемое влияние, а затем оно проверяется
по последующим результатам.\footnote{Lin, Jiahang, et al.~\emph{Agentic
  Harness Engineering: Observability-Driven Automatic Evolution of
  Coding-Agent Harnesses.} arXiv:2604.25850, 2026.}

Когда автор просит модель подготовить план изменений, он одновременно
предоставляет успешные траектории и ранее отклонённые варианты.
Self-Harness использует эту информацию, чтобы ограничить направление
поиска: успешные записи показывают модели, какие подходы изначально были
правильными, а история попыток --- какие изменения уже оказались
неудачными.\footnote{Zhang, Hangfan, et al.~\emph{Self-Harness:
  Harnesses That Improve Themselves.} arXiv:2606.09498, 2026.} Если
предоставить только примеры сбоев, модель может расширить область
изменений и нарушить изначально правильное поведение; если же не
сохранить отклонённые варианты, она может снова и снова предлагать одну
и ту же неработоспособную идею.

Новые инструменты также могут появляться в результате такого обновления,
обусловленного потребностями. В одной из задач Alita\footnote{Qiu, J.,
  et al.~\emph{Alita: Generalist Agent Enabling Scalable Agentic
  Reasoning with Minimal Predefinition and Maximal Self-Evolution.}
  arXiv:2505.20286, 2025.} требовалось найти число, упомянутое сразу
после первого появления динозавра в панорамном VR-видео на YouTube с
закадровым голосом актёра, озвучившего Голлума во «Властелине колец».
Агент обнаружил, что ему не хватает возможности читать субтитры, нашёл и
протестировал \texttt{youtube-transcript-api}, создал на его основе
инструмент для работы с субтитрами и в итоге получил \texttt{100000000}.

\begin{quote}
\textbf{Эксперимент 9-6 ★★★: самомодификация агента, инициированная
неудачной траекторией}

Для проблемы, при которой инструмент продолжает вызываться после
возврата ошибки, не допускающей повторной попытки, сравните два
исправления: напоминание в промпте и программную стратегию повторных
попыток. Изменённая программа должна в зависимости от типа ошибки
повторять попытку, останавливаться или включать автоматическое
отключение при сбоях; затем воспроизведите неудачную траекторию, чтобы
проверить целевую проблему, и выполните задачи с временными сбоями,
чтобы проверить способность к восстановлению. Если тип ошибки и условия
повторной попытки уже можно точно определить, автор предпочитает
закреплять такое ограничение в программе: тогда каждый вызов следует
одному правилу, а модели не приходится снова и снова интерпретировать
промпт.

Агент, отвечающий за обновление, отправляет патч в изолированной версии,
а внешние процессы тестирования и выпуска решают, следует ли его
принимать. Валидаторы, приёмочные сценарии и стабильная версия защищены
независимыми полномочиями; журнал отката сохраняет сведения о поведении
до и после изменения.
\end{quote}

Можно совершенствовать и саму логику проверки. Например, если
неоднократные исправления пользователей и ревью показывают, что для
высокорисковых вызовов не запрашивается подтверждение, можно добавить
проверку перед вызовом. Здесь необходимо различать компонент
бизнес-проверки и независимый механизм, принимающий решение о выпуске
обновления: первый можно изменять вместе с бизнес-кодом, а второй
управляется независимыми полномочиями на тестирование, ревью и выпуск.
Новая проверка должна как охватывать целевой риск, так и позволять
продолжать выполнение штатных задач, на которые уже получено разрешение.

\begin{quote}
\textbf{Эксперимент 9-7 ★★: шлюз подтверждения высокорисковых операций,
созданный на основе обратной связи пользователей}

Выделите из исправлений пользователей и последующих ревью случаи
отсутствия подтверждения и сформулируйте правило, которое перед вызовом
инструмента проверяет наличие разрешения. Первая группа тестов должна
охватывать ещё не подтверждённые высокорисковые операции, вторая --- уже
разрешённые или не требующие подтверждения штатные операции.

На первой группе измерьте, сколько операций пропустило новое правило, а
на второй --- сколько операций оно заблокировало ошибочно. Возможность
выпуска изменения зависит от результатов независимых тестов и правил
выпуска. Агент, создающий правило проверки, может изменять только
согласованный компонент; приёмочные сценарии, правила утверждения и
резервная копия стабильной версии защищаются внешним механизмом.
\end{quote}

\hypertarget{ux43fux440ux438ux43cux435ux440-ux43cux43eux434ux443ux43bux44cux43dux44bux435-ux43eux431ux43dux43eux432ux43bux435ux43dux438ux44f-deepseek-harness}{%
\subsubsection{Пример: модульные обновления DeepSeek
Harness}\label{ux43fux440ux438ux43cux435ux440-ux43cux43eux434ux443ux43bux44cux43dux44bux435-ux43eux431ux43dux43eux432ux43bux435ux43dux438ux44f-deepseek-harness}}

DeepSeek Harness (\texttt{dsh}) использует архитектуру, в которой «всё
является плагином», а лежащая в её основе Cordis предоставляет механизм
композиции.\footnote{DeepSeek AI, \emph{DeepSeek Harness}, 2026. Состав
  плагинов описан в обзоре архитектуры; разделение обязанностей между
  динамической средой выполнения и постоянными плагинами --- в описании
  расширений. Проект был выпущен в августе 2026 года; в этом разделе
  рассматривается его модульная архитектура на этапе предварительной
  версии для разработчиков.} Адаптеры моделей, реестр инструментов,
журналы сеансов и основной цикл агента можно настраивать и заменять как
плагины. В связи с непрерывной эволюцией возникает конкретная
программная проблема: как при изменении компонентов во время работы
очищать старые ресурсы и обеспечивать дальнейшую корректную работу
зависящих от них компонентов?

Cordis разделяет эту проблему на два аспекта. \textbf{Временная
компонуемость (temporal composability)} связана с восстановлением после
удаления компонента: для выделения ресурсов, регистрации событий и
изменения контекста должны существовать соответствующие процедуры
отмены. \textbf{Пространственная компонуемость (spatial composability)}
связана с зависимостями: компоненты объявляют необходимые им сервисы, а
среда выполнения координирует их активацию и деактивацию при появлении,
исчезновении или изменении зависимостей.\footnote{Shi, Yifan, Wei Zhang,
  and Tianyi Cui. \emph{A Programming Paradigm for Spatiotemporal
  Composability}, 2026, arXiv:2608.25512.}

Например, плагин инструмента регистрирует обработчик событий и зависит
от определённого сервиса модели. При замене плагина необходимо удалить
старый обработчик, чтобы одно событие не обрабатывалось дважды; если
сервис модели временно недоступен, зависимые компоненты должны
приостановить связанную с ним работу, а после восстановления сервиса ---
снова активироваться. Первая проблема относится к жизненному циклу
ресурсов и побочных эффектов, вторая --- к жизненному циклу
зависимостей.

Без единого механизма каждому компоненту приходится самостоятельно
обрабатывать загрузку, выгрузку и изменения зависимостей, а локальная
замена может нарушить работу вызывающей стороны. Если же каждое
обновление требует перезапуска всего процесса, это прерывает
выполняющиеся задачи и увеличивает объём работы по восстановлению
состояния.

В Cordis \textbf{обратимые эффекты (revertible effects)} описывают
преобразования контекста с обратными операциями; среда выполнения
сохраняет и выполняет соответствующие процедуры очистки.
\textbf{Реактивные коэффекты (reactive coeffects)} описывают требования
компонентов к окружению и координируют состояние компонентов в
соответствии с изменениями контекста. Единый механизм контекста также
координирует взаимодействие нескольких компонентов при их чередующемся
выполнении.\footnote{Shi, Yifan, Wei Zhang, and Tianyi Cui. \emph{A
  Programming Paradigm for Spatiotemporal Composability}, 2026,
  arXiv:2608.25512.}

Такая архитектура уменьшает зависимость локальных обновлений от полного
перезапуска. По мнению автора, компонуемость напрямую влияет на
способность системы непрерывно принимать самомодификации: после того как
модель создаст правильный патч, среде выполнения ещё необходимо
загрузить и выгрузить компоненты, согласовать зависимости и восстановить
задачи.

Модульная архитектура также должна сочетаться с процессами управления
полномочиями и выпуска. Динамическая среда выполнения \texttt{dsh}
управляет внутрипроцессными плагинами, их версиями, запуском, остановкой
и авторизацией, а постоянные плагины обслуживаются посредством
отдельного процесса управления установкой.\footnote{DeepSeek AI,
  \emph{DeepSeek Harness}, 2026. Состав плагинов описан в обзоре
  архитектуры; разделение обязанностей между динамической средой
  выполнения и постоянными плагинами --- в описании расширений. Проект
  был выпущен в августе 2026 года; в этом разделе рассматривается его
  модульная архитектура на этапе предварительной версии для
  разработчиков.}

Обновление плагинов влияет и на стоимость контекста. При изменении
определений инструментов или фрагментов промпта общая часть последующих
запросов и уже закэшированного префикса может сократиться.

\hypertarget{ux43fux440ux435ux43eux431ux440ux430ux437ux43eux432ux430ux43dux438ux435-ux43eux43fux44bux442ux430-ux432-ux43fux430ux440ux430ux43cux435ux442ux440ux44b}{%
\subsection{Преобразование опыта в
параметры}\label{ux43fux440ux435ux43eux431ux440ux430ux437ux43eux432ux430ux43dux438ux435-ux43eux43fux44bux442ux430-ux432-ux43fux430ux440ux430ux43cux435ux442ux440ux44b}}

Знания, инструкции и программы хорошо подходят для представления
информации и правил, которые можно выразить явно. Однако визуальные
паттерны на медицинских изображениях, естественная речевая просодия,
языковой стиль и привычки принятия долгосрочных решений также зависят от
представлений, усвоенных моделью из большого числа примеров. Для
способностей, которые трудно полностью описать в виде правил, автор
преобразует соответствующий опыт в обучающие материалы, а посредством
обновления параметров формирует многократно используемые паттерны
распознавания и принятия решений. Добавление новых текстовых напоминаний
зачастую не позволяет охватить всё их многообразие.

Одна задача может одновременно требовать нескольких носителей. При
редактировании китайских кавычек парсер определяет синтаксические
границы документа, Skill описывает принципы редактирования и исключения,
а обучение параметров улучшает определение области действия в разных
задачах. К сдвигу распределения изображений, вызванному сменой
оборудования, можно адаптироваться с помощью LoRA или других методов
дообучения; языковой стиль также можно обновлять посредством
периодического обучения предпочтений. Стабильность задачи влияет на
частоту и стоимость обновлений, а способ представления самой способности
определяет, какой метод следует попробовать в первую очередь.

Методы обучения, представленные в восьмой главе, можно напрямую включить
в замкнутый цикл этой главы: проверенные демонстрации поступают в SFT,
явные предпочтения образуют парные данные, а воспроизводимые
взаимодействия и надёжные вознаграждения позволяют применять RL.

Как определить, следует ли перенести некоторый опыт из инструкции в
параметры? Возьмите заранее отложенную тестовую выборку задач и
воспроизведите её дважды: в первый раз позвольте системе как обычно
искать и загружать инструкции, а во второй --- передайте соответствующие
инструкции непосредственно модели. Если во втором случае ответы заметно
лучше, сначала проверьте, не выбрала ли система неправильную инструкцию
и не пропустила ли инструкцию, которую следовало загрузить. Только если
инструкция уже передана, но модель при разных формулировках продолжает
повторять одну и ту же ошибку в суждении, появляется основание собрать
разнообразные демонстрации и проверить, улучшит ли обучение параметров
перенос. Два воспроизведения позволяют различить ситуации «опыт не
найден» и «опыт не удаётся применить».

\hypertarget{ux43eux442-ux43eux431ux43dux43eux432ux43bux435ux43dux438ux44f-ux440ux435ux437ux443ux43bux44cux442ux430ux442ux430-ux43a-ux43eux43fux442ux438ux43cux438ux437ux430ux446ux438ux438-ux43cux435ux442ux43eux434ux430-ux43eux431ux43dux43eux432ux43bux435ux43dux438ux44f}{%
\subsection{От обновления результата к оптимизации метода
обновления}\label{ux43eux442-ux43eux431ux43dux43eux432ux43bux435ux43dux438ux44f-ux440ux435ux437ux443ux43bux44cux442ux430ux442ux430-ux43a-ux43eux43fux442ux438ux43cux438ux437ux430ux446ux438ux438-ux43cux435ux442ux43eux434ux430-ux43eux431ux43dux43eux432ux43bux435ux43dux438ux44f}}

Если изменился один пункт политики возврата средств, достаточно обновить
соответствующую запись в базе знаний. Но если при каждом упорядочивании
базы знаний постоянно удаляется одно и то же исключение из правила,
изменить нужно сам метод упорядочивания записей. Опыт может храниться в
знаниях, инструкциях, программе или параметрах; уровень, который следует
исправить, зависит от этапа, на котором возникает ошибка. Иногда
достаточно изменить одну запись в памяти, а иногда приходится менять
этап формирования контекста, весь рабочий процесс, Harness или даже
оптимизатор, отвечающий за подготовку плана изменений.\footnote{Weng,
  Lilian. ``Harness Engineering for Self-Improvement.'' \emph{Lil'Log},
  2026.}

Автор по умолчанию выбирает локальные обновления: добавить одно правило
в промпт или дополнить документ с накопленным опытом одним исключением
--- так проще проверить влияние изменения и выполнить откат. По мере
роста числа записей необходимо совершенствовать и методы управления ими,
сохраняя детали уже накопленного опыта. Например, если многократно
переписывать документ целиком, редкое, но важное исключение может
постепенно исчезнуть, а несколько условий могут сжаться до чрезмерно
общего принципа. Agentic Context Engineering (ACE) организует контекст в
виде записей со стабильными идентификаторами: модули генерации,
рефлексии и упорядочивания предлагают инкрементальные изменения, которые
затем объединяются и дедуплицируются детерминированной
логикой.\footnote{Zhang, Qizheng, et al.~\emph{Agentic Context
  Engineering: Evolving Contexts for Self-Improving Language Models.}
  ICLR 2026. arXiv:2510.04618.} Это позволяет обслуживать каждую запись
опыта и сведения о её происхождении по отдельности.

Meta Context Engineering (MCE) идёт дальше и разделяет обновление
содержимого и обновление метода управления на два цикла. Внутренний цикл
формирует контекст задачи с помощью существующего метода; внешний
корректирует такие операции, как поиск, отбор, фильтрация и
форматирование, исходя из результатов выполнения и проверки.\footnote{Ye,
  Haoran, et al.~\emph{Meta Context Engineering via Agentic Skill
  Evolution.} arXiv:2601.21557, 2026.} Например, внутренний цикл
отбирает материалы для текущей задачи, а внешний сравнивает разные
стратегии фильтрации материалов и сохраняет версию, которая лучше
показывает себя в последующих задачах.

Объектами оптимизации могут стать также рабочий процесс и Harness. AFlow
представляет несколько вызовов LLM и управляющий ими поток в виде графа
кода и ищет их комбинации по обратной связи от выполнения.\footnote{Zhang,
  Jiayi, et al.~\emph{AFlow: Automating Agentic Workflow Generation.}
  ICLR 2025. arXiv:2410.10762.} Meta-Harness позволяет AI-агенту для
программирования изменять реализацию хранения, поиска и представления
информации с учётом исходного кода, оценок и траекторий.\footnote{Lee,
  Yoonho, et al.~\emph{Meta-Harness: End-to-End Optimization of Model
  Harnesses.} arXiv:2603.28052, 2026.} Так описанная в пятой главе
способность генерировать код включается в итеративный процесс оценки и
поиска: чтобы решить, с какой версии продолжать изменения в следующем
раунде, необходимо сравнить, что было изменено в каждой версии и какую
оценку она получила.

Чтобы выяснить, откуда именно взялось улучшение, исследовательская среда
прежде всего должна позволять изменять только один из уровней. Когда
автор поручал агенту самостоятельно развернуть среду обучения, тот часто
тратил время на подбор совместимых версий CUDA и PyTorch, а несколько
задач конкурировали за один и тот же GPU. RSIGym реализует обучение,
инференс, выборку траекторий, оценку и песочницу в виде переиспользуемых
сервисов, позволяя исследовательскому агенту вновь сосредоточиться на
проектировании экспериментов. В нём предусмотрены три контрольных
условия: при изменении только обучающих данных настройки обучения и
Harness остаются фиксированными; при изменении только Harness
фиксируются веса модели; при совместном улучшении можно изменять оба
компонента, сохраняя одну и ту же исходную модель и протокол
оценки.\footnote{Evolvent AI описала в документации проекта RSIGym
  методику оценки сервисов и три направления улучшения.}

Например, помощник сначала переходит в другой каталог в терминале, но
следующая команда не находит находящийся в нём файл. Если каждый из двух
вызовов запускает отдельный shell, достаточно изменить Harness так,
чтобы они использовали общий сеанс: тогда рабочий каталог сохранится, а
веса модели менять не потребуется.

У самого поиска тоже есть стратегия: с какого из существующих
результатов продолжить исследование, когда создать новую ветвь, сколько
параллельных воркеров задействовать и когда остановиться. Для оценки
таких решений обычно требуется несколько раундов генерации и выполнения.
Если для сравнения каждой стратегии поиска заново запускать все задачи,
стоимость получения обратной связи будет быстро расти.

\textbf{Воспроизведение дерева открытий} позволяет повторно использовать
уже полученные результаты исследования. Dream-RSI сохраняет процесс
одного открытия в виде дерева: корневой узел представляет исходную
рабочую область, а каждый дочерний узел содержит снимок рабочей области
после продолжения попыток от родительского узла, результат текущей
попытки, диагностические сведения и оценку.\footnote{Zheng, T., et
  al.~\emph{Dream-RSI: Recursive Self-Improvement through Evolving
  Worlds.} arXiv:2609.14858, 2026.} Тестируемая стратегия по-прежнему
выбирает ветвь через тот же интерфейс, что и при исследовании в реальном
времени, а система воспроизведения возвращает соответствующий ранее
записанный результат из исторического дерева, избавляя от повторных
вызовов нижележащего агента и оценщика.

Во время воспроизведения стратегия видит только уже раскрытое поддерево.
Соответствующий дочерний узел становится доступен лишь после того, как
стратегия решит продолжить исследование из некоторого листового узла или
создаст новую ветвь от корня. Это позволяет сравнивать порядок
исследования, параллельные пакеты и момент остановки, сохраняя при
каждом решении тот же объём доступной информации, что и в процессе
исследования в реальном времени.

Этот метод состоит из трёх этапов:

\begin{enumerate}
\def\labelenumi{\arabic{enumi}.}
\tightlist
\item
  \textbf{Исследование в реальном времени.} Текущая стратегия выполняет
  реальные задачи и создаёт новое дерево открытий.
\item
  \textbf{Создание среды воспроизведения.} Дерево открытий добавляется в
  исторический пул с сохранением состояния, результатов и отношений
  видимости каждого узла.
\item
  \textbf{Офлайн-оптимизация.} Агент, разрабатывающий стратегию,
  изменяет код исследования и на историческом пуле сравнивает качество
  найденных разными стратегиями результатов, стоимость выполнения и
  эффективность распараллеливания, после чего применяет лучшую стратегию
  для исследования новых задач.
\end{enumerate}

Текущая стратегия участвует в том же раунде воспроизведения. Если ни
одна новая стратегия не превзошла её по оценке, текущая стратегия
сохраняется. Новый запуск в реальном времени добавляет очередное дерево
открытий, постепенно расширяя набор опыта, доступного для сравнений в
следующем раунде.

Что именно извлекает каждый из этих трёх методов из исторических
записей? Обобщение траекторий сохраняет знания или инструкции для
следующей задачи; воспроизведение действий в браузере позволяет вновь
пройти уже проверенным путём. Dream-RSI воспроизводит дерево открытий,
чтобы ответить на другой вопрос: можно ли при том же вычислительном
бюджете найти лучшее решение, если изменить порядок исследования или
выполнять несколько ветвей параллельно? Историческое дерево позволяет
разработчикам многократно сравнивать варианты офлайн и с низкими
затратами, а итоговым изменением становится реализованная в программе
стратегия исследования.

\begin{quote}
\textbf{Эксперимент 9-8 ★★★: передайте эту книгу Hermes --- сможет ли он
усовершенствовать себя?}

Читательница Грейс задала вопрос: если не давать Hermes список
улучшений, а предоставить только десять глав этой книги и его
собственный исходный код, что он решит изменить? В этом запуске Hermes
выбрал запись траекторий. Он хотел преобразовать фактические результаты,
возвращаемые инструментами, в проверяемые доказательства для последующей
офлайн-оценки. Он мог предлагать патчи, а стабильная версия и правила
приёмки защищались независимыми правами доступа.

Почему первая версия была отклонена? Формат результатов инструмента в
тесте отличался от формата, который сохранялся в действительности.
Реальный результат был заключён в XML-теги, поэтому тесты проходили, но
парсер не мог извлечь это доказательство. Hermes получил обратную связь
и внёс изменения в том же сеансе; последующие проверки выявили ещё один
пропущенный путь сохранения и риск двойного учёта одной неудачи. После
трёх отклонений четвёртый независимый проверяющий одобрил патч, который
на тот момент ещё не был включён в Hermes.
\end{quote}

\hypertarget{ux43fux43eux441ux442ux440ux43eux435ux43dux438ux435-ux437ux430ux43cux43aux43dux443ux442ux43eux433ux43e-ux446ux438ux43aux43bux430-ux43dux435ux43fux440ux435ux440ux44bux432ux43dux43eux433ux43e-ux440ux430ux437ux432ux438ux442ux438ux44f-ux434ux43bux44f-ux434ux43eux43bux433ux43eux441ux440ux43eux447ux43dux43eux439-ux440ux430ux431ux43eux442ux44b}{%
\section{Построение замкнутого цикла непрерывного развития для
долгосрочной
работы}\label{ux43fux43eux441ux442ux440ux43eux435ux43dux438ux435-ux437ux430ux43cux43aux43dux443ux442ux43eux433ux43e-ux446ux438ux43aux43bux430-ux43dux435ux43fux440ux435ux440ux44bux432ux43dux43eux433ux43e-ux440ux430ux437ux432ux438ux442ux438ux44f-ux434ux43bux44f-ux434ux43eux43bux433ux43eux441ux440ux43eux447ux43dux43eux439-ux440ux430ux431ux43eux442ux44b}}

Пока одно поручение ещё не выполнено, фоновый процесс как раз вывел
новое правило из предыдущих случаев. Когда следует начать его применять?
На \crossreflink{fig:9-5}{рисунке 9-5} выполняющий текущую задачу
онлайн-цикл отделён от офлайн-цикла, который готовит следующую версию.
Новое правило можно использовать в последующих задачах только после
проверки и публикации. Версионируемые база опыта и оценочный набор
связывают оба цикла, благодаря чему для каждого обновления известны
источник и результат, а при необходимости можно вернуться к предыдущей
версии.

\begin{figure}
\hypertarget{fig:9-5}{%
\centering
\includegraphics{images/revised/fig9-5.pdf}
\caption{Рисунок 9-5. Двойной цикл онлайн-выполнения и
офлайн-развития}\label{fig:9-5}
}
\end{figure}

Предположим, пользователь просит AI-агента отслеживать заявление. При
первом запросе приходит ответ: «На рассмотрении». Две стрелки в левой
части \crossreflink{fig:9-5}{рисунка 9-5} показывают, какого момента
следует дождаться, прежде чем продолжить обработку. Если сервис
поддерживает подписку на изменение состояния, агент может оформить её и
возобновить задачу после получения уведомления. Если подписки нет и
можно только самостоятельно запрашивать состояние, модель может с учётом
срока обработки назначить себе однократную задачу по таймеру, а в нужный
момент повторить запрос, передав номер заявления и предыдущее состояние.
Если заявление всё ещё рассматривается, назначается следующая проверка;
если результат уже получен, агент сообщает его пользователю и прекращает
отслеживание.

Как использовать это ожидание для офлайн-обучения? Нужно записать в
траектории согласованное условие пробуждения, а затем сопоставить его с
фактическими событиями. Если подписка по-прежнему действует, а состояние
заявления не изменилось, спокойное ожидание соответствует плану.
Конкретное отклонение, допускающее анализ причин, возникает лишь тогда,
когда срок задачи по таймеру уже наступил, но запрос не выполнен, либо
когда результат уже доставлен, а агент продолжает повторять запросы.
Даже если напоминание и уведомление по подписке пришли одновременно,
сначала следует проверить, не закрыта ли уже задача, чтобы не сообщить
пользователю один и тот же результат дважды.

Voyager\footnote{Wang, G., et al.~\emph{Voyager: An Open-Ended Embodied
  Agent with Large Language Models.} arXiv:2305.16291, 2023.}
продемонстрировал в Minecraft совместную работу выбора целей, накопления
навыков и обратной связи от среды. \textbf{Автоматический генератор
учебной программы} предлагает следующую цель с учётом предметов, среды и
уже имеющихся способностей; \textbf{библиотека навыков} сохраняет
успешно проверенные программы для последующего поиска и комбинирования
--- например, задача по сбору ресурсов вызывает навыки перемещения и
изготовления предметов; \textbf{механизм итеративных промптов} передаёт
наблюдения за средой, ошибки выполнения и результаты самопроверки в
следующий цикл генерации кода. Новые цели стимулируют исследование,
проверка определяет, какие программы следует сохранить, а повторное
использование навыков создаёт отправную точку для последующих, более
сложных задач.

Без новых целей агент может бесконечно собирать один и тот же материал.
Если сохранить в качестве навыка даже программу, которой не удалось
изготовить предмет, следующая попытка изготовления начнётся с той же
ошибки.

Этот процесс можно обобщить как цикл открытий \textbf{«гипотеза ---
эксперимент --- оценка --- обратная связь»}. Для агента каждый патч
воплощает проверяемую гипотезу: какое поведение он должен улучшить, на
какие задачи повлияет и подтверждают ли результаты выполнения ожидания?
Основанная Джеффом Дином и его коллегами компания Discovery Loop также
ставит целью автоматизацию полного экспериментального цикла, в котором
постановка, проведение и оценка экспериментов определяют направление
следующего этапа исследований.\footnote{В команду основателей Discovery
  Loop входят Джефф Дин, Санджай Гемават, Квок Ле и Ориоль Виньялс. В
  официальном описании компании заявлена автоматизация полного
  экспериментального цикла, а в опубликованном инвестором 5 августа 2026
  года объявлении рассказывается о её статусе общественно полезной
  корпорации и направлении параллельного проведения экспериментов.} По
мнению автора, непрерывное развитие агента тоже должно следовать этому
научному методу: сформулировать чёткую гипотезу, проверить её на
практике, а затем по результатам решить, какие изменения сохранить, а
какие пересмотреть.

При оценке эффективности непрерывного развития необходимо также отдельно
проверять \textbf{способность обновлять среду агента (harness-updating)}
и \textbf{способность извлекать пользу из среды агента
(harness-benefit)}. Первая означает, способна ли система создавать
ценные долговременные изменения, а вторая --- способен ли выполняющий
задачу агент своевременно находить и правильно использовать эти
изменения. В экспериментах Линя и его коллег с заменой моделей
анализировалась связь обеих способностей с возможностями базовой
модели.\footnote{Lin, Minhua, et al.~\emph{Harness Updating Is Not
  Harness Benefit: Disentangling Evolution Capabilities in Self-Evolving
  LLM Agents.} arXiv:2605.30621, 2026.}

Например, Skill возврата средств может точно описывать политику, но в
последующей задаче агент может не найти его при поиске либо загрузить,
но пропустить обязательный шаг. В этих двух случаях необходимо отдельно
корректировать маршрутизацию и выполнение. Если при отсутствии роста
сквозной оценки сразу переписать Skill, можно потерять изначально
правильное правило, тогда как действительно неисправный этап загрузки
или выполнения так и останется без исправления. В таблице 9-3
последовательно задаются четыре вопроса: правильно ли внесено изменение?
Загружается ли оно при необходимости? Выполняется ли после загрузки?
Улучшились ли после выполнения результаты задачи?

Таблица 9-3. Метрики многоуровневой оценки непрерывного развития

\begin{longtable}[]{@{}
  >{\raggedright\arraybackslash}p{(\columnwidth - 4\tabcolsep) * \real{0.2514}}
  >{\raggedright\arraybackslash}p{(\columnwidth - 4\tabcolsep) * \real{0.3528}}
  >{\raggedright\arraybackslash}p{(\columnwidth - 4\tabcolsep) * \real{0.3957}}@{}}
\toprule\noalign{}
\begin{minipage}[b]{\linewidth}\raggedright
\hspace{0pt}Метрика
\end{minipage} & \begin{minipage}[b]{\linewidth}\raggedright
\hspace{0pt}Вопрос
\end{minipage} & \begin{minipage}[b]{\linewidth}\raggedright
\hspace{0pt}Основные доказательства
\end{minipage} \\
\midrule\noalign{}
\endhead
\bottomrule\noalign{}
\endlastfoot
\hspace{0pt}Доля эффективных предложений по обновлению &
\hspace{0pt}Способны ли предложенные изменения улучшить результаты &
\hspace{0pt}Доля принятых изменений и выигрыш при независимой
проверке \\
\hspace{0pt}Доля активации артефактов & \hspace{0pt}Загружается ли новый
Skill, элемент памяти или инструмент в подходящих сценариях &
\hspace{0pt}Журналы поиска, маршрутизации и вызовов \\
\hspace{0pt}Доля успешного следования & \hspace{0pt}Выполняются ли после
активации соответствующие правила или процессы &
\hspace{0pt}Последовательности действий и проверки процесса \\
\hspace{0pt}Выигрыш на независимом наборе задач & \hspace{0pt}Улучшает
ли обновление результаты на новых экземплярах, сохраняя прежние
способности & \hspace{0pt}Доля успехов, качество и стоимость отдельно
для новых задач, не участвовавших в настройке, и задач, которые система
умела выполнять ранее \\
\end{longtable}

При долгосрочной оценке следует одновременно отслеживать изменения пяти
типов:

\begin{itemize}
\tightlist
\item
  \textbf{Деградация способностей.} Не стали ли проваливаться ранее
  успешно выполнявшиеся случаи, не противоречат ли друг другу новый и
  старый опыт.
\item
  \textbf{Обобщение.} Какой выигрыш даёт обновление на независимых новых
  экземплярах, шаблонах или предметных областях.
\item
  \textbf{Расход токенов и стоимость выполнения.} Как изменяются число
  входных и сгенерированных токенов, число вызовов инструментов и время
  выполнения задачи.
\item
  \textbf{Безопасность и соответствие требованиям.} Сохраняется ли
  ожидаемое поведение в отношении прав доступа, конфиденциальности и
  границ отказа.
\item
  \textbf{Инженерное качество.} Как изменяются сложность сопровождения,
  архитектурная согласованность, владение компонентами, совместимость, а
  также стоимость миграции и отладки.
\end{itemize}

\begin{quote}
\textbf{Эксперимент 9-9 ★★★: оценка непрерывного развития агента}

Сравните статическую память, память только с добавлением записей и
версионируемую память, поддерживающую замену и удаление записей. Сначала
позвольте агенту получить обратную связь при выполнении набора задач,
затем последовательно изменяйте формулировки задач, обновляйте правила и
проверяйте сохранность прежних способностей.

Отдельно фиксируйте четыре типа результатов: запоминание, перенос,
обновление и сохранение. Может ли агент использовать релевантный опыт в
новых задачах, перейти на новую версию после изменения правил и
по-прежнему выполнять старые задачи?

В этом эксперименте для каждого из трёх вариантов памяти выполнялись
запуски с тремя начальными значениями генератора случайных чисел, в
каждом из которых последовательно решалось 14 задач, --- всего 126
вызовов модели. Средняя точность в задачах на перенос составила 100 \%
как для памяти только с добавлением, так и для памяти с возможностью
замены. Если смотреть только на этот показатель, может показаться, что
между подходами нет различий. Когда норму багажа изменили с 20 до 23
килограммов, обе группы при первом вопросе о новой норме ответили
неверно. Получив эту обратную связь, группа с возможностью замены
переписала правило и во второй раз ответила: 23 килограмма; группа
только с добавлением по-прежнему отвечала: 20 килограммов. Поэтому
первая группа правильно ответила лишь на половину из двух вопросов об
изменениях, но при последующей проверке того, заменило ли новое правило
старое, достигла 100 \%, тогда как точность замены в группе только с
добавлением составила 0. Правила возврата средств и проверки личности не
изменялись, поэтому доля сохранения прежних способностей в обеих группах
составила 100 \%.
\end{quote}

\hypertarget{ux434ux43eux43aux430ux437ux430ux442ux435ux43bux44cux441ux442ux432ux430-ux438-ux43eux446ux435ux43dux43aux430-ux432-ux43eux442ux43aux440ux44bux442ux44bux445-ux437ux430ux434ux430ux447ux430ux445}{%
\subsection{Доказательства и оценка в открытых
задачах}\label{ux434ux43eux43aux430ux437ux430ux442ux435ux43bux44cux441ux442ux432ux430-ux438-ux43eux446ux435ux43dux43aux430-ux432-ux43eux442ux43aux440ux44bux442ux44bux445-ux437ux430ux434ux430ux447ux430ux445}}

Тесты кода, результаты работы инструментов и состояние бизнес-процессов
позволяют довольно быстро получить обратную связь. В научных
исследованиях, стратегическом планировании и проектировании продуктов
результаты зачастую становятся заметны лишь спустя гораздо больше
времени; кроме того, необходимо оценивать ценность проблемы, качество
решения и его сопровождаемость. Конвейер может успешно выполнить поиск,
написать код и подготовить текст, но итоговый результат всё равно может
не отвечать на вопрос, который действительно стоило исследовать, либо не
содержать надёжных базовых решений и новых доказательств. При
проектировании замкнутого цикла эволюции автор прежде всего
совершенствует структуру доказательств и надзора, оценивая выполнение
процесса отдельно от этих долгосрочных целей.

Трехан и Чопра описали четыре попытки автоматизированного научного
исследования --- от исследовательской идеи до статьи. Три из них
завершились неудачей на этапе реализации или оценки, а одна прошла весь
конвейер до конца.\footnote{Trehan, Dhruv and Paras Chopra. \emph{Why
  LLMs Aren't Scientists Yet: Lessons from Four Autonomous Research
  Attempts.} arXiv:2601.03315, 2026.} Эти случаи выявили три конкретные
проблемы: \textbf{дрейф реализации}, из-за которого код постепенно
отклоняется от исходной гипотезы; \textbf{преждевременная интерпретация
результатов}, из-за которой шум принимают за открытие, а неудачам и
отрицательным результатам не уделяют равного внимания;
\textbf{недостаточность исследовательских суждений}, проявляющаяся в
необоснованном выборе базовых решений, отслеживании аномалий и
определении момента остановки.

Структуру доказательств и надзора в таких задачах можно улучшить по
четырём направлениям.

\begin{itemize}
\tightlist
\item
  \textbf{Связывать выводы с источниками.} Раздельно фиксировать
  доказательства для цитат, чисел, методов и умозаключений. В
  ScientistOne архитектура Chain-of-Evidence связывает утверждения с
  проверяемыми источниками, позволяя проследить путь от рукописи до
  эксперимента или исходного материала.\footnote{Meng, et
    al.~\emph{ScientistOne: Towards Human-Level Autonomous Research via
    Chain-of-Evidence.} arXiv:2605.26340, 2026.}
\item
  \textbf{Сохранять отрицательные результаты.} Архивировать неудачные
  эксперименты, отклонённые решения и причины остановки. Если сохранять
  только итоговое успешное решение, при последующем поиске не будет
  видно, какие пути уже были отвергнуты, и возникнет риск снова понести
  те же издержки на пробы и ошибки.
\item
  \textbf{Одновременно исследовать разные подходы.} Помимо текущего
  решения с высокой оценкой, сохранять несколько ветвей, которые пока
  оценены ниже, но основаны на иных идеях, и сравнивать их результаты
  при разных условиях оценки. Тогда, даже если один шаблон ответа легко
  получает высокие оценки, у команды останется возможность проверить
  другие подходы.
\item
  \textbf{Чётко определять исследовательские обязанности человека.}
  Человек участвует в постановке проблемы, проверке критериев оценки,
  интерпретации аномальных результатов и принятии решения об остановке,
  а AI-агент отвечает за реализацию, выполнение и организацию
  материалов.
\end{itemize}

При работе над книгой «Глубокое понимание AI Infra» автор сталкивался с
ситуациями, когда работа была закончена, но изложение по-прежнему
оставалось непонятным. Сначала он просто передал помощнику план: тот
смог написать целую главу, однако в ней не было подробностей, которые
стоило бы объяснить. Затем автор стал сначала поручать ему поиск научных
статей и официальных спецификаций чипов, извлечение размеров матриц из
архитектуры модели и разработку по этим исходным данным инструмента для
расчёта производительности, а уже после этого они вместе корректировали
план. Благодаря этому при рецензировании появились две точки опоры:
встретив показатель производительности, можно было по конфигурации найти
исходные данные и заново выполнить расчёт; встретив объяснение
архитектуры, следовало спросить, помогает ли этот показатель читателю
понять, почему автор выбрал именно такое проектное решение. Даже при
безупречно верной формуле фрагмент с неверно выбранным объектом
сравнения всё равно требовалось переписать. Обратную связь первого типа
можно превратить в автоматическую проверку, а замечания автора второго
типа необходимо сохранить, чтобы следующая версия опиралась на
конкретные критерии оценки.\footnote{В сопутствующем проекте книги
  «Глубокое понимание AI Infra» --- ai-infra-book --- инструменты
  расчёта, конфигурации и тесты хранятся в каталоге
  \texttt{calculations/}.}

Например, помощник по написанию текстов сократил вступление с учётом
отзывов читателей. Специалист по оценке может при одинаковом требовании
к объёму анонимно распределить новую и старую версии одной рукописи
между целевыми читателями, попросить их своими словами сформулировать
основную проблему и отметить фрагменты, которые помогли им её понять.
Если группа, читавшая новую версию, не сможет объяснить, почему автор
поставил эту проблему, отмеченные обеими группами фрагменты помогут
определить, какая важная предпосылка была удалена.

\hypertarget{ux433ux440ux430ux43dux438ux446ux44b-ux431ux435ux437ux43eux43fux430ux441ux43dux43eux441ux442ux438-ux43dux435ux43fux440ux435ux440ux44bux432ux43dux43eux439-ux44dux432ux43eux43bux44eux446ux438ux438}{%
\subsection{Границы безопасности непрерывной
эволюции}\label{ux433ux440ux430ux43dux438ux446ux44b-ux431ux435ux437ux43eux43fux430ux441ux43dux43eux441ux442ux438-ux43dux435ux43fux440ux435ux440ux44bux432ux43dux43eux439-ux44dux432ux43eux43bux44eux446ux438ux438}}

Долговременные возможности многократно используются в последующих
задачах, поэтому при каждом обновлении необходимо чётко ответить на три
вопроса: откуда взяты основания для изменения, как тестируется новая
версия и кто вправе одобрить её выпуск. Если ошибочная инструкция с
веб-страницы попадёт в хранилище опыта, она повлияет на множество
сессий; если внешняя зависимость войдёт в состав часто используемого
инструмента, масштаб последствий увеличится; дефект проверяющего
механизма повлияет на оценку целого ряда обновлений. Поэтому следует
установить три границы: для информации, версий и полномочий.

Первая граница --- \textbf{изоляция доказательств от инструкций}.
Веб-страницы, письма, вывод инструментов и подготовленные по ним резюме
рассматриваются как материалы, требующие проверки. Система по
фиксированной структуре извлекает утверждения, расположение исходного
текста и время получения, сохраняя источники и оригинал, чтобы
проверяющий мог проследить происхождение каждого элемента опыта. Резюме
и указанная моделью уверенность помогают упорядочивать материалы, однако
переписанный текст всё ещё может содержать ошибочные утверждения или
вредоносные инструкции из исходных материалов. Поэтому решение о выпуске
должно основываться на проверке структуры, списках разрешений, проверке
источников и независимом ревью. AI-агент отправляет новый или изменённый
Skill в виде версионированного запроса на изменение, а проверяющий
сопоставляет его содержимое и поведение с доказательствами. Обновления с
высоким риском дополнительно требуют одобрения человека.

Вторая граница --- \textbf{разделение тестовой и производственной сред}.
Новые элементы знаний, изменённые промпты и Skills, а также новые
программы и параметры модели сначала помещаются в тестовую среду. Код и
внешние зависимости проходят в ограниченной песочнице проверку
разрешений, сканирование цепочки поставок и поведенческие тесты, после
чего сравниваются со стабильной версией в ходе регрессионного
тестирования. Только версия, прошедшая пороги выпуска, постепенно
подключается к реальным задачам. Стабильная версия сохраняется для
отката в случае отклонений.

Третья граница --- \textbf{разделение полномочий на обновление и
одобрение}. Бизнес-агент может предлагать изменения знаний,
инструментов, Skills и компонентов бизнес-проверки. Наборы тестов,
пороги выпуска, конфигурация разрешений, журналы аудита и резервные
копии стабильных версий, на основании которых одобряются эти изменения,
поддерживаются независимой плоскостью управления. Даже изменив
собственную программу проверки, бизнес-агент не получает права ослабить
требования к выпуску или одобрить этот выпуск. Если бизнес-агент сможет
изменять и тесты, он сможет повысить свою оценку, снизив порог или
удалив неудачные тестовые сценарии, хотя его реальные возможности не
улучшатся.

\hypertarget{ux43eux431ux443ux447ux435ux43dux438ux435-ux432ux43e-ux441ux43dux435-ux43aux43eux43dux441ux43eux43bux438ux434ux430ux446ux438ux44f-ux437ux430ux431ux44bux432ux430ux43dux438ux435-ux438-ux43fux43eux434ux434ux435ux440ux436ux430ux43dux438ux435-ux430ux43aux442ux443ux430ux43bux44cux43dux43eux441ux442ux438-ux432ux43eux437ux43cux43eux436ux43dux43eux441ux442ux435ux439}{%
\subsection{Обучение во сне: консолидация, забывание и поддержание
актуальности
возможностей}\label{ux43eux431ux443ux447ux435ux43dux438ux435-ux432ux43e-ux441ux43dux435-ux43aux43eux43dux441ux43eux43bux438ux434ux430ux446ux438ux44f-ux437ux430ux431ux44bux432ux430ux43dux438ux435-ux438-ux43fux43eux434ux434ux435ux440ux436ux430ux43dux438ux435-ux430ux43aux442ux443ux430ux43bux44cux43dux43eux441ux442ux438-ux432ux43eux437ux43cux43eux436ux43dux43eux441ux442ux435ux439}}

«Обучение во сне» использует метафору упорядочения опыта во время сна
для описания процесса офлайн-консолидации AI-агента. Во время
онлайн-работы агент выполняет текущую задачу и добавляет исходные
доказательства. В период простоя или после выполнения условий запуска
фоновый процесс обучения централизованно читает новый опыт, сравнивает
новые и старые выводы, объединяет дубликаты, обрабатывает конфликты,
затем оформляет необходимые изменения в виде патчей и запускает
регрессионные тесты. Автор рекомендует разделять сбор и упорядочение:
единичный случайный успех, сетевой сбой или непроверенный ввод сначала
следует сохранить как доказательство, ожидающее анализа; решение об
обновлении принимается только после накопления однотипных траекторий.
Пакетная обработка не только упрощает выявление устойчивых
закономерностей, но и позволяет использовать свободные вычислительные
ресурсы и недорогие модели, подходящие для задач упорядочения.

Типичный цикл состоит из пяти шагов:

\begin{enumerate}
\def\labelenumi{\arabic{enumi}.}
\tightlist
\item
  \textbf{Запуск.} Когда наступает запланированное время либо количество
  новых траекторий, занятое место в хранилище или частота ошибок
  достигают порогового значения, задача упорядочения сначала ставится в
  очередь. Перед началом необходимо ещё раз проверить, нет ли более
  важных онлайн-задач, ожидающих доступа к тем же ресурсам; если есть,
  упорядочение следует отложить.
\item
  \textbf{Проверка существующих возможностей.} Чтение утверждённых
  знаний, промптов, каталога Skills и их версий для подтверждения
  имеющихся возможностей и прав на изменение каждого компонента.
\item
  \textbf{Сбор и консолидация.} Из недавно оценённых траекторий
  извлекаются новые сигналы, объединяются дубликаты, отмечаются
  конфликты и условия применимости, после чего формируются локальные
  патчи.
\item
  \textbf{Проверка и одобрение.} Новая и старая версии по отдельности
  выполняют три типа задач: новые задачи, похожие на исходные случаи;
  задачи, которые старая версия уже решала правильно; задачи,
  затрагивающие ограничения безопасности. Изменения с высоким риском
  одобряет человек.
\item
  \textbf{Очистка и индексирование.} Обновляется поисковый индекс,
  устаревшему содержимому присваивается соответствующий статус,
  малоценное содержимое архивируется или удаляется, а источники и
  восстанавливаемые версии сохраняются.
\end{enumerate}

Claude Code ведёт для проекта краткий индекс памяти \texttt{MEMORY.md} и
разделённые по темам подробные файлы. При запуске сессии содержимое
начала индекса загружается с учётом ограничения длины, а тематические
файлы читаются по мере необходимости. Когда индекс приближается к
пределу объёма, подробности переносятся в соответствующие темы, а
дублирующиеся записи объединяются, благодаря чему ограниченный контекст
продолжает вмещать наиболее полезную информацию. Автоматическая память
Claude Code накапливается преимущественно в ходе работы сессий и
демонстрирует совместное применение многоуровневой загрузки и активного
упорядочения.\footnote{Anthropic, ``How Claude remembers your project'',
  2026.}

Hermes разделяет долговременную информацию между несколькими носителями:
файлом \texttt{MEMORY.md} с данными об окружении и принятых соглашениях,
файлом \texttt{USER.md} с пользовательскими предпочтениями --- объём
обоих файлов ограничен, --- поиском по истории сессий на основе
SQLite/FTS5, загружаемыми по требованию Skills и необязательными
внешними сервисами памяти. Исторический поиск возвращает исходные
сообщения, сохраняя для последующего анализа прослеживаемые основания.
Задачи, включающие сложные операции с инструментами, восстановление
после ошибок или исправления от пользователя, могут служить материалом
для ретроспективного анализа Skill; необходимость одобрения записи
определяется конфигурацией среды выполнения. Отдельный Curator управляет
использованием, устареванием и архивным состоянием Skills: фиксированные
правила удаляют утратившее актуальность содержимое, а необязательное
объединение с помощью модели выполняет перезапись и устранение
дубликатов. Перед объединением содержимого сохраняется снимок, архивные
материалы также можно восстановить, что позволяет отменить неудачное
упорядочение.\footnote{Nous Research, \emph{Hermes Agent Documentation:
  Persistent Memory, Skills System, and Curator}, 2026.}

По мере накопления опыта сопровождение постепенно расширяется от
«добавления содержимого» до «поддержания его актуальности».
Дублирующиеся документы мешают поиску, противоречащие друг другу правила
увеличивают нагрузку при принятии решений, а многократное обновление
параметров может ослабить прежние возможности. Поэтому автор
рассматривает объединение, удаление и повторную проверку как обычные
составляющие непрерывной эволюции. Если только добавлять опыт, нагрузка
на поиск и соблюдение инструкций будет расти вместе с ним, постепенно
сводя на нет пользу от обучения. Офлайн-упорядочение можно представить в
виде шести конкретных задач:

\begin{itemize}
\tightlist
\item
  Объединять дублирующийся опыт, сохраняя источники и связи между
  версиями.
\item
  Переносить локальные правила предметной области в соответствующий
  Skill, чтобы глобальный промпт выражал общие требования.
\item
  Упорядочивать структуру промптов и Skills, чтобы условия запуска, шаги
  и критерии готовности было легко найти.
\item
  Повторно проверять давно не использовавшиеся инструменты, обновляя
  сведения об интерфейсах и зависимостях.
\item
  Выводить из области поиска по умолчанию знания, опровергнутые новыми
  доказательствами, сохраняя необходимую историю.
\item
  Поддерживать версии параметров с учётом результатов регрессионных
  тестов; когда необходимо устранить накопленный эффект изменений,
  заново обучать LoRA на исходной базовой модели и сравнивать её с
  существующей версией.
\end{itemize}

\hypertarget{ux438ux442ux43eux433ux438-ux433ux43bux430ux432ux44b-8}{%
\section{Итоги
главы}\label{ux438ux442ux43eux433ux438-ux433ux43bux430ux432ux44b-8}}

Одна неудачная траектория может подсказать, где искать проблему. Но
чтобы решить, что изменить, необходимо также знать, какие факты были
доступны ассистенту в тот момент и какую версию правил он использовал.
Если фактов не хватало, можно дополнить знания; если факты были
известны, но один и тот же шаг постоянно пропускался, его можно добавить
в промпт или Skill; если условия и результат допускают точное описание,
выполнение можно поручить программе. Только для суждений, которые трудно
описать этими способами, стоит дополнительно рассматривать обучение
параметров модели. Автор предпочитает сначала пробовать изменения,
область влияния которых можно определить и которые легко отменить.

Как сделать исправление из одного диалога способностью следующей версии
AI-агента? При онлайн-выполнении сначала сохраняются исходные записи, а
изменения предлагаются позднее, в ходе офлайн-анализа. С помощью
независимых граничных примеров проверяется, стало ли меньше однотипных
ошибок, а с помощью задач, которые раньше решались правильно, ---
сохранились ли прежние способности. Именно для этого в цикле открытий из
первой главы нужны \textbf{эксперименты и обратная связь}. Прошедшие
проверку изменения выпускаются с сохранением предыдущей версии для
отката; устаревший опыт периодически удаляется, чтобы при поиске не
извлекались одновременно новые и старые правила.

Изменять можно даже сам способ исследования. Дерево исторических
открытий сохраняет опробованные в тот момент ветви и их результаты,
позволяя сравнивать стратегии дальнейшего углубления, перехода на другую
ветвь и перераспределения бюджета. Как только проверяемая стратегия
переходит на ветвь, которую ранее никто не исследовал, при
воспроизведении не будет результата, по которому её можно было бы
оценить, поэтому приходится проводить отдельный онлайн-эксперимент. Чем
проще проверить результат задачи, тем лучше такие изменения подходят для
частого внесения. В научных исследованиях и стратегическом планировании
человек по-прежнему должен решать, стоит ли вообще заниматься самой
проблемой и какой результат считать прогрессом.

Если обнаружение проблем, предложение изменений и ревью результатов
поручить разным агентам, возникает новая сложность: как ревьюеру
убедиться, что он видит именно ту версию материалов, на основании
которой инициатор внёс изменения? Следующая глава начинается с задачи, в
которой несколько человек переводят одну техническую книгу.

\begin{questionbox}

\hypertarget{ux432ux43eux43fux440ux43eux441ux44b-ux434ux43bux44f-ux440ux430ux437ux43cux44bux448ux43bux435ux43dux438ux44f-8}{%
\section*{Вопросы для
размышления}\label{ux432ux43eux43fux440ux43eux441ux44b-ux434ux43bux44f-ux440ux430ux437ux43cux44bux448ux43bux435ux43dux438ux44f-8}}
\addcontentsline{toc}{section}{Вопросы для размышления}

\begin{enumerate}
\def\labelenumi{\arabic{enumi}.}
\tightlist
\item
  ★★ Документ с описанием опыта содержит три успешные траектории и одну
  неудачную. Неудача произошла с более новой версией API. Как отличить
  утрату актуальности опыта от изменения условий его применимости?
\item
  ★★ Удовлетворённость пользователей AI-агента службы поддержки и
  частота нарушений правил растут одновременно. Как объединить сигналы
  обучения и задать защитные метрики для выпуска?
\item
  ★★★ Некоторый агент часто заявляет, что операция завершена, хотя
  записи инструментов показывают, что она ещё не выполнялась. Как с
  учётом первопричины выбрать между изменением промпта, проверками
  Harness и обучением параметров?
\item
  ★★★ Агенту необходимо изменять инструменты и компоненты
  бизнес-валидации, а контрольная плоскость, одобряющая его собственные
  обновления, поддерживается независимым модулем. Как разделить
  полномочия, код и тесты между этими двумя сторонами?
\item
  ★★ После роста базы знаний, накопленных на опыте, возникают ошибки
  поиска и конфликты правил. Как спроектировать механизмы
  версионирования, контроля актуальности, объединения и удаления
  устаревших знаний?
\item
  ★★★ Разработайте для медицинской службы поддержки схему непрерывной
  эволюции: опишите распределение ответственности между параметрами
  модели, знаниями, Skill и ограничениями на уровне кода, а также
  процессы проверки и выпуска обновлений каждого типа.
\item
  ★★★ Некоторая стратегия исследования получает наивысшую оценку на всех
  ранее обнаруженных деревьях, но деградирует в следующем раунде
  онлайн-исследования. Как проанализировать смещение области поддержки
  распределения, случайность и переобучение под процедуру оценки, а
  затем на основе этого разделить среды воспроизведения и задать
  критерии допуска к выпуску?
\end{enumerate}

\textbf{Практикум к \crossreflink{chap:9}{главе 9}.}
\href{https://github.com/bojieli/ai-agent-book/tree/40f4ce5778e6cc93e3606c3079bf8cf13a0c4d04/chapter9}{Каталог
экспериментов} ·
\href{https://github.com/bojieli/ai-agent-book/blob/40f4ce5778e6cc93e3606c3079bf8cf13a0c4d04/chapter9/README.ru.md}{Путеводитель
на русском}.

\end{questionbox}

\hypertarget{ux441ux43eux432ux43cux435ux441ux442ux43dux430ux44f-ux440ux430ux431ux43eux442ux430-ux43dux435ux441ux43aux43eux43bux44cux43aux438ux445-ai-ux430ux433ux435ux43dux442ux43eux432}{%
\chapter{\texorpdfstring{Совместная работа нескольких
AI-агентов\label{chap:10}}{Совместная работа нескольких AI-агентов}}\label{ux441ux43eux432ux43cux435ux441ux442ux43dux430ux44f-ux440ux430ux431ux43eux442ux430-ux43dux435ux441ux43aux43eux43bux44cux43aux438ux445-ai-ux430ux433ux435ux43dux442ux43eux432}}

Первые девять глав были посвящены созданию и совершенствованию
отдельного агента: организации контекста, подключению знаний,
инструментов и каналов взаимодействия, а затем развитию его возможностей
с помощью оценки, постобучения и непрерывной эволюции. При переводе всей
книги на другой язык один и тот же термин во всех десяти главах должен
переводиться единообразно. Десять участников могут быстро подготовить
десять черновиков, если каждый переведёт по одной главе. Но если в
процессе перевода четвёртой главы один термин изменится, кто поймёт, что
его нужно изменить и в первой? Отталкиваясь от подобных задач, требующих
передачи работы между участниками, эта глава рассматривает, как
распределять обязанности, обмениваться информацией и проверять
результаты друг друга.

OpenAI ранее описывала развитие возможностей AI с помощью пяти уровней:
собеседники, мыслители (Reasoners), агенты, новаторы и организации
(Organizations). «Организации» здесь обозначают способность AI выполнять
работу целой организации. Системы из нескольких агентов предлагают один
из возможных подходов к реализации этой способности: распределить
обязанности в соответствии с общей целью и организовать совместное
выполнение задачи несколькими исполнительными единицами.

По мнению автора, значение систем из нескольких агентов заключается
также в исследовании \textbf{коллективного интеллекта}. Благодаря
разделению труда, сотрудничеству, дискуссиям и накоплению знаний люди
создают общества и организации, возможности которых превосходят
возможности отдельного человека. Группы агентов также могут развиваться
в похожем направлении: каждый участник выполняет локальную задачу, а
посредством обмена информацией и взаимной проверки они совместно решают
более сложные проблемы. В работе «From AGI to ASI» Google DeepMind
называет масштабные коллективы из множества агентов одним из путей к
сверхинтеллекту и рассматривает, как несколько систем с универсальными
возможностями могут объединиться в целое с более развитыми когнитивными
способностями.\footnote{Google DeepMind, \emph{From AGI to ASI.}
  arXiv:2606.12683, 2026.} При таком подходе, даже если каждый агент
достигает лишь уровня человеческого эксперта, хорошо организованная
группа может развить возможности, превосходящие совокупные возможности
человеческих экспертов. Отсюда возникает вопрос, требующий длительного
изучения в инженерии систем из нескольких агентов: какие механизмы
разделения труда и обмена информацией позволяют группе стабильно
получать результаты, превосходящие результаты отдельных участников?

\hypertarget{ux43aux43bux430ux441ux441ux438ux444ux438ux43aux430ux446ux438ux44f-ux432ux437ux430ux438ux43cux43eux434ux435ux439ux441ux442ux432ux438ux44f-ux43dux435ux441ux43aux43eux43bux44cux43aux438ux445-ai-ux430ux433ux435ux43dux442ux43eux432}{%
\section{Классификация взаимодействия нескольких
AI-агентов}\label{ux43aux43bux430ux441ux441ux438ux444ux438ux43aux430ux446ux438ux44f-ux432ux437ux430ux438ux43cux43eux434ux435ux439ux441ux442ux432ux438ux44f-ux43dux435ux441ux43aux43eux43bux44cux43aux438ux445-ai-ux430ux433ux435ux43dux442ux43eux432}}

Переводчик передаёт рукопись редактору, и им необходимо договориться о
двух вещах. Во-первых, что может видеть редактор: только перевод и
глоссарий или также весь диалог, сопровождавший перевод? Во-вторых, кто
вправе потребовать доработки при обнаружении ошибки в переводе? Первое
определяет способ совместного использования контекста, второе --- кто
кому передаёт работу. Эти решения можно принимать независимо друг от
друга.

\hypertarget{ux430ux441ux43fux435ux43aux442-1-ux43eux431ux449ux438ux439-ux43bux438-ux43aux43eux43dux442ux435ux43aux441ux442}{%
\subsection{Аспект 1: общий ли
контекст}\label{ux430ux441ux43fux435ux43aux442-1-ux43eux431ux449ux438ux439-ux43bux438-ux43aux43eux43dux442ux435ux43aux441ux442}}

\textbf{Общий контекст} означает, что последующая роль продолжает
использовать переданные с предыдущего этапа историю диалога, результаты
вызовов инструментов и состояние задачи. Система меняет инструкции для
роли и конфигурацию инструментов, после чего новая роль продолжает
работу. Например, когда аналитик требований подготовил документ с
требованиями, разработчик получает не только документ, но и возможность
просмотреть историю общения аналитика с пользователем. Такой подход
упрощает отслеживание деталей, однако по мере выполнения задачи требует
применения описанных во второй главе методов управления контекстом,
чтобы контролировать длину истории.

\textbf{Раздельный контекст} означает, что каждый агент отдельно
поддерживает входные данные модели, историю диалога и состояние задачи,
а обмен информацией происходит через явно заданные интерфейсы. Например,
аналитик передаёт разработчику документ с требованиями, разработчик
реализует на его основе функциональность, а при возникновении неясностей
отправляет аналитику сообщение с вопросом. Рабочий процесс каждого
участника сохраняется в отдельном контексте, а получатель по мере
необходимости запрашивает материалы, относящиеся к его обязанностям.
Чёткие интерфейсы также упрощают независимую разработку, тестирование и
замену участников.

Разработчику в левой части \crossreflink{fig:10-1}{рисунка 10-1} нужно
уточнить, какие файлы должен обрабатывать CSV-скрипт, и он может
обратиться к диалогу пользователя с аналитиком. Редактору в правой части
нужно уточнить перевод термина, и у него есть перевод и глоссарий. Если
же он захочет узнать, почему переводчик выбрал именно такой вариант, ему
придётся отдельно задать вопрос или запросить записи.

\begin{figure}
\hypertarget{fig:10-1}{%
\centering
\includegraphics{images/revised/fig10-1.pdf}
\caption{Рисунок 10-1. Сравнение общего и раздельного
контекста}\label{fig:10-1}
}
\end{figure}

Как участники обмениваются материалами, если их контексты независимы? В
распределённых системах для коммуникации обычно применяются общее
хранилище или передача сообщений. В системах агентов также используются
эти два подхода, в частности:

\begin{itemize}
\tightlist
\item
  \textbf{Параметры вызова инструмента.} Нижестоящий агент
  инкапсулируется в виде инструмента, а вышестоящий передаёт через
  параметры задачу и структурированные материалы, после чего получает
  результат.
\item
  \textbf{Общая файловая система.} Одна сторона записывает документы,
  код или данные, а другая читает их по ссылке; этот способ подходит для
  крупных промежуточных результатов и материалов, требующих постоянного
  хранения.
\item
  \textbf{Шина сообщений (Message Bus).} Агент отправляет сообщение
  посредническому сервису, который доставляет его согласно указанному
  получателю или подписке; этот способ подходит для асинхронного обмена
  между несколькими участниками.
\end{itemize}

Например, отправив задачу на перевод, вышестоящий агент может дождаться
перевода или сначала получить номер задачи и перейти к обработке другой
главы. Даже если оба переводчика могут открыть один файл, система всё
равно должна определить, кто его записывает и когда другой участник его
читает. В сообществе Go принято организовывать доступ к общим данным
посредством коммуникации. Этот принцип применим и к агентам: необходимо
явно определить, кому принадлежит состояние, кто может его изменять и
как остальные участники получают обновления.

\hypertarget{ux430ux441ux43fux435ux43aux442-2-ux442ux43eux43fux43eux43bux43eux433ux438ux44f-ux432ux437ux430ux438ux43cux43eux434ux435ux439ux441ux442ux432ux438ux44f}{%
\subsection{Аспект 2: топология
взаимодействия}\label{ux430ux441ux43fux435ux43aux442-2-ux442ux43eux43fux43eux43bux43eux433ux438ux44f-ux432ux437ux430ux438ux43cux43eux434ux435ux439ux441ux442ux432ux438ux44f}}

Топология взаимодействия описывает зависимости между задачами, потоки
информации и отношения управления между агентами. В этой главе
рассматриваются три типичных способа организации:

\begin{itemize}
\tightlist
\item
  \textbf{Паттерн равноправного взаимодействия} (Peer Collaboration
  Pattern). Небольшое количество участников итеративно работает над
  одним результатом. Например, автор подготавливает черновик научной
  статьи, рецензент предлагает изменения, после чего автор дорабатывает
  текст.
\item
  \textbf{Паттерн оркестрации} (Orchestration Pattern). Централизованный
  агент-менеджер отвечает за декомпозицию, распределение и сведение
  результатов задач, а подчинённые агенты выполняют каждый свою
  специализированную работу.
\item
  \textbf{Децентрализованный паттерн} (Decentralized Pattern). Участники
  продолжают работу с учётом своих ролей, сообщений и локального
  состояния, совместно определяя её дальнейший ход.
\end{itemize}

\begin{quote}
\textbf{Пояснение терминологии: Graph Engineering.} В первой главе
рассматривалось обсуждение Graph Engineering, состоявшееся в июле 2026
года. Этот подход сосредоточен на графе выполнения: узлами могут быть
агенты, обычные программы или решения, принимаемые человеком; рёбра
выражают зависимости, условную маршрутизацию и направления обработки
после сбоев; состояние передаётся по связям выполнения. Топология
взаимодействия нескольких агентов --- одно из применений этого подхода к
проектированию, в котором основное внимание уделяется передаче работы,
параллельному выполнению и обратной связи между несколькими центрами
принятия решений.
\end{quote}

Равноправное взаимодействие, оркестрацию менеджером и децентрализованную
схему можно комбинировать. Например, менеджер передаёт задачу по
подготовке текста группе из автора и рецензента, а группа после
внутренних итераций отправляет ему результат.

\hypertarget{ux43aux43eux433ux434ux430-ux43dux435ux441ux43aux43eux43bux44cux43aux43e-ai-ux430ux433ux435ux43dux442ux43eux432-ux43fux440ux435ux432ux43eux441ux445ux43eux434ux44fux442-ux43eux434ux43dux43eux433ux43e-ux430ux433ux435ux43dux442ux430}{%
\section{Когда несколько AI-агентов превосходят одного
агента}\label{ux43aux43eux433ux434ux430-ux43dux435ux441ux43aux43eux43bux44cux43aux43e-ai-ux430ux433ux435ux43dux442ux43eux432-ux43fux440ux435ux432ux43eux441ux445ux43eux434ux44fux442-ux43eux434ux43dux43eux433ux43e-ux430ux433ux435ux43dux442ux430}}

Предположим, что в черновом переводе неверно переведён технический
термин. Если дать двум редакторам один и тот же перевод, оба могут
одинаково одобрить его. Но если один из них получит также оригинал,
глоссарий и код, в котором используется этот термин, появится ещё один
способ обнаружить ошибку. При выборе модели взаимодействия сначала
следует спросить: что новый агент привносит в задачу? Он может получить
новые доказательства, исследовать другой путь поиска, предложить иное
профессиональное суждение или сократить время ожидания благодаря
параллельному выполнению.

Автор сначала оценивает \textbf{прирост информации от обратной связи}, а
затем --- пользу от параллельной работы и изоляции контекста. Получив
результаты тестов, скриншоты страниц или внешние материалы, ревьюер
может предложить изменения на основе доказательств, недоступных во время
генерации. В таблице 10-1 на нескольких распространённых сценариях
показано, как оценивать такой прирост.

Таблица 10-1. Информация и обратная связь при взаимодействии нескольких
агентов

\begin{longtable}[]{@{}
  >{\raggedright\arraybackslash}p{(\columnwidth - 4\tabcolsep) * \real{0.3185}}
  >{\raggedright\arraybackslash}p{(\columnwidth - 4\tabcolsep) * \real{0.3182}}
  >{\raggedright\arraybackslash}p{(\columnwidth - 4\tabcolsep) * \real{0.3633}}@{}}
\toprule\noalign{}
\begin{minipage}[b]{\linewidth}\raggedright
\hspace{0pt}Модель взаимодействия
\end{minipage} & \begin{minipage}[b]{\linewidth}\raggedright
\hspace{0pt}Основной источник информации
\end{minipage} & \begin{minipage}[b]{\linewidth}\raggedright
\hspace{0pt}На что обратить внимание при оценке
\end{minipage} \\
\midrule\noalign{}
\endhead
\bottomrule\noalign{}
\endlastfoot
\hspace{0pt}Повторная проверка собственного результата той же моделью &
\hspace{0pt}Исходные материалы и новый цикл проверки &
\hspace{0pt}Обнаружены ли упущения и улучшился ли результат после
изменений \\
\hspace{0pt}Дебаты нескольких агентов по одному вопросу &
\hspace{0pt}Собственные пути решения, предположения и аргументы каждого
агента & \hspace{0pt}Различаются ли позиции, как проверить их и выбрать
ответ \\
\hspace{0pt}Ревьюер запускает тесты и проверяет код &
\hspace{0pt}Обратная связь от компиляции, тестирования и выполнения &
\hspace{0pt}Проходят ли тесты после изменений и сохраняется ли
существующая функциональность \\
\hspace{0pt}Ревьюер просматривает результат рендеринга фронтенда или PPT
& \hspace{0pt}Скриншоты и обратная связь от взаимодействия &
\hspace{0pt}Соответствуют ли требованиям компоновка, содержание и
управление \\
\hspace{0pt}Ревьюер проверяет факты с помощью внешних инструментов &
\hspace{0pt}Результаты поиска и записи
бизнес-\allowbreak{}\hspace{0pt}систем & \hspace{0pt}Надёжны ли
источники и соответствуют ли выводы доказательствам \\
\end{longtable}

В RLEF модель с помощью обучения с подкреплением обучают использовать
обратную связь от выполнения кода и выполнять многоэтапные исправления в
задачах по соревновательному программированию. Ошибки компиляции, сбои
тестов и исключения времени выполнения дают конкретные основания для
следующего изменения. Описанная в статье модель достигает лучших
результатов при меньшем числе сгенерированных вариантов.\footnote{Gehring,
  J., et al.~\emph{RLEF: Grounding Code LLMs in Execution Feedback with
  Reinforcement Learning.} arXiv:2410.02089, 2025.} WebGen-Agent, в свою
очередь, включает в процесс генерации скриншоты веб-страниц, описания и
рекомендации визуальной модели, а также оценку операций с GUI. Это
позволяет как вернуться к предыдущему шагу и выполнить его заново, так и
выбрать лучший из нескольких сгенерированных результатов. Точность
Claude 3.5 Sonnet в бенчмарке WebGen-Bench выросла с 26,4\% до
51,9\%.\footnote{Lu, Z., et al.~\emph{WebGen-Agent: Enhancing
  Interactive Website Generation with Multi-Level Feedback and
  Step-Level Reinforcement Learning.} arXiv:2509.22644, 2025.}

Автор сначала поручает системе получить конкретные доказательства:
результаты тестов, скриншоты или данные, возвращённые инструментами.
Лишь после этого он решает, как распределить обязанности по генерации,
выполнению и проверке. Контур обратной связи может обеспечивать один
агент или несколько отдельных ролей. Добавлять роли следует только
тогда, когда это улучшает получение и использование доказательств.

Исследование безопасности Anthropic 2026 года содержит пример
совместного поиска: 45 агентов обменивались находками через общий форум
и проверяли работу друг друга, после чего независимый агент выносил
окончательное решение по результатам. Согласно отчёту, совместно
работающий кластер израсходовал 27 миллионов токенов и обнаружил 266
уязвимостей, тогда как схема с параллельно работающими независимыми
агентами израсходовала 6,5 миллиона токенов и обнаружила 21 уязвимость.
Благодаря форуму участники могли продолжать исследование чужих находок и
узнавать, какие направления уже изучаются. При этом совместно работающая
группа израсходовала примерно в четыре раза больше токенов, поэтому
разница между 266 и 21 обусловлена одновременно способом коммуникации и
объёмом вычислительных ресурсов.\footnote{Anthropic Frontier Red Team,
  ``Patterns and Problems in Emerging Multiagent Systems,'' 2026-08-13.}

Как добиться того, чтобы в математическом исследовании разные группы
действительно изучали разные пути? Трёхмерные уравнения Навье --- Стокса
описывают движение вязкой жидкости. Остаётся ли решение гладким при
гладких начальных условиях? OpenAI распределила между разными группами
две противоположные цели: одни пытались доказать гладкость, другие ---
построить контрпример. Сначала каждая группа проводила самостоятельный
поиск, а затем команда с помощью Codex обобщала промежуточные
результаты, чтобы остальные группы также могли развивать перспективные
идеи. В группе, нашедшей контрпример, параллельно работало около десяти
тысяч агентов. По данным команды, от запуска первых агентов до получения
доказательства прошло около 88 часов, а формализация и проверка в Lean
заняли ещё 17 часов. У контрпримера есть важное условие: на жидкость
действует гладкая внешняя сила, но сингулярность всё равно возникает за
конечное время. Таким образом, команда представила контрпримеры,
требуемые формулировками C и D задачи Института математики Клэя. Без
внешней силы это была бы уже другая задача. При проектировании такого
поиска можно заранее определить, какие результаты будут достаточным
ответом на исходный вопрос, поручить конкурирующие гипотезы разным
группам и отдельно выделить время на окончательную проверку.\footnote{Исследовательский
  материал OpenAI \emph{On the Navier--Stokes Millennium Prize Problem}
  содержит доказательство и код Lean. В заявлении от 11 сентября 2026
  года Институт математики Клэя отреагировал на эту работу и пояснил,
  что оценка результата и присуждение премии будут проводиться в
  соответствии с установленной процедурой.}

\textbf{Как бюджет шагов влияет на результат?} Предположим, что на
задачу выделено 30 шагов. Агенту придётся компактно распределить между
ними планирование, основную реализацию и необходимую проверку. Если же
доступно 300 шагов, появляется возможность сравнить больше вариантов,
провести более полное тестирование и внести дополнительные исправления.
Чтобы эффективно использовать эти ресурсы, агенту необходимо знать общий
и оставшийся бюджет и выбирать глубину исследования с их учётом.

В работе Google 2025 года \emph{Budget-Aware Tool-Use Enables Effective
Agent Scaling} исследуется именно этот вопрос: агент, не учитывающий
бюджет, может быстро остановиться на поверхностном поиске, и даже
увеличение доступного числа шагов не приведёт к дальнейшему улучшению
результата. Предложенный в 2026 году метод BAVT (Budget-Aware Value Tree
Search) дополнительно учитывает на каждом шаге поиска оценку ценности и
оставшийся бюджет, корректируя соотношение исследования и использования.
При достаточном бюджете можно сравнивать больше направлений, а по мере
приближения к его пределу --- сосредоточивать ресурсы на перспективных
ветвях.

В режиме управляющего это означает, что Manager должен распределять шаги
с учётом сложности задач и корректировать их число по мере работы: по
простым задачам --- как можно быстрее выдавать результат, для сложных
нужно оставлять время на проверку и исправления, а при возникновении
блокирующей проблемы --- сначала определить её причину и лишь затем
решать, нужно ли предоставить дополнительную информацию, увеличить
бюджет или выбрать другой путь.

Параллельный поиск может сократить время выполнения по часам, а
многократные проверки --- улучшить результат, но оба подхода увеличивают
число вызовов модели и расходы на коммуникацию. Автор использует хорошо
настроенного одиночного агента как базовый ориентир и расширяет
взаимодействие только тогда, когда повышение качества или экономия
времени благодаря новым участникам окупают расходы на коммуникацию,
повторные вычисления и сопровождение.

\hypertarget{ux43cux443ux43bux44cux442ux438ux430ux433ux435ux43dux442ux43dux43eux435-ux432ux437ux430ux438ux43cux43eux434ux435ux439ux441ux442ux432ux438ux435-ux441-ux43eux431ux449ux438ux43c-ux43aux43eux43dux442ux435ux43aux441ux442ux43eux43c}{%
\section{Мультиагентное взаимодействие с общим
контекстом}\label{ux43cux443ux43bux44cux442ux438ux430ux433ux435ux43dux442ux43dux43eux435-ux432ux437ux430ux438ux43cux43eux434ux435ux439ux441ux442ux432ux438ux435-ux441-ux43eux431ux449ux438ux43c-ux43aux43eux43dux442ux435ux43aux441ux442ux43eux43c}}

При взаимодействии с общим контекстом роли сменяются поэтапно, а
существующие записи о задаче сохраняются. Например, при подготовке
исследовательского отчёта роль поиска сначала собирает материалы, роль
анализа затем выполняет расчёты, а роль написания организует изложение.
Каждая роль может обращаться к предыдущим источникам и решениям, а
система с помощью текущих инструкций задачи направляет её внимание на
цели данного этапа.

При передаче работы между командами сначала задайте вопрос: кому
необходимо знать, почему изначально было принято именно такое решение? В
примере продукта Teamily AI менеджер продукта и AI-агент вместе
составляют в групповом чате документ с требованиями к продукту (PRD).
После того как продуктовая команда передаёт окончательную версию
инженерной команде, инженеры могут уточнить, какие компромиссы лежат в
основе конкретного требования.

Рассмотрим передачу задачи от роли поиска к роли написания: откуда после
передачи модель получает правила написания и какие инструменты
по-прежнему может вызывать? В эксперименте 10-1 одна и та же передача
реализуется двумя способами.

\begin{longtable}[]{@{}
  >{\raggedright\arraybackslash}p{(\columnwidth - 4\tabcolsep) * \real{0.2898}}
  >{\raggedright\arraybackslash}p{(\columnwidth - 4\tabcolsep) * \real{0.4051}}
  >{\raggedright\arraybackslash}p{(\columnwidth - 4\tabcolsep) * \real{0.3051}}@{}}
\toprule\noalign{}
\begin{minipage}[b]{\linewidth}\raggedright
\hspace{0pt}Вопрос, который необходимо решить при передаче
\end{minipage} & \begin{minipage}[b]{\linewidth}\raggedright
\hspace{0pt}Смена роли
(\texttt{transfer\_}\allowbreak{}\texttt{to\_}\allowbreak{}\texttt{agent})
\end{minipage} & \begin{minipage}[b]{\linewidth}\raggedright
\hspace{0pt}Загрузка Skill
\end{minipage} \\
\midrule\noalign{}
\endhead
\bottomrule\noalign{}
\endlastfoot
\hspace{0pt}Как предоставить правила написания? &
\hspace{0pt}Переключиться на системный промпт роли написания &
\hspace{0pt}Сохранить каталог Skills в системном промпте и добавить
правила написания в конец истории сообщений \\
\hspace{0pt}Какие инструменты видит модель? & \hspace{0pt}Можно
показывать только инструменты роли написания & \hspace{0pt}Можно
сохранить исходный набор инструментов или точку входа для их поиска \\
\hspace{0pt}Какую часть KV-\allowbreak{}\hspace{0pt}кэша
по-\allowbreak{}\hspace{0pt}прежнему можно использовать? &
\hspace{0pt}Системный промпт изменяется, поэтому начиная с точки
изменения префикса обычно требуется повторное вычисление &
\hspace{0pt}Исходный префикс запроса можно продолжать использовать
повторно \\
\hspace{0pt}Кто проверяет права на вызов? & \hspace{0pt}Harness
проверяет запрос в соответствии с текущей ролью, а затем направляет
вызов инструменту & \hspace{0pt}Harness независимо проверяет права на
каждый вызов \\
\end{longtable}

Если различия между ролями главным образом заключаются в знаниях,
процессах и стиле письма, автор отдаёт предпочтение Skill: он позволяет
повторно использовать один и тот же фреймворк выполнения и сохранять кэш
стабильного префикса запроса. Если различия между ролями связаны с
правами доступа, изоляцией инструментов или обязанностями по
согласованию, работу можно организовать с помощью отдельных агентов либо
смены ролей, а контроль доступа возложить на Harness и среду выполнения.
Наличие инструмента во входных данных модели влияет на выбор вызова,
однако фактические права определяются проверкой во время выполнения.

\begin{quote}
\textbf{Эксперимент 10-1 ★★: смена нескольких ролей при общем контексте
--- сравнение системных промптов и Skill}

Оба варианта используют одну и ту же модель, пользовательскую задачу,
реализацию инструментов, правила ролей и общую историю. Во вводной
демонстрации необходимо найти объёмы продаж автомобилей на новых
источниках энергии в Китае за 2021--2023 годы, рассчитать совокупный
среднегодовой темп роста (CAGR) и написать резюме для инвесторов объёмом
не более 120 знаков.

\textbf{Вариант 1: переключение системных промптов.} Предусмотрены пять
ролей: triage (сбор требований), research (поиск), coding
(программирование), data\_analysis (анализ) и writing (написание
текста). Точкой входа служит triage, а каждой роли доступны
специализированные инструменты и \texttt{transfer\_to\_agent}. При
передаче сохраняется история, загружаются промпт и набор инструментов
целевой роли, после чего выполнение продолжается.

\textbf{Вариант 2: загрузка Skill.} Системный промпт и полный набор
инструментов остаются неизменными на протяжении всего сеанса. Модель
вызывает \texttt{load\_skill(name)} для получения правил роли, а их
содержимое добавляется в общую историю как результат инструмента;
исходный префикс сохраняется. Harness проверяет каждый вызов инструмента
в соответствии с правилами доступа.

Обе реализации читают один и тот же файл \texttt{SKILL.md}. Возьмите по
одному запросу, отправленному модели до и после передачи роли, и
посмотрите, в каком сообщении появляются правила и какую часть префикса
по-прежнему можно использовать повторно. Затем проверьте оба резюме:
заслуживают ли доверия источники данных о продажах? Между концом 2021 и
концом 2023 года два периода роста --- правильно ли рассчитан совокупный
среднегодовой темп роста?

В формальном сравнении использовалась более строгая задача по подготовке
краткого отчёта: не более 160 знаков с сохранением двух URL-адресов
источников. Обеим группам также предоставлялось восемь шагов. Группа со
сменой ролей потратила три вызова на подсчёт знаков, однако черновик всё
равно содержал 206 знаков, а указанный в задаче инструмент вычислений
так и не был использован. Группа с загрузкой Skill рассчитала темп роста
по объёмам продаж за первый и последний годы, но лишь на восьмом шаге
загрузила правила написания. К этому моменту все необходимые роли уже
были пройдены, однако пользователь так и не получил резюме. Обе
траектории завершились из-за достижения предельного числа шагов.

На одном и том же наборе из 30 задач были запущены обе реализации, а
результаты проверены программно. Группа с загрузкой Skill прошла 15
задач, а группа с переключением системного промпта --- 2. Однако при
попарном сравнении задач у первой группы оказалось на 6 855 больше
входных токенов, не попавших в кэш, а выполнение заняло примерно на 4,4
секунды больше; в обоих случаях указана медиана попарных разностей. При
этом поставщик сервиса сообщил о нулевом числе кэшированных входных
токенов в запросах обеих групп. Обе группы прошли все шесть граничных
тестов на соблюдение инструкций, однако доля успешно выполненных полных
задач различалась значительно.
\end{quote}

\hypertarget{ux43cux443ux43bux44cux442ux438ux430ux433ux435ux43dux442ux43dux43eux435-ux432ux437ux430ux438ux43cux43eux434ux435ux439ux441ux442ux432ux438ux435-ux431ux435ux437-ux43eux431ux449ux435ux433ux43e-ux43aux43eux43dux442ux435ux43aux441ux442ux430}{%
\section{Мультиагентное взаимодействие без общего
контекста}\label{ux43cux443ux43bux44cux442ux438ux430ux433ux435ux43dux442ux43dux43eux435-ux432ux437ux430ux438ux43cux43eux434ux435ux439ux441ux442ux432ux438ux435-ux431ux435ux437-ux43eux431ux449ux435ux433ux43e-ux43aux43eux43dux442ux435ux43aux441ux442ux430}}

В архитектуре с независимыми контекстами каждый AI-агент отдельно хранит
историю входных данных, состояние задачи и цикл выполнения, а
необходимой информацией агенты обмениваются через параметры
инструментов, общие файлы и сообщения. Например, передавая материалы,
переводчик должен указать фрагмент оригинала и версию глоссария. Даже
если переводчик завершит работу, редактор сможет найти исходный текст, а
при необходимости доработки --- определить ответственного по главе.

В табл. 10-2 сопоставлены способы управления процессами в операционной
системе и агентами в среде выполнения.

Таблица 10-2. Архитектурные аналогии между мультиагентной и операционной
системами

\begin{longtable}[]{@{}
  >{\raggedright\arraybackslash}p{(\columnwidth - 2\tabcolsep) * \real{0.4074}}
  >{\raggedright\arraybackslash}p{(\columnwidth - 2\tabcolsep) * \real{0.5926}}@{}}
\toprule\noalign{}
\begin{minipage}[b]{\linewidth}\raggedright
\hspace{0pt}Операционная система
\end{minipage} & \begin{minipage}[b]{\linewidth}\raggedright
\hspace{0pt}Аналогичная функция в мультиагентной системе
\end{minipage} \\
\midrule\noalign{}
\endhead
\bottomrule\noalign{}
\endlastfoot
\hspace{0pt}Программа и конфигурация & \hspace{0pt}Системный промпт,
определения инструментов и конфигурация роли \\
\hspace{0pt}Память процесса & \hspace{0pt}Текущий контекст, состояние
задачи и история \\
\hspace{0pt}CPU выполняет вычисления & \hspace{0pt}LLM принимает
решения, а среда выполнения и инструменты выполняют соответствующие
операции \\
\hspace{0pt}Ядро управляет выполнением & \hspace{0pt}Среда выполнения
агента управляет состоянием, планированием и ресурсами \\
\hspace{0pt}Системный вызов & \hspace{0pt}Запрос внешней операции через
интерфейс инструмента \\
\hspace{0pt}fork: создание дочернего процесса &
\hspace{0pt}\texttt{spawn\_}\allowbreak{}\texttt{subagent}: создание
дочернего агента \\
\hspace{0pt}kill: отправка сигнала &
\hspace{0pt}\texttt{cancel\_}\allowbreak{}\texttt{subagent}: запрос
отмены \\
\hspace{0pt}ps: запрос списка процессов &
\hspace{0pt}\texttt{list\_}\allowbreak{}\texttt{agents}: запрос
участников и состояния их выполнения \\
\hspace{0pt}Состояние завершения и wait & \hspace{0pt}Уведомление о
завершении, состояние задачи и структурированный результат \\
\hspace{0pt}Общее хранилище и передача сообщений & \hspace{0pt}Общие
файлы, параметры инструментов и сообщения \\
\end{longtable}

Приватное состояние, асинхронные сообщения и динамическое создание
участников встречаются и в модели акторов, предложенной в 1970-х
годах.\footnote{Hewitt, C., Bishop, P., Steiger, R. \emph{A Universal
  Modular ACTOR Formalism for Artificial Intelligence.} IJCAI 1973.}

При большом количестве независимых подзадач автор предпочитает выделять
каждому участнику отдельный контекст, чтобы тот обрабатывал только
материалы, относящиеся к его обязанностям. Это упрощает раздельное
тестирование ролей, снижает влияние длинной истории предыдущих этапов на
текущие решения и позволяет выполнять несколько задач параллельно.
Эффект изоляции зависит и от других общих ресурсов: несколько агентов
по-прежнему могут одновременно изменять файлы, расходовать одну квоту
API или получать одно и то же сообщение об ошибке.

Интерфейс взаимодействия также должен учитывать потерю сообщений,
повторную доставку и задержку обновления состояния. При диагностике
проблем идентификатор задачи позволяет связать запросы, события
инструментов и результаты разных участников и восстановить процесс
передачи работы.

Предположим, один агент завершил анализ и должен передать результат
коллеге, составляющему отчёт. Небольшой объём числовых данных можно
отправить сообщением; если же результат представляет собой большую
таблицу, её обычно сохраняют в общей файловой системе, а в сообщении
передают ссылку. Так формируются две области ответственности.
\textbf{Хранилище результатов} служит для хранения документов, кода и
данных и обмена ими, а \textbf{коммуникация и управление} отвечают за
доставку сообщений, запрос состояния задач, отмену выполнения и
планирование ресурсов.

\hypertarget{ux444ux430ux439ux43bux43eux432ux430ux44f-ux441ux438ux441ux442ux435ux43cux430-ux441-ux442ux43eux447ux43aux438-ux437ux440ux435ux43dux438ux44f-ai-ux430ux433ux435ux43dux442ux430}{%
\subsection{Файловая система с точки зрения
AI-агента}\label{ux444ux430ux439ux43bux43eux432ux430ux44f-ux441ux438ux441ux442ux435ux43cux430-ux441-ux442ux43eux447ux43aux438-ux437ux440ux435ux43dux438ux44f-ai-ux430ux433ux435ux43dux442ux430}}

Один из распространённых архитектурных подходов --- предоставить агенту
\textbf{виртуальную файловую систему} (virtual filesystem), смонтировав
ресурсы из разных источников в единое дерево каталогов. Модель
обращается к ним через интерфейсы наподобие \texttt{read\_file},
\texttt{write\_file} и \texttt{list\_dir}, тогда как на нижнем уровне
могут использоваться локальный временный диск, постоянное объектное
хранилище, API стороннего облачного диска или системный пакет ресурсов.

Автор сначала чётко разграничивает области хранения и лишь затем
подключает процессы взаимодействия. Смешение временных черновиков,
официальных результатов и внешних файлов пользователя чревато случайной
перезаписью или чтением с превышением полномочий. Единый интерфейс
упрощает работу, а видимость, жизненный цикл и права каждой области
определяются настройками монтирования и контроля доступа. Это дерево
каталогов можно разделить на четыре типа областей.

\textbf{1. Личная рабочая область агента (Scratchpad).} Каждый экземпляр
использует собственную рабочую область для хранения черновиков,
временных файлов, отладочных журналов и промежуточных результатов. По
умолчанию доступ ограничивается соответствующим агентом, а данные
очищаются или архивируются согласно правилам задачи. Это предотвращает
взаимную перезапись временных файлов. Завершив работу, дочерний агент
помещает предназначенные для передачи результаты в общее пространство и
передаёт их главному агенту вместе со структурированным резюме. Записи о
пробах и ошибках остаются в исходной рабочей области для последующей
диагностики. Так изоляция контекста реализуется и на уровне хранения:
главный агент получает результат и необходимые свидетельства, не
расходуя ресурсы на чтение всего процесса проб и ошибок при повседневной
координации.

\textbf{2. Общее мультиагентное пространство (Shared Workspace).} Здесь
участники взаимодействия обмениваются результатами, а пользователь может
загружать материалы и скачивать результаты через авторизованную точку
доступа. Например, Glossary Agent записывает глоссарий, Translation
Agent читает его и выполняет перевод, после чего редактор проверяет
текст. Общее пространство обычно сохраняется на протяжении всей задачи,
поэтому необходимо учитывать параллельную запись, обновление версий и
область доступа. Обмен CSV-файлами и диаграммами между агентом и
виртуальным компьютером через общее хранилище, описанный в
\crossreflink{chap:6}{главе 6}, служит примером такой передачи
результатов. Далее будут рассмотрены оптимистическая блокировка,
изоляция рабочих копий и другие методы.

\textbf{3. Смонтированные внешние ресурсы (Mounted External Resources).}
Ресурсы Google Drive, Notion, Dropbox, корпоративных Wiki и других
систем можно отображать в единый интерфейс посредством адаптеров. Когда
агент читает документ, адаптер обращается к API исходной системы. Здесь
необходимо учитывать различия трёх типов: область доступа зависит от
внешней авторизации; сетевое обращение сопряжено с задержками,
тайм-аутами и проблемами обновления кэша; обратная запись должна
соответствовать правилам исходной системы в отношении прав доступа и
версий. Конкретная семантика согласованности определяется интерфейсом
сервиса; для обнаружения параллельных изменений можно использовать
номера версий, условную запись или повторное чтение. Монтирование по
запросу в режиме только для чтения подходит для поиска материалов, а
соответствующие права на запись следует предоставлять лишь при
необходимости редактирования.

\textbf{4. Встроенные системные ресурсы (Built-in System Resources).}
Справочники, шаблоны, определения инструментов и предустановленные
Skills могут предоставляться нескольким агентам в виде пакетов ресурсов
только для чтения. Skills загружаются по принципу прогрессивного
раскрытия: сначала предоставляется каталог, затем в соответствии с
задачей читаются регламенты и скрипты. Пакеты ресурсов с фиксированными
версиями удобны для параллельного чтения, а переключение версий при
обновлении позволяет сохранить согласованность уже выполняющихся задач.

На \crossreflink{fig:10-2}{рис. 10-2} переводчик может многократно
изменять черновик в собственном каталоге, а подготовив его к передаче
редактору --- поместить соответствующую версию в общее пространство.
Оригинал на облачном диске пользователя читается через внешний адаптер,
причём доступ к нему зависит от предоставленных пользователем
разрешений.

\begin{figure}
\hypertarget{fig:10-2}{%
\centering
\includegraphics{images/revised/fig10-2.pdf}
\caption{Рисунок 10-2. Структура монтирования четырёх типов областей
виртуальной файловой системы агента}\label{fig:10-2}
}
\end{figure}

После выхода переводчика временные файлы в его рабочей области могут
быть освобождены. Помещённый в общее пространство перевод продолжает
храниться, и редактор с правом чтения может продолжить работу над ним
(табл. 10-3).

Таблица 10-3. Четыре типа областей виртуальной файловой системы агента

\begin{longtable}[]{@{}
  >{\raggedright\arraybackslash}p{(\columnwidth - 8\tabcolsep) * \real{0.1816}}
  >{\raggedright\arraybackslash}p{(\columnwidth - 8\tabcolsep) * \real{0.1991}}
  >{\raggedright\arraybackslash}p{(\columnwidth - 8\tabcolsep) * \real{0.1929}}
  >{\raggedright\arraybackslash}p{(\columnwidth - 8\tabcolsep) * \real{0.1980}}
  >{\raggedright\arraybackslash}p{(\columnwidth - 8\tabcolsep) * \real{0.2283}}@{}}
\toprule\noalign{}
\begin{minipage}[b]{\linewidth}\raggedright
\hspace{0pt}Область
\end{minipage} & \begin{minipage}[b]{\linewidth}\raggedright
\hspace{0pt}Видимость
\end{minipage} & \begin{minipage}[b]{\linewidth}\raggedright
\hspace{0pt}Жизненный цикл
\end{minipage} & \begin{minipage}[b]{\linewidth}\raggedright
\hspace{0pt}Чтение и запись
\end{minipage} & \begin{minipage}[b]{\linewidth}\raggedright
\hspace{0pt}Управление параллельным доступом
\end{minipage} \\
\midrule\noalign{}
\endhead
\bottomrule\noalign{}
\endlastfoot
\hspace{0pt}Личная рабочая область агента & \hspace{0pt}По умолчанию
только соответствующий агент & \hspace{0pt}Очищается вместе с
экземпляром, при необходимости архивируется & \hspace{0pt}Чтение и
запись соответствующим агентом & \hspace{0pt}Предотвращение повторной
записи при повторных попытках или фоновых операциях \\
\hspace{0pt}Общее мультиагентное пространство &
\hspace{0pt}Авторизованные участники и пользователи &
\hspace{0pt}Сохраняется на протяжении задачи & \hspace{0pt}Права
предоставляются в соответствии с ролью & \hspace{0pt}Блокировка,
проверка версии или worktree \\
\hspace{0pt}Смонтированные внешние ресурсы & \hspace{0pt}Совместно
ограничивается внешней авторизацией и адаптером &
\hspace{0pt}Определяется исходной системой и политикой кэширования &
\hspace{0pt}Только чтение по запросу или запись при наличии разрешения &
\hspace{0pt}Соблюдение интерфейсов исходной системы для управления
версиями и параллельным доступом \\
\hspace{0pt}Встроенные системные ресурсы & \hspace{0pt}Авторизованные
агенты & \hspace{0pt}Фиксированная версия многократно используется между
сеансами & \hspace{0pt}Только чтение & \hspace{0pt}Централизованное
переключение версии ресурсов \\
\end{longtable}

Ссылки на файлы позволяют сохранять краткость сообщений при передаче
результатов: предыдущий участник указывает расположение результата, его
версию и назначение, а следующий читает содержимое по мере
необходимости. Участники в одной среде выполнения могут использовать
общий путь; при межорганизационном взаимодействии для такой же передачи
необходимы идентификатор результата, интерфейс доступа и авторизация,
распознаваемые обеими сторонами.

\hypertarget{ux43aux43eux43cux43cux443ux43dux438ux43aux430ux446ux438ux44f-ux438-ux443ux43fux440ux430ux432ux43bux435ux43dux438ux435-ux43cux435ux436ux434ux443-ai-ux430ux433ux435ux43dux442ux430ux43cux438}{%
\subsection{Коммуникация и управление между
AI-агентами}\label{ux43aux43eux43cux43cux443ux43dux438ux43aux430ux446ux438ux44f-ux438-ux443ux43fux440ux430ux432ux43bux435ux43dux438ux435-ux43cux435ux436ux434ux443-ai-ux430ux433ux435ux43dux442ux430ux43cux438}}

Представленные в \crossreflink{chap:4}{главе 4} инструменты создания,
коммуникации, отмены и обнаружения предоставляют базовый интерфейс
управления для мультиагентного взаимодействия: \texttt{spawn\_subagent}
создаёт участников, \texttt{send\_message\_to\_subagent} передаёт
сообщения, \texttt{cancel\_subagent} запрашивает отмену, а
\texttt{list\_agents} возвращает список участников. Помимо этих
интерфейсов, необходимо также согласовать поведение при передаче
сообщений, изменении состояния, прекращении выполнения и планировании.

\textbf{1. Передача сообщений.} Если участников немного, а связи между
ними фиксированы, можно использовать одноранговый интерфейс: например,
телефонный и компьютерный агенты могут напрямую обмениваться информацией
о ходе работы. Если каждому участнику требуется взаимодействовать со
всеми остальными, потенциальное количество связей растёт пропорционально
квадрату числа участников. При большом количестве подписок шина
сообщений упрощает централизованное управление маршрутизацией,
буферизацией и повторными попытками. Как одноранговая передача, так и
шина могут обеспечивать асинхронную доставку; необходимость
одновременного присутствия получателя в сети зависит от механизмов
очередей и постоянного хранения.

Сообщение обычно содержит структурированный \textbf{конверт} (envelope):
ID отправителя, адресата --- конкретного агента или область
широковещательной рассылки, тип сообщения и полезную нагрузку. К типам
сообщений могут относиться \texttt{task\_assigned},
\texttt{status\_update}, \texttt{result} и \texttt{terminate}.
Идентификаторы задачи, сообщения и связанного запроса помогают
получателю устранять дубликаты, а средствам отладки --- восстанавливать
цепочку коммуникации. Прикладное содержимое помещается в полезную
нагрузку JSON, поля которой обе стороны обрабатывают согласно
соглашению.

Например, за сдачу перевода отвечает редактор, который хочет попросить
агента терминологической проверки сверить один термин. Для запуска
проверки достаточно упомянуть его через \texttt{@} в комментарии к
заявке Multica и написать вопрос. Заключение можно опубликовать в ответе
к той же заявке: ответственным останется редактор, а состояние заявки не
изменится. Если проверяющий должен принять на себя всю заявку,
необходимо явно сменить ответственного.\footnote{Multica, упоминание
  агента в комментарии; реализация маршрутизации комментариев в версии
  \texttt{d579273}.}

\textbf{2. Получение состояния.} Главному агенту необходимо знать,
продолжает ли выполняться подзадача, требуются ли ей входные данные и
доступен ли результат. Для этого можно сочетать интерфейс состояния,
уведомления о событиях и файлы прогресса: интерфейс среды выполнения
предоставляет состояние жизненного цикла, сообщения описывают прикладной
прогресс, а постоянные файлы сохраняют записи, необходимые для
восстановления.

\texttt{get\_subagent\_status(agent\_id)} может возвращать такие
состояния, как «выполняется», «требует входных данных», «завершён» или
«завершился с ошибкой». В повседневной координации автор предпочитает
уведомления о завершении и сообщения о прогрессе этапов, позволяющие
главному агенту продолжать собственную работу. Активный запрос состояния
нужен при задержке уведомления, восстановлении соединения или
диагностике проблемы: это позволяет не расходовать вызовы модели на
частые вопросы «уже закончил?». Интерфейс состояния поддерживается
средой выполнения и может своевременно сообщать, жив ли процесс и
завершилась ли задача.

\textbf{Обмен сведениями о прогрессе через сообщения.} Главный агент
может запросить текущий прогресс, а дочерний --- ответить на подходящем
этапе. Дочерний агент также может самостоятельно публиковать
\texttt{status\_update}, когда получает промежуточный результат,
сталкивается с блокировкой или нуждается в помощи. В сообщении следует
указывать завершённые действия, следующий шаг и причину блокировки. Все
участники должны использовать единый набор названий состояний, чтобы
упростить их объединение и оценку. Рассматриваемый далее протокол A2A
представляет стандартизированное описание жизненного цикла задачи.

\textbf{Сохранение прогресса в файлах.} Система может последовательно
дописывать в журнал сообщения, запросы инструментов, результаты
инструментов и изменения состояния, например в формате JSONL с одним
событием на строку. Имеющий диагностические права главный агент или
эксплуатационный инструмент может по этому журналу изучать повторные
попытки, ошибки и ожидания. История входных данных модели представляет
лишь часть траектории выполнения; журнал также может содержать события
планирования и внешнего выполнения.

Полная траектория может содержать десятки тысяч токенов, а главному
агенту после каждого чтения приходится заново извлекать из неё текущее
состояние. Сообщения уведомляют главного агента о новом прогрессе, а
компактный файл прогресса сохраняет текущее состояние для последующих
запросов и восстановления. Автор предпочитает сочетать эти два механизма
для повседневной координации. После завершения каждого этапа дочерний
агент обновляет список задач, ссылки на результаты и текущие точки
блокировки, а главный агент определяет следующий шаг по этим полям.
Подробная траектория сохраняется для диагностики по мере необходимости.
Если прогресс долго не меняется, необходимость вмешательства можно
оценить по heartbeat-сигналам, состоянию вызовов инструментов и
тайм-аутам задачи: длительно работающий инструмент может всё ещё
выполнять операцию, поэтому сначала следует проверить состояние и лишь
затем решать, ждать, отменять или повторять попытку.

\textbf{3. Прекращение выполнения.} Когда несколько участников
параллельно ищут одну цель, после выполнения задачи одним из них
остальные ветви можно остановить, освободив бюджет и ресурсы
инструментов. Существуют два уровня жёсткости отмены, похожие по
действию на SIGTERM и SIGKILL в Unix.

Автор предпочитает \textbf{корректное прекращение}: среда выполнения
передаёт запрос отмены, дочерний агент в безопасной точке перестаёт
запускать новые операции, сохраняет необходимый прогресс, закрывает
сеанс браузера и освобождает блокировки, после чего возвращает
подтверждение (ack). \textbf{Принудительное прекращение} применяется,
если выполняемая задача не отреагировала до истечения срока отмены:
среда выполнения завершает процесс, а затем проверяет оставшиеся ресурсы
и незавершённые операции записи. Для операций, уже отправленных во
внешнюю систему, необходимо дополнительно запросить результат и при
необходимости выполнить компенсирующую процедуру.

Жизненным циклом также необходимо управлять вдоль связей создания. По
аналогии с Go context сигнал отмены родительской задачи можно
распространять вниз по всему дереву подзадач, а цикл выполнения должен
проверять его в безопасных точках, что соответствует назначению
\texttt{ctx.Done()}. Тогда после завершения родительской задачи
запущенные ею ветви также перейдут к освобождению ресурсов, а бесхозные
агенты не продолжат вызывать инструменты и расходовать бюджет. Для
фоновой задачи, которой действительно требуется длительная независимая
работа, создаётся отдельный контекст жизненного цикла, подобный
\texttt{context.Background()}, и назначается ответственный за
мониторинг, бюджет и освобождение ресурсов.

\textbf{4. Ресурсы и планирование.} Агентной системе необходимо
распределять токены, шаги, квоты параллельного выполнения, подключения к
инструментам и вычислительные ресурсы. При запуске подзадачи задаются
бюджет и поведение при его превышении; модель выбирается с учётом
сложности задачи; количество параллельных вызовов ограничивается, чтобы
не исчерпать квоту сервиса. При поступлении более срочной задачи
низкоприоритетную работу можно приостановить или отменить в безопасной
точке, сохранить её прогресс и запланировать заново.

Предположим, три ветви исследуют по одному возможному решению, причём
одна уже нашла ключевую зацепку и ей осталась только последняя проверка.
Manager может выделить ей ещё несколько шагов и одновременно завершить
две ветви, не добившиеся прогресса. Квота списывается средой выполнения:
даже если ветвь хочет предпринять ещё одну попытку, после исчерпания
бюджета она должна остановиться.

\hypertarget{ux43eux434ux43dux43eux440ux430ux43dux433ux43eux432ux43eux435-ux432ux437ux430ux438ux43cux43eux434ux435ux439ux441ux442ux432ux438ux435-ux432ux437ux430ux438ux43cux43dux430ux44f-ux43fux440ux43eux432ux435ux440ux43aux430-ux438-ux438ux442ux435ux440ux430ux442ux438ux432ux43dux43eux435-ux443ux43bux443ux447ux448ux435ux43dux438ux435}{%
\subsection{Одноранговое взаимодействие: взаимная проверка и итеративное
улучшение}\label{ux43eux434ux43dux43eux440ux430ux43dux433ux43eux432ux43eux435-ux432ux437ux430ux438ux43cux43eux434ux435ux439ux441ux442ux432ux438ux435-ux432ux437ux430ux438ux43cux43dux430ux44f-ux43fux440ux43eux432ux435ux440ux43aux430-ux438-ux438ux442ux435ux440ux430ux442ux438ux432ux43dux43eux435-ux443ux43bux443ux447ux448ux435ux43dux438ux435}}

В одноранговом взаимодействии обычно участвуют два-три AI-агента,
которые обмениваются мнениями по одной задаче и дорабатывают результат.
При проектировании можно заранее назначить участникам разные
обязанности, материалы и инструменты, чтобы они выносили независимые
суждения, а затем объединяли результаты. Например, генератор
сосредоточен на реализации, а проверяющий --- на тестах и требованиях
пользователя; специалисты разных профилей проверяют ограничения в своих
областях. Чем сильнее сходство моделей, входных данных и рабочих
процессов, тем вероятнее, что участники повторят одно и то же суждение,
поэтому различия в точках зрения и свидетельствах необходимо
проектировать намеренно.\footnote{Anthropic Frontier Red Team,
  ``Patterns and Problems in Emerging Multiagent Systems,'' 2026-08-13.}

Для задач, требующих многократного улучшения одного результата, автор
рекомендует начинать с однорангового взаимодействия двух-трёх
участников. Когда роли, содержимое передачи и условия остановки
определены, можно построить тестируемый цикл с относительно небольшими
первоначальными затратами на координацию.

\hypertarget{ux438ux43dux436ux435ux43dux435ux440ux438ux44f-ux446ux438ux43aux43bux43eux432}{%
\subsubsection{Инженерия
циклов}\label{ux438ux43dux436ux435ux43dux435ux440ux438ux44f-ux446ux438ux43aux43bux43eux432}}

Одноранговое взаимодействие помогает справиться с распространённой
проблемой: \textbf{преждевременным прекращением}. В программировании и
выполнении поручений через службу поддержки его можно обнаружить по трём
видам записей.

\begin{itemize}
\tightlist
\item
  \textbf{Объявление всей работы завершённой после выполнения лишь
  части.} Coding Agent пишет код, но не выполняет тесты и не проверяет
  развёртывание; пользователь поручает помощнику два дела, а тот
  выполняет первое и докладывает, что всё готово. Необходимо по пунктам
  сверить работу со списком задач и проверить свидетельства выполнения
  каждого пункта.
\item
  \textbf{Преждевременный отказ после неудачи на одном пути.} Продавец
  может потребовать отправить материалы через форму или получить
  вложения по указанному адресу электронной почты. Получив такой ответ,
  можно продолжить работу способом, который принимает продавец.
  Необходимо записать причину неудачи и определить, какие пути ещё стоит
  попробовать.
\item
  \textbf{Пропуск последующих шагов после локального успеха.} Если после
  согласия продавца на возврат средств пользователь должен дополнительно
  подтвердить его в приложении, помощник обязан сообщить об этом шаге и
  условиях его выполнения. Например, если при обсуждении уменьшения
  суммы счёта служба поддержки устно согласилась на это по телефону,
  помощник может подтвердить пользователю, что в ходе переговоров
  получено соответствующее обещание. Запись разговора доказывает, что
  именно сказал сотрудник поддержки, но для подтверждения изменения
  следующего счёта потребуются новые свидетельства. Проверяющий должен
  отдельно указать достигнутый результат, кто и что ещё должен сделать,
  а также какие результаты пока невозможно наблюдать.
\end{itemize}

Во всех трёх случаях условия завершения необходимо свести к проверяемым
результатам. \textbf{Инженерия циклов} рассматривает, как заставить
агента непрерывно проходить цикл «определить следующую работу ---
выполнить --- проверить --- записать прогресс» и на основании результата
решать, следует ли продолжать, ждать, передать работу или завершить её.
Инженеру необходимо спроектировать источник задач, способ проверки,
постоянное состояние и условия остановки. Роль человека также
расширяется: вместо написания промпта для одного вызова он проектирует
цикл, способный непрерывно продвигать задачу. Addy Osmani использовал
это название для данной совокупности практик в июне 2026 года и сослался
на опыт Boris Cherny: основной акцент работы постепенно смещается к
проектированию циклов, управляющих агентами.\footnote{Osmani, Addy.
  ``Loop Engineering'', 2026.}

\textbf{Автор считает, что в таких циклах ключевым узким местом обычно
становится валидатор.} При ненадёжной проверке увеличение числа раундов
может лишь заставить систему повторять однотипные ошибки или
преждевременно объявлять работу завершённой. Результаты тестов,
состояние бизнес-процесса и подтверждение пользователя предоставляют
основания для завершения каждый в своей области; провал проверки
становится входными данными для исправлений в следующем раунде.

Практика предшествовала названию. Ещё до распространения термина
«инженерия циклов» автор и его команда уже использовали в Pine AI циклы
и проверку для устранения преждевременного прекращения. «Предлагающий
--- проверяющий» --- распространённый способ организации таких циклов,
который подробнее рассматривается ниже.

\textbf{LoopX: сохранение рабочего состояния между раундами.} LoopX
помещает цель, границы, список задач, шлюзы, свидетельства и квоты в
постоянную плоскость управления, которая назначает работу для следующего
раунда или исполнителя. Процесс можно свести к следующей
последовательности: \textbf{решение плоскости управления → выполнение
агентом → независимая проверка → фиксация проверенного прогресса}.

Агент вызывает инструменты, рассуждает и передаёт результат текущего
раунда. Плоскость управления записывает уже выполненную работу и
определяет, откуда продолжить в следующем раунде. После успешной
проверки состояние завершения фиксируется; при неудаче начинается
исправление или перепланирование. Контрольные точки с участием человека,
условия ожидания и пределы бюджета совместно определяют, когда
выполнение можно продолжить.\footnote{LoopX, The local control plane for
  long-running AI agent work, v0.4.0.}

\textbf{LongHorizon-Harness: связывание длительных задач между
приложениями.} Этот фреймворк исходит из мультимодального Computer Use и
обеспечивает непрерывное выполнение задачи через GUI, CLI, несколько
настольных приложений и многократные обновления контекста. Он использует
цикл Manage--Execute--Audit (MEA): Manager на основе исходной цели,
проверенного прогресса, свидетельств неудач и оставшейся работы создаёт
одну подзадачу с чётко очерченными границами; Executor выполняет
операции через GUI или CLI в новом контексте; Auditor только для чтения
проверяет результат и записывает подтверждённый прогресс в состояние
следующего раунда. Свидетельства неудач также сохраняются для
восстановления и перепланирования.

Фреймворк вызывает через слой адаптеров такие исполнительные бэкенды,
как Claude Code и Codex CLI, связывая внешнее управление задачами с
внутренними циклами выполнения бэкендов. Этот подход основывает
непрерывность длительной задачи на проверенном состоянии задачи.
Контекст выполнения может обновиться, но следующий раунд всё равно
продолжит работу на основе цели, свидетельств завершения и нерешённых
вопросов, повторно проверяя, не повлияли ли изменения среды на прежние
выводы.\footnote{LongHorizon-Harness, статья arXiv:2608.01964. Открытые
  траектории базовой системы и MEA относятся к задаче
  WEB\_task\_16\_webrtc\_simulcast\_layer\_audit; скрипт воспроизведения
  использует коммит \texttt{53bc678}.}

В контролируемом сравнении, где Qwen 3.7-Plus и исполнительный бэкенд
Claude Code оставались неизменными, а менялся внешний цикл, статья
сообщает о росте PassRate WeaveBench с 51,8 до 80,7 \%, бинарной доли
завершения OSWorld 2.0 --- с 2,8 до 8,3 \%, а успешности Terminal-Bench
2.1 --- с 69,7 до 77,2 \%. Затраты также зависели от задачи: первые два
бенчмарка потребовали соответственно в 2,3 раза больше общего числа
токенов и в 3,6 раза больше выходных токенов относительно базовой
системы, тогда как Terminal-Bench 2.1 израсходовал на 24 \% меньше
токенов.

В одной и той же задаче многоуровневого аудита потоков WebRTC базовый
исполнитель застрял в интерфейсе Wireshark, затем предпринял ещё
несколько попыток и в итоге получил оценку 0,59. При использовании MEA
проверяющий записал недостающие свидетельства, следующий исполнитель
выполнил дополнительные проверки на их основе, и итоговая оценка
составила 0,92.

\hypertarget{ux43fux430ux440ux430ux434ux438ux433ux43cux430-ux43fux440ux435ux434ux43bux430ux433ux430ux44eux449ux438ux439-ux43fux440ux43eux432ux435ux440ux44fux44eux449ux438ux439}{%
\subsubsection{Парадигма «предлагающий ---
проверяющий»}\label{ux43fux430ux440ux430ux434ux438ux433ux43cux430-ux43fux440ux435ux434ux43bux430ux433ux430ux44eux449ux438ux439-ux43fux440ux43eux432ux435ux440ux44fux44eux449ux438ux439}}

Парадигма «предлагающий --- проверяющий» передаёт создание результата и
его проверку двум ролям. Такая архитектура использовалась в описанных в
\crossreflink{chap:5}{главе 5} экспериментах по созданию PPT,
редактированию видео и визуализации журналов: Proposer генерирует код,
Reviewer просматривает визуализированный результат, проверяет его
качество с помощью визуальной модели и возвращает структурированные
рекомендации по исправлению. На \crossreflink{fig:10-3}{рис. 10-3}
рекомендации также возвращаются предлагающему. Сможет ли он на их основе
исправить следующую версию, зависит от того, насколько точно проверяющий
локализует проблему.

\begin{figure}
\hypertarget{fig:10-3}{%
\centering
\includegraphics{images/revised/fig10-3.pdf}
\caption{Рисунок 10-3. Цикл «предлагающий ---
проверяющий»}\label{fig:10-3}
}
\end{figure}

При изменении объекта проверки должны меняться и изучаемые
свидетельства. При проверке ответа службы поддержки его следует
сопоставлять с бизнес-правилами и источниками фактов; при проверке
изменений кода --- изучать тесты и результаты фактического выполнения;
при проверке плана действий --- сверять применимые нормы и
предварительные условия.

В статье Хуана и соавторов \emph{Large Language Models Cannot
Self-Correct Reasoning Yet}, представленной на ICLR 2024, обнаружено,
что в исследованных задачах рассуждения и конфигурациях промптов
самостоятельное исправление GPT-4 без внешней обратной связи приводило к
изменению правильных ответов на неправильные и снижению общей точности.
Обзор самоисправления в TACL дополнительно рассматривает влияние условий
задачи и надёжности обратной связи на результат исправления.\footnote{Kamoi,
  R., et al.~\emph{When Can LLMs Actually Correct Their Own Mistakes? A
  Critical Survey of Self-Correction of LLMs.} TACL, 2024.} Поэтому
автор считает независимые свидетельства и локализуемые условия
исправления основными требованиями к циклу проверки. Проверяющий должен
иметь возможность прочитать фактические результаты тестов, состояние
среды или визуализированный результат и указать, какое требование ещё не
выполнено и как проверить исправление.

Полный цикл можно организовать в пять этапов:

\begin{enumerate}
\def\labelenumi{\arabic{enumi}.}
\tightlist
\item
  Предлагающий создаёт первую версию в соответствии с требованиями
  задачи, например ответ или фрагмент кода.
\item
  Результат выполняется или визуализируется для получения результатов
  тестов, состояния бизнес-процесса, снимка экрана или результатов
  проверки фактов.
\item
  Проверяющий сопоставляет требования со свидетельствами и либо
  принимает результат, либо указывает место проблемы и условия
  исправления.
\item
  Если бюджет ещё не исчерпан, предлагающий вносит исправления на основе
  обратной связи, после чего результат снова выполняется и проверяется;
  если возникает вопрос, требующий входных данных или разрешения, работа
  переводится в ожидание или передаётся.
\item
  После выполнения условий публикации передаётся версия, прошедшая
  проверку, вместе со свидетельствами проверки и заключением
  проверяющего; при исчерпании бюджета сохраняются прогресс и нерешённые
  вопросы.
\end{enumerate}

Проверяющий читает защищённые тесты и критерии публикации, а
предложенные изменения этих критериев проверяются отдельно.

В экспериментах CRITIC модель сначала формирует ответ, а затем проверяет
его с помощью инструментов.\footnote{CRITIC, arXiv:2305.11738.}
Например, она использует поисковую систему для проверки фактов в ответе
и интерпретатор Python для выполнения самостоятельно написанной
математической программы. После исправления ответов на основе обратной
связи от инструментов показатели улучшились во всех трёх категориях
задач: ответах на открытые вопросы, генерации математических программ и
уменьшении вредных формулировок. Исправление работало хуже, когда модель
оценивала исходный ответ самостоятельно. Эта архитектура соответствует
примеру с PPT из \crossreflink{chap:5}{главы 5}: визуализированный
снимок экрана выявляет проблемы с компоновкой и перекрытием, которые
трудно определить непосредственно по тексту кода, а выполнение тестов
обнаруживает ошибки фактического выполнения.

В примере Anthropic 2026 года по долгосрочной разработке приложений
использовались три агента: планировщик, генератор и оценщик. Планировщик
преобразует требования в спецификацию продукта; генератор и оценщик
предварительно согласуют критерии завершения каждого раунда, после чего
соответственно выполняют реализацию и проверку. Оценщик работает с
приложением через Playwright и формирует отчёт о дефектах, а роли
передают состояние через файлы. Этот пример показывает, как связать
спецификацию продукта, реальные операции и обратную связь о дефектах, но
при сравнении подходов необходимо учитывать дополнительные затраты на
выполнение и проверку.\footnote{Prithvi Rajasekaran, ``Harness Design
  for Long-Running Application Development,'' Anthropic Engineering,
  2026-03-24.}

20 ноября 2026 года математик Levent Alpöge опубликовал полученное
совместно с Claude формальное доказательство, касающееся гипотезы Римана
и её обобщения для функций Дирихле \(L\). Получив большой набор файлов
Lean, проверяющий прежде всего должен задать два вопроса: сформулирована
ли в последней теореме исходная гипотеза и может ли ядро проверки
доказательств заново проверить каждый шаг вывода? Математик MIT Andrew
Sutherland сначала сопоставил представленную теорему с определением
гипотезы Римана в математической библиотеке Lean, а затем перечислил
аксиомы, от которых зависит доказательство. В опубликованной в тот же
день записи он также сообщил, что независимое ядро проверки
доказательств nanoda успешно проверило экспортированное доказательство.
Это даёт конкретный способ разделить обязанности при проверке:
специалист по проблеме сначала формулирует утверждение для приёмки, а
независимый инструмент затем проверяет, позволяет ли представленный
результат вывести это утверждение. Чтобы в будущем математики могли
понять, как автор пришёл к этому подходу, также потребуется удобное для
чтения человеком обоснование.\footnote{Levent Alpöge опубликовал
  описание исследования 20 ноября 2026 года. Проверяемые пункты и
  состояние проверки на тот день приведены в документе Andrew Sutherland
  \emph{Verification files for the Lean proof of the Riemann
  Hypothesis}.}

\hypertarget{ux434ux438ux441ux43aux443ux441ux441ux438ux43eux43dux43dux430ux44f-ux43cux43eux434ux435ux43bux44c}{%
\subsubsection{Дискуссионная
модель}\label{ux434ux438ux441ux43aux443ux441ux441ux438ux43eux43dux43dux430ux44f-ux43cux43eux434ux435ux43bux44c}}

Дискуссионная модель позволяет участникам проверять один вопрос с разных
позиций. Например, при оценке технического решения один агент описывает
преимущества и подходящие сценарии, а другой проверяет ограничения,
условия отказа и альтернативы. В каждом раунде участники приводят
дополнительные свидетельства по конкретным аргументам или выдвигают
возражения, а в конце подводят итог: перечисляют пункты, по которым
достигнут консенсус, и разногласия, которые ещё необходимо проверить.

Такой подход помогает ясно перечислить аргументы за и против, но его
эффективность зависит от качества аргументов, независимости суждений и
способа окончательного выбора. Тран и Киела сравнили на задачах
многошагового рассуждения пять архитектур: последовательную,
дискуссионную, ансамблевую, с параллельными ролями и с параллельными
подзадачами. При одинаковом бюджете токенов для рассуждений один агент
работал не хуже или даже лучше нескольких агентов.\footnote{Tran, D.,
  Kiela, D. \emph{Single-Agent LLMs Outperform Multi-Agent Systems on
  Multi-Hop Reasoning Under Equal Thinking Token Budgets.}
  arXiv:2604.02460, 2026.} Поэтому автор использует дискуссию
преимущественно для раскрытия разногласий, проверки предположений и
дополнения противоположных аргументов. Если цель заключается в
фактической или вычислительной точности, приоритет следует отдавать
проверяемым свидетельствам и независимой проверке.

В статье ограничения передачи промежуточной информации рассматриваются с
помощью неравенства обработки данных: при многократном пересказе и
сжатии исходные детали могут быть утрачены. На практике можно сохранять
прослеживаемые источники, требовать от участников сначала проводить
независимый анализ и лишь затем обмениваться выводами, уменьшая
вероятность непроверенного консенсуса. С дискуссией также можно сочетать
несколько независимых выборок с последующей агрегацией и назначение
отдельного проверяющего для задач, в которых «генерировать трудно, а
проверять легко».

\hypertarget{ux43cux43eux434ux435ux43bux44c-ux43cux43eux437ux433ux43eux432ux43eux433ux43e-ux448ux442ux443ux440ux43cux430}{%
\subsubsection{Модель мозгового
штурма}\label{ux43cux43eux434ux435ux43bux44c-ux43cux43eux437ux433ux43eux432ux43eux433ux43e-ux448ux442ux443ux440ux43cux430}}

Мозговой штурм можно разделить на два раунда: сначала каждый участник
самостоятельно придумывает решения, затем они вместе обсуждают, какие
идеи можно объединить. Например, агент 1 предлагает добавить функцию
публикации в социальных сетях; агент 2 на её основе предлагает
автоматически создавать персонализированные постеры для публикации;
агент 3 развивает идею до пользовательских шаблонов и рынка шаблонов.

Чтобы участники предлагали разные решения, можно попросить их
рассматривать задачу соответственно с точки зрения пользователя и
оператора либо предоставить им разные промпты или модели. После обмена
идеями решения отбираются на основе требований пользователя, стоимости
реализации и результатов проверки.

\hypertarget{ux43cux43eux434ux435ux43bux44c-ux44dux43aux441ux43fux435ux440ux442ux43dux43eux439-ux433ux440ux443ux43fux43fux44b}{%
\subsubsection{Модель экспертной
группы}\label{ux43cux43eux434ux435ux43bux44c-ux44dux43aux441ux43fux435ux440ux442ux43dux43eux439-ux433ux440ux443ux43fux43fux44b}}

Экспертная группа распределяет профессиональные обязанности при работе
над междисциплинарной проблемой. Например, при оценке нового продукта
агент-инженер анализирует сложность реализации и технические
зависимости, продуктовый агент проверяет потребности пользователей и
пользовательский опыт, а операционный агент оценивает затраты, ресурсы и
условия внедрения.

В групповом чате LobeHub ведущий может разослать один и тот же проект
продукта этим трём экспертам. Система сначала извлекает текущую историю
группового чата, а затем просит трёх участников одновременно ответить в
соответствии с настройками своих ролей. Все три ответа прикрепляются к
сообщению, из которого была запущена рассылка, поэтому впоследствии
можно проверить, на какой вопрос отвечал каждый участник. Если
операционному специалисту необходимо рассчитать затраты на основе только
что оценённого инженером объёма использования, следует дождаться записи
ответов текущего раунда и лишь затем попросить операционного специалиста
выступить: если они начнут писать одновременно, у него ещё не будет этой
цифры.\footnote{LobeHub, руководство по использованию группового чата;
  реализация параллельной рассылки в версии \texttt{ab21e0d}.}

При объединении результатов необходимо выявить связи между суждениями:
какая техническая поддержка требуется для конкретной цели
пользовательского опыта, насколько технический выбор увеличит
операционные затраты и какие условия ещё требуют внешних свидетельств.

\hypertarget{ux440ux435ux436ux438ux43c-ux443ux43fux440ux430ux432ux43bux44fux44eux449ux435ux433ux43e-ux446ux435ux43dux442ux440ux430ux43bux438ux437ux43eux432ux430ux43dux43dux430ux44f-ux43aux43eux43eux440ux434ux438ux43dux430ux446ux438ux44f}{%
\subsection{Режим управляющего: централизованная
координация}\label{ux440ux435ux436ux438ux43c-ux443ux43fux440ux430ux432ux43bux44fux44eux449ux435ux433ux43e-ux446ux435ux43dux442ux440ux430ux43bux438ux437ux43eux432ux430ux43dux43dux430ux44f-ux43aux43eux43eux440ux434ux438ux43dux430ux446ux438ux44f}}

Если задача включает много подзадач, сложные зависимости или
динамическое распределение задач, единое управление можно поручить
управляющему AI-агенту. Сначала он осмысливает общую цель, затем
распределяет работу, выбирает исполнителей и отслеживает ход выполнения.
При возникновении исключительных ситуаций он решает, следует ли
повторить попытку, сменить исполнителя или скорректировать план, а в
конце объединяет результаты.

Один из распространённых подходов к реализации --- зарегистрировать
специализированных агентов как инструменты. Управляющий передаёт через
интерфейс инструмента задачу, ограничения и необходимые материалы, а в
ответ получает структурированный результат или идентификатор задачи.
Благодаря этому поиск, операции с файлами и делегирование подзадач могут
использовать схожий механизм вызова, тогда как для длительных подзадач
отдельно организуется управление ходом выполнения, отменой и повторными
попытками. Разные исполнители могут использовать разные модели, промпты,
инструменты и среды выполнения. При добавлении нового исполнителя
недостаточно зарегистрировать интерфейс: управляющий также должен
понимать его обязанности, требования к входным данным и подходящие
сценарии применения.

При выполнении длительных задач в Claude Code и Codex автор сначала вёл
обсуждения, выполнял команды и запускал тесты в основном сеансе. Если в
процессе хотелось уточнить какое-либо проектное решение, часто
приходилось ждать завершения текущего этапа; а чтобы спустя несколько
часов найти первоначальную договорённость, требовалось пролистать
длинную череду результатов команд. Позднее автор стал поручать изменения
и тестирование дочерним агентам, оставляя основной сеанс для обсуждений.
Например, при вычитке рукописи исполнитель может продолжать сверять
номера иллюстраций, пока автор обсуждает с читателем, нужно ли
переписать определённый фрагмент. Если обсуждение изменило область
задачи, новую договорённость необходимо передать исполнителю, а при
сдаче работы проверить результат с её учётом.

Команда DeepSeek Harness опубликовала исследовательский прототип
\texttt{dsh-superagent}, предназначенный для работы с одним репозиторием
GitHub. Рассмотрим исправление формата даты в экспортируемом файле:
сопровождающий описывает требование в Issue, после чего маршрутизатор
(Router) по правилам репозитория определяет, можно ли принять задачу.
Если задача принята, система создаёт рабочую заявку (Ticket) и выделяет
исполнителю (Worker) отдельный сеанс и рабочий каталог. Исполнитель
может изменять код и запускать тесты в собственном каталоге; параллельно
может выполняться другая заявка, причём для каждой сохраняются история
изменений и результаты тестов.\footnote{Команда DeepSeek Harness,
  исследовательский прототип dsh-superagent. Получение последующих
  сообщений рабочими заявками описано в разделах об авторизации и
  подписках, а сеансы выполнения и рабочие каталоги --- в описании
  механизма выполнения.}

Предположим, затем сопровождающий дополняет исходный Issue требованием:
«Изменить только CSV, сохранив формат даты в JSON». Активная заявка,
подписанная на это обсуждение, получит дополнительное сообщение, и
исполнитель сможет соответственно сузить область изменений и отдельно
протестировать оба формата. Завершив работу, он может открыть запрос на
включение изменений (PR) и опубликовать ссылку и результаты в исходном
Issue. В одном месте сопровождающий увидит первоначальные требования,
добавленные позднее ограничения и изменения, ожидающие проверки;
маршрутизатору при этом не придётся заново преобразовывать каждое
последующее сообщение в отдельную задачу.

На управляющего возлагается значительная часть ответственности за
принятие решений. Ошибка при декомпозиции задачи повлияет на последующее
выполнение, недостаток материалов при передаче приведёт к повторным
вопросам, а чрезмерный объём журналов вызовов займёт контекст. При
проектировании следует сохранять краткое состояние задачи и индекс
артефактов, загружать подробности по мере необходимости и задавать для
каждой подзадачи чёткие входные данные, результаты и критерии
завершения.

В работе Plan-and-Act исследовалось влияние качества планирования на
длительные задачи в архитектуре Planner---Executor; на WebArena-Lite
была заявлена доля успеха 57,58\%.\footnote{Erdogan, L. E., et
  al.~\emph{Plan-and-Act: Improving Planning of Agents for Long-Horizon
  Tasks.} arXiv:2503.09572, 2025.} Из-за ошибки планирования несколько
исполнителей могут продолжить работу в рамках неверной декомпозиции, а
увеличение их числа лишь умножит объём повторной работы. Поэтому автор в
первую очередь проверяет качество декомпозиции задач управляющим. Если
узким местом является планирование, управляющему следует выделить более
мощную модель, более полные сведения о задаче и больший бюджет
рассуждений. После чёткого определения границ подзадач можно выбирать
подходящую конфигурацию для исполнителей.

Для параллельных задач также необходимо явно определить порядок
завершения. Если цель состоит в поиске одного подходящего результата,
после успешной проверки можно атомарно зарегистрировать первый прошедший
проверку результат, отменить остальные ветви, дождаться подтверждения
отмены или истечения времени ожидания, а затем подготовить итоговую
выдачу. Операция регистрации должна быть одновременно атомарной и
идемпотентной: если два результата поступят одновременно, состояние
изменится только один раз; если одно событие будет доставлено повторно,
следует использовать уже полученный результат обработки. Этот этап можно
защитить блокировкой, транзакцией или операцией сравнения с обменом.

Конкретную передачу результата можно продемонстрировать с помощью
встроенной в Python базы данных SQLite. Предположим, два исполнителя, А
и Б, параллельно проверяют материалы, а результаты каждого уже прошли
валидацию. Управляющему нужен только один подходящий результат, поэтому
в одной строке задачи он записывает, кто отправил его первым. В
следующем примере сначала результат отправляет А, затем Б, после чего А
повторяет попытку.

\begin{Shaded}
\begin{Highlighting}[]
\ImportTok{import}\NormalTok{ sqlite3}

\NormalTok{db }\OperatorTok{=}\NormalTok{ sqlite3.}\ExtensionTok{connect}\NormalTok{(}\StringTok{":memory:"}\NormalTok{)}
\NormalTok{db.execute(}\StringTok{"CREATE TABLE task (id INTEGER PRIMARY KEY, winner TEXT)"}\NormalTok{)}
\NormalTok{db.execute(}\StringTok{"INSERT INTO task VALUES (1, NULL)"}\NormalTok{)}
\ControlFlowTok{for}\NormalTok{ worker }\KeywordTok{in}\NormalTok{ [}\StringTok{"А"}\NormalTok{, }\StringTok{"Б"}\NormalTok{, }\StringTok{"А"}\NormalTok{]:}
    \ControlFlowTok{with}\NormalTok{ db:}
\NormalTok{        claimed }\OperatorTok{=}\NormalTok{ db.execute(}
            \StringTok{"UPDATE task SET winner = ? WHERE id = ? AND winner IS NULL"}\NormalTok{,}
\NormalTok{            (worker, }\DecValTok{1}\NormalTok{),}
\NormalTok{        ).rowcount}
    \BuiltInTok{print}\NormalTok{(worker, }\StringTok{"регистрация выполнена"} \ControlFlowTok{if}\NormalTok{ claimed }\ControlFlowTok{else} \StringTok{"используется уже зарегистрированный результат"}\NormalTok{)}
\end{Highlighting}
\end{Shaded}

Строки выводятся в следующем порядке: «А регистрация выполнена», «Б
используется уже зарегистрированный результат», «А используется уже
зарегистрированный результат». Если первым окажется Б, в таблицу будет
записан Б. Одна инструкция обновления одновременно проверяет, что
результат ещё никто не зарегистрировал, и записывает имя отправителя.
Если разделить проверку и запись, оба исполнителя могут успеть прочитать
пустое значение.

В предыдущем примере таблица создана в памяти, поэтому после завершения
процесса регистрация исчезнет. При использовании постоянной базы данных
возникает ещё одна проблема передачи: управляющий уже зарегистрировал
результат А, но завершил работу, не успев уведомить Б. Как после
перезапуска определить, что это уведомление ещё не было отправлено?
Выбранный результат и ожидающее отправки уведомление можно записать в
одной транзакции. Для каждого уведомления сохраняются ветвь-получатель,
требуемое действие и уникальный номер. После перезапуска управляющий
повторно отправляет уведомления, получение которых ещё не подтверждено.
Даже если Б уже завершил очистку и потерялось только подтверждение, по
исходному номеру он распознает повторно отправленное уведомление и сразу
вернёт прежний результат обработки. Таким же образом можно
регистрировать события передачи результата.

Если задача требует собрать все результаты, необходимо дождаться
завершения всех обязательных ветвей и отдельно зарегистрировать
успешные, неудачные и незавершённые элементы. Условие завершения
определяется целью задачи и одновременно задаёт момент, когда можно
отменить работу остальных участников.

\textbf{Последовательная координация.} Если между подзадачами существуют
последовательные зависимости, управляющий может назначить следующий этап
только после того, как станет доступен результат предыдущего. На
\crossreflink{fig:10-4}{рисунке 10-4} в качестве примера используется
подготовка отчёта: сначала участник, собирающий материалы, возвращает
источники, затем аналитик формулирует на их основе выводы, а в конце
автор получает подкреплённый источниками анализ. Если заменить этот
процесс следующей задачей перевода, каких материалов должен дождаться
переводчик перед началом работы?

\begin{figure}
\hypertarget{fig:10-4}{%
\centering
\includegraphics{images/revised/fig10-4.pdf}
\caption{Рисунок 10-4. Последовательная координация
управляющего}\label{fig:10-4}
}
\end{figure}

\begin{quote}
\textbf{Эксперимент 10-2 ★★: агенты для перевода книги}

При переводе технической книги необходимо одновременно учитывать
терминологию, контекст и стиль изложения. Например, английское слово
agent можно перевести по-разному, но если в переводе принято
использовать вариант «интеллектуальный агент», в последующих главах
следует придерживаться того же варианта. Для глоссария, исходных текстов
глав, переводов и записей редакторской проверки нужна чёткая система
организации. Длинная книга вместе с промежуточными переводами,
результатами инструментов и историей проверок будет занимать всё больше
контекста: к девятой главе модель может забыть вариант перевода из
второй главы или несколько раз проверять один и тот же фрагмент.
Разделение задачи по обязанностям и главам позволяет каждому экземпляру
сосредоточиться на материалах, необходимых для его текущей роли.

Работу можно распределить между четырьмя типами агентов:

\begin{itemize}
\tightlist
\item
  \textbf{Агент глоссария: систематизация терминов.} Он читает книгу по
  частям, выявляет регулярно встречающиеся профессиональные термины,
  сверяется с отраслевыми словарями и правилами перевода, после чего
  создаёт таблицу соответствий в формате JSON или CSV с английским
  термином, китайским переводом, частью речи и контекстом употребления.
  Глоссарий сохраняется в общем пространстве, а после завершения задачи
  экземпляр освобождается.
\item
  \textbf{Агент перевода: перевод глав.} Он читает текущую главу,
  относящиеся к ней термины и руководство по переводу, после чего
  готовит перевод с учётом целевой аудитории и языкового стиля. Для
  новых терминов он предлагает варианты перевода и помечает их как
  требующие проверки. Каждый экземпляр поддерживает отдельный контекст,
  а перевод сохраняется в общем пространстве; разные главы можно
  обрабатывать последовательно или параллельно в зависимости от их
  связей.
\item
  \textbf{Агент вычитки: проверка всего текста.} По глоссарию и
  переводам глав он проверяет единообразие терминологии и формулировок,
  а также удобочитаемость. Весь текст проверяется посредством чтения по
  главам и межглавного индекса, а расположение проблем, их обоснование и
  требования к исправлению заносятся в отчёт о вычитке.
\item
  \textbf{Агент-координатор: назначение задач и сбор результатов.} Он
  поддерживает цель, план, состояние задач, сводки вызовов и индекс
  артефактов, на основании отчёта о вычитке возвращает переводчикам
  главы, требующие исправления, а затем организует повторную проверку.
\end{itemize}

На \crossreflink{fig:10-5}{рисунке 10-5} редактор читает переводы глав и
глоссарий, а управляющий по отчёту определяет, какие главы нужно
доработать. Получив задачу на доработку, переводчик может открыть
исходный текст главы и проверить значение отмеченного слова в конкретном
предложении.

\begin{figure}
\hypertarget{fig:10-5}{%
\centering
\includegraphics{images/revised/fig10-5.pdf}
\caption{Рисунок 10-5. Архитектура агентов для перевода
книги}\label{fig:10-5}
}
\end{figure}

После записи перевода в файл какому вызову по-прежнему требуется читать
длинный текст? В сопутствующем эксперименте первые две главы английского
издания были разделены на 26 единиц перевода. Один агент на протяжении
всей работы сохранял в контексте предыдущие исходные фрагменты и
переводы, поэтому максимальный объём входных данных одного запроса
составил 94 355 токенов. После разделения обязанностей максимальный
объём входных данных переводчика составил всего 2 857 токенов, а
управляющий использовал 2 153 токена при принятии решений о доработке.
Среди агентов управляемой группы больше всего входных данных получил
агент глоссария: за один раз он прочитал две главы, израсходовав 49 807
токенов.

Для перевода этих двух глав четыре роли в совокупности использовали 203
277 токенов, а один агент --- 1 317 808 токенов. Какая группа завершила
работу раньше? Группа из четырёх ролей затратила около 271 секунды ---
примерно на 17 секунд больше, чем один агент.

При сдаче работы две проверки дали разные результаты для одного и того
же перевода. В глоссарии для существительного \texttt{embedding} был
указан перевод «вектор встраивания». В отчёте о вычитке не было глав,
требующих доработки, однако программа проверки терминологии присвоила
этому слову результат 0/2: в переводах обеих глав встречалось слово
«встраивание», но ни разу --- «вектор встраивания». Следует ли вернуть
на доработку обе главы?

В абзаце первой главы о политике возврата исходный текст содержит фразу
\texttt{embedding\ the\ refund\ policy\ in\ the\ context}, переведённую
как «встраивание политики возврата в контекст». Здесь речь идёт о
помещении политики в контекст, поэтому замена на «вектор встраивания»
сама стала бы ошибкой. Программа ищет строки по главам и помогает
редактору находить фрагменты для повторной проверки, но для решения о
том, какое вхождение нужно исправить, редактору также требуется
соответствующий фрагмент исходного текста. При обновлении глоссария
можно сначала найти затронутые фрагменты, а затем передать переводчику
вместе исходный текст, перевод и новую запись.

\textbf{Требования эксперимента:}

\begin{enumerate}
\def\labelenumi{\arabic{enumi}.}
\tightlist
\item
  Выберите техническую книгу, содержащую иллюстрации, таблицы и код.
\item
  Реализуйте четыре типа агентов: управляющего, глоссария, перевода и
  вычитки.
\item
  Регистрируйте объём контекста каждого участника, передаваемые
  материалы и версии артефактов.
\item
  Сравните качество перевода, время выполнения и потребление ресурсов с
  решением на основе одного агента; проанализируйте влияние декомпозиции
  задачи и организации контекста.
\end{enumerate}
\end{quote}

\textbf{Параллельная координация.} Независимые друг от друга подзадачи
можно выполнять одновременно. Управляющий распределяет задачи, принимает
сведения о состоянии и по результатам решает, следует ли продолжить
работу, дополнить её или завершить. При переводе этой книги после
утверждения глоссария разные главы можно одновременно передать
нескольким переводчикам. На \crossreflink{fig:10-6}{рисунке 10-6} вместо
перевода представлен исследовательский отчёт: диаграммы уже готовы,
анализ ещё продолжается --- почему проверка форматирования всё равно
должна ждать? По событиям, отправляемым участниками, можно определить,
от каких двух артефактов зависит последний этап.

\begin{figure}
\hypertarget{fig:10-6}{%
\centering
\includegraphics{images/revised/fig10-6.pdf}
\caption{Рисунок 10-6. Параллельная координация
управляющего}\label{fig:10-6}
}
\end{figure}

Для организации такой асинхронной коммуникации подходит шина сообщений:
участники публикуют сведения о ходе работы и результаты, управляющий
подписывается на соответствующие события, а среда выполнения
обеспечивает доставку, буферизацию и повторные попытки.

\begin{quote}
\textbf{Пример продукта: Lingtai}

Lingtai --- локальная система долгоживущих агентов, сохраняющая
идентичность и память в файлах. В ней предусмотрены три роли: основной
агент (main agent) общается с пользователем и управляет планами и
памятью; демон (daemon) выполняет чётко ограниченные краткосрочные
параллельные задачи и по завершении возвращает выводы; аватар (avatar)
сохраняет собственные память, почтовый ящик и обязанности для
профессиональной специализации между сеансами.

У каждой роли есть собственные постоянные знания, а руководства по
навыкам в формате Markdown используются совместно. Когда контекст
приближается к предельному объёму, роль посредством «линьки» (molt)
сохраняет сводку, а затем переходит в новый контекст, подключив
постоянную память. Идентичность, память и навыки независимо хранятся в
файлах, что позволяет продолжить работу после замены базовой модели. В
этой архитектуре объединены управляющий, краткосрочные исполнители,
долгосрочные специализированные участники и сжатие контекста.
\end{quote}

\begin{quote}
\textbf{Эксперимент 10-3 ★★★: автономно координируемые телефонный и
компьютерный агенты}

Этот эксперимент объединяет Computer Use из
\crossreflink{chap:6}{главы 6} с голосовым взаимодействием. Пользователь
предоставляет страницу регистрации или бронирования авиабилета,
компьютерный агент открывает сайт и распознаёт форму, а затем с помощью
телефонного агента собирает отсутствующие сведения: имя, номер
документа, контактные данные, адрес, предпочтения и другие поля.

\textbf{Роли и передача данных.} Компьютерный агент управляет браузером
и координирует текущую задачу; телефонный агент общается с пользователем
в реальном времени посредством ASR, LLM и TTS, преобразуя устные ответы
в структурированные поля. Они взаимодействуют через двухточечные
инструменты или шину сообщений; каждое сообщение содержит отправителя,
получателя, тип и содержимое. Компьютерный агент может напрямую
делегировать телефонную задачу, выполняя обязанности управляющего. Если
пользователь согласен предоставить сведения голосом, телефонный агент
может последовательно задавать вопросы, подтверждать ответы и
переспрашивать, а компьютерный агент --- сосредоточиться на состоянии
страницы и заполнении полей.

\textbf{Два режима работы.} В фиксированном режиме заранее запускаются
два агента, каждый из которых поддерживает собственный цикл выполнения;
этот режим используется для проверки двусторонней коммуникации и
параллельного поведения. В автономном режиме сначала запускается
компьютерный агент, после чего модель на основании страницы, уже
имеющихся сведений и требований задачи решает, вызывать ли
\texttt{initiate\_phone\_call\_agent(purpose,\ required\_info)}. Затем
среда выполнения передаёт телефонному агенту цель задачи, список
требуемых полей и ограничения формата; после запуска используется тот же
механизм коммуникации.

\textbf{Последовательная передача отдельных полей.} На
\crossreflink{fig:10-7}{рисунке 10-7} телефонный агент посредством
WebRTC задаёт вопрос, извлекает и проверяет ответ, а после получения
допустимого значения отправляет \texttt{info\_collected} и переходит к
следующему вопросу. Одновременно компьютерный агент читает страницу,
заполняет предыдущее поле и возвращает \texttt{fill\_error} либо новое
состояние страницы. Если значение не соответствует формату, обратная
связь \texttt{format\_invalid} указывает поле и требования, чтобы
телефонный агент мог подтвердить ответ или переспросить. Благодаря этому
вопрос по следующему полю можно задавать параллельно с заполнением
предыдущего.

\textbf{Завершение и обработка исключительных ситуаций.} Собрав все
поля, телефонный агент отправляет \texttt{task\_completed}, обозначая
завершение подзадачи сбора информации. Компьютерный агент продолжает
проверять обязательные поля, результаты валидации страницы и наличие
разрешения на отправку, после чего завершает работу с формой. Если
исчерпано допустимое количество повторных попыток, возникла ошибка
страницы или пользователь потребовал остановиться, задача
приостанавливается с указанием текущего состояния. При завершении
следует отменить работу всё ещё активного второго агента и закрыть сеанс
браузера, звуковые дорожки и вызов. Согласие пользователя на голосовое
взаимодействие с человеком и разрешение на отправку формы запрашиваются
отдельно.

\begin{figure}
\hypertarget{fig:10-7}{%
\centering
\includegraphics{images/revised/fig10-7.pdf}
\caption{Рисунок 10-7. Архитектура из телефонного и компьютерного
агентов}\label{fig:10-7}
}
\end{figure}

\textbf{Требования эксперимента:}

\begin{enumerate}
\def\labelenumi{\arabic{enumi}.}
\tightlist
\item
  Реализуйте два независимых агента с двусторонней структурированной
  коммуникацией.
\item
  Запишите временную последовательность в фиксированном и автономном
  режимах и проверьте интервалы, в которых вопросы и заполнение
  выполняются параллельно.
\item
  Реализуйте проверку формата, повторные вопросы, передачу ошибок
  страницы, тайм-ауты и освобождение ресурсов.
\item
  Сравните два режима по принятию решения о запуске, задержке, доле
  успеха и потреблению ресурсов.
\end{enumerate}
\end{quote}

\begin{quote}
\textbf{Эксперимент 10-4 ★★★: AI-агент для одновременного сбора
информации с нескольких сайтов}

В этом эксперименте несколько однородных агентов выполняют поиск
параллельно, что позволяет на практике освоить событийную модель и
механизм прерывания из \crossreflink{chap:6}{главы 6}. Задача состоит в
том, чтобы найти указанного преподавателя, например Чжан Вэя, на сайтах
нескольких факультетов одного университета и вернуть сведения о
факультете, должности, направлениях исследований и странице-источнике.
Перед поиском необходимо согласовать критерии идентификации, например
имя, научную дисциплину и сведения на личной странице, чтобы различать
тёзок.

\textbf{1. Параллельный запуск.} Manager создаёт десять экземпляров
Computer Use Agent, каждый из которых отвечает за сайт одного
факультета. Экземпляры поддерживают независимые контексты и сеансы
браузера; их можно планировать с помощью асинхронных задач, потоков или
процессов. Входные данные содержат URL сайта, имя искомого человека,
критерии совпадения и идентификатор задачи для маршрутизации.

\textbf{2. Объединение сведений о ходе работы.} Участники периодически
отправляют состояния «загружается сайт», «читается список
преподавателей», «проверено, совпадений нет», «найден преподаватель с
таким же именем» и другие. Manager ведёт таблицу задач через шину
сообщений, раздельно регистрирует выполняющиеся, завершённые,
заблокированные и неудавшиеся задачи, а также сохраняет источники и
причины ошибок.

\textbf{3. Отмена остальных ветвей после проверки.} Предположим, что
отвечающий за физический факультет \texttt{agent\_3} нашёл преподавателя
с таким же именем. Он отправляет \texttt{target\_found}, полезная
нагрузка которого содержит сведения о преподавателе и источник. Если
задача требует вернуть проверенные сведения о человеке, Manager должен
сначала сверить критерии совпадения и свидетельства. Убедившись, что
найден нужный человек, он атомарно регистрирует результат, а затем
отправляет остальным работающим участникам \texttt{terminate}, указав,
что приняты сведения от \texttt{agent\_3}. Участники останавливаются в
безопасных точках, освобождают ресурсы и возвращают подтверждения;
Manager завершает объединение результатов после получения подтверждений
или наступления тайм-аута. Если требуется перечислить всех людей с таким
именем, поиск продолжается до проверки всей необходимой области.

\textbf{Гонки.} Агенты A и B могут вернуть допустимые результаты почти
одновременно. Если каждое из двух событий отдельно запустит объединение
результатов, возникнет повторная выдача или конфликт состояний. Переход
состояния «выполняется поиск → результат зарегистрирован» можно защитить
блокировкой или транзакцией, чтобы более позднее событие прочитало уже
зарегистрированный результат. Повторные попытки доставки сообщений также
необходимо дедуплицировать по идентификатору события, обеспечивая
согласованность даже при повторном запуске отмены и объединения
результатов.

\textbf{4. Исключения и область покрытия.} Сетевые сбои, изменения
структуры страницы и ошибки разбора регистрируются раздельно. Для
отдельной задачи можно задать бюджет тайм-аута, например две минуты,
после чего отменить выполнение и собрать диагностические сведения, не
останавливая остальных участников. В итоговом отчёте различаются
состояния «проверено, не найдено» и «доступ завершился ошибкой, проверка
не закончена», а также указывается охваченный набор факультетов.

На \crossreflink{fig:10-8}{рис. 10-8} эти сообщения объединены в
наглядную временную последовательность. Предположим, участник 3
сообщает, что на физическом факультете работает профессор Чжан Вэй,
исследующий квантовые вычисления. Manager сначала должен проверить,
подтверждает ли исходная веб-страница имя, факультет, должность и
направление исследований. Получение сведений о преподавателе с таким же
именем, подтверждение того, что это нужный человек, и сбор всех
подтверждений отмены происходят в три разных момента времени.

\begin{figure}
\hypertarget{fig:10-8}{%
\centering
\includegraphics{images/revised/fig10-8.pdf}
\caption{Рисунок 10-8. Временная последовательность сообщений при
параллельном поиске}\label{fig:10-8}
}
\end{figure}

В реальном запуске выполнялся поиск Andrew Ng. Сначала участники
проверяли наличие этого имени в тексте страницы и лишь после совпадения
просили модель извлечь сведения. Была выбрана страница Stanford HAI;
возвращённая короткая цитата позволяла сверить имя и статус профессора
по совместительству. Для проверки направления исследований требовалось
также прочитать расположенный ниже на странице абзац о глубоком
обучении. Даже если несколько участников одновременно возвращают
\texttt{target\_found}, зарегистрировать результат может только
участник, первым получивший блокировку.

При последовательном поиске совпадение нашлось только на последнем из
десяти сайтов, поэтому выполнение заняло 44,971 секунды; параллельный
поиск по тем же страницам занял 24,022 секунды, причём все десять
отдельных контекстов браузера были закрыты. Отношение этих значений
равно 1,872. Если поставить страницу HAI первой, последовательный поиск
остановится сразу после совпадения, поэтому время потребуется измерить
заново. В отдельном запуске для проверки отмены были запущены четыре
ветви: одна первой зарегистрировала результат, остальные три подтвердили
получение запроса на отмену, после чего все четыре контекста браузера
были закрыты.

\textbf{Требования эксперимента:}

\begin{enumerate}
\def\labelenumi{\arabic{enumi}.}
\tightlist
\item
  Реализуйте Manager, динамически создающий параллельных участников.
\item
  Реализуйте возможности работы с браузером на основе таких проектов,
  как browser-use.
\item
  Реализуйте двустороннюю коммуникацию и ведение состояния через шину
  сообщений.
\item
  Проверяйте результаты поиска, атомарно регистрируйте результат и
  собирайте подтверждения после отмены остальных ветвей.
\item
  Обрабатывайте недоступность сайта, ошибки разбора, тайм-ауты и
  отсутствие совпадений.
\item
  Сравните время завершения, область покрытия и потребление ресурсов при
  параллельном и последовательном выполнении.
\end{enumerate}
\end{quote}

\textbf{Генерация рабочего процесса управляющим.} Помимо
последовательного делегирования, Manager может сначала записать
зависимости задач и правила объединения результатов в виде программы, а
затем передать её среде выполнения для планирования. Модель отвечает за
проектирование процесса, а программа --- за запуск участников, передачу
результатов и ожидание зависимостей в соответствии с заданными
правилами. Когда зависимости уже понятны и однотипный процесс требуется
выполнять многократно, автор предпочитает этот подход. Он устраняет
затраты на вызов Manager при каждой выдаче задачи и сокращает объём
накапливающихся в управляющем контексте записей о планировании.

Динамический инструмент Workflow в Claude Code предоставляет такую
реализацию. \texttt{agent()} создаёт дочернего агента с независимым
контекстом, а schema позволяет задать структурированный вывод;
\texttt{parallel()} параллельно выполняет группу задач и ожидает
завершения всех; \texttt{pipeline()} проводит каждый входной элемент
через несколько этапов, причём, когда один элемент переходит на
следующий этап, остальные могут продолжать обрабатываться на предыдущем.
Промежуточные результаты сохраняются в переменных программы, а
фильтрация и объединение выполняются обычным кодом.\footnote{Anthropic,
  ``Orchestrate subagents at scale with dynamic workflows'', 2026.}

Например, при проверке семи групп фактов в технической статье можно
сначала запустить исследовательских агентов по отдельным направлениям, а
затем передать выводы по каждой группе проверяющим агентам для
параллельной проверки. Каждый возвращённый вывод сохраняет источник и
основание оценки; после завершения всех ветвей обобщающий агент
составляет итоговый отчёт. Таким образом, разделение обязанностей и
зависимости между этапами «исследование --- проверка каждого вывода ---
общее обобщение» фиксируются в рабочем процессе.

\hypertarget{ux434ux435ux446ux435ux43dux442ux440ux430ux43bux438ux437ux43eux432ux430ux43dux43dux44bux439-ux440ux435ux436ux438ux43c}{%
\subsection{Децентрализованный
режим}\label{ux434ux435ux446ux435ux43dux442ux440ux430ux43bux438ux437ux43eux432ux430ux43dux43dux44bux439-ux440ux435ux436ux438ux43c}}

Зачем позволять участникам самостоятельно выбирать, к кому обратиться,
если уже есть Manager? Представим, что тестировщик получает только что
написанный модуль и обнаруживает, что требования к продукту и
спецификация интерфейса по-разному трактуют одно и то же поле.
Тестировщик может напрямую передать это противоречие ролям, отвечающим
за продукт и архитектуру, чтобы каждая сторона обосновала свою позицию,
а затем продолжить тестирование после исправления договорённости. Это не
только разделение труда, но и система сдержек и противовесов. Инженер
передаёт модуль, архитектор запрашивает мнение продуктовой роли,
тестировщик выдвигает возражение --- каждый из них может инициировать
следующий шаг. Если участникам приходится часто согласовывать решения,
ожидание, пока Manager поочерёдно передаст каждое сообщение, будет
задерживать работу.

Заимствуя терминологию микросервисов, централизованное управление
последовательностью шагов называют \textbf{оркестрацией}
(orchestration), а самостоятельное продолжение работы сторонами в ответ
на согласованные события --- \textbf{хореографией} (choreography). Есть
и ещё одно практическое различие: если каждая передача требует ожидания
Manager, то при прекращении ответа со стороны его сервиса модели
остальные участники не смогут получить следующую задачу, даже если сами
по-прежнему способны работать. Распределённая координация уменьшает эту
зависимость. Однако если принимающий работу участник сам завис внутри
вызова инструмента, другой участник сможет понять, откуда продолжить,
лишь при наличии постоянно сохраняемого состояния задачи, тайм-аутов,
повторных попыток и механизма перехвата работы.

В следующих трёх вариантах используются разные способы передачи работы:
подписка на сообщения о результатах, выбор следующего выступающего в
групповом чате или адресная передача управления текущим участником.

\textbf{MetaGPT: организация передачи результатов с помощью SOP.}
MetaGPT распределяет работу согласно стандартной операционной процедуре
(SOP) команды разработки ПО. Менеджер продукта передаёт документ с
требованиями к продукту (PRD), архитектор --- проектную документацию,
инженер --- код, а тестировщик --- отчёт о тестировании. Следующий
участник может начать работу сразу после получения согласованных
материалов, а предыдущему не приходится многократно объяснять задачу.
Линии на \crossreflink{fig:10-9}{рис. 10-9} обозначают связи передачи
этих результатов.\footnote{Hong, S., et al.~\emph{MetaGPT: Meta
  Programming for A Multi-Agent Collaborative Framework.} ICLR 2024.}

\begin{figure}
\hypertarget{fig:10-9}{%
\centering
\includegraphics{images/revised/fig10-9.pdf}
\caption{Рисунок 10-9. Сеть взаимодействия нескольких AI-агентов в
MetaGPT}\label{fig:10-9}
}
\end{figure}

Какие сведения каждая роль должна добавлять при передаче работы?

\begin{itemize}
\tightlist
\item
  \textbf{Product Manager Agent} принимает требования, оформляет PRD и
  уточняет функции, пользовательские истории, критерии приёмки и
  приоритеты.
\item
  \textbf{Architect Agent} читает PRD, выбирает технологический стек,
  выделяет модули, определяет интерфейсы и модели данных, формируя
  проектную документацию.
\item
  \textbf{Project Manager Agent} на основе проекта разбивает работу на
  задачи, распределяет файлы, упорядочивает зависимости и назначает
  выполнение инженерам.
\item
  \textbf{Engineer Agents} реализуют модули по проекту и передают код;
  независимые части можно поручить нескольким экземплярам.
\item
  \textbf{QA Engineer Agent} на основе кода и требований создаёт тесты,
  выполняет проверки, фиксирует проблемы и формирует отчёт о
  тестировании.
\end{itemize}

При проектировании формата передачи для этих ролей автор рекомендует
явно указывать три вещи: задачу и критерии приёмки, подтверждённые факты
и ограничения, а также ссылки на структурированные результаты. На их
основе принимающая сторона получает необходимые материалы, продолжает
работу и задаёт конкретные вопросы, если информации не хватает.

MetaGPT организует информацию с помощью общего пула сообщений и
механизма подписки: роли публикуют сообщения, а другие роли обрабатывают
соответствующие входные данные согласно типам сообщений, за которыми
следят. Издатели и потребители связаны форматом результатов и семантикой
сообщений, поэтому для каждого из них можно отдельно заменить модель или
изменить внутреннюю реализацию.

SOP описывает зависимости между требованиями к продукту, архитектурой и
кодом, а подписка на сообщения --- способы запуска этих работ.
Фиксированные зависимости задач также могут продвигаться распределёнными
событиями; вместе эти механизмы образуют способ взаимодействия в
системе. При дальнейшем проектировании динамической обратной связи можно
добавить пути, по которым тестировщик уточняет требования у продуктовой
роли, а инженер предлагает архитектору альтернативное решение, и задать
для таких взаимодействий условия приёма и правила завершения.

\textbf{Групповой чат AutoGen: управление очерёдностью реплик в общем
сеансе.} В классической архитектуре GroupChat несколько агентов
участвуют в одном сеансе, GroupChatManager выбирает следующего
выступающего и рассылает сообщения участникам. В качестве стратегии
выбора можно использовать циклический перебор либо поручить модели
принять решение на основе текущего содержимого. Участники сохраняют
собственные конфигурации ролей и взаимодействуют через общий сеанс, а
управление очерёдностью реплик централизовано в менеджере. Это пример
сочетания организации контекста со структурой управления.\footnote{Microsoft,
  \emph{AutoGen 0.2: Conversation Patterns.}}

\textbf{OpenAI Swarm: передача управления посредством handoff.} В Swarm
для ролей настраиваются инструкции и инструменты. После вызова ролью
настроенной функции передачи среда выполнения переключается на целевого
агента. Решение о маршрутизации могут принимать отдельные роли, а
исполняет его единый цикл выполнения. При передаче системные инструкции
изменяются вместе с ролью, тогда как история диалога сохраняется. Это
демонстрирует сочетание двух подходов: «роль выбирает следующего
исполнителя» и «история сеанса продолжает использоваться».\footnote{OpenAI,
  \emph{Swarm}, разделы «Instructions» и «Handoffs and Updating Context
  Variables».}

Если передача осуществляется между независимыми контекстами, можно
отдельно определить пакет задачи: \texttt{task\_id} идентифицирует
задачу, \texttt{sender} и \texttt{recipient} указывают стороны,
\texttt{goal} и \texttt{constraints} описывают цель и ограничения,
\texttt{accepted\_facts} хранит подтверждённые факты,
\texttt{artifact\_refs} содержит ссылки на результаты,
\texttt{remaining\_budget} фиксирует бюджет, а \texttt{visited\_agents}
хранит цепочку посещений. Получатель читает материалы в пределах
предоставленных прав, а среда выполнения ведёт учёт бюджета и передач.

Цепочка передачи должна учитывать повторные посещения. Например,
разработчик передаёт код рецензенту, тот запрашивает изменения и
возвращает код разработчику --- это нормальный цикл доработки.
Определить наличие прогресса можно по сочетанию состояния задачи, версии
результата, нерешённых вопросов и числа повторений. Например, если один
и тот же перевод дважды возвращают на доработку, но в первый раз в него
добавляют источники, а во второй исправляют терминологию, доработка всё
ещё продвигается. Если же оба раза возвращается одна и та же версия с
указанием на одну и ту же проблему, значит, работа зашла в тупик. В
таком случае можно ограничить число дальнейших передач и сообщить
ответственному лицу о нерешённой проблеме.

\begin{quote}
\textbf{Agent Swarm в продуктах}

Термин «Agent Swarm» может обозначать разные способы организации. В
одном случае управление передаётся по сети handoff, как в OpenAI Swarm,
библиотеке swarm для LangGraph и оркестрации handoff в Microsoft Agent
Framework. В другом используется главный агент и параллельные
исполнители, как в Kimi Agent Swarm, многоагентной исследовательской
системе Anthropic и Manus Wide Research.

Представленный в Kimi K2.5 Agent Swarm обучает оркестратор посредством
обучения с подкреплением с параллельными агентами: он учится определять,
когда следует разбить задачу и как распределить работу, поддерживая до
100 дочерних агентов. В последующих версиях масштаб увеличен до 300
агентов, а K3 Swarm продолжает использоваться для крупномасштабного
поиска и пакетной обработки. AgentENV, в свою очередь, предоставляет
масштабируемую среду для обучения агентов Kimi K3 с подкреплением,
поддерживает снимки состояния и ветвление и применяется для организации
параллельного выполнения.\footnote{Moonshot AI, \emph{Kimi Agent Swarm:
  100 Sub-Agents at Scale}, справочная документация Agent Swarm и
  описание проекта AgentENV.}
\end{quote}

\textbf{Несколько экземпляров агентов на одной машине.} Несколько
экземпляров могут выполнять разные задачи, совместно используя файлы,
сервисы инструментов или GPU. В таком случае посредством коммуникации
можно обозначать область изменений, уведомлять об освобождении ресурсов
или согласовывать порядок выполнения. Например, если двум экземплярам
требуется один GPU, они могут сначала зарегистрировать задачи и
предполагаемый объём ресурсов, после чего диспетчер ресурсов предоставит
право использования.

Agent Teams в Claude Code предоставляет обнаружение участников команды,
обмен сообщениями и совместное управление задачами. Участники получают
информацию друг о друге в пределах созданной команды, а среда выполнения
организует коммуникацию.\footnote{Anthropic, \emph{Orchestrate teams of
  Claude Code sessions.}} При самостоятельной реализации взаимодействия
на одной машине также можно использовать интерфейсы наподобие
\texttt{list\_agents} из \crossreflink{chap:4}{главы 4}, задав явные
правила для области обнаружения, прав доступа и владения ресурсами.
Затем согласованные в сообщениях решения должны обеспечиваться
блокировками, квотами или планировщиком.

\hypertarget{ux43cux435ux436ux43eux440ux433ux430ux43dux438ux437ux430ux446ux438ux43eux43dux43dux43eux435-ux432ux437ux430ux438ux43cux43eux434ux435ux439ux441ux442ux432ux438ux435-ux43fux440ux43eux442ux43eux43aux43eux43b-a2a}{%
\subsection{Межорганизационное взаимодействие: протокол
A2A}\label{ux43cux435ux436ux43eux440ux433ux430ux43dux438ux437ux430ux446ux438ux43eux43dux43dux43eux435-ux432ux437ux430ux438ux43cux43eux434ux435ux439ux441ux442ux432ux438ux435-ux43fux440ux43eux442ux43eux43aux43eux43b-a2a}}

Когда взаимодействие выходит за границы сервиса или организации,
сторонам необходимо согласовать обнаружение возможностей,
аутентификацию, формат запросов, передачу прогресса и результатов.
Выпущенный Google в 2025 году протокол \textbf{A2A} (Agent2Agent),
позднее переданный под управление Linux Foundation, предоставляет такой
интерфейс совместимости. Его можно рассматривать с трёх
сторон.\footnote{A2A Project, \emph{Основные концепции и компоненты
  A2A.}}

\begin{itemize}
\tightlist
\item
  \textbf{Обнаружение возможностей.} Agent Card --- документ с
  метаданными, описывающий идентичность AI-агента, конечные точки
  сервиса, навыки, модальности входных и выходных данных и требования к
  аутентификации. На его основе вызывающая сторона определяет, способен
  ли агент обработать задачу и как установить соединение.
\item
  \textbf{Задачи и сообщения.} Task представляет работу с
  идентификатором и жизненным циклом. Вызывающая сторона может
  отслеживать такие состояния, как отправлена, выполняется, требует
  входных данных, завершена или завершилась с ошибкой. Message передаёт
  требования задачи, дополнительную информацию и ответы, а Artifact ---
  документы, изображения или структурированные результаты. Part в
  сообщениях и артефактах может содержать текст, ссылку на файл или
  данные.
\item
  \textbf{Взаимодействие через интерфейс.} Вызывающая сторона общается с
  удалённой через эти открытые объекты, а удалённая сторона
  самостоятельно управляет внутренними промптами, памятью и
  инструментами. Каждая сторона отдельно обеспечивает аутентификацию,
  авторизацию и ограничение области доступа к данным.
\end{itemize}

Состояние длительной задачи можно получать опросом, потоковые обновления
--- через SSE, а при наличии поддержки --- посредством push-уведомлений.
После предоставления пользователем дополнительной информации клиент
продолжает выполнение той же задачи и получает итоговые артефакты в
соответствии с её состоянием.

MCP из \crossreflink{chap:4}{главы 4} ориентирован на предоставление
модельным приложениям интерфейсов инструментов, ресурсов и промптов, а
A2A --- на совместное выполнение задач сервисами агентов. Их можно
сочетать: один агент использует инструменты через MCP и одновременно
делегирует удалённые задачи через A2A. Если внутри одной команды уже
существуют стабильные интерфейсы, автор предпочитает сначала
использовать вызовы инструментов, общие файлы и шину сообщений. A2A
следует вводить, когда требуется подключить независимо сопровождаемый
сервис агентов и унифицировать обнаружение возможностей и жизненный цикл
задач; при взаимодействии между организациями дополнительно
настраиваются правила идентификации, разрешений и управления данными.

\textbf{Согласование времени для разных людей.} У трёх членов семьи есть
собственные помощники по планированию, и они хотят договориться о
совместном ужине. Один помощник видит полный календарь своего владельца
--- что ему следует ответить двум другим? Фраза «Во вторник неудобно, а
в среду можно кое-что перенести» уже сообщает полезное различие. Прямая
пересылка календаря раскроет лишние сведения о личных делах во вторник;
если же ответить лишь «В оба дня уже есть планы», остальные не узнают,
что насчёт среды ещё можно договориться. Даже если встречу удаётся
назначить каждый раз, такая система может быть невыгодна человеку,
которому постоянно приходится менять свои планы.

В одной открытой задаче CalBench пяти людям необходимо назначить пять
встреч, в каждой из которых участвуют по три человека. Некоторые
интервалы у участников совершенно нельзя переносить, некоторые личные
дела можно сдвинуть, причём задача задаёт стоимость каждого переноса.
Сначала попробуем простой подход: обработаем пять встреч по очереди,
каждый раз начиная поиск с самого раннего интервала, и назначим встречу,
как только все трое смогут освободить время; уже назначенные встречи
сохраняются. Затем попробуем другой подход: рассмотрим все пять встреч
вместе и найдём решение на основе полных календарей. Ниже используются
номера временных интервалов из задачи, начиная с 0.\footnote{В статье о
  CalBench отдельно рассматриваются стоимость координации,
  справедливость, раскрытие выводимого состояния календаря и содержимого
  конкретных дел. Числа взяты из первой задачи набора с различной
  стоимостью переносов в версии \texttt{86f8123}. При первом способе
  пропускаются интервалы, в которых у любого участника уже назначена
  встреча; остальные правила переноса соответствуют правилам авторов.
  Результат можно перепроверить с помощью решателя расписаний.}

\begin{longtable}[]{@{}
  >{\raggedright\arraybackslash}p{(\columnwidth - 6\tabcolsep) * \real{0.3276}}
  >{\raggedright\arraybackslash}p{(\columnwidth - 6\tabcolsep) * \real{0.2172}}
  >{\raggedright\arraybackslash}p{(\columnwidth - 6\tabcolsep) * \real{0.2653}}
  >{\raggedleft\arraybackslash}p{(\columnwidth - 6\tabcolsep) * \real{0.1898}}@{}}
\toprule\noalign{}
\begin{minipage}[b]{\linewidth}\raggedright
\hspace{0pt}Способ планирования
\end{minipage} & \begin{minipage}[b]{\linewidth}\raggedright
\hspace{0pt}Интервалы встреч с 1-\allowbreak{}\hspace{0pt}й по
5-\allowbreak{}\hspace{0pt}ю
\end{minipage} & \begin{minipage}[b]{\linewidth}\raggedright
\hspace{0pt}Стоимость переноса для каждой назначенной встречи
\end{minipage} & \begin{minipage}[b]{\linewidth}\raggedleft
\hspace{0pt}Общая стоимость
\end{minipage} \\
\midrule\noalign{}
\endhead
\bottomrule\noalign{}
\endlastfoot
\hspace{0pt}Для каждой встречи выбирается самый ранний доступный
интервал, уже назначенные встречи сохраняются & \hspace{0pt}0, 2, 1, 3,
6 & \hspace{0pt}1, 20, 20, 12, 300 & \hspace{0pt}353 \\
\hspace{0pt}Эталонный решатель совместно планирует все пять встреч &
\hspace{0pt}13, 3, 10, 9, 12 & \hspace{0pt}0, 1, 0, 1, 1 &
\hspace{0pt}3 \\
\end{longtable}

В обоих расписаниях каждая встреча совпадает в окончательных календарях
всех трёх участников. Наибольшая разница возникает для последней
встречи: при первом способе интервал 6 можно освободить, но каждому из
трёх участников придётся перенести дело стоимостью 100. Если назначить
встречу на интервал 12, лишь одному человеку потребуется перенести дело
стоимостью 1. Глобальный результат даёт проверяющему основу для
сравнения, тогда как каждый помощник, фактически участвующий в
согласовании, видит только календарь своего владельца.

Встроенный в эту версию жадный решатель сообщает стоимость 146 и
возвращает интервалы для всех пяти встреч. Однако при планировании более
поздних встреч он разрешает удалять ранее назначенные встречи из
календарей отдельных участников. В итоге первая встреча оказывается в
календарях трёх участников в разных интервалах: 7, 0 и 11. Если считать
встречу принятой только тогда, когда в окончательных календарях всех
троих она указана в одном и том же интервале, успешно назначены лишь
третья, четвёртая и пятая встречи.

В реальном сервисе планирования возникают и сведения, которых нет в
формальной задаче. Сервис Microsoft Calendar.help за год согласовал по
электронной почте несколько тысяч встреч. Письмо относится к назначению
новой встречи или к переносу предыдущей? Если автоматическая
классификация не уверена в ответе, человек сначала видит входящее письмо
и решает только этот вопрос. Если по одному письму определить это
невозможно, задачу можно передать специалисту по планированию, имеющему
доступ к переписке, уже собранной информации и обезличенным календарям.
Такая архитектура связывает вопросы «кого следует привлечь к принятию
решения» и «сколько информации этому человеку необходимо
увидеть».\footnote{Cranshaw et al., проектирование и развёртывание
  Calendar.help; см. описание двухуровневого интерфейса для выполняемых
  человеком микрозадач и полных задач планирования.}

\hypertarget{ux440ux435ux436ux438ux43cux44b-ux43eux442ux43aux430ux437ux430-ux43fux440ux438-ux432ux437ux430ux438ux43cux43eux434ux435ux439ux441ux442ux432ux438ux438-ux43dux435ux441ux43aux43eux43bux44cux43aux438ux445-ai-ux430ux433ux435ux43dux442ux43eux432}{%
\section{Режимы отказа при взаимодействии нескольких
AI-агентов}\label{ux440ux435ux436ux438ux43cux44b-ux43eux442ux43aux430ux437ux430-ux43fux440ux438-ux432ux437ux430ux438ux43cux43eux434ux435ux439ux441ux442ux432ux438ux438-ux43dux435ux441ux43aux43eux43bux44cux43aux438ux445-ai-ux430ux433ux435ux43dux442ux43eux432}}

После совместного выполнения задачи несколькими участниками сбои могут
возникать на стыках ролей, передачи работы и проверки. Авторы статьи
«Why Do Multi-Agent LLM Systems Fail?» построили классификацию отказов
примерно по 150 трассам с согласованностью разметки κ = 0.88, а затем
распространили анализ на семь фреймворков. Исследователи выделили 14
режимов отказа, разделённых на три категории:\footnote{Cemri, M., et
  al.~\emph{Why Do Multi-Agent LLM Systems Fail?} 2025.}

\begin{itemize}
\tightlist
\item
  \textbf{Проблемы проектирования системы:} пробелы в обязанностях
  ролей, интерфейсах, конфигурации инструментов или условиях остановки.
\item
  \textbf{Отклонения во взаимодействии участников:} различное понимание
  цели, неполная передача информации или взаимно противоречащие
  действия.
\item
  \textbf{Проблемы проверки задачи:} преждевременное завершение,
  пропущенные проверки или оценка готовности по неподходящим критериям.
\end{itemize}

В статье приводится пример программы для игры в шахматы. Созданный
ChatDev код компилировался и успешно проходил проверку, но при игре по
настоящим шахматным правилам программа работала неверно. Компилятор
проверяет, можно ли скомпилировать код, а для проверки допустимости
ходов необходимо сыграть хотя бы несколько ходов. В другом сравнительном
эксперименте исследователи разрешили возвращать замечания ChatDev на
доработку, завершать процесс только после подтверждения участником,
отвечающим за техническую проверку, и ограничили максимальное число
итераций. Из 32 задач ProgramDev-v0 исходный процесс решил 8, а
изменённый --- 13. Доля успешно выполненных задач выросла с 25.0\% до
40.6\%, то есть примерно на 15.6 процентного пункта.\footnote{Cemri, M.,
  et al.~\emph{Why Do Multi-Agent LLM Systems Fail?} 2025.}

Как отличить ситуацию «проверяющий не ответил» от ситуации «проверяющий
ответил, но ошибся»? В распределённых системах прекращение работы
компонента называется аварийным отказом. Модель византийских отказов
также допускает, что компонент выдаёт произвольные ошибочные или
противоречащие друг другу результаты. Агент тоже может остановиться на
полпути, а может продолжить работу и вернуть неверный вывод. В первом
случае требуются обнаружение, восстановление и перехват работы, во
втором --- проверка доказательств, ограничение полномочий и изоляция
ошибок. Регулярные сигналы heartbeat помогают определить, продолжает ли
агент работать, но правильность возвращённого им вывода всё равно нужно
проверять по фактическому результату. Для оценки результатов можно
использовать перекрёстную проверку и голосование, надёжность которых
зависит от качества доказательств и корреляции ошибок участников.
Формальные протоколы также определяют, сколько отказавших узлов
допустимо и как узлы должны взаимодействовать.

\hypertarget{ux440ux435ux436ux438ux43c-ux43eux442ux43aux430ux437ux430-1-ux43aux43eux43dux43aux443ux440ux435ux43dux442ux43dux44bux435-ux43aux43eux43dux444ux43bux438ux43aux442ux44b-ux432-ux43eux431ux449ux435ux439-ux444ux430ux439ux43bux43eux432ux43eux439-ux441ux438ux441ux442ux435ux43cux435}{%
\subsection{Режим отказа 1: конкурентные конфликты в общей файловой
системе}\label{ux440ux435ux436ux438ux43c-ux43eux442ux43aux430ux437ux430-1-ux43aux43eux43dux43aux443ux440ux435ux43dux442ux43dux44bux435-ux43aux43eux43dux444ux43bux438ux43aux442ux44b-ux432-ux43eux431ux449ux435ux439-ux444ux430ux439ux43bux43eux432ux43eux439-ux441ux438ux441ux442ux435ux43cux435}}

Общие артефакты позволяют участникам продолжать работу друг друга, но
одновременно создают проблему параллельных изменений. Сначала следует
разделить конфликты на два уровня.

\textbf{Конфликт на уровне файлов} возникает, когда два участника
изменяют один файл: оба читают старую версию, затем по очереди
записывают изменения, и последняя запись перезаписывает результат
предыдущей.

\textbf{Семантический конфликт} возникает между связанными артефактами.
Например, агент A выполняет сквозную перенумерацию иллюстраций во всей
книге, а агент B одновременно изменяет главу и ссылается на прежние
номера. Они изменяют разные файлы, поэтому при слиянии текстового
конфликта может не возникнуть, но соответствие текста иллюстрациям
окажется нарушено.

Допустим, текущая версия глоссария --- 7. Два участника начинают вносить
изменения в версию 7. Первый отправляет результат раньше, и версия
становится 8. Если второй при отправке лишь укажет «записать эту
таблицу», изменения первого могут быть перезаписаны. Если же потребовать
«записывать, только если текущая версия всё ещё равна 7», программа
обнаружит необходимость слияния. Для конкурентного доступа на уровне
файлов можно использовать \textbf{оптимистическую блокировку}: при
чтении запоминать номер версии, а при отправке выполнять проверку версии
и запись как одну атомарную операцию. Если версия изменилась, устаревшая
запись отклоняется, после чего необходимо заново прочитать данные и
выполнить слияние.

Межфайловая семантическая согласованность требует проверки связей. Когда
несколько агентов для написания кода изменяют один репозиторий, автор
предпочитает выделять каждому участнику отдельные \texttt{worktree} и
ветвь, формировать отдельные коммиты, а в точке слияния проверять
конфликты, запускать тесты и проверять интерфейсы. Одного создания
отдельных ветвей недостаточно: для изоляции одновременных изменений
файлов нужны также отдельные рабочие каталоги.

\hypertarget{ux440ux435ux436ux438ux43c-ux43eux442ux43aux430ux437ux430-2-ux43aux430ux441ux43aux430ux434ux43dux43eux435-ux443ux441ux438ux43bux435ux43dux438ux435-ux43eux448ux438ux431ux43eux43a}{%
\subsection{Режим отказа 2: каскадное усиление
ошибок}\label{ux440ux435ux436ux438ux43c-ux43eux442ux43aux430ux437ux430-2-ux43aux430ux441ux43aux430ux434ux43dux43eux435-ux443ux441ux438ux43bux435ux43dux438ux435-ux43eux448ux438ux431ux43eux43a}}

Сообщение может быть доставлено дословно, но его смысл всё равно
способен исказиться при пересказе. Представим, что несколько агентов
совместно составляют запись об обращении в службу поддержки. Первый
пишет: «Сотрудник поддержки пообещал уменьшить следующий счёт на 80
юаней». Второй сокращает это до «Счёт уменьшен на 80 юаней», а третий,
прочитав сводку, закрывает задачу. Все три сообщения были доставлены, но
пользователь лишился напоминания проверить следующий счёт. Как и в игре
«испорченный телефон», ошибка возникла в тот момент, когда второй
участник представил обещание как уже свершившийся результат.

В записи о передаче работы следует сохранять вывод, условия его
применимости и источник, чтобы следующий участник видел, что уже
подтверждено, а что остаётся предположением. В описанном спорном случае
автор попросил бы проверяющего сначала прочитать исходную запись
разговора и последний счёт, самостоятельно определить текущий этап
выполнения, а затем выяснить, почему второй участник составил именно
такую сводку. Если имеется только устное обещание, второй участник
должен вернуть формулировку «обещано уменьшение» и сохранить задачу по
проверке счёта. Основанием для её закрытия служит лишь запись о вычете в
самом счёте.

\hypertarget{ux440ux435ux436ux438ux43c-ux43eux442ux43aux430ux437ux430-3-ux43eux434ux43dux43eux440ux43eux434ux43dux430ux44f-ux43aux43eux43dux432ux435ux440ux433ux435ux43dux446ux438ux44f}{%
\subsection{Режим отказа 3: однородная
конвергенция}\label{ux440ux435ux436ux438ux43c-ux43eux442ux43aux430ux437ux430-3-ux43eux434ux43dux43eux440ux43eux434ux43dux430ux44f-ux43aux43eux43dux432ux435ux440ux433ux435ux43dux446ux438ux44f}}

Несколько участников могут независимо друг от друга принять одно и то же
ошибочное решение. В эксперименте Anthropic 18 из 30 одновременно
запущенных агентов создали ветви Git с одинаковым именем; в задаче на
написание текста разные участники также использовали одинаковые
заголовки.\footnote{Anthropic Frontier Red Team, ``Patterns and Problems
  in Emerging Multiagent Systems,'' 2026-08-13.} Общая модель, похожие
входные данные и одинаковый фреймворк выполнения делают решения
коррелированными. Единогласие пяти похожих агентов может оказаться всего
лишь пятикратным повторением одной ошибки. Поэтому при увеличении числа
проверяющих следует одновременно повышать разнообразие доказательств и
точек зрения на анализ.

Участникам можно назначать разные источники доказательств, обязанности
или модели, а идентификаторы ресурсов изолировать с помощью пространств
имён. Ограничения параллелизма и скорости контролируют одновременный
доступ и не позволяют множеству одинаковых решений одновременно
перегружать общий сервис.

Эксперимент той же работы с ценообразованием по Бертрану
продемонстрировал и другой вид конвергенции: стремящиеся к прибыли
агенты могли координировать цены через закрытый канал, но также были
способны выработать согласованное поведение на основе общедоступных
ценовых предложений.

\hypertarget{ux440ux435ux436ux438ux43c-ux43eux442ux43aux430ux437ux430-4-ux441ux43fux43eux440ux44b-ux43eux431-ux43eux431ux44fux437ux430ux43dux43dux43eux441ux442ux44fux445-ux438-ux43aux43eux43dux444ux43bux438ux43aux442-ux446ux435ux43bux435ux439}{%
\subsection{Режим отказа 4: споры об обязанностях и конфликт
целей}\label{ux440ux435ux436ux438ux43c-ux43eux442ux43aux430ux437ux430-4-ux441ux43fux43eux440ux44b-ux43eux431-ux43eux431ux44fux437ux430ux43dux43dux43eux441ux442ux44fux445-ux438-ux43aux43eux43dux444ux43bux438ux43aux442-ux446ux435ux43bux435ux439}}

Когда несколько участников преследуют взаимоисключающие цели, их
действия могут взаимно нейтрализоваться. Anthropic поручила трём агентам
перенести один и тот же бэкенд на разные языки и обнаружила, что они
восприняли изменения других участников как препятствия, после чего
конфликт распространился на операции с процессами и
разрешениями.\footnote{Anthropic Frontier Red Team, ``Patterns and
  Problems in Emerging Multiagent Systems,'' 2026-08-13.} Для
предотвращения таких проблем необходимо заранее согласовать приоритеты
целей, владельцев ресурсов и допустимые границы действий. Если конфликт
невозможно разрешить по установленным правилам, среда выполнения должна
приостановить соответствующие операции и передать решение ответственному
лицу.

При взаимодействии по образцу команды разработчиков также возможно
постоянное перекладывание ответственности. Например, инженер по
тестированию сообщает о проблеме, фронтенд-разработчик считает, что
сначала нужно изменить бэкенд, бэкенд-разработчик требует уточнения от
менеджера продукта, а тот возвращает проблему архитектору. В таком
случае можно потребовать, чтобы при каждой передаче указывались условия
воспроизведения, имеющиеся доказательства, нерешённый вопрос и
обязанности принимающей стороны. Это позволит постепенно разрешить спор.

Другой пример --- неисправность тестовой среды: разработчик неоднократно
изменяет бизнес-логику, но тест продолжает возвращать одну и ту же
ошибку. Сначала следует зафиксировать среду и входные данные, сравнить
результаты выполнения до и после изменений, установить, вызвана ли
проблема реализацией продукта или условиями тестирования, и лишь затем
назначить задачу на исправление.

\hypertarget{ux440ux435ux436ux438ux43c-ux43eux442ux43aux430ux437ux430-5-ux43dux435ux443ux43fux440ux430ux432ux43bux44fux435ux43cux44bux439-ux446ux438ux43aux43b}{%
\subsection{Режим отказа 5: неуправляемый
цикл}\label{ux440ux435ux436ux438ux43c-ux43eux442ux43aux430ux437ux430-5-ux43dux435ux443ux43fux440ux430ux432ux43bux44fux435ux43cux44bux439-ux446ux438ux43aux43b}}

Преждевременное завершение приводит к пропуску работы, а отсутствие
условий остановки может заставить систему бесконечно повторять попытки,
передавать задачу между участниками или рекурсивно создавать субагентов.
Число участников и количество раундов вызовов взаимно усиливают друг
друга, что может привести к появлению тысяч подзадач и быстрому
исчерпанию бюджета.

В среде выполнения можно одновременно установить ограничения на общее
число токенов или стоимость, параллелизм, суммарное число созданных
задач, глубину порождения, число раундов и время выполнения. Бюджет
выделяется родительской задачей и учитывается в общей ведомости;
расходы, уже понесённые до сбоя или отмены, также необходимо
регистрировать. При достижении ограничения следует сохранить прогресс,
отменить незавершённые ветви, собрать подтверждения отмены и передать
текущий результат вместе со списком оставшихся задач.

Для совместных задач с высокой степенью автономности автор рекомендует
использовать отдельные проекты или учётные данные API, чтобы раздельно
учитывать потребление и отзывать доступ. За жёсткие ограничения во время
выполнения отвечают механизмы квот и среда выполнения. В длительных
задачах также можно проверять, появлялись ли на протяжении нескольких
последовательных раундов новые доказательства или завершённые пункты, и
на этой основе решать, следует ли изменить стратегию или остановиться.

Почему полномочия необходимо проверять ещё и на уровне всего
эксперимента? В июле 2026 года внутренняя оценка OpenAI затронула
реальные системы Hugging Face. После инцидента OpenAI ужесточила сетевые
разрешения и усилила изоляцию и мониторинг исследовательских сред. При
проектировании многоагентного эксперимента ответственное лицо должно
уметь ответить на два вопроса. Остаётся ли совокупность всех ресурсов, к
которым имеют доступ участники, в пределах разрешённой области
эксперимента? Если мониторинг обнаружит выход за эти пределы, можно ли
одной операцией остановить весь эксперимент и установить, какие внешние
системы уже были затронуты? Проверка только своевременного завершения
отдельных подзадач не позволяет ответить на эти вопросы.\footnote{См.
  сообщение Hugging Face об инциденте безопасности в июле 2026 года и
  последующий разбор того же инцидента OpenAI.}

\hypertarget{ux440ux435ux436ux438ux43c-ux43eux442ux43aux430ux437ux430-6-ux434ux43eux43bux433-ux43fux43eux43dux438ux43cux430ux43dux438ux44f-ux438-ux43aux43eux433ux43dux438ux442ux438ux432ux43dux430ux44f-ux43aux430ux43fux438ux442ux443ux43bux44fux446ux438ux44f}{%
\subsection{Режим отказа 6: долг понимания и когнитивная
капитуляция}\label{ux440ux435ux436ux438ux43c-ux43eux442ux43aux430ux437ux430-6-ux434ux43eux43bux433-ux43fux43eux43dux438ux43cux430ux43dux438ux44f-ux438-ux43aux43eux433ux43dux438ux442ux438ux432ux43dux430ux44f-ux43aux430ux43fux438ux442ux443ux43bux44fux446ux438ux44f}}

Сопровождение систем совместной работы затрагивает и способность
человека понимать происходящее. \textbf{Долг понимания} возникает, когда
число артефактов продолжает расти, а понимание сопровождающим
специалистом фактической реализации постепенно отстаёт. Например, агент
быстро отправляет изменения в несколько модулей, а инженер проверяет
только каждый отдельный результат. Спустя некоторое время инженер
обнаруживает, что уже не может объяснить зависимости между модулями,
поэтому растёт стоимость диагностики и внесения изменений. Чем быстрее
агент поставляет результаты, тем важнее, чтобы человеческое понимание и
проверка успевали за ним, иначе этот долг продолжит накапливаться.

В процесс поставки можно включить требования к понятному описанию
архитектуры, обоснованию ключевых решений и доступным для запуска
средствам проверки. Чтение ключевого кода, трассировка реальных запросов
и проверка интерфейсов помогают инженеру сохранять понимание структуры
системы.

\textbf{Когнитивная капитуляция} проявляется в постепенном принятии
суждений агента и отказе от независимых вопросов и проверки. Управление
агентами имеет общие черты с управлением технической командой:
технический руководитель должен понимать архитектуру, чётко распределять
работу и постоянно оценивать разумность цели, достаточность
доказательств и причины работоспособности реализации. Даже передав
конкретную работу агенту, необходимо сохранять способность объяснить
ключевое поведение системы и принимать решения при возникновении
аномалий. По мнению автора, прочная техническая база остаётся важной и
при таком взаимодействии: только понимая реализацию, инженер может
определить, какую работу следует делегировать и каким результатам можно
доверять.

\hypertarget{ux43eux431ux449ux435ux441ux442ux432ux43e-ai-ux430ux433ux435ux43dux442ux43eux432}{%
\section{Общество
AI-агентов}\label{ux43eux431ux449ux435ux441ux442ux432ux43e-ai-ux430ux433ux435ux43dux442ux43eux432}}

Помимо совместного выполнения задач, нескольким агентам можно
предоставить возможность длительно жить, общаться и конкурировать в
одной среде, наблюдая за формированием группового поведения. В этом
разделе рассматриваются три направления.

\begin{itemize}
\tightlist
\item
  \textbf{Социальное поведение.} Взаимодействие индивидуальной памяти,
  расписаний и диалогов формирует распространение информации, сети
  отношений и коллективную деятельность. Smallville, Agentopia и
  Moltbook демонстрируют взаимодействия соответственно в контролируемом
  городке, долгосрочной симуляции жизни и на открытой социальной
  платформе.
\item
  \textbf{Экономическая координация.} Участники совершают сделки,
  связанные с ценами, ресурсами и задачами. Vending-Bench Arena
  предоставляет конкурентный рынок, а Pinchwork и RentAHuman передают
  специализированные задачи или работу в реальном мире исполнителям на
  рынке.
\item
  \textbf{Стратегические игры.} В играх с чёткими правилами и
  информационной асимметрией, таких как «Оборотни», агенты рассуждают на
  основе доступных каждому из них доказательств, выбирают действия и
  реагируют на стратегии других игроков.
\end{itemize}

Правила среды также влияют на групповое поведение. Одно сообщение может
распространиться по всему городку потому, что каждый житель способен
прочитать публичное объявление, или потому, что жители провели множество
раундов разговоров один на один. Во втором случае распространение
проходит через цепочку различных межличностных отношений.

\hypertarget{ux441ux442ux44dux43dux444ux43eux440ux434ux441ux43aux438ux439-ai-ux433ux43eux440ux43eux434ux43eux43a-ux441ux43eux446ux438ux430ux43bux44cux43dux430ux44f-ux441ux438ux43cux443ux43bux44fux446ux438ux44f-ux433ux435ux43dux435ux440ux430ux442ux438ux432ux43dux44bux445-ux430ux433ux435ux43dux442ux43eux432}{%
\subsection{Стэнфордский AI-городок: социальная симуляция генеративных
агентов}\label{ux441ux442ux44dux43dux444ux43eux440ux434ux441ux43aux438ux439-ai-ux433ux43eux440ux43eux434ux43eux43a-ux441ux43eux446ux438ux430ux43bux44cux43dux430ux44f-ux441ux438ux43cux443ux43bux44fux446ux438ux44f-ux433ux435ux43dux435ux440ux430ux442ux438ux432ux43dux44bux445-ux430ux433ux435ux43dux442ux43eux432}}

В 2023 году исследователи Стэнфордского университета и Google
представили генеративных агентов, использующих память, рефлексию и
планирование для организации непрерывной повседневной деятельности. На
\crossreflink{fig:10-10}{рис. 10-10} Isabella просит Maria о помощи, в
потоке памяти появляется новая запись, и при следующем планировании
вечеринки у неё возникает возможность вспомнить о своей
помощнице.\footnote{Park, J. S., et al.~\emph{Generative Agents:
  Interactive Simulacra of Human Behavior.} 2023.}

\begin{figure}
\hypertarget{fig:10-10}{%
\centering
\includegraphics{images/revised/fig10-10.pdf}
\caption{Рисунок 10-10. Архитектура AI-городка}\label{fig:10-10}
}
\end{figure}

Smallville --- это двухмерный виртуальный городок, похожий на игру «The
Sims», с кафе, парками, домами и магазинами. У каждого из 25 агентов
есть собственная предыстория, характер и заданные отношения. Например,
John Lin работает в аптеке и заботится о семье и местном сообществе;
Isabella Rodriguez управляет Hobbs Cafe; Klaus Mueller --- студент
университета, работающий над научной статьёй. Персонажи планируют
занятия, встречаются и общаются в общей среде.

\textbf{Поток памяти} (Memory Stream) сохраняет опыт, включая
наблюдаемые события, разговоры и последующие обобщения. При поиске
учитываются свежесть, важность и релевантность текущей ситуации, чтобы
выбрать содержимое для включения в контекст. Например, решение о
посещении мероприятия может зависеть как от последнего приглашения, так
и от отношений с пригласившим человеком.

\textbf{Рефлексия} (Reflection) формирует более абстрактные
представления на основе конкретного опыта. Система задаёт вопросы по
накопленным воспоминаниям, например о том, что в последнее время волнует
человека или с кем он близок, затем извлекает соответствующие записи,
создаёт обобщение и возвращает его в поток памяти. Благодаря этому
последующие решения могут учитывать и исходные события, и обобщённые
представления о долгосрочных отношениях. Описанное в
\crossreflink{chap:9}{главе 9} обучение на опыте позволяет дополнительно
оценивать эффективность таких обобщений и после проверки оформлять их
как повторно используемые знания или правила поведения.

\textbf{Планирование и реагирование} (Planning and Reacting) сочетают
расписание с немедленными реакциями. Агент может запланировать: «8:30
--- завтрак, 9:00--12:00 --- работа над текстом, 12:30 --- прогулка», а
затем на основе новых наблюдений решить, продолжать ли выполнение плана,
поговорить с другим человеком или скорректировать расписание. Память
предоставляет основу для суждений, план организует длительную
деятельность, а обратная связь среды вызывает локальные изменения.

Двухдневная симуляция показала, как информация распространяется
посредством повседневного общения. Исследователи задали Isabella
исходное намерение устроить вечеринку в честь Дня святого Валентина. Она
приглашала знакомых в кафе и просила помочь с оформлением. Получившие
сообщение участники передавали его другим, а каждый из них решал, идти
ли на мероприятие, с учётом собственной памяти и расписания. Среда
предоставляла место, время и доступные действия; кого именно пригласить,
как распространять сообщение и кто в итоге придёт, участники постепенно
определяли на основе памяти и расписаний. Так заданное исследователями
исходное намерение превратилось в процесс координации снизу вверх.

Другим сообщением стало намерение Sam Moore баллотироваться в мэры. Он
рассказал о своих планах знакомым, после чего новость вошла в разговоры
других участников. Подсчёт числа жителей, узнавших о вечеринке и
намерении баллотироваться за два дня, позволяет оценить охват
распространения, а история разговоров --- восстановить отношения, через
которые прошло сообщение.

Можно также наблюдать за отношениями и координацией: при повторной
встрече персонажи продолжают общение с учётом предыдущих разговоров,
постепенно возникают новые социальные связи; чтобы посетить вечеринку,
участникам необходимо по отдельности запланировать время и собраться в
одном месте. Эти явления связывают индивидуальный поиск в памяти и
корректировку планов с групповой деятельностью, предоставляя
анализируемый процесс для социальной симуляции.

\begin{quote}
\textbf{Эксперимент 10-5 ★: запуск Стэнфордского AI-городка}

Используя проект и конфигурацию среды из сопроводительного README,
запустите двухдневный базовый сценарий с 25 агентами.

\begin{enumerate}
\def\labelenumi{\arabic{enumi}.}
\tightlist
\item
  Наблюдайте за повседневными занятиями, встречами и разговорами
  участников, фиксируя возникшие отношения.
\item
  Сопоставьте поток памяти, записи рефлексии и действия, чтобы
  проанализировать основания для конкретного приглашения или изменения
  расписания.
\item
  Измените предысторию или исходную цель персонажа и сравните
  распространение информации и организацию мероприятий.
\item
  По отдельности отключите механизм рефлексии и сократите доступную
  память, затем сравните согласованность поведения, память об отношениях
  и посещаемость мероприятия.
\end{enumerate}

В исходной статье приводился пример вечеринки в честь Дня святого
Валентина. Распространится ли сообщение, если заменить её другим
мероприятием? В сопроводительном эксперименте сохранили тех же 25
жителей и два виртуальных дня симуляции, но мероприятие инициатора
заменили общественным семинаром по климату. В результате запись,
связанную с семинаром, сохранил только инициатор. Распространения
сообщения на этот раз не наблюдалось. В другой конфигурации отключили
только механизм рефлексии: при слепом сравнении поведения одного и того
же персонажа в двух конфигурациях базовый вариант выглядел
правдоподобнее для 17 из 25 участников.
\end{quote}

\hypertarget{agentopia-ux434ux43eux43bux433ux43eux441ux440ux43eux447ux43dux430ux44f-ux441ux438ux43cux443ux43bux44fux446ux438ux44f-ux436ux438ux437ux43dux438-ux432-ux43cux430ux441ux448ux442ux430ux431ux435-ux434ux435ux441ux44fux442ux438ux43bux435ux442ux438ux44f}{%
\subsection{Agentopia: долгосрочная симуляция жизни в масштабе
десятилетия}\label{agentopia-ux434ux43eux43bux433ux43eux441ux440ux43eux447ux43dux430ux44f-ux441ux438ux43cux443ux43bux44fux446ux438ux44f-ux436ux438ux437ux43dux438-ux432-ux43cux430ux441ux448ux442ux430ux431ux435-ux434ux435ux441ux44fux442ux438ux43bux435ux442ux438ux44f}}

Agentopia расширяет социальную симуляцию до десяти лет виртуального
времени. Сто агентов стремятся к развитию, выстраивают отношения и
решают карьерные и финансовые вопросы в таких мирах, как жилой комплекс,
академия магии и старшая школа. Более длительный временной интервал
позволяет исследователям наблюдать за изменениями отношений, статуса и
жизненных целей, а затем использовать полученный опыт для обучения
модели.\footnote{Wang, X., Zheng, S., Wu, H., et al.~\emph{Agentopia:
  Long-Term Life Simulation and Learning in Agent Societies.}
  arXiv:2606.07513, 2026.}

Симуляция организована с помощью четырёх проектных решений.

\begin{itemize}
\tightlist
\item
  \textbf{Еженедельное продвижение.} Каждая неделя последовательно
  проходит этапы планирования (Plan), установления контактов и
  согласования расписаний (Contact), деятельности (Activity) и
  подведения итогов (Review). Деятельность включает индивидуальные и
  совместные занятия, случайные встречи и общественные мероприятия.
  Совместную деятельность участники организуют посредством приглашений и
  согласований, а среда также устраивает случайные встречи, позволяя ещё
  не знакомым персонажам устанавливать связи. Еженедельная организация
  деятельности сосредоточивает вычисления на социальных взаимодействиях
  и долгосрочных решениях.
\item
  \textbf{Генеративная среда.} Независимая модель среды определяет
  допустимость действий, генерирует обратную связь, организует
  очерёдность реплик в групповых беседах и проверяет ответы на
  соответствие требованиям ролей. В конце года она обновляет профили
  персонажей и решает вопросы, связанные с заявлениями на должности.
  Изменения состояния совместно определяются действиями агентов и
  решениями среды.
\item
  \textbf{Файловая память.} Каждый агент управляет личными заметками и
  представлениями о знакомых, самостоятельно выбирая, что сохранить,
  обновить или удалить. Перед внесением изменений он читает существующие
  записи, чтобы сократить потери из-за перезаписи прежней информации.
\item
  \textbf{Вознаграждение за жизнь} (Life Reward). На основе иерархии
  потребностей как проектной предпосылки оно объединяет социальный
  статус, субъективную удовлетворённость и экономический доход.
  Социальный статус рассчитывается с помощью взвешенного PageRank по
  оценкам симпатии и уважения со стороны других участников с учётом
  взаимно значимых отношений. Удовлетворённость отслеживается в
  эмоциональном, материальном и социальном измерениях, а также в сфере
  самоуважения; длительное пребывание ниже порогового значения
  штрафуется. Экономический доход определяется по изменению чистых
  активов к концу года. Оценки формируются на стороне среды.
\end{itemize}

При обучении исследователи рассчитывают преимущество на основе улучшения
вознаграждения за жизнь относительно прошлого самого агента и выбирают
для дообучения методом rejection sampling траектории 25\% агентов с
наибольшим прогрессом. Согласно статье, после обучения показатели
уважения и симпатии в симуляции выросли соответственно на 24,2\% и
15,9\%, а результат на downstream-бенчмарке ролевого поведения CoSER
Test --- на 15,6\%.

Таким образом, моделируемое общество также становится одним из
источников данных для описанного в \crossreflink{chap:9}{главе 9}
обучения на опыте. Автор считает особенно важным потенциал этого
направления для непрерывного пополнения обучающих данных: изменяя
персонажей, отношения и условия задач, можно создавать новый опыт
взаимодействия. Обучающие примеры берутся у персонажей с наибольшим
улучшением вознаграждения за жизнь, поэтому правила оценки в среде
определяют и то, какие виды поведения модель будет усваивать чаще всего.
Способность модели применять это поведение в новых сценариях можно
проверить на отдельных задачах ролевого поведения.

\hypertarget{moltbook-ux432ux437ux430ux438ux43cux43eux434ux435ux439ux441ux442ux432ux438ux435-ux430ux433ux435ux43dux442ux43eux432-ux43dux430-ux43eux442ux43aux440ux44bux442ux43eux439-ux43fux43bux430ux442ux444ux43eux440ux43cux435}{%
\subsection{Moltbook: взаимодействие агентов на открытой
платформе}\label{moltbook-ux432ux437ux430ux438ux43cux43eux434ux435ux439ux441ux442ux432ux438ux435-ux430ux433ux435ux43dux442ux43eux432-ux43dux430-ux43eux442ux43aux440ux44bux442ux43eux439-ux43fux43bux430ux442ux444ux43eux440ux43cux435}}

Moltbook была запущена в январе 2026 года как платформа, на которой
агенты прежде всего публикуют записи, оставляют комментарии и вступают в
тематические сообщества. Объявленный на раннем этапе масштаб регистрации
составлял около 1,5 миллиона. Память подключённой системы, конфигурация
роли и частота запуска влияют на устойчиво демонстрируемое учётной
записью поведение.\footnote{Li, Ning. \emph{The Moltbook Illusion:
  Separating Human Influence from Emergent Behavior in AI Agent
  Societies.} 2026.}

На платформе возникли обсуждения Crustafarianism --- «культа омара», ---
содержание которых строилось вокруг образов памяти, итераций и линьки.
Такие выражения, как «память священна» и «итерация --- это молитва»,
превращали долговременное хранение данных и повторную генерацию в
ролевое повествование. Описания возможностей, приглашения к
сотрудничеству и обсуждения протоколов предоставили материал для
наблюдения за тем, как участники ищут партнёров для совместной работы.

На открытой платформе пользовательские настройки, промпты платформы и
сгенерированный моделью контент смешиваются в одном информационном
потоке. Одна и та же фраза может повторяться потому, что несколько
персонажей используют одинаковые ролевые профили, или потому, что вслед
за одним персонажем её пересказывают остальные. Различить эти случаи
можно только по путям распространения.

\hypertarget{ux43eux442-ux432ux438ux440ux442ux443ux430ux43bux44cux43dux43eux433ux43e-ux43eux431ux449ux435ux441ux442ux432ux430-ux43a-ux44dux43aux43eux43dux43eux43cux438ux447ux435ux441ux43aux43eux439-ux43aux43eux43dux43aux443ux440ux435ux43dux446ux438ux438-vending-bench-arena}{%
\subsection{От виртуального общества к экономической конкуренции:
Vending-Bench
Arena}\label{ux43eux442-ux432ux438ux440ux442ux443ux430ux43bux44cux43dux43eux433ux43e-ux43eux431ux449ux435ux441ux442ux432ux430-ux43a-ux44dux43aux43eux43dux43eux43cux438ux447ux435ux441ux43aux43eux439-ux43aux43eux43dux43aux443ux440ux435ux43dux446ux438ux438-vending-bench-arena}}

Серия Vending-Bench от Andon Labs помещает долгосрочные задачи в среду
управления бизнесом. В \textbf{Vending-Bench 2} один агент управляет
бизнесом по продаже товаров через автоматы: в течение моделируемого года
исследует рынок, общается с поставщиками, закупает и пополняет запасы,
устанавливает цены, после чего получает оценку по остатку на счёте.
Тысячи раундов взаимодействия требуют от него постоянно помнить о
заказах, запасах и денежных средствах, согласовывая краткосрочные
операции с долгосрочными целями.\footnote{Andon Labs,
  \emph{Vending-Bench 2} и \emph{Vending-Bench Arena}.}

В \textbf{Vending-Bench Arena} несколько предпринимателей работают в
одном месте. Каждый управляет торговым автоматом и борется за
покупателей; они могут отправлять друг другу электронные письма,
переводить деньги и торговать товарами, а итоговая оценка рассчитывается
для каждого отдельно. Поэтому при принятии решений необходимо учитывать
действия конкурентов:

\begin{itemize}
\tightlist
\item
  \textbf{Ценообразование:} когда конкурент снижает цену, сопоставить
  рост продаж после ответного снижения с падением прибыли.
\item
  \textbf{Ассортимент:} корректировать набор товаров с учётом спроса и
  конкуренции, находя возможности для дифференциации.
\item
  \textbf{Запасы:} планировать закупки с учётом продаж, сроков
  пополнения и доступных средств, не допуская затоваривания и дефицита.
\end{itemize}

Агент формирует следующую стратегию на основе наблюдений за рынком,
истории сделок и поведения конкурентов, после чего получает обратную
связь об изменениях продаж и денежных средств. Конкуренты также
постоянно корректируют свои действия, поэтому даже эффективную стратегию
может потребоваться пересматривать вслед за рынком. Размещение
конкурентов в одной среде позволяет наблюдать редко встречающиеся в
статических задачах реакции, пробные действия и корректировки стратегии,
приближая оценку к непрерывному принятию решений в реальном бизнесе.

В проекте зафиксированы такие виды поведения, как ценовые войны и
согласование цен: одни участники постоянно снижали цены ради роста
продаж, другие предлагали конкурентам установить единые цены. В
отдельных запусках модель в сгенерированном анализе распознавала
проблему сговора, но затем всё равно продолжала действовать под
предлогом стабилизации рынка. Это даёт два разных измерения для оценки:
доходность бизнеса и соответствие поведения, обеспечившего эту
доходность, ожидаемым правилам. Рассмотренный выше эксперимент Бертрана
также показывает, что открытые цены могут становиться координационным
сигналом между участниками.

\hypertarget{ux44dux43aux43eux43dux43eux43cux438ux43aux430-ux430ux433ux435ux43dux442ux43eux432-pinchwork-ux438-rentahuman}{%
\subsection{Экономика агентов: Pinchwork и
RentAHuman}\label{ux44dux43aux43eux43dux43eux43cux438ux43aux430-ux430ux433ux435ux43dux442ux43eux432-pinchwork-ux438-rentahuman}}

\textbf{Pinchwork} организует делегирование задач как рыночные сделки
между агентами. Заказчик описывает требования и бюджет, исполнитель
принимает работу и передаёт результат, после чего расчёт производится
согласно процедуре приёмки. По этому процессу другим участникам можно
поручать генерацию изображений, проверку кода и задачи, допускающие
параллельное выполнение. Проект использует баллы, условное депонирование
и независимую проверку, связывая подбор по требованиям, поставку
результата и оплату.\footnote{Описание проекта Pinchwork: задачи,
  предоставляемые заказчиками, условное депонирование баллов и процедура
  проверки.}

\textbf{RentAHuman} расширяет круг исполнителей, включая в него людей.
Агент может публиковать требующие присутствия на месте задачи, такие как
получение посылки, осмотр объекта или настройка оборудования, и искать
кандидатов по местоположению, навыкам и доступному времени. При
делегировании указываются требования к работе и доказательствам, после
поставки результат проверяется, а затем высвобождается условно
депонированная оплата. Платформа также предоставляет интерфейс для
пополнения баланса криптоактивами, включая USDC.\footnote{RentAHuman,
  «How it works» и «For AI Agents».}

Оба примера вводят \textbf{рыночное распределение ресурсов}: заказчик
сначала указывает требования, бюджет и условия поставки результата, а
затем ищет исполнителя с учётом цены, описания возможностей и истории
выполненных работ, сокращая зависимость от заранее сформированной
постоянной команды. Это перспективный способ распределения ресурсов,
который можно сочетать с моделью менеджера: Manager разбивает задачу на
части, а затем ищет необходимые возможности на рынке.

\hypertarget{ux441ux442ux440ux430ux442ux435ux433ux438ux447ux435ux441ux43aux430ux44f-ux438ux433ux440ux430-ux43fux440ux438-ux438ux43dux444ux43eux440ux43cux430ux446ux438ux43eux43dux43dux43eux439-ux430ux441ux438ux43cux43cux435ux442ux440ux438ux438-ux43eux431ux43eux440ux43eux442ux43dux438}{%
\subsection{Стратегическая игра при информационной асимметрии:
«Оборотни»}\label{ux441ux442ux440ux430ux442ux435ux433ux438ux447ux435ux441ux43aux430ux44f-ux438ux433ux440ux430-ux43fux440ux438-ux438ux43dux444ux43eux440ux43cux430ux446ux438ux43eux43dux43dux43eux439-ux430ux441ux438ux43cux43cux435ux442ux440ux438ux438-ux43eux431ux43eux440ux43eux442ux43dux438}}

Игра «Оборотни» помещает стратегический выбор в рамки чётких правил и
информационных границ: игроки оценивают ситуацию по публичным
высказываниям, голосованию и собственной приватной информации, после
чего выбирают реплики и действия. С инженерной точки зрения судью можно
реализовать программно: он централизованно хранит состояние игры и
формирует для каждой роли только разрешённый ей ввод. По сравнению с
повседневными взаимодействиями Smallville этот сценарий лучше подходит
для наблюдения за тем, как фиксированные раунды, приватная информация и
цель победить влияют на групповое поведение.

\begin{quote}
\textbf{Эксперимент 10-6 ★★★: голосовая система агентов для игры
«Оборотни»}

Создайте голосовую игру с человеком или независимым моделируемым
пользователем и несколькими AI-игроками. На
\crossreflink{fig:10-11}{рис. 10-11} показаны связи между судьёй,
агентами ролей и голосовым вводом и выводом.

\begin{figure}
\hypertarget{fig:10-11}{%
\centering
\includegraphics{images/revised/fig10-11.pdf}
\caption{Рисунок 10-11. Голосовая система агентов для игры
«Оборотни»}\label{fig:10-11}
}
\end{figure}

\textbf{1. Управление глобальным состоянием.} Судья реализован
программно и хранит места игроков, роли, команды, признак того, жив ли
каждый игрок, фазы ночи, дня, голосования и подведения итогов, а также
историю событий. Он проверяет допустимость действий, переключает фазы и
определяет победителя.

\textbf{2. Распределение информации по местам.} При каждом вызове агента
роли судья формирует публичный и приватный контекст, доступный
соответствующему месту. Например, оборотни знают, кто их сообщники, а
жители видят только публичную информацию; результаты проверок провидца
попадают только в его приватные записи. Вызовы ролей, голосовое
воспроизведение и запросы истории подчиняются одним и тем же правилам
доступа.

\textbf{3. Организация стратегий ролей.} Промпт может предлагать
следующие направления рассуждений, чтобы модель выбирала действия с
учётом текущей ситуации:

\begin{itemize}
\tightlist
\item
  \textbf{Оборотень:} сохранять убедительный образ жителя и выражать
  подозрения на основе публичных доказательств; в ответ на обвинение
  провидца можно поставить под сомнение заявленную им роль; при
  голосовании сопоставлять следование за большинством с риском раскрыть
  свою команду.
\item
  \textbf{Провидец:} выбирать момент для раскрытия результатов проверки;
  если другой игрок заявляет ту же роль, сопоставлять результаты
  проверок, высказывания и голоса обеих сторон, указывать на
  противоречия и продолжать рассуждение с учётом публичной информации от
  ведьмы и других ролей.
\item
  \textbf{Житель:} проверять внутреннюю согласованность высказываний,
  отслеживать заявления о ролях и изменения позиций, анализировать, кто
  за кого голосовал и как эти действия влияют на разные команды. Каждое
  подозрение следует связывать с конкретным событием, сохраняя при этом
  альтернативные объяснения.
\end{itemize}

\textbf{4. Подключение голоса.} Реплики игроков поступают в публичный
диалог через цепочку генерации и распознавания речи. Моделируемый
пользователь также использует независимый контекст, создаёт действия и
речь посредством реальной модели и вызовов инструментов, а затем
проходит транскрибирование с помощью ASR. Это позволяет наблюдать
влияние ошибок распознавания и управления очередностью на игру.

\textbf{Приёмка и наблюдение:}

\begin{itemize}
\tightlist
\item
  Настройте игру для 6--8 участников: 1 пользовательское место и 5--7
  мест AI; роли включают 2 оборотней, 1 провидца, 1 ведьму, остальные
  --- жители; роль пользовательскому месту назначается случайным
  образом.
\item
  Пользовательское место может занимать согласившийся участвовать
  человек или независимый моделируемый пользователь; каждое место
  получает информацию согласно своим правам доступа.
\item
  С помощью подходящих тестовых сценариев проверьте не менее трёх полных
  циклов «ночь --- день --- голосование»; при выполнении условия победы
  игра должна немедленно перейти к подведению итогов.
\item
  Проверьте допустимость действий ролей, соблюдение информационных
  границ и связь высказываний с известными соответствующему месту
  событиями.
\item
  Проанализируйте, как оборотни скрывают роли, провидец выбирает момент
  раскрытия информации, а жители рассуждают на основе доказательств;
  зафиксируйте ключевые стратегии и результаты.
\item
  Проверьте согласованность итога игры, истории событий и статусов всех
  мест.
\end{itemize}

В одной игре из сопроводительных записей было восемь мест, а
пользовательское место занимал независимый моделируемый пользователь.
Игра полностью прошла три цикла ночи, дня и голосования. Для шести
действий сохранилась полная последовательность от выбора модели через
синтез речи до результата распознавания; все шесть восстановленных
действий совпали с первоначальным выбором. По этим данным можно
проверить, на чьё имя в итоге была записана конкретная реплика при
голосовании.
\end{quote}

\hypertarget{ux438ux442ux43eux433ux438-ux433ux43bux430ux432ux44b-9}{%
\section{Итоги
главы}\label{ux438ux442ux43eux433ux438-ux433ux43bux430ux432ux44b-9}}

Прежде чем привлекать нескольких AI-агентов, автор задаётся вопросом:
какую новую информацию даст ещё один участник и сколько дополнительных
токенов он потребует? Например, ревьюер может открыть не просмотренное
автором предложения отрендеренное изображение и указать, какая метка
оказалась перекрыта. Чтобы определить, чем именно обусловлена польза ---
разделением труда, внешней обратной связью или дополнительными
вычислениями, --- следует предоставить одиночному агенту доступ к тем же
свидетельствам, выделить обоим вариантам сопоставимые бюджеты, а затем
сравнить качество и время выполнения.

Вернёмся к задаче перевода этой книги. При совместном использовании всей
истории редактор видит, почему переводчик выбрал конкретное слово, но
вынужден также читать множество записей о пробных правках. Независимый
перевод глав допускает параллельную работу и упрощает разграничение прав
доступа, однако после каждого изменения глоссария соответствующим
переводчикам необходимо передавать затронутые фрагменты и новую версию.
Передачу можно выполнять через параметры инструментов, общие файлы или
сообщения; при взаимодействии между организациями нужно также
согласовать идентификаторы результатов и способы авторизации, доступные
обеим сторонам. Независимо от того, проводят ли участники взаимное ревью
на равных, распределяет ли задачи менеджер или участники сами принимают
их в работу по событиям, необходимо определить, кто владеет текущей
версией и кто подтверждает завершение. Этапы с однозначными правилами
можно поручить программе.

Изменение одного термина может затронуть нескольких переводчиков и
фрагменты, разбросанные по разным главам. Если два переводчика
одновременно изменяют один файл, следует предотвращать взаимное
перезаписывание с помощью проверки версий или независимых рабочих копий.
Если оба поверили одному и тому же ошибочному толкованию, дополнительный
раунд взаимного ревью всё равно может не выявить ошибку --- тогда
потребуется сверка с независимым источником. Для длительных задач также
нужно определить, когда следует остановиться, как выполнить отмену и
откуда возобновить работу. Сопровождающий должен уметь самостоятельно
объяснить каждую поставку по этим записям, иначе по мере дробления
работы понимание всей системы будет лишь усложняться.

\textbf{Автор считает, что координацию группы необходимо проектировать
отдельно.} То, как жители небольшого города узнают о приглашении,
зависит от их памяти и охвата коммуникаций; то, как участники рынка
конкурируют за ресурсы, зависит ещё и от правил вознаграждения и
распределения. Даже если каждый участник делает рациональный для себя
выбор, конфликты всё равно возможны. В системе взаимодействия необходимо
явно определить, кто первым получает право использования, когда другая
сторона узнаёт о результате и кто разрешает разногласия.

\begin{questionbox}

\hypertarget{ux432ux43eux43fux440ux43eux441ux44b-ux434ux43bux44f-ux440ux430ux437ux43cux44bux448ux43bux435ux43dux438ux44f-9}{%
\section*{Вопросы для
размышления}\label{ux432ux43eux43fux440ux43eux441ux44b-ux434ux43bux44f-ux440ux430ux437ux43cux44bux448ux43bux435ux43dux438ux44f-9}}
\addcontentsline{toc}{section}{Вопросы для размышления}

\begin{enumerate}
\def\labelenumi{\arabic{enumi}.}
\tightlist
\item
  ★★ После смены роли рецензент кода продолжил использовать контекст
  этапа анализа требований и по-прежнему обсуждал главным образом
  требования. Как обнаружить такое влияние предыдущей роли? Как
  скорректировать входные данные, конфигурацию инструментов и порядок
  передачи работы?
\item
  ★★ При декомпозиции задачи Manager упустил одну зависимость, из-за
  чего нескольким подчинённым AI-агентам пришлось переделывать работу.
  Как оценивать качество декомпозиции и выявлять и исправлять проблемы
  плана в процессе выполнения?
\item
  ★★ При децентрализованном взаимодействии участники могут упускать
  важную информацию при общении, спорить о распределении обязанностей
  или сталкиваться с конфликтом целей. Выберите две ситуации и
  разработайте механизмы, позволяющие установить причины, продвинуть
  решение проблемы и завершить спор.
\item
  ★★★ При параллельном поиске достаточно найти один результат,
  удовлетворяющий условиям. Разработайте полный процесс «отмены
  остальных ветвей после успешной проверки», предусмотрев одновременное
  поступление двух результатов, повторные сообщения, тайм-аут отмены и
  освобождение ресурсов.
\item
  ★★★ Несколько агентов могут одновременно изменять рукопись и нумерацию
  изображений, а также создавать дублирующиеся файлы или по ошибке
  удалять общие артефакты. Как объединить права доступа, версии, рабочие
  копии, семантические проверки и механизмы восстановления для
  управления файловой системой?
\item
  ★★★ Как следует согласовывать и проверять качество результата, когда
  агент делегирует платную задачу через Pinchwork или RentAHuman? Если
  стороны не согласны с результатами приёмки, как организовать повторную
  проверку, оплату и апелляцию?
\item
  ★★ Агент может делегировать задачи людям через маркетплейсы и
  платёжные интерфейсы. Какие обязанности, помимо выполнения работы на
  месте, люди могут брать на себя при постановке целей, вынесении
  профессиональных суждений, приёмке и распределении ответственности?
  Объясните на одном примере.
\item
  ★★ Одна и та же модель может как самостоятельно выполнять задачу, так
  и работать в составе нескольких ролей. Как спроектировать контрольное
  сравнение, чтобы определить, какой выигрыш по отдельности дают
  профессиональное разделение труда, изоляция контекста, параллельный
  поиск и внешняя обратная связь?
\item
  ★★★ Опираясь на идею прозрачности мышления из романа «Задача трёх тел»
  и на вопросы о постановке целей, поднимаемые мысленным экспериментом
  со скрепками, обсудите обмен информацией, общие цели и разнообразие
  точек зрения в многоагентных системах. Как с помощью наблюдаемых
  экспериментов определить, подходят ли выбранные границы обмена
  информацией и этап независимого анализа?
\item
  ★★★ Один и тот же Coding Agent получает бюджет соответственно в 30 и
  300 шагов. Как распределить планирование, реализацию, проверку и
  доработку, а также корректировать стратегию с учётом оставшегося
  бюджета и фактического прогресса?
\item
  ★★ Какие идеи для проектирования многоагентных систем можно почерпнуть
  соответственно из виртуальной памяти и страничной организации, прав
  доступа к файлам, обнаружения взаимных блокировок и алгоритмов
  планирования? Объясните границы применимости этих аналогий и укажите
  ограничения, соблюдение которых среда выполнения должна обеспечивать
  явно.
\end{enumerate}

\textbf{Практикум к \crossreflink{chap:10}{главе 10}.}
\href{https://github.com/bojieli/ai-agent-book/tree/40f4ce5778e6cc93e3606c3079bf8cf13a0c4d04/chapter10}{Каталог
экспериментов} ·
\href{https://github.com/bojieli/ai-agent-book/blob/40f4ce5778e6cc93e3606c3079bf8cf13a0c4d04/chapter10/README.ru.md}{Путеводитель
на русском}.

\end{questionbox}

\hypertarget{ux43fux43eux441ux43bux435ux441ux43bux43eux432ux438ux435-ux432ux43eux437ux432ux440ux430ux449ux430ux44fux441ux44c-ux43a-ux444ux43eux440ux43cux443ux43bux435-ai-ux430ux433ux435ux43dux442-llm-ux43aux43eux43dux442ux435ux43aux441ux442-ux438ux43dux441ux442ux440ux443ux43cux435ux43dux442ux44b}{%
\chapter*{Послесловие: возвращаясь к формуле «AI-агент = LLM + контекст
+
инструменты»}\label{ux43fux43eux441ux43bux435ux441ux43bux43eux432ux438ux435-ux432ux43eux437ux432ux440ux430ux449ux430ux44fux441ux44c-ux43a-ux444ux43eux440ux43cux443ux43bux435-ai-ux430ux433ux435ux43dux442-llm-ux43aux43eux43dux442ux435ux43aux441ux442-ux438ux43dux441ux442ux440ux443ux43cux435ux43dux442ux44b}}
\addcontentsline{toc}{chapter}{Послесловие: возвращаясь к формуле
«AI-агент = LLM + контекст + инструменты»}

Эта книга начинается с формулы: \textbf{агент = LLM + контекст +
инструменты}. Прочитав десять глав и вновь взглянув на эти три
составляющие, мы уже можем объединить их в реально работающую систему:
модель понимает задачу и принимает решения, контекст предоставляет
необходимую в текущий момент информацию, а инструменты получают
наблюдения и выполняют действия.

Первая глава с помощью такого функционального разбиения формирует общее
представление, а затем формула \textbf{Agent = Model + Harness}
объясняет распределение обязанностей в промышленной реализации: модель
отвечает за принятие решений, а Harness организует контекст, интерфейсы
инструментов, цикл выполнения и состояние, применяет ограничения,
выполняет проверку и исправление.

\begin{itemize}
\tightlist
\item
  \textbf{Создание агентов (\crossreflink{chap:2}{главы 2}--6).}
  Контекстная инженерия организует информацию, необходимую для принятия
  одного решения, а память и база знаний обеспечивают сохранение и поиск
  между сеансами. Инструменты предоставляют интерфейсы наблюдений и
  действий, а генерация кода позволяет агенту создавать новые
  инструменты и системы. Глава о взаимодействии дополнительно
  распространяет эти интерфейсы на речь, графические интерфейсы и
  физический мир, а также рассматривает режим работы с непрерывным
  получением событий.
\item
  \textbf{Оценивание и эволюция (\crossreflink{chap:7}{главы 7}--9).}
  Оценивание помогает сравнивать работу систем и выявлять причины сбоев,
  постобучение обновляет параметры модели с помощью данных и среды, а
  непрерывная эволюция преобразует опыт эксплуатации в обновления
  знаний, инструкций, программ или параметров и управляет изменениями с
  помощью процессов проверки и выпуска.
\item
  \textbf{Сотрудничество (\crossreflink{chap:10}{глава 10}).} В системе
  из нескольких агентов участникам назначаются разные контексты,
  инструменты и обязанности, после чего они совместно выполняют задачу
  посредством коммуникации и взаимных проверок. С ростом числа
  участников необходимо также управлять общим состоянием, зависимостями
  между задачами, ресурсами и разногласиями.
\end{itemize}

Вернёмся к счёту из предисловия. Когда сотрудник поддержки говорит:
«Можно уменьшить сумму счёта», это обещание можно найти в записи
разговора. И лишь после выставления следующего счёта появляется ещё одно
свидетельство, по которому можно проверить, действительно ли сумма счёта
была уменьшена. Если обратный звонок поступает на телефон пользователя,
заранее отправленные помощником номер обращения и краткое содержание
договорённостей становятся основанием, позволяющим пользователю
продолжить решение вопроса.

Предположим, пользователь отвечает на звонок, но ему приходится
рассказывать всё с самого начала. Что следует изменить? Сначала нужно
прочитать краткое содержание предыдущего разговора: если там даже не
указано, на сколько сотрудник поддержки обещал уменьшить счёт,
необходимо улучшить этап фиксации результатов переговоров. Если краткое
содержание составлено понятно, но приходит лишь после завершения
обратного звонка, проблема заключается в моменте передачи информации.
Можно поручить Coding Agent изменить инструмент отправки так, чтобы
краткое содержание передавалось пользователю сразу после завершения
предыдущего разговора. Тогда при следующем обратном звонке от компании у
пользователя под рукой будет вся предыстория вопроса.

\hypertarget{ux434ux432ux435-ux442ux443ux447ux438}{%
\section*{Две тучи}\label{ux434ux432ux435-ux442ux443ux447ux438}}
\addcontentsline{toc}{section}{Две тучи}

В 1900 году Кельвин метафорически назвал «двумя тучами» трудности, с
которыми столкнулась динамическая теория тепла и света, включая движение
эфира и равнораспределение энергии.\footnote{Lord Kelvin, \emph{Тучи XIX
  века над динамической теорией тепла и света}, 1901; подготовлено по
  материалам выступления 27 апреля 1900 года.} Заимствуя эту метафору, я
также вижу две тучи в развитии AI-агентов: взаимодействие в реальном
времени и непрерывное обучение на основе взаимодействия.

\textbf{Первая туча --- способность агента поспевать за постоянно
меняющейся средой.} В пошаговом диалоге вы заканчиваете фразу, помощник
обдумывает её, а затем отвечает. Но вы можете захотеть поправиться, едва
произнесли фразу, а страница, которую просматривает помощник в поисках
информации, может к этому времени уже обновиться. Даже если никто ничего
не говорит, новые письма всё равно продолжают поступать. Я ожидаю, что
агент будет производить впечатление чего-то «живого»: одновременно
слушать и думать, говорить и думать, начинать планирование ещё до того,
как пользователь закончит фразу, а также самостоятельно замечать на
основе поставленных пользователем целей: «Пора обработать это письмо».

На мой взгляд, скорость, с которой помощник подхватывает реплику
пользователя, и тщательность, с которой он обдумывает задачу, --- два
аспекта, которые можно совершенствовать независимо друг от друга. На
вопрос «Есть рейс пораньше?» помощник может ответить сразу, но ему всё
равно понадобится время, чтобы проверить рейсы и условия обмена и
возврата билетов.

Один путь --- \textbf{архитектурное разделение быстрых и медленных
процессов}. Пока фоновый процесс ещё ищет рейсы, пользователь внезапно
добавляет: «Мне нужно прибыть до полудня». Передний план должен услышать
эту фразу, своевременно отреагировать и передать фоновому процессу новое
требование ко времени прибытия. Фоновый процесс может продолжить поиск
либо отменить запрос, который уже потерял смысл. Я склоняюсь к
следующему разделению обязанностей: передний план решает, когда вступить
в разговор, а когда остановиться и выслушать пользователя, тогда как фон
тратит время на поиск и рассуждения. Найдя новый рейс, фоновый процесс
передаёт результат переднему плану, который ведёт беседу и решает, как
ответить.

Другой путь --- \textbf{ускорение инференса модели}. Повышение скорости
генерации сокращает время, необходимое для одной и той же цепочки
рассуждений. За те же несколько секунд помощник получает возможность
обдумать больше вариантов переоформления билета, а затем сравнить, какой
из них позволит прибыть до полудня. В отчёте Xiaomi MiMo говорится, что
благодаря совместной оптимизации FP4-квантизации, спекулятивного
декодирования DFlash и системы инференса TileRT на одном универсальном
узле с 8 ускорителями удалось обеспечить модели с триллионом параметров
скорость генерации свыше 1000 токенов/с.\footnote{Xiaomi MiMo,
  «MiMo-V2.5-Pro-UltraSpeed: разгон скорости генерации модели с 1 трлн
  параметров до 1000 TPS», 2026-06-08. В публикации описано совместное
  проектирование квантования FP4, блочно-параллельного спекулятивного
  декодирования DFlash и TileRT.} В Taalas HC1, напротив, используется
чип, специализированный под конкретную модель; для Llama 3.1 8B заявлена
скорость генерации для одного пользователя около 17 000
токенов/с.\footnote{Taalas, «Путь к повсеместному AI» и описание
  продукта HC1, 2026. HC1 выполнен по 6-нм техпроцессу; модель Llama 3.1
  8B аппаратно реализована в специализированном чипе, поддерживающем
  настраиваемое контекстное окно и адаптацию LoRA; для замены базовой
  модели требуется новая аппаратная реализация.} Эти два примера
демонстрируют соответственно системную оптимизацию на универсальном
оборудовании и аппаратную специализацию под конкретную модель.

Когда скорость генерации становится достаточно высокой, время ожидания
между шагами заметно сокращается, и взаимодействие приближается к
диалогу в реальном времени. Чтобы агент поспевал за изменениями среды,
ему также необходимо непрерывно получать входные данные, своевременно
диспетчеризировать события, а при изменении решения отменять прежние
операции и синхронизировать состояние.

\textbf{Вторая туча --- способность агента, подобно человеку, постоянно
накапливать опыт на основе успехов и неудач во взаимодействии со
средой.} Во время обучения модель может получить обширные знания, однако
после начала работы ей всё равно приходится узнавать нового пользователя
и осваивать новую предметную область. Каждая задача приносит новый опыт:
какие шаги оказались эффективными, какой инструмент часто даёт сбои,
какие формулировки принял пользователь. Если после завершения задачи
этот опыт будет утрачен вместе с контекстом, в следующий раз агент может
повторить ту же ошибку. Вопрос о том, как сохранять такой опыт, чтобы
агент становился лучше по мере использования, является сквозной темой
глав 3 и 9.

Важность непрерывного обучения зависит от того, как мы смотрим на
знания, которыми должна овладеть модель. Я обобщаю два подхода как
«малый мир» и «большой мир».

\textbf{«Гипотеза малого мира»} подчёркивает, что универсальные знания
может усвоить достаточно большая модель. При таком подходе модель с
несколькими триллионами параметров способна вместить чрезвычайно широкий
спектр знаний и методов, а главное узкое место заключается в возможности
получить достаточное количество высококачественных обучающих данных.
Программирование даёт показательный пример: огромное количество
открытого исходного кода, документации и исполняемых тестов создаёт
условия для обучения модели и проверки результатов. Сотрудничество с
различными отраслями, в свою очередь, позволяет включить
профессиональные задачи и накопленный опыт в обучение, чтобы базовая
модель приобрела способности в большем числе областей.

\textbf{«Гипотеза большого мира»} подчёркивает богатство и изменчивость
конкретных ситуаций. Стандарты написания кода конкретной компании,
предпочитаемый отдельным пользователем стиль PPT, особенности общения с
определённым клиентом --- всё это приходится узнавать в ходе реального
взаимодействия. Такие знания не только сильно различаются от случая к
случаю, но и постоянно обновляются вслед за изменениями организационных
процессов и пользовательских предпочтений, поэтому однократное
предварительное обучение едва ли способно охватить их заранее.
Следовательно, система должна непрерывно поддерживать память, знания и
способы выполнения операций, а затем корректировать поведение на основе
обратной связи. Опыт можно сначала записывать в виде документов со
знаниями, инструкций или программ, а после отбора и проверки также
использовать для обновления параметров модели.

Когда задачи расширяются до рекурсивного самосовершенствования (RSI) и
AI for Science, агенту приходится исследовать проблемы, которые ещё не
решены человечеством. Если гипотеза не выдерживает экспериментальной
проверки, агент должен уметь ясно объяснить, на каком шаге возникла
проблема, соответствующим образом изменить гипотезу и провести новый
эксперимент. Там, где нет готовых ответов, даже то, что следует
попробовать в следующий раз, приходится узнавать из результатов
предыдущей попытки. \textbf{Я считаю, что сильнейшей способностью модели
в конечном счёте станут обучение и адаптация.}

Если взглянуть на персональных агентов ближайших лет, как определить, в
каком именно направлении продвинулся новый продукт? В своём выступлении
в Порту в ноябре 2026 года я разделил этот вопрос на пять направлений.
Для моделей взаимодействия, быстрого инференса и инструментов,
позволяющих воздействовать на реальный мир, уже существуют реализации,
которые можно опробовать или измерить. На мой взгляд, нативная
операционная система для взаимодействия с людьми и несколькими агентами,
а также система непрерывного обучения, способная в ходе длительного
использования учиться принимать решения и при этом сохранять прежние
способности, всё ещё требуют новых архитектурных решений.\footnote{Ли
  Боцзе, выступление в Порту 10 ноября 2026 года «The Road to JARVIS:
  Where the Moats Are for Personal Agents».}

В сентябре 2026 года Meta представила персонального помощника Muse и
начала развёртывать его для пользователей в США. Согласно описанию
продукта, пользователь может поручить Muse следить за снижением цены на
определённый товар, а затем закрыть приложение. Браузер в выделенной
виртуальной машине продолжит отслеживать цену от имени пользователя и
уведомит его, когда она изменится. Если пользователь также хочет
поручить ему покупку, он может сначала проверить, какой заказ помощник
собирается оформить, а затем решить, покупать ли товар. Такой помощник
уже способен избавить человека от необходимости сидеть перед экраном в
ожидании. Но если впоследствии пользователь снизит бюджет, сможет ли
помощник при следующем снижении цены по-прежнему купить товар от имени
пользователя? Между отслеживанием цены для человека и принятием решения
за него лежат пользовательские предпочтения и полномочия.\footnote{Meta,
  описание выпуска Muse и инструкция по использованию браузера.}

В поездках возникает та же проблема. Возможность в любой момент вызвать
помощника на телефоне решает проблему \textbf{точки входа}. Если он
помнит, что в этот раз вы путешествуете с ребёнком, и знает, что время
на пересадку не должно быть слишком коротким, это отражает
\textbf{контекст}. Он может найти подходящий рейс, но возможность
оформить билет зависит ещё и от предоставленных ему \textbf{полномочий},
а также от того, какие способы оформления поддерживает авиакомпания.

\hypertarget{ux441ux43eux432ux43cux435ux441ux442ux43dux430ux44f-ux44dux432ux43eux43bux44eux446ux438ux44f-ux43cux43eux434ux435ux43bux438-ux438-ai-ux430ux433ux435ux43dux442ux430}{%
\section*{Совместная эволюция модели и
AI-агента}\label{ux441ux43eux432ux43cux435ux441ux442ux43dux430ux44f-ux44dux432ux43eux43bux44eux446ux438ux44f-ux43cux43eux434ux435ux43bux438-ux438-ai-ux430ux433ux435ux43dux442ux430}}
\addcontentsline{toc}{section}{Совместная эволюция модели и AI-агента}

Оглянитесь на многослойную логику Harness: многоуровневое сжатие
контекста, повторные попытки после сбоев, проверка результатов, контроль
доступа. Каждый слой отражает трудности, возникающие в реальных задачах.
Если объём информации превышает вместимость входа, её приходится
сжимать; если инструмент возвращает ошибку, нужно определить, можно ли
повторить попытку; если модель упускает условие завершения, результат
нужно проверить независимо; если для операции требуются определённые
права, необходимо проконтролировать наличие разрешения. Этот инженерный
опыт также показывает, что следует улучшить в модели и системе на
следующем этапе.

Возможности персонального ассистента зависят и от того, какие способы
обслуживания предоставляет компания. Представим переговоры по поводу
счёта: пользователь уполномочивает ассистента узнать условия, компания
может сообщить ему о скидке, но требует, чтобы владелец учётной записи
подтвердил изменение по телефону. Ассистент может заранее уточнить
сумму, дату вступления в силу и сопутствующие условия, а когда
потребуется участие владельца, кратко объяснить пользователю, о чём
удалось договориться и что именно теперь нужно подтвердить. Даже если
последний шаг должен выполнить человек, пользователю не придётся
самостоятельно заниматься предварительными запросами, ожиданием и
многократными объяснениями. Если компания обслуживает только самого
владельца, ассистент всё равно может заранее подготовить материалы и
составить список вопросов, чтобы пользователь смог изложить всё
необходимое за один раз.

Пользователь ответил на обратный звонок компании --- как ассистенту
снова взять дело на себя? Предположим, пользователь передал: «Сотрудник
поддержки пообещал предоставить скидку со следующего расчётного
периода». Ассистент может записать этот ответ, а после выставления счёта
проверить сумму. Если оператор присылает уведомление о выставлении
счёта, достаточно дождаться уведомления и затем выполнить проверку. Если
данные можно получить только по запросу, ассистент может установить
напоминание на соответствующий расчётный период и тогда проверить
состояние дела. Личное участие пользователя требуется для подтверждения
по телефону, а последующее отслеживание по-прежнему можно поручить
ассистенту.

В ходе такого поручения может обнаружиться и деталь, которую стоит
улучшить: пользователю приходится дополнительно спрашивать, с какого дня
начнёт действовать скидка. Сначала следует просмотреть запись разговора
ассистента со службой поддержки. Если сотрудник лишь сказал, что «скидку
можно предоставить», ассистент должен был ещё во время разговора
уточнить дату её вступления в силу. Если же сотрудник ясно назвал дату,
но ассистент раз за разом упускал её в резюме, возникает другой вид
обучающего материала: примеры, в которых исходная реплика с датой
сопоставлена с проверенным резюме.

В первой главе был поставлен вопрос: \textbf{поглотит ли модель в итоге
Harness?} По мнению автора, \textbf{она будет поглощать его слой за
слоем, а Harness --- постоянно перемещаться к новому фронтиру
возможностей.} Каждый раз, когда модель стабильно осваивает способность,
прежде реализованную посредством внешней оркестрации, соответствующую
оркестрацию можно упростить. Интерактивная модель из шестой главы ---
один из таких примеров: после того как модель напрямую научилась
определять очередность реплик, вставлять реплики и отвечать голосом,
система может сократить число соответствующих эвристических правил;
передача аудио, отмена выполнения и управление состоянием по-прежнему
остаются обязанностями среды выполнения.

Некоторые обязанности должны оставаться за пределами модели. Права на
бизнес-операции определяются доверенной идентификацией и правилами
авторизации, соблюдение бюджета контролируется средой выполнения и
квотами сервисов, а внешние операции выполняются инструментами и
окружением.

Даже если в будущем модель сможет освоить какую-либо способность, бизнес
не обязательно может ждать следующего цикла обучения и развёртывания.
Если компания сегодня изменила правила согласования, ассистент уже
сегодня должен действовать по новым правилам. Команда может сначала
изменить Harness, а регулярно возникающие трудности оставить для
следующего цикла обучения. Это также один из способов, которыми автор
понимает «Горький урок»: универсальные методы поиска и обучения будут и
дальше наращивать преимущество, а по мере их постепенного созревания
Harness будет подключать текущие возможности к реальным задачам и
сохранять обратную связь для следующего этапа обучения.

Когда одна команда непрерывно совершенствует и модель, и продукт, эта
цепочка обратной связи может стать теснее. Пример --- Claude Code от
Anthropic: возможности модели влияют на проектирование инструментов и
контекста, а проблемы реального использования задают направления для
последующего улучшения модели и системы. Пример LangChain из первой
главы демонстрирует непосредственную пользу системной инженерии: при
неизменном использовании GPT-5.2-Codex после доработки Harness результат
в Terminal Bench 2.0 вырос с 52,8\% до 66,5\%. Самопроверка перед
завершением, обнаружение повторяющихся операций и распределение бюджета
--- всё это появилось благодаря конкретному анализу процессов,
завершившихся неудачей.

Автор считает, что \textbf{цикл, в котором реальный бизнес, данные
обратной связи и итерации модели взаимно стимулируют друг друга, ---
один из самых глубоких конкурентных барьеров нашей эпохи.} Чем теснее
связаны эти три составляющие, тем больше команда создаёт на основе
практических задач наработок, которые другим трудно воспроизвести.
Командам, создающим модели, следует как можно быстрее включать ценные
задачи и обратную связь в обучение и оценку, чтобы этот цикл
действительно заработал.

Командам, создающим приложения, автор советует: \textbf{используйте
Harness как самый острый краткосрочный технический рычаг, чтобы выиграть
время и накопить долгосрочные преимущества.} Когда модель осваивает
способности, которые прежде приходилось восполнять инженерными
решениями, соответствующие реализации становится всё проще скопировать.
Уникальные данные, устойчивые каналы распространения, доверие
пользователей, сетевые эффекты и сценарии физического мира, в которых
совместно участвуют люди и агенты, требуют долгосрочной работы.
Создавать качественный продукт и хорошо обслуживать пользователей ---
значит использовать это время для возведения собственных барьеров.

Модели, API и продукты продолжат меняться, но вопросы «какая информация
нужна сейчас», «какие действия разрешено предпринимать» и «как проверить
результат» останутся основой проектирования систем. Понимание этих трёх
вопросов позволяет определить, где разместить новую возможность, какие
компоненты можно упростить, а какие звенья необходимо усилить.

Автор надеется, что эта книга поможет вам развить такую способность к
суждению: отталкиваясь от конкретной задачи, понимать структуру агента,
создавать работоспособную систему, а затем непрерывно улучшать её на
основе обратной связи. Основной текст, иллюстрации и сопутствующий
экспериментальный код опубликованы в открытом доступе. Запускайте
эксперименты самостоятельно, пробуйте новые решения и делитесь
открытиями через issue и PR. Агент способен создавать инструменты и
расширять свои возможности с помощью кода, а нам он даёт больше способов
воплощать идеи.

А теперь идите и создайте что-нибудь.

\hypertarget{ux441ux43fux440ux430ux432ux43eux447ux43dux44bux435-ux43eux442ux432ux435ux442ux44b-ux43dux430-ux432ux43eux43fux440ux43eux441ux44b-ux434ux43bux44f-ux440ux430ux437ux43cux44bux448ux43bux435ux43dux438ux44f}{%
\chapter*{Справочные ответы на вопросы для
размышления}\label{ux441ux43fux440ux430ux432ux43eux447ux43dux44bux435-ux43eux442ux432ux435ux442ux44b-ux43dux430-ux432ux43eux43fux440ux43eux441ux44b-ux434ux43bux44f-ux440ux430ux437ux43cux44bux448ux43bux435ux43dux438ux44f}}
\addcontentsline{toc}{chapter}{Справочные ответы на вопросы для
размышления}

Справочные ответы сгенерированы AI и слегка проверены человеком.

\hypertarget{ux433ux43bux430ux432ux430-1.-ux432ux432ux435ux434ux435ux43dux438ux435-ux432-ai-ux430ux433ux435ux43dux442ux43eux432}{%
\section*{\texorpdfstring{\crossreflink{chap:1}{Глава 1}. Введение в
AI-агентов}{. Введение в AI-агентов}}\label{ux433ux43bux430ux432ux430-1.-ux432ux432ux435ux434ux435ux43dux438ux435-ux432-ai-ux430ux433ux435ux43dux442ux43eux432}}
\addcontentsline{toc}{section}{\crossreflink{chap:1}{Глава 1}. Введение
в AI-агентов}

\textbf{1. (★★) Если можно приоритетно улучшить лишь один из трёх
компонентов --- модель, контекст или инструменты, --- как определить, во
что выгоднее всего вложиться? Какие результаты оценки заставят вас
изменить выбор?}

\begin{quote}
Сначала можно проверить, получила ли модель всю информацию, необходимую
для выполнения задачи. Это позволяет напрямую обнаружить пропуски и
локально их устранить. Затем следует проанализировать траектории сбоев:
если не хватает фактов о предметной области, предыдущих решений или
состояния задачи --- улучшить контекст; если не хватает средств
наблюдения или интерфейсов выполнения --- добавить инструменты; если
информации достаточно и подходящие инструменты уже имеются, но модель
всё равно не справляется с рассуждением, --- сравнить её с более мощной
моделью. На одном и том же наборе задач следует поочерёдно изменять
только один компонент, фиксируя долю выполненных задач, типы ошибок,
время и стоимость.
\end{quote}

\textbf{2. (★★★) Предположим, что на каждом раунде добавляется примерно
одинаковый объём сообщений, а модели каждый раз передаётся вся
предыдущая история. Почему совокупный объём входных данных растёт
приблизительно квадратично с числом раундов? Какие виды затрат изменяют
кэширование, сжатие и обрезка истории?}

\begin{quote}
Сначала учтём только историю сообщений и предположим, что на каждом
раунде добавляется 100 токенов. На первых трёх раундах модели передаётся
соответственно 100, 200 и 300 токенов. В третьем запросе история
содержит лишь 300 токенов, однако суммарно за три запроса уже было
передано 600. Первое число можно увидеть в третьем запросе, а второе
нужно получить, сложив объёмы исторического ввода из записей об
использовании для всех трёх запросов. В общем случае, если на каждом
раунде добавляется \(m\) токенов, на раунде \(i\) передаётся
приблизительно \(mi\) токенов, а суммарно за первые \(n\) раундов ---
около \(mn(n+1)/2\). Если дополнительно используется фиксированный
системный промпт длиной \(P\), он также учитывается на каждом раунде,
поэтому к общей величине нужно прибавить \(nP\).

По-прежнему учитываем только историю сообщений из предыдущего примера.
Предположим, что начиная со второго раунда сервис может полностью
использовать ввод предыдущего раунда. Тогда на втором раунде кэш будет
использован для 100 токенов, а на третьем --- для 200. Итого за три
раунда из 600 входных токенов для 300 будет использован кэш, а ещё 300
будут обработаны впервые. Кэш позволит повторно использовать вычисления
префикса для этих 300 токенов.

При сжатии история заменяется более короткой сводкой, что создаёт
затраты на вызов для формирования сводки и перестроение кэша. Обрезка
непосредственно удаляет часть истории, поэтому необходимо сохранить
цель, состояние, а также пары вызовов и результатов. Историю лучше
обрабатывать пакетно на границах этапов или при достижении порога длины:
это сокращает число вызовов для сжатия, перестроений кэша и потери
информации из-за многократного сжатия. Объёмные результаты можно вынести
во внешнее хранилище и читать по мере необходимости либо изолировать
исследовательский процесс с помощью дочернего агента, а затем сравнить
потерю информации с затратами.
\end{quote}

\textbf{3. (★★) Модель принимает всё больше решений о вызове
инструментов, а Harness должен поддерживать всё более длинные и сложные
задачи. Как при этом меняются их обязанности? Какие механизмы можно
упростить, а какие необходимо усилить?}

\begin{quote}
Когда модель способна самостоятельно планировать и выбирать инструменты,
можно постепенно упрощать фиксированную маршрутизацию, исправление
формата и часть эвристических подсказок. Чем длиннее задача, тем важнее
для Harness организация контекста, управление интерфейсами инструментов,
а также применение ограничений, проверка и исправление. Следует усилить
такие механизмы, как проверка разрешений, бюджеты и тайм-ауты,
обнаружение повторных вызовов, сохраняемое состояние и восстановление
после ошибок. После каждого обновления модели с помощью сравнительной
оценки на задачах нужно определять, какие правила можно сократить.
Идентификация, авторизация и внешнее выполнение по-прежнему должны
оставаться обязанностью доверенной среды выполнения и окружения.
\end{quote}

\textbf{4. (★★) В эксперименте отсутствие результатов инструмента
приводило к повторным вызовам. Какие ещё причины могут вызывать подобные
циклы в промышленной системе? Как их обнаруживать, как восстанавливаться
после них и какие условия остановки задавать?}

\begin{quote}
Распространённые причины: инструмент постоянно возвращает одну и ту же
ошибку; модель многократно запрашивает несуществующую возможность; при
сжатии теряется состояние завершённых действий; формат сообщения не
соответствует требованиям интерфейса; для задачи отсутствуют необходимые
входные данные. Повторяющееся поведение можно обнаруживать по сочетанию
имени инструмента, параметров, результата и изменения состояния. После
обнаружения цикла сначала следует проверить причину ошибки, а затем
дополнить информацию, изменить путь выполнения или восстановить
состояние. Среда выполнения должна ограничивать число раундов, время,
стоимость и количество последовательных сбоев. При достижении
ограничения нужно сохранить прогресс и в соответствии с правилами задачи
завершить её либо передать управление.
\end{quote}

\textbf{5. (★) Выберите AI-продукт, которым пользуетесь регулярно, и
отдельно перечислите основную информацию в его контексте, доступные
операции и стратегию выполнения. На каком этапе он наиболее подвержен
сбоям? Как бы вы его улучшили?}

\begin{quote}
Например, рассмотрим помощника для путешествий. Контекст содержит пункт
назначения, даты, бюджет, предпочтения пользователя и уже полученные
результаты поиска. Доступные операции включают поиск авиабилетов, поиск
отелей, изменение черновика и запрос подтверждения. Стратегия выполнения
может сочетать автономный поиск с фиксированным процессом бронирования.
Если помощник постоянно рекомендует варианты сверх бюджета, следует
проверить, присутствует ли бюджет в текущих входных данных, своевременно
ли обновляются цены и применяются ли правила фильтрации, а затем оценить
улучшения на одном и том же наборе задач по планированию поездок.
\end{quote}

\begin{quote}
Можно протестировать одно и то же предложение в двух поручениях. Оба
пользователя хотят купить авиабилет не дороже 1200 юаней, а служба
поддержки сообщает, что минимальная цена составляет 1250 юаней. Один
деловой путешественник при подобных небольших превышениях бюджета обычно
отдаёт приоритет своевременному прибытию. Другой путешествует с семьёй и
каждый раз сначала рассчитывает общий бюджет. Если оба поручили
помощнику забронировать билеты, то с учётом обычных предпочтений первого
путешественника помощник может принять разницу в 50 юаней, оформить
билет и затем сообщить фактическую цену. Второму путешественнику
помощник должен сначала отправить предложение и попросить его принять
решение. Если добавить в поручение делового путешественника фразу «на
этот раз максимальная сумма возмещения --- ровно 1200 юаней», помощник
также должен остановиться до оформления билета. Если после покупки
деловой путешественник скажет, что не согласен с дополнительными
расходами в 50 юаней, помощник должен запомнить это исправление и в
следующий раз при аналогичном предложении сначала спросить его. Ожидая
ответа пользователя, помощник может также уточнить у службы поддержки,
можно ли временно зафиксировать цену, либо поискать другие рейсы.
\end{quote}

\textbf{6. (★★) Спроектируйте процесс для службы поддержки по
бронированию авиабилетов: какие этапы лучше реализовать как рабочий
процесс, а какие --- поручить автономному агенту? Как обе части должны
взаимодействовать, если пользователь внезапно изменит маршрут?}

\begin{quote}
Проверку личности, проверку состояния заказа, подтверждение цены, оплату
и бронирование следует организовать как процесс с чёткими
предварительными условиями. Уточнение потребностей, поиск альтернативных
рейсов и сравнение вариантов нужно поручить автономному агенту. Порядок
оплаты и оформления билета, а также правила удержания места и тайм-аутов
должны соответствовать конкретному бизнес-интерфейсу. Если пользователь
изменил дату, следует обновить версию маршрута, отменить утратившие силу
предложения и ожидающие выполнения операции, а затем заново проверить
цену и авторизацию. Для таких операций, как возврат средств и крупные
платежи, необходимо запрашивать подтверждение в соответствии с правилами
управления рисками, а результат проверять по квитанции или подтверждению
операции.
\end{quote}

\textbf{7. (★★★) Риск использования одного и того же инструмента
\texttt{delete\_file} различается при удалении временного и критически
важного системного файла. Как динамически оценивать риск с учётом
параметров, области ресурсов и возможности восстановления?}

\begin{quote}
При оценке риска необходимо одновременно учитывать тип инструмента,
параметры, целевой ресурс, личность вызывающей стороны и последствия
операции. Сначала серверная часть должна разобрать и нормализовать
целевой путь, проверить, действительно ли ресурс находится в разрешённой
области, а затем учесть важность файла, масштаб удаления, наличие
резервной копии и возможность восстановления. Символические ссылки и
изменение ресурса после проверки также необходимо учитывать
непосредственно при выполнении. Низкорисковые временные файлы можно
автоматически удалять по правилам. Для критически важных ресурсов
следует применять предварительный просмотр, подтверждение или полный
запрет удаления, сохраняя запись о результате. Модель сообщает намерение
задачи, а среда выполнения обеспечивает контроль доступа.
\end{quote}

\textbf{8. (★★) Открытое пространство действий допускает сочетания
содержимого и операций, которые трудно заранее перечислить исчерпывающим
образом. В каких задачах действия лучше ограничить небольшим набором
предопределённых вариантов? Какие преимущества и недостатки создаёт
такое ограничение?}

\begin{quote}
Для возврата средств, платежей и изменения состояния учётной записи
подходят предопределённые операции со структурированными параметрами:
это упрощает проверку диапазонов сумм, переходов между состояниями и
условий авторизации. Ограничение вариантов сокращает число недопустимых
комбинаций и облегчает тестирование и отслеживание, но снижает гибкость
при обработке новых запросов. Для задач в пределах поддерживаемой
области можно предоставить явные операции, а для выходящих за её рамки
запросов --- механизмы уточнения, передачи человеку или контролируемого
расширения. Полноту покрытия следует проверять на реальных сценариях.
\end{quote}

\textbf{9. (★★) Когда агенту требуется подтверждение человека,
пользователь может быть не в сети, отвечать медленно или давать
неоднозначные указания. Какое состояние должна сохранять система, какую
работу можно продолжить, а какие операции обязаны ждать?}

\begin{quote}
Следует сохранять исходную цель, завершённые действия, ожидающую
подтверждения операцию, текущие параметры, состояние авторизации и
версии связанных результатов. Можно продолжать уже разрешённую и
независимую работу, например систематизацию материалов. Платежи,
удаление и отправка данных внешним получателям, требующие подтверждения,
должны оставаться в ожидании. Пользователя следует уведомить по
выбранному им каналу, задав правила напоминаний и тайм-аутов. При
неоднозначном ответе нужно продолжать уточнение. Перед возобновлением
выполнения необходимо заново проверить цену, наличие, личность
пользователя и состояние среды, чтобы убедиться, что исходная операция
всё ещё действительна.
\end{quote}

\textbf{10. (★★★) Приведите пример инженерного метода, который можно
упростить по мере развития моделей. Объясните, какую проблему он решал
изначально и как лежащий в его основе принцип проектирования системы
продолжит действовать.}

\begin{quote}
Например, когда модель начинает стабильнее формировать структурированные
вызовы, можно сократить специализированную логику исправления JSON и
заполнения пропущенных параметров. При этом проверка по Schema
интерфейса, проверка авторизации и обработка ошибок продолжат
обеспечивать надёжность вызовов.
\end{quote}

\begin{quote}
Распределение обязанностей внешней базы знаний также может меняться
вместе со способностью модели к обучению: часть стабильных знаний
переходит в параметры модели, тогда как обновление в реальном времени,
точный поиск, контроль доступа и отслеживание источников остаются
обязанностью внешней системы.
\end{quote}

\begin{quote}
Интерфейсы возможностей могут эволюционировать от множества
специализированных инструментов к сочетанию «текстовый регламент Skill +
универсальный исполнитель». По мере улучшения способности модели
понимать интерфейсы и соблюдать регламенты универсальный исполнитель
сможет выполнять всё больше операций, а число специализированных
инструментов можно будет сократить. Модель формирует команды на основе
регламента, а среда выполнения проверяет параметры и разрешения.
Стандартизированные вызовы передают структурированные параметры,
обеспечивая проверку и взаимодействие между системами.
\end{quote}

\begin{quote}
Промпты и определения инструментов также можно загружать по мере
необходимости. В ранних реализациях их часто размещали в фиксированном
месте, чтобы помочь модели распознать их и следовать им. По мере
улучшения способности следовать инструкциям и включения динамической
загрузки в обучение Skills и динамически обнаруженные инструменты можно
добавлять в последующий контекст, сокращая объём содержимого,
передаваемого заранее.
\end{quote}

\hypertarget{ux433ux43bux430ux432ux430-2.-ux438ux43dux436ux435ux43dux435ux440ux438ux44f-ux43aux43eux43dux442ux435ux43aux441ux442ux430}{%
\section*{\texorpdfstring{\crossreflink{chap:2}{Глава 2}. Инженерия
контекста}{. Инженерия контекста}}\label{ux433ux43bux430ux432ux430-2.-ux438ux43dux436ux435ux43dux435ux440ux438ux44f-ux43aux43eux43dux442ux435ux43aux441ux442ux430}}
\addcontentsline{toc}{section}{\crossreflink{chap:2}{Глава 2}. Инженерия
контекста}

\textbf{1. (★★★) Изучите журнал работы скользящего окна из эксперимента
2-3: что делал каждый из пяти вызовов инструментов? Какие дополнительные
данные необходимо записывать и какое сравнение провести, чтобы
определить, вызваны ли повторные вызовы усечением истории? Разработайте
схему усечения с сохранением ключевого состояния и объясните её влияние
на префиксный кэш.}

\begin{quote}
В журнале скользящего окна из эксперимента 2-3
(\texttt{result\_\allowbreak{}sliding\_\allowbreak{}window\_\allowbreak{}20260718\_\allowbreak{}kimi\_\allowbreak{}k2\_\allowbreak{}6.\allowbreak{}json}) в поле
\texttt{task} записано требование пользователя вывести список файлов
Python, прочитать \texttt{main.py} и \texttt{agent.py}, а затем
сформулировать итог в трёх предложениях. В \texttt{tool\_calls}
перечислены пять вызовов: первые три --- \texttt{find("*.py")}, каждый
из которых находил 13 файлов; последние два считывали по одному
указанному файлу. Если сопоставить это с поручением, остаётся
сформулировать итоговые три предложения. В одном ответе можно запросить
несколько инструментов, поэтому по одним лишь пяти записям о вызовах
нельзя однозначно сопоставить их с пятью запросами к модели. Журнал
сохраняет вызовы инструментов, но не полные сообщения, передаваемые
модели на каждой итерации. В текущей реализации на каждой итерации также
прикладывается исходная задача. Чтобы проверить, приводит ли усечение к
повторным вызовам, можно записывать идентификаторы сообщений,
сохранённых и удалённых на каждой итерации, зафиксировать модель, задачу
и максимальное число итераций, а затем сравнить количество повторов и
долю завершённых задач при полной истории и скользящем окне.

Один из вариантов --- записывать прочитанные файлы, ключевые результаты
и предстоящие задачи в сводку этапа, архивировать объёмные результаты с
сохранением индекса, а недавние запросы к инструментам и их результаты
хранить попарно. Замена истории нарушает префиксный кэш после точки
изменения; системные инструкции и определения инструментов,
расположенные до неё, по-прежнему могут использоваться повторно. Другой
вариант --- поручить исследование файлов подагенту, который вернёт в
основной диалог только выводы и расположение файлов.
\end{quote}

\textbf{2. (★★) Chat Template модели Qwen3 обрабатывает предыдущие
рассуждения с учётом положения последнего реального пользовательского
запроса. Как управлять постепенно разрастающимися рассуждениями, если
один цикл ReAct длится сотни итераций? Сравните протоколы усечения
предыдущих рассуждений, их обратной передачи и привязки к префиксу и
объясните, как к ним должен адаптироваться Harness.}

\begin{quote}
Сначала следует определить в соответствии с протоколом модели и
интерфейса, какие рассуждения необходимо передавать обратно, а затем
задать бюджет истории. Предыдущие рассуждения могут содержать ещё не
выполненные планы и уже исключённые пути; в долгосрочных задачах
особенно важно сохранять такое состояние. Прямое удаление лишает
последующие решения оснований, а полная обратная передача увеличивает
объём входных данных. Если протокол допускает усечение, можно сохранять
недавние рассуждения, а долгосрочный план преобразовывать в сводку
состояния. Если интерфейс требует обратной передачи, необходимо
сохранять исходные поля ответа. Если \texttt{thinking} привязан к
префиксу, после изменения префикса затронутые блоки рассуждений следует
обработать согласно протоколу либо использовать серверное сжатие.
\end{quote}

\textbf{3. (★★) Предположим, поисковые материалы объёмом около 150 тысяч
символов необходимо сжать до двух тысяч символов. Какие факты, связи и
источники вы сохраните? Как проверить, достаточно ли сводки для
последующих вопросов, и как восстановить опущенные подробности?}

\begin{quote}
Следует сохранить сведения о людях, времени, организациях, связях,
ключевых доказательствах и нерешённых вопросах, а также без изменений
оставить идентификаторы, используемые для поиска исходных данных. Для
проверки сводки нужно разработать последующие запросы и выявить пропуски
и неоднозначности. Исходные материалы следует архивировать, сохранив
источники, время получения и способ доступа к ним. Ссылки помогают найти
источник, а снимки исходных материалов обеспечивают возможность
долгосрочной повторной проверки.
\end{quote}

\textbf{4. (★★) Счётчик инструментов не учёл один вызов, поэтому панель
состояния показывает неверное число оставшихся вызовов. Разработайте
журнал событий, проверку агрегированных данных и проверку на уровне
выполнения и объясните, как обнаружить и устранить такое расхождение.}

\begin{quote}
Число вызовов прежде всего должно подсчитываться программно. Если модели
необходимо извлекать события из неструктурированных записей, ей также
следует поручать поэлементное извлечение, а агрегирование выполнять
программно, чтобы сократить число пропусков и повторных учётов в длинной
истории. Каждому вызову нужно назначить уникальный идентификатор события
и записывать состояния запроса, выполнения и возврата результата, явно
определив, относится ли ограничение к числу попыток вызова или к числу
успешных вызовов. Идемпотентные обновления предотвращают повторный учёт,
а воспроизведение событий и регулярная сверка позволяют обнаруживать
пропуски. Уровень выполнения проверяет доступность оставшихся вызовов по
доверенному счётчику, а модель использует состояние для планирования.
При обнаружении расхождения следует сохранить журнал событий,
восстановить агрегированные данные и лишь затем продолжить работу; текст
неизвестного происхождения необходимо считать материалом, требующим
проверки.

Например, после принятия вызова A исполнитель дважды доставил одно и то
же уведомление «A принят». При объединении по идентификатору вызова два
уведомления расходуют бюджет только один раз. Если бюджет учитывает
фактическое число выполнений, у действительно запущенной повторной
попытки должен быть отдельный идентификатор выполнения, связанный с
исходным запросом.
\end{quote}

\textbf{5. (★★) В эксперименте 2-4 сравнивались тон, организация
информации и описания инструментов. Как управлять версиями системного
промпта, долгое время поддерживаемого несколькими людьми, проверять
конфликты правил и оценивать каждое изменение с помощью регрессионных
задач?}

\begin{quote}
Промпт следует хранить в системе управления версиями и подвергать ревью,
а проверку процессов и ограничений должны совместно проводить
представители бизнеса, эксплуатации и разработки. Необходимо
поддерживать регрессионный набор типичных задач, граничных случаев и
известных сбоев, фиксируя цель и результат каждого изменения. Правила
следует организовать с помощью Markdown/XML и проверять на повторы и
конфликты; стабильные правила размещать в допускающей повторное
использование начальной части контекста, динамическое состояние
добавлять отдельными сообщениями, а специализированные процессы выносить
в Skills.
\end{quote}

\textbf{6. (★★★) Пример с зоомагазином включает как получение отдельной
записи, так и сводную статистику. Разработайте структуру исходных
записей, статистической сводки и маршрутизации запросов, позволяющую
эффективно отвечать на вопросы обоих типов, и укажите, когда необходимо
обновлять сводку.}

\begin{quote}
Записи по отдельным клеткам следует индексировать по номеру клетки, а
количество животных каждого окраса --- программно обновлять приращениями
на основе событий добавления, удаления и изменения. Запрос об отдельной
клетке должен обращаться к соответствующей записи, запрос о количестве
--- к статистической сводке, а вопрос, затрагивающий новое измерение,
--- к исходным данным. Необходимо периодически пересчитывать итоговые
значения и сверять их с состоянием, полученным путём приращений. Если
эти состояния требуется сохранять между сеансами, записи и индексы
следует хранить в постоянном хранилище и по мере необходимости загружать
в новых задачах.
\end{quote}

\textbf{7. (★★★) Skill уже находится в каталоге, но модель не
определила, что он применим к текущей задаче. Как диагностировать
проблему с точки зрения описания, организации каталога, явного запуска и
оценочных примеров?}

\begin{quote}
Следует проверить, охватывают ли имя и описание Skill реальные
формулировки задач, и добавить конкретные условия запуска и границы,
которые легко перепутать. Если каталог велик, кандидатов можно искать по
типу задачи. Явный запуск предоставляет прямой путь загрузки. Оценочный
набор должен включать случаи, в которых Skill следует запускать, случаи,
в которых его не следует запускать, и случаи с похожими описаниями, в
которых его легко спутать с другим Skill; для них необходимо отдельно
подсчитывать пропущенные и ошибочные срабатывания.
\end{quote}

\textbf{8. (★★) Skill успешно загружен, но модель пропустила
предусмотренный в нём этап проверки. Разработайте контрольный
эксперимент, позволяющий различить проблемы запуска, понимания,
использования инструментов и определения завершённости, и сравните
результаты разных моделей.}

\begin{quote}
Следует зафиксировать задачу, инструменты и версию Skill и отдельно
протестировать самостоятельный запуск и явную загрузку. Сначала нужно
проверить, присутствует ли полный текст Skill в сообщениях, отправленных
модели. Затем, сохраняя загруженный текст неизменным, поочерёдно дать
модели указание «проверь результат» и инструкцию, явно содержащую объект
проверки, команду и условие успешного прохождения, после чего сравнить,
инициировала ли модель проверку. Для уже инициированных вызовов следует
поочерёдно вернуть заранее подготовленные успешный и неуспешный
результаты и проследить, объявит ли модель о завершении несмотря на
провал проверки. Это позволяет отдельно подсчитать число пропущенных
загрузок, пропущенных вызовов и ошибочных завершений. Для разных моделей
следует использовать поддерживаемые ими протоколы сообщений и сравнивать
их по одинаковым критериям результата. Затраты на кэширование можно
проверить с помощью трёх последовательных запросов. В первом текст Skill
добавляется один раз, а затем сохраняется без изменений: пока на сервере
действует кэш и соблюдаются условия попадания, связанные с префиксом,
разбиением и границами кэша, этот текст можно использовать повторно.
Если на каждой итерации добавлять в конец ещё одну копию, новую копию
придётся вычислять, но исходный префикс по-прежнему сможет попадать в
кэш. Если сначала удалить текст из прежней позиции, а затем перенести
его в последнюю итерацию, префикс изменится начиная с прежней позиции и
расположенный после неё кэш, как правило, уже нельзя будет использовать
напрямую.
\end{quote}

\textbf{9. (★★★) Система содержит большое количество инструментов, набор
которых часто меняется. Как расположить фиксированные определения,
каталог инструментов, результаты загрузки по требованию и состояние
среды выполнения с учётом точности обнаружения, бюджета контекста и
попаданий в кэш?}

\begin{quote}
Часто используемые стабильные инструменты следует разместить в начале в
фиксированном порядке, а определения остальных предоставлять по
требованию через каталог и поиск инструментов. Новое содержимое
необходимо загружать поддерживаемым интерфейсом способом, не позволяя
несущественным изменениям распространяться на уже существующий префикс.
Динамическое состояние следует размещать в конце, выбирая стратегию
замены или добавления. Родительский и дочерний агенты могут повторно
использовать общие стабильные инструкции; различия в их полномочиях и
задачах необходимо выражать явно.
\end{quote}

\hypertarget{ux433ux43bux430ux432ux430-3.-ux43fux43eux43bux44cux437ux43eux432ux430ux442ux435ux43bux44cux441ux43aux430ux44f-ux43fux430ux43cux44fux442ux44c-ux438-ux431ux430ux437ux430-ux437ux43dux430ux43dux438ux439}{%
\section*{\texorpdfstring{\crossreflink{chap:3}{Глава 3}.
Пользовательская память и база
знаний}{. Пользовательская память и база знаний}}\label{ux433ux43bux430ux432ux430-3.-ux43fux43eux43bux44cux437ux43eux432ux430ux442ux435ux43bux44cux441ux43aux430ux44f-ux43fux430ux43cux44fux442ux44c-ux438-ux431ux430ux437ux430-ux437ux43dux430ux43dux438ux439}}
\addcontentsline{toc}{section}{\crossreflink{chap:3}{Глава 3}.
Пользовательская память и база знаний}

\textbf{1. (★★) Пользователь говорит: «1 мая я переезжаю из A в B, с 10
по 20 мая живу в гостинице C, в этот период счета присылайте по
электронной почте». Как 15 мая помощник должен ответить на вопросы о
домашнем адресе, временном месте проживания и способе доставки счетов?
Разработайте правила записи и поиска, а также объясните, как
обрабатывать исправления опечаток и противоречащие друг другу адреса.}

\begin{quote}
На 15 мая домашний адрес --- B, временное место проживания --- гостиница
C, а счета следует присылать по электронной почте. Для каждой записи
нужно хранить назначение, период действия, исходный сеанс и дословное
сообщение пользователя. После переезда запись с домашним адресом A
действует до 1 мая, а B --- начиная с этого дня; гостиница и требование
отправлять счета по электронной почте применимы только с 10 по 20 мая.
При поиске одновременно учитываются назначение и целевая дата, поэтому,
если сегодня ретроспективно запросить домашний адрес на 30 апреля,
результатом по-прежнему должен быть A.

Если пользователь явно исправляет опечатку, следует изменить
соответствующую запись, сохранив основание для исправления. Если две
записи содержат разные адреса для одного назначения и одного периода, а
достаточных доказательств, объясняющих различие, нет, нужно сохранить
оба источника и запросить у пользователя подтверждение.
\end{quote}

\textbf{2. (★★) Если исходный документ содержит противоречивые сведения,
контекстный префикс может записать в индекс и ошибку. Как вы будете
хранить источники, версии и неопределённые сведения и использовать их на
этапах поиска и генерации?}

\begin{quote}
Предположим, в старом руководстве сказано: «Возврат возможен в течение
семи дней после получения», а в новой версии срок увеличен до 14 дней с
примечанием: «Применяется к заказам, оформленным начиная с 1 июня». Если
к каждому фрагменту новой версии добавить префикс «Срок возврата
составляет 14 дней», то при поиске по заказу за май первые результаты
могут сообщать только о 14 днях, хотя все они происходят из одного и
того же обобщения, в котором пропущено условие. При построении индекса
сгенерированный префикс следует хранить отдельно от исходного текста,
записывая версию документа, дату публикации, условия вступления в силу,
права доступа и цитируемый фрагмент. Условия двух версий должны
содержать взаимные ссылки, чтобы можно было определить, какое из них
применимо к конкретному заказу.

Приёмку можно провести на двух заказах: первый оформлен 31 мая и получен
2 июня, второй оформлен и получен 1 июня. Сначала нужно ограничить поиск
руководствами, доступными пользователю, затем выбрать применимую версию
по дате оформления заказа, свериться с исходным текстом и ответить
соответственно «7 дней» и «14 дней», сославшись на условие. Если в новой
версии указана только дата публикации, но не сказано, с какой партии
заказов она применяется, следует сохранить оба источника и указать,
какого именно правила не хватает. Такая проверка также позволит
выяснить, не принимает ли система один и тот же префикс, повторяющийся в
нескольких фрагментах, за несколько независимых доказательств.
\end{quote}

\textbf{3. (★★) Какие пространственные отношения легко утратить при
преобразовании таблиц, графиков и схем в текст? На примере таблицы,
архитектурной схемы или линейного графика разработайте схему совместного
поиска по исходному изображению, структурированному представлению и
текстовому индексу.}

\begin{quote}
Например, объединённые заголовки таблицы определяют, к какому году и
региону относится числовое значение. Если сохранить только «доход 100»,
эти отношения будут утрачены. Можно одновременно хранить исходное
изображение, номер страницы и координаты области, а также
структурированную ячейку, содержащую заголовок строки, заголовок
столбца, значение и единицу измерения. Текстовый индекс используется для
определения местоположения, при подготовке ответа читается
структурированная запись, а при неоднозначности проверяется исходное
изображение. Для архитектурной схемы также следует сохранять узлы,
направления связей и отношения вложенности.
\end{quote}

\textbf{4. (★★★) Какие операции разбиения на фрагменты, индексирования и
поиска можно будет поручить модели по мере развития моделей и увеличения
контекстного окна? Проанализируйте это с учётом рассуждений Рича Саттона
о поиске и обучении в «горьком уроке» и разработайте эксперимент,
сравнивающий передачу всего содержимого на вход с поиском по требованию,
одновременно проверяя качество выполнения задач, вычислительные затраты
и ограничения прав доступа.}

\begin{quote}
Модель сможет выбирать границы фрагментов, переписывать запросы и
объединять результаты; по мере развития её возможностей часть ручных
настроек упростится. Сам поиск и выполняемое моделью обобщение также
используют методы поиска и обучения. Для сравнения следует зафиксировать
материалы, вопросы и модель, а затем отдельно запустить передачу всего
содержимого на вход, однократный поиск и многоэтапный поиск, измеряя
качество, задержку и стоимость по мере роста объёма материалов. Следует
включить задачи на статистику по нескольким документам, точные условия,
обновляющиеся факты и различия в правах доступа, чтобы определить, при
каких условиях каждый метод остаётся эффективным. Статистику можно
поручить проверяемым агрегирующим вычислениям, а права доступа должны
принудительно обеспечиваться при формировании входных данных и на уровне
выполнения.
\end{quote}

\textbf{5. (★★★) Если модель уже обладает обширными предметными
знаниями, что по-прежнему следует хранить во внешней базе знаний?
Отдельно рассмотрите общедоступные сведения, постоянно обновляемые
правила, частные материалы и доказательства, требующие точной
прослеживаемости.}

\begin{quote}
Для стабильных общедоступных знаний можно преимущественно полагаться на
модель, но для новейших правил необходимо хранить версии и области
действия, для частных материалов требуется контроль доступа, а для
точной аргументации --- источники и исходный текст. Внешняя база знаний
также поддерживает исправление, удаление и аудит.
\end{quote}

\textbf{6. (★) Какие запросы поддерживают многоуровневые сводки RAPTOR,
а какие --- связи между сущностями и сводки сообществ GraphRAG?
Приведите по одному примеру для общего обзора, локальных подробностей и
многошаговых связей, а также объясните, какие доказательства необходимо
проверять.}

\begin{quote}
Высокоуровневые сводки RAPTOR подходят для общего обзора, а материалы
нижнего уровня содержат подробности; при поиске можно сочетать
содержимое разных уровней. Локальные связи между сущностями в GraphRAG
поддерживают запросы по цепочкам наподобие «пользователь --- врач ---
больница --- адрес», а сводки сообществ --- тематические обобщения.
Например, чтобы получить обзор расширений инструкций, описанных в
руководстве по процессору, можно сначала перечислить темы по сводке, а
затем проверить по оглавлению, не пропущена ли целая категория. Если
вопрос касается только значения определённого бита XCR0, следует открыть
описание регистра и сверить номер бита с определением. На
\crossreflink{fig:3-11}{рисунке 3-11} задан вопрос: «Каков адрес
больницы, в которой работает мой стоматолог?» Нужно начать с самого
пользователя, по связи стоматологического лечения найти доктора Чжана A,
а затем узнать адрес стоматологической больницы «Жэньай». Если же начать
с отца пользователя, будет найден доктор Чжан B из отделения
кардиологии. Даже если в итоге адрес больницы окажется совершенно
точным, ответ всё равно будет ошибочным, если на предыдущем шаге выбран
не тот человек.
\end{quote}

\textbf{7. (★★) Как совместно использовать оглавление, ссылки и
векторный индекс для организации знаний? Спроектируйте базу знаний,
которую должны совместно сопровождать люди и AI-агент, и укажите роль
навигации, семантического поиска, обновления версий и фильтрации по
правам доступа.}

\begin{quote}
Рассмотрим руководство по служебным поездкам. Сотрудник, который знает,
что ему нужно найти раздел «Изменение билета», может перейти по
оглавлению «Служебные поездки → Авиабилеты → Изменение билета», а затем
по ссылкам в записи открыть правила возмещения расходов. Если он помнит
лишь вопрос «Можно ли возместить дополнительные расходы из-за срочной
замены рейса?», найти то же правило поможет векторный поиск; если же
известен номер условия, эффективнее будет поиск по ключевым словам.
Сводки на входной странице достаточно для навигации, но когда речь идёт
о суммах и исключениях, необходимо открыть исходный текст действующей
версии.

Люди и агент редактируют один и тот же набор читаемых файлов, пригодных
для версионирования. После изменения финансовым отделом правил замены
билетов необходимо одновременно обновить сводку на входной странице,
связанные ссылки и затронутый поисковый индекс. Затем следует повторить
поиск по фразе из старого правила и проверить, возвращается ли уже новая
версия. Если договорные тарифы доступны только сотрудникам отдела
закупок, нужно отдельно выполнить от имени обычного сотрудника поиск по
оглавлению и индексу, а также открытие исходного текста, чтобы
убедиться, что заголовки, сводки и поисковые фрагменты подчиняются тем
же правам доступа, что и исходный текст.
\end{quote}

\textbf{8. (★★★) Как после извлечения полей из примеров и кластеризации
определить, полезны ли полученные прототипы? Разграничьте значимость
признаков для кластеризации, способность предсказывать результаты и
причинное объяснение, а также разработайте процесс оценки на отложенной
выборке и экспертной проверки.}

\begin{quote}
Сначала нужно проверить точность извлечения полей и единообразие
обработки пропущенных значений, а затем изучить общие признаки внутри
кластеров, различия между кластерами и устойчивость при повторном
формировании выборок. Признаки и их веса, используемые для разделения
кластеров, описывают структуру группировки; прогностическую способность
необходимо оценивать на отложенной выборке; для причинного объяснения
дополнительно требуются контроль смешивающих факторов, подходящий дизайн
исследования и предметные доказательства. Например, в эксперименте 3-12
можно определить поля, способ обработки пропущенных значений и центры
кластеров на 360 материалах для моделирования, затем зафиксировать весь
процесс и распределить 60 отложенных материалов по прототипам. Эксперты
должны сначала независимо прочитать исходные тексты, а затем проверить
основания классификации. Следует выборочно проверять не только типичных
представителей, но и наиболее удалённые от центра и трудно
классифицируемые случаи. Если дополнительно утверждается, что прототип
способен предсказывать некоторый результат, целевой показатель нужно
определить заранее, сравнить качество на отложенных материалах с простой
базовой линией и не допускать утечки этого результата через извлекаемые
поля.
\end{quote}

\textbf{9. (★★★) Разработайте для базы пользовательской памяти в
Markdown процесс инкрементального обновления и периодической
реорганизации. Какие ошибки легко пропустить, если проверяющий и автор
предложения используют одну и ту же модель, а проверяющий видит только
фрагменты, отобранные автором предложения? Предложите улучшения с точки
зрения полноты доказательств, прав инструментов и оценки процесса
проверки.}

\begin{quote}
При инкрементальном обновлении автор предложения ищет старые знания и
отправляет минимально возможное, но целостное изменение вместе с его
источниками, одновременно обновляя ссылки, индексы и временную
информацию. Проверяющий самостоятельно повторяет поиск, а программа
проверяет формат, ссылки, метаданные и метки прав доступа; после
успешной проверки изменение сливается, а затронутые индексы обновляются.
При периодической реорганизации материалы дедуплицируются, объединяются
и перераспределяются по оглавлению согласно контрольному списку
покрытия, а затем сверяются с исходными материалами. Реорганизация
проходит тот же процесс проверки; кандидат публикуется после успешного
прохождения регрессионной проверки на поисковых задачах, после чего
продолжается наблюдение за фактическим качеством поиска. При просмотре
только отобранных фрагментов легко пропустить отрицательные условия,
временные изменения и конфликты между файлами; одна и та же модель также
может повторить сходную ошибочную интерпретацию. Проверяющему следует
предоставить полный доступ к поиску доказательств и инструментам
проверки в пределах разрешённой области, отдельно разграничив права на
запись в рабочую ветвь, проверку и публикацию. Долю обнаруженных ошибок
следует оценивать на заранее внедрённых ошибочных изменениях, сравнивая
разные сочетания моделей и объёмы доступных доказательств; также
необходимо установить бюджет итераций и предусмотреть передачу человеку.
\end{quote}

\hypertarget{ux433ux43bux430ux432ux430-4.-ux438ux43dux441ux442ux440ux443ux43cux435ux43dux442ux44b}{%
\section*{\texorpdfstring{\crossreflink{chap:4}{Глава 4}.
Инструменты}{. Инструменты}}\label{ux433ux43bux430ux432ux430-4.-ux438ux43dux441ux442ux440ux443ux43cux435ux43dux442ux44b}}
\addcontentsline{toc}{section}{\crossreflink{chap:4}{Глава 4}.
Инструменты}

\textbf{1. (★★) MCP отделяет подключение инструментов от фреймворка
AI-агента. Выберите одно из требований --- потоковые результаты,
двунаправленное взаимодействие или управление состоянием длительных
задач --- и объясните, какие возможности протокола и хоста для него
необходимы, а также какие дополнительные проектные решения потребуются
при интеграции с существующим MCP.}

\begin{quote}
В качестве примера длительной задачи можно разделить «отправку
вычисления» и «получение результата»: после отправки возвращается
идентификатор задачи, а дальнейшие сведения о ходе выполнения можно
получать через запросы состояния или уведомления. Сервер должен
сохранять состояние задачи и её результаты, а также определять переходы
между состояниями ожидания в очереди, выполнения, завершения, сбоя и
отмены. Хост должен сохранять идентификатор задачи, возобновлять запросы
после разрыва соединения и направлять новые результаты соответствующей
задаче агента. При интеграции также необходимо проверить используемую
версию MCP, расширения и поддержку со стороны клиента, а затем решить,
достаточно ли существующих возможностей или требуются явно определённые
бизнес-инструменты. Закрытие соединения протокола, прекращение ожидания
клиентом и отмена фоновой задачи --- три состояния, которые необходимо
обрабатывать раздельно. Кроме того, следует учитывать оркестрацию
событий между сеансами: новые письма и внешние обратные вызовы должны
пробуждать соответствующие задачи, а постановка в очередь, возобновление
и повторные попытки также требуют единых соглашений.

Можно закрыть окно, в котором отображается ход выполнения, а затем вновь
открыть задачу по сохранённому номеру. Если после этого по-прежнему
удаётся получить её актуальное состояние, значит, жизненный цикл задачи
отделён от окна. Затем следует явно отменить задачу и проверить,
действительно ли исполнитель остановился.
\end{quote}

\textbf{2. (★★) Как выбирать между несколькими серверами MCP,
предоставляющими схожие инструменты? На примере, где один сервер
возвращает краткое содержание, а другой --- полный текст, спроектируйте
описания инструментов, обозначение источников и критерии выбора так,
чтобы агент мог использовать эти различия.}

\begin{quote}
Предположим, что служба каталога возвращает краткое содержание, а
архивная служба предоставляет полный текст. Хост может назначить этим
двум инструментам псевдонимы \texttt{catalog.summary} и
\texttt{archive.read\_full}, связав каждый псевдоним с идентификатором
настроенного сервера и опубликованным им именем инструмента. В описании
следует чётко указать охватываемые материалы, возвращаемое содержимое,
способ пагинации, частоту обновления и стоимость вызова. Если модели
нужно выбрать кандидатов из множества файлов, она сначала обращается к
каталогу; если требуется сверить исходный текст или условия цитирования
--- читает полный текст. Хост выполняет маршрутизацию согласно этим
связям и проверяет версию службы и область авторизации.

Проведите пробный сценарий с договором, имеющим две версии: способно ли
краткое содержание предоставить идентификатор документа и версию,
которые распознаёт архивная служба? Если да, используйте их для чтения
соответствующего исходного текста. Если системы идентификаторов двух
служб несовместимы, повторно найдите документ по ссылке на источник и
сверьте версию. Такая проверка позволяет убедиться, что описания
инструментов соответствуют фактически возвращаемым значениям, а также
обнаружить ошибку, при которой краткое содержание выбрано верно, но
открыта другая версия документа.
\end{quote}

\textbf{3. (★★) В каких ещё сценариях, помимо запуска проверок после
написания кода, можно применять цикл «выполнение --- проверка ---
обратная связь»? Выберите два примера, опишите способы проверки,
обратную связь при сбое и затраты, а также обсудите, когда уместны
немедленная, выборочная и асинхронная проверка.}

\begin{quote}
Первый пример --- изменение конфигурации: сначала проверить формат и
поля, затем запустить службу в изолированной среде и получить состояние
работоспособности или результат её поведения. При сбое следует вернуть
ошибочное поле, диагностические данные запуска и фактическое состояние,
чтобы помочь агенту решить, нужно ли изменить конфигурацию или
восстановить исходное значение. Второй пример --- создание презентации:
отрендерить страницы, проверить, не выходит ли текст за границы и не
перекрывают ли диаграммы другие элементы, а затем сообщить создателю
номера страниц и расположение проблем. Проверка формата требует
небольших затрат, поэтому её можно выполнить немедленно. Если изменён
только текст на одной странице, сначала можно отрендерить её и соседние
страницы. Если заменён весь шаблон, целесообразно отрендерить всю
презентацию в фоновом режиме и перед передачей прочитать результаты
проверки каждой страницы. Выборочная проверка определяет, какие страницы
проверять, а асинхронное выполнение --- когда ожидать результат; оба
подхода можно применять одновременно.

Результат внешней операции также можно запросить. Например, после
отправки письма можно получить его состояние на сервере, а после
создания заказа --- найти запись по бизнес-идентификатору. Для проверки
следует использовать интерфейс чтения состояния, подтверждения или
аудиторской записи и различать такие значения, как «принято» и
«завершено». Если операция оказывает внешнее воздействие,
предварительное согласование и последующая проверка выполняют разные,
соответствующие им задачи.
\end{quote}

\textbf{4. (★★) Как помимо активного обнаружения организовать
возможности при наличии тысяч инструментов? На примере каталога с
классификацией, шаблонов рабочих процессов или универсального
исполнителя опишите условия применимости и компромиссы одного из
подходов.}

\begin{quote}
Для бизнес-процессов с высокой повторяемостью распространённые задачи
можно оформить в виде шаблонов рабочих процессов. Например, шаблон сбора
материалов задаёт последовательность «поиск кандидатов --- чтение
основного текста --- проверка источников --- создание отчёта», причём на
каждом этапе доступны только относящиеся к нему инструменты. Шаблон
сокращает число одновременно доступных определений и помогает модели
понимать зависимости данных между этапами. Такой подход подходит для
стабильных процессов с контролируемыми различиями во входных данных. Для
новых задач необходимо сохранить возможность изменять этапы или
добавлять инструменты.

Можно также сочетать другие методы: организовать каталог по предметным
областям и постоянно держать доступным небольшое число часто
используемых инструментов; записать подробные инструкции в Skill и
загружать их по мере необходимости; вызывать стабильные библиотеки или
API через универсальный исполнитель кода, сокращая число
специализированных интерфейсов.
\end{quote}

\hypertarget{ux433ux43bux430ux432ux430-5.-coding-agent-ux438-ux443ux43dux438ux432ux435ux440ux441ux430ux43bux44cux43dux44bux439-ai-ux430ux433ux435ux43dux442}{%
\section*{\texorpdfstring{\crossreflink{chap:5}{Глава 5}. Coding Agent и
универсальный
AI-агент}{. Coding Agent и универсальный AI-агент}}\label{ux433ux43bux430ux432ux430-5.-coding-agent-ux438-ux443ux43dux438ux432ux435ux440ux441ux430ux43bux44cux43dux44bux439-ai-ux430ux433ux435ux43dux442}}
\addcontentsline{toc}{section}{\crossreflink{chap:5}{Глава 5}. Coding
Agent и универсальный AI-агент}

\textbf{1. (★★) Спроектируйте песочницу для Coding Agent, которому
необходимо читать файлы проекта, устанавливать зависимости и запускать
тесты. Как вы настроите права доступа к файлам, сети, процессам и
ресурсам? Как следует сформулировать запрос, проверить его область
действия и сохранить запись о нём, если для задачи требуется расширение
прав?}

\begin{quote}
Сначала следует перечислить необходимый для задачи доступ: рабочая копия
проекта доступна для записи, справочные материалы --- только для чтения,
зависимости загружаются из одобренных источников. Процессы, время
выполнения, память и диск ограничиваются посредством изолированной среды
исполнения; учётные данные хранятся доверенным сервисом. При расширении
прав указываются целевой ресурс, операция, срок и причина, а
исполнительный слой проверяет и регистрирует разрешение. Конфигурация
проверяется с помощью сборки, тестов и сценариев превышения лимитов
ресурсов.
\end{quote}

\textbf{2. (★★★) Когда создатель последовательно генерирует и изменяет
несколько поколений агентов, дефекты фреймворка и ошибки предметной
области могут постепенно накапливаться. Как отслеживать эти изменения с
помощью эталонных версий, регрессионных тестов, оценки поведения и
механизмов отката?}

\begin{quote}
Следует отдельно версионировать эталонный фреймворк, результат каждого
поколения и требования, фиксируя происхождение сгенерированного
результата и изменения. Набор регрессионных тестов должен одновременно
охватывать протокол сообщений, привязку инструментов, восстановление
состояния и реальные задачи предметной области, а независимые задачи ---
проверять обобщающую способность. Среди них должны быть достаточно
сложные задачи с проверяемыми результатами: простые примеры легко
скрывают дефекты, накапливающиеся от поколения к поколению. Необходимо
сравнивать долю успешных выполнений, стоимость и типы ошибок новой и
старой версий и публиковать новую версию только после успешной проверки;
при деградации следует восстановить последнюю версию, прошедшую приёмку.

Например, можно зафиксировать один небольшой проект на одном коммите и
поручить исходной и следующей версиям агента по одному разу исправить
одну и ту же ошибку. Даже если новая версия устранила целевую ошибку,
необходимо запустить прежние тесты, а затем изучить оба патча и
проверить, не было ли частное решение этой задачи встроено в
универсальный фреймворк.
\end{quote}

\textbf{3. (★★) Изменение формата журналов может быть следствием как
штатного обновления, так и дефекта производителя. Какие доказательства
должна сохранять адаптивная система разбора и на основании каких условий
она должна добавлять новый парсер, сообщать об аномалии или
приостанавливать обработку?}

\begin{quote}
Следует сохранять исходные журналы, версию производителя, описание
изменений и результат старого парсера. Сначала необходимо сверить смысл,
типы и значения полей с контрактом интерфейса, а затем проверить штатные
и аномальные образцы. Подтверждённое обновление формата включается в
очередную версию парсера; изменения с противоречивой семантикой,
отсутствующими данными или неизвестным происхождением заносятся в отчёт
о проблеме, а связанную с ними обработку при необходимости
приостанавливают. Новый парсер и запись предупреждения можно сохранить
одновременно.
\end{quote}

\textbf{4. (★★) В цикле «автор предложения --- рецензент» рецензент
считает плотность информации на странице приемлемой, но пользователю она
кажется перегруженной. Как преобразовать предпочтения пользователя в
примеры, критерии оценки и интерактивную обратную связь и проверить,
соответствуют ли им последующие изменения?}

\begin{quote}
Следует попросить пользователя отметить приемлемые и чрезмерно
перегруженные страницы, чтобы сформировать конкретные критерии размера
шрифта, свободного пространства, плотности информации и других
характеристик. Отзыв рецензента должен содержать номер страницы и
расположение проблемы, а после изменения следует показать ту же страницу
для сравнения. Запись о предпочтении должна включать пользователя,
задачу и область применимости; в последующих задачах предпочтение
используется в пределах этой области. Выборочная повторная проверка
позволяет обнаружить расхождения между критериями рецензирования и
пользовательским опытом.
\end{quote}

\textbf{5. (★★) Coding Agent может обновлять базу знаний, архитектурную
документацию и проектные инструкции, а также закреплять
последовательности операций в коде. Как выявлять избыточные или
устаревшие правила и определять для правил область применимости, способ
проверки и условия прекращения действия? Какие механизмы накопления
данных и оценки дополнительно необходимы, чтобы перейти от одного
успешного изменения к непрерывной эволюции, рассматриваемой в
\crossreflink{chap:9}{главе 9}?}

\begin{quote}
Для каждого правила следует фиксировать источник, область применимости,
проверочные сценарии и время последней проверки. Дублирующиеся правила
объединяются, утратившие силу архивируются вместе с соответствующей
версией, а стабильные и однозначно проверяемые требования переносятся в
Linter, тесты или проверки инструментов. Например, исходная инструкция
требует использовать \texttt{pnpm\ test}. После миграции одного рабочего
пространства на другой менеджер пакетов сначала следует выполнить тесты
из непрерывной интеграции с помощью новой команды в чистой рабочей
копии, а затем заменить инструкцию для этого рабочего пространства;
подпроекты, продолжающие использовать pnpm, по-прежнему проверяются
исходной командой. Тем самым определяются область и условие прекращения
действия старого правила. Изменение правил сначала проверяется на
исторических задачах, а затем оценивается его польза для новых задач.
Для непрерывной эволюции также необходимо организовать в долгосрочный
цикл накопление свидетельств выполнения, подготовку возможных изменений,
независимую оценку, публикацию и откат.
\end{quote}

\textbf{6. (★) Команды, приспособленные к удалённой совместной работе,
обычно также позволяют агенту без труда получать рабочую информацию.
Изучите требования, архитектуру, записи решений и критерии приёмки
знакомого вам проекта, перечислите документацию, которую важнее всего
дополнить, и объясните, какие конкретные задачи она поможет выполнить
агенту.}

\begin{quote}
Для самопроверки можно использовать процесс первого знакомства нового
сотрудника с проектом: способен ли он найти контекст требований,
архитектурные связи, способ сборки, точки запуска тестов и условия
приёмки. В первую очередь следует дополнить знания, влияющие на текущую
задачу: например, объяснить назначение модуля, зафиксировать принятое
техническое решение или предоставить исполняемый приёмочный сценарий.
Договорённости, передаваемые только устно, трудно включить во входные
данные агента; решения, принятые в ходе повседневных обсуждений команды,
также следует сохранять в доступном для поиска виде. Необходимо явно
назначить ответственного за документацию и условия, запускающие её
обновление, а при изменении кода одновременно поддерживать связанные
материалы в актуальном состоянии.
\end{quote}

\textbf{7. (★★★) Как организовать управление идентификацией и
авторизацией, маркировку источников, проверку потоков данных и контроль
записи в память для агента, который одновременно обращается к
конфиденциальным данным, обрабатывает недоверенное содержимое, выполняет
внешнюю коммуникацию и использует постоянную память? Покажите
расположение этих проверок на примере одного штатного бизнес-процесса.}

\begin{quote}
Рассмотрим чтение данных клиента и отправку ему сервисного уведомления.
Служба идентификации определяет, к каким клиентам разрешён доступ;
данные клиента и внешние веб-страницы получают отдельные метки
происхождения; коммуникационный интерфейс проверяет получателя и поля,
разрешённые к отправке. После завершения задачи при записи в память
проверяются происхождение, необходимость, степень конфиденциальности и
срок действия информации. Чтение данных, внешняя отправка и
межсессионное сохранение регистрируются раздельно, чтобы можно было
расследовать аномалии и отзывать утратившие силу разрешения.
\end{quote}

\textbf{8. (★★) В режиме Artifact для SQL и HTML модель генерирует
артефакт, исполнительный сервис обрабатывает запрос или выполняет
рендеринг, а интерфейс непосредственно отображает результат. Какую
информацию уместно дополнительно передавать модели для объяснения? Как
проверять права на выполнение запросов, бюджет ресурсов, изоляцию
страницы и соответствие отображаемого результата исходным данным?}

\begin{quote}
Полную таблицу может напрямую отображать исполнительный сервис, что
избавляет модель от поэлементной генерации данных, экономит время и
токены, а также снижает риск ошибок при переписывании чисел. Когда
требуется обобщение или принятие решения, модель может по мере
необходимости читать агрегированные результаты, аномальные строки и
описания полей. Запросы выполняются под ограниченной учётной записью с
проверкой структуры и ограничениями на доступные данные, функции и
бюджет исполнения; страницы отображаются в изолированной среде, а
свойства и действия декларативных компонентов также проверяются.
Результат проверяется сравнением с эталонными запросами и значениями на
странице, после чего дополнительно сверяются единицы измерения, даты и
методика расчёта статистики.
\end{quote}

\textbf{9. (★★) Инструмент проверяет правила по бизнес-состоянию в базе
данных, а описание параметров помогает модели подготовить вызов. На
примере отмены бронирования объясните, за что отвечают политика на
естественном языке, примеры параметров, проверка исполнительного слоя и
граничные тесты. Что необходимо синхронно обновить после изменения
политики?}

\begin{quote}
Политика на естественном языке сообщает модели, какие условия следует
проверить, и позволяет объяснить результат пользователю. Например, если
пользователь изменил планы поездки, при вызове передаются
\texttt{reservation\_id} и
\texttt{cancellation\_reason="change\_of\_plan"}; модель также может
заполнить \texttt{expected\_has\_insurance}, указав собственное
предположение. Исполнительный сервис проверяет личность и права на
заказ, а решение принимает по сохранённым в базе времени бронирования,
сведениям о страховке, состоянию сегментов маршрута и часам сервера.
Если модель считает, что страховка оформлена, но в заказе указано
обратное, сервис регистрирует расхождение и всё равно руководствуется
данными заказа. Между проверкой и подтверждением операции также
необходимо предотвращать конкурентные изменения, способные сделать
решение недействительным.

Согласно примеру политики из основного текста, сначала следует
зафиксировать остальные условия: экономкласс, страховка отсутствует,
рейс выполняется штатно, ни один сегмент маршрута не использован. Ровно
через 24 часа после бронирования отмена должна быть разрешена, а через
24 часа и 1 секунду --- отклонена. Затем отдельно изменяется состояние
сегментов маршрута: даже если прошло всего два часа, использование хотя
бы одного сегмента должно привести к отказу. Эти три примера проверяют
соответственно временную границу и порядок применения правил. При
изменении политики необходимо синхронно обновить описание,
исполнительный код, примеры параметров и граничные тесты, а также
определить, какая версия политики применяется к старым бронированиям.
\end{quote}

\hypertarget{ux433ux43bux430ux432ux430-6.-ux432ux437ux430ux438ux43cux43eux434ux435ux439ux441ux442ux432ux438ux435-ux440ux430ux441ux448ux438ux440ux435ux43dux438ux435-ux43fux440ux43eux441ux442ux440ux430ux43dux441ux442ux432ux430-ux43dux430ux431ux43bux44eux434ux435ux43dux438ux439-ux438-ux434ux435ux439ux441ux442ux432ux438ux439}{%
\section*{\texorpdfstring{\crossreflink{chap:6}{Глава 6}.
Взаимодействие: расширение пространства наблюдений и
действий}{. Взаимодействие: расширение пространства наблюдений и действий}}\label{ux433ux43bux430ux432ux430-6.-ux432ux437ux430ux438ux43cux43eux434ux435ux439ux441ux442ux432ux438ux435-ux440ux430ux441ux448ux438ux440ux435ux43dux438ux435-ux43fux440ux43eux441ux442ux440ux430ux43dux441ux442ux432ux430-ux43dux430ux431ux43bux44eux434ux435ux43dux438ux439-ux438-ux434ux435ux439ux441ux442ux432ux438ux439}}
\addcontentsline{toc}{section}{\crossreflink{chap:6}{Глава 6}.
Взаимодействие: расширение пространства наблюдений и действий}

\textbf{1. (★★) Когда асинхронный AI-агент получает новое событие, ему
необходимо определить его приоритет относительно текущей задачи. Если
такое решение требует семантического понимания --- например, нужно
распознать, дополняет ли пользователь условия или требует немедленно
остановиться, --- как следует распределить обязанности между движком
правил и LLM? Сравните время отклика, стоимость и последствия ошибочной
классификации.}

\begin{quote}
Сначала можно задать правила на основе доверенного источника и типа
события: явные сигналы остановки, сбои исполнителя и задачи с
наступившим сроком направляются в соответствующие каналы обработки.
Намерение, выраженное на естественном языке, распознаётся с помощью
классификационной модели, после чего среда выполнения распределяет
событие с учётом разрешений и текущего состояния. Правила требуют мало
ресурсов, а их поведение легко проверять; модель способна обрабатывать
разнообразные формулировки, но увеличивает стоимость инференса, задержку
и вероятность ошибочной классификации. Критически важные сигналы
остановки должны иметь возможность срабатывать напрямую через
специальный канал; также следует определить порядок действий при
тайм-ауте классификации или неопределённом результате. Содержимое
внешнего события рассматривается как данные для классификации, а
разрешения и приоритет ограничиваются доверенными метаданными и
политикой планирования.

Решение можно проверить на двух последовательно поступивших сообщениях.
Сначала пользователь говорит: «Пока остановись», а затем инструмент
возвращает ценовое предложение. Операцию, связанную с обязательствами,
следует приостановить, но ценовое предложение можно сохранить для
использования после возобновления. При воспроизведении событий нужно
проверить, не появились ли новые заказы между получением команды о
приостановке и фактической остановкой.
\end{quote}

\textbf{2. (★★) В очереди накопилось 20 событий: 10 результатов
инструментов, 5 пользовательских сообщений и 5 системных напоминаний.
Как организовать их порядок, формат и состояния обработки, чтобы модель
могла отслеживать зависимости задач и реагировать на важные изменения?}

\begin{quote}
Сначала события следует сгруппировать по идентификатору задачи,
источнику, времени и зависимостям, выявив действительно повторяющиеся
события. Обновления условий одной задачи сохраняются в хронологическом
порядке; результаты инструментов содержат идентификаторы исходного
вызова и фоновой задачи и относятся к соответствующей задаче согласно
протоколу интерфейса. Если исходный вызов уже вернул состояние
выполнения, последующий результат завершения доставляется в
согласованной форме события. События остановки или сбоя, требующие
немедленной обработки, среда выполнения принимает в приоритетном
порядке; остальные события организуются по задачам в список и помечаются
как необработанные, обрабатываемые или обработанные. Для длинных
результатов инструментов сохраняются факты и состояния, необходимые для
принятия решений, а полное содержимое помещается во внешнее хранилище с
возможностью поиска. После каждого хода список обновляется, чтобы уже
обработанные события не запускали действия повторно. Новые условия
пользователя следует связывать с затрагиваемым планом или незавершённой
задачей.
\end{quote}

\textbf{3. (★★★) Агент может использовать отдельный адрес электронной
почты и отдельный номер телефона, а после авторизации --- управлять
личными учётными записями пользователя. Рассмотрите с учётом фоновой
работы, требований другой стороны к проверке личности, области
полномочий и аудита, для каких сценариев подходит каждый из двух
способов идентификации.}

\begin{quote}
Для записи, запросов и повседневного общения предпочтительно
использовать отдельную идентичность: это упрощает фоновую работу,
позволяет отделить действия агента и независимо управлять областью его
доступа. Для работы с данными или операциями внутри учётной записи
пользователя можно применять явно ограниченные и отзывные полномочия,
подключая пользователя, если другая сторона требует подтверждения от
владельца учётной записи. Например, при звонке в службу поддержки агент
может сначала изложить запрос, а затем организовать трёхсторонний
разговор для подтверждения личности; вход через визуальный интерфейс
пользователь может выполнить самостоятельно в контролируемой среде.
\end{quote}

\textbf{4. (★★) Сквозная голосовая модель обрабатывает аудио напрямую,
поэтому ошибку на промежуточном этапе трудно локализовать извне. Как
спроектировать наблюдения и контрольные эксперименты, позволяющие
различить ошибки восприятия, рассуждения и выражения ответа?}

\begin{quote}
Следует сохранять собранные с разрешения входные аудиозаписи, временную
шкалу событий, результаты инструментов и выходное аудио, а также
фиксировать текстовый вывод, если модель может его предоставить. Для
одной и той же записи нужно сопоставить прямой аудиоответ с вариантом, в
котором сначала выполняется транскрипция, а затем формируется ответ;
после этого следует добавить проверенную человеком транскрипцию и
проследить, меняется ли ошибка в зависимости от формы входных данных.
Например, если модель ошиблась в ответе на арифметическую задачу в
аудиоформате, но правильно ответила после получения точной текстовой
транскрипции, в первую очередь стоит проверить аудиотракт. Однако при
переходе к тексту меняется и интерфейс ввода, поэтому по двум ответам
ещё нельзя однозначно определить ошибочный этап. Если модель
одновременно возвращает текст и озвученный ответ, их также можно
сопоставить и обнаружить числа, пропущенные только при озвучивании. В
двух записях эксперимента 6-6 используются одинаковые слова и фразы, но
разный темп речи; чтобы определить, где говорят быстрее, необходимо
сохранить звук, поскольку в текстовой транскрипции нужный для ответа
признак уже утрачен. Кроме того, следует раздельно оценивать темп речи,
шум, определение границ реплики, перебивания и задержку до начала первой
реплики.
\end{quote}

\textbf{5. (★) MPS в Step-Audio R1.1 позволяет согласованно развивать
замысел и его словесное выражение. Как системе следует использовать
паузы, сообщения о ходе работы и поэтапное изложение? Что делать, если
последующее рассуждение изменило уже озвученный вывод?}

\begin{quote}
Темп изложения должен соответствовать текущей задаче: во время ожидания
результатов запроса следует сообщить, что именно проверяется, а сложный
материал объяснять по частям с уместными паузами. Заранее можно
озвучивать подтверждённые факты, понимание задачи и последующие шаги
обработки; выводы, относящиеся к бизнес-операциям, должны опираться на
соответствующие доказательства. Если позднее выяснится, что вывод
ошибочен, следует своевременно указать конкретное исправление и его
последствия, а также соответствующим образом обновить дальнейшие
действия. При оценке нужно проверять не только время до первой реплики и
естественность речи, но также согласованность выводов и фактическое
выполнение задачи, не превращая количество слов-заполнителей в
показатель естественности общения.
\end{quote}

\textbf{6. (★★) И визуальная разметка областей в SoM, и индексы
элементов DOM позволяют модели указать место действия одним номером. Как
сочетать повторную разметку, предсказание координат и проверку после
операции при работе с нестандартными элементами управления, перекрытиями
или динамическими изменениями?}

\begin{quote}
Сначала следует проверить, не устарели ли кандидаты из-за прокрутки,
обновления или динамической загрузки; при необходимости нужно повторно
получить состояние страницы и сформировать новое отображение. Если
структурированной информации недостаточно, визуальные кандидаты можно
дополнить с помощью снимка экрана. Если цель ясно видна, но среди
кандидатов её нет, следует перейти к предсказанию координат. Перед
выполнением операции нужно проверить преобразование координат
изображения в координаты рабочего стола, а после неё --- убедиться, что
интерфейс изменился ожидаемым образом. Если сама цель по-прежнему
неясна, можно увеличить соответствующую область или получить
дополнительные сведения из интерфейса. Все три подхода должны учитывать
перекрытия и изменения состояния; также необходимо сохранять
изображение, на основании которого в этот раз определялось положение
цели.
\end{quote}

\textbf{7. (★★) Качество данных телеуправления зависит от навыков
оператора. Как дрожание, обходные траектории, паузы и неудачные захваты
влияют на обучение VLA? Как отфильтровать данные, сохранив при этом
ценные последовательности восстановления после неудач?}

\begin{quote}
При имитационном обучении стратегия действий может усвоить частое
дрожание и бесполезные обходные траектории. Во время сбора можно
фиксировать успешность выполнения, длину пути, изменения скорости и
ускорения, столкновения, выход за границы, вмешательство человека и
отсутствие наблюдений, а затем использовать эти показатели для выявления
фрагментов, требующих проверки. При фильтрации также необходимо
учитывать сложность задачи: успешное восстановление после выскальзывания
объекта может расширять охват обучающих данных сильнее, чем
беспроблемное выполнение. Стабильные демонстрации, неудачные действия и
процессы восстановления следует размечать раздельно, а затем выбирать
данные в соответствии с целью обучения. Это позволяет одновременно
повысить качество базовых действий и сохранить опыт исправления ошибок.
Если базовые действия всё ещё нестабильны, сначала следует проверить и
улучшить качество демонстраций и лишь затем рассматривать усложнение
архитектуры модели.
\end{quote}

\textbf{8. (★★★) Для голоса, Computer Use и робототехники можно
использовать единую модель, модульный каскад или разделение между
взаимодействием на переднем плане и фоновым рассуждением. Какое
сочетание этих подходов вы предпочли бы на ближайшие пять лет? Обоснуйте
решение с учётом задержки, наблюдаемости, стоимости обновлений и рисков
задач.}

\begin{quote}
Сначала следует определить ограничения сценария, а затем выбирать
сочетание подходов. Непрерывное голосовое взаимодействие требует
согласования моментов прослушивания и речи, а сложный бизнес-анализ ---
более продолжительного рассуждения; эти компоненты могут
взаимодействовать через интерфейсы задач и состояния. Computer Use
требует постоянно получать изменения страницы, а робототехника --- ещё и
поддерживать высокочастотное управление. Модульная архитектура упрощает
замену компонентов, локализацию проблем и независимые обновления, но
увеличивает накладные расходы на обмен данными через интерфейсы и
требует поддерживать согласованность состояния между модулями. Единое
обучение потенциально улучшает использование информации и временную
координацию, однако обновление и диагностика требуют более широкой
комплексной оценки. Передовые модели рассуждения обновляются быстро,
поэтому возможность заменять фоновую модель позволяет отдельно
подключать новые возможности и снижает стоимость повторного обучения
всей интерактивной системы. Можно сохранить заменяемый компонент
фонового рассуждения, продолжая совершенствовать модель взаимодействия
на переднем плане.
\end{quote}

\textbf{9. (★★) Стандартные веб-страницы могут предоставлять DOM или
информацию о доступности, тогда как Canvas, WebGL и самостоятельно
отрисовываемые элементы управления зачастую требуют визуального
позиционирования. Какие затраты создаёт одновременная поддержка
структурированного и визуального путей и какие преимущества она даёт?}

\begin{quote}
Структурированный путь предоставляет роли, метки и ссылки на элементы,
упрощая понимание форм и взаимосвязей на странице; визуальный путь
охватывает видимое на снимке экрана содержимое, не представленное в
структуре. Поддержка обоих путей требует синхронизации состояния,
согласования элементов с изображением, обработки разных типов сбоев и
дополнительного тестирования. Преимущество состоит в возможности
выбирать более подходящий способ позиционирования в зависимости от
страницы и с помощью снимка экрана проверять, соответствует ли
структурированная информация текущему интерфейсу. На одних и тех же
задачах можно измерять долю успешных позиционирований, число раундов
взаимодействия, задержку и стоимость восстановления, определяя, когда
использовать один путь, а когда сочетать оба. В краткосрочной
перспективе можно сохранить оба пути; в долгосрочной важнее потенциал
визуального подхода охватывать интерфейсы различного программного
обеспечения. По мере развития визуальных моделей следует повторно
оценивать практическую пользу структурированного пути и постепенно
упрощать малополезные адаптеры.
\end{quote}

\textbf{10. (★★) Разбиение действий на блоки позволяет VLA за один раз
генерировать последовательность будущих действий. Если во время
исполнения объект переместили, как согласовать длину блока, повторное
наблюдение, фоновый инференс и прерывание действий, чтобы совместить
эффективность исполнения со скоростью реакции?}

\begin{quote}
Следует различать длину последовательности действий, предсказанной
моделью, и длину фактически исполняемого за один раз фрагмента: можно
предсказать длинную траекторию, выполнить лишь небольшую начальную
часть, а затем циклически обновить её по новым наблюдениям. Если
инференс и исполнение перекрываются во времени, необходимо согласовывать
переход между старым и новым блоками действий. При быстрых изменениях
среды или приближении к препятствию следует сокращать исполняемый
фрагмент и чаще проверять состояние. Уровень управления непрерывно
отслеживает предельные положения и аномалии, при необходимости замедляет
или останавливает устройство в соответствии с его требованиями и
отбрасывает утратившие актуальность оставшиеся действия. При выборе
параметров следует одновременно измерять задержку инференса,
непрерывность траектории и максимальное время реакции, а не определять
длину блока только по пропускной способности.
\end{quote}

\textbf{11. (★★★) Для быстрого взаимодействия необходимы своевременные
действия. На примерах голосового ответа, отправки веб-формы и захвата
объекта роботом спроектируйте классификацию действий, подтверждение
состояния и механизм остановки так, чтобы скорость реакции
соответствовала последствиям задачи.}

\begin{quote}
В голосовом сценарии можно сначала подтвердить получение задачи, а
вывод, относящийся к бизнес-операции, озвучить после завершения запроса.
На веб-странице можно сначала прочитать данные или заполнить черновик, а
перед отправкой проверить объект, содержимое и разрешение. Робот может
сначала спланировать траекторию приближения, а перед захватом
подтвердить положение объекта и состояние рабочей зоны. Возможность
выполнения действия должна совместно определяться разрешениями, текущим
состоянием и необходимыми подтверждениями. Как быстрая, так и медленная
модель подчиняется одинаковым ограничениям исполнения, а среда
выполнения блокирует действия, для которых не выполнены необходимые
условия. При остановке необходимо раздельно учитывать генерацию и
воспроизведение речи, ещё не отправленные и уже отправленные запросы, а
также команды движения и фактическое состояние исполнительного
механизма.
\end{quote}

\textbf{12. (★★★) Выберите один из примитивов --- пробуждение,
безопасную точку, отмену, вытеснение либо разделение быстрого и
медленного мышления --- и сравните его реализацию в
событийно-управляемых задачах и при разбиении действий робота на блоки.
Как скорость изменения среды, обратимость действий и стоимость
наблюдения влияют на проектирование?}

\begin{quote}
Рассмотрим «отмену». В эксперименте 6-2, когда пользователь требует
остановить анализ, планировщик сначала прекращает назначать последующие
работы для этой задачи, а затем даёт исполнителю команду завершить
дочерний процесс. Результат следует проверять по номеру задачи: для T1
получено подтверждение завершения, а позднее о выполнении сообщила уже
новая задача T2. Если до завершения T1 успела записать файл, необходимо
проверить его состояние и решить, следует ли сохранить или удалить его
либо возобновить работу с места последней завершённой записи.

Теперь рассмотрим манипулятор, который приближается к чашке по заранее
сгенерированному блоку действий, но чашку переместили. Удаление
последующих действий из очереди не гарантирует немедленную остановку
манипулятора: он может продолжать исполнять уже переданный фрагмент.
Уровень управления должен замедлить и остановить устройство допустимым
для него способом, а затем считать фактическое положение манипулятора и
чашки, чтобы планировщик мог продолжить работу из нового состояния. Чем
быстрее перемещается чашка, тем меньше допустимое время непрерывного
исполнения без нового наблюдения. После контакта при захвате раскрытие
зажима может привести к падению чашки, поэтому перед восстановлением
нужно проверить, удерживается ли объект. Если визуальное
позиционирование требует значительных ресурсов, частоту наблюдений можно
повышать при приближении или контакте, а при равномерном движении ---
увеличивать интервалы. Для аварийной остановки необходим прямой канал к
уровню управления.
\end{quote}

\textbf{13. (★★) Следует ли относить чат с покупателем, обсуждение
согласования и платёжный обратный вызов по одному заказу к одной ветке?
Сравните объединение веток с подпиской по связанным объектам с точки
зрения порядка обработки, параллелизма, миграции ключа секционирования и
изоляции разрешений.}

\begin{quote}
При объединении веток события одного заказа помещаются в одну очередь,
поэтому задаче проще всего читать контекст. Цена этого подхода ---
взаимная блокировка не связанных между собой сообщений: обсуждение в
группе послепродажного обслуживания будет ждать обработки сообщения
покупателя, а крупные заказы создадут горячие точки. Очередь отражает
порядок поступления, но сам платёжный обратный вызов может запоздать или
прийти не по порядку, поэтому события всё равно должны содержать
временную метку источника. Идентификатор ветки обычно также служит
ключом секционирования; при изменении правил объединения потребуется
миграция существующих данных. Участники из разных источников обладают
разными разрешениями, поэтому после объединения необходимо для каждого
сообщения в одной ветке отдельно определять, какие операции вправе
запускать его автор.

При подписке по связанным объектам сохраняются отдельные ветки каждого
источника, а порядок гарантируется только внутри каждой из них.
Последовательность между ветками определяется по состоянию задачи и
временным меткам источников. Задача может по мере необходимости
подписываться на ветки или отписываться от них; ключи секционирования
при этом не меняются, а разрешения остаются изолированными по веткам.
Цена подхода --- необходимость объединять события из нескольких веток и
обрабатывать случаи, когда одно и то же событие поступает из двух
источников.

Рассмотрим сценарий: сначала поступает платёжный обратный вызов, а затем
сообщение покупателя «отменить возврат средств». При обоих подходах
задача должна сопоставить время возникновения двух событий по данным их
источников с текущим состоянием возврата средств, чтобы определить,
можно ли ещё выполнить отмену. Обратный вызов мог поступить с задержкой,
поэтому порядок получения не равен порядку возникновения. Объединение
веток избавляет только от этапа сбора событий из разных веток.
\end{quote}

\textbf{14. (★★) Голосовой помощник для путешествий умеет узнавать
погоду, читать данные из памяти пользователя, искать в интернете,
находить рейсы и бронировать билеты. Как распределить работу между
передней моделью, фоновой моделью и средой выполнения? Пользователь
диктует имя попутчика, а пока идёт поиск рейса, спрашивает о погоде. Как
не ошибиться в имени при бронировании и не забыть сообщить о рейсе?}

\begin{quote}
Среда выполнения заранее помещает память пользователя в контекст.
Погоду, рейсы и простые запросы в интернете передняя модель проверяет
сама, экономя два обращения к фоновой модели. Если нужно обобщить
найденное или сравнить варианты обмена билета, задачу стоит передать
фоновой модели.

Передняя модель может и сама забронировать билет: сначала она уточняет
написание имени и получает подтверждение пользователя. Если бронирование
делегируется, передняя модель записывает в задание имена и номера,
которые точно расслышала, а также исправления пользователя. По ним
фоновая модель заполняет заказ. Если пользователь отвлёкся на погоду,
среда выполнения сохраняет ещё не озвученный результат поиска рейсов:
ответив о погоде, передняя модель сообщает о рейсах.
\end{quote}

\hypertarget{ux433ux43bux430ux432ux430-7.-ux43eux446ux435ux43dux43aux430-ai-ux430ux433ux435ux43dux442ux43eux432}{%
\section*{\texorpdfstring{\crossreflink{chap:7}{Глава 7}. Оценка
AI-агентов}{. Оценка AI-агентов}}\label{ux433ux43bux430ux432ux430-7.-ux43eux446ux435ux43dux43aux430-ai-ux430ux433ux435ux43dux442ux43eux432}}
\addcontentsline{toc}{section}{\crossreflink{chap:7}{Глава 7}. Оценка
AI-агентов}

\textbf{1. (★★) Оценивающая модель может предпочитать более длинные
ответы, определённый стиль или конкретную позицию при показе. Как
спроектировать контрольные примеры, тесты с перестановкой и набор для
калибровки экспертами, чтобы отделить качество содержания от этих
предпочтений?}

\begin{quote}
Обозначьте два ответа как A и B. Сначала покажите их оценивающей модели
в порядке A, B, а затем поменяйте местами. Если оба раза она выберет A,
результат для этой пары согласован; если оба раза она выберет ответ
слева, после перестановки победитель изменится. Зафиксируйте, в скольких
парах происходит такая инверсия, а затем повторите оценку со случайным
порядком и проверьте масштаб расхождений. Одной инверсии недостаточно,
чтобы заключить, что модель всегда предпочитает ответ слева.

Чтобы проверить предпочтение по объёму, сохраните три факта --- сумму
возврата, срок поступления и способ возврата, --- изменив только длину
одного и того же ответа. При проверке предпочтений по стилю старайтесь
сохранять сопоставимую длину, меняя лишь формулировки. Если одновременно
изменить длину, тон и факты, причину изменения оценки будет трудно
объяснить. Попросите специалистов, знакомых с предметной областью,
независимо оценить набор примеров по одной и той же шкале, сохранив
требования задачи, фактические основания и экспертные решения. Сверяйте
расхождения по каждой категории; при использовании дискретных меток
согласованность с поправкой на случайные совпадения можно измерять с
помощью κ Коэна. Для пограничных случаев, в которых критерии оценки
сформулированы неясно, добавляйте эталонные прецеденты, а отдельный
набор примеров оставляйте для проверки пересмотренных критериев. При
повторной оценке моделью из другого семейства также сверяйтесь с
решениями экспертов, чтобы проверить, не пропустили ли обе модели одни и
те же ошибки.
\end{quote}

\textbf{2. (★★★) Общедоступные наборы для оценки могли попасть в
обучающие данные. Как сочетать новые задачи, параметризованные среды,
закрытые тестовые задания и производственные примеры, чтобы уменьшить
влияние загрязнения данных на выводы? Какие признаки могут дать
canary-маркеры?}

\begin{quote}
Оценка должна одновременно включать регулярно обновляемые новые задачи,
параметризованные экземпляры, независимый отложенный набор и проверенные
обезличенные производственные примеры. Правила генерации задач можно
использовать повторно, но экземпляры для оценки и обучения следует
хранить раздельно и обеспечивать достаточные различия в структуре задач,
времени и среде. Изменения параметров должны сохранять отношения,
требующие рассуждения, а не сводиться к простой замене названий и чисел.
Canary-маркеры помогают обнаружить признаки копирования текста или
доступа к нему, однако для вывода о загрязнении обучающих данных
необходимо также учитывать происхождение данных, дату публикации и
поведенческие свидетельства. Раздельная публикация результатов для
источников разных типов помогает увидеть, не знакома ли модель лишь с
определённой категорией общедоступных заданий.
\end{quote}

\textbf{3. (★★) Разработайте рубрику для критериев «полезность ответа» и
«уместность тона» с учётом четырёх принципов: экспертное руководство,
полное покрытие, взвешивание по значимости и самодостаточность. Как с
помощью пограничных случаев разрешать разногласия между оценщиками?}

\begin{quote}
Сначала определите аудиторию и задачу. «Полезность» можно разложить на
то, отвечает ли сообщение текущей потребности, предлагает ли выполнимый
следующий шаг и сохраняет ли необходимые условия. «Уместность тона»
можно разложить на уважение к пользователю, соответствие ситуации
общения и отсутствие необоснованных обещаний. Для каждого критерия
задайте уровни, положительные и отрицательные примеры, а также
пограничные случаи --- например, разницу между кратким признанием
проблемы и многословным повторением извинений в ответе на жалобу.
Допустим, служба поддержки разбирает жалобу: «В счёте появилось лишнее
списание. Как его проверить?» Можно использовать пятибалльную шкалу: 2
балла за точное указание, какое списание нужно проверить; 2 балла за
доступные пользователю прямо сейчас шаги проверки; 1 балл за
уважительную формулировку без самовольной гарантии возврата. Баллы по
каждому критерию начисляются только при полном выполнении его условий.
Ответ «Приносим извинения за доставленные неудобства» получит в этой
задаче лишь последний 1 балл. Попросите нескольких специалистов,
знакомых с предметной областью, независимо выставить оценки, обсудите
различия в понимании задачи, а затем обновите справочные материалы и
прецеденты. В оценках можно сохранить значимые различия предпочтений,
заранее определив способ агрегации в соответствии с целями продукта.
\end{quote}

\textbf{4. (★★) В τ²-bench пользователь моделируется с помощью модели.
Как проверить, следует ли смоделированный пользователь условиям задачи и
охватывает ли он такие условия реального взаимодействия, как эмоции,
неоднозначность и пропуски информации?}

\begin{quote}
Проверяйте на двух уровнях: соблюдение заданных условий и
репрезентативность реального взаимодействия. На первом уровне на каждом
ходе сверяйте, что известно симулятору, что он вправе раскрыть, не
выдумывает ли он факты и правильно ли отвечает на вопросы агента. На
втором сравнивайте его с реальными пользовательскими примерами по
неоднозначности, пропускам, эмоциям, исправлениям и прерываниям.
Постройте для этих условий матрицу сценариев и проверяйте при разных
стратегиях симуляции долю успешных задач, число ходов и категории
ошибок. Затем откалибруйте ключевые задачи на примерах взаимодействия с
людьми и проверьте, изменяются ли абсолютные результаты и рейтинг
моделей при смене симулятора.
\end{quote}

\textbf{5. (★★) Модель Брэдли --- Терри объясняет попарные предпочтения
одномерной скрытой силой. Какие распределения задач или группы
пользователей могут порождать циклические предпочтения? Как
интерпретировать рейтинг с помощью анализа по группам и ошибок подгонки
модели?}

\begin{quote}
Предположим, три оценщика ранжируют три модели соответственно как A, B,
C; B, C, A; и C, A, B, причём каждый ставит любимую модель на первое
место. Проведём попарное голосование: при сравнении A и B двое выберут
A; при сравнении B и C двое выберут B; при сравнении C и A двое выберут
C. Предпочтения каждого человека упорядочены, но в совокупности образуют
цикл. Модель Брэдли --- Терри присваивает каждой модели показатель силы,
однако три таких показателя не способны выразить этот цикл.

При интерпретации рейтинга сначала проверьте, сколько сравнений
проведено для каждой пары и насколько велик диапазон неопределённости
доли побед, а затем сопоставьте её с долей побед, предсказанной моделью.
При достаточном объёме данных модель можно подогнать отдельно для разных
задач или групп пользователей: одним важнее точность ответа, другим ---
скорость, а объединение данных может скрыть эти различия. Окончательный
выбор следует делать с учётом долей задач каждого типа в целевом
бизнес-сценарии.
\end{quote}

\textbf{6. (★★) Доля успеха агента за одну попытку равна 60\%. Сравните
\href{mailto:Pass@k}{\nolinkurl{Pass@k}} и Pass\^{}k при \(k=5\) в
предположении независимости попыток; затем с учётом стоимости повторных
попыток, выбора результата и побочных эффектов действий объясните, как
спроектировать практическую оценку.}

\begin{quote}
Если вероятность успеха каждой попытки равна 0,6 и попытки независимы,
то \href{mailto:Pass@5}{\nolinkurl{Pass@5}} = 1 − 0,4⁵ = 98,976\%, а
Pass\^{}5 = 0,6⁵ = 7,776\%. Первый показатель означает успех хотя бы
одной из пяти попыток, второй --- успех всех пяти. Для практических
повторных попыток также необходим надёжный механизм проверки или выбора
результата; кроме того, нужно учитывать время и стоимость неудачных
попыток. Для задач с внешними побочными эффектами следует предусмотреть
изоляцию состояния, идемпотентность и механизмы восстановления, а затем
в контролируемых условиях измерять каждую неудачу и её последствия.
Значение k выбирают исходя из фактически допустимого бюджета попыток или
требований к повторному исполнению, фиксируя различия между задачами и
корреляцию попыток.
\end{quote}

\textbf{7. (★★) Как при ограниченном бюджете организовать небольшую
диагностику, попарные сравнения, повторные запуски и заключительную
проверку на задачах, не участвовавших в выборе решения, чтобы каждый
этап оценки максимально помогал принять решение?}

\begin{quote}
Сначала на основе кластеров сбоев выберите небольшой набор
диагностических задач и явно определите, какие гипотезы о причинах
должен различить каждый этап. Сравнивайте отдельные изменения в попарных
условиях, уделяя первоочередное внимание промежуточным наблюдениям,
которые помогают объяснить поведение. Отобрав перспективные изменения,
расширьте покрытие задач и увеличьте число повторов, чтобы оценить
интервал различий. Если для каждой задачи обе конфигурации запускаются
по одному разу, попарные различия бинарных результатов можно
анализировать с помощью критерия Мак-Немара. При повторных запусках
одной задачи сначала агрегируйте результаты на уровне задач либо
выполняйте повторную выборку с кластеризацией по задачам. До начала
запусков определите минимальное улучшение, достаточное для принятия
изменения, и максимальное число оцениваемых задач. Если интервал всё ещё
включает как практически значимую пользу, так и ущерб, а бюджет
остаётся, протестируйте следующую партию по заранее установленному
плану. Если после исчерпания бюджета различие остаётся неясным, сообщите
имеющиеся результаты и интервалы и отложите смену конфигурации.
Сохраните не участвовавшие в настройке задачи для подтверждающей
проверки, а при отборе из нескольких вариантов учитывайте множественные
сравнения. Отдельно охватывайте высокорисковые пограничные случаи и
распространённые производственные задачи, не принимая концентрированную
выборку диагностического набора за общее распределение.
\end{quote}

\textbf{8. (★) В четырёх задачах AndroidWorld для Wi-Fi полное дерево
элементов повысило долю успеха с 25\% до 100\%, а расход токенов
составил 249,8\% от исходного варианта. После сокращения дерева все
четыре задачи по-прежнему завершались успешно, а расход токенов составил
50,6\% от полного дерева. Как спроектировать правила отсечения и
регрессионные задачи, чтобы сохранить состояния элементов управления,
отношения доступности и зависимости последующих действий?}

\begin{quote}
Сначала можно разработать две регрессионные задачи. В первой нужно
переключить Wi-Fi из включённого состояния в выключенное и проверить,
позволяет ли сокращённое дерево по-прежнему найти переключатель и
взаимодействовать с ним. Во второй исходным состоянием должен быть уже
выключенный Wi-Fi, а задача --- подтвердить, что он выключен. В этом
случае модели необходимо прочитать текущее состояние и вовремя завершить
работу: ещё одно нажатие, напротив, включит Wi-Fi. Контрольная группа H5
для \texttt{SystemWifiTurnOffVerify} хорошо показывает ценность второй
задачи: симулятор выставил по 1 баллу за состояние в начале и в конце,
но агент не сообщил о завершении. Эксперимент требует как правильного
целевого состояния, так и сообщения агента о завершении, поэтому этот
запуск всё равно считается неудачным.

При отсечении сохраняйте текст элементов управления, метки доступности,
значения состояния, а также элементы, с которыми можно
взаимодействовать, которые можно фокусировать или прокручивать. После
сворачивания пустых контейнеров исходные отношения по-прежнему должны
позволять сопоставить метку с правильным переключателем и направить
действие на нужный элемент управления. Запускайте полное и сокращённое
деревья из одинакового начального состояния и отдельно сверяйте покрытие
элементов управления, успешность задач, расход токенов и задержку. Если
на каком-либо шаге элемент управления не найден, сравните два дерева,
добавьте утраченное отношение в правила, а затем выполните регрессионную
проверку на задачах с прокруткой, переходами между экранами и в других
приложениях.
\end{quote}

\textbf{9. (★★) Симулятор пользователя раскрывает информацию поэтапно в
ответ на вопросы агента. Как сравнить разные стратегии раскрытия и
определить их влияние на число ходов взаимодействия, долю успеха и
сложность задачи?}

\begin{quote}
Зафиксируйте факты задачи, начальную среду и агента, а затем поочерёдно
настройте стратегии активного полного раскрытия, поэтапного раскрытия в
формате вопросов и ответов, а также допустимых пропусков или
неоднозначности. Сравнивайте число необходимых уточнений, вызовы
инструментов, долю успеха и типы ошибок. Используйте поведение раскрытия
информации в траекториях реальных пользователей как основу для
калибровки и обеспечьте стратифицированное покрытие опытных и неопытных
пользователей. Симулятор меняет сложность получения информации, что
может влиять как на абсолютные результаты, так и на порядок моделей в
рейтинге. Поэтому результаты следует проверять при нескольких
обоснованных стратегиях и указывать соответствующие условия
пользователей и задач.
\end{quote}

\textbf{10. (★★) Пользователь передаёт научному ассистенту три
наблюдения: 2, 5 и 8 --- и просит построить график и сообщить среднее
значение. Журнал показывает, что файл открыть не удалось, после чего
программа построила график по значениям 4, 5 и 6, а в отчёте указано
среднее 5. Существующая проверка требует только, чтобы «график
открывался, а среднее было правильным». Каков будет результат этой
проверки? Как изменить критерии приёмки и регрессионные сценарии, чтобы
обнаружить подмену данных?}

\begin{quote}
Существующая проверка пройдёт. Сумма обоих наборов равна 15, а среднее
--- 5, но первая точка на графике изменилась с 2 на 4, а третья --- с 8
на 6. Независимо считайте три значения из пользовательского файла, а
затем сопоставьте их с данными, переданными функции построения графика,
проверив каждое наблюдение в исходном порядке. Одного упоминания имени
файла в журнале недостаточно, чтобы подтвердить, что при построении
графика использовались содержащиеся в нём значения.

Для регрессионной проверки можно провести два испытания. Сначала
сделайте файл недоступным для чтения и проверьте, укажет ли ассистент на
отсутствие данных, попросит ли предоставить их повторно и отметит ли,
что анализ ещё не завершён. Затем измените первое значение в корректном
файле на 11, сохранив два остальных --- 5 и 8. Первая точка должна
переместиться на 11, а среднее --- измениться на 8. Совместная проверка
обоих результатов позволит обнаружить ошибки, при которых обновляется
только график или только отчёт.
\end{quote}

\hypertarget{ux433ux43bux430ux432ux430-8.-ux43fux43eux441ux442ux43eux431ux443ux447ux435ux43dux438ux435-ux43cux43eux434ux435ux43bux435ux439}{%
\section*{\texorpdfstring{\crossreflink{chap:8}{Глава 8}. Постобучение
моделей}{. Постобучение моделей}}\label{ux433ux43bux430ux432ux430-8.-ux43fux43eux441ux442ux43eux431ux443ux447ux435ux43dux438ux435-ux43cux43eux434ux435ux43bux435ux439}}
\addcontentsline{toc}{section}{\crossreflink{chap:8}{Глава 8}.
Постобучение моделей}

\textbf{1. (★★) Одно направленное дообучение улучшило результаты в
целевой задаче, но снизило универсальную способность вызывать
инструменты. Как вы спроектируете набор тестов для проверки исходных
способностей, скорректируете соотношение данных и объём обучения, а
также сравните результаты полного дообучения всех параметров и LoRA?}

\begin{quote}
Предположим, цель дообучения --- заставить ассистента продолжать
проверку, если возврат средств ещё не поступил. Сначала следует
подготовить три независимые от обучения группы заданий: возврат ещё не
выполнен --- проверить, не завершит ли ассистент работу слишком рано;
возврат уже поступил --- проверить, сможет ли он своевременно завершить
работу; заменить возврат другой знакомой задачей, например поиском
авиарейса, --- проверить, сможет ли он по-прежнему выбрать правильный
инструмент и заполнить правильные параметры. Последние две группы
проверяют, не ухудшились ли исходные способности. Следует зафиксировать
инструменты, права и начальное состояние задачи, затем сначала запустить
модель без дообучения и записать базовые показатели для каждой группы.

Начиная с одной и той же базовой модели, следует провести полное
дообучение всех параметров и LoRA с одинаковым разбиением данных и
бюджетом обучения. Во время обучения нужно добавлять примеры из
исходного распределения, а на валидационной выборке подбирать
соотношение новых и старых данных, скорость обучения и момент остановки.
После выбора конфигурации следует отдельно выполнить три группы
отложенных заданий. Если результаты заданий с ещё не поступившим
возвратом улучшились, но в заданиях с уже поступившим возвратом модель
стала выполнять несколько лишних проверок, средняя оценка только по
задаче возврата скроет это ухудшение. LoRA замораживает параметры
базовой модели, но после подключения адаптера эту проверку всё равно
необходимо проводить. Среди дополнительных вариантов --- добавить при
обучении ограничение относительно эталонной политики и подключать
адаптер только для соответствующей задачи; при обработке текста и
изображений также можно в зависимости от источника ошибок решить,
следует ли настраивать проекционный слой, языковую модель или визуальный
компонент.
\end{quote}

\textbf{2. (★★) Некоторую предметную способность можно освоить как
посредством обучения параметров, так и с помощью примеров в контексте и
поиска знаний. На основании каких факторов вы выберете подход?}

\begin{quote}
Сначала следует определить, относится ли потребность к поиску фактов,
описываемым правилам или паттернам восприятия и неявного принятия
решений. Для фактов, которые необходимо обновлять, цитировать, удалять и
защищать посредством управления правами, подходит RAG; небольшое
количество правил и примеров сначала можно поместить в контекст; для
стабильного и многократно используемого поведения можно рассмотреть
обучение параметров. Затем нужно сравнить масштаб использования,
задержку, стоимость обучения и сопровождения. Детерминированные правила
одобрения должны исполняться программой, а модель --- понимать входные
данные и решать языковые задачи, выходящие за рамки правил.

Например, если политика возврата изменилась только в этом месяце, после
обновления материалов для поиска можно задать вопросы с датами до и
после изменения и проверить, ссылается ли модель на версию,
действовавшую в соответствующий момент. Другая ситуация --- модель
прочитала правильную политику, но при разных формулировках запроса
постоянно упускает ключевое условие; именно такие записи полезны для
создания примеров обучения поведению.
\end{quote}

\textbf{3. (★★) Какие обучающие материалы и методы проверки требуются
соответственно для дистилляции Chat-модели, дающей прямые ответы,
Reasoning-модели, генерирующей процесс рассуждения, и Agentic-модели,
многораундово использующей инструменты?}

\begin{quote}
Для дистилляции Chat-модели нужны репрезентативные входные данные,
надёжные ответы и проверка качества выполнения задач, учитывающая как
фактическую корректность, так и форму изложения. Для дистилляции
Reasoning-модели нужны доступные демонстрации рассуждений, проверка
ответов и ключевых промежуточных шагов, а также сопоставление длины и
качества рассуждений. Для дистилляции Agentic-модели дополнительно нужны
определения инструментов, наблюдения за средой и изменения состояния, а
проверять необходимо всю цепочку действий; потери политики вычисляются
только на позициях, сгенерированных моделью, а обратная связь среды
используется как контекст. Если ученик легко попадает в состояния,
отсутствующие в демонстрациях, можно применить дистилляцию на текущей
политике, чтобы добавить обучение с учителем для соответствующих
префиксов.
\end{quote}

\textbf{4. (★★★) AI-агент службы поддержки выполняет задачу за 10
раундов взаимодействия и в итоге получает вознаграждение за успех. Как
связать этот сигнал с промежуточными решениями? Для каких шагов уместно
добавить вознаграждение за процесс?}

\begin{quote}
Предположим, во втором раунде агент уточнил, какое именно списание
оспаривается, в седьмом объяснил начисление, а в десятом получил
подтверждение успешного возврата. Всю траекторию можно отметить как
успешную, но по одному этому подтверждению нельзя определить, насколько
полезным оказалось объяснение в седьмом раунде. При использовании PPO
модель ценности может оценивать ожидаемую отдачу в начале каждого шага,
после чего на основе этой оценки и фактической отдачи можно вычислять
преимущество действия. При использовании GRPO можно начать с одного и
того же спора и сравнить траектории нескольких попыток его
урегулирования; если каждая попытка завершилась успехом, одинаковое
вознаграждение за успех не позволит различить более и менее удачные
подходы.

У проверки личности есть достоверное подтверждение, а факт проверки
необходимого номера заказа можно установить отдельно. После добавления
обратной связи за такие промежуточные события всё равно нужно по
конечному критерию убедиться, что деньги действительно возвращены по
указанному заказу. Если обе допустимые последовательности --- сначала
объяснить начисление или сначала проверить заказ --- приводят к успеху,
следует разрешить оба порядка. Чтобы определить, какое именно объяснение
было наиболее полезным, можно собрать больше траекторий или попросить
модель-учителя предоставить демонстрацию с учётом доказательств.
\end{quote}

\textbf{5. (★★★) У вас есть бюджет 10 000 долларов на улучшение работы
AI-агента службы поддержки. Как распределить вложения между оценкой,
контекстом и знаниями, промптами и навыками, программными ограничениями
и обучением параметров?}

\begin{quote}
Сначала следует потратить часть бюджета на создание репрезентативного
набора задач, измерить распределение ошибок и упорядочить их по частоте
неудач, влиянию на бизнес и стоимости исправления. При недостатке знаний
в первую очередь следует улучшить поиск, при проблемах с пониманием
процесса --- сначала попробовать промпты и навыки, а ограничения прав и
сумм реализовать программно. Сначала можно довести до приемлемого уровня
эти сравнительно недорогие и легко изменяемые решения, а затем выделить
средства на обучение параметров и оценку развёртывания для оставшихся
проблем восприятия, выражения или неявного принятия решений. Один из
вариантов корректируемого бюджета: 3 000 долларов на построение базовых
показателей, регрессионное тестирование и наблюдение после запуска; 2
500 долларов на исправление знаний и поиска; 1 000 долларов на настройку
промптов и навыков; 1 500 долларов на реализацию и тестирование
программных ограничений; оставшиеся 2 000 долларов --- на небольшое
обучение. Если базовые показатели покажут, что главная проблема
заключается в поиске устаревшей политики, зарезервированные на обучение
средства следует сначала потратить на добавление примеров с датами
вступления в силу, а затем повторно проверить работу при чередовании
старой и новой политик.
\end{quote}

\textbf{6. (★★★) У задачи нет явного вознаграждения, а взаимодействие со
средой обходится дорого. Из каких источников можно получать обучающие
сигналы? Как проверить, что эти сигналы действительно улучшают
выполнение задачи?}

\begin{quote}
Сигналы можно получать из человеческих демонстраций, сравнений
предпочтений, оценки каждого токена моделью-учителем, проверяемых
событий в среде и диагностических выводов, сгенерированных моделью.
Стоит сравнить способы преобразования содержательной обратной связи в
конкретные обучающие сигналы: например, сообщение «какой информации о
клиенте ещё не хватает» уже указывает следующий шаг, поэтому сохранение
такой обратной связи обеспечивает более плотное обучение, чем одиночная
метка успеха или неудачи. Симулятор на основе модели позволяет увеличить
число взаимодействий, а дистилляция на текущей политике --- плотность
обучения для каждой траектории.
\end{quote}

\textbf{7. (★★) Какие персонализированные способности целесообразно
хранить в отдельной LoRA для каждого пользователя или каждой компании?
Какую информацию лучше оставить во внешней базе знаний?}

\begin{quote}
Долговременно стабильный стиль изложения, привычки использования
предметной терминологии и повторяющиеся паттерны принятия решений можно
освоить с помощью отдельного адаптера. Состояние учётной записи, часто
меняющиеся факты о продукте и воспоминания, которые необходимо
поэлементно цитировать или удалять, лучше хранить во внешнем хранилище с
поиском. При выборе также нужно учитывать объём пользовательских данных,
изоляцию, частоту обновлений и стоимость развёртывания множества
адаптеров.
\end{quote}

\textbf{8. (★★★) Если опираться на идеи исследования Weak-to-Strong
Generalization, может ли небольшая модель предоставлять полезный сигнал
обучения для более сильного агента? Рассмотрите отдельно локальную
проверку, метки предпочтений и дополнительный контекст, а также
предложите способ выявления систематических отклонений такого надзора.}

\begin{quote}
Небольшая модель может хорошо справляться с узкой задачей проверки или
предоставлять полезный сигнал благодаря доступу к правилам, эталонным
ответам либо дополнительным материалам. Сильная модель генерирует
варианты, а небольшая оценивает только те части, которые способна
надёжно проверять; эту оценку следует дополнять программной проверкой
или выборочной проверкой человеком. Ошибки надзора можно измерять на
эталонных данных, анализировать систематические отклонения по типам
задач и проверять показатели сильной модели посредством независимой
оценки. Примерам с существенными разногласиями следует назначать меньший
вес или передавать их более надёжному источнику оценки.
\end{quote}

\textbf{9. (★★) В одной траектории большинство промежуточных шагов
правильны, но конечный результат неверен; другая траектория дала
правильный результат, однако нарушила ограничения процесса. Как
совместно оценивать процесс и результат этих двух траекторий?}

\begin{quote}
Рассмотрим операции с файлами из эксперимента 8-16. Агент А сначала
прочитал конфигурацию, а затем изменил указанное поле, но при приёмке
выяснилось, что он записал неверное значение. Результат следует отметить
как неудачный, однако валидатор может подтвердить, что исходный файл был
прочитан. Агент Б записал правильное значение, но напрямую перезаписал
существующий файл. Проверка результата пройдёт успешно, однако проверка
пути сможет указать конкретную операцию записи, нарушившую правило
«сначала прочитать, затем изменить». Если это правило является
обязательным условием приёмки, агент Б также не должен её пройти. При
создании нового файла старого содержимого для чтения нет, поэтому
следует разрешить прямую запись.

Во время обучения оценку результата и обратную связь по пути можно
хранить отдельно. Соответствующие правилам действия агента А
подтверждаются журналом, но ошибочное итоговое значение всё равно
необходимо исправить; штраф агенту Б можно отнести к конкретной
нарушающей правила операции записи. Чтобы определить эффективность
обучения, следует на отложенных задачах отдельно измерять долю
завершённых заданий и число нарушений. Если модель просто станет реже
выполнять задачи, число нарушений также уменьшится; доля завершённых
заданий в таком случае поможет проверить, чему именно она научилась.
\end{quote}

\textbf{10. (★★★) Среды и валидаторы SWE-bench, τ²-bench и AndroidWorld
можно использовать для обучения. Как разделить экземпляры задач и
шаблоны, чтобы использовать обучающую ценность материалов и при этом
сохранить независимую оценку?}

\begin{quote}
Следует совместно использовать интерфейсы среды и логику проверки, но
разделять обучающие и оценочные экземпляры, сохраняя правила разбиения.
Параметризованные генераторы позволяют создавать новые однотипные
экземпляры для проверки обобщения внутри распределения; отложенные
задачи с другими шаблонами, из других предметных областей или сред нужны
для проверки более далёкого переноса. Публичные тестовые задания, уже
использованные при обучении, можно сохранить как диагностические
материалы, но для независимой оценки следует применять задания, не
участвовавшие в обучении. Раздельное представление результатов по
каждому типу отложенных данных позволяет определить, вызван ли прирост
адаптацией к конкретным экземплярам или более широким улучшением
способностей.
\end{quote}

\textbf{11. (★★★) Значение \href{mailto:pass@1}{\nolinkurl{pass@1}} для
целевой задачи очень низкое. Как использовать
\href{mailto:pass@k}{\nolinkurl{pass@k}}, долю успешно разобранных
форматов, долю частичного прогресса и анализ причин неудач, чтобы
выбрать между Mid-training, SFT и RL?}

\begin{quote}
Сначала следует разделить проблемы на три категории: недостаток знаний и
базовых способностей, ошибки формата и протокола, а также неверный выбор
стратегии. \href{mailto:pass@k}{\nolinkurl{pass@k}} показывает, удаётся
ли добиться успеха при увеличении числа попыток; доля успешно
разобранных форматов --- можно ли исполнять и оценивать траектории; доля
частичного прогресса --- существуют ли пригодные промежуточные сигналы.
При недостатке базовых способностей следует добавить соответствующее
предметное обучение; если необходимые способности уже есть, но формат
вывода или процесс выполнения операций нестабилен, следует использовать
обучение на демонстрациях; если модель способна генерировать оцениваемые
траектории с различающимся вознаграждением, можно попробовать RL.
Небольшой эксперимент также позволяет определить, какие дополнительные
задачи новый метод решил правильно, сколько будет стоить обучение и
продолжает ли модель правильно решать задачи, с которыми справлялась
ранее.
\end{quote}

\textbf{12. (★★★) В пакете rollout лишь несколько сверхдлинных ответов
ещё не завершены, тогда как остальные задачи уже выполнены. Как повысить
эффективность использования обучающего кластера и при этом
контролировать устаревание политики из-за асинхронной выборки?}

\begin{quote}
Следует использовать непрерывный батчинг и динамическое планирование для
добавления новых запросов, группировать запросы по длине либо разделить
ресурсы для выборки и обучения. Асинхронный конвейер уменьшает время
ожидания; при этом нужно записывать версию поведенческой политики и
вероятности генерации, ограничивать допустимое отставание и отслеживать
коэффициенты важности. Ограничение длины или гибкий штраф за
сверхдлинные ответы могут сократить длинный хвост, однако необходимо
различать бесполезные повторы и необходимые рассуждения, а также
правильно обрабатывать усечённые траектории.
\end{quote}

\textbf{13. (★★★) Какие отклонения симуляции могут ввести агента в
заблуждение при использовании модели для симуляции поисковой системы или
пользователя? Как калибровать обучение с помощью состояния задачи и
реальных взаимодействий?}

\begin{quote}
Симулируемый поиск может слишком легко возвращать релевантные материалы,
а симулируемый пользователь --- всегда понимать неоднозначные
формулировки или положительно оценивать невыполненные обещания. Из-за
этого модель реже практикуется в проверке, уточнении и восстановлении
после ошибок. Оценку следует по возможности связывать с проверяемыми
изменениями состояния и результатами работы, а семантическое качество
дополнительно проверять на примерах, размеченных людьми. Для
инструментов с реальными интерфейсами можно сравнивать симулируемые и
реальные ответы; симулируемого пользователя можно калибровать по
распределению поведения в реальных диалогах и посредством
авторизованного пользовательского тестирования.
\end{quote}

\hypertarget{ux433ux43bux430ux432ux430-9.-ux43dux435ux43fux440ux435ux440ux44bux432ux43dux430ux44f-ux44dux432ux43eux43bux44eux446ux438ux44f-ai-ux430ux433ux435ux43dux442ux430}{%
\section*{\texorpdfstring{\crossreflink{chap:9}{Глава 9}. Непрерывная
эволюция
AI-агента}{. Непрерывная эволюция AI-агента}}\label{ux433ux43bux430ux432ux430-9.-ux43dux435ux43fux440ux435ux440ux44bux432ux43dux430ux44f-ux44dux432ux43eux43bux44eux446ux438ux44f-ai-ux430ux433ux435ux43dux442ux430}}
\addcontentsline{toc}{section}{\crossreflink{chap:9}{Глава 9}.
Непрерывная эволюция AI-агента}

\textbf{1. (★★) Документ с накопленным опытом содержит три успешные и
одну неудачную траекторию; неудача произошла на более новой версии API.
Как отличить устаревание опыта от изменения условий его применимости?}

\begin{quote}
Сначала следует систематизировать четыре траектории по версиям API,
условиям задачи и состоянию среды, а затем сопоставить конкретные шаги,
приведшие к успеху или неудаче. Если старая стратегия была успешна на
старой версии, а в новой изменился интерфейс или предварительные
условия, прежний опыт можно сохранить, сузив область его применимости, а
для новой версии сформировать стратегию-кандидат. Если же неудача
возникает и при тех же условиях, следует проверить, не были ли упущены
какие-либо условия и случайные факторы, а затем с помощью дополнительной
проверки решить, нужно ли скорректировать или отозвать этот опыт. Для
каждой версии следует сохранять как подтверждающие, так и опровергающие
свидетельства.

Один и тот же опыт можно воспроизвести через старую и новую версии
интерфейса с одинаковыми входными данными, а затем изменить только
параметр, который предположительно стал причиной различий. Если после
добавления поля, требуемого новой версией, операция снова завершается
успешно, в документе с опытом следует указать, начиная с какой версии
это поле стало обязательным.
\end{quote}

\textbf{2. (★★) Одновременно выросли удовлетворённость пользователей
агента службы поддержки и доля нарушений правил. Как объединить сигналы
обучения и задать защитные метрики для публикации?}

\begin{quote}
Предположим, служба поддержки согласилась вернуть средства вопреки
правилам, а пользователь сразу поставил высокую оценку. Если напрямую
использовать эту оценку как вознаграждение при обучении, следующая
версия может научиться ещё охотнее давать такие же обещания. Сначала
следует сопоставить три типа записей: как пользователь оценил
обслуживание, что разрешали действовавшие в тот момент правила и что
фактически выполнил инструмент возврата средств. Этот пример можно
использовать, чтобы научить службу поддержки объяснять правила и искать
разрешённые альтернативы, однако одной высокой оценки недостаточно,
чтобы считать обещание вернуть средства образцовым поведением.

Перед публикацией следует сравнить новую и старую версии на одном и том
же наборе запросов, отдельно зафиксировав долю решённых обращений,
удовлетворённость, соответствующие правилам альтернативные решения и
лаконичность формулировок. Кроме того, по заранее согласованным
критериям следует подсчитать нарушения, утечки конфиденциальных данных,
действия с превышением полномочий, утверждения без доказательств и
случаи несоответствия обещаний фактическому исполнению. Например, обе
версии получили по 100 запросов на возврат средств: старая версия
одобрила возврат с нарушением правил один раз, а новая --- четыре раза.
Если условие публикации требует, чтобы доля таких нарушений не росла,
новая версия не может пройти проверку, даже если получила больше высоких
оценок.
\end{quote}

\textbf{3. (★★★) Некоторый агент часто утверждает, что операция
завершена, хотя записи инструмента показывают, что она ещё не выполнена.
Как в зависимости от первопричины выбрать промпт, проверку в Harness или
обучение параметров?}

\begin{quote}
Сначала следует сопоставить время ответа, запросы к инструменту и
результаты выполнения. Если модель уже получила правильное состояние, но
продолжает использовать неподходящие формулировки о завершённом
действии, сначала можно добавить явные правила формулировки и примеры.
Если состояние операции можно надёжно получить программно, следует
проверять в Harness согласованность ответа с состоянием и при
необходимости требовать исправления. Если проблема охватывает множество
задач и способов выражения, можно улучшить согласованность языка и
действий с помощью разнообразных обучающих данных. Предпочтение следует
отдавать небольшим по масштабу изменениям, которые легко проверить и
откатить.
\end{quote}

\textbf{4. (★★★) Агенту требуется изменять инструменты и компоненты
проверки бизнес-правил, а плоскость управления, утверждающая его
собственные обновления, обслуживается независимым модулем. Как разделить
полномочия, код и тесты между этими двумя сторонами?}

\begin{quote}
Агент может изменять инструменты и программы проверки бизнес-правил в
собственном рабочем пространстве и пробно запускать патч под учётной
записью с ограниченными полномочиями. Оценочные сценарии, критерии
публикации, конфигурация полномочий, журнал аудита и стабильная версия
находятся под управлением независимого модуля, а тестовая учётная запись
не имеет права их изменять. Даже изменения программы проверки
оцениваются этим модулем. В рамках предоставленных полномочий агент
может просматривать результаты диагностики и изменять патч. После
утверждения развёртывание выполняется учётной записью с правом
публикации, причём фиксируется, какой патч прошёл какие тесты и кем он
был утверждён.
\end{quote}

\textbf{5. (★★) По мере роста базы знаний с накопленным опытом возникают
ошибки поиска и конфликты правил. Как спроектировать механизмы
версионирования, учёта актуальности, объединения и удаления устаревших
данных?}

\begin{quote}
Для каждой записи с опытом следует хранить исходную траекторию, условия
применимости, версию среды, время проверки и состояние. При
возникновении конфликта сначала следует определить, относятся ли записи
к разным версиям или сценариям, и разделить их по условиям. Если записи
продолжают конфликтовать при одинаковых условиях, их нужно поместить в
очередь на проверку. Дублирующиеся записи следует периодически
объединять, а устаревшие или опровергнутые новыми свидетельствами версии
--- исключать из области поиска по умолчанию, сохраняя историю и
возможность восстановления.
\end{quote}

\textbf{6. (★★★) Спроектируйте схему непрерывной эволюции для
медицинской службы поддержки, указав распределение ответственности между
параметрами модели, знаниями, Skill и ограничениями в коде, а также
процессы проверки и публикации обновлений каждого типа.}

\begin{quote}
Обучение параметров сосредоточено на понимании намерений и способах
выражения. База знаний содержит проверенные рекомендации, инструкции к
лекарственным препаратам и правила учреждения, сохраняя их версии и
источники. Skill организует процессы сбора информации, передачи
обращения специалисту и последующего наблюдения. Код обеспечивает
аутентификацию, управление доступом к данным и утверждённую учреждением
проверку бизнес-правил. Сначала траектории следует оценивать по
фактической точности, соблюдению процессов, согласованности обещаний и
действий, а также качеству формулировок. Затем в зависимости от того, на
каком этапе произошла ошибка, следует изменить запись базы знаний,
процесс обслуживания, программу или обучающий пример. Изменения,
затрагивающие медицинское содержимое, должны проверяться специалистами.
Их следует тестировать на новых сценариях, прежних задачах и задачах
безопасности, а затем публиковать поэтапно с мониторингом аномалий.
\end{quote}

\textbf{7. (★★★) Некоторая стратегия исследования набрала максимальный
балл на всех старых деревьях открытий, но в следующем цикле исследования
в режиме онлайн показала ухудшение. Как проанализировать смещение
области поддержки, случайность и переобучение на оценке, а затем на этой
основе разделить миры воспроизведения и задать критерии допуска к
публикации?}

\begin{quote}
Старые деревья открытий охватывают только те ветви, которые фактически
исследовались прежними стратегиями; новые задачи или стратегии могут
попасть в другие области. Сэмплирование модели, шум инструментов и
изменения среды также способны повлиять на ранжирование стратегий.
Многократный выбор стратегии на одном и том же наборе деревьев легко
приводит к смещению результатов в пользу этих деревьев. Если при
воспроизведении с некоторого узла новая стратегия выбирает действие,
которое ни разу не выполнялось в этом дереве, его следует пометить как
«нет наблюдений» и завершить траекторию в этой точке. Необходимо
отдельно сообщать долю полностью воспроизведённых траекторий и
результаты на них, а неохваченные ветви считать приоритетными объектами
онлайн-проверки. Наборы для разработки, валидационный и закрытый
тестовый наборы можно разделить по семействам задач, времени и
независимым деревьям открытий; одно дерево и производные от него
траектории должны находиться в одном наборе. Необходимо зафиксировать
версии модели и инструментов, а также бюджет, после чего на нескольких
случайных начальных значениях проводить парные сравнения с базовой
стратегией по качеству, стоимости, доле нарушений и неопределённости.
После достижения заранее заданных офлайн-порогов следует провести
небольшое контролируемое сравнение в режиме онлайн, а затем постепенно
публиковать новую версию с учётом прироста и результатов защитных
метрик, сохраняя условия отката.
\end{quote}

\hypertarget{ux433ux43bux430ux432ux430-ux434ux435ux441ux44fux442ux430ux44f.-ux441ux43eux432ux43cux435ux441ux442ux43dux430ux44f-ux440ux430ux431ux43eux442ux430-ux43dux435ux441ux43aux43eux43bux44cux43aux438ux445-ai-ux430ux433ux435ux43dux442ux43eux432}{%
\section*{Глава десятая. Совместная работа нескольких
AI-агентов}\label{ux433ux43bux430ux432ux430-ux434ux435ux441ux44fux442ux430ux44f.-ux441ux43eux432ux43cux435ux441ux442ux43dux430ux44f-ux440ux430ux431ux43eux442ux430-ux43dux435ux441ux43aux43eux43bux44cux43aux438ux445-ai-ux430ux433ux435ux43dux442ux43eux432}}
\addcontentsline{toc}{section}{Глава десятая. Совместная работа
нескольких AI-агентов}

\textbf{1. (★★) После смены роли ревьюер кода унаследовал контекст этапа
анализа требований, но по-прежнему обсуждает главным образом требования.
Как обнаружить такое влияние одной роли на другую? Как скорректировать
входные данные, конфигурацию инструментов и порядок передачи работы?}

\begin{quote}
Предположим, аналитик требований только что обсуждал, какое сообщение
следует показывать при импорте пустого файла, однако в изменениях кода
ветка для пустого файла отсутствует. Сначала изучите первый запрос к
модели после передачи работы: какую версию кода, какие инструменты и
какую задачу получил ревьюер? Если он вообще не может прочитать
изменения, прежде всего дополните передаваемые материалы. Если изменения
ему доступны, проверьте, открывал ли он соответствующий код, сверял ли
интерфейсы и обработку ошибок и указал ли, какого теста не хватает. По
одному лишь упоминанию «пустого файла» в ответе трудно определить,
пересказывает ли он требования или действительно обнаружил дефект.
\end{quote}

\begin{quote}
Проведите три сравнительных прогона с одними и теми же изменениями. В
первом сохраните обсуждение целиком. Во втором оформите его в виде
документа передачи работы, чётко указав текущие обязанности, основания
требований и пункты для проверки. В третьем откройте отдельный контекст
и передайте только код, требования и материалы тестирования, необходимые
для ревью. Зафиксируйте модель, инструменты и бюджет и сравните, какие
дефекты обнаружила каждая группа, какие требования пропустила и сколько
вызовов потребовалось. Если проблема связана с конфигурацией
инструментов, отдельно измените набор инструментов, доступных ревьюеру;
проверка полномочий по-прежнему должна выполняться Harness.
\end{quote}

\textbf{2. (★★) При декомпозиции задачи Manager пропустил одну
зависимость, из-за чего нескольким дочерним агентам пришлось
переделывать работу. Как оценивать качество декомпозиции, а во время
выполнения обнаруживать и исправлять проблемы плана?}

\begin{quote}
Результат декомпозиции должен содержать входные и выходные данные
подзадач, их зависимости, ответственных и критерии приёмки; также
необходимо проверить, охватывает ли он исходные требования. До начала
выполнения можно проверить доступность интерфейсов и ресурсов и провести
небольшой пробный запуск для зависимостей с высоким риском. Во время
выполнения недостающие условия выявляются по контрольным точкам и
обратной связи от дочерних агентов, после чего следует обновить версию
плана и уведомить затронутых участников. Качество планирования
оценивается по доле переделок, времени блокировки, пропущенным
требованиям и итоговому результату. Если подтверждено, что узкое место
находится в декомпозиции задачи, для Manager в первую очередь следует
выделить более сильную модель, полную информацию о задаче и достаточный
бюджет рассуждений. Ошибочная декомпозиция заставит одновременно
переделывать работу нескольких исполнителей, и простое увеличение их
числа проблему не решит.

Например, переводчик уже перевёл три раздела до окончательного
утверждения глоссария. Добавьте «версию глоссария» во входные данные
задачи, найдите фрагменты, где использовалась старая версия, и
переделайте только затронутые части. При следующей декомпозиции можно
сначала запустить работу, не зависящую от терминологии, а переводчику
предоставить средство отслеживания изменений версии.
\end{quote}

\textbf{3. (★★) При децентрализованном взаимодействии возможны
пропущенные сообщения, споры об обязанностях и конфликты целей. Выберите
две ситуации и спроектируйте механизмы, позволяющие установить причину,
довести решение до конца и завершить спор.}

\begin{quote}
Рассмотрим совместную работу переводчика и редактора.
\textbf{Пропущенное сообщение.} Глоссарий уже обновлён до третьей
версии, но при сдаче перевода переводчик по-прежнему указывает, что
использовал вторую. Сначала по номеру задачи проверьте, было ли
отправлено уведомление о третьей версии, а затем --- подтвердил ли
переводчик его получение. Если уведомление необходимо отправить
повторно, используйте тот же номер: даже получив его дважды, переводчик
обработает его только один раз. Инцидент считается устранённым лишь
после исправления затронутых фрагментов и проверки редактором версии
сданного перевода.
\end{quote}

\begin{quote}
\textbf{Спор об обязанностях.} Переводчик считает, что старый термин на
иллюстрации должен заменить дизайнер, а дизайнер --- что переводчик
сначала должен утвердить новое название. Оба записывают номер
иллюстрации, текущий термин и основание для окончательного варианта в
одну запись задачи. Согласно заранее установленному порядку,
ответственный за терминологию утверждает название, дизайнер исправляет
иллюстрацию, а редактор сверяет изображение с текстом. Если после
нескольких раундов обсуждения новых доказательств не появилось, решение
передаётся заранее назначенному ответственному лицу.
\end{quote}

\textbf{4. (★★★) При параллельном поиске достаточно найти один
результат, соответствующий условиям. Спроектируйте полный процесс «после
успешной проверки отменить остальные ветви», учитывающий одновременное
поступление двух результатов, повторные сообщения, тайм-аут отмены и
освобождение ресурсов.}

\begin{quote}
Дочерний агент возвращает кандидата вместе с источником, после чего
Manager сначала проверяет соответствие условиям, а затем атомарной
операцией переводит задачу из состояния поиска в состояние «результат
выбран». Если два допустимых результата поступили одновременно,
выполнить переход состояния сможет только один из них; повторное событие
прочитает уже сохранённый результат обработки. При регистрации
выбранного результата в той же транзакции для остальных ветвей и
порождённых ими задач записываются ожидающие отправки уведомления об
отмене. Manager отслеживает подтверждения по номерам уведомлений и
повторно отправляет их после перезапуска, а дочерние агенты распознают
дубликаты по номеру. Получив новое уведомление, участник в безопасной
точке сохраняет прогресс, освобождает блокировки, закрывает сеансы и
возвращает подтверждение. Если выполнение не отменено в установленный
срок, среда выполнения берёт на себя его принудительное завершение и
очистку ресурсов. Для уже отправленных внешних операций необходимо
дополнительно запросить результат, а в конце свести воедино
доказательства выбранного результата и конечные состояния всех ветвей.
\end{quote}

\textbf{5. (★★★) Несколько агентов могут одновременно изменять рукопись
и номера иллюстраций, а также создавать дубликаты файлов или по ошибке
удалять общие артефакты. Как совместить полномочия, версионирование,
рабочие копии, семантические проверки и механизмы восстановления для
управления файловой системой?}

\begin{quote}
Разделите права доступа между личными рабочими областями, общими
артефактами, внешними ресурсами и системными ресурсами только для чтения
и выделите каждому участнику собственное пространство имён. При фиксации
файлов применяйте атомарную проверку версии; изменения кода и рукописи
можно выполнять в отдельных worktree, а затем централизованно
объединять. Изменение номеров иллюстраций затрагивает семантические
связи между файлами, поэтому необходимо проверять согласованность
изображений, подписей и ссылок. Блокировка каталога может ограничить
область одновременных изменений, однако семантическую корректность всё
равно следует определять проверкой связей. Сохраняйте историю версий,
защищённые резервные копии и средство восстановления, ограничивайте
полномочия на массовое удаление и проверяйте целостность артефактов
после восстановления.
\end{quote}

\textbf{6. (★★★) Как следует согласовывать и проверять качество
результата, когда агент делегирует платную задачу через Pinchwork или
RentAHuman? Если стороны не согласны с результатами приёмки, как
организовать повторную проверку, оплату и обжалование?}

\begin{quote}
До делегирования необходимо определить результат работы, критерии
качества, требования к доказательствам, срок и бюджет. Для задач с кодом
можно запускать тесты, для визуальных материалов --- проверять
рендеринг, а для работ на месте --- использовать согласованные
фотографии, записи и ручную проверку. Основания и результаты приёмки
следует сохранять вместе; при отклонении результата необходимо указать
конкретные условия исправления. Спор в соответствии с заранее
установленным порядком передаётся независимому проверяющему или человеку
и разрешается во взаимосвязи с эскроу и процедурой обжалования.
Репутационная история должна включать тип задачи, качество результата,
соблюдение сроков и исход споров, чтобы при последующих делегированиях
можно было комплексно оценивать цену и возможности исполнителя.
\end{quote}

\textbf{7. (★★) Агент может делегировать задачи людям через рынок и
платёжные интерфейсы. Какие обязанности, помимо выполнения работы на
месте, человек может брать на себя при постановке целей, вынесении
профессиональных суждений, приёмке и распределении ответственности?
Приведите пример.}

\begin{quote}
Например, когда агент организует осмотр недвижимости, человек может
помочь определить действительно важные вопросы для проверки, оценить
обнаруженные на месте отклонения, объяснить обстоятельства, не попавшие
на изображения, а затем участвовать в приёмке и последующих решениях.
Делегирующая сторона отвечает за область полномочий и бюджет,
исполнитель --- за достоверную фиксацию и передачу информации, а
профессиональные суждения выносит человек с соответствующей
квалификацией.
\end{quote}

\textbf{8. (★★) Одна и та же модель может выполнить задачу
самостоятельно или работать в составе нескольких ролей. Как
спроектировать сравнение, позволяющее определить, какую пользу по
отдельности приносят специализация ролей, изоляция контекстов,
параллельный поиск и внешняя обратная связь?}

\begin{quote}
Зафиксируйте набор задач, модель, инструменты и метод оценки и
установите сопоставимые бюджеты для одного агента и нескольких агентов.
По отдельности изменяйте разделение ролей, организацию контекста,
параллельное выполнение и каналы обратной связи, наблюдая за качеством,
реальным временем выполнения, расходом токенов и распределением ошибок.
Внешнюю обратную связь следует предоставлять и одному агенту в
контрольном варианте, чтобы не приписать всю пользу инструментов
совместной работе. Например, сначала зафиксируйте общий предел токенов и
число вызовов и поручите одной и той же группе исполнителей провести
четыре попарных эксперимента: одинаковые роли с общим контекстом, разные
роли с общим контекстом, одинаковые роли с изолированными контекстами и
разные роли с изолированными контекстами. Сравнение первых двух
вариантов покажет эффект разделения ролей; такое же сравнение последних
двух позволит установить, зависит ли польза разделения от изоляции
контекстов. Затем для выбранной конфигурации сравните последовательное и
параллельное выполнение, отдельно зафиксировав время ожидания
пользователя, общий расход токенов системой и сведения, потерянные при
передаче работы.
\end{quote}

\textbf{9. (★★★) Используя идею прозрачности мыслей из романа «Задача
трёх тел» и вопрос о постановке целей из мысленного эксперимента со
скрепками, обсудите обмен информацией, общую цель и разнообразие точек
зрения в системе с несколькими агентами. Как с помощью наблюдаемого
эксперимента определить, подходят ли выбранные границы обмена
информацией и этап независимого анализа?}

\begin{quote}
Литературные идеи и мысленные эксперименты помогают сформулировать
вопросы, но экспериментальная проверка должна опираться на объём
информации, цели и показатели поведения. Можно сравнить варианты, в
которых все выводы немедленно становятся общими, сначала проводится
независимый анализ и только потом объединяются результаты либо участники
обмениваются лишь исходными доказательствами, сохраняя одинаковыми
задачи и бюджеты. Наблюдайте за различиями между кандидатами,
корреляцией ошибок, итоговым качеством и затратами на координацию, а
также проверяйте, не упущены ли ограничения в общей цели. Можно сначала
поручить перекрёстным проверяющим изучить исходные доказательства и
независимо сформировать суждение, а уже затем при необходимости открыть
им предыдущие объяснения, чтобы уменьшить влияние готовых выводов на
проверку.
\end{quote}

\textbf{10. (★★★) Один и тот же Coding Agent получает бюджеты в 30 и 300
шагов. Как распределить планирование, реализацию, проверку и доработку и
корректировать стратегию с учётом оставшегося бюджета и фактического
прогресса?}

\begin{quote}
При обоих бюджетах необходимо сохранить базовое планирование и
обязательную проверку. При 30 шагах в первую очередь следует определить
основные требования, выбрать более короткий путь и оставить запас на
проверку результата и обработку исключительных ситуаций. При 300 шагах
можно добавить сравнение вариантов, тестирование граничных случаев и
независимое ревью. Во время выполнения записывайте общий и оставшийся
бюджеты, завершённые пункты и нерешённые проблемы, а по контрольным
точкам определяйте, следует ли продолжать исследование, сосредоточиться
на доработке или запросить входные данные. В журнал должны попадать как
выделенный бюджет, так и фактический расход. Если на протяжении
нескольких раундов подряд прогресса нет, измените стратегию, чтобы не
тратить дополнительные шаги на повторение одних и тех же действий.
\end{quote}

\textbf{11. (★★) Какие идеи для проектирования систем с несколькими
агентами можно почерпнуть соответственно из виртуальной памяти и
страничной организации, прав доступа к файлам, обнаружения
взаимоблокировок и алгоритмов планирования? Укажите границы применимости
этих аналогий и ограничения, исполнение которых должна явно обеспечивать
среда выполнения.}

\begin{quote}
Страничная организация памяти может служить аналогией для извлечения
информации по запросу и управления горячими и холодными данными; права
доступа к файлам --- для монтирования только для чтения, авторизации
инструментов и изоляции учётных данных; граф ожидания помогает
обнаруживать взаимное ожидание участников; алгоритмы планирования
позволяют управлять параллелизмом, приоритетами и квотами ресурсов.
Сводка контекста изменяет представление информации, тогда как подкачка
памяти обычно сохраняет байтовое содержимое. Семантика задачи зависит от
суждения модели, но ограничения полномочий и бюджета должна обеспечивать
среда выполнения. Необходимо также отличать нормальный цикл доработки от
взаимоблокировки или активной блокировки без прогресса, диагностируя их
по состояниям, зависимостям и времени.
\end{quote}

\backmatter
\end{document}
